\documentclass[UKenglish]{uiomasterthesis}
\usepackage{{booktabs}}
\usepackage{algpseudocode}
\usepackage{algorithm}
\usepackage{graphicx} 
\usepackage[UKenglish]{uiomasterfp}
\usepackage[nottoc]{tocbibind}
\usepackage[hidelinks]{hyperref}
\usepackage{amsmath}
\usepackage{amssymb}
\usepackage{amsfonts}

\usepackage{amsthm}
\usepackage{bm}
\usepackage{mathtools}
\usepackage[printonlyused]{acronym}
\usepackage{tikz}
\usepackage{pgfplots} 
\usepackage{pgf} 
\usepackage[backend=biber, style=numeric-comp, sorting=none]{biblatex}
\addbibresource{ref.bib}
\usetikzlibrary{arrows, arrows.meta, positioning, calc}
\usepackage{listofitems}

\pgfplotsset{compat=1.18}

\def\mathdefault#1{#1}
\everymath=\expandafter{\the\everymath\displaystyle}

\title{Explainable Reinforcement Learning}
\subtitle{Discovering intent based explanations for heterogeneous cooperative multi agent reinforcement learning agents}
\author{Ada Hatland}
\date{August 2024}


\pagenumbering{roman}

\begin{document}

\uiomasterfp[master, program={Informatics: Robotics and Intelligent Systems},
  color=orange, dept={Department of Informatics}, fac={The Faculty of Mathematics and Natural Sciences},
  supervisors={Dr. Dennis Gro\ss \and Prof. Kyrre Glette\and Dr. Helge Spieker}, image = {images/cat.png}]

\renewcommand*\acffont{\textit}

\section*{Declaration of AI use}
In this thesis generative models have been used for topic suggestion, a small part of figure creation, equations and parts of the code used for creating plots, as well as debugging and documentation. All code suggested has been verified by a human. Generative models have \textbf{not} been used for text generation. I, the author, take full responsibility for the contents of this thesis.


\section*{Acknowledgments}
I would like to thank my external supervisor Dennis Gro\ss, as well as my internal supervisor Kyrre Glette for their guidance and support in writing this thesis. They have been of significant help when discussing how to proceed and what remains on the to-do list. I would also like to thank my cat Nokia for emotional support throughout the writing of this thesis.

\section*{List of Acronyms}
\begin{acronym}[ICANN]
    \acro  {rl}   [RL]   {Reinforcement Learning}
    \acro  {xrl} [XRL] {Explainable Reinforcement Learning}
    \acro  {xdrl} [XDRL] {Explainable Deep Reinforcement Learning}
    \acro  {ai}   [AI]   {Artificial Intelligence}
    \acro  {xai}  [XAI]  {Explainable Artificial Intelligence}
    \acro  {marl}  [MARL]  {Multi Agent Reinforcement Learning}
    \acro  {mdp}  [MDP]  {Markov Decision Process}
    \acro  {dqn}  [DQN]  {Deep Q-Network}
    \acro  {ppo}  [PPO]  {Proximal Policy Optimization}
    \acro  {drl}  [DRL]  {Deep Reinforcement Learning}
    \acro  {rnn}  [RNN]  {Recurrent Neural Network}
    \acro  {lstm}  [LSTM]  {Long Short Term Memory}
    \acro  {mbrl}  [MBRL]  {Model-Based Reinforcement Learning}
    \acro  {llm}  [LLM]  {Large Language Model}
    \acro  {ea}  [EA]  {Evolutionary Algorithm}
    \acro  {nsga}  [NSGA-II]  {Non-Dominated Sorting Genetic Algorithm II}
    \acro  {nn}  [NN]  {Neural Network}
    \acro  {aec}  [AEC]  {Agent Environment Cycle}
    \acro  {cnn}  [CNN]  {Convolutional Neural Network}
    \acro  {drc}  [DRC]  {Deep Repeated ConvLSTM}
    \acro  {pomdp}  [POMDP]  {Partially Observable Markov Decision Process}
\end{acronym}


\abstract
This is the abstract

\lefthyphenmin=1000
\tableofcontents
\listoftables
\listoffigures
\chapter{Introduction}

\pagenumbering{arabic}

Sequential decision-making problems are problems where one decision leads to a new state that requires a new decision to be made. An example of this is autonomously driving from point A to point B through city streets. For these types of problems \ac{rl} systems are used. \ac{rl}, both multi- and single-agent models have seen a significant rise in successful use and applicability in recent years, with models such as AlphaGo\cite{article} and smacv2\cite{ellis2023smacv2}. These models learn by agents performing actions in a given state decided by a policy that lead to a new state, and learning by receiving rewards depending on if the new state is preferable to the previous, the aim is to learn a near optimal policy for achieving a fixed goal\cite{Sutton1998}. The agent uses an observation of the state to choose an action. The environment describes how an action affects the state. See Figure \ref{fig:rl_basics}. In a \ac{marl} system, we would have multiple agents.

\section{Motivation}
\label{sec:motivation}
A significant problem with many deep machine learning models, including \ac{marl} policies, is known as the black box problem\cite{zednik2019solving}. This problem describes how the processing of information is hidden from a human user due to the opacity of the model, and we just have to trust that the model uses relevant data to come to the conclusion it does, which it often doesn't do. For many tasks and models where the impact of the output isn't highly significant this isn't a big issue, but for tasks like autonomous driving we simply cannot use these models without trust for the models processing of data and that the solution given wont hurt anyone. We can considering autonomous driving as a \ac{marl} system if we consider each vehicle as an agent. To solve this we would like to have models that along with an output can provide us with some kind of reasoning for what information is used and how.

All this essentially means, until the black box problem is solved, autopilot will always require a driver behind the wheel\cite{tian2018deeptest}. Despite having high accuracy, precision and recall, a reinforcement learning model might choose a less than preferable action in edge cases or states not well covered by training and test sets, for example driving on snowy roads when all images in the training data is from warmer climates, or combinations that aren't well covered, like a foggy, snowy road with a sharp right turn.
With a way to ask the agent for intent, we could find that it has learned to always expect there to not be a car around the corner, which is not always obvious from looking at the dataset, but if it in a decision relies on the road to be empty to choose a safe route, this is obviously a huge issue. This paper aims to explain how an agent in a \ac{marl} setting decides on an action due to what it expects from other agents in the future. The field working on combating this issue is known is Explainable reinforcement learning.

\section{Problem statement}
\label{sec:problem}
The focus of this paper will be to rework and apply methods for general \ac{xai} to Reinforcement Learning, and expand on \ac{xrl} methods already developed. In particular we will focus on an agent and how the expected future states of other agents will affect it in its decision making process. The environment we will focus on is a \ac{marl} environment known as Knights Archers Zombies \cite{KAZ} made for Python. We use this because it is a cooperative environment where each agent has to consider other agents, both of the same type as itself, and other types. If we used an environment where all the agents are identical, our findings will be less useful for other environments where the agents aren't identical. In many real life scenarios agents will not be identical. If we once again consider autonomous vehicles, most vehicles are different in some way. A very obvious example of this is considering cars and motorcycles. Where because of their size difference, their paths chosen will often be different.

In psychology it is well known that most human based decisions are made with intent\cite{inbook}, and by focusing on the expectation of what future states will look like we can consider this the intent of the agent. If we are accurately able to extract the intent of an \ac{rl} policy we are better able to explain the choice made and we are more comfortable with trusting that the choice made is a good decision. Concretely we hypothesise that by analyzing expected future states and events we could find the intent of the agent, which we could use to ensure the agent acts in a safe or desirable manner. The main question we aim to answer is what information, found in weights or design of the policy, or information extracted by knowing the weights and design, is the most important for separate prediction models to be able to predict future states and actions, in cases where we do not have access to a decent world model, or such a model does not exist.


\section{Scope and limitations}
We aim to construct a framework to describe the intent of the agents in the KAZ environment, and we aim to make it applicable to other environments as well. However, due to time limitations we restrict the framework to only accept PettingZoo example environments or environments created with the PettingZoo environment creation tool. \ac{marl} environments can be divided into two major categories, \ac{aec} and parallel. \ac{aec} environments are environments where the state transition happens after each agent chooses an action, and parallel environments are environments where all the agents choose an action each before the transition happens. See figure \ref{fig:aec}.

Due to limited time and resources the size of policy or value networks, as well as any other \acp{nn} or otherwise computationally expensive algorithm used, will have to be significantly limited, this will impact results in a meaningful way, and therefore we will keep all \acp{nn} comparable sizes, and focus on doing experiments on how other factors than network size impact the results.

We do not have access to any real world datasets, so its hard to ensure the methods we have developed will be directly applicable, however the results we get will still be relevant for later research in the area.
\begin{figure}
\begin{center}

\begin{tikzpicture}[>=Latex, node distance=2.5cm, auto, font=\small]

    % Agent 1 Node
    \node[draw, thick, circle, fill=blue!20, minimum size=2.5cm, align=center] (agent1)
    at (-3,3) {Agent 1 \\ \(\pi_{\theta_1}\)};

    % Transition 1 Node
    \node[draw, thick, circle, fill=green!20, minimum size=2.5cm, align=center] (trans1)
    at (0,6) {State \\ Transition 1};

    % Agent 2 Node
    \node[draw, thick, circle, fill=blue!20, minimum size=2.5cm, align=center] (agent2)
    at (3,3) {Agent 2 \\ \(\pi_{\theta_2}\)};

    % Transition 2 Node
    \node[draw, thick, circle, fill=green!20, minimum size=2.5cm, align=center] (trans2)
    at (0,0) {State \\ Transition 2};

    % Arrows between nodes
    \draw[->, thick] (agent1.north) -- node[left] {\(a^1\)} (trans1.west);
    \draw[->, thick] (trans1.east) -- node[right] {\(o^1\)} (agent2.north);
    \draw[->, thick] (agent2.south) -- node[right] {\(a^2\)} (trans2.east);
    \draw[->, thick] (trans2.west) -- node[left] {\(o^2\)} (agent1.south);

\end{tikzpicture}

\caption{Example \ac{aec} in an environment with 2 agents.}
\label{fig:aec}
\end{center}
\end{figure}

\section{Research methods}
\label{sec:research}
This section describes the methods we use to answer the questions we have posed. The methodology mainly focuses on produced datasets made with a digital \ac{rl} environment, which does mean we have access to a world model, however as we want to answer the research questions without the use of a world model we will only be using the results of simulations to create our methods. The effect of this is that the methods will be compatible with \ac{rl} systems where we do not have access to a world model.

Producing our own dataset has the benefit of no restrictions, other than hardware restrictions, on the datasets size. We can also produce variations of these datasets by altering the simulations used to create them which would in a real world setting often be impossible, e.g. access to gradients during inference.

The datasets produced will be sets of full or partial trajectories in given environments along with other relevant data like integrated gradient, or Shapley values for the observations in the trajectories. We use these datasets to train predictive \acp{nn} with two main types. Event prediction, for example whether an agent encounters a critical state within the next 10 timesteps, which could represent a dangerous situation for an autonomous vehicle, and state prediction, which for example could be the x and y coordinates of a given agent 10 timesteps into the future.

We will first be focusing on including explanations for the observations and observe and analyze the effects on event and state predictors, then we will be using \ac{nn} policies with access to memory and include their representation of memory along with the observation, hidden values for lstm, and attention matrix for GTrXL, and once again observe and analyze effect on the predictors.

\section{Ethical considerations}
\ac{marl} is an important field in development for military purposes\cite{military_marl}, and while this research is not supposed to effectivize warfare its possible that it ends up being relevant for unintended areas of research, including but not limited to military use-cases. This is especially the case because I aim to develop methods for general \ac{marl} use and not a specific area of research, like medicine or sports. It is more or less impossible to restrict access to these methods for specific areas of research.

If this research aids to make \ac{marl} methods applicable to real world settings there is the possibility it risks increasing job displacement issues. This is an issue that affects most \ac{ai} research and can be mitigated by programs to reeducate people whose jobs it might affect, and increasing employment in developing and monitoring \ac{ai} tools.


\section{Main contributions}
As discussed in \ref{sec:motivation}, \ref{sec:problem} and \ref{sec:research}, there are certain specific objectives we aim to reach. In general we aim to expand on current future prediction research.

\begin{itemize}

    \item \textbf{Objective 1: Observation space}\\
        Most current methods assumes the observation space to be image based, and bases research on information found in convolutional layers, either regular \acp{cnn}, or convLSTM layers.

    \item \textbf{Objective 2: Feature importance as input}\\
        Using feature importance from the policy to improve future predictions is a novel idea, and has potential to improve future predictions by a significant margin.

    \item \textbf{Objective 3: Memory as input}\\
        Earlier research has found that including hidden states from the \ac{drc} architecture had a significant effect on accuracy when predicting future states and events. We aim to expand on this and use hidden states from \acp{lstm} layers and observe the effect and analyze how it differs from using the ConvLSTM hidden values, as well as including attention from the GTrXL architecture.

\end{itemize}

\section{Thesis outline}
This thesis will be structured as follows:
\begin{itemize}

    \item \textbf{Chapter 2: Background and related works}\\
        This chapter will introduce important concepts, both \ac{rl} and explainability related, and briefly discuss the history of \ac{xai} methods, and how they are applicable to my use case. We will also consider \ac{xrl} methods that are already developed and how to expand or integrate them into developed methods of our own.

    \item \textbf{Chapter 3: Methodology}\\
        This chapter discusses the specific principles, procedures and techniques used to conduct the \ac{xrl} research done in this thesis, to ensure that the results we get are valid and reliable.

    \item \textbf{Chapter 4: Experiments}\\
        This chapter details the setup, results and analysis of the specific experiments we use to evaluate the hypothesis and answer research questions based on certain evaluation metrics, and baseline comparisons.

    \item \textbf{Chapter 5: Discussion}\\
        This chapter will analyze the results of the experiments and discuss their implication for \ac{xrl} research. It will also contain the limitations of our work and briefly discuss how to expand on any research done.
\end{itemize}


\begin{figure}
\begin{center}
    \hspace{-4cm}
    \begin{tikzpicture}[>=Latex, node distance=2.5cm, auto, font=\small]
    % Agent Node (with Policy and Value Networks)
    \node[draw, thick, rectangle, fill=blue!20, minimum width=3cm, minimum height=2cm, align=center] (agent)
    at (-4,3) {Agent \\ \(\pi_\theta(s)\) \\ \(V_\phi(s)\)};

    % Environment Node
    \node[draw, thick, rectangle, fill=green!20, minimum width=3cm, minimum height=2cm, align=center] (env)
    at (4,3) {Environment};

    % Arrow: Agent to Environment (Action)
    \draw[->, thick] (agent.east) -- node[above] {\(a_t\)} (env.west);

    % Rollout/Trajectory Buffer Node
    \node[draw, thick, rectangle, fill=yellow!20, minimum width=4cm, minimum height=1.5cm, align=center] (buffer)
    at (0,0.5) {Rollout/Trajectory Buffer \\ \((s_t,\, a_t,\, r_t,\, s_{t+1})\)};

    % Arrow: Environment to Rollout Buffer (State & Reward)
    \draw[->, thick] (env.south) .. controls +(0,-1) and +(2,1) .. node[right] {\(s_{t+1},\, r_{t+1}\)} (buffer.east);

    % Arrow: Agent to Rollout Buffer (State, Action, Reward)
    \draw[->, thick] (agent.south) .. controls +(0,-1) and +(-2,1) .. node[left] {\(s_t, a_t, r_t\)} (buffer.west);

    % PPO Update Node
    \node[draw, thick, rectangle, fill=red!20, minimum width=4cm, minimum height=2cm, align=center] (ppo)
    at (0,-3) {PPO Update \\ \(\bullet\) Compute Advantages \\ \(\bullet\) Clip Objective \\ \(\bullet\) Update \(\pi_\theta\) and \(V_\phi\)};

    % Arrow: Rollout Buffer to PPO Update
    \draw[->, thick] (buffer.south) -- node[right] {Optimize PPO Objective} (ppo.north);

    % Arrow: PPO Update to Agent (New Parameters)
    \draw[->, thick] (ppo.west) .. controls +(-7,0) and +(0,0) .. node[right] {New parameters \(\theta, \phi\)} (agent.west);
\end{tikzpicture}
\caption{Diagram showing the \ac{ppo} update algorithm}
\label{fig:rl_basics}
\end{center}
\end{figure}

\medskip
\chapter{Background and Related Works}

\section{Reinforcement Learning Fundamentals}
To make sure all methods are understood well its important to make sure we have a proper understanding of the fundamentals of \ac{rl}, especially \ac{drl} and \ac{marl}. This section will go over the basics needed.


\subsection{Markov Decision Process in Reinforcement Learning}

A \ac{mdp} is a framework used to model decision-making in stochastic environments, like we want to in an \ac{rl} task. It is defined by a tuple with 5 elements:

$$\mathcal{M} = (\mathcal{S}, \mathcal{A}, P, R, \gamma)$$

where \(\mathcal{S}\) is the set of possible states, \(\mathcal{A}\) is the set of possible actions, \(P(s' | s, a)\) is the transition probability function, which defines the probability of moving to state \(s'\) when action \(a\) is taken in state \(s\), \(R(s, a, s')\) is the reward function that provides a reward for transitioning from state \(s\) to \(s'\) via action \(a\), \(\gamma \in [0,1]\) is the discount factor that determines the importance of future rewards.

The objective in an \ac{mdp} is to find a policy \(\pi(s)\) that maximizes the expected cumulative reward:

$$V^\pi(s) = \mathbb{E} \left[ \sum_{t=0}^{\infty} \gamma^t R(s_t, a_t, s_{t+1}) \mid \pi \right]$$

where \(V^\pi(s)\) is the value function that represents the expected return starting from state \(s\) and following policy \(\pi\). See figure \ref{fig:rl_basics}.


\subsection{Markov Decision Process in Multi Agent Reinforcement Learning}
In a \ac{marl} setting agents could have shared or independent \acp{mdp} depending on architecture. There are three main branches of \ac{marl}, cooperative, competitive and mixed. Our focus will be on cooperative \ac{marl}.
In a cooperative \ac{marl} setting the goal is some social welfare function that maximises rewards for the agents, either collective rewards or a mix of collective and individual rewards. In such a setting each agent has \ac{mdp}. Assuming that state space \(\mathcal{S}\) and discount factor \(\gamma\) is shared, which they will be in all my experiments, if two agents have the same reward structure $R$ and action space \(\mathcal{A}\) their objective would be the same, in which case they could share policy \(\pi\) and could therefore share \(\mathcal{M}\).

\subsection{Deep Reinforcement Learning}
Traditional RL methods struggle with scalability due to the fact that they rely on discrete state representations. \ac{drl} uses deep \acp{nn} to approximate the policy function $\pi_\theta(s)$, the value function $V_\phi(s)$, or the Q-function $Q_\theta(s,a)$, making it possible to represent these functions in continuous and complex environments. See figure \ref{fig:neural_network}. Input layer is size of state for all three functions. Output layer is size of action space for $Q_\theta(s,a)$ and $\pi_\theta(s)$. $V_\phi(s)$ only has one output node. In a feed forward \ac{nn} like we have here each node has the value of some activation function $\phi(x)$ where $x$ is the sum of nodes of the previous layer multiplied by their respective weights, usually as well as adding some bias, shown as the connections between the nodes. Common activation functions are $ReLU(x) = max(0,v)$ and $tanh=\frac{e^v-e^v}{e^v+e^v}$ where $v$ is the value before activation.
\begin{figure}[H]
    \begin{center}
    \begin{tikzpicture}[
        x=2.2cm, y=1.4cm, 
        mynode/.style={thick, draw=black, circle, minimum size=22}, % Base style
        inputnode/.style={mynode, fill=blue!35},  % Input layer
        hiddennode/.style={mynode, fill=red!35}, % Hidden layers
        outputnode/.style={mynode, fill=green!35} % Output layer
    ]
    
    \readlist\Nnod{4,5,5,5,3} % Number of nodes per layer
    
    \foreachitem \N \in \Nnod{ % Loop over layers
        \foreach \i [evaluate={\x=\Ncnt; \y=\N/2-\i+0.5; \prev=int(\Ncnt-1);}] in {1,...,\N}{ 
            
            % Choose color based on layer type
            \ifnum\Ncnt=1
                \node[inputnode] (N\Ncnt-\i) at (\x,\y) {}; % Input layer
            \else
                \ifnum\Ncnt=\Nnodlen
                    \node[outputnode] (N\Ncnt-\i) at (\x,\y) {}; % Output layer
                \else
                    \node[hiddennode] (N\Ncnt-\i) at (\x,\y) {}; % Hidden layer
                \fi
            \fi
            
            % Connect to previous layer
            \ifnum\Ncnt>1  
                \foreach \j in {1,...,\Nnod[\prev]}{ 
                    \draw[thick] (N\prev-\j) -- (N\Ncnt-\i); 
                }
            \fi 
        }
    }
    
    \node[align=center, below=1] at (N1-4.90) {Input Layer};
    \node[align=center, below=1] at (N3-5.90) {Hidden Layers};
    \node[align=center, below=1] at (N5-3.90) {Output Layer};
    \end{tikzpicture}
    \end{center}
    \caption{Illustration of an example of a \ac{nn} with three hidden layers}
    \label{fig:neural_network}
\end{figure}


\subsection{Model-Free vs. Model-Based Reinforcement Learning}

Model-based and model-free reinforcement learning differ in how they represent environment transitions. In model-based reinforcement learning, the agent learns or has access to a model of the environment, which includes a transition function \( P(s' | s, a) \) and a reward function \( R(s, a) \) that determines the expected reward for taking action \( a \) in state \( s \) \cite{moerland2022modelbasedreinforcementlearningsurvey}. This allows the agent to simulate future outcomes without taking actions in the real environment. As a result, model-based methods are able to update policies based on imagined rollouts, which increases learning speed compared to executing the real environment transitions. Examples of model-based algorithms are Dyna\cite{10.1145/122344.122377} and DreamerV3\cite{hafner2024masteringdiversedomainsworld}.

In contrast, model-free reinforcement learning does not learn or utilize an explicit model of the environment. Instead, the agent learns by interacting directly with the environment and updating its network weights based on observed rewards. This approach is generally less sample-efficient but is applicable to environments where modeling the transition dynamics isn't feasible. Examples of model-free algorithms include Q-learning, \ac{dqn} \cite{mnih2013playingatarideepreinforcement}, and \ac{ppo} \cite{schulman2017proximalpolicyoptimizationalgorithms}.

Model-based methods enable explicit planning by using the learned environment transitions in methods such as Monte Carlo Tree Search \cite{_wiechowski_2022}, while model-free methods rely on direct experience to optimize behavior. We will be using \ac{ppo} to optimize our algorithms, and will not be using simulations or any other form of explicit planning, ensuring that the methods developed will be applicable in the most amount of environments.

\subsection{Recurrency in Deep Reinforcement Learning Policies}
Recurrency in \ac{nn} is a way to carry over information from previous timesteps to current ones, if we believe that information could be useful. \acp{rnn} are sometimes used when constructing the architecture of \ac{drl} policies \cite{hausknecht2017deeprecurrentqlearningpartially}, mainly for two reasons. The first one is for when your environment can be modeled as a \ac{pomdp}, in which case memory is important as the current observation is not always enough to capture all the information an agent could have. The second reason is for planning, like in the \ac{drc} architecture\cite{guez2019investigationmodelfreeplanning}. See figure \ref{fig:rnn}.

\begin{center}
\begin{figure}
    \hspace{3.5cm}
    \begin{tikzpicture}[>=stealth, auto, node distance=1.5cm and 2cm]

    \begin{scope}[shift={(4,0)}, rotate=90]
  
    % Time step 1
    \node (x1) [draw, rectangle, fill=blue!20, minimum width=1cm, minimum height=0.7cm] {$x_1$};
    \node (h1) [draw, circle, fill=red!20, right=of x1] {$h_1$};
    \node (y1) [draw, rectangle, fill=green!20, right=of h1, minimum width=1cm, minimum height=0.7cm] {$y_1$};

    % Time step 2
    \node (x2) [draw, rectangle, fill=blue!20, below=of x1, minimum width=1cm, minimum height=0.7cm] {$x_2$};
    \node (h2) [draw, circle, fill=red!20, right=of x2] {$h_2$};
    \node (y2) [draw, rectangle, fill=green!20, right=of h2, minimum width=1cm, minimum height=0.7cm] {$y_2$};

    % Time step 3
    \node (x3) [draw, rectangle, fill=blue!20, below=of x2, minimum width=1cm, minimum height=0.7cm] {$x_3$};
    \node (h3) [draw, circle, fill=red!20, right=of x3] {$h_3$};
    \node (y3) [draw, rectangle, fill=green!20, right=of h3, minimum width=1cm, minimum height=0.7cm] {$y_3$};

    % Forward connections within each time step (black)
    \draw[->, thick, black] (x1) -- (h1);
    \draw[->, thick, black] (h1) -- (y1);
  
    \draw[->, thick, black] (x2) -- (h2);
    \draw[->, thick, black] (h2) -- (y2);
  
    \draw[->, thick, black] (x3) -- (h3);
    \draw[->, thick, black] (h3) -- (y3);

    % Recurrent connections between hidden states (orange)
    \draw[->, thick, orange] (h1) -- (h2) node[midway, right=4pt] {$W_h$};
    \draw[->, thick, orange] (h2) -- (h3) node[midway, right=4pt] {$W_h$};
    \end{scope}
\end{tikzpicture}
\caption{Recurrent neural network showing a 3 step time series with hidden values propagating through time.}
\label{fig:rnn}
\end{figure}
\end{center}

\subsection{Backpropagation and gradient descent}
Backpropagation is a training algorithm used to train \acp{dnn} to minimize error. If a network is differentiable you can use backpropagation. First you do a forward pass through the network, i.e. your network calculates a prediction based on an input. Then you calculate loss, the difference between true value and predicted value is passed to some function. MSE is one such function.
\[
\text{MSE} = \frac{1}{n} \sum_{i=1}^{n} (y_i - \hat{y}_i)^2
\]
Where $y$ is true value, $\hat{y}$ is predicted value, $n$ is amount of samples multiplied by outputs. If you are trying to predict $(x,y)$ coordinates, you would have two outputs, if you're trying to predict $x,y$ coordinates in 5 cases you would have $n = 10$.
Then you do a backwards pass, calculating loss with respect to the weights in the network using the chain rule. Then using gradient descent you change the weights to reduce loss.


\section{Approaches to Explainability in Reinforcement Learning}
\subsection{Post Hoc vs. Intrinsically Explainable Models}
This section will discuss benefits and drawbacks of explainable models divided into two main categories, Post hoc models and intrinsically explainable models.

Post hoc explainability refers to finding ways to understand already trained models. Instances where this doesn't pose more of a challenge can often be more efficient as you do not need to construct and train a model, which could often pose an issue both considering time spent and accuracy of the model. Making a model explainable only makes sense if we know it usually makes sensible decisions. There are however issues with making post hoc models.

\subsubsection{Temporal Aspect}
\ac{rl} policies all have a temporal aspect, which means several different actions over a certain period of time might contribute to a singular outcome, this could make it very challenging to pinpoint which actions were made for which outcome, especially if the outcome happens several states after the initial decision was made.

\subsubsection{Nature of Black Box Models}
Many \ac{rl} policies, especially deep \ac{rl} policies, which are the policies we will be working with, have complex inner connections due to the hierarchical structure of learning features in the input, sequences of non linear connections that are very hard to understand without spending a lot of time studying the specific connections learned by a model. It's often very difficult to accurately pinpoint the purpose of a node, especially because we do not know if it even has a purpose at all or it could have a lot of different purposes. Especially in environments with complex observation spaces.


\subsubsection{Lack of Transparency}
Tracing how an observation leads to a state, through the action chosen and the environment, is often challenging in models constructed without this in mind. If we do not know the processing of information of a model it's very hard to explain the rationalisation of an agent.

\subsubsection{Intrinsically Explainable Models}
Intrinsically explainable models have the already stated drawback of relying on careful model construction, which can be challenging and may impact model performance. While these models offer transparency from the get go, they are often less flexible and might not achieve the same level of accuracy as more complex \ac{rl} models. Given that we aim to analyze and compare already trained \ac{rl} policies, we will mostly focus on post hoc explainability methods in this thesis.

\medskip

\subsection{Fidelity of Explainable Artificial Intelligence Methods}
It is important to understand why the current \ac{xai} methods not meant for \ac{rl} applications aren't always useful to us.

\ac{xai} methods not intended for reinforcement learning often provide\\ explanations that don't necessarily represent the inner workings of an \ac{rl} policy, because an \ac{rl} policy has a temporal aspect as well. Broadly these methods can be categorised as feature-based explainers, and they often struggle to fully explain an agents behaviour or specific actions because they cannot capture the future view of an \ac{rl} policy. 

Saliency maps which have been successfully used for classification of images provide explanations about what part of the input was important for the outcome, which is highly relevant for classification tasks, but using the same method for an \ac{rl} policy doesn't sufficiently explain the intent of the agent\cite{atrey2020exploratory}.

Another commonly used \ac{xai} method is model distillation, which works by transferring knowledge from a large model to a smaller, usually interpretable, one, for instance a deep learning network to a decision tree \cite{bastani2019verifiable}. This has use cases in verifyability, but struggles to fully explain the temporal aspect of \ac{rl} policies, and are therefore not sufficient as an explainer.

However, these methods might still prove insightful in conjunction with other intent-based methods, the state in which a decision is made is obviously very relevant to why that particular decision was made. We could perhaps use these methods to answer questions such as "What part of agent As Observation this state, lead it to believe agent B would end up at these coordinates at a later state?"
\medskip

\subsection{Future-Based Explanations}
Next, I will describe in slightly more detail, what is meant by an intent based explainer, like I want to develop, and how to use it.

\subsubsection{Design}
"Towards Future-Based Explanations for Deep RL Network Controllers"\cite{10.1145/3626570.3626607} broadly describes future-based intrinsically explainable methods. Future-based intrinsically explainable methods for \ac{drl} policies often take three inputs, the trajectories experienced by an agent, the environment and the agent. Then, they collect the rewards and interactions and use this information to train an explainer.

During inference, we can then apply the explainer to a state and action to get the expected future consequences of that action. Depending on the architecture we could either get the expected future consequence of any action, or just the one the agent decides on.

\subsubsection{Use Cases}
Designing a \ac{drl} solution requires choosing features in the observation, hyperparameter tuning, policy design and reward function among other things. This is a time, and resource, consuming process. These are usually picked by trial-and-error but could be made easier with assistance from an explainer.

Another, and perhaps more important, use case is safety. If when online, i.e. the chosen action will actually affect the state, we are expecting high likelihood of an unsafe state, we could instead of opting for the chosen action fall back to a known safe action, that could have a lower expected return, i.e. breaking instead of turning a corner if we have limited vision around the corner.


\subsubsection{Definition of Intent}
We define the state transition function $T(s, \textbf{a}, s')$ where $\textbf{a}$ is the set of simultaneous actions made by the set of active agents, in our case, one action per agent. Given $T$ the agent should when prompted output a set of trajectories $\tau = \{(s_0,\textbf{a}_0),(s,\textbf{a}),(s_1,\textbf{a}_1)...(s_n,\textbf{a}_n)\}$ where $s_0$ is the current state, $\textbf{a}_i$ is the set of actions taken in $s_i$, and $s_{i+1}$ is the state reached by $T(s_i,\textbf{a}_i, s_{i+1})$. We will apply a series of methods on $\tau$ to extract intent. One of these methods is discovering counterfactual trajectories, $'\tau$, where the actions made by the agents in $'\tau$ are as similar as possible to the actions made by the agents in $\tau$, but the total reward gained is as low as possible. $'s_0 = s_0$ and given an identical transition function $T$ the goal is to discover which actions are most important to receive the reward. This method is similar to and inspired by ACTER\cite{gajcin2024acter}, see section \ref{sec:acter}. Details in section \ref{sec:counterfactual}.

\section{Explainable Deep Reinforcement Learning: State of the Art and Challenges}
This paper is a comprehensive survey of the most common state of the art methods applicable to \ac{rl}. Section 4 specifically highlights \ac{xdrl} methods. This subsection will focus on these methods\cite{sota}. This section will serve to give an overview of \ac{xrl} methods that already exist, and if and how I use these methods to develop my own.

\subsection{Feature importance}

\subsubsection{LIME and SP-LIME}
LIME and SP-LIME are model-agnostic methods that explain chosen actions in specific states by slightly modifying input features and fitting an interpretable linear surrogate model that approximates the behaviour of the model locally. SP-LIME selects a representative set of these explanations to create a global explanation for how the model selects an action. These methods explain what features are important when deciding on an action in a given state.

\subsubsection{Shapley Additive exPlanations}
Shapley values is a method of computing marginal contributions from cooperative game theory. It assigns a Shapley value to each feature which is the mean of that features marginal contribution over all combinations of features to establish a fair attribution of credit to each player, or feature in the case of machine learning. This method is very computationally expensive and as such, approximations of this method are used instead.

\subsubsection{Layer-wise Relevancy Propagation}
LRP is a method to assign importance to each node in the network by working backwards from the output to the input. It looks at what nodes in the last hidden layer influenced the output and redistributes prediction scores to these nodes, it then looks at what nodes in the layer before influenced the nodes in the last hidden layer, until you have a complete map of how the network nodes affected each other to produce the output. This method helps reveal hidden patterns in model behaviour.

\subsection{Policy explanations}

\subsubsection{HIGHLIGHTS}
This is a intrinsically explainable \ac{rl} algorithm and it works by extracting sub-trajectories with important states, states where a wrong action can lead to significant decreases in reward, to explain how the agent acts in these states. It also provides context in the form of neighbouring states in the important sub-trajectories, in order for users to understand why it chose the trajectory it did.

\subsection{Objective explanations}


\subsection{Outcome explanation}


\section{Relevant Methods}
There are several papers written on \ac{xai} and \ac{xrl} problems. Milani et al.\cite{milani2022survey} categorise \ac{xrl} into three main categories, feature importance (FI), learning process and \ac{mdp} (LPM) and Policy-level (PL). FI considers the immediate context for actions, i.e. what part of the input was important for a single action, LPM considers techniques that highlight which experiences in training were influential for the final policy and PL focuses on long term effects and behaviour. Since we are interested in future states and actions we will look at influential trajectories and transitions within these trajectories. It is important to view these transitions in the context of the trajectory to understand the long term effects and not just immediate, which are of less interest in this paper, if we find the state with a high state importance $I(s)$, $I(s) = max_aQ(s,a)-min_aQ(s,a)$ most similar by some similarity measure to an arbitrary current state we could find the resulting trajectory and expect the agent to intend a similar outcome. There are also ways to convert \ac{rnn} policies to an interpretable models post-hoc, which might be relevant if we use an \ac{rnn}. This paper will explore PL explainability further.

In particular "What did you think would happen? Explaining Agent Behaviour through Intended Outcomes" \cite{yau2020did}, "Explainable Reinforcement Learning via a Causal world model" \cite{yu2024explainable} and "CrystalBox: Future-based explanations for input-driven deep \ac{rl} systems" \cite{patel2024crystalbox} are highly relevant due to the fact that they all describe temporal connections between current actions and future states or actions.

"Are large language models post-hoc explainers" \cite{kroeger2024large} Could be relevant as using other explainers to compare explanations is useful. "ACTER: Diverse and Actionable Counterfactual Sequences for Explaining and Diagnosing \ac{rl} Policies" \cite{gajcin2024acter} 

\subsection{What did you think would happen? Explaining Agent Behaviour through Intended Outcomes}
What did you think would happen describes what an agent expects to happen in future states, and why the current action is chosen based on the future expectations. As stated in the paper, a limitation of their method means it doesn't work well with high dimensionality. The two main difference between this paper and the problem i aim to solve is that they're focusing on an environment with a single agent and that our observation space will be high dimensionality. It uses Q-learning and Markov-chains that train simultaneously with a "belief map" that shows what the agent expects the environment to look like in future states. In the simple examples used in the paper it shows where it believes the taxi should drive and therefore chooses an action to follow this path. This is not directly applicable to my thesis as it's unlikely that Q-learning or Markov chains will be viable for the policy and explainer based on the environment. However the paper is successful in explaining an agents underlying motivations and beliefs about the causal nature of the environment, and using similar methods might be an effective means for making \ac{marl} agents with higher dimensionality understandable from a human perspective. %See Figure \ref{fig:taxi}

\subsection{Explainable Reinforcement Learning via a Causal world model}
Explainable Reinforcement Learning via a Causal world model constructs a sparse model to connect causal relationships, without prior knowledge of the causal relationship, rather than a fully connected one, but they still achieve high accuracy and results comparable to other fully connected \ac{mbrl} policies. Which is important, as there is often a trade off between explainability and performance. Using the same model for both explanations and choosing actions also make the explanations faithful to the intentions of the agent. The paper also describes a novel approach to derive causal chains from the causal influence of the actions, which lets us find the intent of the agent. The paper is successful being applied to \ac{mbrl}, and is also applicable to models with infinite actions spaces, which is a limitation of some other models, see previous sub section.

A limitation of the paper is that it requires a known factorisation of the environment, denoted by $\langle S, A, O, R, P, T, \gamma \rangle$, where S is state space, A is action space, O is observation space, R is reward function, P is probability function for the probability of transitioning from state $s$ to state $s'$ given action $a$, T is the termination condition given the transition tuple $(s,a,o,s')$, and $\gamma$ is the discount factor. Considering we will be working with hand crafted simulations we will have access to all of these, however its not certain that if we depend on this method that our contributions will be applicable to certain other environments where the factorisation is not known.


\subsection{CrystalBox: Future-Based Explanations for Input-Driven Deep Reinforcement Learning Systems}
CrystalBox introduces a model-agnostic, post-hoc, future based explanation for \ac{drl}. It doesn't require altering the controller, and works by decomposing the future rewards into its individual components. While this isn't exactly the kind of explanation we are looking for, it could be a great tool in developing an explainer that considers other agents actions in a multi agent cooperative environment, which is the goal of our paper, because it is post-hoc, and easily deployable. Especially because it was constructed for input-driven environments. The original paper claims it offers high fidelity solutions for both discrete and continuous action spaces. KAZ has a discrete action space but we might do some work with other environments as well.

It's not certain to be useful because it works by decomposing the reward function, and it's not safe to assume the reward function will even be useful to decompose.


\subsection{Are Large Language Models Post Hoc Explainers}
A \ac{llm} is a predictive model that generates text based on a prompt you give it, be it continuing the prompt or responding, and can often give the impression of comprehension of human language and a deeper understanding of the topic at hand.  The paper aims to investigate the question "Can LLMs explain the behaviour of other complex predictive models?" by exploiting the in-context learning (ICL) capabilities of \acp{llm}. ICL allows \acp{llm} to perform well on new tasks by using a few task samples in the prompt. A benefit of using an \ac{llm} as a post-hoc explainer, is that the output given by the model will already be written in natural language and should be understandable by a layman. The paper presumes that the local behaviour of a model is a linear decision boundary, and by using a sample x and perturbing it to x', and presenting both the samples and the perturbation as natural language we could get an explanation from the \ac{llm} for the outcome. With a sufficient number of perturbations in a neighbourhood around x the \ac{llm} is expected to explain the behaviour in this neighbourhood, and rank the features in order of importance.

While the faithfulness of the \ac{llm} as an explainer is on par with other \ac{xai} methods used for classification, meaning that the reasons provided are enough to explain the output, I am sceptical of the fidelity of the \ac{llm} for two reasons. One is the same as for why other \ac{xai} models often struggle with fidelity, the temporal aspect. If applied in the same way as in the paper it would not consider the intent or the past and only what part of the current observation made it make a certain decision. The other is that I am sceptical of claims presented by an \ac{llm} in general as these are all just guesses. Good guesses a lot of the time, but still just guesses. 

We could however potentially change the implementation so it considers the temporal aspect, and this might be a viable post hoc explainer after some more research into prompt engineering.


\subsection{ACTER: Diverse and Actionable Counterfactual Sequences for Explaining and Diagnosing Reinforcement Learning Policies}
\label{sec:acter}
The paper presents ACTER, an algorithm that uses counterfactual sequences with actionable advice on how to avoid failure for an \ac{rl} policy. It does this by using an \ac{ea} known as \ac{nsga} to generate counterfactual sequences that don't lead to failure as close as possible to factual sequences that lead to failure. This paper presents counterfactual sequences and not just actions, which means it also presents how to avoid the state that lead to failure to begin with, which should, if ACTER is implemented correctly, allow us to significantly reduce the amount of times our policy fails. It also offers multiple counterfactuals to allow the end user to decide which counterfactual is preferable to their use case.

There are 4 hypothesises tested by the paper. The last two considers laymen users and are therefore not as interesting to us. The first two however "ACTER can produce counterfactual sequences that prevent failure with lower effort and higher certainty in stochastic environments compared to the baselines." and "ACTER can produce a set of counterfactual sequences that offer more diverse ways of preventing failure compared to the baselines." are partially and fully confirmed respectively. Which means ACTER will likely be a useful tool to explain and diagnose our \ac{rl} policy.


\section{Related Works}
This section describes recent development in the field of planning and future predictions. All the papers described here are directly relevant to the methods i developed.

\subsection{Predicting Future Actions Of Reinforcement Learning Agents}
\label{sec:predicting_actions}
This paper describes a method developed by the author that predicts future actions and events \cite{chung2024predictingfutureactionsreinforcement}. It does this for non planning policies, like Impala, implicit planning policies like \ac{drc}, see section \ref{sec:impl_planning}, and explicit planning policies like MuZero and Thinker.

Non planning policies are policies designed without planning in mind. Pure \ac{nn} policies without recurrency are considered to be non planning, as they have not been found to exhibit planning like behaviour. This paper found that when predicting future states and actions both inner states from implicit planners, and explicit planners improved accuracy over non planning agents.

Explicit planners are agents who simulate future states with a world model before making a decision on what action to take. They rely on the world model being accurate, and so does using them to predict future states and actions. In some environments learning an accurate world model is not feasible, an example of this is autonomous driving.

The paper proposes two methods on predicting future states and actions, simulation based and inner state based. The inner state approach is most effective when working with explicit planners, but also shows significant improvement for implicit planner agents. The simulation based approach performs very well on implicit planner agents, but requires the opportunity to train a decently performing world model. The paper also shows with an ablated world model, the inner state approach performs better.

\subsection{An Investigation of Model Free Planning}
\label{sec:impl_planning}
This paper investigates whether planning like behaviour can emerge from model free \ac{rl}, without any explicit planning methods or inductive biases designed to induce planning behaviour \cite{guez2019investigationmodelfreeplanning}. The authors introduce and explore if the \ac{drc} architecture can learn to implicitly plan through training. World models suffer from scalability issues, and inductive biases require a priori knowledge of the environment. The \ac{drc} architecture is comprised of $3$ layers of ConvLSTMs and each layer is iterated through $3$ times each at each timestep.

This architecture is tested on domains that require planning such as Sokoban, and they test the agents ability to generalize as well as data efficiency. The results of the paper is that not only does the model exhibit planning like behaviour, but it outperforms state of the art methods that use inductive biases or world models. The paper described in section \ref{sec:predicting_actions} finds that using the inner state of implicit planner agents improves predicted future states and actions.


\medskip
\chapter{Methodology}

\section{Critical states}
In a safety critical real world scenario, like autonomous driving, its important to identify when a vehicle is about to make an important decision, a decision where a mistake could lead to important consequences. If we can identify these states we could make efforts to increase the likelihood of correct action in these states, either by taking over as a human, or improving the policy in these scenarios.

\subsection{Maximum logit difference}
The output of the policy, in the case of discrete action space, is a vector with size $n$, where $n$ is the amount of possible actions, without an activation function. This is also known as the logit vector. A state in which there is a high difference between the highest logit and the lowest logit means the agent has learned to avoid the action associated with the low logit and prefers the action with the high logit, which would mean the agent associates the high logit action with high future rewards, and the low logit action with low future rewards. If the logits are close to each other it means the agent is either uncertain about which action to choose or indifferent to the different actions in this state. Given that the policy has been trained well enough that it correctly predicts the optimal action in a state, I consider this state to be critical if the difference between the highest logit and the lowest logit is higher than some chosen threshold.
To compute whether a trajectory contains a critical state I consider the maximum logit difference in each of the states, and assign the maximum of these differences to the trajectory. After computing a value for every trajectory, we consider the median of the values assigned to the trajectories as the threshold for whether a state contains a critical state or not. This means we consider half the trajectories to contain a critical state. The main reason for this is it makes using a baseline of random guesses have expected $50\%$ accuracy. This means if we train a binary classifier with initial observation as input and whether the trajectory contains a critical state as output any improvement over $50\%$ means a network has learned a meaningful connection between the input and the output, as opposed to simply learning that one class is more common than another.

\subsection{Counterfactual sequences}
\label{sec:counterfactual}
In the case of continuous action spaces we cannot look at maximum logit difference, as a policy network doesn't output logits, but rather the actions to choose. To discover counterfactual sequences we use the \ac{ea} known as \ac{nsga}. This algorithm is a significant improvement over earlier multi-objective EAs that use non dominated sorting. It has a complexity of $O(MN^2)$ instead of $O(MN^3)$, elitist approach and a specified sharing parameter, all three of which, many earlier algorithms lacked. \cite{Deb2001AFA} Pseudocode for sorting in algorithm \ref{alg:fnds}. 
We use multi objective search because we want to minimize action change while maximizing total reward change throughout a simulation. This defines our two objectives. This way we can discover what actions or combinations there of have the highest impact on the environment. When working with a discrete action space, that is when we have a finite, usually relatively small, amount of actions to choose from, we define the action objective that we want to minimize as how many actions in a single timestep are different from a predefined sequence found with a trained model. After the timestep with the found actions we again use the model to roll out the rest of the sequence. We compare the altered sequence to the original sequence and measure the total reward both of these sequences receive. We want to maximize this difference. 

\begin{algorithm}
\caption{Fast Non-Dominated Sort}
\label{alg:fnds}
\begin{algorithmic}
    \State $F_i \gets \emptyset$
    \ForAll{$p \in P$}
    \State $S_p \gets \emptyset$
    \State $n_p \gets 0$
        \ForAll{$q \in P$}
            \If{$p \prec q$}
                \State $S_p \gets S_p \cup q$
            \ElsIf{$q \prec p$}
                \State $n_p \gets n_p + 1$
            \EndIf
        \EndFor
        \If{$n_p == 0$}
            \State $F_1 \gets F_1 \cup p$
            \State $p_{rank} \gets 1$
        \EndIf
    \EndFor
    \State $i \gets 1$
    \While{$F_i \neq \emptyset$}
        \State $Q \gets \emptyset$
        \ForAll{$p \in F_i$} 
            \ForAll{$q \in S_p$}
                \State $n_q \gets n_q - 1$
                \If{$n_q \gets 0$}
                    \State $q_{rank} \gets i + 1$
                    \State $Q \gets Q \cup q$
                \EndIf
            \EndFor
        \EndFor
        \State $i \gets i + 1$
        \State $F_i \gets Q$
    \EndWhile
\end{algorithmic}
\end{algorithm}


\section{Feature importance}

\subsection{Shapley values for observations}
The Shapley value is an idea from cooperative game theory where it was used to measure the contribution of each player to the total outcome. It has been adapted to deep reinforcement learning where each feature is considered a player and the policy is considered a game. In our case we initially want to use it to explain actions so we consider that to be the outcome of the game. The Shapley value is calculated by taking each feature and finding the marginal contribution, i.e. seeing how the prediction changes when you include the feature and compare the outcome to when it wasn't included, you do this over all possible coalitions of features. See algorithm \ref{alg:shapley}.
The Shapley value has a computational complexity of $O(2^n)$, and because of this we often use estimates instead of exact Shapley values. There are two main ways of reducing the complexity. The first one is using a surrogate model instead of the policy to reduce the number of features, and the second is reducing the number of coalitions. In our case we reduce the number of coalitions to some number $2^K$ and use K-means clustering to increase representability and reduce variance.
Since we have forward passes of the policy, and each can be considered one example of a game we sample the games and use the average marginal contribution over all the games instead of just one game.
The Shapley value can help us understand how the policy acts in general. We can also use it to get explanations for single observations, however these explanations won't always have high fidelity. We will explore more of how the fidelity of Shapley values for single explanations compare to high fidelity explainability methods like gradient methods. See section \ref{sec:intgrad}

\begin{algorithm}
\caption{Shapley Value Calculation with Baseline Replacement}
\label{alg:shapley}
\begin{algorithmic}
    \State \textbf{Input:} Model $f$, input sample $x$, baseline $x'$, set of features $F$
    \State \textbf{Output:} Shapley values $\phi_f$ for each feature $f \in F$
    \State $\phi_f \gets 0 \; \forall f \in F$ \Comment{Initialize Shapley values}
    \State $n \gets |F|$ \Comment{Number of features}
    
    \ForAll{$f \in F$}
        \ForAll{$S \subseteq F \setminus \{f\}$}
            \State $x_S \gets \texttt{construct}(x, x', S)$ \Comment{Replace features not in $S$ with baseline values}
            \State $x_{S \cup \{f\}} \gets \texttt{construct}(x, x', S \cup \{f\})$
            \State $\Delta \gets f(x_{S \cup \{f\}}) - f(x_S)$ \Comment{Marginal contribution of feature $f$}
            \State $\phi_f \gets \phi_f + \frac{|S|!(n - |S| - 1)!}{n!} \Delta$
        \EndFor
    \EndFor
    \State \Return $\{\phi_f : f \in F\}$
    
    \Function{construct}{$x, x', S$}
        \State \textbf{Input:} Original sample $x$, baseline $x'$, feature subset $S$
        \State \textbf{Output:} Modified sample $x_S$
        \ForAll{$f \in F$}
            \If{$f \in S$}
                \State $x_S[f] \gets x[f]$
            \Else
                \State $x_S[f] \gets x'[f]$ \Comment{Replace with baseline value}
            \EndIf
        \EndFor
        \State \Return $x_S$
    \EndFunction
\end{algorithmic}
\end{algorithm}

\subsection{Integrated Gradients}
\label{sec:intgrad}
By integrating the gradient of the model prediction when going from a specified baseline value as input to our specific observation x we get a high fidelity explanation for what features were important in that specific observation. We integrate from the baseline instead of just getting the gradient in our observation because of the saturation problem. If a feature is important to the action the gradient could be small for the value x, but the integral of the gradient from the baseline to x would still be large. Conversely a gradient for the observation x could be significant in a small area around a feature in x even if that feature isn't as significant as the gradient in x would lead us to believe \cite{sundararajan2017axiomaticattributiondeepnetworks}. See algorithm \ref{alg:intgrad}. Comparing figure \ref{fig:intgrad} with figure \ref{fig:shap} we can see that the values are similar. This means that in the case of this observation, the fidelity of the Shapley explanation is high.

\begin{algorithm}
\caption{Integrated Gradients for Feature Importance in Reinforcement Learning}
\label{alg:intgrad}
\begin{algorithmic}
    \State \textbf{Input:} \ac{rl} policy $\pi_\theta$, input $x$, baseline input $\bar{x}$, number of steps $m$
    \State \textbf{Output:} Integrated gradients $\texttt{IG}(x, \bar{x})$
    \State $\texttt{IG}(x, \bar{x}) \gets 0 \; \forall f \in \texttt{features of } x$ \Comment{Initialize integrated gradients to zero}
    
    \For{$i = 1$ to $m$}
        \State $\alpha_i \gets \frac{i}{m}$
        \State $x_i \gets \bar{x} + \alpha_i \times (x - \bar{x})$
        \State $grad_i \gets \nabla_x \texttt{value}(\pi_\theta, x_i)$ \Comment{Compute the gradient of the value function}
        \State $\texttt{IG}(x, \bar{x}) \gets \texttt{IG}(x, \bar{x}) + grad_i \times \frac{(x - \bar{x})}{m}$
        \Comment{Accumulate gradients scaled by the step size}
    \EndFor
    
    \Function{value}{$\pi_\theta, x$}
        \State \textbf{Input:} \ac{rl} policy $\pi_\theta$, input $x$
        \State \textbf{Output:} Model prediction for input $x$
        
        \State $\texttt{Feed input } x \texttt{ to policy } \pi_\theta$
        \State \Return \texttt{model output}
    \EndFunction
\end{algorithmic}
\end{algorithm}


\subsection{Choosing a feature attribution baseline}
Both the Shapley value attribution and the Integrated Gradients attribution I use require a baseline. A baseline should represents a neutral, or absence of a value, and needs to be from within the domain that a model is trained on, as a models could behave unpredictably on data outside of their domain, In my case, the environments we used has observation spaces with values in $[0,1]$, meaning that for our baseline we need to choose values in this range. There are several ways to choose a baseline, and each baseline will yield different results. If the domain is a set of images, a common baseline to use is black, or $(0,0,0)$ if the images are represented with RGB values\cite{baseline_bird}. Intuitively, this can be understood as \acp{nn} looking for patterns in an image, which disappear if the image is monochrome.
In my case, I do not have images, but tabular data. Using $0$ as baseline values is likely not a good choice, as each value when using tabular data represents an interpretable value and $0$ represents a value the model would process as extreme, instead of neutral. Another possibility is choosing the dataset mean as a baseline, which has been used successfully in other \ac{rl} applications, although not \ac{rl} applications with tabular data that I am aware of \cite{baseline_rl}.


\section{Extracting intent}
In an environment where the agents successfully reaches a desirable state we can assume that at some point in the trajectory an agents intent is to reach this state. We can train a prediction model on such an environment that uses a certain state or observation as input and outputs features of interest some time in the future. In the example of autonomous vehicles the output could be predicted position of the vehicle at a future point in time. If our prediction model is successful, then it could be considered part of an explanation for the intent of an agent.
Additionally, we train different models with different sets of inputs to identify what is important to identify intent. We train one model that takes just the observation as input, one that takes observation and action, and one that takes observation and action one hot encoded. While action as an integer and action one hot encoded contains the same amount of information it could be easier for a network to learn correlations with the additional inputs. Looking at the loss curve during training, as well as final loss, this suspicion is confirmed to be true when dealing with a network architecture and environment like we are currently using, 2 fully connected hidden layers with the hyperbolic tangent as the activation, in the environment simple\_spread\_v3. While the input layer for the three prediction models are slightly different the difference in training time is negligible. 
After 200 epochs the base prediction model reached a evaluation loss of 0.0407, the model with integer action as part of the input reached a evaluation loss of 0.0359 and the model with one-hot encoded action as part of the input reached an evaluation loss of 0.0322. See figure \ref{fig:pred_loss_none}, \ref{fig:pred_loss_action} and \ref{fig:pred_loss_one-hot}. Both the base model and the one-hot action model stagnate long before 200 epochs is reached. If we let the model with integer action as input train for another 50 epochs it does not improve. The test loss for these three models respectively is ... . These graphs and values indicate that while action and an explanation are useful for extracting intent, in the context of our environment and model one-hot values and integrated gradient explanation are the most useful. The difference between using Shapley values and

\section{Temporal information}

\subsection{LSTM values}
An \ac{lstm} cell takes as input a cell state $c_{t-1}$, a hidden state $h_{t-1}$, and an input $X_t$ for each timestep. It outputs a cell state $c_{t}$, a hidden state $h_{t}$ and an output $y_t$. $h_{t-1}$ is concatenated with $X_t$ and then with trainable weight matrices $\textbf{W}$ and biases $b$, $c_t$ and $h_t$ are computed according to the following equations:

\[
f_t = \sigma(W_f [h_{t-1}, x_t] + b_f)
\]

\[
i_t = \sigma(W_i [h_{t-1}, x_t] + b_i)
\]

\[
\tilde{c}_t = \tanh(W_c [h_{t-1}, x_t] + b_c)
\]

\[
c_t = f_t \odot c_{t-1} + i_t \odot \tilde{c}_t
\]

\[
o_t = \sigma(W_o [h_{t-1}, x_t] + b_o)
\]

\[
h_t = o_t \odot \tanh(c_t)
\]

Intuitively, $f_t$ can be viewed as how much of the previous cell state $c_{t-1}$ to forget, $i_t$ can be viewed as how much of the candidate cell state $\tilde{c}_t$ to remember, and $o_t$ can be viewed as how much of the $c_{t}$ to output. Capital letters represent matrices and lowercase letters represent vectors.

\begin{center}
\begin{figure}
\begin{tikzpicture}[>=Stealth, node distance=1.5cm, auto]
\hspace{4cm}
  % LSTM cell as a pastel purple box with rounded corners
  \node[draw, fill=purple!30, rounded corners, minimum width=3cm, minimum height=2cm] (lstm) {LSTM};

  % Incoming arrows on the left: previous cell state (c) and hidden state (h)
  \draw[<-] ([yshift=0.3cm]lstm.west) -- ++(-1,0) node[left] {\(c_{t-1}\)};
  \draw[<-] ([yshift=-0.3cm]lstm.west) -- ++(-1,0) node[left] {\(h_{t-1}\)};

  % Outgoing arrows on the right: new cell state (c) and hidden state (h)
  \draw[->] ([yshift=0.3cm]lstm.east) -- ++(1,0) node[right] {\(c_t\)};
  \draw[->] ([yshift=-0.3cm]lstm.east) -- ++(1,0) node[right] {\(h_t\)};

  \draw[<-] (lstm.south) -- ++(0,-1) node[below] {\(x_t\)};
  \draw[->] (lstm.north) -- ++(0,1) node[above] {\(h_t\)};
\end{tikzpicture}
\caption{Simplified illustration of an \ac{lstm} layer}
\label{fig:lstm}
\end{figure}
\end{center}

As $c_t$ and $h_t$ are helpful when predicting future actions and states \cite{chung2024predictingfutureactionsreinforcement} when considering the \ac{drc} architecture, I hypothesize that the same holds when using a regular \ac{lstm} layer. See figure \ref{fig:lstm}.

\subsection{Frame Stacking}
Frame stacking in \ac{rl} is a technique used to provide an agent with temporal context when tasked with choosing an action by combining the previous $n$ observations into a single input to the policy. In some environments an observation at a single timestep isn't enough to capture the full state. In the case of autonomous vehicles a single observation might not have information on whether a car is standing still or in motion, and the velocity of a car is important to take into account when making an action.
If we have an observation space of 18, and we want to stack the previous 4 frames to feed into a fully connected input layer to a \ac{nn} we would insted of having the input $o_t$, where $o_t$ is a vector of size 18, we would instead get $$s_t = [o_t, o_{t-1}, o_{t-2}, o_{t-3}]$$ which is a $18\times4$ matrix, which we flatten into a single vector of size 72. If we instead of having a fully connected input layer, have a \ac{cnn} or another architecture that takes spatial context into consideration when processing an input, we would not flatten it.
Frame stacking does not give the \ac{nn} architecture memory, however it does provide the network with short term temporal context. I hypothesize that for event and state prediction for an agent, having information about previous observations and their transitions is helpful for predicting how the state will change in the future. In the case of autonomous vehicles its natural to assume that knowing the speed of which a vehicle has been moving will be relevant when trying to predict where it will be $n$ steps into the future.

\section{Statistical tests}
The null hypothesis chosen, $H_0$, is that the method provides no change in expected performance. To confirm that the difference in performance between baseline and our method is statistically significant I performed an independent one-tailed t-test, i.e. I calculated the likelihood that the increase in observed performance, or a higher increase in performance, happened by chance, and not because the true mean of the performance of the developed method is increased from baseline.
We bootstrapped the full dataset several times, i.e. created new datasets by sampling with replacement from the original datasets, and trained my prediction \acp{nn} 
on every dataset. From this we get a distribution of performances that we can use to perform a t-test and get a $p$-value. If this $p$-value is lower than a specified threshold value $\alpha$ I consider the result to be statistically significant. I set $\alpha = 0.05$.


\chapter{Experiments}
\section{General experimental setup}
I have several \acp{nn} for each experiment. The \acp{nn} have two hidden layers of size 64 each, the activation function $tanh$, and outputs either part of the state, $(x,y)$ coordinates for an agent for state prediction, or a single value, whether the trajectory contains a critical state or not, for event prediction. The experiments attempt to test whether my developed methods have an effect on event and state predictions, how impactful these effects are and what contributes to whether the effect is significant or not.I trained them with a scheduler which would divide learning rate by $3$ from an original of $1e-4$, down to a minimum of $1e-6$ whenever the calculated loss on the validation set stopped decreasing.
\subsection{State prediction neural nets and datasets}
For the state prediction I initially trained a \ac{nn}, $state_{base}$ to take the observations made by a type of agent from step $n$, $Obs_n$, and output the $x$ and $y$ coordinates of the same agent at step $n+m$ in a trajectory. The dataset $D_{pos}$ was constructed such that each data point was from a different trajectory $\tau$, to make sure all the data points were independent of each other, in order to make sure the \ac{nn} did not learn any such connection. Each data point is a pairing of values $(Obs_n, (x_{n+m},y_{n+m}))$. The variable $n$ was chosen randomly for each trajectory while $m$ stayed as a constant, in my case $m=10$. I would for example have $50000$ trajectories, each trajectory contains $2$ archers, so we would have a dataset of size $100000$. $100000$ observations and their corresponding positions $10$ steps into the future. We assume that both archers use the same policy as they are homogeneous.

\subsection{Event prediction neural nets and datasets}
I also trained a \ac{nn} for event prediction, $event_{base}$. The event we trained a network to predict was whether an agent would encounter a critical state or not. See section placeholder. The dataset $D_{crit}$ was similar to $D_{pos}$, however instead of positions, it would contain whether the agent would encounter a critical state within the next $5$ steps, $(Obs_n, crit)$, where $crit$ is a binary value, 0, for non critical states, or 1, for critical states.
I consider the accuracy of $event_{base}$ and the average error measured in euclidean distance of $state_{base}$ to be the baselines i compare results against in the first experiment. I will train other neural networks.


\subsection{Baseline performance}
$event_{base}$ got an accuracy of $0.8244$ on the test set and for $state_{base}$ the average distance from target was $0.0409$ on the test set. These \acp{nn} are purely post hoc explainers and used no information about the policy to explain the agents intention. As aforementioned we want to include information about the policy to discover how we can improve this prediction, and better discover the intent of the agent.

 Only The dataset, and therefore also the input layer, was different between the runs. I trained the networks for a total of $200$ epochs, after which improvement was negligible, if at all existent, with a learning rate scheduler that decreased learning rate when loss reached a plateau, which was measured on a validation set.

\section{Action and explanation inclusion}
\subsection{Setup}
I enhanced both the dataset $D_{pos}$ and $D_{crit}$ with a series of other data. For both datasets I included the action, encoded by an integer from $0$ to $n_{acts}$, i.e. the amount of actions an agent has available, taken at $Obs_n$ for each data point, to use as input for an \ac{nn}. I then trained two \acp{nn} with the same architecture as $state_{base}$ and $event_{base}$ respectively, and once again measured accuracy for the event predictor net with included action, and average error for the state predictor net with included action. If the predictions improved I considered the action to be useful on its own for the predictions in the given environment.

\subsubsection{One-hot and integer inclusion}
I then one-hot encoded the action, i.e. encoded the action as a $n_{acts}$ dimensional unit vector with a 1 in the chosen action dimension. This includes the same amount of information as represented as an integer, however I hypothesized that the resulting network would achieve higher performance with one-hot encoded actions as the distance between independent actions when one-hot encoding the actions are all the same, as opposed to when encoding the actions as an integer when the distance between two actions is not constant. For example, between the action encoded as the integer 2 and the action encoded as the integer 3 the Manhattan distance is 1, but between the action encoded as the integer 2 and the action encoded as the integer 4 the Manhattan distance is 2. When one-hot encoding these actions the Manhattan distance between two distinct actions will in all cases be 2. As in the previous cases I also trained a \ac{nn} on the datasets enhanced with one hot encoded actions.

\subsubsection{Integrated Gradients and Shapley Value inclusion}
After creating datasets with actions encoded different ways I also created two feature importance explanations for the action chosen in each of the initial observations. Shapley values and Integrated gradients were the two methods chosen. I decided on 2 feature importance measures that internally uses different methods of calculation, gradients and SHAP, to explore if they would yield different results. The motivation behind including feature importance explanations is that it describes some if the inner processing of the policy, and might therefore elucidate what kind of actions it might take in future states, and therefore include some of the intent. On each of the previous datasets I created new datasets with these two feature importance methods included.\\

\subsubsection{Result presentation}
In total I trained 18 neural networks per environment, one for each dataset constructed, and compared the results for accuracy on the event prediction networks and the results for average error on the state prediction networks. I repeated this experiment for a few different environments to explore if the results vary depending on environment.

All results are presented in tables, with asterisks marking significance, and italics used for baseline value.

\begin{figure}[!ht]
	%% Creator: Matplotlib, PGF backend
%%
%% To include the figure in your LaTeX document, write
%%   \input{<filename>.pgf}
%%
%% Make sure the required packages are loaded in your preamble
%%   \usepackage{pgf}
%%
%% Also ensure that all the required font packages are loaded; for instance,
%% the lmodern package is sometimes necessary when using math font.
%%   \usepackage{lmodern}
%%
%% Figures using additional raster images can only be included by \input if
%% they are in the same directory as the main LaTeX file. For loading figures
%% from other directories you can use the `import` package
%%   \usepackage{import}
%%
%% and then include the figures with
%%   \import{<path to file>}{<filename>.pgf}
%%
%% Matplotlib used the following preamble
%%   \def\mathdefault#1{#1}
%%   \everymath=\expandafter{\the\everymath\displaystyle}
%%   \IfFileExists{scrextend.sty}{
%%     \usepackage[fontsize=10.000000pt]{scrextend}
%%   }{
%%     \renewcommand{\normalsize}{\fontsize{10.000000}{12.000000}\selectfont}
%%     \normalsize
%%   }
%%   
%%   \ifdefined\pdftexversion\else  % non-pdftex case.
%%     \usepackage{fontspec}
%%     \setmainfont{DejaVuSerif.ttf}[Path=\detokenize{/uio/hume/student-u79/adaha/master_thesis/.venv/lib/python3.11/site-packages/matplotlib/mpl-data/fonts/ttf/}]
%%     \setsansfont{DejaVuSans.ttf}[Path=\detokenize{/uio/hume/student-u79/adaha/master_thesis/.venv/lib/python3.11/site-packages/matplotlib/mpl-data/fonts/ttf/}]
%%     \setmonofont{DejaVuSansMono.ttf}[Path=\detokenize{/uio/hume/student-u79/adaha/master_thesis/.venv/lib/python3.11/site-packages/matplotlib/mpl-data/fonts/ttf/}]
%%   \fi
%%   \makeatletter\@ifpackageloaded{underscore}{}{\usepackage[strings]{underscore}}\makeatother
%%
\begingroup%
\makeatletter%
\begin{pgfpicture}%
\pgfpathrectangle{\pgfpointorigin}{\pgfqpoint{6.400000in}{4.800000in}}%
\pgfusepath{use as bounding box, clip}%
\begin{pgfscope}%
\pgfsetbuttcap%
\pgfsetmiterjoin%
\definecolor{currentfill}{rgb}{1.000000,1.000000,1.000000}%
\pgfsetfillcolor{currentfill}%
\pgfsetlinewidth{0.000000pt}%
\definecolor{currentstroke}{rgb}{1.000000,1.000000,1.000000}%
\pgfsetstrokecolor{currentstroke}%
\pgfsetdash{}{0pt}%
\pgfpathmoveto{\pgfqpoint{0.000000in}{0.000000in}}%
\pgfpathlineto{\pgfqpoint{6.400000in}{0.000000in}}%
\pgfpathlineto{\pgfqpoint{6.400000in}{4.800000in}}%
\pgfpathlineto{\pgfqpoint{0.000000in}{4.800000in}}%
\pgfpathlineto{\pgfqpoint{0.000000in}{0.000000in}}%
\pgfpathclose%
\pgfusepath{fill}%
\end{pgfscope}%
\begin{pgfscope}%
\pgfsetbuttcap%
\pgfsetmiterjoin%
\definecolor{currentfill}{rgb}{1.000000,1.000000,1.000000}%
\pgfsetfillcolor{currentfill}%
\pgfsetlinewidth{0.000000pt}%
\definecolor{currentstroke}{rgb}{0.000000,0.000000,0.000000}%
\pgfsetstrokecolor{currentstroke}%
\pgfsetstrokeopacity{0.000000}%
\pgfsetdash{}{0pt}%
\pgfpathmoveto{\pgfqpoint{0.800000in}{0.528000in}}%
\pgfpathlineto{\pgfqpoint{5.760000in}{0.528000in}}%
\pgfpathlineto{\pgfqpoint{5.760000in}{4.224000in}}%
\pgfpathlineto{\pgfqpoint{0.800000in}{4.224000in}}%
\pgfpathlineto{\pgfqpoint{0.800000in}{0.528000in}}%
\pgfpathclose%
\pgfusepath{fill}%
\end{pgfscope}%
\begin{pgfscope}%
\pgfsetbuttcap%
\pgfsetroundjoin%
\definecolor{currentfill}{rgb}{0.000000,0.000000,0.000000}%
\pgfsetfillcolor{currentfill}%
\pgfsetlinewidth{0.803000pt}%
\definecolor{currentstroke}{rgb}{0.000000,0.000000,0.000000}%
\pgfsetstrokecolor{currentstroke}%
\pgfsetdash{}{0pt}%
\pgfsys@defobject{currentmarker}{\pgfqpoint{0.000000in}{-0.048611in}}{\pgfqpoint{0.000000in}{0.000000in}}{%
\pgfpathmoveto{\pgfqpoint{0.000000in}{0.000000in}}%
\pgfpathlineto{\pgfqpoint{0.000000in}{-0.048611in}}%
\pgfusepath{stroke,fill}%
}%
\begin{pgfscope}%
\pgfsys@transformshift{1.002796in}{0.528000in}%
\pgfsys@useobject{currentmarker}{}%
\end{pgfscope}%
\end{pgfscope}%
\begin{pgfscope}%
\definecolor{textcolor}{rgb}{0.000000,0.000000,0.000000}%
\pgfsetstrokecolor{textcolor}%
\pgfsetfillcolor{textcolor}%
\pgftext[x=1.002796in,y=0.430778in,,top]{\color{textcolor}{\sffamily\fontsize{10.000000}{12.000000}\selectfont\catcode`\^=\active\def^{\ifmmode\sp\else\^{}\fi}\catcode`\%=\active\def%{\%}0}}%
\end{pgfscope}%
\begin{pgfscope}%
\pgfsetbuttcap%
\pgfsetroundjoin%
\definecolor{currentfill}{rgb}{0.000000,0.000000,0.000000}%
\pgfsetfillcolor{currentfill}%
\pgfsetlinewidth{0.803000pt}%
\definecolor{currentstroke}{rgb}{0.000000,0.000000,0.000000}%
\pgfsetstrokecolor{currentstroke}%
\pgfsetdash{}{0pt}%
\pgfsys@defobject{currentmarker}{\pgfqpoint{0.000000in}{-0.048611in}}{\pgfqpoint{0.000000in}{0.000000in}}{%
\pgfpathmoveto{\pgfqpoint{0.000000in}{0.000000in}}%
\pgfpathlineto{\pgfqpoint{0.000000in}{-0.048611in}}%
\pgfusepath{stroke,fill}%
}%
\begin{pgfscope}%
\pgfsys@transformshift{1.569265in}{0.528000in}%
\pgfsys@useobject{currentmarker}{}%
\end{pgfscope}%
\end{pgfscope}%
\begin{pgfscope}%
\definecolor{textcolor}{rgb}{0.000000,0.000000,0.000000}%
\pgfsetstrokecolor{textcolor}%
\pgfsetfillcolor{textcolor}%
\pgftext[x=1.569265in,y=0.430778in,,top]{\color{textcolor}{\sffamily\fontsize{10.000000}{12.000000}\selectfont\catcode`\^=\active\def^{\ifmmode\sp\else\^{}\fi}\catcode`\%=\active\def%{\%}25}}%
\end{pgfscope}%
\begin{pgfscope}%
\pgfsetbuttcap%
\pgfsetroundjoin%
\definecolor{currentfill}{rgb}{0.000000,0.000000,0.000000}%
\pgfsetfillcolor{currentfill}%
\pgfsetlinewidth{0.803000pt}%
\definecolor{currentstroke}{rgb}{0.000000,0.000000,0.000000}%
\pgfsetstrokecolor{currentstroke}%
\pgfsetdash{}{0pt}%
\pgfsys@defobject{currentmarker}{\pgfqpoint{0.000000in}{-0.048611in}}{\pgfqpoint{0.000000in}{0.000000in}}{%
\pgfpathmoveto{\pgfqpoint{0.000000in}{0.000000in}}%
\pgfpathlineto{\pgfqpoint{0.000000in}{-0.048611in}}%
\pgfusepath{stroke,fill}%
}%
\begin{pgfscope}%
\pgfsys@transformshift{2.135733in}{0.528000in}%
\pgfsys@useobject{currentmarker}{}%
\end{pgfscope}%
\end{pgfscope}%
\begin{pgfscope}%
\definecolor{textcolor}{rgb}{0.000000,0.000000,0.000000}%
\pgfsetstrokecolor{textcolor}%
\pgfsetfillcolor{textcolor}%
\pgftext[x=2.135733in,y=0.430778in,,top]{\color{textcolor}{\sffamily\fontsize{10.000000}{12.000000}\selectfont\catcode`\^=\active\def^{\ifmmode\sp\else\^{}\fi}\catcode`\%=\active\def%{\%}50}}%
\end{pgfscope}%
\begin{pgfscope}%
\pgfsetbuttcap%
\pgfsetroundjoin%
\definecolor{currentfill}{rgb}{0.000000,0.000000,0.000000}%
\pgfsetfillcolor{currentfill}%
\pgfsetlinewidth{0.803000pt}%
\definecolor{currentstroke}{rgb}{0.000000,0.000000,0.000000}%
\pgfsetstrokecolor{currentstroke}%
\pgfsetdash{}{0pt}%
\pgfsys@defobject{currentmarker}{\pgfqpoint{0.000000in}{-0.048611in}}{\pgfqpoint{0.000000in}{0.000000in}}{%
\pgfpathmoveto{\pgfqpoint{0.000000in}{0.000000in}}%
\pgfpathlineto{\pgfqpoint{0.000000in}{-0.048611in}}%
\pgfusepath{stroke,fill}%
}%
\begin{pgfscope}%
\pgfsys@transformshift{2.702202in}{0.528000in}%
\pgfsys@useobject{currentmarker}{}%
\end{pgfscope}%
\end{pgfscope}%
\begin{pgfscope}%
\definecolor{textcolor}{rgb}{0.000000,0.000000,0.000000}%
\pgfsetstrokecolor{textcolor}%
\pgfsetfillcolor{textcolor}%
\pgftext[x=2.702202in,y=0.430778in,,top]{\color{textcolor}{\sffamily\fontsize{10.000000}{12.000000}\selectfont\catcode`\^=\active\def^{\ifmmode\sp\else\^{}\fi}\catcode`\%=\active\def%{\%}75}}%
\end{pgfscope}%
\begin{pgfscope}%
\pgfsetbuttcap%
\pgfsetroundjoin%
\definecolor{currentfill}{rgb}{0.000000,0.000000,0.000000}%
\pgfsetfillcolor{currentfill}%
\pgfsetlinewidth{0.803000pt}%
\definecolor{currentstroke}{rgb}{0.000000,0.000000,0.000000}%
\pgfsetstrokecolor{currentstroke}%
\pgfsetdash{}{0pt}%
\pgfsys@defobject{currentmarker}{\pgfqpoint{0.000000in}{-0.048611in}}{\pgfqpoint{0.000000in}{0.000000in}}{%
\pgfpathmoveto{\pgfqpoint{0.000000in}{0.000000in}}%
\pgfpathlineto{\pgfqpoint{0.000000in}{-0.048611in}}%
\pgfusepath{stroke,fill}%
}%
\begin{pgfscope}%
\pgfsys@transformshift{3.268671in}{0.528000in}%
\pgfsys@useobject{currentmarker}{}%
\end{pgfscope}%
\end{pgfscope}%
\begin{pgfscope}%
\definecolor{textcolor}{rgb}{0.000000,0.000000,0.000000}%
\pgfsetstrokecolor{textcolor}%
\pgfsetfillcolor{textcolor}%
\pgftext[x=3.268671in,y=0.430778in,,top]{\color{textcolor}{\sffamily\fontsize{10.000000}{12.000000}\selectfont\catcode`\^=\active\def^{\ifmmode\sp\else\^{}\fi}\catcode`\%=\active\def%{\%}100}}%
\end{pgfscope}%
\begin{pgfscope}%
\pgfsetbuttcap%
\pgfsetroundjoin%
\definecolor{currentfill}{rgb}{0.000000,0.000000,0.000000}%
\pgfsetfillcolor{currentfill}%
\pgfsetlinewidth{0.803000pt}%
\definecolor{currentstroke}{rgb}{0.000000,0.000000,0.000000}%
\pgfsetstrokecolor{currentstroke}%
\pgfsetdash{}{0pt}%
\pgfsys@defobject{currentmarker}{\pgfqpoint{0.000000in}{-0.048611in}}{\pgfqpoint{0.000000in}{0.000000in}}{%
\pgfpathmoveto{\pgfqpoint{0.000000in}{0.000000in}}%
\pgfpathlineto{\pgfqpoint{0.000000in}{-0.048611in}}%
\pgfusepath{stroke,fill}%
}%
\begin{pgfscope}%
\pgfsys@transformshift{3.835139in}{0.528000in}%
\pgfsys@useobject{currentmarker}{}%
\end{pgfscope}%
\end{pgfscope}%
\begin{pgfscope}%
\definecolor{textcolor}{rgb}{0.000000,0.000000,0.000000}%
\pgfsetstrokecolor{textcolor}%
\pgfsetfillcolor{textcolor}%
\pgftext[x=3.835139in,y=0.430778in,,top]{\color{textcolor}{\sffamily\fontsize{10.000000}{12.000000}\selectfont\catcode`\^=\active\def^{\ifmmode\sp\else\^{}\fi}\catcode`\%=\active\def%{\%}125}}%
\end{pgfscope}%
\begin{pgfscope}%
\pgfsetbuttcap%
\pgfsetroundjoin%
\definecolor{currentfill}{rgb}{0.000000,0.000000,0.000000}%
\pgfsetfillcolor{currentfill}%
\pgfsetlinewidth{0.803000pt}%
\definecolor{currentstroke}{rgb}{0.000000,0.000000,0.000000}%
\pgfsetstrokecolor{currentstroke}%
\pgfsetdash{}{0pt}%
\pgfsys@defobject{currentmarker}{\pgfqpoint{0.000000in}{-0.048611in}}{\pgfqpoint{0.000000in}{0.000000in}}{%
\pgfpathmoveto{\pgfqpoint{0.000000in}{0.000000in}}%
\pgfpathlineto{\pgfqpoint{0.000000in}{-0.048611in}}%
\pgfusepath{stroke,fill}%
}%
\begin{pgfscope}%
\pgfsys@transformshift{4.401608in}{0.528000in}%
\pgfsys@useobject{currentmarker}{}%
\end{pgfscope}%
\end{pgfscope}%
\begin{pgfscope}%
\definecolor{textcolor}{rgb}{0.000000,0.000000,0.000000}%
\pgfsetstrokecolor{textcolor}%
\pgfsetfillcolor{textcolor}%
\pgftext[x=4.401608in,y=0.430778in,,top]{\color{textcolor}{\sffamily\fontsize{10.000000}{12.000000}\selectfont\catcode`\^=\active\def^{\ifmmode\sp\else\^{}\fi}\catcode`\%=\active\def%{\%}150}}%
\end{pgfscope}%
\begin{pgfscope}%
\pgfsetbuttcap%
\pgfsetroundjoin%
\definecolor{currentfill}{rgb}{0.000000,0.000000,0.000000}%
\pgfsetfillcolor{currentfill}%
\pgfsetlinewidth{0.803000pt}%
\definecolor{currentstroke}{rgb}{0.000000,0.000000,0.000000}%
\pgfsetstrokecolor{currentstroke}%
\pgfsetdash{}{0pt}%
\pgfsys@defobject{currentmarker}{\pgfqpoint{0.000000in}{-0.048611in}}{\pgfqpoint{0.000000in}{0.000000in}}{%
\pgfpathmoveto{\pgfqpoint{0.000000in}{0.000000in}}%
\pgfpathlineto{\pgfqpoint{0.000000in}{-0.048611in}}%
\pgfusepath{stroke,fill}%
}%
\begin{pgfscope}%
\pgfsys@transformshift{4.968077in}{0.528000in}%
\pgfsys@useobject{currentmarker}{}%
\end{pgfscope}%
\end{pgfscope}%
\begin{pgfscope}%
\definecolor{textcolor}{rgb}{0.000000,0.000000,0.000000}%
\pgfsetstrokecolor{textcolor}%
\pgfsetfillcolor{textcolor}%
\pgftext[x=4.968077in,y=0.430778in,,top]{\color{textcolor}{\sffamily\fontsize{10.000000}{12.000000}\selectfont\catcode`\^=\active\def^{\ifmmode\sp\else\^{}\fi}\catcode`\%=\active\def%{\%}175}}%
\end{pgfscope}%
\begin{pgfscope}%
\pgfsetbuttcap%
\pgfsetroundjoin%
\definecolor{currentfill}{rgb}{0.000000,0.000000,0.000000}%
\pgfsetfillcolor{currentfill}%
\pgfsetlinewidth{0.803000pt}%
\definecolor{currentstroke}{rgb}{0.000000,0.000000,0.000000}%
\pgfsetstrokecolor{currentstroke}%
\pgfsetdash{}{0pt}%
\pgfsys@defobject{currentmarker}{\pgfqpoint{0.000000in}{-0.048611in}}{\pgfqpoint{0.000000in}{0.000000in}}{%
\pgfpathmoveto{\pgfqpoint{0.000000in}{0.000000in}}%
\pgfpathlineto{\pgfqpoint{0.000000in}{-0.048611in}}%
\pgfusepath{stroke,fill}%
}%
\begin{pgfscope}%
\pgfsys@transformshift{5.534545in}{0.528000in}%
\pgfsys@useobject{currentmarker}{}%
\end{pgfscope}%
\end{pgfscope}%
\begin{pgfscope}%
\definecolor{textcolor}{rgb}{0.000000,0.000000,0.000000}%
\pgfsetstrokecolor{textcolor}%
\pgfsetfillcolor{textcolor}%
\pgftext[x=5.534545in,y=0.430778in,,top]{\color{textcolor}{\sffamily\fontsize{10.000000}{12.000000}\selectfont\catcode`\^=\active\def^{\ifmmode\sp\else\^{}\fi}\catcode`\%=\active\def%{\%}200}}%
\end{pgfscope}%
\begin{pgfscope}%
\definecolor{textcolor}{rgb}{0.000000,0.000000,0.000000}%
\pgfsetstrokecolor{textcolor}%
\pgfsetfillcolor{textcolor}%
\pgftext[x=3.280000in,y=0.240809in,,top]{\color{textcolor}{\sffamily\fontsize{10.000000}{12.000000}\selectfont\catcode`\^=\active\def^{\ifmmode\sp\else\^{}\fi}\catcode`\%=\active\def%{\%}Epoch}}%
\end{pgfscope}%
\begin{pgfscope}%
\pgfsetbuttcap%
\pgfsetroundjoin%
\definecolor{currentfill}{rgb}{0.000000,0.000000,0.000000}%
\pgfsetfillcolor{currentfill}%
\pgfsetlinewidth{0.803000pt}%
\definecolor{currentstroke}{rgb}{0.000000,0.000000,0.000000}%
\pgfsetstrokecolor{currentstroke}%
\pgfsetdash{}{0pt}%
\pgfsys@defobject{currentmarker}{\pgfqpoint{-0.048611in}{0.000000in}}{\pgfqpoint{-0.000000in}{0.000000in}}{%
\pgfpathmoveto{\pgfqpoint{-0.000000in}{0.000000in}}%
\pgfpathlineto{\pgfqpoint{-0.048611in}{0.000000in}}%
\pgfusepath{stroke,fill}%
}%
\begin{pgfscope}%
\pgfsys@transformshift{0.800000in}{0.833856in}%
\pgfsys@useobject{currentmarker}{}%
\end{pgfscope}%
\end{pgfscope}%
\begin{pgfscope}%
\definecolor{textcolor}{rgb}{0.000000,0.000000,0.000000}%
\pgfsetstrokecolor{textcolor}%
\pgfsetfillcolor{textcolor}%
\pgftext[x=0.393533in, y=0.781094in, left, base]{\color{textcolor}{\sffamily\fontsize{10.000000}{12.000000}\selectfont\catcode`\^=\active\def^{\ifmmode\sp\else\^{}\fi}\catcode`\%=\active\def%{\%}0.40}}%
\end{pgfscope}%
\begin{pgfscope}%
\pgfsetbuttcap%
\pgfsetroundjoin%
\definecolor{currentfill}{rgb}{0.000000,0.000000,0.000000}%
\pgfsetfillcolor{currentfill}%
\pgfsetlinewidth{0.803000pt}%
\definecolor{currentstroke}{rgb}{0.000000,0.000000,0.000000}%
\pgfsetstrokecolor{currentstroke}%
\pgfsetdash{}{0pt}%
\pgfsys@defobject{currentmarker}{\pgfqpoint{-0.048611in}{0.000000in}}{\pgfqpoint{-0.000000in}{0.000000in}}{%
\pgfpathmoveto{\pgfqpoint{-0.000000in}{0.000000in}}%
\pgfpathlineto{\pgfqpoint{-0.048611in}{0.000000in}}%
\pgfusepath{stroke,fill}%
}%
\begin{pgfscope}%
\pgfsys@transformshift{0.800000in}{1.384119in}%
\pgfsys@useobject{currentmarker}{}%
\end{pgfscope}%
\end{pgfscope}%
\begin{pgfscope}%
\definecolor{textcolor}{rgb}{0.000000,0.000000,0.000000}%
\pgfsetstrokecolor{textcolor}%
\pgfsetfillcolor{textcolor}%
\pgftext[x=0.393533in, y=1.331357in, left, base]{\color{textcolor}{\sffamily\fontsize{10.000000}{12.000000}\selectfont\catcode`\^=\active\def^{\ifmmode\sp\else\^{}\fi}\catcode`\%=\active\def%{\%}0.45}}%
\end{pgfscope}%
\begin{pgfscope}%
\pgfsetbuttcap%
\pgfsetroundjoin%
\definecolor{currentfill}{rgb}{0.000000,0.000000,0.000000}%
\pgfsetfillcolor{currentfill}%
\pgfsetlinewidth{0.803000pt}%
\definecolor{currentstroke}{rgb}{0.000000,0.000000,0.000000}%
\pgfsetstrokecolor{currentstroke}%
\pgfsetdash{}{0pt}%
\pgfsys@defobject{currentmarker}{\pgfqpoint{-0.048611in}{0.000000in}}{\pgfqpoint{-0.000000in}{0.000000in}}{%
\pgfpathmoveto{\pgfqpoint{-0.000000in}{0.000000in}}%
\pgfpathlineto{\pgfqpoint{-0.048611in}{0.000000in}}%
\pgfusepath{stroke,fill}%
}%
\begin{pgfscope}%
\pgfsys@transformshift{0.800000in}{1.934382in}%
\pgfsys@useobject{currentmarker}{}%
\end{pgfscope}%
\end{pgfscope}%
\begin{pgfscope}%
\definecolor{textcolor}{rgb}{0.000000,0.000000,0.000000}%
\pgfsetstrokecolor{textcolor}%
\pgfsetfillcolor{textcolor}%
\pgftext[x=0.393533in, y=1.881620in, left, base]{\color{textcolor}{\sffamily\fontsize{10.000000}{12.000000}\selectfont\catcode`\^=\active\def^{\ifmmode\sp\else\^{}\fi}\catcode`\%=\active\def%{\%}0.50}}%
\end{pgfscope}%
\begin{pgfscope}%
\pgfsetbuttcap%
\pgfsetroundjoin%
\definecolor{currentfill}{rgb}{0.000000,0.000000,0.000000}%
\pgfsetfillcolor{currentfill}%
\pgfsetlinewidth{0.803000pt}%
\definecolor{currentstroke}{rgb}{0.000000,0.000000,0.000000}%
\pgfsetstrokecolor{currentstroke}%
\pgfsetdash{}{0pt}%
\pgfsys@defobject{currentmarker}{\pgfqpoint{-0.048611in}{0.000000in}}{\pgfqpoint{-0.000000in}{0.000000in}}{%
\pgfpathmoveto{\pgfqpoint{-0.000000in}{0.000000in}}%
\pgfpathlineto{\pgfqpoint{-0.048611in}{0.000000in}}%
\pgfusepath{stroke,fill}%
}%
\begin{pgfscope}%
\pgfsys@transformshift{0.800000in}{2.484645in}%
\pgfsys@useobject{currentmarker}{}%
\end{pgfscope}%
\end{pgfscope}%
\begin{pgfscope}%
\definecolor{textcolor}{rgb}{0.000000,0.000000,0.000000}%
\pgfsetstrokecolor{textcolor}%
\pgfsetfillcolor{textcolor}%
\pgftext[x=0.393533in, y=2.431883in, left, base]{\color{textcolor}{\sffamily\fontsize{10.000000}{12.000000}\selectfont\catcode`\^=\active\def^{\ifmmode\sp\else\^{}\fi}\catcode`\%=\active\def%{\%}0.55}}%
\end{pgfscope}%
\begin{pgfscope}%
\pgfsetbuttcap%
\pgfsetroundjoin%
\definecolor{currentfill}{rgb}{0.000000,0.000000,0.000000}%
\pgfsetfillcolor{currentfill}%
\pgfsetlinewidth{0.803000pt}%
\definecolor{currentstroke}{rgb}{0.000000,0.000000,0.000000}%
\pgfsetstrokecolor{currentstroke}%
\pgfsetdash{}{0pt}%
\pgfsys@defobject{currentmarker}{\pgfqpoint{-0.048611in}{0.000000in}}{\pgfqpoint{-0.000000in}{0.000000in}}{%
\pgfpathmoveto{\pgfqpoint{-0.000000in}{0.000000in}}%
\pgfpathlineto{\pgfqpoint{-0.048611in}{0.000000in}}%
\pgfusepath{stroke,fill}%
}%
\begin{pgfscope}%
\pgfsys@transformshift{0.800000in}{3.034908in}%
\pgfsys@useobject{currentmarker}{}%
\end{pgfscope}%
\end{pgfscope}%
\begin{pgfscope}%
\definecolor{textcolor}{rgb}{0.000000,0.000000,0.000000}%
\pgfsetstrokecolor{textcolor}%
\pgfsetfillcolor{textcolor}%
\pgftext[x=0.393533in, y=2.982147in, left, base]{\color{textcolor}{\sffamily\fontsize{10.000000}{12.000000}\selectfont\catcode`\^=\active\def^{\ifmmode\sp\else\^{}\fi}\catcode`\%=\active\def%{\%}0.60}}%
\end{pgfscope}%
\begin{pgfscope}%
\pgfsetbuttcap%
\pgfsetroundjoin%
\definecolor{currentfill}{rgb}{0.000000,0.000000,0.000000}%
\pgfsetfillcolor{currentfill}%
\pgfsetlinewidth{0.803000pt}%
\definecolor{currentstroke}{rgb}{0.000000,0.000000,0.000000}%
\pgfsetstrokecolor{currentstroke}%
\pgfsetdash{}{0pt}%
\pgfsys@defobject{currentmarker}{\pgfqpoint{-0.048611in}{0.000000in}}{\pgfqpoint{-0.000000in}{0.000000in}}{%
\pgfpathmoveto{\pgfqpoint{-0.000000in}{0.000000in}}%
\pgfpathlineto{\pgfqpoint{-0.048611in}{0.000000in}}%
\pgfusepath{stroke,fill}%
}%
\begin{pgfscope}%
\pgfsys@transformshift{0.800000in}{3.585171in}%
\pgfsys@useobject{currentmarker}{}%
\end{pgfscope}%
\end{pgfscope}%
\begin{pgfscope}%
\definecolor{textcolor}{rgb}{0.000000,0.000000,0.000000}%
\pgfsetstrokecolor{textcolor}%
\pgfsetfillcolor{textcolor}%
\pgftext[x=0.393533in, y=3.532410in, left, base]{\color{textcolor}{\sffamily\fontsize{10.000000}{12.000000}\selectfont\catcode`\^=\active\def^{\ifmmode\sp\else\^{}\fi}\catcode`\%=\active\def%{\%}0.65}}%
\end{pgfscope}%
\begin{pgfscope}%
\pgfsetbuttcap%
\pgfsetroundjoin%
\definecolor{currentfill}{rgb}{0.000000,0.000000,0.000000}%
\pgfsetfillcolor{currentfill}%
\pgfsetlinewidth{0.803000pt}%
\definecolor{currentstroke}{rgb}{0.000000,0.000000,0.000000}%
\pgfsetstrokecolor{currentstroke}%
\pgfsetdash{}{0pt}%
\pgfsys@defobject{currentmarker}{\pgfqpoint{-0.048611in}{0.000000in}}{\pgfqpoint{-0.000000in}{0.000000in}}{%
\pgfpathmoveto{\pgfqpoint{-0.000000in}{0.000000in}}%
\pgfpathlineto{\pgfqpoint{-0.048611in}{0.000000in}}%
\pgfusepath{stroke,fill}%
}%
\begin{pgfscope}%
\pgfsys@transformshift{0.800000in}{4.135434in}%
\pgfsys@useobject{currentmarker}{}%
\end{pgfscope}%
\end{pgfscope}%
\begin{pgfscope}%
\definecolor{textcolor}{rgb}{0.000000,0.000000,0.000000}%
\pgfsetstrokecolor{textcolor}%
\pgfsetfillcolor{textcolor}%
\pgftext[x=0.393533in, y=4.082673in, left, base]{\color{textcolor}{\sffamily\fontsize{10.000000}{12.000000}\selectfont\catcode`\^=\active\def^{\ifmmode\sp\else\^{}\fi}\catcode`\%=\active\def%{\%}0.70}}%
\end{pgfscope}%
\begin{pgfscope}%
\definecolor{textcolor}{rgb}{0.000000,0.000000,0.000000}%
\pgfsetstrokecolor{textcolor}%
\pgfsetfillcolor{textcolor}%
\pgftext[x=0.337977in,y=2.376000in,,bottom,rotate=90.000000]{\color{textcolor}{\sffamily\fontsize{10.000000}{12.000000}\selectfont\catcode`\^=\active\def^{\ifmmode\sp\else\^{}\fi}\catcode`\%=\active\def%{\%}Loss}}%
\end{pgfscope}%
\begin{pgfscope}%
\pgfpathrectangle{\pgfqpoint{0.800000in}{0.528000in}}{\pgfqpoint{4.960000in}{3.696000in}}%
\pgfusepath{clip}%
\pgfsetrectcap%
\pgfsetroundjoin%
\pgfsetlinewidth{1.505625pt}%
\definecolor{currentstroke}{rgb}{0.121569,0.466667,0.705882}%
\pgfsetstrokecolor{currentstroke}%
\pgfsetdash{}{0pt}%
\pgfpathmoveto{\pgfqpoint{1.025455in}{3.646090in}}%
\pgfpathlineto{\pgfqpoint{1.048113in}{3.316751in}}%
\pgfpathlineto{\pgfqpoint{1.070772in}{3.147608in}}%
\pgfpathlineto{\pgfqpoint{1.093431in}{3.075059in}}%
\pgfpathlineto{\pgfqpoint{1.116090in}{3.009104in}}%
\pgfpathlineto{\pgfqpoint{1.138748in}{2.957928in}}%
\pgfpathlineto{\pgfqpoint{1.161407in}{2.923531in}}%
\pgfpathlineto{\pgfqpoint{1.184066in}{2.886982in}}%
\pgfpathlineto{\pgfqpoint{1.206725in}{2.867509in}}%
\pgfpathlineto{\pgfqpoint{1.229383in}{2.836679in}}%
\pgfpathlineto{\pgfqpoint{1.252042in}{2.811131in}}%
\pgfpathlineto{\pgfqpoint{1.274701in}{2.795475in}}%
\pgfpathlineto{\pgfqpoint{1.297360in}{2.767883in}}%
\pgfpathlineto{\pgfqpoint{1.320018in}{2.747825in}}%
\pgfpathlineto{\pgfqpoint{1.342677in}{2.702942in}}%
\pgfpathlineto{\pgfqpoint{1.365336in}{2.627781in}}%
\pgfpathlineto{\pgfqpoint{1.387995in}{2.527545in}}%
\pgfpathlineto{\pgfqpoint{1.410653in}{2.397779in}}%
\pgfpathlineto{\pgfqpoint{1.455971in}{2.066693in}}%
\pgfpathlineto{\pgfqpoint{1.478630in}{1.918332in}}%
\pgfpathlineto{\pgfqpoint{1.501288in}{1.782035in}}%
\pgfpathlineto{\pgfqpoint{1.523947in}{1.678207in}}%
\pgfpathlineto{\pgfqpoint{1.546606in}{1.613673in}}%
\pgfpathlineto{\pgfqpoint{1.569265in}{1.553824in}}%
\pgfpathlineto{\pgfqpoint{1.591923in}{1.518982in}}%
\pgfpathlineto{\pgfqpoint{1.637241in}{1.457502in}}%
\pgfpathlineto{\pgfqpoint{1.659899in}{1.437342in}}%
\pgfpathlineto{\pgfqpoint{1.682558in}{1.432887in}}%
\pgfpathlineto{\pgfqpoint{1.705217in}{1.424968in}}%
\pgfpathlineto{\pgfqpoint{1.727876in}{1.418983in}}%
\pgfpathlineto{\pgfqpoint{1.750534in}{1.392994in}}%
\pgfpathlineto{\pgfqpoint{1.773193in}{1.400551in}}%
\pgfpathlineto{\pgfqpoint{1.795852in}{1.392070in}}%
\pgfpathlineto{\pgfqpoint{1.818511in}{1.374852in}}%
\pgfpathlineto{\pgfqpoint{1.841169in}{1.392117in}}%
\pgfpathlineto{\pgfqpoint{1.863828in}{1.369300in}}%
\pgfpathlineto{\pgfqpoint{1.886487in}{1.378338in}}%
\pgfpathlineto{\pgfqpoint{1.909146in}{1.355524in}}%
\pgfpathlineto{\pgfqpoint{1.931804in}{1.363897in}}%
\pgfpathlineto{\pgfqpoint{1.954463in}{1.360236in}}%
\pgfpathlineto{\pgfqpoint{1.977122in}{1.358667in}}%
\pgfpathlineto{\pgfqpoint{1.999781in}{1.345921in}}%
\pgfpathlineto{\pgfqpoint{2.022439in}{1.339148in}}%
\pgfpathlineto{\pgfqpoint{2.045098in}{1.336269in}}%
\pgfpathlineto{\pgfqpoint{2.067757in}{1.345409in}}%
\pgfpathlineto{\pgfqpoint{2.090416in}{1.333997in}}%
\pgfpathlineto{\pgfqpoint{2.113074in}{1.327474in}}%
\pgfpathlineto{\pgfqpoint{2.135733in}{1.339803in}}%
\pgfpathlineto{\pgfqpoint{2.158392in}{1.329945in}}%
\pgfpathlineto{\pgfqpoint{2.203709in}{1.312837in}}%
\pgfpathlineto{\pgfqpoint{2.226368in}{1.325118in}}%
\pgfpathlineto{\pgfqpoint{2.249027in}{1.320112in}}%
\pgfpathlineto{\pgfqpoint{2.271686in}{1.316556in}}%
\pgfpathlineto{\pgfqpoint{2.294344in}{1.305167in}}%
\pgfpathlineto{\pgfqpoint{2.317003in}{1.315281in}}%
\pgfpathlineto{\pgfqpoint{2.339662in}{1.311539in}}%
\pgfpathlineto{\pgfqpoint{2.362321in}{1.309087in}}%
\pgfpathlineto{\pgfqpoint{2.384979in}{1.298943in}}%
\pgfpathlineto{\pgfqpoint{2.407638in}{1.322016in}}%
\pgfpathlineto{\pgfqpoint{2.430297in}{1.295086in}}%
\pgfpathlineto{\pgfqpoint{2.452956in}{1.315561in}}%
\pgfpathlineto{\pgfqpoint{2.475614in}{1.283253in}}%
\pgfpathlineto{\pgfqpoint{2.498273in}{1.299417in}}%
\pgfpathlineto{\pgfqpoint{2.520932in}{1.290283in}}%
\pgfpathlineto{\pgfqpoint{2.543591in}{1.278218in}}%
\pgfpathlineto{\pgfqpoint{2.566249in}{1.274648in}}%
\pgfpathlineto{\pgfqpoint{2.588908in}{1.260058in}}%
\pgfpathlineto{\pgfqpoint{2.611567in}{1.261577in}}%
\pgfpathlineto{\pgfqpoint{2.634226in}{1.276150in}}%
\pgfpathlineto{\pgfqpoint{2.656884in}{1.249354in}}%
\pgfpathlineto{\pgfqpoint{2.679543in}{1.265610in}}%
\pgfpathlineto{\pgfqpoint{2.724861in}{1.245832in}}%
\pgfpathlineto{\pgfqpoint{2.747519in}{1.270269in}}%
\pgfpathlineto{\pgfqpoint{2.770178in}{1.250338in}}%
\pgfpathlineto{\pgfqpoint{2.792837in}{1.247128in}}%
\pgfpathlineto{\pgfqpoint{2.815496in}{1.240877in}}%
\pgfpathlineto{\pgfqpoint{2.838154in}{1.238176in}}%
\pgfpathlineto{\pgfqpoint{2.860813in}{1.230538in}}%
\pgfpathlineto{\pgfqpoint{2.883472in}{1.227103in}}%
\pgfpathlineto{\pgfqpoint{2.906131in}{1.234120in}}%
\pgfpathlineto{\pgfqpoint{2.928789in}{1.222840in}}%
\pgfpathlineto{\pgfqpoint{2.951448in}{1.214833in}}%
\pgfpathlineto{\pgfqpoint{2.974107in}{1.225295in}}%
\pgfpathlineto{\pgfqpoint{2.996766in}{1.225850in}}%
\pgfpathlineto{\pgfqpoint{3.019424in}{1.207897in}}%
\pgfpathlineto{\pgfqpoint{3.042083in}{1.218576in}}%
\pgfpathlineto{\pgfqpoint{3.064742in}{1.208954in}}%
\pgfpathlineto{\pgfqpoint{3.087401in}{1.223592in}}%
\pgfpathlineto{\pgfqpoint{3.110059in}{1.203358in}}%
\pgfpathlineto{\pgfqpoint{3.132718in}{1.196760in}}%
\pgfpathlineto{\pgfqpoint{3.155377in}{1.209080in}}%
\pgfpathlineto{\pgfqpoint{3.178036in}{1.188520in}}%
\pgfpathlineto{\pgfqpoint{3.200694in}{1.193126in}}%
\pgfpathlineto{\pgfqpoint{3.223353in}{1.193068in}}%
\pgfpathlineto{\pgfqpoint{3.246012in}{1.175617in}}%
\pgfpathlineto{\pgfqpoint{3.268671in}{1.176908in}}%
\pgfpathlineto{\pgfqpoint{3.291329in}{1.188964in}}%
\pgfpathlineto{\pgfqpoint{3.313988in}{1.177784in}}%
\pgfpathlineto{\pgfqpoint{3.336647in}{1.172817in}}%
\pgfpathlineto{\pgfqpoint{3.359306in}{1.173020in}}%
\pgfpathlineto{\pgfqpoint{3.381964in}{1.163098in}}%
\pgfpathlineto{\pgfqpoint{3.404623in}{1.158882in}}%
\pgfpathlineto{\pgfqpoint{3.427282in}{1.162354in}}%
\pgfpathlineto{\pgfqpoint{3.449941in}{1.163004in}}%
\pgfpathlineto{\pgfqpoint{3.472599in}{1.161265in}}%
\pgfpathlineto{\pgfqpoint{3.495258in}{1.154225in}}%
\pgfpathlineto{\pgfqpoint{3.517917in}{1.140341in}}%
\pgfpathlineto{\pgfqpoint{3.540576in}{1.173674in}}%
\pgfpathlineto{\pgfqpoint{3.563234in}{1.140850in}}%
\pgfpathlineto{\pgfqpoint{3.585893in}{1.146638in}}%
\pgfpathlineto{\pgfqpoint{3.608552in}{1.147971in}}%
\pgfpathlineto{\pgfqpoint{3.631211in}{1.126114in}}%
\pgfpathlineto{\pgfqpoint{3.653869in}{1.133937in}}%
\pgfpathlineto{\pgfqpoint{3.676528in}{1.124021in}}%
\pgfpathlineto{\pgfqpoint{3.699187in}{1.148367in}}%
\pgfpathlineto{\pgfqpoint{3.721846in}{1.126034in}}%
\pgfpathlineto{\pgfqpoint{3.744504in}{1.130817in}}%
\pgfpathlineto{\pgfqpoint{3.767163in}{1.122494in}}%
\pgfpathlineto{\pgfqpoint{3.789822in}{1.106357in}}%
\pgfpathlineto{\pgfqpoint{3.812481in}{1.108213in}}%
\pgfpathlineto{\pgfqpoint{3.835139in}{1.105102in}}%
\pgfpathlineto{\pgfqpoint{3.857798in}{1.107353in}}%
\pgfpathlineto{\pgfqpoint{3.880457in}{1.107272in}}%
\pgfpathlineto{\pgfqpoint{3.903116in}{1.117760in}}%
\pgfpathlineto{\pgfqpoint{3.925774in}{1.103474in}}%
\pgfpathlineto{\pgfqpoint{3.948433in}{1.106673in}}%
\pgfpathlineto{\pgfqpoint{3.971092in}{1.089437in}}%
\pgfpathlineto{\pgfqpoint{3.993751in}{1.102764in}}%
\pgfpathlineto{\pgfqpoint{4.016409in}{1.098403in}}%
\pgfpathlineto{\pgfqpoint{4.039068in}{1.090213in}}%
\pgfpathlineto{\pgfqpoint{4.061727in}{1.105144in}}%
\pgfpathlineto{\pgfqpoint{4.084386in}{1.079537in}}%
\pgfpathlineto{\pgfqpoint{4.107044in}{1.091450in}}%
\pgfpathlineto{\pgfqpoint{4.129703in}{1.083129in}}%
\pgfpathlineto{\pgfqpoint{4.152362in}{1.079302in}}%
\pgfpathlineto{\pgfqpoint{4.175021in}{1.067677in}}%
\pgfpathlineto{\pgfqpoint{4.197679in}{1.072701in}}%
\pgfpathlineto{\pgfqpoint{4.220338in}{1.086010in}}%
\pgfpathlineto{\pgfqpoint{4.242997in}{1.065377in}}%
\pgfpathlineto{\pgfqpoint{4.265656in}{1.084924in}}%
\pgfpathlineto{\pgfqpoint{4.288314in}{1.101791in}}%
\pgfpathlineto{\pgfqpoint{4.310973in}{1.070078in}}%
\pgfpathlineto{\pgfqpoint{4.333632in}{1.054602in}}%
\pgfpathlineto{\pgfqpoint{4.356291in}{1.060780in}}%
\pgfpathlineto{\pgfqpoint{4.378949in}{1.057239in}}%
\pgfpathlineto{\pgfqpoint{4.401608in}{1.067462in}}%
\pgfpathlineto{\pgfqpoint{4.424267in}{1.083334in}}%
\pgfpathlineto{\pgfqpoint{4.446926in}{1.067064in}}%
\pgfpathlineto{\pgfqpoint{4.469584in}{1.058756in}}%
\pgfpathlineto{\pgfqpoint{4.492243in}{1.062653in}}%
\pgfpathlineto{\pgfqpoint{4.514902in}{1.058948in}}%
\pgfpathlineto{\pgfqpoint{4.537561in}{1.047707in}}%
\pgfpathlineto{\pgfqpoint{4.560219in}{1.052525in}}%
\pgfpathlineto{\pgfqpoint{4.582878in}{1.060316in}}%
\pgfpathlineto{\pgfqpoint{4.605537in}{1.040860in}}%
\pgfpathlineto{\pgfqpoint{4.628196in}{1.053316in}}%
\pgfpathlineto{\pgfqpoint{4.650854in}{1.060561in}}%
\pgfpathlineto{\pgfqpoint{4.673513in}{1.043473in}}%
\pgfpathlineto{\pgfqpoint{4.718831in}{1.047014in}}%
\pgfpathlineto{\pgfqpoint{4.741489in}{1.038002in}}%
\pgfpathlineto{\pgfqpoint{4.764148in}{1.035856in}}%
\pgfpathlineto{\pgfqpoint{4.786807in}{1.042970in}}%
\pgfpathlineto{\pgfqpoint{4.809466in}{1.037007in}}%
\pgfpathlineto{\pgfqpoint{4.832124in}{1.038167in}}%
\pgfpathlineto{\pgfqpoint{4.854783in}{1.025563in}}%
\pgfpathlineto{\pgfqpoint{4.877442in}{1.027514in}}%
\pgfpathlineto{\pgfqpoint{4.900101in}{1.039135in}}%
\pgfpathlineto{\pgfqpoint{4.922759in}{1.041890in}}%
\pgfpathlineto{\pgfqpoint{4.945418in}{1.055849in}}%
\pgfpathlineto{\pgfqpoint{4.968077in}{1.032938in}}%
\pgfpathlineto{\pgfqpoint{4.990735in}{1.013241in}}%
\pgfpathlineto{\pgfqpoint{5.013394in}{1.040104in}}%
\pgfpathlineto{\pgfqpoint{5.036053in}{1.019202in}}%
\pgfpathlineto{\pgfqpoint{5.058712in}{1.030864in}}%
\pgfpathlineto{\pgfqpoint{5.081370in}{1.020275in}}%
\pgfpathlineto{\pgfqpoint{5.104029in}{1.027955in}}%
\pgfpathlineto{\pgfqpoint{5.126688in}{1.023274in}}%
\pgfpathlineto{\pgfqpoint{5.149347in}{1.029585in}}%
\pgfpathlineto{\pgfqpoint{5.172005in}{1.037277in}}%
\pgfpathlineto{\pgfqpoint{5.194664in}{1.033608in}}%
\pgfpathlineto{\pgfqpoint{5.217323in}{1.040655in}}%
\pgfpathlineto{\pgfqpoint{5.239982in}{1.034430in}}%
\pgfpathlineto{\pgfqpoint{5.262640in}{1.006769in}}%
\pgfpathlineto{\pgfqpoint{5.285299in}{1.004674in}}%
\pgfpathlineto{\pgfqpoint{5.307958in}{1.009251in}}%
\pgfpathlineto{\pgfqpoint{5.330617in}{1.009351in}}%
\pgfpathlineto{\pgfqpoint{5.353275in}{1.013423in}}%
\pgfpathlineto{\pgfqpoint{5.375934in}{1.010190in}}%
\pgfpathlineto{\pgfqpoint{5.398593in}{1.004303in}}%
\pgfpathlineto{\pgfqpoint{5.421252in}{1.011691in}}%
\pgfpathlineto{\pgfqpoint{5.443910in}{1.004527in}}%
\pgfpathlineto{\pgfqpoint{5.466569in}{1.012506in}}%
\pgfpathlineto{\pgfqpoint{5.489228in}{1.013914in}}%
\pgfpathlineto{\pgfqpoint{5.511887in}{1.007628in}}%
\pgfpathlineto{\pgfqpoint{5.534545in}{1.007526in}}%
\pgfpathlineto{\pgfqpoint{5.534545in}{1.007526in}}%
\pgfusepath{stroke}%
\end{pgfscope}%
\begin{pgfscope}%
\pgfpathrectangle{\pgfqpoint{0.800000in}{0.528000in}}{\pgfqpoint{4.960000in}{3.696000in}}%
\pgfusepath{clip}%
\pgfsetrectcap%
\pgfsetroundjoin%
\pgfsetlinewidth{1.505625pt}%
\definecolor{currentstroke}{rgb}{1.000000,0.498039,0.054902}%
\pgfsetstrokecolor{currentstroke}%
\pgfsetdash{}{0pt}%
\pgfpathmoveto{\pgfqpoint{1.025455in}{4.052679in}}%
\pgfpathlineto{\pgfqpoint{1.048113in}{4.030472in}}%
\pgfpathlineto{\pgfqpoint{1.070772in}{3.954097in}}%
\pgfpathlineto{\pgfqpoint{1.093431in}{3.804777in}}%
\pgfpathlineto{\pgfqpoint{1.116090in}{3.646231in}}%
\pgfpathlineto{\pgfqpoint{1.138748in}{3.531300in}}%
\pgfpathlineto{\pgfqpoint{1.161407in}{3.460978in}}%
\pgfpathlineto{\pgfqpoint{1.184066in}{3.414394in}}%
\pgfpathlineto{\pgfqpoint{1.206725in}{3.380909in}}%
\pgfpathlineto{\pgfqpoint{1.229383in}{3.361856in}}%
\pgfpathlineto{\pgfqpoint{1.252042in}{3.336207in}}%
\pgfpathlineto{\pgfqpoint{1.274701in}{3.323975in}}%
\pgfpathlineto{\pgfqpoint{1.297360in}{3.300776in}}%
\pgfpathlineto{\pgfqpoint{1.320018in}{3.295714in}}%
\pgfpathlineto{\pgfqpoint{1.342677in}{3.269548in}}%
\pgfpathlineto{\pgfqpoint{1.365336in}{3.253098in}}%
\pgfpathlineto{\pgfqpoint{1.387995in}{3.242595in}}%
\pgfpathlineto{\pgfqpoint{1.410653in}{3.228122in}}%
\pgfpathlineto{\pgfqpoint{1.433312in}{3.211579in}}%
\pgfpathlineto{\pgfqpoint{1.455971in}{3.200762in}}%
\pgfpathlineto{\pgfqpoint{1.478630in}{3.187169in}}%
\pgfpathlineto{\pgfqpoint{1.501288in}{3.179999in}}%
\pgfpathlineto{\pgfqpoint{1.569265in}{3.145493in}}%
\pgfpathlineto{\pgfqpoint{1.591923in}{3.122241in}}%
\pgfpathlineto{\pgfqpoint{1.614582in}{3.111951in}}%
\pgfpathlineto{\pgfqpoint{1.637241in}{3.099285in}}%
\pgfpathlineto{\pgfqpoint{1.659899in}{3.094134in}}%
\pgfpathlineto{\pgfqpoint{1.682558in}{3.087369in}}%
\pgfpathlineto{\pgfqpoint{1.705217in}{3.064570in}}%
\pgfpathlineto{\pgfqpoint{1.727876in}{3.067895in}}%
\pgfpathlineto{\pgfqpoint{1.750534in}{3.043177in}}%
\pgfpathlineto{\pgfqpoint{1.773193in}{3.034923in}}%
\pgfpathlineto{\pgfqpoint{1.795852in}{3.019659in}}%
\pgfpathlineto{\pgfqpoint{1.818511in}{3.023157in}}%
\pgfpathlineto{\pgfqpoint{1.841169in}{3.001855in}}%
\pgfpathlineto{\pgfqpoint{1.863828in}{2.996814in}}%
\pgfpathlineto{\pgfqpoint{1.886487in}{2.988499in}}%
\pgfpathlineto{\pgfqpoint{1.909146in}{2.985667in}}%
\pgfpathlineto{\pgfqpoint{1.931804in}{2.966821in}}%
\pgfpathlineto{\pgfqpoint{1.954463in}{2.969320in}}%
\pgfpathlineto{\pgfqpoint{1.977122in}{2.952545in}}%
\pgfpathlineto{\pgfqpoint{1.999781in}{2.956262in}}%
\pgfpathlineto{\pgfqpoint{2.022439in}{2.936565in}}%
\pgfpathlineto{\pgfqpoint{2.045098in}{2.934377in}}%
\pgfpathlineto{\pgfqpoint{2.067757in}{2.928526in}}%
\pgfpathlineto{\pgfqpoint{2.090416in}{2.927089in}}%
\pgfpathlineto{\pgfqpoint{2.113074in}{2.931486in}}%
\pgfpathlineto{\pgfqpoint{2.135733in}{2.906503in}}%
\pgfpathlineto{\pgfqpoint{2.158392in}{2.901603in}}%
\pgfpathlineto{\pgfqpoint{2.181051in}{2.931180in}}%
\pgfpathlineto{\pgfqpoint{2.203709in}{2.892586in}}%
\pgfpathlineto{\pgfqpoint{2.226368in}{2.889207in}}%
\pgfpathlineto{\pgfqpoint{2.249027in}{2.903301in}}%
\pgfpathlineto{\pgfqpoint{2.271686in}{2.883161in}}%
\pgfpathlineto{\pgfqpoint{2.294344in}{2.880970in}}%
\pgfpathlineto{\pgfqpoint{2.317003in}{2.863774in}}%
\pgfpathlineto{\pgfqpoint{2.339662in}{2.859921in}}%
\pgfpathlineto{\pgfqpoint{2.362321in}{2.865141in}}%
\pgfpathlineto{\pgfqpoint{2.384979in}{2.857301in}}%
\pgfpathlineto{\pgfqpoint{2.407638in}{2.853876in}}%
\pgfpathlineto{\pgfqpoint{2.430297in}{2.856910in}}%
\pgfpathlineto{\pgfqpoint{2.452956in}{2.842357in}}%
\pgfpathlineto{\pgfqpoint{2.475614in}{2.832644in}}%
\pgfpathlineto{\pgfqpoint{2.498273in}{2.859446in}}%
\pgfpathlineto{\pgfqpoint{2.520932in}{2.820978in}}%
\pgfpathlineto{\pgfqpoint{2.543591in}{2.853091in}}%
\pgfpathlineto{\pgfqpoint{2.566249in}{2.822884in}}%
\pgfpathlineto{\pgfqpoint{2.588908in}{2.817554in}}%
\pgfpathlineto{\pgfqpoint{2.611567in}{2.823970in}}%
\pgfpathlineto{\pgfqpoint{2.634226in}{2.818624in}}%
\pgfpathlineto{\pgfqpoint{2.656884in}{2.814929in}}%
\pgfpathlineto{\pgfqpoint{2.679543in}{2.794606in}}%
\pgfpathlineto{\pgfqpoint{2.702202in}{2.792424in}}%
\pgfpathlineto{\pgfqpoint{2.724861in}{2.788648in}}%
\pgfpathlineto{\pgfqpoint{2.747519in}{2.786529in}}%
\pgfpathlineto{\pgfqpoint{2.770178in}{2.780283in}}%
\pgfpathlineto{\pgfqpoint{2.792837in}{2.780377in}}%
\pgfpathlineto{\pgfqpoint{2.815496in}{2.775106in}}%
\pgfpathlineto{\pgfqpoint{2.838154in}{2.766364in}}%
\pgfpathlineto{\pgfqpoint{2.860813in}{2.767119in}}%
\pgfpathlineto{\pgfqpoint{2.883472in}{2.756663in}}%
\pgfpathlineto{\pgfqpoint{2.906131in}{2.762386in}}%
\pgfpathlineto{\pgfqpoint{2.928789in}{2.779394in}}%
\pgfpathlineto{\pgfqpoint{2.951448in}{2.744225in}}%
\pgfpathlineto{\pgfqpoint{2.974107in}{2.736840in}}%
\pgfpathlineto{\pgfqpoint{2.996766in}{2.739083in}}%
\pgfpathlineto{\pgfqpoint{3.019424in}{2.747754in}}%
\pgfpathlineto{\pgfqpoint{3.042083in}{2.733171in}}%
\pgfpathlineto{\pgfqpoint{3.064742in}{2.736248in}}%
\pgfpathlineto{\pgfqpoint{3.087401in}{2.722583in}}%
\pgfpathlineto{\pgfqpoint{3.110059in}{2.710394in}}%
\pgfpathlineto{\pgfqpoint{3.132718in}{2.729140in}}%
\pgfpathlineto{\pgfqpoint{3.155377in}{2.705875in}}%
\pgfpathlineto{\pgfqpoint{3.178036in}{2.704566in}}%
\pgfpathlineto{\pgfqpoint{3.200694in}{2.699613in}}%
\pgfpathlineto{\pgfqpoint{3.223353in}{2.705201in}}%
\pgfpathlineto{\pgfqpoint{3.246012in}{2.686649in}}%
\pgfpathlineto{\pgfqpoint{3.268671in}{2.695509in}}%
\pgfpathlineto{\pgfqpoint{3.291329in}{2.676649in}}%
\pgfpathlineto{\pgfqpoint{3.313988in}{2.676069in}}%
\pgfpathlineto{\pgfqpoint{3.336647in}{2.690018in}}%
\pgfpathlineto{\pgfqpoint{3.359306in}{2.673531in}}%
\pgfpathlineto{\pgfqpoint{3.381964in}{2.666774in}}%
\pgfpathlineto{\pgfqpoint{3.404623in}{2.677851in}}%
\pgfpathlineto{\pgfqpoint{3.427282in}{2.651868in}}%
\pgfpathlineto{\pgfqpoint{3.449941in}{2.654148in}}%
\pgfpathlineto{\pgfqpoint{3.472599in}{2.647883in}}%
\pgfpathlineto{\pgfqpoint{3.495258in}{2.656901in}}%
\pgfpathlineto{\pgfqpoint{3.517917in}{2.646447in}}%
\pgfpathlineto{\pgfqpoint{3.540576in}{2.653501in}}%
\pgfpathlineto{\pgfqpoint{3.563234in}{2.638858in}}%
\pgfpathlineto{\pgfqpoint{3.585893in}{2.651454in}}%
\pgfpathlineto{\pgfqpoint{3.608552in}{2.626957in}}%
\pgfpathlineto{\pgfqpoint{3.653869in}{2.624992in}}%
\pgfpathlineto{\pgfqpoint{3.676528in}{2.613869in}}%
\pgfpathlineto{\pgfqpoint{3.699187in}{2.612734in}}%
\pgfpathlineto{\pgfqpoint{3.721846in}{2.613097in}}%
\pgfpathlineto{\pgfqpoint{3.744504in}{2.637535in}}%
\pgfpathlineto{\pgfqpoint{3.767163in}{2.618919in}}%
\pgfpathlineto{\pgfqpoint{3.789822in}{2.603313in}}%
\pgfpathlineto{\pgfqpoint{3.812481in}{2.621660in}}%
\pgfpathlineto{\pgfqpoint{3.835139in}{2.600293in}}%
\pgfpathlineto{\pgfqpoint{3.857798in}{2.612588in}}%
\pgfpathlineto{\pgfqpoint{3.880457in}{2.600595in}}%
\pgfpathlineto{\pgfqpoint{3.903116in}{2.583286in}}%
\pgfpathlineto{\pgfqpoint{3.925774in}{2.616522in}}%
\pgfpathlineto{\pgfqpoint{3.948433in}{2.583937in}}%
\pgfpathlineto{\pgfqpoint{3.971092in}{2.581817in}}%
\pgfpathlineto{\pgfqpoint{3.993751in}{2.595370in}}%
\pgfpathlineto{\pgfqpoint{4.016409in}{2.578295in}}%
\pgfpathlineto{\pgfqpoint{4.039068in}{2.574268in}}%
\pgfpathlineto{\pgfqpoint{4.061727in}{2.577458in}}%
\pgfpathlineto{\pgfqpoint{4.084386in}{2.570424in}}%
\pgfpathlineto{\pgfqpoint{4.107044in}{2.579272in}}%
\pgfpathlineto{\pgfqpoint{4.129703in}{2.592816in}}%
\pgfpathlineto{\pgfqpoint{4.152362in}{2.564994in}}%
\pgfpathlineto{\pgfqpoint{4.175021in}{2.566408in}}%
\pgfpathlineto{\pgfqpoint{4.197679in}{2.552566in}}%
\pgfpathlineto{\pgfqpoint{4.242997in}{2.549306in}}%
\pgfpathlineto{\pgfqpoint{4.265656in}{2.545591in}}%
\pgfpathlineto{\pgfqpoint{4.288314in}{2.548525in}}%
\pgfpathlineto{\pgfqpoint{4.310973in}{2.545421in}}%
\pgfpathlineto{\pgfqpoint{4.356291in}{2.528807in}}%
\pgfpathlineto{\pgfqpoint{4.378949in}{2.546115in}}%
\pgfpathlineto{\pgfqpoint{4.401608in}{2.523781in}}%
\pgfpathlineto{\pgfqpoint{4.424267in}{2.546443in}}%
\pgfpathlineto{\pgfqpoint{4.446926in}{2.537891in}}%
\pgfpathlineto{\pgfqpoint{4.469584in}{2.535134in}}%
\pgfpathlineto{\pgfqpoint{4.492243in}{2.518448in}}%
\pgfpathlineto{\pgfqpoint{4.537561in}{2.511995in}}%
\pgfpathlineto{\pgfqpoint{4.560219in}{2.519722in}}%
\pgfpathlineto{\pgfqpoint{4.582878in}{2.499354in}}%
\pgfpathlineto{\pgfqpoint{4.605537in}{2.499509in}}%
\pgfpathlineto{\pgfqpoint{4.628196in}{2.517664in}}%
\pgfpathlineto{\pgfqpoint{4.650854in}{2.493520in}}%
\pgfpathlineto{\pgfqpoint{4.673513in}{2.504560in}}%
\pgfpathlineto{\pgfqpoint{4.696172in}{2.499551in}}%
\pgfpathlineto{\pgfqpoint{4.718831in}{2.498920in}}%
\pgfpathlineto{\pgfqpoint{4.741489in}{2.487993in}}%
\pgfpathlineto{\pgfqpoint{4.764148in}{2.483787in}}%
\pgfpathlineto{\pgfqpoint{4.786807in}{2.482811in}}%
\pgfpathlineto{\pgfqpoint{4.809466in}{2.507893in}}%
\pgfpathlineto{\pgfqpoint{4.832124in}{2.482087in}}%
\pgfpathlineto{\pgfqpoint{4.854783in}{2.482618in}}%
\pgfpathlineto{\pgfqpoint{4.877442in}{2.467286in}}%
\pgfpathlineto{\pgfqpoint{4.900101in}{2.468442in}}%
\pgfpathlineto{\pgfqpoint{4.922759in}{2.471915in}}%
\pgfpathlineto{\pgfqpoint{4.945418in}{2.462575in}}%
\pgfpathlineto{\pgfqpoint{4.968077in}{2.470234in}}%
\pgfpathlineto{\pgfqpoint{4.990735in}{2.463232in}}%
\pgfpathlineto{\pgfqpoint{5.013394in}{2.460162in}}%
\pgfpathlineto{\pgfqpoint{5.036053in}{2.452576in}}%
\pgfpathlineto{\pgfqpoint{5.058712in}{2.457333in}}%
\pgfpathlineto{\pgfqpoint{5.081370in}{2.447474in}}%
\pgfpathlineto{\pgfqpoint{5.104029in}{2.446567in}}%
\pgfpathlineto{\pgfqpoint{5.126688in}{2.440614in}}%
\pgfpathlineto{\pgfqpoint{5.149347in}{2.444242in}}%
\pgfpathlineto{\pgfqpoint{5.194664in}{2.434251in}}%
\pgfpathlineto{\pgfqpoint{5.217323in}{2.442179in}}%
\pgfpathlineto{\pgfqpoint{5.239982in}{2.433163in}}%
\pgfpathlineto{\pgfqpoint{5.262640in}{2.445824in}}%
\pgfpathlineto{\pgfqpoint{5.285299in}{2.436916in}}%
\pgfpathlineto{\pgfqpoint{5.307958in}{2.444297in}}%
\pgfpathlineto{\pgfqpoint{5.330617in}{2.421105in}}%
\pgfpathlineto{\pgfqpoint{5.353275in}{2.411478in}}%
\pgfpathlineto{\pgfqpoint{5.375934in}{2.416368in}}%
\pgfpathlineto{\pgfqpoint{5.421252in}{2.405666in}}%
\pgfpathlineto{\pgfqpoint{5.443910in}{2.402438in}}%
\pgfpathlineto{\pgfqpoint{5.466569in}{2.409599in}}%
\pgfpathlineto{\pgfqpoint{5.489228in}{2.404517in}}%
\pgfpathlineto{\pgfqpoint{5.511887in}{2.403290in}}%
\pgfpathlineto{\pgfqpoint{5.534545in}{2.403503in}}%
\pgfpathlineto{\pgfqpoint{5.534545in}{2.403503in}}%
\pgfusepath{stroke}%
\end{pgfscope}%
\begin{pgfscope}%
\pgfpathrectangle{\pgfqpoint{0.800000in}{0.528000in}}{\pgfqpoint{4.960000in}{3.696000in}}%
\pgfusepath{clip}%
\pgfsetrectcap%
\pgfsetroundjoin%
\pgfsetlinewidth{1.505625pt}%
\definecolor{currentstroke}{rgb}{0.172549,0.627451,0.172549}%
\pgfsetstrokecolor{currentstroke}%
\pgfsetdash{}{0pt}%
\pgfpathmoveto{\pgfqpoint{1.025455in}{3.041976in}}%
\pgfpathlineto{\pgfqpoint{1.048113in}{2.756325in}}%
\pgfpathlineto{\pgfqpoint{1.070772in}{2.650460in}}%
\pgfpathlineto{\pgfqpoint{1.093431in}{2.569946in}}%
\pgfpathlineto{\pgfqpoint{1.161407in}{2.415498in}}%
\pgfpathlineto{\pgfqpoint{1.184066in}{2.361606in}}%
\pgfpathlineto{\pgfqpoint{1.206725in}{2.299440in}}%
\pgfpathlineto{\pgfqpoint{1.229383in}{2.224739in}}%
\pgfpathlineto{\pgfqpoint{1.252042in}{2.167497in}}%
\pgfpathlineto{\pgfqpoint{1.274701in}{2.114583in}}%
\pgfpathlineto{\pgfqpoint{1.297360in}{2.036222in}}%
\pgfpathlineto{\pgfqpoint{1.320018in}{1.971910in}}%
\pgfpathlineto{\pgfqpoint{1.365336in}{1.823008in}}%
\pgfpathlineto{\pgfqpoint{1.387995in}{1.740054in}}%
\pgfpathlineto{\pgfqpoint{1.410653in}{1.662437in}}%
\pgfpathlineto{\pgfqpoint{1.433312in}{1.589165in}}%
\pgfpathlineto{\pgfqpoint{1.478630in}{1.424904in}}%
\pgfpathlineto{\pgfqpoint{1.501288in}{1.377361in}}%
\pgfpathlineto{\pgfqpoint{1.523947in}{1.318632in}}%
\pgfpathlineto{\pgfqpoint{1.546606in}{1.286580in}}%
\pgfpathlineto{\pgfqpoint{1.569265in}{1.268993in}}%
\pgfpathlineto{\pgfqpoint{1.591923in}{1.228793in}}%
\pgfpathlineto{\pgfqpoint{1.614582in}{1.207002in}}%
\pgfpathlineto{\pgfqpoint{1.637241in}{1.190814in}}%
\pgfpathlineto{\pgfqpoint{1.659899in}{1.147829in}}%
\pgfpathlineto{\pgfqpoint{1.682558in}{1.129487in}}%
\pgfpathlineto{\pgfqpoint{1.705217in}{1.101057in}}%
\pgfpathlineto{\pgfqpoint{1.727876in}{1.090051in}}%
\pgfpathlineto{\pgfqpoint{1.750534in}{1.076191in}}%
\pgfpathlineto{\pgfqpoint{1.773193in}{1.072186in}}%
\pgfpathlineto{\pgfqpoint{1.795852in}{1.059659in}}%
\pgfpathlineto{\pgfqpoint{1.818511in}{1.038144in}}%
\pgfpathlineto{\pgfqpoint{1.841169in}{1.049766in}}%
\pgfpathlineto{\pgfqpoint{1.863828in}{1.037641in}}%
\pgfpathlineto{\pgfqpoint{1.886487in}{1.006187in}}%
\pgfpathlineto{\pgfqpoint{1.909146in}{0.994721in}}%
\pgfpathlineto{\pgfqpoint{1.931804in}{0.990813in}}%
\pgfpathlineto{\pgfqpoint{1.954463in}{0.988267in}}%
\pgfpathlineto{\pgfqpoint{1.977122in}{0.967896in}}%
\pgfpathlineto{\pgfqpoint{1.999781in}{0.996681in}}%
\pgfpathlineto{\pgfqpoint{2.022439in}{0.967863in}}%
\pgfpathlineto{\pgfqpoint{2.045098in}{0.959822in}}%
\pgfpathlineto{\pgfqpoint{2.067757in}{0.958441in}}%
\pgfpathlineto{\pgfqpoint{2.090416in}{0.946374in}}%
\pgfpathlineto{\pgfqpoint{2.113074in}{0.943320in}}%
\pgfpathlineto{\pgfqpoint{2.135733in}{0.951741in}}%
\pgfpathlineto{\pgfqpoint{2.158392in}{0.929312in}}%
\pgfpathlineto{\pgfqpoint{2.181051in}{0.923092in}}%
\pgfpathlineto{\pgfqpoint{2.203709in}{0.915324in}}%
\pgfpathlineto{\pgfqpoint{2.226368in}{0.920061in}}%
\pgfpathlineto{\pgfqpoint{2.249027in}{0.931717in}}%
\pgfpathlineto{\pgfqpoint{2.271686in}{0.913583in}}%
\pgfpathlineto{\pgfqpoint{2.294344in}{0.898462in}}%
\pgfpathlineto{\pgfqpoint{2.317003in}{0.908765in}}%
\pgfpathlineto{\pgfqpoint{2.339662in}{0.903992in}}%
\pgfpathlineto{\pgfqpoint{2.362321in}{0.926154in}}%
\pgfpathlineto{\pgfqpoint{2.384979in}{0.883716in}}%
\pgfpathlineto{\pgfqpoint{2.407638in}{0.892799in}}%
\pgfpathlineto{\pgfqpoint{2.452956in}{0.871593in}}%
\pgfpathlineto{\pgfqpoint{2.475614in}{0.885860in}}%
\pgfpathlineto{\pgfqpoint{2.498273in}{0.866358in}}%
\pgfpathlineto{\pgfqpoint{2.520932in}{0.885856in}}%
\pgfpathlineto{\pgfqpoint{2.543591in}{0.887078in}}%
\pgfpathlineto{\pgfqpoint{2.566249in}{0.896336in}}%
\pgfpathlineto{\pgfqpoint{2.588908in}{0.877307in}}%
\pgfpathlineto{\pgfqpoint{2.611567in}{0.852847in}}%
\pgfpathlineto{\pgfqpoint{2.634226in}{0.852649in}}%
\pgfpathlineto{\pgfqpoint{2.656884in}{0.860352in}}%
\pgfpathlineto{\pgfqpoint{2.679543in}{0.847857in}}%
\pgfpathlineto{\pgfqpoint{2.702202in}{0.837531in}}%
\pgfpathlineto{\pgfqpoint{2.724861in}{0.849403in}}%
\pgfpathlineto{\pgfqpoint{2.747519in}{0.849087in}}%
\pgfpathlineto{\pgfqpoint{2.770178in}{0.836348in}}%
\pgfpathlineto{\pgfqpoint{2.792837in}{0.868317in}}%
\pgfpathlineto{\pgfqpoint{2.815496in}{0.831347in}}%
\pgfpathlineto{\pgfqpoint{2.838154in}{0.828857in}}%
\pgfpathlineto{\pgfqpoint{2.860813in}{0.834157in}}%
\pgfpathlineto{\pgfqpoint{2.883472in}{0.830231in}}%
\pgfpathlineto{\pgfqpoint{2.906131in}{0.842245in}}%
\pgfpathlineto{\pgfqpoint{2.928789in}{0.841430in}}%
\pgfpathlineto{\pgfqpoint{2.951448in}{0.838524in}}%
\pgfpathlineto{\pgfqpoint{2.974107in}{0.829864in}}%
\pgfpathlineto{\pgfqpoint{2.996766in}{0.831706in}}%
\pgfpathlineto{\pgfqpoint{3.019424in}{0.813760in}}%
\pgfpathlineto{\pgfqpoint{3.042083in}{0.851102in}}%
\pgfpathlineto{\pgfqpoint{3.064742in}{0.822970in}}%
\pgfpathlineto{\pgfqpoint{3.087401in}{0.812808in}}%
\pgfpathlineto{\pgfqpoint{3.110059in}{0.807655in}}%
\pgfpathlineto{\pgfqpoint{3.132718in}{0.808579in}}%
\pgfpathlineto{\pgfqpoint{3.155377in}{0.802844in}}%
\pgfpathlineto{\pgfqpoint{3.178036in}{0.835490in}}%
\pgfpathlineto{\pgfqpoint{3.200694in}{0.809244in}}%
\pgfpathlineto{\pgfqpoint{3.223353in}{0.812627in}}%
\pgfpathlineto{\pgfqpoint{3.246012in}{0.798047in}}%
\pgfpathlineto{\pgfqpoint{3.268671in}{0.789591in}}%
\pgfpathlineto{\pgfqpoint{3.291329in}{0.792630in}}%
\pgfpathlineto{\pgfqpoint{3.313988in}{0.796858in}}%
\pgfpathlineto{\pgfqpoint{3.336647in}{0.816536in}}%
\pgfpathlineto{\pgfqpoint{3.359306in}{0.782490in}}%
\pgfpathlineto{\pgfqpoint{3.381964in}{0.808956in}}%
\pgfpathlineto{\pgfqpoint{3.404623in}{0.806102in}}%
\pgfpathlineto{\pgfqpoint{3.427282in}{0.791433in}}%
\pgfpathlineto{\pgfqpoint{3.449941in}{0.795713in}}%
\pgfpathlineto{\pgfqpoint{3.472599in}{0.780870in}}%
\pgfpathlineto{\pgfqpoint{3.495258in}{0.789959in}}%
\pgfpathlineto{\pgfqpoint{3.517917in}{0.810880in}}%
\pgfpathlineto{\pgfqpoint{3.540576in}{0.788315in}}%
\pgfpathlineto{\pgfqpoint{3.563234in}{0.776090in}}%
\pgfpathlineto{\pgfqpoint{3.608552in}{0.778032in}}%
\pgfpathlineto{\pgfqpoint{3.631211in}{0.770653in}}%
\pgfpathlineto{\pgfqpoint{3.653869in}{0.767723in}}%
\pgfpathlineto{\pgfqpoint{3.676528in}{0.766077in}}%
\pgfpathlineto{\pgfqpoint{3.699187in}{0.775104in}}%
\pgfpathlineto{\pgfqpoint{3.721846in}{0.771955in}}%
\pgfpathlineto{\pgfqpoint{3.744504in}{0.762785in}}%
\pgfpathlineto{\pgfqpoint{3.767163in}{0.775850in}}%
\pgfpathlineto{\pgfqpoint{3.789822in}{0.762473in}}%
\pgfpathlineto{\pgfqpoint{3.812481in}{0.776083in}}%
\pgfpathlineto{\pgfqpoint{3.835139in}{0.774984in}}%
\pgfpathlineto{\pgfqpoint{3.857798in}{0.767421in}}%
\pgfpathlineto{\pgfqpoint{3.880457in}{0.774636in}}%
\pgfpathlineto{\pgfqpoint{3.903116in}{0.762399in}}%
\pgfpathlineto{\pgfqpoint{3.925774in}{0.762302in}}%
\pgfpathlineto{\pgfqpoint{3.948433in}{0.754777in}}%
\pgfpathlineto{\pgfqpoint{3.971092in}{0.752807in}}%
\pgfpathlineto{\pgfqpoint{3.993751in}{0.756321in}}%
\pgfpathlineto{\pgfqpoint{4.016409in}{0.753103in}}%
\pgfpathlineto{\pgfqpoint{4.039068in}{0.770670in}}%
\pgfpathlineto{\pgfqpoint{4.061727in}{0.760860in}}%
\pgfpathlineto{\pgfqpoint{4.084386in}{0.773771in}}%
\pgfpathlineto{\pgfqpoint{4.107044in}{0.758552in}}%
\pgfpathlineto{\pgfqpoint{4.129703in}{0.748553in}}%
\pgfpathlineto{\pgfqpoint{4.152362in}{0.755222in}}%
\pgfpathlineto{\pgfqpoint{4.175021in}{0.763609in}}%
\pgfpathlineto{\pgfqpoint{4.197679in}{0.751446in}}%
\pgfpathlineto{\pgfqpoint{4.220338in}{0.754047in}}%
\pgfpathlineto{\pgfqpoint{4.242997in}{0.749686in}}%
\pgfpathlineto{\pgfqpoint{4.265656in}{0.738428in}}%
\pgfpathlineto{\pgfqpoint{4.288314in}{0.741794in}}%
\pgfpathlineto{\pgfqpoint{4.310973in}{0.760465in}}%
\pgfpathlineto{\pgfqpoint{4.333632in}{0.752278in}}%
\pgfpathlineto{\pgfqpoint{4.356291in}{0.733360in}}%
\pgfpathlineto{\pgfqpoint{4.378949in}{0.747479in}}%
\pgfpathlineto{\pgfqpoint{4.401608in}{0.742795in}}%
\pgfpathlineto{\pgfqpoint{4.424267in}{0.743944in}}%
\pgfpathlineto{\pgfqpoint{4.446926in}{0.738775in}}%
\pgfpathlineto{\pgfqpoint{4.469584in}{0.746697in}}%
\pgfpathlineto{\pgfqpoint{4.492243in}{0.731863in}}%
\pgfpathlineto{\pgfqpoint{4.514902in}{0.725138in}}%
\pgfpathlineto{\pgfqpoint{4.537561in}{0.735976in}}%
\pgfpathlineto{\pgfqpoint{4.560219in}{0.739373in}}%
\pgfpathlineto{\pgfqpoint{4.605537in}{0.736739in}}%
\pgfpathlineto{\pgfqpoint{4.628196in}{0.739955in}}%
\pgfpathlineto{\pgfqpoint{4.650854in}{0.744638in}}%
\pgfpathlineto{\pgfqpoint{4.673513in}{0.725563in}}%
\pgfpathlineto{\pgfqpoint{4.696172in}{0.722314in}}%
\pgfpathlineto{\pgfqpoint{4.741489in}{0.736942in}}%
\pgfpathlineto{\pgfqpoint{4.764148in}{0.734555in}}%
\pgfpathlineto{\pgfqpoint{4.786807in}{0.742912in}}%
\pgfpathlineto{\pgfqpoint{4.809466in}{0.749521in}}%
\pgfpathlineto{\pgfqpoint{4.832124in}{0.728219in}}%
\pgfpathlineto{\pgfqpoint{4.854783in}{0.735578in}}%
\pgfpathlineto{\pgfqpoint{4.877442in}{0.717373in}}%
\pgfpathlineto{\pgfqpoint{4.900101in}{0.721138in}}%
\pgfpathlineto{\pgfqpoint{4.922759in}{0.752175in}}%
\pgfpathlineto{\pgfqpoint{4.945418in}{0.719543in}}%
\pgfpathlineto{\pgfqpoint{4.968077in}{0.726197in}}%
\pgfpathlineto{\pgfqpoint{4.990735in}{0.727616in}}%
\pgfpathlineto{\pgfqpoint{5.013394in}{0.719825in}}%
\pgfpathlineto{\pgfqpoint{5.036053in}{0.717845in}}%
\pgfpathlineto{\pgfqpoint{5.058712in}{0.728883in}}%
\pgfpathlineto{\pgfqpoint{5.081370in}{0.736116in}}%
\pgfpathlineto{\pgfqpoint{5.104029in}{0.722972in}}%
\pgfpathlineto{\pgfqpoint{5.126688in}{0.736705in}}%
\pgfpathlineto{\pgfqpoint{5.149347in}{0.701443in}}%
\pgfpathlineto{\pgfqpoint{5.172005in}{0.705361in}}%
\pgfpathlineto{\pgfqpoint{5.194664in}{0.704759in}}%
\pgfpathlineto{\pgfqpoint{5.217323in}{0.706384in}}%
\pgfpathlineto{\pgfqpoint{5.239982in}{0.702562in}}%
\pgfpathlineto{\pgfqpoint{5.262640in}{0.705622in}}%
\pgfpathlineto{\pgfqpoint{5.285299in}{0.703120in}}%
\pgfpathlineto{\pgfqpoint{5.307958in}{0.702588in}}%
\pgfpathlineto{\pgfqpoint{5.330617in}{0.705239in}}%
\pgfpathlineto{\pgfqpoint{5.353275in}{0.704760in}}%
\pgfpathlineto{\pgfqpoint{5.375934in}{0.707428in}}%
\pgfpathlineto{\pgfqpoint{5.398593in}{0.704119in}}%
\pgfpathlineto{\pgfqpoint{5.421252in}{0.699352in}}%
\pgfpathlineto{\pgfqpoint{5.466569in}{0.701096in}}%
\pgfpathlineto{\pgfqpoint{5.489228in}{0.697775in}}%
\pgfpathlineto{\pgfqpoint{5.511887in}{0.699715in}}%
\pgfpathlineto{\pgfqpoint{5.534545in}{0.699329in}}%
\pgfpathlineto{\pgfqpoint{5.534545in}{0.699329in}}%
\pgfusepath{stroke}%
\end{pgfscope}%
\begin{pgfscope}%
\pgfpathrectangle{\pgfqpoint{0.800000in}{0.528000in}}{\pgfqpoint{4.960000in}{3.696000in}}%
\pgfusepath{clip}%
\pgfsetrectcap%
\pgfsetroundjoin%
\pgfsetlinewidth{1.505625pt}%
\definecolor{currentstroke}{rgb}{0.839216,0.152941,0.156863}%
\pgfsetstrokecolor{currentstroke}%
\pgfsetdash{}{0pt}%
\pgfpathmoveto{\pgfqpoint{1.025455in}{3.560907in}}%
\pgfpathlineto{\pgfqpoint{1.048113in}{3.269264in}}%
\pgfpathlineto{\pgfqpoint{1.070772in}{3.136862in}}%
\pgfpathlineto{\pgfqpoint{1.093431in}{3.060639in}}%
\pgfpathlineto{\pgfqpoint{1.116090in}{2.998719in}}%
\pgfpathlineto{\pgfqpoint{1.161407in}{2.899312in}}%
\pgfpathlineto{\pgfqpoint{1.184066in}{2.854795in}}%
\pgfpathlineto{\pgfqpoint{1.206725in}{2.781384in}}%
\pgfpathlineto{\pgfqpoint{1.229383in}{2.696480in}}%
\pgfpathlineto{\pgfqpoint{1.252042in}{2.604498in}}%
\pgfpathlineto{\pgfqpoint{1.297360in}{2.386187in}}%
\pgfpathlineto{\pgfqpoint{1.320018in}{2.254460in}}%
\pgfpathlineto{\pgfqpoint{1.342677in}{2.134235in}}%
\pgfpathlineto{\pgfqpoint{1.365336in}{1.997321in}}%
\pgfpathlineto{\pgfqpoint{1.387995in}{1.880942in}}%
\pgfpathlineto{\pgfqpoint{1.410653in}{1.791040in}}%
\pgfpathlineto{\pgfqpoint{1.433312in}{1.685690in}}%
\pgfpathlineto{\pgfqpoint{1.455971in}{1.607209in}}%
\pgfpathlineto{\pgfqpoint{1.478630in}{1.544365in}}%
\pgfpathlineto{\pgfqpoint{1.501288in}{1.492036in}}%
\pgfpathlineto{\pgfqpoint{1.523947in}{1.459514in}}%
\pgfpathlineto{\pgfqpoint{1.546606in}{1.453848in}}%
\pgfpathlineto{\pgfqpoint{1.569265in}{1.415365in}}%
\pgfpathlineto{\pgfqpoint{1.591923in}{1.400167in}}%
\pgfpathlineto{\pgfqpoint{1.614582in}{1.387185in}}%
\pgfpathlineto{\pgfqpoint{1.637241in}{1.371081in}}%
\pgfpathlineto{\pgfqpoint{1.659899in}{1.368643in}}%
\pgfpathlineto{\pgfqpoint{1.682558in}{1.342287in}}%
\pgfpathlineto{\pgfqpoint{1.705217in}{1.346773in}}%
\pgfpathlineto{\pgfqpoint{1.727876in}{1.335688in}}%
\pgfpathlineto{\pgfqpoint{1.750534in}{1.328473in}}%
\pgfpathlineto{\pgfqpoint{1.773193in}{1.313084in}}%
\pgfpathlineto{\pgfqpoint{1.795852in}{1.320206in}}%
\pgfpathlineto{\pgfqpoint{1.818511in}{1.324373in}}%
\pgfpathlineto{\pgfqpoint{1.841169in}{1.310942in}}%
\pgfpathlineto{\pgfqpoint{1.863828in}{1.306416in}}%
\pgfpathlineto{\pgfqpoint{1.886487in}{1.288079in}}%
\pgfpathlineto{\pgfqpoint{1.909146in}{1.286146in}}%
\pgfpathlineto{\pgfqpoint{1.931804in}{1.288238in}}%
\pgfpathlineto{\pgfqpoint{1.954463in}{1.275407in}}%
\pgfpathlineto{\pgfqpoint{1.977122in}{1.270542in}}%
\pgfpathlineto{\pgfqpoint{1.999781in}{1.270027in}}%
\pgfpathlineto{\pgfqpoint{2.022439in}{1.249370in}}%
\pgfpathlineto{\pgfqpoint{2.045098in}{1.245796in}}%
\pgfpathlineto{\pgfqpoint{2.067757in}{1.252222in}}%
\pgfpathlineto{\pgfqpoint{2.090416in}{1.249195in}}%
\pgfpathlineto{\pgfqpoint{2.113074in}{1.227802in}}%
\pgfpathlineto{\pgfqpoint{2.135733in}{1.231997in}}%
\pgfpathlineto{\pgfqpoint{2.158392in}{1.233174in}}%
\pgfpathlineto{\pgfqpoint{2.181051in}{1.227995in}}%
\pgfpathlineto{\pgfqpoint{2.203709in}{1.211350in}}%
\pgfpathlineto{\pgfqpoint{2.226368in}{1.212998in}}%
\pgfpathlineto{\pgfqpoint{2.249027in}{1.217324in}}%
\pgfpathlineto{\pgfqpoint{2.271686in}{1.197945in}}%
\pgfpathlineto{\pgfqpoint{2.294344in}{1.200223in}}%
\pgfpathlineto{\pgfqpoint{2.317003in}{1.190427in}}%
\pgfpathlineto{\pgfqpoint{2.339662in}{1.187388in}}%
\pgfpathlineto{\pgfqpoint{2.362321in}{1.176359in}}%
\pgfpathlineto{\pgfqpoint{2.384979in}{1.189841in}}%
\pgfpathlineto{\pgfqpoint{2.407638in}{1.196175in}}%
\pgfpathlineto{\pgfqpoint{2.452956in}{1.171014in}}%
\pgfpathlineto{\pgfqpoint{2.475614in}{1.171925in}}%
\pgfpathlineto{\pgfqpoint{2.498273in}{1.164083in}}%
\pgfpathlineto{\pgfqpoint{2.520932in}{1.166648in}}%
\pgfpathlineto{\pgfqpoint{2.543591in}{1.172723in}}%
\pgfpathlineto{\pgfqpoint{2.566249in}{1.146297in}}%
\pgfpathlineto{\pgfqpoint{2.588908in}{1.157438in}}%
\pgfpathlineto{\pgfqpoint{2.611567in}{1.148443in}}%
\pgfpathlineto{\pgfqpoint{2.634226in}{1.148905in}}%
\pgfpathlineto{\pgfqpoint{2.656884in}{1.156532in}}%
\pgfpathlineto{\pgfqpoint{2.702202in}{1.118289in}}%
\pgfpathlineto{\pgfqpoint{2.724861in}{1.117653in}}%
\pgfpathlineto{\pgfqpoint{2.747519in}{1.141137in}}%
\pgfpathlineto{\pgfqpoint{2.770178in}{1.120575in}}%
\pgfpathlineto{\pgfqpoint{2.792837in}{1.122689in}}%
\pgfpathlineto{\pgfqpoint{2.815496in}{1.107691in}}%
\pgfpathlineto{\pgfqpoint{2.838154in}{1.111066in}}%
\pgfpathlineto{\pgfqpoint{2.860813in}{1.102201in}}%
\pgfpathlineto{\pgfqpoint{2.883472in}{1.090670in}}%
\pgfpathlineto{\pgfqpoint{2.928789in}{1.106365in}}%
\pgfpathlineto{\pgfqpoint{2.951448in}{1.083624in}}%
\pgfpathlineto{\pgfqpoint{2.974107in}{1.090815in}}%
\pgfpathlineto{\pgfqpoint{2.996766in}{1.078642in}}%
\pgfpathlineto{\pgfqpoint{3.019424in}{1.099471in}}%
\pgfpathlineto{\pgfqpoint{3.042083in}{1.073344in}}%
\pgfpathlineto{\pgfqpoint{3.064742in}{1.089327in}}%
\pgfpathlineto{\pgfqpoint{3.110059in}{1.061572in}}%
\pgfpathlineto{\pgfqpoint{3.132718in}{1.075664in}}%
\pgfpathlineto{\pgfqpoint{3.155377in}{1.052184in}}%
\pgfpathlineto{\pgfqpoint{3.178036in}{1.061672in}}%
\pgfpathlineto{\pgfqpoint{3.200694in}{1.056890in}}%
\pgfpathlineto{\pgfqpoint{3.223353in}{1.044303in}}%
\pgfpathlineto{\pgfqpoint{3.246012in}{1.042973in}}%
\pgfpathlineto{\pgfqpoint{3.268671in}{1.038172in}}%
\pgfpathlineto{\pgfqpoint{3.291329in}{1.058338in}}%
\pgfpathlineto{\pgfqpoint{3.313988in}{1.033796in}}%
\pgfpathlineto{\pgfqpoint{3.336647in}{1.044572in}}%
\pgfpathlineto{\pgfqpoint{3.359306in}{1.026197in}}%
\pgfpathlineto{\pgfqpoint{3.381964in}{1.015103in}}%
\pgfpathlineto{\pgfqpoint{3.404623in}{1.024674in}}%
\pgfpathlineto{\pgfqpoint{3.427282in}{1.002970in}}%
\pgfpathlineto{\pgfqpoint{3.449941in}{1.001434in}}%
\pgfpathlineto{\pgfqpoint{3.472599in}{1.005106in}}%
\pgfpathlineto{\pgfqpoint{3.495258in}{0.998064in}}%
\pgfpathlineto{\pgfqpoint{3.517917in}{0.992422in}}%
\pgfpathlineto{\pgfqpoint{3.540576in}{0.981390in}}%
\pgfpathlineto{\pgfqpoint{3.563234in}{0.984920in}}%
\pgfpathlineto{\pgfqpoint{3.585893in}{0.997497in}}%
\pgfpathlineto{\pgfqpoint{3.608552in}{0.979636in}}%
\pgfpathlineto{\pgfqpoint{3.631211in}{0.987777in}}%
\pgfpathlineto{\pgfqpoint{3.653869in}{0.989364in}}%
\pgfpathlineto{\pgfqpoint{3.676528in}{0.985951in}}%
\pgfpathlineto{\pgfqpoint{3.699187in}{0.965794in}}%
\pgfpathlineto{\pgfqpoint{3.721846in}{0.961297in}}%
\pgfpathlineto{\pgfqpoint{3.744504in}{0.970860in}}%
\pgfpathlineto{\pgfqpoint{3.767163in}{0.948240in}}%
\pgfpathlineto{\pgfqpoint{3.789822in}{0.957042in}}%
\pgfpathlineto{\pgfqpoint{3.812481in}{0.946162in}}%
\pgfpathlineto{\pgfqpoint{3.835139in}{0.959058in}}%
\pgfpathlineto{\pgfqpoint{3.857798in}{0.945718in}}%
\pgfpathlineto{\pgfqpoint{3.903116in}{0.948080in}}%
\pgfpathlineto{\pgfqpoint{3.925774in}{0.932683in}}%
\pgfpathlineto{\pgfqpoint{3.948433in}{0.933569in}}%
\pgfpathlineto{\pgfqpoint{3.971092in}{0.929739in}}%
\pgfpathlineto{\pgfqpoint{3.993751in}{0.929035in}}%
\pgfpathlineto{\pgfqpoint{4.016409in}{0.948869in}}%
\pgfpathlineto{\pgfqpoint{4.039068in}{0.922643in}}%
\pgfpathlineto{\pgfqpoint{4.061727in}{0.929893in}}%
\pgfpathlineto{\pgfqpoint{4.084386in}{0.919822in}}%
\pgfpathlineto{\pgfqpoint{4.107044in}{0.927658in}}%
\pgfpathlineto{\pgfqpoint{4.129703in}{0.913004in}}%
\pgfpathlineto{\pgfqpoint{4.152362in}{0.917758in}}%
\pgfpathlineto{\pgfqpoint{4.175021in}{0.925190in}}%
\pgfpathlineto{\pgfqpoint{4.197679in}{0.914716in}}%
\pgfpathlineto{\pgfqpoint{4.220338in}{0.932484in}}%
\pgfpathlineto{\pgfqpoint{4.242997in}{0.923687in}}%
\pgfpathlineto{\pgfqpoint{4.265656in}{0.922671in}}%
\pgfpathlineto{\pgfqpoint{4.288314in}{0.936779in}}%
\pgfpathlineto{\pgfqpoint{4.310973in}{0.912751in}}%
\pgfpathlineto{\pgfqpoint{4.333632in}{0.913360in}}%
\pgfpathlineto{\pgfqpoint{4.356291in}{0.894450in}}%
\pgfpathlineto{\pgfqpoint{4.378949in}{0.900938in}}%
\pgfpathlineto{\pgfqpoint{4.401608in}{0.902299in}}%
\pgfpathlineto{\pgfqpoint{4.424267in}{0.899626in}}%
\pgfpathlineto{\pgfqpoint{4.446926in}{0.914849in}}%
\pgfpathlineto{\pgfqpoint{4.469584in}{0.894153in}}%
\pgfpathlineto{\pgfqpoint{4.492243in}{0.896992in}}%
\pgfpathlineto{\pgfqpoint{4.514902in}{0.895725in}}%
\pgfpathlineto{\pgfqpoint{4.537561in}{0.889942in}}%
\pgfpathlineto{\pgfqpoint{4.560219in}{0.889320in}}%
\pgfpathlineto{\pgfqpoint{4.582878in}{0.877706in}}%
\pgfpathlineto{\pgfqpoint{4.605537in}{0.889848in}}%
\pgfpathlineto{\pgfqpoint{4.628196in}{0.891830in}}%
\pgfpathlineto{\pgfqpoint{4.650854in}{0.891821in}}%
\pgfpathlineto{\pgfqpoint{4.673513in}{0.877417in}}%
\pgfpathlineto{\pgfqpoint{4.696172in}{0.895176in}}%
\pgfpathlineto{\pgfqpoint{4.718831in}{0.880000in}}%
\pgfpathlineto{\pgfqpoint{4.741489in}{0.878401in}}%
\pgfpathlineto{\pgfqpoint{4.764148in}{0.879898in}}%
\pgfpathlineto{\pgfqpoint{4.786807in}{0.882784in}}%
\pgfpathlineto{\pgfqpoint{4.809466in}{0.880738in}}%
\pgfpathlineto{\pgfqpoint{4.832124in}{0.873576in}}%
\pgfpathlineto{\pgfqpoint{4.854783in}{0.880708in}}%
\pgfpathlineto{\pgfqpoint{4.877442in}{0.864131in}}%
\pgfpathlineto{\pgfqpoint{4.900101in}{0.866266in}}%
\pgfpathlineto{\pgfqpoint{4.922759in}{0.881572in}}%
\pgfpathlineto{\pgfqpoint{4.945418in}{0.870936in}}%
\pgfpathlineto{\pgfqpoint{4.968077in}{0.857320in}}%
\pgfpathlineto{\pgfqpoint{4.990735in}{0.869497in}}%
\pgfpathlineto{\pgfqpoint{5.013394in}{0.871823in}}%
\pgfpathlineto{\pgfqpoint{5.036053in}{0.888206in}}%
\pgfpathlineto{\pgfqpoint{5.058712in}{0.861316in}}%
\pgfpathlineto{\pgfqpoint{5.081370in}{0.866004in}}%
\pgfpathlineto{\pgfqpoint{5.104029in}{0.853787in}}%
\pgfpathlineto{\pgfqpoint{5.126688in}{0.854571in}}%
\pgfpathlineto{\pgfqpoint{5.149347in}{0.874896in}}%
\pgfpathlineto{\pgfqpoint{5.172005in}{0.859814in}}%
\pgfpathlineto{\pgfqpoint{5.194664in}{0.852170in}}%
\pgfpathlineto{\pgfqpoint{5.217323in}{0.859416in}}%
\pgfpathlineto{\pgfqpoint{5.239982in}{0.857505in}}%
\pgfpathlineto{\pgfqpoint{5.285299in}{0.861058in}}%
\pgfpathlineto{\pgfqpoint{5.307958in}{0.886492in}}%
\pgfpathlineto{\pgfqpoint{5.330617in}{0.846427in}}%
\pgfpathlineto{\pgfqpoint{5.353275in}{0.857381in}}%
\pgfpathlineto{\pgfqpoint{5.375934in}{0.862442in}}%
\pgfpathlineto{\pgfqpoint{5.398593in}{0.842146in}}%
\pgfpathlineto{\pgfqpoint{5.421252in}{0.856597in}}%
\pgfpathlineto{\pgfqpoint{5.443910in}{0.837296in}}%
\pgfpathlineto{\pgfqpoint{5.466569in}{0.838963in}}%
\pgfpathlineto{\pgfqpoint{5.489228in}{0.844201in}}%
\pgfpathlineto{\pgfqpoint{5.511887in}{0.838025in}}%
\pgfpathlineto{\pgfqpoint{5.534545in}{0.842681in}}%
\pgfpathlineto{\pgfqpoint{5.534545in}{0.842681in}}%
\pgfusepath{stroke}%
\end{pgfscope}%
\begin{pgfscope}%
\pgfpathrectangle{\pgfqpoint{0.800000in}{0.528000in}}{\pgfqpoint{4.960000in}{3.696000in}}%
\pgfusepath{clip}%
\pgfsetrectcap%
\pgfsetroundjoin%
\pgfsetlinewidth{1.505625pt}%
\definecolor{currentstroke}{rgb}{0.580392,0.403922,0.741176}%
\pgfsetstrokecolor{currentstroke}%
\pgfsetdash{}{0pt}%
\pgfpathmoveto{\pgfqpoint{1.025455in}{4.056000in}}%
\pgfpathlineto{\pgfqpoint{1.048113in}{4.047234in}}%
\pgfpathlineto{\pgfqpoint{1.070772in}{4.028364in}}%
\pgfpathlineto{\pgfqpoint{1.093431in}{3.973613in}}%
\pgfpathlineto{\pgfqpoint{1.116090in}{3.872648in}}%
\pgfpathlineto{\pgfqpoint{1.161407in}{3.616688in}}%
\pgfpathlineto{\pgfqpoint{1.184066in}{3.501861in}}%
\pgfpathlineto{\pgfqpoint{1.206725in}{3.408693in}}%
\pgfpathlineto{\pgfqpoint{1.229383in}{3.332505in}}%
\pgfpathlineto{\pgfqpoint{1.252042in}{3.286182in}}%
\pgfpathlineto{\pgfqpoint{1.274701in}{3.255044in}}%
\pgfpathlineto{\pgfqpoint{1.297360in}{3.231445in}}%
\pgfpathlineto{\pgfqpoint{1.320018in}{3.209993in}}%
\pgfpathlineto{\pgfqpoint{1.342677in}{3.183776in}}%
\pgfpathlineto{\pgfqpoint{1.365336in}{3.164411in}}%
\pgfpathlineto{\pgfqpoint{1.387995in}{3.155851in}}%
\pgfpathlineto{\pgfqpoint{1.410653in}{3.127659in}}%
\pgfpathlineto{\pgfqpoint{1.433312in}{3.121435in}}%
\pgfpathlineto{\pgfqpoint{1.455971in}{3.078198in}}%
\pgfpathlineto{\pgfqpoint{1.478630in}{3.057865in}}%
\pgfpathlineto{\pgfqpoint{1.501288in}{3.042699in}}%
\pgfpathlineto{\pgfqpoint{1.546606in}{2.994600in}}%
\pgfpathlineto{\pgfqpoint{1.569265in}{2.973521in}}%
\pgfpathlineto{\pgfqpoint{1.591923in}{2.956530in}}%
\pgfpathlineto{\pgfqpoint{1.614582in}{2.927350in}}%
\pgfpathlineto{\pgfqpoint{1.637241in}{2.914168in}}%
\pgfpathlineto{\pgfqpoint{1.659899in}{2.897946in}}%
\pgfpathlineto{\pgfqpoint{1.682558in}{2.869242in}}%
\pgfpathlineto{\pgfqpoint{1.705217in}{2.854387in}}%
\pgfpathlineto{\pgfqpoint{1.727876in}{2.828310in}}%
\pgfpathlineto{\pgfqpoint{1.750534in}{2.808516in}}%
\pgfpathlineto{\pgfqpoint{1.773193in}{2.795915in}}%
\pgfpathlineto{\pgfqpoint{1.795852in}{2.770204in}}%
\pgfpathlineto{\pgfqpoint{1.818511in}{2.747814in}}%
\pgfpathlineto{\pgfqpoint{1.841169in}{2.734299in}}%
\pgfpathlineto{\pgfqpoint{1.863828in}{2.703252in}}%
\pgfpathlineto{\pgfqpoint{1.886487in}{2.682737in}}%
\pgfpathlineto{\pgfqpoint{1.909146in}{2.665075in}}%
\pgfpathlineto{\pgfqpoint{1.931804in}{2.649481in}}%
\pgfpathlineto{\pgfqpoint{1.954463in}{2.628591in}}%
\pgfpathlineto{\pgfqpoint{1.977122in}{2.615938in}}%
\pgfpathlineto{\pgfqpoint{1.999781in}{2.587292in}}%
\pgfpathlineto{\pgfqpoint{2.045098in}{2.564993in}}%
\pgfpathlineto{\pgfqpoint{2.067757in}{2.541488in}}%
\pgfpathlineto{\pgfqpoint{2.090416in}{2.527900in}}%
\pgfpathlineto{\pgfqpoint{2.113074in}{2.512129in}}%
\pgfpathlineto{\pgfqpoint{2.135733in}{2.492528in}}%
\pgfpathlineto{\pgfqpoint{2.158392in}{2.479125in}}%
\pgfpathlineto{\pgfqpoint{2.181051in}{2.470010in}}%
\pgfpathlineto{\pgfqpoint{2.203709in}{2.450244in}}%
\pgfpathlineto{\pgfqpoint{2.226368in}{2.450116in}}%
\pgfpathlineto{\pgfqpoint{2.249027in}{2.433312in}}%
\pgfpathlineto{\pgfqpoint{2.271686in}{2.422086in}}%
\pgfpathlineto{\pgfqpoint{2.294344in}{2.404046in}}%
\pgfpathlineto{\pgfqpoint{2.317003in}{2.390885in}}%
\pgfpathlineto{\pgfqpoint{2.339662in}{2.387436in}}%
\pgfpathlineto{\pgfqpoint{2.362321in}{2.374909in}}%
\pgfpathlineto{\pgfqpoint{2.384979in}{2.367593in}}%
\pgfpathlineto{\pgfqpoint{2.407638in}{2.357479in}}%
\pgfpathlineto{\pgfqpoint{2.430297in}{2.364871in}}%
\pgfpathlineto{\pgfqpoint{2.452956in}{2.337706in}}%
\pgfpathlineto{\pgfqpoint{2.475614in}{2.334566in}}%
\pgfpathlineto{\pgfqpoint{2.520932in}{2.314795in}}%
\pgfpathlineto{\pgfqpoint{2.543591in}{2.310699in}}%
\pgfpathlineto{\pgfqpoint{2.566249in}{2.308324in}}%
\pgfpathlineto{\pgfqpoint{2.588908in}{2.293703in}}%
\pgfpathlineto{\pgfqpoint{2.611567in}{2.297681in}}%
\pgfpathlineto{\pgfqpoint{2.634226in}{2.289194in}}%
\pgfpathlineto{\pgfqpoint{2.656884in}{2.275432in}}%
\pgfpathlineto{\pgfqpoint{2.679543in}{2.280077in}}%
\pgfpathlineto{\pgfqpoint{2.702202in}{2.265627in}}%
\pgfpathlineto{\pgfqpoint{2.724861in}{2.284831in}}%
\pgfpathlineto{\pgfqpoint{2.747519in}{2.257792in}}%
\pgfpathlineto{\pgfqpoint{2.770178in}{2.247484in}}%
\pgfpathlineto{\pgfqpoint{2.792837in}{2.251013in}}%
\pgfpathlineto{\pgfqpoint{2.815496in}{2.244615in}}%
\pgfpathlineto{\pgfqpoint{2.838154in}{2.233169in}}%
\pgfpathlineto{\pgfqpoint{2.860813in}{2.230791in}}%
\pgfpathlineto{\pgfqpoint{2.883472in}{2.226974in}}%
\pgfpathlineto{\pgfqpoint{2.928789in}{2.225309in}}%
\pgfpathlineto{\pgfqpoint{2.951448in}{2.219331in}}%
\pgfpathlineto{\pgfqpoint{2.974107in}{2.233802in}}%
\pgfpathlineto{\pgfqpoint{2.996766in}{2.220481in}}%
\pgfpathlineto{\pgfqpoint{3.042083in}{2.209304in}}%
\pgfpathlineto{\pgfqpoint{3.064742in}{2.205896in}}%
\pgfpathlineto{\pgfqpoint{3.087401in}{2.195029in}}%
\pgfpathlineto{\pgfqpoint{3.132718in}{2.195757in}}%
\pgfpathlineto{\pgfqpoint{3.155377in}{2.194593in}}%
\pgfpathlineto{\pgfqpoint{3.178036in}{2.177273in}}%
\pgfpathlineto{\pgfqpoint{3.200694in}{2.184133in}}%
\pgfpathlineto{\pgfqpoint{3.223353in}{2.180511in}}%
\pgfpathlineto{\pgfqpoint{3.246012in}{2.178363in}}%
\pgfpathlineto{\pgfqpoint{3.268671in}{2.171841in}}%
\pgfpathlineto{\pgfqpoint{3.291329in}{2.179311in}}%
\pgfpathlineto{\pgfqpoint{3.313988in}{2.168013in}}%
\pgfpathlineto{\pgfqpoint{3.336647in}{2.166797in}}%
\pgfpathlineto{\pgfqpoint{3.359306in}{2.176544in}}%
\pgfpathlineto{\pgfqpoint{3.381964in}{2.169143in}}%
\pgfpathlineto{\pgfqpoint{3.404623in}{2.157174in}}%
\pgfpathlineto{\pgfqpoint{3.427282in}{2.150594in}}%
\pgfpathlineto{\pgfqpoint{3.449941in}{2.153090in}}%
\pgfpathlineto{\pgfqpoint{3.472599in}{2.154172in}}%
\pgfpathlineto{\pgfqpoint{3.517917in}{2.133946in}}%
\pgfpathlineto{\pgfqpoint{3.540576in}{2.133625in}}%
\pgfpathlineto{\pgfqpoint{3.563234in}{2.146079in}}%
\pgfpathlineto{\pgfqpoint{3.585893in}{2.129313in}}%
\pgfpathlineto{\pgfqpoint{3.608552in}{2.127435in}}%
\pgfpathlineto{\pgfqpoint{3.631211in}{2.136889in}}%
\pgfpathlineto{\pgfqpoint{3.653869in}{2.129025in}}%
\pgfpathlineto{\pgfqpoint{3.676528in}{2.132867in}}%
\pgfpathlineto{\pgfqpoint{3.699187in}{2.118134in}}%
\pgfpathlineto{\pgfqpoint{3.721846in}{2.126062in}}%
\pgfpathlineto{\pgfqpoint{3.744504in}{2.109898in}}%
\pgfpathlineto{\pgfqpoint{3.767163in}{2.103122in}}%
\pgfpathlineto{\pgfqpoint{3.789822in}{2.108482in}}%
\pgfpathlineto{\pgfqpoint{3.812481in}{2.112487in}}%
\pgfpathlineto{\pgfqpoint{3.835139in}{2.105339in}}%
\pgfpathlineto{\pgfqpoint{3.857798in}{2.105074in}}%
\pgfpathlineto{\pgfqpoint{3.880457in}{2.100333in}}%
\pgfpathlineto{\pgfqpoint{3.903116in}{2.092914in}}%
\pgfpathlineto{\pgfqpoint{3.925774in}{2.114655in}}%
\pgfpathlineto{\pgfqpoint{3.948433in}{2.093700in}}%
\pgfpathlineto{\pgfqpoint{3.971092in}{2.087155in}}%
\pgfpathlineto{\pgfqpoint{4.016409in}{2.083201in}}%
\pgfpathlineto{\pgfqpoint{4.039068in}{2.093364in}}%
\pgfpathlineto{\pgfqpoint{4.061727in}{2.069396in}}%
\pgfpathlineto{\pgfqpoint{4.084386in}{2.080519in}}%
\pgfpathlineto{\pgfqpoint{4.107044in}{2.072313in}}%
\pgfpathlineto{\pgfqpoint{4.129703in}{2.059827in}}%
\pgfpathlineto{\pgfqpoint{4.152362in}{2.067149in}}%
\pgfpathlineto{\pgfqpoint{4.175021in}{2.063331in}}%
\pgfpathlineto{\pgfqpoint{4.242997in}{2.048427in}}%
\pgfpathlineto{\pgfqpoint{4.265656in}{2.064628in}}%
\pgfpathlineto{\pgfqpoint{4.288314in}{2.056170in}}%
\pgfpathlineto{\pgfqpoint{4.310973in}{2.060344in}}%
\pgfpathlineto{\pgfqpoint{4.333632in}{2.049228in}}%
\pgfpathlineto{\pgfqpoint{4.356291in}{2.044829in}}%
\pgfpathlineto{\pgfqpoint{4.401608in}{2.049005in}}%
\pgfpathlineto{\pgfqpoint{4.424267in}{2.030905in}}%
\pgfpathlineto{\pgfqpoint{4.446926in}{2.034941in}}%
\pgfpathlineto{\pgfqpoint{4.469584in}{2.029880in}}%
\pgfpathlineto{\pgfqpoint{4.492243in}{2.028289in}}%
\pgfpathlineto{\pgfqpoint{4.514902in}{2.021734in}}%
\pgfpathlineto{\pgfqpoint{4.537561in}{2.029342in}}%
\pgfpathlineto{\pgfqpoint{4.582878in}{2.033418in}}%
\pgfpathlineto{\pgfqpoint{4.605537in}{2.021289in}}%
\pgfpathlineto{\pgfqpoint{4.628196in}{2.021985in}}%
\pgfpathlineto{\pgfqpoint{4.650854in}{2.013455in}}%
\pgfpathlineto{\pgfqpoint{4.673513in}{2.012238in}}%
\pgfpathlineto{\pgfqpoint{4.696172in}{2.008730in}}%
\pgfpathlineto{\pgfqpoint{4.718831in}{2.020277in}}%
\pgfpathlineto{\pgfqpoint{4.741489in}{2.011815in}}%
\pgfpathlineto{\pgfqpoint{4.764148in}{2.011532in}}%
\pgfpathlineto{\pgfqpoint{4.786807in}{1.996838in}}%
\pgfpathlineto{\pgfqpoint{4.809466in}{1.999760in}}%
\pgfpathlineto{\pgfqpoint{4.832124in}{1.997000in}}%
\pgfpathlineto{\pgfqpoint{4.854783in}{1.996860in}}%
\pgfpathlineto{\pgfqpoint{4.877442in}{2.004906in}}%
\pgfpathlineto{\pgfqpoint{4.900101in}{1.996621in}}%
\pgfpathlineto{\pgfqpoint{4.922759in}{1.990485in}}%
\pgfpathlineto{\pgfqpoint{4.945418in}{1.994533in}}%
\pgfpathlineto{\pgfqpoint{4.968077in}{1.996512in}}%
\pgfpathlineto{\pgfqpoint{4.990735in}{2.001511in}}%
\pgfpathlineto{\pgfqpoint{5.013394in}{1.989492in}}%
\pgfpathlineto{\pgfqpoint{5.036053in}{1.994557in}}%
\pgfpathlineto{\pgfqpoint{5.058712in}{1.976232in}}%
\pgfpathlineto{\pgfqpoint{5.081370in}{1.975102in}}%
\pgfpathlineto{\pgfqpoint{5.104029in}{1.997740in}}%
\pgfpathlineto{\pgfqpoint{5.126688in}{1.978078in}}%
\pgfpathlineto{\pgfqpoint{5.149347in}{1.971329in}}%
\pgfpathlineto{\pgfqpoint{5.172005in}{1.977102in}}%
\pgfpathlineto{\pgfqpoint{5.194664in}{1.967864in}}%
\pgfpathlineto{\pgfqpoint{5.217323in}{1.976162in}}%
\pgfpathlineto{\pgfqpoint{5.239982in}{1.965613in}}%
\pgfpathlineto{\pgfqpoint{5.262640in}{1.978513in}}%
\pgfpathlineto{\pgfqpoint{5.285299in}{1.959454in}}%
\pgfpathlineto{\pgfqpoint{5.307958in}{1.974158in}}%
\pgfpathlineto{\pgfqpoint{5.330617in}{1.980074in}}%
\pgfpathlineto{\pgfqpoint{5.353275in}{1.967686in}}%
\pgfpathlineto{\pgfqpoint{5.375934in}{1.961370in}}%
\pgfpathlineto{\pgfqpoint{5.398593in}{1.966404in}}%
\pgfpathlineto{\pgfqpoint{5.421252in}{1.958487in}}%
\pgfpathlineto{\pgfqpoint{5.443910in}{1.968199in}}%
\pgfpathlineto{\pgfqpoint{5.466569in}{1.952309in}}%
\pgfpathlineto{\pgfqpoint{5.489228in}{1.951956in}}%
\pgfpathlineto{\pgfqpoint{5.534545in}{1.968705in}}%
\pgfpathlineto{\pgfqpoint{5.534545in}{1.968705in}}%
\pgfusepath{stroke}%
\end{pgfscope}%
\begin{pgfscope}%
\pgfpathrectangle{\pgfqpoint{0.800000in}{0.528000in}}{\pgfqpoint{4.960000in}{3.696000in}}%
\pgfusepath{clip}%
\pgfsetrectcap%
\pgfsetroundjoin%
\pgfsetlinewidth{1.505625pt}%
\definecolor{currentstroke}{rgb}{0.549020,0.337255,0.294118}%
\pgfsetstrokecolor{currentstroke}%
\pgfsetdash{}{0pt}%
\pgfpathmoveto{\pgfqpoint{1.025455in}{3.111747in}}%
\pgfpathlineto{\pgfqpoint{1.048113in}{2.738895in}}%
\pgfpathlineto{\pgfqpoint{1.070772in}{2.632610in}}%
\pgfpathlineto{\pgfqpoint{1.093431in}{2.593706in}}%
\pgfpathlineto{\pgfqpoint{1.116090in}{2.562720in}}%
\pgfpathlineto{\pgfqpoint{1.138748in}{2.517979in}}%
\pgfpathlineto{\pgfqpoint{1.161407in}{2.488401in}}%
\pgfpathlineto{\pgfqpoint{1.184066in}{2.463372in}}%
\pgfpathlineto{\pgfqpoint{1.206725in}{2.444434in}}%
\pgfpathlineto{\pgfqpoint{1.229383in}{2.419986in}}%
\pgfpathlineto{\pgfqpoint{1.252042in}{2.362507in}}%
\pgfpathlineto{\pgfqpoint{1.274701in}{2.325118in}}%
\pgfpathlineto{\pgfqpoint{1.297360in}{2.263632in}}%
\pgfpathlineto{\pgfqpoint{1.320018in}{2.194372in}}%
\pgfpathlineto{\pgfqpoint{1.342677in}{2.098809in}}%
\pgfpathlineto{\pgfqpoint{1.365336in}{1.988527in}}%
\pgfpathlineto{\pgfqpoint{1.433312in}{1.550198in}}%
\pgfpathlineto{\pgfqpoint{1.455971in}{1.424655in}}%
\pgfpathlineto{\pgfqpoint{1.478630in}{1.322242in}}%
\pgfpathlineto{\pgfqpoint{1.501288in}{1.252345in}}%
\pgfpathlineto{\pgfqpoint{1.523947in}{1.186097in}}%
\pgfpathlineto{\pgfqpoint{1.546606in}{1.131145in}}%
\pgfpathlineto{\pgfqpoint{1.569265in}{1.097999in}}%
\pgfpathlineto{\pgfqpoint{1.591923in}{1.058112in}}%
\pgfpathlineto{\pgfqpoint{1.614582in}{1.044185in}}%
\pgfpathlineto{\pgfqpoint{1.637241in}{1.020425in}}%
\pgfpathlineto{\pgfqpoint{1.659899in}{0.998862in}}%
\pgfpathlineto{\pgfqpoint{1.682558in}{0.972484in}}%
\pgfpathlineto{\pgfqpoint{1.705217in}{0.956763in}}%
\pgfpathlineto{\pgfqpoint{1.727876in}{0.974797in}}%
\pgfpathlineto{\pgfqpoint{1.750534in}{0.940605in}}%
\pgfpathlineto{\pgfqpoint{1.773193in}{0.945968in}}%
\pgfpathlineto{\pgfqpoint{1.795852in}{0.935266in}}%
\pgfpathlineto{\pgfqpoint{1.818511in}{0.919363in}}%
\pgfpathlineto{\pgfqpoint{1.841169in}{0.913779in}}%
\pgfpathlineto{\pgfqpoint{1.863828in}{0.890136in}}%
\pgfpathlineto{\pgfqpoint{1.886487in}{0.905861in}}%
\pgfpathlineto{\pgfqpoint{1.909146in}{0.880132in}}%
\pgfpathlineto{\pgfqpoint{1.931804in}{0.875946in}}%
\pgfpathlineto{\pgfqpoint{1.954463in}{0.875410in}}%
\pgfpathlineto{\pgfqpoint{1.977122in}{0.886653in}}%
\pgfpathlineto{\pgfqpoint{1.999781in}{0.859762in}}%
\pgfpathlineto{\pgfqpoint{2.022439in}{0.846965in}}%
\pgfpathlineto{\pgfqpoint{2.045098in}{0.856246in}}%
\pgfpathlineto{\pgfqpoint{2.067757in}{0.861118in}}%
\pgfpathlineto{\pgfqpoint{2.090416in}{0.847921in}}%
\pgfpathlineto{\pgfqpoint{2.113074in}{0.836766in}}%
\pgfpathlineto{\pgfqpoint{2.135733in}{0.834555in}}%
\pgfpathlineto{\pgfqpoint{2.158392in}{0.841463in}}%
\pgfpathlineto{\pgfqpoint{2.181051in}{0.841381in}}%
\pgfpathlineto{\pgfqpoint{2.203709in}{0.822168in}}%
\pgfpathlineto{\pgfqpoint{2.226368in}{0.832636in}}%
\pgfpathlineto{\pgfqpoint{2.249027in}{0.819190in}}%
\pgfpathlineto{\pgfqpoint{2.271686in}{0.820267in}}%
\pgfpathlineto{\pgfqpoint{2.294344in}{0.810365in}}%
\pgfpathlineto{\pgfqpoint{2.317003in}{0.824417in}}%
\pgfpathlineto{\pgfqpoint{2.339662in}{0.807893in}}%
\pgfpathlineto{\pgfqpoint{2.362321in}{0.808659in}}%
\pgfpathlineto{\pgfqpoint{2.384979in}{0.796102in}}%
\pgfpathlineto{\pgfqpoint{2.407638in}{0.799980in}}%
\pgfpathlineto{\pgfqpoint{2.430297in}{0.795088in}}%
\pgfpathlineto{\pgfqpoint{2.452956in}{0.781887in}}%
\pgfpathlineto{\pgfqpoint{2.475614in}{0.791717in}}%
\pgfpathlineto{\pgfqpoint{2.520932in}{0.786510in}}%
\pgfpathlineto{\pgfqpoint{2.543591in}{0.788272in}}%
\pgfpathlineto{\pgfqpoint{2.566249in}{0.772082in}}%
\pgfpathlineto{\pgfqpoint{2.588908in}{0.771144in}}%
\pgfpathlineto{\pgfqpoint{2.611567in}{0.780268in}}%
\pgfpathlineto{\pgfqpoint{2.634226in}{0.765467in}}%
\pgfpathlineto{\pgfqpoint{2.656884in}{0.768560in}}%
\pgfpathlineto{\pgfqpoint{2.679543in}{0.761320in}}%
\pgfpathlineto{\pgfqpoint{2.702202in}{0.787003in}}%
\pgfpathlineto{\pgfqpoint{2.724861in}{0.776338in}}%
\pgfpathlineto{\pgfqpoint{2.747519in}{0.758065in}}%
\pgfpathlineto{\pgfqpoint{2.770178in}{0.766526in}}%
\pgfpathlineto{\pgfqpoint{2.792837in}{0.763921in}}%
\pgfpathlineto{\pgfqpoint{2.815496in}{0.777368in}}%
\pgfpathlineto{\pgfqpoint{2.838154in}{0.771499in}}%
\pgfpathlineto{\pgfqpoint{2.860813in}{0.749992in}}%
\pgfpathlineto{\pgfqpoint{2.883472in}{0.763204in}}%
\pgfpathlineto{\pgfqpoint{2.906131in}{0.745296in}}%
\pgfpathlineto{\pgfqpoint{2.951448in}{0.779044in}}%
\pgfpathlineto{\pgfqpoint{2.974107in}{0.738572in}}%
\pgfpathlineto{\pgfqpoint{2.996766in}{0.757456in}}%
\pgfpathlineto{\pgfqpoint{3.019424in}{0.760816in}}%
\pgfpathlineto{\pgfqpoint{3.042083in}{0.747890in}}%
\pgfpathlineto{\pgfqpoint{3.064742in}{0.739268in}}%
\pgfpathlineto{\pgfqpoint{3.087401in}{0.766184in}}%
\pgfpathlineto{\pgfqpoint{3.110059in}{0.726410in}}%
\pgfpathlineto{\pgfqpoint{3.132718in}{0.731860in}}%
\pgfpathlineto{\pgfqpoint{3.155377in}{0.757101in}}%
\pgfpathlineto{\pgfqpoint{3.178036in}{0.746250in}}%
\pgfpathlineto{\pgfqpoint{3.200694in}{0.739048in}}%
\pgfpathlineto{\pgfqpoint{3.223353in}{0.737709in}}%
\pgfpathlineto{\pgfqpoint{3.246012in}{0.732976in}}%
\pgfpathlineto{\pgfqpoint{3.268671in}{0.730304in}}%
\pgfpathlineto{\pgfqpoint{3.291329in}{0.750583in}}%
\pgfpathlineto{\pgfqpoint{3.313988in}{0.735349in}}%
\pgfpathlineto{\pgfqpoint{3.336647in}{0.746509in}}%
\pgfpathlineto{\pgfqpoint{3.359306in}{0.734893in}}%
\pgfpathlineto{\pgfqpoint{3.381964in}{0.708840in}}%
\pgfpathlineto{\pgfqpoint{3.404623in}{0.708722in}}%
\pgfpathlineto{\pgfqpoint{3.427282in}{0.719101in}}%
\pgfpathlineto{\pgfqpoint{3.449941in}{0.710050in}}%
\pgfpathlineto{\pgfqpoint{3.472599in}{0.707240in}}%
\pgfpathlineto{\pgfqpoint{3.495258in}{0.711214in}}%
\pgfpathlineto{\pgfqpoint{3.517917in}{0.707391in}}%
\pgfpathlineto{\pgfqpoint{3.540576in}{0.710660in}}%
\pgfpathlineto{\pgfqpoint{3.563234in}{0.711962in}}%
\pgfpathlineto{\pgfqpoint{3.585893in}{0.711786in}}%
\pgfpathlineto{\pgfqpoint{3.608552in}{0.706026in}}%
\pgfpathlineto{\pgfqpoint{3.631211in}{0.712586in}}%
\pgfpathlineto{\pgfqpoint{3.653869in}{0.708930in}}%
\pgfpathlineto{\pgfqpoint{3.676528in}{0.707528in}}%
\pgfpathlineto{\pgfqpoint{3.699187in}{0.710741in}}%
\pgfpathlineto{\pgfqpoint{3.744504in}{0.707282in}}%
\pgfpathlineto{\pgfqpoint{3.767163in}{0.711513in}}%
\pgfpathlineto{\pgfqpoint{3.789822in}{0.708670in}}%
\pgfpathlineto{\pgfqpoint{3.812481in}{0.715320in}}%
\pgfpathlineto{\pgfqpoint{3.835139in}{0.706300in}}%
\pgfpathlineto{\pgfqpoint{3.857798in}{0.704970in}}%
\pgfpathlineto{\pgfqpoint{3.880457in}{0.709888in}}%
\pgfpathlineto{\pgfqpoint{3.903116in}{0.706543in}}%
\pgfpathlineto{\pgfqpoint{3.925774in}{0.710572in}}%
\pgfpathlineto{\pgfqpoint{3.948433in}{0.707180in}}%
\pgfpathlineto{\pgfqpoint{3.971092in}{0.705024in}}%
\pgfpathlineto{\pgfqpoint{3.993751in}{0.708142in}}%
\pgfpathlineto{\pgfqpoint{4.039068in}{0.706271in}}%
\pgfpathlineto{\pgfqpoint{4.084386in}{0.702191in}}%
\pgfpathlineto{\pgfqpoint{4.107044in}{0.710482in}}%
\pgfpathlineto{\pgfqpoint{4.129703in}{0.705023in}}%
\pgfpathlineto{\pgfqpoint{4.175021in}{0.709640in}}%
\pgfpathlineto{\pgfqpoint{4.197679in}{0.700836in}}%
\pgfpathlineto{\pgfqpoint{4.220338in}{0.706863in}}%
\pgfpathlineto{\pgfqpoint{4.242997in}{0.703746in}}%
\pgfpathlineto{\pgfqpoint{4.265656in}{0.703696in}}%
\pgfpathlineto{\pgfqpoint{4.288314in}{0.706907in}}%
\pgfpathlineto{\pgfqpoint{4.310973in}{0.702000in}}%
\pgfpathlineto{\pgfqpoint{4.356291in}{0.705932in}}%
\pgfpathlineto{\pgfqpoint{4.378949in}{0.705343in}}%
\pgfpathlineto{\pgfqpoint{4.401608in}{0.703520in}}%
\pgfpathlineto{\pgfqpoint{4.446926in}{0.704605in}}%
\pgfpathlineto{\pgfqpoint{4.469584in}{0.696079in}}%
\pgfpathlineto{\pgfqpoint{4.537561in}{0.697488in}}%
\pgfpathlineto{\pgfqpoint{4.582878in}{0.697165in}}%
\pgfpathlineto{\pgfqpoint{4.605537in}{0.699731in}}%
\pgfpathlineto{\pgfqpoint{4.628196in}{0.697515in}}%
\pgfpathlineto{\pgfqpoint{4.650854in}{0.703421in}}%
\pgfpathlineto{\pgfqpoint{4.673513in}{0.696867in}}%
\pgfpathlineto{\pgfqpoint{4.718831in}{0.699239in}}%
\pgfpathlineto{\pgfqpoint{4.741489in}{0.697176in}}%
\pgfpathlineto{\pgfqpoint{4.809466in}{0.696207in}}%
\pgfpathlineto{\pgfqpoint{4.877442in}{0.696650in}}%
\pgfpathlineto{\pgfqpoint{5.534545in}{0.696500in}}%
\pgfpathlineto{\pgfqpoint{5.534545in}{0.696500in}}%
\pgfusepath{stroke}%
\end{pgfscope}%
\begin{pgfscope}%
\pgfpathrectangle{\pgfqpoint{0.800000in}{0.528000in}}{\pgfqpoint{4.960000in}{3.696000in}}%
\pgfusepath{clip}%
\pgfsetrectcap%
\pgfsetroundjoin%
\pgfsetlinewidth{1.505625pt}%
\definecolor{currentstroke}{rgb}{0.890196,0.466667,0.760784}%
\pgfsetstrokecolor{currentstroke}%
\pgfsetdash{}{0pt}%
\pgfpathmoveto{\pgfqpoint{1.025455in}{3.654559in}}%
\pgfpathlineto{\pgfqpoint{1.048113in}{3.281113in}}%
\pgfpathlineto{\pgfqpoint{1.070772in}{3.094072in}}%
\pgfpathlineto{\pgfqpoint{1.093431in}{3.014840in}}%
\pgfpathlineto{\pgfqpoint{1.138748in}{2.920434in}}%
\pgfpathlineto{\pgfqpoint{1.161407in}{2.885644in}}%
\pgfpathlineto{\pgfqpoint{1.184066in}{2.848708in}}%
\pgfpathlineto{\pgfqpoint{1.206725in}{2.821440in}}%
\pgfpathlineto{\pgfqpoint{1.229383in}{2.779497in}}%
\pgfpathlineto{\pgfqpoint{1.252042in}{2.728162in}}%
\pgfpathlineto{\pgfqpoint{1.274701in}{2.672204in}}%
\pgfpathlineto{\pgfqpoint{1.297360in}{2.578634in}}%
\pgfpathlineto{\pgfqpoint{1.320018in}{2.450354in}}%
\pgfpathlineto{\pgfqpoint{1.342677in}{2.296623in}}%
\pgfpathlineto{\pgfqpoint{1.387995in}{1.872713in}}%
\pgfpathlineto{\pgfqpoint{1.410653in}{1.688286in}}%
\pgfpathlineto{\pgfqpoint{1.433312in}{1.538850in}}%
\pgfpathlineto{\pgfqpoint{1.455971in}{1.427839in}}%
\pgfpathlineto{\pgfqpoint{1.478630in}{1.355097in}}%
\pgfpathlineto{\pgfqpoint{1.501288in}{1.322003in}}%
\pgfpathlineto{\pgfqpoint{1.523947in}{1.277649in}}%
\pgfpathlineto{\pgfqpoint{1.546606in}{1.230564in}}%
\pgfpathlineto{\pgfqpoint{1.569265in}{1.218831in}}%
\pgfpathlineto{\pgfqpoint{1.591923in}{1.214922in}}%
\pgfpathlineto{\pgfqpoint{1.614582in}{1.187591in}}%
\pgfpathlineto{\pgfqpoint{1.637241in}{1.181133in}}%
\pgfpathlineto{\pgfqpoint{1.659899in}{1.167937in}}%
\pgfpathlineto{\pgfqpoint{1.682558in}{1.150332in}}%
\pgfpathlineto{\pgfqpoint{1.705217in}{1.191775in}}%
\pgfpathlineto{\pgfqpoint{1.727876in}{1.141522in}}%
\pgfpathlineto{\pgfqpoint{1.750534in}{1.146055in}}%
\pgfpathlineto{\pgfqpoint{1.773193in}{1.126774in}}%
\pgfpathlineto{\pgfqpoint{1.795852in}{1.129423in}}%
\pgfpathlineto{\pgfqpoint{1.818511in}{1.121061in}}%
\pgfpathlineto{\pgfqpoint{1.841169in}{1.147737in}}%
\pgfpathlineto{\pgfqpoint{1.863828in}{1.132202in}}%
\pgfpathlineto{\pgfqpoint{1.931804in}{1.128213in}}%
\pgfpathlineto{\pgfqpoint{1.954463in}{1.118124in}}%
\pgfpathlineto{\pgfqpoint{1.977122in}{1.122743in}}%
\pgfpathlineto{\pgfqpoint{1.999781in}{1.125532in}}%
\pgfpathlineto{\pgfqpoint{2.022439in}{1.130701in}}%
\pgfpathlineto{\pgfqpoint{2.045098in}{1.115736in}}%
\pgfpathlineto{\pgfqpoint{2.067757in}{1.125637in}}%
\pgfpathlineto{\pgfqpoint{2.090416in}{1.107504in}}%
\pgfpathlineto{\pgfqpoint{2.158392in}{1.116313in}}%
\pgfpathlineto{\pgfqpoint{2.181051in}{1.125269in}}%
\pgfpathlineto{\pgfqpoint{2.203709in}{1.106055in}}%
\pgfpathlineto{\pgfqpoint{2.226368in}{1.104370in}}%
\pgfpathlineto{\pgfqpoint{2.249027in}{1.113097in}}%
\pgfpathlineto{\pgfqpoint{2.271686in}{1.123050in}}%
\pgfpathlineto{\pgfqpoint{2.294344in}{1.106596in}}%
\pgfpathlineto{\pgfqpoint{2.317003in}{1.132328in}}%
\pgfpathlineto{\pgfqpoint{2.339662in}{1.116764in}}%
\pgfpathlineto{\pgfqpoint{2.362321in}{1.112009in}}%
\pgfpathlineto{\pgfqpoint{2.384979in}{1.118689in}}%
\pgfpathlineto{\pgfqpoint{2.407638in}{1.112052in}}%
\pgfpathlineto{\pgfqpoint{2.430297in}{1.107650in}}%
\pgfpathlineto{\pgfqpoint{2.452956in}{1.111997in}}%
\pgfpathlineto{\pgfqpoint{2.475614in}{1.095592in}}%
\pgfpathlineto{\pgfqpoint{2.498273in}{1.100794in}}%
\pgfpathlineto{\pgfqpoint{2.520932in}{1.089177in}}%
\pgfpathlineto{\pgfqpoint{2.566249in}{1.110434in}}%
\pgfpathlineto{\pgfqpoint{2.588908in}{1.114060in}}%
\pgfpathlineto{\pgfqpoint{2.611567in}{1.111378in}}%
\pgfpathlineto{\pgfqpoint{2.634226in}{1.097934in}}%
\pgfpathlineto{\pgfqpoint{2.656884in}{1.105947in}}%
\pgfpathlineto{\pgfqpoint{2.679543in}{1.108451in}}%
\pgfpathlineto{\pgfqpoint{2.702202in}{1.103755in}}%
\pgfpathlineto{\pgfqpoint{2.724861in}{1.096882in}}%
\pgfpathlineto{\pgfqpoint{2.747519in}{1.105997in}}%
\pgfpathlineto{\pgfqpoint{2.770178in}{1.095044in}}%
\pgfpathlineto{\pgfqpoint{2.792837in}{1.080234in}}%
\pgfpathlineto{\pgfqpoint{2.815496in}{1.086882in}}%
\pgfpathlineto{\pgfqpoint{2.838154in}{1.077832in}}%
\pgfpathlineto{\pgfqpoint{2.860813in}{1.086263in}}%
\pgfpathlineto{\pgfqpoint{2.906131in}{1.088762in}}%
\pgfpathlineto{\pgfqpoint{2.928789in}{1.085732in}}%
\pgfpathlineto{\pgfqpoint{2.996766in}{1.084677in}}%
\pgfpathlineto{\pgfqpoint{3.019424in}{1.081853in}}%
\pgfpathlineto{\pgfqpoint{3.042083in}{1.085046in}}%
\pgfpathlineto{\pgfqpoint{3.087401in}{1.084763in}}%
\pgfpathlineto{\pgfqpoint{3.110059in}{1.078849in}}%
\pgfpathlineto{\pgfqpoint{3.132718in}{1.079502in}}%
\pgfpathlineto{\pgfqpoint{3.155377in}{1.077963in}}%
\pgfpathlineto{\pgfqpoint{3.178036in}{1.080521in}}%
\pgfpathlineto{\pgfqpoint{3.200694in}{1.081195in}}%
\pgfpathlineto{\pgfqpoint{3.223353in}{1.078785in}}%
\pgfpathlineto{\pgfqpoint{3.291329in}{1.080020in}}%
\pgfpathlineto{\pgfqpoint{3.313988in}{1.082341in}}%
\pgfpathlineto{\pgfqpoint{3.381964in}{1.079461in}}%
\pgfpathlineto{\pgfqpoint{3.495258in}{1.079547in}}%
\pgfpathlineto{\pgfqpoint{3.585893in}{1.079392in}}%
\pgfpathlineto{\pgfqpoint{4.016409in}{1.079911in}}%
\pgfpathlineto{\pgfqpoint{5.239982in}{1.080361in}}%
\pgfpathlineto{\pgfqpoint{5.534545in}{1.080490in}}%
\pgfpathlineto{\pgfqpoint{5.534545in}{1.080490in}}%
\pgfusepath{stroke}%
\end{pgfscope}%
\begin{pgfscope}%
\pgfpathrectangle{\pgfqpoint{0.800000in}{0.528000in}}{\pgfqpoint{4.960000in}{3.696000in}}%
\pgfusepath{clip}%
\pgfsetrectcap%
\pgfsetroundjoin%
\pgfsetlinewidth{1.505625pt}%
\definecolor{currentstroke}{rgb}{0.498039,0.498039,0.498039}%
\pgfsetstrokecolor{currentstroke}%
\pgfsetdash{}{0pt}%
\pgfpathmoveto{\pgfqpoint{1.025455in}{4.055231in}}%
\pgfpathlineto{\pgfqpoint{1.048113in}{4.041314in}}%
\pgfpathlineto{\pgfqpoint{1.070772in}{4.014919in}}%
\pgfpathlineto{\pgfqpoint{1.093431in}{3.964758in}}%
\pgfpathlineto{\pgfqpoint{1.116090in}{3.866495in}}%
\pgfpathlineto{\pgfqpoint{1.161407in}{3.615049in}}%
\pgfpathlineto{\pgfqpoint{1.184066in}{3.525015in}}%
\pgfpathlineto{\pgfqpoint{1.206725in}{3.460446in}}%
\pgfpathlineto{\pgfqpoint{1.229383in}{3.411261in}}%
\pgfpathlineto{\pgfqpoint{1.252042in}{3.370699in}}%
\pgfpathlineto{\pgfqpoint{1.297360in}{3.318473in}}%
\pgfpathlineto{\pgfqpoint{1.320018in}{3.294331in}}%
\pgfpathlineto{\pgfqpoint{1.342677in}{3.281600in}}%
\pgfpathlineto{\pgfqpoint{1.365336in}{3.254384in}}%
\pgfpathlineto{\pgfqpoint{1.387995in}{3.244286in}}%
\pgfpathlineto{\pgfqpoint{1.433312in}{3.209661in}}%
\pgfpathlineto{\pgfqpoint{1.455971in}{3.189185in}}%
\pgfpathlineto{\pgfqpoint{1.478630in}{3.182265in}}%
\pgfpathlineto{\pgfqpoint{1.501288in}{3.150443in}}%
\pgfpathlineto{\pgfqpoint{1.546606in}{3.117743in}}%
\pgfpathlineto{\pgfqpoint{1.569265in}{3.107907in}}%
\pgfpathlineto{\pgfqpoint{1.591923in}{3.091349in}}%
\pgfpathlineto{\pgfqpoint{1.614582in}{3.071749in}}%
\pgfpathlineto{\pgfqpoint{1.637241in}{3.061565in}}%
\pgfpathlineto{\pgfqpoint{1.659899in}{3.049407in}}%
\pgfpathlineto{\pgfqpoint{1.682558in}{3.029851in}}%
\pgfpathlineto{\pgfqpoint{1.705217in}{3.023135in}}%
\pgfpathlineto{\pgfqpoint{1.727876in}{3.006191in}}%
\pgfpathlineto{\pgfqpoint{1.750534in}{3.005415in}}%
\pgfpathlineto{\pgfqpoint{1.795852in}{2.967033in}}%
\pgfpathlineto{\pgfqpoint{1.818511in}{2.966049in}}%
\pgfpathlineto{\pgfqpoint{1.841169in}{2.952038in}}%
\pgfpathlineto{\pgfqpoint{1.863828in}{2.944464in}}%
\pgfpathlineto{\pgfqpoint{1.886487in}{2.935415in}}%
\pgfpathlineto{\pgfqpoint{1.909146in}{2.920552in}}%
\pgfpathlineto{\pgfqpoint{1.931804in}{2.913336in}}%
\pgfpathlineto{\pgfqpoint{1.954463in}{2.898968in}}%
\pgfpathlineto{\pgfqpoint{1.977122in}{2.895129in}}%
\pgfpathlineto{\pgfqpoint{1.999781in}{2.889295in}}%
\pgfpathlineto{\pgfqpoint{2.022439in}{2.878688in}}%
\pgfpathlineto{\pgfqpoint{2.045098in}{2.871734in}}%
\pgfpathlineto{\pgfqpoint{2.067757in}{2.877875in}}%
\pgfpathlineto{\pgfqpoint{2.090416in}{2.863572in}}%
\pgfpathlineto{\pgfqpoint{2.113074in}{2.852698in}}%
\pgfpathlineto{\pgfqpoint{2.135733in}{2.852095in}}%
\pgfpathlineto{\pgfqpoint{2.181051in}{2.845292in}}%
\pgfpathlineto{\pgfqpoint{2.226368in}{2.841208in}}%
\pgfpathlineto{\pgfqpoint{2.249027in}{2.825389in}}%
\pgfpathlineto{\pgfqpoint{2.271686in}{2.828179in}}%
\pgfpathlineto{\pgfqpoint{2.294344in}{2.821150in}}%
\pgfpathlineto{\pgfqpoint{2.317003in}{2.807831in}}%
\pgfpathlineto{\pgfqpoint{2.339662in}{2.813085in}}%
\pgfpathlineto{\pgfqpoint{2.362321in}{2.811163in}}%
\pgfpathlineto{\pgfqpoint{2.384979in}{2.797317in}}%
\pgfpathlineto{\pgfqpoint{2.407638in}{2.802267in}}%
\pgfpathlineto{\pgfqpoint{2.452956in}{2.798712in}}%
\pgfpathlineto{\pgfqpoint{2.475614in}{2.782616in}}%
\pgfpathlineto{\pgfqpoint{2.498273in}{2.779790in}}%
\pgfpathlineto{\pgfqpoint{2.520932in}{2.775547in}}%
\pgfpathlineto{\pgfqpoint{2.543591in}{2.782754in}}%
\pgfpathlineto{\pgfqpoint{2.566249in}{2.762909in}}%
\pgfpathlineto{\pgfqpoint{2.588908in}{2.767632in}}%
\pgfpathlineto{\pgfqpoint{2.611567in}{2.754834in}}%
\pgfpathlineto{\pgfqpoint{2.634226in}{2.753822in}}%
\pgfpathlineto{\pgfqpoint{2.679543in}{2.747506in}}%
\pgfpathlineto{\pgfqpoint{2.702202in}{2.751274in}}%
\pgfpathlineto{\pgfqpoint{2.724861in}{2.740340in}}%
\pgfpathlineto{\pgfqpoint{2.747519in}{2.730728in}}%
\pgfpathlineto{\pgfqpoint{2.770178in}{2.733855in}}%
\pgfpathlineto{\pgfqpoint{2.792837in}{2.739018in}}%
\pgfpathlineto{\pgfqpoint{2.815496in}{2.724484in}}%
\pgfpathlineto{\pgfqpoint{2.838154in}{2.721618in}}%
\pgfpathlineto{\pgfqpoint{2.860813in}{2.715842in}}%
\pgfpathlineto{\pgfqpoint{2.883472in}{2.714631in}}%
\pgfpathlineto{\pgfqpoint{2.906131in}{2.706501in}}%
\pgfpathlineto{\pgfqpoint{2.928789in}{2.702385in}}%
\pgfpathlineto{\pgfqpoint{2.951448in}{2.706052in}}%
\pgfpathlineto{\pgfqpoint{2.974107in}{2.696335in}}%
\pgfpathlineto{\pgfqpoint{2.996766in}{2.696479in}}%
\pgfpathlineto{\pgfqpoint{3.019424in}{2.686901in}}%
\pgfpathlineto{\pgfqpoint{3.042083in}{2.697805in}}%
\pgfpathlineto{\pgfqpoint{3.064742in}{2.684178in}}%
\pgfpathlineto{\pgfqpoint{3.087401in}{2.695313in}}%
\pgfpathlineto{\pgfqpoint{3.110059in}{2.676366in}}%
\pgfpathlineto{\pgfqpoint{3.132718in}{2.674505in}}%
\pgfpathlineto{\pgfqpoint{3.155377in}{2.676508in}}%
\pgfpathlineto{\pgfqpoint{3.178036in}{2.668445in}}%
\pgfpathlineto{\pgfqpoint{3.200694in}{2.666230in}}%
\pgfpathlineto{\pgfqpoint{3.223353in}{2.672912in}}%
\pgfpathlineto{\pgfqpoint{3.246012in}{2.661185in}}%
\pgfpathlineto{\pgfqpoint{3.268671in}{2.663352in}}%
\pgfpathlineto{\pgfqpoint{3.291329in}{2.654665in}}%
\pgfpathlineto{\pgfqpoint{3.313988in}{2.659538in}}%
\pgfpathlineto{\pgfqpoint{3.336647in}{2.656766in}}%
\pgfpathlineto{\pgfqpoint{3.359306in}{2.647396in}}%
\pgfpathlineto{\pgfqpoint{3.381964in}{2.647936in}}%
\pgfpathlineto{\pgfqpoint{3.404623in}{2.654444in}}%
\pgfpathlineto{\pgfqpoint{3.427282in}{2.640803in}}%
\pgfpathlineto{\pgfqpoint{3.449941in}{2.633367in}}%
\pgfpathlineto{\pgfqpoint{3.472599in}{2.629566in}}%
\pgfpathlineto{\pgfqpoint{3.495258in}{2.638317in}}%
\pgfpathlineto{\pgfqpoint{3.517917in}{2.627872in}}%
\pgfpathlineto{\pgfqpoint{3.540576in}{2.629947in}}%
\pgfpathlineto{\pgfqpoint{3.563234in}{2.624067in}}%
\pgfpathlineto{\pgfqpoint{3.585893in}{2.628006in}}%
\pgfpathlineto{\pgfqpoint{3.608552in}{2.623565in}}%
\pgfpathlineto{\pgfqpoint{3.631211in}{2.622053in}}%
\pgfpathlineto{\pgfqpoint{3.653869in}{2.623804in}}%
\pgfpathlineto{\pgfqpoint{3.676528in}{2.617532in}}%
\pgfpathlineto{\pgfqpoint{3.699187in}{2.607799in}}%
\pgfpathlineto{\pgfqpoint{3.721846in}{2.608609in}}%
\pgfpathlineto{\pgfqpoint{3.744504in}{2.605558in}}%
\pgfpathlineto{\pgfqpoint{3.767163in}{2.604937in}}%
\pgfpathlineto{\pgfqpoint{3.789822in}{2.602548in}}%
\pgfpathlineto{\pgfqpoint{3.812481in}{2.605674in}}%
\pgfpathlineto{\pgfqpoint{3.835139in}{2.602863in}}%
\pgfpathlineto{\pgfqpoint{3.857798in}{2.593216in}}%
\pgfpathlineto{\pgfqpoint{3.880457in}{2.593577in}}%
\pgfpathlineto{\pgfqpoint{3.903116in}{2.591785in}}%
\pgfpathlineto{\pgfqpoint{3.925774in}{2.587351in}}%
\pgfpathlineto{\pgfqpoint{3.948433in}{2.589263in}}%
\pgfpathlineto{\pgfqpoint{3.971092in}{2.581247in}}%
\pgfpathlineto{\pgfqpoint{3.993751in}{2.582373in}}%
\pgfpathlineto{\pgfqpoint{4.039068in}{2.574479in}}%
\pgfpathlineto{\pgfqpoint{4.129703in}{2.571706in}}%
\pgfpathlineto{\pgfqpoint{4.152362in}{2.565878in}}%
\pgfpathlineto{\pgfqpoint{4.175021in}{2.564404in}}%
\pgfpathlineto{\pgfqpoint{4.197679in}{2.564347in}}%
\pgfpathlineto{\pgfqpoint{4.220338in}{2.568041in}}%
\pgfpathlineto{\pgfqpoint{4.242997in}{2.564872in}}%
\pgfpathlineto{\pgfqpoint{4.265656in}{2.556021in}}%
\pgfpathlineto{\pgfqpoint{4.288314in}{2.554052in}}%
\pgfpathlineto{\pgfqpoint{4.310973in}{2.560496in}}%
\pgfpathlineto{\pgfqpoint{4.333632in}{2.549690in}}%
\pgfpathlineto{\pgfqpoint{4.356291in}{2.551999in}}%
\pgfpathlineto{\pgfqpoint{4.378949in}{2.549648in}}%
\pgfpathlineto{\pgfqpoint{4.424267in}{2.551312in}}%
\pgfpathlineto{\pgfqpoint{4.469584in}{2.543162in}}%
\pgfpathlineto{\pgfqpoint{4.514902in}{2.537283in}}%
\pgfpathlineto{\pgfqpoint{4.537561in}{2.528913in}}%
\pgfpathlineto{\pgfqpoint{4.560219in}{2.530855in}}%
\pgfpathlineto{\pgfqpoint{4.582878in}{2.526482in}}%
\pgfpathlineto{\pgfqpoint{4.605537in}{2.530241in}}%
\pgfpathlineto{\pgfqpoint{4.628196in}{2.519560in}}%
\pgfpathlineto{\pgfqpoint{4.650854in}{2.520141in}}%
\pgfpathlineto{\pgfqpoint{4.673513in}{2.527137in}}%
\pgfpathlineto{\pgfqpoint{4.696172in}{2.512131in}}%
\pgfpathlineto{\pgfqpoint{4.718831in}{2.516621in}}%
\pgfpathlineto{\pgfqpoint{4.741489in}{2.519386in}}%
\pgfpathlineto{\pgfqpoint{4.764148in}{2.507295in}}%
\pgfpathlineto{\pgfqpoint{4.786807in}{2.506304in}}%
\pgfpathlineto{\pgfqpoint{4.809466in}{2.509456in}}%
\pgfpathlineto{\pgfqpoint{4.832124in}{2.513855in}}%
\pgfpathlineto{\pgfqpoint{4.854783in}{2.502926in}}%
\pgfpathlineto{\pgfqpoint{4.877442in}{2.507880in}}%
\pgfpathlineto{\pgfqpoint{4.900101in}{2.499316in}}%
\pgfpathlineto{\pgfqpoint{4.922759in}{2.493448in}}%
\pgfpathlineto{\pgfqpoint{4.945418in}{2.497267in}}%
\pgfpathlineto{\pgfqpoint{4.968077in}{2.485870in}}%
\pgfpathlineto{\pgfqpoint{4.990735in}{2.493297in}}%
\pgfpathlineto{\pgfqpoint{5.013394in}{2.485265in}}%
\pgfpathlineto{\pgfqpoint{5.036053in}{2.494031in}}%
\pgfpathlineto{\pgfqpoint{5.058712in}{2.484141in}}%
\pgfpathlineto{\pgfqpoint{5.081370in}{2.477634in}}%
\pgfpathlineto{\pgfqpoint{5.104029in}{2.475178in}}%
\pgfpathlineto{\pgfqpoint{5.149347in}{2.474790in}}%
\pgfpathlineto{\pgfqpoint{5.172005in}{2.479357in}}%
\pgfpathlineto{\pgfqpoint{5.194664in}{2.472461in}}%
\pgfpathlineto{\pgfqpoint{5.217323in}{2.468417in}}%
\pgfpathlineto{\pgfqpoint{5.239982in}{2.460187in}}%
\pgfpathlineto{\pgfqpoint{5.262640in}{2.460647in}}%
\pgfpathlineto{\pgfqpoint{5.285299in}{2.458681in}}%
\pgfpathlineto{\pgfqpoint{5.307958in}{2.469136in}}%
\pgfpathlineto{\pgfqpoint{5.330617in}{2.468121in}}%
\pgfpathlineto{\pgfqpoint{5.353275in}{2.451171in}}%
\pgfpathlineto{\pgfqpoint{5.375934in}{2.449739in}}%
\pgfpathlineto{\pgfqpoint{5.398593in}{2.450383in}}%
\pgfpathlineto{\pgfqpoint{5.421252in}{2.447099in}}%
\pgfpathlineto{\pgfqpoint{5.466569in}{2.444488in}}%
\pgfpathlineto{\pgfqpoint{5.489228in}{2.439833in}}%
\pgfpathlineto{\pgfqpoint{5.511887in}{2.442482in}}%
\pgfpathlineto{\pgfqpoint{5.534545in}{2.439569in}}%
\pgfpathlineto{\pgfqpoint{5.534545in}{2.439569in}}%
\pgfusepath{stroke}%
\end{pgfscope}%
\begin{pgfscope}%
\pgfpathrectangle{\pgfqpoint{0.800000in}{0.528000in}}{\pgfqpoint{4.960000in}{3.696000in}}%
\pgfusepath{clip}%
\pgfsetrectcap%
\pgfsetroundjoin%
\pgfsetlinewidth{1.505625pt}%
\definecolor{currentstroke}{rgb}{0.737255,0.741176,0.133333}%
\pgfsetstrokecolor{currentstroke}%
\pgfsetdash{}{0pt}%
\pgfpathmoveto{\pgfqpoint{1.025455in}{3.283278in}}%
\pgfpathlineto{\pgfqpoint{1.048113in}{2.847525in}}%
\pgfpathlineto{\pgfqpoint{1.070772in}{2.716059in}}%
\pgfpathlineto{\pgfqpoint{1.093431in}{2.650960in}}%
\pgfpathlineto{\pgfqpoint{1.116090in}{2.596536in}}%
\pgfpathlineto{\pgfqpoint{1.138748in}{2.557386in}}%
\pgfpathlineto{\pgfqpoint{1.161407in}{2.544670in}}%
\pgfpathlineto{\pgfqpoint{1.184066in}{2.524332in}}%
\pgfpathlineto{\pgfqpoint{1.206725in}{2.479444in}}%
\pgfpathlineto{\pgfqpoint{1.229383in}{2.463523in}}%
\pgfpathlineto{\pgfqpoint{1.252042in}{2.421800in}}%
\pgfpathlineto{\pgfqpoint{1.274701in}{2.389650in}}%
\pgfpathlineto{\pgfqpoint{1.297360in}{2.367235in}}%
\pgfpathlineto{\pgfqpoint{1.320018in}{2.330924in}}%
\pgfpathlineto{\pgfqpoint{1.342677in}{2.303818in}}%
\pgfpathlineto{\pgfqpoint{1.365336in}{2.280942in}}%
\pgfpathlineto{\pgfqpoint{1.387995in}{2.238702in}}%
\pgfpathlineto{\pgfqpoint{1.410653in}{2.180089in}}%
\pgfpathlineto{\pgfqpoint{1.433312in}{2.133638in}}%
\pgfpathlineto{\pgfqpoint{1.455971in}{2.050368in}}%
\pgfpathlineto{\pgfqpoint{1.478630in}{1.959086in}}%
\pgfpathlineto{\pgfqpoint{1.523947in}{1.724715in}}%
\pgfpathlineto{\pgfqpoint{1.546606in}{1.592376in}}%
\pgfpathlineto{\pgfqpoint{1.569265in}{1.487222in}}%
\pgfpathlineto{\pgfqpoint{1.614582in}{1.305996in}}%
\pgfpathlineto{\pgfqpoint{1.637241in}{1.286332in}}%
\pgfpathlineto{\pgfqpoint{1.682558in}{1.181082in}}%
\pgfpathlineto{\pgfqpoint{1.727876in}{1.137144in}}%
\pgfpathlineto{\pgfqpoint{1.750534in}{1.145460in}}%
\pgfpathlineto{\pgfqpoint{1.773193in}{1.089209in}}%
\pgfpathlineto{\pgfqpoint{1.795852in}{1.068362in}}%
\pgfpathlineto{\pgfqpoint{1.818511in}{1.052811in}}%
\pgfpathlineto{\pgfqpoint{1.841169in}{1.031419in}}%
\pgfpathlineto{\pgfqpoint{1.863828in}{1.027703in}}%
\pgfpathlineto{\pgfqpoint{1.886487in}{1.007378in}}%
\pgfpathlineto{\pgfqpoint{1.909146in}{1.031891in}}%
\pgfpathlineto{\pgfqpoint{1.931804in}{0.996273in}}%
\pgfpathlineto{\pgfqpoint{1.977122in}{0.970743in}}%
\pgfpathlineto{\pgfqpoint{1.999781in}{0.964849in}}%
\pgfpathlineto{\pgfqpoint{2.022439in}{0.972221in}}%
\pgfpathlineto{\pgfqpoint{2.067757in}{0.941756in}}%
\pgfpathlineto{\pgfqpoint{2.090416in}{0.930231in}}%
\pgfpathlineto{\pgfqpoint{2.113074in}{0.960050in}}%
\pgfpathlineto{\pgfqpoint{2.135733in}{0.917300in}}%
\pgfpathlineto{\pgfqpoint{2.158392in}{0.928795in}}%
\pgfpathlineto{\pgfqpoint{2.181051in}{0.926906in}}%
\pgfpathlineto{\pgfqpoint{2.203709in}{0.896077in}}%
\pgfpathlineto{\pgfqpoint{2.226368in}{0.896796in}}%
\pgfpathlineto{\pgfqpoint{2.249027in}{0.899773in}}%
\pgfpathlineto{\pgfqpoint{2.271686in}{0.879455in}}%
\pgfpathlineto{\pgfqpoint{2.294344in}{0.876963in}}%
\pgfpathlineto{\pgfqpoint{2.317003in}{0.881221in}}%
\pgfpathlineto{\pgfqpoint{2.339662in}{0.877850in}}%
\pgfpathlineto{\pgfqpoint{2.362321in}{0.882627in}}%
\pgfpathlineto{\pgfqpoint{2.384979in}{0.864383in}}%
\pgfpathlineto{\pgfqpoint{2.407638in}{0.873609in}}%
\pgfpathlineto{\pgfqpoint{2.430297in}{0.868797in}}%
\pgfpathlineto{\pgfqpoint{2.452956in}{0.853333in}}%
\pgfpathlineto{\pgfqpoint{2.498273in}{0.861923in}}%
\pgfpathlineto{\pgfqpoint{2.520932in}{0.859378in}}%
\pgfpathlineto{\pgfqpoint{2.543591in}{0.851789in}}%
\pgfpathlineto{\pgfqpoint{2.566249in}{0.859509in}}%
\pgfpathlineto{\pgfqpoint{2.588908in}{0.849655in}}%
\pgfpathlineto{\pgfqpoint{2.611567in}{0.856212in}}%
\pgfpathlineto{\pgfqpoint{2.634226in}{0.855917in}}%
\pgfpathlineto{\pgfqpoint{2.656884in}{0.828744in}}%
\pgfpathlineto{\pgfqpoint{2.679543in}{0.858682in}}%
\pgfpathlineto{\pgfqpoint{2.702202in}{0.837766in}}%
\pgfpathlineto{\pgfqpoint{2.724861in}{0.830234in}}%
\pgfpathlineto{\pgfqpoint{2.747519in}{0.829954in}}%
\pgfpathlineto{\pgfqpoint{2.770178in}{0.817350in}}%
\pgfpathlineto{\pgfqpoint{2.792837in}{0.821844in}}%
\pgfpathlineto{\pgfqpoint{2.815496in}{0.833845in}}%
\pgfpathlineto{\pgfqpoint{2.838154in}{0.817950in}}%
\pgfpathlineto{\pgfqpoint{2.860813in}{0.805470in}}%
\pgfpathlineto{\pgfqpoint{2.883472in}{0.809981in}}%
\pgfpathlineto{\pgfqpoint{2.906131in}{0.812951in}}%
\pgfpathlineto{\pgfqpoint{2.928789in}{0.802265in}}%
\pgfpathlineto{\pgfqpoint{2.974107in}{0.796750in}}%
\pgfpathlineto{\pgfqpoint{2.996766in}{0.823661in}}%
\pgfpathlineto{\pgfqpoint{3.019424in}{0.808702in}}%
\pgfpathlineto{\pgfqpoint{3.042083in}{0.818270in}}%
\pgfpathlineto{\pgfqpoint{3.064742in}{0.806273in}}%
\pgfpathlineto{\pgfqpoint{3.087401in}{0.813258in}}%
\pgfpathlineto{\pgfqpoint{3.110059in}{0.826645in}}%
\pgfpathlineto{\pgfqpoint{3.132718in}{0.794239in}}%
\pgfpathlineto{\pgfqpoint{3.155377in}{0.798011in}}%
\pgfpathlineto{\pgfqpoint{3.178036in}{0.787927in}}%
\pgfpathlineto{\pgfqpoint{3.223353in}{0.796242in}}%
\pgfpathlineto{\pgfqpoint{3.246012in}{0.809314in}}%
\pgfpathlineto{\pgfqpoint{3.268671in}{0.785806in}}%
\pgfpathlineto{\pgfqpoint{3.291329in}{0.786155in}}%
\pgfpathlineto{\pgfqpoint{3.313988in}{0.801720in}}%
\pgfpathlineto{\pgfqpoint{3.336647in}{0.774943in}}%
\pgfpathlineto{\pgfqpoint{3.359306in}{0.791268in}}%
\pgfpathlineto{\pgfqpoint{3.381964in}{0.791802in}}%
\pgfpathlineto{\pgfqpoint{3.404623in}{0.781474in}}%
\pgfpathlineto{\pgfqpoint{3.427282in}{0.785057in}}%
\pgfpathlineto{\pgfqpoint{3.449941in}{0.799516in}}%
\pgfpathlineto{\pgfqpoint{3.472599in}{0.780423in}}%
\pgfpathlineto{\pgfqpoint{3.495258in}{0.785258in}}%
\pgfpathlineto{\pgfqpoint{3.517917in}{0.786237in}}%
\pgfpathlineto{\pgfqpoint{3.540576in}{0.783031in}}%
\pgfpathlineto{\pgfqpoint{3.563234in}{0.776171in}}%
\pgfpathlineto{\pgfqpoint{3.585893in}{0.809946in}}%
\pgfpathlineto{\pgfqpoint{3.608552in}{0.758652in}}%
\pgfpathlineto{\pgfqpoint{3.631211in}{0.770520in}}%
\pgfpathlineto{\pgfqpoint{3.653869in}{0.756256in}}%
\pgfpathlineto{\pgfqpoint{3.676528in}{0.758461in}}%
\pgfpathlineto{\pgfqpoint{3.699187in}{0.766851in}}%
\pgfpathlineto{\pgfqpoint{3.721846in}{0.767517in}}%
\pgfpathlineto{\pgfqpoint{3.744504in}{0.764938in}}%
\pgfpathlineto{\pgfqpoint{3.767163in}{0.760573in}}%
\pgfpathlineto{\pgfqpoint{3.789822in}{0.757405in}}%
\pgfpathlineto{\pgfqpoint{3.812481in}{0.764382in}}%
\pgfpathlineto{\pgfqpoint{3.835139in}{0.754483in}}%
\pgfpathlineto{\pgfqpoint{3.857798in}{0.757382in}}%
\pgfpathlineto{\pgfqpoint{3.880457in}{0.762438in}}%
\pgfpathlineto{\pgfqpoint{3.903116in}{0.760077in}}%
\pgfpathlineto{\pgfqpoint{3.925774in}{0.754115in}}%
\pgfpathlineto{\pgfqpoint{3.948433in}{0.756722in}}%
\pgfpathlineto{\pgfqpoint{3.971092in}{0.757477in}}%
\pgfpathlineto{\pgfqpoint{3.993751in}{0.760311in}}%
\pgfpathlineto{\pgfqpoint{4.016409in}{0.756242in}}%
\pgfpathlineto{\pgfqpoint{4.039068in}{0.760428in}}%
\pgfpathlineto{\pgfqpoint{4.061727in}{0.755420in}}%
\pgfpathlineto{\pgfqpoint{4.084386in}{0.755623in}}%
\pgfpathlineto{\pgfqpoint{4.107044in}{0.751528in}}%
\pgfpathlineto{\pgfqpoint{4.129703in}{0.752047in}}%
\pgfpathlineto{\pgfqpoint{4.152362in}{0.754232in}}%
\pgfpathlineto{\pgfqpoint{4.220338in}{0.750927in}}%
\pgfpathlineto{\pgfqpoint{4.265656in}{0.752048in}}%
\pgfpathlineto{\pgfqpoint{4.288314in}{0.751042in}}%
\pgfpathlineto{\pgfqpoint{4.333632in}{0.752659in}}%
\pgfpathlineto{\pgfqpoint{4.356291in}{0.750612in}}%
\pgfpathlineto{\pgfqpoint{4.378949in}{0.753179in}}%
\pgfpathlineto{\pgfqpoint{4.401608in}{0.751082in}}%
\pgfpathlineto{\pgfqpoint{4.424267in}{0.755953in}}%
\pgfpathlineto{\pgfqpoint{4.446926in}{0.750328in}}%
\pgfpathlineto{\pgfqpoint{4.469584in}{0.752915in}}%
\pgfpathlineto{\pgfqpoint{4.492243in}{0.751015in}}%
\pgfpathlineto{\pgfqpoint{4.514902in}{0.753491in}}%
\pgfpathlineto{\pgfqpoint{4.537561in}{0.751388in}}%
\pgfpathlineto{\pgfqpoint{4.605537in}{0.752857in}}%
\pgfpathlineto{\pgfqpoint{4.628196in}{0.750366in}}%
\pgfpathlineto{\pgfqpoint{4.650854in}{0.751966in}}%
\pgfpathlineto{\pgfqpoint{4.673513in}{0.750924in}}%
\pgfpathlineto{\pgfqpoint{4.696172in}{0.750997in}}%
\pgfpathlineto{\pgfqpoint{4.718831in}{0.749939in}}%
\pgfpathlineto{\pgfqpoint{4.809466in}{0.750518in}}%
\pgfpathlineto{\pgfqpoint{4.877442in}{0.750255in}}%
\pgfpathlineto{\pgfqpoint{4.900101in}{0.750878in}}%
\pgfpathlineto{\pgfqpoint{4.945418in}{0.750356in}}%
\pgfpathlineto{\pgfqpoint{5.534545in}{0.750418in}}%
\pgfpathlineto{\pgfqpoint{5.534545in}{0.750418in}}%
\pgfusepath{stroke}%
\end{pgfscope}%
\begin{pgfscope}%
\pgfpathrectangle{\pgfqpoint{0.800000in}{0.528000in}}{\pgfqpoint{4.960000in}{3.696000in}}%
\pgfusepath{clip}%
\pgfsetrectcap%
\pgfsetroundjoin%
\pgfsetlinewidth{1.505625pt}%
\definecolor{currentstroke}{rgb}{0.090196,0.745098,0.811765}%
\pgfsetstrokecolor{currentstroke}%
\pgfsetdash{}{0pt}%
\pgfpathmoveto{\pgfqpoint{1.025455in}{3.721129in}}%
\pgfpathlineto{\pgfqpoint{1.048113in}{3.439800in}}%
\pgfpathlineto{\pgfqpoint{1.070772in}{3.218311in}}%
\pgfpathlineto{\pgfqpoint{1.093431in}{3.129574in}}%
\pgfpathlineto{\pgfqpoint{1.116090in}{3.051858in}}%
\pgfpathlineto{\pgfqpoint{1.138748in}{3.016580in}}%
\pgfpathlineto{\pgfqpoint{1.161407in}{2.977465in}}%
\pgfpathlineto{\pgfqpoint{1.184066in}{2.942533in}}%
\pgfpathlineto{\pgfqpoint{1.206725in}{2.920515in}}%
\pgfpathlineto{\pgfqpoint{1.229383in}{2.884322in}}%
\pgfpathlineto{\pgfqpoint{1.252042in}{2.852048in}}%
\pgfpathlineto{\pgfqpoint{1.274701in}{2.827566in}}%
\pgfpathlineto{\pgfqpoint{1.297360in}{2.789390in}}%
\pgfpathlineto{\pgfqpoint{1.320018in}{2.754045in}}%
\pgfpathlineto{\pgfqpoint{1.342677in}{2.705243in}}%
\pgfpathlineto{\pgfqpoint{1.365336in}{2.631638in}}%
\pgfpathlineto{\pgfqpoint{1.387995in}{2.526269in}}%
\pgfpathlineto{\pgfqpoint{1.410653in}{2.393498in}}%
\pgfpathlineto{\pgfqpoint{1.433312in}{2.229458in}}%
\pgfpathlineto{\pgfqpoint{1.455971in}{2.074747in}}%
\pgfpathlineto{\pgfqpoint{1.478630in}{1.906353in}}%
\pgfpathlineto{\pgfqpoint{1.501288in}{1.770917in}}%
\pgfpathlineto{\pgfqpoint{1.523947in}{1.673765in}}%
\pgfpathlineto{\pgfqpoint{1.546606in}{1.606736in}}%
\pgfpathlineto{\pgfqpoint{1.569265in}{1.554694in}}%
\pgfpathlineto{\pgfqpoint{1.591923in}{1.515115in}}%
\pgfpathlineto{\pgfqpoint{1.614582in}{1.480449in}}%
\pgfpathlineto{\pgfqpoint{1.637241in}{1.455973in}}%
\pgfpathlineto{\pgfqpoint{1.659899in}{1.465862in}}%
\pgfpathlineto{\pgfqpoint{1.682558in}{1.428093in}}%
\pgfpathlineto{\pgfqpoint{1.705217in}{1.436263in}}%
\pgfpathlineto{\pgfqpoint{1.727876in}{1.410020in}}%
\pgfpathlineto{\pgfqpoint{1.750534in}{1.402271in}}%
\pgfpathlineto{\pgfqpoint{1.773193in}{1.386239in}}%
\pgfpathlineto{\pgfqpoint{1.795852in}{1.392105in}}%
\pgfpathlineto{\pgfqpoint{1.818511in}{1.379020in}}%
\pgfpathlineto{\pgfqpoint{1.841169in}{1.359815in}}%
\pgfpathlineto{\pgfqpoint{1.863828in}{1.348731in}}%
\pgfpathlineto{\pgfqpoint{1.886487in}{1.369217in}}%
\pgfpathlineto{\pgfqpoint{1.909146in}{1.347212in}}%
\pgfpathlineto{\pgfqpoint{1.931804in}{1.346290in}}%
\pgfpathlineto{\pgfqpoint{1.954463in}{1.338766in}}%
\pgfpathlineto{\pgfqpoint{1.999781in}{1.335749in}}%
\pgfpathlineto{\pgfqpoint{2.022439in}{1.337875in}}%
\pgfpathlineto{\pgfqpoint{2.045098in}{1.326538in}}%
\pgfpathlineto{\pgfqpoint{2.067757in}{1.332521in}}%
\pgfpathlineto{\pgfqpoint{2.090416in}{1.320565in}}%
\pgfpathlineto{\pgfqpoint{2.113074in}{1.328470in}}%
\pgfpathlineto{\pgfqpoint{2.135733in}{1.311906in}}%
\pgfpathlineto{\pgfqpoint{2.158392in}{1.317829in}}%
\pgfpathlineto{\pgfqpoint{2.181051in}{1.317482in}}%
\pgfpathlineto{\pgfqpoint{2.203709in}{1.297018in}}%
\pgfpathlineto{\pgfqpoint{2.226368in}{1.290718in}}%
\pgfpathlineto{\pgfqpoint{2.249027in}{1.309266in}}%
\pgfpathlineto{\pgfqpoint{2.271686in}{1.339874in}}%
\pgfpathlineto{\pgfqpoint{2.294344in}{1.273950in}}%
\pgfpathlineto{\pgfqpoint{2.317003in}{1.298594in}}%
\pgfpathlineto{\pgfqpoint{2.339662in}{1.280474in}}%
\pgfpathlineto{\pgfqpoint{2.362321in}{1.275048in}}%
\pgfpathlineto{\pgfqpoint{2.384979in}{1.257260in}}%
\pgfpathlineto{\pgfqpoint{2.407638in}{1.269397in}}%
\pgfpathlineto{\pgfqpoint{2.430297in}{1.263236in}}%
\pgfpathlineto{\pgfqpoint{2.452956in}{1.267747in}}%
\pgfpathlineto{\pgfqpoint{2.475614in}{1.248319in}}%
\pgfpathlineto{\pgfqpoint{2.498273in}{1.244010in}}%
\pgfpathlineto{\pgfqpoint{2.520932in}{1.245772in}}%
\pgfpathlineto{\pgfqpoint{2.543591in}{1.274869in}}%
\pgfpathlineto{\pgfqpoint{2.566249in}{1.247749in}}%
\pgfpathlineto{\pgfqpoint{2.588908in}{1.229551in}}%
\pgfpathlineto{\pgfqpoint{2.611567in}{1.228631in}}%
\pgfpathlineto{\pgfqpoint{2.634226in}{1.246684in}}%
\pgfpathlineto{\pgfqpoint{2.656884in}{1.212776in}}%
\pgfpathlineto{\pgfqpoint{2.679543in}{1.204421in}}%
\pgfpathlineto{\pgfqpoint{2.702202in}{1.206106in}}%
\pgfpathlineto{\pgfqpoint{2.724861in}{1.190290in}}%
\pgfpathlineto{\pgfqpoint{2.747519in}{1.186746in}}%
\pgfpathlineto{\pgfqpoint{2.770178in}{1.198263in}}%
\pgfpathlineto{\pgfqpoint{2.792837in}{1.181997in}}%
\pgfpathlineto{\pgfqpoint{2.815496in}{1.198410in}}%
\pgfpathlineto{\pgfqpoint{2.838154in}{1.175146in}}%
\pgfpathlineto{\pgfqpoint{2.860813in}{1.164504in}}%
\pgfpathlineto{\pgfqpoint{2.883472in}{1.170343in}}%
\pgfpathlineto{\pgfqpoint{2.906131in}{1.150620in}}%
\pgfpathlineto{\pgfqpoint{2.928789in}{1.147944in}}%
\pgfpathlineto{\pgfqpoint{2.951448in}{1.156272in}}%
\pgfpathlineto{\pgfqpoint{2.974107in}{1.130489in}}%
\pgfpathlineto{\pgfqpoint{2.996766in}{1.143541in}}%
\pgfpathlineto{\pgfqpoint{3.019424in}{1.153576in}}%
\pgfpathlineto{\pgfqpoint{3.042083in}{1.124686in}}%
\pgfpathlineto{\pgfqpoint{3.064742in}{1.124024in}}%
\pgfpathlineto{\pgfqpoint{3.087401in}{1.128589in}}%
\pgfpathlineto{\pgfqpoint{3.110059in}{1.128075in}}%
\pgfpathlineto{\pgfqpoint{3.132718in}{1.121029in}}%
\pgfpathlineto{\pgfqpoint{3.155377in}{1.122243in}}%
\pgfpathlineto{\pgfqpoint{3.178036in}{1.129895in}}%
\pgfpathlineto{\pgfqpoint{3.200694in}{1.115573in}}%
\pgfpathlineto{\pgfqpoint{3.223353in}{1.106694in}}%
\pgfpathlineto{\pgfqpoint{3.246012in}{1.093416in}}%
\pgfpathlineto{\pgfqpoint{3.268671in}{1.124816in}}%
\pgfpathlineto{\pgfqpoint{3.291329in}{1.096203in}}%
\pgfpathlineto{\pgfqpoint{3.313988in}{1.095974in}}%
\pgfpathlineto{\pgfqpoint{3.336647in}{1.104626in}}%
\pgfpathlineto{\pgfqpoint{3.359306in}{1.093496in}}%
\pgfpathlineto{\pgfqpoint{3.381964in}{1.105073in}}%
\pgfpathlineto{\pgfqpoint{3.404623in}{1.079446in}}%
\pgfpathlineto{\pgfqpoint{3.427282in}{1.089571in}}%
\pgfpathlineto{\pgfqpoint{3.449941in}{1.082245in}}%
\pgfpathlineto{\pgfqpoint{3.472599in}{1.060268in}}%
\pgfpathlineto{\pgfqpoint{3.495258in}{1.070508in}}%
\pgfpathlineto{\pgfqpoint{3.517917in}{1.065425in}}%
\pgfpathlineto{\pgfqpoint{3.540576in}{1.058213in}}%
\pgfpathlineto{\pgfqpoint{3.563234in}{1.059357in}}%
\pgfpathlineto{\pgfqpoint{3.585893in}{1.051366in}}%
\pgfpathlineto{\pgfqpoint{3.631211in}{1.052449in}}%
\pgfpathlineto{\pgfqpoint{3.653869in}{1.049295in}}%
\pgfpathlineto{\pgfqpoint{3.676528in}{1.042556in}}%
\pgfpathlineto{\pgfqpoint{3.699187in}{1.063474in}}%
\pgfpathlineto{\pgfqpoint{3.721846in}{1.050866in}}%
\pgfpathlineto{\pgfqpoint{3.744504in}{1.055211in}}%
\pgfpathlineto{\pgfqpoint{3.767163in}{1.043468in}}%
\pgfpathlineto{\pgfqpoint{3.789822in}{1.039939in}}%
\pgfpathlineto{\pgfqpoint{3.812481in}{1.044533in}}%
\pgfpathlineto{\pgfqpoint{3.835139in}{1.034378in}}%
\pgfpathlineto{\pgfqpoint{3.857798in}{1.038270in}}%
\pgfpathlineto{\pgfqpoint{3.880457in}{1.025425in}}%
\pgfpathlineto{\pgfqpoint{3.903116in}{1.027695in}}%
\pgfpathlineto{\pgfqpoint{3.925774in}{1.031289in}}%
\pgfpathlineto{\pgfqpoint{3.948433in}{1.028705in}}%
\pgfpathlineto{\pgfqpoint{3.971092in}{1.027568in}}%
\pgfpathlineto{\pgfqpoint{3.993751in}{1.021708in}}%
\pgfpathlineto{\pgfqpoint{4.016409in}{1.043621in}}%
\pgfpathlineto{\pgfqpoint{4.039068in}{1.041781in}}%
\pgfpathlineto{\pgfqpoint{4.061727in}{1.023535in}}%
\pgfpathlineto{\pgfqpoint{4.084386in}{1.027146in}}%
\pgfpathlineto{\pgfqpoint{4.107044in}{1.022232in}}%
\pgfpathlineto{\pgfqpoint{4.129703in}{1.031927in}}%
\pgfpathlineto{\pgfqpoint{4.152362in}{1.012144in}}%
\pgfpathlineto{\pgfqpoint{4.175021in}{1.012383in}}%
\pgfpathlineto{\pgfqpoint{4.197679in}{1.026134in}}%
\pgfpathlineto{\pgfqpoint{4.220338in}{1.011029in}}%
\pgfpathlineto{\pgfqpoint{4.242997in}{0.998922in}}%
\pgfpathlineto{\pgfqpoint{4.265656in}{1.037719in}}%
\pgfpathlineto{\pgfqpoint{4.288314in}{1.005146in}}%
\pgfpathlineto{\pgfqpoint{4.310973in}{1.005031in}}%
\pgfpathlineto{\pgfqpoint{4.333632in}{1.027444in}}%
\pgfpathlineto{\pgfqpoint{4.356291in}{1.002147in}}%
\pgfpathlineto{\pgfqpoint{4.378949in}{0.992786in}}%
\pgfpathlineto{\pgfqpoint{4.401608in}{1.007968in}}%
\pgfpathlineto{\pgfqpoint{4.424267in}{0.989815in}}%
\pgfpathlineto{\pgfqpoint{4.446926in}{1.001776in}}%
\pgfpathlineto{\pgfqpoint{4.469584in}{1.004256in}}%
\pgfpathlineto{\pgfqpoint{4.492243in}{1.021318in}}%
\pgfpathlineto{\pgfqpoint{4.514902in}{0.991834in}}%
\pgfpathlineto{\pgfqpoint{4.537561in}{0.991459in}}%
\pgfpathlineto{\pgfqpoint{4.560219in}{0.993670in}}%
\pgfpathlineto{\pgfqpoint{4.582878in}{0.985760in}}%
\pgfpathlineto{\pgfqpoint{4.605537in}{0.990790in}}%
\pgfpathlineto{\pgfqpoint{4.628196in}{0.999437in}}%
\pgfpathlineto{\pgfqpoint{4.673513in}{1.001117in}}%
\pgfpathlineto{\pgfqpoint{4.696172in}{1.000885in}}%
\pgfpathlineto{\pgfqpoint{4.718831in}{0.980596in}}%
\pgfpathlineto{\pgfqpoint{4.741489in}{1.007593in}}%
\pgfpathlineto{\pgfqpoint{4.764148in}{0.989763in}}%
\pgfpathlineto{\pgfqpoint{4.786807in}{0.995728in}}%
\pgfpathlineto{\pgfqpoint{4.809466in}{0.965999in}}%
\pgfpathlineto{\pgfqpoint{4.832124in}{0.982323in}}%
\pgfpathlineto{\pgfqpoint{4.854783in}{0.974497in}}%
\pgfpathlineto{\pgfqpoint{4.877442in}{0.986324in}}%
\pgfpathlineto{\pgfqpoint{4.900101in}{0.974475in}}%
\pgfpathlineto{\pgfqpoint{4.922759in}{0.970353in}}%
\pgfpathlineto{\pgfqpoint{4.945418in}{0.972562in}}%
\pgfpathlineto{\pgfqpoint{4.968077in}{0.986500in}}%
\pgfpathlineto{\pgfqpoint{4.990735in}{0.990564in}}%
\pgfpathlineto{\pgfqpoint{5.013394in}{0.973139in}}%
\pgfpathlineto{\pgfqpoint{5.036053in}{0.985978in}}%
\pgfpathlineto{\pgfqpoint{5.058712in}{0.977638in}}%
\pgfpathlineto{\pgfqpoint{5.081370in}{0.956681in}}%
\pgfpathlineto{\pgfqpoint{5.104029in}{0.952847in}}%
\pgfpathlineto{\pgfqpoint{5.126688in}{0.955643in}}%
\pgfpathlineto{\pgfqpoint{5.149347in}{0.952461in}}%
\pgfpathlineto{\pgfqpoint{5.172005in}{0.965782in}}%
\pgfpathlineto{\pgfqpoint{5.194664in}{0.952009in}}%
\pgfpathlineto{\pgfqpoint{5.217323in}{0.959763in}}%
\pgfpathlineto{\pgfqpoint{5.239982in}{0.956431in}}%
\pgfpathlineto{\pgfqpoint{5.262640in}{0.951730in}}%
\pgfpathlineto{\pgfqpoint{5.285299in}{0.955745in}}%
\pgfpathlineto{\pgfqpoint{5.307958in}{0.958623in}}%
\pgfpathlineto{\pgfqpoint{5.330617in}{0.954395in}}%
\pgfpathlineto{\pgfqpoint{5.353275in}{0.955142in}}%
\pgfpathlineto{\pgfqpoint{5.375934in}{0.957670in}}%
\pgfpathlineto{\pgfqpoint{5.398593in}{0.953394in}}%
\pgfpathlineto{\pgfqpoint{5.421252in}{0.961144in}}%
\pgfpathlineto{\pgfqpoint{5.466569in}{0.951982in}}%
\pgfpathlineto{\pgfqpoint{5.511887in}{0.951522in}}%
\pgfpathlineto{\pgfqpoint{5.534545in}{0.954655in}}%
\pgfpathlineto{\pgfqpoint{5.534545in}{0.954655in}}%
\pgfusepath{stroke}%
\end{pgfscope}%
\begin{pgfscope}%
\pgfpathrectangle{\pgfqpoint{0.800000in}{0.528000in}}{\pgfqpoint{4.960000in}{3.696000in}}%
\pgfusepath{clip}%
\pgfsetrectcap%
\pgfsetroundjoin%
\pgfsetlinewidth{1.505625pt}%
\definecolor{currentstroke}{rgb}{0.121569,0.466667,0.705882}%
\pgfsetstrokecolor{currentstroke}%
\pgfsetdash{}{0pt}%
\pgfpathmoveto{\pgfqpoint{1.025455in}{4.040654in}}%
\pgfpathlineto{\pgfqpoint{1.048113in}{3.998832in}}%
\pgfpathlineto{\pgfqpoint{1.070772in}{3.903703in}}%
\pgfpathlineto{\pgfqpoint{1.093431in}{3.753504in}}%
\pgfpathlineto{\pgfqpoint{1.116090in}{3.577085in}}%
\pgfpathlineto{\pgfqpoint{1.138748in}{3.478681in}}%
\pgfpathlineto{\pgfqpoint{1.161407in}{3.423929in}}%
\pgfpathlineto{\pgfqpoint{1.184066in}{3.383350in}}%
\pgfpathlineto{\pgfqpoint{1.206725in}{3.360698in}}%
\pgfpathlineto{\pgfqpoint{1.252042in}{3.318487in}}%
\pgfpathlineto{\pgfqpoint{1.274701in}{3.303742in}}%
\pgfpathlineto{\pgfqpoint{1.297360in}{3.286510in}}%
\pgfpathlineto{\pgfqpoint{1.320018in}{3.270947in}}%
\pgfpathlineto{\pgfqpoint{1.342677in}{3.259206in}}%
\pgfpathlineto{\pgfqpoint{1.365336in}{3.244532in}}%
\pgfpathlineto{\pgfqpoint{1.410653in}{3.229091in}}%
\pgfpathlineto{\pgfqpoint{1.433312in}{3.205344in}}%
\pgfpathlineto{\pgfqpoint{1.478630in}{3.194055in}}%
\pgfpathlineto{\pgfqpoint{1.501288in}{3.174448in}}%
\pgfpathlineto{\pgfqpoint{1.523947in}{3.173748in}}%
\pgfpathlineto{\pgfqpoint{1.546606in}{3.156348in}}%
\pgfpathlineto{\pgfqpoint{1.569265in}{3.145766in}}%
\pgfpathlineto{\pgfqpoint{1.591923in}{3.130501in}}%
\pgfpathlineto{\pgfqpoint{1.614582in}{3.121529in}}%
\pgfpathlineto{\pgfqpoint{1.637241in}{3.106838in}}%
\pgfpathlineto{\pgfqpoint{1.659899in}{3.096158in}}%
\pgfpathlineto{\pgfqpoint{1.682558in}{3.089515in}}%
\pgfpathlineto{\pgfqpoint{1.705217in}{3.079406in}}%
\pgfpathlineto{\pgfqpoint{1.727876in}{3.077184in}}%
\pgfpathlineto{\pgfqpoint{1.750534in}{3.063338in}}%
\pgfpathlineto{\pgfqpoint{1.773193in}{3.060216in}}%
\pgfpathlineto{\pgfqpoint{1.795852in}{3.042184in}}%
\pgfpathlineto{\pgfqpoint{1.818511in}{3.040744in}}%
\pgfpathlineto{\pgfqpoint{1.841169in}{3.035608in}}%
\pgfpathlineto{\pgfqpoint{1.863828in}{3.018827in}}%
\pgfpathlineto{\pgfqpoint{1.886487in}{3.018131in}}%
\pgfpathlineto{\pgfqpoint{1.909146in}{2.996505in}}%
\pgfpathlineto{\pgfqpoint{1.954463in}{2.988554in}}%
\pgfpathlineto{\pgfqpoint{1.977122in}{2.977978in}}%
\pgfpathlineto{\pgfqpoint{1.999781in}{2.978121in}}%
\pgfpathlineto{\pgfqpoint{2.022439in}{2.967154in}}%
\pgfpathlineto{\pgfqpoint{2.045098in}{2.974195in}}%
\pgfpathlineto{\pgfqpoint{2.067757in}{2.958820in}}%
\pgfpathlineto{\pgfqpoint{2.090416in}{2.961536in}}%
\pgfpathlineto{\pgfqpoint{2.113074in}{2.945999in}}%
\pgfpathlineto{\pgfqpoint{2.135733in}{2.943248in}}%
\pgfpathlineto{\pgfqpoint{2.158392in}{2.936693in}}%
\pgfpathlineto{\pgfqpoint{2.181051in}{2.934742in}}%
\pgfpathlineto{\pgfqpoint{2.203709in}{2.923317in}}%
\pgfpathlineto{\pgfqpoint{2.226368in}{2.927968in}}%
\pgfpathlineto{\pgfqpoint{2.249027in}{2.916389in}}%
\pgfpathlineto{\pgfqpoint{2.271686in}{2.909437in}}%
\pgfpathlineto{\pgfqpoint{2.294344in}{2.911930in}}%
\pgfpathlineto{\pgfqpoint{2.317003in}{2.908114in}}%
\pgfpathlineto{\pgfqpoint{2.339662in}{2.913879in}}%
\pgfpathlineto{\pgfqpoint{2.362321in}{2.893167in}}%
\pgfpathlineto{\pgfqpoint{2.384979in}{2.899316in}}%
\pgfpathlineto{\pgfqpoint{2.407638in}{2.883855in}}%
\pgfpathlineto{\pgfqpoint{2.430297in}{2.880927in}}%
\pgfpathlineto{\pgfqpoint{2.452956in}{2.903784in}}%
\pgfpathlineto{\pgfqpoint{2.475614in}{2.873014in}}%
\pgfpathlineto{\pgfqpoint{2.498273in}{2.884392in}}%
\pgfpathlineto{\pgfqpoint{2.520932in}{2.863927in}}%
\pgfpathlineto{\pgfqpoint{2.543591in}{2.859034in}}%
\pgfpathlineto{\pgfqpoint{2.566249in}{2.859769in}}%
\pgfpathlineto{\pgfqpoint{2.588908in}{2.856988in}}%
\pgfpathlineto{\pgfqpoint{2.611567in}{2.866682in}}%
\pgfpathlineto{\pgfqpoint{2.634226in}{2.839984in}}%
\pgfpathlineto{\pgfqpoint{2.656884in}{2.838244in}}%
\pgfpathlineto{\pgfqpoint{2.702202in}{2.840844in}}%
\pgfpathlineto{\pgfqpoint{2.747519in}{2.833966in}}%
\pgfpathlineto{\pgfqpoint{2.770178in}{2.819630in}}%
\pgfpathlineto{\pgfqpoint{2.792837in}{2.832206in}}%
\pgfpathlineto{\pgfqpoint{2.815496in}{2.816750in}}%
\pgfpathlineto{\pgfqpoint{2.838154in}{2.808016in}}%
\pgfpathlineto{\pgfqpoint{2.860813in}{2.814992in}}%
\pgfpathlineto{\pgfqpoint{2.883472in}{2.793373in}}%
\pgfpathlineto{\pgfqpoint{2.906131in}{2.793327in}}%
\pgfpathlineto{\pgfqpoint{2.928789in}{2.797029in}}%
\pgfpathlineto{\pgfqpoint{2.951448in}{2.785130in}}%
\pgfpathlineto{\pgfqpoint{2.974107in}{2.780877in}}%
\pgfpathlineto{\pgfqpoint{2.996766in}{2.798554in}}%
\pgfpathlineto{\pgfqpoint{3.019424in}{2.772210in}}%
\pgfpathlineto{\pgfqpoint{3.042083in}{2.770542in}}%
\pgfpathlineto{\pgfqpoint{3.064742in}{2.773951in}}%
\pgfpathlineto{\pgfqpoint{3.087401in}{2.760030in}}%
\pgfpathlineto{\pgfqpoint{3.110059in}{2.765841in}}%
\pgfpathlineto{\pgfqpoint{3.132718in}{2.754470in}}%
\pgfpathlineto{\pgfqpoint{3.155377in}{2.750337in}}%
\pgfpathlineto{\pgfqpoint{3.178036in}{2.771466in}}%
\pgfpathlineto{\pgfqpoint{3.200694in}{2.752456in}}%
\pgfpathlineto{\pgfqpoint{3.223353in}{2.756605in}}%
\pgfpathlineto{\pgfqpoint{3.246012in}{2.747773in}}%
\pgfpathlineto{\pgfqpoint{3.268671in}{2.745497in}}%
\pgfpathlineto{\pgfqpoint{3.291329in}{2.747638in}}%
\pgfpathlineto{\pgfqpoint{3.313988in}{2.736845in}}%
\pgfpathlineto{\pgfqpoint{3.336647in}{2.744529in}}%
\pgfpathlineto{\pgfqpoint{3.359306in}{2.731337in}}%
\pgfpathlineto{\pgfqpoint{3.381964in}{2.725645in}}%
\pgfpathlineto{\pgfqpoint{3.404623in}{2.761014in}}%
\pgfpathlineto{\pgfqpoint{3.427282in}{2.725620in}}%
\pgfpathlineto{\pgfqpoint{3.449941in}{2.725790in}}%
\pgfpathlineto{\pgfqpoint{3.472599in}{2.723834in}}%
\pgfpathlineto{\pgfqpoint{3.495258in}{2.727393in}}%
\pgfpathlineto{\pgfqpoint{3.517917in}{2.726186in}}%
\pgfpathlineto{\pgfqpoint{3.540576in}{2.710252in}}%
\pgfpathlineto{\pgfqpoint{3.563234in}{2.737632in}}%
\pgfpathlineto{\pgfqpoint{3.585893in}{2.708340in}}%
\pgfpathlineto{\pgfqpoint{3.608552in}{2.715160in}}%
\pgfpathlineto{\pgfqpoint{3.631211in}{2.704125in}}%
\pgfpathlineto{\pgfqpoint{3.653869in}{2.700879in}}%
\pgfpathlineto{\pgfqpoint{3.676528in}{2.693531in}}%
\pgfpathlineto{\pgfqpoint{3.699187in}{2.702392in}}%
\pgfpathlineto{\pgfqpoint{3.721846in}{2.696278in}}%
\pgfpathlineto{\pgfqpoint{3.744504in}{2.708395in}}%
\pgfpathlineto{\pgfqpoint{3.767163in}{2.694232in}}%
\pgfpathlineto{\pgfqpoint{3.789822in}{2.693219in}}%
\pgfpathlineto{\pgfqpoint{3.812481in}{2.682535in}}%
\pgfpathlineto{\pgfqpoint{3.835139in}{2.687555in}}%
\pgfpathlineto{\pgfqpoint{3.857798in}{2.683393in}}%
\pgfpathlineto{\pgfqpoint{3.903116in}{2.670940in}}%
\pgfpathlineto{\pgfqpoint{3.925774in}{2.675590in}}%
\pgfpathlineto{\pgfqpoint{3.948433in}{2.674302in}}%
\pgfpathlineto{\pgfqpoint{3.971092in}{2.665768in}}%
\pgfpathlineto{\pgfqpoint{3.993751in}{2.679328in}}%
\pgfpathlineto{\pgfqpoint{4.016409in}{2.673516in}}%
\pgfpathlineto{\pgfqpoint{4.039068in}{2.665281in}}%
\pgfpathlineto{\pgfqpoint{4.061727in}{2.663496in}}%
\pgfpathlineto{\pgfqpoint{4.084386in}{2.651949in}}%
\pgfpathlineto{\pgfqpoint{4.107044in}{2.666633in}}%
\pgfpathlineto{\pgfqpoint{4.129703in}{2.648977in}}%
\pgfpathlineto{\pgfqpoint{4.152362in}{2.662672in}}%
\pgfpathlineto{\pgfqpoint{4.175021in}{2.648175in}}%
\pgfpathlineto{\pgfqpoint{4.197679in}{2.664343in}}%
\pgfpathlineto{\pgfqpoint{4.220338in}{2.644890in}}%
\pgfpathlineto{\pgfqpoint{4.265656in}{2.638203in}}%
\pgfpathlineto{\pgfqpoint{4.288314in}{2.638690in}}%
\pgfpathlineto{\pgfqpoint{4.310973in}{2.635817in}}%
\pgfpathlineto{\pgfqpoint{4.333632in}{2.638472in}}%
\pgfpathlineto{\pgfqpoint{4.356291in}{2.636406in}}%
\pgfpathlineto{\pgfqpoint{4.378949in}{2.644422in}}%
\pgfpathlineto{\pgfqpoint{4.401608in}{2.623119in}}%
\pgfpathlineto{\pgfqpoint{4.424267in}{2.626825in}}%
\pgfpathlineto{\pgfqpoint{4.446926in}{2.636559in}}%
\pgfpathlineto{\pgfqpoint{4.469584in}{2.623529in}}%
\pgfpathlineto{\pgfqpoint{4.492243in}{2.633190in}}%
\pgfpathlineto{\pgfqpoint{4.514902in}{2.610911in}}%
\pgfpathlineto{\pgfqpoint{4.537561in}{2.624517in}}%
\pgfpathlineto{\pgfqpoint{4.560219in}{2.612490in}}%
\pgfpathlineto{\pgfqpoint{4.582878in}{2.607483in}}%
\pgfpathlineto{\pgfqpoint{4.605537in}{2.611291in}}%
\pgfpathlineto{\pgfqpoint{4.650854in}{2.603114in}}%
\pgfpathlineto{\pgfqpoint{4.673513in}{2.597500in}}%
\pgfpathlineto{\pgfqpoint{4.718831in}{2.594776in}}%
\pgfpathlineto{\pgfqpoint{4.741489in}{2.594965in}}%
\pgfpathlineto{\pgfqpoint{4.764148in}{2.586175in}}%
\pgfpathlineto{\pgfqpoint{4.786807in}{2.578859in}}%
\pgfpathlineto{\pgfqpoint{4.809466in}{2.582533in}}%
\pgfpathlineto{\pgfqpoint{4.832124in}{2.578628in}}%
\pgfpathlineto{\pgfqpoint{4.854783in}{2.570620in}}%
\pgfpathlineto{\pgfqpoint{4.877442in}{2.572554in}}%
\pgfpathlineto{\pgfqpoint{4.900101in}{2.564689in}}%
\pgfpathlineto{\pgfqpoint{4.945418in}{2.562343in}}%
\pgfpathlineto{\pgfqpoint{4.968077in}{2.549933in}}%
\pgfpathlineto{\pgfqpoint{4.990735in}{2.549726in}}%
\pgfpathlineto{\pgfqpoint{5.013394in}{2.565718in}}%
\pgfpathlineto{\pgfqpoint{5.036053in}{2.548470in}}%
\pgfpathlineto{\pgfqpoint{5.058712in}{2.536895in}}%
\pgfpathlineto{\pgfqpoint{5.081370in}{2.536951in}}%
\pgfpathlineto{\pgfqpoint{5.104029in}{2.531562in}}%
\pgfpathlineto{\pgfqpoint{5.126688in}{2.530845in}}%
\pgfpathlineto{\pgfqpoint{5.149347in}{2.517897in}}%
\pgfpathlineto{\pgfqpoint{5.172005in}{2.526427in}}%
\pgfpathlineto{\pgfqpoint{5.194664in}{2.518621in}}%
\pgfpathlineto{\pgfqpoint{5.217323in}{2.518807in}}%
\pgfpathlineto{\pgfqpoint{5.239982in}{2.507695in}}%
\pgfpathlineto{\pgfqpoint{5.262640in}{2.507261in}}%
\pgfpathlineto{\pgfqpoint{5.285299in}{2.509836in}}%
\pgfpathlineto{\pgfqpoint{5.307958in}{2.501756in}}%
\pgfpathlineto{\pgfqpoint{5.330617in}{2.489190in}}%
\pgfpathlineto{\pgfqpoint{5.353275in}{2.491688in}}%
\pgfpathlineto{\pgfqpoint{5.375934in}{2.498813in}}%
\pgfpathlineto{\pgfqpoint{5.398593in}{2.481652in}}%
\pgfpathlineto{\pgfqpoint{5.421252in}{2.474961in}}%
\pgfpathlineto{\pgfqpoint{5.443910in}{2.472501in}}%
\pgfpathlineto{\pgfqpoint{5.466569in}{2.464151in}}%
\pgfpathlineto{\pgfqpoint{5.489228in}{2.468817in}}%
\pgfpathlineto{\pgfqpoint{5.511887in}{2.454782in}}%
\pgfpathlineto{\pgfqpoint{5.534545in}{2.458417in}}%
\pgfpathlineto{\pgfqpoint{5.534545in}{2.458417in}}%
\pgfusepath{stroke}%
\end{pgfscope}%
\begin{pgfscope}%
\pgfpathrectangle{\pgfqpoint{0.800000in}{0.528000in}}{\pgfqpoint{4.960000in}{3.696000in}}%
\pgfusepath{clip}%
\pgfsetrectcap%
\pgfsetroundjoin%
\pgfsetlinewidth{1.505625pt}%
\definecolor{currentstroke}{rgb}{1.000000,0.498039,0.054902}%
\pgfsetstrokecolor{currentstroke}%
\pgfsetdash{}{0pt}%
\pgfpathmoveto{\pgfqpoint{1.025455in}{3.099476in}}%
\pgfpathlineto{\pgfqpoint{1.048113in}{2.785047in}}%
\pgfpathlineto{\pgfqpoint{1.070772in}{2.657195in}}%
\pgfpathlineto{\pgfqpoint{1.093431in}{2.575763in}}%
\pgfpathlineto{\pgfqpoint{1.138748in}{2.469362in}}%
\pgfpathlineto{\pgfqpoint{1.161407in}{2.407321in}}%
\pgfpathlineto{\pgfqpoint{1.184066in}{2.363281in}}%
\pgfpathlineto{\pgfqpoint{1.206725in}{2.300292in}}%
\pgfpathlineto{\pgfqpoint{1.274701in}{2.099421in}}%
\pgfpathlineto{\pgfqpoint{1.320018in}{1.945329in}}%
\pgfpathlineto{\pgfqpoint{1.342677in}{1.852053in}}%
\pgfpathlineto{\pgfqpoint{1.365336in}{1.748283in}}%
\pgfpathlineto{\pgfqpoint{1.410653in}{1.561206in}}%
\pgfpathlineto{\pgfqpoint{1.501288in}{1.279164in}}%
\pgfpathlineto{\pgfqpoint{1.523947in}{1.233908in}}%
\pgfpathlineto{\pgfqpoint{1.546606in}{1.177063in}}%
\pgfpathlineto{\pgfqpoint{1.569265in}{1.153307in}}%
\pgfpathlineto{\pgfqpoint{1.591923in}{1.118739in}}%
\pgfpathlineto{\pgfqpoint{1.614582in}{1.074366in}}%
\pgfpathlineto{\pgfqpoint{1.637241in}{1.048334in}}%
\pgfpathlineto{\pgfqpoint{1.659899in}{1.035370in}}%
\pgfpathlineto{\pgfqpoint{1.682558in}{1.034592in}}%
\pgfpathlineto{\pgfqpoint{1.705217in}{1.016039in}}%
\pgfpathlineto{\pgfqpoint{1.727876in}{0.986789in}}%
\pgfpathlineto{\pgfqpoint{1.750534in}{0.989110in}}%
\pgfpathlineto{\pgfqpoint{1.773193in}{0.962538in}}%
\pgfpathlineto{\pgfqpoint{1.795852in}{0.990810in}}%
\pgfpathlineto{\pgfqpoint{1.818511in}{0.941386in}}%
\pgfpathlineto{\pgfqpoint{1.841169in}{0.940271in}}%
\pgfpathlineto{\pgfqpoint{1.863828in}{0.923008in}}%
\pgfpathlineto{\pgfqpoint{1.886487in}{0.955900in}}%
\pgfpathlineto{\pgfqpoint{1.909146in}{0.917412in}}%
\pgfpathlineto{\pgfqpoint{1.931804in}{0.911654in}}%
\pgfpathlineto{\pgfqpoint{1.954463in}{0.909566in}}%
\pgfpathlineto{\pgfqpoint{1.977122in}{0.921631in}}%
\pgfpathlineto{\pgfqpoint{1.999781in}{0.883588in}}%
\pgfpathlineto{\pgfqpoint{2.022439in}{0.900825in}}%
\pgfpathlineto{\pgfqpoint{2.045098in}{0.885384in}}%
\pgfpathlineto{\pgfqpoint{2.067757in}{0.871955in}}%
\pgfpathlineto{\pgfqpoint{2.113074in}{0.879390in}}%
\pgfpathlineto{\pgfqpoint{2.135733in}{0.875985in}}%
\pgfpathlineto{\pgfqpoint{2.158392in}{0.856766in}}%
\pgfpathlineto{\pgfqpoint{2.181051in}{0.862341in}}%
\pgfpathlineto{\pgfqpoint{2.203709in}{0.855020in}}%
\pgfpathlineto{\pgfqpoint{2.226368in}{0.866905in}}%
\pgfpathlineto{\pgfqpoint{2.249027in}{0.860336in}}%
\pgfpathlineto{\pgfqpoint{2.271686in}{0.850861in}}%
\pgfpathlineto{\pgfqpoint{2.294344in}{0.840060in}}%
\pgfpathlineto{\pgfqpoint{2.317003in}{0.846165in}}%
\pgfpathlineto{\pgfqpoint{2.339662in}{0.834143in}}%
\pgfpathlineto{\pgfqpoint{2.362321in}{0.841265in}}%
\pgfpathlineto{\pgfqpoint{2.384979in}{0.825802in}}%
\pgfpathlineto{\pgfqpoint{2.430297in}{0.821232in}}%
\pgfpathlineto{\pgfqpoint{2.452956in}{0.829015in}}%
\pgfpathlineto{\pgfqpoint{2.475614in}{0.823953in}}%
\pgfpathlineto{\pgfqpoint{2.498273in}{0.810101in}}%
\pgfpathlineto{\pgfqpoint{2.520932in}{0.837335in}}%
\pgfpathlineto{\pgfqpoint{2.543591in}{0.831344in}}%
\pgfpathlineto{\pgfqpoint{2.566249in}{0.811536in}}%
\pgfpathlineto{\pgfqpoint{2.588908in}{0.814074in}}%
\pgfpathlineto{\pgfqpoint{2.634226in}{0.795134in}}%
\pgfpathlineto{\pgfqpoint{2.656884in}{0.792413in}}%
\pgfpathlineto{\pgfqpoint{2.679543in}{0.813829in}}%
\pgfpathlineto{\pgfqpoint{2.702202in}{0.804558in}}%
\pgfpathlineto{\pgfqpoint{2.724861in}{0.782185in}}%
\pgfpathlineto{\pgfqpoint{2.747519in}{0.783658in}}%
\pgfpathlineto{\pgfqpoint{2.770178in}{0.803142in}}%
\pgfpathlineto{\pgfqpoint{2.792837in}{0.795279in}}%
\pgfpathlineto{\pgfqpoint{2.838154in}{0.795479in}}%
\pgfpathlineto{\pgfqpoint{2.860813in}{0.774551in}}%
\pgfpathlineto{\pgfqpoint{2.883472in}{0.804361in}}%
\pgfpathlineto{\pgfqpoint{2.906131in}{0.785611in}}%
\pgfpathlineto{\pgfqpoint{2.928789in}{0.778802in}}%
\pgfpathlineto{\pgfqpoint{2.951448in}{0.777478in}}%
\pgfpathlineto{\pgfqpoint{2.974107in}{0.785440in}}%
\pgfpathlineto{\pgfqpoint{3.042083in}{0.797174in}}%
\pgfpathlineto{\pgfqpoint{3.064742in}{0.786150in}}%
\pgfpathlineto{\pgfqpoint{3.087401in}{0.765819in}}%
\pgfpathlineto{\pgfqpoint{3.110059in}{0.762780in}}%
\pgfpathlineto{\pgfqpoint{3.132718in}{0.757049in}}%
\pgfpathlineto{\pgfqpoint{3.155377in}{0.779137in}}%
\pgfpathlineto{\pgfqpoint{3.178036in}{0.761299in}}%
\pgfpathlineto{\pgfqpoint{3.200694in}{0.748798in}}%
\pgfpathlineto{\pgfqpoint{3.223353in}{0.756296in}}%
\pgfpathlineto{\pgfqpoint{3.246012in}{0.753340in}}%
\pgfpathlineto{\pgfqpoint{3.268671in}{0.782277in}}%
\pgfpathlineto{\pgfqpoint{3.291329in}{0.765614in}}%
\pgfpathlineto{\pgfqpoint{3.313988in}{0.755048in}}%
\pgfpathlineto{\pgfqpoint{3.336647in}{0.752664in}}%
\pgfpathlineto{\pgfqpoint{3.359306in}{0.756905in}}%
\pgfpathlineto{\pgfqpoint{3.381964in}{0.754894in}}%
\pgfpathlineto{\pgfqpoint{3.404623in}{0.745036in}}%
\pgfpathlineto{\pgfqpoint{3.427282in}{0.756702in}}%
\pgfpathlineto{\pgfqpoint{3.449941in}{0.749512in}}%
\pgfpathlineto{\pgfqpoint{3.472599in}{0.743505in}}%
\pgfpathlineto{\pgfqpoint{3.495258in}{0.750146in}}%
\pgfpathlineto{\pgfqpoint{3.517917in}{0.746930in}}%
\pgfpathlineto{\pgfqpoint{3.540576in}{0.739340in}}%
\pgfpathlineto{\pgfqpoint{3.563234in}{0.746012in}}%
\pgfpathlineto{\pgfqpoint{3.585893in}{0.762609in}}%
\pgfpathlineto{\pgfqpoint{3.608552in}{0.740543in}}%
\pgfpathlineto{\pgfqpoint{3.631211in}{0.748050in}}%
\pgfpathlineto{\pgfqpoint{3.653869in}{0.751865in}}%
\pgfpathlineto{\pgfqpoint{3.676528in}{0.727107in}}%
\pgfpathlineto{\pgfqpoint{3.699187in}{0.744987in}}%
\pgfpathlineto{\pgfqpoint{3.721846in}{0.740335in}}%
\pgfpathlineto{\pgfqpoint{3.744504in}{0.750375in}}%
\pgfpathlineto{\pgfqpoint{3.767163in}{0.735929in}}%
\pgfpathlineto{\pgfqpoint{3.789822in}{0.735323in}}%
\pgfpathlineto{\pgfqpoint{3.812481in}{0.728652in}}%
\pgfpathlineto{\pgfqpoint{3.835139in}{0.737354in}}%
\pgfpathlineto{\pgfqpoint{3.857798in}{0.722478in}}%
\pgfpathlineto{\pgfqpoint{3.880457in}{0.742052in}}%
\pgfpathlineto{\pgfqpoint{3.903116in}{0.740131in}}%
\pgfpathlineto{\pgfqpoint{3.925774in}{0.729776in}}%
\pgfpathlineto{\pgfqpoint{3.948433in}{0.755554in}}%
\pgfpathlineto{\pgfqpoint{3.971092in}{0.726428in}}%
\pgfpathlineto{\pgfqpoint{3.993751in}{0.743000in}}%
\pgfpathlineto{\pgfqpoint{4.016409in}{0.738153in}}%
\pgfpathlineto{\pgfqpoint{4.039068in}{0.751336in}}%
\pgfpathlineto{\pgfqpoint{4.061727in}{0.737004in}}%
\pgfpathlineto{\pgfqpoint{4.084386in}{0.724075in}}%
\pgfpathlineto{\pgfqpoint{4.107044in}{0.726977in}}%
\pgfpathlineto{\pgfqpoint{4.129703in}{0.711275in}}%
\pgfpathlineto{\pgfqpoint{4.152362in}{0.715714in}}%
\pgfpathlineto{\pgfqpoint{4.175021in}{0.706088in}}%
\pgfpathlineto{\pgfqpoint{4.197679in}{0.713620in}}%
\pgfpathlineto{\pgfqpoint{4.220338in}{0.711271in}}%
\pgfpathlineto{\pgfqpoint{4.242997in}{0.712395in}}%
\pgfpathlineto{\pgfqpoint{4.265656in}{0.705000in}}%
\pgfpathlineto{\pgfqpoint{4.288314in}{0.707773in}}%
\pgfpathlineto{\pgfqpoint{4.310973in}{0.712994in}}%
\pgfpathlineto{\pgfqpoint{4.333632in}{0.713075in}}%
\pgfpathlineto{\pgfqpoint{4.356291in}{0.706028in}}%
\pgfpathlineto{\pgfqpoint{4.378949in}{0.712894in}}%
\pgfpathlineto{\pgfqpoint{4.401608in}{0.710974in}}%
\pgfpathlineto{\pgfqpoint{4.424267in}{0.707091in}}%
\pgfpathlineto{\pgfqpoint{4.469584in}{0.714699in}}%
\pgfpathlineto{\pgfqpoint{4.492243in}{0.709475in}}%
\pgfpathlineto{\pgfqpoint{4.514902in}{0.713200in}}%
\pgfpathlineto{\pgfqpoint{4.537561in}{0.705756in}}%
\pgfpathlineto{\pgfqpoint{4.560219in}{0.703462in}}%
\pgfpathlineto{\pgfqpoint{4.673513in}{0.705067in}}%
\pgfpathlineto{\pgfqpoint{4.696172in}{0.703294in}}%
\pgfpathlineto{\pgfqpoint{4.741489in}{0.703974in}}%
\pgfpathlineto{\pgfqpoint{4.764148in}{0.702844in}}%
\pgfpathlineto{\pgfqpoint{4.786807in}{0.705543in}}%
\pgfpathlineto{\pgfqpoint{4.809466in}{0.703428in}}%
\pgfpathlineto{\pgfqpoint{4.854783in}{0.704122in}}%
\pgfpathlineto{\pgfqpoint{4.877442in}{0.702429in}}%
\pgfpathlineto{\pgfqpoint{4.900101in}{0.704337in}}%
\pgfpathlineto{\pgfqpoint{4.922759in}{0.702093in}}%
\pgfpathlineto{\pgfqpoint{4.945418in}{0.703708in}}%
\pgfpathlineto{\pgfqpoint{4.968077in}{0.703953in}}%
\pgfpathlineto{\pgfqpoint{4.990735in}{0.702965in}}%
\pgfpathlineto{\pgfqpoint{5.013394in}{0.705636in}}%
\pgfpathlineto{\pgfqpoint{5.036053in}{0.702613in}}%
\pgfpathlineto{\pgfqpoint{5.081370in}{0.704542in}}%
\pgfpathlineto{\pgfqpoint{5.172005in}{0.703840in}}%
\pgfpathlineto{\pgfqpoint{5.217323in}{0.702809in}}%
\pgfpathlineto{\pgfqpoint{5.398593in}{0.702794in}}%
\pgfpathlineto{\pgfqpoint{5.534545in}{0.703044in}}%
\pgfpathlineto{\pgfqpoint{5.534545in}{0.703044in}}%
\pgfusepath{stroke}%
\end{pgfscope}%
\begin{pgfscope}%
\pgfsetrectcap%
\pgfsetmiterjoin%
\pgfsetlinewidth{0.803000pt}%
\definecolor{currentstroke}{rgb}{0.000000,0.000000,0.000000}%
\pgfsetstrokecolor{currentstroke}%
\pgfsetdash{}{0pt}%
\pgfpathmoveto{\pgfqpoint{0.800000in}{0.528000in}}%
\pgfpathlineto{\pgfqpoint{0.800000in}{4.224000in}}%
\pgfusepath{stroke}%
\end{pgfscope}%
\begin{pgfscope}%
\pgfsetrectcap%
\pgfsetmiterjoin%
\pgfsetlinewidth{0.803000pt}%
\definecolor{currentstroke}{rgb}{0.000000,0.000000,0.000000}%
\pgfsetstrokecolor{currentstroke}%
\pgfsetdash{}{0pt}%
\pgfpathmoveto{\pgfqpoint{5.760000in}{0.528000in}}%
\pgfpathlineto{\pgfqpoint{5.760000in}{4.224000in}}%
\pgfusepath{stroke}%
\end{pgfscope}%
\begin{pgfscope}%
\pgfsetrectcap%
\pgfsetmiterjoin%
\pgfsetlinewidth{0.803000pt}%
\definecolor{currentstroke}{rgb}{0.000000,0.000000,0.000000}%
\pgfsetstrokecolor{currentstroke}%
\pgfsetdash{}{0pt}%
\pgfpathmoveto{\pgfqpoint{0.800000in}{0.528000in}}%
\pgfpathlineto{\pgfqpoint{5.760000in}{0.528000in}}%
\pgfusepath{stroke}%
\end{pgfscope}%
\begin{pgfscope}%
\pgfsetrectcap%
\pgfsetmiterjoin%
\pgfsetlinewidth{0.803000pt}%
\definecolor{currentstroke}{rgb}{0.000000,0.000000,0.000000}%
\pgfsetstrokecolor{currentstroke}%
\pgfsetdash{}{0pt}%
\pgfpathmoveto{\pgfqpoint{0.800000in}{4.224000in}}%
\pgfpathlineto{\pgfqpoint{5.760000in}{4.224000in}}%
\pgfusepath{stroke}%
\end{pgfscope}%
\begin{pgfscope}%
\definecolor{textcolor}{rgb}{0.000000,0.000000,0.000000}%
\pgfsetstrokecolor{textcolor}%
\pgfsetfillcolor{textcolor}%
\pgftext[x=3.280000in,y=4.307333in,,base]{\color{textcolor}{\sffamily\fontsize{12.000000}{14.400000}\selectfont\catcode`\^=\active\def^{\ifmmode\sp\else\^{}\fi}\catcode`\%=\active\def%{\%}Loss on evaluation set during training}}%
\end{pgfscope}%
\end{pgfpicture}%
\makeatother%
\endgroup%

	\label{fig:pred_loss_action_ig} 
\end{figure}

\subsection{Results}

The result which deviates the most from the baseline in table \ref{tab:state_simpl} is the drop in error of $1 - \frac{0.288748}{0.302494} = 1 - 0.954558 \approx 4.54\%$ when including both actions encoded as one-hot values as well as Shapley values as input for the prediction model, compared to when neither action nor any form of feature importance was included. See table \ref{tab:state_simpl}.

\begin{center}
\captionof{table}{Average error for state prediction on the Simple Spread environment without frame stacking}
\label{tab:state_simpl}
\begin{tabular}{lrrr}
\toprule
 & No explanation & Integrated Gradients & Shapley Values \\
\midrule
No action & 0.302494 & 0.301515 & 0.304524 \\
Integer & 0.298232 & 0.301110 & 0.302363 \\
One-hot & 0.296869 & 0.297798 & 0.288748\\
\bottomrule
\end{tabular}
\end{center}


When considering the event predictions for the Simple Spread environment the inclusion of explanations make a noticeably increased difference. With misclassifications being reduced from $1 - 0.700933 = 0.299067$ to $1-0.742933=0.257067$ when including actions encoded as one hot vectors. A decrease of $1 - \frac{0.257067}{0.299067} = 1 - 0.859563 \approx 14.04\%$. Which is the smallest improvement of the $7$ methods that improved the predictions by more than $2\%$. See table \ref{tab:crit_simpl}.

\begin{center}
\captionof{table}{Accuracy for event prediction on the Simple Spread environment without frame stacking, mean over 10 computations}
\label{tab:crit_simpl}
\begin{tabular}{lrrr}
\toprule
 & No explanation & Integrated Gradients & Shapley Values \\
\midrule
    No Action & \textit{0.711730} & 0.806637*** & 0.793757*** \\
Integer & 0.709863 & 0.810210*** & 0.800030*** \\
One-hot & 0.745763*** & 0.811403*** & 0.792647*** \\
\bottomrule
\multicolumn{3}{l}{\textsuperscript{***}$p<0.001$, 
  \textsuperscript{**}$p<0.01$, 
  \textsuperscript{*}$p<0.05$}
\end{tabular}
\end{center}


In the case of state predictions on the Knights Archers Zombies environment the additions of feature importance does not seem to have any important impact as none of the improvements. The inclusion of actions seem to have made an impact on the accuracy of the state prediction, though this decrease in error might just stem from learned knowledge of how an action changes the state, rather than knowledge of how to extract intent from an action. See table \ref{tab:state_kaz}.

\begin{center}
\captionof{table}{Average error, euclidean distance, for state prediction on the Knights Archers Zombies environment}
\label{tab:state_kaz}
\begin{tabular}{lrrr}
\toprule
 & No explanation & Integrated Gradients & Shapley Values \\
\midrule
No Action & 0.040901 & 0.040928 & 0.040728 \\
Integer & 0.039768 & 0.040216 & 0.039505 \\
One-hot & 0.038924 & 0.039241 & 0.038873 \\
\bottomrule
\end{tabular}
\end{center}

In the case of event predictions on the Knights Archers Zombies environment none of the additions seem to have any important impact as none of the improvements. As none of the changes are higher than $2\%$ we do not consider these results relevant. See table \ref{tab:event_kaz}.

\begin{center}
\captionof{table}{Accuracy for event prediction on the Knights Archers Zombies environment}
\label{tab:event_kaz}
\begin{tabular}{lrrr}
\toprule
 & No explanation & Integrated Gradients & Shapley Values \\
\midrule
No action & 0.824400 & 0.829300 & 0.820650 \\
Integer & 0.822700 & 0.828850 & 0.819250 \\
One-hot & 0.824800 & 0.826700 & 0.822950 \\
\bottomrule
\end{tabular}
\end{center}

\subsection{Analysis}
to-do, statistical tests

\section{Memory inclusion}

\subsection{Setup}

There are two main ways of including temporal information in reinforcement learning policies. The first is to include more than just the last observation as input, known as frame stacking. Another design I implemented, that also includes temporal information for the policy is a design that includes an \ac{lstm} layer in that takes input from the second hidden layer and outputs to the last fully connected layer.

We once again construct datasets which contain pairs on the form $(obs, pos)$ and $(obs, crit)$ respectively. For the observations for the policy using an \ac{lstm} we also include the hidden values and cell states, $(h,c)$, as part of the input for the event prediction and state prediction \acp{nn}, to include the temporal information used by the policy for the prediction \acp{nn}, similarly to in the case with frame stacking.

I will once again train 18 neural networks for each environment and see how including the hidden values and cell states of the \ac{lstm} affects the predictions.

\subsection{Results}
When we include frame stacking in the Simple Spread environment in our case we include the previous 3 observations as well so the total observation space has size $18\times 4 = 72$. This auxiliary information requires a new policy to be trained as the representation of the environment is different from the case without frame stacking. The baseline results are unchanged from the case without frame stacking, and all other improvements are less pronounced. See table \ref{tab:state_simpl_stack}.

\begin{center}
\captionof{table}{Average error for state prediction on the Simple Spread environment with frame stacking}
\label{tab:state_simpl_stack}
\begin{tabular}{lrrr}
\toprule
 & No explanation & Integrated Gradients & Shapley Values \\
\midrule
No action & 0.303008 & 0.304275 & 0.298301 \\
Integer & 0.289100 & 0.305222 & 0.291256 \\
One-hot & 0.305099 & 0.305001 & 0.298962 \\
\bottomrule
\end{tabular}
\end{center}

For event prediction on the Simple Spread environment with frame stacking the results are improved for no explanation included, and improved in some cases when considering explanations as well, from the results in the case without frame stacking. See table \ref{tab:event_simpl_stack}.


\begin{center}
\captionof{table}{Average accuracy for event prediction on the Simple Spread environment with frame stacking}
\label{tab:event_simpl_stack}
\begin{tabular}{lrrr}
\toprule
 & No explanation & Integrated Gradients & Shapley Values \\
\midrule
No action & 0.737900 & 0.80300 & 0.756633 \\
Integer & 0.724767 & 0.804433 & 0.763100 \\
One-hot & 0.766567 & 0.812067 & 0.780700 \\
\bottomrule
\end{tabular}
\end{center}

\begin{center}
\captionof{table}{Average error for state prediction on the Simple Spread environment with \ac{lstm} output and cell states}
\label{tab:state_simpl_lstm}
\begin{tabular}{lrrr}
\toprule
 & No explanation & Integrated Gradients & Shapley Value \\
\midrule
No action & 0.306358 & 0.304827 & 0.305336 \\
Integer & 0.306486 & 0.304529 & 0.306255 \\
One-hot & 0.307398 & 0.305968 & 0.305619 \\
\bottomrule
\end{tabular}
\end{center}

\begin{center}
\captionof{table}{Average accuracy for event prediction on the Simple Spread environment with \ac{lstm} output and cell states}
\label{tab:event_simpl_lstm}
\begin{tabular}{lrrr}
\toprule
 & No explanation & Integrated Gradients & Shapley Values \\
\midrule
No action & 0.730967 & 0.775767 & 0.739200 \\
Integer & 0.715100 & 0.778400 & 0.741933 \\
One-hot & 0.728867 & 0.779600 & 0.742567 \\
\bottomrule
\end{tabular}
\end{center}

\subsection{Analysis}

\begin{figure}[!ht]
	%% Creator: Matplotlib, PGF backend
%%
%% To include the figure in your LaTeX document, write
%%   \input{<filename>.pgf}
%%
%% Make sure the required packages are loaded in your preamble
%%   \usepackage{pgf}
%%
%% Also ensure that all the required font packages are loaded; for instance,
%% the lmodern package is sometimes necessary when using math font.
%%   \usepackage{lmodern}
%%
%% Figures using additional raster images can only be included by \input if
%% they are in the same directory as the main LaTeX file. For loading figures
%% from other directories you can use the `import` package
%%   \usepackage{import}
%%
%% and then include the figures with
%%   \import{<path to file>}{<filename>.pgf}
%%
%% Matplotlib used the following preamble
%%   \def\mathdefault#1{#1}
%%   \everymath=\expandafter{\the\everymath\displaystyle}
%%   \IfFileExists{scrextend.sty}{
%%     \usepackage[fontsize=10.000000pt]{scrextend}
%%   }{
%%     \renewcommand{\normalsize}{\fontsize{10.000000}{12.000000}\selectfont}
%%     \normalsize
%%   }
%%   
%%   \makeatletter\@ifpackageloaded{underscore}{}{\usepackage[strings]{underscore}}\makeatother
%%
\begingroup%
\makeatletter%
\begin{pgfpicture}%
\pgfpathrectangle{\pgfpointorigin}{\pgfqpoint{6.000000in}{8.000000in}}%
\pgfusepath{use as bounding box, clip}%
\begin{pgfscope}%
\pgfsetbuttcap%
\pgfsetmiterjoin%
\definecolor{currentfill}{rgb}{1.000000,1.000000,1.000000}%
\pgfsetfillcolor{currentfill}%
\pgfsetlinewidth{0.000000pt}%
\definecolor{currentstroke}{rgb}{1.000000,1.000000,1.000000}%
\pgfsetstrokecolor{currentstroke}%
\pgfsetdash{}{0pt}%
\pgfpathmoveto{\pgfqpoint{0.000000in}{0.000000in}}%
\pgfpathlineto{\pgfqpoint{6.000000in}{0.000000in}}%
\pgfpathlineto{\pgfqpoint{6.000000in}{8.000000in}}%
\pgfpathlineto{\pgfqpoint{0.000000in}{8.000000in}}%
\pgfpathlineto{\pgfqpoint{0.000000in}{0.000000in}}%
\pgfpathclose%
\pgfusepath{fill}%
\end{pgfscope}%
\begin{pgfscope}%
\pgfsetbuttcap%
\pgfsetmiterjoin%
\definecolor{currentfill}{rgb}{1.000000,1.000000,1.000000}%
\pgfsetfillcolor{currentfill}%
\pgfsetlinewidth{0.000000pt}%
\definecolor{currentstroke}{rgb}{0.000000,0.000000,0.000000}%
\pgfsetstrokecolor{currentstroke}%
\pgfsetstrokeopacity{0.000000}%
\pgfsetdash{}{0pt}%
\pgfpathmoveto{\pgfqpoint{1.414736in}{0.548769in}}%
\pgfpathlineto{\pgfqpoint{4.912483in}{0.548769in}}%
\pgfpathlineto{\pgfqpoint{4.912483in}{7.801955in}}%
\pgfpathlineto{\pgfqpoint{1.414736in}{7.801955in}}%
\pgfpathlineto{\pgfqpoint{1.414736in}{0.548769in}}%
\pgfpathclose%
\pgfusepath{fill}%
\end{pgfscope}%
\begin{pgfscope}%
\pgfpathrectangle{\pgfqpoint{1.414736in}{0.548769in}}{\pgfqpoint{3.497747in}{7.253186in}}%
\pgfusepath{clip}%
\pgfsetbuttcap%
\pgfsetroundjoin%
\definecolor{currentfill}{rgb}{0.974417,0.903590,0.130215}%
\pgfsetfillcolor{currentfill}%
\pgfsetlinewidth{1.003750pt}%
\definecolor{currentstroke}{rgb}{0.974417,0.903590,0.130215}%
\pgfsetstrokecolor{currentstroke}%
\pgfsetdash{}{0pt}%
\pgfpathmoveto{\pgfqpoint{3.421232in}{0.856499in}}%
\pgfpathcurveto{\pgfqpoint{3.427056in}{0.856499in}}{\pgfqpoint{3.432643in}{0.858813in}}{\pgfqpoint{3.436761in}{0.862931in}}%
\pgfpathcurveto{\pgfqpoint{3.440879in}{0.867049in}}{\pgfqpoint{3.443193in}{0.872635in}}{\pgfqpoint{3.443193in}{0.878459in}}%
\pgfpathcurveto{\pgfqpoint{3.443193in}{0.884283in}}{\pgfqpoint{3.440879in}{0.889869in}}{\pgfqpoint{3.436761in}{0.893987in}}%
\pgfpathcurveto{\pgfqpoint{3.432643in}{0.898105in}}{\pgfqpoint{3.427056in}{0.900419in}}{\pgfqpoint{3.421232in}{0.900419in}}%
\pgfpathcurveto{\pgfqpoint{3.415408in}{0.900419in}}{\pgfqpoint{3.409822in}{0.898105in}}{\pgfqpoint{3.405704in}{0.893987in}}%
\pgfpathcurveto{\pgfqpoint{3.401586in}{0.889869in}}{\pgfqpoint{3.399272in}{0.884283in}}{\pgfqpoint{3.399272in}{0.878459in}}%
\pgfpathcurveto{\pgfqpoint{3.399272in}{0.872635in}}{\pgfqpoint{3.401586in}{0.867049in}}{\pgfqpoint{3.405704in}{0.862931in}}%
\pgfpathcurveto{\pgfqpoint{3.409822in}{0.858813in}}{\pgfqpoint{3.415408in}{0.856499in}}{\pgfqpoint{3.421232in}{0.856499in}}%
\pgfpathlineto{\pgfqpoint{3.421232in}{0.856499in}}%
\pgfpathclose%
\pgfusepath{stroke,fill}%
\end{pgfscope}%
\begin{pgfscope}%
\pgfpathrectangle{\pgfqpoint{1.414736in}{0.548769in}}{\pgfqpoint{3.497747in}{7.253186in}}%
\pgfusepath{clip}%
\pgfsetbuttcap%
\pgfsetroundjoin%
\definecolor{currentfill}{rgb}{0.440137,0.811138,0.340967}%
\pgfsetfillcolor{currentfill}%
\pgfsetlinewidth{1.003750pt}%
\definecolor{currentstroke}{rgb}{0.440137,0.811138,0.340967}%
\pgfsetstrokecolor{currentstroke}%
\pgfsetdash{}{0pt}%
\pgfpathmoveto{\pgfqpoint{3.718776in}{1.327485in}}%
\pgfpathcurveto{\pgfqpoint{3.724600in}{1.327485in}}{\pgfqpoint{3.730186in}{1.329799in}}{\pgfqpoint{3.734305in}{1.333917in}}%
\pgfpathcurveto{\pgfqpoint{3.738423in}{1.338035in}}{\pgfqpoint{3.740737in}{1.343621in}}{\pgfqpoint{3.740737in}{1.349445in}}%
\pgfpathcurveto{\pgfqpoint{3.740737in}{1.355269in}}{\pgfqpoint{3.738423in}{1.360855in}}{\pgfqpoint{3.734305in}{1.364973in}}%
\pgfpathcurveto{\pgfqpoint{3.730186in}{1.369092in}}{\pgfqpoint{3.724600in}{1.371405in}}{\pgfqpoint{3.718776in}{1.371405in}}%
\pgfpathcurveto{\pgfqpoint{3.712952in}{1.371405in}}{\pgfqpoint{3.707366in}{1.369092in}}{\pgfqpoint{3.703248in}{1.364973in}}%
\pgfpathcurveto{\pgfqpoint{3.699130in}{1.360855in}}{\pgfqpoint{3.696816in}{1.355269in}}{\pgfqpoint{3.696816in}{1.349445in}}%
\pgfpathcurveto{\pgfqpoint{3.696816in}{1.343621in}}{\pgfqpoint{3.699130in}{1.338035in}}{\pgfqpoint{3.703248in}{1.333917in}}%
\pgfpathcurveto{\pgfqpoint{3.707366in}{1.329799in}}{\pgfqpoint{3.712952in}{1.327485in}}{\pgfqpoint{3.718776in}{1.327485in}}%
\pgfpathlineto{\pgfqpoint{3.718776in}{1.327485in}}%
\pgfpathclose%
\pgfusepath{stroke,fill}%
\end{pgfscope}%
\begin{pgfscope}%
\pgfpathrectangle{\pgfqpoint{1.414736in}{0.548769in}}{\pgfqpoint{3.497747in}{7.253186in}}%
\pgfusepath{clip}%
\pgfsetbuttcap%
\pgfsetroundjoin%
\definecolor{currentfill}{rgb}{0.412913,0.803041,0.357269}%
\pgfsetfillcolor{currentfill}%
\pgfsetlinewidth{1.003750pt}%
\definecolor{currentstroke}{rgb}{0.412913,0.803041,0.357269}%
\pgfsetstrokecolor{currentstroke}%
\pgfsetdash{}{0pt}%
\pgfpathmoveto{\pgfqpoint{3.499015in}{1.798471in}}%
\pgfpathcurveto{\pgfqpoint{3.504839in}{1.798471in}}{\pgfqpoint{3.510425in}{1.800785in}}{\pgfqpoint{3.514543in}{1.804903in}}%
\pgfpathcurveto{\pgfqpoint{3.518661in}{1.809021in}}{\pgfqpoint{3.520975in}{1.814607in}}{\pgfqpoint{3.520975in}{1.820431in}}%
\pgfpathcurveto{\pgfqpoint{3.520975in}{1.826255in}}{\pgfqpoint{3.518661in}{1.831841in}}{\pgfqpoint{3.514543in}{1.835960in}}%
\pgfpathcurveto{\pgfqpoint{3.510425in}{1.840078in}}{\pgfqpoint{3.504839in}{1.842392in}}{\pgfqpoint{3.499015in}{1.842392in}}%
\pgfpathcurveto{\pgfqpoint{3.493191in}{1.842392in}}{\pgfqpoint{3.487605in}{1.840078in}}{\pgfqpoint{3.483487in}{1.835960in}}%
\pgfpathcurveto{\pgfqpoint{3.479368in}{1.831841in}}{\pgfqpoint{3.477055in}{1.826255in}}{\pgfqpoint{3.477055in}{1.820431in}}%
\pgfpathcurveto{\pgfqpoint{3.477055in}{1.814607in}}{\pgfqpoint{3.479368in}{1.809021in}}{\pgfqpoint{3.483487in}{1.804903in}}%
\pgfpathcurveto{\pgfqpoint{3.487605in}{1.800785in}}{\pgfqpoint{3.493191in}{1.798471in}}{\pgfqpoint{3.499015in}{1.798471in}}%
\pgfpathlineto{\pgfqpoint{3.499015in}{1.798471in}}%
\pgfpathclose%
\pgfusepath{stroke,fill}%
\end{pgfscope}%
\begin{pgfscope}%
\pgfpathrectangle{\pgfqpoint{1.414736in}{0.548769in}}{\pgfqpoint{3.497747in}{7.253186in}}%
\pgfusepath{clip}%
\pgfsetbuttcap%
\pgfsetroundjoin%
\definecolor{currentfill}{rgb}{0.993248,0.906157,0.143936}%
\pgfsetfillcolor{currentfill}%
\pgfsetlinewidth{1.003750pt}%
\definecolor{currentstroke}{rgb}{0.993248,0.906157,0.143936}%
\pgfsetstrokecolor{currentstroke}%
\pgfsetdash{}{0pt}%
\pgfpathmoveto{\pgfqpoint{3.852525in}{2.269457in}}%
\pgfpathcurveto{\pgfqpoint{3.858349in}{2.269457in}}{\pgfqpoint{3.863935in}{2.271771in}}{\pgfqpoint{3.868054in}{2.275889in}}%
\pgfpathcurveto{\pgfqpoint{3.872172in}{2.280007in}}{\pgfqpoint{3.874486in}{2.285593in}}{\pgfqpoint{3.874486in}{2.291417in}}%
\pgfpathcurveto{\pgfqpoint{3.874486in}{2.297241in}}{\pgfqpoint{3.872172in}{2.302828in}}{\pgfqpoint{3.868054in}{2.306946in}}%
\pgfpathcurveto{\pgfqpoint{3.863935in}{2.311064in}}{\pgfqpoint{3.858349in}{2.313378in}}{\pgfqpoint{3.852525in}{2.313378in}}%
\pgfpathcurveto{\pgfqpoint{3.846701in}{2.313378in}}{\pgfqpoint{3.841115in}{2.311064in}}{\pgfqpoint{3.836997in}{2.306946in}}%
\pgfpathcurveto{\pgfqpoint{3.832879in}{2.302828in}}{\pgfqpoint{3.830565in}{2.297241in}}{\pgfqpoint{3.830565in}{2.291417in}}%
\pgfpathcurveto{\pgfqpoint{3.830565in}{2.285593in}}{\pgfqpoint{3.832879in}{2.280007in}}{\pgfqpoint{3.836997in}{2.275889in}}%
\pgfpathcurveto{\pgfqpoint{3.841115in}{2.271771in}}{\pgfqpoint{3.846701in}{2.269457in}}{\pgfqpoint{3.852525in}{2.269457in}}%
\pgfpathlineto{\pgfqpoint{3.852525in}{2.269457in}}%
\pgfpathclose%
\pgfusepath{stroke,fill}%
\end{pgfscope}%
\begin{pgfscope}%
\pgfpathrectangle{\pgfqpoint{1.414736in}{0.548769in}}{\pgfqpoint{3.497747in}{7.253186in}}%
\pgfusepath{clip}%
\pgfsetbuttcap%
\pgfsetroundjoin%
\definecolor{currentfill}{rgb}{0.232815,0.732247,0.459277}%
\pgfsetfillcolor{currentfill}%
\pgfsetlinewidth{1.003750pt}%
\definecolor{currentstroke}{rgb}{0.232815,0.732247,0.459277}%
\pgfsetstrokecolor{currentstroke}%
\pgfsetdash{}{0pt}%
\pgfpathmoveto{\pgfqpoint{3.889990in}{2.740443in}}%
\pgfpathcurveto{\pgfqpoint{3.895814in}{2.740443in}}{\pgfqpoint{3.901400in}{2.742757in}}{\pgfqpoint{3.905518in}{2.746875in}}%
\pgfpathcurveto{\pgfqpoint{3.909636in}{2.750993in}}{\pgfqpoint{3.911950in}{2.756580in}}{\pgfqpoint{3.911950in}{2.762404in}}%
\pgfpathcurveto{\pgfqpoint{3.911950in}{2.768227in}}{\pgfqpoint{3.909636in}{2.773814in}}{\pgfqpoint{3.905518in}{2.777932in}}%
\pgfpathcurveto{\pgfqpoint{3.901400in}{2.782050in}}{\pgfqpoint{3.895814in}{2.784364in}}{\pgfqpoint{3.889990in}{2.784364in}}%
\pgfpathcurveto{\pgfqpoint{3.884166in}{2.784364in}}{\pgfqpoint{3.878580in}{2.782050in}}{\pgfqpoint{3.874462in}{2.777932in}}%
\pgfpathcurveto{\pgfqpoint{3.870343in}{2.773814in}}{\pgfqpoint{3.868030in}{2.768227in}}{\pgfqpoint{3.868030in}{2.762404in}}%
\pgfpathcurveto{\pgfqpoint{3.868030in}{2.756580in}}{\pgfqpoint{3.870343in}{2.750993in}}{\pgfqpoint{3.874462in}{2.746875in}}%
\pgfpathcurveto{\pgfqpoint{3.878580in}{2.742757in}}{\pgfqpoint{3.884166in}{2.740443in}}{\pgfqpoint{3.889990in}{2.740443in}}%
\pgfpathlineto{\pgfqpoint{3.889990in}{2.740443in}}%
\pgfpathclose%
\pgfusepath{stroke,fill}%
\end{pgfscope}%
\begin{pgfscope}%
\pgfpathrectangle{\pgfqpoint{1.414736in}{0.548769in}}{\pgfqpoint{3.497747in}{7.253186in}}%
\pgfusepath{clip}%
\pgfsetbuttcap%
\pgfsetroundjoin%
\definecolor{currentfill}{rgb}{0.515992,0.831158,0.294279}%
\pgfsetfillcolor{currentfill}%
\pgfsetlinewidth{1.003750pt}%
\definecolor{currentstroke}{rgb}{0.515992,0.831158,0.294279}%
\pgfsetstrokecolor{currentstroke}%
\pgfsetdash{}{0pt}%
\pgfpathmoveto{\pgfqpoint{3.297483in}{3.211429in}}%
\pgfpathcurveto{\pgfqpoint{3.303307in}{3.211429in}}{\pgfqpoint{3.308893in}{3.213743in}}{\pgfqpoint{3.313011in}{3.217861in}}%
\pgfpathcurveto{\pgfqpoint{3.317129in}{3.221980in}}{\pgfqpoint{3.319443in}{3.227566in}}{\pgfqpoint{3.319443in}{3.233390in}}%
\pgfpathcurveto{\pgfqpoint{3.319443in}{3.239214in}}{\pgfqpoint{3.317129in}{3.244800in}}{\pgfqpoint{3.313011in}{3.248918in}}%
\pgfpathcurveto{\pgfqpoint{3.308893in}{3.253036in}}{\pgfqpoint{3.303307in}{3.255350in}}{\pgfqpoint{3.297483in}{3.255350in}}%
\pgfpathcurveto{\pgfqpoint{3.291659in}{3.255350in}}{\pgfqpoint{3.286073in}{3.253036in}}{\pgfqpoint{3.281955in}{3.248918in}}%
\pgfpathcurveto{\pgfqpoint{3.277837in}{3.244800in}}{\pgfqpoint{3.275523in}{3.239214in}}{\pgfqpoint{3.275523in}{3.233390in}}%
\pgfpathcurveto{\pgfqpoint{3.275523in}{3.227566in}}{\pgfqpoint{3.277837in}{3.221980in}}{\pgfqpoint{3.281955in}{3.217861in}}%
\pgfpathcurveto{\pgfqpoint{3.286073in}{3.213743in}}{\pgfqpoint{3.291659in}{3.211429in}}{\pgfqpoint{3.297483in}{3.211429in}}%
\pgfpathlineto{\pgfqpoint{3.297483in}{3.211429in}}%
\pgfpathclose%
\pgfusepath{stroke,fill}%
\end{pgfscope}%
\begin{pgfscope}%
\pgfpathrectangle{\pgfqpoint{1.414736in}{0.548769in}}{\pgfqpoint{3.497747in}{7.253186in}}%
\pgfusepath{clip}%
\pgfsetbuttcap%
\pgfsetroundjoin%
\definecolor{currentfill}{rgb}{0.896320,0.893616,0.096335}%
\pgfsetfillcolor{currentfill}%
\pgfsetlinewidth{1.003750pt}%
\definecolor{currentstroke}{rgb}{0.896320,0.893616,0.096335}%
\pgfsetstrokecolor{currentstroke}%
\pgfsetdash{}{0pt}%
\pgfpathmoveto{\pgfqpoint{3.870898in}{3.682415in}}%
\pgfpathcurveto{\pgfqpoint{3.876722in}{3.682415in}}{\pgfqpoint{3.882308in}{3.684729in}}{\pgfqpoint{3.886426in}{3.688847in}}%
\pgfpathcurveto{\pgfqpoint{3.890544in}{3.692966in}}{\pgfqpoint{3.892858in}{3.698552in}}{\pgfqpoint{3.892858in}{3.704376in}}%
\pgfpathcurveto{\pgfqpoint{3.892858in}{3.710200in}}{\pgfqpoint{3.890544in}{3.715786in}}{\pgfqpoint{3.886426in}{3.719904in}}%
\pgfpathcurveto{\pgfqpoint{3.882308in}{3.724022in}}{\pgfqpoint{3.876722in}{3.726336in}}{\pgfqpoint{3.870898in}{3.726336in}}%
\pgfpathcurveto{\pgfqpoint{3.865074in}{3.726336in}}{\pgfqpoint{3.859488in}{3.724022in}}{\pgfqpoint{3.855369in}{3.719904in}}%
\pgfpathcurveto{\pgfqpoint{3.851251in}{3.715786in}}{\pgfqpoint{3.848937in}{3.710200in}}{\pgfqpoint{3.848937in}{3.704376in}}%
\pgfpathcurveto{\pgfqpoint{3.848937in}{3.698552in}}{\pgfqpoint{3.851251in}{3.692966in}}{\pgfqpoint{3.855369in}{3.688847in}}%
\pgfpathcurveto{\pgfqpoint{3.859488in}{3.684729in}}{\pgfqpoint{3.865074in}{3.682415in}}{\pgfqpoint{3.870898in}{3.682415in}}%
\pgfpathlineto{\pgfqpoint{3.870898in}{3.682415in}}%
\pgfpathclose%
\pgfusepath{stroke,fill}%
\end{pgfscope}%
\begin{pgfscope}%
\pgfpathrectangle{\pgfqpoint{1.414736in}{0.548769in}}{\pgfqpoint{3.497747in}{7.253186in}}%
\pgfusepath{clip}%
\pgfsetbuttcap%
\pgfsetroundjoin%
\definecolor{currentfill}{rgb}{0.993248,0.906157,0.143936}%
\pgfsetfillcolor{currentfill}%
\pgfsetlinewidth{1.003750pt}%
\definecolor{currentstroke}{rgb}{0.993248,0.906157,0.143936}%
\pgfsetstrokecolor{currentstroke}%
\pgfsetdash{}{0pt}%
\pgfpathmoveto{\pgfqpoint{3.181930in}{4.153402in}}%
\pgfpathcurveto{\pgfqpoint{3.187754in}{4.153402in}}{\pgfqpoint{3.193340in}{4.155715in}}{\pgfqpoint{3.197458in}{4.159834in}}%
\pgfpathcurveto{\pgfqpoint{3.201577in}{4.163952in}}{\pgfqpoint{3.203890in}{4.169538in}}{\pgfqpoint{3.203890in}{4.175362in}}%
\pgfpathcurveto{\pgfqpoint{3.203890in}{4.181186in}}{\pgfqpoint{3.201577in}{4.186772in}}{\pgfqpoint{3.197458in}{4.190890in}}%
\pgfpathcurveto{\pgfqpoint{3.193340in}{4.195008in}}{\pgfqpoint{3.187754in}{4.197322in}}{\pgfqpoint{3.181930in}{4.197322in}}%
\pgfpathcurveto{\pgfqpoint{3.176106in}{4.197322in}}{\pgfqpoint{3.170520in}{4.195008in}}{\pgfqpoint{3.166402in}{4.190890in}}%
\pgfpathcurveto{\pgfqpoint{3.162284in}{4.186772in}}{\pgfqpoint{3.159970in}{4.181186in}}{\pgfqpoint{3.159970in}{4.175362in}}%
\pgfpathcurveto{\pgfqpoint{3.159970in}{4.169538in}}{\pgfqpoint{3.162284in}{4.163952in}}{\pgfqpoint{3.166402in}{4.159834in}}%
\pgfpathcurveto{\pgfqpoint{3.170520in}{4.155715in}}{\pgfqpoint{3.176106in}{4.153402in}}{\pgfqpoint{3.181930in}{4.153402in}}%
\pgfpathlineto{\pgfqpoint{3.181930in}{4.153402in}}%
\pgfpathclose%
\pgfusepath{stroke,fill}%
\end{pgfscope}%
\begin{pgfscope}%
\pgfpathrectangle{\pgfqpoint{1.414736in}{0.548769in}}{\pgfqpoint{3.497747in}{7.253186in}}%
\pgfusepath{clip}%
\pgfsetbuttcap%
\pgfsetroundjoin%
\definecolor{currentfill}{rgb}{0.626579,0.854645,0.223353}%
\pgfsetfillcolor{currentfill}%
\pgfsetlinewidth{1.003750pt}%
\definecolor{currentstroke}{rgb}{0.626579,0.854645,0.223353}%
\pgfsetstrokecolor{currentstroke}%
\pgfsetdash{}{0pt}%
\pgfpathmoveto{\pgfqpoint{3.261342in}{4.624388in}}%
\pgfpathcurveto{\pgfqpoint{3.267166in}{4.624388in}}{\pgfqpoint{3.272752in}{4.626702in}}{\pgfqpoint{3.276870in}{4.630820in}}%
\pgfpathcurveto{\pgfqpoint{3.280988in}{4.634938in}}{\pgfqpoint{3.283302in}{4.640524in}}{\pgfqpoint{3.283302in}{4.646348in}}%
\pgfpathcurveto{\pgfqpoint{3.283302in}{4.652172in}}{\pgfqpoint{3.280988in}{4.657758in}}{\pgfqpoint{3.276870in}{4.661876in}}%
\pgfpathcurveto{\pgfqpoint{3.272752in}{4.665994in}}{\pgfqpoint{3.267166in}{4.668308in}}{\pgfqpoint{3.261342in}{4.668308in}}%
\pgfpathcurveto{\pgfqpoint{3.255518in}{4.668308in}}{\pgfqpoint{3.249932in}{4.665994in}}{\pgfqpoint{3.245813in}{4.661876in}}%
\pgfpathcurveto{\pgfqpoint{3.241695in}{4.657758in}}{\pgfqpoint{3.239381in}{4.652172in}}{\pgfqpoint{3.239381in}{4.646348in}}%
\pgfpathcurveto{\pgfqpoint{3.239381in}{4.640524in}}{\pgfqpoint{3.241695in}{4.634938in}}{\pgfqpoint{3.245813in}{4.630820in}}%
\pgfpathcurveto{\pgfqpoint{3.249932in}{4.626702in}}{\pgfqpoint{3.255518in}{4.624388in}}{\pgfqpoint{3.261342in}{4.624388in}}%
\pgfpathlineto{\pgfqpoint{3.261342in}{4.624388in}}%
\pgfpathclose%
\pgfusepath{stroke,fill}%
\end{pgfscope}%
\begin{pgfscope}%
\pgfpathrectangle{\pgfqpoint{1.414736in}{0.548769in}}{\pgfqpoint{3.497747in}{7.253186in}}%
\pgfusepath{clip}%
\pgfsetbuttcap%
\pgfsetroundjoin%
\definecolor{currentfill}{rgb}{0.283072,0.130895,0.449241}%
\pgfsetfillcolor{currentfill}%
\pgfsetlinewidth{1.003750pt}%
\definecolor{currentstroke}{rgb}{0.283072,0.130895,0.449241}%
\pgfsetstrokecolor{currentstroke}%
\pgfsetdash{}{0pt}%
\pgfpathmoveto{\pgfqpoint{3.237127in}{5.095374in}}%
\pgfpathcurveto{\pgfqpoint{3.242951in}{5.095374in}}{\pgfqpoint{3.248537in}{5.097688in}}{\pgfqpoint{3.252655in}{5.101806in}}%
\pgfpathcurveto{\pgfqpoint{3.256773in}{5.105924in}}{\pgfqpoint{3.259087in}{5.111510in}}{\pgfqpoint{3.259087in}{5.117334in}}%
\pgfpathcurveto{\pgfqpoint{3.259087in}{5.123158in}}{\pgfqpoint{3.256773in}{5.128744in}}{\pgfqpoint{3.252655in}{5.132862in}}%
\pgfpathcurveto{\pgfqpoint{3.248537in}{5.136980in}}{\pgfqpoint{3.242951in}{5.139294in}}{\pgfqpoint{3.237127in}{5.139294in}}%
\pgfpathcurveto{\pgfqpoint{3.231303in}{5.139294in}}{\pgfqpoint{3.225717in}{5.136980in}}{\pgfqpoint{3.221599in}{5.132862in}}%
\pgfpathcurveto{\pgfqpoint{3.217481in}{5.128744in}}{\pgfqpoint{3.215167in}{5.123158in}}{\pgfqpoint{3.215167in}{5.117334in}}%
\pgfpathcurveto{\pgfqpoint{3.215167in}{5.111510in}}{\pgfqpoint{3.217481in}{5.105924in}}{\pgfqpoint{3.221599in}{5.101806in}}%
\pgfpathcurveto{\pgfqpoint{3.225717in}{5.097688in}}{\pgfqpoint{3.231303in}{5.095374in}}{\pgfqpoint{3.237127in}{5.095374in}}%
\pgfpathlineto{\pgfqpoint{3.237127in}{5.095374in}}%
\pgfpathclose%
\pgfusepath{stroke,fill}%
\end{pgfscope}%
\begin{pgfscope}%
\pgfpathrectangle{\pgfqpoint{1.414736in}{0.548769in}}{\pgfqpoint{3.497747in}{7.253186in}}%
\pgfusepath{clip}%
\pgfsetbuttcap%
\pgfsetroundjoin%
\definecolor{currentfill}{rgb}{0.187231,0.414746,0.556547}%
\pgfsetfillcolor{currentfill}%
\pgfsetlinewidth{1.003750pt}%
\definecolor{currentstroke}{rgb}{0.187231,0.414746,0.556547}%
\pgfsetstrokecolor{currentstroke}%
\pgfsetdash{}{0pt}%
\pgfpathmoveto{\pgfqpoint{3.937575in}{5.566360in}}%
\pgfpathcurveto{\pgfqpoint{3.943399in}{5.566360in}}{\pgfqpoint{3.948985in}{5.568674in}}{\pgfqpoint{3.953103in}{5.572792in}}%
\pgfpathcurveto{\pgfqpoint{3.957222in}{5.576910in}}{\pgfqpoint{3.959535in}{5.582496in}}{\pgfqpoint{3.959535in}{5.588320in}}%
\pgfpathcurveto{\pgfqpoint{3.959535in}{5.594144in}}{\pgfqpoint{3.957222in}{5.599730in}}{\pgfqpoint{3.953103in}{5.603848in}}%
\pgfpathcurveto{\pgfqpoint{3.948985in}{5.607967in}}{\pgfqpoint{3.943399in}{5.610280in}}{\pgfqpoint{3.937575in}{5.610280in}}%
\pgfpathcurveto{\pgfqpoint{3.931751in}{5.610280in}}{\pgfqpoint{3.926165in}{5.607967in}}{\pgfqpoint{3.922047in}{5.603848in}}%
\pgfpathcurveto{\pgfqpoint{3.917929in}{5.599730in}}{\pgfqpoint{3.915615in}{5.594144in}}{\pgfqpoint{3.915615in}{5.588320in}}%
\pgfpathcurveto{\pgfqpoint{3.915615in}{5.582496in}}{\pgfqpoint{3.917929in}{5.576910in}}{\pgfqpoint{3.922047in}{5.572792in}}%
\pgfpathcurveto{\pgfqpoint{3.926165in}{5.568674in}}{\pgfqpoint{3.931751in}{5.566360in}}{\pgfqpoint{3.937575in}{5.566360in}}%
\pgfpathlineto{\pgfqpoint{3.937575in}{5.566360in}}%
\pgfpathclose%
\pgfusepath{stroke,fill}%
\end{pgfscope}%
\begin{pgfscope}%
\pgfpathrectangle{\pgfqpoint{1.414736in}{0.548769in}}{\pgfqpoint{3.497747in}{7.253186in}}%
\pgfusepath{clip}%
\pgfsetbuttcap%
\pgfsetroundjoin%
\definecolor{currentfill}{rgb}{0.153894,0.680203,0.504172}%
\pgfsetfillcolor{currentfill}%
\pgfsetlinewidth{1.003750pt}%
\definecolor{currentstroke}{rgb}{0.153894,0.680203,0.504172}%
\pgfsetstrokecolor{currentstroke}%
\pgfsetdash{}{0pt}%
\pgfpathmoveto{\pgfqpoint{3.450411in}{6.037346in}}%
\pgfpathcurveto{\pgfqpoint{3.456235in}{6.037346in}}{\pgfqpoint{3.461821in}{6.039660in}}{\pgfqpoint{3.465939in}{6.043778in}}%
\pgfpathcurveto{\pgfqpoint{3.470057in}{6.047896in}}{\pgfqpoint{3.472371in}{6.053482in}}{\pgfqpoint{3.472371in}{6.059306in}}%
\pgfpathcurveto{\pgfqpoint{3.472371in}{6.065130in}}{\pgfqpoint{3.470057in}{6.070716in}}{\pgfqpoint{3.465939in}{6.074835in}}%
\pgfpathcurveto{\pgfqpoint{3.461821in}{6.078953in}}{\pgfqpoint{3.456235in}{6.081267in}}{\pgfqpoint{3.450411in}{6.081267in}}%
\pgfpathcurveto{\pgfqpoint{3.444587in}{6.081267in}}{\pgfqpoint{3.439001in}{6.078953in}}{\pgfqpoint{3.434883in}{6.074835in}}%
\pgfpathcurveto{\pgfqpoint{3.430765in}{6.070716in}}{\pgfqpoint{3.428451in}{6.065130in}}{\pgfqpoint{3.428451in}{6.059306in}}%
\pgfpathcurveto{\pgfqpoint{3.428451in}{6.053482in}}{\pgfqpoint{3.430765in}{6.047896in}}{\pgfqpoint{3.434883in}{6.043778in}}%
\pgfpathcurveto{\pgfqpoint{3.439001in}{6.039660in}}{\pgfqpoint{3.444587in}{6.037346in}}{\pgfqpoint{3.450411in}{6.037346in}}%
\pgfpathlineto{\pgfqpoint{3.450411in}{6.037346in}}%
\pgfpathclose%
\pgfusepath{stroke,fill}%
\end{pgfscope}%
\begin{pgfscope}%
\pgfpathrectangle{\pgfqpoint{1.414736in}{0.548769in}}{\pgfqpoint{3.497747in}{7.253186in}}%
\pgfusepath{clip}%
\pgfsetbuttcap%
\pgfsetroundjoin%
\definecolor{currentfill}{rgb}{0.575563,0.844566,0.256415}%
\pgfsetfillcolor{currentfill}%
\pgfsetlinewidth{1.003750pt}%
\definecolor{currentstroke}{rgb}{0.575563,0.844566,0.256415}%
\pgfsetstrokecolor{currentstroke}%
\pgfsetdash{}{0pt}%
\pgfpathmoveto{\pgfqpoint{4.075476in}{6.508332in}}%
\pgfpathcurveto{\pgfqpoint{4.081300in}{6.508332in}}{\pgfqpoint{4.086886in}{6.510646in}}{\pgfqpoint{4.091005in}{6.514764in}}%
\pgfpathcurveto{\pgfqpoint{4.095123in}{6.518882in}}{\pgfqpoint{4.097437in}{6.524468in}}{\pgfqpoint{4.097437in}{6.530292in}}%
\pgfpathcurveto{\pgfqpoint{4.097437in}{6.536116in}}{\pgfqpoint{4.095123in}{6.541702in}}{\pgfqpoint{4.091005in}{6.545821in}}%
\pgfpathcurveto{\pgfqpoint{4.086886in}{6.549939in}}{\pgfqpoint{4.081300in}{6.552253in}}{\pgfqpoint{4.075476in}{6.552253in}}%
\pgfpathcurveto{\pgfqpoint{4.069652in}{6.552253in}}{\pgfqpoint{4.064066in}{6.549939in}}{\pgfqpoint{4.059948in}{6.545821in}}%
\pgfpathcurveto{\pgfqpoint{4.055830in}{6.541702in}}{\pgfqpoint{4.053516in}{6.536116in}}{\pgfqpoint{4.053516in}{6.530292in}}%
\pgfpathcurveto{\pgfqpoint{4.053516in}{6.524468in}}{\pgfqpoint{4.055830in}{6.518882in}}{\pgfqpoint{4.059948in}{6.514764in}}%
\pgfpathcurveto{\pgfqpoint{4.064066in}{6.510646in}}{\pgfqpoint{4.069652in}{6.508332in}}{\pgfqpoint{4.075476in}{6.508332in}}%
\pgfpathlineto{\pgfqpoint{4.075476in}{6.508332in}}%
\pgfpathclose%
\pgfusepath{stroke,fill}%
\end{pgfscope}%
\begin{pgfscope}%
\pgfpathrectangle{\pgfqpoint{1.414736in}{0.548769in}}{\pgfqpoint{3.497747in}{7.253186in}}%
\pgfusepath{clip}%
\pgfsetbuttcap%
\pgfsetroundjoin%
\definecolor{currentfill}{rgb}{0.246070,0.738910,0.452024}%
\pgfsetfillcolor{currentfill}%
\pgfsetlinewidth{1.003750pt}%
\definecolor{currentstroke}{rgb}{0.246070,0.738910,0.452024}%
\pgfsetstrokecolor{currentstroke}%
\pgfsetdash{}{0pt}%
\pgfpathmoveto{\pgfqpoint{3.424865in}{6.979318in}}%
\pgfpathcurveto{\pgfqpoint{3.430688in}{6.979318in}}{\pgfqpoint{3.436275in}{6.981632in}}{\pgfqpoint{3.440393in}{6.985750in}}%
\pgfpathcurveto{\pgfqpoint{3.444511in}{6.989868in}}{\pgfqpoint{3.446825in}{6.995455in}}{\pgfqpoint{3.446825in}{7.001278in}}%
\pgfpathcurveto{\pgfqpoint{3.446825in}{7.007102in}}{\pgfqpoint{3.444511in}{7.012689in}}{\pgfqpoint{3.440393in}{7.016807in}}%
\pgfpathcurveto{\pgfqpoint{3.436275in}{7.020925in}}{\pgfqpoint{3.430688in}{7.023239in}}{\pgfqpoint{3.424865in}{7.023239in}}%
\pgfpathcurveto{\pgfqpoint{3.419041in}{7.023239in}}{\pgfqpoint{3.413454in}{7.020925in}}{\pgfqpoint{3.409336in}{7.016807in}}%
\pgfpathcurveto{\pgfqpoint{3.405218in}{7.012689in}}{\pgfqpoint{3.402904in}{7.007102in}}{\pgfqpoint{3.402904in}{7.001278in}}%
\pgfpathcurveto{\pgfqpoint{3.402904in}{6.995455in}}{\pgfqpoint{3.405218in}{6.989868in}}{\pgfqpoint{3.409336in}{6.985750in}}%
\pgfpathcurveto{\pgfqpoint{3.413454in}{6.981632in}}{\pgfqpoint{3.419041in}{6.979318in}}{\pgfqpoint{3.424865in}{6.979318in}}%
\pgfpathlineto{\pgfqpoint{3.424865in}{6.979318in}}%
\pgfpathclose%
\pgfusepath{stroke,fill}%
\end{pgfscope}%
\begin{pgfscope}%
\pgfpathrectangle{\pgfqpoint{1.414736in}{0.548769in}}{\pgfqpoint{3.497747in}{7.253186in}}%
\pgfusepath{clip}%
\pgfsetbuttcap%
\pgfsetroundjoin%
\definecolor{currentfill}{rgb}{0.214000,0.722114,0.469588}%
\pgfsetfillcolor{currentfill}%
\pgfsetlinewidth{1.003750pt}%
\definecolor{currentstroke}{rgb}{0.214000,0.722114,0.469588}%
\pgfsetstrokecolor{currentstroke}%
\pgfsetdash{}{0pt}%
\pgfpathmoveto{\pgfqpoint{3.901538in}{7.450304in}}%
\pgfpathcurveto{\pgfqpoint{3.907362in}{7.450304in}}{\pgfqpoint{3.912948in}{7.452618in}}{\pgfqpoint{3.917067in}{7.456736in}}%
\pgfpathcurveto{\pgfqpoint{3.921185in}{7.460854in}}{\pgfqpoint{3.923499in}{7.466441in}}{\pgfqpoint{3.923499in}{7.472265in}}%
\pgfpathcurveto{\pgfqpoint{3.923499in}{7.478089in}}{\pgfqpoint{3.921185in}{7.483675in}}{\pgfqpoint{3.917067in}{7.487793in}}%
\pgfpathcurveto{\pgfqpoint{3.912948in}{7.491911in}}{\pgfqpoint{3.907362in}{7.494225in}}{\pgfqpoint{3.901538in}{7.494225in}}%
\pgfpathcurveto{\pgfqpoint{3.895714in}{7.494225in}}{\pgfqpoint{3.890128in}{7.491911in}}{\pgfqpoint{3.886010in}{7.487793in}}%
\pgfpathcurveto{\pgfqpoint{3.881892in}{7.483675in}}{\pgfqpoint{3.879578in}{7.478089in}}{\pgfqpoint{3.879578in}{7.472265in}}%
\pgfpathcurveto{\pgfqpoint{3.879578in}{7.466441in}}{\pgfqpoint{3.881892in}{7.460854in}}{\pgfqpoint{3.886010in}{7.456736in}}%
\pgfpathcurveto{\pgfqpoint{3.890128in}{7.452618in}}{\pgfqpoint{3.895714in}{7.450304in}}{\pgfqpoint{3.901538in}{7.450304in}}%
\pgfpathlineto{\pgfqpoint{3.901538in}{7.450304in}}%
\pgfpathclose%
\pgfusepath{stroke,fill}%
\end{pgfscope}%
\begin{pgfscope}%
\pgfpathrectangle{\pgfqpoint{1.414736in}{0.548769in}}{\pgfqpoint{3.497747in}{7.253186in}}%
\pgfusepath{clip}%
\pgfsetbuttcap%
\pgfsetroundjoin%
\definecolor{currentfill}{rgb}{0.993248,0.906157,0.143936}%
\pgfsetfillcolor{currentfill}%
\pgfsetlinewidth{1.003750pt}%
\definecolor{currentstroke}{rgb}{0.993248,0.906157,0.143936}%
\pgfsetstrokecolor{currentstroke}%
\pgfsetdash{}{0pt}%
\pgfpathmoveto{\pgfqpoint{3.634590in}{0.856499in}}%
\pgfpathcurveto{\pgfqpoint{3.640414in}{0.856499in}}{\pgfqpoint{3.646000in}{0.858813in}}{\pgfqpoint{3.650118in}{0.862931in}}%
\pgfpathcurveto{\pgfqpoint{3.654236in}{0.867049in}}{\pgfqpoint{3.656550in}{0.872635in}}{\pgfqpoint{3.656550in}{0.878459in}}%
\pgfpathcurveto{\pgfqpoint{3.656550in}{0.884283in}}{\pgfqpoint{3.654236in}{0.889869in}}{\pgfqpoint{3.650118in}{0.893987in}}%
\pgfpathcurveto{\pgfqpoint{3.646000in}{0.898105in}}{\pgfqpoint{3.640414in}{0.900419in}}{\pgfqpoint{3.634590in}{0.900419in}}%
\pgfpathcurveto{\pgfqpoint{3.628766in}{0.900419in}}{\pgfqpoint{3.623180in}{0.898105in}}{\pgfqpoint{3.619062in}{0.893987in}}%
\pgfpathcurveto{\pgfqpoint{3.614943in}{0.889869in}}{\pgfqpoint{3.612630in}{0.884283in}}{\pgfqpoint{3.612630in}{0.878459in}}%
\pgfpathcurveto{\pgfqpoint{3.612630in}{0.872635in}}{\pgfqpoint{3.614943in}{0.867049in}}{\pgfqpoint{3.619062in}{0.862931in}}%
\pgfpathcurveto{\pgfqpoint{3.623180in}{0.858813in}}{\pgfqpoint{3.628766in}{0.856499in}}{\pgfqpoint{3.634590in}{0.856499in}}%
\pgfpathlineto{\pgfqpoint{3.634590in}{0.856499in}}%
\pgfpathclose%
\pgfusepath{stroke,fill}%
\end{pgfscope}%
\begin{pgfscope}%
\pgfpathrectangle{\pgfqpoint{1.414736in}{0.548769in}}{\pgfqpoint{3.497747in}{7.253186in}}%
\pgfusepath{clip}%
\pgfsetbuttcap%
\pgfsetroundjoin%
\definecolor{currentfill}{rgb}{0.668054,0.861999,0.196293}%
\pgfsetfillcolor{currentfill}%
\pgfsetlinewidth{1.003750pt}%
\definecolor{currentstroke}{rgb}{0.668054,0.861999,0.196293}%
\pgfsetstrokecolor{currentstroke}%
\pgfsetdash{}{0pt}%
\pgfpathmoveto{\pgfqpoint{3.478725in}{1.327485in}}%
\pgfpathcurveto{\pgfqpoint{3.484549in}{1.327485in}}{\pgfqpoint{3.490135in}{1.329799in}}{\pgfqpoint{3.494253in}{1.333917in}}%
\pgfpathcurveto{\pgfqpoint{3.498371in}{1.338035in}}{\pgfqpoint{3.500685in}{1.343621in}}{\pgfqpoint{3.500685in}{1.349445in}}%
\pgfpathcurveto{\pgfqpoint{3.500685in}{1.355269in}}{\pgfqpoint{3.498371in}{1.360855in}}{\pgfqpoint{3.494253in}{1.364973in}}%
\pgfpathcurveto{\pgfqpoint{3.490135in}{1.369092in}}{\pgfqpoint{3.484549in}{1.371405in}}{\pgfqpoint{3.478725in}{1.371405in}}%
\pgfpathcurveto{\pgfqpoint{3.472901in}{1.371405in}}{\pgfqpoint{3.467315in}{1.369092in}}{\pgfqpoint{3.463197in}{1.364973in}}%
\pgfpathcurveto{\pgfqpoint{3.459079in}{1.360855in}}{\pgfqpoint{3.456765in}{1.355269in}}{\pgfqpoint{3.456765in}{1.349445in}}%
\pgfpathcurveto{\pgfqpoint{3.456765in}{1.343621in}}{\pgfqpoint{3.459079in}{1.338035in}}{\pgfqpoint{3.463197in}{1.333917in}}%
\pgfpathcurveto{\pgfqpoint{3.467315in}{1.329799in}}{\pgfqpoint{3.472901in}{1.327485in}}{\pgfqpoint{3.478725in}{1.327485in}}%
\pgfpathlineto{\pgfqpoint{3.478725in}{1.327485in}}%
\pgfpathclose%
\pgfusepath{stroke,fill}%
\end{pgfscope}%
\begin{pgfscope}%
\pgfpathrectangle{\pgfqpoint{1.414736in}{0.548769in}}{\pgfqpoint{3.497747in}{7.253186in}}%
\pgfusepath{clip}%
\pgfsetbuttcap%
\pgfsetroundjoin%
\definecolor{currentfill}{rgb}{0.412913,0.803041,0.357269}%
\pgfsetfillcolor{currentfill}%
\pgfsetlinewidth{1.003750pt}%
\definecolor{currentstroke}{rgb}{0.412913,0.803041,0.357269}%
\pgfsetstrokecolor{currentstroke}%
\pgfsetdash{}{0pt}%
\pgfpathmoveto{\pgfqpoint{3.410962in}{1.798471in}}%
\pgfpathcurveto{\pgfqpoint{3.416786in}{1.798471in}}{\pgfqpoint{3.422372in}{1.800785in}}{\pgfqpoint{3.426490in}{1.804903in}}%
\pgfpathcurveto{\pgfqpoint{3.430609in}{1.809021in}}{\pgfqpoint{3.432922in}{1.814607in}}{\pgfqpoint{3.432922in}{1.820431in}}%
\pgfpathcurveto{\pgfqpoint{3.432922in}{1.826255in}}{\pgfqpoint{3.430609in}{1.831841in}}{\pgfqpoint{3.426490in}{1.835960in}}%
\pgfpathcurveto{\pgfqpoint{3.422372in}{1.840078in}}{\pgfqpoint{3.416786in}{1.842392in}}{\pgfqpoint{3.410962in}{1.842392in}}%
\pgfpathcurveto{\pgfqpoint{3.405138in}{1.842392in}}{\pgfqpoint{3.399552in}{1.840078in}}{\pgfqpoint{3.395434in}{1.835960in}}%
\pgfpathcurveto{\pgfqpoint{3.391316in}{1.831841in}}{\pgfqpoint{3.389002in}{1.826255in}}{\pgfqpoint{3.389002in}{1.820431in}}%
\pgfpathcurveto{\pgfqpoint{3.389002in}{1.814607in}}{\pgfqpoint{3.391316in}{1.809021in}}{\pgfqpoint{3.395434in}{1.804903in}}%
\pgfpathcurveto{\pgfqpoint{3.399552in}{1.800785in}}{\pgfqpoint{3.405138in}{1.798471in}}{\pgfqpoint{3.410962in}{1.798471in}}%
\pgfpathlineto{\pgfqpoint{3.410962in}{1.798471in}}%
\pgfpathclose%
\pgfusepath{stroke,fill}%
\end{pgfscope}%
\begin{pgfscope}%
\pgfpathrectangle{\pgfqpoint{1.414736in}{0.548769in}}{\pgfqpoint{3.497747in}{7.253186in}}%
\pgfusepath{clip}%
\pgfsetbuttcap%
\pgfsetroundjoin%
\definecolor{currentfill}{rgb}{0.993248,0.906157,0.143936}%
\pgfsetfillcolor{currentfill}%
\pgfsetlinewidth{1.003750pt}%
\definecolor{currentstroke}{rgb}{0.993248,0.906157,0.143936}%
\pgfsetstrokecolor{currentstroke}%
\pgfsetdash{}{0pt}%
\pgfpathmoveto{\pgfqpoint{3.336650in}{2.269457in}}%
\pgfpathcurveto{\pgfqpoint{3.342474in}{2.269457in}}{\pgfqpoint{3.348060in}{2.271771in}}{\pgfqpoint{3.352178in}{2.275889in}}%
\pgfpathcurveto{\pgfqpoint{3.356296in}{2.280007in}}{\pgfqpoint{3.358610in}{2.285593in}}{\pgfqpoint{3.358610in}{2.291417in}}%
\pgfpathcurveto{\pgfqpoint{3.358610in}{2.297241in}}{\pgfqpoint{3.356296in}{2.302828in}}{\pgfqpoint{3.352178in}{2.306946in}}%
\pgfpathcurveto{\pgfqpoint{3.348060in}{2.311064in}}{\pgfqpoint{3.342474in}{2.313378in}}{\pgfqpoint{3.336650in}{2.313378in}}%
\pgfpathcurveto{\pgfqpoint{3.330826in}{2.313378in}}{\pgfqpoint{3.325240in}{2.311064in}}{\pgfqpoint{3.321121in}{2.306946in}}%
\pgfpathcurveto{\pgfqpoint{3.317003in}{2.302828in}}{\pgfqpoint{3.314689in}{2.297241in}}{\pgfqpoint{3.314689in}{2.291417in}}%
\pgfpathcurveto{\pgfqpoint{3.314689in}{2.285593in}}{\pgfqpoint{3.317003in}{2.280007in}}{\pgfqpoint{3.321121in}{2.275889in}}%
\pgfpathcurveto{\pgfqpoint{3.325240in}{2.271771in}}{\pgfqpoint{3.330826in}{2.269457in}}{\pgfqpoint{3.336650in}{2.269457in}}%
\pgfpathlineto{\pgfqpoint{3.336650in}{2.269457in}}%
\pgfpathclose%
\pgfusepath{stroke,fill}%
\end{pgfscope}%
\begin{pgfscope}%
\pgfpathrectangle{\pgfqpoint{1.414736in}{0.548769in}}{\pgfqpoint{3.497747in}{7.253186in}}%
\pgfusepath{clip}%
\pgfsetbuttcap%
\pgfsetroundjoin%
\definecolor{currentfill}{rgb}{0.232815,0.732247,0.459277}%
\pgfsetfillcolor{currentfill}%
\pgfsetlinewidth{1.003750pt}%
\definecolor{currentstroke}{rgb}{0.232815,0.732247,0.459277}%
\pgfsetstrokecolor{currentstroke}%
\pgfsetdash{}{0pt}%
\pgfpathmoveto{\pgfqpoint{3.862348in}{2.740443in}}%
\pgfpathcurveto{\pgfqpoint{3.868172in}{2.740443in}}{\pgfqpoint{3.873758in}{2.742757in}}{\pgfqpoint{3.877876in}{2.746875in}}%
\pgfpathcurveto{\pgfqpoint{3.881994in}{2.750993in}}{\pgfqpoint{3.884308in}{2.756580in}}{\pgfqpoint{3.884308in}{2.762404in}}%
\pgfpathcurveto{\pgfqpoint{3.884308in}{2.768227in}}{\pgfqpoint{3.881994in}{2.773814in}}{\pgfqpoint{3.877876in}{2.777932in}}%
\pgfpathcurveto{\pgfqpoint{3.873758in}{2.782050in}}{\pgfqpoint{3.868172in}{2.784364in}}{\pgfqpoint{3.862348in}{2.784364in}}%
\pgfpathcurveto{\pgfqpoint{3.856524in}{2.784364in}}{\pgfqpoint{3.850938in}{2.782050in}}{\pgfqpoint{3.846820in}{2.777932in}}%
\pgfpathcurveto{\pgfqpoint{3.842701in}{2.773814in}}{\pgfqpoint{3.840387in}{2.768227in}}{\pgfqpoint{3.840387in}{2.762404in}}%
\pgfpathcurveto{\pgfqpoint{3.840387in}{2.756580in}}{\pgfqpoint{3.842701in}{2.750993in}}{\pgfqpoint{3.846820in}{2.746875in}}%
\pgfpathcurveto{\pgfqpoint{3.850938in}{2.742757in}}{\pgfqpoint{3.856524in}{2.740443in}}{\pgfqpoint{3.862348in}{2.740443in}}%
\pgfpathlineto{\pgfqpoint{3.862348in}{2.740443in}}%
\pgfpathclose%
\pgfusepath{stroke,fill}%
\end{pgfscope}%
\begin{pgfscope}%
\pgfpathrectangle{\pgfqpoint{1.414736in}{0.548769in}}{\pgfqpoint{3.497747in}{7.253186in}}%
\pgfusepath{clip}%
\pgfsetbuttcap%
\pgfsetroundjoin%
\definecolor{currentfill}{rgb}{0.496615,0.826376,0.306377}%
\pgfsetfillcolor{currentfill}%
\pgfsetlinewidth{1.003750pt}%
\definecolor{currentstroke}{rgb}{0.496615,0.826376,0.306377}%
\pgfsetstrokecolor{currentstroke}%
\pgfsetdash{}{0pt}%
\pgfpathmoveto{\pgfqpoint{3.403781in}{3.211429in}}%
\pgfpathcurveto{\pgfqpoint{3.409605in}{3.211429in}}{\pgfqpoint{3.415191in}{3.213743in}}{\pgfqpoint{3.419309in}{3.217861in}}%
\pgfpathcurveto{\pgfqpoint{3.423428in}{3.221980in}}{\pgfqpoint{3.425741in}{3.227566in}}{\pgfqpoint{3.425741in}{3.233390in}}%
\pgfpathcurveto{\pgfqpoint{3.425741in}{3.239214in}}{\pgfqpoint{3.423428in}{3.244800in}}{\pgfqpoint{3.419309in}{3.248918in}}%
\pgfpathcurveto{\pgfqpoint{3.415191in}{3.253036in}}{\pgfqpoint{3.409605in}{3.255350in}}{\pgfqpoint{3.403781in}{3.255350in}}%
\pgfpathcurveto{\pgfqpoint{3.397957in}{3.255350in}}{\pgfqpoint{3.392371in}{3.253036in}}{\pgfqpoint{3.388253in}{3.248918in}}%
\pgfpathcurveto{\pgfqpoint{3.384135in}{3.244800in}}{\pgfqpoint{3.381821in}{3.239214in}}{\pgfqpoint{3.381821in}{3.233390in}}%
\pgfpathcurveto{\pgfqpoint{3.381821in}{3.227566in}}{\pgfqpoint{3.384135in}{3.221980in}}{\pgfqpoint{3.388253in}{3.217861in}}%
\pgfpathcurveto{\pgfqpoint{3.392371in}{3.213743in}}{\pgfqpoint{3.397957in}{3.211429in}}{\pgfqpoint{3.403781in}{3.211429in}}%
\pgfpathlineto{\pgfqpoint{3.403781in}{3.211429in}}%
\pgfpathclose%
\pgfusepath{stroke,fill}%
\end{pgfscope}%
\begin{pgfscope}%
\pgfpathrectangle{\pgfqpoint{1.414736in}{0.548769in}}{\pgfqpoint{3.497747in}{7.253186in}}%
\pgfusepath{clip}%
\pgfsetbuttcap%
\pgfsetroundjoin%
\definecolor{currentfill}{rgb}{0.896320,0.893616,0.096335}%
\pgfsetfillcolor{currentfill}%
\pgfsetlinewidth{1.003750pt}%
\definecolor{currentstroke}{rgb}{0.896320,0.893616,0.096335}%
\pgfsetstrokecolor{currentstroke}%
\pgfsetdash{}{0pt}%
\pgfpathmoveto{\pgfqpoint{3.587786in}{3.682415in}}%
\pgfpathcurveto{\pgfqpoint{3.593610in}{3.682415in}}{\pgfqpoint{3.599196in}{3.684729in}}{\pgfqpoint{3.603314in}{3.688847in}}%
\pgfpathcurveto{\pgfqpoint{3.607432in}{3.692966in}}{\pgfqpoint{3.609746in}{3.698552in}}{\pgfqpoint{3.609746in}{3.704376in}}%
\pgfpathcurveto{\pgfqpoint{3.609746in}{3.710200in}}{\pgfqpoint{3.607432in}{3.715786in}}{\pgfqpoint{3.603314in}{3.719904in}}%
\pgfpathcurveto{\pgfqpoint{3.599196in}{3.724022in}}{\pgfqpoint{3.593610in}{3.726336in}}{\pgfqpoint{3.587786in}{3.726336in}}%
\pgfpathcurveto{\pgfqpoint{3.581962in}{3.726336in}}{\pgfqpoint{3.576376in}{3.724022in}}{\pgfqpoint{3.572258in}{3.719904in}}%
\pgfpathcurveto{\pgfqpoint{3.568140in}{3.715786in}}{\pgfqpoint{3.565826in}{3.710200in}}{\pgfqpoint{3.565826in}{3.704376in}}%
\pgfpathcurveto{\pgfqpoint{3.565826in}{3.698552in}}{\pgfqpoint{3.568140in}{3.692966in}}{\pgfqpoint{3.572258in}{3.688847in}}%
\pgfpathcurveto{\pgfqpoint{3.576376in}{3.684729in}}{\pgfqpoint{3.581962in}{3.682415in}}{\pgfqpoint{3.587786in}{3.682415in}}%
\pgfpathlineto{\pgfqpoint{3.587786in}{3.682415in}}%
\pgfpathclose%
\pgfusepath{stroke,fill}%
\end{pgfscope}%
\begin{pgfscope}%
\pgfpathrectangle{\pgfqpoint{1.414736in}{0.548769in}}{\pgfqpoint{3.497747in}{7.253186in}}%
\pgfusepath{clip}%
\pgfsetbuttcap%
\pgfsetroundjoin%
\definecolor{currentfill}{rgb}{0.993248,0.906157,0.143936}%
\pgfsetfillcolor{currentfill}%
\pgfsetlinewidth{1.003750pt}%
\definecolor{currentstroke}{rgb}{0.993248,0.906157,0.143936}%
\pgfsetstrokecolor{currentstroke}%
\pgfsetdash{}{0pt}%
\pgfpathmoveto{\pgfqpoint{3.308240in}{4.153402in}}%
\pgfpathcurveto{\pgfqpoint{3.314064in}{4.153402in}}{\pgfqpoint{3.319650in}{4.155715in}}{\pgfqpoint{3.323768in}{4.159834in}}%
\pgfpathcurveto{\pgfqpoint{3.327886in}{4.163952in}}{\pgfqpoint{3.330200in}{4.169538in}}{\pgfqpoint{3.330200in}{4.175362in}}%
\pgfpathcurveto{\pgfqpoint{3.330200in}{4.181186in}}{\pgfqpoint{3.327886in}{4.186772in}}{\pgfqpoint{3.323768in}{4.190890in}}%
\pgfpathcurveto{\pgfqpoint{3.319650in}{4.195008in}}{\pgfqpoint{3.314064in}{4.197322in}}{\pgfqpoint{3.308240in}{4.197322in}}%
\pgfpathcurveto{\pgfqpoint{3.302416in}{4.197322in}}{\pgfqpoint{3.296830in}{4.195008in}}{\pgfqpoint{3.292712in}{4.190890in}}%
\pgfpathcurveto{\pgfqpoint{3.288593in}{4.186772in}}{\pgfqpoint{3.286280in}{4.181186in}}{\pgfqpoint{3.286280in}{4.175362in}}%
\pgfpathcurveto{\pgfqpoint{3.286280in}{4.169538in}}{\pgfqpoint{3.288593in}{4.163952in}}{\pgfqpoint{3.292712in}{4.159834in}}%
\pgfpathcurveto{\pgfqpoint{3.296830in}{4.155715in}}{\pgfqpoint{3.302416in}{4.153402in}}{\pgfqpoint{3.308240in}{4.153402in}}%
\pgfpathlineto{\pgfqpoint{3.308240in}{4.153402in}}%
\pgfpathclose%
\pgfusepath{stroke,fill}%
\end{pgfscope}%
\begin{pgfscope}%
\pgfpathrectangle{\pgfqpoint{1.414736in}{0.548769in}}{\pgfqpoint{3.497747in}{7.253186in}}%
\pgfusepath{clip}%
\pgfsetbuttcap%
\pgfsetroundjoin%
\definecolor{currentfill}{rgb}{0.955300,0.901065,0.118128}%
\pgfsetfillcolor{currentfill}%
\pgfsetlinewidth{1.003750pt}%
\definecolor{currentstroke}{rgb}{0.955300,0.901065,0.118128}%
\pgfsetstrokecolor{currentstroke}%
\pgfsetdash{}{0pt}%
\pgfpathmoveto{\pgfqpoint{3.811003in}{4.624388in}}%
\pgfpathcurveto{\pgfqpoint{3.816827in}{4.624388in}}{\pgfqpoint{3.822414in}{4.626702in}}{\pgfqpoint{3.826532in}{4.630820in}}%
\pgfpathcurveto{\pgfqpoint{3.830650in}{4.634938in}}{\pgfqpoint{3.832964in}{4.640524in}}{\pgfqpoint{3.832964in}{4.646348in}}%
\pgfpathcurveto{\pgfqpoint{3.832964in}{4.652172in}}{\pgfqpoint{3.830650in}{4.657758in}}{\pgfqpoint{3.826532in}{4.661876in}}%
\pgfpathcurveto{\pgfqpoint{3.822414in}{4.665994in}}{\pgfqpoint{3.816827in}{4.668308in}}{\pgfqpoint{3.811003in}{4.668308in}}%
\pgfpathcurveto{\pgfqpoint{3.805179in}{4.668308in}}{\pgfqpoint{3.799593in}{4.665994in}}{\pgfqpoint{3.795475in}{4.661876in}}%
\pgfpathcurveto{\pgfqpoint{3.791357in}{4.657758in}}{\pgfqpoint{3.789043in}{4.652172in}}{\pgfqpoint{3.789043in}{4.646348in}}%
\pgfpathcurveto{\pgfqpoint{3.789043in}{4.640524in}}{\pgfqpoint{3.791357in}{4.634938in}}{\pgfqpoint{3.795475in}{4.630820in}}%
\pgfpathcurveto{\pgfqpoint{3.799593in}{4.626702in}}{\pgfqpoint{3.805179in}{4.624388in}}{\pgfqpoint{3.811003in}{4.624388in}}%
\pgfpathlineto{\pgfqpoint{3.811003in}{4.624388in}}%
\pgfpathclose%
\pgfusepath{stroke,fill}%
\end{pgfscope}%
\begin{pgfscope}%
\pgfpathrectangle{\pgfqpoint{1.414736in}{0.548769in}}{\pgfqpoint{3.497747in}{7.253186in}}%
\pgfusepath{clip}%
\pgfsetbuttcap%
\pgfsetroundjoin%
\definecolor{currentfill}{rgb}{0.157729,0.485932,0.558013}%
\pgfsetfillcolor{currentfill}%
\pgfsetlinewidth{1.003750pt}%
\definecolor{currentstroke}{rgb}{0.157729,0.485932,0.558013}%
\pgfsetstrokecolor{currentstroke}%
\pgfsetdash{}{0pt}%
\pgfpathmoveto{\pgfqpoint{3.663887in}{5.095374in}}%
\pgfpathcurveto{\pgfqpoint{3.669710in}{5.095374in}}{\pgfqpoint{3.675297in}{5.097688in}}{\pgfqpoint{3.679415in}{5.101806in}}%
\pgfpathcurveto{\pgfqpoint{3.683533in}{5.105924in}}{\pgfqpoint{3.685847in}{5.111510in}}{\pgfqpoint{3.685847in}{5.117334in}}%
\pgfpathcurveto{\pgfqpoint{3.685847in}{5.123158in}}{\pgfqpoint{3.683533in}{5.128744in}}{\pgfqpoint{3.679415in}{5.132862in}}%
\pgfpathcurveto{\pgfqpoint{3.675297in}{5.136980in}}{\pgfqpoint{3.669710in}{5.139294in}}{\pgfqpoint{3.663887in}{5.139294in}}%
\pgfpathcurveto{\pgfqpoint{3.658063in}{5.139294in}}{\pgfqpoint{3.652476in}{5.136980in}}{\pgfqpoint{3.648358in}{5.132862in}}%
\pgfpathcurveto{\pgfqpoint{3.644240in}{5.128744in}}{\pgfqpoint{3.641926in}{5.123158in}}{\pgfqpoint{3.641926in}{5.117334in}}%
\pgfpathcurveto{\pgfqpoint{3.641926in}{5.111510in}}{\pgfqpoint{3.644240in}{5.105924in}}{\pgfqpoint{3.648358in}{5.101806in}}%
\pgfpathcurveto{\pgfqpoint{3.652476in}{5.097688in}}{\pgfqpoint{3.658063in}{5.095374in}}{\pgfqpoint{3.663887in}{5.095374in}}%
\pgfpathlineto{\pgfqpoint{3.663887in}{5.095374in}}%
\pgfpathclose%
\pgfusepath{stroke,fill}%
\end{pgfscope}%
\begin{pgfscope}%
\pgfpathrectangle{\pgfqpoint{1.414736in}{0.548769in}}{\pgfqpoint{3.497747in}{7.253186in}}%
\pgfusepath{clip}%
\pgfsetbuttcap%
\pgfsetroundjoin%
\definecolor{currentfill}{rgb}{0.192357,0.403199,0.555836}%
\pgfsetfillcolor{currentfill}%
\pgfsetlinewidth{1.003750pt}%
\definecolor{currentstroke}{rgb}{0.192357,0.403199,0.555836}%
\pgfsetstrokecolor{currentstroke}%
\pgfsetdash{}{0pt}%
\pgfpathmoveto{\pgfqpoint{3.558866in}{5.566360in}}%
\pgfpathcurveto{\pgfqpoint{3.564690in}{5.566360in}}{\pgfqpoint{3.570276in}{5.568674in}}{\pgfqpoint{3.574394in}{5.572792in}}%
\pgfpathcurveto{\pgfqpoint{3.578512in}{5.576910in}}{\pgfqpoint{3.580826in}{5.582496in}}{\pgfqpoint{3.580826in}{5.588320in}}%
\pgfpathcurveto{\pgfqpoint{3.580826in}{5.594144in}}{\pgfqpoint{3.578512in}{5.599730in}}{\pgfqpoint{3.574394in}{5.603848in}}%
\pgfpathcurveto{\pgfqpoint{3.570276in}{5.607967in}}{\pgfqpoint{3.564690in}{5.610280in}}{\pgfqpoint{3.558866in}{5.610280in}}%
\pgfpathcurveto{\pgfqpoint{3.553042in}{5.610280in}}{\pgfqpoint{3.547456in}{5.607967in}}{\pgfqpoint{3.543338in}{5.603848in}}%
\pgfpathcurveto{\pgfqpoint{3.539219in}{5.599730in}}{\pgfqpoint{3.536905in}{5.594144in}}{\pgfqpoint{3.536905in}{5.588320in}}%
\pgfpathcurveto{\pgfqpoint{3.536905in}{5.582496in}}{\pgfqpoint{3.539219in}{5.576910in}}{\pgfqpoint{3.543338in}{5.572792in}}%
\pgfpathcurveto{\pgfqpoint{3.547456in}{5.568674in}}{\pgfqpoint{3.553042in}{5.566360in}}{\pgfqpoint{3.558866in}{5.566360in}}%
\pgfpathlineto{\pgfqpoint{3.558866in}{5.566360in}}%
\pgfpathclose%
\pgfusepath{stroke,fill}%
\end{pgfscope}%
\begin{pgfscope}%
\pgfpathrectangle{\pgfqpoint{1.414736in}{0.548769in}}{\pgfqpoint{3.497747in}{7.253186in}}%
\pgfusepath{clip}%
\pgfsetbuttcap%
\pgfsetroundjoin%
\definecolor{currentfill}{rgb}{0.153894,0.680203,0.504172}%
\pgfsetfillcolor{currentfill}%
\pgfsetlinewidth{1.003750pt}%
\definecolor{currentstroke}{rgb}{0.153894,0.680203,0.504172}%
\pgfsetstrokecolor{currentstroke}%
\pgfsetdash{}{0pt}%
\pgfpathmoveto{\pgfqpoint{3.268132in}{6.037346in}}%
\pgfpathcurveto{\pgfqpoint{3.273956in}{6.037346in}}{\pgfqpoint{3.279542in}{6.039660in}}{\pgfqpoint{3.283660in}{6.043778in}}%
\pgfpathcurveto{\pgfqpoint{3.287778in}{6.047896in}}{\pgfqpoint{3.290092in}{6.053482in}}{\pgfqpoint{3.290092in}{6.059306in}}%
\pgfpathcurveto{\pgfqpoint{3.290092in}{6.065130in}}{\pgfqpoint{3.287778in}{6.070716in}}{\pgfqpoint{3.283660in}{6.074835in}}%
\pgfpathcurveto{\pgfqpoint{3.279542in}{6.078953in}}{\pgfqpoint{3.273956in}{6.081267in}}{\pgfqpoint{3.268132in}{6.081267in}}%
\pgfpathcurveto{\pgfqpoint{3.262308in}{6.081267in}}{\pgfqpoint{3.256722in}{6.078953in}}{\pgfqpoint{3.252604in}{6.074835in}}%
\pgfpathcurveto{\pgfqpoint{3.248486in}{6.070716in}}{\pgfqpoint{3.246172in}{6.065130in}}{\pgfqpoint{3.246172in}{6.059306in}}%
\pgfpathcurveto{\pgfqpoint{3.246172in}{6.053482in}}{\pgfqpoint{3.248486in}{6.047896in}}{\pgfqpoint{3.252604in}{6.043778in}}%
\pgfpathcurveto{\pgfqpoint{3.256722in}{6.039660in}}{\pgfqpoint{3.262308in}{6.037346in}}{\pgfqpoint{3.268132in}{6.037346in}}%
\pgfpathlineto{\pgfqpoint{3.268132in}{6.037346in}}%
\pgfpathclose%
\pgfusepath{stroke,fill}%
\end{pgfscope}%
\begin{pgfscope}%
\pgfpathrectangle{\pgfqpoint{1.414736in}{0.548769in}}{\pgfqpoint{3.497747in}{7.253186in}}%
\pgfusepath{clip}%
\pgfsetbuttcap%
\pgfsetroundjoin%
\definecolor{currentfill}{rgb}{0.575563,0.844566,0.256415}%
\pgfsetfillcolor{currentfill}%
\pgfsetlinewidth{1.003750pt}%
\definecolor{currentstroke}{rgb}{0.575563,0.844566,0.256415}%
\pgfsetstrokecolor{currentstroke}%
\pgfsetdash{}{0pt}%
\pgfpathmoveto{\pgfqpoint{3.919999in}{6.508332in}}%
\pgfpathcurveto{\pgfqpoint{3.925823in}{6.508332in}}{\pgfqpoint{3.931409in}{6.510646in}}{\pgfqpoint{3.935527in}{6.514764in}}%
\pgfpathcurveto{\pgfqpoint{3.939645in}{6.518882in}}{\pgfqpoint{3.941959in}{6.524468in}}{\pgfqpoint{3.941959in}{6.530292in}}%
\pgfpathcurveto{\pgfqpoint{3.941959in}{6.536116in}}{\pgfqpoint{3.939645in}{6.541702in}}{\pgfqpoint{3.935527in}{6.545821in}}%
\pgfpathcurveto{\pgfqpoint{3.931409in}{6.549939in}}{\pgfqpoint{3.925823in}{6.552253in}}{\pgfqpoint{3.919999in}{6.552253in}}%
\pgfpathcurveto{\pgfqpoint{3.914175in}{6.552253in}}{\pgfqpoint{3.908589in}{6.549939in}}{\pgfqpoint{3.904471in}{6.545821in}}%
\pgfpathcurveto{\pgfqpoint{3.900352in}{6.541702in}}{\pgfqpoint{3.898039in}{6.536116in}}{\pgfqpoint{3.898039in}{6.530292in}}%
\pgfpathcurveto{\pgfqpoint{3.898039in}{6.524468in}}{\pgfqpoint{3.900352in}{6.518882in}}{\pgfqpoint{3.904471in}{6.514764in}}%
\pgfpathcurveto{\pgfqpoint{3.908589in}{6.510646in}}{\pgfqpoint{3.914175in}{6.508332in}}{\pgfqpoint{3.919999in}{6.508332in}}%
\pgfpathlineto{\pgfqpoint{3.919999in}{6.508332in}}%
\pgfpathclose%
\pgfusepath{stroke,fill}%
\end{pgfscope}%
\begin{pgfscope}%
\pgfpathrectangle{\pgfqpoint{1.414736in}{0.548769in}}{\pgfqpoint{3.497747in}{7.253186in}}%
\pgfusepath{clip}%
\pgfsetbuttcap%
\pgfsetroundjoin%
\definecolor{currentfill}{rgb}{0.866013,0.889868,0.095953}%
\pgfsetfillcolor{currentfill}%
\pgfsetlinewidth{1.003750pt}%
\definecolor{currentstroke}{rgb}{0.866013,0.889868,0.095953}%
\pgfsetstrokecolor{currentstroke}%
\pgfsetdash{}{0pt}%
\pgfpathmoveto{\pgfqpoint{3.542565in}{6.979318in}}%
\pgfpathcurveto{\pgfqpoint{3.548389in}{6.979318in}}{\pgfqpoint{3.553975in}{6.981632in}}{\pgfqpoint{3.558093in}{6.985750in}}%
\pgfpathcurveto{\pgfqpoint{3.562211in}{6.989868in}}{\pgfqpoint{3.564525in}{6.995455in}}{\pgfqpoint{3.564525in}{7.001278in}}%
\pgfpathcurveto{\pgfqpoint{3.564525in}{7.007102in}}{\pgfqpoint{3.562211in}{7.012689in}}{\pgfqpoint{3.558093in}{7.016807in}}%
\pgfpathcurveto{\pgfqpoint{3.553975in}{7.020925in}}{\pgfqpoint{3.548389in}{7.023239in}}{\pgfqpoint{3.542565in}{7.023239in}}%
\pgfpathcurveto{\pgfqpoint{3.536741in}{7.023239in}}{\pgfqpoint{3.531155in}{7.020925in}}{\pgfqpoint{3.527036in}{7.016807in}}%
\pgfpathcurveto{\pgfqpoint{3.522918in}{7.012689in}}{\pgfqpoint{3.520604in}{7.007102in}}{\pgfqpoint{3.520604in}{7.001278in}}%
\pgfpathcurveto{\pgfqpoint{3.520604in}{6.995455in}}{\pgfqpoint{3.522918in}{6.989868in}}{\pgfqpoint{3.527036in}{6.985750in}}%
\pgfpathcurveto{\pgfqpoint{3.531155in}{6.981632in}}{\pgfqpoint{3.536741in}{6.979318in}}{\pgfqpoint{3.542565in}{6.979318in}}%
\pgfpathlineto{\pgfqpoint{3.542565in}{6.979318in}}%
\pgfpathclose%
\pgfusepath{stroke,fill}%
\end{pgfscope}%
\begin{pgfscope}%
\pgfpathrectangle{\pgfqpoint{1.414736in}{0.548769in}}{\pgfqpoint{3.497747in}{7.253186in}}%
\pgfusepath{clip}%
\pgfsetbuttcap%
\pgfsetroundjoin%
\definecolor{currentfill}{rgb}{0.688944,0.865448,0.182725}%
\pgfsetfillcolor{currentfill}%
\pgfsetlinewidth{1.003750pt}%
\definecolor{currentstroke}{rgb}{0.688944,0.865448,0.182725}%
\pgfsetstrokecolor{currentstroke}%
\pgfsetdash{}{0pt}%
\pgfpathmoveto{\pgfqpoint{3.613635in}{7.450304in}}%
\pgfpathcurveto{\pgfqpoint{3.619459in}{7.450304in}}{\pgfqpoint{3.625045in}{7.452618in}}{\pgfqpoint{3.629163in}{7.456736in}}%
\pgfpathcurveto{\pgfqpoint{3.633281in}{7.460854in}}{\pgfqpoint{3.635595in}{7.466441in}}{\pgfqpoint{3.635595in}{7.472265in}}%
\pgfpathcurveto{\pgfqpoint{3.635595in}{7.478089in}}{\pgfqpoint{3.633281in}{7.483675in}}{\pgfqpoint{3.629163in}{7.487793in}}%
\pgfpathcurveto{\pgfqpoint{3.625045in}{7.491911in}}{\pgfqpoint{3.619459in}{7.494225in}}{\pgfqpoint{3.613635in}{7.494225in}}%
\pgfpathcurveto{\pgfqpoint{3.607811in}{7.494225in}}{\pgfqpoint{3.602225in}{7.491911in}}{\pgfqpoint{3.598106in}{7.487793in}}%
\pgfpathcurveto{\pgfqpoint{3.593988in}{7.483675in}}{\pgfqpoint{3.591674in}{7.478089in}}{\pgfqpoint{3.591674in}{7.472265in}}%
\pgfpathcurveto{\pgfqpoint{3.591674in}{7.466441in}}{\pgfqpoint{3.593988in}{7.460854in}}{\pgfqpoint{3.598106in}{7.456736in}}%
\pgfpathcurveto{\pgfqpoint{3.602225in}{7.452618in}}{\pgfqpoint{3.607811in}{7.450304in}}{\pgfqpoint{3.613635in}{7.450304in}}%
\pgfpathlineto{\pgfqpoint{3.613635in}{7.450304in}}%
\pgfpathclose%
\pgfusepath{stroke,fill}%
\end{pgfscope}%
\begin{pgfscope}%
\pgfpathrectangle{\pgfqpoint{1.414736in}{0.548769in}}{\pgfqpoint{3.497747in}{7.253186in}}%
\pgfusepath{clip}%
\pgfsetbuttcap%
\pgfsetroundjoin%
\definecolor{currentfill}{rgb}{0.926106,0.897330,0.104071}%
\pgfsetfillcolor{currentfill}%
\pgfsetlinewidth{1.003750pt}%
\definecolor{currentstroke}{rgb}{0.926106,0.897330,0.104071}%
\pgfsetstrokecolor{currentstroke}%
\pgfsetdash{}{0pt}%
\pgfpathmoveto{\pgfqpoint{3.752792in}{0.856499in}}%
\pgfpathcurveto{\pgfqpoint{3.758616in}{0.856499in}}{\pgfqpoint{3.764202in}{0.858813in}}{\pgfqpoint{3.768320in}{0.862931in}}%
\pgfpathcurveto{\pgfqpoint{3.772438in}{0.867049in}}{\pgfqpoint{3.774752in}{0.872635in}}{\pgfqpoint{3.774752in}{0.878459in}}%
\pgfpathcurveto{\pgfqpoint{3.774752in}{0.884283in}}{\pgfqpoint{3.772438in}{0.889869in}}{\pgfqpoint{3.768320in}{0.893987in}}%
\pgfpathcurveto{\pgfqpoint{3.764202in}{0.898105in}}{\pgfqpoint{3.758616in}{0.900419in}}{\pgfqpoint{3.752792in}{0.900419in}}%
\pgfpathcurveto{\pgfqpoint{3.746968in}{0.900419in}}{\pgfqpoint{3.741382in}{0.898105in}}{\pgfqpoint{3.737263in}{0.893987in}}%
\pgfpathcurveto{\pgfqpoint{3.733145in}{0.889869in}}{\pgfqpoint{3.730831in}{0.884283in}}{\pgfqpoint{3.730831in}{0.878459in}}%
\pgfpathcurveto{\pgfqpoint{3.730831in}{0.872635in}}{\pgfqpoint{3.733145in}{0.867049in}}{\pgfqpoint{3.737263in}{0.862931in}}%
\pgfpathcurveto{\pgfqpoint{3.741382in}{0.858813in}}{\pgfqpoint{3.746968in}{0.856499in}}{\pgfqpoint{3.752792in}{0.856499in}}%
\pgfpathlineto{\pgfqpoint{3.752792in}{0.856499in}}%
\pgfpathclose%
\pgfusepath{stroke,fill}%
\end{pgfscope}%
\begin{pgfscope}%
\pgfpathrectangle{\pgfqpoint{1.414736in}{0.548769in}}{\pgfqpoint{3.497747in}{7.253186in}}%
\pgfusepath{clip}%
\pgfsetbuttcap%
\pgfsetroundjoin%
\definecolor{currentfill}{rgb}{0.678489,0.863742,0.189503}%
\pgfsetfillcolor{currentfill}%
\pgfsetlinewidth{1.003750pt}%
\definecolor{currentstroke}{rgb}{0.678489,0.863742,0.189503}%
\pgfsetstrokecolor{currentstroke}%
\pgfsetdash{}{0pt}%
\pgfpathmoveto{\pgfqpoint{3.454497in}{1.327485in}}%
\pgfpathcurveto{\pgfqpoint{3.460321in}{1.327485in}}{\pgfqpoint{3.465907in}{1.329799in}}{\pgfqpoint{3.470025in}{1.333917in}}%
\pgfpathcurveto{\pgfqpoint{3.474143in}{1.338035in}}{\pgfqpoint{3.476457in}{1.343621in}}{\pgfqpoint{3.476457in}{1.349445in}}%
\pgfpathcurveto{\pgfqpoint{3.476457in}{1.355269in}}{\pgfqpoint{3.474143in}{1.360855in}}{\pgfqpoint{3.470025in}{1.364973in}}%
\pgfpathcurveto{\pgfqpoint{3.465907in}{1.369092in}}{\pgfqpoint{3.460321in}{1.371405in}}{\pgfqpoint{3.454497in}{1.371405in}}%
\pgfpathcurveto{\pgfqpoint{3.448673in}{1.371405in}}{\pgfqpoint{3.443087in}{1.369092in}}{\pgfqpoint{3.438968in}{1.364973in}}%
\pgfpathcurveto{\pgfqpoint{3.434850in}{1.360855in}}{\pgfqpoint{3.432536in}{1.355269in}}{\pgfqpoint{3.432536in}{1.349445in}}%
\pgfpathcurveto{\pgfqpoint{3.432536in}{1.343621in}}{\pgfqpoint{3.434850in}{1.338035in}}{\pgfqpoint{3.438968in}{1.333917in}}%
\pgfpathcurveto{\pgfqpoint{3.443087in}{1.329799in}}{\pgfqpoint{3.448673in}{1.327485in}}{\pgfqpoint{3.454497in}{1.327485in}}%
\pgfpathlineto{\pgfqpoint{3.454497in}{1.327485in}}%
\pgfpathclose%
\pgfusepath{stroke,fill}%
\end{pgfscope}%
\begin{pgfscope}%
\pgfpathrectangle{\pgfqpoint{1.414736in}{0.548769in}}{\pgfqpoint{3.497747in}{7.253186in}}%
\pgfusepath{clip}%
\pgfsetbuttcap%
\pgfsetroundjoin%
\definecolor{currentfill}{rgb}{0.699415,0.867117,0.175971}%
\pgfsetfillcolor{currentfill}%
\pgfsetlinewidth{1.003750pt}%
\definecolor{currentstroke}{rgb}{0.699415,0.867117,0.175971}%
\pgfsetstrokecolor{currentstroke}%
\pgfsetdash{}{0pt}%
\pgfpathmoveto{\pgfqpoint{3.747288in}{1.798471in}}%
\pgfpathcurveto{\pgfqpoint{3.753112in}{1.798471in}}{\pgfqpoint{3.758698in}{1.800785in}}{\pgfqpoint{3.762816in}{1.804903in}}%
\pgfpathcurveto{\pgfqpoint{3.766934in}{1.809021in}}{\pgfqpoint{3.769248in}{1.814607in}}{\pgfqpoint{3.769248in}{1.820431in}}%
\pgfpathcurveto{\pgfqpoint{3.769248in}{1.826255in}}{\pgfqpoint{3.766934in}{1.831841in}}{\pgfqpoint{3.762816in}{1.835960in}}%
\pgfpathcurveto{\pgfqpoint{3.758698in}{1.840078in}}{\pgfqpoint{3.753112in}{1.842392in}}{\pgfqpoint{3.747288in}{1.842392in}}%
\pgfpathcurveto{\pgfqpoint{3.741464in}{1.842392in}}{\pgfqpoint{3.735878in}{1.840078in}}{\pgfqpoint{3.731760in}{1.835960in}}%
\pgfpathcurveto{\pgfqpoint{3.727642in}{1.831841in}}{\pgfqpoint{3.725328in}{1.826255in}}{\pgfqpoint{3.725328in}{1.820431in}}%
\pgfpathcurveto{\pgfqpoint{3.725328in}{1.814607in}}{\pgfqpoint{3.727642in}{1.809021in}}{\pgfqpoint{3.731760in}{1.804903in}}%
\pgfpathcurveto{\pgfqpoint{3.735878in}{1.800785in}}{\pgfqpoint{3.741464in}{1.798471in}}{\pgfqpoint{3.747288in}{1.798471in}}%
\pgfpathlineto{\pgfqpoint{3.747288in}{1.798471in}}%
\pgfpathclose%
\pgfusepath{stroke,fill}%
\end{pgfscope}%
\begin{pgfscope}%
\pgfpathrectangle{\pgfqpoint{1.414736in}{0.548769in}}{\pgfqpoint{3.497747in}{7.253186in}}%
\pgfusepath{clip}%
\pgfsetbuttcap%
\pgfsetroundjoin%
\definecolor{currentfill}{rgb}{0.993248,0.906157,0.143936}%
\pgfsetfillcolor{currentfill}%
\pgfsetlinewidth{1.003750pt}%
\definecolor{currentstroke}{rgb}{0.993248,0.906157,0.143936}%
\pgfsetstrokecolor{currentstroke}%
\pgfsetdash{}{0pt}%
\pgfpathmoveto{\pgfqpoint{3.796393in}{2.269457in}}%
\pgfpathcurveto{\pgfqpoint{3.802217in}{2.269457in}}{\pgfqpoint{3.807803in}{2.271771in}}{\pgfqpoint{3.811921in}{2.275889in}}%
\pgfpathcurveto{\pgfqpoint{3.816039in}{2.280007in}}{\pgfqpoint{3.818353in}{2.285593in}}{\pgfqpoint{3.818353in}{2.291417in}}%
\pgfpathcurveto{\pgfqpoint{3.818353in}{2.297241in}}{\pgfqpoint{3.816039in}{2.302828in}}{\pgfqpoint{3.811921in}{2.306946in}}%
\pgfpathcurveto{\pgfqpoint{3.807803in}{2.311064in}}{\pgfqpoint{3.802217in}{2.313378in}}{\pgfqpoint{3.796393in}{2.313378in}}%
\pgfpathcurveto{\pgfqpoint{3.790569in}{2.313378in}}{\pgfqpoint{3.784983in}{2.311064in}}{\pgfqpoint{3.780864in}{2.306946in}}%
\pgfpathcurveto{\pgfqpoint{3.776746in}{2.302828in}}{\pgfqpoint{3.774432in}{2.297241in}}{\pgfqpoint{3.774432in}{2.291417in}}%
\pgfpathcurveto{\pgfqpoint{3.774432in}{2.285593in}}{\pgfqpoint{3.776746in}{2.280007in}}{\pgfqpoint{3.780864in}{2.275889in}}%
\pgfpathcurveto{\pgfqpoint{3.784983in}{2.271771in}}{\pgfqpoint{3.790569in}{2.269457in}}{\pgfqpoint{3.796393in}{2.269457in}}%
\pgfpathlineto{\pgfqpoint{3.796393in}{2.269457in}}%
\pgfpathclose%
\pgfusepath{stroke,fill}%
\end{pgfscope}%
\begin{pgfscope}%
\pgfpathrectangle{\pgfqpoint{1.414736in}{0.548769in}}{\pgfqpoint{3.497747in}{7.253186in}}%
\pgfusepath{clip}%
\pgfsetbuttcap%
\pgfsetroundjoin%
\definecolor{currentfill}{rgb}{0.252899,0.742211,0.448284}%
\pgfsetfillcolor{currentfill}%
\pgfsetlinewidth{1.003750pt}%
\definecolor{currentstroke}{rgb}{0.252899,0.742211,0.448284}%
\pgfsetstrokecolor{currentstroke}%
\pgfsetdash{}{0pt}%
\pgfpathmoveto{\pgfqpoint{3.771522in}{2.740443in}}%
\pgfpathcurveto{\pgfqpoint{3.777346in}{2.740443in}}{\pgfqpoint{3.782932in}{2.742757in}}{\pgfqpoint{3.787050in}{2.746875in}}%
\pgfpathcurveto{\pgfqpoint{3.791168in}{2.750993in}}{\pgfqpoint{3.793482in}{2.756580in}}{\pgfqpoint{3.793482in}{2.762404in}}%
\pgfpathcurveto{\pgfqpoint{3.793482in}{2.768227in}}{\pgfqpoint{3.791168in}{2.773814in}}{\pgfqpoint{3.787050in}{2.777932in}}%
\pgfpathcurveto{\pgfqpoint{3.782932in}{2.782050in}}{\pgfqpoint{3.777346in}{2.784364in}}{\pgfqpoint{3.771522in}{2.784364in}}%
\pgfpathcurveto{\pgfqpoint{3.765698in}{2.784364in}}{\pgfqpoint{3.760112in}{2.782050in}}{\pgfqpoint{3.755994in}{2.777932in}}%
\pgfpathcurveto{\pgfqpoint{3.751875in}{2.773814in}}{\pgfqpoint{3.749562in}{2.768227in}}{\pgfqpoint{3.749562in}{2.762404in}}%
\pgfpathcurveto{\pgfqpoint{3.749562in}{2.756580in}}{\pgfqpoint{3.751875in}{2.750993in}}{\pgfqpoint{3.755994in}{2.746875in}}%
\pgfpathcurveto{\pgfqpoint{3.760112in}{2.742757in}}{\pgfqpoint{3.765698in}{2.740443in}}{\pgfqpoint{3.771522in}{2.740443in}}%
\pgfpathlineto{\pgfqpoint{3.771522in}{2.740443in}}%
\pgfpathclose%
\pgfusepath{stroke,fill}%
\end{pgfscope}%
\begin{pgfscope}%
\pgfpathrectangle{\pgfqpoint{1.414736in}{0.548769in}}{\pgfqpoint{3.497747in}{7.253186in}}%
\pgfusepath{clip}%
\pgfsetbuttcap%
\pgfsetroundjoin%
\definecolor{currentfill}{rgb}{0.876168,0.891125,0.095250}%
\pgfsetfillcolor{currentfill}%
\pgfsetlinewidth{1.003750pt}%
\definecolor{currentstroke}{rgb}{0.876168,0.891125,0.095250}%
\pgfsetstrokecolor{currentstroke}%
\pgfsetdash{}{0pt}%
\pgfpathmoveto{\pgfqpoint{3.509472in}{3.211429in}}%
\pgfpathcurveto{\pgfqpoint{3.515296in}{3.211429in}}{\pgfqpoint{3.520882in}{3.213743in}}{\pgfqpoint{3.525001in}{3.217861in}}%
\pgfpathcurveto{\pgfqpoint{3.529119in}{3.221980in}}{\pgfqpoint{3.531433in}{3.227566in}}{\pgfqpoint{3.531433in}{3.233390in}}%
\pgfpathcurveto{\pgfqpoint{3.531433in}{3.239214in}}{\pgfqpoint{3.529119in}{3.244800in}}{\pgfqpoint{3.525001in}{3.248918in}}%
\pgfpathcurveto{\pgfqpoint{3.520882in}{3.253036in}}{\pgfqpoint{3.515296in}{3.255350in}}{\pgfqpoint{3.509472in}{3.255350in}}%
\pgfpathcurveto{\pgfqpoint{3.503648in}{3.255350in}}{\pgfqpoint{3.498062in}{3.253036in}}{\pgfqpoint{3.493944in}{3.248918in}}%
\pgfpathcurveto{\pgfqpoint{3.489826in}{3.244800in}}{\pgfqpoint{3.487512in}{3.239214in}}{\pgfqpoint{3.487512in}{3.233390in}}%
\pgfpathcurveto{\pgfqpoint{3.487512in}{3.227566in}}{\pgfqpoint{3.489826in}{3.221980in}}{\pgfqpoint{3.493944in}{3.217861in}}%
\pgfpathcurveto{\pgfqpoint{3.498062in}{3.213743in}}{\pgfqpoint{3.503648in}{3.211429in}}{\pgfqpoint{3.509472in}{3.211429in}}%
\pgfpathlineto{\pgfqpoint{3.509472in}{3.211429in}}%
\pgfpathclose%
\pgfusepath{stroke,fill}%
\end{pgfscope}%
\begin{pgfscope}%
\pgfpathrectangle{\pgfqpoint{1.414736in}{0.548769in}}{\pgfqpoint{3.497747in}{7.253186in}}%
\pgfusepath{clip}%
\pgfsetbuttcap%
\pgfsetroundjoin%
\definecolor{currentfill}{rgb}{0.876168,0.891125,0.095250}%
\pgfsetfillcolor{currentfill}%
\pgfsetlinewidth{1.003750pt}%
\definecolor{currentstroke}{rgb}{0.876168,0.891125,0.095250}%
\pgfsetstrokecolor{currentstroke}%
\pgfsetdash{}{0pt}%
\pgfpathmoveto{\pgfqpoint{3.487113in}{3.682415in}}%
\pgfpathcurveto{\pgfqpoint{3.492937in}{3.682415in}}{\pgfqpoint{3.498523in}{3.684729in}}{\pgfqpoint{3.502641in}{3.688847in}}%
\pgfpathcurveto{\pgfqpoint{3.506760in}{3.692966in}}{\pgfqpoint{3.509073in}{3.698552in}}{\pgfqpoint{3.509073in}{3.704376in}}%
\pgfpathcurveto{\pgfqpoint{3.509073in}{3.710200in}}{\pgfqpoint{3.506760in}{3.715786in}}{\pgfqpoint{3.502641in}{3.719904in}}%
\pgfpathcurveto{\pgfqpoint{3.498523in}{3.724022in}}{\pgfqpoint{3.492937in}{3.726336in}}{\pgfqpoint{3.487113in}{3.726336in}}%
\pgfpathcurveto{\pgfqpoint{3.481289in}{3.726336in}}{\pgfqpoint{3.475703in}{3.724022in}}{\pgfqpoint{3.471585in}{3.719904in}}%
\pgfpathcurveto{\pgfqpoint{3.467467in}{3.715786in}}{\pgfqpoint{3.465153in}{3.710200in}}{\pgfqpoint{3.465153in}{3.704376in}}%
\pgfpathcurveto{\pgfqpoint{3.465153in}{3.698552in}}{\pgfqpoint{3.467467in}{3.692966in}}{\pgfqpoint{3.471585in}{3.688847in}}%
\pgfpathcurveto{\pgfqpoint{3.475703in}{3.684729in}}{\pgfqpoint{3.481289in}{3.682415in}}{\pgfqpoint{3.487113in}{3.682415in}}%
\pgfpathlineto{\pgfqpoint{3.487113in}{3.682415in}}%
\pgfpathclose%
\pgfusepath{stroke,fill}%
\end{pgfscope}%
\begin{pgfscope}%
\pgfpathrectangle{\pgfqpoint{1.414736in}{0.548769in}}{\pgfqpoint{3.497747in}{7.253186in}}%
\pgfusepath{clip}%
\pgfsetbuttcap%
\pgfsetroundjoin%
\definecolor{currentfill}{rgb}{0.606045,0.850733,0.236712}%
\pgfsetfillcolor{currentfill}%
\pgfsetlinewidth{1.003750pt}%
\definecolor{currentstroke}{rgb}{0.606045,0.850733,0.236712}%
\pgfsetstrokecolor{currentstroke}%
\pgfsetdash{}{0pt}%
\pgfpathmoveto{\pgfqpoint{3.675151in}{4.153402in}}%
\pgfpathcurveto{\pgfqpoint{3.680975in}{4.153402in}}{\pgfqpoint{3.686561in}{4.155715in}}{\pgfqpoint{3.690679in}{4.159834in}}%
\pgfpathcurveto{\pgfqpoint{3.694797in}{4.163952in}}{\pgfqpoint{3.697111in}{4.169538in}}{\pgfqpoint{3.697111in}{4.175362in}}%
\pgfpathcurveto{\pgfqpoint{3.697111in}{4.181186in}}{\pgfqpoint{3.694797in}{4.186772in}}{\pgfqpoint{3.690679in}{4.190890in}}%
\pgfpathcurveto{\pgfqpoint{3.686561in}{4.195008in}}{\pgfqpoint{3.680975in}{4.197322in}}{\pgfqpoint{3.675151in}{4.197322in}}%
\pgfpathcurveto{\pgfqpoint{3.669327in}{4.197322in}}{\pgfqpoint{3.663741in}{4.195008in}}{\pgfqpoint{3.659623in}{4.190890in}}%
\pgfpathcurveto{\pgfqpoint{3.655505in}{4.186772in}}{\pgfqpoint{3.653191in}{4.181186in}}{\pgfqpoint{3.653191in}{4.175362in}}%
\pgfpathcurveto{\pgfqpoint{3.653191in}{4.169538in}}{\pgfqpoint{3.655505in}{4.163952in}}{\pgfqpoint{3.659623in}{4.159834in}}%
\pgfpathcurveto{\pgfqpoint{3.663741in}{4.155715in}}{\pgfqpoint{3.669327in}{4.153402in}}{\pgfqpoint{3.675151in}{4.153402in}}%
\pgfpathlineto{\pgfqpoint{3.675151in}{4.153402in}}%
\pgfpathclose%
\pgfusepath{stroke,fill}%
\end{pgfscope}%
\begin{pgfscope}%
\pgfpathrectangle{\pgfqpoint{1.414736in}{0.548769in}}{\pgfqpoint{3.497747in}{7.253186in}}%
\pgfusepath{clip}%
\pgfsetbuttcap%
\pgfsetroundjoin%
\definecolor{currentfill}{rgb}{0.876168,0.891125,0.095250}%
\pgfsetfillcolor{currentfill}%
\pgfsetlinewidth{1.003750pt}%
\definecolor{currentstroke}{rgb}{0.876168,0.891125,0.095250}%
\pgfsetstrokecolor{currentstroke}%
\pgfsetdash{}{0pt}%
\pgfpathmoveto{\pgfqpoint{3.812410in}{4.624388in}}%
\pgfpathcurveto{\pgfqpoint{3.818234in}{4.624388in}}{\pgfqpoint{3.823820in}{4.626702in}}{\pgfqpoint{3.827938in}{4.630820in}}%
\pgfpathcurveto{\pgfqpoint{3.832056in}{4.634938in}}{\pgfqpoint{3.834370in}{4.640524in}}{\pgfqpoint{3.834370in}{4.646348in}}%
\pgfpathcurveto{\pgfqpoint{3.834370in}{4.652172in}}{\pgfqpoint{3.832056in}{4.657758in}}{\pgfqpoint{3.827938in}{4.661876in}}%
\pgfpathcurveto{\pgfqpoint{3.823820in}{4.665994in}}{\pgfqpoint{3.818234in}{4.668308in}}{\pgfqpoint{3.812410in}{4.668308in}}%
\pgfpathcurveto{\pgfqpoint{3.806586in}{4.668308in}}{\pgfqpoint{3.801000in}{4.665994in}}{\pgfqpoint{3.796882in}{4.661876in}}%
\pgfpathcurveto{\pgfqpoint{3.792763in}{4.657758in}}{\pgfqpoint{3.790450in}{4.652172in}}{\pgfqpoint{3.790450in}{4.646348in}}%
\pgfpathcurveto{\pgfqpoint{3.790450in}{4.640524in}}{\pgfqpoint{3.792763in}{4.634938in}}{\pgfqpoint{3.796882in}{4.630820in}}%
\pgfpathcurveto{\pgfqpoint{3.801000in}{4.626702in}}{\pgfqpoint{3.806586in}{4.624388in}}{\pgfqpoint{3.812410in}{4.624388in}}%
\pgfpathlineto{\pgfqpoint{3.812410in}{4.624388in}}%
\pgfpathclose%
\pgfusepath{stroke,fill}%
\end{pgfscope}%
\begin{pgfscope}%
\pgfpathrectangle{\pgfqpoint{1.414736in}{0.548769in}}{\pgfqpoint{3.497747in}{7.253186in}}%
\pgfusepath{clip}%
\pgfsetbuttcap%
\pgfsetroundjoin%
\definecolor{currentfill}{rgb}{0.185783,0.704891,0.485273}%
\pgfsetfillcolor{currentfill}%
\pgfsetlinewidth{1.003750pt}%
\definecolor{currentstroke}{rgb}{0.185783,0.704891,0.485273}%
\pgfsetstrokecolor{currentstroke}%
\pgfsetdash{}{0pt}%
\pgfpathmoveto{\pgfqpoint{3.574773in}{5.095374in}}%
\pgfpathcurveto{\pgfqpoint{3.580597in}{5.095374in}}{\pgfqpoint{3.586183in}{5.097688in}}{\pgfqpoint{3.590301in}{5.101806in}}%
\pgfpathcurveto{\pgfqpoint{3.594419in}{5.105924in}}{\pgfqpoint{3.596733in}{5.111510in}}{\pgfqpoint{3.596733in}{5.117334in}}%
\pgfpathcurveto{\pgfqpoint{3.596733in}{5.123158in}}{\pgfqpoint{3.594419in}{5.128744in}}{\pgfqpoint{3.590301in}{5.132862in}}%
\pgfpathcurveto{\pgfqpoint{3.586183in}{5.136980in}}{\pgfqpoint{3.580597in}{5.139294in}}{\pgfqpoint{3.574773in}{5.139294in}}%
\pgfpathcurveto{\pgfqpoint{3.568949in}{5.139294in}}{\pgfqpoint{3.563363in}{5.136980in}}{\pgfqpoint{3.559244in}{5.132862in}}%
\pgfpathcurveto{\pgfqpoint{3.555126in}{5.128744in}}{\pgfqpoint{3.552812in}{5.123158in}}{\pgfqpoint{3.552812in}{5.117334in}}%
\pgfpathcurveto{\pgfqpoint{3.552812in}{5.111510in}}{\pgfqpoint{3.555126in}{5.105924in}}{\pgfqpoint{3.559244in}{5.101806in}}%
\pgfpathcurveto{\pgfqpoint{3.563363in}{5.097688in}}{\pgfqpoint{3.568949in}{5.095374in}}{\pgfqpoint{3.574773in}{5.095374in}}%
\pgfpathlineto{\pgfqpoint{3.574773in}{5.095374in}}%
\pgfpathclose%
\pgfusepath{stroke,fill}%
\end{pgfscope}%
\begin{pgfscope}%
\pgfpathrectangle{\pgfqpoint{1.414736in}{0.548769in}}{\pgfqpoint{3.497747in}{7.253186in}}%
\pgfusepath{clip}%
\pgfsetbuttcap%
\pgfsetroundjoin%
\definecolor{currentfill}{rgb}{0.136408,0.541173,0.554483}%
\pgfsetfillcolor{currentfill}%
\pgfsetlinewidth{1.003750pt}%
\definecolor{currentstroke}{rgb}{0.136408,0.541173,0.554483}%
\pgfsetstrokecolor{currentstroke}%
\pgfsetdash{}{0pt}%
\pgfpathmoveto{\pgfqpoint{3.692594in}{5.566360in}}%
\pgfpathcurveto{\pgfqpoint{3.698417in}{5.566360in}}{\pgfqpoint{3.704004in}{5.568674in}}{\pgfqpoint{3.708122in}{5.572792in}}%
\pgfpathcurveto{\pgfqpoint{3.712240in}{5.576910in}}{\pgfqpoint{3.714554in}{5.582496in}}{\pgfqpoint{3.714554in}{5.588320in}}%
\pgfpathcurveto{\pgfqpoint{3.714554in}{5.594144in}}{\pgfqpoint{3.712240in}{5.599730in}}{\pgfqpoint{3.708122in}{5.603848in}}%
\pgfpathcurveto{\pgfqpoint{3.704004in}{5.607967in}}{\pgfqpoint{3.698417in}{5.610280in}}{\pgfqpoint{3.692594in}{5.610280in}}%
\pgfpathcurveto{\pgfqpoint{3.686770in}{5.610280in}}{\pgfqpoint{3.681183in}{5.607967in}}{\pgfqpoint{3.677065in}{5.603848in}}%
\pgfpathcurveto{\pgfqpoint{3.672947in}{5.599730in}}{\pgfqpoint{3.670633in}{5.594144in}}{\pgfqpoint{3.670633in}{5.588320in}}%
\pgfpathcurveto{\pgfqpoint{3.670633in}{5.582496in}}{\pgfqpoint{3.672947in}{5.576910in}}{\pgfqpoint{3.677065in}{5.572792in}}%
\pgfpathcurveto{\pgfqpoint{3.681183in}{5.568674in}}{\pgfqpoint{3.686770in}{5.566360in}}{\pgfqpoint{3.692594in}{5.566360in}}%
\pgfpathlineto{\pgfqpoint{3.692594in}{5.566360in}}%
\pgfpathclose%
\pgfusepath{stroke,fill}%
\end{pgfscope}%
\begin{pgfscope}%
\pgfpathrectangle{\pgfqpoint{1.414736in}{0.548769in}}{\pgfqpoint{3.497747in}{7.253186in}}%
\pgfusepath{clip}%
\pgfsetbuttcap%
\pgfsetroundjoin%
\definecolor{currentfill}{rgb}{0.162016,0.687316,0.499129}%
\pgfsetfillcolor{currentfill}%
\pgfsetlinewidth{1.003750pt}%
\definecolor{currentstroke}{rgb}{0.162016,0.687316,0.499129}%
\pgfsetstrokecolor{currentstroke}%
\pgfsetdash{}{0pt}%
\pgfpathmoveto{\pgfqpoint{3.434914in}{6.037346in}}%
\pgfpathcurveto{\pgfqpoint{3.440738in}{6.037346in}}{\pgfqpoint{3.446324in}{6.039660in}}{\pgfqpoint{3.450442in}{6.043778in}}%
\pgfpathcurveto{\pgfqpoint{3.454560in}{6.047896in}}{\pgfqpoint{3.456874in}{6.053482in}}{\pgfqpoint{3.456874in}{6.059306in}}%
\pgfpathcurveto{\pgfqpoint{3.456874in}{6.065130in}}{\pgfqpoint{3.454560in}{6.070716in}}{\pgfqpoint{3.450442in}{6.074835in}}%
\pgfpathcurveto{\pgfqpoint{3.446324in}{6.078953in}}{\pgfqpoint{3.440738in}{6.081267in}}{\pgfqpoint{3.434914in}{6.081267in}}%
\pgfpathcurveto{\pgfqpoint{3.429090in}{6.081267in}}{\pgfqpoint{3.423504in}{6.078953in}}{\pgfqpoint{3.419386in}{6.074835in}}%
\pgfpathcurveto{\pgfqpoint{3.415268in}{6.070716in}}{\pgfqpoint{3.412954in}{6.065130in}}{\pgfqpoint{3.412954in}{6.059306in}}%
\pgfpathcurveto{\pgfqpoint{3.412954in}{6.053482in}}{\pgfqpoint{3.415268in}{6.047896in}}{\pgfqpoint{3.419386in}{6.043778in}}%
\pgfpathcurveto{\pgfqpoint{3.423504in}{6.039660in}}{\pgfqpoint{3.429090in}{6.037346in}}{\pgfqpoint{3.434914in}{6.037346in}}%
\pgfpathlineto{\pgfqpoint{3.434914in}{6.037346in}}%
\pgfpathclose%
\pgfusepath{stroke,fill}%
\end{pgfscope}%
\begin{pgfscope}%
\pgfpathrectangle{\pgfqpoint{1.414736in}{0.548769in}}{\pgfqpoint{3.497747in}{7.253186in}}%
\pgfusepath{clip}%
\pgfsetbuttcap%
\pgfsetroundjoin%
\definecolor{currentfill}{rgb}{0.595839,0.848717,0.243329}%
\pgfsetfillcolor{currentfill}%
\pgfsetlinewidth{1.003750pt}%
\definecolor{currentstroke}{rgb}{0.595839,0.848717,0.243329}%
\pgfsetstrokecolor{currentstroke}%
\pgfsetdash{}{0pt}%
\pgfpathmoveto{\pgfqpoint{3.792863in}{6.508332in}}%
\pgfpathcurveto{\pgfqpoint{3.798687in}{6.508332in}}{\pgfqpoint{3.804273in}{6.510646in}}{\pgfqpoint{3.808391in}{6.514764in}}%
\pgfpathcurveto{\pgfqpoint{3.812509in}{6.518882in}}{\pgfqpoint{3.814823in}{6.524468in}}{\pgfqpoint{3.814823in}{6.530292in}}%
\pgfpathcurveto{\pgfqpoint{3.814823in}{6.536116in}}{\pgfqpoint{3.812509in}{6.541702in}}{\pgfqpoint{3.808391in}{6.545821in}}%
\pgfpathcurveto{\pgfqpoint{3.804273in}{6.549939in}}{\pgfqpoint{3.798687in}{6.552253in}}{\pgfqpoint{3.792863in}{6.552253in}}%
\pgfpathcurveto{\pgfqpoint{3.787039in}{6.552253in}}{\pgfqpoint{3.781453in}{6.549939in}}{\pgfqpoint{3.777334in}{6.545821in}}%
\pgfpathcurveto{\pgfqpoint{3.773216in}{6.541702in}}{\pgfqpoint{3.770902in}{6.536116in}}{\pgfqpoint{3.770902in}{6.530292in}}%
\pgfpathcurveto{\pgfqpoint{3.770902in}{6.524468in}}{\pgfqpoint{3.773216in}{6.518882in}}{\pgfqpoint{3.777334in}{6.514764in}}%
\pgfpathcurveto{\pgfqpoint{3.781453in}{6.510646in}}{\pgfqpoint{3.787039in}{6.508332in}}{\pgfqpoint{3.792863in}{6.508332in}}%
\pgfpathlineto{\pgfqpoint{3.792863in}{6.508332in}}%
\pgfpathclose%
\pgfusepath{stroke,fill}%
\end{pgfscope}%
\begin{pgfscope}%
\pgfpathrectangle{\pgfqpoint{1.414736in}{0.548769in}}{\pgfqpoint{3.497747in}{7.253186in}}%
\pgfusepath{clip}%
\pgfsetbuttcap%
\pgfsetroundjoin%
\definecolor{currentfill}{rgb}{0.906311,0.894855,0.098125}%
\pgfsetfillcolor{currentfill}%
\pgfsetlinewidth{1.003750pt}%
\definecolor{currentstroke}{rgb}{0.906311,0.894855,0.098125}%
\pgfsetstrokecolor{currentstroke}%
\pgfsetdash{}{0pt}%
\pgfpathmoveto{\pgfqpoint{3.792156in}{6.979318in}}%
\pgfpathcurveto{\pgfqpoint{3.797980in}{6.979318in}}{\pgfqpoint{3.803566in}{6.981632in}}{\pgfqpoint{3.807684in}{6.985750in}}%
\pgfpathcurveto{\pgfqpoint{3.811802in}{6.989868in}}{\pgfqpoint{3.814116in}{6.995455in}}{\pgfqpoint{3.814116in}{7.001278in}}%
\pgfpathcurveto{\pgfqpoint{3.814116in}{7.007102in}}{\pgfqpoint{3.811802in}{7.012689in}}{\pgfqpoint{3.807684in}{7.016807in}}%
\pgfpathcurveto{\pgfqpoint{3.803566in}{7.020925in}}{\pgfqpoint{3.797980in}{7.023239in}}{\pgfqpoint{3.792156in}{7.023239in}}%
\pgfpathcurveto{\pgfqpoint{3.786332in}{7.023239in}}{\pgfqpoint{3.780746in}{7.020925in}}{\pgfqpoint{3.776628in}{7.016807in}}%
\pgfpathcurveto{\pgfqpoint{3.772509in}{7.012689in}}{\pgfqpoint{3.770196in}{7.007102in}}{\pgfqpoint{3.770196in}{7.001278in}}%
\pgfpathcurveto{\pgfqpoint{3.770196in}{6.995455in}}{\pgfqpoint{3.772509in}{6.989868in}}{\pgfqpoint{3.776628in}{6.985750in}}%
\pgfpathcurveto{\pgfqpoint{3.780746in}{6.981632in}}{\pgfqpoint{3.786332in}{6.979318in}}{\pgfqpoint{3.792156in}{6.979318in}}%
\pgfpathlineto{\pgfqpoint{3.792156in}{6.979318in}}%
\pgfpathclose%
\pgfusepath{stroke,fill}%
\end{pgfscope}%
\begin{pgfscope}%
\pgfpathrectangle{\pgfqpoint{1.414736in}{0.548769in}}{\pgfqpoint{3.497747in}{7.253186in}}%
\pgfusepath{clip}%
\pgfsetbuttcap%
\pgfsetroundjoin%
\definecolor{currentfill}{rgb}{0.468053,0.818921,0.323998}%
\pgfsetfillcolor{currentfill}%
\pgfsetlinewidth{1.003750pt}%
\definecolor{currentstroke}{rgb}{0.468053,0.818921,0.323998}%
\pgfsetstrokecolor{currentstroke}%
\pgfsetdash{}{0pt}%
\pgfpathmoveto{\pgfqpoint{3.299398in}{7.450304in}}%
\pgfpathcurveto{\pgfqpoint{3.305221in}{7.450304in}}{\pgfqpoint{3.310808in}{7.452618in}}{\pgfqpoint{3.314926in}{7.456736in}}%
\pgfpathcurveto{\pgfqpoint{3.319044in}{7.460854in}}{\pgfqpoint{3.321358in}{7.466441in}}{\pgfqpoint{3.321358in}{7.472265in}}%
\pgfpathcurveto{\pgfqpoint{3.321358in}{7.478089in}}{\pgfqpoint{3.319044in}{7.483675in}}{\pgfqpoint{3.314926in}{7.487793in}}%
\pgfpathcurveto{\pgfqpoint{3.310808in}{7.491911in}}{\pgfqpoint{3.305221in}{7.494225in}}{\pgfqpoint{3.299398in}{7.494225in}}%
\pgfpathcurveto{\pgfqpoint{3.293574in}{7.494225in}}{\pgfqpoint{3.287987in}{7.491911in}}{\pgfqpoint{3.283869in}{7.487793in}}%
\pgfpathcurveto{\pgfqpoint{3.279751in}{7.483675in}}{\pgfqpoint{3.277437in}{7.478089in}}{\pgfqpoint{3.277437in}{7.472265in}}%
\pgfpathcurveto{\pgfqpoint{3.277437in}{7.466441in}}{\pgfqpoint{3.279751in}{7.460854in}}{\pgfqpoint{3.283869in}{7.456736in}}%
\pgfpathcurveto{\pgfqpoint{3.287987in}{7.452618in}}{\pgfqpoint{3.293574in}{7.450304in}}{\pgfqpoint{3.299398in}{7.450304in}}%
\pgfpathlineto{\pgfqpoint{3.299398in}{7.450304in}}%
\pgfpathclose%
\pgfusepath{stroke,fill}%
\end{pgfscope}%
\begin{pgfscope}%
\pgfpathrectangle{\pgfqpoint{1.414736in}{0.548769in}}{\pgfqpoint{3.497747in}{7.253186in}}%
\pgfusepath{clip}%
\pgfsetbuttcap%
\pgfsetroundjoin%
\definecolor{currentfill}{rgb}{0.993248,0.906157,0.143936}%
\pgfsetfillcolor{currentfill}%
\pgfsetlinewidth{1.003750pt}%
\definecolor{currentstroke}{rgb}{0.993248,0.906157,0.143936}%
\pgfsetstrokecolor{currentstroke}%
\pgfsetdash{}{0pt}%
\pgfpathmoveto{\pgfqpoint{4.054983in}{0.856499in}}%
\pgfpathcurveto{\pgfqpoint{4.060807in}{0.856499in}}{\pgfqpoint{4.066393in}{0.858813in}}{\pgfqpoint{4.070512in}{0.862931in}}%
\pgfpathcurveto{\pgfqpoint{4.074630in}{0.867049in}}{\pgfqpoint{4.076944in}{0.872635in}}{\pgfqpoint{4.076944in}{0.878459in}}%
\pgfpathcurveto{\pgfqpoint{4.076944in}{0.884283in}}{\pgfqpoint{4.074630in}{0.889869in}}{\pgfqpoint{4.070512in}{0.893987in}}%
\pgfpathcurveto{\pgfqpoint{4.066393in}{0.898105in}}{\pgfqpoint{4.060807in}{0.900419in}}{\pgfqpoint{4.054983in}{0.900419in}}%
\pgfpathcurveto{\pgfqpoint{4.049159in}{0.900419in}}{\pgfqpoint{4.043573in}{0.898105in}}{\pgfqpoint{4.039455in}{0.893987in}}%
\pgfpathcurveto{\pgfqpoint{4.035337in}{0.889869in}}{\pgfqpoint{4.033023in}{0.884283in}}{\pgfqpoint{4.033023in}{0.878459in}}%
\pgfpathcurveto{\pgfqpoint{4.033023in}{0.872635in}}{\pgfqpoint{4.035337in}{0.867049in}}{\pgfqpoint{4.039455in}{0.862931in}}%
\pgfpathcurveto{\pgfqpoint{4.043573in}{0.858813in}}{\pgfqpoint{4.049159in}{0.856499in}}{\pgfqpoint{4.054983in}{0.856499in}}%
\pgfpathlineto{\pgfqpoint{4.054983in}{0.856499in}}%
\pgfpathclose%
\pgfusepath{stroke,fill}%
\end{pgfscope}%
\begin{pgfscope}%
\pgfpathrectangle{\pgfqpoint{1.414736in}{0.548769in}}{\pgfqpoint{3.497747in}{7.253186in}}%
\pgfusepath{clip}%
\pgfsetbuttcap%
\pgfsetroundjoin%
\definecolor{currentfill}{rgb}{0.824940,0.884720,0.106217}%
\pgfsetfillcolor{currentfill}%
\pgfsetlinewidth{1.003750pt}%
\definecolor{currentstroke}{rgb}{0.824940,0.884720,0.106217}%
\pgfsetstrokecolor{currentstroke}%
\pgfsetdash{}{0pt}%
\pgfpathmoveto{\pgfqpoint{3.229888in}{1.327485in}}%
\pgfpathcurveto{\pgfqpoint{3.235712in}{1.327485in}}{\pgfqpoint{3.241298in}{1.329799in}}{\pgfqpoint{3.245416in}{1.333917in}}%
\pgfpathcurveto{\pgfqpoint{3.249534in}{1.338035in}}{\pgfqpoint{3.251848in}{1.343621in}}{\pgfqpoint{3.251848in}{1.349445in}}%
\pgfpathcurveto{\pgfqpoint{3.251848in}{1.355269in}}{\pgfqpoint{3.249534in}{1.360855in}}{\pgfqpoint{3.245416in}{1.364973in}}%
\pgfpathcurveto{\pgfqpoint{3.241298in}{1.369092in}}{\pgfqpoint{3.235712in}{1.371405in}}{\pgfqpoint{3.229888in}{1.371405in}}%
\pgfpathcurveto{\pgfqpoint{3.224064in}{1.371405in}}{\pgfqpoint{3.218478in}{1.369092in}}{\pgfqpoint{3.214360in}{1.364973in}}%
\pgfpathcurveto{\pgfqpoint{3.210242in}{1.360855in}}{\pgfqpoint{3.207928in}{1.355269in}}{\pgfqpoint{3.207928in}{1.349445in}}%
\pgfpathcurveto{\pgfqpoint{3.207928in}{1.343621in}}{\pgfqpoint{3.210242in}{1.338035in}}{\pgfqpoint{3.214360in}{1.333917in}}%
\pgfpathcurveto{\pgfqpoint{3.218478in}{1.329799in}}{\pgfqpoint{3.224064in}{1.327485in}}{\pgfqpoint{3.229888in}{1.327485in}}%
\pgfpathlineto{\pgfqpoint{3.229888in}{1.327485in}}%
\pgfpathclose%
\pgfusepath{stroke,fill}%
\end{pgfscope}%
\begin{pgfscope}%
\pgfpathrectangle{\pgfqpoint{1.414736in}{0.548769in}}{\pgfqpoint{3.497747in}{7.253186in}}%
\pgfusepath{clip}%
\pgfsetbuttcap%
\pgfsetroundjoin%
\definecolor{currentfill}{rgb}{0.657642,0.860219,0.203082}%
\pgfsetfillcolor{currentfill}%
\pgfsetlinewidth{1.003750pt}%
\definecolor{currentstroke}{rgb}{0.657642,0.860219,0.203082}%
\pgfsetstrokecolor{currentstroke}%
\pgfsetdash{}{0pt}%
\pgfpathmoveto{\pgfqpoint{3.254166in}{1.798471in}}%
\pgfpathcurveto{\pgfqpoint{3.259990in}{1.798471in}}{\pgfqpoint{3.265576in}{1.800785in}}{\pgfqpoint{3.269694in}{1.804903in}}%
\pgfpathcurveto{\pgfqpoint{3.273813in}{1.809021in}}{\pgfqpoint{3.276126in}{1.814607in}}{\pgfqpoint{3.276126in}{1.820431in}}%
\pgfpathcurveto{\pgfqpoint{3.276126in}{1.826255in}}{\pgfqpoint{3.273813in}{1.831841in}}{\pgfqpoint{3.269694in}{1.835960in}}%
\pgfpathcurveto{\pgfqpoint{3.265576in}{1.840078in}}{\pgfqpoint{3.259990in}{1.842392in}}{\pgfqpoint{3.254166in}{1.842392in}}%
\pgfpathcurveto{\pgfqpoint{3.248342in}{1.842392in}}{\pgfqpoint{3.242756in}{1.840078in}}{\pgfqpoint{3.238638in}{1.835960in}}%
\pgfpathcurveto{\pgfqpoint{3.234520in}{1.831841in}}{\pgfqpoint{3.232206in}{1.826255in}}{\pgfqpoint{3.232206in}{1.820431in}}%
\pgfpathcurveto{\pgfqpoint{3.232206in}{1.814607in}}{\pgfqpoint{3.234520in}{1.809021in}}{\pgfqpoint{3.238638in}{1.804903in}}%
\pgfpathcurveto{\pgfqpoint{3.242756in}{1.800785in}}{\pgfqpoint{3.248342in}{1.798471in}}{\pgfqpoint{3.254166in}{1.798471in}}%
\pgfpathlineto{\pgfqpoint{3.254166in}{1.798471in}}%
\pgfpathclose%
\pgfusepath{stroke,fill}%
\end{pgfscope}%
\begin{pgfscope}%
\pgfpathrectangle{\pgfqpoint{1.414736in}{0.548769in}}{\pgfqpoint{3.497747in}{7.253186in}}%
\pgfusepath{clip}%
\pgfsetbuttcap%
\pgfsetroundjoin%
\definecolor{currentfill}{rgb}{0.993248,0.906157,0.143936}%
\pgfsetfillcolor{currentfill}%
\pgfsetlinewidth{1.003750pt}%
\definecolor{currentstroke}{rgb}{0.993248,0.906157,0.143936}%
\pgfsetstrokecolor{currentstroke}%
\pgfsetdash{}{0pt}%
\pgfpathmoveto{\pgfqpoint{3.524031in}{2.269457in}}%
\pgfpathcurveto{\pgfqpoint{3.529855in}{2.269457in}}{\pgfqpoint{3.535441in}{2.271771in}}{\pgfqpoint{3.539560in}{2.275889in}}%
\pgfpathcurveto{\pgfqpoint{3.543678in}{2.280007in}}{\pgfqpoint{3.545992in}{2.285593in}}{\pgfqpoint{3.545992in}{2.291417in}}%
\pgfpathcurveto{\pgfqpoint{3.545992in}{2.297241in}}{\pgfqpoint{3.543678in}{2.302828in}}{\pgfqpoint{3.539560in}{2.306946in}}%
\pgfpathcurveto{\pgfqpoint{3.535441in}{2.311064in}}{\pgfqpoint{3.529855in}{2.313378in}}{\pgfqpoint{3.524031in}{2.313378in}}%
\pgfpathcurveto{\pgfqpoint{3.518207in}{2.313378in}}{\pgfqpoint{3.512621in}{2.311064in}}{\pgfqpoint{3.508503in}{2.306946in}}%
\pgfpathcurveto{\pgfqpoint{3.504385in}{2.302828in}}{\pgfqpoint{3.502071in}{2.297241in}}{\pgfqpoint{3.502071in}{2.291417in}}%
\pgfpathcurveto{\pgfqpoint{3.502071in}{2.285593in}}{\pgfqpoint{3.504385in}{2.280007in}}{\pgfqpoint{3.508503in}{2.275889in}}%
\pgfpathcurveto{\pgfqpoint{3.512621in}{2.271771in}}{\pgfqpoint{3.518207in}{2.269457in}}{\pgfqpoint{3.524031in}{2.269457in}}%
\pgfpathlineto{\pgfqpoint{3.524031in}{2.269457in}}%
\pgfpathclose%
\pgfusepath{stroke,fill}%
\end{pgfscope}%
\begin{pgfscope}%
\pgfpathrectangle{\pgfqpoint{1.414736in}{0.548769in}}{\pgfqpoint{3.497747in}{7.253186in}}%
\pgfusepath{clip}%
\pgfsetbuttcap%
\pgfsetroundjoin%
\definecolor{currentfill}{rgb}{0.449368,0.813768,0.335384}%
\pgfsetfillcolor{currentfill}%
\pgfsetlinewidth{1.003750pt}%
\definecolor{currentstroke}{rgb}{0.449368,0.813768,0.335384}%
\pgfsetstrokecolor{currentstroke}%
\pgfsetdash{}{0pt}%
\pgfpathmoveto{\pgfqpoint{3.532540in}{2.740443in}}%
\pgfpathcurveto{\pgfqpoint{3.538364in}{2.740443in}}{\pgfqpoint{3.543950in}{2.742757in}}{\pgfqpoint{3.548068in}{2.746875in}}%
\pgfpathcurveto{\pgfqpoint{3.552186in}{2.750993in}}{\pgfqpoint{3.554500in}{2.756580in}}{\pgfqpoint{3.554500in}{2.762404in}}%
\pgfpathcurveto{\pgfqpoint{3.554500in}{2.768227in}}{\pgfqpoint{3.552186in}{2.773814in}}{\pgfqpoint{3.548068in}{2.777932in}}%
\pgfpathcurveto{\pgfqpoint{3.543950in}{2.782050in}}{\pgfqpoint{3.538364in}{2.784364in}}{\pgfqpoint{3.532540in}{2.784364in}}%
\pgfpathcurveto{\pgfqpoint{3.526716in}{2.784364in}}{\pgfqpoint{3.521130in}{2.782050in}}{\pgfqpoint{3.517012in}{2.777932in}}%
\pgfpathcurveto{\pgfqpoint{3.512893in}{2.773814in}}{\pgfqpoint{3.510580in}{2.768227in}}{\pgfqpoint{3.510580in}{2.762404in}}%
\pgfpathcurveto{\pgfqpoint{3.510580in}{2.756580in}}{\pgfqpoint{3.512893in}{2.750993in}}{\pgfqpoint{3.517012in}{2.746875in}}%
\pgfpathcurveto{\pgfqpoint{3.521130in}{2.742757in}}{\pgfqpoint{3.526716in}{2.740443in}}{\pgfqpoint{3.532540in}{2.740443in}}%
\pgfpathlineto{\pgfqpoint{3.532540in}{2.740443in}}%
\pgfpathclose%
\pgfusepath{stroke,fill}%
\end{pgfscope}%
\begin{pgfscope}%
\pgfpathrectangle{\pgfqpoint{1.414736in}{0.548769in}}{\pgfqpoint{3.497747in}{7.253186in}}%
\pgfusepath{clip}%
\pgfsetbuttcap%
\pgfsetroundjoin%
\definecolor{currentfill}{rgb}{0.916242,0.896091,0.100717}%
\pgfsetfillcolor{currentfill}%
\pgfsetlinewidth{1.003750pt}%
\definecolor{currentstroke}{rgb}{0.916242,0.896091,0.100717}%
\pgfsetstrokecolor{currentstroke}%
\pgfsetdash{}{0pt}%
\pgfpathmoveto{\pgfqpoint{3.793311in}{3.211429in}}%
\pgfpathcurveto{\pgfqpoint{3.799135in}{3.211429in}}{\pgfqpoint{3.804721in}{3.213743in}}{\pgfqpoint{3.808839in}{3.217861in}}%
\pgfpathcurveto{\pgfqpoint{3.812957in}{3.221980in}}{\pgfqpoint{3.815271in}{3.227566in}}{\pgfqpoint{3.815271in}{3.233390in}}%
\pgfpathcurveto{\pgfqpoint{3.815271in}{3.239214in}}{\pgfqpoint{3.812957in}{3.244800in}}{\pgfqpoint{3.808839in}{3.248918in}}%
\pgfpathcurveto{\pgfqpoint{3.804721in}{3.253036in}}{\pgfqpoint{3.799135in}{3.255350in}}{\pgfqpoint{3.793311in}{3.255350in}}%
\pgfpathcurveto{\pgfqpoint{3.787487in}{3.255350in}}{\pgfqpoint{3.781901in}{3.253036in}}{\pgfqpoint{3.777783in}{3.248918in}}%
\pgfpathcurveto{\pgfqpoint{3.773665in}{3.244800in}}{\pgfqpoint{3.771351in}{3.239214in}}{\pgfqpoint{3.771351in}{3.233390in}}%
\pgfpathcurveto{\pgfqpoint{3.771351in}{3.227566in}}{\pgfqpoint{3.773665in}{3.221980in}}{\pgfqpoint{3.777783in}{3.217861in}}%
\pgfpathcurveto{\pgfqpoint{3.781901in}{3.213743in}}{\pgfqpoint{3.787487in}{3.211429in}}{\pgfqpoint{3.793311in}{3.211429in}}%
\pgfpathlineto{\pgfqpoint{3.793311in}{3.211429in}}%
\pgfpathclose%
\pgfusepath{stroke,fill}%
\end{pgfscope}%
\begin{pgfscope}%
\pgfpathrectangle{\pgfqpoint{1.414736in}{0.548769in}}{\pgfqpoint{3.497747in}{7.253186in}}%
\pgfusepath{clip}%
\pgfsetbuttcap%
\pgfsetroundjoin%
\definecolor{currentfill}{rgb}{0.993248,0.906157,0.143936}%
\pgfsetfillcolor{currentfill}%
\pgfsetlinewidth{1.003750pt}%
\definecolor{currentstroke}{rgb}{0.993248,0.906157,0.143936}%
\pgfsetstrokecolor{currentstroke}%
\pgfsetdash{}{0pt}%
\pgfpathmoveto{\pgfqpoint{3.728679in}{3.682415in}}%
\pgfpathcurveto{\pgfqpoint{3.734503in}{3.682415in}}{\pgfqpoint{3.740089in}{3.684729in}}{\pgfqpoint{3.744207in}{3.688847in}}%
\pgfpathcurveto{\pgfqpoint{3.748325in}{3.692966in}}{\pgfqpoint{3.750639in}{3.698552in}}{\pgfqpoint{3.750639in}{3.704376in}}%
\pgfpathcurveto{\pgfqpoint{3.750639in}{3.710200in}}{\pgfqpoint{3.748325in}{3.715786in}}{\pgfqpoint{3.744207in}{3.719904in}}%
\pgfpathcurveto{\pgfqpoint{3.740089in}{3.724022in}}{\pgfqpoint{3.734503in}{3.726336in}}{\pgfqpoint{3.728679in}{3.726336in}}%
\pgfpathcurveto{\pgfqpoint{3.722855in}{3.726336in}}{\pgfqpoint{3.717269in}{3.724022in}}{\pgfqpoint{3.713151in}{3.719904in}}%
\pgfpathcurveto{\pgfqpoint{3.709032in}{3.715786in}}{\pgfqpoint{3.706719in}{3.710200in}}{\pgfqpoint{3.706719in}{3.704376in}}%
\pgfpathcurveto{\pgfqpoint{3.706719in}{3.698552in}}{\pgfqpoint{3.709032in}{3.692966in}}{\pgfqpoint{3.713151in}{3.688847in}}%
\pgfpathcurveto{\pgfqpoint{3.717269in}{3.684729in}}{\pgfqpoint{3.722855in}{3.682415in}}{\pgfqpoint{3.728679in}{3.682415in}}%
\pgfpathlineto{\pgfqpoint{3.728679in}{3.682415in}}%
\pgfpathclose%
\pgfusepath{stroke,fill}%
\end{pgfscope}%
\begin{pgfscope}%
\pgfpathrectangle{\pgfqpoint{1.414736in}{0.548769in}}{\pgfqpoint{3.497747in}{7.253186in}}%
\pgfusepath{clip}%
\pgfsetbuttcap%
\pgfsetroundjoin%
\definecolor{currentfill}{rgb}{0.993248,0.906157,0.143936}%
\pgfsetfillcolor{currentfill}%
\pgfsetlinewidth{1.003750pt}%
\definecolor{currentstroke}{rgb}{0.993248,0.906157,0.143936}%
\pgfsetstrokecolor{currentstroke}%
\pgfsetdash{}{0pt}%
\pgfpathmoveto{\pgfqpoint{3.383299in}{4.153402in}}%
\pgfpathcurveto{\pgfqpoint{3.389123in}{4.153402in}}{\pgfqpoint{3.394709in}{4.155715in}}{\pgfqpoint{3.398827in}{4.159834in}}%
\pgfpathcurveto{\pgfqpoint{3.402945in}{4.163952in}}{\pgfqpoint{3.405259in}{4.169538in}}{\pgfqpoint{3.405259in}{4.175362in}}%
\pgfpathcurveto{\pgfqpoint{3.405259in}{4.181186in}}{\pgfqpoint{3.402945in}{4.186772in}}{\pgfqpoint{3.398827in}{4.190890in}}%
\pgfpathcurveto{\pgfqpoint{3.394709in}{4.195008in}}{\pgfqpoint{3.389123in}{4.197322in}}{\pgfqpoint{3.383299in}{4.197322in}}%
\pgfpathcurveto{\pgfqpoint{3.377475in}{4.197322in}}{\pgfqpoint{3.371889in}{4.195008in}}{\pgfqpoint{3.367771in}{4.190890in}}%
\pgfpathcurveto{\pgfqpoint{3.363653in}{4.186772in}}{\pgfqpoint{3.361339in}{4.181186in}}{\pgfqpoint{3.361339in}{4.175362in}}%
\pgfpathcurveto{\pgfqpoint{3.361339in}{4.169538in}}{\pgfqpoint{3.363653in}{4.163952in}}{\pgfqpoint{3.367771in}{4.159834in}}%
\pgfpathcurveto{\pgfqpoint{3.371889in}{4.155715in}}{\pgfqpoint{3.377475in}{4.153402in}}{\pgfqpoint{3.383299in}{4.153402in}}%
\pgfpathlineto{\pgfqpoint{3.383299in}{4.153402in}}%
\pgfpathclose%
\pgfusepath{stroke,fill}%
\end{pgfscope}%
\begin{pgfscope}%
\pgfpathrectangle{\pgfqpoint{1.414736in}{0.548769in}}{\pgfqpoint{3.497747in}{7.253186in}}%
\pgfusepath{clip}%
\pgfsetbuttcap%
\pgfsetroundjoin%
\definecolor{currentfill}{rgb}{0.993248,0.906157,0.143936}%
\pgfsetfillcolor{currentfill}%
\pgfsetlinewidth{1.003750pt}%
\definecolor{currentstroke}{rgb}{0.993248,0.906157,0.143936}%
\pgfsetstrokecolor{currentstroke}%
\pgfsetdash{}{0pt}%
\pgfpathmoveto{\pgfqpoint{4.404536in}{4.624388in}}%
\pgfpathcurveto{\pgfqpoint{4.410360in}{4.624388in}}{\pgfqpoint{4.415946in}{4.626702in}}{\pgfqpoint{4.420064in}{4.630820in}}%
\pgfpathcurveto{\pgfqpoint{4.424182in}{4.634938in}}{\pgfqpoint{4.426496in}{4.640524in}}{\pgfqpoint{4.426496in}{4.646348in}}%
\pgfpathcurveto{\pgfqpoint{4.426496in}{4.652172in}}{\pgfqpoint{4.424182in}{4.657758in}}{\pgfqpoint{4.420064in}{4.661876in}}%
\pgfpathcurveto{\pgfqpoint{4.415946in}{4.665994in}}{\pgfqpoint{4.410360in}{4.668308in}}{\pgfqpoint{4.404536in}{4.668308in}}%
\pgfpathcurveto{\pgfqpoint{4.398712in}{4.668308in}}{\pgfqpoint{4.393126in}{4.665994in}}{\pgfqpoint{4.389008in}{4.661876in}}%
\pgfpathcurveto{\pgfqpoint{4.384890in}{4.657758in}}{\pgfqpoint{4.382576in}{4.652172in}}{\pgfqpoint{4.382576in}{4.646348in}}%
\pgfpathcurveto{\pgfqpoint{4.382576in}{4.640524in}}{\pgfqpoint{4.384890in}{4.634938in}}{\pgfqpoint{4.389008in}{4.630820in}}%
\pgfpathcurveto{\pgfqpoint{4.393126in}{4.626702in}}{\pgfqpoint{4.398712in}{4.624388in}}{\pgfqpoint{4.404536in}{4.624388in}}%
\pgfpathlineto{\pgfqpoint{4.404536in}{4.624388in}}%
\pgfpathclose%
\pgfusepath{stroke,fill}%
\end{pgfscope}%
\begin{pgfscope}%
\pgfpathrectangle{\pgfqpoint{1.414736in}{0.548769in}}{\pgfqpoint{3.497747in}{7.253186in}}%
\pgfusepath{clip}%
\pgfsetbuttcap%
\pgfsetroundjoin%
\definecolor{currentfill}{rgb}{0.783315,0.879285,0.125405}%
\pgfsetfillcolor{currentfill}%
\pgfsetlinewidth{1.003750pt}%
\definecolor{currentstroke}{rgb}{0.783315,0.879285,0.125405}%
\pgfsetstrokecolor{currentstroke}%
\pgfsetdash{}{0pt}%
\pgfpathmoveto{\pgfqpoint{3.202994in}{5.095374in}}%
\pgfpathcurveto{\pgfqpoint{3.208818in}{5.095374in}}{\pgfqpoint{3.214404in}{5.097688in}}{\pgfqpoint{3.218522in}{5.101806in}}%
\pgfpathcurveto{\pgfqpoint{3.222640in}{5.105924in}}{\pgfqpoint{3.224954in}{5.111510in}}{\pgfqpoint{3.224954in}{5.117334in}}%
\pgfpathcurveto{\pgfqpoint{3.224954in}{5.123158in}}{\pgfqpoint{3.222640in}{5.128744in}}{\pgfqpoint{3.218522in}{5.132862in}}%
\pgfpathcurveto{\pgfqpoint{3.214404in}{5.136980in}}{\pgfqpoint{3.208818in}{5.139294in}}{\pgfqpoint{3.202994in}{5.139294in}}%
\pgfpathcurveto{\pgfqpoint{3.197170in}{5.139294in}}{\pgfqpoint{3.191583in}{5.136980in}}{\pgfqpoint{3.187465in}{5.132862in}}%
\pgfpathcurveto{\pgfqpoint{3.183347in}{5.128744in}}{\pgfqpoint{3.181033in}{5.123158in}}{\pgfqpoint{3.181033in}{5.117334in}}%
\pgfpathcurveto{\pgfqpoint{3.181033in}{5.111510in}}{\pgfqpoint{3.183347in}{5.105924in}}{\pgfqpoint{3.187465in}{5.101806in}}%
\pgfpathcurveto{\pgfqpoint{3.191583in}{5.097688in}}{\pgfqpoint{3.197170in}{5.095374in}}{\pgfqpoint{3.202994in}{5.095374in}}%
\pgfpathlineto{\pgfqpoint{3.202994in}{5.095374in}}%
\pgfpathclose%
\pgfusepath{stroke,fill}%
\end{pgfscope}%
\begin{pgfscope}%
\pgfpathrectangle{\pgfqpoint{1.414736in}{0.548769in}}{\pgfqpoint{3.497747in}{7.253186in}}%
\pgfusepath{clip}%
\pgfsetbuttcap%
\pgfsetroundjoin%
\definecolor{currentfill}{rgb}{0.185783,0.704891,0.485273}%
\pgfsetfillcolor{currentfill}%
\pgfsetlinewidth{1.003750pt}%
\definecolor{currentstroke}{rgb}{0.185783,0.704891,0.485273}%
\pgfsetstrokecolor{currentstroke}%
\pgfsetdash{}{0pt}%
\pgfpathmoveto{\pgfqpoint{3.052048in}{5.566360in}}%
\pgfpathcurveto{\pgfqpoint{3.057872in}{5.566360in}}{\pgfqpoint{3.063458in}{5.568674in}}{\pgfqpoint{3.067576in}{5.572792in}}%
\pgfpathcurveto{\pgfqpoint{3.071694in}{5.576910in}}{\pgfqpoint{3.074008in}{5.582496in}}{\pgfqpoint{3.074008in}{5.588320in}}%
\pgfpathcurveto{\pgfqpoint{3.074008in}{5.594144in}}{\pgfqpoint{3.071694in}{5.599730in}}{\pgfqpoint{3.067576in}{5.603848in}}%
\pgfpathcurveto{\pgfqpoint{3.063458in}{5.607967in}}{\pgfqpoint{3.057872in}{5.610280in}}{\pgfqpoint{3.052048in}{5.610280in}}%
\pgfpathcurveto{\pgfqpoint{3.046224in}{5.610280in}}{\pgfqpoint{3.040638in}{5.607967in}}{\pgfqpoint{3.036520in}{5.603848in}}%
\pgfpathcurveto{\pgfqpoint{3.032402in}{5.599730in}}{\pgfqpoint{3.030088in}{5.594144in}}{\pgfqpoint{3.030088in}{5.588320in}}%
\pgfpathcurveto{\pgfqpoint{3.030088in}{5.582496in}}{\pgfqpoint{3.032402in}{5.576910in}}{\pgfqpoint{3.036520in}{5.572792in}}%
\pgfpathcurveto{\pgfqpoint{3.040638in}{5.568674in}}{\pgfqpoint{3.046224in}{5.566360in}}{\pgfqpoint{3.052048in}{5.566360in}}%
\pgfpathlineto{\pgfqpoint{3.052048in}{5.566360in}}%
\pgfpathclose%
\pgfusepath{stroke,fill}%
\end{pgfscope}%
\begin{pgfscope}%
\pgfpathrectangle{\pgfqpoint{1.414736in}{0.548769in}}{\pgfqpoint{3.497747in}{7.253186in}}%
\pgfusepath{clip}%
\pgfsetbuttcap%
\pgfsetroundjoin%
\definecolor{currentfill}{rgb}{0.993248,0.906157,0.143936}%
\pgfsetfillcolor{currentfill}%
\pgfsetlinewidth{1.003750pt}%
\definecolor{currentstroke}{rgb}{0.993248,0.906157,0.143936}%
\pgfsetstrokecolor{currentstroke}%
\pgfsetdash{}{0pt}%
\pgfpathmoveto{\pgfqpoint{3.061412in}{6.037346in}}%
\pgfpathcurveto{\pgfqpoint{3.067236in}{6.037346in}}{\pgfqpoint{3.072822in}{6.039660in}}{\pgfqpoint{3.076941in}{6.043778in}}%
\pgfpathcurveto{\pgfqpoint{3.081059in}{6.047896in}}{\pgfqpoint{3.083373in}{6.053482in}}{\pgfqpoint{3.083373in}{6.059306in}}%
\pgfpathcurveto{\pgfqpoint{3.083373in}{6.065130in}}{\pgfqpoint{3.081059in}{6.070716in}}{\pgfqpoint{3.076941in}{6.074835in}}%
\pgfpathcurveto{\pgfqpoint{3.072822in}{6.078953in}}{\pgfqpoint{3.067236in}{6.081267in}}{\pgfqpoint{3.061412in}{6.081267in}}%
\pgfpathcurveto{\pgfqpoint{3.055588in}{6.081267in}}{\pgfqpoint{3.050002in}{6.078953in}}{\pgfqpoint{3.045884in}{6.074835in}}%
\pgfpathcurveto{\pgfqpoint{3.041766in}{6.070716in}}{\pgfqpoint{3.039452in}{6.065130in}}{\pgfqpoint{3.039452in}{6.059306in}}%
\pgfpathcurveto{\pgfqpoint{3.039452in}{6.053482in}}{\pgfqpoint{3.041766in}{6.047896in}}{\pgfqpoint{3.045884in}{6.043778in}}%
\pgfpathcurveto{\pgfqpoint{3.050002in}{6.039660in}}{\pgfqpoint{3.055588in}{6.037346in}}{\pgfqpoint{3.061412in}{6.037346in}}%
\pgfpathlineto{\pgfqpoint{3.061412in}{6.037346in}}%
\pgfpathclose%
\pgfusepath{stroke,fill}%
\end{pgfscope}%
\begin{pgfscope}%
\pgfpathrectangle{\pgfqpoint{1.414736in}{0.548769in}}{\pgfqpoint{3.497747in}{7.253186in}}%
\pgfusepath{clip}%
\pgfsetbuttcap%
\pgfsetroundjoin%
\definecolor{currentfill}{rgb}{0.993248,0.906157,0.143936}%
\pgfsetfillcolor{currentfill}%
\pgfsetlinewidth{1.003750pt}%
\definecolor{currentstroke}{rgb}{0.993248,0.906157,0.143936}%
\pgfsetstrokecolor{currentstroke}%
\pgfsetdash{}{0pt}%
\pgfpathmoveto{\pgfqpoint{3.442846in}{6.508332in}}%
\pgfpathcurveto{\pgfqpoint{3.448670in}{6.508332in}}{\pgfqpoint{3.454256in}{6.510646in}}{\pgfqpoint{3.458374in}{6.514764in}}%
\pgfpathcurveto{\pgfqpoint{3.462492in}{6.518882in}}{\pgfqpoint{3.464806in}{6.524468in}}{\pgfqpoint{3.464806in}{6.530292in}}%
\pgfpathcurveto{\pgfqpoint{3.464806in}{6.536116in}}{\pgfqpoint{3.462492in}{6.541702in}}{\pgfqpoint{3.458374in}{6.545821in}}%
\pgfpathcurveto{\pgfqpoint{3.454256in}{6.549939in}}{\pgfqpoint{3.448670in}{6.552253in}}{\pgfqpoint{3.442846in}{6.552253in}}%
\pgfpathcurveto{\pgfqpoint{3.437022in}{6.552253in}}{\pgfqpoint{3.431436in}{6.549939in}}{\pgfqpoint{3.427317in}{6.545821in}}%
\pgfpathcurveto{\pgfqpoint{3.423199in}{6.541702in}}{\pgfqpoint{3.420885in}{6.536116in}}{\pgfqpoint{3.420885in}{6.530292in}}%
\pgfpathcurveto{\pgfqpoint{3.420885in}{6.524468in}}{\pgfqpoint{3.423199in}{6.518882in}}{\pgfqpoint{3.427317in}{6.514764in}}%
\pgfpathcurveto{\pgfqpoint{3.431436in}{6.510646in}}{\pgfqpoint{3.437022in}{6.508332in}}{\pgfqpoint{3.442846in}{6.508332in}}%
\pgfpathlineto{\pgfqpoint{3.442846in}{6.508332in}}%
\pgfpathclose%
\pgfusepath{stroke,fill}%
\end{pgfscope}%
\begin{pgfscope}%
\pgfpathrectangle{\pgfqpoint{1.414736in}{0.548769in}}{\pgfqpoint{3.497747in}{7.253186in}}%
\pgfusepath{clip}%
\pgfsetbuttcap%
\pgfsetroundjoin%
\definecolor{currentfill}{rgb}{0.993248,0.906157,0.143936}%
\pgfsetfillcolor{currentfill}%
\pgfsetlinewidth{1.003750pt}%
\definecolor{currentstroke}{rgb}{0.993248,0.906157,0.143936}%
\pgfsetstrokecolor{currentstroke}%
\pgfsetdash{}{0pt}%
\pgfpathmoveto{\pgfqpoint{2.774205in}{6.979318in}}%
\pgfpathcurveto{\pgfqpoint{2.780029in}{6.979318in}}{\pgfqpoint{2.785615in}{6.981632in}}{\pgfqpoint{2.789734in}{6.985750in}}%
\pgfpathcurveto{\pgfqpoint{2.793852in}{6.989868in}}{\pgfqpoint{2.796166in}{6.995455in}}{\pgfqpoint{2.796166in}{7.001278in}}%
\pgfpathcurveto{\pgfqpoint{2.796166in}{7.007102in}}{\pgfqpoint{2.793852in}{7.012689in}}{\pgfqpoint{2.789734in}{7.016807in}}%
\pgfpathcurveto{\pgfqpoint{2.785615in}{7.020925in}}{\pgfqpoint{2.780029in}{7.023239in}}{\pgfqpoint{2.774205in}{7.023239in}}%
\pgfpathcurveto{\pgfqpoint{2.768381in}{7.023239in}}{\pgfqpoint{2.762795in}{7.020925in}}{\pgfqpoint{2.758677in}{7.016807in}}%
\pgfpathcurveto{\pgfqpoint{2.754559in}{7.012689in}}{\pgfqpoint{2.752245in}{7.007102in}}{\pgfqpoint{2.752245in}{7.001278in}}%
\pgfpathcurveto{\pgfqpoint{2.752245in}{6.995455in}}{\pgfqpoint{2.754559in}{6.989868in}}{\pgfqpoint{2.758677in}{6.985750in}}%
\pgfpathcurveto{\pgfqpoint{2.762795in}{6.981632in}}{\pgfqpoint{2.768381in}{6.979318in}}{\pgfqpoint{2.774205in}{6.979318in}}%
\pgfpathlineto{\pgfqpoint{2.774205in}{6.979318in}}%
\pgfpathclose%
\pgfusepath{stroke,fill}%
\end{pgfscope}%
\begin{pgfscope}%
\pgfpathrectangle{\pgfqpoint{1.414736in}{0.548769in}}{\pgfqpoint{3.497747in}{7.253186in}}%
\pgfusepath{clip}%
\pgfsetbuttcap%
\pgfsetroundjoin%
\definecolor{currentfill}{rgb}{0.993248,0.906157,0.143936}%
\pgfsetfillcolor{currentfill}%
\pgfsetlinewidth{1.003750pt}%
\definecolor{currentstroke}{rgb}{0.993248,0.906157,0.143936}%
\pgfsetstrokecolor{currentstroke}%
\pgfsetdash{}{0pt}%
\pgfpathmoveto{\pgfqpoint{3.589930in}{7.450304in}}%
\pgfpathcurveto{\pgfqpoint{3.595754in}{7.450304in}}{\pgfqpoint{3.601340in}{7.452618in}}{\pgfqpoint{3.605459in}{7.456736in}}%
\pgfpathcurveto{\pgfqpoint{3.609577in}{7.460854in}}{\pgfqpoint{3.611891in}{7.466441in}}{\pgfqpoint{3.611891in}{7.472265in}}%
\pgfpathcurveto{\pgfqpoint{3.611891in}{7.478089in}}{\pgfqpoint{3.609577in}{7.483675in}}{\pgfqpoint{3.605459in}{7.487793in}}%
\pgfpathcurveto{\pgfqpoint{3.601340in}{7.491911in}}{\pgfqpoint{3.595754in}{7.494225in}}{\pgfqpoint{3.589930in}{7.494225in}}%
\pgfpathcurveto{\pgfqpoint{3.584106in}{7.494225in}}{\pgfqpoint{3.578520in}{7.491911in}}{\pgfqpoint{3.574402in}{7.487793in}}%
\pgfpathcurveto{\pgfqpoint{3.570284in}{7.483675in}}{\pgfqpoint{3.567970in}{7.478089in}}{\pgfqpoint{3.567970in}{7.472265in}}%
\pgfpathcurveto{\pgfqpoint{3.567970in}{7.466441in}}{\pgfqpoint{3.570284in}{7.460854in}}{\pgfqpoint{3.574402in}{7.456736in}}%
\pgfpathcurveto{\pgfqpoint{3.578520in}{7.452618in}}{\pgfqpoint{3.584106in}{7.450304in}}{\pgfqpoint{3.589930in}{7.450304in}}%
\pgfpathlineto{\pgfqpoint{3.589930in}{7.450304in}}%
\pgfpathclose%
\pgfusepath{stroke,fill}%
\end{pgfscope}%
\begin{pgfscope}%
\pgfpathrectangle{\pgfqpoint{1.414736in}{0.548769in}}{\pgfqpoint{3.497747in}{7.253186in}}%
\pgfusepath{clip}%
\pgfsetbuttcap%
\pgfsetroundjoin%
\definecolor{currentfill}{rgb}{0.993248,0.906157,0.143936}%
\pgfsetfillcolor{currentfill}%
\pgfsetlinewidth{1.003750pt}%
\definecolor{currentstroke}{rgb}{0.993248,0.906157,0.143936}%
\pgfsetstrokecolor{currentstroke}%
\pgfsetdash{}{0pt}%
\pgfpathmoveto{\pgfqpoint{3.506077in}{0.856499in}}%
\pgfpathcurveto{\pgfqpoint{3.511901in}{0.856499in}}{\pgfqpoint{3.517487in}{0.858813in}}{\pgfqpoint{3.521606in}{0.862931in}}%
\pgfpathcurveto{\pgfqpoint{3.525724in}{0.867049in}}{\pgfqpoint{3.528038in}{0.872635in}}{\pgfqpoint{3.528038in}{0.878459in}}%
\pgfpathcurveto{\pgfqpoint{3.528038in}{0.884283in}}{\pgfqpoint{3.525724in}{0.889869in}}{\pgfqpoint{3.521606in}{0.893987in}}%
\pgfpathcurveto{\pgfqpoint{3.517487in}{0.898105in}}{\pgfqpoint{3.511901in}{0.900419in}}{\pgfqpoint{3.506077in}{0.900419in}}%
\pgfpathcurveto{\pgfqpoint{3.500253in}{0.900419in}}{\pgfqpoint{3.494667in}{0.898105in}}{\pgfqpoint{3.490549in}{0.893987in}}%
\pgfpathcurveto{\pgfqpoint{3.486431in}{0.889869in}}{\pgfqpoint{3.484117in}{0.884283in}}{\pgfqpoint{3.484117in}{0.878459in}}%
\pgfpathcurveto{\pgfqpoint{3.484117in}{0.872635in}}{\pgfqpoint{3.486431in}{0.867049in}}{\pgfqpoint{3.490549in}{0.862931in}}%
\pgfpathcurveto{\pgfqpoint{3.494667in}{0.858813in}}{\pgfqpoint{3.500253in}{0.856499in}}{\pgfqpoint{3.506077in}{0.856499in}}%
\pgfpathlineto{\pgfqpoint{3.506077in}{0.856499in}}%
\pgfpathclose%
\pgfusepath{stroke,fill}%
\end{pgfscope}%
\begin{pgfscope}%
\pgfpathrectangle{\pgfqpoint{1.414736in}{0.548769in}}{\pgfqpoint{3.497747in}{7.253186in}}%
\pgfusepath{clip}%
\pgfsetbuttcap%
\pgfsetroundjoin%
\definecolor{currentfill}{rgb}{0.458674,0.816363,0.329727}%
\pgfsetfillcolor{currentfill}%
\pgfsetlinewidth{1.003750pt}%
\definecolor{currentstroke}{rgb}{0.458674,0.816363,0.329727}%
\pgfsetstrokecolor{currentstroke}%
\pgfsetdash{}{0pt}%
\pgfpathmoveto{\pgfqpoint{3.664220in}{1.327485in}}%
\pgfpathcurveto{\pgfqpoint{3.670044in}{1.327485in}}{\pgfqpoint{3.675630in}{1.329799in}}{\pgfqpoint{3.679748in}{1.333917in}}%
\pgfpathcurveto{\pgfqpoint{3.683866in}{1.338035in}}{\pgfqpoint{3.686180in}{1.343621in}}{\pgfqpoint{3.686180in}{1.349445in}}%
\pgfpathcurveto{\pgfqpoint{3.686180in}{1.355269in}}{\pgfqpoint{3.683866in}{1.360855in}}{\pgfqpoint{3.679748in}{1.364973in}}%
\pgfpathcurveto{\pgfqpoint{3.675630in}{1.369092in}}{\pgfqpoint{3.670044in}{1.371405in}}{\pgfqpoint{3.664220in}{1.371405in}}%
\pgfpathcurveto{\pgfqpoint{3.658396in}{1.371405in}}{\pgfqpoint{3.652810in}{1.369092in}}{\pgfqpoint{3.648692in}{1.364973in}}%
\pgfpathcurveto{\pgfqpoint{3.644574in}{1.360855in}}{\pgfqpoint{3.642260in}{1.355269in}}{\pgfqpoint{3.642260in}{1.349445in}}%
\pgfpathcurveto{\pgfqpoint{3.642260in}{1.343621in}}{\pgfqpoint{3.644574in}{1.338035in}}{\pgfqpoint{3.648692in}{1.333917in}}%
\pgfpathcurveto{\pgfqpoint{3.652810in}{1.329799in}}{\pgfqpoint{3.658396in}{1.327485in}}{\pgfqpoint{3.664220in}{1.327485in}}%
\pgfpathlineto{\pgfqpoint{3.664220in}{1.327485in}}%
\pgfpathclose%
\pgfusepath{stroke,fill}%
\end{pgfscope}%
\begin{pgfscope}%
\pgfpathrectangle{\pgfqpoint{1.414736in}{0.548769in}}{\pgfqpoint{3.497747in}{7.253186in}}%
\pgfusepath{clip}%
\pgfsetbuttcap%
\pgfsetroundjoin%
\definecolor{currentfill}{rgb}{0.430983,0.808473,0.346476}%
\pgfsetfillcolor{currentfill}%
\pgfsetlinewidth{1.003750pt}%
\definecolor{currentstroke}{rgb}{0.430983,0.808473,0.346476}%
\pgfsetstrokecolor{currentstroke}%
\pgfsetdash{}{0pt}%
\pgfpathmoveto{\pgfqpoint{3.421607in}{1.798471in}}%
\pgfpathcurveto{\pgfqpoint{3.427431in}{1.798471in}}{\pgfqpoint{3.433017in}{1.800785in}}{\pgfqpoint{3.437135in}{1.804903in}}%
\pgfpathcurveto{\pgfqpoint{3.441253in}{1.809021in}}{\pgfqpoint{3.443567in}{1.814607in}}{\pgfqpoint{3.443567in}{1.820431in}}%
\pgfpathcurveto{\pgfqpoint{3.443567in}{1.826255in}}{\pgfqpoint{3.441253in}{1.831841in}}{\pgfqpoint{3.437135in}{1.835960in}}%
\pgfpathcurveto{\pgfqpoint{3.433017in}{1.840078in}}{\pgfqpoint{3.427431in}{1.842392in}}{\pgfqpoint{3.421607in}{1.842392in}}%
\pgfpathcurveto{\pgfqpoint{3.415783in}{1.842392in}}{\pgfqpoint{3.410197in}{1.840078in}}{\pgfqpoint{3.406079in}{1.835960in}}%
\pgfpathcurveto{\pgfqpoint{3.401960in}{1.831841in}}{\pgfqpoint{3.399647in}{1.826255in}}{\pgfqpoint{3.399647in}{1.820431in}}%
\pgfpathcurveto{\pgfqpoint{3.399647in}{1.814607in}}{\pgfqpoint{3.401960in}{1.809021in}}{\pgfqpoint{3.406079in}{1.804903in}}%
\pgfpathcurveto{\pgfqpoint{3.410197in}{1.800785in}}{\pgfqpoint{3.415783in}{1.798471in}}{\pgfqpoint{3.421607in}{1.798471in}}%
\pgfpathlineto{\pgfqpoint{3.421607in}{1.798471in}}%
\pgfpathclose%
\pgfusepath{stroke,fill}%
\end{pgfscope}%
\begin{pgfscope}%
\pgfpathrectangle{\pgfqpoint{1.414736in}{0.548769in}}{\pgfqpoint{3.497747in}{7.253186in}}%
\pgfusepath{clip}%
\pgfsetbuttcap%
\pgfsetroundjoin%
\definecolor{currentfill}{rgb}{0.824940,0.884720,0.106217}%
\pgfsetfillcolor{currentfill}%
\pgfsetlinewidth{1.003750pt}%
\definecolor{currentstroke}{rgb}{0.824940,0.884720,0.106217}%
\pgfsetstrokecolor{currentstroke}%
\pgfsetdash{}{0pt}%
\pgfpathmoveto{\pgfqpoint{3.354258in}{2.269457in}}%
\pgfpathcurveto{\pgfqpoint{3.360082in}{2.269457in}}{\pgfqpoint{3.365668in}{2.271771in}}{\pgfqpoint{3.369786in}{2.275889in}}%
\pgfpathcurveto{\pgfqpoint{3.373904in}{2.280007in}}{\pgfqpoint{3.376218in}{2.285593in}}{\pgfqpoint{3.376218in}{2.291417in}}%
\pgfpathcurveto{\pgfqpoint{3.376218in}{2.297241in}}{\pgfqpoint{3.373904in}{2.302828in}}{\pgfqpoint{3.369786in}{2.306946in}}%
\pgfpathcurveto{\pgfqpoint{3.365668in}{2.311064in}}{\pgfqpoint{3.360082in}{2.313378in}}{\pgfqpoint{3.354258in}{2.313378in}}%
\pgfpathcurveto{\pgfqpoint{3.348434in}{2.313378in}}{\pgfqpoint{3.342848in}{2.311064in}}{\pgfqpoint{3.338730in}{2.306946in}}%
\pgfpathcurveto{\pgfqpoint{3.334611in}{2.302828in}}{\pgfqpoint{3.332298in}{2.297241in}}{\pgfqpoint{3.332298in}{2.291417in}}%
\pgfpathcurveto{\pgfqpoint{3.332298in}{2.285593in}}{\pgfqpoint{3.334611in}{2.280007in}}{\pgfqpoint{3.338730in}{2.275889in}}%
\pgfpathcurveto{\pgfqpoint{3.342848in}{2.271771in}}{\pgfqpoint{3.348434in}{2.269457in}}{\pgfqpoint{3.354258in}{2.269457in}}%
\pgfpathlineto{\pgfqpoint{3.354258in}{2.269457in}}%
\pgfpathclose%
\pgfusepath{stroke,fill}%
\end{pgfscope}%
\begin{pgfscope}%
\pgfpathrectangle{\pgfqpoint{1.414736in}{0.548769in}}{\pgfqpoint{3.497747in}{7.253186in}}%
\pgfusepath{clip}%
\pgfsetbuttcap%
\pgfsetroundjoin%
\definecolor{currentfill}{rgb}{0.202219,0.715272,0.476084}%
\pgfsetfillcolor{currentfill}%
\pgfsetlinewidth{1.003750pt}%
\definecolor{currentstroke}{rgb}{0.202219,0.715272,0.476084}%
\pgfsetstrokecolor{currentstroke}%
\pgfsetdash{}{0pt}%
\pgfpathmoveto{\pgfqpoint{4.001701in}{2.740443in}}%
\pgfpathcurveto{\pgfqpoint{4.007525in}{2.740443in}}{\pgfqpoint{4.013111in}{2.742757in}}{\pgfqpoint{4.017230in}{2.746875in}}%
\pgfpathcurveto{\pgfqpoint{4.021348in}{2.750993in}}{\pgfqpoint{4.023662in}{2.756580in}}{\pgfqpoint{4.023662in}{2.762404in}}%
\pgfpathcurveto{\pgfqpoint{4.023662in}{2.768227in}}{\pgfqpoint{4.021348in}{2.773814in}}{\pgfqpoint{4.017230in}{2.777932in}}%
\pgfpathcurveto{\pgfqpoint{4.013111in}{2.782050in}}{\pgfqpoint{4.007525in}{2.784364in}}{\pgfqpoint{4.001701in}{2.784364in}}%
\pgfpathcurveto{\pgfqpoint{3.995877in}{2.784364in}}{\pgfqpoint{3.990291in}{2.782050in}}{\pgfqpoint{3.986173in}{2.777932in}}%
\pgfpathcurveto{\pgfqpoint{3.982055in}{2.773814in}}{\pgfqpoint{3.979741in}{2.768227in}}{\pgfqpoint{3.979741in}{2.762404in}}%
\pgfpathcurveto{\pgfqpoint{3.979741in}{2.756580in}}{\pgfqpoint{3.982055in}{2.750993in}}{\pgfqpoint{3.986173in}{2.746875in}}%
\pgfpathcurveto{\pgfqpoint{3.990291in}{2.742757in}}{\pgfqpoint{3.995877in}{2.740443in}}{\pgfqpoint{4.001701in}{2.740443in}}%
\pgfpathlineto{\pgfqpoint{4.001701in}{2.740443in}}%
\pgfpathclose%
\pgfusepath{stroke,fill}%
\end{pgfscope}%
\begin{pgfscope}%
\pgfpathrectangle{\pgfqpoint{1.414736in}{0.548769in}}{\pgfqpoint{3.497747in}{7.253186in}}%
\pgfusepath{clip}%
\pgfsetbuttcap%
\pgfsetroundjoin%
\definecolor{currentfill}{rgb}{0.575563,0.844566,0.256415}%
\pgfsetfillcolor{currentfill}%
\pgfsetlinewidth{1.003750pt}%
\definecolor{currentstroke}{rgb}{0.575563,0.844566,0.256415}%
\pgfsetstrokecolor{currentstroke}%
\pgfsetdash{}{0pt}%
\pgfpathmoveto{\pgfqpoint{3.363145in}{3.211429in}}%
\pgfpathcurveto{\pgfqpoint{3.368969in}{3.211429in}}{\pgfqpoint{3.374555in}{3.213743in}}{\pgfqpoint{3.378673in}{3.217861in}}%
\pgfpathcurveto{\pgfqpoint{3.382792in}{3.221980in}}{\pgfqpoint{3.385105in}{3.227566in}}{\pgfqpoint{3.385105in}{3.233390in}}%
\pgfpathcurveto{\pgfqpoint{3.385105in}{3.239214in}}{\pgfqpoint{3.382792in}{3.244800in}}{\pgfqpoint{3.378673in}{3.248918in}}%
\pgfpathcurveto{\pgfqpoint{3.374555in}{3.253036in}}{\pgfqpoint{3.368969in}{3.255350in}}{\pgfqpoint{3.363145in}{3.255350in}}%
\pgfpathcurveto{\pgfqpoint{3.357321in}{3.255350in}}{\pgfqpoint{3.351735in}{3.253036in}}{\pgfqpoint{3.347617in}{3.248918in}}%
\pgfpathcurveto{\pgfqpoint{3.343499in}{3.244800in}}{\pgfqpoint{3.341185in}{3.239214in}}{\pgfqpoint{3.341185in}{3.233390in}}%
\pgfpathcurveto{\pgfqpoint{3.341185in}{3.227566in}}{\pgfqpoint{3.343499in}{3.221980in}}{\pgfqpoint{3.347617in}{3.217861in}}%
\pgfpathcurveto{\pgfqpoint{3.351735in}{3.213743in}}{\pgfqpoint{3.357321in}{3.211429in}}{\pgfqpoint{3.363145in}{3.211429in}}%
\pgfpathlineto{\pgfqpoint{3.363145in}{3.211429in}}%
\pgfpathclose%
\pgfusepath{stroke,fill}%
\end{pgfscope}%
\begin{pgfscope}%
\pgfpathrectangle{\pgfqpoint{1.414736in}{0.548769in}}{\pgfqpoint{3.497747in}{7.253186in}}%
\pgfusepath{clip}%
\pgfsetbuttcap%
\pgfsetroundjoin%
\definecolor{currentfill}{rgb}{0.964894,0.902323,0.123941}%
\pgfsetfillcolor{currentfill}%
\pgfsetlinewidth{1.003750pt}%
\definecolor{currentstroke}{rgb}{0.964894,0.902323,0.123941}%
\pgfsetstrokecolor{currentstroke}%
\pgfsetdash{}{0pt}%
\pgfpathmoveto{\pgfqpoint{3.593418in}{3.682415in}}%
\pgfpathcurveto{\pgfqpoint{3.599242in}{3.682415in}}{\pgfqpoint{3.604828in}{3.684729in}}{\pgfqpoint{3.608946in}{3.688847in}}%
\pgfpathcurveto{\pgfqpoint{3.613064in}{3.692966in}}{\pgfqpoint{3.615378in}{3.698552in}}{\pgfqpoint{3.615378in}{3.704376in}}%
\pgfpathcurveto{\pgfqpoint{3.615378in}{3.710200in}}{\pgfqpoint{3.613064in}{3.715786in}}{\pgfqpoint{3.608946in}{3.719904in}}%
\pgfpathcurveto{\pgfqpoint{3.604828in}{3.724022in}}{\pgfqpoint{3.599242in}{3.726336in}}{\pgfqpoint{3.593418in}{3.726336in}}%
\pgfpathcurveto{\pgfqpoint{3.587594in}{3.726336in}}{\pgfqpoint{3.582008in}{3.724022in}}{\pgfqpoint{3.577890in}{3.719904in}}%
\pgfpathcurveto{\pgfqpoint{3.573771in}{3.715786in}}{\pgfqpoint{3.571458in}{3.710200in}}{\pgfqpoint{3.571458in}{3.704376in}}%
\pgfpathcurveto{\pgfqpoint{3.571458in}{3.698552in}}{\pgfqpoint{3.573771in}{3.692966in}}{\pgfqpoint{3.577890in}{3.688847in}}%
\pgfpathcurveto{\pgfqpoint{3.582008in}{3.684729in}}{\pgfqpoint{3.587594in}{3.682415in}}{\pgfqpoint{3.593418in}{3.682415in}}%
\pgfpathlineto{\pgfqpoint{3.593418in}{3.682415in}}%
\pgfpathclose%
\pgfusepath{stroke,fill}%
\end{pgfscope}%
\begin{pgfscope}%
\pgfpathrectangle{\pgfqpoint{1.414736in}{0.548769in}}{\pgfqpoint{3.497747in}{7.253186in}}%
\pgfusepath{clip}%
\pgfsetbuttcap%
\pgfsetroundjoin%
\definecolor{currentfill}{rgb}{0.993248,0.906157,0.143936}%
\pgfsetfillcolor{currentfill}%
\pgfsetlinewidth{1.003750pt}%
\definecolor{currentstroke}{rgb}{0.993248,0.906157,0.143936}%
\pgfsetstrokecolor{currentstroke}%
\pgfsetdash{}{0pt}%
\pgfpathmoveto{\pgfqpoint{3.815026in}{4.153402in}}%
\pgfpathcurveto{\pgfqpoint{3.820850in}{4.153402in}}{\pgfqpoint{3.826436in}{4.155715in}}{\pgfqpoint{3.830554in}{4.159834in}}%
\pgfpathcurveto{\pgfqpoint{3.834672in}{4.163952in}}{\pgfqpoint{3.836986in}{4.169538in}}{\pgfqpoint{3.836986in}{4.175362in}}%
\pgfpathcurveto{\pgfqpoint{3.836986in}{4.181186in}}{\pgfqpoint{3.834672in}{4.186772in}}{\pgfqpoint{3.830554in}{4.190890in}}%
\pgfpathcurveto{\pgfqpoint{3.826436in}{4.195008in}}{\pgfqpoint{3.820850in}{4.197322in}}{\pgfqpoint{3.815026in}{4.197322in}}%
\pgfpathcurveto{\pgfqpoint{3.809202in}{4.197322in}}{\pgfqpoint{3.803616in}{4.195008in}}{\pgfqpoint{3.799498in}{4.190890in}}%
\pgfpathcurveto{\pgfqpoint{3.795380in}{4.186772in}}{\pgfqpoint{3.793066in}{4.181186in}}{\pgfqpoint{3.793066in}{4.175362in}}%
\pgfpathcurveto{\pgfqpoint{3.793066in}{4.169538in}}{\pgfqpoint{3.795380in}{4.163952in}}{\pgfqpoint{3.799498in}{4.159834in}}%
\pgfpathcurveto{\pgfqpoint{3.803616in}{4.155715in}}{\pgfqpoint{3.809202in}{4.153402in}}{\pgfqpoint{3.815026in}{4.153402in}}%
\pgfpathlineto{\pgfqpoint{3.815026in}{4.153402in}}%
\pgfpathclose%
\pgfusepath{stroke,fill}%
\end{pgfscope}%
\begin{pgfscope}%
\pgfpathrectangle{\pgfqpoint{1.414736in}{0.548769in}}{\pgfqpoint{3.497747in}{7.253186in}}%
\pgfusepath{clip}%
\pgfsetbuttcap%
\pgfsetroundjoin%
\definecolor{currentfill}{rgb}{0.804182,0.882046,0.114965}%
\pgfsetfillcolor{currentfill}%
\pgfsetlinewidth{1.003750pt}%
\definecolor{currentstroke}{rgb}{0.804182,0.882046,0.114965}%
\pgfsetstrokecolor{currentstroke}%
\pgfsetdash{}{0pt}%
\pgfpathmoveto{\pgfqpoint{3.588250in}{4.624388in}}%
\pgfpathcurveto{\pgfqpoint{3.594074in}{4.624388in}}{\pgfqpoint{3.599660in}{4.626702in}}{\pgfqpoint{3.603778in}{4.630820in}}%
\pgfpathcurveto{\pgfqpoint{3.607896in}{4.634938in}}{\pgfqpoint{3.610210in}{4.640524in}}{\pgfqpoint{3.610210in}{4.646348in}}%
\pgfpathcurveto{\pgfqpoint{3.610210in}{4.652172in}}{\pgfqpoint{3.607896in}{4.657758in}}{\pgfqpoint{3.603778in}{4.661876in}}%
\pgfpathcurveto{\pgfqpoint{3.599660in}{4.665994in}}{\pgfqpoint{3.594074in}{4.668308in}}{\pgfqpoint{3.588250in}{4.668308in}}%
\pgfpathcurveto{\pgfqpoint{3.582426in}{4.668308in}}{\pgfqpoint{3.576840in}{4.665994in}}{\pgfqpoint{3.572722in}{4.661876in}}%
\pgfpathcurveto{\pgfqpoint{3.568604in}{4.657758in}}{\pgfqpoint{3.566290in}{4.652172in}}{\pgfqpoint{3.566290in}{4.646348in}}%
\pgfpathcurveto{\pgfqpoint{3.566290in}{4.640524in}}{\pgfqpoint{3.568604in}{4.634938in}}{\pgfqpoint{3.572722in}{4.630820in}}%
\pgfpathcurveto{\pgfqpoint{3.576840in}{4.626702in}}{\pgfqpoint{3.582426in}{4.624388in}}{\pgfqpoint{3.588250in}{4.624388in}}%
\pgfpathlineto{\pgfqpoint{3.588250in}{4.624388in}}%
\pgfpathclose%
\pgfusepath{stroke,fill}%
\end{pgfscope}%
\begin{pgfscope}%
\pgfpathrectangle{\pgfqpoint{1.414736in}{0.548769in}}{\pgfqpoint{3.497747in}{7.253186in}}%
\pgfusepath{clip}%
\pgfsetbuttcap%
\pgfsetroundjoin%
\definecolor{currentfill}{rgb}{0.280894,0.078907,0.402329}%
\pgfsetfillcolor{currentfill}%
\pgfsetlinewidth{1.003750pt}%
\definecolor{currentstroke}{rgb}{0.280894,0.078907,0.402329}%
\pgfsetstrokecolor{currentstroke}%
\pgfsetdash{}{0pt}%
\pgfpathmoveto{\pgfqpoint{3.359988in}{5.095374in}}%
\pgfpathcurveto{\pgfqpoint{3.365812in}{5.095374in}}{\pgfqpoint{3.371398in}{5.097688in}}{\pgfqpoint{3.375516in}{5.101806in}}%
\pgfpathcurveto{\pgfqpoint{3.379635in}{5.105924in}}{\pgfqpoint{3.381948in}{5.111510in}}{\pgfqpoint{3.381948in}{5.117334in}}%
\pgfpathcurveto{\pgfqpoint{3.381948in}{5.123158in}}{\pgfqpoint{3.379635in}{5.128744in}}{\pgfqpoint{3.375516in}{5.132862in}}%
\pgfpathcurveto{\pgfqpoint{3.371398in}{5.136980in}}{\pgfqpoint{3.365812in}{5.139294in}}{\pgfqpoint{3.359988in}{5.139294in}}%
\pgfpathcurveto{\pgfqpoint{3.354164in}{5.139294in}}{\pgfqpoint{3.348578in}{5.136980in}}{\pgfqpoint{3.344460in}{5.132862in}}%
\pgfpathcurveto{\pgfqpoint{3.340342in}{5.128744in}}{\pgfqpoint{3.338028in}{5.123158in}}{\pgfqpoint{3.338028in}{5.117334in}}%
\pgfpathcurveto{\pgfqpoint{3.338028in}{5.111510in}}{\pgfqpoint{3.340342in}{5.105924in}}{\pgfqpoint{3.344460in}{5.101806in}}%
\pgfpathcurveto{\pgfqpoint{3.348578in}{5.097688in}}{\pgfqpoint{3.354164in}{5.095374in}}{\pgfqpoint{3.359988in}{5.095374in}}%
\pgfpathlineto{\pgfqpoint{3.359988in}{5.095374in}}%
\pgfpathclose%
\pgfusepath{stroke,fill}%
\end{pgfscope}%
\begin{pgfscope}%
\pgfpathrectangle{\pgfqpoint{1.414736in}{0.548769in}}{\pgfqpoint{3.497747in}{7.253186in}}%
\pgfusepath{clip}%
\pgfsetbuttcap%
\pgfsetroundjoin%
\definecolor{currentfill}{rgb}{0.216210,0.351535,0.550627}%
\pgfsetfillcolor{currentfill}%
\pgfsetlinewidth{1.003750pt}%
\definecolor{currentstroke}{rgb}{0.216210,0.351535,0.550627}%
\pgfsetstrokecolor{currentstroke}%
\pgfsetdash{}{0pt}%
\pgfpathmoveto{\pgfqpoint{3.713161in}{5.566360in}}%
\pgfpathcurveto{\pgfqpoint{3.718985in}{5.566360in}}{\pgfqpoint{3.724571in}{5.568674in}}{\pgfqpoint{3.728689in}{5.572792in}}%
\pgfpathcurveto{\pgfqpoint{3.732807in}{5.576910in}}{\pgfqpoint{3.735121in}{5.582496in}}{\pgfqpoint{3.735121in}{5.588320in}}%
\pgfpathcurveto{\pgfqpoint{3.735121in}{5.594144in}}{\pgfqpoint{3.732807in}{5.599730in}}{\pgfqpoint{3.728689in}{5.603848in}}%
\pgfpathcurveto{\pgfqpoint{3.724571in}{5.607967in}}{\pgfqpoint{3.718985in}{5.610280in}}{\pgfqpoint{3.713161in}{5.610280in}}%
\pgfpathcurveto{\pgfqpoint{3.707337in}{5.610280in}}{\pgfqpoint{3.701751in}{5.607967in}}{\pgfqpoint{3.697632in}{5.603848in}}%
\pgfpathcurveto{\pgfqpoint{3.693514in}{5.599730in}}{\pgfqpoint{3.691200in}{5.594144in}}{\pgfqpoint{3.691200in}{5.588320in}}%
\pgfpathcurveto{\pgfqpoint{3.691200in}{5.582496in}}{\pgfqpoint{3.693514in}{5.576910in}}{\pgfqpoint{3.697632in}{5.572792in}}%
\pgfpathcurveto{\pgfqpoint{3.701751in}{5.568674in}}{\pgfqpoint{3.707337in}{5.566360in}}{\pgfqpoint{3.713161in}{5.566360in}}%
\pgfpathlineto{\pgfqpoint{3.713161in}{5.566360in}}%
\pgfpathclose%
\pgfusepath{stroke,fill}%
\end{pgfscope}%
\begin{pgfscope}%
\pgfpathrectangle{\pgfqpoint{1.414736in}{0.548769in}}{\pgfqpoint{3.497747in}{7.253186in}}%
\pgfusepath{clip}%
\pgfsetbuttcap%
\pgfsetroundjoin%
\definecolor{currentfill}{rgb}{0.135066,0.544853,0.554029}%
\pgfsetfillcolor{currentfill}%
\pgfsetlinewidth{1.003750pt}%
\definecolor{currentstroke}{rgb}{0.135066,0.544853,0.554029}%
\pgfsetstrokecolor{currentstroke}%
\pgfsetdash{}{0pt}%
\pgfpathmoveto{\pgfqpoint{3.897681in}{6.037346in}}%
\pgfpathcurveto{\pgfqpoint{3.903505in}{6.037346in}}{\pgfqpoint{3.909091in}{6.039660in}}{\pgfqpoint{3.913210in}{6.043778in}}%
\pgfpathcurveto{\pgfqpoint{3.917328in}{6.047896in}}{\pgfqpoint{3.919642in}{6.053482in}}{\pgfqpoint{3.919642in}{6.059306in}}%
\pgfpathcurveto{\pgfqpoint{3.919642in}{6.065130in}}{\pgfqpoint{3.917328in}{6.070716in}}{\pgfqpoint{3.913210in}{6.074835in}}%
\pgfpathcurveto{\pgfqpoint{3.909091in}{6.078953in}}{\pgfqpoint{3.903505in}{6.081267in}}{\pgfqpoint{3.897681in}{6.081267in}}%
\pgfpathcurveto{\pgfqpoint{3.891857in}{6.081267in}}{\pgfqpoint{3.886271in}{6.078953in}}{\pgfqpoint{3.882153in}{6.074835in}}%
\pgfpathcurveto{\pgfqpoint{3.878035in}{6.070716in}}{\pgfqpoint{3.875721in}{6.065130in}}{\pgfqpoint{3.875721in}{6.059306in}}%
\pgfpathcurveto{\pgfqpoint{3.875721in}{6.053482in}}{\pgfqpoint{3.878035in}{6.047896in}}{\pgfqpoint{3.882153in}{6.043778in}}%
\pgfpathcurveto{\pgfqpoint{3.886271in}{6.039660in}}{\pgfqpoint{3.891857in}{6.037346in}}{\pgfqpoint{3.897681in}{6.037346in}}%
\pgfpathlineto{\pgfqpoint{3.897681in}{6.037346in}}%
\pgfpathclose%
\pgfusepath{stroke,fill}%
\end{pgfscope}%
\begin{pgfscope}%
\pgfpathrectangle{\pgfqpoint{1.414736in}{0.548769in}}{\pgfqpoint{3.497747in}{7.253186in}}%
\pgfusepath{clip}%
\pgfsetbuttcap%
\pgfsetroundjoin%
\definecolor{currentfill}{rgb}{0.993248,0.906157,0.143936}%
\pgfsetfillcolor{currentfill}%
\pgfsetlinewidth{1.003750pt}%
\definecolor{currentstroke}{rgb}{0.993248,0.906157,0.143936}%
\pgfsetstrokecolor{currentstroke}%
\pgfsetdash{}{0pt}%
\pgfpathmoveto{\pgfqpoint{3.916956in}{6.508332in}}%
\pgfpathcurveto{\pgfqpoint{3.922780in}{6.508332in}}{\pgfqpoint{3.928366in}{6.510646in}}{\pgfqpoint{3.932484in}{6.514764in}}%
\pgfpathcurveto{\pgfqpoint{3.936602in}{6.518882in}}{\pgfqpoint{3.938916in}{6.524468in}}{\pgfqpoint{3.938916in}{6.530292in}}%
\pgfpathcurveto{\pgfqpoint{3.938916in}{6.536116in}}{\pgfqpoint{3.936602in}{6.541702in}}{\pgfqpoint{3.932484in}{6.545821in}}%
\pgfpathcurveto{\pgfqpoint{3.928366in}{6.549939in}}{\pgfqpoint{3.922780in}{6.552253in}}{\pgfqpoint{3.916956in}{6.552253in}}%
\pgfpathcurveto{\pgfqpoint{3.911132in}{6.552253in}}{\pgfqpoint{3.905546in}{6.549939in}}{\pgfqpoint{3.901427in}{6.545821in}}%
\pgfpathcurveto{\pgfqpoint{3.897309in}{6.541702in}}{\pgfqpoint{3.894995in}{6.536116in}}{\pgfqpoint{3.894995in}{6.530292in}}%
\pgfpathcurveto{\pgfqpoint{3.894995in}{6.524468in}}{\pgfqpoint{3.897309in}{6.518882in}}{\pgfqpoint{3.901427in}{6.514764in}}%
\pgfpathcurveto{\pgfqpoint{3.905546in}{6.510646in}}{\pgfqpoint{3.911132in}{6.508332in}}{\pgfqpoint{3.916956in}{6.508332in}}%
\pgfpathlineto{\pgfqpoint{3.916956in}{6.508332in}}%
\pgfpathclose%
\pgfusepath{stroke,fill}%
\end{pgfscope}%
\begin{pgfscope}%
\pgfpathrectangle{\pgfqpoint{1.414736in}{0.548769in}}{\pgfqpoint{3.497747in}{7.253186in}}%
\pgfusepath{clip}%
\pgfsetbuttcap%
\pgfsetroundjoin%
\definecolor{currentfill}{rgb}{0.352360,0.783011,0.392636}%
\pgfsetfillcolor{currentfill}%
\pgfsetlinewidth{1.003750pt}%
\definecolor{currentstroke}{rgb}{0.352360,0.783011,0.392636}%
\pgfsetstrokecolor{currentstroke}%
\pgfsetdash{}{0pt}%
\pgfpathmoveto{\pgfqpoint{3.560666in}{6.979318in}}%
\pgfpathcurveto{\pgfqpoint{3.566490in}{6.979318in}}{\pgfqpoint{3.572076in}{6.981632in}}{\pgfqpoint{3.576194in}{6.985750in}}%
\pgfpathcurveto{\pgfqpoint{3.580312in}{6.989868in}}{\pgfqpoint{3.582626in}{6.995455in}}{\pgfqpoint{3.582626in}{7.001278in}}%
\pgfpathcurveto{\pgfqpoint{3.582626in}{7.007102in}}{\pgfqpoint{3.580312in}{7.012689in}}{\pgfqpoint{3.576194in}{7.016807in}}%
\pgfpathcurveto{\pgfqpoint{3.572076in}{7.020925in}}{\pgfqpoint{3.566490in}{7.023239in}}{\pgfqpoint{3.560666in}{7.023239in}}%
\pgfpathcurveto{\pgfqpoint{3.554842in}{7.023239in}}{\pgfqpoint{3.549256in}{7.020925in}}{\pgfqpoint{3.545138in}{7.016807in}}%
\pgfpathcurveto{\pgfqpoint{3.541019in}{7.012689in}}{\pgfqpoint{3.538706in}{7.007102in}}{\pgfqpoint{3.538706in}{7.001278in}}%
\pgfpathcurveto{\pgfqpoint{3.538706in}{6.995455in}}{\pgfqpoint{3.541019in}{6.989868in}}{\pgfqpoint{3.545138in}{6.985750in}}%
\pgfpathcurveto{\pgfqpoint{3.549256in}{6.981632in}}{\pgfqpoint{3.554842in}{6.979318in}}{\pgfqpoint{3.560666in}{6.979318in}}%
\pgfpathlineto{\pgfqpoint{3.560666in}{6.979318in}}%
\pgfpathclose%
\pgfusepath{stroke,fill}%
\end{pgfscope}%
\begin{pgfscope}%
\pgfpathrectangle{\pgfqpoint{1.414736in}{0.548769in}}{\pgfqpoint{3.497747in}{7.253186in}}%
\pgfusepath{clip}%
\pgfsetbuttcap%
\pgfsetroundjoin%
\definecolor{currentfill}{rgb}{0.866013,0.889868,0.095953}%
\pgfsetfillcolor{currentfill}%
\pgfsetlinewidth{1.003750pt}%
\definecolor{currentstroke}{rgb}{0.866013,0.889868,0.095953}%
\pgfsetstrokecolor{currentstroke}%
\pgfsetdash{}{0pt}%
\pgfpathmoveto{\pgfqpoint{3.443270in}{7.450304in}}%
\pgfpathcurveto{\pgfqpoint{3.449094in}{7.450304in}}{\pgfqpoint{3.454680in}{7.452618in}}{\pgfqpoint{3.458798in}{7.456736in}}%
\pgfpathcurveto{\pgfqpoint{3.462917in}{7.460854in}}{\pgfqpoint{3.465230in}{7.466441in}}{\pgfqpoint{3.465230in}{7.472265in}}%
\pgfpathcurveto{\pgfqpoint{3.465230in}{7.478089in}}{\pgfqpoint{3.462917in}{7.483675in}}{\pgfqpoint{3.458798in}{7.487793in}}%
\pgfpathcurveto{\pgfqpoint{3.454680in}{7.491911in}}{\pgfqpoint{3.449094in}{7.494225in}}{\pgfqpoint{3.443270in}{7.494225in}}%
\pgfpathcurveto{\pgfqpoint{3.437446in}{7.494225in}}{\pgfqpoint{3.431860in}{7.491911in}}{\pgfqpoint{3.427742in}{7.487793in}}%
\pgfpathcurveto{\pgfqpoint{3.423624in}{7.483675in}}{\pgfqpoint{3.421310in}{7.478089in}}{\pgfqpoint{3.421310in}{7.472265in}}%
\pgfpathcurveto{\pgfqpoint{3.421310in}{7.466441in}}{\pgfqpoint{3.423624in}{7.460854in}}{\pgfqpoint{3.427742in}{7.456736in}}%
\pgfpathcurveto{\pgfqpoint{3.431860in}{7.452618in}}{\pgfqpoint{3.437446in}{7.450304in}}{\pgfqpoint{3.443270in}{7.450304in}}%
\pgfpathlineto{\pgfqpoint{3.443270in}{7.450304in}}%
\pgfpathclose%
\pgfusepath{stroke,fill}%
\end{pgfscope}%
\begin{pgfscope}%
\pgfpathrectangle{\pgfqpoint{1.414736in}{0.548769in}}{\pgfqpoint{3.497747in}{7.253186in}}%
\pgfusepath{clip}%
\pgfsetbuttcap%
\pgfsetroundjoin%
\definecolor{currentfill}{rgb}{0.993248,0.906157,0.143936}%
\pgfsetfillcolor{currentfill}%
\pgfsetlinewidth{1.003750pt}%
\definecolor{currentstroke}{rgb}{0.993248,0.906157,0.143936}%
\pgfsetstrokecolor{currentstroke}%
\pgfsetdash{}{0pt}%
\pgfpathmoveto{\pgfqpoint{3.538333in}{0.856499in}}%
\pgfpathcurveto{\pgfqpoint{3.544157in}{0.856499in}}{\pgfqpoint{3.549743in}{0.858813in}}{\pgfqpoint{3.553861in}{0.862931in}}%
\pgfpathcurveto{\pgfqpoint{3.557979in}{0.867049in}}{\pgfqpoint{3.560293in}{0.872635in}}{\pgfqpoint{3.560293in}{0.878459in}}%
\pgfpathcurveto{\pgfqpoint{3.560293in}{0.884283in}}{\pgfqpoint{3.557979in}{0.889869in}}{\pgfqpoint{3.553861in}{0.893987in}}%
\pgfpathcurveto{\pgfqpoint{3.549743in}{0.898105in}}{\pgfqpoint{3.544157in}{0.900419in}}{\pgfqpoint{3.538333in}{0.900419in}}%
\pgfpathcurveto{\pgfqpoint{3.532509in}{0.900419in}}{\pgfqpoint{3.526923in}{0.898105in}}{\pgfqpoint{3.522805in}{0.893987in}}%
\pgfpathcurveto{\pgfqpoint{3.518686in}{0.889869in}}{\pgfqpoint{3.516373in}{0.884283in}}{\pgfqpoint{3.516373in}{0.878459in}}%
\pgfpathcurveto{\pgfqpoint{3.516373in}{0.872635in}}{\pgfqpoint{3.518686in}{0.867049in}}{\pgfqpoint{3.522805in}{0.862931in}}%
\pgfpathcurveto{\pgfqpoint{3.526923in}{0.858813in}}{\pgfqpoint{3.532509in}{0.856499in}}{\pgfqpoint{3.538333in}{0.856499in}}%
\pgfpathlineto{\pgfqpoint{3.538333in}{0.856499in}}%
\pgfpathclose%
\pgfusepath{stroke,fill}%
\end{pgfscope}%
\begin{pgfscope}%
\pgfpathrectangle{\pgfqpoint{1.414736in}{0.548769in}}{\pgfqpoint{3.497747in}{7.253186in}}%
\pgfusepath{clip}%
\pgfsetbuttcap%
\pgfsetroundjoin%
\definecolor{currentfill}{rgb}{0.535621,0.835785,0.281908}%
\pgfsetfillcolor{currentfill}%
\pgfsetlinewidth{1.003750pt}%
\definecolor{currentstroke}{rgb}{0.535621,0.835785,0.281908}%
\pgfsetstrokecolor{currentstroke}%
\pgfsetdash{}{0pt}%
\pgfpathmoveto{\pgfqpoint{3.510423in}{1.327485in}}%
\pgfpathcurveto{\pgfqpoint{3.516246in}{1.327485in}}{\pgfqpoint{3.521833in}{1.329799in}}{\pgfqpoint{3.525951in}{1.333917in}}%
\pgfpathcurveto{\pgfqpoint{3.530069in}{1.338035in}}{\pgfqpoint{3.532383in}{1.343621in}}{\pgfqpoint{3.532383in}{1.349445in}}%
\pgfpathcurveto{\pgfqpoint{3.532383in}{1.355269in}}{\pgfqpoint{3.530069in}{1.360855in}}{\pgfqpoint{3.525951in}{1.364973in}}%
\pgfpathcurveto{\pgfqpoint{3.521833in}{1.369092in}}{\pgfqpoint{3.516246in}{1.371405in}}{\pgfqpoint{3.510423in}{1.371405in}}%
\pgfpathcurveto{\pgfqpoint{3.504599in}{1.371405in}}{\pgfqpoint{3.499012in}{1.369092in}}{\pgfqpoint{3.494894in}{1.364973in}}%
\pgfpathcurveto{\pgfqpoint{3.490776in}{1.360855in}}{\pgfqpoint{3.488462in}{1.355269in}}{\pgfqpoint{3.488462in}{1.349445in}}%
\pgfpathcurveto{\pgfqpoint{3.488462in}{1.343621in}}{\pgfqpoint{3.490776in}{1.338035in}}{\pgfqpoint{3.494894in}{1.333917in}}%
\pgfpathcurveto{\pgfqpoint{3.499012in}{1.329799in}}{\pgfqpoint{3.504599in}{1.327485in}}{\pgfqpoint{3.510423in}{1.327485in}}%
\pgfpathlineto{\pgfqpoint{3.510423in}{1.327485in}}%
\pgfpathclose%
\pgfusepath{stroke,fill}%
\end{pgfscope}%
\begin{pgfscope}%
\pgfpathrectangle{\pgfqpoint{1.414736in}{0.548769in}}{\pgfqpoint{3.497747in}{7.253186in}}%
\pgfusepath{clip}%
\pgfsetbuttcap%
\pgfsetroundjoin%
\definecolor{currentfill}{rgb}{0.699415,0.867117,0.175971}%
\pgfsetfillcolor{currentfill}%
\pgfsetlinewidth{1.003750pt}%
\definecolor{currentstroke}{rgb}{0.699415,0.867117,0.175971}%
\pgfsetstrokecolor{currentstroke}%
\pgfsetdash{}{0pt}%
\pgfpathmoveto{\pgfqpoint{3.665672in}{1.798471in}}%
\pgfpathcurveto{\pgfqpoint{3.671496in}{1.798471in}}{\pgfqpoint{3.677082in}{1.800785in}}{\pgfqpoint{3.681200in}{1.804903in}}%
\pgfpathcurveto{\pgfqpoint{3.685319in}{1.809021in}}{\pgfqpoint{3.687633in}{1.814607in}}{\pgfqpoint{3.687633in}{1.820431in}}%
\pgfpathcurveto{\pgfqpoint{3.687633in}{1.826255in}}{\pgfqpoint{3.685319in}{1.831841in}}{\pgfqpoint{3.681200in}{1.835960in}}%
\pgfpathcurveto{\pgfqpoint{3.677082in}{1.840078in}}{\pgfqpoint{3.671496in}{1.842392in}}{\pgfqpoint{3.665672in}{1.842392in}}%
\pgfpathcurveto{\pgfqpoint{3.659848in}{1.842392in}}{\pgfqpoint{3.654262in}{1.840078in}}{\pgfqpoint{3.650144in}{1.835960in}}%
\pgfpathcurveto{\pgfqpoint{3.646026in}{1.831841in}}{\pgfqpoint{3.643712in}{1.826255in}}{\pgfqpoint{3.643712in}{1.820431in}}%
\pgfpathcurveto{\pgfqpoint{3.643712in}{1.814607in}}{\pgfqpoint{3.646026in}{1.809021in}}{\pgfqpoint{3.650144in}{1.804903in}}%
\pgfpathcurveto{\pgfqpoint{3.654262in}{1.800785in}}{\pgfqpoint{3.659848in}{1.798471in}}{\pgfqpoint{3.665672in}{1.798471in}}%
\pgfpathlineto{\pgfqpoint{3.665672in}{1.798471in}}%
\pgfpathclose%
\pgfusepath{stroke,fill}%
\end{pgfscope}%
\begin{pgfscope}%
\pgfpathrectangle{\pgfqpoint{1.414736in}{0.548769in}}{\pgfqpoint{3.497747in}{7.253186in}}%
\pgfusepath{clip}%
\pgfsetbuttcap%
\pgfsetroundjoin%
\definecolor{currentfill}{rgb}{0.993248,0.906157,0.143936}%
\pgfsetfillcolor{currentfill}%
\pgfsetlinewidth{1.003750pt}%
\definecolor{currentstroke}{rgb}{0.993248,0.906157,0.143936}%
\pgfsetstrokecolor{currentstroke}%
\pgfsetdash{}{0pt}%
\pgfpathmoveto{\pgfqpoint{3.383288in}{2.269457in}}%
\pgfpathcurveto{\pgfqpoint{3.389112in}{2.269457in}}{\pgfqpoint{3.394698in}{2.271771in}}{\pgfqpoint{3.398816in}{2.275889in}}%
\pgfpathcurveto{\pgfqpoint{3.402934in}{2.280007in}}{\pgfqpoint{3.405248in}{2.285593in}}{\pgfqpoint{3.405248in}{2.291417in}}%
\pgfpathcurveto{\pgfqpoint{3.405248in}{2.297241in}}{\pgfqpoint{3.402934in}{2.302828in}}{\pgfqpoint{3.398816in}{2.306946in}}%
\pgfpathcurveto{\pgfqpoint{3.394698in}{2.311064in}}{\pgfqpoint{3.389112in}{2.313378in}}{\pgfqpoint{3.383288in}{2.313378in}}%
\pgfpathcurveto{\pgfqpoint{3.377464in}{2.313378in}}{\pgfqpoint{3.371878in}{2.311064in}}{\pgfqpoint{3.367760in}{2.306946in}}%
\pgfpathcurveto{\pgfqpoint{3.363642in}{2.302828in}}{\pgfqpoint{3.361328in}{2.297241in}}{\pgfqpoint{3.361328in}{2.291417in}}%
\pgfpathcurveto{\pgfqpoint{3.361328in}{2.285593in}}{\pgfqpoint{3.363642in}{2.280007in}}{\pgfqpoint{3.367760in}{2.275889in}}%
\pgfpathcurveto{\pgfqpoint{3.371878in}{2.271771in}}{\pgfqpoint{3.377464in}{2.269457in}}{\pgfqpoint{3.383288in}{2.269457in}}%
\pgfpathlineto{\pgfqpoint{3.383288in}{2.269457in}}%
\pgfpathclose%
\pgfusepath{stroke,fill}%
\end{pgfscope}%
\begin{pgfscope}%
\pgfpathrectangle{\pgfqpoint{1.414736in}{0.548769in}}{\pgfqpoint{3.497747in}{7.253186in}}%
\pgfusepath{clip}%
\pgfsetbuttcap%
\pgfsetroundjoin%
\definecolor{currentfill}{rgb}{0.751884,0.874951,0.143228}%
\pgfsetfillcolor{currentfill}%
\pgfsetlinewidth{1.003750pt}%
\definecolor{currentstroke}{rgb}{0.751884,0.874951,0.143228}%
\pgfsetstrokecolor{currentstroke}%
\pgfsetdash{}{0pt}%
\pgfpathmoveto{\pgfqpoint{3.270179in}{2.740443in}}%
\pgfpathcurveto{\pgfqpoint{3.276003in}{2.740443in}}{\pgfqpoint{3.281589in}{2.742757in}}{\pgfqpoint{3.285707in}{2.746875in}}%
\pgfpathcurveto{\pgfqpoint{3.289826in}{2.750993in}}{\pgfqpoint{3.292139in}{2.756580in}}{\pgfqpoint{3.292139in}{2.762404in}}%
\pgfpathcurveto{\pgfqpoint{3.292139in}{2.768227in}}{\pgfqpoint{3.289826in}{2.773814in}}{\pgfqpoint{3.285707in}{2.777932in}}%
\pgfpathcurveto{\pgfqpoint{3.281589in}{2.782050in}}{\pgfqpoint{3.276003in}{2.784364in}}{\pgfqpoint{3.270179in}{2.784364in}}%
\pgfpathcurveto{\pgfqpoint{3.264355in}{2.784364in}}{\pgfqpoint{3.258769in}{2.782050in}}{\pgfqpoint{3.254651in}{2.777932in}}%
\pgfpathcurveto{\pgfqpoint{3.250533in}{2.773814in}}{\pgfqpoint{3.248219in}{2.768227in}}{\pgfqpoint{3.248219in}{2.762404in}}%
\pgfpathcurveto{\pgfqpoint{3.248219in}{2.756580in}}{\pgfqpoint{3.250533in}{2.750993in}}{\pgfqpoint{3.254651in}{2.746875in}}%
\pgfpathcurveto{\pgfqpoint{3.258769in}{2.742757in}}{\pgfqpoint{3.264355in}{2.740443in}}{\pgfqpoint{3.270179in}{2.740443in}}%
\pgfpathlineto{\pgfqpoint{3.270179in}{2.740443in}}%
\pgfpathclose%
\pgfusepath{stroke,fill}%
\end{pgfscope}%
\begin{pgfscope}%
\pgfpathrectangle{\pgfqpoint{1.414736in}{0.548769in}}{\pgfqpoint{3.497747in}{7.253186in}}%
\pgfusepath{clip}%
\pgfsetbuttcap%
\pgfsetroundjoin%
\definecolor{currentfill}{rgb}{0.525776,0.833491,0.288127}%
\pgfsetfillcolor{currentfill}%
\pgfsetlinewidth{1.003750pt}%
\definecolor{currentstroke}{rgb}{0.525776,0.833491,0.288127}%
\pgfsetstrokecolor{currentstroke}%
\pgfsetdash{}{0pt}%
\pgfpathmoveto{\pgfqpoint{3.702247in}{3.211429in}}%
\pgfpathcurveto{\pgfqpoint{3.708071in}{3.211429in}}{\pgfqpoint{3.713657in}{3.213743in}}{\pgfqpoint{3.717775in}{3.217861in}}%
\pgfpathcurveto{\pgfqpoint{3.721893in}{3.221980in}}{\pgfqpoint{3.724207in}{3.227566in}}{\pgfqpoint{3.724207in}{3.233390in}}%
\pgfpathcurveto{\pgfqpoint{3.724207in}{3.239214in}}{\pgfqpoint{3.721893in}{3.244800in}}{\pgfqpoint{3.717775in}{3.248918in}}%
\pgfpathcurveto{\pgfqpoint{3.713657in}{3.253036in}}{\pgfqpoint{3.708071in}{3.255350in}}{\pgfqpoint{3.702247in}{3.255350in}}%
\pgfpathcurveto{\pgfqpoint{3.696423in}{3.255350in}}{\pgfqpoint{3.690837in}{3.253036in}}{\pgfqpoint{3.686718in}{3.248918in}}%
\pgfpathcurveto{\pgfqpoint{3.682600in}{3.244800in}}{\pgfqpoint{3.680286in}{3.239214in}}{\pgfqpoint{3.680286in}{3.233390in}}%
\pgfpathcurveto{\pgfqpoint{3.680286in}{3.227566in}}{\pgfqpoint{3.682600in}{3.221980in}}{\pgfqpoint{3.686718in}{3.217861in}}%
\pgfpathcurveto{\pgfqpoint{3.690837in}{3.213743in}}{\pgfqpoint{3.696423in}{3.211429in}}{\pgfqpoint{3.702247in}{3.211429in}}%
\pgfpathlineto{\pgfqpoint{3.702247in}{3.211429in}}%
\pgfpathclose%
\pgfusepath{stroke,fill}%
\end{pgfscope}%
\begin{pgfscope}%
\pgfpathrectangle{\pgfqpoint{1.414736in}{0.548769in}}{\pgfqpoint{3.497747in}{7.253186in}}%
\pgfusepath{clip}%
\pgfsetbuttcap%
\pgfsetroundjoin%
\definecolor{currentfill}{rgb}{0.993248,0.906157,0.143936}%
\pgfsetfillcolor{currentfill}%
\pgfsetlinewidth{1.003750pt}%
\definecolor{currentstroke}{rgb}{0.993248,0.906157,0.143936}%
\pgfsetstrokecolor{currentstroke}%
\pgfsetdash{}{0pt}%
\pgfpathmoveto{\pgfqpoint{3.365412in}{3.682415in}}%
\pgfpathcurveto{\pgfqpoint{3.371236in}{3.682415in}}{\pgfqpoint{3.376822in}{3.684729in}}{\pgfqpoint{3.380941in}{3.688847in}}%
\pgfpathcurveto{\pgfqpoint{3.385059in}{3.692966in}}{\pgfqpoint{3.387373in}{3.698552in}}{\pgfqpoint{3.387373in}{3.704376in}}%
\pgfpathcurveto{\pgfqpoint{3.387373in}{3.710200in}}{\pgfqpoint{3.385059in}{3.715786in}}{\pgfqpoint{3.380941in}{3.719904in}}%
\pgfpathcurveto{\pgfqpoint{3.376822in}{3.724022in}}{\pgfqpoint{3.371236in}{3.726336in}}{\pgfqpoint{3.365412in}{3.726336in}}%
\pgfpathcurveto{\pgfqpoint{3.359588in}{3.726336in}}{\pgfqpoint{3.354002in}{3.724022in}}{\pgfqpoint{3.349884in}{3.719904in}}%
\pgfpathcurveto{\pgfqpoint{3.345766in}{3.715786in}}{\pgfqpoint{3.343452in}{3.710200in}}{\pgfqpoint{3.343452in}{3.704376in}}%
\pgfpathcurveto{\pgfqpoint{3.343452in}{3.698552in}}{\pgfqpoint{3.345766in}{3.692966in}}{\pgfqpoint{3.349884in}{3.688847in}}%
\pgfpathcurveto{\pgfqpoint{3.354002in}{3.684729in}}{\pgfqpoint{3.359588in}{3.682415in}}{\pgfqpoint{3.365412in}{3.682415in}}%
\pgfpathlineto{\pgfqpoint{3.365412in}{3.682415in}}%
\pgfpathclose%
\pgfusepath{stroke,fill}%
\end{pgfscope}%
\begin{pgfscope}%
\pgfpathrectangle{\pgfqpoint{1.414736in}{0.548769in}}{\pgfqpoint{3.497747in}{7.253186in}}%
\pgfusepath{clip}%
\pgfsetbuttcap%
\pgfsetroundjoin%
\definecolor{currentfill}{rgb}{0.709898,0.868751,0.169257}%
\pgfsetfillcolor{currentfill}%
\pgfsetlinewidth{1.003750pt}%
\definecolor{currentstroke}{rgb}{0.709898,0.868751,0.169257}%
\pgfsetstrokecolor{currentstroke}%
\pgfsetdash{}{0pt}%
\pgfpathmoveto{\pgfqpoint{3.659460in}{4.153402in}}%
\pgfpathcurveto{\pgfqpoint{3.665283in}{4.153402in}}{\pgfqpoint{3.670870in}{4.155715in}}{\pgfqpoint{3.674988in}{4.159834in}}%
\pgfpathcurveto{\pgfqpoint{3.679106in}{4.163952in}}{\pgfqpoint{3.681420in}{4.169538in}}{\pgfqpoint{3.681420in}{4.175362in}}%
\pgfpathcurveto{\pgfqpoint{3.681420in}{4.181186in}}{\pgfqpoint{3.679106in}{4.186772in}}{\pgfqpoint{3.674988in}{4.190890in}}%
\pgfpathcurveto{\pgfqpoint{3.670870in}{4.195008in}}{\pgfqpoint{3.665283in}{4.197322in}}{\pgfqpoint{3.659460in}{4.197322in}}%
\pgfpathcurveto{\pgfqpoint{3.653636in}{4.197322in}}{\pgfqpoint{3.648049in}{4.195008in}}{\pgfqpoint{3.643931in}{4.190890in}}%
\pgfpathcurveto{\pgfqpoint{3.639813in}{4.186772in}}{\pgfqpoint{3.637499in}{4.181186in}}{\pgfqpoint{3.637499in}{4.175362in}}%
\pgfpathcurveto{\pgfqpoint{3.637499in}{4.169538in}}{\pgfqpoint{3.639813in}{4.163952in}}{\pgfqpoint{3.643931in}{4.159834in}}%
\pgfpathcurveto{\pgfqpoint{3.648049in}{4.155715in}}{\pgfqpoint{3.653636in}{4.153402in}}{\pgfqpoint{3.659460in}{4.153402in}}%
\pgfpathlineto{\pgfqpoint{3.659460in}{4.153402in}}%
\pgfpathclose%
\pgfusepath{stroke,fill}%
\end{pgfscope}%
\begin{pgfscope}%
\pgfpathrectangle{\pgfqpoint{1.414736in}{0.548769in}}{\pgfqpoint{3.497747in}{7.253186in}}%
\pgfusepath{clip}%
\pgfsetbuttcap%
\pgfsetroundjoin%
\definecolor{currentfill}{rgb}{0.804182,0.882046,0.114965}%
\pgfsetfillcolor{currentfill}%
\pgfsetlinewidth{1.003750pt}%
\definecolor{currentstroke}{rgb}{0.804182,0.882046,0.114965}%
\pgfsetstrokecolor{currentstroke}%
\pgfsetdash{}{0pt}%
\pgfpathmoveto{\pgfqpoint{3.574103in}{4.624388in}}%
\pgfpathcurveto{\pgfqpoint{3.579927in}{4.624388in}}{\pgfqpoint{3.585513in}{4.626702in}}{\pgfqpoint{3.589631in}{4.630820in}}%
\pgfpathcurveto{\pgfqpoint{3.593749in}{4.634938in}}{\pgfqpoint{3.596063in}{4.640524in}}{\pgfqpoint{3.596063in}{4.646348in}}%
\pgfpathcurveto{\pgfqpoint{3.596063in}{4.652172in}}{\pgfqpoint{3.593749in}{4.657758in}}{\pgfqpoint{3.589631in}{4.661876in}}%
\pgfpathcurveto{\pgfqpoint{3.585513in}{4.665994in}}{\pgfqpoint{3.579927in}{4.668308in}}{\pgfqpoint{3.574103in}{4.668308in}}%
\pgfpathcurveto{\pgfqpoint{3.568279in}{4.668308in}}{\pgfqpoint{3.562693in}{4.665994in}}{\pgfqpoint{3.558575in}{4.661876in}}%
\pgfpathcurveto{\pgfqpoint{3.554457in}{4.657758in}}{\pgfqpoint{3.552143in}{4.652172in}}{\pgfqpoint{3.552143in}{4.646348in}}%
\pgfpathcurveto{\pgfqpoint{3.552143in}{4.640524in}}{\pgfqpoint{3.554457in}{4.634938in}}{\pgfqpoint{3.558575in}{4.630820in}}%
\pgfpathcurveto{\pgfqpoint{3.562693in}{4.626702in}}{\pgfqpoint{3.568279in}{4.624388in}}{\pgfqpoint{3.574103in}{4.624388in}}%
\pgfpathlineto{\pgfqpoint{3.574103in}{4.624388in}}%
\pgfpathclose%
\pgfusepath{stroke,fill}%
\end{pgfscope}%
\begin{pgfscope}%
\pgfpathrectangle{\pgfqpoint{1.414736in}{0.548769in}}{\pgfqpoint{3.497747in}{7.253186in}}%
\pgfusepath{clip}%
\pgfsetbuttcap%
\pgfsetroundjoin%
\definecolor{currentfill}{rgb}{0.125394,0.574318,0.549086}%
\pgfsetfillcolor{currentfill}%
\pgfsetlinewidth{1.003750pt}%
\definecolor{currentstroke}{rgb}{0.125394,0.574318,0.549086}%
\pgfsetstrokecolor{currentstroke}%
\pgfsetdash{}{0pt}%
\pgfpathmoveto{\pgfqpoint{3.867795in}{5.095374in}}%
\pgfpathcurveto{\pgfqpoint{3.873619in}{5.095374in}}{\pgfqpoint{3.879205in}{5.097688in}}{\pgfqpoint{3.883323in}{5.101806in}}%
\pgfpathcurveto{\pgfqpoint{3.887441in}{5.105924in}}{\pgfqpoint{3.889755in}{5.111510in}}{\pgfqpoint{3.889755in}{5.117334in}}%
\pgfpathcurveto{\pgfqpoint{3.889755in}{5.123158in}}{\pgfqpoint{3.887441in}{5.128744in}}{\pgfqpoint{3.883323in}{5.132862in}}%
\pgfpathcurveto{\pgfqpoint{3.879205in}{5.136980in}}{\pgfqpoint{3.873619in}{5.139294in}}{\pgfqpoint{3.867795in}{5.139294in}}%
\pgfpathcurveto{\pgfqpoint{3.861971in}{5.139294in}}{\pgfqpoint{3.856385in}{5.136980in}}{\pgfqpoint{3.852267in}{5.132862in}}%
\pgfpathcurveto{\pgfqpoint{3.848148in}{5.128744in}}{\pgfqpoint{3.845835in}{5.123158in}}{\pgfqpoint{3.845835in}{5.117334in}}%
\pgfpathcurveto{\pgfqpoint{3.845835in}{5.111510in}}{\pgfqpoint{3.848148in}{5.105924in}}{\pgfqpoint{3.852267in}{5.101806in}}%
\pgfpathcurveto{\pgfqpoint{3.856385in}{5.097688in}}{\pgfqpoint{3.861971in}{5.095374in}}{\pgfqpoint{3.867795in}{5.095374in}}%
\pgfpathlineto{\pgfqpoint{3.867795in}{5.095374in}}%
\pgfpathclose%
\pgfusepath{stroke,fill}%
\end{pgfscope}%
\begin{pgfscope}%
\pgfpathrectangle{\pgfqpoint{1.414736in}{0.548769in}}{\pgfqpoint{3.497747in}{7.253186in}}%
\pgfusepath{clip}%
\pgfsetbuttcap%
\pgfsetroundjoin%
\definecolor{currentfill}{rgb}{0.163625,0.471133,0.558148}%
\pgfsetfillcolor{currentfill}%
\pgfsetlinewidth{1.003750pt}%
\definecolor{currentstroke}{rgb}{0.163625,0.471133,0.558148}%
\pgfsetstrokecolor{currentstroke}%
\pgfsetdash{}{0pt}%
\pgfpathmoveto{\pgfqpoint{3.789875in}{5.566360in}}%
\pgfpathcurveto{\pgfqpoint{3.795699in}{5.566360in}}{\pgfqpoint{3.801285in}{5.568674in}}{\pgfqpoint{3.805403in}{5.572792in}}%
\pgfpathcurveto{\pgfqpoint{3.809521in}{5.576910in}}{\pgfqpoint{3.811835in}{5.582496in}}{\pgfqpoint{3.811835in}{5.588320in}}%
\pgfpathcurveto{\pgfqpoint{3.811835in}{5.594144in}}{\pgfqpoint{3.809521in}{5.599730in}}{\pgfqpoint{3.805403in}{5.603848in}}%
\pgfpathcurveto{\pgfqpoint{3.801285in}{5.607967in}}{\pgfqpoint{3.795699in}{5.610280in}}{\pgfqpoint{3.789875in}{5.610280in}}%
\pgfpathcurveto{\pgfqpoint{3.784051in}{5.610280in}}{\pgfqpoint{3.778465in}{5.607967in}}{\pgfqpoint{3.774347in}{5.603848in}}%
\pgfpathcurveto{\pgfqpoint{3.770229in}{5.599730in}}{\pgfqpoint{3.767915in}{5.594144in}}{\pgfqpoint{3.767915in}{5.588320in}}%
\pgfpathcurveto{\pgfqpoint{3.767915in}{5.582496in}}{\pgfqpoint{3.770229in}{5.576910in}}{\pgfqpoint{3.774347in}{5.572792in}}%
\pgfpathcurveto{\pgfqpoint{3.778465in}{5.568674in}}{\pgfqpoint{3.784051in}{5.566360in}}{\pgfqpoint{3.789875in}{5.566360in}}%
\pgfpathlineto{\pgfqpoint{3.789875in}{5.566360in}}%
\pgfpathclose%
\pgfusepath{stroke,fill}%
\end{pgfscope}%
\begin{pgfscope}%
\pgfpathrectangle{\pgfqpoint{1.414736in}{0.548769in}}{\pgfqpoint{3.497747in}{7.253186in}}%
\pgfusepath{clip}%
\pgfsetbuttcap%
\pgfsetroundjoin%
\definecolor{currentfill}{rgb}{0.535621,0.835785,0.281908}%
\pgfsetfillcolor{currentfill}%
\pgfsetlinewidth{1.003750pt}%
\definecolor{currentstroke}{rgb}{0.535621,0.835785,0.281908}%
\pgfsetstrokecolor{currentstroke}%
\pgfsetdash{}{0pt}%
\pgfpathmoveto{\pgfqpoint{3.329225in}{6.037346in}}%
\pgfpathcurveto{\pgfqpoint{3.335048in}{6.037346in}}{\pgfqpoint{3.340635in}{6.039660in}}{\pgfqpoint{3.344753in}{6.043778in}}%
\pgfpathcurveto{\pgfqpoint{3.348871in}{6.047896in}}{\pgfqpoint{3.351185in}{6.053482in}}{\pgfqpoint{3.351185in}{6.059306in}}%
\pgfpathcurveto{\pgfqpoint{3.351185in}{6.065130in}}{\pgfqpoint{3.348871in}{6.070716in}}{\pgfqpoint{3.344753in}{6.074835in}}%
\pgfpathcurveto{\pgfqpoint{3.340635in}{6.078953in}}{\pgfqpoint{3.335048in}{6.081267in}}{\pgfqpoint{3.329225in}{6.081267in}}%
\pgfpathcurveto{\pgfqpoint{3.323401in}{6.081267in}}{\pgfqpoint{3.317814in}{6.078953in}}{\pgfqpoint{3.313696in}{6.074835in}}%
\pgfpathcurveto{\pgfqpoint{3.309578in}{6.070716in}}{\pgfqpoint{3.307264in}{6.065130in}}{\pgfqpoint{3.307264in}{6.059306in}}%
\pgfpathcurveto{\pgfqpoint{3.307264in}{6.053482in}}{\pgfqpoint{3.309578in}{6.047896in}}{\pgfqpoint{3.313696in}{6.043778in}}%
\pgfpathcurveto{\pgfqpoint{3.317814in}{6.039660in}}{\pgfqpoint{3.323401in}{6.037346in}}{\pgfqpoint{3.329225in}{6.037346in}}%
\pgfpathlineto{\pgfqpoint{3.329225in}{6.037346in}}%
\pgfpathclose%
\pgfusepath{stroke,fill}%
\end{pgfscope}%
\begin{pgfscope}%
\pgfpathrectangle{\pgfqpoint{1.414736in}{0.548769in}}{\pgfqpoint{3.497747in}{7.253186in}}%
\pgfusepath{clip}%
\pgfsetbuttcap%
\pgfsetroundjoin%
\definecolor{currentfill}{rgb}{0.783315,0.879285,0.125405}%
\pgfsetfillcolor{currentfill}%
\pgfsetlinewidth{1.003750pt}%
\definecolor{currentstroke}{rgb}{0.783315,0.879285,0.125405}%
\pgfsetstrokecolor{currentstroke}%
\pgfsetdash{}{0pt}%
\pgfpathmoveto{\pgfqpoint{3.723479in}{6.508332in}}%
\pgfpathcurveto{\pgfqpoint{3.729303in}{6.508332in}}{\pgfqpoint{3.734889in}{6.510646in}}{\pgfqpoint{3.739007in}{6.514764in}}%
\pgfpathcurveto{\pgfqpoint{3.743125in}{6.518882in}}{\pgfqpoint{3.745439in}{6.524468in}}{\pgfqpoint{3.745439in}{6.530292in}}%
\pgfpathcurveto{\pgfqpoint{3.745439in}{6.536116in}}{\pgfqpoint{3.743125in}{6.541702in}}{\pgfqpoint{3.739007in}{6.545821in}}%
\pgfpathcurveto{\pgfqpoint{3.734889in}{6.549939in}}{\pgfqpoint{3.729303in}{6.552253in}}{\pgfqpoint{3.723479in}{6.552253in}}%
\pgfpathcurveto{\pgfqpoint{3.717655in}{6.552253in}}{\pgfqpoint{3.712069in}{6.549939in}}{\pgfqpoint{3.707951in}{6.545821in}}%
\pgfpathcurveto{\pgfqpoint{3.703833in}{6.541702in}}{\pgfqpoint{3.701519in}{6.536116in}}{\pgfqpoint{3.701519in}{6.530292in}}%
\pgfpathcurveto{\pgfqpoint{3.701519in}{6.524468in}}{\pgfqpoint{3.703833in}{6.518882in}}{\pgfqpoint{3.707951in}{6.514764in}}%
\pgfpathcurveto{\pgfqpoint{3.712069in}{6.510646in}}{\pgfqpoint{3.717655in}{6.508332in}}{\pgfqpoint{3.723479in}{6.508332in}}%
\pgfpathlineto{\pgfqpoint{3.723479in}{6.508332in}}%
\pgfpathclose%
\pgfusepath{stroke,fill}%
\end{pgfscope}%
\begin{pgfscope}%
\pgfpathrectangle{\pgfqpoint{1.414736in}{0.548769in}}{\pgfqpoint{3.497747in}{7.253186in}}%
\pgfusepath{clip}%
\pgfsetbuttcap%
\pgfsetroundjoin%
\definecolor{currentfill}{rgb}{0.626579,0.854645,0.223353}%
\pgfsetfillcolor{currentfill}%
\pgfsetlinewidth{1.003750pt}%
\definecolor{currentstroke}{rgb}{0.626579,0.854645,0.223353}%
\pgfsetstrokecolor{currentstroke}%
\pgfsetdash{}{0pt}%
\pgfpathmoveto{\pgfqpoint{4.457437in}{6.979318in}}%
\pgfpathcurveto{\pgfqpoint{4.463261in}{6.979318in}}{\pgfqpoint{4.468847in}{6.981632in}}{\pgfqpoint{4.472965in}{6.985750in}}%
\pgfpathcurveto{\pgfqpoint{4.477083in}{6.989868in}}{\pgfqpoint{4.479397in}{6.995455in}}{\pgfqpoint{4.479397in}{7.001278in}}%
\pgfpathcurveto{\pgfqpoint{4.479397in}{7.007102in}}{\pgfqpoint{4.477083in}{7.012689in}}{\pgfqpoint{4.472965in}{7.016807in}}%
\pgfpathcurveto{\pgfqpoint{4.468847in}{7.020925in}}{\pgfqpoint{4.463261in}{7.023239in}}{\pgfqpoint{4.457437in}{7.023239in}}%
\pgfpathcurveto{\pgfqpoint{4.451613in}{7.023239in}}{\pgfqpoint{4.446027in}{7.020925in}}{\pgfqpoint{4.441909in}{7.016807in}}%
\pgfpathcurveto{\pgfqpoint{4.437791in}{7.012689in}}{\pgfqpoint{4.435477in}{7.007102in}}{\pgfqpoint{4.435477in}{7.001278in}}%
\pgfpathcurveto{\pgfqpoint{4.435477in}{6.995455in}}{\pgfqpoint{4.437791in}{6.989868in}}{\pgfqpoint{4.441909in}{6.985750in}}%
\pgfpathcurveto{\pgfqpoint{4.446027in}{6.981632in}}{\pgfqpoint{4.451613in}{6.979318in}}{\pgfqpoint{4.457437in}{6.979318in}}%
\pgfpathlineto{\pgfqpoint{4.457437in}{6.979318in}}%
\pgfpathclose%
\pgfusepath{stroke,fill}%
\end{pgfscope}%
\begin{pgfscope}%
\pgfpathrectangle{\pgfqpoint{1.414736in}{0.548769in}}{\pgfqpoint{3.497747in}{7.253186in}}%
\pgfusepath{clip}%
\pgfsetbuttcap%
\pgfsetroundjoin%
\definecolor{currentfill}{rgb}{0.430983,0.808473,0.346476}%
\pgfsetfillcolor{currentfill}%
\pgfsetlinewidth{1.003750pt}%
\definecolor{currentstroke}{rgb}{0.430983,0.808473,0.346476}%
\pgfsetstrokecolor{currentstroke}%
\pgfsetdash{}{0pt}%
\pgfpathmoveto{\pgfqpoint{3.838915in}{7.450304in}}%
\pgfpathcurveto{\pgfqpoint{3.844739in}{7.450304in}}{\pgfqpoint{3.850325in}{7.452618in}}{\pgfqpoint{3.854443in}{7.456736in}}%
\pgfpathcurveto{\pgfqpoint{3.858562in}{7.460854in}}{\pgfqpoint{3.860876in}{7.466441in}}{\pgfqpoint{3.860876in}{7.472265in}}%
\pgfpathcurveto{\pgfqpoint{3.860876in}{7.478089in}}{\pgfqpoint{3.858562in}{7.483675in}}{\pgfqpoint{3.854443in}{7.487793in}}%
\pgfpathcurveto{\pgfqpoint{3.850325in}{7.491911in}}{\pgfqpoint{3.844739in}{7.494225in}}{\pgfqpoint{3.838915in}{7.494225in}}%
\pgfpathcurveto{\pgfqpoint{3.833091in}{7.494225in}}{\pgfqpoint{3.827505in}{7.491911in}}{\pgfqpoint{3.823387in}{7.487793in}}%
\pgfpathcurveto{\pgfqpoint{3.819269in}{7.483675in}}{\pgfqpoint{3.816955in}{7.478089in}}{\pgfqpoint{3.816955in}{7.472265in}}%
\pgfpathcurveto{\pgfqpoint{3.816955in}{7.466441in}}{\pgfqpoint{3.819269in}{7.460854in}}{\pgfqpoint{3.823387in}{7.456736in}}%
\pgfpathcurveto{\pgfqpoint{3.827505in}{7.452618in}}{\pgfqpoint{3.833091in}{7.450304in}}{\pgfqpoint{3.838915in}{7.450304in}}%
\pgfpathlineto{\pgfqpoint{3.838915in}{7.450304in}}%
\pgfpathclose%
\pgfusepath{stroke,fill}%
\end{pgfscope}%
\begin{pgfscope}%
\pgfpathrectangle{\pgfqpoint{1.414736in}{0.548769in}}{\pgfqpoint{3.497747in}{7.253186in}}%
\pgfusepath{clip}%
\pgfsetbuttcap%
\pgfsetroundjoin%
\definecolor{currentfill}{rgb}{0.866013,0.889868,0.095953}%
\pgfsetfillcolor{currentfill}%
\pgfsetlinewidth{1.003750pt}%
\definecolor{currentstroke}{rgb}{0.866013,0.889868,0.095953}%
\pgfsetstrokecolor{currentstroke}%
\pgfsetdash{}{0pt}%
\pgfpathmoveto{\pgfqpoint{3.403630in}{0.856499in}}%
\pgfpathcurveto{\pgfqpoint{3.409454in}{0.856499in}}{\pgfqpoint{3.415040in}{0.858813in}}{\pgfqpoint{3.419158in}{0.862931in}}%
\pgfpathcurveto{\pgfqpoint{3.423277in}{0.867049in}}{\pgfqpoint{3.425590in}{0.872635in}}{\pgfqpoint{3.425590in}{0.878459in}}%
\pgfpathcurveto{\pgfqpoint{3.425590in}{0.884283in}}{\pgfqpoint{3.423277in}{0.889869in}}{\pgfqpoint{3.419158in}{0.893987in}}%
\pgfpathcurveto{\pgfqpoint{3.415040in}{0.898105in}}{\pgfqpoint{3.409454in}{0.900419in}}{\pgfqpoint{3.403630in}{0.900419in}}%
\pgfpathcurveto{\pgfqpoint{3.397806in}{0.900419in}}{\pgfqpoint{3.392220in}{0.898105in}}{\pgfqpoint{3.388102in}{0.893987in}}%
\pgfpathcurveto{\pgfqpoint{3.383984in}{0.889869in}}{\pgfqpoint{3.381670in}{0.884283in}}{\pgfqpoint{3.381670in}{0.878459in}}%
\pgfpathcurveto{\pgfqpoint{3.381670in}{0.872635in}}{\pgfqpoint{3.383984in}{0.867049in}}{\pgfqpoint{3.388102in}{0.862931in}}%
\pgfpathcurveto{\pgfqpoint{3.392220in}{0.858813in}}{\pgfqpoint{3.397806in}{0.856499in}}{\pgfqpoint{3.403630in}{0.856499in}}%
\pgfpathlineto{\pgfqpoint{3.403630in}{0.856499in}}%
\pgfpathclose%
\pgfusepath{stroke,fill}%
\end{pgfscope}%
\begin{pgfscope}%
\pgfpathrectangle{\pgfqpoint{1.414736in}{0.548769in}}{\pgfqpoint{3.497747in}{7.253186in}}%
\pgfusepath{clip}%
\pgfsetbuttcap%
\pgfsetroundjoin%
\definecolor{currentfill}{rgb}{0.335885,0.777018,0.402049}%
\pgfsetfillcolor{currentfill}%
\pgfsetlinewidth{1.003750pt}%
\definecolor{currentstroke}{rgb}{0.335885,0.777018,0.402049}%
\pgfsetstrokecolor{currentstroke}%
\pgfsetdash{}{0pt}%
\pgfpathmoveto{\pgfqpoint{3.972421in}{1.327485in}}%
\pgfpathcurveto{\pgfqpoint{3.978245in}{1.327485in}}{\pgfqpoint{3.983831in}{1.329799in}}{\pgfqpoint{3.987949in}{1.333917in}}%
\pgfpathcurveto{\pgfqpoint{3.992067in}{1.338035in}}{\pgfqpoint{3.994381in}{1.343621in}}{\pgfqpoint{3.994381in}{1.349445in}}%
\pgfpathcurveto{\pgfqpoint{3.994381in}{1.355269in}}{\pgfqpoint{3.992067in}{1.360855in}}{\pgfqpoint{3.987949in}{1.364973in}}%
\pgfpathcurveto{\pgfqpoint{3.983831in}{1.369092in}}{\pgfqpoint{3.978245in}{1.371405in}}{\pgfqpoint{3.972421in}{1.371405in}}%
\pgfpathcurveto{\pgfqpoint{3.966597in}{1.371405in}}{\pgfqpoint{3.961011in}{1.369092in}}{\pgfqpoint{3.956893in}{1.364973in}}%
\pgfpathcurveto{\pgfqpoint{3.952774in}{1.360855in}}{\pgfqpoint{3.950460in}{1.355269in}}{\pgfqpoint{3.950460in}{1.349445in}}%
\pgfpathcurveto{\pgfqpoint{3.950460in}{1.343621in}}{\pgfqpoint{3.952774in}{1.338035in}}{\pgfqpoint{3.956893in}{1.333917in}}%
\pgfpathcurveto{\pgfqpoint{3.961011in}{1.329799in}}{\pgfqpoint{3.966597in}{1.327485in}}{\pgfqpoint{3.972421in}{1.327485in}}%
\pgfpathlineto{\pgfqpoint{3.972421in}{1.327485in}}%
\pgfpathclose%
\pgfusepath{stroke,fill}%
\end{pgfscope}%
\begin{pgfscope}%
\pgfpathrectangle{\pgfqpoint{1.414736in}{0.548769in}}{\pgfqpoint{3.497747in}{7.253186in}}%
\pgfusepath{clip}%
\pgfsetbuttcap%
\pgfsetroundjoin%
\definecolor{currentfill}{rgb}{0.412913,0.803041,0.357269}%
\pgfsetfillcolor{currentfill}%
\pgfsetlinewidth{1.003750pt}%
\definecolor{currentstroke}{rgb}{0.412913,0.803041,0.357269}%
\pgfsetstrokecolor{currentstroke}%
\pgfsetdash{}{0pt}%
\pgfpathmoveto{\pgfqpoint{2.897513in}{1.798471in}}%
\pgfpathcurveto{\pgfqpoint{2.903337in}{1.798471in}}{\pgfqpoint{2.908924in}{1.800785in}}{\pgfqpoint{2.913042in}{1.804903in}}%
\pgfpathcurveto{\pgfqpoint{2.917160in}{1.809021in}}{\pgfqpoint{2.919474in}{1.814607in}}{\pgfqpoint{2.919474in}{1.820431in}}%
\pgfpathcurveto{\pgfqpoint{2.919474in}{1.826255in}}{\pgfqpoint{2.917160in}{1.831841in}}{\pgfqpoint{2.913042in}{1.835960in}}%
\pgfpathcurveto{\pgfqpoint{2.908924in}{1.840078in}}{\pgfqpoint{2.903337in}{1.842392in}}{\pgfqpoint{2.897513in}{1.842392in}}%
\pgfpathcurveto{\pgfqpoint{2.891690in}{1.842392in}}{\pgfqpoint{2.886103in}{1.840078in}}{\pgfqpoint{2.881985in}{1.835960in}}%
\pgfpathcurveto{\pgfqpoint{2.877867in}{1.831841in}}{\pgfqpoint{2.875553in}{1.826255in}}{\pgfqpoint{2.875553in}{1.820431in}}%
\pgfpathcurveto{\pgfqpoint{2.875553in}{1.814607in}}{\pgfqpoint{2.877867in}{1.809021in}}{\pgfqpoint{2.881985in}{1.804903in}}%
\pgfpathcurveto{\pgfqpoint{2.886103in}{1.800785in}}{\pgfqpoint{2.891690in}{1.798471in}}{\pgfqpoint{2.897513in}{1.798471in}}%
\pgfpathlineto{\pgfqpoint{2.897513in}{1.798471in}}%
\pgfpathclose%
\pgfusepath{stroke,fill}%
\end{pgfscope}%
\begin{pgfscope}%
\pgfpathrectangle{\pgfqpoint{1.414736in}{0.548769in}}{\pgfqpoint{3.497747in}{7.253186in}}%
\pgfusepath{clip}%
\pgfsetbuttcap%
\pgfsetroundjoin%
\definecolor{currentfill}{rgb}{0.866013,0.889868,0.095953}%
\pgfsetfillcolor{currentfill}%
\pgfsetlinewidth{1.003750pt}%
\definecolor{currentstroke}{rgb}{0.866013,0.889868,0.095953}%
\pgfsetstrokecolor{currentstroke}%
\pgfsetdash{}{0pt}%
\pgfpathmoveto{\pgfqpoint{3.850423in}{2.269457in}}%
\pgfpathcurveto{\pgfqpoint{3.856247in}{2.269457in}}{\pgfqpoint{3.861833in}{2.271771in}}{\pgfqpoint{3.865951in}{2.275889in}}%
\pgfpathcurveto{\pgfqpoint{3.870069in}{2.280007in}}{\pgfqpoint{3.872383in}{2.285593in}}{\pgfqpoint{3.872383in}{2.291417in}}%
\pgfpathcurveto{\pgfqpoint{3.872383in}{2.297241in}}{\pgfqpoint{3.870069in}{2.302828in}}{\pgfqpoint{3.865951in}{2.306946in}}%
\pgfpathcurveto{\pgfqpoint{3.861833in}{2.311064in}}{\pgfqpoint{3.856247in}{2.313378in}}{\pgfqpoint{3.850423in}{2.313378in}}%
\pgfpathcurveto{\pgfqpoint{3.844599in}{2.313378in}}{\pgfqpoint{3.839013in}{2.311064in}}{\pgfqpoint{3.834895in}{2.306946in}}%
\pgfpathcurveto{\pgfqpoint{3.830777in}{2.302828in}}{\pgfqpoint{3.828463in}{2.297241in}}{\pgfqpoint{3.828463in}{2.291417in}}%
\pgfpathcurveto{\pgfqpoint{3.828463in}{2.285593in}}{\pgfqpoint{3.830777in}{2.280007in}}{\pgfqpoint{3.834895in}{2.275889in}}%
\pgfpathcurveto{\pgfqpoint{3.839013in}{2.271771in}}{\pgfqpoint{3.844599in}{2.269457in}}{\pgfqpoint{3.850423in}{2.269457in}}%
\pgfpathlineto{\pgfqpoint{3.850423in}{2.269457in}}%
\pgfpathclose%
\pgfusepath{stroke,fill}%
\end{pgfscope}%
\begin{pgfscope}%
\pgfpathrectangle{\pgfqpoint{1.414736in}{0.548769in}}{\pgfqpoint{3.497747in}{7.253186in}}%
\pgfusepath{clip}%
\pgfsetbuttcap%
\pgfsetroundjoin%
\definecolor{currentfill}{rgb}{0.647257,0.858400,0.209861}%
\pgfsetfillcolor{currentfill}%
\pgfsetlinewidth{1.003750pt}%
\definecolor{currentstroke}{rgb}{0.647257,0.858400,0.209861}%
\pgfsetstrokecolor{currentstroke}%
\pgfsetdash{}{0pt}%
\pgfpathmoveto{\pgfqpoint{3.481565in}{2.740443in}}%
\pgfpathcurveto{\pgfqpoint{3.487389in}{2.740443in}}{\pgfqpoint{3.492975in}{2.742757in}}{\pgfqpoint{3.497093in}{2.746875in}}%
\pgfpathcurveto{\pgfqpoint{3.501211in}{2.750993in}}{\pgfqpoint{3.503525in}{2.756580in}}{\pgfqpoint{3.503525in}{2.762404in}}%
\pgfpathcurveto{\pgfqpoint{3.503525in}{2.768227in}}{\pgfqpoint{3.501211in}{2.773814in}}{\pgfqpoint{3.497093in}{2.777932in}}%
\pgfpathcurveto{\pgfqpoint{3.492975in}{2.782050in}}{\pgfqpoint{3.487389in}{2.784364in}}{\pgfqpoint{3.481565in}{2.784364in}}%
\pgfpathcurveto{\pgfqpoint{3.475741in}{2.784364in}}{\pgfqpoint{3.470155in}{2.782050in}}{\pgfqpoint{3.466037in}{2.777932in}}%
\pgfpathcurveto{\pgfqpoint{3.461918in}{2.773814in}}{\pgfqpoint{3.459605in}{2.768227in}}{\pgfqpoint{3.459605in}{2.762404in}}%
\pgfpathcurveto{\pgfqpoint{3.459605in}{2.756580in}}{\pgfqpoint{3.461918in}{2.750993in}}{\pgfqpoint{3.466037in}{2.746875in}}%
\pgfpathcurveto{\pgfqpoint{3.470155in}{2.742757in}}{\pgfqpoint{3.475741in}{2.740443in}}{\pgfqpoint{3.481565in}{2.740443in}}%
\pgfpathlineto{\pgfqpoint{3.481565in}{2.740443in}}%
\pgfpathclose%
\pgfusepath{stroke,fill}%
\end{pgfscope}%
\begin{pgfscope}%
\pgfpathrectangle{\pgfqpoint{1.414736in}{0.548769in}}{\pgfqpoint{3.497747in}{7.253186in}}%
\pgfusepath{clip}%
\pgfsetbuttcap%
\pgfsetroundjoin%
\definecolor{currentfill}{rgb}{0.751884,0.874951,0.143228}%
\pgfsetfillcolor{currentfill}%
\pgfsetlinewidth{1.003750pt}%
\definecolor{currentstroke}{rgb}{0.751884,0.874951,0.143228}%
\pgfsetstrokecolor{currentstroke}%
\pgfsetdash{}{0pt}%
\pgfpathmoveto{\pgfqpoint{3.768312in}{3.211429in}}%
\pgfpathcurveto{\pgfqpoint{3.774136in}{3.211429in}}{\pgfqpoint{3.779722in}{3.213743in}}{\pgfqpoint{3.783840in}{3.217861in}}%
\pgfpathcurveto{\pgfqpoint{3.787959in}{3.221980in}}{\pgfqpoint{3.790272in}{3.227566in}}{\pgfqpoint{3.790272in}{3.233390in}}%
\pgfpathcurveto{\pgfqpoint{3.790272in}{3.239214in}}{\pgfqpoint{3.787959in}{3.244800in}}{\pgfqpoint{3.783840in}{3.248918in}}%
\pgfpathcurveto{\pgfqpoint{3.779722in}{3.253036in}}{\pgfqpoint{3.774136in}{3.255350in}}{\pgfqpoint{3.768312in}{3.255350in}}%
\pgfpathcurveto{\pgfqpoint{3.762488in}{3.255350in}}{\pgfqpoint{3.756902in}{3.253036in}}{\pgfqpoint{3.752784in}{3.248918in}}%
\pgfpathcurveto{\pgfqpoint{3.748666in}{3.244800in}}{\pgfqpoint{3.746352in}{3.239214in}}{\pgfqpoint{3.746352in}{3.233390in}}%
\pgfpathcurveto{\pgfqpoint{3.746352in}{3.227566in}}{\pgfqpoint{3.748666in}{3.221980in}}{\pgfqpoint{3.752784in}{3.217861in}}%
\pgfpathcurveto{\pgfqpoint{3.756902in}{3.213743in}}{\pgfqpoint{3.762488in}{3.211429in}}{\pgfqpoint{3.768312in}{3.211429in}}%
\pgfpathlineto{\pgfqpoint{3.768312in}{3.211429in}}%
\pgfpathclose%
\pgfusepath{stroke,fill}%
\end{pgfscope}%
\begin{pgfscope}%
\pgfpathrectangle{\pgfqpoint{1.414736in}{0.548769in}}{\pgfqpoint{3.497747in}{7.253186in}}%
\pgfusepath{clip}%
\pgfsetbuttcap%
\pgfsetroundjoin%
\definecolor{currentfill}{rgb}{0.804182,0.882046,0.114965}%
\pgfsetfillcolor{currentfill}%
\pgfsetlinewidth{1.003750pt}%
\definecolor{currentstroke}{rgb}{0.804182,0.882046,0.114965}%
\pgfsetstrokecolor{currentstroke}%
\pgfsetdash{}{0pt}%
\pgfpathmoveto{\pgfqpoint{3.416299in}{3.682415in}}%
\pgfpathcurveto{\pgfqpoint{3.422123in}{3.682415in}}{\pgfqpoint{3.427710in}{3.684729in}}{\pgfqpoint{3.431828in}{3.688847in}}%
\pgfpathcurveto{\pgfqpoint{3.435946in}{3.692966in}}{\pgfqpoint{3.438260in}{3.698552in}}{\pgfqpoint{3.438260in}{3.704376in}}%
\pgfpathcurveto{\pgfqpoint{3.438260in}{3.710200in}}{\pgfqpoint{3.435946in}{3.715786in}}{\pgfqpoint{3.431828in}{3.719904in}}%
\pgfpathcurveto{\pgfqpoint{3.427710in}{3.724022in}}{\pgfqpoint{3.422123in}{3.726336in}}{\pgfqpoint{3.416299in}{3.726336in}}%
\pgfpathcurveto{\pgfqpoint{3.410476in}{3.726336in}}{\pgfqpoint{3.404889in}{3.724022in}}{\pgfqpoint{3.400771in}{3.719904in}}%
\pgfpathcurveto{\pgfqpoint{3.396653in}{3.715786in}}{\pgfqpoint{3.394339in}{3.710200in}}{\pgfqpoint{3.394339in}{3.704376in}}%
\pgfpathcurveto{\pgfqpoint{3.394339in}{3.698552in}}{\pgfqpoint{3.396653in}{3.692966in}}{\pgfqpoint{3.400771in}{3.688847in}}%
\pgfpathcurveto{\pgfqpoint{3.404889in}{3.684729in}}{\pgfqpoint{3.410476in}{3.682415in}}{\pgfqpoint{3.416299in}{3.682415in}}%
\pgfpathlineto{\pgfqpoint{3.416299in}{3.682415in}}%
\pgfpathclose%
\pgfusepath{stroke,fill}%
\end{pgfscope}%
\begin{pgfscope}%
\pgfpathrectangle{\pgfqpoint{1.414736in}{0.548769in}}{\pgfqpoint{3.497747in}{7.253186in}}%
\pgfusepath{clip}%
\pgfsetbuttcap%
\pgfsetroundjoin%
\definecolor{currentfill}{rgb}{0.993248,0.906157,0.143936}%
\pgfsetfillcolor{currentfill}%
\pgfsetlinewidth{1.003750pt}%
\definecolor{currentstroke}{rgb}{0.993248,0.906157,0.143936}%
\pgfsetstrokecolor{currentstroke}%
\pgfsetdash{}{0pt}%
\pgfpathmoveto{\pgfqpoint{3.428498in}{4.153402in}}%
\pgfpathcurveto{\pgfqpoint{3.434322in}{4.153402in}}{\pgfqpoint{3.439908in}{4.155715in}}{\pgfqpoint{3.444026in}{4.159834in}}%
\pgfpathcurveto{\pgfqpoint{3.448144in}{4.163952in}}{\pgfqpoint{3.450458in}{4.169538in}}{\pgfqpoint{3.450458in}{4.175362in}}%
\pgfpathcurveto{\pgfqpoint{3.450458in}{4.181186in}}{\pgfqpoint{3.448144in}{4.186772in}}{\pgfqpoint{3.444026in}{4.190890in}}%
\pgfpathcurveto{\pgfqpoint{3.439908in}{4.195008in}}{\pgfqpoint{3.434322in}{4.197322in}}{\pgfqpoint{3.428498in}{4.197322in}}%
\pgfpathcurveto{\pgfqpoint{3.422674in}{4.197322in}}{\pgfqpoint{3.417088in}{4.195008in}}{\pgfqpoint{3.412970in}{4.190890in}}%
\pgfpathcurveto{\pgfqpoint{3.408852in}{4.186772in}}{\pgfqpoint{3.406538in}{4.181186in}}{\pgfqpoint{3.406538in}{4.175362in}}%
\pgfpathcurveto{\pgfqpoint{3.406538in}{4.169538in}}{\pgfqpoint{3.408852in}{4.163952in}}{\pgfqpoint{3.412970in}{4.159834in}}%
\pgfpathcurveto{\pgfqpoint{3.417088in}{4.155715in}}{\pgfqpoint{3.422674in}{4.153402in}}{\pgfqpoint{3.428498in}{4.153402in}}%
\pgfpathlineto{\pgfqpoint{3.428498in}{4.153402in}}%
\pgfpathclose%
\pgfusepath{stroke,fill}%
\end{pgfscope}%
\begin{pgfscope}%
\pgfpathrectangle{\pgfqpoint{1.414736in}{0.548769in}}{\pgfqpoint{3.497747in}{7.253186in}}%
\pgfusepath{clip}%
\pgfsetbuttcap%
\pgfsetroundjoin%
\definecolor{currentfill}{rgb}{0.585678,0.846661,0.249897}%
\pgfsetfillcolor{currentfill}%
\pgfsetlinewidth{1.003750pt}%
\definecolor{currentstroke}{rgb}{0.585678,0.846661,0.249897}%
\pgfsetstrokecolor{currentstroke}%
\pgfsetdash{}{0pt}%
\pgfpathmoveto{\pgfqpoint{2.946899in}{4.624388in}}%
\pgfpathcurveto{\pgfqpoint{2.952723in}{4.624388in}}{\pgfqpoint{2.958310in}{4.626702in}}{\pgfqpoint{2.962428in}{4.630820in}}%
\pgfpathcurveto{\pgfqpoint{2.966546in}{4.634938in}}{\pgfqpoint{2.968860in}{4.640524in}}{\pgfqpoint{2.968860in}{4.646348in}}%
\pgfpathcurveto{\pgfqpoint{2.968860in}{4.652172in}}{\pgfqpoint{2.966546in}{4.657758in}}{\pgfqpoint{2.962428in}{4.661876in}}%
\pgfpathcurveto{\pgfqpoint{2.958310in}{4.665994in}}{\pgfqpoint{2.952723in}{4.668308in}}{\pgfqpoint{2.946899in}{4.668308in}}%
\pgfpathcurveto{\pgfqpoint{2.941076in}{4.668308in}}{\pgfqpoint{2.935489in}{4.665994in}}{\pgfqpoint{2.931371in}{4.661876in}}%
\pgfpathcurveto{\pgfqpoint{2.927253in}{4.657758in}}{\pgfqpoint{2.924939in}{4.652172in}}{\pgfqpoint{2.924939in}{4.646348in}}%
\pgfpathcurveto{\pgfqpoint{2.924939in}{4.640524in}}{\pgfqpoint{2.927253in}{4.634938in}}{\pgfqpoint{2.931371in}{4.630820in}}%
\pgfpathcurveto{\pgfqpoint{2.935489in}{4.626702in}}{\pgfqpoint{2.941076in}{4.624388in}}{\pgfqpoint{2.946899in}{4.624388in}}%
\pgfpathlineto{\pgfqpoint{2.946899in}{4.624388in}}%
\pgfpathclose%
\pgfusepath{stroke,fill}%
\end{pgfscope}%
\begin{pgfscope}%
\pgfpathrectangle{\pgfqpoint{1.414736in}{0.548769in}}{\pgfqpoint{3.497747in}{7.253186in}}%
\pgfusepath{clip}%
\pgfsetbuttcap%
\pgfsetroundjoin%
\definecolor{currentfill}{rgb}{0.282290,0.145912,0.461510}%
\pgfsetfillcolor{currentfill}%
\pgfsetlinewidth{1.003750pt}%
\definecolor{currentstroke}{rgb}{0.282290,0.145912,0.461510}%
\pgfsetstrokecolor{currentstroke}%
\pgfsetdash{}{0pt}%
\pgfpathmoveto{\pgfqpoint{3.279309in}{5.095374in}}%
\pgfpathcurveto{\pgfqpoint{3.285133in}{5.095374in}}{\pgfqpoint{3.290719in}{5.097688in}}{\pgfqpoint{3.294837in}{5.101806in}}%
\pgfpathcurveto{\pgfqpoint{3.298956in}{5.105924in}}{\pgfqpoint{3.301269in}{5.111510in}}{\pgfqpoint{3.301269in}{5.117334in}}%
\pgfpathcurveto{\pgfqpoint{3.301269in}{5.123158in}}{\pgfqpoint{3.298956in}{5.128744in}}{\pgfqpoint{3.294837in}{5.132862in}}%
\pgfpathcurveto{\pgfqpoint{3.290719in}{5.136980in}}{\pgfqpoint{3.285133in}{5.139294in}}{\pgfqpoint{3.279309in}{5.139294in}}%
\pgfpathcurveto{\pgfqpoint{3.273485in}{5.139294in}}{\pgfqpoint{3.267899in}{5.136980in}}{\pgfqpoint{3.263781in}{5.132862in}}%
\pgfpathcurveto{\pgfqpoint{3.259663in}{5.128744in}}{\pgfqpoint{3.257349in}{5.123158in}}{\pgfqpoint{3.257349in}{5.117334in}}%
\pgfpathcurveto{\pgfqpoint{3.257349in}{5.111510in}}{\pgfqpoint{3.259663in}{5.105924in}}{\pgfqpoint{3.263781in}{5.101806in}}%
\pgfpathcurveto{\pgfqpoint{3.267899in}{5.097688in}}{\pgfqpoint{3.273485in}{5.095374in}}{\pgfqpoint{3.279309in}{5.095374in}}%
\pgfpathlineto{\pgfqpoint{3.279309in}{5.095374in}}%
\pgfpathclose%
\pgfusepath{stroke,fill}%
\end{pgfscope}%
\begin{pgfscope}%
\pgfpathrectangle{\pgfqpoint{1.414736in}{0.548769in}}{\pgfqpoint{3.497747in}{7.253186in}}%
\pgfusepath{clip}%
\pgfsetbuttcap%
\pgfsetroundjoin%
\definecolor{currentfill}{rgb}{0.195860,0.395433,0.555276}%
\pgfsetfillcolor{currentfill}%
\pgfsetlinewidth{1.003750pt}%
\definecolor{currentstroke}{rgb}{0.195860,0.395433,0.555276}%
\pgfsetstrokecolor{currentstroke}%
\pgfsetdash{}{0pt}%
\pgfpathmoveto{\pgfqpoint{4.103111in}{5.566360in}}%
\pgfpathcurveto{\pgfqpoint{4.108935in}{5.566360in}}{\pgfqpoint{4.114521in}{5.568674in}}{\pgfqpoint{4.118639in}{5.572792in}}%
\pgfpathcurveto{\pgfqpoint{4.122757in}{5.576910in}}{\pgfqpoint{4.125071in}{5.582496in}}{\pgfqpoint{4.125071in}{5.588320in}}%
\pgfpathcurveto{\pgfqpoint{4.125071in}{5.594144in}}{\pgfqpoint{4.122757in}{5.599730in}}{\pgfqpoint{4.118639in}{5.603848in}}%
\pgfpathcurveto{\pgfqpoint{4.114521in}{5.607967in}}{\pgfqpoint{4.108935in}{5.610280in}}{\pgfqpoint{4.103111in}{5.610280in}}%
\pgfpathcurveto{\pgfqpoint{4.097287in}{5.610280in}}{\pgfqpoint{4.091701in}{5.607967in}}{\pgfqpoint{4.087582in}{5.603848in}}%
\pgfpathcurveto{\pgfqpoint{4.083464in}{5.599730in}}{\pgfqpoint{4.081150in}{5.594144in}}{\pgfqpoint{4.081150in}{5.588320in}}%
\pgfpathcurveto{\pgfqpoint{4.081150in}{5.582496in}}{\pgfqpoint{4.083464in}{5.576910in}}{\pgfqpoint{4.087582in}{5.572792in}}%
\pgfpathcurveto{\pgfqpoint{4.091701in}{5.568674in}}{\pgfqpoint{4.097287in}{5.566360in}}{\pgfqpoint{4.103111in}{5.566360in}}%
\pgfpathlineto{\pgfqpoint{4.103111in}{5.566360in}}%
\pgfpathclose%
\pgfusepath{stroke,fill}%
\end{pgfscope}%
\begin{pgfscope}%
\pgfpathrectangle{\pgfqpoint{1.414736in}{0.548769in}}{\pgfqpoint{3.497747in}{7.253186in}}%
\pgfusepath{clip}%
\pgfsetbuttcap%
\pgfsetroundjoin%
\definecolor{currentfill}{rgb}{0.140210,0.665859,0.513427}%
\pgfsetfillcolor{currentfill}%
\pgfsetlinewidth{1.003750pt}%
\definecolor{currentstroke}{rgb}{0.140210,0.665859,0.513427}%
\pgfsetstrokecolor{currentstroke}%
\pgfsetdash{}{0pt}%
\pgfpathmoveto{\pgfqpoint{3.878364in}{6.037346in}}%
\pgfpathcurveto{\pgfqpoint{3.884188in}{6.037346in}}{\pgfqpoint{3.889774in}{6.039660in}}{\pgfqpoint{3.893892in}{6.043778in}}%
\pgfpathcurveto{\pgfqpoint{3.898010in}{6.047896in}}{\pgfqpoint{3.900324in}{6.053482in}}{\pgfqpoint{3.900324in}{6.059306in}}%
\pgfpathcurveto{\pgfqpoint{3.900324in}{6.065130in}}{\pgfqpoint{3.898010in}{6.070716in}}{\pgfqpoint{3.893892in}{6.074835in}}%
\pgfpathcurveto{\pgfqpoint{3.889774in}{6.078953in}}{\pgfqpoint{3.884188in}{6.081267in}}{\pgfqpoint{3.878364in}{6.081267in}}%
\pgfpathcurveto{\pgfqpoint{3.872540in}{6.081267in}}{\pgfqpoint{3.866954in}{6.078953in}}{\pgfqpoint{3.862836in}{6.074835in}}%
\pgfpathcurveto{\pgfqpoint{3.858717in}{6.070716in}}{\pgfqpoint{3.856404in}{6.065130in}}{\pgfqpoint{3.856404in}{6.059306in}}%
\pgfpathcurveto{\pgfqpoint{3.856404in}{6.053482in}}{\pgfqpoint{3.858717in}{6.047896in}}{\pgfqpoint{3.862836in}{6.043778in}}%
\pgfpathcurveto{\pgfqpoint{3.866954in}{6.039660in}}{\pgfqpoint{3.872540in}{6.037346in}}{\pgfqpoint{3.878364in}{6.037346in}}%
\pgfpathlineto{\pgfqpoint{3.878364in}{6.037346in}}%
\pgfpathclose%
\pgfusepath{stroke,fill}%
\end{pgfscope}%
\begin{pgfscope}%
\pgfpathrectangle{\pgfqpoint{1.414736in}{0.548769in}}{\pgfqpoint{3.497747in}{7.253186in}}%
\pgfusepath{clip}%
\pgfsetbuttcap%
\pgfsetroundjoin%
\definecolor{currentfill}{rgb}{0.993248,0.906157,0.143936}%
\pgfsetfillcolor{currentfill}%
\pgfsetlinewidth{1.003750pt}%
\definecolor{currentstroke}{rgb}{0.993248,0.906157,0.143936}%
\pgfsetstrokecolor{currentstroke}%
\pgfsetdash{}{0pt}%
\pgfpathmoveto{\pgfqpoint{3.141189in}{6.508332in}}%
\pgfpathcurveto{\pgfqpoint{3.147013in}{6.508332in}}{\pgfqpoint{3.152600in}{6.510646in}}{\pgfqpoint{3.156718in}{6.514764in}}%
\pgfpathcurveto{\pgfqpoint{3.160836in}{6.518882in}}{\pgfqpoint{3.163150in}{6.524468in}}{\pgfqpoint{3.163150in}{6.530292in}}%
\pgfpathcurveto{\pgfqpoint{3.163150in}{6.536116in}}{\pgfqpoint{3.160836in}{6.541702in}}{\pgfqpoint{3.156718in}{6.545821in}}%
\pgfpathcurveto{\pgfqpoint{3.152600in}{6.549939in}}{\pgfqpoint{3.147013in}{6.552253in}}{\pgfqpoint{3.141189in}{6.552253in}}%
\pgfpathcurveto{\pgfqpoint{3.135366in}{6.552253in}}{\pgfqpoint{3.129779in}{6.549939in}}{\pgfqpoint{3.125661in}{6.545821in}}%
\pgfpathcurveto{\pgfqpoint{3.121543in}{6.541702in}}{\pgfqpoint{3.119229in}{6.536116in}}{\pgfqpoint{3.119229in}{6.530292in}}%
\pgfpathcurveto{\pgfqpoint{3.119229in}{6.524468in}}{\pgfqpoint{3.121543in}{6.518882in}}{\pgfqpoint{3.125661in}{6.514764in}}%
\pgfpathcurveto{\pgfqpoint{3.129779in}{6.510646in}}{\pgfqpoint{3.135366in}{6.508332in}}{\pgfqpoint{3.141189in}{6.508332in}}%
\pgfpathlineto{\pgfqpoint{3.141189in}{6.508332in}}%
\pgfpathclose%
\pgfusepath{stroke,fill}%
\end{pgfscope}%
\begin{pgfscope}%
\pgfpathrectangle{\pgfqpoint{1.414736in}{0.548769in}}{\pgfqpoint{3.497747in}{7.253186in}}%
\pgfusepath{clip}%
\pgfsetbuttcap%
\pgfsetroundjoin%
\definecolor{currentfill}{rgb}{0.906311,0.894855,0.098125}%
\pgfsetfillcolor{currentfill}%
\pgfsetlinewidth{1.003750pt}%
\definecolor{currentstroke}{rgb}{0.906311,0.894855,0.098125}%
\pgfsetstrokecolor{currentstroke}%
\pgfsetdash{}{0pt}%
\pgfpathmoveto{\pgfqpoint{4.257179in}{6.979318in}}%
\pgfpathcurveto{\pgfqpoint{4.263003in}{6.979318in}}{\pgfqpoint{4.268589in}{6.981632in}}{\pgfqpoint{4.272707in}{6.985750in}}%
\pgfpathcurveto{\pgfqpoint{4.276826in}{6.989868in}}{\pgfqpoint{4.279139in}{6.995455in}}{\pgfqpoint{4.279139in}{7.001278in}}%
\pgfpathcurveto{\pgfqpoint{4.279139in}{7.007102in}}{\pgfqpoint{4.276826in}{7.012689in}}{\pgfqpoint{4.272707in}{7.016807in}}%
\pgfpathcurveto{\pgfqpoint{4.268589in}{7.020925in}}{\pgfqpoint{4.263003in}{7.023239in}}{\pgfqpoint{4.257179in}{7.023239in}}%
\pgfpathcurveto{\pgfqpoint{4.251355in}{7.023239in}}{\pgfqpoint{4.245769in}{7.020925in}}{\pgfqpoint{4.241651in}{7.016807in}}%
\pgfpathcurveto{\pgfqpoint{4.237533in}{7.012689in}}{\pgfqpoint{4.235219in}{7.007102in}}{\pgfqpoint{4.235219in}{7.001278in}}%
\pgfpathcurveto{\pgfqpoint{4.235219in}{6.995455in}}{\pgfqpoint{4.237533in}{6.989868in}}{\pgfqpoint{4.241651in}{6.985750in}}%
\pgfpathcurveto{\pgfqpoint{4.245769in}{6.981632in}}{\pgfqpoint{4.251355in}{6.979318in}}{\pgfqpoint{4.257179in}{6.979318in}}%
\pgfpathlineto{\pgfqpoint{4.257179in}{6.979318in}}%
\pgfpathclose%
\pgfusepath{stroke,fill}%
\end{pgfscope}%
\begin{pgfscope}%
\pgfpathrectangle{\pgfqpoint{1.414736in}{0.548769in}}{\pgfqpoint{3.497747in}{7.253186in}}%
\pgfusepath{clip}%
\pgfsetbuttcap%
\pgfsetroundjoin%
\definecolor{currentfill}{rgb}{0.688944,0.865448,0.182725}%
\pgfsetfillcolor{currentfill}%
\pgfsetlinewidth{1.003750pt}%
\definecolor{currentstroke}{rgb}{0.688944,0.865448,0.182725}%
\pgfsetstrokecolor{currentstroke}%
\pgfsetdash{}{0pt}%
\pgfpathmoveto{\pgfqpoint{3.998279in}{7.450304in}}%
\pgfpathcurveto{\pgfqpoint{4.004103in}{7.450304in}}{\pgfqpoint{4.009690in}{7.452618in}}{\pgfqpoint{4.013808in}{7.456736in}}%
\pgfpathcurveto{\pgfqpoint{4.017926in}{7.460854in}}{\pgfqpoint{4.020240in}{7.466441in}}{\pgfqpoint{4.020240in}{7.472265in}}%
\pgfpathcurveto{\pgfqpoint{4.020240in}{7.478089in}}{\pgfqpoint{4.017926in}{7.483675in}}{\pgfqpoint{4.013808in}{7.487793in}}%
\pgfpathcurveto{\pgfqpoint{4.009690in}{7.491911in}}{\pgfqpoint{4.004103in}{7.494225in}}{\pgfqpoint{3.998279in}{7.494225in}}%
\pgfpathcurveto{\pgfqpoint{3.992455in}{7.494225in}}{\pgfqpoint{3.986869in}{7.491911in}}{\pgfqpoint{3.982751in}{7.487793in}}%
\pgfpathcurveto{\pgfqpoint{3.978633in}{7.483675in}}{\pgfqpoint{3.976319in}{7.478089in}}{\pgfqpoint{3.976319in}{7.472265in}}%
\pgfpathcurveto{\pgfqpoint{3.976319in}{7.466441in}}{\pgfqpoint{3.978633in}{7.460854in}}{\pgfqpoint{3.982751in}{7.456736in}}%
\pgfpathcurveto{\pgfqpoint{3.986869in}{7.452618in}}{\pgfqpoint{3.992455in}{7.450304in}}{\pgfqpoint{3.998279in}{7.450304in}}%
\pgfpathlineto{\pgfqpoint{3.998279in}{7.450304in}}%
\pgfpathclose%
\pgfusepath{stroke,fill}%
\end{pgfscope}%
\begin{pgfscope}%
\pgfpathrectangle{\pgfqpoint{1.414736in}{0.548769in}}{\pgfqpoint{3.497747in}{7.253186in}}%
\pgfusepath{clip}%
\pgfsetbuttcap%
\pgfsetroundjoin%
\definecolor{currentfill}{rgb}{0.964894,0.902323,0.123941}%
\pgfsetfillcolor{currentfill}%
\pgfsetlinewidth{1.003750pt}%
\definecolor{currentstroke}{rgb}{0.964894,0.902323,0.123941}%
\pgfsetstrokecolor{currentstroke}%
\pgfsetdash{}{0pt}%
\pgfpathmoveto{\pgfqpoint{3.472309in}{0.856499in}}%
\pgfpathcurveto{\pgfqpoint{3.478133in}{0.856499in}}{\pgfqpoint{3.483719in}{0.858813in}}{\pgfqpoint{3.487837in}{0.862931in}}%
\pgfpathcurveto{\pgfqpoint{3.491955in}{0.867049in}}{\pgfqpoint{3.494269in}{0.872635in}}{\pgfqpoint{3.494269in}{0.878459in}}%
\pgfpathcurveto{\pgfqpoint{3.494269in}{0.884283in}}{\pgfqpoint{3.491955in}{0.889869in}}{\pgfqpoint{3.487837in}{0.893987in}}%
\pgfpathcurveto{\pgfqpoint{3.483719in}{0.898105in}}{\pgfqpoint{3.478133in}{0.900419in}}{\pgfqpoint{3.472309in}{0.900419in}}%
\pgfpathcurveto{\pgfqpoint{3.466485in}{0.900419in}}{\pgfqpoint{3.460898in}{0.898105in}}{\pgfqpoint{3.456780in}{0.893987in}}%
\pgfpathcurveto{\pgfqpoint{3.452662in}{0.889869in}}{\pgfqpoint{3.450348in}{0.884283in}}{\pgfqpoint{3.450348in}{0.878459in}}%
\pgfpathcurveto{\pgfqpoint{3.450348in}{0.872635in}}{\pgfqpoint{3.452662in}{0.867049in}}{\pgfqpoint{3.456780in}{0.862931in}}%
\pgfpathcurveto{\pgfqpoint{3.460898in}{0.858813in}}{\pgfqpoint{3.466485in}{0.856499in}}{\pgfqpoint{3.472309in}{0.856499in}}%
\pgfpathlineto{\pgfqpoint{3.472309in}{0.856499in}}%
\pgfpathclose%
\pgfusepath{stroke,fill}%
\end{pgfscope}%
\begin{pgfscope}%
\pgfpathrectangle{\pgfqpoint{1.414736in}{0.548769in}}{\pgfqpoint{3.497747in}{7.253186in}}%
\pgfusepath{clip}%
\pgfsetbuttcap%
\pgfsetroundjoin%
\definecolor{currentfill}{rgb}{0.412913,0.803041,0.357269}%
\pgfsetfillcolor{currentfill}%
\pgfsetlinewidth{1.003750pt}%
\definecolor{currentstroke}{rgb}{0.412913,0.803041,0.357269}%
\pgfsetstrokecolor{currentstroke}%
\pgfsetdash{}{0pt}%
\pgfpathmoveto{\pgfqpoint{3.777909in}{1.327485in}}%
\pgfpathcurveto{\pgfqpoint{3.783733in}{1.327485in}}{\pgfqpoint{3.789319in}{1.329799in}}{\pgfqpoint{3.793437in}{1.333917in}}%
\pgfpathcurveto{\pgfqpoint{3.797555in}{1.338035in}}{\pgfqpoint{3.799869in}{1.343621in}}{\pgfqpoint{3.799869in}{1.349445in}}%
\pgfpathcurveto{\pgfqpoint{3.799869in}{1.355269in}}{\pgfqpoint{3.797555in}{1.360855in}}{\pgfqpoint{3.793437in}{1.364973in}}%
\pgfpathcurveto{\pgfqpoint{3.789319in}{1.369092in}}{\pgfqpoint{3.783733in}{1.371405in}}{\pgfqpoint{3.777909in}{1.371405in}}%
\pgfpathcurveto{\pgfqpoint{3.772085in}{1.371405in}}{\pgfqpoint{3.766499in}{1.369092in}}{\pgfqpoint{3.762381in}{1.364973in}}%
\pgfpathcurveto{\pgfqpoint{3.758263in}{1.360855in}}{\pgfqpoint{3.755949in}{1.355269in}}{\pgfqpoint{3.755949in}{1.349445in}}%
\pgfpathcurveto{\pgfqpoint{3.755949in}{1.343621in}}{\pgfqpoint{3.758263in}{1.338035in}}{\pgfqpoint{3.762381in}{1.333917in}}%
\pgfpathcurveto{\pgfqpoint{3.766499in}{1.329799in}}{\pgfqpoint{3.772085in}{1.327485in}}{\pgfqpoint{3.777909in}{1.327485in}}%
\pgfpathlineto{\pgfqpoint{3.777909in}{1.327485in}}%
\pgfpathclose%
\pgfusepath{stroke,fill}%
\end{pgfscope}%
\begin{pgfscope}%
\pgfpathrectangle{\pgfqpoint{1.414736in}{0.548769in}}{\pgfqpoint{3.497747in}{7.253186in}}%
\pgfusepath{clip}%
\pgfsetbuttcap%
\pgfsetroundjoin%
\definecolor{currentfill}{rgb}{0.430983,0.808473,0.346476}%
\pgfsetfillcolor{currentfill}%
\pgfsetlinewidth{1.003750pt}%
\definecolor{currentstroke}{rgb}{0.430983,0.808473,0.346476}%
\pgfsetstrokecolor{currentstroke}%
\pgfsetdash{}{0pt}%
\pgfpathmoveto{\pgfqpoint{3.377050in}{1.798471in}}%
\pgfpathcurveto{\pgfqpoint{3.382874in}{1.798471in}}{\pgfqpoint{3.388460in}{1.800785in}}{\pgfqpoint{3.392579in}{1.804903in}}%
\pgfpathcurveto{\pgfqpoint{3.396697in}{1.809021in}}{\pgfqpoint{3.399011in}{1.814607in}}{\pgfqpoint{3.399011in}{1.820431in}}%
\pgfpathcurveto{\pgfqpoint{3.399011in}{1.826255in}}{\pgfqpoint{3.396697in}{1.831841in}}{\pgfqpoint{3.392579in}{1.835960in}}%
\pgfpathcurveto{\pgfqpoint{3.388460in}{1.840078in}}{\pgfqpoint{3.382874in}{1.842392in}}{\pgfqpoint{3.377050in}{1.842392in}}%
\pgfpathcurveto{\pgfqpoint{3.371226in}{1.842392in}}{\pgfqpoint{3.365640in}{1.840078in}}{\pgfqpoint{3.361522in}{1.835960in}}%
\pgfpathcurveto{\pgfqpoint{3.357404in}{1.831841in}}{\pgfqpoint{3.355090in}{1.826255in}}{\pgfqpoint{3.355090in}{1.820431in}}%
\pgfpathcurveto{\pgfqpoint{3.355090in}{1.814607in}}{\pgfqpoint{3.357404in}{1.809021in}}{\pgfqpoint{3.361522in}{1.804903in}}%
\pgfpathcurveto{\pgfqpoint{3.365640in}{1.800785in}}{\pgfqpoint{3.371226in}{1.798471in}}{\pgfqpoint{3.377050in}{1.798471in}}%
\pgfpathlineto{\pgfqpoint{3.377050in}{1.798471in}}%
\pgfpathclose%
\pgfusepath{stroke,fill}%
\end{pgfscope}%
\begin{pgfscope}%
\pgfpathrectangle{\pgfqpoint{1.414736in}{0.548769in}}{\pgfqpoint{3.497747in}{7.253186in}}%
\pgfusepath{clip}%
\pgfsetbuttcap%
\pgfsetroundjoin%
\definecolor{currentfill}{rgb}{0.935904,0.898570,0.108131}%
\pgfsetfillcolor{currentfill}%
\pgfsetlinewidth{1.003750pt}%
\definecolor{currentstroke}{rgb}{0.935904,0.898570,0.108131}%
\pgfsetstrokecolor{currentstroke}%
\pgfsetdash{}{0pt}%
\pgfpathmoveto{\pgfqpoint{3.410261in}{2.269457in}}%
\pgfpathcurveto{\pgfqpoint{3.416085in}{2.269457in}}{\pgfqpoint{3.421671in}{2.271771in}}{\pgfqpoint{3.425789in}{2.275889in}}%
\pgfpathcurveto{\pgfqpoint{3.429907in}{2.280007in}}{\pgfqpoint{3.432221in}{2.285593in}}{\pgfqpoint{3.432221in}{2.291417in}}%
\pgfpathcurveto{\pgfqpoint{3.432221in}{2.297241in}}{\pgfqpoint{3.429907in}{2.302828in}}{\pgfqpoint{3.425789in}{2.306946in}}%
\pgfpathcurveto{\pgfqpoint{3.421671in}{2.311064in}}{\pgfqpoint{3.416085in}{2.313378in}}{\pgfqpoint{3.410261in}{2.313378in}}%
\pgfpathcurveto{\pgfqpoint{3.404437in}{2.313378in}}{\pgfqpoint{3.398851in}{2.311064in}}{\pgfqpoint{3.394732in}{2.306946in}}%
\pgfpathcurveto{\pgfqpoint{3.390614in}{2.302828in}}{\pgfqpoint{3.388300in}{2.297241in}}{\pgfqpoint{3.388300in}{2.291417in}}%
\pgfpathcurveto{\pgfqpoint{3.388300in}{2.285593in}}{\pgfqpoint{3.390614in}{2.280007in}}{\pgfqpoint{3.394732in}{2.275889in}}%
\pgfpathcurveto{\pgfqpoint{3.398851in}{2.271771in}}{\pgfqpoint{3.404437in}{2.269457in}}{\pgfqpoint{3.410261in}{2.269457in}}%
\pgfpathlineto{\pgfqpoint{3.410261in}{2.269457in}}%
\pgfpathclose%
\pgfusepath{stroke,fill}%
\end{pgfscope}%
\begin{pgfscope}%
\pgfpathrectangle{\pgfqpoint{1.414736in}{0.548769in}}{\pgfqpoint{3.497747in}{7.253186in}}%
\pgfusepath{clip}%
\pgfsetbuttcap%
\pgfsetroundjoin%
\definecolor{currentfill}{rgb}{0.239374,0.735588,0.455688}%
\pgfsetfillcolor{currentfill}%
\pgfsetlinewidth{1.003750pt}%
\definecolor{currentstroke}{rgb}{0.239374,0.735588,0.455688}%
\pgfsetstrokecolor{currentstroke}%
\pgfsetdash{}{0pt}%
\pgfpathmoveto{\pgfqpoint{3.871751in}{2.740443in}}%
\pgfpathcurveto{\pgfqpoint{3.877575in}{2.740443in}}{\pgfqpoint{3.883161in}{2.742757in}}{\pgfqpoint{3.887280in}{2.746875in}}%
\pgfpathcurveto{\pgfqpoint{3.891398in}{2.750993in}}{\pgfqpoint{3.893712in}{2.756580in}}{\pgfqpoint{3.893712in}{2.762404in}}%
\pgfpathcurveto{\pgfqpoint{3.893712in}{2.768227in}}{\pgfqpoint{3.891398in}{2.773814in}}{\pgfqpoint{3.887280in}{2.777932in}}%
\pgfpathcurveto{\pgfqpoint{3.883161in}{2.782050in}}{\pgfqpoint{3.877575in}{2.784364in}}{\pgfqpoint{3.871751in}{2.784364in}}%
\pgfpathcurveto{\pgfqpoint{3.865927in}{2.784364in}}{\pgfqpoint{3.860341in}{2.782050in}}{\pgfqpoint{3.856223in}{2.777932in}}%
\pgfpathcurveto{\pgfqpoint{3.852105in}{2.773814in}}{\pgfqpoint{3.849791in}{2.768227in}}{\pgfqpoint{3.849791in}{2.762404in}}%
\pgfpathcurveto{\pgfqpoint{3.849791in}{2.756580in}}{\pgfqpoint{3.852105in}{2.750993in}}{\pgfqpoint{3.856223in}{2.746875in}}%
\pgfpathcurveto{\pgfqpoint{3.860341in}{2.742757in}}{\pgfqpoint{3.865927in}{2.740443in}}{\pgfqpoint{3.871751in}{2.740443in}}%
\pgfpathlineto{\pgfqpoint{3.871751in}{2.740443in}}%
\pgfpathclose%
\pgfusepath{stroke,fill}%
\end{pgfscope}%
\begin{pgfscope}%
\pgfpathrectangle{\pgfqpoint{1.414736in}{0.548769in}}{\pgfqpoint{3.497747in}{7.253186in}}%
\pgfusepath{clip}%
\pgfsetbuttcap%
\pgfsetroundjoin%
\definecolor{currentfill}{rgb}{0.595839,0.848717,0.243329}%
\pgfsetfillcolor{currentfill}%
\pgfsetlinewidth{1.003750pt}%
\definecolor{currentstroke}{rgb}{0.595839,0.848717,0.243329}%
\pgfsetstrokecolor{currentstroke}%
\pgfsetdash{}{0pt}%
\pgfpathmoveto{\pgfqpoint{3.356287in}{3.211429in}}%
\pgfpathcurveto{\pgfqpoint{3.362111in}{3.211429in}}{\pgfqpoint{3.367697in}{3.213743in}}{\pgfqpoint{3.371815in}{3.217861in}}%
\pgfpathcurveto{\pgfqpoint{3.375933in}{3.221980in}}{\pgfqpoint{3.378247in}{3.227566in}}{\pgfqpoint{3.378247in}{3.233390in}}%
\pgfpathcurveto{\pgfqpoint{3.378247in}{3.239214in}}{\pgfqpoint{3.375933in}{3.244800in}}{\pgfqpoint{3.371815in}{3.248918in}}%
\pgfpathcurveto{\pgfqpoint{3.367697in}{3.253036in}}{\pgfqpoint{3.362111in}{3.255350in}}{\pgfqpoint{3.356287in}{3.255350in}}%
\pgfpathcurveto{\pgfqpoint{3.350463in}{3.255350in}}{\pgfqpoint{3.344877in}{3.253036in}}{\pgfqpoint{3.340758in}{3.248918in}}%
\pgfpathcurveto{\pgfqpoint{3.336640in}{3.244800in}}{\pgfqpoint{3.334326in}{3.239214in}}{\pgfqpoint{3.334326in}{3.233390in}}%
\pgfpathcurveto{\pgfqpoint{3.334326in}{3.227566in}}{\pgfqpoint{3.336640in}{3.221980in}}{\pgfqpoint{3.340758in}{3.217861in}}%
\pgfpathcurveto{\pgfqpoint{3.344877in}{3.213743in}}{\pgfqpoint{3.350463in}{3.211429in}}{\pgfqpoint{3.356287in}{3.211429in}}%
\pgfpathlineto{\pgfqpoint{3.356287in}{3.211429in}}%
\pgfpathclose%
\pgfusepath{stroke,fill}%
\end{pgfscope}%
\begin{pgfscope}%
\pgfpathrectangle{\pgfqpoint{1.414736in}{0.548769in}}{\pgfqpoint{3.497747in}{7.253186in}}%
\pgfusepath{clip}%
\pgfsetbuttcap%
\pgfsetroundjoin%
\definecolor{currentfill}{rgb}{0.845561,0.887322,0.099702}%
\pgfsetfillcolor{currentfill}%
\pgfsetlinewidth{1.003750pt}%
\definecolor{currentstroke}{rgb}{0.845561,0.887322,0.099702}%
\pgfsetstrokecolor{currentstroke}%
\pgfsetdash{}{0pt}%
\pgfpathmoveto{\pgfqpoint{3.924055in}{3.682415in}}%
\pgfpathcurveto{\pgfqpoint{3.929879in}{3.682415in}}{\pgfqpoint{3.935465in}{3.684729in}}{\pgfqpoint{3.939583in}{3.688847in}}%
\pgfpathcurveto{\pgfqpoint{3.943701in}{3.692966in}}{\pgfqpoint{3.946015in}{3.698552in}}{\pgfqpoint{3.946015in}{3.704376in}}%
\pgfpathcurveto{\pgfqpoint{3.946015in}{3.710200in}}{\pgfqpoint{3.943701in}{3.715786in}}{\pgfqpoint{3.939583in}{3.719904in}}%
\pgfpathcurveto{\pgfqpoint{3.935465in}{3.724022in}}{\pgfqpoint{3.929879in}{3.726336in}}{\pgfqpoint{3.924055in}{3.726336in}}%
\pgfpathcurveto{\pgfqpoint{3.918231in}{3.726336in}}{\pgfqpoint{3.912645in}{3.724022in}}{\pgfqpoint{3.908527in}{3.719904in}}%
\pgfpathcurveto{\pgfqpoint{3.904409in}{3.715786in}}{\pgfqpoint{3.902095in}{3.710200in}}{\pgfqpoint{3.902095in}{3.704376in}}%
\pgfpathcurveto{\pgfqpoint{3.902095in}{3.698552in}}{\pgfqpoint{3.904409in}{3.692966in}}{\pgfqpoint{3.908527in}{3.688847in}}%
\pgfpathcurveto{\pgfqpoint{3.912645in}{3.684729in}}{\pgfqpoint{3.918231in}{3.682415in}}{\pgfqpoint{3.924055in}{3.682415in}}%
\pgfpathlineto{\pgfqpoint{3.924055in}{3.682415in}}%
\pgfpathclose%
\pgfusepath{stroke,fill}%
\end{pgfscope}%
\begin{pgfscope}%
\pgfpathrectangle{\pgfqpoint{1.414736in}{0.548769in}}{\pgfqpoint{3.497747in}{7.253186in}}%
\pgfusepath{clip}%
\pgfsetbuttcap%
\pgfsetroundjoin%
\definecolor{currentfill}{rgb}{0.993248,0.906157,0.143936}%
\pgfsetfillcolor{currentfill}%
\pgfsetlinewidth{1.003750pt}%
\definecolor{currentstroke}{rgb}{0.993248,0.906157,0.143936}%
\pgfsetstrokecolor{currentstroke}%
\pgfsetdash{}{0pt}%
\pgfpathmoveto{\pgfqpoint{3.474955in}{4.153402in}}%
\pgfpathcurveto{\pgfqpoint{3.480779in}{4.153402in}}{\pgfqpoint{3.486365in}{4.155715in}}{\pgfqpoint{3.490483in}{4.159834in}}%
\pgfpathcurveto{\pgfqpoint{3.494602in}{4.163952in}}{\pgfqpoint{3.496915in}{4.169538in}}{\pgfqpoint{3.496915in}{4.175362in}}%
\pgfpathcurveto{\pgfqpoint{3.496915in}{4.181186in}}{\pgfqpoint{3.494602in}{4.186772in}}{\pgfqpoint{3.490483in}{4.190890in}}%
\pgfpathcurveto{\pgfqpoint{3.486365in}{4.195008in}}{\pgfqpoint{3.480779in}{4.197322in}}{\pgfqpoint{3.474955in}{4.197322in}}%
\pgfpathcurveto{\pgfqpoint{3.469131in}{4.197322in}}{\pgfqpoint{3.463545in}{4.195008in}}{\pgfqpoint{3.459427in}{4.190890in}}%
\pgfpathcurveto{\pgfqpoint{3.455309in}{4.186772in}}{\pgfqpoint{3.452995in}{4.181186in}}{\pgfqpoint{3.452995in}{4.175362in}}%
\pgfpathcurveto{\pgfqpoint{3.452995in}{4.169538in}}{\pgfqpoint{3.455309in}{4.163952in}}{\pgfqpoint{3.459427in}{4.159834in}}%
\pgfpathcurveto{\pgfqpoint{3.463545in}{4.155715in}}{\pgfqpoint{3.469131in}{4.153402in}}{\pgfqpoint{3.474955in}{4.153402in}}%
\pgfpathlineto{\pgfqpoint{3.474955in}{4.153402in}}%
\pgfpathclose%
\pgfusepath{stroke,fill}%
\end{pgfscope}%
\begin{pgfscope}%
\pgfpathrectangle{\pgfqpoint{1.414736in}{0.548769in}}{\pgfqpoint{3.497747in}{7.253186in}}%
\pgfusepath{clip}%
\pgfsetbuttcap%
\pgfsetroundjoin%
\definecolor{currentfill}{rgb}{0.730889,0.871916,0.156029}%
\pgfsetfillcolor{currentfill}%
\pgfsetlinewidth{1.003750pt}%
\definecolor{currentstroke}{rgb}{0.730889,0.871916,0.156029}%
\pgfsetstrokecolor{currentstroke}%
\pgfsetdash{}{0pt}%
\pgfpathmoveto{\pgfqpoint{3.424805in}{4.624388in}}%
\pgfpathcurveto{\pgfqpoint{3.430629in}{4.624388in}}{\pgfqpoint{3.436216in}{4.626702in}}{\pgfqpoint{3.440334in}{4.630820in}}%
\pgfpathcurveto{\pgfqpoint{3.444452in}{4.634938in}}{\pgfqpoint{3.446766in}{4.640524in}}{\pgfqpoint{3.446766in}{4.646348in}}%
\pgfpathcurveto{\pgfqpoint{3.446766in}{4.652172in}}{\pgfqpoint{3.444452in}{4.657758in}}{\pgfqpoint{3.440334in}{4.661876in}}%
\pgfpathcurveto{\pgfqpoint{3.436216in}{4.665994in}}{\pgfqpoint{3.430629in}{4.668308in}}{\pgfqpoint{3.424805in}{4.668308in}}%
\pgfpathcurveto{\pgfqpoint{3.418981in}{4.668308in}}{\pgfqpoint{3.413395in}{4.665994in}}{\pgfqpoint{3.409277in}{4.661876in}}%
\pgfpathcurveto{\pgfqpoint{3.405159in}{4.657758in}}{\pgfqpoint{3.402845in}{4.652172in}}{\pgfqpoint{3.402845in}{4.646348in}}%
\pgfpathcurveto{\pgfqpoint{3.402845in}{4.640524in}}{\pgfqpoint{3.405159in}{4.634938in}}{\pgfqpoint{3.409277in}{4.630820in}}%
\pgfpathcurveto{\pgfqpoint{3.413395in}{4.626702in}}{\pgfqpoint{3.418981in}{4.624388in}}{\pgfqpoint{3.424805in}{4.624388in}}%
\pgfpathlineto{\pgfqpoint{3.424805in}{4.624388in}}%
\pgfpathclose%
\pgfusepath{stroke,fill}%
\end{pgfscope}%
\begin{pgfscope}%
\pgfpathrectangle{\pgfqpoint{1.414736in}{0.548769in}}{\pgfqpoint{3.497747in}{7.253186in}}%
\pgfusepath{clip}%
\pgfsetbuttcap%
\pgfsetroundjoin%
\definecolor{currentfill}{rgb}{0.272594,0.025563,0.353093}%
\pgfsetfillcolor{currentfill}%
\pgfsetlinewidth{1.003750pt}%
\definecolor{currentstroke}{rgb}{0.272594,0.025563,0.353093}%
\pgfsetstrokecolor{currentstroke}%
\pgfsetdash{}{0pt}%
\pgfpathmoveto{\pgfqpoint{3.275622in}{5.095374in}}%
\pgfpathcurveto{\pgfqpoint{3.281446in}{5.095374in}}{\pgfqpoint{3.287032in}{5.097688in}}{\pgfqpoint{3.291150in}{5.101806in}}%
\pgfpathcurveto{\pgfqpoint{3.295268in}{5.105924in}}{\pgfqpoint{3.297582in}{5.111510in}}{\pgfqpoint{3.297582in}{5.117334in}}%
\pgfpathcurveto{\pgfqpoint{3.297582in}{5.123158in}}{\pgfqpoint{3.295268in}{5.128744in}}{\pgfqpoint{3.291150in}{5.132862in}}%
\pgfpathcurveto{\pgfqpoint{3.287032in}{5.136980in}}{\pgfqpoint{3.281446in}{5.139294in}}{\pgfqpoint{3.275622in}{5.139294in}}%
\pgfpathcurveto{\pgfqpoint{3.269798in}{5.139294in}}{\pgfqpoint{3.264212in}{5.136980in}}{\pgfqpoint{3.260094in}{5.132862in}}%
\pgfpathcurveto{\pgfqpoint{3.255975in}{5.128744in}}{\pgfqpoint{3.253662in}{5.123158in}}{\pgfqpoint{3.253662in}{5.117334in}}%
\pgfpathcurveto{\pgfqpoint{3.253662in}{5.111510in}}{\pgfqpoint{3.255975in}{5.105924in}}{\pgfqpoint{3.260094in}{5.101806in}}%
\pgfpathcurveto{\pgfqpoint{3.264212in}{5.097688in}}{\pgfqpoint{3.269798in}{5.095374in}}{\pgfqpoint{3.275622in}{5.095374in}}%
\pgfpathlineto{\pgfqpoint{3.275622in}{5.095374in}}%
\pgfpathclose%
\pgfusepath{stroke,fill}%
\end{pgfscope}%
\begin{pgfscope}%
\pgfpathrectangle{\pgfqpoint{1.414736in}{0.548769in}}{\pgfqpoint{3.497747in}{7.253186in}}%
\pgfusepath{clip}%
\pgfsetbuttcap%
\pgfsetroundjoin%
\definecolor{currentfill}{rgb}{0.218130,0.347432,0.550038}%
\pgfsetfillcolor{currentfill}%
\pgfsetlinewidth{1.003750pt}%
\definecolor{currentstroke}{rgb}{0.218130,0.347432,0.550038}%
\pgfsetstrokecolor{currentstroke}%
\pgfsetdash{}{0pt}%
\pgfpathmoveto{\pgfqpoint{3.764263in}{5.566360in}}%
\pgfpathcurveto{\pgfqpoint{3.770087in}{5.566360in}}{\pgfqpoint{3.775674in}{5.568674in}}{\pgfqpoint{3.779792in}{5.572792in}}%
\pgfpathcurveto{\pgfqpoint{3.783910in}{5.576910in}}{\pgfqpoint{3.786224in}{5.582496in}}{\pgfqpoint{3.786224in}{5.588320in}}%
\pgfpathcurveto{\pgfqpoint{3.786224in}{5.594144in}}{\pgfqpoint{3.783910in}{5.599730in}}{\pgfqpoint{3.779792in}{5.603848in}}%
\pgfpathcurveto{\pgfqpoint{3.775674in}{5.607967in}}{\pgfqpoint{3.770087in}{5.610280in}}{\pgfqpoint{3.764263in}{5.610280in}}%
\pgfpathcurveto{\pgfqpoint{3.758439in}{5.610280in}}{\pgfqpoint{3.752853in}{5.607967in}}{\pgfqpoint{3.748735in}{5.603848in}}%
\pgfpathcurveto{\pgfqpoint{3.744617in}{5.599730in}}{\pgfqpoint{3.742303in}{5.594144in}}{\pgfqpoint{3.742303in}{5.588320in}}%
\pgfpathcurveto{\pgfqpoint{3.742303in}{5.582496in}}{\pgfqpoint{3.744617in}{5.576910in}}{\pgfqpoint{3.748735in}{5.572792in}}%
\pgfpathcurveto{\pgfqpoint{3.752853in}{5.568674in}}{\pgfqpoint{3.758439in}{5.566360in}}{\pgfqpoint{3.764263in}{5.566360in}}%
\pgfpathlineto{\pgfqpoint{3.764263in}{5.566360in}}%
\pgfpathclose%
\pgfusepath{stroke,fill}%
\end{pgfscope}%
\begin{pgfscope}%
\pgfpathrectangle{\pgfqpoint{1.414736in}{0.548769in}}{\pgfqpoint{3.497747in}{7.253186in}}%
\pgfusepath{clip}%
\pgfsetbuttcap%
\pgfsetroundjoin%
\definecolor{currentfill}{rgb}{0.146180,0.515413,0.556823}%
\pgfsetfillcolor{currentfill}%
\pgfsetlinewidth{1.003750pt}%
\definecolor{currentstroke}{rgb}{0.146180,0.515413,0.556823}%
\pgfsetstrokecolor{currentstroke}%
\pgfsetdash{}{0pt}%
\pgfpathmoveto{\pgfqpoint{3.787747in}{6.037346in}}%
\pgfpathcurveto{\pgfqpoint{3.793571in}{6.037346in}}{\pgfqpoint{3.799157in}{6.039660in}}{\pgfqpoint{3.803276in}{6.043778in}}%
\pgfpathcurveto{\pgfqpoint{3.807394in}{6.047896in}}{\pgfqpoint{3.809708in}{6.053482in}}{\pgfqpoint{3.809708in}{6.059306in}}%
\pgfpathcurveto{\pgfqpoint{3.809708in}{6.065130in}}{\pgfqpoint{3.807394in}{6.070716in}}{\pgfqpoint{3.803276in}{6.074835in}}%
\pgfpathcurveto{\pgfqpoint{3.799157in}{6.078953in}}{\pgfqpoint{3.793571in}{6.081267in}}{\pgfqpoint{3.787747in}{6.081267in}}%
\pgfpathcurveto{\pgfqpoint{3.781923in}{6.081267in}}{\pgfqpoint{3.776337in}{6.078953in}}{\pgfqpoint{3.772219in}{6.074835in}}%
\pgfpathcurveto{\pgfqpoint{3.768101in}{6.070716in}}{\pgfqpoint{3.765787in}{6.065130in}}{\pgfqpoint{3.765787in}{6.059306in}}%
\pgfpathcurveto{\pgfqpoint{3.765787in}{6.053482in}}{\pgfqpoint{3.768101in}{6.047896in}}{\pgfqpoint{3.772219in}{6.043778in}}%
\pgfpathcurveto{\pgfqpoint{3.776337in}{6.039660in}}{\pgfqpoint{3.781923in}{6.037346in}}{\pgfqpoint{3.787747in}{6.037346in}}%
\pgfpathlineto{\pgfqpoint{3.787747in}{6.037346in}}%
\pgfpathclose%
\pgfusepath{stroke,fill}%
\end{pgfscope}%
\begin{pgfscope}%
\pgfpathrectangle{\pgfqpoint{1.414736in}{0.548769in}}{\pgfqpoint{3.497747in}{7.253186in}}%
\pgfusepath{clip}%
\pgfsetbuttcap%
\pgfsetroundjoin%
\definecolor{currentfill}{rgb}{0.993248,0.906157,0.143936}%
\pgfsetfillcolor{currentfill}%
\pgfsetlinewidth{1.003750pt}%
\definecolor{currentstroke}{rgb}{0.993248,0.906157,0.143936}%
\pgfsetstrokecolor{currentstroke}%
\pgfsetdash{}{0pt}%
\pgfpathmoveto{\pgfqpoint{3.835927in}{6.508332in}}%
\pgfpathcurveto{\pgfqpoint{3.841751in}{6.508332in}}{\pgfqpoint{3.847338in}{6.510646in}}{\pgfqpoint{3.851456in}{6.514764in}}%
\pgfpathcurveto{\pgfqpoint{3.855574in}{6.518882in}}{\pgfqpoint{3.857888in}{6.524468in}}{\pgfqpoint{3.857888in}{6.530292in}}%
\pgfpathcurveto{\pgfqpoint{3.857888in}{6.536116in}}{\pgfqpoint{3.855574in}{6.541702in}}{\pgfqpoint{3.851456in}{6.545821in}}%
\pgfpathcurveto{\pgfqpoint{3.847338in}{6.549939in}}{\pgfqpoint{3.841751in}{6.552253in}}{\pgfqpoint{3.835927in}{6.552253in}}%
\pgfpathcurveto{\pgfqpoint{3.830103in}{6.552253in}}{\pgfqpoint{3.824517in}{6.549939in}}{\pgfqpoint{3.820399in}{6.545821in}}%
\pgfpathcurveto{\pgfqpoint{3.816281in}{6.541702in}}{\pgfqpoint{3.813967in}{6.536116in}}{\pgfqpoint{3.813967in}{6.530292in}}%
\pgfpathcurveto{\pgfqpoint{3.813967in}{6.524468in}}{\pgfqpoint{3.816281in}{6.518882in}}{\pgfqpoint{3.820399in}{6.514764in}}%
\pgfpathcurveto{\pgfqpoint{3.824517in}{6.510646in}}{\pgfqpoint{3.830103in}{6.508332in}}{\pgfqpoint{3.835927in}{6.508332in}}%
\pgfpathlineto{\pgfqpoint{3.835927in}{6.508332in}}%
\pgfpathclose%
\pgfusepath{stroke,fill}%
\end{pgfscope}%
\begin{pgfscope}%
\pgfpathrectangle{\pgfqpoint{1.414736in}{0.548769in}}{\pgfqpoint{3.497747in}{7.253186in}}%
\pgfusepath{clip}%
\pgfsetbuttcap%
\pgfsetroundjoin%
\definecolor{currentfill}{rgb}{0.814576,0.883393,0.110347}%
\pgfsetfillcolor{currentfill}%
\pgfsetlinewidth{1.003750pt}%
\definecolor{currentstroke}{rgb}{0.814576,0.883393,0.110347}%
\pgfsetstrokecolor{currentstroke}%
\pgfsetdash{}{0pt}%
\pgfpathmoveto{\pgfqpoint{3.121100in}{6.979318in}}%
\pgfpathcurveto{\pgfqpoint{3.126923in}{6.979318in}}{\pgfqpoint{3.132510in}{6.981632in}}{\pgfqpoint{3.136628in}{6.985750in}}%
\pgfpathcurveto{\pgfqpoint{3.140746in}{6.989868in}}{\pgfqpoint{3.143060in}{6.995455in}}{\pgfqpoint{3.143060in}{7.001278in}}%
\pgfpathcurveto{\pgfqpoint{3.143060in}{7.007102in}}{\pgfqpoint{3.140746in}{7.012689in}}{\pgfqpoint{3.136628in}{7.016807in}}%
\pgfpathcurveto{\pgfqpoint{3.132510in}{7.020925in}}{\pgfqpoint{3.126923in}{7.023239in}}{\pgfqpoint{3.121100in}{7.023239in}}%
\pgfpathcurveto{\pgfqpoint{3.115276in}{7.023239in}}{\pgfqpoint{3.109689in}{7.020925in}}{\pgfqpoint{3.105571in}{7.016807in}}%
\pgfpathcurveto{\pgfqpoint{3.101453in}{7.012689in}}{\pgfqpoint{3.099139in}{7.007102in}}{\pgfqpoint{3.099139in}{7.001278in}}%
\pgfpathcurveto{\pgfqpoint{3.099139in}{6.995455in}}{\pgfqpoint{3.101453in}{6.989868in}}{\pgfqpoint{3.105571in}{6.985750in}}%
\pgfpathcurveto{\pgfqpoint{3.109689in}{6.981632in}}{\pgfqpoint{3.115276in}{6.979318in}}{\pgfqpoint{3.121100in}{6.979318in}}%
\pgfpathlineto{\pgfqpoint{3.121100in}{6.979318in}}%
\pgfpathclose%
\pgfusepath{stroke,fill}%
\end{pgfscope}%
\begin{pgfscope}%
\pgfpathrectangle{\pgfqpoint{1.414736in}{0.548769in}}{\pgfqpoint{3.497747in}{7.253186in}}%
\pgfusepath{clip}%
\pgfsetbuttcap%
\pgfsetroundjoin%
\definecolor{currentfill}{rgb}{0.845561,0.887322,0.099702}%
\pgfsetfillcolor{currentfill}%
\pgfsetlinewidth{1.003750pt}%
\definecolor{currentstroke}{rgb}{0.845561,0.887322,0.099702}%
\pgfsetstrokecolor{currentstroke}%
\pgfsetdash{}{0pt}%
\pgfpathmoveto{\pgfqpoint{4.023066in}{7.450304in}}%
\pgfpathcurveto{\pgfqpoint{4.028890in}{7.450304in}}{\pgfqpoint{4.034476in}{7.452618in}}{\pgfqpoint{4.038594in}{7.456736in}}%
\pgfpathcurveto{\pgfqpoint{4.042712in}{7.460854in}}{\pgfqpoint{4.045026in}{7.466441in}}{\pgfqpoint{4.045026in}{7.472265in}}%
\pgfpathcurveto{\pgfqpoint{4.045026in}{7.478089in}}{\pgfqpoint{4.042712in}{7.483675in}}{\pgfqpoint{4.038594in}{7.487793in}}%
\pgfpathcurveto{\pgfqpoint{4.034476in}{7.491911in}}{\pgfqpoint{4.028890in}{7.494225in}}{\pgfqpoint{4.023066in}{7.494225in}}%
\pgfpathcurveto{\pgfqpoint{4.017242in}{7.494225in}}{\pgfqpoint{4.011655in}{7.491911in}}{\pgfqpoint{4.007537in}{7.487793in}}%
\pgfpathcurveto{\pgfqpoint{4.003419in}{7.483675in}}{\pgfqpoint{4.001105in}{7.478089in}}{\pgfqpoint{4.001105in}{7.472265in}}%
\pgfpathcurveto{\pgfqpoint{4.001105in}{7.466441in}}{\pgfqpoint{4.003419in}{7.460854in}}{\pgfqpoint{4.007537in}{7.456736in}}%
\pgfpathcurveto{\pgfqpoint{4.011655in}{7.452618in}}{\pgfqpoint{4.017242in}{7.450304in}}{\pgfqpoint{4.023066in}{7.450304in}}%
\pgfpathlineto{\pgfqpoint{4.023066in}{7.450304in}}%
\pgfpathclose%
\pgfusepath{stroke,fill}%
\end{pgfscope}%
\begin{pgfscope}%
\pgfpathrectangle{\pgfqpoint{1.414736in}{0.548769in}}{\pgfqpoint{3.497747in}{7.253186in}}%
\pgfusepath{clip}%
\pgfsetbuttcap%
\pgfsetroundjoin%
\definecolor{currentfill}{rgb}{0.699415,0.867117,0.175971}%
\pgfsetfillcolor{currentfill}%
\pgfsetlinewidth{1.003750pt}%
\definecolor{currentstroke}{rgb}{0.699415,0.867117,0.175971}%
\pgfsetstrokecolor{currentstroke}%
\pgfsetdash{}{0pt}%
\pgfpathmoveto{\pgfqpoint{3.643165in}{0.856499in}}%
\pgfpathcurveto{\pgfqpoint{3.648989in}{0.856499in}}{\pgfqpoint{3.654575in}{0.858813in}}{\pgfqpoint{3.658693in}{0.862931in}}%
\pgfpathcurveto{\pgfqpoint{3.662811in}{0.867049in}}{\pgfqpoint{3.665125in}{0.872635in}}{\pgfqpoint{3.665125in}{0.878459in}}%
\pgfpathcurveto{\pgfqpoint{3.665125in}{0.884283in}}{\pgfqpoint{3.662811in}{0.889869in}}{\pgfqpoint{3.658693in}{0.893987in}}%
\pgfpathcurveto{\pgfqpoint{3.654575in}{0.898105in}}{\pgfqpoint{3.648989in}{0.900419in}}{\pgfqpoint{3.643165in}{0.900419in}}%
\pgfpathcurveto{\pgfqpoint{3.637341in}{0.900419in}}{\pgfqpoint{3.631754in}{0.898105in}}{\pgfqpoint{3.627636in}{0.893987in}}%
\pgfpathcurveto{\pgfqpoint{3.623518in}{0.889869in}}{\pgfqpoint{3.621204in}{0.884283in}}{\pgfqpoint{3.621204in}{0.878459in}}%
\pgfpathcurveto{\pgfqpoint{3.621204in}{0.872635in}}{\pgfqpoint{3.623518in}{0.867049in}}{\pgfqpoint{3.627636in}{0.862931in}}%
\pgfpathcurveto{\pgfqpoint{3.631754in}{0.858813in}}{\pgfqpoint{3.637341in}{0.856499in}}{\pgfqpoint{3.643165in}{0.856499in}}%
\pgfpathlineto{\pgfqpoint{3.643165in}{0.856499in}}%
\pgfpathclose%
\pgfusepath{stroke,fill}%
\end{pgfscope}%
\begin{pgfscope}%
\pgfpathrectangle{\pgfqpoint{1.414736in}{0.548769in}}{\pgfqpoint{3.497747in}{7.253186in}}%
\pgfusepath{clip}%
\pgfsetbuttcap%
\pgfsetroundjoin%
\definecolor{currentfill}{rgb}{0.595839,0.848717,0.243329}%
\pgfsetfillcolor{currentfill}%
\pgfsetlinewidth{1.003750pt}%
\definecolor{currentstroke}{rgb}{0.595839,0.848717,0.243329}%
\pgfsetstrokecolor{currentstroke}%
\pgfsetdash{}{0pt}%
\pgfpathmoveto{\pgfqpoint{3.438806in}{1.327485in}}%
\pgfpathcurveto{\pgfqpoint{3.444630in}{1.327485in}}{\pgfqpoint{3.450216in}{1.329799in}}{\pgfqpoint{3.454334in}{1.333917in}}%
\pgfpathcurveto{\pgfqpoint{3.458452in}{1.338035in}}{\pgfqpoint{3.460766in}{1.343621in}}{\pgfqpoint{3.460766in}{1.349445in}}%
\pgfpathcurveto{\pgfqpoint{3.460766in}{1.355269in}}{\pgfqpoint{3.458452in}{1.360855in}}{\pgfqpoint{3.454334in}{1.364973in}}%
\pgfpathcurveto{\pgfqpoint{3.450216in}{1.369092in}}{\pgfqpoint{3.444630in}{1.371405in}}{\pgfqpoint{3.438806in}{1.371405in}}%
\pgfpathcurveto{\pgfqpoint{3.432982in}{1.371405in}}{\pgfqpoint{3.427396in}{1.369092in}}{\pgfqpoint{3.423277in}{1.364973in}}%
\pgfpathcurveto{\pgfqpoint{3.419159in}{1.360855in}}{\pgfqpoint{3.416845in}{1.355269in}}{\pgfqpoint{3.416845in}{1.349445in}}%
\pgfpathcurveto{\pgfqpoint{3.416845in}{1.343621in}}{\pgfqpoint{3.419159in}{1.338035in}}{\pgfqpoint{3.423277in}{1.333917in}}%
\pgfpathcurveto{\pgfqpoint{3.427396in}{1.329799in}}{\pgfqpoint{3.432982in}{1.327485in}}{\pgfqpoint{3.438806in}{1.327485in}}%
\pgfpathlineto{\pgfqpoint{3.438806in}{1.327485in}}%
\pgfpathclose%
\pgfusepath{stroke,fill}%
\end{pgfscope}%
\begin{pgfscope}%
\pgfpathrectangle{\pgfqpoint{1.414736in}{0.548769in}}{\pgfqpoint{3.497747in}{7.253186in}}%
\pgfusepath{clip}%
\pgfsetbuttcap%
\pgfsetroundjoin%
\definecolor{currentfill}{rgb}{0.699415,0.867117,0.175971}%
\pgfsetfillcolor{currentfill}%
\pgfsetlinewidth{1.003750pt}%
\definecolor{currentstroke}{rgb}{0.699415,0.867117,0.175971}%
\pgfsetstrokecolor{currentstroke}%
\pgfsetdash{}{0pt}%
\pgfpathmoveto{\pgfqpoint{3.744897in}{1.798471in}}%
\pgfpathcurveto{\pgfqpoint{3.750721in}{1.798471in}}{\pgfqpoint{3.756307in}{1.800785in}}{\pgfqpoint{3.760425in}{1.804903in}}%
\pgfpathcurveto{\pgfqpoint{3.764543in}{1.809021in}}{\pgfqpoint{3.766857in}{1.814607in}}{\pgfqpoint{3.766857in}{1.820431in}}%
\pgfpathcurveto{\pgfqpoint{3.766857in}{1.826255in}}{\pgfqpoint{3.764543in}{1.831841in}}{\pgfqpoint{3.760425in}{1.835960in}}%
\pgfpathcurveto{\pgfqpoint{3.756307in}{1.840078in}}{\pgfqpoint{3.750721in}{1.842392in}}{\pgfqpoint{3.744897in}{1.842392in}}%
\pgfpathcurveto{\pgfqpoint{3.739073in}{1.842392in}}{\pgfqpoint{3.733487in}{1.840078in}}{\pgfqpoint{3.729369in}{1.835960in}}%
\pgfpathcurveto{\pgfqpoint{3.725250in}{1.831841in}}{\pgfqpoint{3.722937in}{1.826255in}}{\pgfqpoint{3.722937in}{1.820431in}}%
\pgfpathcurveto{\pgfqpoint{3.722937in}{1.814607in}}{\pgfqpoint{3.725250in}{1.809021in}}{\pgfqpoint{3.729369in}{1.804903in}}%
\pgfpathcurveto{\pgfqpoint{3.733487in}{1.800785in}}{\pgfqpoint{3.739073in}{1.798471in}}{\pgfqpoint{3.744897in}{1.798471in}}%
\pgfpathlineto{\pgfqpoint{3.744897in}{1.798471in}}%
\pgfpathclose%
\pgfusepath{stroke,fill}%
\end{pgfscope}%
\begin{pgfscope}%
\pgfpathrectangle{\pgfqpoint{1.414736in}{0.548769in}}{\pgfqpoint{3.497747in}{7.253186in}}%
\pgfusepath{clip}%
\pgfsetbuttcap%
\pgfsetroundjoin%
\definecolor{currentfill}{rgb}{0.993248,0.906157,0.143936}%
\pgfsetfillcolor{currentfill}%
\pgfsetlinewidth{1.003750pt}%
\definecolor{currentstroke}{rgb}{0.993248,0.906157,0.143936}%
\pgfsetstrokecolor{currentstroke}%
\pgfsetdash{}{0pt}%
\pgfpathmoveto{\pgfqpoint{3.824624in}{2.269457in}}%
\pgfpathcurveto{\pgfqpoint{3.830448in}{2.269457in}}{\pgfqpoint{3.836034in}{2.271771in}}{\pgfqpoint{3.840152in}{2.275889in}}%
\pgfpathcurveto{\pgfqpoint{3.844270in}{2.280007in}}{\pgfqpoint{3.846584in}{2.285593in}}{\pgfqpoint{3.846584in}{2.291417in}}%
\pgfpathcurveto{\pgfqpoint{3.846584in}{2.297241in}}{\pgfqpoint{3.844270in}{2.302828in}}{\pgfqpoint{3.840152in}{2.306946in}}%
\pgfpathcurveto{\pgfqpoint{3.836034in}{2.311064in}}{\pgfqpoint{3.830448in}{2.313378in}}{\pgfqpoint{3.824624in}{2.313378in}}%
\pgfpathcurveto{\pgfqpoint{3.818800in}{2.313378in}}{\pgfqpoint{3.813213in}{2.311064in}}{\pgfqpoint{3.809095in}{2.306946in}}%
\pgfpathcurveto{\pgfqpoint{3.804977in}{2.302828in}}{\pgfqpoint{3.802663in}{2.297241in}}{\pgfqpoint{3.802663in}{2.291417in}}%
\pgfpathcurveto{\pgfqpoint{3.802663in}{2.285593in}}{\pgfqpoint{3.804977in}{2.280007in}}{\pgfqpoint{3.809095in}{2.275889in}}%
\pgfpathcurveto{\pgfqpoint{3.813213in}{2.271771in}}{\pgfqpoint{3.818800in}{2.269457in}}{\pgfqpoint{3.824624in}{2.269457in}}%
\pgfpathlineto{\pgfqpoint{3.824624in}{2.269457in}}%
\pgfpathclose%
\pgfusepath{stroke,fill}%
\end{pgfscope}%
\begin{pgfscope}%
\pgfpathrectangle{\pgfqpoint{1.414736in}{0.548769in}}{\pgfqpoint{3.497747in}{7.253186in}}%
\pgfusepath{clip}%
\pgfsetbuttcap%
\pgfsetroundjoin%
\definecolor{currentfill}{rgb}{0.239374,0.735588,0.455688}%
\pgfsetfillcolor{currentfill}%
\pgfsetlinewidth{1.003750pt}%
\definecolor{currentstroke}{rgb}{0.239374,0.735588,0.455688}%
\pgfsetstrokecolor{currentstroke}%
\pgfsetdash{}{0pt}%
\pgfpathmoveto{\pgfqpoint{3.810814in}{2.740443in}}%
\pgfpathcurveto{\pgfqpoint{3.816638in}{2.740443in}}{\pgfqpoint{3.822224in}{2.742757in}}{\pgfqpoint{3.826342in}{2.746875in}}%
\pgfpathcurveto{\pgfqpoint{3.830460in}{2.750993in}}{\pgfqpoint{3.832774in}{2.756580in}}{\pgfqpoint{3.832774in}{2.762404in}}%
\pgfpathcurveto{\pgfqpoint{3.832774in}{2.768227in}}{\pgfqpoint{3.830460in}{2.773814in}}{\pgfqpoint{3.826342in}{2.777932in}}%
\pgfpathcurveto{\pgfqpoint{3.822224in}{2.782050in}}{\pgfqpoint{3.816638in}{2.784364in}}{\pgfqpoint{3.810814in}{2.784364in}}%
\pgfpathcurveto{\pgfqpoint{3.804990in}{2.784364in}}{\pgfqpoint{3.799404in}{2.782050in}}{\pgfqpoint{3.795286in}{2.777932in}}%
\pgfpathcurveto{\pgfqpoint{3.791168in}{2.773814in}}{\pgfqpoint{3.788854in}{2.768227in}}{\pgfqpoint{3.788854in}{2.762404in}}%
\pgfpathcurveto{\pgfqpoint{3.788854in}{2.756580in}}{\pgfqpoint{3.791168in}{2.750993in}}{\pgfqpoint{3.795286in}{2.746875in}}%
\pgfpathcurveto{\pgfqpoint{3.799404in}{2.742757in}}{\pgfqpoint{3.804990in}{2.740443in}}{\pgfqpoint{3.810814in}{2.740443in}}%
\pgfpathlineto{\pgfqpoint{3.810814in}{2.740443in}}%
\pgfpathclose%
\pgfusepath{stroke,fill}%
\end{pgfscope}%
\begin{pgfscope}%
\pgfpathrectangle{\pgfqpoint{1.414736in}{0.548769in}}{\pgfqpoint{3.497747in}{7.253186in}}%
\pgfusepath{clip}%
\pgfsetbuttcap%
\pgfsetroundjoin%
\definecolor{currentfill}{rgb}{0.993248,0.906157,0.143936}%
\pgfsetfillcolor{currentfill}%
\pgfsetlinewidth{1.003750pt}%
\definecolor{currentstroke}{rgb}{0.993248,0.906157,0.143936}%
\pgfsetstrokecolor{currentstroke}%
\pgfsetdash{}{0pt}%
\pgfpathmoveto{\pgfqpoint{3.639384in}{3.211429in}}%
\pgfpathcurveto{\pgfqpoint{3.645208in}{3.211429in}}{\pgfqpoint{3.650794in}{3.213743in}}{\pgfqpoint{3.654913in}{3.217861in}}%
\pgfpathcurveto{\pgfqpoint{3.659031in}{3.221980in}}{\pgfqpoint{3.661345in}{3.227566in}}{\pgfqpoint{3.661345in}{3.233390in}}%
\pgfpathcurveto{\pgfqpoint{3.661345in}{3.239214in}}{\pgfqpoint{3.659031in}{3.244800in}}{\pgfqpoint{3.654913in}{3.248918in}}%
\pgfpathcurveto{\pgfqpoint{3.650794in}{3.253036in}}{\pgfqpoint{3.645208in}{3.255350in}}{\pgfqpoint{3.639384in}{3.255350in}}%
\pgfpathcurveto{\pgfqpoint{3.633560in}{3.255350in}}{\pgfqpoint{3.627974in}{3.253036in}}{\pgfqpoint{3.623856in}{3.248918in}}%
\pgfpathcurveto{\pgfqpoint{3.619738in}{3.244800in}}{\pgfqpoint{3.617424in}{3.239214in}}{\pgfqpoint{3.617424in}{3.233390in}}%
\pgfpathcurveto{\pgfqpoint{3.617424in}{3.227566in}}{\pgfqpoint{3.619738in}{3.221980in}}{\pgfqpoint{3.623856in}{3.217861in}}%
\pgfpathcurveto{\pgfqpoint{3.627974in}{3.213743in}}{\pgfqpoint{3.633560in}{3.211429in}}{\pgfqpoint{3.639384in}{3.211429in}}%
\pgfpathlineto{\pgfqpoint{3.639384in}{3.211429in}}%
\pgfpathclose%
\pgfusepath{stroke,fill}%
\end{pgfscope}%
\begin{pgfscope}%
\pgfpathrectangle{\pgfqpoint{1.414736in}{0.548769in}}{\pgfqpoint{3.497747in}{7.253186in}}%
\pgfusepath{clip}%
\pgfsetbuttcap%
\pgfsetroundjoin%
\definecolor{currentfill}{rgb}{0.855810,0.888601,0.097452}%
\pgfsetfillcolor{currentfill}%
\pgfsetlinewidth{1.003750pt}%
\definecolor{currentstroke}{rgb}{0.855810,0.888601,0.097452}%
\pgfsetstrokecolor{currentstroke}%
\pgfsetdash{}{0pt}%
\pgfpathmoveto{\pgfqpoint{3.585882in}{3.682415in}}%
\pgfpathcurveto{\pgfqpoint{3.591706in}{3.682415in}}{\pgfqpoint{3.597292in}{3.684729in}}{\pgfqpoint{3.601410in}{3.688847in}}%
\pgfpathcurveto{\pgfqpoint{3.605528in}{3.692966in}}{\pgfqpoint{3.607842in}{3.698552in}}{\pgfqpoint{3.607842in}{3.704376in}}%
\pgfpathcurveto{\pgfqpoint{3.607842in}{3.710200in}}{\pgfqpoint{3.605528in}{3.715786in}}{\pgfqpoint{3.601410in}{3.719904in}}%
\pgfpathcurveto{\pgfqpoint{3.597292in}{3.724022in}}{\pgfqpoint{3.591706in}{3.726336in}}{\pgfqpoint{3.585882in}{3.726336in}}%
\pgfpathcurveto{\pgfqpoint{3.580058in}{3.726336in}}{\pgfqpoint{3.574472in}{3.724022in}}{\pgfqpoint{3.570354in}{3.719904in}}%
\pgfpathcurveto{\pgfqpoint{3.566236in}{3.715786in}}{\pgfqpoint{3.563922in}{3.710200in}}{\pgfqpoint{3.563922in}{3.704376in}}%
\pgfpathcurveto{\pgfqpoint{3.563922in}{3.698552in}}{\pgfqpoint{3.566236in}{3.692966in}}{\pgfqpoint{3.570354in}{3.688847in}}%
\pgfpathcurveto{\pgfqpoint{3.574472in}{3.684729in}}{\pgfqpoint{3.580058in}{3.682415in}}{\pgfqpoint{3.585882in}{3.682415in}}%
\pgfpathlineto{\pgfqpoint{3.585882in}{3.682415in}}%
\pgfpathclose%
\pgfusepath{stroke,fill}%
\end{pgfscope}%
\begin{pgfscope}%
\pgfpathrectangle{\pgfqpoint{1.414736in}{0.548769in}}{\pgfqpoint{3.497747in}{7.253186in}}%
\pgfusepath{clip}%
\pgfsetbuttcap%
\pgfsetroundjoin%
\definecolor{currentfill}{rgb}{0.412913,0.803041,0.357269}%
\pgfsetfillcolor{currentfill}%
\pgfsetlinewidth{1.003750pt}%
\definecolor{currentstroke}{rgb}{0.412913,0.803041,0.357269}%
\pgfsetstrokecolor{currentstroke}%
\pgfsetdash{}{0pt}%
\pgfpathmoveto{\pgfqpoint{3.266470in}{4.153402in}}%
\pgfpathcurveto{\pgfqpoint{3.272294in}{4.153402in}}{\pgfqpoint{3.277880in}{4.155715in}}{\pgfqpoint{3.281999in}{4.159834in}}%
\pgfpathcurveto{\pgfqpoint{3.286117in}{4.163952in}}{\pgfqpoint{3.288431in}{4.169538in}}{\pgfqpoint{3.288431in}{4.175362in}}%
\pgfpathcurveto{\pgfqpoint{3.288431in}{4.181186in}}{\pgfqpoint{3.286117in}{4.186772in}}{\pgfqpoint{3.281999in}{4.190890in}}%
\pgfpathcurveto{\pgfqpoint{3.277880in}{4.195008in}}{\pgfqpoint{3.272294in}{4.197322in}}{\pgfqpoint{3.266470in}{4.197322in}}%
\pgfpathcurveto{\pgfqpoint{3.260646in}{4.197322in}}{\pgfqpoint{3.255060in}{4.195008in}}{\pgfqpoint{3.250942in}{4.190890in}}%
\pgfpathcurveto{\pgfqpoint{3.246824in}{4.186772in}}{\pgfqpoint{3.244510in}{4.181186in}}{\pgfqpoint{3.244510in}{4.175362in}}%
\pgfpathcurveto{\pgfqpoint{3.244510in}{4.169538in}}{\pgfqpoint{3.246824in}{4.163952in}}{\pgfqpoint{3.250942in}{4.159834in}}%
\pgfpathcurveto{\pgfqpoint{3.255060in}{4.155715in}}{\pgfqpoint{3.260646in}{4.153402in}}{\pgfqpoint{3.266470in}{4.153402in}}%
\pgfpathlineto{\pgfqpoint{3.266470in}{4.153402in}}%
\pgfpathclose%
\pgfusepath{stroke,fill}%
\end{pgfscope}%
\begin{pgfscope}%
\pgfpathrectangle{\pgfqpoint{1.414736in}{0.548769in}}{\pgfqpoint{3.497747in}{7.253186in}}%
\pgfusepath{clip}%
\pgfsetbuttcap%
\pgfsetroundjoin%
\definecolor{currentfill}{rgb}{0.751884,0.874951,0.143228}%
\pgfsetfillcolor{currentfill}%
\pgfsetlinewidth{1.003750pt}%
\definecolor{currentstroke}{rgb}{0.751884,0.874951,0.143228}%
\pgfsetstrokecolor{currentstroke}%
\pgfsetdash{}{0pt}%
\pgfpathmoveto{\pgfqpoint{3.681680in}{4.624388in}}%
\pgfpathcurveto{\pgfqpoint{3.687504in}{4.624388in}}{\pgfqpoint{3.693090in}{4.626702in}}{\pgfqpoint{3.697208in}{4.630820in}}%
\pgfpathcurveto{\pgfqpoint{3.701326in}{4.634938in}}{\pgfqpoint{3.703640in}{4.640524in}}{\pgfqpoint{3.703640in}{4.646348in}}%
\pgfpathcurveto{\pgfqpoint{3.703640in}{4.652172in}}{\pgfqpoint{3.701326in}{4.657758in}}{\pgfqpoint{3.697208in}{4.661876in}}%
\pgfpathcurveto{\pgfqpoint{3.693090in}{4.665994in}}{\pgfqpoint{3.687504in}{4.668308in}}{\pgfqpoint{3.681680in}{4.668308in}}%
\pgfpathcurveto{\pgfqpoint{3.675856in}{4.668308in}}{\pgfqpoint{3.670270in}{4.665994in}}{\pgfqpoint{3.666152in}{4.661876in}}%
\pgfpathcurveto{\pgfqpoint{3.662034in}{4.657758in}}{\pgfqpoint{3.659720in}{4.652172in}}{\pgfqpoint{3.659720in}{4.646348in}}%
\pgfpathcurveto{\pgfqpoint{3.659720in}{4.640524in}}{\pgfqpoint{3.662034in}{4.634938in}}{\pgfqpoint{3.666152in}{4.630820in}}%
\pgfpathcurveto{\pgfqpoint{3.670270in}{4.626702in}}{\pgfqpoint{3.675856in}{4.624388in}}{\pgfqpoint{3.681680in}{4.624388in}}%
\pgfpathlineto{\pgfqpoint{3.681680in}{4.624388in}}%
\pgfpathclose%
\pgfusepath{stroke,fill}%
\end{pgfscope}%
\begin{pgfscope}%
\pgfpathrectangle{\pgfqpoint{1.414736in}{0.548769in}}{\pgfqpoint{3.497747in}{7.253186in}}%
\pgfusepath{clip}%
\pgfsetbuttcap%
\pgfsetroundjoin%
\definecolor{currentfill}{rgb}{0.143303,0.669459,0.511215}%
\pgfsetfillcolor{currentfill}%
\pgfsetlinewidth{1.003750pt}%
\definecolor{currentstroke}{rgb}{0.143303,0.669459,0.511215}%
\pgfsetstrokecolor{currentstroke}%
\pgfsetdash{}{0pt}%
\pgfpathmoveto{\pgfqpoint{4.077384in}{5.095374in}}%
\pgfpathcurveto{\pgfqpoint{4.083208in}{5.095374in}}{\pgfqpoint{4.088794in}{5.097688in}}{\pgfqpoint{4.092912in}{5.101806in}}%
\pgfpathcurveto{\pgfqpoint{4.097030in}{5.105924in}}{\pgfqpoint{4.099344in}{5.111510in}}{\pgfqpoint{4.099344in}{5.117334in}}%
\pgfpathcurveto{\pgfqpoint{4.099344in}{5.123158in}}{\pgfqpoint{4.097030in}{5.128744in}}{\pgfqpoint{4.092912in}{5.132862in}}%
\pgfpathcurveto{\pgfqpoint{4.088794in}{5.136980in}}{\pgfqpoint{4.083208in}{5.139294in}}{\pgfqpoint{4.077384in}{5.139294in}}%
\pgfpathcurveto{\pgfqpoint{4.071560in}{5.139294in}}{\pgfqpoint{4.065974in}{5.136980in}}{\pgfqpoint{4.061856in}{5.132862in}}%
\pgfpathcurveto{\pgfqpoint{4.057738in}{5.128744in}}{\pgfqpoint{4.055424in}{5.123158in}}{\pgfqpoint{4.055424in}{5.117334in}}%
\pgfpathcurveto{\pgfqpoint{4.055424in}{5.111510in}}{\pgfqpoint{4.057738in}{5.105924in}}{\pgfqpoint{4.061856in}{5.101806in}}%
\pgfpathcurveto{\pgfqpoint{4.065974in}{5.097688in}}{\pgfqpoint{4.071560in}{5.095374in}}{\pgfqpoint{4.077384in}{5.095374in}}%
\pgfpathlineto{\pgfqpoint{4.077384in}{5.095374in}}%
\pgfpathclose%
\pgfusepath{stroke,fill}%
\end{pgfscope}%
\begin{pgfscope}%
\pgfpathrectangle{\pgfqpoint{1.414736in}{0.548769in}}{\pgfqpoint{3.497747in}{7.253186in}}%
\pgfusepath{clip}%
\pgfsetbuttcap%
\pgfsetroundjoin%
\definecolor{currentfill}{rgb}{0.121831,0.589055,0.545623}%
\pgfsetfillcolor{currentfill}%
\pgfsetlinewidth{1.003750pt}%
\definecolor{currentstroke}{rgb}{0.121831,0.589055,0.545623}%
\pgfsetstrokecolor{currentstroke}%
\pgfsetdash{}{0pt}%
\pgfpathmoveto{\pgfqpoint{3.152198in}{5.566360in}}%
\pgfpathcurveto{\pgfqpoint{3.158022in}{5.566360in}}{\pgfqpoint{3.163608in}{5.568674in}}{\pgfqpoint{3.167726in}{5.572792in}}%
\pgfpathcurveto{\pgfqpoint{3.171844in}{5.576910in}}{\pgfqpoint{3.174158in}{5.582496in}}{\pgfqpoint{3.174158in}{5.588320in}}%
\pgfpathcurveto{\pgfqpoint{3.174158in}{5.594144in}}{\pgfqpoint{3.171844in}{5.599730in}}{\pgfqpoint{3.167726in}{5.603848in}}%
\pgfpathcurveto{\pgfqpoint{3.163608in}{5.607967in}}{\pgfqpoint{3.158022in}{5.610280in}}{\pgfqpoint{3.152198in}{5.610280in}}%
\pgfpathcurveto{\pgfqpoint{3.146374in}{5.610280in}}{\pgfqpoint{3.140788in}{5.607967in}}{\pgfqpoint{3.136669in}{5.603848in}}%
\pgfpathcurveto{\pgfqpoint{3.132551in}{5.599730in}}{\pgfqpoint{3.130237in}{5.594144in}}{\pgfqpoint{3.130237in}{5.588320in}}%
\pgfpathcurveto{\pgfqpoint{3.130237in}{5.582496in}}{\pgfqpoint{3.132551in}{5.576910in}}{\pgfqpoint{3.136669in}{5.572792in}}%
\pgfpathcurveto{\pgfqpoint{3.140788in}{5.568674in}}{\pgfqpoint{3.146374in}{5.566360in}}{\pgfqpoint{3.152198in}{5.566360in}}%
\pgfpathlineto{\pgfqpoint{3.152198in}{5.566360in}}%
\pgfpathclose%
\pgfusepath{stroke,fill}%
\end{pgfscope}%
\begin{pgfscope}%
\pgfpathrectangle{\pgfqpoint{1.414736in}{0.548769in}}{\pgfqpoint{3.497747in}{7.253186in}}%
\pgfusepath{clip}%
\pgfsetbuttcap%
\pgfsetroundjoin%
\definecolor{currentfill}{rgb}{0.157851,0.683765,0.501686}%
\pgfsetfillcolor{currentfill}%
\pgfsetlinewidth{1.003750pt}%
\definecolor{currentstroke}{rgb}{0.157851,0.683765,0.501686}%
\pgfsetstrokecolor{currentstroke}%
\pgfsetdash{}{0pt}%
\pgfpathmoveto{\pgfqpoint{3.853246in}{6.037346in}}%
\pgfpathcurveto{\pgfqpoint{3.859070in}{6.037346in}}{\pgfqpoint{3.864657in}{6.039660in}}{\pgfqpoint{3.868775in}{6.043778in}}%
\pgfpathcurveto{\pgfqpoint{3.872893in}{6.047896in}}{\pgfqpoint{3.875207in}{6.053482in}}{\pgfqpoint{3.875207in}{6.059306in}}%
\pgfpathcurveto{\pgfqpoint{3.875207in}{6.065130in}}{\pgfqpoint{3.872893in}{6.070716in}}{\pgfqpoint{3.868775in}{6.074835in}}%
\pgfpathcurveto{\pgfqpoint{3.864657in}{6.078953in}}{\pgfqpoint{3.859070in}{6.081267in}}{\pgfqpoint{3.853246in}{6.081267in}}%
\pgfpathcurveto{\pgfqpoint{3.847423in}{6.081267in}}{\pgfqpoint{3.841836in}{6.078953in}}{\pgfqpoint{3.837718in}{6.074835in}}%
\pgfpathcurveto{\pgfqpoint{3.833600in}{6.070716in}}{\pgfqpoint{3.831286in}{6.065130in}}{\pgfqpoint{3.831286in}{6.059306in}}%
\pgfpathcurveto{\pgfqpoint{3.831286in}{6.053482in}}{\pgfqpoint{3.833600in}{6.047896in}}{\pgfqpoint{3.837718in}{6.043778in}}%
\pgfpathcurveto{\pgfqpoint{3.841836in}{6.039660in}}{\pgfqpoint{3.847423in}{6.037346in}}{\pgfqpoint{3.853246in}{6.037346in}}%
\pgfpathlineto{\pgfqpoint{3.853246in}{6.037346in}}%
\pgfpathclose%
\pgfusepath{stroke,fill}%
\end{pgfscope}%
\begin{pgfscope}%
\pgfpathrectangle{\pgfqpoint{1.414736in}{0.548769in}}{\pgfqpoint{3.497747in}{7.253186in}}%
\pgfusepath{clip}%
\pgfsetbuttcap%
\pgfsetroundjoin%
\definecolor{currentfill}{rgb}{0.477504,0.821444,0.318195}%
\pgfsetfillcolor{currentfill}%
\pgfsetlinewidth{1.003750pt}%
\definecolor{currentstroke}{rgb}{0.477504,0.821444,0.318195}%
\pgfsetstrokecolor{currentstroke}%
\pgfsetdash{}{0pt}%
\pgfpathmoveto{\pgfqpoint{3.758808in}{6.508332in}}%
\pgfpathcurveto{\pgfqpoint{3.764632in}{6.508332in}}{\pgfqpoint{3.770218in}{6.510646in}}{\pgfqpoint{3.774336in}{6.514764in}}%
\pgfpathcurveto{\pgfqpoint{3.778454in}{6.518882in}}{\pgfqpoint{3.780768in}{6.524468in}}{\pgfqpoint{3.780768in}{6.530292in}}%
\pgfpathcurveto{\pgfqpoint{3.780768in}{6.536116in}}{\pgfqpoint{3.778454in}{6.541702in}}{\pgfqpoint{3.774336in}{6.545821in}}%
\pgfpathcurveto{\pgfqpoint{3.770218in}{6.549939in}}{\pgfqpoint{3.764632in}{6.552253in}}{\pgfqpoint{3.758808in}{6.552253in}}%
\pgfpathcurveto{\pgfqpoint{3.752984in}{6.552253in}}{\pgfqpoint{3.747398in}{6.549939in}}{\pgfqpoint{3.743279in}{6.545821in}}%
\pgfpathcurveto{\pgfqpoint{3.739161in}{6.541702in}}{\pgfqpoint{3.736847in}{6.536116in}}{\pgfqpoint{3.736847in}{6.530292in}}%
\pgfpathcurveto{\pgfqpoint{3.736847in}{6.524468in}}{\pgfqpoint{3.739161in}{6.518882in}}{\pgfqpoint{3.743279in}{6.514764in}}%
\pgfpathcurveto{\pgfqpoint{3.747398in}{6.510646in}}{\pgfqpoint{3.752984in}{6.508332in}}{\pgfqpoint{3.758808in}{6.508332in}}%
\pgfpathlineto{\pgfqpoint{3.758808in}{6.508332in}}%
\pgfpathclose%
\pgfusepath{stroke,fill}%
\end{pgfscope}%
\begin{pgfscope}%
\pgfpathrectangle{\pgfqpoint{1.414736in}{0.548769in}}{\pgfqpoint{3.497747in}{7.253186in}}%
\pgfusepath{clip}%
\pgfsetbuttcap%
\pgfsetroundjoin%
\definecolor{currentfill}{rgb}{0.545524,0.838039,0.275626}%
\pgfsetfillcolor{currentfill}%
\pgfsetlinewidth{1.003750pt}%
\definecolor{currentstroke}{rgb}{0.545524,0.838039,0.275626}%
\pgfsetstrokecolor{currentstroke}%
\pgfsetdash{}{0pt}%
\pgfpathmoveto{\pgfqpoint{3.882549in}{6.979318in}}%
\pgfpathcurveto{\pgfqpoint{3.888373in}{6.979318in}}{\pgfqpoint{3.893959in}{6.981632in}}{\pgfqpoint{3.898077in}{6.985750in}}%
\pgfpathcurveto{\pgfqpoint{3.902195in}{6.989868in}}{\pgfqpoint{3.904509in}{6.995455in}}{\pgfqpoint{3.904509in}{7.001278in}}%
\pgfpathcurveto{\pgfqpoint{3.904509in}{7.007102in}}{\pgfqpoint{3.902195in}{7.012689in}}{\pgfqpoint{3.898077in}{7.016807in}}%
\pgfpathcurveto{\pgfqpoint{3.893959in}{7.020925in}}{\pgfqpoint{3.888373in}{7.023239in}}{\pgfqpoint{3.882549in}{7.023239in}}%
\pgfpathcurveto{\pgfqpoint{3.876725in}{7.023239in}}{\pgfqpoint{3.871139in}{7.020925in}}{\pgfqpoint{3.867020in}{7.016807in}}%
\pgfpathcurveto{\pgfqpoint{3.862902in}{7.012689in}}{\pgfqpoint{3.860588in}{7.007102in}}{\pgfqpoint{3.860588in}{7.001278in}}%
\pgfpathcurveto{\pgfqpoint{3.860588in}{6.995455in}}{\pgfqpoint{3.862902in}{6.989868in}}{\pgfqpoint{3.867020in}{6.985750in}}%
\pgfpathcurveto{\pgfqpoint{3.871139in}{6.981632in}}{\pgfqpoint{3.876725in}{6.979318in}}{\pgfqpoint{3.882549in}{6.979318in}}%
\pgfpathlineto{\pgfqpoint{3.882549in}{6.979318in}}%
\pgfpathclose%
\pgfusepath{stroke,fill}%
\end{pgfscope}%
\begin{pgfscope}%
\pgfpathrectangle{\pgfqpoint{1.414736in}{0.548769in}}{\pgfqpoint{3.497747in}{7.253186in}}%
\pgfusepath{clip}%
\pgfsetbuttcap%
\pgfsetroundjoin%
\definecolor{currentfill}{rgb}{0.412913,0.803041,0.357269}%
\pgfsetfillcolor{currentfill}%
\pgfsetlinewidth{1.003750pt}%
\definecolor{currentstroke}{rgb}{0.412913,0.803041,0.357269}%
\pgfsetstrokecolor{currentstroke}%
\pgfsetdash{}{0pt}%
\pgfpathmoveto{\pgfqpoint{3.181265in}{7.450304in}}%
\pgfpathcurveto{\pgfqpoint{3.187089in}{7.450304in}}{\pgfqpoint{3.192675in}{7.452618in}}{\pgfqpoint{3.196793in}{7.456736in}}%
\pgfpathcurveto{\pgfqpoint{3.200911in}{7.460854in}}{\pgfqpoint{3.203225in}{7.466441in}}{\pgfqpoint{3.203225in}{7.472265in}}%
\pgfpathcurveto{\pgfqpoint{3.203225in}{7.478089in}}{\pgfqpoint{3.200911in}{7.483675in}}{\pgfqpoint{3.196793in}{7.487793in}}%
\pgfpathcurveto{\pgfqpoint{3.192675in}{7.491911in}}{\pgfqpoint{3.187089in}{7.494225in}}{\pgfqpoint{3.181265in}{7.494225in}}%
\pgfpathcurveto{\pgfqpoint{3.175441in}{7.494225in}}{\pgfqpoint{3.169855in}{7.491911in}}{\pgfqpoint{3.165737in}{7.487793in}}%
\pgfpathcurveto{\pgfqpoint{3.161618in}{7.483675in}}{\pgfqpoint{3.159304in}{7.478089in}}{\pgfqpoint{3.159304in}{7.472265in}}%
\pgfpathcurveto{\pgfqpoint{3.159304in}{7.466441in}}{\pgfqpoint{3.161618in}{7.460854in}}{\pgfqpoint{3.165737in}{7.456736in}}%
\pgfpathcurveto{\pgfqpoint{3.169855in}{7.452618in}}{\pgfqpoint{3.175441in}{7.450304in}}{\pgfqpoint{3.181265in}{7.450304in}}%
\pgfpathlineto{\pgfqpoint{3.181265in}{7.450304in}}%
\pgfpathclose%
\pgfusepath{stroke,fill}%
\end{pgfscope}%
\begin{pgfscope}%
\pgfpathrectangle{\pgfqpoint{1.414736in}{0.548769in}}{\pgfqpoint{3.497747in}{7.253186in}}%
\pgfusepath{clip}%
\pgfsetbuttcap%
\pgfsetroundjoin%
\definecolor{currentfill}{rgb}{0.993248,0.906157,0.143936}%
\pgfsetfillcolor{currentfill}%
\pgfsetlinewidth{1.003750pt}%
\definecolor{currentstroke}{rgb}{0.993248,0.906157,0.143936}%
\pgfsetstrokecolor{currentstroke}%
\pgfsetdash{}{0pt}%
\pgfpathmoveto{\pgfqpoint{3.584128in}{0.856499in}}%
\pgfpathcurveto{\pgfqpoint{3.589952in}{0.856499in}}{\pgfqpoint{3.595538in}{0.858813in}}{\pgfqpoint{3.599656in}{0.862931in}}%
\pgfpathcurveto{\pgfqpoint{3.603775in}{0.867049in}}{\pgfqpoint{3.606088in}{0.872635in}}{\pgfqpoint{3.606088in}{0.878459in}}%
\pgfpathcurveto{\pgfqpoint{3.606088in}{0.884283in}}{\pgfqpoint{3.603775in}{0.889869in}}{\pgfqpoint{3.599656in}{0.893987in}}%
\pgfpathcurveto{\pgfqpoint{3.595538in}{0.898105in}}{\pgfqpoint{3.589952in}{0.900419in}}{\pgfqpoint{3.584128in}{0.900419in}}%
\pgfpathcurveto{\pgfqpoint{3.578304in}{0.900419in}}{\pgfqpoint{3.572718in}{0.898105in}}{\pgfqpoint{3.568600in}{0.893987in}}%
\pgfpathcurveto{\pgfqpoint{3.564482in}{0.889869in}}{\pgfqpoint{3.562168in}{0.884283in}}{\pgfqpoint{3.562168in}{0.878459in}}%
\pgfpathcurveto{\pgfqpoint{3.562168in}{0.872635in}}{\pgfqpoint{3.564482in}{0.867049in}}{\pgfqpoint{3.568600in}{0.862931in}}%
\pgfpathcurveto{\pgfqpoint{3.572718in}{0.858813in}}{\pgfqpoint{3.578304in}{0.856499in}}{\pgfqpoint{3.584128in}{0.856499in}}%
\pgfpathlineto{\pgfqpoint{3.584128in}{0.856499in}}%
\pgfpathclose%
\pgfusepath{stroke,fill}%
\end{pgfscope}%
\begin{pgfscope}%
\pgfpathrectangle{\pgfqpoint{1.414736in}{0.548769in}}{\pgfqpoint{3.497747in}{7.253186in}}%
\pgfusepath{clip}%
\pgfsetbuttcap%
\pgfsetroundjoin%
\definecolor{currentfill}{rgb}{0.506271,0.828786,0.300362}%
\pgfsetfillcolor{currentfill}%
\pgfsetlinewidth{1.003750pt}%
\definecolor{currentstroke}{rgb}{0.506271,0.828786,0.300362}%
\pgfsetstrokecolor{currentstroke}%
\pgfsetdash{}{0pt}%
\pgfpathmoveto{\pgfqpoint{3.607731in}{1.327485in}}%
\pgfpathcurveto{\pgfqpoint{3.613555in}{1.327485in}}{\pgfqpoint{3.619141in}{1.329799in}}{\pgfqpoint{3.623259in}{1.333917in}}%
\pgfpathcurveto{\pgfqpoint{3.627377in}{1.338035in}}{\pgfqpoint{3.629691in}{1.343621in}}{\pgfqpoint{3.629691in}{1.349445in}}%
\pgfpathcurveto{\pgfqpoint{3.629691in}{1.355269in}}{\pgfqpoint{3.627377in}{1.360855in}}{\pgfqpoint{3.623259in}{1.364973in}}%
\pgfpathcurveto{\pgfqpoint{3.619141in}{1.369092in}}{\pgfqpoint{3.613555in}{1.371405in}}{\pgfqpoint{3.607731in}{1.371405in}}%
\pgfpathcurveto{\pgfqpoint{3.601907in}{1.371405in}}{\pgfqpoint{3.596321in}{1.369092in}}{\pgfqpoint{3.592203in}{1.364973in}}%
\pgfpathcurveto{\pgfqpoint{3.588085in}{1.360855in}}{\pgfqpoint{3.585771in}{1.355269in}}{\pgfqpoint{3.585771in}{1.349445in}}%
\pgfpathcurveto{\pgfqpoint{3.585771in}{1.343621in}}{\pgfqpoint{3.588085in}{1.338035in}}{\pgfqpoint{3.592203in}{1.333917in}}%
\pgfpathcurveto{\pgfqpoint{3.596321in}{1.329799in}}{\pgfqpoint{3.601907in}{1.327485in}}{\pgfqpoint{3.607731in}{1.327485in}}%
\pgfpathlineto{\pgfqpoint{3.607731in}{1.327485in}}%
\pgfpathclose%
\pgfusepath{stroke,fill}%
\end{pgfscope}%
\begin{pgfscope}%
\pgfpathrectangle{\pgfqpoint{1.414736in}{0.548769in}}{\pgfqpoint{3.497747in}{7.253186in}}%
\pgfusepath{clip}%
\pgfsetbuttcap%
\pgfsetroundjoin%
\definecolor{currentfill}{rgb}{0.699415,0.867117,0.175971}%
\pgfsetfillcolor{currentfill}%
\pgfsetlinewidth{1.003750pt}%
\definecolor{currentstroke}{rgb}{0.699415,0.867117,0.175971}%
\pgfsetstrokecolor{currentstroke}%
\pgfsetdash{}{0pt}%
\pgfpathmoveto{\pgfqpoint{3.685954in}{1.798471in}}%
\pgfpathcurveto{\pgfqpoint{3.691778in}{1.798471in}}{\pgfqpoint{3.697364in}{1.800785in}}{\pgfqpoint{3.701482in}{1.804903in}}%
\pgfpathcurveto{\pgfqpoint{3.705600in}{1.809021in}}{\pgfqpoint{3.707914in}{1.814607in}}{\pgfqpoint{3.707914in}{1.820431in}}%
\pgfpathcurveto{\pgfqpoint{3.707914in}{1.826255in}}{\pgfqpoint{3.705600in}{1.831841in}}{\pgfqpoint{3.701482in}{1.835960in}}%
\pgfpathcurveto{\pgfqpoint{3.697364in}{1.840078in}}{\pgfqpoint{3.691778in}{1.842392in}}{\pgfqpoint{3.685954in}{1.842392in}}%
\pgfpathcurveto{\pgfqpoint{3.680130in}{1.842392in}}{\pgfqpoint{3.674544in}{1.840078in}}{\pgfqpoint{3.670426in}{1.835960in}}%
\pgfpathcurveto{\pgfqpoint{3.666308in}{1.831841in}}{\pgfqpoint{3.663994in}{1.826255in}}{\pgfqpoint{3.663994in}{1.820431in}}%
\pgfpathcurveto{\pgfqpoint{3.663994in}{1.814607in}}{\pgfqpoint{3.666308in}{1.809021in}}{\pgfqpoint{3.670426in}{1.804903in}}%
\pgfpathcurveto{\pgfqpoint{3.674544in}{1.800785in}}{\pgfqpoint{3.680130in}{1.798471in}}{\pgfqpoint{3.685954in}{1.798471in}}%
\pgfpathlineto{\pgfqpoint{3.685954in}{1.798471in}}%
\pgfpathclose%
\pgfusepath{stroke,fill}%
\end{pgfscope}%
\begin{pgfscope}%
\pgfpathrectangle{\pgfqpoint{1.414736in}{0.548769in}}{\pgfqpoint{3.497747in}{7.253186in}}%
\pgfusepath{clip}%
\pgfsetbuttcap%
\pgfsetroundjoin%
\definecolor{currentfill}{rgb}{0.983868,0.904867,0.136897}%
\pgfsetfillcolor{currentfill}%
\pgfsetlinewidth{1.003750pt}%
\definecolor{currentstroke}{rgb}{0.983868,0.904867,0.136897}%
\pgfsetstrokecolor{currentstroke}%
\pgfsetdash{}{0pt}%
\pgfpathmoveto{\pgfqpoint{3.226483in}{2.269457in}}%
\pgfpathcurveto{\pgfqpoint{3.232307in}{2.269457in}}{\pgfqpoint{3.237894in}{2.271771in}}{\pgfqpoint{3.242012in}{2.275889in}}%
\pgfpathcurveto{\pgfqpoint{3.246130in}{2.280007in}}{\pgfqpoint{3.248444in}{2.285593in}}{\pgfqpoint{3.248444in}{2.291417in}}%
\pgfpathcurveto{\pgfqpoint{3.248444in}{2.297241in}}{\pgfqpoint{3.246130in}{2.302828in}}{\pgfqpoint{3.242012in}{2.306946in}}%
\pgfpathcurveto{\pgfqpoint{3.237894in}{2.311064in}}{\pgfqpoint{3.232307in}{2.313378in}}{\pgfqpoint{3.226483in}{2.313378in}}%
\pgfpathcurveto{\pgfqpoint{3.220659in}{2.313378in}}{\pgfqpoint{3.215073in}{2.311064in}}{\pgfqpoint{3.210955in}{2.306946in}}%
\pgfpathcurveto{\pgfqpoint{3.206837in}{2.302828in}}{\pgfqpoint{3.204523in}{2.297241in}}{\pgfqpoint{3.204523in}{2.291417in}}%
\pgfpathcurveto{\pgfqpoint{3.204523in}{2.285593in}}{\pgfqpoint{3.206837in}{2.280007in}}{\pgfqpoint{3.210955in}{2.275889in}}%
\pgfpathcurveto{\pgfqpoint{3.215073in}{2.271771in}}{\pgfqpoint{3.220659in}{2.269457in}}{\pgfqpoint{3.226483in}{2.269457in}}%
\pgfpathlineto{\pgfqpoint{3.226483in}{2.269457in}}%
\pgfpathclose%
\pgfusepath{stroke,fill}%
\end{pgfscope}%
\begin{pgfscope}%
\pgfpathrectangle{\pgfqpoint{1.414736in}{0.548769in}}{\pgfqpoint{3.497747in}{7.253186in}}%
\pgfusepath{clip}%
\pgfsetbuttcap%
\pgfsetroundjoin%
\definecolor{currentfill}{rgb}{0.239374,0.735588,0.455688}%
\pgfsetfillcolor{currentfill}%
\pgfsetlinewidth{1.003750pt}%
\definecolor{currentstroke}{rgb}{0.239374,0.735588,0.455688}%
\pgfsetstrokecolor{currentstroke}%
\pgfsetdash{}{0pt}%
\pgfpathmoveto{\pgfqpoint{3.850067in}{2.740443in}}%
\pgfpathcurveto{\pgfqpoint{3.855891in}{2.740443in}}{\pgfqpoint{3.861478in}{2.742757in}}{\pgfqpoint{3.865596in}{2.746875in}}%
\pgfpathcurveto{\pgfqpoint{3.869714in}{2.750993in}}{\pgfqpoint{3.872028in}{2.756580in}}{\pgfqpoint{3.872028in}{2.762404in}}%
\pgfpathcurveto{\pgfqpoint{3.872028in}{2.768227in}}{\pgfqpoint{3.869714in}{2.773814in}}{\pgfqpoint{3.865596in}{2.777932in}}%
\pgfpathcurveto{\pgfqpoint{3.861478in}{2.782050in}}{\pgfqpoint{3.855891in}{2.784364in}}{\pgfqpoint{3.850067in}{2.784364in}}%
\pgfpathcurveto{\pgfqpoint{3.844244in}{2.784364in}}{\pgfqpoint{3.838657in}{2.782050in}}{\pgfqpoint{3.834539in}{2.777932in}}%
\pgfpathcurveto{\pgfqpoint{3.830421in}{2.773814in}}{\pgfqpoint{3.828107in}{2.768227in}}{\pgfqpoint{3.828107in}{2.762404in}}%
\pgfpathcurveto{\pgfqpoint{3.828107in}{2.756580in}}{\pgfqpoint{3.830421in}{2.750993in}}{\pgfqpoint{3.834539in}{2.746875in}}%
\pgfpathcurveto{\pgfqpoint{3.838657in}{2.742757in}}{\pgfqpoint{3.844244in}{2.740443in}}{\pgfqpoint{3.850067in}{2.740443in}}%
\pgfpathlineto{\pgfqpoint{3.850067in}{2.740443in}}%
\pgfpathclose%
\pgfusepath{stroke,fill}%
\end{pgfscope}%
\begin{pgfscope}%
\pgfpathrectangle{\pgfqpoint{1.414736in}{0.548769in}}{\pgfqpoint{3.497747in}{7.253186in}}%
\pgfusepath{clip}%
\pgfsetbuttcap%
\pgfsetroundjoin%
\definecolor{currentfill}{rgb}{0.616293,0.852709,0.230052}%
\pgfsetfillcolor{currentfill}%
\pgfsetlinewidth{1.003750pt}%
\definecolor{currentstroke}{rgb}{0.616293,0.852709,0.230052}%
\pgfsetstrokecolor{currentstroke}%
\pgfsetdash{}{0pt}%
\pgfpathmoveto{\pgfqpoint{3.418720in}{3.211429in}}%
\pgfpathcurveto{\pgfqpoint{3.424544in}{3.211429in}}{\pgfqpoint{3.430130in}{3.213743in}}{\pgfqpoint{3.434248in}{3.217861in}}%
\pgfpathcurveto{\pgfqpoint{3.438367in}{3.221980in}}{\pgfqpoint{3.440680in}{3.227566in}}{\pgfqpoint{3.440680in}{3.233390in}}%
\pgfpathcurveto{\pgfqpoint{3.440680in}{3.239214in}}{\pgfqpoint{3.438367in}{3.244800in}}{\pgfqpoint{3.434248in}{3.248918in}}%
\pgfpathcurveto{\pgfqpoint{3.430130in}{3.253036in}}{\pgfqpoint{3.424544in}{3.255350in}}{\pgfqpoint{3.418720in}{3.255350in}}%
\pgfpathcurveto{\pgfqpoint{3.412896in}{3.255350in}}{\pgfqpoint{3.407310in}{3.253036in}}{\pgfqpoint{3.403192in}{3.248918in}}%
\pgfpathcurveto{\pgfqpoint{3.399074in}{3.244800in}}{\pgfqpoint{3.396760in}{3.239214in}}{\pgfqpoint{3.396760in}{3.233390in}}%
\pgfpathcurveto{\pgfqpoint{3.396760in}{3.227566in}}{\pgfqpoint{3.399074in}{3.221980in}}{\pgfqpoint{3.403192in}{3.217861in}}%
\pgfpathcurveto{\pgfqpoint{3.407310in}{3.213743in}}{\pgfqpoint{3.412896in}{3.211429in}}{\pgfqpoint{3.418720in}{3.211429in}}%
\pgfpathlineto{\pgfqpoint{3.418720in}{3.211429in}}%
\pgfpathclose%
\pgfusepath{stroke,fill}%
\end{pgfscope}%
\begin{pgfscope}%
\pgfpathrectangle{\pgfqpoint{1.414736in}{0.548769in}}{\pgfqpoint{3.497747in}{7.253186in}}%
\pgfusepath{clip}%
\pgfsetbuttcap%
\pgfsetroundjoin%
\definecolor{currentfill}{rgb}{0.993248,0.906157,0.143936}%
\pgfsetfillcolor{currentfill}%
\pgfsetlinewidth{1.003750pt}%
\definecolor{currentstroke}{rgb}{0.993248,0.906157,0.143936}%
\pgfsetstrokecolor{currentstroke}%
\pgfsetdash{}{0pt}%
\pgfpathmoveto{\pgfqpoint{3.535670in}{3.682415in}}%
\pgfpathcurveto{\pgfqpoint{3.541494in}{3.682415in}}{\pgfqpoint{3.547080in}{3.684729in}}{\pgfqpoint{3.551198in}{3.688847in}}%
\pgfpathcurveto{\pgfqpoint{3.555316in}{3.692966in}}{\pgfqpoint{3.557630in}{3.698552in}}{\pgfqpoint{3.557630in}{3.704376in}}%
\pgfpathcurveto{\pgfqpoint{3.557630in}{3.710200in}}{\pgfqpoint{3.555316in}{3.715786in}}{\pgfqpoint{3.551198in}{3.719904in}}%
\pgfpathcurveto{\pgfqpoint{3.547080in}{3.724022in}}{\pgfqpoint{3.541494in}{3.726336in}}{\pgfqpoint{3.535670in}{3.726336in}}%
\pgfpathcurveto{\pgfqpoint{3.529846in}{3.726336in}}{\pgfqpoint{3.524260in}{3.724022in}}{\pgfqpoint{3.520142in}{3.719904in}}%
\pgfpathcurveto{\pgfqpoint{3.516024in}{3.715786in}}{\pgfqpoint{3.513710in}{3.710200in}}{\pgfqpoint{3.513710in}{3.704376in}}%
\pgfpathcurveto{\pgfqpoint{3.513710in}{3.698552in}}{\pgfqpoint{3.516024in}{3.692966in}}{\pgfqpoint{3.520142in}{3.688847in}}%
\pgfpathcurveto{\pgfqpoint{3.524260in}{3.684729in}}{\pgfqpoint{3.529846in}{3.682415in}}{\pgfqpoint{3.535670in}{3.682415in}}%
\pgfpathlineto{\pgfqpoint{3.535670in}{3.682415in}}%
\pgfpathclose%
\pgfusepath{stroke,fill}%
\end{pgfscope}%
\begin{pgfscope}%
\pgfpathrectangle{\pgfqpoint{1.414736in}{0.548769in}}{\pgfqpoint{3.497747in}{7.253186in}}%
\pgfusepath{clip}%
\pgfsetbuttcap%
\pgfsetroundjoin%
\definecolor{currentfill}{rgb}{0.266941,0.748751,0.440573}%
\pgfsetfillcolor{currentfill}%
\pgfsetlinewidth{1.003750pt}%
\definecolor{currentstroke}{rgb}{0.266941,0.748751,0.440573}%
\pgfsetstrokecolor{currentstroke}%
\pgfsetdash{}{0pt}%
\pgfpathmoveto{\pgfqpoint{3.645490in}{4.153402in}}%
\pgfpathcurveto{\pgfqpoint{3.651314in}{4.153402in}}{\pgfqpoint{3.656900in}{4.155715in}}{\pgfqpoint{3.661018in}{4.159834in}}%
\pgfpathcurveto{\pgfqpoint{3.665136in}{4.163952in}}{\pgfqpoint{3.667450in}{4.169538in}}{\pgfqpoint{3.667450in}{4.175362in}}%
\pgfpathcurveto{\pgfqpoint{3.667450in}{4.181186in}}{\pgfqpoint{3.665136in}{4.186772in}}{\pgfqpoint{3.661018in}{4.190890in}}%
\pgfpathcurveto{\pgfqpoint{3.656900in}{4.195008in}}{\pgfqpoint{3.651314in}{4.197322in}}{\pgfqpoint{3.645490in}{4.197322in}}%
\pgfpathcurveto{\pgfqpoint{3.639666in}{4.197322in}}{\pgfqpoint{3.634080in}{4.195008in}}{\pgfqpoint{3.629961in}{4.190890in}}%
\pgfpathcurveto{\pgfqpoint{3.625843in}{4.186772in}}{\pgfqpoint{3.623529in}{4.181186in}}{\pgfqpoint{3.623529in}{4.175362in}}%
\pgfpathcurveto{\pgfqpoint{3.623529in}{4.169538in}}{\pgfqpoint{3.625843in}{4.163952in}}{\pgfqpoint{3.629961in}{4.159834in}}%
\pgfpathcurveto{\pgfqpoint{3.634080in}{4.155715in}}{\pgfqpoint{3.639666in}{4.153402in}}{\pgfqpoint{3.645490in}{4.153402in}}%
\pgfpathlineto{\pgfqpoint{3.645490in}{4.153402in}}%
\pgfpathclose%
\pgfusepath{stroke,fill}%
\end{pgfscope}%
\begin{pgfscope}%
\pgfpathrectangle{\pgfqpoint{1.414736in}{0.548769in}}{\pgfqpoint{3.497747in}{7.253186in}}%
\pgfusepath{clip}%
\pgfsetbuttcap%
\pgfsetroundjoin%
\definecolor{currentfill}{rgb}{0.835270,0.886029,0.102646}%
\pgfsetfillcolor{currentfill}%
\pgfsetlinewidth{1.003750pt}%
\definecolor{currentstroke}{rgb}{0.835270,0.886029,0.102646}%
\pgfsetstrokecolor{currentstroke}%
\pgfsetdash{}{0pt}%
\pgfpathmoveto{\pgfqpoint{3.637447in}{4.624388in}}%
\pgfpathcurveto{\pgfqpoint{3.643271in}{4.624388in}}{\pgfqpoint{3.648857in}{4.626702in}}{\pgfqpoint{3.652975in}{4.630820in}}%
\pgfpathcurveto{\pgfqpoint{3.657093in}{4.634938in}}{\pgfqpoint{3.659407in}{4.640524in}}{\pgfqpoint{3.659407in}{4.646348in}}%
\pgfpathcurveto{\pgfqpoint{3.659407in}{4.652172in}}{\pgfqpoint{3.657093in}{4.657758in}}{\pgfqpoint{3.652975in}{4.661876in}}%
\pgfpathcurveto{\pgfqpoint{3.648857in}{4.665994in}}{\pgfqpoint{3.643271in}{4.668308in}}{\pgfqpoint{3.637447in}{4.668308in}}%
\pgfpathcurveto{\pgfqpoint{3.631623in}{4.668308in}}{\pgfqpoint{3.626036in}{4.665994in}}{\pgfqpoint{3.621918in}{4.661876in}}%
\pgfpathcurveto{\pgfqpoint{3.617800in}{4.657758in}}{\pgfqpoint{3.615486in}{4.652172in}}{\pgfqpoint{3.615486in}{4.646348in}}%
\pgfpathcurveto{\pgfqpoint{3.615486in}{4.640524in}}{\pgfqpoint{3.617800in}{4.634938in}}{\pgfqpoint{3.621918in}{4.630820in}}%
\pgfpathcurveto{\pgfqpoint{3.626036in}{4.626702in}}{\pgfqpoint{3.631623in}{4.624388in}}{\pgfqpoint{3.637447in}{4.624388in}}%
\pgfpathlineto{\pgfqpoint{3.637447in}{4.624388in}}%
\pgfpathclose%
\pgfusepath{stroke,fill}%
\end{pgfscope}%
\begin{pgfscope}%
\pgfpathrectangle{\pgfqpoint{1.414736in}{0.548769in}}{\pgfqpoint{3.497747in}{7.253186in}}%
\pgfusepath{clip}%
\pgfsetbuttcap%
\pgfsetroundjoin%
\definecolor{currentfill}{rgb}{0.279566,0.067836,0.391917}%
\pgfsetfillcolor{currentfill}%
\pgfsetlinewidth{1.003750pt}%
\definecolor{currentstroke}{rgb}{0.279566,0.067836,0.391917}%
\pgfsetstrokecolor{currentstroke}%
\pgfsetdash{}{0pt}%
\pgfpathmoveto{\pgfqpoint{3.689083in}{5.095374in}}%
\pgfpathcurveto{\pgfqpoint{3.694907in}{5.095374in}}{\pgfqpoint{3.700493in}{5.097688in}}{\pgfqpoint{3.704611in}{5.101806in}}%
\pgfpathcurveto{\pgfqpoint{3.708730in}{5.105924in}}{\pgfqpoint{3.711043in}{5.111510in}}{\pgfqpoint{3.711043in}{5.117334in}}%
\pgfpathcurveto{\pgfqpoint{3.711043in}{5.123158in}}{\pgfqpoint{3.708730in}{5.128744in}}{\pgfqpoint{3.704611in}{5.132862in}}%
\pgfpathcurveto{\pgfqpoint{3.700493in}{5.136980in}}{\pgfqpoint{3.694907in}{5.139294in}}{\pgfqpoint{3.689083in}{5.139294in}}%
\pgfpathcurveto{\pgfqpoint{3.683259in}{5.139294in}}{\pgfqpoint{3.677673in}{5.136980in}}{\pgfqpoint{3.673555in}{5.132862in}}%
\pgfpathcurveto{\pgfqpoint{3.669437in}{5.128744in}}{\pgfqpoint{3.667123in}{5.123158in}}{\pgfqpoint{3.667123in}{5.117334in}}%
\pgfpathcurveto{\pgfqpoint{3.667123in}{5.111510in}}{\pgfqpoint{3.669437in}{5.105924in}}{\pgfqpoint{3.673555in}{5.101806in}}%
\pgfpathcurveto{\pgfqpoint{3.677673in}{5.097688in}}{\pgfqpoint{3.683259in}{5.095374in}}{\pgfqpoint{3.689083in}{5.095374in}}%
\pgfpathlineto{\pgfqpoint{3.689083in}{5.095374in}}%
\pgfpathclose%
\pgfusepath{stroke,fill}%
\end{pgfscope}%
\begin{pgfscope}%
\pgfpathrectangle{\pgfqpoint{1.414736in}{0.548769in}}{\pgfqpoint{3.497747in}{7.253186in}}%
\pgfusepath{clip}%
\pgfsetbuttcap%
\pgfsetroundjoin%
\definecolor{currentfill}{rgb}{0.241237,0.296485,0.539709}%
\pgfsetfillcolor{currentfill}%
\pgfsetlinewidth{1.003750pt}%
\definecolor{currentstroke}{rgb}{0.241237,0.296485,0.539709}%
\pgfsetstrokecolor{currentstroke}%
\pgfsetdash{}{0pt}%
\pgfpathmoveto{\pgfqpoint{3.994223in}{5.566360in}}%
\pgfpathcurveto{\pgfqpoint{4.000047in}{5.566360in}}{\pgfqpoint{4.005633in}{5.568674in}}{\pgfqpoint{4.009751in}{5.572792in}}%
\pgfpathcurveto{\pgfqpoint{4.013869in}{5.576910in}}{\pgfqpoint{4.016183in}{5.582496in}}{\pgfqpoint{4.016183in}{5.588320in}}%
\pgfpathcurveto{\pgfqpoint{4.016183in}{5.594144in}}{\pgfqpoint{4.013869in}{5.599730in}}{\pgfqpoint{4.009751in}{5.603848in}}%
\pgfpathcurveto{\pgfqpoint{4.005633in}{5.607967in}}{\pgfqpoint{4.000047in}{5.610280in}}{\pgfqpoint{3.994223in}{5.610280in}}%
\pgfpathcurveto{\pgfqpoint{3.988399in}{5.610280in}}{\pgfqpoint{3.982813in}{5.607967in}}{\pgfqpoint{3.978694in}{5.603848in}}%
\pgfpathcurveto{\pgfqpoint{3.974576in}{5.599730in}}{\pgfqpoint{3.972262in}{5.594144in}}{\pgfqpoint{3.972262in}{5.588320in}}%
\pgfpathcurveto{\pgfqpoint{3.972262in}{5.582496in}}{\pgfqpoint{3.974576in}{5.576910in}}{\pgfqpoint{3.978694in}{5.572792in}}%
\pgfpathcurveto{\pgfqpoint{3.982813in}{5.568674in}}{\pgfqpoint{3.988399in}{5.566360in}}{\pgfqpoint{3.994223in}{5.566360in}}%
\pgfpathlineto{\pgfqpoint{3.994223in}{5.566360in}}%
\pgfpathclose%
\pgfusepath{stroke,fill}%
\end{pgfscope}%
\begin{pgfscope}%
\pgfpathrectangle{\pgfqpoint{1.414736in}{0.548769in}}{\pgfqpoint{3.497747in}{7.253186in}}%
\pgfusepath{clip}%
\pgfsetbuttcap%
\pgfsetroundjoin%
\definecolor{currentfill}{rgb}{0.128087,0.647749,0.523491}%
\pgfsetfillcolor{currentfill}%
\pgfsetlinewidth{1.003750pt}%
\definecolor{currentstroke}{rgb}{0.128087,0.647749,0.523491}%
\pgfsetstrokecolor{currentstroke}%
\pgfsetdash{}{0pt}%
\pgfpathmoveto{\pgfqpoint{3.457792in}{6.037346in}}%
\pgfpathcurveto{\pgfqpoint{3.463616in}{6.037346in}}{\pgfqpoint{3.469202in}{6.039660in}}{\pgfqpoint{3.473320in}{6.043778in}}%
\pgfpathcurveto{\pgfqpoint{3.477438in}{6.047896in}}{\pgfqpoint{3.479752in}{6.053482in}}{\pgfqpoint{3.479752in}{6.059306in}}%
\pgfpathcurveto{\pgfqpoint{3.479752in}{6.065130in}}{\pgfqpoint{3.477438in}{6.070716in}}{\pgfqpoint{3.473320in}{6.074835in}}%
\pgfpathcurveto{\pgfqpoint{3.469202in}{6.078953in}}{\pgfqpoint{3.463616in}{6.081267in}}{\pgfqpoint{3.457792in}{6.081267in}}%
\pgfpathcurveto{\pgfqpoint{3.451968in}{6.081267in}}{\pgfqpoint{3.446382in}{6.078953in}}{\pgfqpoint{3.442264in}{6.074835in}}%
\pgfpathcurveto{\pgfqpoint{3.438145in}{6.070716in}}{\pgfqpoint{3.435832in}{6.065130in}}{\pgfqpoint{3.435832in}{6.059306in}}%
\pgfpathcurveto{\pgfqpoint{3.435832in}{6.053482in}}{\pgfqpoint{3.438145in}{6.047896in}}{\pgfqpoint{3.442264in}{6.043778in}}%
\pgfpathcurveto{\pgfqpoint{3.446382in}{6.039660in}}{\pgfqpoint{3.451968in}{6.037346in}}{\pgfqpoint{3.457792in}{6.037346in}}%
\pgfpathlineto{\pgfqpoint{3.457792in}{6.037346in}}%
\pgfpathclose%
\pgfusepath{stroke,fill}%
\end{pgfscope}%
\begin{pgfscope}%
\pgfpathrectangle{\pgfqpoint{1.414736in}{0.548769in}}{\pgfqpoint{3.497747in}{7.253186in}}%
\pgfusepath{clip}%
\pgfsetbuttcap%
\pgfsetroundjoin%
\definecolor{currentfill}{rgb}{0.699415,0.867117,0.175971}%
\pgfsetfillcolor{currentfill}%
\pgfsetlinewidth{1.003750pt}%
\definecolor{currentstroke}{rgb}{0.699415,0.867117,0.175971}%
\pgfsetstrokecolor{currentstroke}%
\pgfsetdash{}{0pt}%
\pgfpathmoveto{\pgfqpoint{3.661765in}{6.508332in}}%
\pgfpathcurveto{\pgfqpoint{3.667589in}{6.508332in}}{\pgfqpoint{3.673175in}{6.510646in}}{\pgfqpoint{3.677293in}{6.514764in}}%
\pgfpathcurveto{\pgfqpoint{3.681411in}{6.518882in}}{\pgfqpoint{3.683725in}{6.524468in}}{\pgfqpoint{3.683725in}{6.530292in}}%
\pgfpathcurveto{\pgfqpoint{3.683725in}{6.536116in}}{\pgfqpoint{3.681411in}{6.541702in}}{\pgfqpoint{3.677293in}{6.545821in}}%
\pgfpathcurveto{\pgfqpoint{3.673175in}{6.549939in}}{\pgfqpoint{3.667589in}{6.552253in}}{\pgfqpoint{3.661765in}{6.552253in}}%
\pgfpathcurveto{\pgfqpoint{3.655941in}{6.552253in}}{\pgfqpoint{3.650355in}{6.549939in}}{\pgfqpoint{3.646237in}{6.545821in}}%
\pgfpathcurveto{\pgfqpoint{3.642119in}{6.541702in}}{\pgfqpoint{3.639805in}{6.536116in}}{\pgfqpoint{3.639805in}{6.530292in}}%
\pgfpathcurveto{\pgfqpoint{3.639805in}{6.524468in}}{\pgfqpoint{3.642119in}{6.518882in}}{\pgfqpoint{3.646237in}{6.514764in}}%
\pgfpathcurveto{\pgfqpoint{3.650355in}{6.510646in}}{\pgfqpoint{3.655941in}{6.508332in}}{\pgfqpoint{3.661765in}{6.508332in}}%
\pgfpathlineto{\pgfqpoint{3.661765in}{6.508332in}}%
\pgfpathclose%
\pgfusepath{stroke,fill}%
\end{pgfscope}%
\begin{pgfscope}%
\pgfpathrectangle{\pgfqpoint{1.414736in}{0.548769in}}{\pgfqpoint{3.497747in}{7.253186in}}%
\pgfusepath{clip}%
\pgfsetbuttcap%
\pgfsetroundjoin%
\definecolor{currentfill}{rgb}{0.202219,0.715272,0.476084}%
\pgfsetfillcolor{currentfill}%
\pgfsetlinewidth{1.003750pt}%
\definecolor{currentstroke}{rgb}{0.202219,0.715272,0.476084}%
\pgfsetstrokecolor{currentstroke}%
\pgfsetdash{}{0pt}%
\pgfpathmoveto{\pgfqpoint{4.242317in}{6.979318in}}%
\pgfpathcurveto{\pgfqpoint{4.248141in}{6.979318in}}{\pgfqpoint{4.253727in}{6.981632in}}{\pgfqpoint{4.257845in}{6.985750in}}%
\pgfpathcurveto{\pgfqpoint{4.261963in}{6.989868in}}{\pgfqpoint{4.264277in}{6.995455in}}{\pgfqpoint{4.264277in}{7.001278in}}%
\pgfpathcurveto{\pgfqpoint{4.264277in}{7.007102in}}{\pgfqpoint{4.261963in}{7.012689in}}{\pgfqpoint{4.257845in}{7.016807in}}%
\pgfpathcurveto{\pgfqpoint{4.253727in}{7.020925in}}{\pgfqpoint{4.248141in}{7.023239in}}{\pgfqpoint{4.242317in}{7.023239in}}%
\pgfpathcurveto{\pgfqpoint{4.236493in}{7.023239in}}{\pgfqpoint{4.230907in}{7.020925in}}{\pgfqpoint{4.226789in}{7.016807in}}%
\pgfpathcurveto{\pgfqpoint{4.222671in}{7.012689in}}{\pgfqpoint{4.220357in}{7.007102in}}{\pgfqpoint{4.220357in}{7.001278in}}%
\pgfpathcurveto{\pgfqpoint{4.220357in}{6.995455in}}{\pgfqpoint{4.222671in}{6.989868in}}{\pgfqpoint{4.226789in}{6.985750in}}%
\pgfpathcurveto{\pgfqpoint{4.230907in}{6.981632in}}{\pgfqpoint{4.236493in}{6.979318in}}{\pgfqpoint{4.242317in}{6.979318in}}%
\pgfpathlineto{\pgfqpoint{4.242317in}{6.979318in}}%
\pgfpathclose%
\pgfusepath{stroke,fill}%
\end{pgfscope}%
\begin{pgfscope}%
\pgfpathrectangle{\pgfqpoint{1.414736in}{0.548769in}}{\pgfqpoint{3.497747in}{7.253186in}}%
\pgfusepath{clip}%
\pgfsetbuttcap%
\pgfsetroundjoin%
\definecolor{currentfill}{rgb}{0.120638,0.625828,0.533488}%
\pgfsetfillcolor{currentfill}%
\pgfsetlinewidth{1.003750pt}%
\definecolor{currentstroke}{rgb}{0.120638,0.625828,0.533488}%
\pgfsetstrokecolor{currentstroke}%
\pgfsetdash{}{0pt}%
\pgfpathmoveto{\pgfqpoint{3.905353in}{7.450304in}}%
\pgfpathcurveto{\pgfqpoint{3.911177in}{7.450304in}}{\pgfqpoint{3.916763in}{7.452618in}}{\pgfqpoint{3.920882in}{7.456736in}}%
\pgfpathcurveto{\pgfqpoint{3.925000in}{7.460854in}}{\pgfqpoint{3.927314in}{7.466441in}}{\pgfqpoint{3.927314in}{7.472265in}}%
\pgfpathcurveto{\pgfqpoint{3.927314in}{7.478089in}}{\pgfqpoint{3.925000in}{7.483675in}}{\pgfqpoint{3.920882in}{7.487793in}}%
\pgfpathcurveto{\pgfqpoint{3.916763in}{7.491911in}}{\pgfqpoint{3.911177in}{7.494225in}}{\pgfqpoint{3.905353in}{7.494225in}}%
\pgfpathcurveto{\pgfqpoint{3.899529in}{7.494225in}}{\pgfqpoint{3.893943in}{7.491911in}}{\pgfqpoint{3.889825in}{7.487793in}}%
\pgfpathcurveto{\pgfqpoint{3.885707in}{7.483675in}}{\pgfqpoint{3.883393in}{7.478089in}}{\pgfqpoint{3.883393in}{7.472265in}}%
\pgfpathcurveto{\pgfqpoint{3.883393in}{7.466441in}}{\pgfqpoint{3.885707in}{7.460854in}}{\pgfqpoint{3.889825in}{7.456736in}}%
\pgfpathcurveto{\pgfqpoint{3.893943in}{7.452618in}}{\pgfqpoint{3.899529in}{7.450304in}}{\pgfqpoint{3.905353in}{7.450304in}}%
\pgfpathlineto{\pgfqpoint{3.905353in}{7.450304in}}%
\pgfpathclose%
\pgfusepath{stroke,fill}%
\end{pgfscope}%
\begin{pgfscope}%
\pgfpathrectangle{\pgfqpoint{1.414736in}{0.548769in}}{\pgfqpoint{3.497747in}{7.253186in}}%
\pgfusepath{clip}%
\pgfsetbuttcap%
\pgfsetroundjoin%
\definecolor{currentfill}{rgb}{0.751884,0.874951,0.143228}%
\pgfsetfillcolor{currentfill}%
\pgfsetlinewidth{1.003750pt}%
\definecolor{currentstroke}{rgb}{0.751884,0.874951,0.143228}%
\pgfsetstrokecolor{currentstroke}%
\pgfsetdash{}{0pt}%
\pgfpathmoveto{\pgfqpoint{3.714262in}{0.856499in}}%
\pgfpathcurveto{\pgfqpoint{3.720086in}{0.856499in}}{\pgfqpoint{3.725673in}{0.858813in}}{\pgfqpoint{3.729791in}{0.862931in}}%
\pgfpathcurveto{\pgfqpoint{3.733909in}{0.867049in}}{\pgfqpoint{3.736223in}{0.872635in}}{\pgfqpoint{3.736223in}{0.878459in}}%
\pgfpathcurveto{\pgfqpoint{3.736223in}{0.884283in}}{\pgfqpoint{3.733909in}{0.889869in}}{\pgfqpoint{3.729791in}{0.893987in}}%
\pgfpathcurveto{\pgfqpoint{3.725673in}{0.898105in}}{\pgfqpoint{3.720086in}{0.900419in}}{\pgfqpoint{3.714262in}{0.900419in}}%
\pgfpathcurveto{\pgfqpoint{3.708439in}{0.900419in}}{\pgfqpoint{3.702852in}{0.898105in}}{\pgfqpoint{3.698734in}{0.893987in}}%
\pgfpathcurveto{\pgfqpoint{3.694616in}{0.889869in}}{\pgfqpoint{3.692302in}{0.884283in}}{\pgfqpoint{3.692302in}{0.878459in}}%
\pgfpathcurveto{\pgfqpoint{3.692302in}{0.872635in}}{\pgfqpoint{3.694616in}{0.867049in}}{\pgfqpoint{3.698734in}{0.862931in}}%
\pgfpathcurveto{\pgfqpoint{3.702852in}{0.858813in}}{\pgfqpoint{3.708439in}{0.856499in}}{\pgfqpoint{3.714262in}{0.856499in}}%
\pgfpathlineto{\pgfqpoint{3.714262in}{0.856499in}}%
\pgfpathclose%
\pgfusepath{stroke,fill}%
\end{pgfscope}%
\begin{pgfscope}%
\pgfpathrectangle{\pgfqpoint{1.414736in}{0.548769in}}{\pgfqpoint{3.497747in}{7.253186in}}%
\pgfusepath{clip}%
\pgfsetbuttcap%
\pgfsetroundjoin%
\definecolor{currentfill}{rgb}{0.606045,0.850733,0.236712}%
\pgfsetfillcolor{currentfill}%
\pgfsetlinewidth{1.003750pt}%
\definecolor{currentstroke}{rgb}{0.606045,0.850733,0.236712}%
\pgfsetstrokecolor{currentstroke}%
\pgfsetdash{}{0pt}%
\pgfpathmoveto{\pgfqpoint{3.463204in}{1.327485in}}%
\pgfpathcurveto{\pgfqpoint{3.469028in}{1.327485in}}{\pgfqpoint{3.474614in}{1.329799in}}{\pgfqpoint{3.478732in}{1.333917in}}%
\pgfpathcurveto{\pgfqpoint{3.482850in}{1.338035in}}{\pgfqpoint{3.485164in}{1.343621in}}{\pgfqpoint{3.485164in}{1.349445in}}%
\pgfpathcurveto{\pgfqpoint{3.485164in}{1.355269in}}{\pgfqpoint{3.482850in}{1.360855in}}{\pgfqpoint{3.478732in}{1.364973in}}%
\pgfpathcurveto{\pgfqpoint{3.474614in}{1.369092in}}{\pgfqpoint{3.469028in}{1.371405in}}{\pgfqpoint{3.463204in}{1.371405in}}%
\pgfpathcurveto{\pgfqpoint{3.457380in}{1.371405in}}{\pgfqpoint{3.451794in}{1.369092in}}{\pgfqpoint{3.447675in}{1.364973in}}%
\pgfpathcurveto{\pgfqpoint{3.443557in}{1.360855in}}{\pgfqpoint{3.441243in}{1.355269in}}{\pgfqpoint{3.441243in}{1.349445in}}%
\pgfpathcurveto{\pgfqpoint{3.441243in}{1.343621in}}{\pgfqpoint{3.443557in}{1.338035in}}{\pgfqpoint{3.447675in}{1.333917in}}%
\pgfpathcurveto{\pgfqpoint{3.451794in}{1.329799in}}{\pgfqpoint{3.457380in}{1.327485in}}{\pgfqpoint{3.463204in}{1.327485in}}%
\pgfpathlineto{\pgfqpoint{3.463204in}{1.327485in}}%
\pgfpathclose%
\pgfusepath{stroke,fill}%
\end{pgfscope}%
\begin{pgfscope}%
\pgfpathrectangle{\pgfqpoint{1.414736in}{0.548769in}}{\pgfqpoint{3.497747in}{7.253186in}}%
\pgfusepath{clip}%
\pgfsetbuttcap%
\pgfsetroundjoin%
\definecolor{currentfill}{rgb}{0.699415,0.867117,0.175971}%
\pgfsetfillcolor{currentfill}%
\pgfsetlinewidth{1.003750pt}%
\definecolor{currentstroke}{rgb}{0.699415,0.867117,0.175971}%
\pgfsetstrokecolor{currentstroke}%
\pgfsetdash{}{0pt}%
\pgfpathmoveto{\pgfqpoint{3.748223in}{1.798471in}}%
\pgfpathcurveto{\pgfqpoint{3.754047in}{1.798471in}}{\pgfqpoint{3.759633in}{1.800785in}}{\pgfqpoint{3.763751in}{1.804903in}}%
\pgfpathcurveto{\pgfqpoint{3.767869in}{1.809021in}}{\pgfqpoint{3.770183in}{1.814607in}}{\pgfqpoint{3.770183in}{1.820431in}}%
\pgfpathcurveto{\pgfqpoint{3.770183in}{1.826255in}}{\pgfqpoint{3.767869in}{1.831841in}}{\pgfqpoint{3.763751in}{1.835960in}}%
\pgfpathcurveto{\pgfqpoint{3.759633in}{1.840078in}}{\pgfqpoint{3.754047in}{1.842392in}}{\pgfqpoint{3.748223in}{1.842392in}}%
\pgfpathcurveto{\pgfqpoint{3.742399in}{1.842392in}}{\pgfqpoint{3.736813in}{1.840078in}}{\pgfqpoint{3.732695in}{1.835960in}}%
\pgfpathcurveto{\pgfqpoint{3.728577in}{1.831841in}}{\pgfqpoint{3.726263in}{1.826255in}}{\pgfqpoint{3.726263in}{1.820431in}}%
\pgfpathcurveto{\pgfqpoint{3.726263in}{1.814607in}}{\pgfqpoint{3.728577in}{1.809021in}}{\pgfqpoint{3.732695in}{1.804903in}}%
\pgfpathcurveto{\pgfqpoint{3.736813in}{1.800785in}}{\pgfqpoint{3.742399in}{1.798471in}}{\pgfqpoint{3.748223in}{1.798471in}}%
\pgfpathlineto{\pgfqpoint{3.748223in}{1.798471in}}%
\pgfpathclose%
\pgfusepath{stroke,fill}%
\end{pgfscope}%
\begin{pgfscope}%
\pgfpathrectangle{\pgfqpoint{1.414736in}{0.548769in}}{\pgfqpoint{3.497747in}{7.253186in}}%
\pgfusepath{clip}%
\pgfsetbuttcap%
\pgfsetroundjoin%
\definecolor{currentfill}{rgb}{0.993248,0.906157,0.143936}%
\pgfsetfillcolor{currentfill}%
\pgfsetlinewidth{1.003750pt}%
\definecolor{currentstroke}{rgb}{0.993248,0.906157,0.143936}%
\pgfsetstrokecolor{currentstroke}%
\pgfsetdash{}{0pt}%
\pgfpathmoveto{\pgfqpoint{3.802653in}{2.269457in}}%
\pgfpathcurveto{\pgfqpoint{3.808477in}{2.269457in}}{\pgfqpoint{3.814063in}{2.271771in}}{\pgfqpoint{3.818181in}{2.275889in}}%
\pgfpathcurveto{\pgfqpoint{3.822299in}{2.280007in}}{\pgfqpoint{3.824613in}{2.285593in}}{\pgfqpoint{3.824613in}{2.291417in}}%
\pgfpathcurveto{\pgfqpoint{3.824613in}{2.297241in}}{\pgfqpoint{3.822299in}{2.302828in}}{\pgfqpoint{3.818181in}{2.306946in}}%
\pgfpathcurveto{\pgfqpoint{3.814063in}{2.311064in}}{\pgfqpoint{3.808477in}{2.313378in}}{\pgfqpoint{3.802653in}{2.313378in}}%
\pgfpathcurveto{\pgfqpoint{3.796829in}{2.313378in}}{\pgfqpoint{3.791243in}{2.311064in}}{\pgfqpoint{3.787125in}{2.306946in}}%
\pgfpathcurveto{\pgfqpoint{3.783007in}{2.302828in}}{\pgfqpoint{3.780693in}{2.297241in}}{\pgfqpoint{3.780693in}{2.291417in}}%
\pgfpathcurveto{\pgfqpoint{3.780693in}{2.285593in}}{\pgfqpoint{3.783007in}{2.280007in}}{\pgfqpoint{3.787125in}{2.275889in}}%
\pgfpathcurveto{\pgfqpoint{3.791243in}{2.271771in}}{\pgfqpoint{3.796829in}{2.269457in}}{\pgfqpoint{3.802653in}{2.269457in}}%
\pgfpathlineto{\pgfqpoint{3.802653in}{2.269457in}}%
\pgfpathclose%
\pgfusepath{stroke,fill}%
\end{pgfscope}%
\begin{pgfscope}%
\pgfpathrectangle{\pgfqpoint{1.414736in}{0.548769in}}{\pgfqpoint{3.497747in}{7.253186in}}%
\pgfusepath{clip}%
\pgfsetbuttcap%
\pgfsetroundjoin%
\definecolor{currentfill}{rgb}{0.246070,0.738910,0.452024}%
\pgfsetfillcolor{currentfill}%
\pgfsetlinewidth{1.003750pt}%
\definecolor{currentstroke}{rgb}{0.246070,0.738910,0.452024}%
\pgfsetstrokecolor{currentstroke}%
\pgfsetdash{}{0pt}%
\pgfpathmoveto{\pgfqpoint{3.796622in}{2.740443in}}%
\pgfpathcurveto{\pgfqpoint{3.802446in}{2.740443in}}{\pgfqpoint{3.808032in}{2.742757in}}{\pgfqpoint{3.812151in}{2.746875in}}%
\pgfpathcurveto{\pgfqpoint{3.816269in}{2.750993in}}{\pgfqpoint{3.818583in}{2.756580in}}{\pgfqpoint{3.818583in}{2.762404in}}%
\pgfpathcurveto{\pgfqpoint{3.818583in}{2.768227in}}{\pgfqpoint{3.816269in}{2.773814in}}{\pgfqpoint{3.812151in}{2.777932in}}%
\pgfpathcurveto{\pgfqpoint{3.808032in}{2.782050in}}{\pgfqpoint{3.802446in}{2.784364in}}{\pgfqpoint{3.796622in}{2.784364in}}%
\pgfpathcurveto{\pgfqpoint{3.790798in}{2.784364in}}{\pgfqpoint{3.785212in}{2.782050in}}{\pgfqpoint{3.781094in}{2.777932in}}%
\pgfpathcurveto{\pgfqpoint{3.776976in}{2.773814in}}{\pgfqpoint{3.774662in}{2.768227in}}{\pgfqpoint{3.774662in}{2.762404in}}%
\pgfpathcurveto{\pgfqpoint{3.774662in}{2.756580in}}{\pgfqpoint{3.776976in}{2.750993in}}{\pgfqpoint{3.781094in}{2.746875in}}%
\pgfpathcurveto{\pgfqpoint{3.785212in}{2.742757in}}{\pgfqpoint{3.790798in}{2.740443in}}{\pgfqpoint{3.796622in}{2.740443in}}%
\pgfpathlineto{\pgfqpoint{3.796622in}{2.740443in}}%
\pgfpathclose%
\pgfusepath{stroke,fill}%
\end{pgfscope}%
\begin{pgfscope}%
\pgfpathrectangle{\pgfqpoint{1.414736in}{0.548769in}}{\pgfqpoint{3.497747in}{7.253186in}}%
\pgfusepath{clip}%
\pgfsetbuttcap%
\pgfsetroundjoin%
\definecolor{currentfill}{rgb}{0.916242,0.896091,0.100717}%
\pgfsetfillcolor{currentfill}%
\pgfsetlinewidth{1.003750pt}%
\definecolor{currentstroke}{rgb}{0.916242,0.896091,0.100717}%
\pgfsetstrokecolor{currentstroke}%
\pgfsetdash{}{0pt}%
\pgfpathmoveto{\pgfqpoint{3.448359in}{3.211429in}}%
\pgfpathcurveto{\pgfqpoint{3.454183in}{3.211429in}}{\pgfqpoint{3.459769in}{3.213743in}}{\pgfqpoint{3.463887in}{3.217861in}}%
\pgfpathcurveto{\pgfqpoint{3.468005in}{3.221980in}}{\pgfqpoint{3.470319in}{3.227566in}}{\pgfqpoint{3.470319in}{3.233390in}}%
\pgfpathcurveto{\pgfqpoint{3.470319in}{3.239214in}}{\pgfqpoint{3.468005in}{3.244800in}}{\pgfqpoint{3.463887in}{3.248918in}}%
\pgfpathcurveto{\pgfqpoint{3.459769in}{3.253036in}}{\pgfqpoint{3.454183in}{3.255350in}}{\pgfqpoint{3.448359in}{3.255350in}}%
\pgfpathcurveto{\pgfqpoint{3.442535in}{3.255350in}}{\pgfqpoint{3.436949in}{3.253036in}}{\pgfqpoint{3.432831in}{3.248918in}}%
\pgfpathcurveto{\pgfqpoint{3.428713in}{3.244800in}}{\pgfqpoint{3.426399in}{3.239214in}}{\pgfqpoint{3.426399in}{3.233390in}}%
\pgfpathcurveto{\pgfqpoint{3.426399in}{3.227566in}}{\pgfqpoint{3.428713in}{3.221980in}}{\pgfqpoint{3.432831in}{3.217861in}}%
\pgfpathcurveto{\pgfqpoint{3.436949in}{3.213743in}}{\pgfqpoint{3.442535in}{3.211429in}}{\pgfqpoint{3.448359in}{3.211429in}}%
\pgfpathlineto{\pgfqpoint{3.448359in}{3.211429in}}%
\pgfpathclose%
\pgfusepath{stroke,fill}%
\end{pgfscope}%
\begin{pgfscope}%
\pgfpathrectangle{\pgfqpoint{1.414736in}{0.548769in}}{\pgfqpoint{3.497747in}{7.253186in}}%
\pgfusepath{clip}%
\pgfsetbuttcap%
\pgfsetroundjoin%
\definecolor{currentfill}{rgb}{0.876168,0.891125,0.095250}%
\pgfsetfillcolor{currentfill}%
\pgfsetlinewidth{1.003750pt}%
\definecolor{currentstroke}{rgb}{0.876168,0.891125,0.095250}%
\pgfsetstrokecolor{currentstroke}%
\pgfsetdash{}{0pt}%
\pgfpathmoveto{\pgfqpoint{3.406412in}{3.682415in}}%
\pgfpathcurveto{\pgfqpoint{3.412236in}{3.682415in}}{\pgfqpoint{3.417822in}{3.684729in}}{\pgfqpoint{3.421940in}{3.688847in}}%
\pgfpathcurveto{\pgfqpoint{3.426058in}{3.692966in}}{\pgfqpoint{3.428372in}{3.698552in}}{\pgfqpoint{3.428372in}{3.704376in}}%
\pgfpathcurveto{\pgfqpoint{3.428372in}{3.710200in}}{\pgfqpoint{3.426058in}{3.715786in}}{\pgfqpoint{3.421940in}{3.719904in}}%
\pgfpathcurveto{\pgfqpoint{3.417822in}{3.724022in}}{\pgfqpoint{3.412236in}{3.726336in}}{\pgfqpoint{3.406412in}{3.726336in}}%
\pgfpathcurveto{\pgfqpoint{3.400588in}{3.726336in}}{\pgfqpoint{3.395002in}{3.724022in}}{\pgfqpoint{3.390883in}{3.719904in}}%
\pgfpathcurveto{\pgfqpoint{3.386765in}{3.715786in}}{\pgfqpoint{3.384451in}{3.710200in}}{\pgfqpoint{3.384451in}{3.704376in}}%
\pgfpathcurveto{\pgfqpoint{3.384451in}{3.698552in}}{\pgfqpoint{3.386765in}{3.692966in}}{\pgfqpoint{3.390883in}{3.688847in}}%
\pgfpathcurveto{\pgfqpoint{3.395002in}{3.684729in}}{\pgfqpoint{3.400588in}{3.682415in}}{\pgfqpoint{3.406412in}{3.682415in}}%
\pgfpathlineto{\pgfqpoint{3.406412in}{3.682415in}}%
\pgfpathclose%
\pgfusepath{stroke,fill}%
\end{pgfscope}%
\begin{pgfscope}%
\pgfpathrectangle{\pgfqpoint{1.414736in}{0.548769in}}{\pgfqpoint{3.497747in}{7.253186in}}%
\pgfusepath{clip}%
\pgfsetbuttcap%
\pgfsetroundjoin%
\definecolor{currentfill}{rgb}{0.487026,0.823929,0.312321}%
\pgfsetfillcolor{currentfill}%
\pgfsetlinewidth{1.003750pt}%
\definecolor{currentstroke}{rgb}{0.487026,0.823929,0.312321}%
\pgfsetstrokecolor{currentstroke}%
\pgfsetdash{}{0pt}%
\pgfpathmoveto{\pgfqpoint{3.440672in}{4.153402in}}%
\pgfpathcurveto{\pgfqpoint{3.446496in}{4.153402in}}{\pgfqpoint{3.452082in}{4.155715in}}{\pgfqpoint{3.456201in}{4.159834in}}%
\pgfpathcurveto{\pgfqpoint{3.460319in}{4.163952in}}{\pgfqpoint{3.462633in}{4.169538in}}{\pgfqpoint{3.462633in}{4.175362in}}%
\pgfpathcurveto{\pgfqpoint{3.462633in}{4.181186in}}{\pgfqpoint{3.460319in}{4.186772in}}{\pgfqpoint{3.456201in}{4.190890in}}%
\pgfpathcurveto{\pgfqpoint{3.452082in}{4.195008in}}{\pgfqpoint{3.446496in}{4.197322in}}{\pgfqpoint{3.440672in}{4.197322in}}%
\pgfpathcurveto{\pgfqpoint{3.434848in}{4.197322in}}{\pgfqpoint{3.429262in}{4.195008in}}{\pgfqpoint{3.425144in}{4.190890in}}%
\pgfpathcurveto{\pgfqpoint{3.421026in}{4.186772in}}{\pgfqpoint{3.418712in}{4.181186in}}{\pgfqpoint{3.418712in}{4.175362in}}%
\pgfpathcurveto{\pgfqpoint{3.418712in}{4.169538in}}{\pgfqpoint{3.421026in}{4.163952in}}{\pgfqpoint{3.425144in}{4.159834in}}%
\pgfpathcurveto{\pgfqpoint{3.429262in}{4.155715in}}{\pgfqpoint{3.434848in}{4.153402in}}{\pgfqpoint{3.440672in}{4.153402in}}%
\pgfpathlineto{\pgfqpoint{3.440672in}{4.153402in}}%
\pgfpathclose%
\pgfusepath{stroke,fill}%
\end{pgfscope}%
\begin{pgfscope}%
\pgfpathrectangle{\pgfqpoint{1.414736in}{0.548769in}}{\pgfqpoint{3.497747in}{7.253186in}}%
\pgfusepath{clip}%
\pgfsetbuttcap%
\pgfsetroundjoin%
\definecolor{currentfill}{rgb}{0.804182,0.882046,0.114965}%
\pgfsetfillcolor{currentfill}%
\pgfsetlinewidth{1.003750pt}%
\definecolor{currentstroke}{rgb}{0.804182,0.882046,0.114965}%
\pgfsetstrokecolor{currentstroke}%
\pgfsetdash{}{0pt}%
\pgfpathmoveto{\pgfqpoint{3.821430in}{4.624388in}}%
\pgfpathcurveto{\pgfqpoint{3.827253in}{4.624388in}}{\pgfqpoint{3.832840in}{4.626702in}}{\pgfqpoint{3.836958in}{4.630820in}}%
\pgfpathcurveto{\pgfqpoint{3.841076in}{4.634938in}}{\pgfqpoint{3.843390in}{4.640524in}}{\pgfqpoint{3.843390in}{4.646348in}}%
\pgfpathcurveto{\pgfqpoint{3.843390in}{4.652172in}}{\pgfqpoint{3.841076in}{4.657758in}}{\pgfqpoint{3.836958in}{4.661876in}}%
\pgfpathcurveto{\pgfqpoint{3.832840in}{4.665994in}}{\pgfqpoint{3.827253in}{4.668308in}}{\pgfqpoint{3.821430in}{4.668308in}}%
\pgfpathcurveto{\pgfqpoint{3.815606in}{4.668308in}}{\pgfqpoint{3.810019in}{4.665994in}}{\pgfqpoint{3.805901in}{4.661876in}}%
\pgfpathcurveto{\pgfqpoint{3.801783in}{4.657758in}}{\pgfqpoint{3.799469in}{4.652172in}}{\pgfqpoint{3.799469in}{4.646348in}}%
\pgfpathcurveto{\pgfqpoint{3.799469in}{4.640524in}}{\pgfqpoint{3.801783in}{4.634938in}}{\pgfqpoint{3.805901in}{4.630820in}}%
\pgfpathcurveto{\pgfqpoint{3.810019in}{4.626702in}}{\pgfqpoint{3.815606in}{4.624388in}}{\pgfqpoint{3.821430in}{4.624388in}}%
\pgfpathlineto{\pgfqpoint{3.821430in}{4.624388in}}%
\pgfpathclose%
\pgfusepath{stroke,fill}%
\end{pgfscope}%
\begin{pgfscope}%
\pgfpathrectangle{\pgfqpoint{1.414736in}{0.548769in}}{\pgfqpoint{3.497747in}{7.253186in}}%
\pgfusepath{clip}%
\pgfsetbuttcap%
\pgfsetroundjoin%
\definecolor{currentfill}{rgb}{0.202219,0.715272,0.476084}%
\pgfsetfillcolor{currentfill}%
\pgfsetlinewidth{1.003750pt}%
\definecolor{currentstroke}{rgb}{0.202219,0.715272,0.476084}%
\pgfsetstrokecolor{currentstroke}%
\pgfsetdash{}{0pt}%
\pgfpathmoveto{\pgfqpoint{3.971876in}{5.095374in}}%
\pgfpathcurveto{\pgfqpoint{3.977700in}{5.095374in}}{\pgfqpoint{3.983286in}{5.097688in}}{\pgfqpoint{3.987405in}{5.101806in}}%
\pgfpathcurveto{\pgfqpoint{3.991523in}{5.105924in}}{\pgfqpoint{3.993837in}{5.111510in}}{\pgfqpoint{3.993837in}{5.117334in}}%
\pgfpathcurveto{\pgfqpoint{3.993837in}{5.123158in}}{\pgfqpoint{3.991523in}{5.128744in}}{\pgfqpoint{3.987405in}{5.132862in}}%
\pgfpathcurveto{\pgfqpoint{3.983286in}{5.136980in}}{\pgfqpoint{3.977700in}{5.139294in}}{\pgfqpoint{3.971876in}{5.139294in}}%
\pgfpathcurveto{\pgfqpoint{3.966052in}{5.139294in}}{\pgfqpoint{3.960466in}{5.136980in}}{\pgfqpoint{3.956348in}{5.132862in}}%
\pgfpathcurveto{\pgfqpoint{3.952230in}{5.128744in}}{\pgfqpoint{3.949916in}{5.123158in}}{\pgfqpoint{3.949916in}{5.117334in}}%
\pgfpathcurveto{\pgfqpoint{3.949916in}{5.111510in}}{\pgfqpoint{3.952230in}{5.105924in}}{\pgfqpoint{3.956348in}{5.101806in}}%
\pgfpathcurveto{\pgfqpoint{3.960466in}{5.097688in}}{\pgfqpoint{3.966052in}{5.095374in}}{\pgfqpoint{3.971876in}{5.095374in}}%
\pgfpathlineto{\pgfqpoint{3.971876in}{5.095374in}}%
\pgfpathclose%
\pgfusepath{stroke,fill}%
\end{pgfscope}%
\begin{pgfscope}%
\pgfpathrectangle{\pgfqpoint{1.414736in}{0.548769in}}{\pgfqpoint{3.497747in}{7.253186in}}%
\pgfusepath{clip}%
\pgfsetbuttcap%
\pgfsetroundjoin%
\definecolor{currentfill}{rgb}{0.119512,0.607464,0.540218}%
\pgfsetfillcolor{currentfill}%
\pgfsetlinewidth{1.003750pt}%
\definecolor{currentstroke}{rgb}{0.119512,0.607464,0.540218}%
\pgfsetstrokecolor{currentstroke}%
\pgfsetdash{}{0pt}%
\pgfpathmoveto{\pgfqpoint{3.204842in}{5.566360in}}%
\pgfpathcurveto{\pgfqpoint{3.210666in}{5.566360in}}{\pgfqpoint{3.216252in}{5.568674in}}{\pgfqpoint{3.220370in}{5.572792in}}%
\pgfpathcurveto{\pgfqpoint{3.224488in}{5.576910in}}{\pgfqpoint{3.226802in}{5.582496in}}{\pgfqpoint{3.226802in}{5.588320in}}%
\pgfpathcurveto{\pgfqpoint{3.226802in}{5.594144in}}{\pgfqpoint{3.224488in}{5.599730in}}{\pgfqpoint{3.220370in}{5.603848in}}%
\pgfpathcurveto{\pgfqpoint{3.216252in}{5.607967in}}{\pgfqpoint{3.210666in}{5.610280in}}{\pgfqpoint{3.204842in}{5.610280in}}%
\pgfpathcurveto{\pgfqpoint{3.199018in}{5.610280in}}{\pgfqpoint{3.193432in}{5.607967in}}{\pgfqpoint{3.189314in}{5.603848in}}%
\pgfpathcurveto{\pgfqpoint{3.185195in}{5.599730in}}{\pgfqpoint{3.182882in}{5.594144in}}{\pgfqpoint{3.182882in}{5.588320in}}%
\pgfpathcurveto{\pgfqpoint{3.182882in}{5.582496in}}{\pgfqpoint{3.185195in}{5.576910in}}{\pgfqpoint{3.189314in}{5.572792in}}%
\pgfpathcurveto{\pgfqpoint{3.193432in}{5.568674in}}{\pgfqpoint{3.199018in}{5.566360in}}{\pgfqpoint{3.204842in}{5.566360in}}%
\pgfpathlineto{\pgfqpoint{3.204842in}{5.566360in}}%
\pgfpathclose%
\pgfusepath{stroke,fill}%
\end{pgfscope}%
\begin{pgfscope}%
\pgfpathrectangle{\pgfqpoint{1.414736in}{0.548769in}}{\pgfqpoint{3.497747in}{7.253186in}}%
\pgfusepath{clip}%
\pgfsetbuttcap%
\pgfsetroundjoin%
\definecolor{currentfill}{rgb}{0.226397,0.728888,0.462789}%
\pgfsetfillcolor{currentfill}%
\pgfsetlinewidth{1.003750pt}%
\definecolor{currentstroke}{rgb}{0.226397,0.728888,0.462789}%
\pgfsetstrokecolor{currentstroke}%
\pgfsetdash{}{0pt}%
\pgfpathmoveto{\pgfqpoint{3.961991in}{6.037346in}}%
\pgfpathcurveto{\pgfqpoint{3.967815in}{6.037346in}}{\pgfqpoint{3.973401in}{6.039660in}}{\pgfqpoint{3.977519in}{6.043778in}}%
\pgfpathcurveto{\pgfqpoint{3.981637in}{6.047896in}}{\pgfqpoint{3.983951in}{6.053482in}}{\pgfqpoint{3.983951in}{6.059306in}}%
\pgfpathcurveto{\pgfqpoint{3.983951in}{6.065130in}}{\pgfqpoint{3.981637in}{6.070716in}}{\pgfqpoint{3.977519in}{6.074835in}}%
\pgfpathcurveto{\pgfqpoint{3.973401in}{6.078953in}}{\pgfqpoint{3.967815in}{6.081267in}}{\pgfqpoint{3.961991in}{6.081267in}}%
\pgfpathcurveto{\pgfqpoint{3.956167in}{6.081267in}}{\pgfqpoint{3.950581in}{6.078953in}}{\pgfqpoint{3.946463in}{6.074835in}}%
\pgfpathcurveto{\pgfqpoint{3.942345in}{6.070716in}}{\pgfqpoint{3.940031in}{6.065130in}}{\pgfqpoint{3.940031in}{6.059306in}}%
\pgfpathcurveto{\pgfqpoint{3.940031in}{6.053482in}}{\pgfqpoint{3.942345in}{6.047896in}}{\pgfqpoint{3.946463in}{6.043778in}}%
\pgfpathcurveto{\pgfqpoint{3.950581in}{6.039660in}}{\pgfqpoint{3.956167in}{6.037346in}}{\pgfqpoint{3.961991in}{6.037346in}}%
\pgfpathlineto{\pgfqpoint{3.961991in}{6.037346in}}%
\pgfpathclose%
\pgfusepath{stroke,fill}%
\end{pgfscope}%
\begin{pgfscope}%
\pgfpathrectangle{\pgfqpoint{1.414736in}{0.548769in}}{\pgfqpoint{3.497747in}{7.253186in}}%
\pgfusepath{clip}%
\pgfsetbuttcap%
\pgfsetroundjoin%
\definecolor{currentfill}{rgb}{0.545524,0.838039,0.275626}%
\pgfsetfillcolor{currentfill}%
\pgfsetlinewidth{1.003750pt}%
\definecolor{currentstroke}{rgb}{0.545524,0.838039,0.275626}%
\pgfsetstrokecolor{currentstroke}%
\pgfsetdash{}{0pt}%
\pgfpathmoveto{\pgfqpoint{3.742488in}{6.508332in}}%
\pgfpathcurveto{\pgfqpoint{3.748312in}{6.508332in}}{\pgfqpoint{3.753898in}{6.510646in}}{\pgfqpoint{3.758016in}{6.514764in}}%
\pgfpathcurveto{\pgfqpoint{3.762134in}{6.518882in}}{\pgfqpoint{3.764448in}{6.524468in}}{\pgfqpoint{3.764448in}{6.530292in}}%
\pgfpathcurveto{\pgfqpoint{3.764448in}{6.536116in}}{\pgfqpoint{3.762134in}{6.541702in}}{\pgfqpoint{3.758016in}{6.545821in}}%
\pgfpathcurveto{\pgfqpoint{3.753898in}{6.549939in}}{\pgfqpoint{3.748312in}{6.552253in}}{\pgfqpoint{3.742488in}{6.552253in}}%
\pgfpathcurveto{\pgfqpoint{3.736664in}{6.552253in}}{\pgfqpoint{3.731078in}{6.549939in}}{\pgfqpoint{3.726960in}{6.545821in}}%
\pgfpathcurveto{\pgfqpoint{3.722841in}{6.541702in}}{\pgfqpoint{3.720528in}{6.536116in}}{\pgfqpoint{3.720528in}{6.530292in}}%
\pgfpathcurveto{\pgfqpoint{3.720528in}{6.524468in}}{\pgfqpoint{3.722841in}{6.518882in}}{\pgfqpoint{3.726960in}{6.514764in}}%
\pgfpathcurveto{\pgfqpoint{3.731078in}{6.510646in}}{\pgfqpoint{3.736664in}{6.508332in}}{\pgfqpoint{3.742488in}{6.508332in}}%
\pgfpathlineto{\pgfqpoint{3.742488in}{6.508332in}}%
\pgfpathclose%
\pgfusepath{stroke,fill}%
\end{pgfscope}%
\begin{pgfscope}%
\pgfpathrectangle{\pgfqpoint{1.414736in}{0.548769in}}{\pgfqpoint{3.497747in}{7.253186in}}%
\pgfusepath{clip}%
\pgfsetbuttcap%
\pgfsetroundjoin%
\definecolor{currentfill}{rgb}{0.585678,0.846661,0.249897}%
\pgfsetfillcolor{currentfill}%
\pgfsetlinewidth{1.003750pt}%
\definecolor{currentstroke}{rgb}{0.585678,0.846661,0.249897}%
\pgfsetstrokecolor{currentstroke}%
\pgfsetdash{}{0pt}%
\pgfpathmoveto{\pgfqpoint{3.814783in}{6.979318in}}%
\pgfpathcurveto{\pgfqpoint{3.820607in}{6.979318in}}{\pgfqpoint{3.826193in}{6.981632in}}{\pgfqpoint{3.830311in}{6.985750in}}%
\pgfpathcurveto{\pgfqpoint{3.834430in}{6.989868in}}{\pgfqpoint{3.836744in}{6.995455in}}{\pgfqpoint{3.836744in}{7.001278in}}%
\pgfpathcurveto{\pgfqpoint{3.836744in}{7.007102in}}{\pgfqpoint{3.834430in}{7.012689in}}{\pgfqpoint{3.830311in}{7.016807in}}%
\pgfpathcurveto{\pgfqpoint{3.826193in}{7.020925in}}{\pgfqpoint{3.820607in}{7.023239in}}{\pgfqpoint{3.814783in}{7.023239in}}%
\pgfpathcurveto{\pgfqpoint{3.808959in}{7.023239in}}{\pgfqpoint{3.803373in}{7.020925in}}{\pgfqpoint{3.799255in}{7.016807in}}%
\pgfpathcurveto{\pgfqpoint{3.795137in}{7.012689in}}{\pgfqpoint{3.792823in}{7.007102in}}{\pgfqpoint{3.792823in}{7.001278in}}%
\pgfpathcurveto{\pgfqpoint{3.792823in}{6.995455in}}{\pgfqpoint{3.795137in}{6.989868in}}{\pgfqpoint{3.799255in}{6.985750in}}%
\pgfpathcurveto{\pgfqpoint{3.803373in}{6.981632in}}{\pgfqpoint{3.808959in}{6.979318in}}{\pgfqpoint{3.814783in}{6.979318in}}%
\pgfpathlineto{\pgfqpoint{3.814783in}{6.979318in}}%
\pgfpathclose%
\pgfusepath{stroke,fill}%
\end{pgfscope}%
\begin{pgfscope}%
\pgfpathrectangle{\pgfqpoint{1.414736in}{0.548769in}}{\pgfqpoint{3.497747in}{7.253186in}}%
\pgfusepath{clip}%
\pgfsetbuttcap%
\pgfsetroundjoin%
\definecolor{currentfill}{rgb}{0.430983,0.808473,0.346476}%
\pgfsetfillcolor{currentfill}%
\pgfsetlinewidth{1.003750pt}%
\definecolor{currentstroke}{rgb}{0.430983,0.808473,0.346476}%
\pgfsetstrokecolor{currentstroke}%
\pgfsetdash{}{0pt}%
\pgfpathmoveto{\pgfqpoint{3.187237in}{7.450304in}}%
\pgfpathcurveto{\pgfqpoint{3.193061in}{7.450304in}}{\pgfqpoint{3.198647in}{7.452618in}}{\pgfqpoint{3.202765in}{7.456736in}}%
\pgfpathcurveto{\pgfqpoint{3.206883in}{7.460854in}}{\pgfqpoint{3.209197in}{7.466441in}}{\pgfqpoint{3.209197in}{7.472265in}}%
\pgfpathcurveto{\pgfqpoint{3.209197in}{7.478089in}}{\pgfqpoint{3.206883in}{7.483675in}}{\pgfqpoint{3.202765in}{7.487793in}}%
\pgfpathcurveto{\pgfqpoint{3.198647in}{7.491911in}}{\pgfqpoint{3.193061in}{7.494225in}}{\pgfqpoint{3.187237in}{7.494225in}}%
\pgfpathcurveto{\pgfqpoint{3.181413in}{7.494225in}}{\pgfqpoint{3.175827in}{7.491911in}}{\pgfqpoint{3.171708in}{7.487793in}}%
\pgfpathcurveto{\pgfqpoint{3.167590in}{7.483675in}}{\pgfqpoint{3.165276in}{7.478089in}}{\pgfqpoint{3.165276in}{7.472265in}}%
\pgfpathcurveto{\pgfqpoint{3.165276in}{7.466441in}}{\pgfqpoint{3.167590in}{7.460854in}}{\pgfqpoint{3.171708in}{7.456736in}}%
\pgfpathcurveto{\pgfqpoint{3.175827in}{7.452618in}}{\pgfqpoint{3.181413in}{7.450304in}}{\pgfqpoint{3.187237in}{7.450304in}}%
\pgfpathlineto{\pgfqpoint{3.187237in}{7.450304in}}%
\pgfpathclose%
\pgfusepath{stroke,fill}%
\end{pgfscope}%
\begin{pgfscope}%
\pgfpathrectangle{\pgfqpoint{1.414736in}{0.548769in}}{\pgfqpoint{3.497747in}{7.253186in}}%
\pgfusepath{clip}%
\pgfsetbuttcap%
\pgfsetroundjoin%
\definecolor{currentfill}{rgb}{0.993248,0.906157,0.143936}%
\pgfsetfillcolor{currentfill}%
\pgfsetlinewidth{1.003750pt}%
\definecolor{currentstroke}{rgb}{0.993248,0.906157,0.143936}%
\pgfsetstrokecolor{currentstroke}%
\pgfsetdash{}{0pt}%
\pgfpathmoveto{\pgfqpoint{3.873602in}{0.856499in}}%
\pgfpathcurveto{\pgfqpoint{3.879426in}{0.856499in}}{\pgfqpoint{3.885012in}{0.858813in}}{\pgfqpoint{3.889130in}{0.862931in}}%
\pgfpathcurveto{\pgfqpoint{3.893248in}{0.867049in}}{\pgfqpoint{3.895562in}{0.872635in}}{\pgfqpoint{3.895562in}{0.878459in}}%
\pgfpathcurveto{\pgfqpoint{3.895562in}{0.884283in}}{\pgfqpoint{3.893248in}{0.889869in}}{\pgfqpoint{3.889130in}{0.893987in}}%
\pgfpathcurveto{\pgfqpoint{3.885012in}{0.898105in}}{\pgfqpoint{3.879426in}{0.900419in}}{\pgfqpoint{3.873602in}{0.900419in}}%
\pgfpathcurveto{\pgfqpoint{3.867778in}{0.900419in}}{\pgfqpoint{3.862192in}{0.898105in}}{\pgfqpoint{3.858074in}{0.893987in}}%
\pgfpathcurveto{\pgfqpoint{3.853955in}{0.889869in}}{\pgfqpoint{3.851642in}{0.884283in}}{\pgfqpoint{3.851642in}{0.878459in}}%
\pgfpathcurveto{\pgfqpoint{3.851642in}{0.872635in}}{\pgfqpoint{3.853955in}{0.867049in}}{\pgfqpoint{3.858074in}{0.862931in}}%
\pgfpathcurveto{\pgfqpoint{3.862192in}{0.858813in}}{\pgfqpoint{3.867778in}{0.856499in}}{\pgfqpoint{3.873602in}{0.856499in}}%
\pgfpathlineto{\pgfqpoint{3.873602in}{0.856499in}}%
\pgfpathclose%
\pgfusepath{stroke,fill}%
\end{pgfscope}%
\begin{pgfscope}%
\pgfpathrectangle{\pgfqpoint{1.414736in}{0.548769in}}{\pgfqpoint{3.497747in}{7.253186in}}%
\pgfusepath{clip}%
\pgfsetbuttcap%
\pgfsetroundjoin%
\definecolor{currentfill}{rgb}{0.636902,0.856542,0.216620}%
\pgfsetfillcolor{currentfill}%
\pgfsetlinewidth{1.003750pt}%
\definecolor{currentstroke}{rgb}{0.636902,0.856542,0.216620}%
\pgfsetstrokecolor{currentstroke}%
\pgfsetdash{}{0pt}%
\pgfpathmoveto{\pgfqpoint{3.382630in}{1.327485in}}%
\pgfpathcurveto{\pgfqpoint{3.388454in}{1.327485in}}{\pgfqpoint{3.394040in}{1.329799in}}{\pgfqpoint{3.398158in}{1.333917in}}%
\pgfpathcurveto{\pgfqpoint{3.402277in}{1.338035in}}{\pgfqpoint{3.404590in}{1.343621in}}{\pgfqpoint{3.404590in}{1.349445in}}%
\pgfpathcurveto{\pgfqpoint{3.404590in}{1.355269in}}{\pgfqpoint{3.402277in}{1.360855in}}{\pgfqpoint{3.398158in}{1.364973in}}%
\pgfpathcurveto{\pgfqpoint{3.394040in}{1.369092in}}{\pgfqpoint{3.388454in}{1.371405in}}{\pgfqpoint{3.382630in}{1.371405in}}%
\pgfpathcurveto{\pgfqpoint{3.376806in}{1.371405in}}{\pgfqpoint{3.371220in}{1.369092in}}{\pgfqpoint{3.367102in}{1.364973in}}%
\pgfpathcurveto{\pgfqpoint{3.362984in}{1.360855in}}{\pgfqpoint{3.360670in}{1.355269in}}{\pgfqpoint{3.360670in}{1.349445in}}%
\pgfpathcurveto{\pgfqpoint{3.360670in}{1.343621in}}{\pgfqpoint{3.362984in}{1.338035in}}{\pgfqpoint{3.367102in}{1.333917in}}%
\pgfpathcurveto{\pgfqpoint{3.371220in}{1.329799in}}{\pgfqpoint{3.376806in}{1.327485in}}{\pgfqpoint{3.382630in}{1.327485in}}%
\pgfpathlineto{\pgfqpoint{3.382630in}{1.327485in}}%
\pgfpathclose%
\pgfusepath{stroke,fill}%
\end{pgfscope}%
\begin{pgfscope}%
\pgfpathrectangle{\pgfqpoint{1.414736in}{0.548769in}}{\pgfqpoint{3.497747in}{7.253186in}}%
\pgfusepath{clip}%
\pgfsetbuttcap%
\pgfsetroundjoin%
\definecolor{currentfill}{rgb}{0.412913,0.803041,0.357269}%
\pgfsetfillcolor{currentfill}%
\pgfsetlinewidth{1.003750pt}%
\definecolor{currentstroke}{rgb}{0.412913,0.803041,0.357269}%
\pgfsetstrokecolor{currentstroke}%
\pgfsetdash{}{0pt}%
\pgfpathmoveto{\pgfqpoint{3.511846in}{1.798471in}}%
\pgfpathcurveto{\pgfqpoint{3.517670in}{1.798471in}}{\pgfqpoint{3.523256in}{1.800785in}}{\pgfqpoint{3.527374in}{1.804903in}}%
\pgfpathcurveto{\pgfqpoint{3.531492in}{1.809021in}}{\pgfqpoint{3.533806in}{1.814607in}}{\pgfqpoint{3.533806in}{1.820431in}}%
\pgfpathcurveto{\pgfqpoint{3.533806in}{1.826255in}}{\pgfqpoint{3.531492in}{1.831841in}}{\pgfqpoint{3.527374in}{1.835960in}}%
\pgfpathcurveto{\pgfqpoint{3.523256in}{1.840078in}}{\pgfqpoint{3.517670in}{1.842392in}}{\pgfqpoint{3.511846in}{1.842392in}}%
\pgfpathcurveto{\pgfqpoint{3.506022in}{1.842392in}}{\pgfqpoint{3.500436in}{1.840078in}}{\pgfqpoint{3.496317in}{1.835960in}}%
\pgfpathcurveto{\pgfqpoint{3.492199in}{1.831841in}}{\pgfqpoint{3.489885in}{1.826255in}}{\pgfqpoint{3.489885in}{1.820431in}}%
\pgfpathcurveto{\pgfqpoint{3.489885in}{1.814607in}}{\pgfqpoint{3.492199in}{1.809021in}}{\pgfqpoint{3.496317in}{1.804903in}}%
\pgfpathcurveto{\pgfqpoint{3.500436in}{1.800785in}}{\pgfqpoint{3.506022in}{1.798471in}}{\pgfqpoint{3.511846in}{1.798471in}}%
\pgfpathlineto{\pgfqpoint{3.511846in}{1.798471in}}%
\pgfpathclose%
\pgfusepath{stroke,fill}%
\end{pgfscope}%
\begin{pgfscope}%
\pgfpathrectangle{\pgfqpoint{1.414736in}{0.548769in}}{\pgfqpoint{3.497747in}{7.253186in}}%
\pgfusepath{clip}%
\pgfsetbuttcap%
\pgfsetroundjoin%
\definecolor{currentfill}{rgb}{0.993248,0.906157,0.143936}%
\pgfsetfillcolor{currentfill}%
\pgfsetlinewidth{1.003750pt}%
\definecolor{currentstroke}{rgb}{0.993248,0.906157,0.143936}%
\pgfsetstrokecolor{currentstroke}%
\pgfsetdash{}{0pt}%
\pgfpathmoveto{\pgfqpoint{3.133681in}{2.269457in}}%
\pgfpathcurveto{\pgfqpoint{3.139505in}{2.269457in}}{\pgfqpoint{3.145091in}{2.271771in}}{\pgfqpoint{3.149210in}{2.275889in}}%
\pgfpathcurveto{\pgfqpoint{3.153328in}{2.280007in}}{\pgfqpoint{3.155642in}{2.285593in}}{\pgfqpoint{3.155642in}{2.291417in}}%
\pgfpathcurveto{\pgfqpoint{3.155642in}{2.297241in}}{\pgfqpoint{3.153328in}{2.302828in}}{\pgfqpoint{3.149210in}{2.306946in}}%
\pgfpathcurveto{\pgfqpoint{3.145091in}{2.311064in}}{\pgfqpoint{3.139505in}{2.313378in}}{\pgfqpoint{3.133681in}{2.313378in}}%
\pgfpathcurveto{\pgfqpoint{3.127857in}{2.313378in}}{\pgfqpoint{3.122271in}{2.311064in}}{\pgfqpoint{3.118153in}{2.306946in}}%
\pgfpathcurveto{\pgfqpoint{3.114035in}{2.302828in}}{\pgfqpoint{3.111721in}{2.297241in}}{\pgfqpoint{3.111721in}{2.291417in}}%
\pgfpathcurveto{\pgfqpoint{3.111721in}{2.285593in}}{\pgfqpoint{3.114035in}{2.280007in}}{\pgfqpoint{3.118153in}{2.275889in}}%
\pgfpathcurveto{\pgfqpoint{3.122271in}{2.271771in}}{\pgfqpoint{3.127857in}{2.269457in}}{\pgfqpoint{3.133681in}{2.269457in}}%
\pgfpathlineto{\pgfqpoint{3.133681in}{2.269457in}}%
\pgfpathclose%
\pgfusepath{stroke,fill}%
\end{pgfscope}%
\begin{pgfscope}%
\pgfpathrectangle{\pgfqpoint{1.414736in}{0.548769in}}{\pgfqpoint{3.497747in}{7.253186in}}%
\pgfusepath{clip}%
\pgfsetbuttcap%
\pgfsetroundjoin%
\definecolor{currentfill}{rgb}{0.657642,0.860219,0.203082}%
\pgfsetfillcolor{currentfill}%
\pgfsetlinewidth{1.003750pt}%
\definecolor{currentstroke}{rgb}{0.657642,0.860219,0.203082}%
\pgfsetstrokecolor{currentstroke}%
\pgfsetdash{}{0pt}%
\pgfpathmoveto{\pgfqpoint{3.488884in}{2.740443in}}%
\pgfpathcurveto{\pgfqpoint{3.494707in}{2.740443in}}{\pgfqpoint{3.500294in}{2.742757in}}{\pgfqpoint{3.504412in}{2.746875in}}%
\pgfpathcurveto{\pgfqpoint{3.508530in}{2.750993in}}{\pgfqpoint{3.510844in}{2.756580in}}{\pgfqpoint{3.510844in}{2.762404in}}%
\pgfpathcurveto{\pgfqpoint{3.510844in}{2.768227in}}{\pgfqpoint{3.508530in}{2.773814in}}{\pgfqpoint{3.504412in}{2.777932in}}%
\pgfpathcurveto{\pgfqpoint{3.500294in}{2.782050in}}{\pgfqpoint{3.494707in}{2.784364in}}{\pgfqpoint{3.488884in}{2.784364in}}%
\pgfpathcurveto{\pgfqpoint{3.483060in}{2.784364in}}{\pgfqpoint{3.477473in}{2.782050in}}{\pgfqpoint{3.473355in}{2.777932in}}%
\pgfpathcurveto{\pgfqpoint{3.469237in}{2.773814in}}{\pgfqpoint{3.466923in}{2.768227in}}{\pgfqpoint{3.466923in}{2.762404in}}%
\pgfpathcurveto{\pgfqpoint{3.466923in}{2.756580in}}{\pgfqpoint{3.469237in}{2.750993in}}{\pgfqpoint{3.473355in}{2.746875in}}%
\pgfpathcurveto{\pgfqpoint{3.477473in}{2.742757in}}{\pgfqpoint{3.483060in}{2.740443in}}{\pgfqpoint{3.488884in}{2.740443in}}%
\pgfpathlineto{\pgfqpoint{3.488884in}{2.740443in}}%
\pgfpathclose%
\pgfusepath{stroke,fill}%
\end{pgfscope}%
\begin{pgfscope}%
\pgfpathrectangle{\pgfqpoint{1.414736in}{0.548769in}}{\pgfqpoint{3.497747in}{7.253186in}}%
\pgfusepath{clip}%
\pgfsetbuttcap%
\pgfsetroundjoin%
\definecolor{currentfill}{rgb}{0.955300,0.901065,0.118128}%
\pgfsetfillcolor{currentfill}%
\pgfsetlinewidth{1.003750pt}%
\definecolor{currentstroke}{rgb}{0.955300,0.901065,0.118128}%
\pgfsetstrokecolor{currentstroke}%
\pgfsetdash{}{0pt}%
\pgfpathmoveto{\pgfqpoint{3.965606in}{3.211429in}}%
\pgfpathcurveto{\pgfqpoint{3.971429in}{3.211429in}}{\pgfqpoint{3.977016in}{3.213743in}}{\pgfqpoint{3.981134in}{3.217861in}}%
\pgfpathcurveto{\pgfqpoint{3.985252in}{3.221980in}}{\pgfqpoint{3.987566in}{3.227566in}}{\pgfqpoint{3.987566in}{3.233390in}}%
\pgfpathcurveto{\pgfqpoint{3.987566in}{3.239214in}}{\pgfqpoint{3.985252in}{3.244800in}}{\pgfqpoint{3.981134in}{3.248918in}}%
\pgfpathcurveto{\pgfqpoint{3.977016in}{3.253036in}}{\pgfqpoint{3.971429in}{3.255350in}}{\pgfqpoint{3.965606in}{3.255350in}}%
\pgfpathcurveto{\pgfqpoint{3.959782in}{3.255350in}}{\pgfqpoint{3.954195in}{3.253036in}}{\pgfqpoint{3.950077in}{3.248918in}}%
\pgfpathcurveto{\pgfqpoint{3.945959in}{3.244800in}}{\pgfqpoint{3.943645in}{3.239214in}}{\pgfqpoint{3.943645in}{3.233390in}}%
\pgfpathcurveto{\pgfqpoint{3.943645in}{3.227566in}}{\pgfqpoint{3.945959in}{3.221980in}}{\pgfqpoint{3.950077in}{3.217861in}}%
\pgfpathcurveto{\pgfqpoint{3.954195in}{3.213743in}}{\pgfqpoint{3.959782in}{3.211429in}}{\pgfqpoint{3.965606in}{3.211429in}}%
\pgfpathlineto{\pgfqpoint{3.965606in}{3.211429in}}%
\pgfpathclose%
\pgfusepath{stroke,fill}%
\end{pgfscope}%
\begin{pgfscope}%
\pgfpathrectangle{\pgfqpoint{1.414736in}{0.548769in}}{\pgfqpoint{3.497747in}{7.253186in}}%
\pgfusepath{clip}%
\pgfsetbuttcap%
\pgfsetroundjoin%
\definecolor{currentfill}{rgb}{0.616293,0.852709,0.230052}%
\pgfsetfillcolor{currentfill}%
\pgfsetlinewidth{1.003750pt}%
\definecolor{currentstroke}{rgb}{0.616293,0.852709,0.230052}%
\pgfsetstrokecolor{currentstroke}%
\pgfsetdash{}{0pt}%
\pgfpathmoveto{\pgfqpoint{3.917497in}{3.682415in}}%
\pgfpathcurveto{\pgfqpoint{3.923321in}{3.682415in}}{\pgfqpoint{3.928908in}{3.684729in}}{\pgfqpoint{3.933026in}{3.688847in}}%
\pgfpathcurveto{\pgfqpoint{3.937144in}{3.692966in}}{\pgfqpoint{3.939458in}{3.698552in}}{\pgfqpoint{3.939458in}{3.704376in}}%
\pgfpathcurveto{\pgfqpoint{3.939458in}{3.710200in}}{\pgfqpoint{3.937144in}{3.715786in}}{\pgfqpoint{3.933026in}{3.719904in}}%
\pgfpathcurveto{\pgfqpoint{3.928908in}{3.724022in}}{\pgfqpoint{3.923321in}{3.726336in}}{\pgfqpoint{3.917497in}{3.726336in}}%
\pgfpathcurveto{\pgfqpoint{3.911673in}{3.726336in}}{\pgfqpoint{3.906087in}{3.724022in}}{\pgfqpoint{3.901969in}{3.719904in}}%
\pgfpathcurveto{\pgfqpoint{3.897851in}{3.715786in}}{\pgfqpoint{3.895537in}{3.710200in}}{\pgfqpoint{3.895537in}{3.704376in}}%
\pgfpathcurveto{\pgfqpoint{3.895537in}{3.698552in}}{\pgfqpoint{3.897851in}{3.692966in}}{\pgfqpoint{3.901969in}{3.688847in}}%
\pgfpathcurveto{\pgfqpoint{3.906087in}{3.684729in}}{\pgfqpoint{3.911673in}{3.682415in}}{\pgfqpoint{3.917497in}{3.682415in}}%
\pgfpathlineto{\pgfqpoint{3.917497in}{3.682415in}}%
\pgfpathclose%
\pgfusepath{stroke,fill}%
\end{pgfscope}%
\begin{pgfscope}%
\pgfpathrectangle{\pgfqpoint{1.414736in}{0.548769in}}{\pgfqpoint{3.497747in}{7.253186in}}%
\pgfusepath{clip}%
\pgfsetbuttcap%
\pgfsetroundjoin%
\definecolor{currentfill}{rgb}{0.926106,0.897330,0.104071}%
\pgfsetfillcolor{currentfill}%
\pgfsetlinewidth{1.003750pt}%
\definecolor{currentstroke}{rgb}{0.926106,0.897330,0.104071}%
\pgfsetstrokecolor{currentstroke}%
\pgfsetdash{}{0pt}%
\pgfpathmoveto{\pgfqpoint{2.827642in}{4.153402in}}%
\pgfpathcurveto{\pgfqpoint{2.833466in}{4.153402in}}{\pgfqpoint{2.839052in}{4.155715in}}{\pgfqpoint{2.843170in}{4.159834in}}%
\pgfpathcurveto{\pgfqpoint{2.847288in}{4.163952in}}{\pgfqpoint{2.849602in}{4.169538in}}{\pgfqpoint{2.849602in}{4.175362in}}%
\pgfpathcurveto{\pgfqpoint{2.849602in}{4.181186in}}{\pgfqpoint{2.847288in}{4.186772in}}{\pgfqpoint{2.843170in}{4.190890in}}%
\pgfpathcurveto{\pgfqpoint{2.839052in}{4.195008in}}{\pgfqpoint{2.833466in}{4.197322in}}{\pgfqpoint{2.827642in}{4.197322in}}%
\pgfpathcurveto{\pgfqpoint{2.821818in}{4.197322in}}{\pgfqpoint{2.816232in}{4.195008in}}{\pgfqpoint{2.812113in}{4.190890in}}%
\pgfpathcurveto{\pgfqpoint{2.807995in}{4.186772in}}{\pgfqpoint{2.805681in}{4.181186in}}{\pgfqpoint{2.805681in}{4.175362in}}%
\pgfpathcurveto{\pgfqpoint{2.805681in}{4.169538in}}{\pgfqpoint{2.807995in}{4.163952in}}{\pgfqpoint{2.812113in}{4.159834in}}%
\pgfpathcurveto{\pgfqpoint{2.816232in}{4.155715in}}{\pgfqpoint{2.821818in}{4.153402in}}{\pgfqpoint{2.827642in}{4.153402in}}%
\pgfpathlineto{\pgfqpoint{2.827642in}{4.153402in}}%
\pgfpathclose%
\pgfusepath{stroke,fill}%
\end{pgfscope}%
\begin{pgfscope}%
\pgfpathrectangle{\pgfqpoint{1.414736in}{0.548769in}}{\pgfqpoint{3.497747in}{7.253186in}}%
\pgfusepath{clip}%
\pgfsetbuttcap%
\pgfsetroundjoin%
\definecolor{currentfill}{rgb}{0.866013,0.889868,0.095953}%
\pgfsetfillcolor{currentfill}%
\pgfsetlinewidth{1.003750pt}%
\definecolor{currentstroke}{rgb}{0.866013,0.889868,0.095953}%
\pgfsetstrokecolor{currentstroke}%
\pgfsetdash{}{0pt}%
\pgfpathmoveto{\pgfqpoint{3.803366in}{4.624388in}}%
\pgfpathcurveto{\pgfqpoint{3.809189in}{4.624388in}}{\pgfqpoint{3.814776in}{4.626702in}}{\pgfqpoint{3.818894in}{4.630820in}}%
\pgfpathcurveto{\pgfqpoint{3.823012in}{4.634938in}}{\pgfqpoint{3.825326in}{4.640524in}}{\pgfqpoint{3.825326in}{4.646348in}}%
\pgfpathcurveto{\pgfqpoint{3.825326in}{4.652172in}}{\pgfqpoint{3.823012in}{4.657758in}}{\pgfqpoint{3.818894in}{4.661876in}}%
\pgfpathcurveto{\pgfqpoint{3.814776in}{4.665994in}}{\pgfqpoint{3.809189in}{4.668308in}}{\pgfqpoint{3.803366in}{4.668308in}}%
\pgfpathcurveto{\pgfqpoint{3.797542in}{4.668308in}}{\pgfqpoint{3.791955in}{4.665994in}}{\pgfqpoint{3.787837in}{4.661876in}}%
\pgfpathcurveto{\pgfqpoint{3.783719in}{4.657758in}}{\pgfqpoint{3.781405in}{4.652172in}}{\pgfqpoint{3.781405in}{4.646348in}}%
\pgfpathcurveto{\pgfqpoint{3.781405in}{4.640524in}}{\pgfqpoint{3.783719in}{4.634938in}}{\pgfqpoint{3.787837in}{4.630820in}}%
\pgfpathcurveto{\pgfqpoint{3.791955in}{4.626702in}}{\pgfqpoint{3.797542in}{4.624388in}}{\pgfqpoint{3.803366in}{4.624388in}}%
\pgfpathlineto{\pgfqpoint{3.803366in}{4.624388in}}%
\pgfpathclose%
\pgfusepath{stroke,fill}%
\end{pgfscope}%
\begin{pgfscope}%
\pgfpathrectangle{\pgfqpoint{1.414736in}{0.548769in}}{\pgfqpoint{3.497747in}{7.253186in}}%
\pgfusepath{clip}%
\pgfsetbuttcap%
\pgfsetroundjoin%
\definecolor{currentfill}{rgb}{0.165117,0.467423,0.558141}%
\pgfsetfillcolor{currentfill}%
\pgfsetlinewidth{1.003750pt}%
\definecolor{currentstroke}{rgb}{0.165117,0.467423,0.558141}%
\pgfsetstrokecolor{currentstroke}%
\pgfsetdash{}{0pt}%
\pgfpathmoveto{\pgfqpoint{4.682758in}{5.095374in}}%
\pgfpathcurveto{\pgfqpoint{4.688582in}{5.095374in}}{\pgfqpoint{4.694168in}{5.097688in}}{\pgfqpoint{4.698286in}{5.101806in}}%
\pgfpathcurveto{\pgfqpoint{4.702404in}{5.105924in}}{\pgfqpoint{4.704718in}{5.111510in}}{\pgfqpoint{4.704718in}{5.117334in}}%
\pgfpathcurveto{\pgfqpoint{4.704718in}{5.123158in}}{\pgfqpoint{4.702404in}{5.128744in}}{\pgfqpoint{4.698286in}{5.132862in}}%
\pgfpathcurveto{\pgfqpoint{4.694168in}{5.136980in}}{\pgfqpoint{4.688582in}{5.139294in}}{\pgfqpoint{4.682758in}{5.139294in}}%
\pgfpathcurveto{\pgfqpoint{4.676934in}{5.139294in}}{\pgfqpoint{4.671348in}{5.136980in}}{\pgfqpoint{4.667229in}{5.132862in}}%
\pgfpathcurveto{\pgfqpoint{4.663111in}{5.128744in}}{\pgfqpoint{4.660797in}{5.123158in}}{\pgfqpoint{4.660797in}{5.117334in}}%
\pgfpathcurveto{\pgfqpoint{4.660797in}{5.111510in}}{\pgfqpoint{4.663111in}{5.105924in}}{\pgfqpoint{4.667229in}{5.101806in}}%
\pgfpathcurveto{\pgfqpoint{4.671348in}{5.097688in}}{\pgfqpoint{4.676934in}{5.095374in}}{\pgfqpoint{4.682758in}{5.095374in}}%
\pgfpathlineto{\pgfqpoint{4.682758in}{5.095374in}}%
\pgfpathclose%
\pgfusepath{stroke,fill}%
\end{pgfscope}%
\begin{pgfscope}%
\pgfpathrectangle{\pgfqpoint{1.414736in}{0.548769in}}{\pgfqpoint{3.497747in}{7.253186in}}%
\pgfusepath{clip}%
\pgfsetbuttcap%
\pgfsetroundjoin%
\definecolor{currentfill}{rgb}{0.166617,0.463708,0.558119}%
\pgfsetfillcolor{currentfill}%
\pgfsetlinewidth{1.003750pt}%
\definecolor{currentstroke}{rgb}{0.166617,0.463708,0.558119}%
\pgfsetstrokecolor{currentstroke}%
\pgfsetdash{}{0pt}%
\pgfpathmoveto{\pgfqpoint{3.575352in}{5.566360in}}%
\pgfpathcurveto{\pgfqpoint{3.581176in}{5.566360in}}{\pgfqpoint{3.586762in}{5.568674in}}{\pgfqpoint{3.590880in}{5.572792in}}%
\pgfpathcurveto{\pgfqpoint{3.594998in}{5.576910in}}{\pgfqpoint{3.597312in}{5.582496in}}{\pgfqpoint{3.597312in}{5.588320in}}%
\pgfpathcurveto{\pgfqpoint{3.597312in}{5.594144in}}{\pgfqpoint{3.594998in}{5.599730in}}{\pgfqpoint{3.590880in}{5.603848in}}%
\pgfpathcurveto{\pgfqpoint{3.586762in}{5.607967in}}{\pgfqpoint{3.581176in}{5.610280in}}{\pgfqpoint{3.575352in}{5.610280in}}%
\pgfpathcurveto{\pgfqpoint{3.569528in}{5.610280in}}{\pgfqpoint{3.563942in}{5.607967in}}{\pgfqpoint{3.559824in}{5.603848in}}%
\pgfpathcurveto{\pgfqpoint{3.555706in}{5.599730in}}{\pgfqpoint{3.553392in}{5.594144in}}{\pgfqpoint{3.553392in}{5.588320in}}%
\pgfpathcurveto{\pgfqpoint{3.553392in}{5.582496in}}{\pgfqpoint{3.555706in}{5.576910in}}{\pgfqpoint{3.559824in}{5.572792in}}%
\pgfpathcurveto{\pgfqpoint{3.563942in}{5.568674in}}{\pgfqpoint{3.569528in}{5.566360in}}{\pgfqpoint{3.575352in}{5.566360in}}%
\pgfpathlineto{\pgfqpoint{3.575352in}{5.566360in}}%
\pgfpathclose%
\pgfusepath{stroke,fill}%
\end{pgfscope}%
\begin{pgfscope}%
\pgfpathrectangle{\pgfqpoint{1.414736in}{0.548769in}}{\pgfqpoint{3.497747in}{7.253186in}}%
\pgfusepath{clip}%
\pgfsetbuttcap%
\pgfsetroundjoin%
\definecolor{currentfill}{rgb}{0.824940,0.884720,0.106217}%
\pgfsetfillcolor{currentfill}%
\pgfsetlinewidth{1.003750pt}%
\definecolor{currentstroke}{rgb}{0.824940,0.884720,0.106217}%
\pgfsetstrokecolor{currentstroke}%
\pgfsetdash{}{0pt}%
\pgfpathmoveto{\pgfqpoint{3.031322in}{6.037346in}}%
\pgfpathcurveto{\pgfqpoint{3.037146in}{6.037346in}}{\pgfqpoint{3.042732in}{6.039660in}}{\pgfqpoint{3.046851in}{6.043778in}}%
\pgfpathcurveto{\pgfqpoint{3.050969in}{6.047896in}}{\pgfqpoint{3.053283in}{6.053482in}}{\pgfqpoint{3.053283in}{6.059306in}}%
\pgfpathcurveto{\pgfqpoint{3.053283in}{6.065130in}}{\pgfqpoint{3.050969in}{6.070716in}}{\pgfqpoint{3.046851in}{6.074835in}}%
\pgfpathcurveto{\pgfqpoint{3.042732in}{6.078953in}}{\pgfqpoint{3.037146in}{6.081267in}}{\pgfqpoint{3.031322in}{6.081267in}}%
\pgfpathcurveto{\pgfqpoint{3.025498in}{6.081267in}}{\pgfqpoint{3.019912in}{6.078953in}}{\pgfqpoint{3.015794in}{6.074835in}}%
\pgfpathcurveto{\pgfqpoint{3.011676in}{6.070716in}}{\pgfqpoint{3.009362in}{6.065130in}}{\pgfqpoint{3.009362in}{6.059306in}}%
\pgfpathcurveto{\pgfqpoint{3.009362in}{6.053482in}}{\pgfqpoint{3.011676in}{6.047896in}}{\pgfqpoint{3.015794in}{6.043778in}}%
\pgfpathcurveto{\pgfqpoint{3.019912in}{6.039660in}}{\pgfqpoint{3.025498in}{6.037346in}}{\pgfqpoint{3.031322in}{6.037346in}}%
\pgfpathlineto{\pgfqpoint{3.031322in}{6.037346in}}%
\pgfpathclose%
\pgfusepath{stroke,fill}%
\end{pgfscope}%
\begin{pgfscope}%
\pgfpathrectangle{\pgfqpoint{1.414736in}{0.548769in}}{\pgfqpoint{3.497747in}{7.253186in}}%
\pgfusepath{clip}%
\pgfsetbuttcap%
\pgfsetroundjoin%
\definecolor{currentfill}{rgb}{0.993248,0.906157,0.143936}%
\pgfsetfillcolor{currentfill}%
\pgfsetlinewidth{1.003750pt}%
\definecolor{currentstroke}{rgb}{0.993248,0.906157,0.143936}%
\pgfsetstrokecolor{currentstroke}%
\pgfsetdash{}{0pt}%
\pgfpathmoveto{\pgfqpoint{3.925664in}{6.508332in}}%
\pgfpathcurveto{\pgfqpoint{3.931488in}{6.508332in}}{\pgfqpoint{3.937074in}{6.510646in}}{\pgfqpoint{3.941192in}{6.514764in}}%
\pgfpathcurveto{\pgfqpoint{3.945311in}{6.518882in}}{\pgfqpoint{3.947624in}{6.524468in}}{\pgfqpoint{3.947624in}{6.530292in}}%
\pgfpathcurveto{\pgfqpoint{3.947624in}{6.536116in}}{\pgfqpoint{3.945311in}{6.541702in}}{\pgfqpoint{3.941192in}{6.545821in}}%
\pgfpathcurveto{\pgfqpoint{3.937074in}{6.549939in}}{\pgfqpoint{3.931488in}{6.552253in}}{\pgfqpoint{3.925664in}{6.552253in}}%
\pgfpathcurveto{\pgfqpoint{3.919840in}{6.552253in}}{\pgfqpoint{3.914254in}{6.549939in}}{\pgfqpoint{3.910136in}{6.545821in}}%
\pgfpathcurveto{\pgfqpoint{3.906018in}{6.541702in}}{\pgfqpoint{3.903704in}{6.536116in}}{\pgfqpoint{3.903704in}{6.530292in}}%
\pgfpathcurveto{\pgfqpoint{3.903704in}{6.524468in}}{\pgfqpoint{3.906018in}{6.518882in}}{\pgfqpoint{3.910136in}{6.514764in}}%
\pgfpathcurveto{\pgfqpoint{3.914254in}{6.510646in}}{\pgfqpoint{3.919840in}{6.508332in}}{\pgfqpoint{3.925664in}{6.508332in}}%
\pgfpathlineto{\pgfqpoint{3.925664in}{6.508332in}}%
\pgfpathclose%
\pgfusepath{stroke,fill}%
\end{pgfscope}%
\begin{pgfscope}%
\pgfpathrectangle{\pgfqpoint{1.414736in}{0.548769in}}{\pgfqpoint{3.497747in}{7.253186in}}%
\pgfusepath{clip}%
\pgfsetbuttcap%
\pgfsetroundjoin%
\definecolor{currentfill}{rgb}{0.741388,0.873449,0.149561}%
\pgfsetfillcolor{currentfill}%
\pgfsetlinewidth{1.003750pt}%
\definecolor{currentstroke}{rgb}{0.741388,0.873449,0.149561}%
\pgfsetstrokecolor{currentstroke}%
\pgfsetdash{}{0pt}%
\pgfpathmoveto{\pgfqpoint{3.227858in}{6.979318in}}%
\pgfpathcurveto{\pgfqpoint{3.233682in}{6.979318in}}{\pgfqpoint{3.239268in}{6.981632in}}{\pgfqpoint{3.243386in}{6.985750in}}%
\pgfpathcurveto{\pgfqpoint{3.247504in}{6.989868in}}{\pgfqpoint{3.249818in}{6.995455in}}{\pgfqpoint{3.249818in}{7.001278in}}%
\pgfpathcurveto{\pgfqpoint{3.249818in}{7.007102in}}{\pgfqpoint{3.247504in}{7.012689in}}{\pgfqpoint{3.243386in}{7.016807in}}%
\pgfpathcurveto{\pgfqpoint{3.239268in}{7.020925in}}{\pgfqpoint{3.233682in}{7.023239in}}{\pgfqpoint{3.227858in}{7.023239in}}%
\pgfpathcurveto{\pgfqpoint{3.222034in}{7.023239in}}{\pgfqpoint{3.216448in}{7.020925in}}{\pgfqpoint{3.212329in}{7.016807in}}%
\pgfpathcurveto{\pgfqpoint{3.208211in}{7.012689in}}{\pgfqpoint{3.205897in}{7.007102in}}{\pgfqpoint{3.205897in}{7.001278in}}%
\pgfpathcurveto{\pgfqpoint{3.205897in}{6.995455in}}{\pgfqpoint{3.208211in}{6.989868in}}{\pgfqpoint{3.212329in}{6.985750in}}%
\pgfpathcurveto{\pgfqpoint{3.216448in}{6.981632in}}{\pgfqpoint{3.222034in}{6.979318in}}{\pgfqpoint{3.227858in}{6.979318in}}%
\pgfpathlineto{\pgfqpoint{3.227858in}{6.979318in}}%
\pgfpathclose%
\pgfusepath{stroke,fill}%
\end{pgfscope}%
\begin{pgfscope}%
\pgfpathrectangle{\pgfqpoint{1.414736in}{0.548769in}}{\pgfqpoint{3.497747in}{7.253186in}}%
\pgfusepath{clip}%
\pgfsetbuttcap%
\pgfsetroundjoin%
\definecolor{currentfill}{rgb}{0.983868,0.904867,0.136897}%
\pgfsetfillcolor{currentfill}%
\pgfsetlinewidth{1.003750pt}%
\definecolor{currentstroke}{rgb}{0.983868,0.904867,0.136897}%
\pgfsetstrokecolor{currentstroke}%
\pgfsetdash{}{0pt}%
\pgfpathmoveto{\pgfqpoint{2.578998in}{7.450304in}}%
\pgfpathcurveto{\pgfqpoint{2.584822in}{7.450304in}}{\pgfqpoint{2.590408in}{7.452618in}}{\pgfqpoint{2.594526in}{7.456736in}}%
\pgfpathcurveto{\pgfqpoint{2.598644in}{7.460854in}}{\pgfqpoint{2.600958in}{7.466441in}}{\pgfqpoint{2.600958in}{7.472265in}}%
\pgfpathcurveto{\pgfqpoint{2.600958in}{7.478089in}}{\pgfqpoint{2.598644in}{7.483675in}}{\pgfqpoint{2.594526in}{7.487793in}}%
\pgfpathcurveto{\pgfqpoint{2.590408in}{7.491911in}}{\pgfqpoint{2.584822in}{7.494225in}}{\pgfqpoint{2.578998in}{7.494225in}}%
\pgfpathcurveto{\pgfqpoint{2.573174in}{7.494225in}}{\pgfqpoint{2.567588in}{7.491911in}}{\pgfqpoint{2.563470in}{7.487793in}}%
\pgfpathcurveto{\pgfqpoint{2.559351in}{7.483675in}}{\pgfqpoint{2.557038in}{7.478089in}}{\pgfqpoint{2.557038in}{7.472265in}}%
\pgfpathcurveto{\pgfqpoint{2.557038in}{7.466441in}}{\pgfqpoint{2.559351in}{7.460854in}}{\pgfqpoint{2.563470in}{7.456736in}}%
\pgfpathcurveto{\pgfqpoint{2.567588in}{7.452618in}}{\pgfqpoint{2.573174in}{7.450304in}}{\pgfqpoint{2.578998in}{7.450304in}}%
\pgfpathlineto{\pgfqpoint{2.578998in}{7.450304in}}%
\pgfpathclose%
\pgfusepath{stroke,fill}%
\end{pgfscope}%
\begin{pgfscope}%
\pgfpathrectangle{\pgfqpoint{1.414736in}{0.548769in}}{\pgfqpoint{3.497747in}{7.253186in}}%
\pgfusepath{clip}%
\pgfsetbuttcap%
\pgfsetroundjoin%
\definecolor{currentfill}{rgb}{0.699415,0.867117,0.175971}%
\pgfsetfillcolor{currentfill}%
\pgfsetlinewidth{1.003750pt}%
\definecolor{currentstroke}{rgb}{0.699415,0.867117,0.175971}%
\pgfsetstrokecolor{currentstroke}%
\pgfsetdash{}{0pt}%
\pgfpathmoveto{\pgfqpoint{3.915900in}{0.856499in}}%
\pgfpathcurveto{\pgfqpoint{3.921724in}{0.856499in}}{\pgfqpoint{3.927310in}{0.858813in}}{\pgfqpoint{3.931428in}{0.862931in}}%
\pgfpathcurveto{\pgfqpoint{3.935546in}{0.867049in}}{\pgfqpoint{3.937860in}{0.872635in}}{\pgfqpoint{3.937860in}{0.878459in}}%
\pgfpathcurveto{\pgfqpoint{3.937860in}{0.884283in}}{\pgfqpoint{3.935546in}{0.889869in}}{\pgfqpoint{3.931428in}{0.893987in}}%
\pgfpathcurveto{\pgfqpoint{3.927310in}{0.898105in}}{\pgfqpoint{3.921724in}{0.900419in}}{\pgfqpoint{3.915900in}{0.900419in}}%
\pgfpathcurveto{\pgfqpoint{3.910076in}{0.900419in}}{\pgfqpoint{3.904490in}{0.898105in}}{\pgfqpoint{3.900372in}{0.893987in}}%
\pgfpathcurveto{\pgfqpoint{3.896254in}{0.889869in}}{\pgfqpoint{3.893940in}{0.884283in}}{\pgfqpoint{3.893940in}{0.878459in}}%
\pgfpathcurveto{\pgfqpoint{3.893940in}{0.872635in}}{\pgfqpoint{3.896254in}{0.867049in}}{\pgfqpoint{3.900372in}{0.862931in}}%
\pgfpathcurveto{\pgfqpoint{3.904490in}{0.858813in}}{\pgfqpoint{3.910076in}{0.856499in}}{\pgfqpoint{3.915900in}{0.856499in}}%
\pgfpathlineto{\pgfqpoint{3.915900in}{0.856499in}}%
\pgfpathclose%
\pgfusepath{stroke,fill}%
\end{pgfscope}%
\begin{pgfscope}%
\pgfpathrectangle{\pgfqpoint{1.414736in}{0.548769in}}{\pgfqpoint{3.497747in}{7.253186in}}%
\pgfusepath{clip}%
\pgfsetbuttcap%
\pgfsetroundjoin%
\definecolor{currentfill}{rgb}{0.699415,0.867117,0.175971}%
\pgfsetfillcolor{currentfill}%
\pgfsetlinewidth{1.003750pt}%
\definecolor{currentstroke}{rgb}{0.699415,0.867117,0.175971}%
\pgfsetstrokecolor{currentstroke}%
\pgfsetdash{}{0pt}%
\pgfpathmoveto{\pgfqpoint{3.240579in}{1.327485in}}%
\pgfpathcurveto{\pgfqpoint{3.246403in}{1.327485in}}{\pgfqpoint{3.251989in}{1.329799in}}{\pgfqpoint{3.256107in}{1.333917in}}%
\pgfpathcurveto{\pgfqpoint{3.260226in}{1.338035in}}{\pgfqpoint{3.262539in}{1.343621in}}{\pgfqpoint{3.262539in}{1.349445in}}%
\pgfpathcurveto{\pgfqpoint{3.262539in}{1.355269in}}{\pgfqpoint{3.260226in}{1.360855in}}{\pgfqpoint{3.256107in}{1.364973in}}%
\pgfpathcurveto{\pgfqpoint{3.251989in}{1.369092in}}{\pgfqpoint{3.246403in}{1.371405in}}{\pgfqpoint{3.240579in}{1.371405in}}%
\pgfpathcurveto{\pgfqpoint{3.234755in}{1.371405in}}{\pgfqpoint{3.229169in}{1.369092in}}{\pgfqpoint{3.225051in}{1.364973in}}%
\pgfpathcurveto{\pgfqpoint{3.220933in}{1.360855in}}{\pgfqpoint{3.218619in}{1.355269in}}{\pgfqpoint{3.218619in}{1.349445in}}%
\pgfpathcurveto{\pgfqpoint{3.218619in}{1.343621in}}{\pgfqpoint{3.220933in}{1.338035in}}{\pgfqpoint{3.225051in}{1.333917in}}%
\pgfpathcurveto{\pgfqpoint{3.229169in}{1.329799in}}{\pgfqpoint{3.234755in}{1.327485in}}{\pgfqpoint{3.240579in}{1.327485in}}%
\pgfpathlineto{\pgfqpoint{3.240579in}{1.327485in}}%
\pgfpathclose%
\pgfusepath{stroke,fill}%
\end{pgfscope}%
\begin{pgfscope}%
\pgfpathrectangle{\pgfqpoint{1.414736in}{0.548769in}}{\pgfqpoint{3.497747in}{7.253186in}}%
\pgfusepath{clip}%
\pgfsetbuttcap%
\pgfsetroundjoin%
\definecolor{currentfill}{rgb}{0.699415,0.867117,0.175971}%
\pgfsetfillcolor{currentfill}%
\pgfsetlinewidth{1.003750pt}%
\definecolor{currentstroke}{rgb}{0.699415,0.867117,0.175971}%
\pgfsetstrokecolor{currentstroke}%
\pgfsetdash{}{0pt}%
\pgfpathmoveto{\pgfqpoint{3.842450in}{1.798471in}}%
\pgfpathcurveto{\pgfqpoint{3.848274in}{1.798471in}}{\pgfqpoint{3.853860in}{1.800785in}}{\pgfqpoint{3.857978in}{1.804903in}}%
\pgfpathcurveto{\pgfqpoint{3.862096in}{1.809021in}}{\pgfqpoint{3.864410in}{1.814607in}}{\pgfqpoint{3.864410in}{1.820431in}}%
\pgfpathcurveto{\pgfqpoint{3.864410in}{1.826255in}}{\pgfqpoint{3.862096in}{1.831841in}}{\pgfqpoint{3.857978in}{1.835960in}}%
\pgfpathcurveto{\pgfqpoint{3.853860in}{1.840078in}}{\pgfqpoint{3.848274in}{1.842392in}}{\pgfqpoint{3.842450in}{1.842392in}}%
\pgfpathcurveto{\pgfqpoint{3.836626in}{1.842392in}}{\pgfqpoint{3.831039in}{1.840078in}}{\pgfqpoint{3.826921in}{1.835960in}}%
\pgfpathcurveto{\pgfqpoint{3.822803in}{1.831841in}}{\pgfqpoint{3.820489in}{1.826255in}}{\pgfqpoint{3.820489in}{1.820431in}}%
\pgfpathcurveto{\pgfqpoint{3.820489in}{1.814607in}}{\pgfqpoint{3.822803in}{1.809021in}}{\pgfqpoint{3.826921in}{1.804903in}}%
\pgfpathcurveto{\pgfqpoint{3.831039in}{1.800785in}}{\pgfqpoint{3.836626in}{1.798471in}}{\pgfqpoint{3.842450in}{1.798471in}}%
\pgfpathlineto{\pgfqpoint{3.842450in}{1.798471in}}%
\pgfpathclose%
\pgfusepath{stroke,fill}%
\end{pgfscope}%
\begin{pgfscope}%
\pgfpathrectangle{\pgfqpoint{1.414736in}{0.548769in}}{\pgfqpoint{3.497747in}{7.253186in}}%
\pgfusepath{clip}%
\pgfsetbuttcap%
\pgfsetroundjoin%
\definecolor{currentfill}{rgb}{0.699415,0.867117,0.175971}%
\pgfsetfillcolor{currentfill}%
\pgfsetlinewidth{1.003750pt}%
\definecolor{currentstroke}{rgb}{0.699415,0.867117,0.175971}%
\pgfsetstrokecolor{currentstroke}%
\pgfsetdash{}{0pt}%
\pgfpathmoveto{\pgfqpoint{3.932279in}{2.269457in}}%
\pgfpathcurveto{\pgfqpoint{3.938103in}{2.269457in}}{\pgfqpoint{3.943689in}{2.271771in}}{\pgfqpoint{3.947807in}{2.275889in}}%
\pgfpathcurveto{\pgfqpoint{3.951925in}{2.280007in}}{\pgfqpoint{3.954239in}{2.285593in}}{\pgfqpoint{3.954239in}{2.291417in}}%
\pgfpathcurveto{\pgfqpoint{3.954239in}{2.297241in}}{\pgfqpoint{3.951925in}{2.302828in}}{\pgfqpoint{3.947807in}{2.306946in}}%
\pgfpathcurveto{\pgfqpoint{3.943689in}{2.311064in}}{\pgfqpoint{3.938103in}{2.313378in}}{\pgfqpoint{3.932279in}{2.313378in}}%
\pgfpathcurveto{\pgfqpoint{3.926455in}{2.313378in}}{\pgfqpoint{3.920869in}{2.311064in}}{\pgfqpoint{3.916751in}{2.306946in}}%
\pgfpathcurveto{\pgfqpoint{3.912632in}{2.302828in}}{\pgfqpoint{3.910319in}{2.297241in}}{\pgfqpoint{3.910319in}{2.291417in}}%
\pgfpathcurveto{\pgfqpoint{3.910319in}{2.285593in}}{\pgfqpoint{3.912632in}{2.280007in}}{\pgfqpoint{3.916751in}{2.275889in}}%
\pgfpathcurveto{\pgfqpoint{3.920869in}{2.271771in}}{\pgfqpoint{3.926455in}{2.269457in}}{\pgfqpoint{3.932279in}{2.269457in}}%
\pgfpathlineto{\pgfqpoint{3.932279in}{2.269457in}}%
\pgfpathclose%
\pgfusepath{stroke,fill}%
\end{pgfscope}%
\begin{pgfscope}%
\pgfpathrectangle{\pgfqpoint{1.414736in}{0.548769in}}{\pgfqpoint{3.497747in}{7.253186in}}%
\pgfusepath{clip}%
\pgfsetbuttcap%
\pgfsetroundjoin%
\definecolor{currentfill}{rgb}{0.699415,0.867117,0.175971}%
\pgfsetfillcolor{currentfill}%
\pgfsetlinewidth{1.003750pt}%
\definecolor{currentstroke}{rgb}{0.699415,0.867117,0.175971}%
\pgfsetstrokecolor{currentstroke}%
\pgfsetdash{}{0pt}%
\pgfpathmoveto{\pgfqpoint{3.286941in}{2.740443in}}%
\pgfpathcurveto{\pgfqpoint{3.292765in}{2.740443in}}{\pgfqpoint{3.298351in}{2.742757in}}{\pgfqpoint{3.302469in}{2.746875in}}%
\pgfpathcurveto{\pgfqpoint{3.306587in}{2.750993in}}{\pgfqpoint{3.308901in}{2.756580in}}{\pgfqpoint{3.308901in}{2.762404in}}%
\pgfpathcurveto{\pgfqpoint{3.308901in}{2.768227in}}{\pgfqpoint{3.306587in}{2.773814in}}{\pgfqpoint{3.302469in}{2.777932in}}%
\pgfpathcurveto{\pgfqpoint{3.298351in}{2.782050in}}{\pgfqpoint{3.292765in}{2.784364in}}{\pgfqpoint{3.286941in}{2.784364in}}%
\pgfpathcurveto{\pgfqpoint{3.281117in}{2.784364in}}{\pgfqpoint{3.275531in}{2.782050in}}{\pgfqpoint{3.271413in}{2.777932in}}%
\pgfpathcurveto{\pgfqpoint{3.267294in}{2.773814in}}{\pgfqpoint{3.264980in}{2.768227in}}{\pgfqpoint{3.264980in}{2.762404in}}%
\pgfpathcurveto{\pgfqpoint{3.264980in}{2.756580in}}{\pgfqpoint{3.267294in}{2.750993in}}{\pgfqpoint{3.271413in}{2.746875in}}%
\pgfpathcurveto{\pgfqpoint{3.275531in}{2.742757in}}{\pgfqpoint{3.281117in}{2.740443in}}{\pgfqpoint{3.286941in}{2.740443in}}%
\pgfpathlineto{\pgfqpoint{3.286941in}{2.740443in}}%
\pgfpathclose%
\pgfusepath{stroke,fill}%
\end{pgfscope}%
\begin{pgfscope}%
\pgfpathrectangle{\pgfqpoint{1.414736in}{0.548769in}}{\pgfqpoint{3.497747in}{7.253186in}}%
\pgfusepath{clip}%
\pgfsetbuttcap%
\pgfsetroundjoin%
\definecolor{currentfill}{rgb}{0.699415,0.867117,0.175971}%
\pgfsetfillcolor{currentfill}%
\pgfsetlinewidth{1.003750pt}%
\definecolor{currentstroke}{rgb}{0.699415,0.867117,0.175971}%
\pgfsetstrokecolor{currentstroke}%
\pgfsetdash{}{0pt}%
\pgfpathmoveto{\pgfqpoint{3.687129in}{3.211429in}}%
\pgfpathcurveto{\pgfqpoint{3.692953in}{3.211429in}}{\pgfqpoint{3.698539in}{3.213743in}}{\pgfqpoint{3.702657in}{3.217861in}}%
\pgfpathcurveto{\pgfqpoint{3.706775in}{3.221980in}}{\pgfqpoint{3.709089in}{3.227566in}}{\pgfqpoint{3.709089in}{3.233390in}}%
\pgfpathcurveto{\pgfqpoint{3.709089in}{3.239214in}}{\pgfqpoint{3.706775in}{3.244800in}}{\pgfqpoint{3.702657in}{3.248918in}}%
\pgfpathcurveto{\pgfqpoint{3.698539in}{3.253036in}}{\pgfqpoint{3.692953in}{3.255350in}}{\pgfqpoint{3.687129in}{3.255350in}}%
\pgfpathcurveto{\pgfqpoint{3.681305in}{3.255350in}}{\pgfqpoint{3.675719in}{3.253036in}}{\pgfqpoint{3.671601in}{3.248918in}}%
\pgfpathcurveto{\pgfqpoint{3.667483in}{3.244800in}}{\pgfqpoint{3.665169in}{3.239214in}}{\pgfqpoint{3.665169in}{3.233390in}}%
\pgfpathcurveto{\pgfqpoint{3.665169in}{3.227566in}}{\pgfqpoint{3.667483in}{3.221980in}}{\pgfqpoint{3.671601in}{3.217861in}}%
\pgfpathcurveto{\pgfqpoint{3.675719in}{3.213743in}}{\pgfqpoint{3.681305in}{3.211429in}}{\pgfqpoint{3.687129in}{3.211429in}}%
\pgfpathlineto{\pgfqpoint{3.687129in}{3.211429in}}%
\pgfpathclose%
\pgfusepath{stroke,fill}%
\end{pgfscope}%
\begin{pgfscope}%
\pgfpathrectangle{\pgfqpoint{1.414736in}{0.548769in}}{\pgfqpoint{3.497747in}{7.253186in}}%
\pgfusepath{clip}%
\pgfsetbuttcap%
\pgfsetroundjoin%
\definecolor{currentfill}{rgb}{0.699415,0.867117,0.175971}%
\pgfsetfillcolor{currentfill}%
\pgfsetlinewidth{1.003750pt}%
\definecolor{currentstroke}{rgb}{0.699415,0.867117,0.175971}%
\pgfsetstrokecolor{currentstroke}%
\pgfsetdash{}{0pt}%
\pgfpathmoveto{\pgfqpoint{4.753495in}{3.682415in}}%
\pgfpathcurveto{\pgfqpoint{4.759319in}{3.682415in}}{\pgfqpoint{4.764905in}{3.684729in}}{\pgfqpoint{4.769023in}{3.688847in}}%
\pgfpathcurveto{\pgfqpoint{4.773141in}{3.692966in}}{\pgfqpoint{4.775455in}{3.698552in}}{\pgfqpoint{4.775455in}{3.704376in}}%
\pgfpathcurveto{\pgfqpoint{4.775455in}{3.710200in}}{\pgfqpoint{4.773141in}{3.715786in}}{\pgfqpoint{4.769023in}{3.719904in}}%
\pgfpathcurveto{\pgfqpoint{4.764905in}{3.724022in}}{\pgfqpoint{4.759319in}{3.726336in}}{\pgfqpoint{4.753495in}{3.726336in}}%
\pgfpathcurveto{\pgfqpoint{4.747671in}{3.726336in}}{\pgfqpoint{4.742085in}{3.724022in}}{\pgfqpoint{4.737966in}{3.719904in}}%
\pgfpathcurveto{\pgfqpoint{4.733848in}{3.715786in}}{\pgfqpoint{4.731534in}{3.710200in}}{\pgfqpoint{4.731534in}{3.704376in}}%
\pgfpathcurveto{\pgfqpoint{4.731534in}{3.698552in}}{\pgfqpoint{4.733848in}{3.692966in}}{\pgfqpoint{4.737966in}{3.688847in}}%
\pgfpathcurveto{\pgfqpoint{4.742085in}{3.684729in}}{\pgfqpoint{4.747671in}{3.682415in}}{\pgfqpoint{4.753495in}{3.682415in}}%
\pgfpathlineto{\pgfqpoint{4.753495in}{3.682415in}}%
\pgfpathclose%
\pgfusepath{stroke,fill}%
\end{pgfscope}%
\begin{pgfscope}%
\pgfpathrectangle{\pgfqpoint{1.414736in}{0.548769in}}{\pgfqpoint{3.497747in}{7.253186in}}%
\pgfusepath{clip}%
\pgfsetbuttcap%
\pgfsetroundjoin%
\definecolor{currentfill}{rgb}{0.699415,0.867117,0.175971}%
\pgfsetfillcolor{currentfill}%
\pgfsetlinewidth{1.003750pt}%
\definecolor{currentstroke}{rgb}{0.699415,0.867117,0.175971}%
\pgfsetstrokecolor{currentstroke}%
\pgfsetdash{}{0pt}%
\pgfpathmoveto{\pgfqpoint{4.030088in}{4.153402in}}%
\pgfpathcurveto{\pgfqpoint{4.035912in}{4.153402in}}{\pgfqpoint{4.041498in}{4.155715in}}{\pgfqpoint{4.045616in}{4.159834in}}%
\pgfpathcurveto{\pgfqpoint{4.049734in}{4.163952in}}{\pgfqpoint{4.052048in}{4.169538in}}{\pgfqpoint{4.052048in}{4.175362in}}%
\pgfpathcurveto{\pgfqpoint{4.052048in}{4.181186in}}{\pgfqpoint{4.049734in}{4.186772in}}{\pgfqpoint{4.045616in}{4.190890in}}%
\pgfpathcurveto{\pgfqpoint{4.041498in}{4.195008in}}{\pgfqpoint{4.035912in}{4.197322in}}{\pgfqpoint{4.030088in}{4.197322in}}%
\pgfpathcurveto{\pgfqpoint{4.024264in}{4.197322in}}{\pgfqpoint{4.018678in}{4.195008in}}{\pgfqpoint{4.014560in}{4.190890in}}%
\pgfpathcurveto{\pgfqpoint{4.010442in}{4.186772in}}{\pgfqpoint{4.008128in}{4.181186in}}{\pgfqpoint{4.008128in}{4.175362in}}%
\pgfpathcurveto{\pgfqpoint{4.008128in}{4.169538in}}{\pgfqpoint{4.010442in}{4.163952in}}{\pgfqpoint{4.014560in}{4.159834in}}%
\pgfpathcurveto{\pgfqpoint{4.018678in}{4.155715in}}{\pgfqpoint{4.024264in}{4.153402in}}{\pgfqpoint{4.030088in}{4.153402in}}%
\pgfpathlineto{\pgfqpoint{4.030088in}{4.153402in}}%
\pgfpathclose%
\pgfusepath{stroke,fill}%
\end{pgfscope}%
\begin{pgfscope}%
\pgfpathrectangle{\pgfqpoint{1.414736in}{0.548769in}}{\pgfqpoint{3.497747in}{7.253186in}}%
\pgfusepath{clip}%
\pgfsetbuttcap%
\pgfsetroundjoin%
\definecolor{currentfill}{rgb}{0.699415,0.867117,0.175971}%
\pgfsetfillcolor{currentfill}%
\pgfsetlinewidth{1.003750pt}%
\definecolor{currentstroke}{rgb}{0.699415,0.867117,0.175971}%
\pgfsetstrokecolor{currentstroke}%
\pgfsetdash{}{0pt}%
\pgfpathmoveto{\pgfqpoint{3.730682in}{4.624388in}}%
\pgfpathcurveto{\pgfqpoint{3.736506in}{4.624388in}}{\pgfqpoint{3.742093in}{4.626702in}}{\pgfqpoint{3.746211in}{4.630820in}}%
\pgfpathcurveto{\pgfqpoint{3.750329in}{4.634938in}}{\pgfqpoint{3.752643in}{4.640524in}}{\pgfqpoint{3.752643in}{4.646348in}}%
\pgfpathcurveto{\pgfqpoint{3.752643in}{4.652172in}}{\pgfqpoint{3.750329in}{4.657758in}}{\pgfqpoint{3.746211in}{4.661876in}}%
\pgfpathcurveto{\pgfqpoint{3.742093in}{4.665994in}}{\pgfqpoint{3.736506in}{4.668308in}}{\pgfqpoint{3.730682in}{4.668308in}}%
\pgfpathcurveto{\pgfqpoint{3.724858in}{4.668308in}}{\pgfqpoint{3.719272in}{4.665994in}}{\pgfqpoint{3.715154in}{4.661876in}}%
\pgfpathcurveto{\pgfqpoint{3.711036in}{4.657758in}}{\pgfqpoint{3.708722in}{4.652172in}}{\pgfqpoint{3.708722in}{4.646348in}}%
\pgfpathcurveto{\pgfqpoint{3.708722in}{4.640524in}}{\pgfqpoint{3.711036in}{4.634938in}}{\pgfqpoint{3.715154in}{4.630820in}}%
\pgfpathcurveto{\pgfqpoint{3.719272in}{4.626702in}}{\pgfqpoint{3.724858in}{4.624388in}}{\pgfqpoint{3.730682in}{4.624388in}}%
\pgfpathlineto{\pgfqpoint{3.730682in}{4.624388in}}%
\pgfpathclose%
\pgfusepath{stroke,fill}%
\end{pgfscope}%
\begin{pgfscope}%
\pgfpathrectangle{\pgfqpoint{1.414736in}{0.548769in}}{\pgfqpoint{3.497747in}{7.253186in}}%
\pgfusepath{clip}%
\pgfsetbuttcap%
\pgfsetroundjoin%
\definecolor{currentfill}{rgb}{0.699415,0.867117,0.175971}%
\pgfsetfillcolor{currentfill}%
\pgfsetlinewidth{1.003750pt}%
\definecolor{currentstroke}{rgb}{0.699415,0.867117,0.175971}%
\pgfsetstrokecolor{currentstroke}%
\pgfsetdash{}{0pt}%
\pgfpathmoveto{\pgfqpoint{2.920101in}{5.095374in}}%
\pgfpathcurveto{\pgfqpoint{2.925925in}{5.095374in}}{\pgfqpoint{2.931511in}{5.097688in}}{\pgfqpoint{2.935629in}{5.101806in}}%
\pgfpathcurveto{\pgfqpoint{2.939747in}{5.105924in}}{\pgfqpoint{2.942061in}{5.111510in}}{\pgfqpoint{2.942061in}{5.117334in}}%
\pgfpathcurveto{\pgfqpoint{2.942061in}{5.123158in}}{\pgfqpoint{2.939747in}{5.128744in}}{\pgfqpoint{2.935629in}{5.132862in}}%
\pgfpathcurveto{\pgfqpoint{2.931511in}{5.136980in}}{\pgfqpoint{2.925925in}{5.139294in}}{\pgfqpoint{2.920101in}{5.139294in}}%
\pgfpathcurveto{\pgfqpoint{2.914277in}{5.139294in}}{\pgfqpoint{2.908691in}{5.136980in}}{\pgfqpoint{2.904572in}{5.132862in}}%
\pgfpathcurveto{\pgfqpoint{2.900454in}{5.128744in}}{\pgfqpoint{2.898140in}{5.123158in}}{\pgfqpoint{2.898140in}{5.117334in}}%
\pgfpathcurveto{\pgfqpoint{2.898140in}{5.111510in}}{\pgfqpoint{2.900454in}{5.105924in}}{\pgfqpoint{2.904572in}{5.101806in}}%
\pgfpathcurveto{\pgfqpoint{2.908691in}{5.097688in}}{\pgfqpoint{2.914277in}{5.095374in}}{\pgfqpoint{2.920101in}{5.095374in}}%
\pgfpathlineto{\pgfqpoint{2.920101in}{5.095374in}}%
\pgfpathclose%
\pgfusepath{stroke,fill}%
\end{pgfscope}%
\begin{pgfscope}%
\pgfpathrectangle{\pgfqpoint{1.414736in}{0.548769in}}{\pgfqpoint{3.497747in}{7.253186in}}%
\pgfusepath{clip}%
\pgfsetbuttcap%
\pgfsetroundjoin%
\definecolor{currentfill}{rgb}{0.699415,0.867117,0.175971}%
\pgfsetfillcolor{currentfill}%
\pgfsetlinewidth{1.003750pt}%
\definecolor{currentstroke}{rgb}{0.699415,0.867117,0.175971}%
\pgfsetstrokecolor{currentstroke}%
\pgfsetdash{}{0pt}%
\pgfpathmoveto{\pgfqpoint{1.573725in}{5.566360in}}%
\pgfpathcurveto{\pgfqpoint{1.579548in}{5.566360in}}{\pgfqpoint{1.585135in}{5.568674in}}{\pgfqpoint{1.589253in}{5.572792in}}%
\pgfpathcurveto{\pgfqpoint{1.593371in}{5.576910in}}{\pgfqpoint{1.595685in}{5.582496in}}{\pgfqpoint{1.595685in}{5.588320in}}%
\pgfpathcurveto{\pgfqpoint{1.595685in}{5.594144in}}{\pgfqpoint{1.593371in}{5.599730in}}{\pgfqpoint{1.589253in}{5.603848in}}%
\pgfpathcurveto{\pgfqpoint{1.585135in}{5.607967in}}{\pgfqpoint{1.579548in}{5.610280in}}{\pgfqpoint{1.573725in}{5.610280in}}%
\pgfpathcurveto{\pgfqpoint{1.567901in}{5.610280in}}{\pgfqpoint{1.562314in}{5.607967in}}{\pgfqpoint{1.558196in}{5.603848in}}%
\pgfpathcurveto{\pgfqpoint{1.554078in}{5.599730in}}{\pgfqpoint{1.551764in}{5.594144in}}{\pgfqpoint{1.551764in}{5.588320in}}%
\pgfpathcurveto{\pgfqpoint{1.551764in}{5.582496in}}{\pgfqpoint{1.554078in}{5.576910in}}{\pgfqpoint{1.558196in}{5.572792in}}%
\pgfpathcurveto{\pgfqpoint{1.562314in}{5.568674in}}{\pgfqpoint{1.567901in}{5.566360in}}{\pgfqpoint{1.573725in}{5.566360in}}%
\pgfpathlineto{\pgfqpoint{1.573725in}{5.566360in}}%
\pgfpathclose%
\pgfusepath{stroke,fill}%
\end{pgfscope}%
\begin{pgfscope}%
\pgfpathrectangle{\pgfqpoint{1.414736in}{0.548769in}}{\pgfqpoint{3.497747in}{7.253186in}}%
\pgfusepath{clip}%
\pgfsetbuttcap%
\pgfsetroundjoin%
\definecolor{currentfill}{rgb}{0.699415,0.867117,0.175971}%
\pgfsetfillcolor{currentfill}%
\pgfsetlinewidth{1.003750pt}%
\definecolor{currentstroke}{rgb}{0.699415,0.867117,0.175971}%
\pgfsetstrokecolor{currentstroke}%
\pgfsetdash{}{0pt}%
\pgfpathmoveto{\pgfqpoint{3.877156in}{6.037346in}}%
\pgfpathcurveto{\pgfqpoint{3.882980in}{6.037346in}}{\pgfqpoint{3.888567in}{6.039660in}}{\pgfqpoint{3.892685in}{6.043778in}}%
\pgfpathcurveto{\pgfqpoint{3.896803in}{6.047896in}}{\pgfqpoint{3.899117in}{6.053482in}}{\pgfqpoint{3.899117in}{6.059306in}}%
\pgfpathcurveto{\pgfqpoint{3.899117in}{6.065130in}}{\pgfqpoint{3.896803in}{6.070716in}}{\pgfqpoint{3.892685in}{6.074835in}}%
\pgfpathcurveto{\pgfqpoint{3.888567in}{6.078953in}}{\pgfqpoint{3.882980in}{6.081267in}}{\pgfqpoint{3.877156in}{6.081267in}}%
\pgfpathcurveto{\pgfqpoint{3.871333in}{6.081267in}}{\pgfqpoint{3.865746in}{6.078953in}}{\pgfqpoint{3.861628in}{6.074835in}}%
\pgfpathcurveto{\pgfqpoint{3.857510in}{6.070716in}}{\pgfqpoint{3.855196in}{6.065130in}}{\pgfqpoint{3.855196in}{6.059306in}}%
\pgfpathcurveto{\pgfqpoint{3.855196in}{6.053482in}}{\pgfqpoint{3.857510in}{6.047896in}}{\pgfqpoint{3.861628in}{6.043778in}}%
\pgfpathcurveto{\pgfqpoint{3.865746in}{6.039660in}}{\pgfqpoint{3.871333in}{6.037346in}}{\pgfqpoint{3.877156in}{6.037346in}}%
\pgfpathlineto{\pgfqpoint{3.877156in}{6.037346in}}%
\pgfpathclose%
\pgfusepath{stroke,fill}%
\end{pgfscope}%
\begin{pgfscope}%
\pgfpathrectangle{\pgfqpoint{1.414736in}{0.548769in}}{\pgfqpoint{3.497747in}{7.253186in}}%
\pgfusepath{clip}%
\pgfsetbuttcap%
\pgfsetroundjoin%
\definecolor{currentfill}{rgb}{0.699415,0.867117,0.175971}%
\pgfsetfillcolor{currentfill}%
\pgfsetlinewidth{1.003750pt}%
\definecolor{currentstroke}{rgb}{0.699415,0.867117,0.175971}%
\pgfsetstrokecolor{currentstroke}%
\pgfsetdash{}{0pt}%
\pgfpathmoveto{\pgfqpoint{4.176053in}{6.508332in}}%
\pgfpathcurveto{\pgfqpoint{4.181877in}{6.508332in}}{\pgfqpoint{4.187464in}{6.510646in}}{\pgfqpoint{4.191582in}{6.514764in}}%
\pgfpathcurveto{\pgfqpoint{4.195700in}{6.518882in}}{\pgfqpoint{4.198014in}{6.524468in}}{\pgfqpoint{4.198014in}{6.530292in}}%
\pgfpathcurveto{\pgfqpoint{4.198014in}{6.536116in}}{\pgfqpoint{4.195700in}{6.541702in}}{\pgfqpoint{4.191582in}{6.545821in}}%
\pgfpathcurveto{\pgfqpoint{4.187464in}{6.549939in}}{\pgfqpoint{4.181877in}{6.552253in}}{\pgfqpoint{4.176053in}{6.552253in}}%
\pgfpathcurveto{\pgfqpoint{4.170230in}{6.552253in}}{\pgfqpoint{4.164643in}{6.549939in}}{\pgfqpoint{4.160525in}{6.545821in}}%
\pgfpathcurveto{\pgfqpoint{4.156407in}{6.541702in}}{\pgfqpoint{4.154093in}{6.536116in}}{\pgfqpoint{4.154093in}{6.530292in}}%
\pgfpathcurveto{\pgfqpoint{4.154093in}{6.524468in}}{\pgfqpoint{4.156407in}{6.518882in}}{\pgfqpoint{4.160525in}{6.514764in}}%
\pgfpathcurveto{\pgfqpoint{4.164643in}{6.510646in}}{\pgfqpoint{4.170230in}{6.508332in}}{\pgfqpoint{4.176053in}{6.508332in}}%
\pgfpathlineto{\pgfqpoint{4.176053in}{6.508332in}}%
\pgfpathclose%
\pgfusepath{stroke,fill}%
\end{pgfscope}%
\begin{pgfscope}%
\pgfpathrectangle{\pgfqpoint{1.414736in}{0.548769in}}{\pgfqpoint{3.497747in}{7.253186in}}%
\pgfusepath{clip}%
\pgfsetbuttcap%
\pgfsetroundjoin%
\definecolor{currentfill}{rgb}{0.699415,0.867117,0.175971}%
\pgfsetfillcolor{currentfill}%
\pgfsetlinewidth{1.003750pt}%
\definecolor{currentstroke}{rgb}{0.699415,0.867117,0.175971}%
\pgfsetstrokecolor{currentstroke}%
\pgfsetdash{}{0pt}%
\pgfpathmoveto{\pgfqpoint{4.605088in}{6.979318in}}%
\pgfpathcurveto{\pgfqpoint{4.610912in}{6.979318in}}{\pgfqpoint{4.616498in}{6.981632in}}{\pgfqpoint{4.620616in}{6.985750in}}%
\pgfpathcurveto{\pgfqpoint{4.624734in}{6.989868in}}{\pgfqpoint{4.627048in}{6.995455in}}{\pgfqpoint{4.627048in}{7.001278in}}%
\pgfpathcurveto{\pgfqpoint{4.627048in}{7.007102in}}{\pgfqpoint{4.624734in}{7.012689in}}{\pgfqpoint{4.620616in}{7.016807in}}%
\pgfpathcurveto{\pgfqpoint{4.616498in}{7.020925in}}{\pgfqpoint{4.610912in}{7.023239in}}{\pgfqpoint{4.605088in}{7.023239in}}%
\pgfpathcurveto{\pgfqpoint{4.599264in}{7.023239in}}{\pgfqpoint{4.593678in}{7.020925in}}{\pgfqpoint{4.589560in}{7.016807in}}%
\pgfpathcurveto{\pgfqpoint{4.585442in}{7.012689in}}{\pgfqpoint{4.583128in}{7.007102in}}{\pgfqpoint{4.583128in}{7.001278in}}%
\pgfpathcurveto{\pgfqpoint{4.583128in}{6.995455in}}{\pgfqpoint{4.585442in}{6.989868in}}{\pgfqpoint{4.589560in}{6.985750in}}%
\pgfpathcurveto{\pgfqpoint{4.593678in}{6.981632in}}{\pgfqpoint{4.599264in}{6.979318in}}{\pgfqpoint{4.605088in}{6.979318in}}%
\pgfpathlineto{\pgfqpoint{4.605088in}{6.979318in}}%
\pgfpathclose%
\pgfusepath{stroke,fill}%
\end{pgfscope}%
\begin{pgfscope}%
\pgfpathrectangle{\pgfqpoint{1.414736in}{0.548769in}}{\pgfqpoint{3.497747in}{7.253186in}}%
\pgfusepath{clip}%
\pgfsetbuttcap%
\pgfsetroundjoin%
\definecolor{currentfill}{rgb}{0.699415,0.867117,0.175971}%
\pgfsetfillcolor{currentfill}%
\pgfsetlinewidth{1.003750pt}%
\definecolor{currentstroke}{rgb}{0.699415,0.867117,0.175971}%
\pgfsetstrokecolor{currentstroke}%
\pgfsetdash{}{0pt}%
\pgfpathmoveto{\pgfqpoint{2.101439in}{7.450304in}}%
\pgfpathcurveto{\pgfqpoint{2.107263in}{7.450304in}}{\pgfqpoint{2.112849in}{7.452618in}}{\pgfqpoint{2.116967in}{7.456736in}}%
\pgfpathcurveto{\pgfqpoint{2.121085in}{7.460854in}}{\pgfqpoint{2.123399in}{7.466441in}}{\pgfqpoint{2.123399in}{7.472265in}}%
\pgfpathcurveto{\pgfqpoint{2.123399in}{7.478089in}}{\pgfqpoint{2.121085in}{7.483675in}}{\pgfqpoint{2.116967in}{7.487793in}}%
\pgfpathcurveto{\pgfqpoint{2.112849in}{7.491911in}}{\pgfqpoint{2.107263in}{7.494225in}}{\pgfqpoint{2.101439in}{7.494225in}}%
\pgfpathcurveto{\pgfqpoint{2.095615in}{7.494225in}}{\pgfqpoint{2.090029in}{7.491911in}}{\pgfqpoint{2.085910in}{7.487793in}}%
\pgfpathcurveto{\pgfqpoint{2.081792in}{7.483675in}}{\pgfqpoint{2.079478in}{7.478089in}}{\pgfqpoint{2.079478in}{7.472265in}}%
\pgfpathcurveto{\pgfqpoint{2.079478in}{7.466441in}}{\pgfqpoint{2.081792in}{7.460854in}}{\pgfqpoint{2.085910in}{7.456736in}}%
\pgfpathcurveto{\pgfqpoint{2.090029in}{7.452618in}}{\pgfqpoint{2.095615in}{7.450304in}}{\pgfqpoint{2.101439in}{7.450304in}}%
\pgfpathlineto{\pgfqpoint{2.101439in}{7.450304in}}%
\pgfpathclose%
\pgfusepath{stroke,fill}%
\end{pgfscope}%
\begin{pgfscope}%
\pgfpathrectangle{\pgfqpoint{1.414736in}{0.548769in}}{\pgfqpoint{3.497747in}{7.253186in}}%
\pgfusepath{clip}%
\pgfsetbuttcap%
\pgfsetroundjoin%
\definecolor{currentfill}{rgb}{0.699415,0.867117,0.175971}%
\pgfsetfillcolor{currentfill}%
\pgfsetlinewidth{1.003750pt}%
\definecolor{currentstroke}{rgb}{0.699415,0.867117,0.175971}%
\pgfsetstrokecolor{currentstroke}%
\pgfsetdash{}{0pt}%
\pgfpathmoveto{\pgfqpoint{3.965226in}{0.856499in}}%
\pgfpathcurveto{\pgfqpoint{3.971050in}{0.856499in}}{\pgfqpoint{3.976636in}{0.858813in}}{\pgfqpoint{3.980754in}{0.862931in}}%
\pgfpathcurveto{\pgfqpoint{3.984872in}{0.867049in}}{\pgfqpoint{3.987186in}{0.872635in}}{\pgfqpoint{3.987186in}{0.878459in}}%
\pgfpathcurveto{\pgfqpoint{3.987186in}{0.884283in}}{\pgfqpoint{3.984872in}{0.889869in}}{\pgfqpoint{3.980754in}{0.893987in}}%
\pgfpathcurveto{\pgfqpoint{3.976636in}{0.898105in}}{\pgfqpoint{3.971050in}{0.900419in}}{\pgfqpoint{3.965226in}{0.900419in}}%
\pgfpathcurveto{\pgfqpoint{3.959402in}{0.900419in}}{\pgfqpoint{3.953816in}{0.898105in}}{\pgfqpoint{3.949698in}{0.893987in}}%
\pgfpathcurveto{\pgfqpoint{3.945579in}{0.889869in}}{\pgfqpoint{3.943266in}{0.884283in}}{\pgfqpoint{3.943266in}{0.878459in}}%
\pgfpathcurveto{\pgfqpoint{3.943266in}{0.872635in}}{\pgfqpoint{3.945579in}{0.867049in}}{\pgfqpoint{3.949698in}{0.862931in}}%
\pgfpathcurveto{\pgfqpoint{3.953816in}{0.858813in}}{\pgfqpoint{3.959402in}{0.856499in}}{\pgfqpoint{3.965226in}{0.856499in}}%
\pgfpathlineto{\pgfqpoint{3.965226in}{0.856499in}}%
\pgfpathclose%
\pgfusepath{stroke,fill}%
\end{pgfscope}%
\begin{pgfscope}%
\pgfpathrectangle{\pgfqpoint{1.414736in}{0.548769in}}{\pgfqpoint{3.497747in}{7.253186in}}%
\pgfusepath{clip}%
\pgfsetbuttcap%
\pgfsetroundjoin%
\definecolor{currentfill}{rgb}{0.699415,0.867117,0.175971}%
\pgfsetfillcolor{currentfill}%
\pgfsetlinewidth{1.003750pt}%
\definecolor{currentstroke}{rgb}{0.699415,0.867117,0.175971}%
\pgfsetstrokecolor{currentstroke}%
\pgfsetdash{}{0pt}%
\pgfpathmoveto{\pgfqpoint{3.368049in}{1.327485in}}%
\pgfpathcurveto{\pgfqpoint{3.373873in}{1.327485in}}{\pgfqpoint{3.379459in}{1.329799in}}{\pgfqpoint{3.383577in}{1.333917in}}%
\pgfpathcurveto{\pgfqpoint{3.387695in}{1.338035in}}{\pgfqpoint{3.390009in}{1.343621in}}{\pgfqpoint{3.390009in}{1.349445in}}%
\pgfpathcurveto{\pgfqpoint{3.390009in}{1.355269in}}{\pgfqpoint{3.387695in}{1.360855in}}{\pgfqpoint{3.383577in}{1.364973in}}%
\pgfpathcurveto{\pgfqpoint{3.379459in}{1.369092in}}{\pgfqpoint{3.373873in}{1.371405in}}{\pgfqpoint{3.368049in}{1.371405in}}%
\pgfpathcurveto{\pgfqpoint{3.362225in}{1.371405in}}{\pgfqpoint{3.356639in}{1.369092in}}{\pgfqpoint{3.352520in}{1.364973in}}%
\pgfpathcurveto{\pgfqpoint{3.348402in}{1.360855in}}{\pgfqpoint{3.346088in}{1.355269in}}{\pgfqpoint{3.346088in}{1.349445in}}%
\pgfpathcurveto{\pgfqpoint{3.346088in}{1.343621in}}{\pgfqpoint{3.348402in}{1.338035in}}{\pgfqpoint{3.352520in}{1.333917in}}%
\pgfpathcurveto{\pgfqpoint{3.356639in}{1.329799in}}{\pgfqpoint{3.362225in}{1.327485in}}{\pgfqpoint{3.368049in}{1.327485in}}%
\pgfpathlineto{\pgfqpoint{3.368049in}{1.327485in}}%
\pgfpathclose%
\pgfusepath{stroke,fill}%
\end{pgfscope}%
\begin{pgfscope}%
\pgfpathrectangle{\pgfqpoint{1.414736in}{0.548769in}}{\pgfqpoint{3.497747in}{7.253186in}}%
\pgfusepath{clip}%
\pgfsetbuttcap%
\pgfsetroundjoin%
\definecolor{currentfill}{rgb}{0.699415,0.867117,0.175971}%
\pgfsetfillcolor{currentfill}%
\pgfsetlinewidth{1.003750pt}%
\definecolor{currentstroke}{rgb}{0.699415,0.867117,0.175971}%
\pgfsetstrokecolor{currentstroke}%
\pgfsetdash{}{0pt}%
\pgfpathmoveto{\pgfqpoint{3.762013in}{1.798471in}}%
\pgfpathcurveto{\pgfqpoint{3.767837in}{1.798471in}}{\pgfqpoint{3.773423in}{1.800785in}}{\pgfqpoint{3.777541in}{1.804903in}}%
\pgfpathcurveto{\pgfqpoint{3.781659in}{1.809021in}}{\pgfqpoint{3.783973in}{1.814607in}}{\pgfqpoint{3.783973in}{1.820431in}}%
\pgfpathcurveto{\pgfqpoint{3.783973in}{1.826255in}}{\pgfqpoint{3.781659in}{1.831841in}}{\pgfqpoint{3.777541in}{1.835960in}}%
\pgfpathcurveto{\pgfqpoint{3.773423in}{1.840078in}}{\pgfqpoint{3.767837in}{1.842392in}}{\pgfqpoint{3.762013in}{1.842392in}}%
\pgfpathcurveto{\pgfqpoint{3.756189in}{1.842392in}}{\pgfqpoint{3.750603in}{1.840078in}}{\pgfqpoint{3.746485in}{1.835960in}}%
\pgfpathcurveto{\pgfqpoint{3.742367in}{1.831841in}}{\pgfqpoint{3.740053in}{1.826255in}}{\pgfqpoint{3.740053in}{1.820431in}}%
\pgfpathcurveto{\pgfqpoint{3.740053in}{1.814607in}}{\pgfqpoint{3.742367in}{1.809021in}}{\pgfqpoint{3.746485in}{1.804903in}}%
\pgfpathcurveto{\pgfqpoint{3.750603in}{1.800785in}}{\pgfqpoint{3.756189in}{1.798471in}}{\pgfqpoint{3.762013in}{1.798471in}}%
\pgfpathlineto{\pgfqpoint{3.762013in}{1.798471in}}%
\pgfpathclose%
\pgfusepath{stroke,fill}%
\end{pgfscope}%
\begin{pgfscope}%
\pgfpathrectangle{\pgfqpoint{1.414736in}{0.548769in}}{\pgfqpoint{3.497747in}{7.253186in}}%
\pgfusepath{clip}%
\pgfsetbuttcap%
\pgfsetroundjoin%
\definecolor{currentfill}{rgb}{0.699415,0.867117,0.175971}%
\pgfsetfillcolor{currentfill}%
\pgfsetlinewidth{1.003750pt}%
\definecolor{currentstroke}{rgb}{0.699415,0.867117,0.175971}%
\pgfsetstrokecolor{currentstroke}%
\pgfsetdash{}{0pt}%
\pgfpathmoveto{\pgfqpoint{4.051201in}{2.269457in}}%
\pgfpathcurveto{\pgfqpoint{4.057025in}{2.269457in}}{\pgfqpoint{4.062611in}{2.271771in}}{\pgfqpoint{4.066730in}{2.275889in}}%
\pgfpathcurveto{\pgfqpoint{4.070848in}{2.280007in}}{\pgfqpoint{4.073162in}{2.285593in}}{\pgfqpoint{4.073162in}{2.291417in}}%
\pgfpathcurveto{\pgfqpoint{4.073162in}{2.297241in}}{\pgfqpoint{4.070848in}{2.302828in}}{\pgfqpoint{4.066730in}{2.306946in}}%
\pgfpathcurveto{\pgfqpoint{4.062611in}{2.311064in}}{\pgfqpoint{4.057025in}{2.313378in}}{\pgfqpoint{4.051201in}{2.313378in}}%
\pgfpathcurveto{\pgfqpoint{4.045377in}{2.313378in}}{\pgfqpoint{4.039791in}{2.311064in}}{\pgfqpoint{4.035673in}{2.306946in}}%
\pgfpathcurveto{\pgfqpoint{4.031555in}{2.302828in}}{\pgfqpoint{4.029241in}{2.297241in}}{\pgfqpoint{4.029241in}{2.291417in}}%
\pgfpathcurveto{\pgfqpoint{4.029241in}{2.285593in}}{\pgfqpoint{4.031555in}{2.280007in}}{\pgfqpoint{4.035673in}{2.275889in}}%
\pgfpathcurveto{\pgfqpoint{4.039791in}{2.271771in}}{\pgfqpoint{4.045377in}{2.269457in}}{\pgfqpoint{4.051201in}{2.269457in}}%
\pgfpathlineto{\pgfqpoint{4.051201in}{2.269457in}}%
\pgfpathclose%
\pgfusepath{stroke,fill}%
\end{pgfscope}%
\begin{pgfscope}%
\pgfpathrectangle{\pgfqpoint{1.414736in}{0.548769in}}{\pgfqpoint{3.497747in}{7.253186in}}%
\pgfusepath{clip}%
\pgfsetbuttcap%
\pgfsetroundjoin%
\definecolor{currentfill}{rgb}{0.699415,0.867117,0.175971}%
\pgfsetfillcolor{currentfill}%
\pgfsetlinewidth{1.003750pt}%
\definecolor{currentstroke}{rgb}{0.699415,0.867117,0.175971}%
\pgfsetstrokecolor{currentstroke}%
\pgfsetdash{}{0pt}%
\pgfpathmoveto{\pgfqpoint{3.292253in}{2.740443in}}%
\pgfpathcurveto{\pgfqpoint{3.298077in}{2.740443in}}{\pgfqpoint{3.303663in}{2.742757in}}{\pgfqpoint{3.307781in}{2.746875in}}%
\pgfpathcurveto{\pgfqpoint{3.311899in}{2.750993in}}{\pgfqpoint{3.314213in}{2.756580in}}{\pgfqpoint{3.314213in}{2.762404in}}%
\pgfpathcurveto{\pgfqpoint{3.314213in}{2.768227in}}{\pgfqpoint{3.311899in}{2.773814in}}{\pgfqpoint{3.307781in}{2.777932in}}%
\pgfpathcurveto{\pgfqpoint{3.303663in}{2.782050in}}{\pgfqpoint{3.298077in}{2.784364in}}{\pgfqpoint{3.292253in}{2.784364in}}%
\pgfpathcurveto{\pgfqpoint{3.286429in}{2.784364in}}{\pgfqpoint{3.280843in}{2.782050in}}{\pgfqpoint{3.276725in}{2.777932in}}%
\pgfpathcurveto{\pgfqpoint{3.272606in}{2.773814in}}{\pgfqpoint{3.270293in}{2.768227in}}{\pgfqpoint{3.270293in}{2.762404in}}%
\pgfpathcurveto{\pgfqpoint{3.270293in}{2.756580in}}{\pgfqpoint{3.272606in}{2.750993in}}{\pgfqpoint{3.276725in}{2.746875in}}%
\pgfpathcurveto{\pgfqpoint{3.280843in}{2.742757in}}{\pgfqpoint{3.286429in}{2.740443in}}{\pgfqpoint{3.292253in}{2.740443in}}%
\pgfpathlineto{\pgfqpoint{3.292253in}{2.740443in}}%
\pgfpathclose%
\pgfusepath{stroke,fill}%
\end{pgfscope}%
\begin{pgfscope}%
\pgfpathrectangle{\pgfqpoint{1.414736in}{0.548769in}}{\pgfqpoint{3.497747in}{7.253186in}}%
\pgfusepath{clip}%
\pgfsetbuttcap%
\pgfsetroundjoin%
\definecolor{currentfill}{rgb}{0.699415,0.867117,0.175971}%
\pgfsetfillcolor{currentfill}%
\pgfsetlinewidth{1.003750pt}%
\definecolor{currentstroke}{rgb}{0.699415,0.867117,0.175971}%
\pgfsetstrokecolor{currentstroke}%
\pgfsetdash{}{0pt}%
\pgfpathmoveto{\pgfqpoint{3.630692in}{3.211429in}}%
\pgfpathcurveto{\pgfqpoint{3.636516in}{3.211429in}}{\pgfqpoint{3.642102in}{3.213743in}}{\pgfqpoint{3.646221in}{3.217861in}}%
\pgfpathcurveto{\pgfqpoint{3.650339in}{3.221980in}}{\pgfqpoint{3.652653in}{3.227566in}}{\pgfqpoint{3.652653in}{3.233390in}}%
\pgfpathcurveto{\pgfqpoint{3.652653in}{3.239214in}}{\pgfqpoint{3.650339in}{3.244800in}}{\pgfqpoint{3.646221in}{3.248918in}}%
\pgfpathcurveto{\pgfqpoint{3.642102in}{3.253036in}}{\pgfqpoint{3.636516in}{3.255350in}}{\pgfqpoint{3.630692in}{3.255350in}}%
\pgfpathcurveto{\pgfqpoint{3.624868in}{3.255350in}}{\pgfqpoint{3.619282in}{3.253036in}}{\pgfqpoint{3.615164in}{3.248918in}}%
\pgfpathcurveto{\pgfqpoint{3.611046in}{3.244800in}}{\pgfqpoint{3.608732in}{3.239214in}}{\pgfqpoint{3.608732in}{3.233390in}}%
\pgfpathcurveto{\pgfqpoint{3.608732in}{3.227566in}}{\pgfqpoint{3.611046in}{3.221980in}}{\pgfqpoint{3.615164in}{3.217861in}}%
\pgfpathcurveto{\pgfqpoint{3.619282in}{3.213743in}}{\pgfqpoint{3.624868in}{3.211429in}}{\pgfqpoint{3.630692in}{3.211429in}}%
\pgfpathlineto{\pgfqpoint{3.630692in}{3.211429in}}%
\pgfpathclose%
\pgfusepath{stroke,fill}%
\end{pgfscope}%
\begin{pgfscope}%
\pgfpathrectangle{\pgfqpoint{1.414736in}{0.548769in}}{\pgfqpoint{3.497747in}{7.253186in}}%
\pgfusepath{clip}%
\pgfsetbuttcap%
\pgfsetroundjoin%
\definecolor{currentfill}{rgb}{0.699415,0.867117,0.175971}%
\pgfsetfillcolor{currentfill}%
\pgfsetlinewidth{1.003750pt}%
\definecolor{currentstroke}{rgb}{0.699415,0.867117,0.175971}%
\pgfsetstrokecolor{currentstroke}%
\pgfsetdash{}{0pt}%
\pgfpathmoveto{\pgfqpoint{4.595575in}{3.682415in}}%
\pgfpathcurveto{\pgfqpoint{4.601399in}{3.682415in}}{\pgfqpoint{4.606985in}{3.684729in}}{\pgfqpoint{4.611103in}{3.688847in}}%
\pgfpathcurveto{\pgfqpoint{4.615221in}{3.692966in}}{\pgfqpoint{4.617535in}{3.698552in}}{\pgfqpoint{4.617535in}{3.704376in}}%
\pgfpathcurveto{\pgfqpoint{4.617535in}{3.710200in}}{\pgfqpoint{4.615221in}{3.715786in}}{\pgfqpoint{4.611103in}{3.719904in}}%
\pgfpathcurveto{\pgfqpoint{4.606985in}{3.724022in}}{\pgfqpoint{4.601399in}{3.726336in}}{\pgfqpoint{4.595575in}{3.726336in}}%
\pgfpathcurveto{\pgfqpoint{4.589751in}{3.726336in}}{\pgfqpoint{4.584165in}{3.724022in}}{\pgfqpoint{4.580047in}{3.719904in}}%
\pgfpathcurveto{\pgfqpoint{4.575929in}{3.715786in}}{\pgfqpoint{4.573615in}{3.710200in}}{\pgfqpoint{4.573615in}{3.704376in}}%
\pgfpathcurveto{\pgfqpoint{4.573615in}{3.698552in}}{\pgfqpoint{4.575929in}{3.692966in}}{\pgfqpoint{4.580047in}{3.688847in}}%
\pgfpathcurveto{\pgfqpoint{4.584165in}{3.684729in}}{\pgfqpoint{4.589751in}{3.682415in}}{\pgfqpoint{4.595575in}{3.682415in}}%
\pgfpathlineto{\pgfqpoint{4.595575in}{3.682415in}}%
\pgfpathclose%
\pgfusepath{stroke,fill}%
\end{pgfscope}%
\begin{pgfscope}%
\pgfpathrectangle{\pgfqpoint{1.414736in}{0.548769in}}{\pgfqpoint{3.497747in}{7.253186in}}%
\pgfusepath{clip}%
\pgfsetbuttcap%
\pgfsetroundjoin%
\definecolor{currentfill}{rgb}{0.699415,0.867117,0.175971}%
\pgfsetfillcolor{currentfill}%
\pgfsetlinewidth{1.003750pt}%
\definecolor{currentstroke}{rgb}{0.699415,0.867117,0.175971}%
\pgfsetstrokecolor{currentstroke}%
\pgfsetdash{}{0pt}%
\pgfpathmoveto{\pgfqpoint{3.947326in}{4.153402in}}%
\pgfpathcurveto{\pgfqpoint{3.953150in}{4.153402in}}{\pgfqpoint{3.958736in}{4.155715in}}{\pgfqpoint{3.962854in}{4.159834in}}%
\pgfpathcurveto{\pgfqpoint{3.966972in}{4.163952in}}{\pgfqpoint{3.969286in}{4.169538in}}{\pgfqpoint{3.969286in}{4.175362in}}%
\pgfpathcurveto{\pgfqpoint{3.969286in}{4.181186in}}{\pgfqpoint{3.966972in}{4.186772in}}{\pgfqpoint{3.962854in}{4.190890in}}%
\pgfpathcurveto{\pgfqpoint{3.958736in}{4.195008in}}{\pgfqpoint{3.953150in}{4.197322in}}{\pgfqpoint{3.947326in}{4.197322in}}%
\pgfpathcurveto{\pgfqpoint{3.941502in}{4.197322in}}{\pgfqpoint{3.935916in}{4.195008in}}{\pgfqpoint{3.931798in}{4.190890in}}%
\pgfpathcurveto{\pgfqpoint{3.927679in}{4.186772in}}{\pgfqpoint{3.925366in}{4.181186in}}{\pgfqpoint{3.925366in}{4.175362in}}%
\pgfpathcurveto{\pgfqpoint{3.925366in}{4.169538in}}{\pgfqpoint{3.927679in}{4.163952in}}{\pgfqpoint{3.931798in}{4.159834in}}%
\pgfpathcurveto{\pgfqpoint{3.935916in}{4.155715in}}{\pgfqpoint{3.941502in}{4.153402in}}{\pgfqpoint{3.947326in}{4.153402in}}%
\pgfpathlineto{\pgfqpoint{3.947326in}{4.153402in}}%
\pgfpathclose%
\pgfusepath{stroke,fill}%
\end{pgfscope}%
\begin{pgfscope}%
\pgfpathrectangle{\pgfqpoint{1.414736in}{0.548769in}}{\pgfqpoint{3.497747in}{7.253186in}}%
\pgfusepath{clip}%
\pgfsetbuttcap%
\pgfsetroundjoin%
\definecolor{currentfill}{rgb}{0.699415,0.867117,0.175971}%
\pgfsetfillcolor{currentfill}%
\pgfsetlinewidth{1.003750pt}%
\definecolor{currentstroke}{rgb}{0.699415,0.867117,0.175971}%
\pgfsetstrokecolor{currentstroke}%
\pgfsetdash{}{0pt}%
\pgfpathmoveto{\pgfqpoint{3.713663in}{4.624388in}}%
\pgfpathcurveto{\pgfqpoint{3.719487in}{4.624388in}}{\pgfqpoint{3.725073in}{4.626702in}}{\pgfqpoint{3.729191in}{4.630820in}}%
\pgfpathcurveto{\pgfqpoint{3.733309in}{4.634938in}}{\pgfqpoint{3.735623in}{4.640524in}}{\pgfqpoint{3.735623in}{4.646348in}}%
\pgfpathcurveto{\pgfqpoint{3.735623in}{4.652172in}}{\pgfqpoint{3.733309in}{4.657758in}}{\pgfqpoint{3.729191in}{4.661876in}}%
\pgfpathcurveto{\pgfqpoint{3.725073in}{4.665994in}}{\pgfqpoint{3.719487in}{4.668308in}}{\pgfqpoint{3.713663in}{4.668308in}}%
\pgfpathcurveto{\pgfqpoint{3.707839in}{4.668308in}}{\pgfqpoint{3.702253in}{4.665994in}}{\pgfqpoint{3.698135in}{4.661876in}}%
\pgfpathcurveto{\pgfqpoint{3.694017in}{4.657758in}}{\pgfqpoint{3.691703in}{4.652172in}}{\pgfqpoint{3.691703in}{4.646348in}}%
\pgfpathcurveto{\pgfqpoint{3.691703in}{4.640524in}}{\pgfqpoint{3.694017in}{4.634938in}}{\pgfqpoint{3.698135in}{4.630820in}}%
\pgfpathcurveto{\pgfqpoint{3.702253in}{4.626702in}}{\pgfqpoint{3.707839in}{4.624388in}}{\pgfqpoint{3.713663in}{4.624388in}}%
\pgfpathlineto{\pgfqpoint{3.713663in}{4.624388in}}%
\pgfpathclose%
\pgfusepath{stroke,fill}%
\end{pgfscope}%
\begin{pgfscope}%
\pgfpathrectangle{\pgfqpoint{1.414736in}{0.548769in}}{\pgfqpoint{3.497747in}{7.253186in}}%
\pgfusepath{clip}%
\pgfsetbuttcap%
\pgfsetroundjoin%
\definecolor{currentfill}{rgb}{0.699415,0.867117,0.175971}%
\pgfsetfillcolor{currentfill}%
\pgfsetlinewidth{1.003750pt}%
\definecolor{currentstroke}{rgb}{0.699415,0.867117,0.175971}%
\pgfsetstrokecolor{currentstroke}%
\pgfsetdash{}{0pt}%
\pgfpathmoveto{\pgfqpoint{3.084646in}{5.095374in}}%
\pgfpathcurveto{\pgfqpoint{3.090469in}{5.095374in}}{\pgfqpoint{3.096056in}{5.097688in}}{\pgfqpoint{3.100174in}{5.101806in}}%
\pgfpathcurveto{\pgfqpoint{3.104292in}{5.105924in}}{\pgfqpoint{3.106606in}{5.111510in}}{\pgfqpoint{3.106606in}{5.117334in}}%
\pgfpathcurveto{\pgfqpoint{3.106606in}{5.123158in}}{\pgfqpoint{3.104292in}{5.128744in}}{\pgfqpoint{3.100174in}{5.132862in}}%
\pgfpathcurveto{\pgfqpoint{3.096056in}{5.136980in}}{\pgfqpoint{3.090469in}{5.139294in}}{\pgfqpoint{3.084646in}{5.139294in}}%
\pgfpathcurveto{\pgfqpoint{3.078822in}{5.139294in}}{\pgfqpoint{3.073235in}{5.136980in}}{\pgfqpoint{3.069117in}{5.132862in}}%
\pgfpathcurveto{\pgfqpoint{3.064999in}{5.128744in}}{\pgfqpoint{3.062685in}{5.123158in}}{\pgfqpoint{3.062685in}{5.117334in}}%
\pgfpathcurveto{\pgfqpoint{3.062685in}{5.111510in}}{\pgfqpoint{3.064999in}{5.105924in}}{\pgfqpoint{3.069117in}{5.101806in}}%
\pgfpathcurveto{\pgfqpoint{3.073235in}{5.097688in}}{\pgfqpoint{3.078822in}{5.095374in}}{\pgfqpoint{3.084646in}{5.095374in}}%
\pgfpathlineto{\pgfqpoint{3.084646in}{5.095374in}}%
\pgfpathclose%
\pgfusepath{stroke,fill}%
\end{pgfscope}%
\begin{pgfscope}%
\pgfpathrectangle{\pgfqpoint{1.414736in}{0.548769in}}{\pgfqpoint{3.497747in}{7.253186in}}%
\pgfusepath{clip}%
\pgfsetbuttcap%
\pgfsetroundjoin%
\definecolor{currentfill}{rgb}{0.699415,0.867117,0.175971}%
\pgfsetfillcolor{currentfill}%
\pgfsetlinewidth{1.003750pt}%
\definecolor{currentstroke}{rgb}{0.699415,0.867117,0.175971}%
\pgfsetstrokecolor{currentstroke}%
\pgfsetdash{}{0pt}%
\pgfpathmoveto{\pgfqpoint{1.820167in}{5.566360in}}%
\pgfpathcurveto{\pgfqpoint{1.825990in}{5.566360in}}{\pgfqpoint{1.831577in}{5.568674in}}{\pgfqpoint{1.835695in}{5.572792in}}%
\pgfpathcurveto{\pgfqpoint{1.839813in}{5.576910in}}{\pgfqpoint{1.842127in}{5.582496in}}{\pgfqpoint{1.842127in}{5.588320in}}%
\pgfpathcurveto{\pgfqpoint{1.842127in}{5.594144in}}{\pgfqpoint{1.839813in}{5.599730in}}{\pgfqpoint{1.835695in}{5.603848in}}%
\pgfpathcurveto{\pgfqpoint{1.831577in}{5.607967in}}{\pgfqpoint{1.825990in}{5.610280in}}{\pgfqpoint{1.820167in}{5.610280in}}%
\pgfpathcurveto{\pgfqpoint{1.814343in}{5.610280in}}{\pgfqpoint{1.808756in}{5.607967in}}{\pgfqpoint{1.804638in}{5.603848in}}%
\pgfpathcurveto{\pgfqpoint{1.800520in}{5.599730in}}{\pgfqpoint{1.798206in}{5.594144in}}{\pgfqpoint{1.798206in}{5.588320in}}%
\pgfpathcurveto{\pgfqpoint{1.798206in}{5.582496in}}{\pgfqpoint{1.800520in}{5.576910in}}{\pgfqpoint{1.804638in}{5.572792in}}%
\pgfpathcurveto{\pgfqpoint{1.808756in}{5.568674in}}{\pgfqpoint{1.814343in}{5.566360in}}{\pgfqpoint{1.820167in}{5.566360in}}%
\pgfpathlineto{\pgfqpoint{1.820167in}{5.566360in}}%
\pgfpathclose%
\pgfusepath{stroke,fill}%
\end{pgfscope}%
\begin{pgfscope}%
\pgfpathrectangle{\pgfqpoint{1.414736in}{0.548769in}}{\pgfqpoint{3.497747in}{7.253186in}}%
\pgfusepath{clip}%
\pgfsetbuttcap%
\pgfsetroundjoin%
\definecolor{currentfill}{rgb}{0.699415,0.867117,0.175971}%
\pgfsetfillcolor{currentfill}%
\pgfsetlinewidth{1.003750pt}%
\definecolor{currentstroke}{rgb}{0.699415,0.867117,0.175971}%
\pgfsetstrokecolor{currentstroke}%
\pgfsetdash{}{0pt}%
\pgfpathmoveto{\pgfqpoint{3.668342in}{6.037346in}}%
\pgfpathcurveto{\pgfqpoint{3.674166in}{6.037346in}}{\pgfqpoint{3.679752in}{6.039660in}}{\pgfqpoint{3.683870in}{6.043778in}}%
\pgfpathcurveto{\pgfqpoint{3.687988in}{6.047896in}}{\pgfqpoint{3.690302in}{6.053482in}}{\pgfqpoint{3.690302in}{6.059306in}}%
\pgfpathcurveto{\pgfqpoint{3.690302in}{6.065130in}}{\pgfqpoint{3.687988in}{6.070716in}}{\pgfqpoint{3.683870in}{6.074835in}}%
\pgfpathcurveto{\pgfqpoint{3.679752in}{6.078953in}}{\pgfqpoint{3.674166in}{6.081267in}}{\pgfqpoint{3.668342in}{6.081267in}}%
\pgfpathcurveto{\pgfqpoint{3.662518in}{6.081267in}}{\pgfqpoint{3.656932in}{6.078953in}}{\pgfqpoint{3.652813in}{6.074835in}}%
\pgfpathcurveto{\pgfqpoint{3.648695in}{6.070716in}}{\pgfqpoint{3.646381in}{6.065130in}}{\pgfqpoint{3.646381in}{6.059306in}}%
\pgfpathcurveto{\pgfqpoint{3.646381in}{6.053482in}}{\pgfqpoint{3.648695in}{6.047896in}}{\pgfqpoint{3.652813in}{6.043778in}}%
\pgfpathcurveto{\pgfqpoint{3.656932in}{6.039660in}}{\pgfqpoint{3.662518in}{6.037346in}}{\pgfqpoint{3.668342in}{6.037346in}}%
\pgfpathlineto{\pgfqpoint{3.668342in}{6.037346in}}%
\pgfpathclose%
\pgfusepath{stroke,fill}%
\end{pgfscope}%
\begin{pgfscope}%
\pgfpathrectangle{\pgfqpoint{1.414736in}{0.548769in}}{\pgfqpoint{3.497747in}{7.253186in}}%
\pgfusepath{clip}%
\pgfsetbuttcap%
\pgfsetroundjoin%
\definecolor{currentfill}{rgb}{0.699415,0.867117,0.175971}%
\pgfsetfillcolor{currentfill}%
\pgfsetlinewidth{1.003750pt}%
\definecolor{currentstroke}{rgb}{0.699415,0.867117,0.175971}%
\pgfsetstrokecolor{currentstroke}%
\pgfsetdash{}{0pt}%
\pgfpathmoveto{\pgfqpoint{3.921705in}{6.508332in}}%
\pgfpathcurveto{\pgfqpoint{3.927529in}{6.508332in}}{\pgfqpoint{3.933115in}{6.510646in}}{\pgfqpoint{3.937233in}{6.514764in}}%
\pgfpathcurveto{\pgfqpoint{3.941351in}{6.518882in}}{\pgfqpoint{3.943665in}{6.524468in}}{\pgfqpoint{3.943665in}{6.530292in}}%
\pgfpathcurveto{\pgfqpoint{3.943665in}{6.536116in}}{\pgfqpoint{3.941351in}{6.541702in}}{\pgfqpoint{3.937233in}{6.545821in}}%
\pgfpathcurveto{\pgfqpoint{3.933115in}{6.549939in}}{\pgfqpoint{3.927529in}{6.552253in}}{\pgfqpoint{3.921705in}{6.552253in}}%
\pgfpathcurveto{\pgfqpoint{3.915881in}{6.552253in}}{\pgfqpoint{3.910295in}{6.549939in}}{\pgfqpoint{3.906176in}{6.545821in}}%
\pgfpathcurveto{\pgfqpoint{3.902058in}{6.541702in}}{\pgfqpoint{3.899744in}{6.536116in}}{\pgfqpoint{3.899744in}{6.530292in}}%
\pgfpathcurveto{\pgfqpoint{3.899744in}{6.524468in}}{\pgfqpoint{3.902058in}{6.518882in}}{\pgfqpoint{3.906176in}{6.514764in}}%
\pgfpathcurveto{\pgfqpoint{3.910295in}{6.510646in}}{\pgfqpoint{3.915881in}{6.508332in}}{\pgfqpoint{3.921705in}{6.508332in}}%
\pgfpathlineto{\pgfqpoint{3.921705in}{6.508332in}}%
\pgfpathclose%
\pgfusepath{stroke,fill}%
\end{pgfscope}%
\begin{pgfscope}%
\pgfpathrectangle{\pgfqpoint{1.414736in}{0.548769in}}{\pgfqpoint{3.497747in}{7.253186in}}%
\pgfusepath{clip}%
\pgfsetbuttcap%
\pgfsetroundjoin%
\definecolor{currentfill}{rgb}{0.699415,0.867117,0.175971}%
\pgfsetfillcolor{currentfill}%
\pgfsetlinewidth{1.003750pt}%
\definecolor{currentstroke}{rgb}{0.699415,0.867117,0.175971}%
\pgfsetstrokecolor{currentstroke}%
\pgfsetdash{}{0pt}%
\pgfpathmoveto{\pgfqpoint{4.447505in}{6.979318in}}%
\pgfpathcurveto{\pgfqpoint{4.453329in}{6.979318in}}{\pgfqpoint{4.458916in}{6.981632in}}{\pgfqpoint{4.463034in}{6.985750in}}%
\pgfpathcurveto{\pgfqpoint{4.467152in}{6.989868in}}{\pgfqpoint{4.469466in}{6.995455in}}{\pgfqpoint{4.469466in}{7.001278in}}%
\pgfpathcurveto{\pgfqpoint{4.469466in}{7.007102in}}{\pgfqpoint{4.467152in}{7.012689in}}{\pgfqpoint{4.463034in}{7.016807in}}%
\pgfpathcurveto{\pgfqpoint{4.458916in}{7.020925in}}{\pgfqpoint{4.453329in}{7.023239in}}{\pgfqpoint{4.447505in}{7.023239in}}%
\pgfpathcurveto{\pgfqpoint{4.441682in}{7.023239in}}{\pgfqpoint{4.436095in}{7.020925in}}{\pgfqpoint{4.431977in}{7.016807in}}%
\pgfpathcurveto{\pgfqpoint{4.427859in}{7.012689in}}{\pgfqpoint{4.425545in}{7.007102in}}{\pgfqpoint{4.425545in}{7.001278in}}%
\pgfpathcurveto{\pgfqpoint{4.425545in}{6.995455in}}{\pgfqpoint{4.427859in}{6.989868in}}{\pgfqpoint{4.431977in}{6.985750in}}%
\pgfpathcurveto{\pgfqpoint{4.436095in}{6.981632in}}{\pgfqpoint{4.441682in}{6.979318in}}{\pgfqpoint{4.447505in}{6.979318in}}%
\pgfpathlineto{\pgfqpoint{4.447505in}{6.979318in}}%
\pgfpathclose%
\pgfusepath{stroke,fill}%
\end{pgfscope}%
\begin{pgfscope}%
\pgfpathrectangle{\pgfqpoint{1.414736in}{0.548769in}}{\pgfqpoint{3.497747in}{7.253186in}}%
\pgfusepath{clip}%
\pgfsetbuttcap%
\pgfsetroundjoin%
\definecolor{currentfill}{rgb}{0.699415,0.867117,0.175971}%
\pgfsetfillcolor{currentfill}%
\pgfsetlinewidth{1.003750pt}%
\definecolor{currentstroke}{rgb}{0.699415,0.867117,0.175971}%
\pgfsetstrokecolor{currentstroke}%
\pgfsetdash{}{0pt}%
\pgfpathmoveto{\pgfqpoint{2.051060in}{7.450304in}}%
\pgfpathcurveto{\pgfqpoint{2.056884in}{7.450304in}}{\pgfqpoint{2.062470in}{7.452618in}}{\pgfqpoint{2.066588in}{7.456736in}}%
\pgfpathcurveto{\pgfqpoint{2.070706in}{7.460854in}}{\pgfqpoint{2.073020in}{7.466441in}}{\pgfqpoint{2.073020in}{7.472265in}}%
\pgfpathcurveto{\pgfqpoint{2.073020in}{7.478089in}}{\pgfqpoint{2.070706in}{7.483675in}}{\pgfqpoint{2.066588in}{7.487793in}}%
\pgfpathcurveto{\pgfqpoint{2.062470in}{7.491911in}}{\pgfqpoint{2.056884in}{7.494225in}}{\pgfqpoint{2.051060in}{7.494225in}}%
\pgfpathcurveto{\pgfqpoint{2.045236in}{7.494225in}}{\pgfqpoint{2.039650in}{7.491911in}}{\pgfqpoint{2.035532in}{7.487793in}}%
\pgfpathcurveto{\pgfqpoint{2.031414in}{7.483675in}}{\pgfqpoint{2.029100in}{7.478089in}}{\pgfqpoint{2.029100in}{7.472265in}}%
\pgfpathcurveto{\pgfqpoint{2.029100in}{7.466441in}}{\pgfqpoint{2.031414in}{7.460854in}}{\pgfqpoint{2.035532in}{7.456736in}}%
\pgfpathcurveto{\pgfqpoint{2.039650in}{7.452618in}}{\pgfqpoint{2.045236in}{7.450304in}}{\pgfqpoint{2.051060in}{7.450304in}}%
\pgfpathlineto{\pgfqpoint{2.051060in}{7.450304in}}%
\pgfpathclose%
\pgfusepath{stroke,fill}%
\end{pgfscope}%
\begin{pgfscope}%
\pgfpathrectangle{\pgfqpoint{1.414736in}{0.548769in}}{\pgfqpoint{3.497747in}{7.253186in}}%
\pgfusepath{clip}%
\pgfsetbuttcap%
\pgfsetroundjoin%
\definecolor{currentfill}{rgb}{0.916242,0.896091,0.100717}%
\pgfsetfillcolor{currentfill}%
\pgfsetlinewidth{1.003750pt}%
\definecolor{currentstroke}{rgb}{0.916242,0.896091,0.100717}%
\pgfsetstrokecolor{currentstroke}%
\pgfsetdash{}{0pt}%
\pgfpathmoveto{\pgfqpoint{3.610309in}{0.856499in}}%
\pgfpathcurveto{\pgfqpoint{3.616133in}{0.856499in}}{\pgfqpoint{3.621719in}{0.858813in}}{\pgfqpoint{3.625837in}{0.862931in}}%
\pgfpathcurveto{\pgfqpoint{3.629955in}{0.867049in}}{\pgfqpoint{3.632269in}{0.872635in}}{\pgfqpoint{3.632269in}{0.878459in}}%
\pgfpathcurveto{\pgfqpoint{3.632269in}{0.884283in}}{\pgfqpoint{3.629955in}{0.889869in}}{\pgfqpoint{3.625837in}{0.893987in}}%
\pgfpathcurveto{\pgfqpoint{3.621719in}{0.898105in}}{\pgfqpoint{3.616133in}{0.900419in}}{\pgfqpoint{3.610309in}{0.900419in}}%
\pgfpathcurveto{\pgfqpoint{3.604485in}{0.900419in}}{\pgfqpoint{3.598899in}{0.898105in}}{\pgfqpoint{3.594780in}{0.893987in}}%
\pgfpathcurveto{\pgfqpoint{3.590662in}{0.889869in}}{\pgfqpoint{3.588348in}{0.884283in}}{\pgfqpoint{3.588348in}{0.878459in}}%
\pgfpathcurveto{\pgfqpoint{3.588348in}{0.872635in}}{\pgfqpoint{3.590662in}{0.867049in}}{\pgfqpoint{3.594780in}{0.862931in}}%
\pgfpathcurveto{\pgfqpoint{3.598899in}{0.858813in}}{\pgfqpoint{3.604485in}{0.856499in}}{\pgfqpoint{3.610309in}{0.856499in}}%
\pgfpathlineto{\pgfqpoint{3.610309in}{0.856499in}}%
\pgfpathclose%
\pgfusepath{stroke,fill}%
\end{pgfscope}%
\begin{pgfscope}%
\pgfpathrectangle{\pgfqpoint{1.414736in}{0.548769in}}{\pgfqpoint{3.497747in}{7.253186in}}%
\pgfusepath{clip}%
\pgfsetbuttcap%
\pgfsetroundjoin%
\definecolor{currentfill}{rgb}{0.616293,0.852709,0.230052}%
\pgfsetfillcolor{currentfill}%
\pgfsetlinewidth{1.003750pt}%
\definecolor{currentstroke}{rgb}{0.616293,0.852709,0.230052}%
\pgfsetstrokecolor{currentstroke}%
\pgfsetdash{}{0pt}%
\pgfpathmoveto{\pgfqpoint{3.461072in}{1.327485in}}%
\pgfpathcurveto{\pgfqpoint{3.466896in}{1.327485in}}{\pgfqpoint{3.472482in}{1.329799in}}{\pgfqpoint{3.476600in}{1.333917in}}%
\pgfpathcurveto{\pgfqpoint{3.480718in}{1.338035in}}{\pgfqpoint{3.483032in}{1.343621in}}{\pgfqpoint{3.483032in}{1.349445in}}%
\pgfpathcurveto{\pgfqpoint{3.483032in}{1.355269in}}{\pgfqpoint{3.480718in}{1.360855in}}{\pgfqpoint{3.476600in}{1.364973in}}%
\pgfpathcurveto{\pgfqpoint{3.472482in}{1.369092in}}{\pgfqpoint{3.466896in}{1.371405in}}{\pgfqpoint{3.461072in}{1.371405in}}%
\pgfpathcurveto{\pgfqpoint{3.455248in}{1.371405in}}{\pgfqpoint{3.449662in}{1.369092in}}{\pgfqpoint{3.445544in}{1.364973in}}%
\pgfpathcurveto{\pgfqpoint{3.441425in}{1.360855in}}{\pgfqpoint{3.439112in}{1.355269in}}{\pgfqpoint{3.439112in}{1.349445in}}%
\pgfpathcurveto{\pgfqpoint{3.439112in}{1.343621in}}{\pgfqpoint{3.441425in}{1.338035in}}{\pgfqpoint{3.445544in}{1.333917in}}%
\pgfpathcurveto{\pgfqpoint{3.449662in}{1.329799in}}{\pgfqpoint{3.455248in}{1.327485in}}{\pgfqpoint{3.461072in}{1.327485in}}%
\pgfpathlineto{\pgfqpoint{3.461072in}{1.327485in}}%
\pgfpathclose%
\pgfusepath{stroke,fill}%
\end{pgfscope}%
\begin{pgfscope}%
\pgfpathrectangle{\pgfqpoint{1.414736in}{0.548769in}}{\pgfqpoint{3.497747in}{7.253186in}}%
\pgfusepath{clip}%
\pgfsetbuttcap%
\pgfsetroundjoin%
\definecolor{currentfill}{rgb}{0.699415,0.867117,0.175971}%
\pgfsetfillcolor{currentfill}%
\pgfsetlinewidth{1.003750pt}%
\definecolor{currentstroke}{rgb}{0.699415,0.867117,0.175971}%
\pgfsetstrokecolor{currentstroke}%
\pgfsetdash{}{0pt}%
\pgfpathmoveto{\pgfqpoint{3.726923in}{1.798471in}}%
\pgfpathcurveto{\pgfqpoint{3.732747in}{1.798471in}}{\pgfqpoint{3.738333in}{1.800785in}}{\pgfqpoint{3.742451in}{1.804903in}}%
\pgfpathcurveto{\pgfqpoint{3.746569in}{1.809021in}}{\pgfqpoint{3.748883in}{1.814607in}}{\pgfqpoint{3.748883in}{1.820431in}}%
\pgfpathcurveto{\pgfqpoint{3.748883in}{1.826255in}}{\pgfqpoint{3.746569in}{1.831841in}}{\pgfqpoint{3.742451in}{1.835960in}}%
\pgfpathcurveto{\pgfqpoint{3.738333in}{1.840078in}}{\pgfqpoint{3.732747in}{1.842392in}}{\pgfqpoint{3.726923in}{1.842392in}}%
\pgfpathcurveto{\pgfqpoint{3.721099in}{1.842392in}}{\pgfqpoint{3.715513in}{1.840078in}}{\pgfqpoint{3.711395in}{1.835960in}}%
\pgfpathcurveto{\pgfqpoint{3.707276in}{1.831841in}}{\pgfqpoint{3.704963in}{1.826255in}}{\pgfqpoint{3.704963in}{1.820431in}}%
\pgfpathcurveto{\pgfqpoint{3.704963in}{1.814607in}}{\pgfqpoint{3.707276in}{1.809021in}}{\pgfqpoint{3.711395in}{1.804903in}}%
\pgfpathcurveto{\pgfqpoint{3.715513in}{1.800785in}}{\pgfqpoint{3.721099in}{1.798471in}}{\pgfqpoint{3.726923in}{1.798471in}}%
\pgfpathlineto{\pgfqpoint{3.726923in}{1.798471in}}%
\pgfpathclose%
\pgfusepath{stroke,fill}%
\end{pgfscope}%
\begin{pgfscope}%
\pgfpathrectangle{\pgfqpoint{1.414736in}{0.548769in}}{\pgfqpoint{3.497747in}{7.253186in}}%
\pgfusepath{clip}%
\pgfsetbuttcap%
\pgfsetroundjoin%
\definecolor{currentfill}{rgb}{0.993248,0.906157,0.143936}%
\pgfsetfillcolor{currentfill}%
\pgfsetlinewidth{1.003750pt}%
\definecolor{currentstroke}{rgb}{0.993248,0.906157,0.143936}%
\pgfsetstrokecolor{currentstroke}%
\pgfsetdash{}{0pt}%
\pgfpathmoveto{\pgfqpoint{3.779541in}{2.269457in}}%
\pgfpathcurveto{\pgfqpoint{3.785365in}{2.269457in}}{\pgfqpoint{3.790952in}{2.271771in}}{\pgfqpoint{3.795070in}{2.275889in}}%
\pgfpathcurveto{\pgfqpoint{3.799188in}{2.280007in}}{\pgfqpoint{3.801502in}{2.285593in}}{\pgfqpoint{3.801502in}{2.291417in}}%
\pgfpathcurveto{\pgfqpoint{3.801502in}{2.297241in}}{\pgfqpoint{3.799188in}{2.302828in}}{\pgfqpoint{3.795070in}{2.306946in}}%
\pgfpathcurveto{\pgfqpoint{3.790952in}{2.311064in}}{\pgfqpoint{3.785365in}{2.313378in}}{\pgfqpoint{3.779541in}{2.313378in}}%
\pgfpathcurveto{\pgfqpoint{3.773718in}{2.313378in}}{\pgfqpoint{3.768131in}{2.311064in}}{\pgfqpoint{3.764013in}{2.306946in}}%
\pgfpathcurveto{\pgfqpoint{3.759895in}{2.302828in}}{\pgfqpoint{3.757581in}{2.297241in}}{\pgfqpoint{3.757581in}{2.291417in}}%
\pgfpathcurveto{\pgfqpoint{3.757581in}{2.285593in}}{\pgfqpoint{3.759895in}{2.280007in}}{\pgfqpoint{3.764013in}{2.275889in}}%
\pgfpathcurveto{\pgfqpoint{3.768131in}{2.271771in}}{\pgfqpoint{3.773718in}{2.269457in}}{\pgfqpoint{3.779541in}{2.269457in}}%
\pgfpathlineto{\pgfqpoint{3.779541in}{2.269457in}}%
\pgfpathclose%
\pgfusepath{stroke,fill}%
\end{pgfscope}%
\begin{pgfscope}%
\pgfpathrectangle{\pgfqpoint{1.414736in}{0.548769in}}{\pgfqpoint{3.497747in}{7.253186in}}%
\pgfusepath{clip}%
\pgfsetbuttcap%
\pgfsetroundjoin%
\definecolor{currentfill}{rgb}{0.311925,0.767822,0.415586}%
\pgfsetfillcolor{currentfill}%
\pgfsetlinewidth{1.003750pt}%
\definecolor{currentstroke}{rgb}{0.311925,0.767822,0.415586}%
\pgfsetstrokecolor{currentstroke}%
\pgfsetdash{}{0pt}%
\pgfpathmoveto{\pgfqpoint{3.711804in}{2.740443in}}%
\pgfpathcurveto{\pgfqpoint{3.717627in}{2.740443in}}{\pgfqpoint{3.723214in}{2.742757in}}{\pgfqpoint{3.727332in}{2.746875in}}%
\pgfpathcurveto{\pgfqpoint{3.731450in}{2.750993in}}{\pgfqpoint{3.733764in}{2.756580in}}{\pgfqpoint{3.733764in}{2.762404in}}%
\pgfpathcurveto{\pgfqpoint{3.733764in}{2.768227in}}{\pgfqpoint{3.731450in}{2.773814in}}{\pgfqpoint{3.727332in}{2.777932in}}%
\pgfpathcurveto{\pgfqpoint{3.723214in}{2.782050in}}{\pgfqpoint{3.717627in}{2.784364in}}{\pgfqpoint{3.711804in}{2.784364in}}%
\pgfpathcurveto{\pgfqpoint{3.705980in}{2.784364in}}{\pgfqpoint{3.700393in}{2.782050in}}{\pgfqpoint{3.696275in}{2.777932in}}%
\pgfpathcurveto{\pgfqpoint{3.692157in}{2.773814in}}{\pgfqpoint{3.689843in}{2.768227in}}{\pgfqpoint{3.689843in}{2.762404in}}%
\pgfpathcurveto{\pgfqpoint{3.689843in}{2.756580in}}{\pgfqpoint{3.692157in}{2.750993in}}{\pgfqpoint{3.696275in}{2.746875in}}%
\pgfpathcurveto{\pgfqpoint{3.700393in}{2.742757in}}{\pgfqpoint{3.705980in}{2.740443in}}{\pgfqpoint{3.711804in}{2.740443in}}%
\pgfpathlineto{\pgfqpoint{3.711804in}{2.740443in}}%
\pgfpathclose%
\pgfusepath{stroke,fill}%
\end{pgfscope}%
\begin{pgfscope}%
\pgfpathrectangle{\pgfqpoint{1.414736in}{0.548769in}}{\pgfqpoint{3.497747in}{7.253186in}}%
\pgfusepath{clip}%
\pgfsetbuttcap%
\pgfsetroundjoin%
\definecolor{currentfill}{rgb}{0.772852,0.877868,0.131109}%
\pgfsetfillcolor{currentfill}%
\pgfsetlinewidth{1.003750pt}%
\definecolor{currentstroke}{rgb}{0.772852,0.877868,0.131109}%
\pgfsetstrokecolor{currentstroke}%
\pgfsetdash{}{0pt}%
\pgfpathmoveto{\pgfqpoint{3.372103in}{3.211429in}}%
\pgfpathcurveto{\pgfqpoint{3.377927in}{3.211429in}}{\pgfqpoint{3.383513in}{3.213743in}}{\pgfqpoint{3.387631in}{3.217861in}}%
\pgfpathcurveto{\pgfqpoint{3.391749in}{3.221980in}}{\pgfqpoint{3.394063in}{3.227566in}}{\pgfqpoint{3.394063in}{3.233390in}}%
\pgfpathcurveto{\pgfqpoint{3.394063in}{3.239214in}}{\pgfqpoint{3.391749in}{3.244800in}}{\pgfqpoint{3.387631in}{3.248918in}}%
\pgfpathcurveto{\pgfqpoint{3.383513in}{3.253036in}}{\pgfqpoint{3.377927in}{3.255350in}}{\pgfqpoint{3.372103in}{3.255350in}}%
\pgfpathcurveto{\pgfqpoint{3.366279in}{3.255350in}}{\pgfqpoint{3.360693in}{3.253036in}}{\pgfqpoint{3.356575in}{3.248918in}}%
\pgfpathcurveto{\pgfqpoint{3.352457in}{3.244800in}}{\pgfqpoint{3.350143in}{3.239214in}}{\pgfqpoint{3.350143in}{3.233390in}}%
\pgfpathcurveto{\pgfqpoint{3.350143in}{3.227566in}}{\pgfqpoint{3.352457in}{3.221980in}}{\pgfqpoint{3.356575in}{3.217861in}}%
\pgfpathcurveto{\pgfqpoint{3.360693in}{3.213743in}}{\pgfqpoint{3.366279in}{3.211429in}}{\pgfqpoint{3.372103in}{3.211429in}}%
\pgfpathlineto{\pgfqpoint{3.372103in}{3.211429in}}%
\pgfpathclose%
\pgfusepath{stroke,fill}%
\end{pgfscope}%
\begin{pgfscope}%
\pgfpathrectangle{\pgfqpoint{1.414736in}{0.548769in}}{\pgfqpoint{3.497747in}{7.253186in}}%
\pgfusepath{clip}%
\pgfsetbuttcap%
\pgfsetroundjoin%
\definecolor{currentfill}{rgb}{0.993248,0.906157,0.143936}%
\pgfsetfillcolor{currentfill}%
\pgfsetlinewidth{1.003750pt}%
\definecolor{currentstroke}{rgb}{0.993248,0.906157,0.143936}%
\pgfsetstrokecolor{currentstroke}%
\pgfsetdash{}{0pt}%
\pgfpathmoveto{\pgfqpoint{3.352442in}{3.682415in}}%
\pgfpathcurveto{\pgfqpoint{3.358265in}{3.682415in}}{\pgfqpoint{3.363852in}{3.684729in}}{\pgfqpoint{3.367970in}{3.688847in}}%
\pgfpathcurveto{\pgfqpoint{3.372088in}{3.692966in}}{\pgfqpoint{3.374402in}{3.698552in}}{\pgfqpoint{3.374402in}{3.704376in}}%
\pgfpathcurveto{\pgfqpoint{3.374402in}{3.710200in}}{\pgfqpoint{3.372088in}{3.715786in}}{\pgfqpoint{3.367970in}{3.719904in}}%
\pgfpathcurveto{\pgfqpoint{3.363852in}{3.724022in}}{\pgfqpoint{3.358265in}{3.726336in}}{\pgfqpoint{3.352442in}{3.726336in}}%
\pgfpathcurveto{\pgfqpoint{3.346618in}{3.726336in}}{\pgfqpoint{3.341031in}{3.724022in}}{\pgfqpoint{3.336913in}{3.719904in}}%
\pgfpathcurveto{\pgfqpoint{3.332795in}{3.715786in}}{\pgfqpoint{3.330481in}{3.710200in}}{\pgfqpoint{3.330481in}{3.704376in}}%
\pgfpathcurveto{\pgfqpoint{3.330481in}{3.698552in}}{\pgfqpoint{3.332795in}{3.692966in}}{\pgfqpoint{3.336913in}{3.688847in}}%
\pgfpathcurveto{\pgfqpoint{3.341031in}{3.684729in}}{\pgfqpoint{3.346618in}{3.682415in}}{\pgfqpoint{3.352442in}{3.682415in}}%
\pgfpathlineto{\pgfqpoint{3.352442in}{3.682415in}}%
\pgfpathclose%
\pgfusepath{stroke,fill}%
\end{pgfscope}%
\begin{pgfscope}%
\pgfpathrectangle{\pgfqpoint{1.414736in}{0.548769in}}{\pgfqpoint{3.497747in}{7.253186in}}%
\pgfusepath{clip}%
\pgfsetbuttcap%
\pgfsetroundjoin%
\definecolor{currentfill}{rgb}{0.993248,0.906157,0.143936}%
\pgfsetfillcolor{currentfill}%
\pgfsetlinewidth{1.003750pt}%
\definecolor{currentstroke}{rgb}{0.993248,0.906157,0.143936}%
\pgfsetstrokecolor{currentstroke}%
\pgfsetdash{}{0pt}%
\pgfpathmoveto{\pgfqpoint{4.302329in}{4.153402in}}%
\pgfpathcurveto{\pgfqpoint{4.308153in}{4.153402in}}{\pgfqpoint{4.313739in}{4.155715in}}{\pgfqpoint{4.317857in}{4.159834in}}%
\pgfpathcurveto{\pgfqpoint{4.321975in}{4.163952in}}{\pgfqpoint{4.324289in}{4.169538in}}{\pgfqpoint{4.324289in}{4.175362in}}%
\pgfpathcurveto{\pgfqpoint{4.324289in}{4.181186in}}{\pgfqpoint{4.321975in}{4.186772in}}{\pgfqpoint{4.317857in}{4.190890in}}%
\pgfpathcurveto{\pgfqpoint{4.313739in}{4.195008in}}{\pgfqpoint{4.308153in}{4.197322in}}{\pgfqpoint{4.302329in}{4.197322in}}%
\pgfpathcurveto{\pgfqpoint{4.296505in}{4.197322in}}{\pgfqpoint{4.290919in}{4.195008in}}{\pgfqpoint{4.286801in}{4.190890in}}%
\pgfpathcurveto{\pgfqpoint{4.282682in}{4.186772in}}{\pgfqpoint{4.280369in}{4.181186in}}{\pgfqpoint{4.280369in}{4.175362in}}%
\pgfpathcurveto{\pgfqpoint{4.280369in}{4.169538in}}{\pgfqpoint{4.282682in}{4.163952in}}{\pgfqpoint{4.286801in}{4.159834in}}%
\pgfpathcurveto{\pgfqpoint{4.290919in}{4.155715in}}{\pgfqpoint{4.296505in}{4.153402in}}{\pgfqpoint{4.302329in}{4.153402in}}%
\pgfpathlineto{\pgfqpoint{4.302329in}{4.153402in}}%
\pgfpathclose%
\pgfusepath{stroke,fill}%
\end{pgfscope}%
\begin{pgfscope}%
\pgfpathrectangle{\pgfqpoint{1.414736in}{0.548769in}}{\pgfqpoint{3.497747in}{7.253186in}}%
\pgfusepath{clip}%
\pgfsetbuttcap%
\pgfsetroundjoin%
\definecolor{currentfill}{rgb}{0.824940,0.884720,0.106217}%
\pgfsetfillcolor{currentfill}%
\pgfsetlinewidth{1.003750pt}%
\definecolor{currentstroke}{rgb}{0.824940,0.884720,0.106217}%
\pgfsetstrokecolor{currentstroke}%
\pgfsetdash{}{0pt}%
\pgfpathmoveto{\pgfqpoint{3.600728in}{4.624388in}}%
\pgfpathcurveto{\pgfqpoint{3.606552in}{4.624388in}}{\pgfqpoint{3.612138in}{4.626702in}}{\pgfqpoint{3.616257in}{4.630820in}}%
\pgfpathcurveto{\pgfqpoint{3.620375in}{4.634938in}}{\pgfqpoint{3.622689in}{4.640524in}}{\pgfqpoint{3.622689in}{4.646348in}}%
\pgfpathcurveto{\pgfqpoint{3.622689in}{4.652172in}}{\pgfqpoint{3.620375in}{4.657758in}}{\pgfqpoint{3.616257in}{4.661876in}}%
\pgfpathcurveto{\pgfqpoint{3.612138in}{4.665994in}}{\pgfqpoint{3.606552in}{4.668308in}}{\pgfqpoint{3.600728in}{4.668308in}}%
\pgfpathcurveto{\pgfqpoint{3.594904in}{4.668308in}}{\pgfqpoint{3.589318in}{4.665994in}}{\pgfqpoint{3.585200in}{4.661876in}}%
\pgfpathcurveto{\pgfqpoint{3.581082in}{4.657758in}}{\pgfqpoint{3.578768in}{4.652172in}}{\pgfqpoint{3.578768in}{4.646348in}}%
\pgfpathcurveto{\pgfqpoint{3.578768in}{4.640524in}}{\pgfqpoint{3.581082in}{4.634938in}}{\pgfqpoint{3.585200in}{4.630820in}}%
\pgfpathcurveto{\pgfqpoint{3.589318in}{4.626702in}}{\pgfqpoint{3.594904in}{4.624388in}}{\pgfqpoint{3.600728in}{4.624388in}}%
\pgfpathlineto{\pgfqpoint{3.600728in}{4.624388in}}%
\pgfpathclose%
\pgfusepath{stroke,fill}%
\end{pgfscope}%
\begin{pgfscope}%
\pgfpathrectangle{\pgfqpoint{1.414736in}{0.548769in}}{\pgfqpoint{3.497747in}{7.253186in}}%
\pgfusepath{clip}%
\pgfsetbuttcap%
\pgfsetroundjoin%
\definecolor{currentfill}{rgb}{0.151918,0.500685,0.557587}%
\pgfsetfillcolor{currentfill}%
\pgfsetlinewidth{1.003750pt}%
\definecolor{currentstroke}{rgb}{0.151918,0.500685,0.557587}%
\pgfsetstrokecolor{currentstroke}%
\pgfsetdash{}{0pt}%
\pgfpathmoveto{\pgfqpoint{3.351270in}{5.095374in}}%
\pgfpathcurveto{\pgfqpoint{3.357094in}{5.095374in}}{\pgfqpoint{3.362680in}{5.097688in}}{\pgfqpoint{3.366798in}{5.101806in}}%
\pgfpathcurveto{\pgfqpoint{3.370916in}{5.105924in}}{\pgfqpoint{3.373230in}{5.111510in}}{\pgfqpoint{3.373230in}{5.117334in}}%
\pgfpathcurveto{\pgfqpoint{3.373230in}{5.123158in}}{\pgfqpoint{3.370916in}{5.128744in}}{\pgfqpoint{3.366798in}{5.132862in}}%
\pgfpathcurveto{\pgfqpoint{3.362680in}{5.136980in}}{\pgfqpoint{3.357094in}{5.139294in}}{\pgfqpoint{3.351270in}{5.139294in}}%
\pgfpathcurveto{\pgfqpoint{3.345446in}{5.139294in}}{\pgfqpoint{3.339860in}{5.136980in}}{\pgfqpoint{3.335741in}{5.132862in}}%
\pgfpathcurveto{\pgfqpoint{3.331623in}{5.128744in}}{\pgfqpoint{3.329309in}{5.123158in}}{\pgfqpoint{3.329309in}{5.117334in}}%
\pgfpathcurveto{\pgfqpoint{3.329309in}{5.111510in}}{\pgfqpoint{3.331623in}{5.105924in}}{\pgfqpoint{3.335741in}{5.101806in}}%
\pgfpathcurveto{\pgfqpoint{3.339860in}{5.097688in}}{\pgfqpoint{3.345446in}{5.095374in}}{\pgfqpoint{3.351270in}{5.095374in}}%
\pgfpathlineto{\pgfqpoint{3.351270in}{5.095374in}}%
\pgfpathclose%
\pgfusepath{stroke,fill}%
\end{pgfscope}%
\begin{pgfscope}%
\pgfpathrectangle{\pgfqpoint{1.414736in}{0.548769in}}{\pgfqpoint{3.497747in}{7.253186in}}%
\pgfusepath{clip}%
\pgfsetbuttcap%
\pgfsetroundjoin%
\definecolor{currentfill}{rgb}{0.199430,0.387607,0.554642}%
\pgfsetfillcolor{currentfill}%
\pgfsetlinewidth{1.003750pt}%
\definecolor{currentstroke}{rgb}{0.199430,0.387607,0.554642}%
\pgfsetstrokecolor{currentstroke}%
\pgfsetdash{}{0pt}%
\pgfpathmoveto{\pgfqpoint{3.890848in}{5.566360in}}%
\pgfpathcurveto{\pgfqpoint{3.896672in}{5.566360in}}{\pgfqpoint{3.902258in}{5.568674in}}{\pgfqpoint{3.906377in}{5.572792in}}%
\pgfpathcurveto{\pgfqpoint{3.910495in}{5.576910in}}{\pgfqpoint{3.912809in}{5.582496in}}{\pgfqpoint{3.912809in}{5.588320in}}%
\pgfpathcurveto{\pgfqpoint{3.912809in}{5.594144in}}{\pgfqpoint{3.910495in}{5.599730in}}{\pgfqpoint{3.906377in}{5.603848in}}%
\pgfpathcurveto{\pgfqpoint{3.902258in}{5.607967in}}{\pgfqpoint{3.896672in}{5.610280in}}{\pgfqpoint{3.890848in}{5.610280in}}%
\pgfpathcurveto{\pgfqpoint{3.885024in}{5.610280in}}{\pgfqpoint{3.879438in}{5.607967in}}{\pgfqpoint{3.875320in}{5.603848in}}%
\pgfpathcurveto{\pgfqpoint{3.871202in}{5.599730in}}{\pgfqpoint{3.868888in}{5.594144in}}{\pgfqpoint{3.868888in}{5.588320in}}%
\pgfpathcurveto{\pgfqpoint{3.868888in}{5.582496in}}{\pgfqpoint{3.871202in}{5.576910in}}{\pgfqpoint{3.875320in}{5.572792in}}%
\pgfpathcurveto{\pgfqpoint{3.879438in}{5.568674in}}{\pgfqpoint{3.885024in}{5.566360in}}{\pgfqpoint{3.890848in}{5.566360in}}%
\pgfpathlineto{\pgfqpoint{3.890848in}{5.566360in}}%
\pgfpathclose%
\pgfusepath{stroke,fill}%
\end{pgfscope}%
\begin{pgfscope}%
\pgfpathrectangle{\pgfqpoint{1.414736in}{0.548769in}}{\pgfqpoint{3.497747in}{7.253186in}}%
\pgfusepath{clip}%
\pgfsetbuttcap%
\pgfsetroundjoin%
\definecolor{currentfill}{rgb}{0.259857,0.745492,0.444467}%
\pgfsetfillcolor{currentfill}%
\pgfsetlinewidth{1.003750pt}%
\definecolor{currentstroke}{rgb}{0.259857,0.745492,0.444467}%
\pgfsetstrokecolor{currentstroke}%
\pgfsetdash{}{0pt}%
\pgfpathmoveto{\pgfqpoint{3.443429in}{6.037346in}}%
\pgfpathcurveto{\pgfqpoint{3.449253in}{6.037346in}}{\pgfqpoint{3.454839in}{6.039660in}}{\pgfqpoint{3.458957in}{6.043778in}}%
\pgfpathcurveto{\pgfqpoint{3.463075in}{6.047896in}}{\pgfqpoint{3.465389in}{6.053482in}}{\pgfqpoint{3.465389in}{6.059306in}}%
\pgfpathcurveto{\pgfqpoint{3.465389in}{6.065130in}}{\pgfqpoint{3.463075in}{6.070716in}}{\pgfqpoint{3.458957in}{6.074835in}}%
\pgfpathcurveto{\pgfqpoint{3.454839in}{6.078953in}}{\pgfqpoint{3.449253in}{6.081267in}}{\pgfqpoint{3.443429in}{6.081267in}}%
\pgfpathcurveto{\pgfqpoint{3.437605in}{6.081267in}}{\pgfqpoint{3.432019in}{6.078953in}}{\pgfqpoint{3.427900in}{6.074835in}}%
\pgfpathcurveto{\pgfqpoint{3.423782in}{6.070716in}}{\pgfqpoint{3.421468in}{6.065130in}}{\pgfqpoint{3.421468in}{6.059306in}}%
\pgfpathcurveto{\pgfqpoint{3.421468in}{6.053482in}}{\pgfqpoint{3.423782in}{6.047896in}}{\pgfqpoint{3.427900in}{6.043778in}}%
\pgfpathcurveto{\pgfqpoint{3.432019in}{6.039660in}}{\pgfqpoint{3.437605in}{6.037346in}}{\pgfqpoint{3.443429in}{6.037346in}}%
\pgfpathlineto{\pgfqpoint{3.443429in}{6.037346in}}%
\pgfpathclose%
\pgfusepath{stroke,fill}%
\end{pgfscope}%
\begin{pgfscope}%
\pgfpathrectangle{\pgfqpoint{1.414736in}{0.548769in}}{\pgfqpoint{3.497747in}{7.253186in}}%
\pgfusepath{clip}%
\pgfsetbuttcap%
\pgfsetroundjoin%
\definecolor{currentfill}{rgb}{0.993248,0.906157,0.143936}%
\pgfsetfillcolor{currentfill}%
\pgfsetlinewidth{1.003750pt}%
\definecolor{currentstroke}{rgb}{0.993248,0.906157,0.143936}%
\pgfsetstrokecolor{currentstroke}%
\pgfsetdash{}{0pt}%
\pgfpathmoveto{\pgfqpoint{3.327235in}{6.508332in}}%
\pgfpathcurveto{\pgfqpoint{3.333059in}{6.508332in}}{\pgfqpoint{3.338645in}{6.510646in}}{\pgfqpoint{3.342763in}{6.514764in}}%
\pgfpathcurveto{\pgfqpoint{3.346881in}{6.518882in}}{\pgfqpoint{3.349195in}{6.524468in}}{\pgfqpoint{3.349195in}{6.530292in}}%
\pgfpathcurveto{\pgfqpoint{3.349195in}{6.536116in}}{\pgfqpoint{3.346881in}{6.541702in}}{\pgfqpoint{3.342763in}{6.545821in}}%
\pgfpathcurveto{\pgfqpoint{3.338645in}{6.549939in}}{\pgfqpoint{3.333059in}{6.552253in}}{\pgfqpoint{3.327235in}{6.552253in}}%
\pgfpathcurveto{\pgfqpoint{3.321411in}{6.552253in}}{\pgfqpoint{3.315825in}{6.549939in}}{\pgfqpoint{3.311706in}{6.545821in}}%
\pgfpathcurveto{\pgfqpoint{3.307588in}{6.541702in}}{\pgfqpoint{3.305274in}{6.536116in}}{\pgfqpoint{3.305274in}{6.530292in}}%
\pgfpathcurveto{\pgfqpoint{3.305274in}{6.524468in}}{\pgfqpoint{3.307588in}{6.518882in}}{\pgfqpoint{3.311706in}{6.514764in}}%
\pgfpathcurveto{\pgfqpoint{3.315825in}{6.510646in}}{\pgfqpoint{3.321411in}{6.508332in}}{\pgfqpoint{3.327235in}{6.508332in}}%
\pgfpathlineto{\pgfqpoint{3.327235in}{6.508332in}}%
\pgfpathclose%
\pgfusepath{stroke,fill}%
\end{pgfscope}%
\begin{pgfscope}%
\pgfpathrectangle{\pgfqpoint{1.414736in}{0.548769in}}{\pgfqpoint{3.497747in}{7.253186in}}%
\pgfusepath{clip}%
\pgfsetbuttcap%
\pgfsetroundjoin%
\definecolor{currentfill}{rgb}{0.835270,0.886029,0.102646}%
\pgfsetfillcolor{currentfill}%
\pgfsetlinewidth{1.003750pt}%
\definecolor{currentstroke}{rgb}{0.835270,0.886029,0.102646}%
\pgfsetstrokecolor{currentstroke}%
\pgfsetdash{}{0pt}%
\pgfpathmoveto{\pgfqpoint{3.827451in}{6.979318in}}%
\pgfpathcurveto{\pgfqpoint{3.833274in}{6.979318in}}{\pgfqpoint{3.838861in}{6.981632in}}{\pgfqpoint{3.842979in}{6.985750in}}%
\pgfpathcurveto{\pgfqpoint{3.847097in}{6.989868in}}{\pgfqpoint{3.849411in}{6.995455in}}{\pgfqpoint{3.849411in}{7.001278in}}%
\pgfpathcurveto{\pgfqpoint{3.849411in}{7.007102in}}{\pgfqpoint{3.847097in}{7.012689in}}{\pgfqpoint{3.842979in}{7.016807in}}%
\pgfpathcurveto{\pgfqpoint{3.838861in}{7.020925in}}{\pgfqpoint{3.833274in}{7.023239in}}{\pgfqpoint{3.827451in}{7.023239in}}%
\pgfpathcurveto{\pgfqpoint{3.821627in}{7.023239in}}{\pgfqpoint{3.816040in}{7.020925in}}{\pgfqpoint{3.811922in}{7.016807in}}%
\pgfpathcurveto{\pgfqpoint{3.807804in}{7.012689in}}{\pgfqpoint{3.805490in}{7.007102in}}{\pgfqpoint{3.805490in}{7.001278in}}%
\pgfpathcurveto{\pgfqpoint{3.805490in}{6.995455in}}{\pgfqpoint{3.807804in}{6.989868in}}{\pgfqpoint{3.811922in}{6.985750in}}%
\pgfpathcurveto{\pgfqpoint{3.816040in}{6.981632in}}{\pgfqpoint{3.821627in}{6.979318in}}{\pgfqpoint{3.827451in}{6.979318in}}%
\pgfpathlineto{\pgfqpoint{3.827451in}{6.979318in}}%
\pgfpathclose%
\pgfusepath{stroke,fill}%
\end{pgfscope}%
\begin{pgfscope}%
\pgfpathrectangle{\pgfqpoint{1.414736in}{0.548769in}}{\pgfqpoint{3.497747in}{7.253186in}}%
\pgfusepath{clip}%
\pgfsetbuttcap%
\pgfsetroundjoin%
\definecolor{currentfill}{rgb}{0.699415,0.867117,0.175971}%
\pgfsetfillcolor{currentfill}%
\pgfsetlinewidth{1.003750pt}%
\definecolor{currentstroke}{rgb}{0.699415,0.867117,0.175971}%
\pgfsetstrokecolor{currentstroke}%
\pgfsetdash{}{0pt}%
\pgfpathmoveto{\pgfqpoint{4.179723in}{7.450304in}}%
\pgfpathcurveto{\pgfqpoint{4.185547in}{7.450304in}}{\pgfqpoint{4.191133in}{7.452618in}}{\pgfqpoint{4.195251in}{7.456736in}}%
\pgfpathcurveto{\pgfqpoint{4.199369in}{7.460854in}}{\pgfqpoint{4.201683in}{7.466441in}}{\pgfqpoint{4.201683in}{7.472265in}}%
\pgfpathcurveto{\pgfqpoint{4.201683in}{7.478089in}}{\pgfqpoint{4.199369in}{7.483675in}}{\pgfqpoint{4.195251in}{7.487793in}}%
\pgfpathcurveto{\pgfqpoint{4.191133in}{7.491911in}}{\pgfqpoint{4.185547in}{7.494225in}}{\pgfqpoint{4.179723in}{7.494225in}}%
\pgfpathcurveto{\pgfqpoint{4.173899in}{7.494225in}}{\pgfqpoint{4.168313in}{7.491911in}}{\pgfqpoint{4.164195in}{7.487793in}}%
\pgfpathcurveto{\pgfqpoint{4.160077in}{7.483675in}}{\pgfqpoint{4.157763in}{7.478089in}}{\pgfqpoint{4.157763in}{7.472265in}}%
\pgfpathcurveto{\pgfqpoint{4.157763in}{7.466441in}}{\pgfqpoint{4.160077in}{7.460854in}}{\pgfqpoint{4.164195in}{7.456736in}}%
\pgfpathcurveto{\pgfqpoint{4.168313in}{7.452618in}}{\pgfqpoint{4.173899in}{7.450304in}}{\pgfqpoint{4.179723in}{7.450304in}}%
\pgfpathlineto{\pgfqpoint{4.179723in}{7.450304in}}%
\pgfpathclose%
\pgfusepath{stroke,fill}%
\end{pgfscope}%
\begin{pgfscope}%
\pgfpathrectangle{\pgfqpoint{1.414736in}{0.548769in}}{\pgfqpoint{3.497747in}{7.253186in}}%
\pgfusepath{clip}%
\pgfsetbuttcap%
\pgfsetroundjoin%
\definecolor{currentfill}{rgb}{0.772852,0.877868,0.131109}%
\pgfsetfillcolor{currentfill}%
\pgfsetlinewidth{1.003750pt}%
\definecolor{currentstroke}{rgb}{0.772852,0.877868,0.131109}%
\pgfsetstrokecolor{currentstroke}%
\pgfsetdash{}{0pt}%
\pgfpathmoveto{\pgfqpoint{3.526710in}{0.856499in}}%
\pgfpathcurveto{\pgfqpoint{3.532533in}{0.856499in}}{\pgfqpoint{3.538120in}{0.858813in}}{\pgfqpoint{3.542238in}{0.862931in}}%
\pgfpathcurveto{\pgfqpoint{3.546356in}{0.867049in}}{\pgfqpoint{3.548670in}{0.872635in}}{\pgfqpoint{3.548670in}{0.878459in}}%
\pgfpathcurveto{\pgfqpoint{3.548670in}{0.884283in}}{\pgfqpoint{3.546356in}{0.889869in}}{\pgfqpoint{3.542238in}{0.893987in}}%
\pgfpathcurveto{\pgfqpoint{3.538120in}{0.898105in}}{\pgfqpoint{3.532533in}{0.900419in}}{\pgfqpoint{3.526710in}{0.900419in}}%
\pgfpathcurveto{\pgfqpoint{3.520886in}{0.900419in}}{\pgfqpoint{3.515299in}{0.898105in}}{\pgfqpoint{3.511181in}{0.893987in}}%
\pgfpathcurveto{\pgfqpoint{3.507063in}{0.889869in}}{\pgfqpoint{3.504749in}{0.884283in}}{\pgfqpoint{3.504749in}{0.878459in}}%
\pgfpathcurveto{\pgfqpoint{3.504749in}{0.872635in}}{\pgfqpoint{3.507063in}{0.867049in}}{\pgfqpoint{3.511181in}{0.862931in}}%
\pgfpathcurveto{\pgfqpoint{3.515299in}{0.858813in}}{\pgfqpoint{3.520886in}{0.856499in}}{\pgfqpoint{3.526710in}{0.856499in}}%
\pgfpathlineto{\pgfqpoint{3.526710in}{0.856499in}}%
\pgfpathclose%
\pgfusepath{stroke,fill}%
\end{pgfscope}%
\begin{pgfscope}%
\pgfpathrectangle{\pgfqpoint{1.414736in}{0.548769in}}{\pgfqpoint{3.497747in}{7.253186in}}%
\pgfusepath{clip}%
\pgfsetbuttcap%
\pgfsetroundjoin%
\definecolor{currentfill}{rgb}{0.185783,0.704891,0.485273}%
\pgfsetfillcolor{currentfill}%
\pgfsetlinewidth{1.003750pt}%
\definecolor{currentstroke}{rgb}{0.185783,0.704891,0.485273}%
\pgfsetstrokecolor{currentstroke}%
\pgfsetdash{}{0pt}%
\pgfpathmoveto{\pgfqpoint{4.005432in}{1.327485in}}%
\pgfpathcurveto{\pgfqpoint{4.011256in}{1.327485in}}{\pgfqpoint{4.016842in}{1.329799in}}{\pgfqpoint{4.020961in}{1.333917in}}%
\pgfpathcurveto{\pgfqpoint{4.025079in}{1.338035in}}{\pgfqpoint{4.027393in}{1.343621in}}{\pgfqpoint{4.027393in}{1.349445in}}%
\pgfpathcurveto{\pgfqpoint{4.027393in}{1.355269in}}{\pgfqpoint{4.025079in}{1.360855in}}{\pgfqpoint{4.020961in}{1.364973in}}%
\pgfpathcurveto{\pgfqpoint{4.016842in}{1.369092in}}{\pgfqpoint{4.011256in}{1.371405in}}{\pgfqpoint{4.005432in}{1.371405in}}%
\pgfpathcurveto{\pgfqpoint{3.999608in}{1.371405in}}{\pgfqpoint{3.994022in}{1.369092in}}{\pgfqpoint{3.989904in}{1.364973in}}%
\pgfpathcurveto{\pgfqpoint{3.985786in}{1.360855in}}{\pgfqpoint{3.983472in}{1.355269in}}{\pgfqpoint{3.983472in}{1.349445in}}%
\pgfpathcurveto{\pgfqpoint{3.983472in}{1.343621in}}{\pgfqpoint{3.985786in}{1.338035in}}{\pgfqpoint{3.989904in}{1.333917in}}%
\pgfpathcurveto{\pgfqpoint{3.994022in}{1.329799in}}{\pgfqpoint{3.999608in}{1.327485in}}{\pgfqpoint{4.005432in}{1.327485in}}%
\pgfpathlineto{\pgfqpoint{4.005432in}{1.327485in}}%
\pgfpathclose%
\pgfusepath{stroke,fill}%
\end{pgfscope}%
\begin{pgfscope}%
\pgfpathrectangle{\pgfqpoint{1.414736in}{0.548769in}}{\pgfqpoint{3.497747in}{7.253186in}}%
\pgfusepath{clip}%
\pgfsetbuttcap%
\pgfsetroundjoin%
\definecolor{currentfill}{rgb}{0.626579,0.854645,0.223353}%
\pgfsetfillcolor{currentfill}%
\pgfsetlinewidth{1.003750pt}%
\definecolor{currentstroke}{rgb}{0.626579,0.854645,0.223353}%
\pgfsetstrokecolor{currentstroke}%
\pgfsetdash{}{0pt}%
\pgfpathmoveto{\pgfqpoint{3.073459in}{1.798471in}}%
\pgfpathcurveto{\pgfqpoint{3.079283in}{1.798471in}}{\pgfqpoint{3.084869in}{1.800785in}}{\pgfqpoint{3.088987in}{1.804903in}}%
\pgfpathcurveto{\pgfqpoint{3.093105in}{1.809021in}}{\pgfqpoint{3.095419in}{1.814607in}}{\pgfqpoint{3.095419in}{1.820431in}}%
\pgfpathcurveto{\pgfqpoint{3.095419in}{1.826255in}}{\pgfqpoint{3.093105in}{1.831841in}}{\pgfqpoint{3.088987in}{1.835960in}}%
\pgfpathcurveto{\pgfqpoint{3.084869in}{1.840078in}}{\pgfqpoint{3.079283in}{1.842392in}}{\pgfqpoint{3.073459in}{1.842392in}}%
\pgfpathcurveto{\pgfqpoint{3.067635in}{1.842392in}}{\pgfqpoint{3.062049in}{1.840078in}}{\pgfqpoint{3.057931in}{1.835960in}}%
\pgfpathcurveto{\pgfqpoint{3.053812in}{1.831841in}}{\pgfqpoint{3.051499in}{1.826255in}}{\pgfqpoint{3.051499in}{1.820431in}}%
\pgfpathcurveto{\pgfqpoint{3.051499in}{1.814607in}}{\pgfqpoint{3.053812in}{1.809021in}}{\pgfqpoint{3.057931in}{1.804903in}}%
\pgfpathcurveto{\pgfqpoint{3.062049in}{1.800785in}}{\pgfqpoint{3.067635in}{1.798471in}}{\pgfqpoint{3.073459in}{1.798471in}}%
\pgfpathlineto{\pgfqpoint{3.073459in}{1.798471in}}%
\pgfpathclose%
\pgfusepath{stroke,fill}%
\end{pgfscope}%
\begin{pgfscope}%
\pgfpathrectangle{\pgfqpoint{1.414736in}{0.548769in}}{\pgfqpoint{3.497747in}{7.253186in}}%
\pgfusepath{clip}%
\pgfsetbuttcap%
\pgfsetroundjoin%
\definecolor{currentfill}{rgb}{0.352360,0.783011,0.392636}%
\pgfsetfillcolor{currentfill}%
\pgfsetlinewidth{1.003750pt}%
\definecolor{currentstroke}{rgb}{0.352360,0.783011,0.392636}%
\pgfsetstrokecolor{currentstroke}%
\pgfsetdash{}{0pt}%
\pgfpathmoveto{\pgfqpoint{3.567685in}{2.269457in}}%
\pgfpathcurveto{\pgfqpoint{3.573509in}{2.269457in}}{\pgfqpoint{3.579095in}{2.271771in}}{\pgfqpoint{3.583213in}{2.275889in}}%
\pgfpathcurveto{\pgfqpoint{3.587331in}{2.280007in}}{\pgfqpoint{3.589645in}{2.285593in}}{\pgfqpoint{3.589645in}{2.291417in}}%
\pgfpathcurveto{\pgfqpoint{3.589645in}{2.297241in}}{\pgfqpoint{3.587331in}{2.302828in}}{\pgfqpoint{3.583213in}{2.306946in}}%
\pgfpathcurveto{\pgfqpoint{3.579095in}{2.311064in}}{\pgfqpoint{3.573509in}{2.313378in}}{\pgfqpoint{3.567685in}{2.313378in}}%
\pgfpathcurveto{\pgfqpoint{3.561861in}{2.313378in}}{\pgfqpoint{3.556275in}{2.311064in}}{\pgfqpoint{3.552156in}{2.306946in}}%
\pgfpathcurveto{\pgfqpoint{3.548038in}{2.302828in}}{\pgfqpoint{3.545724in}{2.297241in}}{\pgfqpoint{3.545724in}{2.291417in}}%
\pgfpathcurveto{\pgfqpoint{3.545724in}{2.285593in}}{\pgfqpoint{3.548038in}{2.280007in}}{\pgfqpoint{3.552156in}{2.275889in}}%
\pgfpathcurveto{\pgfqpoint{3.556275in}{2.271771in}}{\pgfqpoint{3.561861in}{2.269457in}}{\pgfqpoint{3.567685in}{2.269457in}}%
\pgfpathlineto{\pgfqpoint{3.567685in}{2.269457in}}%
\pgfpathclose%
\pgfusepath{stroke,fill}%
\end{pgfscope}%
\begin{pgfscope}%
\pgfpathrectangle{\pgfqpoint{1.414736in}{0.548769in}}{\pgfqpoint{3.497747in}{7.253186in}}%
\pgfusepath{clip}%
\pgfsetbuttcap%
\pgfsetroundjoin%
\definecolor{currentfill}{rgb}{0.430983,0.808473,0.346476}%
\pgfsetfillcolor{currentfill}%
\pgfsetlinewidth{1.003750pt}%
\definecolor{currentstroke}{rgb}{0.430983,0.808473,0.346476}%
\pgfsetstrokecolor{currentstroke}%
\pgfsetdash{}{0pt}%
\pgfpathmoveto{\pgfqpoint{3.685025in}{2.740443in}}%
\pgfpathcurveto{\pgfqpoint{3.690849in}{2.740443in}}{\pgfqpoint{3.696435in}{2.742757in}}{\pgfqpoint{3.700553in}{2.746875in}}%
\pgfpathcurveto{\pgfqpoint{3.704671in}{2.750993in}}{\pgfqpoint{3.706985in}{2.756580in}}{\pgfqpoint{3.706985in}{2.762404in}}%
\pgfpathcurveto{\pgfqpoint{3.706985in}{2.768227in}}{\pgfqpoint{3.704671in}{2.773814in}}{\pgfqpoint{3.700553in}{2.777932in}}%
\pgfpathcurveto{\pgfqpoint{3.696435in}{2.782050in}}{\pgfqpoint{3.690849in}{2.784364in}}{\pgfqpoint{3.685025in}{2.784364in}}%
\pgfpathcurveto{\pgfqpoint{3.679201in}{2.784364in}}{\pgfqpoint{3.673615in}{2.782050in}}{\pgfqpoint{3.669496in}{2.777932in}}%
\pgfpathcurveto{\pgfqpoint{3.665378in}{2.773814in}}{\pgfqpoint{3.663064in}{2.768227in}}{\pgfqpoint{3.663064in}{2.762404in}}%
\pgfpathcurveto{\pgfqpoint{3.663064in}{2.756580in}}{\pgfqpoint{3.665378in}{2.750993in}}{\pgfqpoint{3.669496in}{2.746875in}}%
\pgfpathcurveto{\pgfqpoint{3.673615in}{2.742757in}}{\pgfqpoint{3.679201in}{2.740443in}}{\pgfqpoint{3.685025in}{2.740443in}}%
\pgfpathlineto{\pgfqpoint{3.685025in}{2.740443in}}%
\pgfpathclose%
\pgfusepath{stroke,fill}%
\end{pgfscope}%
\begin{pgfscope}%
\pgfpathrectangle{\pgfqpoint{1.414736in}{0.548769in}}{\pgfqpoint{3.497747in}{7.253186in}}%
\pgfusepath{clip}%
\pgfsetbuttcap%
\pgfsetroundjoin%
\definecolor{currentfill}{rgb}{0.606045,0.850733,0.236712}%
\pgfsetfillcolor{currentfill}%
\pgfsetlinewidth{1.003750pt}%
\definecolor{currentstroke}{rgb}{0.606045,0.850733,0.236712}%
\pgfsetstrokecolor{currentstroke}%
\pgfsetdash{}{0pt}%
\pgfpathmoveto{\pgfqpoint{3.832433in}{3.211429in}}%
\pgfpathcurveto{\pgfqpoint{3.838256in}{3.211429in}}{\pgfqpoint{3.843843in}{3.213743in}}{\pgfqpoint{3.847961in}{3.217861in}}%
\pgfpathcurveto{\pgfqpoint{3.852079in}{3.221980in}}{\pgfqpoint{3.854393in}{3.227566in}}{\pgfqpoint{3.854393in}{3.233390in}}%
\pgfpathcurveto{\pgfqpoint{3.854393in}{3.239214in}}{\pgfqpoint{3.852079in}{3.244800in}}{\pgfqpoint{3.847961in}{3.248918in}}%
\pgfpathcurveto{\pgfqpoint{3.843843in}{3.253036in}}{\pgfqpoint{3.838256in}{3.255350in}}{\pgfqpoint{3.832433in}{3.255350in}}%
\pgfpathcurveto{\pgfqpoint{3.826609in}{3.255350in}}{\pgfqpoint{3.821022in}{3.253036in}}{\pgfqpoint{3.816904in}{3.248918in}}%
\pgfpathcurveto{\pgfqpoint{3.812786in}{3.244800in}}{\pgfqpoint{3.810472in}{3.239214in}}{\pgfqpoint{3.810472in}{3.233390in}}%
\pgfpathcurveto{\pgfqpoint{3.810472in}{3.227566in}}{\pgfqpoint{3.812786in}{3.221980in}}{\pgfqpoint{3.816904in}{3.217861in}}%
\pgfpathcurveto{\pgfqpoint{3.821022in}{3.213743in}}{\pgfqpoint{3.826609in}{3.211429in}}{\pgfqpoint{3.832433in}{3.211429in}}%
\pgfpathlineto{\pgfqpoint{3.832433in}{3.211429in}}%
\pgfpathclose%
\pgfusepath{stroke,fill}%
\end{pgfscope}%
\begin{pgfscope}%
\pgfpathrectangle{\pgfqpoint{1.414736in}{0.548769in}}{\pgfqpoint{3.497747in}{7.253186in}}%
\pgfusepath{clip}%
\pgfsetbuttcap%
\pgfsetroundjoin%
\definecolor{currentfill}{rgb}{0.955300,0.901065,0.118128}%
\pgfsetfillcolor{currentfill}%
\pgfsetlinewidth{1.003750pt}%
\definecolor{currentstroke}{rgb}{0.955300,0.901065,0.118128}%
\pgfsetstrokecolor{currentstroke}%
\pgfsetdash{}{0pt}%
\pgfpathmoveto{\pgfqpoint{3.517342in}{3.682415in}}%
\pgfpathcurveto{\pgfqpoint{3.523165in}{3.682415in}}{\pgfqpoint{3.528752in}{3.684729in}}{\pgfqpoint{3.532870in}{3.688847in}}%
\pgfpathcurveto{\pgfqpoint{3.536988in}{3.692966in}}{\pgfqpoint{3.539302in}{3.698552in}}{\pgfqpoint{3.539302in}{3.704376in}}%
\pgfpathcurveto{\pgfqpoint{3.539302in}{3.710200in}}{\pgfqpoint{3.536988in}{3.715786in}}{\pgfqpoint{3.532870in}{3.719904in}}%
\pgfpathcurveto{\pgfqpoint{3.528752in}{3.724022in}}{\pgfqpoint{3.523165in}{3.726336in}}{\pgfqpoint{3.517342in}{3.726336in}}%
\pgfpathcurveto{\pgfqpoint{3.511518in}{3.726336in}}{\pgfqpoint{3.505931in}{3.724022in}}{\pgfqpoint{3.501813in}{3.719904in}}%
\pgfpathcurveto{\pgfqpoint{3.497695in}{3.715786in}}{\pgfqpoint{3.495381in}{3.710200in}}{\pgfqpoint{3.495381in}{3.704376in}}%
\pgfpathcurveto{\pgfqpoint{3.495381in}{3.698552in}}{\pgfqpoint{3.497695in}{3.692966in}}{\pgfqpoint{3.501813in}{3.688847in}}%
\pgfpathcurveto{\pgfqpoint{3.505931in}{3.684729in}}{\pgfqpoint{3.511518in}{3.682415in}}{\pgfqpoint{3.517342in}{3.682415in}}%
\pgfpathlineto{\pgfqpoint{3.517342in}{3.682415in}}%
\pgfpathclose%
\pgfusepath{stroke,fill}%
\end{pgfscope}%
\begin{pgfscope}%
\pgfpathrectangle{\pgfqpoint{1.414736in}{0.548769in}}{\pgfqpoint{3.497747in}{7.253186in}}%
\pgfusepath{clip}%
\pgfsetbuttcap%
\pgfsetroundjoin%
\definecolor{currentfill}{rgb}{0.993248,0.906157,0.143936}%
\pgfsetfillcolor{currentfill}%
\pgfsetlinewidth{1.003750pt}%
\definecolor{currentstroke}{rgb}{0.993248,0.906157,0.143936}%
\pgfsetstrokecolor{currentstroke}%
\pgfsetdash{}{0pt}%
\pgfpathmoveto{\pgfqpoint{4.330373in}{4.153402in}}%
\pgfpathcurveto{\pgfqpoint{4.336197in}{4.153402in}}{\pgfqpoint{4.341783in}{4.155715in}}{\pgfqpoint{4.345901in}{4.159834in}}%
\pgfpathcurveto{\pgfqpoint{4.350019in}{4.163952in}}{\pgfqpoint{4.352333in}{4.169538in}}{\pgfqpoint{4.352333in}{4.175362in}}%
\pgfpathcurveto{\pgfqpoint{4.352333in}{4.181186in}}{\pgfqpoint{4.350019in}{4.186772in}}{\pgfqpoint{4.345901in}{4.190890in}}%
\pgfpathcurveto{\pgfqpoint{4.341783in}{4.195008in}}{\pgfqpoint{4.336197in}{4.197322in}}{\pgfqpoint{4.330373in}{4.197322in}}%
\pgfpathcurveto{\pgfqpoint{4.324549in}{4.197322in}}{\pgfqpoint{4.318963in}{4.195008in}}{\pgfqpoint{4.314845in}{4.190890in}}%
\pgfpathcurveto{\pgfqpoint{4.310727in}{4.186772in}}{\pgfqpoint{4.308413in}{4.181186in}}{\pgfqpoint{4.308413in}{4.175362in}}%
\pgfpathcurveto{\pgfqpoint{4.308413in}{4.169538in}}{\pgfqpoint{4.310727in}{4.163952in}}{\pgfqpoint{4.314845in}{4.159834in}}%
\pgfpathcurveto{\pgfqpoint{4.318963in}{4.155715in}}{\pgfqpoint{4.324549in}{4.153402in}}{\pgfqpoint{4.330373in}{4.153402in}}%
\pgfpathlineto{\pgfqpoint{4.330373in}{4.153402in}}%
\pgfpathclose%
\pgfusepath{stroke,fill}%
\end{pgfscope}%
\begin{pgfscope}%
\pgfpathrectangle{\pgfqpoint{1.414736in}{0.548769in}}{\pgfqpoint{3.497747in}{7.253186in}}%
\pgfusepath{clip}%
\pgfsetbuttcap%
\pgfsetroundjoin%
\definecolor{currentfill}{rgb}{0.575563,0.844566,0.256415}%
\pgfsetfillcolor{currentfill}%
\pgfsetlinewidth{1.003750pt}%
\definecolor{currentstroke}{rgb}{0.575563,0.844566,0.256415}%
\pgfsetstrokecolor{currentstroke}%
\pgfsetdash{}{0pt}%
\pgfpathmoveto{\pgfqpoint{3.185270in}{4.624388in}}%
\pgfpathcurveto{\pgfqpoint{3.191094in}{4.624388in}}{\pgfqpoint{3.196680in}{4.626702in}}{\pgfqpoint{3.200798in}{4.630820in}}%
\pgfpathcurveto{\pgfqpoint{3.204916in}{4.634938in}}{\pgfqpoint{3.207230in}{4.640524in}}{\pgfqpoint{3.207230in}{4.646348in}}%
\pgfpathcurveto{\pgfqpoint{3.207230in}{4.652172in}}{\pgfqpoint{3.204916in}{4.657758in}}{\pgfqpoint{3.200798in}{4.661876in}}%
\pgfpathcurveto{\pgfqpoint{3.196680in}{4.665994in}}{\pgfqpoint{3.191094in}{4.668308in}}{\pgfqpoint{3.185270in}{4.668308in}}%
\pgfpathcurveto{\pgfqpoint{3.179446in}{4.668308in}}{\pgfqpoint{3.173860in}{4.665994in}}{\pgfqpoint{3.169742in}{4.661876in}}%
\pgfpathcurveto{\pgfqpoint{3.165623in}{4.657758in}}{\pgfqpoint{3.163310in}{4.652172in}}{\pgfqpoint{3.163310in}{4.646348in}}%
\pgfpathcurveto{\pgfqpoint{3.163310in}{4.640524in}}{\pgfqpoint{3.165623in}{4.634938in}}{\pgfqpoint{3.169742in}{4.630820in}}%
\pgfpathcurveto{\pgfqpoint{3.173860in}{4.626702in}}{\pgfqpoint{3.179446in}{4.624388in}}{\pgfqpoint{3.185270in}{4.624388in}}%
\pgfpathlineto{\pgfqpoint{3.185270in}{4.624388in}}%
\pgfpathclose%
\pgfusepath{stroke,fill}%
\end{pgfscope}%
\begin{pgfscope}%
\pgfpathrectangle{\pgfqpoint{1.414736in}{0.548769in}}{\pgfqpoint{3.497747in}{7.253186in}}%
\pgfusepath{clip}%
\pgfsetbuttcap%
\pgfsetroundjoin%
\definecolor{currentfill}{rgb}{0.282327,0.094955,0.417331}%
\pgfsetfillcolor{currentfill}%
\pgfsetlinewidth{1.003750pt}%
\definecolor{currentstroke}{rgb}{0.282327,0.094955,0.417331}%
\pgfsetstrokecolor{currentstroke}%
\pgfsetdash{}{0pt}%
\pgfpathmoveto{\pgfqpoint{3.113327in}{5.095374in}}%
\pgfpathcurveto{\pgfqpoint{3.119151in}{5.095374in}}{\pgfqpoint{3.124737in}{5.097688in}}{\pgfqpoint{3.128856in}{5.101806in}}%
\pgfpathcurveto{\pgfqpoint{3.132974in}{5.105924in}}{\pgfqpoint{3.135288in}{5.111510in}}{\pgfqpoint{3.135288in}{5.117334in}}%
\pgfpathcurveto{\pgfqpoint{3.135288in}{5.123158in}}{\pgfqpoint{3.132974in}{5.128744in}}{\pgfqpoint{3.128856in}{5.132862in}}%
\pgfpathcurveto{\pgfqpoint{3.124737in}{5.136980in}}{\pgfqpoint{3.119151in}{5.139294in}}{\pgfqpoint{3.113327in}{5.139294in}}%
\pgfpathcurveto{\pgfqpoint{3.107503in}{5.139294in}}{\pgfqpoint{3.101917in}{5.136980in}}{\pgfqpoint{3.097799in}{5.132862in}}%
\pgfpathcurveto{\pgfqpoint{3.093681in}{5.128744in}}{\pgfqpoint{3.091367in}{5.123158in}}{\pgfqpoint{3.091367in}{5.117334in}}%
\pgfpathcurveto{\pgfqpoint{3.091367in}{5.111510in}}{\pgfqpoint{3.093681in}{5.105924in}}{\pgfqpoint{3.097799in}{5.101806in}}%
\pgfpathcurveto{\pgfqpoint{3.101917in}{5.097688in}}{\pgfqpoint{3.107503in}{5.095374in}}{\pgfqpoint{3.113327in}{5.095374in}}%
\pgfpathlineto{\pgfqpoint{3.113327in}{5.095374in}}%
\pgfpathclose%
\pgfusepath{stroke,fill}%
\end{pgfscope}%
\begin{pgfscope}%
\pgfpathrectangle{\pgfqpoint{1.414736in}{0.548769in}}{\pgfqpoint{3.497747in}{7.253186in}}%
\pgfusepath{clip}%
\pgfsetbuttcap%
\pgfsetroundjoin%
\definecolor{currentfill}{rgb}{0.235526,0.309527,0.542944}%
\pgfsetfillcolor{currentfill}%
\pgfsetlinewidth{1.003750pt}%
\definecolor{currentstroke}{rgb}{0.235526,0.309527,0.542944}%
\pgfsetstrokecolor{currentstroke}%
\pgfsetdash{}{0pt}%
\pgfpathmoveto{\pgfqpoint{4.017059in}{5.566360in}}%
\pgfpathcurveto{\pgfqpoint{4.022883in}{5.566360in}}{\pgfqpoint{4.028469in}{5.568674in}}{\pgfqpoint{4.032587in}{5.572792in}}%
\pgfpathcurveto{\pgfqpoint{4.036705in}{5.576910in}}{\pgfqpoint{4.039019in}{5.582496in}}{\pgfqpoint{4.039019in}{5.588320in}}%
\pgfpathcurveto{\pgfqpoint{4.039019in}{5.594144in}}{\pgfqpoint{4.036705in}{5.599730in}}{\pgfqpoint{4.032587in}{5.603848in}}%
\pgfpathcurveto{\pgfqpoint{4.028469in}{5.607967in}}{\pgfqpoint{4.022883in}{5.610280in}}{\pgfqpoint{4.017059in}{5.610280in}}%
\pgfpathcurveto{\pgfqpoint{4.011235in}{5.610280in}}{\pgfqpoint{4.005649in}{5.607967in}}{\pgfqpoint{4.001531in}{5.603848in}}%
\pgfpathcurveto{\pgfqpoint{3.997413in}{5.599730in}}{\pgfqpoint{3.995099in}{5.594144in}}{\pgfqpoint{3.995099in}{5.588320in}}%
\pgfpathcurveto{\pgfqpoint{3.995099in}{5.582496in}}{\pgfqpoint{3.997413in}{5.576910in}}{\pgfqpoint{4.001531in}{5.572792in}}%
\pgfpathcurveto{\pgfqpoint{4.005649in}{5.568674in}}{\pgfqpoint{4.011235in}{5.566360in}}{\pgfqpoint{4.017059in}{5.566360in}}%
\pgfpathlineto{\pgfqpoint{4.017059in}{5.566360in}}%
\pgfpathclose%
\pgfusepath{stroke,fill}%
\end{pgfscope}%
\begin{pgfscope}%
\pgfpathrectangle{\pgfqpoint{1.414736in}{0.548769in}}{\pgfqpoint{3.497747in}{7.253186in}}%
\pgfusepath{clip}%
\pgfsetbuttcap%
\pgfsetroundjoin%
\definecolor{currentfill}{rgb}{0.162142,0.474838,0.558140}%
\pgfsetfillcolor{currentfill}%
\pgfsetlinewidth{1.003750pt}%
\definecolor{currentstroke}{rgb}{0.162142,0.474838,0.558140}%
\pgfsetstrokecolor{currentstroke}%
\pgfsetdash{}{0pt}%
\pgfpathmoveto{\pgfqpoint{4.131293in}{6.037346in}}%
\pgfpathcurveto{\pgfqpoint{4.137117in}{6.037346in}}{\pgfqpoint{4.142703in}{6.039660in}}{\pgfqpoint{4.146821in}{6.043778in}}%
\pgfpathcurveto{\pgfqpoint{4.150939in}{6.047896in}}{\pgfqpoint{4.153253in}{6.053482in}}{\pgfqpoint{4.153253in}{6.059306in}}%
\pgfpathcurveto{\pgfqpoint{4.153253in}{6.065130in}}{\pgfqpoint{4.150939in}{6.070716in}}{\pgfqpoint{4.146821in}{6.074835in}}%
\pgfpathcurveto{\pgfqpoint{4.142703in}{6.078953in}}{\pgfqpoint{4.137117in}{6.081267in}}{\pgfqpoint{4.131293in}{6.081267in}}%
\pgfpathcurveto{\pgfqpoint{4.125469in}{6.081267in}}{\pgfqpoint{4.119883in}{6.078953in}}{\pgfqpoint{4.115764in}{6.074835in}}%
\pgfpathcurveto{\pgfqpoint{4.111646in}{6.070716in}}{\pgfqpoint{4.109332in}{6.065130in}}{\pgfqpoint{4.109332in}{6.059306in}}%
\pgfpathcurveto{\pgfqpoint{4.109332in}{6.053482in}}{\pgfqpoint{4.111646in}{6.047896in}}{\pgfqpoint{4.115764in}{6.043778in}}%
\pgfpathcurveto{\pgfqpoint{4.119883in}{6.039660in}}{\pgfqpoint{4.125469in}{6.037346in}}{\pgfqpoint{4.131293in}{6.037346in}}%
\pgfpathlineto{\pgfqpoint{4.131293in}{6.037346in}}%
\pgfpathclose%
\pgfusepath{stroke,fill}%
\end{pgfscope}%
\begin{pgfscope}%
\pgfpathrectangle{\pgfqpoint{1.414736in}{0.548769in}}{\pgfqpoint{3.497747in}{7.253186in}}%
\pgfusepath{clip}%
\pgfsetbuttcap%
\pgfsetroundjoin%
\definecolor{currentfill}{rgb}{0.993248,0.906157,0.143936}%
\pgfsetfillcolor{currentfill}%
\pgfsetlinewidth{1.003750pt}%
\definecolor{currentstroke}{rgb}{0.993248,0.906157,0.143936}%
\pgfsetstrokecolor{currentstroke}%
\pgfsetdash{}{0pt}%
\pgfpathmoveto{\pgfqpoint{3.076022in}{6.508332in}}%
\pgfpathcurveto{\pgfqpoint{3.081846in}{6.508332in}}{\pgfqpoint{3.087432in}{6.510646in}}{\pgfqpoint{3.091550in}{6.514764in}}%
\pgfpathcurveto{\pgfqpoint{3.095668in}{6.518882in}}{\pgfqpoint{3.097982in}{6.524468in}}{\pgfqpoint{3.097982in}{6.530292in}}%
\pgfpathcurveto{\pgfqpoint{3.097982in}{6.536116in}}{\pgfqpoint{3.095668in}{6.541702in}}{\pgfqpoint{3.091550in}{6.545821in}}%
\pgfpathcurveto{\pgfqpoint{3.087432in}{6.549939in}}{\pgfqpoint{3.081846in}{6.552253in}}{\pgfqpoint{3.076022in}{6.552253in}}%
\pgfpathcurveto{\pgfqpoint{3.070198in}{6.552253in}}{\pgfqpoint{3.064612in}{6.549939in}}{\pgfqpoint{3.060494in}{6.545821in}}%
\pgfpathcurveto{\pgfqpoint{3.056376in}{6.541702in}}{\pgfqpoint{3.054062in}{6.536116in}}{\pgfqpoint{3.054062in}{6.530292in}}%
\pgfpathcurveto{\pgfqpoint{3.054062in}{6.524468in}}{\pgfqpoint{3.056376in}{6.518882in}}{\pgfqpoint{3.060494in}{6.514764in}}%
\pgfpathcurveto{\pgfqpoint{3.064612in}{6.510646in}}{\pgfqpoint{3.070198in}{6.508332in}}{\pgfqpoint{3.076022in}{6.508332in}}%
\pgfpathlineto{\pgfqpoint{3.076022in}{6.508332in}}%
\pgfpathclose%
\pgfusepath{stroke,fill}%
\end{pgfscope}%
\begin{pgfscope}%
\pgfpathrectangle{\pgfqpoint{1.414736in}{0.548769in}}{\pgfqpoint{3.497747in}{7.253186in}}%
\pgfusepath{clip}%
\pgfsetbuttcap%
\pgfsetroundjoin%
\definecolor{currentfill}{rgb}{0.214000,0.722114,0.469588}%
\pgfsetfillcolor{currentfill}%
\pgfsetlinewidth{1.003750pt}%
\definecolor{currentstroke}{rgb}{0.214000,0.722114,0.469588}%
\pgfsetstrokecolor{currentstroke}%
\pgfsetdash{}{0pt}%
\pgfpathmoveto{\pgfqpoint{2.843691in}{6.979318in}}%
\pgfpathcurveto{\pgfqpoint{2.849515in}{6.979318in}}{\pgfqpoint{2.855101in}{6.981632in}}{\pgfqpoint{2.859219in}{6.985750in}}%
\pgfpathcurveto{\pgfqpoint{2.863337in}{6.989868in}}{\pgfqpoint{2.865651in}{6.995455in}}{\pgfqpoint{2.865651in}{7.001278in}}%
\pgfpathcurveto{\pgfqpoint{2.865651in}{7.007102in}}{\pgfqpoint{2.863337in}{7.012689in}}{\pgfqpoint{2.859219in}{7.016807in}}%
\pgfpathcurveto{\pgfqpoint{2.855101in}{7.020925in}}{\pgfqpoint{2.849515in}{7.023239in}}{\pgfqpoint{2.843691in}{7.023239in}}%
\pgfpathcurveto{\pgfqpoint{2.837867in}{7.023239in}}{\pgfqpoint{2.832281in}{7.020925in}}{\pgfqpoint{2.828163in}{7.016807in}}%
\pgfpathcurveto{\pgfqpoint{2.824045in}{7.012689in}}{\pgfqpoint{2.821731in}{7.007102in}}{\pgfqpoint{2.821731in}{7.001278in}}%
\pgfpathcurveto{\pgfqpoint{2.821731in}{6.995455in}}{\pgfqpoint{2.824045in}{6.989868in}}{\pgfqpoint{2.828163in}{6.985750in}}%
\pgfpathcurveto{\pgfqpoint{2.832281in}{6.981632in}}{\pgfqpoint{2.837867in}{6.979318in}}{\pgfqpoint{2.843691in}{6.979318in}}%
\pgfpathlineto{\pgfqpoint{2.843691in}{6.979318in}}%
\pgfpathclose%
\pgfusepath{stroke,fill}%
\end{pgfscope}%
\begin{pgfscope}%
\pgfpathrectangle{\pgfqpoint{1.414736in}{0.548769in}}{\pgfqpoint{3.497747in}{7.253186in}}%
\pgfusepath{clip}%
\pgfsetbuttcap%
\pgfsetroundjoin%
\definecolor{currentfill}{rgb}{0.993248,0.906157,0.143936}%
\pgfsetfillcolor{currentfill}%
\pgfsetlinewidth{1.003750pt}%
\definecolor{currentstroke}{rgb}{0.993248,0.906157,0.143936}%
\pgfsetstrokecolor{currentstroke}%
\pgfsetdash{}{0pt}%
\pgfpathmoveto{\pgfqpoint{4.563427in}{7.450304in}}%
\pgfpathcurveto{\pgfqpoint{4.569251in}{7.450304in}}{\pgfqpoint{4.574837in}{7.452618in}}{\pgfqpoint{4.578955in}{7.456736in}}%
\pgfpathcurveto{\pgfqpoint{4.583073in}{7.460854in}}{\pgfqpoint{4.585387in}{7.466441in}}{\pgfqpoint{4.585387in}{7.472265in}}%
\pgfpathcurveto{\pgfqpoint{4.585387in}{7.478089in}}{\pgfqpoint{4.583073in}{7.483675in}}{\pgfqpoint{4.578955in}{7.487793in}}%
\pgfpathcurveto{\pgfqpoint{4.574837in}{7.491911in}}{\pgfqpoint{4.569251in}{7.494225in}}{\pgfqpoint{4.563427in}{7.494225in}}%
\pgfpathcurveto{\pgfqpoint{4.557603in}{7.494225in}}{\pgfqpoint{4.552017in}{7.491911in}}{\pgfqpoint{4.547899in}{7.487793in}}%
\pgfpathcurveto{\pgfqpoint{4.543781in}{7.483675in}}{\pgfqpoint{4.541467in}{7.478089in}}{\pgfqpoint{4.541467in}{7.472265in}}%
\pgfpathcurveto{\pgfqpoint{4.541467in}{7.466441in}}{\pgfqpoint{4.543781in}{7.460854in}}{\pgfqpoint{4.547899in}{7.456736in}}%
\pgfpathcurveto{\pgfqpoint{4.552017in}{7.452618in}}{\pgfqpoint{4.557603in}{7.450304in}}{\pgfqpoint{4.563427in}{7.450304in}}%
\pgfpathlineto{\pgfqpoint{4.563427in}{7.450304in}}%
\pgfpathclose%
\pgfusepath{stroke,fill}%
\end{pgfscope}%
\begin{pgfscope}%
\pgfpathrectangle{\pgfqpoint{1.414736in}{0.548769in}}{\pgfqpoint{3.497747in}{7.253186in}}%
\pgfusepath{clip}%
\pgfsetbuttcap%
\pgfsetroundjoin%
\definecolor{currentfill}{rgb}{0.906311,0.894855,0.098125}%
\pgfsetfillcolor{currentfill}%
\pgfsetlinewidth{1.003750pt}%
\definecolor{currentstroke}{rgb}{0.906311,0.894855,0.098125}%
\pgfsetstrokecolor{currentstroke}%
\pgfsetdash{}{0pt}%
\pgfpathmoveto{\pgfqpoint{3.405731in}{0.856499in}}%
\pgfpathcurveto{\pgfqpoint{3.411555in}{0.856499in}}{\pgfqpoint{3.417141in}{0.858813in}}{\pgfqpoint{3.421260in}{0.862931in}}%
\pgfpathcurveto{\pgfqpoint{3.425378in}{0.867049in}}{\pgfqpoint{3.427692in}{0.872635in}}{\pgfqpoint{3.427692in}{0.878459in}}%
\pgfpathcurveto{\pgfqpoint{3.427692in}{0.884283in}}{\pgfqpoint{3.425378in}{0.889869in}}{\pgfqpoint{3.421260in}{0.893987in}}%
\pgfpathcurveto{\pgfqpoint{3.417141in}{0.898105in}}{\pgfqpoint{3.411555in}{0.900419in}}{\pgfqpoint{3.405731in}{0.900419in}}%
\pgfpathcurveto{\pgfqpoint{3.399907in}{0.900419in}}{\pgfqpoint{3.394321in}{0.898105in}}{\pgfqpoint{3.390203in}{0.893987in}}%
\pgfpathcurveto{\pgfqpoint{3.386085in}{0.889869in}}{\pgfqpoint{3.383771in}{0.884283in}}{\pgfqpoint{3.383771in}{0.878459in}}%
\pgfpathcurveto{\pgfqpoint{3.383771in}{0.872635in}}{\pgfqpoint{3.386085in}{0.867049in}}{\pgfqpoint{3.390203in}{0.862931in}}%
\pgfpathcurveto{\pgfqpoint{3.394321in}{0.858813in}}{\pgfqpoint{3.399907in}{0.856499in}}{\pgfqpoint{3.405731in}{0.856499in}}%
\pgfpathlineto{\pgfqpoint{3.405731in}{0.856499in}}%
\pgfpathclose%
\pgfusepath{stroke,fill}%
\end{pgfscope}%
\begin{pgfscope}%
\pgfpathrectangle{\pgfqpoint{1.414736in}{0.548769in}}{\pgfqpoint{3.497747in}{7.253186in}}%
\pgfusepath{clip}%
\pgfsetbuttcap%
\pgfsetroundjoin%
\definecolor{currentfill}{rgb}{0.377779,0.791781,0.377939}%
\pgfsetfillcolor{currentfill}%
\pgfsetlinewidth{1.003750pt}%
\definecolor{currentstroke}{rgb}{0.377779,0.791781,0.377939}%
\pgfsetstrokecolor{currentstroke}%
\pgfsetdash{}{0pt}%
\pgfpathmoveto{\pgfqpoint{3.839924in}{1.327485in}}%
\pgfpathcurveto{\pgfqpoint{3.845748in}{1.327485in}}{\pgfqpoint{3.851334in}{1.329799in}}{\pgfqpoint{3.855452in}{1.333917in}}%
\pgfpathcurveto{\pgfqpoint{3.859570in}{1.338035in}}{\pgfqpoint{3.861884in}{1.343621in}}{\pgfqpoint{3.861884in}{1.349445in}}%
\pgfpathcurveto{\pgfqpoint{3.861884in}{1.355269in}}{\pgfqpoint{3.859570in}{1.360855in}}{\pgfqpoint{3.855452in}{1.364973in}}%
\pgfpathcurveto{\pgfqpoint{3.851334in}{1.369092in}}{\pgfqpoint{3.845748in}{1.371405in}}{\pgfqpoint{3.839924in}{1.371405in}}%
\pgfpathcurveto{\pgfqpoint{3.834100in}{1.371405in}}{\pgfqpoint{3.828514in}{1.369092in}}{\pgfqpoint{3.824395in}{1.364973in}}%
\pgfpathcurveto{\pgfqpoint{3.820277in}{1.360855in}}{\pgfqpoint{3.817963in}{1.355269in}}{\pgfqpoint{3.817963in}{1.349445in}}%
\pgfpathcurveto{\pgfqpoint{3.817963in}{1.343621in}}{\pgfqpoint{3.820277in}{1.338035in}}{\pgfqpoint{3.824395in}{1.333917in}}%
\pgfpathcurveto{\pgfqpoint{3.828514in}{1.329799in}}{\pgfqpoint{3.834100in}{1.327485in}}{\pgfqpoint{3.839924in}{1.327485in}}%
\pgfpathlineto{\pgfqpoint{3.839924in}{1.327485in}}%
\pgfpathclose%
\pgfusepath{stroke,fill}%
\end{pgfscope}%
\begin{pgfscope}%
\pgfpathrectangle{\pgfqpoint{1.414736in}{0.548769in}}{\pgfqpoint{3.497747in}{7.253186in}}%
\pgfusepath{clip}%
\pgfsetbuttcap%
\pgfsetroundjoin%
\definecolor{currentfill}{rgb}{0.412913,0.803041,0.357269}%
\pgfsetfillcolor{currentfill}%
\pgfsetlinewidth{1.003750pt}%
\definecolor{currentstroke}{rgb}{0.412913,0.803041,0.357269}%
\pgfsetstrokecolor{currentstroke}%
\pgfsetdash{}{0pt}%
\pgfpathmoveto{\pgfqpoint{3.462748in}{1.798471in}}%
\pgfpathcurveto{\pgfqpoint{3.468572in}{1.798471in}}{\pgfqpoint{3.474158in}{1.800785in}}{\pgfqpoint{3.478276in}{1.804903in}}%
\pgfpathcurveto{\pgfqpoint{3.482395in}{1.809021in}}{\pgfqpoint{3.484708in}{1.814607in}}{\pgfqpoint{3.484708in}{1.820431in}}%
\pgfpathcurveto{\pgfqpoint{3.484708in}{1.826255in}}{\pgfqpoint{3.482395in}{1.831841in}}{\pgfqpoint{3.478276in}{1.835960in}}%
\pgfpathcurveto{\pgfqpoint{3.474158in}{1.840078in}}{\pgfqpoint{3.468572in}{1.842392in}}{\pgfqpoint{3.462748in}{1.842392in}}%
\pgfpathcurveto{\pgfqpoint{3.456924in}{1.842392in}}{\pgfqpoint{3.451338in}{1.840078in}}{\pgfqpoint{3.447220in}{1.835960in}}%
\pgfpathcurveto{\pgfqpoint{3.443102in}{1.831841in}}{\pgfqpoint{3.440788in}{1.826255in}}{\pgfqpoint{3.440788in}{1.820431in}}%
\pgfpathcurveto{\pgfqpoint{3.440788in}{1.814607in}}{\pgfqpoint{3.443102in}{1.809021in}}{\pgfqpoint{3.447220in}{1.804903in}}%
\pgfpathcurveto{\pgfqpoint{3.451338in}{1.800785in}}{\pgfqpoint{3.456924in}{1.798471in}}{\pgfqpoint{3.462748in}{1.798471in}}%
\pgfpathlineto{\pgfqpoint{3.462748in}{1.798471in}}%
\pgfpathclose%
\pgfusepath{stroke,fill}%
\end{pgfscope}%
\begin{pgfscope}%
\pgfpathrectangle{\pgfqpoint{1.414736in}{0.548769in}}{\pgfqpoint{3.497747in}{7.253186in}}%
\pgfusepath{clip}%
\pgfsetbuttcap%
\pgfsetroundjoin%
\definecolor{currentfill}{rgb}{0.440137,0.811138,0.340967}%
\pgfsetfillcolor{currentfill}%
\pgfsetlinewidth{1.003750pt}%
\definecolor{currentstroke}{rgb}{0.440137,0.811138,0.340967}%
\pgfsetstrokecolor{currentstroke}%
\pgfsetdash{}{0pt}%
\pgfpathmoveto{\pgfqpoint{3.703658in}{2.269457in}}%
\pgfpathcurveto{\pgfqpoint{3.709482in}{2.269457in}}{\pgfqpoint{3.715068in}{2.271771in}}{\pgfqpoint{3.719186in}{2.275889in}}%
\pgfpathcurveto{\pgfqpoint{3.723304in}{2.280007in}}{\pgfqpoint{3.725618in}{2.285593in}}{\pgfqpoint{3.725618in}{2.291417in}}%
\pgfpathcurveto{\pgfqpoint{3.725618in}{2.297241in}}{\pgfqpoint{3.723304in}{2.302828in}}{\pgfqpoint{3.719186in}{2.306946in}}%
\pgfpathcurveto{\pgfqpoint{3.715068in}{2.311064in}}{\pgfqpoint{3.709482in}{2.313378in}}{\pgfqpoint{3.703658in}{2.313378in}}%
\pgfpathcurveto{\pgfqpoint{3.697834in}{2.313378in}}{\pgfqpoint{3.692248in}{2.311064in}}{\pgfqpoint{3.688130in}{2.306946in}}%
\pgfpathcurveto{\pgfqpoint{3.684012in}{2.302828in}}{\pgfqpoint{3.681698in}{2.297241in}}{\pgfqpoint{3.681698in}{2.291417in}}%
\pgfpathcurveto{\pgfqpoint{3.681698in}{2.285593in}}{\pgfqpoint{3.684012in}{2.280007in}}{\pgfqpoint{3.688130in}{2.275889in}}%
\pgfpathcurveto{\pgfqpoint{3.692248in}{2.271771in}}{\pgfqpoint{3.697834in}{2.269457in}}{\pgfqpoint{3.703658in}{2.269457in}}%
\pgfpathlineto{\pgfqpoint{3.703658in}{2.269457in}}%
\pgfpathclose%
\pgfusepath{stroke,fill}%
\end{pgfscope}%
\begin{pgfscope}%
\pgfpathrectangle{\pgfqpoint{1.414736in}{0.548769in}}{\pgfqpoint{3.497747in}{7.253186in}}%
\pgfusepath{clip}%
\pgfsetbuttcap%
\pgfsetroundjoin%
\definecolor{currentfill}{rgb}{0.214000,0.722114,0.469588}%
\pgfsetfillcolor{currentfill}%
\pgfsetlinewidth{1.003750pt}%
\definecolor{currentstroke}{rgb}{0.214000,0.722114,0.469588}%
\pgfsetstrokecolor{currentstroke}%
\pgfsetdash{}{0pt}%
\pgfpathmoveto{\pgfqpoint{3.786335in}{2.740443in}}%
\pgfpathcurveto{\pgfqpoint{3.792159in}{2.740443in}}{\pgfqpoint{3.797745in}{2.742757in}}{\pgfqpoint{3.801863in}{2.746875in}}%
\pgfpathcurveto{\pgfqpoint{3.805981in}{2.750993in}}{\pgfqpoint{3.808295in}{2.756580in}}{\pgfqpoint{3.808295in}{2.762404in}}%
\pgfpathcurveto{\pgfqpoint{3.808295in}{2.768227in}}{\pgfqpoint{3.805981in}{2.773814in}}{\pgfqpoint{3.801863in}{2.777932in}}%
\pgfpathcurveto{\pgfqpoint{3.797745in}{2.782050in}}{\pgfqpoint{3.792159in}{2.784364in}}{\pgfqpoint{3.786335in}{2.784364in}}%
\pgfpathcurveto{\pgfqpoint{3.780511in}{2.784364in}}{\pgfqpoint{3.774925in}{2.782050in}}{\pgfqpoint{3.770807in}{2.777932in}}%
\pgfpathcurveto{\pgfqpoint{3.766689in}{2.773814in}}{\pgfqpoint{3.764375in}{2.768227in}}{\pgfqpoint{3.764375in}{2.762404in}}%
\pgfpathcurveto{\pgfqpoint{3.764375in}{2.756580in}}{\pgfqpoint{3.766689in}{2.750993in}}{\pgfqpoint{3.770807in}{2.746875in}}%
\pgfpathcurveto{\pgfqpoint{3.774925in}{2.742757in}}{\pgfqpoint{3.780511in}{2.740443in}}{\pgfqpoint{3.786335in}{2.740443in}}%
\pgfpathlineto{\pgfqpoint{3.786335in}{2.740443in}}%
\pgfpathclose%
\pgfusepath{stroke,fill}%
\end{pgfscope}%
\begin{pgfscope}%
\pgfpathrectangle{\pgfqpoint{1.414736in}{0.548769in}}{\pgfqpoint{3.497747in}{7.253186in}}%
\pgfusepath{clip}%
\pgfsetbuttcap%
\pgfsetroundjoin%
\definecolor{currentfill}{rgb}{0.535621,0.835785,0.281908}%
\pgfsetfillcolor{currentfill}%
\pgfsetlinewidth{1.003750pt}%
\definecolor{currentstroke}{rgb}{0.535621,0.835785,0.281908}%
\pgfsetstrokecolor{currentstroke}%
\pgfsetdash{}{0pt}%
\pgfpathmoveto{\pgfqpoint{3.361670in}{3.211429in}}%
\pgfpathcurveto{\pgfqpoint{3.367494in}{3.211429in}}{\pgfqpoint{3.373080in}{3.213743in}}{\pgfqpoint{3.377198in}{3.217861in}}%
\pgfpathcurveto{\pgfqpoint{3.381316in}{3.221980in}}{\pgfqpoint{3.383630in}{3.227566in}}{\pgfqpoint{3.383630in}{3.233390in}}%
\pgfpathcurveto{\pgfqpoint{3.383630in}{3.239214in}}{\pgfqpoint{3.381316in}{3.244800in}}{\pgfqpoint{3.377198in}{3.248918in}}%
\pgfpathcurveto{\pgfqpoint{3.373080in}{3.253036in}}{\pgfqpoint{3.367494in}{3.255350in}}{\pgfqpoint{3.361670in}{3.255350in}}%
\pgfpathcurveto{\pgfqpoint{3.355846in}{3.255350in}}{\pgfqpoint{3.350260in}{3.253036in}}{\pgfqpoint{3.346141in}{3.248918in}}%
\pgfpathcurveto{\pgfqpoint{3.342023in}{3.244800in}}{\pgfqpoint{3.339709in}{3.239214in}}{\pgfqpoint{3.339709in}{3.233390in}}%
\pgfpathcurveto{\pgfqpoint{3.339709in}{3.227566in}}{\pgfqpoint{3.342023in}{3.221980in}}{\pgfqpoint{3.346141in}{3.217861in}}%
\pgfpathcurveto{\pgfqpoint{3.350260in}{3.213743in}}{\pgfqpoint{3.355846in}{3.211429in}}{\pgfqpoint{3.361670in}{3.211429in}}%
\pgfpathlineto{\pgfqpoint{3.361670in}{3.211429in}}%
\pgfpathclose%
\pgfusepath{stroke,fill}%
\end{pgfscope}%
\begin{pgfscope}%
\pgfpathrectangle{\pgfqpoint{1.414736in}{0.548769in}}{\pgfqpoint{3.497747in}{7.253186in}}%
\pgfusepath{clip}%
\pgfsetbuttcap%
\pgfsetroundjoin%
\definecolor{currentfill}{rgb}{0.709898,0.868751,0.169257}%
\pgfsetfillcolor{currentfill}%
\pgfsetlinewidth{1.003750pt}%
\definecolor{currentstroke}{rgb}{0.709898,0.868751,0.169257}%
\pgfsetstrokecolor{currentstroke}%
\pgfsetdash{}{0pt}%
\pgfpathmoveto{\pgfqpoint{3.913947in}{3.682415in}}%
\pgfpathcurveto{\pgfqpoint{3.919771in}{3.682415in}}{\pgfqpoint{3.925357in}{3.684729in}}{\pgfqpoint{3.929475in}{3.688847in}}%
\pgfpathcurveto{\pgfqpoint{3.933594in}{3.692966in}}{\pgfqpoint{3.935907in}{3.698552in}}{\pgfqpoint{3.935907in}{3.704376in}}%
\pgfpathcurveto{\pgfqpoint{3.935907in}{3.710200in}}{\pgfqpoint{3.933594in}{3.715786in}}{\pgfqpoint{3.929475in}{3.719904in}}%
\pgfpathcurveto{\pgfqpoint{3.925357in}{3.724022in}}{\pgfqpoint{3.919771in}{3.726336in}}{\pgfqpoint{3.913947in}{3.726336in}}%
\pgfpathcurveto{\pgfqpoint{3.908123in}{3.726336in}}{\pgfqpoint{3.902537in}{3.724022in}}{\pgfqpoint{3.898419in}{3.719904in}}%
\pgfpathcurveto{\pgfqpoint{3.894301in}{3.715786in}}{\pgfqpoint{3.891987in}{3.710200in}}{\pgfqpoint{3.891987in}{3.704376in}}%
\pgfpathcurveto{\pgfqpoint{3.891987in}{3.698552in}}{\pgfqpoint{3.894301in}{3.692966in}}{\pgfqpoint{3.898419in}{3.688847in}}%
\pgfpathcurveto{\pgfqpoint{3.902537in}{3.684729in}}{\pgfqpoint{3.908123in}{3.682415in}}{\pgfqpoint{3.913947in}{3.682415in}}%
\pgfpathlineto{\pgfqpoint{3.913947in}{3.682415in}}%
\pgfpathclose%
\pgfusepath{stroke,fill}%
\end{pgfscope}%
\begin{pgfscope}%
\pgfpathrectangle{\pgfqpoint{1.414736in}{0.548769in}}{\pgfqpoint{3.497747in}{7.253186in}}%
\pgfusepath{clip}%
\pgfsetbuttcap%
\pgfsetroundjoin%
\definecolor{currentfill}{rgb}{0.993248,0.906157,0.143936}%
\pgfsetfillcolor{currentfill}%
\pgfsetlinewidth{1.003750pt}%
\definecolor{currentstroke}{rgb}{0.993248,0.906157,0.143936}%
\pgfsetstrokecolor{currentstroke}%
\pgfsetdash{}{0pt}%
\pgfpathmoveto{\pgfqpoint{3.494121in}{4.153402in}}%
\pgfpathcurveto{\pgfqpoint{3.499945in}{4.153402in}}{\pgfqpoint{3.505532in}{4.155715in}}{\pgfqpoint{3.509650in}{4.159834in}}%
\pgfpathcurveto{\pgfqpoint{3.513768in}{4.163952in}}{\pgfqpoint{3.516082in}{4.169538in}}{\pgfqpoint{3.516082in}{4.175362in}}%
\pgfpathcurveto{\pgfqpoint{3.516082in}{4.181186in}}{\pgfqpoint{3.513768in}{4.186772in}}{\pgfqpoint{3.509650in}{4.190890in}}%
\pgfpathcurveto{\pgfqpoint{3.505532in}{4.195008in}}{\pgfqpoint{3.499945in}{4.197322in}}{\pgfqpoint{3.494121in}{4.197322in}}%
\pgfpathcurveto{\pgfqpoint{3.488298in}{4.197322in}}{\pgfqpoint{3.482711in}{4.195008in}}{\pgfqpoint{3.478593in}{4.190890in}}%
\pgfpathcurveto{\pgfqpoint{3.474475in}{4.186772in}}{\pgfqpoint{3.472161in}{4.181186in}}{\pgfqpoint{3.472161in}{4.175362in}}%
\pgfpathcurveto{\pgfqpoint{3.472161in}{4.169538in}}{\pgfqpoint{3.474475in}{4.163952in}}{\pgfqpoint{3.478593in}{4.159834in}}%
\pgfpathcurveto{\pgfqpoint{3.482711in}{4.155715in}}{\pgfqpoint{3.488298in}{4.153402in}}{\pgfqpoint{3.494121in}{4.153402in}}%
\pgfpathlineto{\pgfqpoint{3.494121in}{4.153402in}}%
\pgfpathclose%
\pgfusepath{stroke,fill}%
\end{pgfscope}%
\begin{pgfscope}%
\pgfpathrectangle{\pgfqpoint{1.414736in}{0.548769in}}{\pgfqpoint{3.497747in}{7.253186in}}%
\pgfusepath{clip}%
\pgfsetbuttcap%
\pgfsetroundjoin%
\definecolor{currentfill}{rgb}{0.595839,0.848717,0.243329}%
\pgfsetfillcolor{currentfill}%
\pgfsetlinewidth{1.003750pt}%
\definecolor{currentstroke}{rgb}{0.595839,0.848717,0.243329}%
\pgfsetstrokecolor{currentstroke}%
\pgfsetdash{}{0pt}%
\pgfpathmoveto{\pgfqpoint{3.090964in}{4.624388in}}%
\pgfpathcurveto{\pgfqpoint{3.096788in}{4.624388in}}{\pgfqpoint{3.102374in}{4.626702in}}{\pgfqpoint{3.106492in}{4.630820in}}%
\pgfpathcurveto{\pgfqpoint{3.110610in}{4.634938in}}{\pgfqpoint{3.112924in}{4.640524in}}{\pgfqpoint{3.112924in}{4.646348in}}%
\pgfpathcurveto{\pgfqpoint{3.112924in}{4.652172in}}{\pgfqpoint{3.110610in}{4.657758in}}{\pgfqpoint{3.106492in}{4.661876in}}%
\pgfpathcurveto{\pgfqpoint{3.102374in}{4.665994in}}{\pgfqpoint{3.096788in}{4.668308in}}{\pgfqpoint{3.090964in}{4.668308in}}%
\pgfpathcurveto{\pgfqpoint{3.085140in}{4.668308in}}{\pgfqpoint{3.079554in}{4.665994in}}{\pgfqpoint{3.075436in}{4.661876in}}%
\pgfpathcurveto{\pgfqpoint{3.071318in}{4.657758in}}{\pgfqpoint{3.069004in}{4.652172in}}{\pgfqpoint{3.069004in}{4.646348in}}%
\pgfpathcurveto{\pgfqpoint{3.069004in}{4.640524in}}{\pgfqpoint{3.071318in}{4.634938in}}{\pgfqpoint{3.075436in}{4.630820in}}%
\pgfpathcurveto{\pgfqpoint{3.079554in}{4.626702in}}{\pgfqpoint{3.085140in}{4.624388in}}{\pgfqpoint{3.090964in}{4.624388in}}%
\pgfpathlineto{\pgfqpoint{3.090964in}{4.624388in}}%
\pgfpathclose%
\pgfusepath{stroke,fill}%
\end{pgfscope}%
\begin{pgfscope}%
\pgfpathrectangle{\pgfqpoint{1.414736in}{0.548769in}}{\pgfqpoint{3.497747in}{7.253186in}}%
\pgfusepath{clip}%
\pgfsetbuttcap%
\pgfsetroundjoin%
\definecolor{currentfill}{rgb}{0.280868,0.160771,0.472899}%
\pgfsetfillcolor{currentfill}%
\pgfsetlinewidth{1.003750pt}%
\definecolor{currentstroke}{rgb}{0.280868,0.160771,0.472899}%
\pgfsetstrokecolor{currentstroke}%
\pgfsetdash{}{0pt}%
\pgfpathmoveto{\pgfqpoint{3.284220in}{5.095374in}}%
\pgfpathcurveto{\pgfqpoint{3.290044in}{5.095374in}}{\pgfqpoint{3.295630in}{5.097688in}}{\pgfqpoint{3.299748in}{5.101806in}}%
\pgfpathcurveto{\pgfqpoint{3.303867in}{5.105924in}}{\pgfqpoint{3.306180in}{5.111510in}}{\pgfqpoint{3.306180in}{5.117334in}}%
\pgfpathcurveto{\pgfqpoint{3.306180in}{5.123158in}}{\pgfqpoint{3.303867in}{5.128744in}}{\pgfqpoint{3.299748in}{5.132862in}}%
\pgfpathcurveto{\pgfqpoint{3.295630in}{5.136980in}}{\pgfqpoint{3.290044in}{5.139294in}}{\pgfqpoint{3.284220in}{5.139294in}}%
\pgfpathcurveto{\pgfqpoint{3.278396in}{5.139294in}}{\pgfqpoint{3.272810in}{5.136980in}}{\pgfqpoint{3.268692in}{5.132862in}}%
\pgfpathcurveto{\pgfqpoint{3.264574in}{5.128744in}}{\pgfqpoint{3.262260in}{5.123158in}}{\pgfqpoint{3.262260in}{5.117334in}}%
\pgfpathcurveto{\pgfqpoint{3.262260in}{5.111510in}}{\pgfqpoint{3.264574in}{5.105924in}}{\pgfqpoint{3.268692in}{5.101806in}}%
\pgfpathcurveto{\pgfqpoint{3.272810in}{5.097688in}}{\pgfqpoint{3.278396in}{5.095374in}}{\pgfqpoint{3.284220in}{5.095374in}}%
\pgfpathlineto{\pgfqpoint{3.284220in}{5.095374in}}%
\pgfpathclose%
\pgfusepath{stroke,fill}%
\end{pgfscope}%
\begin{pgfscope}%
\pgfpathrectangle{\pgfqpoint{1.414736in}{0.548769in}}{\pgfqpoint{3.497747in}{7.253186in}}%
\pgfusepath{clip}%
\pgfsetbuttcap%
\pgfsetroundjoin%
\definecolor{currentfill}{rgb}{0.168126,0.459988,0.558082}%
\pgfsetfillcolor{currentfill}%
\pgfsetlinewidth{1.003750pt}%
\definecolor{currentstroke}{rgb}{0.168126,0.459988,0.558082}%
\pgfsetstrokecolor{currentstroke}%
\pgfsetdash{}{0pt}%
\pgfpathmoveto{\pgfqpoint{3.668951in}{5.566360in}}%
\pgfpathcurveto{\pgfqpoint{3.674774in}{5.566360in}}{\pgfqpoint{3.680361in}{5.568674in}}{\pgfqpoint{3.684479in}{5.572792in}}%
\pgfpathcurveto{\pgfqpoint{3.688597in}{5.576910in}}{\pgfqpoint{3.690911in}{5.582496in}}{\pgfqpoint{3.690911in}{5.588320in}}%
\pgfpathcurveto{\pgfqpoint{3.690911in}{5.594144in}}{\pgfqpoint{3.688597in}{5.599730in}}{\pgfqpoint{3.684479in}{5.603848in}}%
\pgfpathcurveto{\pgfqpoint{3.680361in}{5.607967in}}{\pgfqpoint{3.674774in}{5.610280in}}{\pgfqpoint{3.668951in}{5.610280in}}%
\pgfpathcurveto{\pgfqpoint{3.663127in}{5.610280in}}{\pgfqpoint{3.657540in}{5.607967in}}{\pgfqpoint{3.653422in}{5.603848in}}%
\pgfpathcurveto{\pgfqpoint{3.649304in}{5.599730in}}{\pgfqpoint{3.646990in}{5.594144in}}{\pgfqpoint{3.646990in}{5.588320in}}%
\pgfpathcurveto{\pgfqpoint{3.646990in}{5.582496in}}{\pgfqpoint{3.649304in}{5.576910in}}{\pgfqpoint{3.653422in}{5.572792in}}%
\pgfpathcurveto{\pgfqpoint{3.657540in}{5.568674in}}{\pgfqpoint{3.663127in}{5.566360in}}{\pgfqpoint{3.668951in}{5.566360in}}%
\pgfpathlineto{\pgfqpoint{3.668951in}{5.566360in}}%
\pgfpathclose%
\pgfusepath{stroke,fill}%
\end{pgfscope}%
\begin{pgfscope}%
\pgfpathrectangle{\pgfqpoint{1.414736in}{0.548769in}}{\pgfqpoint{3.497747in}{7.253186in}}%
\pgfusepath{clip}%
\pgfsetbuttcap%
\pgfsetroundjoin%
\definecolor{currentfill}{rgb}{0.177423,0.437527,0.557565}%
\pgfsetfillcolor{currentfill}%
\pgfsetlinewidth{1.003750pt}%
\definecolor{currentstroke}{rgb}{0.177423,0.437527,0.557565}%
\pgfsetstrokecolor{currentstroke}%
\pgfsetdash{}{0pt}%
\pgfpathmoveto{\pgfqpoint{3.891065in}{6.037346in}}%
\pgfpathcurveto{\pgfqpoint{3.896889in}{6.037346in}}{\pgfqpoint{3.902475in}{6.039660in}}{\pgfqpoint{3.906593in}{6.043778in}}%
\pgfpathcurveto{\pgfqpoint{3.910711in}{6.047896in}}{\pgfqpoint{3.913025in}{6.053482in}}{\pgfqpoint{3.913025in}{6.059306in}}%
\pgfpathcurveto{\pgfqpoint{3.913025in}{6.065130in}}{\pgfqpoint{3.910711in}{6.070716in}}{\pgfqpoint{3.906593in}{6.074835in}}%
\pgfpathcurveto{\pgfqpoint{3.902475in}{6.078953in}}{\pgfqpoint{3.896889in}{6.081267in}}{\pgfqpoint{3.891065in}{6.081267in}}%
\pgfpathcurveto{\pgfqpoint{3.885241in}{6.081267in}}{\pgfqpoint{3.879655in}{6.078953in}}{\pgfqpoint{3.875536in}{6.074835in}}%
\pgfpathcurveto{\pgfqpoint{3.871418in}{6.070716in}}{\pgfqpoint{3.869104in}{6.065130in}}{\pgfqpoint{3.869104in}{6.059306in}}%
\pgfpathcurveto{\pgfqpoint{3.869104in}{6.053482in}}{\pgfqpoint{3.871418in}{6.047896in}}{\pgfqpoint{3.875536in}{6.043778in}}%
\pgfpathcurveto{\pgfqpoint{3.879655in}{6.039660in}}{\pgfqpoint{3.885241in}{6.037346in}}{\pgfqpoint{3.891065in}{6.037346in}}%
\pgfpathlineto{\pgfqpoint{3.891065in}{6.037346in}}%
\pgfpathclose%
\pgfusepath{stroke,fill}%
\end{pgfscope}%
\begin{pgfscope}%
\pgfpathrectangle{\pgfqpoint{1.414736in}{0.548769in}}{\pgfqpoint{3.497747in}{7.253186in}}%
\pgfusepath{clip}%
\pgfsetbuttcap%
\pgfsetroundjoin%
\definecolor{currentfill}{rgb}{0.866013,0.889868,0.095953}%
\pgfsetfillcolor{currentfill}%
\pgfsetlinewidth{1.003750pt}%
\definecolor{currentstroke}{rgb}{0.866013,0.889868,0.095953}%
\pgfsetstrokecolor{currentstroke}%
\pgfsetdash{}{0pt}%
\pgfpathmoveto{\pgfqpoint{3.999622in}{6.508332in}}%
\pgfpathcurveto{\pgfqpoint{4.005446in}{6.508332in}}{\pgfqpoint{4.011032in}{6.510646in}}{\pgfqpoint{4.015150in}{6.514764in}}%
\pgfpathcurveto{\pgfqpoint{4.019268in}{6.518882in}}{\pgfqpoint{4.021582in}{6.524468in}}{\pgfqpoint{4.021582in}{6.530292in}}%
\pgfpathcurveto{\pgfqpoint{4.021582in}{6.536116in}}{\pgfqpoint{4.019268in}{6.541702in}}{\pgfqpoint{4.015150in}{6.545821in}}%
\pgfpathcurveto{\pgfqpoint{4.011032in}{6.549939in}}{\pgfqpoint{4.005446in}{6.552253in}}{\pgfqpoint{3.999622in}{6.552253in}}%
\pgfpathcurveto{\pgfqpoint{3.993798in}{6.552253in}}{\pgfqpoint{3.988212in}{6.549939in}}{\pgfqpoint{3.984093in}{6.545821in}}%
\pgfpathcurveto{\pgfqpoint{3.979975in}{6.541702in}}{\pgfqpoint{3.977661in}{6.536116in}}{\pgfqpoint{3.977661in}{6.530292in}}%
\pgfpathcurveto{\pgfqpoint{3.977661in}{6.524468in}}{\pgfqpoint{3.979975in}{6.518882in}}{\pgfqpoint{3.984093in}{6.514764in}}%
\pgfpathcurveto{\pgfqpoint{3.988212in}{6.510646in}}{\pgfqpoint{3.993798in}{6.508332in}}{\pgfqpoint{3.999622in}{6.508332in}}%
\pgfpathlineto{\pgfqpoint{3.999622in}{6.508332in}}%
\pgfpathclose%
\pgfusepath{stroke,fill}%
\end{pgfscope}%
\begin{pgfscope}%
\pgfpathrectangle{\pgfqpoint{1.414736in}{0.548769in}}{\pgfqpoint{3.497747in}{7.253186in}}%
\pgfusepath{clip}%
\pgfsetbuttcap%
\pgfsetroundjoin%
\definecolor{currentfill}{rgb}{0.730889,0.871916,0.156029}%
\pgfsetfillcolor{currentfill}%
\pgfsetlinewidth{1.003750pt}%
\definecolor{currentstroke}{rgb}{0.730889,0.871916,0.156029}%
\pgfsetstrokecolor{currentstroke}%
\pgfsetdash{}{0pt}%
\pgfpathmoveto{\pgfqpoint{3.691996in}{6.979318in}}%
\pgfpathcurveto{\pgfqpoint{3.697820in}{6.979318in}}{\pgfqpoint{3.703406in}{6.981632in}}{\pgfqpoint{3.707524in}{6.985750in}}%
\pgfpathcurveto{\pgfqpoint{3.711642in}{6.989868in}}{\pgfqpoint{3.713956in}{6.995455in}}{\pgfqpoint{3.713956in}{7.001278in}}%
\pgfpathcurveto{\pgfqpoint{3.713956in}{7.007102in}}{\pgfqpoint{3.711642in}{7.012689in}}{\pgfqpoint{3.707524in}{7.016807in}}%
\pgfpathcurveto{\pgfqpoint{3.703406in}{7.020925in}}{\pgfqpoint{3.697820in}{7.023239in}}{\pgfqpoint{3.691996in}{7.023239in}}%
\pgfpathcurveto{\pgfqpoint{3.686172in}{7.023239in}}{\pgfqpoint{3.680586in}{7.020925in}}{\pgfqpoint{3.676468in}{7.016807in}}%
\pgfpathcurveto{\pgfqpoint{3.672350in}{7.012689in}}{\pgfqpoint{3.670036in}{7.007102in}}{\pgfqpoint{3.670036in}{7.001278in}}%
\pgfpathcurveto{\pgfqpoint{3.670036in}{6.995455in}}{\pgfqpoint{3.672350in}{6.989868in}}{\pgfqpoint{3.676468in}{6.985750in}}%
\pgfpathcurveto{\pgfqpoint{3.680586in}{6.981632in}}{\pgfqpoint{3.686172in}{6.979318in}}{\pgfqpoint{3.691996in}{6.979318in}}%
\pgfpathlineto{\pgfqpoint{3.691996in}{6.979318in}}%
\pgfpathclose%
\pgfusepath{stroke,fill}%
\end{pgfscope}%
\begin{pgfscope}%
\pgfpathrectangle{\pgfqpoint{1.414736in}{0.548769in}}{\pgfqpoint{3.497747in}{7.253186in}}%
\pgfusepath{clip}%
\pgfsetbuttcap%
\pgfsetroundjoin%
\definecolor{currentfill}{rgb}{0.636902,0.856542,0.216620}%
\pgfsetfillcolor{currentfill}%
\pgfsetlinewidth{1.003750pt}%
\definecolor{currentstroke}{rgb}{0.636902,0.856542,0.216620}%
\pgfsetstrokecolor{currentstroke}%
\pgfsetdash{}{0pt}%
\pgfpathmoveto{\pgfqpoint{3.437661in}{7.450304in}}%
\pgfpathcurveto{\pgfqpoint{3.443485in}{7.450304in}}{\pgfqpoint{3.449071in}{7.452618in}}{\pgfqpoint{3.453189in}{7.456736in}}%
\pgfpathcurveto{\pgfqpoint{3.457308in}{7.460854in}}{\pgfqpoint{3.459621in}{7.466441in}}{\pgfqpoint{3.459621in}{7.472265in}}%
\pgfpathcurveto{\pgfqpoint{3.459621in}{7.478089in}}{\pgfqpoint{3.457308in}{7.483675in}}{\pgfqpoint{3.453189in}{7.487793in}}%
\pgfpathcurveto{\pgfqpoint{3.449071in}{7.491911in}}{\pgfqpoint{3.443485in}{7.494225in}}{\pgfqpoint{3.437661in}{7.494225in}}%
\pgfpathcurveto{\pgfqpoint{3.431837in}{7.494225in}}{\pgfqpoint{3.426251in}{7.491911in}}{\pgfqpoint{3.422133in}{7.487793in}}%
\pgfpathcurveto{\pgfqpoint{3.418015in}{7.483675in}}{\pgfqpoint{3.415701in}{7.478089in}}{\pgfqpoint{3.415701in}{7.472265in}}%
\pgfpathcurveto{\pgfqpoint{3.415701in}{7.466441in}}{\pgfqpoint{3.418015in}{7.460854in}}{\pgfqpoint{3.422133in}{7.456736in}}%
\pgfpathcurveto{\pgfqpoint{3.426251in}{7.452618in}}{\pgfqpoint{3.431837in}{7.450304in}}{\pgfqpoint{3.437661in}{7.450304in}}%
\pgfpathlineto{\pgfqpoint{3.437661in}{7.450304in}}%
\pgfpathclose%
\pgfusepath{stroke,fill}%
\end{pgfscope}%
\begin{pgfscope}%
\pgfpathrectangle{\pgfqpoint{1.414736in}{0.548769in}}{\pgfqpoint{3.497747in}{7.253186in}}%
\pgfusepath{clip}%
\pgfsetbuttcap%
\pgfsetroundjoin%
\definecolor{currentfill}{rgb}{0.983868,0.904867,0.136897}%
\pgfsetfillcolor{currentfill}%
\pgfsetlinewidth{1.003750pt}%
\definecolor{currentstroke}{rgb}{0.983868,0.904867,0.136897}%
\pgfsetstrokecolor{currentstroke}%
\pgfsetdash{}{0pt}%
\pgfpathmoveto{\pgfqpoint{3.568748in}{0.856499in}}%
\pgfpathcurveto{\pgfqpoint{3.574572in}{0.856499in}}{\pgfqpoint{3.580158in}{0.858813in}}{\pgfqpoint{3.584276in}{0.862931in}}%
\pgfpathcurveto{\pgfqpoint{3.588395in}{0.867049in}}{\pgfqpoint{3.590708in}{0.872635in}}{\pgfqpoint{3.590708in}{0.878459in}}%
\pgfpathcurveto{\pgfqpoint{3.590708in}{0.884283in}}{\pgfqpoint{3.588395in}{0.889869in}}{\pgfqpoint{3.584276in}{0.893987in}}%
\pgfpathcurveto{\pgfqpoint{3.580158in}{0.898105in}}{\pgfqpoint{3.574572in}{0.900419in}}{\pgfqpoint{3.568748in}{0.900419in}}%
\pgfpathcurveto{\pgfqpoint{3.562924in}{0.900419in}}{\pgfqpoint{3.557338in}{0.898105in}}{\pgfqpoint{3.553220in}{0.893987in}}%
\pgfpathcurveto{\pgfqpoint{3.549102in}{0.889869in}}{\pgfqpoint{3.546788in}{0.884283in}}{\pgfqpoint{3.546788in}{0.878459in}}%
\pgfpathcurveto{\pgfqpoint{3.546788in}{0.872635in}}{\pgfqpoint{3.549102in}{0.867049in}}{\pgfqpoint{3.553220in}{0.862931in}}%
\pgfpathcurveto{\pgfqpoint{3.557338in}{0.858813in}}{\pgfqpoint{3.562924in}{0.856499in}}{\pgfqpoint{3.568748in}{0.856499in}}%
\pgfpathlineto{\pgfqpoint{3.568748in}{0.856499in}}%
\pgfpathclose%
\pgfusepath{stroke,fill}%
\end{pgfscope}%
\begin{pgfscope}%
\pgfpathrectangle{\pgfqpoint{1.414736in}{0.548769in}}{\pgfqpoint{3.497747in}{7.253186in}}%
\pgfusepath{clip}%
\pgfsetbuttcap%
\pgfsetroundjoin%
\definecolor{currentfill}{rgb}{0.688944,0.865448,0.182725}%
\pgfsetfillcolor{currentfill}%
\pgfsetlinewidth{1.003750pt}%
\definecolor{currentstroke}{rgb}{0.688944,0.865448,0.182725}%
\pgfsetstrokecolor{currentstroke}%
\pgfsetdash{}{0pt}%
\pgfpathmoveto{\pgfqpoint{3.419685in}{1.327485in}}%
\pgfpathcurveto{\pgfqpoint{3.425509in}{1.327485in}}{\pgfqpoint{3.431095in}{1.329799in}}{\pgfqpoint{3.435213in}{1.333917in}}%
\pgfpathcurveto{\pgfqpoint{3.439331in}{1.338035in}}{\pgfqpoint{3.441645in}{1.343621in}}{\pgfqpoint{3.441645in}{1.349445in}}%
\pgfpathcurveto{\pgfqpoint{3.441645in}{1.355269in}}{\pgfqpoint{3.439331in}{1.360855in}}{\pgfqpoint{3.435213in}{1.364973in}}%
\pgfpathcurveto{\pgfqpoint{3.431095in}{1.369092in}}{\pgfqpoint{3.425509in}{1.371405in}}{\pgfqpoint{3.419685in}{1.371405in}}%
\pgfpathcurveto{\pgfqpoint{3.413861in}{1.371405in}}{\pgfqpoint{3.408275in}{1.369092in}}{\pgfqpoint{3.404157in}{1.364973in}}%
\pgfpathcurveto{\pgfqpoint{3.400039in}{1.360855in}}{\pgfqpoint{3.397725in}{1.355269in}}{\pgfqpoint{3.397725in}{1.349445in}}%
\pgfpathcurveto{\pgfqpoint{3.397725in}{1.343621in}}{\pgfqpoint{3.400039in}{1.338035in}}{\pgfqpoint{3.404157in}{1.333917in}}%
\pgfpathcurveto{\pgfqpoint{3.408275in}{1.329799in}}{\pgfqpoint{3.413861in}{1.327485in}}{\pgfqpoint{3.419685in}{1.327485in}}%
\pgfpathlineto{\pgfqpoint{3.419685in}{1.327485in}}%
\pgfpathclose%
\pgfusepath{stroke,fill}%
\end{pgfscope}%
\begin{pgfscope}%
\pgfpathrectangle{\pgfqpoint{1.414736in}{0.548769in}}{\pgfqpoint{3.497747in}{7.253186in}}%
\pgfusepath{clip}%
\pgfsetbuttcap%
\pgfsetroundjoin%
\definecolor{currentfill}{rgb}{0.699415,0.867117,0.175971}%
\pgfsetfillcolor{currentfill}%
\pgfsetlinewidth{1.003750pt}%
\definecolor{currentstroke}{rgb}{0.699415,0.867117,0.175971}%
\pgfsetstrokecolor{currentstroke}%
\pgfsetdash{}{0pt}%
\pgfpathmoveto{\pgfqpoint{3.753807in}{1.798471in}}%
\pgfpathcurveto{\pgfqpoint{3.759631in}{1.798471in}}{\pgfqpoint{3.765217in}{1.800785in}}{\pgfqpoint{3.769335in}{1.804903in}}%
\pgfpathcurveto{\pgfqpoint{3.773453in}{1.809021in}}{\pgfqpoint{3.775767in}{1.814607in}}{\pgfqpoint{3.775767in}{1.820431in}}%
\pgfpathcurveto{\pgfqpoint{3.775767in}{1.826255in}}{\pgfqpoint{3.773453in}{1.831841in}}{\pgfqpoint{3.769335in}{1.835960in}}%
\pgfpathcurveto{\pgfqpoint{3.765217in}{1.840078in}}{\pgfqpoint{3.759631in}{1.842392in}}{\pgfqpoint{3.753807in}{1.842392in}}%
\pgfpathcurveto{\pgfqpoint{3.747983in}{1.842392in}}{\pgfqpoint{3.742396in}{1.840078in}}{\pgfqpoint{3.738278in}{1.835960in}}%
\pgfpathcurveto{\pgfqpoint{3.734160in}{1.831841in}}{\pgfqpoint{3.731846in}{1.826255in}}{\pgfqpoint{3.731846in}{1.820431in}}%
\pgfpathcurveto{\pgfqpoint{3.731846in}{1.814607in}}{\pgfqpoint{3.734160in}{1.809021in}}{\pgfqpoint{3.738278in}{1.804903in}}%
\pgfpathcurveto{\pgfqpoint{3.742396in}{1.800785in}}{\pgfqpoint{3.747983in}{1.798471in}}{\pgfqpoint{3.753807in}{1.798471in}}%
\pgfpathlineto{\pgfqpoint{3.753807in}{1.798471in}}%
\pgfpathclose%
\pgfusepath{stroke,fill}%
\end{pgfscope}%
\begin{pgfscope}%
\pgfpathrectangle{\pgfqpoint{1.414736in}{0.548769in}}{\pgfqpoint{3.497747in}{7.253186in}}%
\pgfusepath{clip}%
\pgfsetbuttcap%
\pgfsetroundjoin%
\definecolor{currentfill}{rgb}{0.993248,0.906157,0.143936}%
\pgfsetfillcolor{currentfill}%
\pgfsetlinewidth{1.003750pt}%
\definecolor{currentstroke}{rgb}{0.993248,0.906157,0.143936}%
\pgfsetstrokecolor{currentstroke}%
\pgfsetdash{}{0pt}%
\pgfpathmoveto{\pgfqpoint{3.522702in}{2.269457in}}%
\pgfpathcurveto{\pgfqpoint{3.528526in}{2.269457in}}{\pgfqpoint{3.534112in}{2.271771in}}{\pgfqpoint{3.538230in}{2.275889in}}%
\pgfpathcurveto{\pgfqpoint{3.542348in}{2.280007in}}{\pgfqpoint{3.544662in}{2.285593in}}{\pgfqpoint{3.544662in}{2.291417in}}%
\pgfpathcurveto{\pgfqpoint{3.544662in}{2.297241in}}{\pgfqpoint{3.542348in}{2.302828in}}{\pgfqpoint{3.538230in}{2.306946in}}%
\pgfpathcurveto{\pgfqpoint{3.534112in}{2.311064in}}{\pgfqpoint{3.528526in}{2.313378in}}{\pgfqpoint{3.522702in}{2.313378in}}%
\pgfpathcurveto{\pgfqpoint{3.516878in}{2.313378in}}{\pgfqpoint{3.511292in}{2.311064in}}{\pgfqpoint{3.507174in}{2.306946in}}%
\pgfpathcurveto{\pgfqpoint{3.503056in}{2.302828in}}{\pgfqpoint{3.500742in}{2.297241in}}{\pgfqpoint{3.500742in}{2.291417in}}%
\pgfpathcurveto{\pgfqpoint{3.500742in}{2.285593in}}{\pgfqpoint{3.503056in}{2.280007in}}{\pgfqpoint{3.507174in}{2.275889in}}%
\pgfpathcurveto{\pgfqpoint{3.511292in}{2.271771in}}{\pgfqpoint{3.516878in}{2.269457in}}{\pgfqpoint{3.522702in}{2.269457in}}%
\pgfpathlineto{\pgfqpoint{3.522702in}{2.269457in}}%
\pgfpathclose%
\pgfusepath{stroke,fill}%
\end{pgfscope}%
\begin{pgfscope}%
\pgfpathrectangle{\pgfqpoint{1.414736in}{0.548769in}}{\pgfqpoint{3.497747in}{7.253186in}}%
\pgfusepath{clip}%
\pgfsetbuttcap%
\pgfsetroundjoin%
\definecolor{currentfill}{rgb}{0.352360,0.783011,0.392636}%
\pgfsetfillcolor{currentfill}%
\pgfsetlinewidth{1.003750pt}%
\definecolor{currentstroke}{rgb}{0.352360,0.783011,0.392636}%
\pgfsetstrokecolor{currentstroke}%
\pgfsetdash{}{0pt}%
\pgfpathmoveto{\pgfqpoint{3.216253in}{2.740443in}}%
\pgfpathcurveto{\pgfqpoint{3.222077in}{2.740443in}}{\pgfqpoint{3.227664in}{2.742757in}}{\pgfqpoint{3.231782in}{2.746875in}}%
\pgfpathcurveto{\pgfqpoint{3.235900in}{2.750993in}}{\pgfqpoint{3.238214in}{2.756580in}}{\pgfqpoint{3.238214in}{2.762404in}}%
\pgfpathcurveto{\pgfqpoint{3.238214in}{2.768227in}}{\pgfqpoint{3.235900in}{2.773814in}}{\pgfqpoint{3.231782in}{2.777932in}}%
\pgfpathcurveto{\pgfqpoint{3.227664in}{2.782050in}}{\pgfqpoint{3.222077in}{2.784364in}}{\pgfqpoint{3.216253in}{2.784364in}}%
\pgfpathcurveto{\pgfqpoint{3.210429in}{2.784364in}}{\pgfqpoint{3.204843in}{2.782050in}}{\pgfqpoint{3.200725in}{2.777932in}}%
\pgfpathcurveto{\pgfqpoint{3.196607in}{2.773814in}}{\pgfqpoint{3.194293in}{2.768227in}}{\pgfqpoint{3.194293in}{2.762404in}}%
\pgfpathcurveto{\pgfqpoint{3.194293in}{2.756580in}}{\pgfqpoint{3.196607in}{2.750993in}}{\pgfqpoint{3.200725in}{2.746875in}}%
\pgfpathcurveto{\pgfqpoint{3.204843in}{2.742757in}}{\pgfqpoint{3.210429in}{2.740443in}}{\pgfqpoint{3.216253in}{2.740443in}}%
\pgfpathlineto{\pgfqpoint{3.216253in}{2.740443in}}%
\pgfpathclose%
\pgfusepath{stroke,fill}%
\end{pgfscope}%
\begin{pgfscope}%
\pgfpathrectangle{\pgfqpoint{1.414736in}{0.548769in}}{\pgfqpoint{3.497747in}{7.253186in}}%
\pgfusepath{clip}%
\pgfsetbuttcap%
\pgfsetroundjoin%
\definecolor{currentfill}{rgb}{0.545524,0.838039,0.275626}%
\pgfsetfillcolor{currentfill}%
\pgfsetlinewidth{1.003750pt}%
\definecolor{currentstroke}{rgb}{0.545524,0.838039,0.275626}%
\pgfsetstrokecolor{currentstroke}%
\pgfsetdash{}{0pt}%
\pgfpathmoveto{\pgfqpoint{3.798420in}{3.211429in}}%
\pgfpathcurveto{\pgfqpoint{3.804244in}{3.211429in}}{\pgfqpoint{3.809830in}{3.213743in}}{\pgfqpoint{3.813949in}{3.217861in}}%
\pgfpathcurveto{\pgfqpoint{3.818067in}{3.221980in}}{\pgfqpoint{3.820381in}{3.227566in}}{\pgfqpoint{3.820381in}{3.233390in}}%
\pgfpathcurveto{\pgfqpoint{3.820381in}{3.239214in}}{\pgfqpoint{3.818067in}{3.244800in}}{\pgfqpoint{3.813949in}{3.248918in}}%
\pgfpathcurveto{\pgfqpoint{3.809830in}{3.253036in}}{\pgfqpoint{3.804244in}{3.255350in}}{\pgfqpoint{3.798420in}{3.255350in}}%
\pgfpathcurveto{\pgfqpoint{3.792596in}{3.255350in}}{\pgfqpoint{3.787010in}{3.253036in}}{\pgfqpoint{3.782892in}{3.248918in}}%
\pgfpathcurveto{\pgfqpoint{3.778774in}{3.244800in}}{\pgfqpoint{3.776460in}{3.239214in}}{\pgfqpoint{3.776460in}{3.233390in}}%
\pgfpathcurveto{\pgfqpoint{3.776460in}{3.227566in}}{\pgfqpoint{3.778774in}{3.221980in}}{\pgfqpoint{3.782892in}{3.217861in}}%
\pgfpathcurveto{\pgfqpoint{3.787010in}{3.213743in}}{\pgfqpoint{3.792596in}{3.211429in}}{\pgfqpoint{3.798420in}{3.211429in}}%
\pgfpathlineto{\pgfqpoint{3.798420in}{3.211429in}}%
\pgfpathclose%
\pgfusepath{stroke,fill}%
\end{pgfscope}%
\begin{pgfscope}%
\pgfpathrectangle{\pgfqpoint{1.414736in}{0.548769in}}{\pgfqpoint{3.497747in}{7.253186in}}%
\pgfusepath{clip}%
\pgfsetbuttcap%
\pgfsetroundjoin%
\definecolor{currentfill}{rgb}{0.993248,0.906157,0.143936}%
\pgfsetfillcolor{currentfill}%
\pgfsetlinewidth{1.003750pt}%
\definecolor{currentstroke}{rgb}{0.993248,0.906157,0.143936}%
\pgfsetstrokecolor{currentstroke}%
\pgfsetdash{}{0pt}%
\pgfpathmoveto{\pgfqpoint{3.433738in}{3.682415in}}%
\pgfpathcurveto{\pgfqpoint{3.439562in}{3.682415in}}{\pgfqpoint{3.445148in}{3.684729in}}{\pgfqpoint{3.449266in}{3.688847in}}%
\pgfpathcurveto{\pgfqpoint{3.453385in}{3.692966in}}{\pgfqpoint{3.455698in}{3.698552in}}{\pgfqpoint{3.455698in}{3.704376in}}%
\pgfpathcurveto{\pgfqpoint{3.455698in}{3.710200in}}{\pgfqpoint{3.453385in}{3.715786in}}{\pgfqpoint{3.449266in}{3.719904in}}%
\pgfpathcurveto{\pgfqpoint{3.445148in}{3.724022in}}{\pgfqpoint{3.439562in}{3.726336in}}{\pgfqpoint{3.433738in}{3.726336in}}%
\pgfpathcurveto{\pgfqpoint{3.427914in}{3.726336in}}{\pgfqpoint{3.422328in}{3.724022in}}{\pgfqpoint{3.418210in}{3.719904in}}%
\pgfpathcurveto{\pgfqpoint{3.414092in}{3.715786in}}{\pgfqpoint{3.411778in}{3.710200in}}{\pgfqpoint{3.411778in}{3.704376in}}%
\pgfpathcurveto{\pgfqpoint{3.411778in}{3.698552in}}{\pgfqpoint{3.414092in}{3.692966in}}{\pgfqpoint{3.418210in}{3.688847in}}%
\pgfpathcurveto{\pgfqpoint{3.422328in}{3.684729in}}{\pgfqpoint{3.427914in}{3.682415in}}{\pgfqpoint{3.433738in}{3.682415in}}%
\pgfpathlineto{\pgfqpoint{3.433738in}{3.682415in}}%
\pgfpathclose%
\pgfusepath{stroke,fill}%
\end{pgfscope}%
\begin{pgfscope}%
\pgfpathrectangle{\pgfqpoint{1.414736in}{0.548769in}}{\pgfqpoint{3.497747in}{7.253186in}}%
\pgfusepath{clip}%
\pgfsetbuttcap%
\pgfsetroundjoin%
\definecolor{currentfill}{rgb}{0.657642,0.860219,0.203082}%
\pgfsetfillcolor{currentfill}%
\pgfsetlinewidth{1.003750pt}%
\definecolor{currentstroke}{rgb}{0.657642,0.860219,0.203082}%
\pgfsetstrokecolor{currentstroke}%
\pgfsetdash{}{0pt}%
\pgfpathmoveto{\pgfqpoint{3.640091in}{4.153402in}}%
\pgfpathcurveto{\pgfqpoint{3.645915in}{4.153402in}}{\pgfqpoint{3.651501in}{4.155715in}}{\pgfqpoint{3.655619in}{4.159834in}}%
\pgfpathcurveto{\pgfqpoint{3.659737in}{4.163952in}}{\pgfqpoint{3.662051in}{4.169538in}}{\pgfqpoint{3.662051in}{4.175362in}}%
\pgfpathcurveto{\pgfqpoint{3.662051in}{4.181186in}}{\pgfqpoint{3.659737in}{4.186772in}}{\pgfqpoint{3.655619in}{4.190890in}}%
\pgfpathcurveto{\pgfqpoint{3.651501in}{4.195008in}}{\pgfqpoint{3.645915in}{4.197322in}}{\pgfqpoint{3.640091in}{4.197322in}}%
\pgfpathcurveto{\pgfqpoint{3.634267in}{4.197322in}}{\pgfqpoint{3.628681in}{4.195008in}}{\pgfqpoint{3.624562in}{4.190890in}}%
\pgfpathcurveto{\pgfqpoint{3.620444in}{4.186772in}}{\pgfqpoint{3.618130in}{4.181186in}}{\pgfqpoint{3.618130in}{4.175362in}}%
\pgfpathcurveto{\pgfqpoint{3.618130in}{4.169538in}}{\pgfqpoint{3.620444in}{4.163952in}}{\pgfqpoint{3.624562in}{4.159834in}}%
\pgfpathcurveto{\pgfqpoint{3.628681in}{4.155715in}}{\pgfqpoint{3.634267in}{4.153402in}}{\pgfqpoint{3.640091in}{4.153402in}}%
\pgfpathlineto{\pgfqpoint{3.640091in}{4.153402in}}%
\pgfpathclose%
\pgfusepath{stroke,fill}%
\end{pgfscope}%
\begin{pgfscope}%
\pgfpathrectangle{\pgfqpoint{1.414736in}{0.548769in}}{\pgfqpoint{3.497747in}{7.253186in}}%
\pgfusepath{clip}%
\pgfsetbuttcap%
\pgfsetroundjoin%
\definecolor{currentfill}{rgb}{0.699415,0.867117,0.175971}%
\pgfsetfillcolor{currentfill}%
\pgfsetlinewidth{1.003750pt}%
\definecolor{currentstroke}{rgb}{0.699415,0.867117,0.175971}%
\pgfsetstrokecolor{currentstroke}%
\pgfsetdash{}{0pt}%
\pgfpathmoveto{\pgfqpoint{3.886444in}{4.624388in}}%
\pgfpathcurveto{\pgfqpoint{3.892268in}{4.624388in}}{\pgfqpoint{3.897854in}{4.626702in}}{\pgfqpoint{3.901972in}{4.630820in}}%
\pgfpathcurveto{\pgfqpoint{3.906090in}{4.634938in}}{\pgfqpoint{3.908404in}{4.640524in}}{\pgfqpoint{3.908404in}{4.646348in}}%
\pgfpathcurveto{\pgfqpoint{3.908404in}{4.652172in}}{\pgfqpoint{3.906090in}{4.657758in}}{\pgfqpoint{3.901972in}{4.661876in}}%
\pgfpathcurveto{\pgfqpoint{3.897854in}{4.665994in}}{\pgfqpoint{3.892268in}{4.668308in}}{\pgfqpoint{3.886444in}{4.668308in}}%
\pgfpathcurveto{\pgfqpoint{3.880620in}{4.668308in}}{\pgfqpoint{3.875034in}{4.665994in}}{\pgfqpoint{3.870916in}{4.661876in}}%
\pgfpathcurveto{\pgfqpoint{3.866798in}{4.657758in}}{\pgfqpoint{3.864484in}{4.652172in}}{\pgfqpoint{3.864484in}{4.646348in}}%
\pgfpathcurveto{\pgfqpoint{3.864484in}{4.640524in}}{\pgfqpoint{3.866798in}{4.634938in}}{\pgfqpoint{3.870916in}{4.630820in}}%
\pgfpathcurveto{\pgfqpoint{3.875034in}{4.626702in}}{\pgfqpoint{3.880620in}{4.624388in}}{\pgfqpoint{3.886444in}{4.624388in}}%
\pgfpathlineto{\pgfqpoint{3.886444in}{4.624388in}}%
\pgfpathclose%
\pgfusepath{stroke,fill}%
\end{pgfscope}%
\begin{pgfscope}%
\pgfpathrectangle{\pgfqpoint{1.414736in}{0.548769in}}{\pgfqpoint{3.497747in}{7.253186in}}%
\pgfusepath{clip}%
\pgfsetbuttcap%
\pgfsetroundjoin%
\definecolor{currentfill}{rgb}{0.120565,0.596422,0.543611}%
\pgfsetfillcolor{currentfill}%
\pgfsetlinewidth{1.003750pt}%
\definecolor{currentstroke}{rgb}{0.120565,0.596422,0.543611}%
\pgfsetstrokecolor{currentstroke}%
\pgfsetdash{}{0pt}%
\pgfpathmoveto{\pgfqpoint{3.419707in}{5.095374in}}%
\pgfpathcurveto{\pgfqpoint{3.425531in}{5.095374in}}{\pgfqpoint{3.431117in}{5.097688in}}{\pgfqpoint{3.435236in}{5.101806in}}%
\pgfpathcurveto{\pgfqpoint{3.439354in}{5.105924in}}{\pgfqpoint{3.441668in}{5.111510in}}{\pgfqpoint{3.441668in}{5.117334in}}%
\pgfpathcurveto{\pgfqpoint{3.441668in}{5.123158in}}{\pgfqpoint{3.439354in}{5.128744in}}{\pgfqpoint{3.435236in}{5.132862in}}%
\pgfpathcurveto{\pgfqpoint{3.431117in}{5.136980in}}{\pgfqpoint{3.425531in}{5.139294in}}{\pgfqpoint{3.419707in}{5.139294in}}%
\pgfpathcurveto{\pgfqpoint{3.413883in}{5.139294in}}{\pgfqpoint{3.408297in}{5.136980in}}{\pgfqpoint{3.404179in}{5.132862in}}%
\pgfpathcurveto{\pgfqpoint{3.400061in}{5.128744in}}{\pgfqpoint{3.397747in}{5.123158in}}{\pgfqpoint{3.397747in}{5.117334in}}%
\pgfpathcurveto{\pgfqpoint{3.397747in}{5.111510in}}{\pgfqpoint{3.400061in}{5.105924in}}{\pgfqpoint{3.404179in}{5.101806in}}%
\pgfpathcurveto{\pgfqpoint{3.408297in}{5.097688in}}{\pgfqpoint{3.413883in}{5.095374in}}{\pgfqpoint{3.419707in}{5.095374in}}%
\pgfpathlineto{\pgfqpoint{3.419707in}{5.095374in}}%
\pgfpathclose%
\pgfusepath{stroke,fill}%
\end{pgfscope}%
\begin{pgfscope}%
\pgfpathrectangle{\pgfqpoint{1.414736in}{0.548769in}}{\pgfqpoint{3.497747in}{7.253186in}}%
\pgfusepath{clip}%
\pgfsetbuttcap%
\pgfsetroundjoin%
\definecolor{currentfill}{rgb}{0.159194,0.482237,0.558073}%
\pgfsetfillcolor{currentfill}%
\pgfsetlinewidth{1.003750pt}%
\definecolor{currentstroke}{rgb}{0.159194,0.482237,0.558073}%
\pgfsetstrokecolor{currentstroke}%
\pgfsetdash{}{0pt}%
\pgfpathmoveto{\pgfqpoint{3.550398in}{5.566360in}}%
\pgfpathcurveto{\pgfqpoint{3.556222in}{5.566360in}}{\pgfqpoint{3.561808in}{5.568674in}}{\pgfqpoint{3.565926in}{5.572792in}}%
\pgfpathcurveto{\pgfqpoint{3.570044in}{5.576910in}}{\pgfqpoint{3.572358in}{5.582496in}}{\pgfqpoint{3.572358in}{5.588320in}}%
\pgfpathcurveto{\pgfqpoint{3.572358in}{5.594144in}}{\pgfqpoint{3.570044in}{5.599730in}}{\pgfqpoint{3.565926in}{5.603848in}}%
\pgfpathcurveto{\pgfqpoint{3.561808in}{5.607967in}}{\pgfqpoint{3.556222in}{5.610280in}}{\pgfqpoint{3.550398in}{5.610280in}}%
\pgfpathcurveto{\pgfqpoint{3.544574in}{5.610280in}}{\pgfqpoint{3.538988in}{5.607967in}}{\pgfqpoint{3.534870in}{5.603848in}}%
\pgfpathcurveto{\pgfqpoint{3.530751in}{5.599730in}}{\pgfqpoint{3.528438in}{5.594144in}}{\pgfqpoint{3.528438in}{5.588320in}}%
\pgfpathcurveto{\pgfqpoint{3.528438in}{5.582496in}}{\pgfqpoint{3.530751in}{5.576910in}}{\pgfqpoint{3.534870in}{5.572792in}}%
\pgfpathcurveto{\pgfqpoint{3.538988in}{5.568674in}}{\pgfqpoint{3.544574in}{5.566360in}}{\pgfqpoint{3.550398in}{5.566360in}}%
\pgfpathlineto{\pgfqpoint{3.550398in}{5.566360in}}%
\pgfpathclose%
\pgfusepath{stroke,fill}%
\end{pgfscope}%
\begin{pgfscope}%
\pgfpathrectangle{\pgfqpoint{1.414736in}{0.548769in}}{\pgfqpoint{3.497747in}{7.253186in}}%
\pgfusepath{clip}%
\pgfsetbuttcap%
\pgfsetroundjoin%
\definecolor{currentfill}{rgb}{0.506271,0.828786,0.300362}%
\pgfsetfillcolor{currentfill}%
\pgfsetlinewidth{1.003750pt}%
\definecolor{currentstroke}{rgb}{0.506271,0.828786,0.300362}%
\pgfsetstrokecolor{currentstroke}%
\pgfsetdash{}{0pt}%
\pgfpathmoveto{\pgfqpoint{3.624761in}{6.037346in}}%
\pgfpathcurveto{\pgfqpoint{3.630585in}{6.037346in}}{\pgfqpoint{3.636172in}{6.039660in}}{\pgfqpoint{3.640290in}{6.043778in}}%
\pgfpathcurveto{\pgfqpoint{3.644408in}{6.047896in}}{\pgfqpoint{3.646722in}{6.053482in}}{\pgfqpoint{3.646722in}{6.059306in}}%
\pgfpathcurveto{\pgfqpoint{3.646722in}{6.065130in}}{\pgfqpoint{3.644408in}{6.070716in}}{\pgfqpoint{3.640290in}{6.074835in}}%
\pgfpathcurveto{\pgfqpoint{3.636172in}{6.078953in}}{\pgfqpoint{3.630585in}{6.081267in}}{\pgfqpoint{3.624761in}{6.081267in}}%
\pgfpathcurveto{\pgfqpoint{3.618937in}{6.081267in}}{\pgfqpoint{3.613351in}{6.078953in}}{\pgfqpoint{3.609233in}{6.074835in}}%
\pgfpathcurveto{\pgfqpoint{3.605115in}{6.070716in}}{\pgfqpoint{3.602801in}{6.065130in}}{\pgfqpoint{3.602801in}{6.059306in}}%
\pgfpathcurveto{\pgfqpoint{3.602801in}{6.053482in}}{\pgfqpoint{3.605115in}{6.047896in}}{\pgfqpoint{3.609233in}{6.043778in}}%
\pgfpathcurveto{\pgfqpoint{3.613351in}{6.039660in}}{\pgfqpoint{3.618937in}{6.037346in}}{\pgfqpoint{3.624761in}{6.037346in}}%
\pgfpathlineto{\pgfqpoint{3.624761in}{6.037346in}}%
\pgfpathclose%
\pgfusepath{stroke,fill}%
\end{pgfscope}%
\begin{pgfscope}%
\pgfpathrectangle{\pgfqpoint{1.414736in}{0.548769in}}{\pgfqpoint{3.497747in}{7.253186in}}%
\pgfusepath{clip}%
\pgfsetbuttcap%
\pgfsetroundjoin%
\definecolor{currentfill}{rgb}{0.993248,0.906157,0.143936}%
\pgfsetfillcolor{currentfill}%
\pgfsetlinewidth{1.003750pt}%
\definecolor{currentstroke}{rgb}{0.993248,0.906157,0.143936}%
\pgfsetstrokecolor{currentstroke}%
\pgfsetdash{}{0pt}%
\pgfpathmoveto{\pgfqpoint{3.574810in}{6.508332in}}%
\pgfpathcurveto{\pgfqpoint{3.580634in}{6.508332in}}{\pgfqpoint{3.586220in}{6.510646in}}{\pgfqpoint{3.590338in}{6.514764in}}%
\pgfpathcurveto{\pgfqpoint{3.594456in}{6.518882in}}{\pgfqpoint{3.596770in}{6.524468in}}{\pgfqpoint{3.596770in}{6.530292in}}%
\pgfpathcurveto{\pgfqpoint{3.596770in}{6.536116in}}{\pgfqpoint{3.594456in}{6.541702in}}{\pgfqpoint{3.590338in}{6.545821in}}%
\pgfpathcurveto{\pgfqpoint{3.586220in}{6.549939in}}{\pgfqpoint{3.580634in}{6.552253in}}{\pgfqpoint{3.574810in}{6.552253in}}%
\pgfpathcurveto{\pgfqpoint{3.568986in}{6.552253in}}{\pgfqpoint{3.563400in}{6.549939in}}{\pgfqpoint{3.559282in}{6.545821in}}%
\pgfpathcurveto{\pgfqpoint{3.555164in}{6.541702in}}{\pgfqpoint{3.552850in}{6.536116in}}{\pgfqpoint{3.552850in}{6.530292in}}%
\pgfpathcurveto{\pgfqpoint{3.552850in}{6.524468in}}{\pgfqpoint{3.555164in}{6.518882in}}{\pgfqpoint{3.559282in}{6.514764in}}%
\pgfpathcurveto{\pgfqpoint{3.563400in}{6.510646in}}{\pgfqpoint{3.568986in}{6.508332in}}{\pgfqpoint{3.574810in}{6.508332in}}%
\pgfpathlineto{\pgfqpoint{3.574810in}{6.508332in}}%
\pgfpathclose%
\pgfusepath{stroke,fill}%
\end{pgfscope}%
\begin{pgfscope}%
\pgfpathrectangle{\pgfqpoint{1.414736in}{0.548769in}}{\pgfqpoint{3.497747in}{7.253186in}}%
\pgfusepath{clip}%
\pgfsetbuttcap%
\pgfsetroundjoin%
\definecolor{currentfill}{rgb}{0.678489,0.863742,0.189503}%
\pgfsetfillcolor{currentfill}%
\pgfsetlinewidth{1.003750pt}%
\definecolor{currentstroke}{rgb}{0.678489,0.863742,0.189503}%
\pgfsetstrokecolor{currentstroke}%
\pgfsetdash{}{0pt}%
\pgfpathmoveto{\pgfqpoint{4.169822in}{6.979318in}}%
\pgfpathcurveto{\pgfqpoint{4.175646in}{6.979318in}}{\pgfqpoint{4.181232in}{6.981632in}}{\pgfqpoint{4.185350in}{6.985750in}}%
\pgfpathcurveto{\pgfqpoint{4.189468in}{6.989868in}}{\pgfqpoint{4.191782in}{6.995455in}}{\pgfqpoint{4.191782in}{7.001278in}}%
\pgfpathcurveto{\pgfqpoint{4.191782in}{7.007102in}}{\pgfqpoint{4.189468in}{7.012689in}}{\pgfqpoint{4.185350in}{7.016807in}}%
\pgfpathcurveto{\pgfqpoint{4.181232in}{7.020925in}}{\pgfqpoint{4.175646in}{7.023239in}}{\pgfqpoint{4.169822in}{7.023239in}}%
\pgfpathcurveto{\pgfqpoint{4.163998in}{7.023239in}}{\pgfqpoint{4.158412in}{7.020925in}}{\pgfqpoint{4.154294in}{7.016807in}}%
\pgfpathcurveto{\pgfqpoint{4.150175in}{7.012689in}}{\pgfqpoint{4.147862in}{7.007102in}}{\pgfqpoint{4.147862in}{7.001278in}}%
\pgfpathcurveto{\pgfqpoint{4.147862in}{6.995455in}}{\pgfqpoint{4.150175in}{6.989868in}}{\pgfqpoint{4.154294in}{6.985750in}}%
\pgfpathcurveto{\pgfqpoint{4.158412in}{6.981632in}}{\pgfqpoint{4.163998in}{6.979318in}}{\pgfqpoint{4.169822in}{6.979318in}}%
\pgfpathlineto{\pgfqpoint{4.169822in}{6.979318in}}%
\pgfpathclose%
\pgfusepath{stroke,fill}%
\end{pgfscope}%
\begin{pgfscope}%
\pgfpathrectangle{\pgfqpoint{1.414736in}{0.548769in}}{\pgfqpoint{3.497747in}{7.253186in}}%
\pgfusepath{clip}%
\pgfsetbuttcap%
\pgfsetroundjoin%
\definecolor{currentfill}{rgb}{0.246070,0.738910,0.452024}%
\pgfsetfillcolor{currentfill}%
\pgfsetlinewidth{1.003750pt}%
\definecolor{currentstroke}{rgb}{0.246070,0.738910,0.452024}%
\pgfsetstrokecolor{currentstroke}%
\pgfsetdash{}{0pt}%
\pgfpathmoveto{\pgfqpoint{3.648592in}{7.450304in}}%
\pgfpathcurveto{\pgfqpoint{3.654416in}{7.450304in}}{\pgfqpoint{3.660002in}{7.452618in}}{\pgfqpoint{3.664120in}{7.456736in}}%
\pgfpathcurveto{\pgfqpoint{3.668239in}{7.460854in}}{\pgfqpoint{3.670552in}{7.466441in}}{\pgfqpoint{3.670552in}{7.472265in}}%
\pgfpathcurveto{\pgfqpoint{3.670552in}{7.478089in}}{\pgfqpoint{3.668239in}{7.483675in}}{\pgfqpoint{3.664120in}{7.487793in}}%
\pgfpathcurveto{\pgfqpoint{3.660002in}{7.491911in}}{\pgfqpoint{3.654416in}{7.494225in}}{\pgfqpoint{3.648592in}{7.494225in}}%
\pgfpathcurveto{\pgfqpoint{3.642768in}{7.494225in}}{\pgfqpoint{3.637182in}{7.491911in}}{\pgfqpoint{3.633064in}{7.487793in}}%
\pgfpathcurveto{\pgfqpoint{3.628946in}{7.483675in}}{\pgfqpoint{3.626632in}{7.478089in}}{\pgfqpoint{3.626632in}{7.472265in}}%
\pgfpathcurveto{\pgfqpoint{3.626632in}{7.466441in}}{\pgfqpoint{3.628946in}{7.460854in}}{\pgfqpoint{3.633064in}{7.456736in}}%
\pgfpathcurveto{\pgfqpoint{3.637182in}{7.452618in}}{\pgfqpoint{3.642768in}{7.450304in}}{\pgfqpoint{3.648592in}{7.450304in}}%
\pgfpathlineto{\pgfqpoint{3.648592in}{7.450304in}}%
\pgfpathclose%
\pgfusepath{stroke,fill}%
\end{pgfscope}%
\begin{pgfscope}%
\pgfpathrectangle{\pgfqpoint{1.414736in}{0.548769in}}{\pgfqpoint{3.497747in}{7.253186in}}%
\pgfusepath{clip}%
\pgfsetbuttcap%
\pgfsetroundjoin%
\definecolor{currentfill}{rgb}{0.916242,0.896091,0.100717}%
\pgfsetfillcolor{currentfill}%
\pgfsetlinewidth{1.003750pt}%
\definecolor{currentstroke}{rgb}{0.916242,0.896091,0.100717}%
\pgfsetstrokecolor{currentstroke}%
\pgfsetdash{}{0pt}%
\pgfpathmoveto{\pgfqpoint{3.547305in}{0.856499in}}%
\pgfpathcurveto{\pgfqpoint{3.553129in}{0.856499in}}{\pgfqpoint{3.558715in}{0.858813in}}{\pgfqpoint{3.562833in}{0.862931in}}%
\pgfpathcurveto{\pgfqpoint{3.566951in}{0.867049in}}{\pgfqpoint{3.569265in}{0.872635in}}{\pgfqpoint{3.569265in}{0.878459in}}%
\pgfpathcurveto{\pgfqpoint{3.569265in}{0.884283in}}{\pgfqpoint{3.566951in}{0.889869in}}{\pgfqpoint{3.562833in}{0.893987in}}%
\pgfpathcurveto{\pgfqpoint{3.558715in}{0.898105in}}{\pgfqpoint{3.553129in}{0.900419in}}{\pgfqpoint{3.547305in}{0.900419in}}%
\pgfpathcurveto{\pgfqpoint{3.541481in}{0.900419in}}{\pgfqpoint{3.535895in}{0.898105in}}{\pgfqpoint{3.531777in}{0.893987in}}%
\pgfpathcurveto{\pgfqpoint{3.527659in}{0.889869in}}{\pgfqpoint{3.525345in}{0.884283in}}{\pgfqpoint{3.525345in}{0.878459in}}%
\pgfpathcurveto{\pgfqpoint{3.525345in}{0.872635in}}{\pgfqpoint{3.527659in}{0.867049in}}{\pgfqpoint{3.531777in}{0.862931in}}%
\pgfpathcurveto{\pgfqpoint{3.535895in}{0.858813in}}{\pgfqpoint{3.541481in}{0.856499in}}{\pgfqpoint{3.547305in}{0.856499in}}%
\pgfpathlineto{\pgfqpoint{3.547305in}{0.856499in}}%
\pgfpathclose%
\pgfusepath{stroke,fill}%
\end{pgfscope}%
\begin{pgfscope}%
\pgfpathrectangle{\pgfqpoint{1.414736in}{0.548769in}}{\pgfqpoint{3.497747in}{7.253186in}}%
\pgfusepath{clip}%
\pgfsetbuttcap%
\pgfsetroundjoin%
\definecolor{currentfill}{rgb}{0.506271,0.828786,0.300362}%
\pgfsetfillcolor{currentfill}%
\pgfsetlinewidth{1.003750pt}%
\definecolor{currentstroke}{rgb}{0.506271,0.828786,0.300362}%
\pgfsetstrokecolor{currentstroke}%
\pgfsetdash{}{0pt}%
\pgfpathmoveto{\pgfqpoint{3.366951in}{1.327485in}}%
\pgfpathcurveto{\pgfqpoint{3.372775in}{1.327485in}}{\pgfqpoint{3.378362in}{1.329799in}}{\pgfqpoint{3.382480in}{1.333917in}}%
\pgfpathcurveto{\pgfqpoint{3.386598in}{1.338035in}}{\pgfqpoint{3.388912in}{1.343621in}}{\pgfqpoint{3.388912in}{1.349445in}}%
\pgfpathcurveto{\pgfqpoint{3.388912in}{1.355269in}}{\pgfqpoint{3.386598in}{1.360855in}}{\pgfqpoint{3.382480in}{1.364973in}}%
\pgfpathcurveto{\pgfqpoint{3.378362in}{1.369092in}}{\pgfqpoint{3.372775in}{1.371405in}}{\pgfqpoint{3.366951in}{1.371405in}}%
\pgfpathcurveto{\pgfqpoint{3.361127in}{1.371405in}}{\pgfqpoint{3.355541in}{1.369092in}}{\pgfqpoint{3.351423in}{1.364973in}}%
\pgfpathcurveto{\pgfqpoint{3.347305in}{1.360855in}}{\pgfqpoint{3.344991in}{1.355269in}}{\pgfqpoint{3.344991in}{1.349445in}}%
\pgfpathcurveto{\pgfqpoint{3.344991in}{1.343621in}}{\pgfqpoint{3.347305in}{1.338035in}}{\pgfqpoint{3.351423in}{1.333917in}}%
\pgfpathcurveto{\pgfqpoint{3.355541in}{1.329799in}}{\pgfqpoint{3.361127in}{1.327485in}}{\pgfqpoint{3.366951in}{1.327485in}}%
\pgfpathlineto{\pgfqpoint{3.366951in}{1.327485in}}%
\pgfpathclose%
\pgfusepath{stroke,fill}%
\end{pgfscope}%
\begin{pgfscope}%
\pgfpathrectangle{\pgfqpoint{1.414736in}{0.548769in}}{\pgfqpoint{3.497747in}{7.253186in}}%
\pgfusepath{clip}%
\pgfsetbuttcap%
\pgfsetroundjoin%
\definecolor{currentfill}{rgb}{0.935904,0.898570,0.108131}%
\pgfsetfillcolor{currentfill}%
\pgfsetlinewidth{1.003750pt}%
\definecolor{currentstroke}{rgb}{0.935904,0.898570,0.108131}%
\pgfsetstrokecolor{currentstroke}%
\pgfsetdash{}{0pt}%
\pgfpathmoveto{\pgfqpoint{4.012867in}{1.798471in}}%
\pgfpathcurveto{\pgfqpoint{4.018691in}{1.798471in}}{\pgfqpoint{4.024277in}{1.800785in}}{\pgfqpoint{4.028395in}{1.804903in}}%
\pgfpathcurveto{\pgfqpoint{4.032513in}{1.809021in}}{\pgfqpoint{4.034827in}{1.814607in}}{\pgfqpoint{4.034827in}{1.820431in}}%
\pgfpathcurveto{\pgfqpoint{4.034827in}{1.826255in}}{\pgfqpoint{4.032513in}{1.831841in}}{\pgfqpoint{4.028395in}{1.835960in}}%
\pgfpathcurveto{\pgfqpoint{4.024277in}{1.840078in}}{\pgfqpoint{4.018691in}{1.842392in}}{\pgfqpoint{4.012867in}{1.842392in}}%
\pgfpathcurveto{\pgfqpoint{4.007043in}{1.842392in}}{\pgfqpoint{4.001457in}{1.840078in}}{\pgfqpoint{3.997338in}{1.835960in}}%
\pgfpathcurveto{\pgfqpoint{3.993220in}{1.831841in}}{\pgfqpoint{3.990906in}{1.826255in}}{\pgfqpoint{3.990906in}{1.820431in}}%
\pgfpathcurveto{\pgfqpoint{3.990906in}{1.814607in}}{\pgfqpoint{3.993220in}{1.809021in}}{\pgfqpoint{3.997338in}{1.804903in}}%
\pgfpathcurveto{\pgfqpoint{4.001457in}{1.800785in}}{\pgfqpoint{4.007043in}{1.798471in}}{\pgfqpoint{4.012867in}{1.798471in}}%
\pgfpathlineto{\pgfqpoint{4.012867in}{1.798471in}}%
\pgfpathclose%
\pgfusepath{stroke,fill}%
\end{pgfscope}%
\begin{pgfscope}%
\pgfpathrectangle{\pgfqpoint{1.414736in}{0.548769in}}{\pgfqpoint{3.497747in}{7.253186in}}%
\pgfusepath{clip}%
\pgfsetbuttcap%
\pgfsetroundjoin%
\definecolor{currentfill}{rgb}{0.993248,0.906157,0.143936}%
\pgfsetfillcolor{currentfill}%
\pgfsetlinewidth{1.003750pt}%
\definecolor{currentstroke}{rgb}{0.993248,0.906157,0.143936}%
\pgfsetstrokecolor{currentstroke}%
\pgfsetdash{}{0pt}%
\pgfpathmoveto{\pgfqpoint{3.687099in}{2.269457in}}%
\pgfpathcurveto{\pgfqpoint{3.692923in}{2.269457in}}{\pgfqpoint{3.698509in}{2.271771in}}{\pgfqpoint{3.702627in}{2.275889in}}%
\pgfpathcurveto{\pgfqpoint{3.706745in}{2.280007in}}{\pgfqpoint{3.709059in}{2.285593in}}{\pgfqpoint{3.709059in}{2.291417in}}%
\pgfpathcurveto{\pgfqpoint{3.709059in}{2.297241in}}{\pgfqpoint{3.706745in}{2.302828in}}{\pgfqpoint{3.702627in}{2.306946in}}%
\pgfpathcurveto{\pgfqpoint{3.698509in}{2.311064in}}{\pgfqpoint{3.692923in}{2.313378in}}{\pgfqpoint{3.687099in}{2.313378in}}%
\pgfpathcurveto{\pgfqpoint{3.681275in}{2.313378in}}{\pgfqpoint{3.675689in}{2.311064in}}{\pgfqpoint{3.671570in}{2.306946in}}%
\pgfpathcurveto{\pgfqpoint{3.667452in}{2.302828in}}{\pgfqpoint{3.665138in}{2.297241in}}{\pgfqpoint{3.665138in}{2.291417in}}%
\pgfpathcurveto{\pgfqpoint{3.665138in}{2.285593in}}{\pgfqpoint{3.667452in}{2.280007in}}{\pgfqpoint{3.671570in}{2.275889in}}%
\pgfpathcurveto{\pgfqpoint{3.675689in}{2.271771in}}{\pgfqpoint{3.681275in}{2.269457in}}{\pgfqpoint{3.687099in}{2.269457in}}%
\pgfpathlineto{\pgfqpoint{3.687099in}{2.269457in}}%
\pgfpathclose%
\pgfusepath{stroke,fill}%
\end{pgfscope}%
\begin{pgfscope}%
\pgfpathrectangle{\pgfqpoint{1.414736in}{0.548769in}}{\pgfqpoint{3.497747in}{7.253186in}}%
\pgfusepath{clip}%
\pgfsetbuttcap%
\pgfsetroundjoin%
\definecolor{currentfill}{rgb}{0.762373,0.876424,0.137064}%
\pgfsetfillcolor{currentfill}%
\pgfsetlinewidth{1.003750pt}%
\definecolor{currentstroke}{rgb}{0.762373,0.876424,0.137064}%
\pgfsetstrokecolor{currentstroke}%
\pgfsetdash{}{0pt}%
\pgfpathmoveto{\pgfqpoint{2.615621in}{2.740443in}}%
\pgfpathcurveto{\pgfqpoint{2.621445in}{2.740443in}}{\pgfqpoint{2.627031in}{2.742757in}}{\pgfqpoint{2.631149in}{2.746875in}}%
\pgfpathcurveto{\pgfqpoint{2.635267in}{2.750993in}}{\pgfqpoint{2.637581in}{2.756580in}}{\pgfqpoint{2.637581in}{2.762404in}}%
\pgfpathcurveto{\pgfqpoint{2.637581in}{2.768227in}}{\pgfqpoint{2.635267in}{2.773814in}}{\pgfqpoint{2.631149in}{2.777932in}}%
\pgfpathcurveto{\pgfqpoint{2.627031in}{2.782050in}}{\pgfqpoint{2.621445in}{2.784364in}}{\pgfqpoint{2.615621in}{2.784364in}}%
\pgfpathcurveto{\pgfqpoint{2.609797in}{2.784364in}}{\pgfqpoint{2.604211in}{2.782050in}}{\pgfqpoint{2.600092in}{2.777932in}}%
\pgfpathcurveto{\pgfqpoint{2.595974in}{2.773814in}}{\pgfqpoint{2.593660in}{2.768227in}}{\pgfqpoint{2.593660in}{2.762404in}}%
\pgfpathcurveto{\pgfqpoint{2.593660in}{2.756580in}}{\pgfqpoint{2.595974in}{2.750993in}}{\pgfqpoint{2.600092in}{2.746875in}}%
\pgfpathcurveto{\pgfqpoint{2.604211in}{2.742757in}}{\pgfqpoint{2.609797in}{2.740443in}}{\pgfqpoint{2.615621in}{2.740443in}}%
\pgfpathlineto{\pgfqpoint{2.615621in}{2.740443in}}%
\pgfpathclose%
\pgfusepath{stroke,fill}%
\end{pgfscope}%
\begin{pgfscope}%
\pgfpathrectangle{\pgfqpoint{1.414736in}{0.548769in}}{\pgfqpoint{3.497747in}{7.253186in}}%
\pgfusepath{clip}%
\pgfsetbuttcap%
\pgfsetroundjoin%
\definecolor{currentfill}{rgb}{0.876168,0.891125,0.095250}%
\pgfsetfillcolor{currentfill}%
\pgfsetlinewidth{1.003750pt}%
\definecolor{currentstroke}{rgb}{0.876168,0.891125,0.095250}%
\pgfsetstrokecolor{currentstroke}%
\pgfsetdash{}{0pt}%
\pgfpathmoveto{\pgfqpoint{4.279805in}{3.211429in}}%
\pgfpathcurveto{\pgfqpoint{4.285629in}{3.211429in}}{\pgfqpoint{4.291215in}{3.213743in}}{\pgfqpoint{4.295333in}{3.217861in}}%
\pgfpathcurveto{\pgfqpoint{4.299451in}{3.221980in}}{\pgfqpoint{4.301765in}{3.227566in}}{\pgfqpoint{4.301765in}{3.233390in}}%
\pgfpathcurveto{\pgfqpoint{4.301765in}{3.239214in}}{\pgfqpoint{4.299451in}{3.244800in}}{\pgfqpoint{4.295333in}{3.248918in}}%
\pgfpathcurveto{\pgfqpoint{4.291215in}{3.253036in}}{\pgfqpoint{4.285629in}{3.255350in}}{\pgfqpoint{4.279805in}{3.255350in}}%
\pgfpathcurveto{\pgfqpoint{4.273981in}{3.255350in}}{\pgfqpoint{4.268395in}{3.253036in}}{\pgfqpoint{4.264277in}{3.248918in}}%
\pgfpathcurveto{\pgfqpoint{4.260159in}{3.244800in}}{\pgfqpoint{4.257845in}{3.239214in}}{\pgfqpoint{4.257845in}{3.233390in}}%
\pgfpathcurveto{\pgfqpoint{4.257845in}{3.227566in}}{\pgfqpoint{4.260159in}{3.221980in}}{\pgfqpoint{4.264277in}{3.217861in}}%
\pgfpathcurveto{\pgfqpoint{4.268395in}{3.213743in}}{\pgfqpoint{4.273981in}{3.211429in}}{\pgfqpoint{4.279805in}{3.211429in}}%
\pgfpathlineto{\pgfqpoint{4.279805in}{3.211429in}}%
\pgfpathclose%
\pgfusepath{stroke,fill}%
\end{pgfscope}%
\begin{pgfscope}%
\pgfpathrectangle{\pgfqpoint{1.414736in}{0.548769in}}{\pgfqpoint{3.497747in}{7.253186in}}%
\pgfusepath{clip}%
\pgfsetbuttcap%
\pgfsetroundjoin%
\definecolor{currentfill}{rgb}{0.993248,0.906157,0.143936}%
\pgfsetfillcolor{currentfill}%
\pgfsetlinewidth{1.003750pt}%
\definecolor{currentstroke}{rgb}{0.993248,0.906157,0.143936}%
\pgfsetstrokecolor{currentstroke}%
\pgfsetdash{}{0pt}%
\pgfpathmoveto{\pgfqpoint{3.403912in}{3.682415in}}%
\pgfpathcurveto{\pgfqpoint{3.409736in}{3.682415in}}{\pgfqpoint{3.415322in}{3.684729in}}{\pgfqpoint{3.419441in}{3.688847in}}%
\pgfpathcurveto{\pgfqpoint{3.423559in}{3.692966in}}{\pgfqpoint{3.425873in}{3.698552in}}{\pgfqpoint{3.425873in}{3.704376in}}%
\pgfpathcurveto{\pgfqpoint{3.425873in}{3.710200in}}{\pgfqpoint{3.423559in}{3.715786in}}{\pgfqpoint{3.419441in}{3.719904in}}%
\pgfpathcurveto{\pgfqpoint{3.415322in}{3.724022in}}{\pgfqpoint{3.409736in}{3.726336in}}{\pgfqpoint{3.403912in}{3.726336in}}%
\pgfpathcurveto{\pgfqpoint{3.398088in}{3.726336in}}{\pgfqpoint{3.392502in}{3.724022in}}{\pgfqpoint{3.388384in}{3.719904in}}%
\pgfpathcurveto{\pgfqpoint{3.384266in}{3.715786in}}{\pgfqpoint{3.381952in}{3.710200in}}{\pgfqpoint{3.381952in}{3.704376in}}%
\pgfpathcurveto{\pgfqpoint{3.381952in}{3.698552in}}{\pgfqpoint{3.384266in}{3.692966in}}{\pgfqpoint{3.388384in}{3.688847in}}%
\pgfpathcurveto{\pgfqpoint{3.392502in}{3.684729in}}{\pgfqpoint{3.398088in}{3.682415in}}{\pgfqpoint{3.403912in}{3.682415in}}%
\pgfpathlineto{\pgfqpoint{3.403912in}{3.682415in}}%
\pgfpathclose%
\pgfusepath{stroke,fill}%
\end{pgfscope}%
\begin{pgfscope}%
\pgfpathrectangle{\pgfqpoint{1.414736in}{0.548769in}}{\pgfqpoint{3.497747in}{7.253186in}}%
\pgfusepath{clip}%
\pgfsetbuttcap%
\pgfsetroundjoin%
\definecolor{currentfill}{rgb}{0.993248,0.906157,0.143936}%
\pgfsetfillcolor{currentfill}%
\pgfsetlinewidth{1.003750pt}%
\definecolor{currentstroke}{rgb}{0.993248,0.906157,0.143936}%
\pgfsetstrokecolor{currentstroke}%
\pgfsetdash{}{0pt}%
\pgfpathmoveto{\pgfqpoint{3.319219in}{4.153402in}}%
\pgfpathcurveto{\pgfqpoint{3.325042in}{4.153402in}}{\pgfqpoint{3.330629in}{4.155715in}}{\pgfqpoint{3.334747in}{4.159834in}}%
\pgfpathcurveto{\pgfqpoint{3.338865in}{4.163952in}}{\pgfqpoint{3.341179in}{4.169538in}}{\pgfqpoint{3.341179in}{4.175362in}}%
\pgfpathcurveto{\pgfqpoint{3.341179in}{4.181186in}}{\pgfqpoint{3.338865in}{4.186772in}}{\pgfqpoint{3.334747in}{4.190890in}}%
\pgfpathcurveto{\pgfqpoint{3.330629in}{4.195008in}}{\pgfqpoint{3.325042in}{4.197322in}}{\pgfqpoint{3.319219in}{4.197322in}}%
\pgfpathcurveto{\pgfqpoint{3.313395in}{4.197322in}}{\pgfqpoint{3.307808in}{4.195008in}}{\pgfqpoint{3.303690in}{4.190890in}}%
\pgfpathcurveto{\pgfqpoint{3.299572in}{4.186772in}}{\pgfqpoint{3.297258in}{4.181186in}}{\pgfqpoint{3.297258in}{4.175362in}}%
\pgfpathcurveto{\pgfqpoint{3.297258in}{4.169538in}}{\pgfqpoint{3.299572in}{4.163952in}}{\pgfqpoint{3.303690in}{4.159834in}}%
\pgfpathcurveto{\pgfqpoint{3.307808in}{4.155715in}}{\pgfqpoint{3.313395in}{4.153402in}}{\pgfqpoint{3.319219in}{4.153402in}}%
\pgfpathlineto{\pgfqpoint{3.319219in}{4.153402in}}%
\pgfpathclose%
\pgfusepath{stroke,fill}%
\end{pgfscope}%
\begin{pgfscope}%
\pgfpathrectangle{\pgfqpoint{1.414736in}{0.548769in}}{\pgfqpoint{3.497747in}{7.253186in}}%
\pgfusepath{clip}%
\pgfsetbuttcap%
\pgfsetroundjoin%
\definecolor{currentfill}{rgb}{0.565498,0.842430,0.262877}%
\pgfsetfillcolor{currentfill}%
\pgfsetlinewidth{1.003750pt}%
\definecolor{currentstroke}{rgb}{0.565498,0.842430,0.262877}%
\pgfsetstrokecolor{currentstroke}%
\pgfsetdash{}{0pt}%
\pgfpathmoveto{\pgfqpoint{4.020209in}{4.624388in}}%
\pgfpathcurveto{\pgfqpoint{4.026033in}{4.624388in}}{\pgfqpoint{4.031620in}{4.626702in}}{\pgfqpoint{4.035738in}{4.630820in}}%
\pgfpathcurveto{\pgfqpoint{4.039856in}{4.634938in}}{\pgfqpoint{4.042170in}{4.640524in}}{\pgfqpoint{4.042170in}{4.646348in}}%
\pgfpathcurveto{\pgfqpoint{4.042170in}{4.652172in}}{\pgfqpoint{4.039856in}{4.657758in}}{\pgfqpoint{4.035738in}{4.661876in}}%
\pgfpathcurveto{\pgfqpoint{4.031620in}{4.665994in}}{\pgfqpoint{4.026033in}{4.668308in}}{\pgfqpoint{4.020209in}{4.668308in}}%
\pgfpathcurveto{\pgfqpoint{4.014386in}{4.668308in}}{\pgfqpoint{4.008799in}{4.665994in}}{\pgfqpoint{4.004681in}{4.661876in}}%
\pgfpathcurveto{\pgfqpoint{4.000563in}{4.657758in}}{\pgfqpoint{3.998249in}{4.652172in}}{\pgfqpoint{3.998249in}{4.646348in}}%
\pgfpathcurveto{\pgfqpoint{3.998249in}{4.640524in}}{\pgfqpoint{4.000563in}{4.634938in}}{\pgfqpoint{4.004681in}{4.630820in}}%
\pgfpathcurveto{\pgfqpoint{4.008799in}{4.626702in}}{\pgfqpoint{4.014386in}{4.624388in}}{\pgfqpoint{4.020209in}{4.624388in}}%
\pgfpathlineto{\pgfqpoint{4.020209in}{4.624388in}}%
\pgfpathclose%
\pgfusepath{stroke,fill}%
\end{pgfscope}%
\begin{pgfscope}%
\pgfpathrectangle{\pgfqpoint{1.414736in}{0.548769in}}{\pgfqpoint{3.497747in}{7.253186in}}%
\pgfusepath{clip}%
\pgfsetbuttcap%
\pgfsetroundjoin%
\definecolor{currentfill}{rgb}{0.126326,0.644107,0.525311}%
\pgfsetfillcolor{currentfill}%
\pgfsetlinewidth{1.003750pt}%
\definecolor{currentstroke}{rgb}{0.126326,0.644107,0.525311}%
\pgfsetstrokecolor{currentstroke}%
\pgfsetdash{}{0pt}%
\pgfpathmoveto{\pgfqpoint{3.517176in}{5.095374in}}%
\pgfpathcurveto{\pgfqpoint{3.523000in}{5.095374in}}{\pgfqpoint{3.528586in}{5.097688in}}{\pgfqpoint{3.532704in}{5.101806in}}%
\pgfpathcurveto{\pgfqpoint{3.536822in}{5.105924in}}{\pgfqpoint{3.539136in}{5.111510in}}{\pgfqpoint{3.539136in}{5.117334in}}%
\pgfpathcurveto{\pgfqpoint{3.539136in}{5.123158in}}{\pgfqpoint{3.536822in}{5.128744in}}{\pgfqpoint{3.532704in}{5.132862in}}%
\pgfpathcurveto{\pgfqpoint{3.528586in}{5.136980in}}{\pgfqpoint{3.523000in}{5.139294in}}{\pgfqpoint{3.517176in}{5.139294in}}%
\pgfpathcurveto{\pgfqpoint{3.511352in}{5.139294in}}{\pgfqpoint{3.505766in}{5.136980in}}{\pgfqpoint{3.501648in}{5.132862in}}%
\pgfpathcurveto{\pgfqpoint{3.497530in}{5.128744in}}{\pgfqpoint{3.495216in}{5.123158in}}{\pgfqpoint{3.495216in}{5.117334in}}%
\pgfpathcurveto{\pgfqpoint{3.495216in}{5.111510in}}{\pgfqpoint{3.497530in}{5.105924in}}{\pgfqpoint{3.501648in}{5.101806in}}%
\pgfpathcurveto{\pgfqpoint{3.505766in}{5.097688in}}{\pgfqpoint{3.511352in}{5.095374in}}{\pgfqpoint{3.517176in}{5.095374in}}%
\pgfpathlineto{\pgfqpoint{3.517176in}{5.095374in}}%
\pgfpathclose%
\pgfusepath{stroke,fill}%
\end{pgfscope}%
\begin{pgfscope}%
\pgfpathrectangle{\pgfqpoint{1.414736in}{0.548769in}}{\pgfqpoint{3.497747in}{7.253186in}}%
\pgfusepath{clip}%
\pgfsetbuttcap%
\pgfsetroundjoin%
\definecolor{currentfill}{rgb}{0.154815,0.493313,0.557840}%
\pgfsetfillcolor{currentfill}%
\pgfsetlinewidth{1.003750pt}%
\definecolor{currentstroke}{rgb}{0.154815,0.493313,0.557840}%
\pgfsetstrokecolor{currentstroke}%
\pgfsetdash{}{0pt}%
\pgfpathmoveto{\pgfqpoint{3.738442in}{5.566360in}}%
\pgfpathcurveto{\pgfqpoint{3.744266in}{5.566360in}}{\pgfqpoint{3.749852in}{5.568674in}}{\pgfqpoint{3.753971in}{5.572792in}}%
\pgfpathcurveto{\pgfqpoint{3.758089in}{5.576910in}}{\pgfqpoint{3.760403in}{5.582496in}}{\pgfqpoint{3.760403in}{5.588320in}}%
\pgfpathcurveto{\pgfqpoint{3.760403in}{5.594144in}}{\pgfqpoint{3.758089in}{5.599730in}}{\pgfqpoint{3.753971in}{5.603848in}}%
\pgfpathcurveto{\pgfqpoint{3.749852in}{5.607967in}}{\pgfqpoint{3.744266in}{5.610280in}}{\pgfqpoint{3.738442in}{5.610280in}}%
\pgfpathcurveto{\pgfqpoint{3.732618in}{5.610280in}}{\pgfqpoint{3.727032in}{5.607967in}}{\pgfqpoint{3.722914in}{5.603848in}}%
\pgfpathcurveto{\pgfqpoint{3.718796in}{5.599730in}}{\pgfqpoint{3.716482in}{5.594144in}}{\pgfqpoint{3.716482in}{5.588320in}}%
\pgfpathcurveto{\pgfqpoint{3.716482in}{5.582496in}}{\pgfqpoint{3.718796in}{5.576910in}}{\pgfqpoint{3.722914in}{5.572792in}}%
\pgfpathcurveto{\pgfqpoint{3.727032in}{5.568674in}}{\pgfqpoint{3.732618in}{5.566360in}}{\pgfqpoint{3.738442in}{5.566360in}}%
\pgfpathlineto{\pgfqpoint{3.738442in}{5.566360in}}%
\pgfpathclose%
\pgfusepath{stroke,fill}%
\end{pgfscope}%
\begin{pgfscope}%
\pgfpathrectangle{\pgfqpoint{1.414736in}{0.548769in}}{\pgfqpoint{3.497747in}{7.253186in}}%
\pgfusepath{clip}%
\pgfsetbuttcap%
\pgfsetroundjoin%
\definecolor{currentfill}{rgb}{0.955300,0.901065,0.118128}%
\pgfsetfillcolor{currentfill}%
\pgfsetlinewidth{1.003750pt}%
\definecolor{currentstroke}{rgb}{0.955300,0.901065,0.118128}%
\pgfsetstrokecolor{currentstroke}%
\pgfsetdash{}{0pt}%
\pgfpathmoveto{\pgfqpoint{4.504169in}{6.037346in}}%
\pgfpathcurveto{\pgfqpoint{4.509993in}{6.037346in}}{\pgfqpoint{4.515579in}{6.039660in}}{\pgfqpoint{4.519697in}{6.043778in}}%
\pgfpathcurveto{\pgfqpoint{4.523815in}{6.047896in}}{\pgfqpoint{4.526129in}{6.053482in}}{\pgfqpoint{4.526129in}{6.059306in}}%
\pgfpathcurveto{\pgfqpoint{4.526129in}{6.065130in}}{\pgfqpoint{4.523815in}{6.070716in}}{\pgfqpoint{4.519697in}{6.074835in}}%
\pgfpathcurveto{\pgfqpoint{4.515579in}{6.078953in}}{\pgfqpoint{4.509993in}{6.081267in}}{\pgfqpoint{4.504169in}{6.081267in}}%
\pgfpathcurveto{\pgfqpoint{4.498345in}{6.081267in}}{\pgfqpoint{4.492759in}{6.078953in}}{\pgfqpoint{4.488641in}{6.074835in}}%
\pgfpathcurveto{\pgfqpoint{4.484523in}{6.070716in}}{\pgfqpoint{4.482209in}{6.065130in}}{\pgfqpoint{4.482209in}{6.059306in}}%
\pgfpathcurveto{\pgfqpoint{4.482209in}{6.053482in}}{\pgfqpoint{4.484523in}{6.047896in}}{\pgfqpoint{4.488641in}{6.043778in}}%
\pgfpathcurveto{\pgfqpoint{4.492759in}{6.039660in}}{\pgfqpoint{4.498345in}{6.037346in}}{\pgfqpoint{4.504169in}{6.037346in}}%
\pgfpathlineto{\pgfqpoint{4.504169in}{6.037346in}}%
\pgfpathclose%
\pgfusepath{stroke,fill}%
\end{pgfscope}%
\begin{pgfscope}%
\pgfpathrectangle{\pgfqpoint{1.414736in}{0.548769in}}{\pgfqpoint{3.497747in}{7.253186in}}%
\pgfusepath{clip}%
\pgfsetbuttcap%
\pgfsetroundjoin%
\definecolor{currentfill}{rgb}{0.993248,0.906157,0.143936}%
\pgfsetfillcolor{currentfill}%
\pgfsetlinewidth{1.003750pt}%
\definecolor{currentstroke}{rgb}{0.993248,0.906157,0.143936}%
\pgfsetstrokecolor{currentstroke}%
\pgfsetdash{}{0pt}%
\pgfpathmoveto{\pgfqpoint{2.670778in}{6.508332in}}%
\pgfpathcurveto{\pgfqpoint{2.676601in}{6.508332in}}{\pgfqpoint{2.682188in}{6.510646in}}{\pgfqpoint{2.686306in}{6.514764in}}%
\pgfpathcurveto{\pgfqpoint{2.690424in}{6.518882in}}{\pgfqpoint{2.692738in}{6.524468in}}{\pgfqpoint{2.692738in}{6.530292in}}%
\pgfpathcurveto{\pgfqpoint{2.692738in}{6.536116in}}{\pgfqpoint{2.690424in}{6.541702in}}{\pgfqpoint{2.686306in}{6.545821in}}%
\pgfpathcurveto{\pgfqpoint{2.682188in}{6.549939in}}{\pgfqpoint{2.676601in}{6.552253in}}{\pgfqpoint{2.670778in}{6.552253in}}%
\pgfpathcurveto{\pgfqpoint{2.664954in}{6.552253in}}{\pgfqpoint{2.659367in}{6.549939in}}{\pgfqpoint{2.655249in}{6.545821in}}%
\pgfpathcurveto{\pgfqpoint{2.651131in}{6.541702in}}{\pgfqpoint{2.648817in}{6.536116in}}{\pgfqpoint{2.648817in}{6.530292in}}%
\pgfpathcurveto{\pgfqpoint{2.648817in}{6.524468in}}{\pgfqpoint{2.651131in}{6.518882in}}{\pgfqpoint{2.655249in}{6.514764in}}%
\pgfpathcurveto{\pgfqpoint{2.659367in}{6.510646in}}{\pgfqpoint{2.664954in}{6.508332in}}{\pgfqpoint{2.670778in}{6.508332in}}%
\pgfpathlineto{\pgfqpoint{2.670778in}{6.508332in}}%
\pgfpathclose%
\pgfusepath{stroke,fill}%
\end{pgfscope}%
\begin{pgfscope}%
\pgfpathrectangle{\pgfqpoint{1.414736in}{0.548769in}}{\pgfqpoint{3.497747in}{7.253186in}}%
\pgfusepath{clip}%
\pgfsetbuttcap%
\pgfsetroundjoin%
\definecolor{currentfill}{rgb}{0.876168,0.891125,0.095250}%
\pgfsetfillcolor{currentfill}%
\pgfsetlinewidth{1.003750pt}%
\definecolor{currentstroke}{rgb}{0.876168,0.891125,0.095250}%
\pgfsetstrokecolor{currentstroke}%
\pgfsetdash{}{0pt}%
\pgfpathmoveto{\pgfqpoint{3.871097in}{6.979318in}}%
\pgfpathcurveto{\pgfqpoint{3.876921in}{6.979318in}}{\pgfqpoint{3.882507in}{6.981632in}}{\pgfqpoint{3.886626in}{6.985750in}}%
\pgfpathcurveto{\pgfqpoint{3.890744in}{6.989868in}}{\pgfqpoint{3.893058in}{6.995455in}}{\pgfqpoint{3.893058in}{7.001278in}}%
\pgfpathcurveto{\pgfqpoint{3.893058in}{7.007102in}}{\pgfqpoint{3.890744in}{7.012689in}}{\pgfqpoint{3.886626in}{7.016807in}}%
\pgfpathcurveto{\pgfqpoint{3.882507in}{7.020925in}}{\pgfqpoint{3.876921in}{7.023239in}}{\pgfqpoint{3.871097in}{7.023239in}}%
\pgfpathcurveto{\pgfqpoint{3.865273in}{7.023239in}}{\pgfqpoint{3.859687in}{7.020925in}}{\pgfqpoint{3.855569in}{7.016807in}}%
\pgfpathcurveto{\pgfqpoint{3.851451in}{7.012689in}}{\pgfqpoint{3.849137in}{7.007102in}}{\pgfqpoint{3.849137in}{7.001278in}}%
\pgfpathcurveto{\pgfqpoint{3.849137in}{6.995455in}}{\pgfqpoint{3.851451in}{6.989868in}}{\pgfqpoint{3.855569in}{6.985750in}}%
\pgfpathcurveto{\pgfqpoint{3.859687in}{6.981632in}}{\pgfqpoint{3.865273in}{6.979318in}}{\pgfqpoint{3.871097in}{6.979318in}}%
\pgfpathlineto{\pgfqpoint{3.871097in}{6.979318in}}%
\pgfpathclose%
\pgfusepath{stroke,fill}%
\end{pgfscope}%
\begin{pgfscope}%
\pgfpathrectangle{\pgfqpoint{1.414736in}{0.548769in}}{\pgfqpoint{3.497747in}{7.253186in}}%
\pgfusepath{clip}%
\pgfsetbuttcap%
\pgfsetroundjoin%
\definecolor{currentfill}{rgb}{0.783315,0.879285,0.125405}%
\pgfsetfillcolor{currentfill}%
\pgfsetlinewidth{1.003750pt}%
\definecolor{currentstroke}{rgb}{0.783315,0.879285,0.125405}%
\pgfsetstrokecolor{currentstroke}%
\pgfsetdash{}{0pt}%
\pgfpathmoveto{\pgfqpoint{3.765595in}{7.450304in}}%
\pgfpathcurveto{\pgfqpoint{3.771419in}{7.450304in}}{\pgfqpoint{3.777005in}{7.452618in}}{\pgfqpoint{3.781123in}{7.456736in}}%
\pgfpathcurveto{\pgfqpoint{3.785242in}{7.460854in}}{\pgfqpoint{3.787555in}{7.466441in}}{\pgfqpoint{3.787555in}{7.472265in}}%
\pgfpathcurveto{\pgfqpoint{3.787555in}{7.478089in}}{\pgfqpoint{3.785242in}{7.483675in}}{\pgfqpoint{3.781123in}{7.487793in}}%
\pgfpathcurveto{\pgfqpoint{3.777005in}{7.491911in}}{\pgfqpoint{3.771419in}{7.494225in}}{\pgfqpoint{3.765595in}{7.494225in}}%
\pgfpathcurveto{\pgfqpoint{3.759771in}{7.494225in}}{\pgfqpoint{3.754185in}{7.491911in}}{\pgfqpoint{3.750067in}{7.487793in}}%
\pgfpathcurveto{\pgfqpoint{3.745949in}{7.483675in}}{\pgfqpoint{3.743635in}{7.478089in}}{\pgfqpoint{3.743635in}{7.472265in}}%
\pgfpathcurveto{\pgfqpoint{3.743635in}{7.466441in}}{\pgfqpoint{3.745949in}{7.460854in}}{\pgfqpoint{3.750067in}{7.456736in}}%
\pgfpathcurveto{\pgfqpoint{3.754185in}{7.452618in}}{\pgfqpoint{3.759771in}{7.450304in}}{\pgfqpoint{3.765595in}{7.450304in}}%
\pgfpathlineto{\pgfqpoint{3.765595in}{7.450304in}}%
\pgfpathclose%
\pgfusepath{stroke,fill}%
\end{pgfscope}%
\begin{pgfscope}%
\pgfpathrectangle{\pgfqpoint{1.414736in}{0.548769in}}{\pgfqpoint{3.497747in}{7.253186in}}%
\pgfusepath{clip}%
\pgfsetbuttcap%
\pgfsetroundjoin%
\definecolor{currentfill}{rgb}{0.993248,0.906157,0.143936}%
\pgfsetfillcolor{currentfill}%
\pgfsetlinewidth{1.003750pt}%
\definecolor{currentstroke}{rgb}{0.993248,0.906157,0.143936}%
\pgfsetstrokecolor{currentstroke}%
\pgfsetdash{}{0pt}%
\pgfpathmoveto{\pgfqpoint{3.645354in}{0.856499in}}%
\pgfpathcurveto{\pgfqpoint{3.651178in}{0.856499in}}{\pgfqpoint{3.656764in}{0.858813in}}{\pgfqpoint{3.660882in}{0.862931in}}%
\pgfpathcurveto{\pgfqpoint{3.665000in}{0.867049in}}{\pgfqpoint{3.667314in}{0.872635in}}{\pgfqpoint{3.667314in}{0.878459in}}%
\pgfpathcurveto{\pgfqpoint{3.667314in}{0.884283in}}{\pgfqpoint{3.665000in}{0.889869in}}{\pgfqpoint{3.660882in}{0.893987in}}%
\pgfpathcurveto{\pgfqpoint{3.656764in}{0.898105in}}{\pgfqpoint{3.651178in}{0.900419in}}{\pgfqpoint{3.645354in}{0.900419in}}%
\pgfpathcurveto{\pgfqpoint{3.639530in}{0.900419in}}{\pgfqpoint{3.633944in}{0.898105in}}{\pgfqpoint{3.629826in}{0.893987in}}%
\pgfpathcurveto{\pgfqpoint{3.625707in}{0.889869in}}{\pgfqpoint{3.623393in}{0.884283in}}{\pgfqpoint{3.623393in}{0.878459in}}%
\pgfpathcurveto{\pgfqpoint{3.623393in}{0.872635in}}{\pgfqpoint{3.625707in}{0.867049in}}{\pgfqpoint{3.629826in}{0.862931in}}%
\pgfpathcurveto{\pgfqpoint{3.633944in}{0.858813in}}{\pgfqpoint{3.639530in}{0.856499in}}{\pgfqpoint{3.645354in}{0.856499in}}%
\pgfpathlineto{\pgfqpoint{3.645354in}{0.856499in}}%
\pgfpathclose%
\pgfusepath{stroke,fill}%
\end{pgfscope}%
\begin{pgfscope}%
\pgfpathrectangle{\pgfqpoint{1.414736in}{0.548769in}}{\pgfqpoint{3.497747in}{7.253186in}}%
\pgfusepath{clip}%
\pgfsetbuttcap%
\pgfsetroundjoin%
\definecolor{currentfill}{rgb}{0.783315,0.879285,0.125405}%
\pgfsetfillcolor{currentfill}%
\pgfsetlinewidth{1.003750pt}%
\definecolor{currentstroke}{rgb}{0.783315,0.879285,0.125405}%
\pgfsetstrokecolor{currentstroke}%
\pgfsetdash{}{0pt}%
\pgfpathmoveto{\pgfqpoint{3.310363in}{1.327485in}}%
\pgfpathcurveto{\pgfqpoint{3.316187in}{1.327485in}}{\pgfqpoint{3.321773in}{1.329799in}}{\pgfqpoint{3.325892in}{1.333917in}}%
\pgfpathcurveto{\pgfqpoint{3.330010in}{1.338035in}}{\pgfqpoint{3.332324in}{1.343621in}}{\pgfqpoint{3.332324in}{1.349445in}}%
\pgfpathcurveto{\pgfqpoint{3.332324in}{1.355269in}}{\pgfqpoint{3.330010in}{1.360855in}}{\pgfqpoint{3.325892in}{1.364973in}}%
\pgfpathcurveto{\pgfqpoint{3.321773in}{1.369092in}}{\pgfqpoint{3.316187in}{1.371405in}}{\pgfqpoint{3.310363in}{1.371405in}}%
\pgfpathcurveto{\pgfqpoint{3.304539in}{1.371405in}}{\pgfqpoint{3.298953in}{1.369092in}}{\pgfqpoint{3.294835in}{1.364973in}}%
\pgfpathcurveto{\pgfqpoint{3.290717in}{1.360855in}}{\pgfqpoint{3.288403in}{1.355269in}}{\pgfqpoint{3.288403in}{1.349445in}}%
\pgfpathcurveto{\pgfqpoint{3.288403in}{1.343621in}}{\pgfqpoint{3.290717in}{1.338035in}}{\pgfqpoint{3.294835in}{1.333917in}}%
\pgfpathcurveto{\pgfqpoint{3.298953in}{1.329799in}}{\pgfqpoint{3.304539in}{1.327485in}}{\pgfqpoint{3.310363in}{1.327485in}}%
\pgfpathlineto{\pgfqpoint{3.310363in}{1.327485in}}%
\pgfpathclose%
\pgfusepath{stroke,fill}%
\end{pgfscope}%
\begin{pgfscope}%
\pgfpathrectangle{\pgfqpoint{1.414736in}{0.548769in}}{\pgfqpoint{3.497747in}{7.253186in}}%
\pgfusepath{clip}%
\pgfsetbuttcap%
\pgfsetroundjoin%
\definecolor{currentfill}{rgb}{0.412913,0.803041,0.357269}%
\pgfsetfillcolor{currentfill}%
\pgfsetlinewidth{1.003750pt}%
\definecolor{currentstroke}{rgb}{0.412913,0.803041,0.357269}%
\pgfsetstrokecolor{currentstroke}%
\pgfsetdash{}{0pt}%
\pgfpathmoveto{\pgfqpoint{3.382724in}{1.798471in}}%
\pgfpathcurveto{\pgfqpoint{3.388547in}{1.798471in}}{\pgfqpoint{3.394134in}{1.800785in}}{\pgfqpoint{3.398252in}{1.804903in}}%
\pgfpathcurveto{\pgfqpoint{3.402370in}{1.809021in}}{\pgfqpoint{3.404684in}{1.814607in}}{\pgfqpoint{3.404684in}{1.820431in}}%
\pgfpathcurveto{\pgfqpoint{3.404684in}{1.826255in}}{\pgfqpoint{3.402370in}{1.831841in}}{\pgfqpoint{3.398252in}{1.835960in}}%
\pgfpathcurveto{\pgfqpoint{3.394134in}{1.840078in}}{\pgfqpoint{3.388547in}{1.842392in}}{\pgfqpoint{3.382724in}{1.842392in}}%
\pgfpathcurveto{\pgfqpoint{3.376900in}{1.842392in}}{\pgfqpoint{3.371313in}{1.840078in}}{\pgfqpoint{3.367195in}{1.835960in}}%
\pgfpathcurveto{\pgfqpoint{3.363077in}{1.831841in}}{\pgfqpoint{3.360763in}{1.826255in}}{\pgfqpoint{3.360763in}{1.820431in}}%
\pgfpathcurveto{\pgfqpoint{3.360763in}{1.814607in}}{\pgfqpoint{3.363077in}{1.809021in}}{\pgfqpoint{3.367195in}{1.804903in}}%
\pgfpathcurveto{\pgfqpoint{3.371313in}{1.800785in}}{\pgfqpoint{3.376900in}{1.798471in}}{\pgfqpoint{3.382724in}{1.798471in}}%
\pgfpathlineto{\pgfqpoint{3.382724in}{1.798471in}}%
\pgfpathclose%
\pgfusepath{stroke,fill}%
\end{pgfscope}%
\begin{pgfscope}%
\pgfpathrectangle{\pgfqpoint{1.414736in}{0.548769in}}{\pgfqpoint{3.497747in}{7.253186in}}%
\pgfusepath{clip}%
\pgfsetbuttcap%
\pgfsetroundjoin%
\definecolor{currentfill}{rgb}{0.993248,0.906157,0.143936}%
\pgfsetfillcolor{currentfill}%
\pgfsetlinewidth{1.003750pt}%
\definecolor{currentstroke}{rgb}{0.993248,0.906157,0.143936}%
\pgfsetstrokecolor{currentstroke}%
\pgfsetdash{}{0pt}%
\pgfpathmoveto{\pgfqpoint{3.494880in}{2.269457in}}%
\pgfpathcurveto{\pgfqpoint{3.500704in}{2.269457in}}{\pgfqpoint{3.506291in}{2.271771in}}{\pgfqpoint{3.510409in}{2.275889in}}%
\pgfpathcurveto{\pgfqpoint{3.514527in}{2.280007in}}{\pgfqpoint{3.516841in}{2.285593in}}{\pgfqpoint{3.516841in}{2.291417in}}%
\pgfpathcurveto{\pgfqpoint{3.516841in}{2.297241in}}{\pgfqpoint{3.514527in}{2.302828in}}{\pgfqpoint{3.510409in}{2.306946in}}%
\pgfpathcurveto{\pgfqpoint{3.506291in}{2.311064in}}{\pgfqpoint{3.500704in}{2.313378in}}{\pgfqpoint{3.494880in}{2.313378in}}%
\pgfpathcurveto{\pgfqpoint{3.489057in}{2.313378in}}{\pgfqpoint{3.483470in}{2.311064in}}{\pgfqpoint{3.479352in}{2.306946in}}%
\pgfpathcurveto{\pgfqpoint{3.475234in}{2.302828in}}{\pgfqpoint{3.472920in}{2.297241in}}{\pgfqpoint{3.472920in}{2.291417in}}%
\pgfpathcurveto{\pgfqpoint{3.472920in}{2.285593in}}{\pgfqpoint{3.475234in}{2.280007in}}{\pgfqpoint{3.479352in}{2.275889in}}%
\pgfpathcurveto{\pgfqpoint{3.483470in}{2.271771in}}{\pgfqpoint{3.489057in}{2.269457in}}{\pgfqpoint{3.494880in}{2.269457in}}%
\pgfpathlineto{\pgfqpoint{3.494880in}{2.269457in}}%
\pgfpathclose%
\pgfusepath{stroke,fill}%
\end{pgfscope}%
\begin{pgfscope}%
\pgfpathrectangle{\pgfqpoint{1.414736in}{0.548769in}}{\pgfqpoint{3.497747in}{7.253186in}}%
\pgfusepath{clip}%
\pgfsetbuttcap%
\pgfsetroundjoin%
\definecolor{currentfill}{rgb}{0.575563,0.844566,0.256415}%
\pgfsetfillcolor{currentfill}%
\pgfsetlinewidth{1.003750pt}%
\definecolor{currentstroke}{rgb}{0.575563,0.844566,0.256415}%
\pgfsetstrokecolor{currentstroke}%
\pgfsetdash{}{0pt}%
\pgfpathmoveto{\pgfqpoint{3.416781in}{2.740443in}}%
\pgfpathcurveto{\pgfqpoint{3.422605in}{2.740443in}}{\pgfqpoint{3.428191in}{2.742757in}}{\pgfqpoint{3.432309in}{2.746875in}}%
\pgfpathcurveto{\pgfqpoint{3.436427in}{2.750993in}}{\pgfqpoint{3.438741in}{2.756580in}}{\pgfqpoint{3.438741in}{2.762404in}}%
\pgfpathcurveto{\pgfqpoint{3.438741in}{2.768227in}}{\pgfqpoint{3.436427in}{2.773814in}}{\pgfqpoint{3.432309in}{2.777932in}}%
\pgfpathcurveto{\pgfqpoint{3.428191in}{2.782050in}}{\pgfqpoint{3.422605in}{2.784364in}}{\pgfqpoint{3.416781in}{2.784364in}}%
\pgfpathcurveto{\pgfqpoint{3.410957in}{2.784364in}}{\pgfqpoint{3.405371in}{2.782050in}}{\pgfqpoint{3.401252in}{2.777932in}}%
\pgfpathcurveto{\pgfqpoint{3.397134in}{2.773814in}}{\pgfqpoint{3.394820in}{2.768227in}}{\pgfqpoint{3.394820in}{2.762404in}}%
\pgfpathcurveto{\pgfqpoint{3.394820in}{2.756580in}}{\pgfqpoint{3.397134in}{2.750993in}}{\pgfqpoint{3.401252in}{2.746875in}}%
\pgfpathcurveto{\pgfqpoint{3.405371in}{2.742757in}}{\pgfqpoint{3.410957in}{2.740443in}}{\pgfqpoint{3.416781in}{2.740443in}}%
\pgfpathlineto{\pgfqpoint{3.416781in}{2.740443in}}%
\pgfpathclose%
\pgfusepath{stroke,fill}%
\end{pgfscope}%
\begin{pgfscope}%
\pgfpathrectangle{\pgfqpoint{1.414736in}{0.548769in}}{\pgfqpoint{3.497747in}{7.253186in}}%
\pgfusepath{clip}%
\pgfsetbuttcap%
\pgfsetroundjoin%
\definecolor{currentfill}{rgb}{0.699415,0.867117,0.175971}%
\pgfsetfillcolor{currentfill}%
\pgfsetlinewidth{1.003750pt}%
\definecolor{currentstroke}{rgb}{0.699415,0.867117,0.175971}%
\pgfsetstrokecolor{currentstroke}%
\pgfsetdash{}{0pt}%
\pgfpathmoveto{\pgfqpoint{3.532887in}{3.211429in}}%
\pgfpathcurveto{\pgfqpoint{3.538711in}{3.211429in}}{\pgfqpoint{3.544297in}{3.213743in}}{\pgfqpoint{3.548415in}{3.217861in}}%
\pgfpathcurveto{\pgfqpoint{3.552533in}{3.221980in}}{\pgfqpoint{3.554847in}{3.227566in}}{\pgfqpoint{3.554847in}{3.233390in}}%
\pgfpathcurveto{\pgfqpoint{3.554847in}{3.239214in}}{\pgfqpoint{3.552533in}{3.244800in}}{\pgfqpoint{3.548415in}{3.248918in}}%
\pgfpathcurveto{\pgfqpoint{3.544297in}{3.253036in}}{\pgfqpoint{3.538711in}{3.255350in}}{\pgfqpoint{3.532887in}{3.255350in}}%
\pgfpathcurveto{\pgfqpoint{3.527063in}{3.255350in}}{\pgfqpoint{3.521477in}{3.253036in}}{\pgfqpoint{3.517359in}{3.248918in}}%
\pgfpathcurveto{\pgfqpoint{3.513241in}{3.244800in}}{\pgfqpoint{3.510927in}{3.239214in}}{\pgfqpoint{3.510927in}{3.233390in}}%
\pgfpathcurveto{\pgfqpoint{3.510927in}{3.227566in}}{\pgfqpoint{3.513241in}{3.221980in}}{\pgfqpoint{3.517359in}{3.217861in}}%
\pgfpathcurveto{\pgfqpoint{3.521477in}{3.213743in}}{\pgfqpoint{3.527063in}{3.211429in}}{\pgfqpoint{3.532887in}{3.211429in}}%
\pgfpathlineto{\pgfqpoint{3.532887in}{3.211429in}}%
\pgfpathclose%
\pgfusepath{stroke,fill}%
\end{pgfscope}%
\begin{pgfscope}%
\pgfpathrectangle{\pgfqpoint{1.414736in}{0.548769in}}{\pgfqpoint{3.497747in}{7.253186in}}%
\pgfusepath{clip}%
\pgfsetbuttcap%
\pgfsetroundjoin%
\definecolor{currentfill}{rgb}{0.993248,0.906157,0.143936}%
\pgfsetfillcolor{currentfill}%
\pgfsetlinewidth{1.003750pt}%
\definecolor{currentstroke}{rgb}{0.993248,0.906157,0.143936}%
\pgfsetstrokecolor{currentstroke}%
\pgfsetdash{}{0pt}%
\pgfpathmoveto{\pgfqpoint{3.706379in}{3.682415in}}%
\pgfpathcurveto{\pgfqpoint{3.712203in}{3.682415in}}{\pgfqpoint{3.717789in}{3.684729in}}{\pgfqpoint{3.721907in}{3.688847in}}%
\pgfpathcurveto{\pgfqpoint{3.726025in}{3.692966in}}{\pgfqpoint{3.728339in}{3.698552in}}{\pgfqpoint{3.728339in}{3.704376in}}%
\pgfpathcurveto{\pgfqpoint{3.728339in}{3.710200in}}{\pgfqpoint{3.726025in}{3.715786in}}{\pgfqpoint{3.721907in}{3.719904in}}%
\pgfpathcurveto{\pgfqpoint{3.717789in}{3.724022in}}{\pgfqpoint{3.712203in}{3.726336in}}{\pgfqpoint{3.706379in}{3.726336in}}%
\pgfpathcurveto{\pgfqpoint{3.700555in}{3.726336in}}{\pgfqpoint{3.694969in}{3.724022in}}{\pgfqpoint{3.690851in}{3.719904in}}%
\pgfpathcurveto{\pgfqpoint{3.686733in}{3.715786in}}{\pgfqpoint{3.684419in}{3.710200in}}{\pgfqpoint{3.684419in}{3.704376in}}%
\pgfpathcurveto{\pgfqpoint{3.684419in}{3.698552in}}{\pgfqpoint{3.686733in}{3.692966in}}{\pgfqpoint{3.690851in}{3.688847in}}%
\pgfpathcurveto{\pgfqpoint{3.694969in}{3.684729in}}{\pgfqpoint{3.700555in}{3.682415in}}{\pgfqpoint{3.706379in}{3.682415in}}%
\pgfpathlineto{\pgfqpoint{3.706379in}{3.682415in}}%
\pgfpathclose%
\pgfusepath{stroke,fill}%
\end{pgfscope}%
\begin{pgfscope}%
\pgfpathrectangle{\pgfqpoint{1.414736in}{0.548769in}}{\pgfqpoint{3.497747in}{7.253186in}}%
\pgfusepath{clip}%
\pgfsetbuttcap%
\pgfsetroundjoin%
\definecolor{currentfill}{rgb}{0.993248,0.906157,0.143936}%
\pgfsetfillcolor{currentfill}%
\pgfsetlinewidth{1.003750pt}%
\definecolor{currentstroke}{rgb}{0.993248,0.906157,0.143936}%
\pgfsetstrokecolor{currentstroke}%
\pgfsetdash{}{0pt}%
\pgfpathmoveto{\pgfqpoint{3.729880in}{4.153402in}}%
\pgfpathcurveto{\pgfqpoint{3.735704in}{4.153402in}}{\pgfqpoint{3.741291in}{4.155715in}}{\pgfqpoint{3.745409in}{4.159834in}}%
\pgfpathcurveto{\pgfqpoint{3.749527in}{4.163952in}}{\pgfqpoint{3.751841in}{4.169538in}}{\pgfqpoint{3.751841in}{4.175362in}}%
\pgfpathcurveto{\pgfqpoint{3.751841in}{4.181186in}}{\pgfqpoint{3.749527in}{4.186772in}}{\pgfqpoint{3.745409in}{4.190890in}}%
\pgfpathcurveto{\pgfqpoint{3.741291in}{4.195008in}}{\pgfqpoint{3.735704in}{4.197322in}}{\pgfqpoint{3.729880in}{4.197322in}}%
\pgfpathcurveto{\pgfqpoint{3.724057in}{4.197322in}}{\pgfqpoint{3.718470in}{4.195008in}}{\pgfqpoint{3.714352in}{4.190890in}}%
\pgfpathcurveto{\pgfqpoint{3.710234in}{4.186772in}}{\pgfqpoint{3.707920in}{4.181186in}}{\pgfqpoint{3.707920in}{4.175362in}}%
\pgfpathcurveto{\pgfqpoint{3.707920in}{4.169538in}}{\pgfqpoint{3.710234in}{4.163952in}}{\pgfqpoint{3.714352in}{4.159834in}}%
\pgfpathcurveto{\pgfqpoint{3.718470in}{4.155715in}}{\pgfqpoint{3.724057in}{4.153402in}}{\pgfqpoint{3.729880in}{4.153402in}}%
\pgfpathlineto{\pgfqpoint{3.729880in}{4.153402in}}%
\pgfpathclose%
\pgfusepath{stroke,fill}%
\end{pgfscope}%
\begin{pgfscope}%
\pgfpathrectangle{\pgfqpoint{1.414736in}{0.548769in}}{\pgfqpoint{3.497747in}{7.253186in}}%
\pgfusepath{clip}%
\pgfsetbuttcap%
\pgfsetroundjoin%
\definecolor{currentfill}{rgb}{0.993248,0.906157,0.143936}%
\pgfsetfillcolor{currentfill}%
\pgfsetlinewidth{1.003750pt}%
\definecolor{currentstroke}{rgb}{0.993248,0.906157,0.143936}%
\pgfsetstrokecolor{currentstroke}%
\pgfsetdash{}{0pt}%
\pgfpathmoveto{\pgfqpoint{4.026782in}{4.624388in}}%
\pgfpathcurveto{\pgfqpoint{4.032606in}{4.624388in}}{\pgfqpoint{4.038192in}{4.626702in}}{\pgfqpoint{4.042310in}{4.630820in}}%
\pgfpathcurveto{\pgfqpoint{4.046428in}{4.634938in}}{\pgfqpoint{4.048742in}{4.640524in}}{\pgfqpoint{4.048742in}{4.646348in}}%
\pgfpathcurveto{\pgfqpoint{4.048742in}{4.652172in}}{\pgfqpoint{4.046428in}{4.657758in}}{\pgfqpoint{4.042310in}{4.661876in}}%
\pgfpathcurveto{\pgfqpoint{4.038192in}{4.665994in}}{\pgfqpoint{4.032606in}{4.668308in}}{\pgfqpoint{4.026782in}{4.668308in}}%
\pgfpathcurveto{\pgfqpoint{4.020958in}{4.668308in}}{\pgfqpoint{4.015372in}{4.665994in}}{\pgfqpoint{4.011254in}{4.661876in}}%
\pgfpathcurveto{\pgfqpoint{4.007136in}{4.657758in}}{\pgfqpoint{4.004822in}{4.652172in}}{\pgfqpoint{4.004822in}{4.646348in}}%
\pgfpathcurveto{\pgfqpoint{4.004822in}{4.640524in}}{\pgfqpoint{4.007136in}{4.634938in}}{\pgfqpoint{4.011254in}{4.630820in}}%
\pgfpathcurveto{\pgfqpoint{4.015372in}{4.626702in}}{\pgfqpoint{4.020958in}{4.624388in}}{\pgfqpoint{4.026782in}{4.624388in}}%
\pgfpathlineto{\pgfqpoint{4.026782in}{4.624388in}}%
\pgfpathclose%
\pgfusepath{stroke,fill}%
\end{pgfscope}%
\begin{pgfscope}%
\pgfpathrectangle{\pgfqpoint{1.414736in}{0.548769in}}{\pgfqpoint{3.497747in}{7.253186in}}%
\pgfusepath{clip}%
\pgfsetbuttcap%
\pgfsetroundjoin%
\definecolor{currentfill}{rgb}{0.335885,0.777018,0.402049}%
\pgfsetfillcolor{currentfill}%
\pgfsetlinewidth{1.003750pt}%
\definecolor{currentstroke}{rgb}{0.335885,0.777018,0.402049}%
\pgfsetstrokecolor{currentstroke}%
\pgfsetdash{}{0pt}%
\pgfpathmoveto{\pgfqpoint{3.794447in}{5.095374in}}%
\pgfpathcurveto{\pgfqpoint{3.800271in}{5.095374in}}{\pgfqpoint{3.805857in}{5.097688in}}{\pgfqpoint{3.809975in}{5.101806in}}%
\pgfpathcurveto{\pgfqpoint{3.814093in}{5.105924in}}{\pgfqpoint{3.816407in}{5.111510in}}{\pgfqpoint{3.816407in}{5.117334in}}%
\pgfpathcurveto{\pgfqpoint{3.816407in}{5.123158in}}{\pgfqpoint{3.814093in}{5.128744in}}{\pgfqpoint{3.809975in}{5.132862in}}%
\pgfpathcurveto{\pgfqpoint{3.805857in}{5.136980in}}{\pgfqpoint{3.800271in}{5.139294in}}{\pgfqpoint{3.794447in}{5.139294in}}%
\pgfpathcurveto{\pgfqpoint{3.788623in}{5.139294in}}{\pgfqpoint{3.783037in}{5.136980in}}{\pgfqpoint{3.778918in}{5.132862in}}%
\pgfpathcurveto{\pgfqpoint{3.774800in}{5.128744in}}{\pgfqpoint{3.772486in}{5.123158in}}{\pgfqpoint{3.772486in}{5.117334in}}%
\pgfpathcurveto{\pgfqpoint{3.772486in}{5.111510in}}{\pgfqpoint{3.774800in}{5.105924in}}{\pgfqpoint{3.778918in}{5.101806in}}%
\pgfpathcurveto{\pgfqpoint{3.783037in}{5.097688in}}{\pgfqpoint{3.788623in}{5.095374in}}{\pgfqpoint{3.794447in}{5.095374in}}%
\pgfpathlineto{\pgfqpoint{3.794447in}{5.095374in}}%
\pgfpathclose%
\pgfusepath{stroke,fill}%
\end{pgfscope}%
\begin{pgfscope}%
\pgfpathrectangle{\pgfqpoint{1.414736in}{0.548769in}}{\pgfqpoint{3.497747in}{7.253186in}}%
\pgfusepath{clip}%
\pgfsetbuttcap%
\pgfsetroundjoin%
\definecolor{currentfill}{rgb}{0.128729,0.563265,0.551229}%
\pgfsetfillcolor{currentfill}%
\pgfsetlinewidth{1.003750pt}%
\definecolor{currentstroke}{rgb}{0.128729,0.563265,0.551229}%
\pgfsetstrokecolor{currentstroke}%
\pgfsetdash{}{0pt}%
\pgfpathmoveto{\pgfqpoint{3.538595in}{5.566360in}}%
\pgfpathcurveto{\pgfqpoint{3.544419in}{5.566360in}}{\pgfqpoint{3.550005in}{5.568674in}}{\pgfqpoint{3.554124in}{5.572792in}}%
\pgfpathcurveto{\pgfqpoint{3.558242in}{5.576910in}}{\pgfqpoint{3.560556in}{5.582496in}}{\pgfqpoint{3.560556in}{5.588320in}}%
\pgfpathcurveto{\pgfqpoint{3.560556in}{5.594144in}}{\pgfqpoint{3.558242in}{5.599730in}}{\pgfqpoint{3.554124in}{5.603848in}}%
\pgfpathcurveto{\pgfqpoint{3.550005in}{5.607967in}}{\pgfqpoint{3.544419in}{5.610280in}}{\pgfqpoint{3.538595in}{5.610280in}}%
\pgfpathcurveto{\pgfqpoint{3.532771in}{5.610280in}}{\pgfqpoint{3.527185in}{5.607967in}}{\pgfqpoint{3.523067in}{5.603848in}}%
\pgfpathcurveto{\pgfqpoint{3.518949in}{5.599730in}}{\pgfqpoint{3.516635in}{5.594144in}}{\pgfqpoint{3.516635in}{5.588320in}}%
\pgfpathcurveto{\pgfqpoint{3.516635in}{5.582496in}}{\pgfqpoint{3.518949in}{5.576910in}}{\pgfqpoint{3.523067in}{5.572792in}}%
\pgfpathcurveto{\pgfqpoint{3.527185in}{5.568674in}}{\pgfqpoint{3.532771in}{5.566360in}}{\pgfqpoint{3.538595in}{5.566360in}}%
\pgfpathlineto{\pgfqpoint{3.538595in}{5.566360in}}%
\pgfpathclose%
\pgfusepath{stroke,fill}%
\end{pgfscope}%
\begin{pgfscope}%
\pgfpathrectangle{\pgfqpoint{1.414736in}{0.548769in}}{\pgfqpoint{3.497747in}{7.253186in}}%
\pgfusepath{clip}%
\pgfsetbuttcap%
\pgfsetroundjoin%
\definecolor{currentfill}{rgb}{0.993248,0.906157,0.143936}%
\pgfsetfillcolor{currentfill}%
\pgfsetlinewidth{1.003750pt}%
\definecolor{currentstroke}{rgb}{0.993248,0.906157,0.143936}%
\pgfsetstrokecolor{currentstroke}%
\pgfsetdash{}{0pt}%
\pgfpathmoveto{\pgfqpoint{3.295820in}{6.037346in}}%
\pgfpathcurveto{\pgfqpoint{3.301644in}{6.037346in}}{\pgfqpoint{3.307230in}{6.039660in}}{\pgfqpoint{3.311348in}{6.043778in}}%
\pgfpathcurveto{\pgfqpoint{3.315466in}{6.047896in}}{\pgfqpoint{3.317780in}{6.053482in}}{\pgfqpoint{3.317780in}{6.059306in}}%
\pgfpathcurveto{\pgfqpoint{3.317780in}{6.065130in}}{\pgfqpoint{3.315466in}{6.070716in}}{\pgfqpoint{3.311348in}{6.074835in}}%
\pgfpathcurveto{\pgfqpoint{3.307230in}{6.078953in}}{\pgfqpoint{3.301644in}{6.081267in}}{\pgfqpoint{3.295820in}{6.081267in}}%
\pgfpathcurveto{\pgfqpoint{3.289996in}{6.081267in}}{\pgfqpoint{3.284410in}{6.078953in}}{\pgfqpoint{3.280292in}{6.074835in}}%
\pgfpathcurveto{\pgfqpoint{3.276173in}{6.070716in}}{\pgfqpoint{3.273860in}{6.065130in}}{\pgfqpoint{3.273860in}{6.059306in}}%
\pgfpathcurveto{\pgfqpoint{3.273860in}{6.053482in}}{\pgfqpoint{3.276173in}{6.047896in}}{\pgfqpoint{3.280292in}{6.043778in}}%
\pgfpathcurveto{\pgfqpoint{3.284410in}{6.039660in}}{\pgfqpoint{3.289996in}{6.037346in}}{\pgfqpoint{3.295820in}{6.037346in}}%
\pgfpathlineto{\pgfqpoint{3.295820in}{6.037346in}}%
\pgfpathclose%
\pgfusepath{stroke,fill}%
\end{pgfscope}%
\begin{pgfscope}%
\pgfpathrectangle{\pgfqpoint{1.414736in}{0.548769in}}{\pgfqpoint{3.497747in}{7.253186in}}%
\pgfusepath{clip}%
\pgfsetbuttcap%
\pgfsetroundjoin%
\definecolor{currentfill}{rgb}{0.993248,0.906157,0.143936}%
\pgfsetfillcolor{currentfill}%
\pgfsetlinewidth{1.003750pt}%
\definecolor{currentstroke}{rgb}{0.993248,0.906157,0.143936}%
\pgfsetstrokecolor{currentstroke}%
\pgfsetdash{}{0pt}%
\pgfpathmoveto{\pgfqpoint{3.932917in}{6.508332in}}%
\pgfpathcurveto{\pgfqpoint{3.938741in}{6.508332in}}{\pgfqpoint{3.944327in}{6.510646in}}{\pgfqpoint{3.948445in}{6.514764in}}%
\pgfpathcurveto{\pgfqpoint{3.952563in}{6.518882in}}{\pgfqpoint{3.954877in}{6.524468in}}{\pgfqpoint{3.954877in}{6.530292in}}%
\pgfpathcurveto{\pgfqpoint{3.954877in}{6.536116in}}{\pgfqpoint{3.952563in}{6.541702in}}{\pgfqpoint{3.948445in}{6.545821in}}%
\pgfpathcurveto{\pgfqpoint{3.944327in}{6.549939in}}{\pgfqpoint{3.938741in}{6.552253in}}{\pgfqpoint{3.932917in}{6.552253in}}%
\pgfpathcurveto{\pgfqpoint{3.927093in}{6.552253in}}{\pgfqpoint{3.921507in}{6.549939in}}{\pgfqpoint{3.917389in}{6.545821in}}%
\pgfpathcurveto{\pgfqpoint{3.913271in}{6.541702in}}{\pgfqpoint{3.910957in}{6.536116in}}{\pgfqpoint{3.910957in}{6.530292in}}%
\pgfpathcurveto{\pgfqpoint{3.910957in}{6.524468in}}{\pgfqpoint{3.913271in}{6.518882in}}{\pgfqpoint{3.917389in}{6.514764in}}%
\pgfpathcurveto{\pgfqpoint{3.921507in}{6.510646in}}{\pgfqpoint{3.927093in}{6.508332in}}{\pgfqpoint{3.932917in}{6.508332in}}%
\pgfpathlineto{\pgfqpoint{3.932917in}{6.508332in}}%
\pgfpathclose%
\pgfusepath{stroke,fill}%
\end{pgfscope}%
\begin{pgfscope}%
\pgfpathrectangle{\pgfqpoint{1.414736in}{0.548769in}}{\pgfqpoint{3.497747in}{7.253186in}}%
\pgfusepath{clip}%
\pgfsetbuttcap%
\pgfsetroundjoin%
\definecolor{currentfill}{rgb}{0.993248,0.906157,0.143936}%
\pgfsetfillcolor{currentfill}%
\pgfsetlinewidth{1.003750pt}%
\definecolor{currentstroke}{rgb}{0.993248,0.906157,0.143936}%
\pgfsetstrokecolor{currentstroke}%
\pgfsetdash{}{0pt}%
\pgfpathmoveto{\pgfqpoint{3.481732in}{6.979318in}}%
\pgfpathcurveto{\pgfqpoint{3.487556in}{6.979318in}}{\pgfqpoint{3.493142in}{6.981632in}}{\pgfqpoint{3.497260in}{6.985750in}}%
\pgfpathcurveto{\pgfqpoint{3.501378in}{6.989868in}}{\pgfqpoint{3.503692in}{6.995455in}}{\pgfqpoint{3.503692in}{7.001278in}}%
\pgfpathcurveto{\pgfqpoint{3.503692in}{7.007102in}}{\pgfqpoint{3.501378in}{7.012689in}}{\pgfqpoint{3.497260in}{7.016807in}}%
\pgfpathcurveto{\pgfqpoint{3.493142in}{7.020925in}}{\pgfqpoint{3.487556in}{7.023239in}}{\pgfqpoint{3.481732in}{7.023239in}}%
\pgfpathcurveto{\pgfqpoint{3.475908in}{7.023239in}}{\pgfqpoint{3.470322in}{7.020925in}}{\pgfqpoint{3.466204in}{7.016807in}}%
\pgfpathcurveto{\pgfqpoint{3.462085in}{7.012689in}}{\pgfqpoint{3.459772in}{7.007102in}}{\pgfqpoint{3.459772in}{7.001278in}}%
\pgfpathcurveto{\pgfqpoint{3.459772in}{6.995455in}}{\pgfqpoint{3.462085in}{6.989868in}}{\pgfqpoint{3.466204in}{6.985750in}}%
\pgfpathcurveto{\pgfqpoint{3.470322in}{6.981632in}}{\pgfqpoint{3.475908in}{6.979318in}}{\pgfqpoint{3.481732in}{6.979318in}}%
\pgfpathlineto{\pgfqpoint{3.481732in}{6.979318in}}%
\pgfpathclose%
\pgfusepath{stroke,fill}%
\end{pgfscope}%
\begin{pgfscope}%
\pgfpathrectangle{\pgfqpoint{1.414736in}{0.548769in}}{\pgfqpoint{3.497747in}{7.253186in}}%
\pgfusepath{clip}%
\pgfsetbuttcap%
\pgfsetroundjoin%
\definecolor{currentfill}{rgb}{0.835270,0.886029,0.102646}%
\pgfsetfillcolor{currentfill}%
\pgfsetlinewidth{1.003750pt}%
\definecolor{currentstroke}{rgb}{0.835270,0.886029,0.102646}%
\pgfsetstrokecolor{currentstroke}%
\pgfsetdash{}{0pt}%
\pgfpathmoveto{\pgfqpoint{4.013835in}{7.450304in}}%
\pgfpathcurveto{\pgfqpoint{4.019659in}{7.450304in}}{\pgfqpoint{4.025245in}{7.452618in}}{\pgfqpoint{4.029364in}{7.456736in}}%
\pgfpathcurveto{\pgfqpoint{4.033482in}{7.460854in}}{\pgfqpoint{4.035796in}{7.466441in}}{\pgfqpoint{4.035796in}{7.472265in}}%
\pgfpathcurveto{\pgfqpoint{4.035796in}{7.478089in}}{\pgfqpoint{4.033482in}{7.483675in}}{\pgfqpoint{4.029364in}{7.487793in}}%
\pgfpathcurveto{\pgfqpoint{4.025245in}{7.491911in}}{\pgfqpoint{4.019659in}{7.494225in}}{\pgfqpoint{4.013835in}{7.494225in}}%
\pgfpathcurveto{\pgfqpoint{4.008011in}{7.494225in}}{\pgfqpoint{4.002425in}{7.491911in}}{\pgfqpoint{3.998307in}{7.487793in}}%
\pgfpathcurveto{\pgfqpoint{3.994189in}{7.483675in}}{\pgfqpoint{3.991875in}{7.478089in}}{\pgfqpoint{3.991875in}{7.472265in}}%
\pgfpathcurveto{\pgfqpoint{3.991875in}{7.466441in}}{\pgfqpoint{3.994189in}{7.460854in}}{\pgfqpoint{3.998307in}{7.456736in}}%
\pgfpathcurveto{\pgfqpoint{4.002425in}{7.452618in}}{\pgfqpoint{4.008011in}{7.450304in}}{\pgfqpoint{4.013835in}{7.450304in}}%
\pgfpathlineto{\pgfqpoint{4.013835in}{7.450304in}}%
\pgfpathclose%
\pgfusepath{stroke,fill}%
\end{pgfscope}%
\begin{pgfscope}%
\pgfpathrectangle{\pgfqpoint{1.414736in}{0.548769in}}{\pgfqpoint{3.497747in}{7.253186in}}%
\pgfusepath{clip}%
\pgfsetbuttcap%
\pgfsetroundjoin%
\definecolor{currentfill}{rgb}{0.993248,0.906157,0.143936}%
\pgfsetfillcolor{currentfill}%
\pgfsetlinewidth{1.003750pt}%
\definecolor{currentstroke}{rgb}{0.993248,0.906157,0.143936}%
\pgfsetstrokecolor{currentstroke}%
\pgfsetdash{}{0pt}%
\pgfpathmoveto{\pgfqpoint{3.493854in}{0.856499in}}%
\pgfpathcurveto{\pgfqpoint{3.499678in}{0.856499in}}{\pgfqpoint{3.505264in}{0.858813in}}{\pgfqpoint{3.509382in}{0.862931in}}%
\pgfpathcurveto{\pgfqpoint{3.513500in}{0.867049in}}{\pgfqpoint{3.515814in}{0.872635in}}{\pgfqpoint{3.515814in}{0.878459in}}%
\pgfpathcurveto{\pgfqpoint{3.515814in}{0.884283in}}{\pgfqpoint{3.513500in}{0.889869in}}{\pgfqpoint{3.509382in}{0.893987in}}%
\pgfpathcurveto{\pgfqpoint{3.505264in}{0.898105in}}{\pgfqpoint{3.499678in}{0.900419in}}{\pgfqpoint{3.493854in}{0.900419in}}%
\pgfpathcurveto{\pgfqpoint{3.488030in}{0.900419in}}{\pgfqpoint{3.482444in}{0.898105in}}{\pgfqpoint{3.478326in}{0.893987in}}%
\pgfpathcurveto{\pgfqpoint{3.474208in}{0.889869in}}{\pgfqpoint{3.471894in}{0.884283in}}{\pgfqpoint{3.471894in}{0.878459in}}%
\pgfpathcurveto{\pgfqpoint{3.471894in}{0.872635in}}{\pgfqpoint{3.474208in}{0.867049in}}{\pgfqpoint{3.478326in}{0.862931in}}%
\pgfpathcurveto{\pgfqpoint{3.482444in}{0.858813in}}{\pgfqpoint{3.488030in}{0.856499in}}{\pgfqpoint{3.493854in}{0.856499in}}%
\pgfpathlineto{\pgfqpoint{3.493854in}{0.856499in}}%
\pgfpathclose%
\pgfusepath{stroke,fill}%
\end{pgfscope}%
\begin{pgfscope}%
\pgfpathrectangle{\pgfqpoint{1.414736in}{0.548769in}}{\pgfqpoint{3.497747in}{7.253186in}}%
\pgfusepath{clip}%
\pgfsetbuttcap%
\pgfsetroundjoin%
\definecolor{currentfill}{rgb}{0.458674,0.816363,0.329727}%
\pgfsetfillcolor{currentfill}%
\pgfsetlinewidth{1.003750pt}%
\definecolor{currentstroke}{rgb}{0.458674,0.816363,0.329727}%
\pgfsetstrokecolor{currentstroke}%
\pgfsetdash{}{0pt}%
\pgfpathmoveto{\pgfqpoint{3.508191in}{1.327485in}}%
\pgfpathcurveto{\pgfqpoint{3.514015in}{1.327485in}}{\pgfqpoint{3.519602in}{1.329799in}}{\pgfqpoint{3.523720in}{1.333917in}}%
\pgfpathcurveto{\pgfqpoint{3.527838in}{1.338035in}}{\pgfqpoint{3.530152in}{1.343621in}}{\pgfqpoint{3.530152in}{1.349445in}}%
\pgfpathcurveto{\pgfqpoint{3.530152in}{1.355269in}}{\pgfqpoint{3.527838in}{1.360855in}}{\pgfqpoint{3.523720in}{1.364973in}}%
\pgfpathcurveto{\pgfqpoint{3.519602in}{1.369092in}}{\pgfqpoint{3.514015in}{1.371405in}}{\pgfqpoint{3.508191in}{1.371405in}}%
\pgfpathcurveto{\pgfqpoint{3.502368in}{1.371405in}}{\pgfqpoint{3.496781in}{1.369092in}}{\pgfqpoint{3.492663in}{1.364973in}}%
\pgfpathcurveto{\pgfqpoint{3.488545in}{1.360855in}}{\pgfqpoint{3.486231in}{1.355269in}}{\pgfqpoint{3.486231in}{1.349445in}}%
\pgfpathcurveto{\pgfqpoint{3.486231in}{1.343621in}}{\pgfqpoint{3.488545in}{1.338035in}}{\pgfqpoint{3.492663in}{1.333917in}}%
\pgfpathcurveto{\pgfqpoint{3.496781in}{1.329799in}}{\pgfqpoint{3.502368in}{1.327485in}}{\pgfqpoint{3.508191in}{1.327485in}}%
\pgfpathlineto{\pgfqpoint{3.508191in}{1.327485in}}%
\pgfpathclose%
\pgfusepath{stroke,fill}%
\end{pgfscope}%
\begin{pgfscope}%
\pgfpathrectangle{\pgfqpoint{1.414736in}{0.548769in}}{\pgfqpoint{3.497747in}{7.253186in}}%
\pgfusepath{clip}%
\pgfsetbuttcap%
\pgfsetroundjoin%
\definecolor{currentfill}{rgb}{0.964894,0.902323,0.123941}%
\pgfsetfillcolor{currentfill}%
\pgfsetlinewidth{1.003750pt}%
\definecolor{currentstroke}{rgb}{0.964894,0.902323,0.123941}%
\pgfsetstrokecolor{currentstroke}%
\pgfsetdash{}{0pt}%
\pgfpathmoveto{\pgfqpoint{3.946322in}{1.798471in}}%
\pgfpathcurveto{\pgfqpoint{3.952146in}{1.798471in}}{\pgfqpoint{3.957732in}{1.800785in}}{\pgfqpoint{3.961850in}{1.804903in}}%
\pgfpathcurveto{\pgfqpoint{3.965969in}{1.809021in}}{\pgfqpoint{3.968282in}{1.814607in}}{\pgfqpoint{3.968282in}{1.820431in}}%
\pgfpathcurveto{\pgfqpoint{3.968282in}{1.826255in}}{\pgfqpoint{3.965969in}{1.831841in}}{\pgfqpoint{3.961850in}{1.835960in}}%
\pgfpathcurveto{\pgfqpoint{3.957732in}{1.840078in}}{\pgfqpoint{3.952146in}{1.842392in}}{\pgfqpoint{3.946322in}{1.842392in}}%
\pgfpathcurveto{\pgfqpoint{3.940498in}{1.842392in}}{\pgfqpoint{3.934912in}{1.840078in}}{\pgfqpoint{3.930794in}{1.835960in}}%
\pgfpathcurveto{\pgfqpoint{3.926676in}{1.831841in}}{\pgfqpoint{3.924362in}{1.826255in}}{\pgfqpoint{3.924362in}{1.820431in}}%
\pgfpathcurveto{\pgfqpoint{3.924362in}{1.814607in}}{\pgfqpoint{3.926676in}{1.809021in}}{\pgfqpoint{3.930794in}{1.804903in}}%
\pgfpathcurveto{\pgfqpoint{3.934912in}{1.800785in}}{\pgfqpoint{3.940498in}{1.798471in}}{\pgfqpoint{3.946322in}{1.798471in}}%
\pgfpathlineto{\pgfqpoint{3.946322in}{1.798471in}}%
\pgfpathclose%
\pgfusepath{stroke,fill}%
\end{pgfscope}%
\begin{pgfscope}%
\pgfpathrectangle{\pgfqpoint{1.414736in}{0.548769in}}{\pgfqpoint{3.497747in}{7.253186in}}%
\pgfusepath{clip}%
\pgfsetbuttcap%
\pgfsetroundjoin%
\definecolor{currentfill}{rgb}{0.993248,0.906157,0.143936}%
\pgfsetfillcolor{currentfill}%
\pgfsetlinewidth{1.003750pt}%
\definecolor{currentstroke}{rgb}{0.993248,0.906157,0.143936}%
\pgfsetstrokecolor{currentstroke}%
\pgfsetdash{}{0pt}%
\pgfpathmoveto{\pgfqpoint{3.705821in}{2.269457in}}%
\pgfpathcurveto{\pgfqpoint{3.711645in}{2.269457in}}{\pgfqpoint{3.717231in}{2.271771in}}{\pgfqpoint{3.721349in}{2.275889in}}%
\pgfpathcurveto{\pgfqpoint{3.725468in}{2.280007in}}{\pgfqpoint{3.727781in}{2.285593in}}{\pgfqpoint{3.727781in}{2.291417in}}%
\pgfpathcurveto{\pgfqpoint{3.727781in}{2.297241in}}{\pgfqpoint{3.725468in}{2.302828in}}{\pgfqpoint{3.721349in}{2.306946in}}%
\pgfpathcurveto{\pgfqpoint{3.717231in}{2.311064in}}{\pgfqpoint{3.711645in}{2.313378in}}{\pgfqpoint{3.705821in}{2.313378in}}%
\pgfpathcurveto{\pgfqpoint{3.699997in}{2.313378in}}{\pgfqpoint{3.694411in}{2.311064in}}{\pgfqpoint{3.690293in}{2.306946in}}%
\pgfpathcurveto{\pgfqpoint{3.686175in}{2.302828in}}{\pgfqpoint{3.683861in}{2.297241in}}{\pgfqpoint{3.683861in}{2.291417in}}%
\pgfpathcurveto{\pgfqpoint{3.683861in}{2.285593in}}{\pgfqpoint{3.686175in}{2.280007in}}{\pgfqpoint{3.690293in}{2.275889in}}%
\pgfpathcurveto{\pgfqpoint{3.694411in}{2.271771in}}{\pgfqpoint{3.699997in}{2.269457in}}{\pgfqpoint{3.705821in}{2.269457in}}%
\pgfpathlineto{\pgfqpoint{3.705821in}{2.269457in}}%
\pgfpathclose%
\pgfusepath{stroke,fill}%
\end{pgfscope}%
\begin{pgfscope}%
\pgfpathrectangle{\pgfqpoint{1.414736in}{0.548769in}}{\pgfqpoint{3.497747in}{7.253186in}}%
\pgfusepath{clip}%
\pgfsetbuttcap%
\pgfsetroundjoin%
\definecolor{currentfill}{rgb}{0.720391,0.870350,0.162603}%
\pgfsetfillcolor{currentfill}%
\pgfsetlinewidth{1.003750pt}%
\definecolor{currentstroke}{rgb}{0.720391,0.870350,0.162603}%
\pgfsetstrokecolor{currentstroke}%
\pgfsetdash{}{0pt}%
\pgfpathmoveto{\pgfqpoint{3.223583in}{2.740443in}}%
\pgfpathcurveto{\pgfqpoint{3.229407in}{2.740443in}}{\pgfqpoint{3.234993in}{2.742757in}}{\pgfqpoint{3.239112in}{2.746875in}}%
\pgfpathcurveto{\pgfqpoint{3.243230in}{2.750993in}}{\pgfqpoint{3.245544in}{2.756580in}}{\pgfqpoint{3.245544in}{2.762404in}}%
\pgfpathcurveto{\pgfqpoint{3.245544in}{2.768227in}}{\pgfqpoint{3.243230in}{2.773814in}}{\pgfqpoint{3.239112in}{2.777932in}}%
\pgfpathcurveto{\pgfqpoint{3.234993in}{2.782050in}}{\pgfqpoint{3.229407in}{2.784364in}}{\pgfqpoint{3.223583in}{2.784364in}}%
\pgfpathcurveto{\pgfqpoint{3.217759in}{2.784364in}}{\pgfqpoint{3.212173in}{2.782050in}}{\pgfqpoint{3.208055in}{2.777932in}}%
\pgfpathcurveto{\pgfqpoint{3.203937in}{2.773814in}}{\pgfqpoint{3.201623in}{2.768227in}}{\pgfqpoint{3.201623in}{2.762404in}}%
\pgfpathcurveto{\pgfqpoint{3.201623in}{2.756580in}}{\pgfqpoint{3.203937in}{2.750993in}}{\pgfqpoint{3.208055in}{2.746875in}}%
\pgfpathcurveto{\pgfqpoint{3.212173in}{2.742757in}}{\pgfqpoint{3.217759in}{2.740443in}}{\pgfqpoint{3.223583in}{2.740443in}}%
\pgfpathlineto{\pgfqpoint{3.223583in}{2.740443in}}%
\pgfpathclose%
\pgfusepath{stroke,fill}%
\end{pgfscope}%
\begin{pgfscope}%
\pgfpathrectangle{\pgfqpoint{1.414736in}{0.548769in}}{\pgfqpoint{3.497747in}{7.253186in}}%
\pgfusepath{clip}%
\pgfsetbuttcap%
\pgfsetroundjoin%
\definecolor{currentfill}{rgb}{0.311925,0.767822,0.415586}%
\pgfsetfillcolor{currentfill}%
\pgfsetlinewidth{1.003750pt}%
\definecolor{currentstroke}{rgb}{0.311925,0.767822,0.415586}%
\pgfsetstrokecolor{currentstroke}%
\pgfsetdash{}{0pt}%
\pgfpathmoveto{\pgfqpoint{4.285735in}{3.211429in}}%
\pgfpathcurveto{\pgfqpoint{4.291559in}{3.211429in}}{\pgfqpoint{4.297145in}{3.213743in}}{\pgfqpoint{4.301263in}{3.217861in}}%
\pgfpathcurveto{\pgfqpoint{4.305381in}{3.221980in}}{\pgfqpoint{4.307695in}{3.227566in}}{\pgfqpoint{4.307695in}{3.233390in}}%
\pgfpathcurveto{\pgfqpoint{4.307695in}{3.239214in}}{\pgfqpoint{4.305381in}{3.244800in}}{\pgfqpoint{4.301263in}{3.248918in}}%
\pgfpathcurveto{\pgfqpoint{4.297145in}{3.253036in}}{\pgfqpoint{4.291559in}{3.255350in}}{\pgfqpoint{4.285735in}{3.255350in}}%
\pgfpathcurveto{\pgfqpoint{4.279911in}{3.255350in}}{\pgfqpoint{4.274325in}{3.253036in}}{\pgfqpoint{4.270206in}{3.248918in}}%
\pgfpathcurveto{\pgfqpoint{4.266088in}{3.244800in}}{\pgfqpoint{4.263774in}{3.239214in}}{\pgfqpoint{4.263774in}{3.233390in}}%
\pgfpathcurveto{\pgfqpoint{4.263774in}{3.227566in}}{\pgfqpoint{4.266088in}{3.221980in}}{\pgfqpoint{4.270206in}{3.217861in}}%
\pgfpathcurveto{\pgfqpoint{4.274325in}{3.213743in}}{\pgfqpoint{4.279911in}{3.211429in}}{\pgfqpoint{4.285735in}{3.211429in}}%
\pgfpathlineto{\pgfqpoint{4.285735in}{3.211429in}}%
\pgfpathclose%
\pgfusepath{stroke,fill}%
\end{pgfscope}%
\begin{pgfscope}%
\pgfpathrectangle{\pgfqpoint{1.414736in}{0.548769in}}{\pgfqpoint{3.497747in}{7.253186in}}%
\pgfusepath{clip}%
\pgfsetbuttcap%
\pgfsetroundjoin%
\definecolor{currentfill}{rgb}{0.993248,0.906157,0.143936}%
\pgfsetfillcolor{currentfill}%
\pgfsetlinewidth{1.003750pt}%
\definecolor{currentstroke}{rgb}{0.993248,0.906157,0.143936}%
\pgfsetstrokecolor{currentstroke}%
\pgfsetdash{}{0pt}%
\pgfpathmoveto{\pgfqpoint{3.218625in}{3.682415in}}%
\pgfpathcurveto{\pgfqpoint{3.224449in}{3.682415in}}{\pgfqpoint{3.230035in}{3.684729in}}{\pgfqpoint{3.234153in}{3.688847in}}%
\pgfpathcurveto{\pgfqpoint{3.238272in}{3.692966in}}{\pgfqpoint{3.240585in}{3.698552in}}{\pgfqpoint{3.240585in}{3.704376in}}%
\pgfpathcurveto{\pgfqpoint{3.240585in}{3.710200in}}{\pgfqpoint{3.238272in}{3.715786in}}{\pgfqpoint{3.234153in}{3.719904in}}%
\pgfpathcurveto{\pgfqpoint{3.230035in}{3.724022in}}{\pgfqpoint{3.224449in}{3.726336in}}{\pgfqpoint{3.218625in}{3.726336in}}%
\pgfpathcurveto{\pgfqpoint{3.212801in}{3.726336in}}{\pgfqpoint{3.207215in}{3.724022in}}{\pgfqpoint{3.203097in}{3.719904in}}%
\pgfpathcurveto{\pgfqpoint{3.198979in}{3.715786in}}{\pgfqpoint{3.196665in}{3.710200in}}{\pgfqpoint{3.196665in}{3.704376in}}%
\pgfpathcurveto{\pgfqpoint{3.196665in}{3.698552in}}{\pgfqpoint{3.198979in}{3.692966in}}{\pgfqpoint{3.203097in}{3.688847in}}%
\pgfpathcurveto{\pgfqpoint{3.207215in}{3.684729in}}{\pgfqpoint{3.212801in}{3.682415in}}{\pgfqpoint{3.218625in}{3.682415in}}%
\pgfpathlineto{\pgfqpoint{3.218625in}{3.682415in}}%
\pgfpathclose%
\pgfusepath{stroke,fill}%
\end{pgfscope}%
\begin{pgfscope}%
\pgfpathrectangle{\pgfqpoint{1.414736in}{0.548769in}}{\pgfqpoint{3.497747in}{7.253186in}}%
\pgfusepath{clip}%
\pgfsetbuttcap%
\pgfsetroundjoin%
\definecolor{currentfill}{rgb}{0.974417,0.903590,0.130215}%
\pgfsetfillcolor{currentfill}%
\pgfsetlinewidth{1.003750pt}%
\definecolor{currentstroke}{rgb}{0.974417,0.903590,0.130215}%
\pgfsetstrokecolor{currentstroke}%
\pgfsetdash{}{0pt}%
\pgfpathmoveto{\pgfqpoint{3.492594in}{4.153402in}}%
\pgfpathcurveto{\pgfqpoint{3.498418in}{4.153402in}}{\pgfqpoint{3.504004in}{4.155715in}}{\pgfqpoint{3.508122in}{4.159834in}}%
\pgfpathcurveto{\pgfqpoint{3.512240in}{4.163952in}}{\pgfqpoint{3.514554in}{4.169538in}}{\pgfqpoint{3.514554in}{4.175362in}}%
\pgfpathcurveto{\pgfqpoint{3.514554in}{4.181186in}}{\pgfqpoint{3.512240in}{4.186772in}}{\pgfqpoint{3.508122in}{4.190890in}}%
\pgfpathcurveto{\pgfqpoint{3.504004in}{4.195008in}}{\pgfqpoint{3.498418in}{4.197322in}}{\pgfqpoint{3.492594in}{4.197322in}}%
\pgfpathcurveto{\pgfqpoint{3.486770in}{4.197322in}}{\pgfqpoint{3.481184in}{4.195008in}}{\pgfqpoint{3.477066in}{4.190890in}}%
\pgfpathcurveto{\pgfqpoint{3.472947in}{4.186772in}}{\pgfqpoint{3.470633in}{4.181186in}}{\pgfqpoint{3.470633in}{4.175362in}}%
\pgfpathcurveto{\pgfqpoint{3.470633in}{4.169538in}}{\pgfqpoint{3.472947in}{4.163952in}}{\pgfqpoint{3.477066in}{4.159834in}}%
\pgfpathcurveto{\pgfqpoint{3.481184in}{4.155715in}}{\pgfqpoint{3.486770in}{4.153402in}}{\pgfqpoint{3.492594in}{4.153402in}}%
\pgfpathlineto{\pgfqpoint{3.492594in}{4.153402in}}%
\pgfpathclose%
\pgfusepath{stroke,fill}%
\end{pgfscope}%
\begin{pgfscope}%
\pgfpathrectangle{\pgfqpoint{1.414736in}{0.548769in}}{\pgfqpoint{3.497747in}{7.253186in}}%
\pgfusepath{clip}%
\pgfsetbuttcap%
\pgfsetroundjoin%
\definecolor{currentfill}{rgb}{0.804182,0.882046,0.114965}%
\pgfsetfillcolor{currentfill}%
\pgfsetlinewidth{1.003750pt}%
\definecolor{currentstroke}{rgb}{0.804182,0.882046,0.114965}%
\pgfsetstrokecolor{currentstroke}%
\pgfsetdash{}{0pt}%
\pgfpathmoveto{\pgfqpoint{3.443732in}{4.624388in}}%
\pgfpathcurveto{\pgfqpoint{3.449556in}{4.624388in}}{\pgfqpoint{3.455142in}{4.626702in}}{\pgfqpoint{3.459260in}{4.630820in}}%
\pgfpathcurveto{\pgfqpoint{3.463379in}{4.634938in}}{\pgfqpoint{3.465692in}{4.640524in}}{\pgfqpoint{3.465692in}{4.646348in}}%
\pgfpathcurveto{\pgfqpoint{3.465692in}{4.652172in}}{\pgfqpoint{3.463379in}{4.657758in}}{\pgfqpoint{3.459260in}{4.661876in}}%
\pgfpathcurveto{\pgfqpoint{3.455142in}{4.665994in}}{\pgfqpoint{3.449556in}{4.668308in}}{\pgfqpoint{3.443732in}{4.668308in}}%
\pgfpathcurveto{\pgfqpoint{3.437908in}{4.668308in}}{\pgfqpoint{3.432322in}{4.665994in}}{\pgfqpoint{3.428204in}{4.661876in}}%
\pgfpathcurveto{\pgfqpoint{3.424086in}{4.657758in}}{\pgfqpoint{3.421772in}{4.652172in}}{\pgfqpoint{3.421772in}{4.646348in}}%
\pgfpathcurveto{\pgfqpoint{3.421772in}{4.640524in}}{\pgfqpoint{3.424086in}{4.634938in}}{\pgfqpoint{3.428204in}{4.630820in}}%
\pgfpathcurveto{\pgfqpoint{3.432322in}{4.626702in}}{\pgfqpoint{3.437908in}{4.624388in}}{\pgfqpoint{3.443732in}{4.624388in}}%
\pgfpathlineto{\pgfqpoint{3.443732in}{4.624388in}}%
\pgfpathclose%
\pgfusepath{stroke,fill}%
\end{pgfscope}%
\begin{pgfscope}%
\pgfpathrectangle{\pgfqpoint{1.414736in}{0.548769in}}{\pgfqpoint{3.497747in}{7.253186in}}%
\pgfusepath{clip}%
\pgfsetbuttcap%
\pgfsetroundjoin%
\definecolor{currentfill}{rgb}{0.122312,0.633153,0.530398}%
\pgfsetfillcolor{currentfill}%
\pgfsetlinewidth{1.003750pt}%
\definecolor{currentstroke}{rgb}{0.122312,0.633153,0.530398}%
\pgfsetstrokecolor{currentstroke}%
\pgfsetdash{}{0pt}%
\pgfpathmoveto{\pgfqpoint{3.608250in}{5.095374in}}%
\pgfpathcurveto{\pgfqpoint{3.614074in}{5.095374in}}{\pgfqpoint{3.619660in}{5.097688in}}{\pgfqpoint{3.623778in}{5.101806in}}%
\pgfpathcurveto{\pgfqpoint{3.627896in}{5.105924in}}{\pgfqpoint{3.630210in}{5.111510in}}{\pgfqpoint{3.630210in}{5.117334in}}%
\pgfpathcurveto{\pgfqpoint{3.630210in}{5.123158in}}{\pgfqpoint{3.627896in}{5.128744in}}{\pgfqpoint{3.623778in}{5.132862in}}%
\pgfpathcurveto{\pgfqpoint{3.619660in}{5.136980in}}{\pgfqpoint{3.614074in}{5.139294in}}{\pgfqpoint{3.608250in}{5.139294in}}%
\pgfpathcurveto{\pgfqpoint{3.602426in}{5.139294in}}{\pgfqpoint{3.596840in}{5.136980in}}{\pgfqpoint{3.592721in}{5.132862in}}%
\pgfpathcurveto{\pgfqpoint{3.588603in}{5.128744in}}{\pgfqpoint{3.586289in}{5.123158in}}{\pgfqpoint{3.586289in}{5.117334in}}%
\pgfpathcurveto{\pgfqpoint{3.586289in}{5.111510in}}{\pgfqpoint{3.588603in}{5.105924in}}{\pgfqpoint{3.592721in}{5.101806in}}%
\pgfpathcurveto{\pgfqpoint{3.596840in}{5.097688in}}{\pgfqpoint{3.602426in}{5.095374in}}{\pgfqpoint{3.608250in}{5.095374in}}%
\pgfpathlineto{\pgfqpoint{3.608250in}{5.095374in}}%
\pgfpathclose%
\pgfusepath{stroke,fill}%
\end{pgfscope}%
\begin{pgfscope}%
\pgfpathrectangle{\pgfqpoint{1.414736in}{0.548769in}}{\pgfqpoint{3.497747in}{7.253186in}}%
\pgfusepath{clip}%
\pgfsetbuttcap%
\pgfsetroundjoin%
\definecolor{currentfill}{rgb}{0.135066,0.544853,0.554029}%
\pgfsetfillcolor{currentfill}%
\pgfsetlinewidth{1.003750pt}%
\definecolor{currentstroke}{rgb}{0.135066,0.544853,0.554029}%
\pgfsetstrokecolor{currentstroke}%
\pgfsetdash{}{0pt}%
\pgfpathmoveto{\pgfqpoint{3.556117in}{5.566360in}}%
\pgfpathcurveto{\pgfqpoint{3.561941in}{5.566360in}}{\pgfqpoint{3.567528in}{5.568674in}}{\pgfqpoint{3.571646in}{5.572792in}}%
\pgfpathcurveto{\pgfqpoint{3.575764in}{5.576910in}}{\pgfqpoint{3.578078in}{5.582496in}}{\pgfqpoint{3.578078in}{5.588320in}}%
\pgfpathcurveto{\pgfqpoint{3.578078in}{5.594144in}}{\pgfqpoint{3.575764in}{5.599730in}}{\pgfqpoint{3.571646in}{5.603848in}}%
\pgfpathcurveto{\pgfqpoint{3.567528in}{5.607967in}}{\pgfqpoint{3.561941in}{5.610280in}}{\pgfqpoint{3.556117in}{5.610280in}}%
\pgfpathcurveto{\pgfqpoint{3.550294in}{5.610280in}}{\pgfqpoint{3.544707in}{5.607967in}}{\pgfqpoint{3.540589in}{5.603848in}}%
\pgfpathcurveto{\pgfqpoint{3.536471in}{5.599730in}}{\pgfqpoint{3.534157in}{5.594144in}}{\pgfqpoint{3.534157in}{5.588320in}}%
\pgfpathcurveto{\pgfqpoint{3.534157in}{5.582496in}}{\pgfqpoint{3.536471in}{5.576910in}}{\pgfqpoint{3.540589in}{5.572792in}}%
\pgfpathcurveto{\pgfqpoint{3.544707in}{5.568674in}}{\pgfqpoint{3.550294in}{5.566360in}}{\pgfqpoint{3.556117in}{5.566360in}}%
\pgfpathlineto{\pgfqpoint{3.556117in}{5.566360in}}%
\pgfpathclose%
\pgfusepath{stroke,fill}%
\end{pgfscope}%
\begin{pgfscope}%
\pgfpathrectangle{\pgfqpoint{1.414736in}{0.548769in}}{\pgfqpoint{3.497747in}{7.253186in}}%
\pgfusepath{clip}%
\pgfsetbuttcap%
\pgfsetroundjoin%
\definecolor{currentfill}{rgb}{0.762373,0.876424,0.137064}%
\pgfsetfillcolor{currentfill}%
\pgfsetlinewidth{1.003750pt}%
\definecolor{currentstroke}{rgb}{0.762373,0.876424,0.137064}%
\pgfsetstrokecolor{currentstroke}%
\pgfsetdash{}{0pt}%
\pgfpathmoveto{\pgfqpoint{3.009790in}{6.037346in}}%
\pgfpathcurveto{\pgfqpoint{3.015614in}{6.037346in}}{\pgfqpoint{3.021200in}{6.039660in}}{\pgfqpoint{3.025318in}{6.043778in}}%
\pgfpathcurveto{\pgfqpoint{3.029436in}{6.047896in}}{\pgfqpoint{3.031750in}{6.053482in}}{\pgfqpoint{3.031750in}{6.059306in}}%
\pgfpathcurveto{\pgfqpoint{3.031750in}{6.065130in}}{\pgfqpoint{3.029436in}{6.070716in}}{\pgfqpoint{3.025318in}{6.074835in}}%
\pgfpathcurveto{\pgfqpoint{3.021200in}{6.078953in}}{\pgfqpoint{3.015614in}{6.081267in}}{\pgfqpoint{3.009790in}{6.081267in}}%
\pgfpathcurveto{\pgfqpoint{3.003966in}{6.081267in}}{\pgfqpoint{2.998380in}{6.078953in}}{\pgfqpoint{2.994262in}{6.074835in}}%
\pgfpathcurveto{\pgfqpoint{2.990143in}{6.070716in}}{\pgfqpoint{2.987830in}{6.065130in}}{\pgfqpoint{2.987830in}{6.059306in}}%
\pgfpathcurveto{\pgfqpoint{2.987830in}{6.053482in}}{\pgfqpoint{2.990143in}{6.047896in}}{\pgfqpoint{2.994262in}{6.043778in}}%
\pgfpathcurveto{\pgfqpoint{2.998380in}{6.039660in}}{\pgfqpoint{3.003966in}{6.037346in}}{\pgfqpoint{3.009790in}{6.037346in}}%
\pgfpathlineto{\pgfqpoint{3.009790in}{6.037346in}}%
\pgfpathclose%
\pgfusepath{stroke,fill}%
\end{pgfscope}%
\begin{pgfscope}%
\pgfpathrectangle{\pgfqpoint{1.414736in}{0.548769in}}{\pgfqpoint{3.497747in}{7.253186in}}%
\pgfusepath{clip}%
\pgfsetbuttcap%
\pgfsetroundjoin%
\definecolor{currentfill}{rgb}{0.993248,0.906157,0.143936}%
\pgfsetfillcolor{currentfill}%
\pgfsetlinewidth{1.003750pt}%
\definecolor{currentstroke}{rgb}{0.993248,0.906157,0.143936}%
\pgfsetstrokecolor{currentstroke}%
\pgfsetdash{}{0pt}%
\pgfpathmoveto{\pgfqpoint{3.114733in}{6.508332in}}%
\pgfpathcurveto{\pgfqpoint{3.120557in}{6.508332in}}{\pgfqpoint{3.126143in}{6.510646in}}{\pgfqpoint{3.130261in}{6.514764in}}%
\pgfpathcurveto{\pgfqpoint{3.134379in}{6.518882in}}{\pgfqpoint{3.136693in}{6.524468in}}{\pgfqpoint{3.136693in}{6.530292in}}%
\pgfpathcurveto{\pgfqpoint{3.136693in}{6.536116in}}{\pgfqpoint{3.134379in}{6.541702in}}{\pgfqpoint{3.130261in}{6.545821in}}%
\pgfpathcurveto{\pgfqpoint{3.126143in}{6.549939in}}{\pgfqpoint{3.120557in}{6.552253in}}{\pgfqpoint{3.114733in}{6.552253in}}%
\pgfpathcurveto{\pgfqpoint{3.108909in}{6.552253in}}{\pgfqpoint{3.103323in}{6.549939in}}{\pgfqpoint{3.099205in}{6.545821in}}%
\pgfpathcurveto{\pgfqpoint{3.095086in}{6.541702in}}{\pgfqpoint{3.092773in}{6.536116in}}{\pgfqpoint{3.092773in}{6.530292in}}%
\pgfpathcurveto{\pgfqpoint{3.092773in}{6.524468in}}{\pgfqpoint{3.095086in}{6.518882in}}{\pgfqpoint{3.099205in}{6.514764in}}%
\pgfpathcurveto{\pgfqpoint{3.103323in}{6.510646in}}{\pgfqpoint{3.108909in}{6.508332in}}{\pgfqpoint{3.114733in}{6.508332in}}%
\pgfpathlineto{\pgfqpoint{3.114733in}{6.508332in}}%
\pgfpathclose%
\pgfusepath{stroke,fill}%
\end{pgfscope}%
\begin{pgfscope}%
\pgfpathrectangle{\pgfqpoint{1.414736in}{0.548769in}}{\pgfqpoint{3.497747in}{7.253186in}}%
\pgfusepath{clip}%
\pgfsetbuttcap%
\pgfsetroundjoin%
\definecolor{currentfill}{rgb}{0.804182,0.882046,0.114965}%
\pgfsetfillcolor{currentfill}%
\pgfsetlinewidth{1.003750pt}%
\definecolor{currentstroke}{rgb}{0.804182,0.882046,0.114965}%
\pgfsetstrokecolor{currentstroke}%
\pgfsetdash{}{0pt}%
\pgfpathmoveto{\pgfqpoint{3.773479in}{6.979318in}}%
\pgfpathcurveto{\pgfqpoint{3.779303in}{6.979318in}}{\pgfqpoint{3.784889in}{6.981632in}}{\pgfqpoint{3.789007in}{6.985750in}}%
\pgfpathcurveto{\pgfqpoint{3.793126in}{6.989868in}}{\pgfqpoint{3.795439in}{6.995455in}}{\pgfqpoint{3.795439in}{7.001278in}}%
\pgfpathcurveto{\pgfqpoint{3.795439in}{7.007102in}}{\pgfqpoint{3.793126in}{7.012689in}}{\pgfqpoint{3.789007in}{7.016807in}}%
\pgfpathcurveto{\pgfqpoint{3.784889in}{7.020925in}}{\pgfqpoint{3.779303in}{7.023239in}}{\pgfqpoint{3.773479in}{7.023239in}}%
\pgfpathcurveto{\pgfqpoint{3.767655in}{7.023239in}}{\pgfqpoint{3.762069in}{7.020925in}}{\pgfqpoint{3.757951in}{7.016807in}}%
\pgfpathcurveto{\pgfqpoint{3.753833in}{7.012689in}}{\pgfqpoint{3.751519in}{7.007102in}}{\pgfqpoint{3.751519in}{7.001278in}}%
\pgfpathcurveto{\pgfqpoint{3.751519in}{6.995455in}}{\pgfqpoint{3.753833in}{6.989868in}}{\pgfqpoint{3.757951in}{6.985750in}}%
\pgfpathcurveto{\pgfqpoint{3.762069in}{6.981632in}}{\pgfqpoint{3.767655in}{6.979318in}}{\pgfqpoint{3.773479in}{6.979318in}}%
\pgfpathlineto{\pgfqpoint{3.773479in}{6.979318in}}%
\pgfpathclose%
\pgfusepath{stroke,fill}%
\end{pgfscope}%
\begin{pgfscope}%
\pgfpathrectangle{\pgfqpoint{1.414736in}{0.548769in}}{\pgfqpoint{3.497747in}{7.253186in}}%
\pgfusepath{clip}%
\pgfsetbuttcap%
\pgfsetroundjoin%
\definecolor{currentfill}{rgb}{0.458674,0.816363,0.329727}%
\pgfsetfillcolor{currentfill}%
\pgfsetlinewidth{1.003750pt}%
\definecolor{currentstroke}{rgb}{0.458674,0.816363,0.329727}%
\pgfsetstrokecolor{currentstroke}%
\pgfsetdash{}{0pt}%
\pgfpathmoveto{\pgfqpoint{3.680436in}{7.450304in}}%
\pgfpathcurveto{\pgfqpoint{3.686260in}{7.450304in}}{\pgfqpoint{3.691846in}{7.452618in}}{\pgfqpoint{3.695964in}{7.456736in}}%
\pgfpathcurveto{\pgfqpoint{3.700082in}{7.460854in}}{\pgfqpoint{3.702396in}{7.466441in}}{\pgfqpoint{3.702396in}{7.472265in}}%
\pgfpathcurveto{\pgfqpoint{3.702396in}{7.478089in}}{\pgfqpoint{3.700082in}{7.483675in}}{\pgfqpoint{3.695964in}{7.487793in}}%
\pgfpathcurveto{\pgfqpoint{3.691846in}{7.491911in}}{\pgfqpoint{3.686260in}{7.494225in}}{\pgfqpoint{3.680436in}{7.494225in}}%
\pgfpathcurveto{\pgfqpoint{3.674612in}{7.494225in}}{\pgfqpoint{3.669026in}{7.491911in}}{\pgfqpoint{3.664908in}{7.487793in}}%
\pgfpathcurveto{\pgfqpoint{3.660790in}{7.483675in}}{\pgfqpoint{3.658476in}{7.478089in}}{\pgfqpoint{3.658476in}{7.472265in}}%
\pgfpathcurveto{\pgfqpoint{3.658476in}{7.466441in}}{\pgfqpoint{3.660790in}{7.460854in}}{\pgfqpoint{3.664908in}{7.456736in}}%
\pgfpathcurveto{\pgfqpoint{3.669026in}{7.452618in}}{\pgfqpoint{3.674612in}{7.450304in}}{\pgfqpoint{3.680436in}{7.450304in}}%
\pgfpathlineto{\pgfqpoint{3.680436in}{7.450304in}}%
\pgfpathclose%
\pgfusepath{stroke,fill}%
\end{pgfscope}%
\begin{pgfscope}%
\pgfpathrectangle{\pgfqpoint{1.414736in}{0.548769in}}{\pgfqpoint{3.497747in}{7.253186in}}%
\pgfusepath{clip}%
\pgfsetbuttcap%
\pgfsetroundjoin%
\definecolor{currentfill}{rgb}{0.955300,0.901065,0.118128}%
\pgfsetfillcolor{currentfill}%
\pgfsetlinewidth{1.003750pt}%
\definecolor{currentstroke}{rgb}{0.955300,0.901065,0.118128}%
\pgfsetstrokecolor{currentstroke}%
\pgfsetdash{}{0pt}%
\pgfpathmoveto{\pgfqpoint{3.440712in}{0.856499in}}%
\pgfpathcurveto{\pgfqpoint{3.446536in}{0.856499in}}{\pgfqpoint{3.452122in}{0.858813in}}{\pgfqpoint{3.456240in}{0.862931in}}%
\pgfpathcurveto{\pgfqpoint{3.460359in}{0.867049in}}{\pgfqpoint{3.462672in}{0.872635in}}{\pgfqpoint{3.462672in}{0.878459in}}%
\pgfpathcurveto{\pgfqpoint{3.462672in}{0.884283in}}{\pgfqpoint{3.460359in}{0.889869in}}{\pgfqpoint{3.456240in}{0.893987in}}%
\pgfpathcurveto{\pgfqpoint{3.452122in}{0.898105in}}{\pgfqpoint{3.446536in}{0.900419in}}{\pgfqpoint{3.440712in}{0.900419in}}%
\pgfpathcurveto{\pgfqpoint{3.434888in}{0.900419in}}{\pgfqpoint{3.429302in}{0.898105in}}{\pgfqpoint{3.425184in}{0.893987in}}%
\pgfpathcurveto{\pgfqpoint{3.421066in}{0.889869in}}{\pgfqpoint{3.418752in}{0.884283in}}{\pgfqpoint{3.418752in}{0.878459in}}%
\pgfpathcurveto{\pgfqpoint{3.418752in}{0.872635in}}{\pgfqpoint{3.421066in}{0.867049in}}{\pgfqpoint{3.425184in}{0.862931in}}%
\pgfpathcurveto{\pgfqpoint{3.429302in}{0.858813in}}{\pgfqpoint{3.434888in}{0.856499in}}{\pgfqpoint{3.440712in}{0.856499in}}%
\pgfpathlineto{\pgfqpoint{3.440712in}{0.856499in}}%
\pgfpathclose%
\pgfusepath{stroke,fill}%
\end{pgfscope}%
\begin{pgfscope}%
\pgfpathrectangle{\pgfqpoint{1.414736in}{0.548769in}}{\pgfqpoint{3.497747in}{7.253186in}}%
\pgfusepath{clip}%
\pgfsetbuttcap%
\pgfsetroundjoin%
\definecolor{currentfill}{rgb}{0.327796,0.773980,0.406640}%
\pgfsetfillcolor{currentfill}%
\pgfsetlinewidth{1.003750pt}%
\definecolor{currentstroke}{rgb}{0.327796,0.773980,0.406640}%
\pgfsetstrokecolor{currentstroke}%
\pgfsetdash{}{0pt}%
\pgfpathmoveto{\pgfqpoint{4.046051in}{1.327485in}}%
\pgfpathcurveto{\pgfqpoint{4.051874in}{1.327485in}}{\pgfqpoint{4.057461in}{1.329799in}}{\pgfqpoint{4.061579in}{1.333917in}}%
\pgfpathcurveto{\pgfqpoint{4.065697in}{1.338035in}}{\pgfqpoint{4.068011in}{1.343621in}}{\pgfqpoint{4.068011in}{1.349445in}}%
\pgfpathcurveto{\pgfqpoint{4.068011in}{1.355269in}}{\pgfqpoint{4.065697in}{1.360855in}}{\pgfqpoint{4.061579in}{1.364973in}}%
\pgfpathcurveto{\pgfqpoint{4.057461in}{1.369092in}}{\pgfqpoint{4.051874in}{1.371405in}}{\pgfqpoint{4.046051in}{1.371405in}}%
\pgfpathcurveto{\pgfqpoint{4.040227in}{1.371405in}}{\pgfqpoint{4.034640in}{1.369092in}}{\pgfqpoint{4.030522in}{1.364973in}}%
\pgfpathcurveto{\pgfqpoint{4.026404in}{1.360855in}}{\pgfqpoint{4.024090in}{1.355269in}}{\pgfqpoint{4.024090in}{1.349445in}}%
\pgfpathcurveto{\pgfqpoint{4.024090in}{1.343621in}}{\pgfqpoint{4.026404in}{1.338035in}}{\pgfqpoint{4.030522in}{1.333917in}}%
\pgfpathcurveto{\pgfqpoint{4.034640in}{1.329799in}}{\pgfqpoint{4.040227in}{1.327485in}}{\pgfqpoint{4.046051in}{1.327485in}}%
\pgfpathlineto{\pgfqpoint{4.046051in}{1.327485in}}%
\pgfpathclose%
\pgfusepath{stroke,fill}%
\end{pgfscope}%
\begin{pgfscope}%
\pgfpathrectangle{\pgfqpoint{1.414736in}{0.548769in}}{\pgfqpoint{3.497747in}{7.253186in}}%
\pgfusepath{clip}%
\pgfsetbuttcap%
\pgfsetroundjoin%
\definecolor{currentfill}{rgb}{0.430983,0.808473,0.346476}%
\pgfsetfillcolor{currentfill}%
\pgfsetlinewidth{1.003750pt}%
\definecolor{currentstroke}{rgb}{0.430983,0.808473,0.346476}%
\pgfsetstrokecolor{currentstroke}%
\pgfsetdash{}{0pt}%
\pgfpathmoveto{\pgfqpoint{3.456027in}{1.798471in}}%
\pgfpathcurveto{\pgfqpoint{3.461851in}{1.798471in}}{\pgfqpoint{3.467437in}{1.800785in}}{\pgfqpoint{3.471555in}{1.804903in}}%
\pgfpathcurveto{\pgfqpoint{3.475673in}{1.809021in}}{\pgfqpoint{3.477987in}{1.814607in}}{\pgfqpoint{3.477987in}{1.820431in}}%
\pgfpathcurveto{\pgfqpoint{3.477987in}{1.826255in}}{\pgfqpoint{3.475673in}{1.831841in}}{\pgfqpoint{3.471555in}{1.835960in}}%
\pgfpathcurveto{\pgfqpoint{3.467437in}{1.840078in}}{\pgfqpoint{3.461851in}{1.842392in}}{\pgfqpoint{3.456027in}{1.842392in}}%
\pgfpathcurveto{\pgfqpoint{3.450203in}{1.842392in}}{\pgfqpoint{3.444617in}{1.840078in}}{\pgfqpoint{3.440499in}{1.835960in}}%
\pgfpathcurveto{\pgfqpoint{3.436380in}{1.831841in}}{\pgfqpoint{3.434067in}{1.826255in}}{\pgfqpoint{3.434067in}{1.820431in}}%
\pgfpathcurveto{\pgfqpoint{3.434067in}{1.814607in}}{\pgfqpoint{3.436380in}{1.809021in}}{\pgfqpoint{3.440499in}{1.804903in}}%
\pgfpathcurveto{\pgfqpoint{3.444617in}{1.800785in}}{\pgfqpoint{3.450203in}{1.798471in}}{\pgfqpoint{3.456027in}{1.798471in}}%
\pgfpathlineto{\pgfqpoint{3.456027in}{1.798471in}}%
\pgfpathclose%
\pgfusepath{stroke,fill}%
\end{pgfscope}%
\begin{pgfscope}%
\pgfpathrectangle{\pgfqpoint{1.414736in}{0.548769in}}{\pgfqpoint{3.497747in}{7.253186in}}%
\pgfusepath{clip}%
\pgfsetbuttcap%
\pgfsetroundjoin%
\definecolor{currentfill}{rgb}{0.535621,0.835785,0.281908}%
\pgfsetfillcolor{currentfill}%
\pgfsetlinewidth{1.003750pt}%
\definecolor{currentstroke}{rgb}{0.535621,0.835785,0.281908}%
\pgfsetstrokecolor{currentstroke}%
\pgfsetdash{}{0pt}%
\pgfpathmoveto{\pgfqpoint{3.391521in}{2.269457in}}%
\pgfpathcurveto{\pgfqpoint{3.397345in}{2.269457in}}{\pgfqpoint{3.402931in}{2.271771in}}{\pgfqpoint{3.407049in}{2.275889in}}%
\pgfpathcurveto{\pgfqpoint{3.411167in}{2.280007in}}{\pgfqpoint{3.413481in}{2.285593in}}{\pgfqpoint{3.413481in}{2.291417in}}%
\pgfpathcurveto{\pgfqpoint{3.413481in}{2.297241in}}{\pgfqpoint{3.411167in}{2.302828in}}{\pgfqpoint{3.407049in}{2.306946in}}%
\pgfpathcurveto{\pgfqpoint{3.402931in}{2.311064in}}{\pgfqpoint{3.397345in}{2.313378in}}{\pgfqpoint{3.391521in}{2.313378in}}%
\pgfpathcurveto{\pgfqpoint{3.385697in}{2.313378in}}{\pgfqpoint{3.380111in}{2.311064in}}{\pgfqpoint{3.375993in}{2.306946in}}%
\pgfpathcurveto{\pgfqpoint{3.371875in}{2.302828in}}{\pgfqpoint{3.369561in}{2.297241in}}{\pgfqpoint{3.369561in}{2.291417in}}%
\pgfpathcurveto{\pgfqpoint{3.369561in}{2.285593in}}{\pgfqpoint{3.371875in}{2.280007in}}{\pgfqpoint{3.375993in}{2.275889in}}%
\pgfpathcurveto{\pgfqpoint{3.380111in}{2.271771in}}{\pgfqpoint{3.385697in}{2.269457in}}{\pgfqpoint{3.391521in}{2.269457in}}%
\pgfpathlineto{\pgfqpoint{3.391521in}{2.269457in}}%
\pgfpathclose%
\pgfusepath{stroke,fill}%
\end{pgfscope}%
\begin{pgfscope}%
\pgfpathrectangle{\pgfqpoint{1.414736in}{0.548769in}}{\pgfqpoint{3.497747in}{7.253186in}}%
\pgfusepath{clip}%
\pgfsetbuttcap%
\pgfsetroundjoin%
\definecolor{currentfill}{rgb}{0.352360,0.783011,0.392636}%
\pgfsetfillcolor{currentfill}%
\pgfsetlinewidth{1.003750pt}%
\definecolor{currentstroke}{rgb}{0.352360,0.783011,0.392636}%
\pgfsetstrokecolor{currentstroke}%
\pgfsetdash{}{0pt}%
\pgfpathmoveto{\pgfqpoint{3.752707in}{2.740443in}}%
\pgfpathcurveto{\pgfqpoint{3.758531in}{2.740443in}}{\pgfqpoint{3.764117in}{2.742757in}}{\pgfqpoint{3.768235in}{2.746875in}}%
\pgfpathcurveto{\pgfqpoint{3.772354in}{2.750993in}}{\pgfqpoint{3.774667in}{2.756580in}}{\pgfqpoint{3.774667in}{2.762404in}}%
\pgfpathcurveto{\pgfqpoint{3.774667in}{2.768227in}}{\pgfqpoint{3.772354in}{2.773814in}}{\pgfqpoint{3.768235in}{2.777932in}}%
\pgfpathcurveto{\pgfqpoint{3.764117in}{2.782050in}}{\pgfqpoint{3.758531in}{2.784364in}}{\pgfqpoint{3.752707in}{2.784364in}}%
\pgfpathcurveto{\pgfqpoint{3.746883in}{2.784364in}}{\pgfqpoint{3.741297in}{2.782050in}}{\pgfqpoint{3.737179in}{2.777932in}}%
\pgfpathcurveto{\pgfqpoint{3.733061in}{2.773814in}}{\pgfqpoint{3.730747in}{2.768227in}}{\pgfqpoint{3.730747in}{2.762404in}}%
\pgfpathcurveto{\pgfqpoint{3.730747in}{2.756580in}}{\pgfqpoint{3.733061in}{2.750993in}}{\pgfqpoint{3.737179in}{2.746875in}}%
\pgfpathcurveto{\pgfqpoint{3.741297in}{2.742757in}}{\pgfqpoint{3.746883in}{2.740443in}}{\pgfqpoint{3.752707in}{2.740443in}}%
\pgfpathlineto{\pgfqpoint{3.752707in}{2.740443in}}%
\pgfpathclose%
\pgfusepath{stroke,fill}%
\end{pgfscope}%
\begin{pgfscope}%
\pgfpathrectangle{\pgfqpoint{1.414736in}{0.548769in}}{\pgfqpoint{3.497747in}{7.253186in}}%
\pgfusepath{clip}%
\pgfsetbuttcap%
\pgfsetroundjoin%
\definecolor{currentfill}{rgb}{0.678489,0.863742,0.189503}%
\pgfsetfillcolor{currentfill}%
\pgfsetlinewidth{1.003750pt}%
\definecolor{currentstroke}{rgb}{0.678489,0.863742,0.189503}%
\pgfsetstrokecolor{currentstroke}%
\pgfsetdash{}{0pt}%
\pgfpathmoveto{\pgfqpoint{3.300481in}{3.211429in}}%
\pgfpathcurveto{\pgfqpoint{3.306304in}{3.211429in}}{\pgfqpoint{3.311891in}{3.213743in}}{\pgfqpoint{3.316009in}{3.217861in}}%
\pgfpathcurveto{\pgfqpoint{3.320127in}{3.221980in}}{\pgfqpoint{3.322441in}{3.227566in}}{\pgfqpoint{3.322441in}{3.233390in}}%
\pgfpathcurveto{\pgfqpoint{3.322441in}{3.239214in}}{\pgfqpoint{3.320127in}{3.244800in}}{\pgfqpoint{3.316009in}{3.248918in}}%
\pgfpathcurveto{\pgfqpoint{3.311891in}{3.253036in}}{\pgfqpoint{3.306304in}{3.255350in}}{\pgfqpoint{3.300481in}{3.255350in}}%
\pgfpathcurveto{\pgfqpoint{3.294657in}{3.255350in}}{\pgfqpoint{3.289070in}{3.253036in}}{\pgfqpoint{3.284952in}{3.248918in}}%
\pgfpathcurveto{\pgfqpoint{3.280834in}{3.244800in}}{\pgfqpoint{3.278520in}{3.239214in}}{\pgfqpoint{3.278520in}{3.233390in}}%
\pgfpathcurveto{\pgfqpoint{3.278520in}{3.227566in}}{\pgfqpoint{3.280834in}{3.221980in}}{\pgfqpoint{3.284952in}{3.217861in}}%
\pgfpathcurveto{\pgfqpoint{3.289070in}{3.213743in}}{\pgfqpoint{3.294657in}{3.211429in}}{\pgfqpoint{3.300481in}{3.211429in}}%
\pgfpathlineto{\pgfqpoint{3.300481in}{3.211429in}}%
\pgfpathclose%
\pgfusepath{stroke,fill}%
\end{pgfscope}%
\begin{pgfscope}%
\pgfpathrectangle{\pgfqpoint{1.414736in}{0.548769in}}{\pgfqpoint{3.497747in}{7.253186in}}%
\pgfusepath{clip}%
\pgfsetbuttcap%
\pgfsetroundjoin%
\definecolor{currentfill}{rgb}{0.896320,0.893616,0.096335}%
\pgfsetfillcolor{currentfill}%
\pgfsetlinewidth{1.003750pt}%
\definecolor{currentstroke}{rgb}{0.896320,0.893616,0.096335}%
\pgfsetstrokecolor{currentstroke}%
\pgfsetdash{}{0pt}%
\pgfpathmoveto{\pgfqpoint{3.753336in}{3.682415in}}%
\pgfpathcurveto{\pgfqpoint{3.759159in}{3.682415in}}{\pgfqpoint{3.764746in}{3.684729in}}{\pgfqpoint{3.768864in}{3.688847in}}%
\pgfpathcurveto{\pgfqpoint{3.772982in}{3.692966in}}{\pgfqpoint{3.775296in}{3.698552in}}{\pgfqpoint{3.775296in}{3.704376in}}%
\pgfpathcurveto{\pgfqpoint{3.775296in}{3.710200in}}{\pgfqpoint{3.772982in}{3.715786in}}{\pgfqpoint{3.768864in}{3.719904in}}%
\pgfpathcurveto{\pgfqpoint{3.764746in}{3.724022in}}{\pgfqpoint{3.759159in}{3.726336in}}{\pgfqpoint{3.753336in}{3.726336in}}%
\pgfpathcurveto{\pgfqpoint{3.747512in}{3.726336in}}{\pgfqpoint{3.741925in}{3.724022in}}{\pgfqpoint{3.737807in}{3.719904in}}%
\pgfpathcurveto{\pgfqpoint{3.733689in}{3.715786in}}{\pgfqpoint{3.731375in}{3.710200in}}{\pgfqpoint{3.731375in}{3.704376in}}%
\pgfpathcurveto{\pgfqpoint{3.731375in}{3.698552in}}{\pgfqpoint{3.733689in}{3.692966in}}{\pgfqpoint{3.737807in}{3.688847in}}%
\pgfpathcurveto{\pgfqpoint{3.741925in}{3.684729in}}{\pgfqpoint{3.747512in}{3.682415in}}{\pgfqpoint{3.753336in}{3.682415in}}%
\pgfpathlineto{\pgfqpoint{3.753336in}{3.682415in}}%
\pgfpathclose%
\pgfusepath{stroke,fill}%
\end{pgfscope}%
\begin{pgfscope}%
\pgfpathrectangle{\pgfqpoint{1.414736in}{0.548769in}}{\pgfqpoint{3.497747in}{7.253186in}}%
\pgfusepath{clip}%
\pgfsetbuttcap%
\pgfsetroundjoin%
\definecolor{currentfill}{rgb}{0.993248,0.906157,0.143936}%
\pgfsetfillcolor{currentfill}%
\pgfsetlinewidth{1.003750pt}%
\definecolor{currentstroke}{rgb}{0.993248,0.906157,0.143936}%
\pgfsetstrokecolor{currentstroke}%
\pgfsetdash{}{0pt}%
\pgfpathmoveto{\pgfqpoint{3.287848in}{4.153402in}}%
\pgfpathcurveto{\pgfqpoint{3.293672in}{4.153402in}}{\pgfqpoint{3.299259in}{4.155715in}}{\pgfqpoint{3.303377in}{4.159834in}}%
\pgfpathcurveto{\pgfqpoint{3.307495in}{4.163952in}}{\pgfqpoint{3.309809in}{4.169538in}}{\pgfqpoint{3.309809in}{4.175362in}}%
\pgfpathcurveto{\pgfqpoint{3.309809in}{4.181186in}}{\pgfqpoint{3.307495in}{4.186772in}}{\pgfqpoint{3.303377in}{4.190890in}}%
\pgfpathcurveto{\pgfqpoint{3.299259in}{4.195008in}}{\pgfqpoint{3.293672in}{4.197322in}}{\pgfqpoint{3.287848in}{4.197322in}}%
\pgfpathcurveto{\pgfqpoint{3.282025in}{4.197322in}}{\pgfqpoint{3.276438in}{4.195008in}}{\pgfqpoint{3.272320in}{4.190890in}}%
\pgfpathcurveto{\pgfqpoint{3.268202in}{4.186772in}}{\pgfqpoint{3.265888in}{4.181186in}}{\pgfqpoint{3.265888in}{4.175362in}}%
\pgfpathcurveto{\pgfqpoint{3.265888in}{4.169538in}}{\pgfqpoint{3.268202in}{4.163952in}}{\pgfqpoint{3.272320in}{4.159834in}}%
\pgfpathcurveto{\pgfqpoint{3.276438in}{4.155715in}}{\pgfqpoint{3.282025in}{4.153402in}}{\pgfqpoint{3.287848in}{4.153402in}}%
\pgfpathlineto{\pgfqpoint{3.287848in}{4.153402in}}%
\pgfpathclose%
\pgfusepath{stroke,fill}%
\end{pgfscope}%
\begin{pgfscope}%
\pgfpathrectangle{\pgfqpoint{1.414736in}{0.548769in}}{\pgfqpoint{3.497747in}{7.253186in}}%
\pgfusepath{clip}%
\pgfsetbuttcap%
\pgfsetroundjoin%
\definecolor{currentfill}{rgb}{0.647257,0.858400,0.209861}%
\pgfsetfillcolor{currentfill}%
\pgfsetlinewidth{1.003750pt}%
\definecolor{currentstroke}{rgb}{0.647257,0.858400,0.209861}%
\pgfsetstrokecolor{currentstroke}%
\pgfsetdash{}{0pt}%
\pgfpathmoveto{\pgfqpoint{3.109445in}{4.624388in}}%
\pgfpathcurveto{\pgfqpoint{3.115269in}{4.624388in}}{\pgfqpoint{3.120855in}{4.626702in}}{\pgfqpoint{3.124973in}{4.630820in}}%
\pgfpathcurveto{\pgfqpoint{3.129091in}{4.634938in}}{\pgfqpoint{3.131405in}{4.640524in}}{\pgfqpoint{3.131405in}{4.646348in}}%
\pgfpathcurveto{\pgfqpoint{3.131405in}{4.652172in}}{\pgfqpoint{3.129091in}{4.657758in}}{\pgfqpoint{3.124973in}{4.661876in}}%
\pgfpathcurveto{\pgfqpoint{3.120855in}{4.665994in}}{\pgfqpoint{3.115269in}{4.668308in}}{\pgfqpoint{3.109445in}{4.668308in}}%
\pgfpathcurveto{\pgfqpoint{3.103621in}{4.668308in}}{\pgfqpoint{3.098035in}{4.665994in}}{\pgfqpoint{3.093916in}{4.661876in}}%
\pgfpathcurveto{\pgfqpoint{3.089798in}{4.657758in}}{\pgfqpoint{3.087484in}{4.652172in}}{\pgfqpoint{3.087484in}{4.646348in}}%
\pgfpathcurveto{\pgfqpoint{3.087484in}{4.640524in}}{\pgfqpoint{3.089798in}{4.634938in}}{\pgfqpoint{3.093916in}{4.630820in}}%
\pgfpathcurveto{\pgfqpoint{3.098035in}{4.626702in}}{\pgfqpoint{3.103621in}{4.624388in}}{\pgfqpoint{3.109445in}{4.624388in}}%
\pgfpathlineto{\pgfqpoint{3.109445in}{4.624388in}}%
\pgfpathclose%
\pgfusepath{stroke,fill}%
\end{pgfscope}%
\begin{pgfscope}%
\pgfpathrectangle{\pgfqpoint{1.414736in}{0.548769in}}{\pgfqpoint{3.497747in}{7.253186in}}%
\pgfusepath{clip}%
\pgfsetbuttcap%
\pgfsetroundjoin%
\definecolor{currentfill}{rgb}{0.267004,0.004874,0.329415}%
\pgfsetfillcolor{currentfill}%
\pgfsetlinewidth{1.003750pt}%
\definecolor{currentstroke}{rgb}{0.267004,0.004874,0.329415}%
\pgfsetstrokecolor{currentstroke}%
\pgfsetdash{}{0pt}%
\pgfpathmoveto{\pgfqpoint{3.267932in}{5.095374in}}%
\pgfpathcurveto{\pgfqpoint{3.273756in}{5.095374in}}{\pgfqpoint{3.279342in}{5.097688in}}{\pgfqpoint{3.283460in}{5.101806in}}%
\pgfpathcurveto{\pgfqpoint{3.287578in}{5.105924in}}{\pgfqpoint{3.289892in}{5.111510in}}{\pgfqpoint{3.289892in}{5.117334in}}%
\pgfpathcurveto{\pgfqpoint{3.289892in}{5.123158in}}{\pgfqpoint{3.287578in}{5.128744in}}{\pgfqpoint{3.283460in}{5.132862in}}%
\pgfpathcurveto{\pgfqpoint{3.279342in}{5.136980in}}{\pgfqpoint{3.273756in}{5.139294in}}{\pgfqpoint{3.267932in}{5.139294in}}%
\pgfpathcurveto{\pgfqpoint{3.262108in}{5.139294in}}{\pgfqpoint{3.256522in}{5.136980in}}{\pgfqpoint{3.252403in}{5.132862in}}%
\pgfpathcurveto{\pgfqpoint{3.248285in}{5.128744in}}{\pgfqpoint{3.245971in}{5.123158in}}{\pgfqpoint{3.245971in}{5.117334in}}%
\pgfpathcurveto{\pgfqpoint{3.245971in}{5.111510in}}{\pgfqpoint{3.248285in}{5.105924in}}{\pgfqpoint{3.252403in}{5.101806in}}%
\pgfpathcurveto{\pgfqpoint{3.256522in}{5.097688in}}{\pgfqpoint{3.262108in}{5.095374in}}{\pgfqpoint{3.267932in}{5.095374in}}%
\pgfpathlineto{\pgfqpoint{3.267932in}{5.095374in}}%
\pgfpathclose%
\pgfusepath{stroke,fill}%
\end{pgfscope}%
\begin{pgfscope}%
\pgfpathrectangle{\pgfqpoint{1.414736in}{0.548769in}}{\pgfqpoint{3.497747in}{7.253186in}}%
\pgfusepath{clip}%
\pgfsetbuttcap%
\pgfsetroundjoin%
\definecolor{currentfill}{rgb}{0.165117,0.467423,0.558141}%
\pgfsetfillcolor{currentfill}%
\pgfsetlinewidth{1.003750pt}%
\definecolor{currentstroke}{rgb}{0.165117,0.467423,0.558141}%
\pgfsetstrokecolor{currentstroke}%
\pgfsetdash{}{0pt}%
\pgfpathmoveto{\pgfqpoint{3.664896in}{5.566360in}}%
\pgfpathcurveto{\pgfqpoint{3.670720in}{5.566360in}}{\pgfqpoint{3.676306in}{5.568674in}}{\pgfqpoint{3.680424in}{5.572792in}}%
\pgfpathcurveto{\pgfqpoint{3.684542in}{5.576910in}}{\pgfqpoint{3.686856in}{5.582496in}}{\pgfqpoint{3.686856in}{5.588320in}}%
\pgfpathcurveto{\pgfqpoint{3.686856in}{5.594144in}}{\pgfqpoint{3.684542in}{5.599730in}}{\pgfqpoint{3.680424in}{5.603848in}}%
\pgfpathcurveto{\pgfqpoint{3.676306in}{5.607967in}}{\pgfqpoint{3.670720in}{5.610280in}}{\pgfqpoint{3.664896in}{5.610280in}}%
\pgfpathcurveto{\pgfqpoint{3.659072in}{5.610280in}}{\pgfqpoint{3.653486in}{5.607967in}}{\pgfqpoint{3.649368in}{5.603848in}}%
\pgfpathcurveto{\pgfqpoint{3.645250in}{5.599730in}}{\pgfqpoint{3.642936in}{5.594144in}}{\pgfqpoint{3.642936in}{5.588320in}}%
\pgfpathcurveto{\pgfqpoint{3.642936in}{5.582496in}}{\pgfqpoint{3.645250in}{5.576910in}}{\pgfqpoint{3.649368in}{5.572792in}}%
\pgfpathcurveto{\pgfqpoint{3.653486in}{5.568674in}}{\pgfqpoint{3.659072in}{5.566360in}}{\pgfqpoint{3.664896in}{5.566360in}}%
\pgfpathlineto{\pgfqpoint{3.664896in}{5.566360in}}%
\pgfpathclose%
\pgfusepath{stroke,fill}%
\end{pgfscope}%
\begin{pgfscope}%
\pgfpathrectangle{\pgfqpoint{1.414736in}{0.548769in}}{\pgfqpoint{3.497747in}{7.253186in}}%
\pgfusepath{clip}%
\pgfsetbuttcap%
\pgfsetroundjoin%
\definecolor{currentfill}{rgb}{0.121831,0.589055,0.545623}%
\pgfsetfillcolor{currentfill}%
\pgfsetlinewidth{1.003750pt}%
\definecolor{currentstroke}{rgb}{0.121831,0.589055,0.545623}%
\pgfsetstrokecolor{currentstroke}%
\pgfsetdash{}{0pt}%
\pgfpathmoveto{\pgfqpoint{4.380750in}{6.037346in}}%
\pgfpathcurveto{\pgfqpoint{4.386574in}{6.037346in}}{\pgfqpoint{4.392160in}{6.039660in}}{\pgfqpoint{4.396279in}{6.043778in}}%
\pgfpathcurveto{\pgfqpoint{4.400397in}{6.047896in}}{\pgfqpoint{4.402711in}{6.053482in}}{\pgfqpoint{4.402711in}{6.059306in}}%
\pgfpathcurveto{\pgfqpoint{4.402711in}{6.065130in}}{\pgfqpoint{4.400397in}{6.070716in}}{\pgfqpoint{4.396279in}{6.074835in}}%
\pgfpathcurveto{\pgfqpoint{4.392160in}{6.078953in}}{\pgfqpoint{4.386574in}{6.081267in}}{\pgfqpoint{4.380750in}{6.081267in}}%
\pgfpathcurveto{\pgfqpoint{4.374926in}{6.081267in}}{\pgfqpoint{4.369340in}{6.078953in}}{\pgfqpoint{4.365222in}{6.074835in}}%
\pgfpathcurveto{\pgfqpoint{4.361104in}{6.070716in}}{\pgfqpoint{4.358790in}{6.065130in}}{\pgfqpoint{4.358790in}{6.059306in}}%
\pgfpathcurveto{\pgfqpoint{4.358790in}{6.053482in}}{\pgfqpoint{4.361104in}{6.047896in}}{\pgfqpoint{4.365222in}{6.043778in}}%
\pgfpathcurveto{\pgfqpoint{4.369340in}{6.039660in}}{\pgfqpoint{4.374926in}{6.037346in}}{\pgfqpoint{4.380750in}{6.037346in}}%
\pgfpathlineto{\pgfqpoint{4.380750in}{6.037346in}}%
\pgfpathclose%
\pgfusepath{stroke,fill}%
\end{pgfscope}%
\begin{pgfscope}%
\pgfpathrectangle{\pgfqpoint{1.414736in}{0.548769in}}{\pgfqpoint{3.497747in}{7.253186in}}%
\pgfusepath{clip}%
\pgfsetbuttcap%
\pgfsetroundjoin%
\definecolor{currentfill}{rgb}{0.993248,0.906157,0.143936}%
\pgfsetfillcolor{currentfill}%
\pgfsetlinewidth{1.003750pt}%
\definecolor{currentstroke}{rgb}{0.993248,0.906157,0.143936}%
\pgfsetstrokecolor{currentstroke}%
\pgfsetdash{}{0pt}%
\pgfpathmoveto{\pgfqpoint{4.325978in}{6.508332in}}%
\pgfpathcurveto{\pgfqpoint{4.331802in}{6.508332in}}{\pgfqpoint{4.337388in}{6.510646in}}{\pgfqpoint{4.341506in}{6.514764in}}%
\pgfpathcurveto{\pgfqpoint{4.345624in}{6.518882in}}{\pgfqpoint{4.347938in}{6.524468in}}{\pgfqpoint{4.347938in}{6.530292in}}%
\pgfpathcurveto{\pgfqpoint{4.347938in}{6.536116in}}{\pgfqpoint{4.345624in}{6.541702in}}{\pgfqpoint{4.341506in}{6.545821in}}%
\pgfpathcurveto{\pgfqpoint{4.337388in}{6.549939in}}{\pgfqpoint{4.331802in}{6.552253in}}{\pgfqpoint{4.325978in}{6.552253in}}%
\pgfpathcurveto{\pgfqpoint{4.320154in}{6.552253in}}{\pgfqpoint{4.314568in}{6.549939in}}{\pgfqpoint{4.310450in}{6.545821in}}%
\pgfpathcurveto{\pgfqpoint{4.306332in}{6.541702in}}{\pgfqpoint{4.304018in}{6.536116in}}{\pgfqpoint{4.304018in}{6.530292in}}%
\pgfpathcurveto{\pgfqpoint{4.304018in}{6.524468in}}{\pgfqpoint{4.306332in}{6.518882in}}{\pgfqpoint{4.310450in}{6.514764in}}%
\pgfpathcurveto{\pgfqpoint{4.314568in}{6.510646in}}{\pgfqpoint{4.320154in}{6.508332in}}{\pgfqpoint{4.325978in}{6.508332in}}%
\pgfpathlineto{\pgfqpoint{4.325978in}{6.508332in}}%
\pgfpathclose%
\pgfusepath{stroke,fill}%
\end{pgfscope}%
\begin{pgfscope}%
\pgfpathrectangle{\pgfqpoint{1.414736in}{0.548769in}}{\pgfqpoint{3.497747in}{7.253186in}}%
\pgfusepath{clip}%
\pgfsetbuttcap%
\pgfsetroundjoin%
\definecolor{currentfill}{rgb}{0.772852,0.877868,0.131109}%
\pgfsetfillcolor{currentfill}%
\pgfsetlinewidth{1.003750pt}%
\definecolor{currentstroke}{rgb}{0.772852,0.877868,0.131109}%
\pgfsetstrokecolor{currentstroke}%
\pgfsetdash{}{0pt}%
\pgfpathmoveto{\pgfqpoint{3.120002in}{6.979318in}}%
\pgfpathcurveto{\pgfqpoint{3.125826in}{6.979318in}}{\pgfqpoint{3.131412in}{6.981632in}}{\pgfqpoint{3.135530in}{6.985750in}}%
\pgfpathcurveto{\pgfqpoint{3.139648in}{6.989868in}}{\pgfqpoint{3.141962in}{6.995455in}}{\pgfqpoint{3.141962in}{7.001278in}}%
\pgfpathcurveto{\pgfqpoint{3.141962in}{7.007102in}}{\pgfqpoint{3.139648in}{7.012689in}}{\pgfqpoint{3.135530in}{7.016807in}}%
\pgfpathcurveto{\pgfqpoint{3.131412in}{7.020925in}}{\pgfqpoint{3.125826in}{7.023239in}}{\pgfqpoint{3.120002in}{7.023239in}}%
\pgfpathcurveto{\pgfqpoint{3.114178in}{7.023239in}}{\pgfqpoint{3.108592in}{7.020925in}}{\pgfqpoint{3.104473in}{7.016807in}}%
\pgfpathcurveto{\pgfqpoint{3.100355in}{7.012689in}}{\pgfqpoint{3.098041in}{7.007102in}}{\pgfqpoint{3.098041in}{7.001278in}}%
\pgfpathcurveto{\pgfqpoint{3.098041in}{6.995455in}}{\pgfqpoint{3.100355in}{6.989868in}}{\pgfqpoint{3.104473in}{6.985750in}}%
\pgfpathcurveto{\pgfqpoint{3.108592in}{6.981632in}}{\pgfqpoint{3.114178in}{6.979318in}}{\pgfqpoint{3.120002in}{6.979318in}}%
\pgfpathlineto{\pgfqpoint{3.120002in}{6.979318in}}%
\pgfpathclose%
\pgfusepath{stroke,fill}%
\end{pgfscope}%
\begin{pgfscope}%
\pgfpathrectangle{\pgfqpoint{1.414736in}{0.548769in}}{\pgfqpoint{3.497747in}{7.253186in}}%
\pgfusepath{clip}%
\pgfsetbuttcap%
\pgfsetroundjoin%
\definecolor{currentfill}{rgb}{0.783315,0.879285,0.125405}%
\pgfsetfillcolor{currentfill}%
\pgfsetlinewidth{1.003750pt}%
\definecolor{currentstroke}{rgb}{0.783315,0.879285,0.125405}%
\pgfsetstrokecolor{currentstroke}%
\pgfsetdash{}{0pt}%
\pgfpathmoveto{\pgfqpoint{3.619837in}{7.450304in}}%
\pgfpathcurveto{\pgfqpoint{3.625661in}{7.450304in}}{\pgfqpoint{3.631247in}{7.452618in}}{\pgfqpoint{3.635365in}{7.456736in}}%
\pgfpathcurveto{\pgfqpoint{3.639483in}{7.460854in}}{\pgfqpoint{3.641797in}{7.466441in}}{\pgfqpoint{3.641797in}{7.472265in}}%
\pgfpathcurveto{\pgfqpoint{3.641797in}{7.478089in}}{\pgfqpoint{3.639483in}{7.483675in}}{\pgfqpoint{3.635365in}{7.487793in}}%
\pgfpathcurveto{\pgfqpoint{3.631247in}{7.491911in}}{\pgfqpoint{3.625661in}{7.494225in}}{\pgfqpoint{3.619837in}{7.494225in}}%
\pgfpathcurveto{\pgfqpoint{3.614013in}{7.494225in}}{\pgfqpoint{3.608427in}{7.491911in}}{\pgfqpoint{3.604309in}{7.487793in}}%
\pgfpathcurveto{\pgfqpoint{3.600191in}{7.483675in}}{\pgfqpoint{3.597877in}{7.478089in}}{\pgfqpoint{3.597877in}{7.472265in}}%
\pgfpathcurveto{\pgfqpoint{3.597877in}{7.466441in}}{\pgfqpoint{3.600191in}{7.460854in}}{\pgfqpoint{3.604309in}{7.456736in}}%
\pgfpathcurveto{\pgfqpoint{3.608427in}{7.452618in}}{\pgfqpoint{3.614013in}{7.450304in}}{\pgfqpoint{3.619837in}{7.450304in}}%
\pgfpathlineto{\pgfqpoint{3.619837in}{7.450304in}}%
\pgfpathclose%
\pgfusepath{stroke,fill}%
\end{pgfscope}%
\begin{pgfscope}%
\pgfpathrectangle{\pgfqpoint{1.414736in}{0.548769in}}{\pgfqpoint{3.497747in}{7.253186in}}%
\pgfusepath{clip}%
\pgfsetbuttcap%
\pgfsetroundjoin%
\definecolor{currentfill}{rgb}{0.926106,0.897330,0.104071}%
\pgfsetfillcolor{currentfill}%
\pgfsetlinewidth{1.003750pt}%
\definecolor{currentstroke}{rgb}{0.926106,0.897330,0.104071}%
\pgfsetstrokecolor{currentstroke}%
\pgfsetdash{}{0pt}%
\pgfpathmoveto{\pgfqpoint{3.786544in}{0.856499in}}%
\pgfpathcurveto{\pgfqpoint{3.792368in}{0.856499in}}{\pgfqpoint{3.797954in}{0.858813in}}{\pgfqpoint{3.802072in}{0.862931in}}%
\pgfpathcurveto{\pgfqpoint{3.806190in}{0.867049in}}{\pgfqpoint{3.808504in}{0.872635in}}{\pgfqpoint{3.808504in}{0.878459in}}%
\pgfpathcurveto{\pgfqpoint{3.808504in}{0.884283in}}{\pgfqpoint{3.806190in}{0.889869in}}{\pgfqpoint{3.802072in}{0.893987in}}%
\pgfpathcurveto{\pgfqpoint{3.797954in}{0.898105in}}{\pgfqpoint{3.792368in}{0.900419in}}{\pgfqpoint{3.786544in}{0.900419in}}%
\pgfpathcurveto{\pgfqpoint{3.780720in}{0.900419in}}{\pgfqpoint{3.775134in}{0.898105in}}{\pgfqpoint{3.771016in}{0.893987in}}%
\pgfpathcurveto{\pgfqpoint{3.766897in}{0.889869in}}{\pgfqpoint{3.764584in}{0.884283in}}{\pgfqpoint{3.764584in}{0.878459in}}%
\pgfpathcurveto{\pgfqpoint{3.764584in}{0.872635in}}{\pgfqpoint{3.766897in}{0.867049in}}{\pgfqpoint{3.771016in}{0.862931in}}%
\pgfpathcurveto{\pgfqpoint{3.775134in}{0.858813in}}{\pgfqpoint{3.780720in}{0.856499in}}{\pgfqpoint{3.786544in}{0.856499in}}%
\pgfpathlineto{\pgfqpoint{3.786544in}{0.856499in}}%
\pgfpathclose%
\pgfusepath{stroke,fill}%
\end{pgfscope}%
\begin{pgfscope}%
\pgfpathrectangle{\pgfqpoint{1.414736in}{0.548769in}}{\pgfqpoint{3.497747in}{7.253186in}}%
\pgfusepath{clip}%
\pgfsetbuttcap%
\pgfsetroundjoin%
\definecolor{currentfill}{rgb}{0.678489,0.863742,0.189503}%
\pgfsetfillcolor{currentfill}%
\pgfsetlinewidth{1.003750pt}%
\definecolor{currentstroke}{rgb}{0.678489,0.863742,0.189503}%
\pgfsetstrokecolor{currentstroke}%
\pgfsetdash{}{0pt}%
\pgfpathmoveto{\pgfqpoint{3.441935in}{1.327485in}}%
\pgfpathcurveto{\pgfqpoint{3.447758in}{1.327485in}}{\pgfqpoint{3.453345in}{1.329799in}}{\pgfqpoint{3.457463in}{1.333917in}}%
\pgfpathcurveto{\pgfqpoint{3.461581in}{1.338035in}}{\pgfqpoint{3.463895in}{1.343621in}}{\pgfqpoint{3.463895in}{1.349445in}}%
\pgfpathcurveto{\pgfqpoint{3.463895in}{1.355269in}}{\pgfqpoint{3.461581in}{1.360855in}}{\pgfqpoint{3.457463in}{1.364973in}}%
\pgfpathcurveto{\pgfqpoint{3.453345in}{1.369092in}}{\pgfqpoint{3.447758in}{1.371405in}}{\pgfqpoint{3.441935in}{1.371405in}}%
\pgfpathcurveto{\pgfqpoint{3.436111in}{1.371405in}}{\pgfqpoint{3.430524in}{1.369092in}}{\pgfqpoint{3.426406in}{1.364973in}}%
\pgfpathcurveto{\pgfqpoint{3.422288in}{1.360855in}}{\pgfqpoint{3.419974in}{1.355269in}}{\pgfqpoint{3.419974in}{1.349445in}}%
\pgfpathcurveto{\pgfqpoint{3.419974in}{1.343621in}}{\pgfqpoint{3.422288in}{1.338035in}}{\pgfqpoint{3.426406in}{1.333917in}}%
\pgfpathcurveto{\pgfqpoint{3.430524in}{1.329799in}}{\pgfqpoint{3.436111in}{1.327485in}}{\pgfqpoint{3.441935in}{1.327485in}}%
\pgfpathlineto{\pgfqpoint{3.441935in}{1.327485in}}%
\pgfpathclose%
\pgfusepath{stroke,fill}%
\end{pgfscope}%
\begin{pgfscope}%
\pgfpathrectangle{\pgfqpoint{1.414736in}{0.548769in}}{\pgfqpoint{3.497747in}{7.253186in}}%
\pgfusepath{clip}%
\pgfsetbuttcap%
\pgfsetroundjoin%
\definecolor{currentfill}{rgb}{0.699415,0.867117,0.175971}%
\pgfsetfillcolor{currentfill}%
\pgfsetlinewidth{1.003750pt}%
\definecolor{currentstroke}{rgb}{0.699415,0.867117,0.175971}%
\pgfsetstrokecolor{currentstroke}%
\pgfsetdash{}{0pt}%
\pgfpathmoveto{\pgfqpoint{3.752733in}{1.798471in}}%
\pgfpathcurveto{\pgfqpoint{3.758557in}{1.798471in}}{\pgfqpoint{3.764143in}{1.800785in}}{\pgfqpoint{3.768261in}{1.804903in}}%
\pgfpathcurveto{\pgfqpoint{3.772379in}{1.809021in}}{\pgfqpoint{3.774693in}{1.814607in}}{\pgfqpoint{3.774693in}{1.820431in}}%
\pgfpathcurveto{\pgfqpoint{3.774693in}{1.826255in}}{\pgfqpoint{3.772379in}{1.831841in}}{\pgfqpoint{3.768261in}{1.835960in}}%
\pgfpathcurveto{\pgfqpoint{3.764143in}{1.840078in}}{\pgfqpoint{3.758557in}{1.842392in}}{\pgfqpoint{3.752733in}{1.842392in}}%
\pgfpathcurveto{\pgfqpoint{3.746909in}{1.842392in}}{\pgfqpoint{3.741323in}{1.840078in}}{\pgfqpoint{3.737205in}{1.835960in}}%
\pgfpathcurveto{\pgfqpoint{3.733087in}{1.831841in}}{\pgfqpoint{3.730773in}{1.826255in}}{\pgfqpoint{3.730773in}{1.820431in}}%
\pgfpathcurveto{\pgfqpoint{3.730773in}{1.814607in}}{\pgfqpoint{3.733087in}{1.809021in}}{\pgfqpoint{3.737205in}{1.804903in}}%
\pgfpathcurveto{\pgfqpoint{3.741323in}{1.800785in}}{\pgfqpoint{3.746909in}{1.798471in}}{\pgfqpoint{3.752733in}{1.798471in}}%
\pgfpathlineto{\pgfqpoint{3.752733in}{1.798471in}}%
\pgfpathclose%
\pgfusepath{stroke,fill}%
\end{pgfscope}%
\begin{pgfscope}%
\pgfpathrectangle{\pgfqpoint{1.414736in}{0.548769in}}{\pgfqpoint{3.497747in}{7.253186in}}%
\pgfusepath{clip}%
\pgfsetbuttcap%
\pgfsetroundjoin%
\definecolor{currentfill}{rgb}{0.993248,0.906157,0.143936}%
\pgfsetfillcolor{currentfill}%
\pgfsetlinewidth{1.003750pt}%
\definecolor{currentstroke}{rgb}{0.993248,0.906157,0.143936}%
\pgfsetstrokecolor{currentstroke}%
\pgfsetdash{}{0pt}%
\pgfpathmoveto{\pgfqpoint{3.753678in}{2.269457in}}%
\pgfpathcurveto{\pgfqpoint{3.759502in}{2.269457in}}{\pgfqpoint{3.765088in}{2.271771in}}{\pgfqpoint{3.769206in}{2.275889in}}%
\pgfpathcurveto{\pgfqpoint{3.773324in}{2.280007in}}{\pgfqpoint{3.775638in}{2.285593in}}{\pgfqpoint{3.775638in}{2.291417in}}%
\pgfpathcurveto{\pgfqpoint{3.775638in}{2.297241in}}{\pgfqpoint{3.773324in}{2.302828in}}{\pgfqpoint{3.769206in}{2.306946in}}%
\pgfpathcurveto{\pgfqpoint{3.765088in}{2.311064in}}{\pgfqpoint{3.759502in}{2.313378in}}{\pgfqpoint{3.753678in}{2.313378in}}%
\pgfpathcurveto{\pgfqpoint{3.747854in}{2.313378in}}{\pgfqpoint{3.742268in}{2.311064in}}{\pgfqpoint{3.738149in}{2.306946in}}%
\pgfpathcurveto{\pgfqpoint{3.734031in}{2.302828in}}{\pgfqpoint{3.731717in}{2.297241in}}{\pgfqpoint{3.731717in}{2.291417in}}%
\pgfpathcurveto{\pgfqpoint{3.731717in}{2.285593in}}{\pgfqpoint{3.734031in}{2.280007in}}{\pgfqpoint{3.738149in}{2.275889in}}%
\pgfpathcurveto{\pgfqpoint{3.742268in}{2.271771in}}{\pgfqpoint{3.747854in}{2.269457in}}{\pgfqpoint{3.753678in}{2.269457in}}%
\pgfpathlineto{\pgfqpoint{3.753678in}{2.269457in}}%
\pgfpathclose%
\pgfusepath{stroke,fill}%
\end{pgfscope}%
\begin{pgfscope}%
\pgfpathrectangle{\pgfqpoint{1.414736in}{0.548769in}}{\pgfqpoint{3.497747in}{7.253186in}}%
\pgfusepath{clip}%
\pgfsetbuttcap%
\pgfsetroundjoin%
\definecolor{currentfill}{rgb}{0.252899,0.742211,0.448284}%
\pgfsetfillcolor{currentfill}%
\pgfsetlinewidth{1.003750pt}%
\definecolor{currentstroke}{rgb}{0.252899,0.742211,0.448284}%
\pgfsetstrokecolor{currentstroke}%
\pgfsetdash{}{0pt}%
\pgfpathmoveto{\pgfqpoint{3.790386in}{2.740443in}}%
\pgfpathcurveto{\pgfqpoint{3.796210in}{2.740443in}}{\pgfqpoint{3.801796in}{2.742757in}}{\pgfqpoint{3.805914in}{2.746875in}}%
\pgfpathcurveto{\pgfqpoint{3.810032in}{2.750993in}}{\pgfqpoint{3.812346in}{2.756580in}}{\pgfqpoint{3.812346in}{2.762404in}}%
\pgfpathcurveto{\pgfqpoint{3.812346in}{2.768227in}}{\pgfqpoint{3.810032in}{2.773814in}}{\pgfqpoint{3.805914in}{2.777932in}}%
\pgfpathcurveto{\pgfqpoint{3.801796in}{2.782050in}}{\pgfqpoint{3.796210in}{2.784364in}}{\pgfqpoint{3.790386in}{2.784364in}}%
\pgfpathcurveto{\pgfqpoint{3.784562in}{2.784364in}}{\pgfqpoint{3.778976in}{2.782050in}}{\pgfqpoint{3.774857in}{2.777932in}}%
\pgfpathcurveto{\pgfqpoint{3.770739in}{2.773814in}}{\pgfqpoint{3.768425in}{2.768227in}}{\pgfqpoint{3.768425in}{2.762404in}}%
\pgfpathcurveto{\pgfqpoint{3.768425in}{2.756580in}}{\pgfqpoint{3.770739in}{2.750993in}}{\pgfqpoint{3.774857in}{2.746875in}}%
\pgfpathcurveto{\pgfqpoint{3.778976in}{2.742757in}}{\pgfqpoint{3.784562in}{2.740443in}}{\pgfqpoint{3.790386in}{2.740443in}}%
\pgfpathlineto{\pgfqpoint{3.790386in}{2.740443in}}%
\pgfpathclose%
\pgfusepath{stroke,fill}%
\end{pgfscope}%
\begin{pgfscope}%
\pgfpathrectangle{\pgfqpoint{1.414736in}{0.548769in}}{\pgfqpoint{3.497747in}{7.253186in}}%
\pgfusepath{clip}%
\pgfsetbuttcap%
\pgfsetroundjoin%
\definecolor{currentfill}{rgb}{0.876168,0.891125,0.095250}%
\pgfsetfillcolor{currentfill}%
\pgfsetlinewidth{1.003750pt}%
\definecolor{currentstroke}{rgb}{0.876168,0.891125,0.095250}%
\pgfsetstrokecolor{currentstroke}%
\pgfsetdash{}{0pt}%
\pgfpathmoveto{\pgfqpoint{3.521503in}{3.211429in}}%
\pgfpathcurveto{\pgfqpoint{3.527327in}{3.211429in}}{\pgfqpoint{3.532913in}{3.213743in}}{\pgfqpoint{3.537031in}{3.217861in}}%
\pgfpathcurveto{\pgfqpoint{3.541149in}{3.221980in}}{\pgfqpoint{3.543463in}{3.227566in}}{\pgfqpoint{3.543463in}{3.233390in}}%
\pgfpathcurveto{\pgfqpoint{3.543463in}{3.239214in}}{\pgfqpoint{3.541149in}{3.244800in}}{\pgfqpoint{3.537031in}{3.248918in}}%
\pgfpathcurveto{\pgfqpoint{3.532913in}{3.253036in}}{\pgfqpoint{3.527327in}{3.255350in}}{\pgfqpoint{3.521503in}{3.255350in}}%
\pgfpathcurveto{\pgfqpoint{3.515679in}{3.255350in}}{\pgfqpoint{3.510092in}{3.253036in}}{\pgfqpoint{3.505974in}{3.248918in}}%
\pgfpathcurveto{\pgfqpoint{3.501856in}{3.244800in}}{\pgfqpoint{3.499542in}{3.239214in}}{\pgfqpoint{3.499542in}{3.233390in}}%
\pgfpathcurveto{\pgfqpoint{3.499542in}{3.227566in}}{\pgfqpoint{3.501856in}{3.221980in}}{\pgfqpoint{3.505974in}{3.217861in}}%
\pgfpathcurveto{\pgfqpoint{3.510092in}{3.213743in}}{\pgfqpoint{3.515679in}{3.211429in}}{\pgfqpoint{3.521503in}{3.211429in}}%
\pgfpathlineto{\pgfqpoint{3.521503in}{3.211429in}}%
\pgfpathclose%
\pgfusepath{stroke,fill}%
\end{pgfscope}%
\begin{pgfscope}%
\pgfpathrectangle{\pgfqpoint{1.414736in}{0.548769in}}{\pgfqpoint{3.497747in}{7.253186in}}%
\pgfusepath{clip}%
\pgfsetbuttcap%
\pgfsetroundjoin%
\definecolor{currentfill}{rgb}{0.876168,0.891125,0.095250}%
\pgfsetfillcolor{currentfill}%
\pgfsetlinewidth{1.003750pt}%
\definecolor{currentstroke}{rgb}{0.876168,0.891125,0.095250}%
\pgfsetstrokecolor{currentstroke}%
\pgfsetdash{}{0pt}%
\pgfpathmoveto{\pgfqpoint{3.443747in}{3.682415in}}%
\pgfpathcurveto{\pgfqpoint{3.449571in}{3.682415in}}{\pgfqpoint{3.455158in}{3.684729in}}{\pgfqpoint{3.459276in}{3.688847in}}%
\pgfpathcurveto{\pgfqpoint{3.463394in}{3.692966in}}{\pgfqpoint{3.465708in}{3.698552in}}{\pgfqpoint{3.465708in}{3.704376in}}%
\pgfpathcurveto{\pgfqpoint{3.465708in}{3.710200in}}{\pgfqpoint{3.463394in}{3.715786in}}{\pgfqpoint{3.459276in}{3.719904in}}%
\pgfpathcurveto{\pgfqpoint{3.455158in}{3.724022in}}{\pgfqpoint{3.449571in}{3.726336in}}{\pgfqpoint{3.443747in}{3.726336in}}%
\pgfpathcurveto{\pgfqpoint{3.437924in}{3.726336in}}{\pgfqpoint{3.432337in}{3.724022in}}{\pgfqpoint{3.428219in}{3.719904in}}%
\pgfpathcurveto{\pgfqpoint{3.424101in}{3.715786in}}{\pgfqpoint{3.421787in}{3.710200in}}{\pgfqpoint{3.421787in}{3.704376in}}%
\pgfpathcurveto{\pgfqpoint{3.421787in}{3.698552in}}{\pgfqpoint{3.424101in}{3.692966in}}{\pgfqpoint{3.428219in}{3.688847in}}%
\pgfpathcurveto{\pgfqpoint{3.432337in}{3.684729in}}{\pgfqpoint{3.437924in}{3.682415in}}{\pgfqpoint{3.443747in}{3.682415in}}%
\pgfpathlineto{\pgfqpoint{3.443747in}{3.682415in}}%
\pgfpathclose%
\pgfusepath{stroke,fill}%
\end{pgfscope}%
\begin{pgfscope}%
\pgfpathrectangle{\pgfqpoint{1.414736in}{0.548769in}}{\pgfqpoint{3.497747in}{7.253186in}}%
\pgfusepath{clip}%
\pgfsetbuttcap%
\pgfsetroundjoin%
\definecolor{currentfill}{rgb}{0.606045,0.850733,0.236712}%
\pgfsetfillcolor{currentfill}%
\pgfsetlinewidth{1.003750pt}%
\definecolor{currentstroke}{rgb}{0.606045,0.850733,0.236712}%
\pgfsetstrokecolor{currentstroke}%
\pgfsetdash{}{0pt}%
\pgfpathmoveto{\pgfqpoint{3.716793in}{4.153402in}}%
\pgfpathcurveto{\pgfqpoint{3.722617in}{4.153402in}}{\pgfqpoint{3.728204in}{4.155715in}}{\pgfqpoint{3.732322in}{4.159834in}}%
\pgfpathcurveto{\pgfqpoint{3.736440in}{4.163952in}}{\pgfqpoint{3.738754in}{4.169538in}}{\pgfqpoint{3.738754in}{4.175362in}}%
\pgfpathcurveto{\pgfqpoint{3.738754in}{4.181186in}}{\pgfqpoint{3.736440in}{4.186772in}}{\pgfqpoint{3.732322in}{4.190890in}}%
\pgfpathcurveto{\pgfqpoint{3.728204in}{4.195008in}}{\pgfqpoint{3.722617in}{4.197322in}}{\pgfqpoint{3.716793in}{4.197322in}}%
\pgfpathcurveto{\pgfqpoint{3.710969in}{4.197322in}}{\pgfqpoint{3.705383in}{4.195008in}}{\pgfqpoint{3.701265in}{4.190890in}}%
\pgfpathcurveto{\pgfqpoint{3.697147in}{4.186772in}}{\pgfqpoint{3.694833in}{4.181186in}}{\pgfqpoint{3.694833in}{4.175362in}}%
\pgfpathcurveto{\pgfqpoint{3.694833in}{4.169538in}}{\pgfqpoint{3.697147in}{4.163952in}}{\pgfqpoint{3.701265in}{4.159834in}}%
\pgfpathcurveto{\pgfqpoint{3.705383in}{4.155715in}}{\pgfqpoint{3.710969in}{4.153402in}}{\pgfqpoint{3.716793in}{4.153402in}}%
\pgfpathlineto{\pgfqpoint{3.716793in}{4.153402in}}%
\pgfpathclose%
\pgfusepath{stroke,fill}%
\end{pgfscope}%
\begin{pgfscope}%
\pgfpathrectangle{\pgfqpoint{1.414736in}{0.548769in}}{\pgfqpoint{3.497747in}{7.253186in}}%
\pgfusepath{clip}%
\pgfsetbuttcap%
\pgfsetroundjoin%
\definecolor{currentfill}{rgb}{0.876168,0.891125,0.095250}%
\pgfsetfillcolor{currentfill}%
\pgfsetlinewidth{1.003750pt}%
\definecolor{currentstroke}{rgb}{0.876168,0.891125,0.095250}%
\pgfsetstrokecolor{currentstroke}%
\pgfsetdash{}{0pt}%
\pgfpathmoveto{\pgfqpoint{3.809361in}{4.624388in}}%
\pgfpathcurveto{\pgfqpoint{3.815185in}{4.624388in}}{\pgfqpoint{3.820771in}{4.626702in}}{\pgfqpoint{3.824889in}{4.630820in}}%
\pgfpathcurveto{\pgfqpoint{3.829007in}{4.634938in}}{\pgfqpoint{3.831321in}{4.640524in}}{\pgfqpoint{3.831321in}{4.646348in}}%
\pgfpathcurveto{\pgfqpoint{3.831321in}{4.652172in}}{\pgfqpoint{3.829007in}{4.657758in}}{\pgfqpoint{3.824889in}{4.661876in}}%
\pgfpathcurveto{\pgfqpoint{3.820771in}{4.665994in}}{\pgfqpoint{3.815185in}{4.668308in}}{\pgfqpoint{3.809361in}{4.668308in}}%
\pgfpathcurveto{\pgfqpoint{3.803537in}{4.668308in}}{\pgfqpoint{3.797951in}{4.665994in}}{\pgfqpoint{3.793833in}{4.661876in}}%
\pgfpathcurveto{\pgfqpoint{3.789715in}{4.657758in}}{\pgfqpoint{3.787401in}{4.652172in}}{\pgfqpoint{3.787401in}{4.646348in}}%
\pgfpathcurveto{\pgfqpoint{3.787401in}{4.640524in}}{\pgfqpoint{3.789715in}{4.634938in}}{\pgfqpoint{3.793833in}{4.630820in}}%
\pgfpathcurveto{\pgfqpoint{3.797951in}{4.626702in}}{\pgfqpoint{3.803537in}{4.624388in}}{\pgfqpoint{3.809361in}{4.624388in}}%
\pgfpathlineto{\pgfqpoint{3.809361in}{4.624388in}}%
\pgfpathclose%
\pgfusepath{stroke,fill}%
\end{pgfscope}%
\begin{pgfscope}%
\pgfpathrectangle{\pgfqpoint{1.414736in}{0.548769in}}{\pgfqpoint{3.497747in}{7.253186in}}%
\pgfusepath{clip}%
\pgfsetbuttcap%
\pgfsetroundjoin%
\definecolor{currentfill}{rgb}{0.185783,0.704891,0.485273}%
\pgfsetfillcolor{currentfill}%
\pgfsetlinewidth{1.003750pt}%
\definecolor{currentstroke}{rgb}{0.185783,0.704891,0.485273}%
\pgfsetstrokecolor{currentstroke}%
\pgfsetdash{}{0pt}%
\pgfpathmoveto{\pgfqpoint{3.585991in}{5.095374in}}%
\pgfpathcurveto{\pgfqpoint{3.591815in}{5.095374in}}{\pgfqpoint{3.597401in}{5.097688in}}{\pgfqpoint{3.601519in}{5.101806in}}%
\pgfpathcurveto{\pgfqpoint{3.605637in}{5.105924in}}{\pgfqpoint{3.607951in}{5.111510in}}{\pgfqpoint{3.607951in}{5.117334in}}%
\pgfpathcurveto{\pgfqpoint{3.607951in}{5.123158in}}{\pgfqpoint{3.605637in}{5.128744in}}{\pgfqpoint{3.601519in}{5.132862in}}%
\pgfpathcurveto{\pgfqpoint{3.597401in}{5.136980in}}{\pgfqpoint{3.591815in}{5.139294in}}{\pgfqpoint{3.585991in}{5.139294in}}%
\pgfpathcurveto{\pgfqpoint{3.580167in}{5.139294in}}{\pgfqpoint{3.574581in}{5.136980in}}{\pgfqpoint{3.570463in}{5.132862in}}%
\pgfpathcurveto{\pgfqpoint{3.566345in}{5.128744in}}{\pgfqpoint{3.564031in}{5.123158in}}{\pgfqpoint{3.564031in}{5.117334in}}%
\pgfpathcurveto{\pgfqpoint{3.564031in}{5.111510in}}{\pgfqpoint{3.566345in}{5.105924in}}{\pgfqpoint{3.570463in}{5.101806in}}%
\pgfpathcurveto{\pgfqpoint{3.574581in}{5.097688in}}{\pgfqpoint{3.580167in}{5.095374in}}{\pgfqpoint{3.585991in}{5.095374in}}%
\pgfpathlineto{\pgfqpoint{3.585991in}{5.095374in}}%
\pgfpathclose%
\pgfusepath{stroke,fill}%
\end{pgfscope}%
\begin{pgfscope}%
\pgfpathrectangle{\pgfqpoint{1.414736in}{0.548769in}}{\pgfqpoint{3.497747in}{7.253186in}}%
\pgfusepath{clip}%
\pgfsetbuttcap%
\pgfsetroundjoin%
\definecolor{currentfill}{rgb}{0.136408,0.541173,0.554483}%
\pgfsetfillcolor{currentfill}%
\pgfsetlinewidth{1.003750pt}%
\definecolor{currentstroke}{rgb}{0.136408,0.541173,0.554483}%
\pgfsetstrokecolor{currentstroke}%
\pgfsetdash{}{0pt}%
\pgfpathmoveto{\pgfqpoint{3.657799in}{5.566360in}}%
\pgfpathcurveto{\pgfqpoint{3.663623in}{5.566360in}}{\pgfqpoint{3.669209in}{5.568674in}}{\pgfqpoint{3.673327in}{5.572792in}}%
\pgfpathcurveto{\pgfqpoint{3.677445in}{5.576910in}}{\pgfqpoint{3.679759in}{5.582496in}}{\pgfqpoint{3.679759in}{5.588320in}}%
\pgfpathcurveto{\pgfqpoint{3.679759in}{5.594144in}}{\pgfqpoint{3.677445in}{5.599730in}}{\pgfqpoint{3.673327in}{5.603848in}}%
\pgfpathcurveto{\pgfqpoint{3.669209in}{5.607967in}}{\pgfqpoint{3.663623in}{5.610280in}}{\pgfqpoint{3.657799in}{5.610280in}}%
\pgfpathcurveto{\pgfqpoint{3.651975in}{5.610280in}}{\pgfqpoint{3.646389in}{5.607967in}}{\pgfqpoint{3.642270in}{5.603848in}}%
\pgfpathcurveto{\pgfqpoint{3.638152in}{5.599730in}}{\pgfqpoint{3.635838in}{5.594144in}}{\pgfqpoint{3.635838in}{5.588320in}}%
\pgfpathcurveto{\pgfqpoint{3.635838in}{5.582496in}}{\pgfqpoint{3.638152in}{5.576910in}}{\pgfqpoint{3.642270in}{5.572792in}}%
\pgfpathcurveto{\pgfqpoint{3.646389in}{5.568674in}}{\pgfqpoint{3.651975in}{5.566360in}}{\pgfqpoint{3.657799in}{5.566360in}}%
\pgfpathlineto{\pgfqpoint{3.657799in}{5.566360in}}%
\pgfpathclose%
\pgfusepath{stroke,fill}%
\end{pgfscope}%
\begin{pgfscope}%
\pgfpathrectangle{\pgfqpoint{1.414736in}{0.548769in}}{\pgfqpoint{3.497747in}{7.253186in}}%
\pgfusepath{clip}%
\pgfsetbuttcap%
\pgfsetroundjoin%
\definecolor{currentfill}{rgb}{0.175707,0.697900,0.491033}%
\pgfsetfillcolor{currentfill}%
\pgfsetlinewidth{1.003750pt}%
\definecolor{currentstroke}{rgb}{0.175707,0.697900,0.491033}%
\pgfsetstrokecolor{currentstroke}%
\pgfsetdash{}{0pt}%
\pgfpathmoveto{\pgfqpoint{3.374252in}{6.037346in}}%
\pgfpathcurveto{\pgfqpoint{3.380076in}{6.037346in}}{\pgfqpoint{3.385662in}{6.039660in}}{\pgfqpoint{3.389780in}{6.043778in}}%
\pgfpathcurveto{\pgfqpoint{3.393898in}{6.047896in}}{\pgfqpoint{3.396212in}{6.053482in}}{\pgfqpoint{3.396212in}{6.059306in}}%
\pgfpathcurveto{\pgfqpoint{3.396212in}{6.065130in}}{\pgfqpoint{3.393898in}{6.070716in}}{\pgfqpoint{3.389780in}{6.074835in}}%
\pgfpathcurveto{\pgfqpoint{3.385662in}{6.078953in}}{\pgfqpoint{3.380076in}{6.081267in}}{\pgfqpoint{3.374252in}{6.081267in}}%
\pgfpathcurveto{\pgfqpoint{3.368428in}{6.081267in}}{\pgfqpoint{3.362842in}{6.078953in}}{\pgfqpoint{3.358724in}{6.074835in}}%
\pgfpathcurveto{\pgfqpoint{3.354606in}{6.070716in}}{\pgfqpoint{3.352292in}{6.065130in}}{\pgfqpoint{3.352292in}{6.059306in}}%
\pgfpathcurveto{\pgfqpoint{3.352292in}{6.053482in}}{\pgfqpoint{3.354606in}{6.047896in}}{\pgfqpoint{3.358724in}{6.043778in}}%
\pgfpathcurveto{\pgfqpoint{3.362842in}{6.039660in}}{\pgfqpoint{3.368428in}{6.037346in}}{\pgfqpoint{3.374252in}{6.037346in}}%
\pgfpathlineto{\pgfqpoint{3.374252in}{6.037346in}}%
\pgfpathclose%
\pgfusepath{stroke,fill}%
\end{pgfscope}%
\begin{pgfscope}%
\pgfpathrectangle{\pgfqpoint{1.414736in}{0.548769in}}{\pgfqpoint{3.497747in}{7.253186in}}%
\pgfusepath{clip}%
\pgfsetbuttcap%
\pgfsetroundjoin%
\definecolor{currentfill}{rgb}{0.595839,0.848717,0.243329}%
\pgfsetfillcolor{currentfill}%
\pgfsetlinewidth{1.003750pt}%
\definecolor{currentstroke}{rgb}{0.595839,0.848717,0.243329}%
\pgfsetstrokecolor{currentstroke}%
\pgfsetdash{}{0pt}%
\pgfpathmoveto{\pgfqpoint{3.835806in}{6.508332in}}%
\pgfpathcurveto{\pgfqpoint{3.841630in}{6.508332in}}{\pgfqpoint{3.847216in}{6.510646in}}{\pgfqpoint{3.851334in}{6.514764in}}%
\pgfpathcurveto{\pgfqpoint{3.855452in}{6.518882in}}{\pgfqpoint{3.857766in}{6.524468in}}{\pgfqpoint{3.857766in}{6.530292in}}%
\pgfpathcurveto{\pgfqpoint{3.857766in}{6.536116in}}{\pgfqpoint{3.855452in}{6.541702in}}{\pgfqpoint{3.851334in}{6.545821in}}%
\pgfpathcurveto{\pgfqpoint{3.847216in}{6.549939in}}{\pgfqpoint{3.841630in}{6.552253in}}{\pgfqpoint{3.835806in}{6.552253in}}%
\pgfpathcurveto{\pgfqpoint{3.829982in}{6.552253in}}{\pgfqpoint{3.824396in}{6.549939in}}{\pgfqpoint{3.820277in}{6.545821in}}%
\pgfpathcurveto{\pgfqpoint{3.816159in}{6.541702in}}{\pgfqpoint{3.813845in}{6.536116in}}{\pgfqpoint{3.813845in}{6.530292in}}%
\pgfpathcurveto{\pgfqpoint{3.813845in}{6.524468in}}{\pgfqpoint{3.816159in}{6.518882in}}{\pgfqpoint{3.820277in}{6.514764in}}%
\pgfpathcurveto{\pgfqpoint{3.824396in}{6.510646in}}{\pgfqpoint{3.829982in}{6.508332in}}{\pgfqpoint{3.835806in}{6.508332in}}%
\pgfpathlineto{\pgfqpoint{3.835806in}{6.508332in}}%
\pgfpathclose%
\pgfusepath{stroke,fill}%
\end{pgfscope}%
\begin{pgfscope}%
\pgfpathrectangle{\pgfqpoint{1.414736in}{0.548769in}}{\pgfqpoint{3.497747in}{7.253186in}}%
\pgfusepath{clip}%
\pgfsetbuttcap%
\pgfsetroundjoin%
\definecolor{currentfill}{rgb}{0.906311,0.894855,0.098125}%
\pgfsetfillcolor{currentfill}%
\pgfsetlinewidth{1.003750pt}%
\definecolor{currentstroke}{rgb}{0.906311,0.894855,0.098125}%
\pgfsetstrokecolor{currentstroke}%
\pgfsetdash{}{0pt}%
\pgfpathmoveto{\pgfqpoint{3.770768in}{6.979318in}}%
\pgfpathcurveto{\pgfqpoint{3.776592in}{6.979318in}}{\pgfqpoint{3.782178in}{6.981632in}}{\pgfqpoint{3.786296in}{6.985750in}}%
\pgfpathcurveto{\pgfqpoint{3.790414in}{6.989868in}}{\pgfqpoint{3.792728in}{6.995455in}}{\pgfqpoint{3.792728in}{7.001278in}}%
\pgfpathcurveto{\pgfqpoint{3.792728in}{7.007102in}}{\pgfqpoint{3.790414in}{7.012689in}}{\pgfqpoint{3.786296in}{7.016807in}}%
\pgfpathcurveto{\pgfqpoint{3.782178in}{7.020925in}}{\pgfqpoint{3.776592in}{7.023239in}}{\pgfqpoint{3.770768in}{7.023239in}}%
\pgfpathcurveto{\pgfqpoint{3.764944in}{7.023239in}}{\pgfqpoint{3.759358in}{7.020925in}}{\pgfqpoint{3.755240in}{7.016807in}}%
\pgfpathcurveto{\pgfqpoint{3.751121in}{7.012689in}}{\pgfqpoint{3.748808in}{7.007102in}}{\pgfqpoint{3.748808in}{7.001278in}}%
\pgfpathcurveto{\pgfqpoint{3.748808in}{6.995455in}}{\pgfqpoint{3.751121in}{6.989868in}}{\pgfqpoint{3.755240in}{6.985750in}}%
\pgfpathcurveto{\pgfqpoint{3.759358in}{6.981632in}}{\pgfqpoint{3.764944in}{6.979318in}}{\pgfqpoint{3.770768in}{6.979318in}}%
\pgfpathlineto{\pgfqpoint{3.770768in}{6.979318in}}%
\pgfpathclose%
\pgfusepath{stroke,fill}%
\end{pgfscope}%
\begin{pgfscope}%
\pgfpathrectangle{\pgfqpoint{1.414736in}{0.548769in}}{\pgfqpoint{3.497747in}{7.253186in}}%
\pgfusepath{clip}%
\pgfsetbuttcap%
\pgfsetroundjoin%
\definecolor{currentfill}{rgb}{0.468053,0.818921,0.323998}%
\pgfsetfillcolor{currentfill}%
\pgfsetlinewidth{1.003750pt}%
\definecolor{currentstroke}{rgb}{0.468053,0.818921,0.323998}%
\pgfsetstrokecolor{currentstroke}%
\pgfsetdash{}{0pt}%
\pgfpathmoveto{\pgfqpoint{3.226632in}{7.450304in}}%
\pgfpathcurveto{\pgfqpoint{3.232456in}{7.450304in}}{\pgfqpoint{3.238042in}{7.452618in}}{\pgfqpoint{3.242160in}{7.456736in}}%
\pgfpathcurveto{\pgfqpoint{3.246278in}{7.460854in}}{\pgfqpoint{3.248592in}{7.466441in}}{\pgfqpoint{3.248592in}{7.472265in}}%
\pgfpathcurveto{\pgfqpoint{3.248592in}{7.478089in}}{\pgfqpoint{3.246278in}{7.483675in}}{\pgfqpoint{3.242160in}{7.487793in}}%
\pgfpathcurveto{\pgfqpoint{3.238042in}{7.491911in}}{\pgfqpoint{3.232456in}{7.494225in}}{\pgfqpoint{3.226632in}{7.494225in}}%
\pgfpathcurveto{\pgfqpoint{3.220808in}{7.494225in}}{\pgfqpoint{3.215222in}{7.491911in}}{\pgfqpoint{3.211104in}{7.487793in}}%
\pgfpathcurveto{\pgfqpoint{3.206986in}{7.483675in}}{\pgfqpoint{3.204672in}{7.478089in}}{\pgfqpoint{3.204672in}{7.472265in}}%
\pgfpathcurveto{\pgfqpoint{3.204672in}{7.466441in}}{\pgfqpoint{3.206986in}{7.460854in}}{\pgfqpoint{3.211104in}{7.456736in}}%
\pgfpathcurveto{\pgfqpoint{3.215222in}{7.452618in}}{\pgfqpoint{3.220808in}{7.450304in}}{\pgfqpoint{3.226632in}{7.450304in}}%
\pgfpathlineto{\pgfqpoint{3.226632in}{7.450304in}}%
\pgfpathclose%
\pgfusepath{stroke,fill}%
\end{pgfscope}%
\begin{pgfscope}%
\pgfpathrectangle{\pgfqpoint{1.414736in}{0.548769in}}{\pgfqpoint{3.497747in}{7.253186in}}%
\pgfusepath{clip}%
\pgfsetbuttcap%
\pgfsetroundjoin%
\definecolor{currentfill}{rgb}{0.993248,0.906157,0.143936}%
\pgfsetfillcolor{currentfill}%
\pgfsetlinewidth{1.003750pt}%
\definecolor{currentstroke}{rgb}{0.993248,0.906157,0.143936}%
\pgfsetstrokecolor{currentstroke}%
\pgfsetdash{}{0pt}%
\pgfpathmoveto{\pgfqpoint{3.741850in}{0.856499in}}%
\pgfpathcurveto{\pgfqpoint{3.747674in}{0.856499in}}{\pgfqpoint{3.753260in}{0.858813in}}{\pgfqpoint{3.757378in}{0.862931in}}%
\pgfpathcurveto{\pgfqpoint{3.761497in}{0.867049in}}{\pgfqpoint{3.763810in}{0.872635in}}{\pgfqpoint{3.763810in}{0.878459in}}%
\pgfpathcurveto{\pgfqpoint{3.763810in}{0.884283in}}{\pgfqpoint{3.761497in}{0.889869in}}{\pgfqpoint{3.757378in}{0.893987in}}%
\pgfpathcurveto{\pgfqpoint{3.753260in}{0.898105in}}{\pgfqpoint{3.747674in}{0.900419in}}{\pgfqpoint{3.741850in}{0.900419in}}%
\pgfpathcurveto{\pgfqpoint{3.736026in}{0.900419in}}{\pgfqpoint{3.730440in}{0.898105in}}{\pgfqpoint{3.726322in}{0.893987in}}%
\pgfpathcurveto{\pgfqpoint{3.722204in}{0.889869in}}{\pgfqpoint{3.719890in}{0.884283in}}{\pgfqpoint{3.719890in}{0.878459in}}%
\pgfpathcurveto{\pgfqpoint{3.719890in}{0.872635in}}{\pgfqpoint{3.722204in}{0.867049in}}{\pgfqpoint{3.726322in}{0.862931in}}%
\pgfpathcurveto{\pgfqpoint{3.730440in}{0.858813in}}{\pgfqpoint{3.736026in}{0.856499in}}{\pgfqpoint{3.741850in}{0.856499in}}%
\pgfpathlineto{\pgfqpoint{3.741850in}{0.856499in}}%
\pgfpathclose%
\pgfusepath{stroke,fill}%
\end{pgfscope}%
\begin{pgfscope}%
\pgfpathrectangle{\pgfqpoint{1.414736in}{0.548769in}}{\pgfqpoint{3.497747in}{7.253186in}}%
\pgfusepath{clip}%
\pgfsetbuttcap%
\pgfsetroundjoin%
\definecolor{currentfill}{rgb}{0.668054,0.861999,0.196293}%
\pgfsetfillcolor{currentfill}%
\pgfsetlinewidth{1.003750pt}%
\definecolor{currentstroke}{rgb}{0.668054,0.861999,0.196293}%
\pgfsetstrokecolor{currentstroke}%
\pgfsetdash{}{0pt}%
\pgfpathmoveto{\pgfqpoint{3.402656in}{1.327485in}}%
\pgfpathcurveto{\pgfqpoint{3.408480in}{1.327485in}}{\pgfqpoint{3.414066in}{1.329799in}}{\pgfqpoint{3.418184in}{1.333917in}}%
\pgfpathcurveto{\pgfqpoint{3.422303in}{1.338035in}}{\pgfqpoint{3.424617in}{1.343621in}}{\pgfqpoint{3.424617in}{1.349445in}}%
\pgfpathcurveto{\pgfqpoint{3.424617in}{1.355269in}}{\pgfqpoint{3.422303in}{1.360855in}}{\pgfqpoint{3.418184in}{1.364973in}}%
\pgfpathcurveto{\pgfqpoint{3.414066in}{1.369092in}}{\pgfqpoint{3.408480in}{1.371405in}}{\pgfqpoint{3.402656in}{1.371405in}}%
\pgfpathcurveto{\pgfqpoint{3.396832in}{1.371405in}}{\pgfqpoint{3.391246in}{1.369092in}}{\pgfqpoint{3.387128in}{1.364973in}}%
\pgfpathcurveto{\pgfqpoint{3.383010in}{1.360855in}}{\pgfqpoint{3.380696in}{1.355269in}}{\pgfqpoint{3.380696in}{1.349445in}}%
\pgfpathcurveto{\pgfqpoint{3.380696in}{1.343621in}}{\pgfqpoint{3.383010in}{1.338035in}}{\pgfqpoint{3.387128in}{1.333917in}}%
\pgfpathcurveto{\pgfqpoint{3.391246in}{1.329799in}}{\pgfqpoint{3.396832in}{1.327485in}}{\pgfqpoint{3.402656in}{1.327485in}}%
\pgfpathlineto{\pgfqpoint{3.402656in}{1.327485in}}%
\pgfpathclose%
\pgfusepath{stroke,fill}%
\end{pgfscope}%
\begin{pgfscope}%
\pgfpathrectangle{\pgfqpoint{1.414736in}{0.548769in}}{\pgfqpoint{3.497747in}{7.253186in}}%
\pgfusepath{clip}%
\pgfsetbuttcap%
\pgfsetroundjoin%
\definecolor{currentfill}{rgb}{0.404001,0.800275,0.362552}%
\pgfsetfillcolor{currentfill}%
\pgfsetlinewidth{1.003750pt}%
\definecolor{currentstroke}{rgb}{0.404001,0.800275,0.362552}%
\pgfsetstrokecolor{currentstroke}%
\pgfsetdash{}{0pt}%
\pgfpathmoveto{\pgfqpoint{3.013765in}{1.798471in}}%
\pgfpathcurveto{\pgfqpoint{3.019589in}{1.798471in}}{\pgfqpoint{3.025175in}{1.800785in}}{\pgfqpoint{3.029293in}{1.804903in}}%
\pgfpathcurveto{\pgfqpoint{3.033411in}{1.809021in}}{\pgfqpoint{3.035725in}{1.814607in}}{\pgfqpoint{3.035725in}{1.820431in}}%
\pgfpathcurveto{\pgfqpoint{3.035725in}{1.826255in}}{\pgfqpoint{3.033411in}{1.831841in}}{\pgfqpoint{3.029293in}{1.835960in}}%
\pgfpathcurveto{\pgfqpoint{3.025175in}{1.840078in}}{\pgfqpoint{3.019589in}{1.842392in}}{\pgfqpoint{3.013765in}{1.842392in}}%
\pgfpathcurveto{\pgfqpoint{3.007941in}{1.842392in}}{\pgfqpoint{3.002355in}{1.840078in}}{\pgfqpoint{2.998237in}{1.835960in}}%
\pgfpathcurveto{\pgfqpoint{2.994119in}{1.831841in}}{\pgfqpoint{2.991805in}{1.826255in}}{\pgfqpoint{2.991805in}{1.820431in}}%
\pgfpathcurveto{\pgfqpoint{2.991805in}{1.814607in}}{\pgfqpoint{2.994119in}{1.809021in}}{\pgfqpoint{2.998237in}{1.804903in}}%
\pgfpathcurveto{\pgfqpoint{3.002355in}{1.800785in}}{\pgfqpoint{3.007941in}{1.798471in}}{\pgfqpoint{3.013765in}{1.798471in}}%
\pgfpathlineto{\pgfqpoint{3.013765in}{1.798471in}}%
\pgfpathclose%
\pgfusepath{stroke,fill}%
\end{pgfscope}%
\begin{pgfscope}%
\pgfpathrectangle{\pgfqpoint{1.414736in}{0.548769in}}{\pgfqpoint{3.497747in}{7.253186in}}%
\pgfusepath{clip}%
\pgfsetbuttcap%
\pgfsetroundjoin%
\definecolor{currentfill}{rgb}{0.993248,0.906157,0.143936}%
\pgfsetfillcolor{currentfill}%
\pgfsetlinewidth{1.003750pt}%
\definecolor{currentstroke}{rgb}{0.993248,0.906157,0.143936}%
\pgfsetstrokecolor{currentstroke}%
\pgfsetdash{}{0pt}%
\pgfpathmoveto{\pgfqpoint{3.039277in}{2.269457in}}%
\pgfpathcurveto{\pgfqpoint{3.045100in}{2.269457in}}{\pgfqpoint{3.050687in}{2.271771in}}{\pgfqpoint{3.054805in}{2.275889in}}%
\pgfpathcurveto{\pgfqpoint{3.058923in}{2.280007in}}{\pgfqpoint{3.061237in}{2.285593in}}{\pgfqpoint{3.061237in}{2.291417in}}%
\pgfpathcurveto{\pgfqpoint{3.061237in}{2.297241in}}{\pgfqpoint{3.058923in}{2.302828in}}{\pgfqpoint{3.054805in}{2.306946in}}%
\pgfpathcurveto{\pgfqpoint{3.050687in}{2.311064in}}{\pgfqpoint{3.045100in}{2.313378in}}{\pgfqpoint{3.039277in}{2.313378in}}%
\pgfpathcurveto{\pgfqpoint{3.033453in}{2.313378in}}{\pgfqpoint{3.027866in}{2.311064in}}{\pgfqpoint{3.023748in}{2.306946in}}%
\pgfpathcurveto{\pgfqpoint{3.019630in}{2.302828in}}{\pgfqpoint{3.017316in}{2.297241in}}{\pgfqpoint{3.017316in}{2.291417in}}%
\pgfpathcurveto{\pgfqpoint{3.017316in}{2.285593in}}{\pgfqpoint{3.019630in}{2.280007in}}{\pgfqpoint{3.023748in}{2.275889in}}%
\pgfpathcurveto{\pgfqpoint{3.027866in}{2.271771in}}{\pgfqpoint{3.033453in}{2.269457in}}{\pgfqpoint{3.039277in}{2.269457in}}%
\pgfpathlineto{\pgfqpoint{3.039277in}{2.269457in}}%
\pgfpathclose%
\pgfusepath{stroke,fill}%
\end{pgfscope}%
\begin{pgfscope}%
\pgfpathrectangle{\pgfqpoint{1.414736in}{0.548769in}}{\pgfqpoint{3.497747in}{7.253186in}}%
\pgfusepath{clip}%
\pgfsetbuttcap%
\pgfsetroundjoin%
\definecolor{currentfill}{rgb}{0.281477,0.755203,0.432552}%
\pgfsetfillcolor{currentfill}%
\pgfsetlinewidth{1.003750pt}%
\definecolor{currentstroke}{rgb}{0.281477,0.755203,0.432552}%
\pgfsetstrokecolor{currentstroke}%
\pgfsetdash{}{0pt}%
\pgfpathmoveto{\pgfqpoint{3.874831in}{2.740443in}}%
\pgfpathcurveto{\pgfqpoint{3.880655in}{2.740443in}}{\pgfqpoint{3.886241in}{2.742757in}}{\pgfqpoint{3.890360in}{2.746875in}}%
\pgfpathcurveto{\pgfqpoint{3.894478in}{2.750993in}}{\pgfqpoint{3.896792in}{2.756580in}}{\pgfqpoint{3.896792in}{2.762404in}}%
\pgfpathcurveto{\pgfqpoint{3.896792in}{2.768227in}}{\pgfqpoint{3.894478in}{2.773814in}}{\pgfqpoint{3.890360in}{2.777932in}}%
\pgfpathcurveto{\pgfqpoint{3.886241in}{2.782050in}}{\pgfqpoint{3.880655in}{2.784364in}}{\pgfqpoint{3.874831in}{2.784364in}}%
\pgfpathcurveto{\pgfqpoint{3.869007in}{2.784364in}}{\pgfqpoint{3.863421in}{2.782050in}}{\pgfqpoint{3.859303in}{2.777932in}}%
\pgfpathcurveto{\pgfqpoint{3.855185in}{2.773814in}}{\pgfqpoint{3.852871in}{2.768227in}}{\pgfqpoint{3.852871in}{2.762404in}}%
\pgfpathcurveto{\pgfqpoint{3.852871in}{2.756580in}}{\pgfqpoint{3.855185in}{2.750993in}}{\pgfqpoint{3.859303in}{2.746875in}}%
\pgfpathcurveto{\pgfqpoint{3.863421in}{2.742757in}}{\pgfqpoint{3.869007in}{2.740443in}}{\pgfqpoint{3.874831in}{2.740443in}}%
\pgfpathlineto{\pgfqpoint{3.874831in}{2.740443in}}%
\pgfpathclose%
\pgfusepath{stroke,fill}%
\end{pgfscope}%
\begin{pgfscope}%
\pgfpathrectangle{\pgfqpoint{1.414736in}{0.548769in}}{\pgfqpoint{3.497747in}{7.253186in}}%
\pgfusepath{clip}%
\pgfsetbuttcap%
\pgfsetroundjoin%
\definecolor{currentfill}{rgb}{0.616293,0.852709,0.230052}%
\pgfsetfillcolor{currentfill}%
\pgfsetlinewidth{1.003750pt}%
\definecolor{currentstroke}{rgb}{0.616293,0.852709,0.230052}%
\pgfsetstrokecolor{currentstroke}%
\pgfsetdash{}{0pt}%
\pgfpathmoveto{\pgfqpoint{3.227701in}{3.211429in}}%
\pgfpathcurveto{\pgfqpoint{3.233525in}{3.211429in}}{\pgfqpoint{3.239111in}{3.213743in}}{\pgfqpoint{3.243229in}{3.217861in}}%
\pgfpathcurveto{\pgfqpoint{3.247347in}{3.221980in}}{\pgfqpoint{3.249661in}{3.227566in}}{\pgfqpoint{3.249661in}{3.233390in}}%
\pgfpathcurveto{\pgfqpoint{3.249661in}{3.239214in}}{\pgfqpoint{3.247347in}{3.244800in}}{\pgfqpoint{3.243229in}{3.248918in}}%
\pgfpathcurveto{\pgfqpoint{3.239111in}{3.253036in}}{\pgfqpoint{3.233525in}{3.255350in}}{\pgfqpoint{3.227701in}{3.255350in}}%
\pgfpathcurveto{\pgfqpoint{3.221877in}{3.255350in}}{\pgfqpoint{3.216291in}{3.253036in}}{\pgfqpoint{3.212172in}{3.248918in}}%
\pgfpathcurveto{\pgfqpoint{3.208054in}{3.244800in}}{\pgfqpoint{3.205740in}{3.239214in}}{\pgfqpoint{3.205740in}{3.233390in}}%
\pgfpathcurveto{\pgfqpoint{3.205740in}{3.227566in}}{\pgfqpoint{3.208054in}{3.221980in}}{\pgfqpoint{3.212172in}{3.217861in}}%
\pgfpathcurveto{\pgfqpoint{3.216291in}{3.213743in}}{\pgfqpoint{3.221877in}{3.211429in}}{\pgfqpoint{3.227701in}{3.211429in}}%
\pgfpathlineto{\pgfqpoint{3.227701in}{3.211429in}}%
\pgfpathclose%
\pgfusepath{stroke,fill}%
\end{pgfscope}%
\begin{pgfscope}%
\pgfpathrectangle{\pgfqpoint{1.414736in}{0.548769in}}{\pgfqpoint{3.497747in}{7.253186in}}%
\pgfusepath{clip}%
\pgfsetbuttcap%
\pgfsetroundjoin%
\definecolor{currentfill}{rgb}{0.657642,0.860219,0.203082}%
\pgfsetfillcolor{currentfill}%
\pgfsetlinewidth{1.003750pt}%
\definecolor{currentstroke}{rgb}{0.657642,0.860219,0.203082}%
\pgfsetstrokecolor{currentstroke}%
\pgfsetdash{}{0pt}%
\pgfpathmoveto{\pgfqpoint{3.470547in}{3.682415in}}%
\pgfpathcurveto{\pgfqpoint{3.476371in}{3.682415in}}{\pgfqpoint{3.481957in}{3.684729in}}{\pgfqpoint{3.486075in}{3.688847in}}%
\pgfpathcurveto{\pgfqpoint{3.490193in}{3.692966in}}{\pgfqpoint{3.492507in}{3.698552in}}{\pgfqpoint{3.492507in}{3.704376in}}%
\pgfpathcurveto{\pgfqpoint{3.492507in}{3.710200in}}{\pgfqpoint{3.490193in}{3.715786in}}{\pgfqpoint{3.486075in}{3.719904in}}%
\pgfpathcurveto{\pgfqpoint{3.481957in}{3.724022in}}{\pgfqpoint{3.476371in}{3.726336in}}{\pgfqpoint{3.470547in}{3.726336in}}%
\pgfpathcurveto{\pgfqpoint{3.464723in}{3.726336in}}{\pgfqpoint{3.459137in}{3.724022in}}{\pgfqpoint{3.455018in}{3.719904in}}%
\pgfpathcurveto{\pgfqpoint{3.450900in}{3.715786in}}{\pgfqpoint{3.448586in}{3.710200in}}{\pgfqpoint{3.448586in}{3.704376in}}%
\pgfpathcurveto{\pgfqpoint{3.448586in}{3.698552in}}{\pgfqpoint{3.450900in}{3.692966in}}{\pgfqpoint{3.455018in}{3.688847in}}%
\pgfpathcurveto{\pgfqpoint{3.459137in}{3.684729in}}{\pgfqpoint{3.464723in}{3.682415in}}{\pgfqpoint{3.470547in}{3.682415in}}%
\pgfpathlineto{\pgfqpoint{3.470547in}{3.682415in}}%
\pgfpathclose%
\pgfusepath{stroke,fill}%
\end{pgfscope}%
\begin{pgfscope}%
\pgfpathrectangle{\pgfqpoint{1.414736in}{0.548769in}}{\pgfqpoint{3.497747in}{7.253186in}}%
\pgfusepath{clip}%
\pgfsetbuttcap%
\pgfsetroundjoin%
\definecolor{currentfill}{rgb}{0.845561,0.887322,0.099702}%
\pgfsetfillcolor{currentfill}%
\pgfsetlinewidth{1.003750pt}%
\definecolor{currentstroke}{rgb}{0.845561,0.887322,0.099702}%
\pgfsetstrokecolor{currentstroke}%
\pgfsetdash{}{0pt}%
\pgfpathmoveto{\pgfqpoint{3.223262in}{4.153402in}}%
\pgfpathcurveto{\pgfqpoint{3.229086in}{4.153402in}}{\pgfqpoint{3.234673in}{4.155715in}}{\pgfqpoint{3.238791in}{4.159834in}}%
\pgfpathcurveto{\pgfqpoint{3.242909in}{4.163952in}}{\pgfqpoint{3.245223in}{4.169538in}}{\pgfqpoint{3.245223in}{4.175362in}}%
\pgfpathcurveto{\pgfqpoint{3.245223in}{4.181186in}}{\pgfqpoint{3.242909in}{4.186772in}}{\pgfqpoint{3.238791in}{4.190890in}}%
\pgfpathcurveto{\pgfqpoint{3.234673in}{4.195008in}}{\pgfqpoint{3.229086in}{4.197322in}}{\pgfqpoint{3.223262in}{4.197322in}}%
\pgfpathcurveto{\pgfqpoint{3.217439in}{4.197322in}}{\pgfqpoint{3.211852in}{4.195008in}}{\pgfqpoint{3.207734in}{4.190890in}}%
\pgfpathcurveto{\pgfqpoint{3.203616in}{4.186772in}}{\pgfqpoint{3.201302in}{4.181186in}}{\pgfqpoint{3.201302in}{4.175362in}}%
\pgfpathcurveto{\pgfqpoint{3.201302in}{4.169538in}}{\pgfqpoint{3.203616in}{4.163952in}}{\pgfqpoint{3.207734in}{4.159834in}}%
\pgfpathcurveto{\pgfqpoint{3.211852in}{4.155715in}}{\pgfqpoint{3.217439in}{4.153402in}}{\pgfqpoint{3.223262in}{4.153402in}}%
\pgfpathlineto{\pgfqpoint{3.223262in}{4.153402in}}%
\pgfpathclose%
\pgfusepath{stroke,fill}%
\end{pgfscope}%
\begin{pgfscope}%
\pgfpathrectangle{\pgfqpoint{1.414736in}{0.548769in}}{\pgfqpoint{3.497747in}{7.253186in}}%
\pgfusepath{clip}%
\pgfsetbuttcap%
\pgfsetroundjoin%
\definecolor{currentfill}{rgb}{0.945636,0.899815,0.112838}%
\pgfsetfillcolor{currentfill}%
\pgfsetlinewidth{1.003750pt}%
\definecolor{currentstroke}{rgb}{0.945636,0.899815,0.112838}%
\pgfsetstrokecolor{currentstroke}%
\pgfsetdash{}{0pt}%
\pgfpathmoveto{\pgfqpoint{3.960652in}{4.624388in}}%
\pgfpathcurveto{\pgfqpoint{3.966476in}{4.624388in}}{\pgfqpoint{3.972062in}{4.626702in}}{\pgfqpoint{3.976180in}{4.630820in}}%
\pgfpathcurveto{\pgfqpoint{3.980298in}{4.634938in}}{\pgfqpoint{3.982612in}{4.640524in}}{\pgfqpoint{3.982612in}{4.646348in}}%
\pgfpathcurveto{\pgfqpoint{3.982612in}{4.652172in}}{\pgfqpoint{3.980298in}{4.657758in}}{\pgfqpoint{3.976180in}{4.661876in}}%
\pgfpathcurveto{\pgfqpoint{3.972062in}{4.665994in}}{\pgfqpoint{3.966476in}{4.668308in}}{\pgfqpoint{3.960652in}{4.668308in}}%
\pgfpathcurveto{\pgfqpoint{3.954828in}{4.668308in}}{\pgfqpoint{3.949242in}{4.665994in}}{\pgfqpoint{3.945123in}{4.661876in}}%
\pgfpathcurveto{\pgfqpoint{3.941005in}{4.657758in}}{\pgfqpoint{3.938691in}{4.652172in}}{\pgfqpoint{3.938691in}{4.646348in}}%
\pgfpathcurveto{\pgfqpoint{3.938691in}{4.640524in}}{\pgfqpoint{3.941005in}{4.634938in}}{\pgfqpoint{3.945123in}{4.630820in}}%
\pgfpathcurveto{\pgfqpoint{3.949242in}{4.626702in}}{\pgfqpoint{3.954828in}{4.624388in}}{\pgfqpoint{3.960652in}{4.624388in}}%
\pgfpathlineto{\pgfqpoint{3.960652in}{4.624388in}}%
\pgfpathclose%
\pgfusepath{stroke,fill}%
\end{pgfscope}%
\begin{pgfscope}%
\pgfpathrectangle{\pgfqpoint{1.414736in}{0.548769in}}{\pgfqpoint{3.497747in}{7.253186in}}%
\pgfusepath{clip}%
\pgfsetbuttcap%
\pgfsetroundjoin%
\definecolor{currentfill}{rgb}{0.248629,0.278775,0.534556}%
\pgfsetfillcolor{currentfill}%
\pgfsetlinewidth{1.003750pt}%
\definecolor{currentstroke}{rgb}{0.248629,0.278775,0.534556}%
\pgfsetstrokecolor{currentstroke}%
\pgfsetdash{}{0pt}%
\pgfpathmoveto{\pgfqpoint{4.018987in}{5.095374in}}%
\pgfpathcurveto{\pgfqpoint{4.024811in}{5.095374in}}{\pgfqpoint{4.030397in}{5.097688in}}{\pgfqpoint{4.034515in}{5.101806in}}%
\pgfpathcurveto{\pgfqpoint{4.038633in}{5.105924in}}{\pgfqpoint{4.040947in}{5.111510in}}{\pgfqpoint{4.040947in}{5.117334in}}%
\pgfpathcurveto{\pgfqpoint{4.040947in}{5.123158in}}{\pgfqpoint{4.038633in}{5.128744in}}{\pgfqpoint{4.034515in}{5.132862in}}%
\pgfpathcurveto{\pgfqpoint{4.030397in}{5.136980in}}{\pgfqpoint{4.024811in}{5.139294in}}{\pgfqpoint{4.018987in}{5.139294in}}%
\pgfpathcurveto{\pgfqpoint{4.013163in}{5.139294in}}{\pgfqpoint{4.007577in}{5.136980in}}{\pgfqpoint{4.003459in}{5.132862in}}%
\pgfpathcurveto{\pgfqpoint{3.999341in}{5.128744in}}{\pgfqpoint{3.997027in}{5.123158in}}{\pgfqpoint{3.997027in}{5.117334in}}%
\pgfpathcurveto{\pgfqpoint{3.997027in}{5.111510in}}{\pgfqpoint{3.999341in}{5.105924in}}{\pgfqpoint{4.003459in}{5.101806in}}%
\pgfpathcurveto{\pgfqpoint{4.007577in}{5.097688in}}{\pgfqpoint{4.013163in}{5.095374in}}{\pgfqpoint{4.018987in}{5.095374in}}%
\pgfpathlineto{\pgfqpoint{4.018987in}{5.095374in}}%
\pgfpathclose%
\pgfusepath{stroke,fill}%
\end{pgfscope}%
\begin{pgfscope}%
\pgfpathrectangle{\pgfqpoint{1.414736in}{0.548769in}}{\pgfqpoint{3.497747in}{7.253186in}}%
\pgfusepath{clip}%
\pgfsetbuttcap%
\pgfsetroundjoin%
\definecolor{currentfill}{rgb}{0.172719,0.448791,0.557885}%
\pgfsetfillcolor{currentfill}%
\pgfsetlinewidth{1.003750pt}%
\definecolor{currentstroke}{rgb}{0.172719,0.448791,0.557885}%
\pgfsetstrokecolor{currentstroke}%
\pgfsetdash{}{0pt}%
\pgfpathmoveto{\pgfqpoint{3.672450in}{5.566360in}}%
\pgfpathcurveto{\pgfqpoint{3.678274in}{5.566360in}}{\pgfqpoint{3.683860in}{5.568674in}}{\pgfqpoint{3.687978in}{5.572792in}}%
\pgfpathcurveto{\pgfqpoint{3.692096in}{5.576910in}}{\pgfqpoint{3.694410in}{5.582496in}}{\pgfqpoint{3.694410in}{5.588320in}}%
\pgfpathcurveto{\pgfqpoint{3.694410in}{5.594144in}}{\pgfqpoint{3.692096in}{5.599730in}}{\pgfqpoint{3.687978in}{5.603848in}}%
\pgfpathcurveto{\pgfqpoint{3.683860in}{5.607967in}}{\pgfqpoint{3.678274in}{5.610280in}}{\pgfqpoint{3.672450in}{5.610280in}}%
\pgfpathcurveto{\pgfqpoint{3.666626in}{5.610280in}}{\pgfqpoint{3.661040in}{5.607967in}}{\pgfqpoint{3.656922in}{5.603848in}}%
\pgfpathcurveto{\pgfqpoint{3.652804in}{5.599730in}}{\pgfqpoint{3.650490in}{5.594144in}}{\pgfqpoint{3.650490in}{5.588320in}}%
\pgfpathcurveto{\pgfqpoint{3.650490in}{5.582496in}}{\pgfqpoint{3.652804in}{5.576910in}}{\pgfqpoint{3.656922in}{5.572792in}}%
\pgfpathcurveto{\pgfqpoint{3.661040in}{5.568674in}}{\pgfqpoint{3.666626in}{5.566360in}}{\pgfqpoint{3.672450in}{5.566360in}}%
\pgfpathlineto{\pgfqpoint{3.672450in}{5.566360in}}%
\pgfpathclose%
\pgfusepath{stroke,fill}%
\end{pgfscope}%
\begin{pgfscope}%
\pgfpathrectangle{\pgfqpoint{1.414736in}{0.548769in}}{\pgfqpoint{3.497747in}{7.253186in}}%
\pgfusepath{clip}%
\pgfsetbuttcap%
\pgfsetroundjoin%
\definecolor{currentfill}{rgb}{0.266941,0.748751,0.440573}%
\pgfsetfillcolor{currentfill}%
\pgfsetlinewidth{1.003750pt}%
\definecolor{currentstroke}{rgb}{0.266941,0.748751,0.440573}%
\pgfsetstrokecolor{currentstroke}%
\pgfsetdash{}{0pt}%
\pgfpathmoveto{\pgfqpoint{3.461396in}{6.037346in}}%
\pgfpathcurveto{\pgfqpoint{3.467220in}{6.037346in}}{\pgfqpoint{3.472806in}{6.039660in}}{\pgfqpoint{3.476924in}{6.043778in}}%
\pgfpathcurveto{\pgfqpoint{3.481042in}{6.047896in}}{\pgfqpoint{3.483356in}{6.053482in}}{\pgfqpoint{3.483356in}{6.059306in}}%
\pgfpathcurveto{\pgfqpoint{3.483356in}{6.065130in}}{\pgfqpoint{3.481042in}{6.070716in}}{\pgfqpoint{3.476924in}{6.074835in}}%
\pgfpathcurveto{\pgfqpoint{3.472806in}{6.078953in}}{\pgfqpoint{3.467220in}{6.081267in}}{\pgfqpoint{3.461396in}{6.081267in}}%
\pgfpathcurveto{\pgfqpoint{3.455572in}{6.081267in}}{\pgfqpoint{3.449986in}{6.078953in}}{\pgfqpoint{3.445867in}{6.074835in}}%
\pgfpathcurveto{\pgfqpoint{3.441749in}{6.070716in}}{\pgfqpoint{3.439435in}{6.065130in}}{\pgfqpoint{3.439435in}{6.059306in}}%
\pgfpathcurveto{\pgfqpoint{3.439435in}{6.053482in}}{\pgfqpoint{3.441749in}{6.047896in}}{\pgfqpoint{3.445867in}{6.043778in}}%
\pgfpathcurveto{\pgfqpoint{3.449986in}{6.039660in}}{\pgfqpoint{3.455572in}{6.037346in}}{\pgfqpoint{3.461396in}{6.037346in}}%
\pgfpathlineto{\pgfqpoint{3.461396in}{6.037346in}}%
\pgfpathclose%
\pgfusepath{stroke,fill}%
\end{pgfscope}%
\begin{pgfscope}%
\pgfpathrectangle{\pgfqpoint{1.414736in}{0.548769in}}{\pgfqpoint{3.497747in}{7.253186in}}%
\pgfusepath{clip}%
\pgfsetbuttcap%
\pgfsetroundjoin%
\definecolor{currentfill}{rgb}{0.935904,0.898570,0.108131}%
\pgfsetfillcolor{currentfill}%
\pgfsetlinewidth{1.003750pt}%
\definecolor{currentstroke}{rgb}{0.935904,0.898570,0.108131}%
\pgfsetstrokecolor{currentstroke}%
\pgfsetdash{}{0pt}%
\pgfpathmoveto{\pgfqpoint{3.708731in}{6.508332in}}%
\pgfpathcurveto{\pgfqpoint{3.714555in}{6.508332in}}{\pgfqpoint{3.720141in}{6.510646in}}{\pgfqpoint{3.724259in}{6.514764in}}%
\pgfpathcurveto{\pgfqpoint{3.728377in}{6.518882in}}{\pgfqpoint{3.730691in}{6.524468in}}{\pgfqpoint{3.730691in}{6.530292in}}%
\pgfpathcurveto{\pgfqpoint{3.730691in}{6.536116in}}{\pgfqpoint{3.728377in}{6.541702in}}{\pgfqpoint{3.724259in}{6.545821in}}%
\pgfpathcurveto{\pgfqpoint{3.720141in}{6.549939in}}{\pgfqpoint{3.714555in}{6.552253in}}{\pgfqpoint{3.708731in}{6.552253in}}%
\pgfpathcurveto{\pgfqpoint{3.702907in}{6.552253in}}{\pgfqpoint{3.697321in}{6.549939in}}{\pgfqpoint{3.693203in}{6.545821in}}%
\pgfpathcurveto{\pgfqpoint{3.689085in}{6.541702in}}{\pgfqpoint{3.686771in}{6.536116in}}{\pgfqpoint{3.686771in}{6.530292in}}%
\pgfpathcurveto{\pgfqpoint{3.686771in}{6.524468in}}{\pgfqpoint{3.689085in}{6.518882in}}{\pgfqpoint{3.693203in}{6.514764in}}%
\pgfpathcurveto{\pgfqpoint{3.697321in}{6.510646in}}{\pgfqpoint{3.702907in}{6.508332in}}{\pgfqpoint{3.708731in}{6.508332in}}%
\pgfpathlineto{\pgfqpoint{3.708731in}{6.508332in}}%
\pgfpathclose%
\pgfusepath{stroke,fill}%
\end{pgfscope}%
\begin{pgfscope}%
\pgfpathrectangle{\pgfqpoint{1.414736in}{0.548769in}}{\pgfqpoint{3.497747in}{7.253186in}}%
\pgfusepath{clip}%
\pgfsetbuttcap%
\pgfsetroundjoin%
\definecolor{currentfill}{rgb}{0.926106,0.897330,0.104071}%
\pgfsetfillcolor{currentfill}%
\pgfsetlinewidth{1.003750pt}%
\definecolor{currentstroke}{rgb}{0.926106,0.897330,0.104071}%
\pgfsetstrokecolor{currentstroke}%
\pgfsetdash{}{0pt}%
\pgfpathmoveto{\pgfqpoint{4.080520in}{6.979318in}}%
\pgfpathcurveto{\pgfqpoint{4.086344in}{6.979318in}}{\pgfqpoint{4.091930in}{6.981632in}}{\pgfqpoint{4.096048in}{6.985750in}}%
\pgfpathcurveto{\pgfqpoint{4.100167in}{6.989868in}}{\pgfqpoint{4.102480in}{6.995455in}}{\pgfqpoint{4.102480in}{7.001278in}}%
\pgfpathcurveto{\pgfqpoint{4.102480in}{7.007102in}}{\pgfqpoint{4.100167in}{7.012689in}}{\pgfqpoint{4.096048in}{7.016807in}}%
\pgfpathcurveto{\pgfqpoint{4.091930in}{7.020925in}}{\pgfqpoint{4.086344in}{7.023239in}}{\pgfqpoint{4.080520in}{7.023239in}}%
\pgfpathcurveto{\pgfqpoint{4.074696in}{7.023239in}}{\pgfqpoint{4.069110in}{7.020925in}}{\pgfqpoint{4.064992in}{7.016807in}}%
\pgfpathcurveto{\pgfqpoint{4.060874in}{7.012689in}}{\pgfqpoint{4.058560in}{7.007102in}}{\pgfqpoint{4.058560in}{7.001278in}}%
\pgfpathcurveto{\pgfqpoint{4.058560in}{6.995455in}}{\pgfqpoint{4.060874in}{6.989868in}}{\pgfqpoint{4.064992in}{6.985750in}}%
\pgfpathcurveto{\pgfqpoint{4.069110in}{6.981632in}}{\pgfqpoint{4.074696in}{6.979318in}}{\pgfqpoint{4.080520in}{6.979318in}}%
\pgfpathlineto{\pgfqpoint{4.080520in}{6.979318in}}%
\pgfpathclose%
\pgfusepath{stroke,fill}%
\end{pgfscope}%
\begin{pgfscope}%
\pgfpathrectangle{\pgfqpoint{1.414736in}{0.548769in}}{\pgfqpoint{3.497747in}{7.253186in}}%
\pgfusepath{clip}%
\pgfsetbuttcap%
\pgfsetroundjoin%
\definecolor{currentfill}{rgb}{0.993248,0.906157,0.143936}%
\pgfsetfillcolor{currentfill}%
\pgfsetlinewidth{1.003750pt}%
\definecolor{currentstroke}{rgb}{0.993248,0.906157,0.143936}%
\pgfsetstrokecolor{currentstroke}%
\pgfsetdash{}{0pt}%
\pgfpathmoveto{\pgfqpoint{3.097506in}{7.450304in}}%
\pgfpathcurveto{\pgfqpoint{3.103330in}{7.450304in}}{\pgfqpoint{3.108917in}{7.452618in}}{\pgfqpoint{3.113035in}{7.456736in}}%
\pgfpathcurveto{\pgfqpoint{3.117153in}{7.460854in}}{\pgfqpoint{3.119467in}{7.466441in}}{\pgfqpoint{3.119467in}{7.472265in}}%
\pgfpathcurveto{\pgfqpoint{3.119467in}{7.478089in}}{\pgfqpoint{3.117153in}{7.483675in}}{\pgfqpoint{3.113035in}{7.487793in}}%
\pgfpathcurveto{\pgfqpoint{3.108917in}{7.491911in}}{\pgfqpoint{3.103330in}{7.494225in}}{\pgfqpoint{3.097506in}{7.494225in}}%
\pgfpathcurveto{\pgfqpoint{3.091683in}{7.494225in}}{\pgfqpoint{3.086096in}{7.491911in}}{\pgfqpoint{3.081978in}{7.487793in}}%
\pgfpathcurveto{\pgfqpoint{3.077860in}{7.483675in}}{\pgfqpoint{3.075546in}{7.478089in}}{\pgfqpoint{3.075546in}{7.472265in}}%
\pgfpathcurveto{\pgfqpoint{3.075546in}{7.466441in}}{\pgfqpoint{3.077860in}{7.460854in}}{\pgfqpoint{3.081978in}{7.456736in}}%
\pgfpathcurveto{\pgfqpoint{3.086096in}{7.452618in}}{\pgfqpoint{3.091683in}{7.450304in}}{\pgfqpoint{3.097506in}{7.450304in}}%
\pgfpathlineto{\pgfqpoint{3.097506in}{7.450304in}}%
\pgfpathclose%
\pgfusepath{stroke,fill}%
\end{pgfscope}%
\begin{pgfscope}%
\pgfpathrectangle{\pgfqpoint{1.414736in}{0.548769in}}{\pgfqpoint{3.497747in}{7.253186in}}%
\pgfusepath{clip}%
\pgfsetbuttcap%
\pgfsetroundjoin%
\definecolor{currentfill}{rgb}{0.993248,0.906157,0.143936}%
\pgfsetfillcolor{currentfill}%
\pgfsetlinewidth{1.003750pt}%
\definecolor{currentstroke}{rgb}{0.993248,0.906157,0.143936}%
\pgfsetstrokecolor{currentstroke}%
\pgfsetdash{}{0pt}%
\pgfpathmoveto{\pgfqpoint{3.422331in}{0.856499in}}%
\pgfpathcurveto{\pgfqpoint{3.428155in}{0.856499in}}{\pgfqpoint{3.433741in}{0.858813in}}{\pgfqpoint{3.437860in}{0.862931in}}%
\pgfpathcurveto{\pgfqpoint{3.441978in}{0.867049in}}{\pgfqpoint{3.444292in}{0.872635in}}{\pgfqpoint{3.444292in}{0.878459in}}%
\pgfpathcurveto{\pgfqpoint{3.444292in}{0.884283in}}{\pgfqpoint{3.441978in}{0.889869in}}{\pgfqpoint{3.437860in}{0.893987in}}%
\pgfpathcurveto{\pgfqpoint{3.433741in}{0.898105in}}{\pgfqpoint{3.428155in}{0.900419in}}{\pgfqpoint{3.422331in}{0.900419in}}%
\pgfpathcurveto{\pgfqpoint{3.416507in}{0.900419in}}{\pgfqpoint{3.410921in}{0.898105in}}{\pgfqpoint{3.406803in}{0.893987in}}%
\pgfpathcurveto{\pgfqpoint{3.402685in}{0.889869in}}{\pgfqpoint{3.400371in}{0.884283in}}{\pgfqpoint{3.400371in}{0.878459in}}%
\pgfpathcurveto{\pgfqpoint{3.400371in}{0.872635in}}{\pgfqpoint{3.402685in}{0.867049in}}{\pgfqpoint{3.406803in}{0.862931in}}%
\pgfpathcurveto{\pgfqpoint{3.410921in}{0.858813in}}{\pgfqpoint{3.416507in}{0.856499in}}{\pgfqpoint{3.422331in}{0.856499in}}%
\pgfpathlineto{\pgfqpoint{3.422331in}{0.856499in}}%
\pgfpathclose%
\pgfusepath{stroke,fill}%
\end{pgfscope}%
\begin{pgfscope}%
\pgfpathrectangle{\pgfqpoint{1.414736in}{0.548769in}}{\pgfqpoint{3.497747in}{7.253186in}}%
\pgfusepath{clip}%
\pgfsetbuttcap%
\pgfsetroundjoin%
\definecolor{currentfill}{rgb}{0.525776,0.833491,0.288127}%
\pgfsetfillcolor{currentfill}%
\pgfsetlinewidth{1.003750pt}%
\definecolor{currentstroke}{rgb}{0.525776,0.833491,0.288127}%
\pgfsetstrokecolor{currentstroke}%
\pgfsetdash{}{0pt}%
\pgfpathmoveto{\pgfqpoint{3.590720in}{1.327485in}}%
\pgfpathcurveto{\pgfqpoint{3.596544in}{1.327485in}}{\pgfqpoint{3.602130in}{1.329799in}}{\pgfqpoint{3.606248in}{1.333917in}}%
\pgfpathcurveto{\pgfqpoint{3.610366in}{1.338035in}}{\pgfqpoint{3.612680in}{1.343621in}}{\pgfqpoint{3.612680in}{1.349445in}}%
\pgfpathcurveto{\pgfqpoint{3.612680in}{1.355269in}}{\pgfqpoint{3.610366in}{1.360855in}}{\pgfqpoint{3.606248in}{1.364973in}}%
\pgfpathcurveto{\pgfqpoint{3.602130in}{1.369092in}}{\pgfqpoint{3.596544in}{1.371405in}}{\pgfqpoint{3.590720in}{1.371405in}}%
\pgfpathcurveto{\pgfqpoint{3.584896in}{1.371405in}}{\pgfqpoint{3.579310in}{1.369092in}}{\pgfqpoint{3.575192in}{1.364973in}}%
\pgfpathcurveto{\pgfqpoint{3.571074in}{1.360855in}}{\pgfqpoint{3.568760in}{1.355269in}}{\pgfqpoint{3.568760in}{1.349445in}}%
\pgfpathcurveto{\pgfqpoint{3.568760in}{1.343621in}}{\pgfqpoint{3.571074in}{1.338035in}}{\pgfqpoint{3.575192in}{1.333917in}}%
\pgfpathcurveto{\pgfqpoint{3.579310in}{1.329799in}}{\pgfqpoint{3.584896in}{1.327485in}}{\pgfqpoint{3.590720in}{1.327485in}}%
\pgfpathlineto{\pgfqpoint{3.590720in}{1.327485in}}%
\pgfpathclose%
\pgfusepath{stroke,fill}%
\end{pgfscope}%
\begin{pgfscope}%
\pgfpathrectangle{\pgfqpoint{1.414736in}{0.548769in}}{\pgfqpoint{3.497747in}{7.253186in}}%
\pgfusepath{clip}%
\pgfsetbuttcap%
\pgfsetroundjoin%
\definecolor{currentfill}{rgb}{0.430983,0.808473,0.346476}%
\pgfsetfillcolor{currentfill}%
\pgfsetlinewidth{1.003750pt}%
\definecolor{currentstroke}{rgb}{0.430983,0.808473,0.346476}%
\pgfsetstrokecolor{currentstroke}%
\pgfsetdash{}{0pt}%
\pgfpathmoveto{\pgfqpoint{3.432217in}{1.798471in}}%
\pgfpathcurveto{\pgfqpoint{3.438041in}{1.798471in}}{\pgfqpoint{3.443627in}{1.800785in}}{\pgfqpoint{3.447745in}{1.804903in}}%
\pgfpathcurveto{\pgfqpoint{3.451863in}{1.809021in}}{\pgfqpoint{3.454177in}{1.814607in}}{\pgfqpoint{3.454177in}{1.820431in}}%
\pgfpathcurveto{\pgfqpoint{3.454177in}{1.826255in}}{\pgfqpoint{3.451863in}{1.831841in}}{\pgfqpoint{3.447745in}{1.835960in}}%
\pgfpathcurveto{\pgfqpoint{3.443627in}{1.840078in}}{\pgfqpoint{3.438041in}{1.842392in}}{\pgfqpoint{3.432217in}{1.842392in}}%
\pgfpathcurveto{\pgfqpoint{3.426393in}{1.842392in}}{\pgfqpoint{3.420807in}{1.840078in}}{\pgfqpoint{3.416689in}{1.835960in}}%
\pgfpathcurveto{\pgfqpoint{3.412571in}{1.831841in}}{\pgfqpoint{3.410257in}{1.826255in}}{\pgfqpoint{3.410257in}{1.820431in}}%
\pgfpathcurveto{\pgfqpoint{3.410257in}{1.814607in}}{\pgfqpoint{3.412571in}{1.809021in}}{\pgfqpoint{3.416689in}{1.804903in}}%
\pgfpathcurveto{\pgfqpoint{3.420807in}{1.800785in}}{\pgfqpoint{3.426393in}{1.798471in}}{\pgfqpoint{3.432217in}{1.798471in}}%
\pgfpathlineto{\pgfqpoint{3.432217in}{1.798471in}}%
\pgfpathclose%
\pgfusepath{stroke,fill}%
\end{pgfscope}%
\begin{pgfscope}%
\pgfpathrectangle{\pgfqpoint{1.414736in}{0.548769in}}{\pgfqpoint{3.497747in}{7.253186in}}%
\pgfusepath{clip}%
\pgfsetbuttcap%
\pgfsetroundjoin%
\definecolor{currentfill}{rgb}{0.974417,0.903590,0.130215}%
\pgfsetfillcolor{currentfill}%
\pgfsetlinewidth{1.003750pt}%
\definecolor{currentstroke}{rgb}{0.974417,0.903590,0.130215}%
\pgfsetstrokecolor{currentstroke}%
\pgfsetdash{}{0pt}%
\pgfpathmoveto{\pgfqpoint{3.129114in}{2.269457in}}%
\pgfpathcurveto{\pgfqpoint{3.134938in}{2.269457in}}{\pgfqpoint{3.140524in}{2.271771in}}{\pgfqpoint{3.144642in}{2.275889in}}%
\pgfpathcurveto{\pgfqpoint{3.148760in}{2.280007in}}{\pgfqpoint{3.151074in}{2.285593in}}{\pgfqpoint{3.151074in}{2.291417in}}%
\pgfpathcurveto{\pgfqpoint{3.151074in}{2.297241in}}{\pgfqpoint{3.148760in}{2.302828in}}{\pgfqpoint{3.144642in}{2.306946in}}%
\pgfpathcurveto{\pgfqpoint{3.140524in}{2.311064in}}{\pgfqpoint{3.134938in}{2.313378in}}{\pgfqpoint{3.129114in}{2.313378in}}%
\pgfpathcurveto{\pgfqpoint{3.123290in}{2.313378in}}{\pgfqpoint{3.117704in}{2.311064in}}{\pgfqpoint{3.113586in}{2.306946in}}%
\pgfpathcurveto{\pgfqpoint{3.109468in}{2.302828in}}{\pgfqpoint{3.107154in}{2.297241in}}{\pgfqpoint{3.107154in}{2.291417in}}%
\pgfpathcurveto{\pgfqpoint{3.107154in}{2.285593in}}{\pgfqpoint{3.109468in}{2.280007in}}{\pgfqpoint{3.113586in}{2.275889in}}%
\pgfpathcurveto{\pgfqpoint{3.117704in}{2.271771in}}{\pgfqpoint{3.123290in}{2.269457in}}{\pgfqpoint{3.129114in}{2.269457in}}%
\pgfpathlineto{\pgfqpoint{3.129114in}{2.269457in}}%
\pgfpathclose%
\pgfusepath{stroke,fill}%
\end{pgfscope}%
\begin{pgfscope}%
\pgfpathrectangle{\pgfqpoint{1.414736in}{0.548769in}}{\pgfqpoint{3.497747in}{7.253186in}}%
\pgfusepath{clip}%
\pgfsetbuttcap%
\pgfsetroundjoin%
\definecolor{currentfill}{rgb}{0.616293,0.852709,0.230052}%
\pgfsetfillcolor{currentfill}%
\pgfsetlinewidth{1.003750pt}%
\definecolor{currentstroke}{rgb}{0.616293,0.852709,0.230052}%
\pgfsetstrokecolor{currentstroke}%
\pgfsetdash{}{0pt}%
\pgfpathmoveto{\pgfqpoint{3.532147in}{2.740443in}}%
\pgfpathcurveto{\pgfqpoint{3.537970in}{2.740443in}}{\pgfqpoint{3.543557in}{2.742757in}}{\pgfqpoint{3.547675in}{2.746875in}}%
\pgfpathcurveto{\pgfqpoint{3.551793in}{2.750993in}}{\pgfqpoint{3.554107in}{2.756580in}}{\pgfqpoint{3.554107in}{2.762404in}}%
\pgfpathcurveto{\pgfqpoint{3.554107in}{2.768227in}}{\pgfqpoint{3.551793in}{2.773814in}}{\pgfqpoint{3.547675in}{2.777932in}}%
\pgfpathcurveto{\pgfqpoint{3.543557in}{2.782050in}}{\pgfqpoint{3.537970in}{2.784364in}}{\pgfqpoint{3.532147in}{2.784364in}}%
\pgfpathcurveto{\pgfqpoint{3.526323in}{2.784364in}}{\pgfqpoint{3.520736in}{2.782050in}}{\pgfqpoint{3.516618in}{2.777932in}}%
\pgfpathcurveto{\pgfqpoint{3.512500in}{2.773814in}}{\pgfqpoint{3.510186in}{2.768227in}}{\pgfqpoint{3.510186in}{2.762404in}}%
\pgfpathcurveto{\pgfqpoint{3.510186in}{2.756580in}}{\pgfqpoint{3.512500in}{2.750993in}}{\pgfqpoint{3.516618in}{2.746875in}}%
\pgfpathcurveto{\pgfqpoint{3.520736in}{2.742757in}}{\pgfqpoint{3.526323in}{2.740443in}}{\pgfqpoint{3.532147in}{2.740443in}}%
\pgfpathlineto{\pgfqpoint{3.532147in}{2.740443in}}%
\pgfpathclose%
\pgfusepath{stroke,fill}%
\end{pgfscope}%
\begin{pgfscope}%
\pgfpathrectangle{\pgfqpoint{1.414736in}{0.548769in}}{\pgfqpoint{3.497747in}{7.253186in}}%
\pgfusepath{clip}%
\pgfsetbuttcap%
\pgfsetroundjoin%
\definecolor{currentfill}{rgb}{0.751884,0.874951,0.143228}%
\pgfsetfillcolor{currentfill}%
\pgfsetlinewidth{1.003750pt}%
\definecolor{currentstroke}{rgb}{0.751884,0.874951,0.143228}%
\pgfsetstrokecolor{currentstroke}%
\pgfsetdash{}{0pt}%
\pgfpathmoveto{\pgfqpoint{4.503129in}{3.211429in}}%
\pgfpathcurveto{\pgfqpoint{4.508953in}{3.211429in}}{\pgfqpoint{4.514539in}{3.213743in}}{\pgfqpoint{4.518657in}{3.217861in}}%
\pgfpathcurveto{\pgfqpoint{4.522775in}{3.221980in}}{\pgfqpoint{4.525089in}{3.227566in}}{\pgfqpoint{4.525089in}{3.233390in}}%
\pgfpathcurveto{\pgfqpoint{4.525089in}{3.239214in}}{\pgfqpoint{4.522775in}{3.244800in}}{\pgfqpoint{4.518657in}{3.248918in}}%
\pgfpathcurveto{\pgfqpoint{4.514539in}{3.253036in}}{\pgfqpoint{4.508953in}{3.255350in}}{\pgfqpoint{4.503129in}{3.255350in}}%
\pgfpathcurveto{\pgfqpoint{4.497305in}{3.255350in}}{\pgfqpoint{4.491719in}{3.253036in}}{\pgfqpoint{4.487601in}{3.248918in}}%
\pgfpathcurveto{\pgfqpoint{4.483483in}{3.244800in}}{\pgfqpoint{4.481169in}{3.239214in}}{\pgfqpoint{4.481169in}{3.233390in}}%
\pgfpathcurveto{\pgfqpoint{4.481169in}{3.227566in}}{\pgfqpoint{4.483483in}{3.221980in}}{\pgfqpoint{4.487601in}{3.217861in}}%
\pgfpathcurveto{\pgfqpoint{4.491719in}{3.213743in}}{\pgfqpoint{4.497305in}{3.211429in}}{\pgfqpoint{4.503129in}{3.211429in}}%
\pgfpathlineto{\pgfqpoint{4.503129in}{3.211429in}}%
\pgfpathclose%
\pgfusepath{stroke,fill}%
\end{pgfscope}%
\begin{pgfscope}%
\pgfpathrectangle{\pgfqpoint{1.414736in}{0.548769in}}{\pgfqpoint{3.497747in}{7.253186in}}%
\pgfusepath{clip}%
\pgfsetbuttcap%
\pgfsetroundjoin%
\definecolor{currentfill}{rgb}{0.993248,0.906157,0.143936}%
\pgfsetfillcolor{currentfill}%
\pgfsetlinewidth{1.003750pt}%
\definecolor{currentstroke}{rgb}{0.993248,0.906157,0.143936}%
\pgfsetstrokecolor{currentstroke}%
\pgfsetdash{}{0pt}%
\pgfpathmoveto{\pgfqpoint{2.890294in}{3.682415in}}%
\pgfpathcurveto{\pgfqpoint{2.896117in}{3.682415in}}{\pgfqpoint{2.901704in}{3.684729in}}{\pgfqpoint{2.905822in}{3.688847in}}%
\pgfpathcurveto{\pgfqpoint{2.909940in}{3.692966in}}{\pgfqpoint{2.912254in}{3.698552in}}{\pgfqpoint{2.912254in}{3.704376in}}%
\pgfpathcurveto{\pgfqpoint{2.912254in}{3.710200in}}{\pgfqpoint{2.909940in}{3.715786in}}{\pgfqpoint{2.905822in}{3.719904in}}%
\pgfpathcurveto{\pgfqpoint{2.901704in}{3.724022in}}{\pgfqpoint{2.896117in}{3.726336in}}{\pgfqpoint{2.890294in}{3.726336in}}%
\pgfpathcurveto{\pgfqpoint{2.884470in}{3.726336in}}{\pgfqpoint{2.878883in}{3.724022in}}{\pgfqpoint{2.874765in}{3.719904in}}%
\pgfpathcurveto{\pgfqpoint{2.870647in}{3.715786in}}{\pgfqpoint{2.868333in}{3.710200in}}{\pgfqpoint{2.868333in}{3.704376in}}%
\pgfpathcurveto{\pgfqpoint{2.868333in}{3.698552in}}{\pgfqpoint{2.870647in}{3.692966in}}{\pgfqpoint{2.874765in}{3.688847in}}%
\pgfpathcurveto{\pgfqpoint{2.878883in}{3.684729in}}{\pgfqpoint{2.884470in}{3.682415in}}{\pgfqpoint{2.890294in}{3.682415in}}%
\pgfpathlineto{\pgfqpoint{2.890294in}{3.682415in}}%
\pgfpathclose%
\pgfusepath{stroke,fill}%
\end{pgfscope}%
\begin{pgfscope}%
\pgfpathrectangle{\pgfqpoint{1.414736in}{0.548769in}}{\pgfqpoint{3.497747in}{7.253186in}}%
\pgfusepath{clip}%
\pgfsetbuttcap%
\pgfsetroundjoin%
\definecolor{currentfill}{rgb}{0.720391,0.870350,0.162603}%
\pgfsetfillcolor{currentfill}%
\pgfsetlinewidth{1.003750pt}%
\definecolor{currentstroke}{rgb}{0.720391,0.870350,0.162603}%
\pgfsetstrokecolor{currentstroke}%
\pgfsetdash{}{0pt}%
\pgfpathmoveto{\pgfqpoint{2.941250in}{4.153402in}}%
\pgfpathcurveto{\pgfqpoint{2.947074in}{4.153402in}}{\pgfqpoint{2.952660in}{4.155715in}}{\pgfqpoint{2.956778in}{4.159834in}}%
\pgfpathcurveto{\pgfqpoint{2.960896in}{4.163952in}}{\pgfqpoint{2.963210in}{4.169538in}}{\pgfqpoint{2.963210in}{4.175362in}}%
\pgfpathcurveto{\pgfqpoint{2.963210in}{4.181186in}}{\pgfqpoint{2.960896in}{4.186772in}}{\pgfqpoint{2.956778in}{4.190890in}}%
\pgfpathcurveto{\pgfqpoint{2.952660in}{4.195008in}}{\pgfqpoint{2.947074in}{4.197322in}}{\pgfqpoint{2.941250in}{4.197322in}}%
\pgfpathcurveto{\pgfqpoint{2.935426in}{4.197322in}}{\pgfqpoint{2.929840in}{4.195008in}}{\pgfqpoint{2.925722in}{4.190890in}}%
\pgfpathcurveto{\pgfqpoint{2.921604in}{4.186772in}}{\pgfqpoint{2.919290in}{4.181186in}}{\pgfqpoint{2.919290in}{4.175362in}}%
\pgfpathcurveto{\pgfqpoint{2.919290in}{4.169538in}}{\pgfqpoint{2.921604in}{4.163952in}}{\pgfqpoint{2.925722in}{4.159834in}}%
\pgfpathcurveto{\pgfqpoint{2.929840in}{4.155715in}}{\pgfqpoint{2.935426in}{4.153402in}}{\pgfqpoint{2.941250in}{4.153402in}}%
\pgfpathlineto{\pgfqpoint{2.941250in}{4.153402in}}%
\pgfpathclose%
\pgfusepath{stroke,fill}%
\end{pgfscope}%
\begin{pgfscope}%
\pgfpathrectangle{\pgfqpoint{1.414736in}{0.548769in}}{\pgfqpoint{3.497747in}{7.253186in}}%
\pgfusepath{clip}%
\pgfsetbuttcap%
\pgfsetroundjoin%
\definecolor{currentfill}{rgb}{0.668054,0.861999,0.196293}%
\pgfsetfillcolor{currentfill}%
\pgfsetlinewidth{1.003750pt}%
\definecolor{currentstroke}{rgb}{0.668054,0.861999,0.196293}%
\pgfsetstrokecolor{currentstroke}%
\pgfsetdash{}{0pt}%
\pgfpathmoveto{\pgfqpoint{3.279126in}{4.624388in}}%
\pgfpathcurveto{\pgfqpoint{3.284950in}{4.624388in}}{\pgfqpoint{3.290536in}{4.626702in}}{\pgfqpoint{3.294654in}{4.630820in}}%
\pgfpathcurveto{\pgfqpoint{3.298773in}{4.634938in}}{\pgfqpoint{3.301086in}{4.640524in}}{\pgfqpoint{3.301086in}{4.646348in}}%
\pgfpathcurveto{\pgfqpoint{3.301086in}{4.652172in}}{\pgfqpoint{3.298773in}{4.657758in}}{\pgfqpoint{3.294654in}{4.661876in}}%
\pgfpathcurveto{\pgfqpoint{3.290536in}{4.665994in}}{\pgfqpoint{3.284950in}{4.668308in}}{\pgfqpoint{3.279126in}{4.668308in}}%
\pgfpathcurveto{\pgfqpoint{3.273302in}{4.668308in}}{\pgfqpoint{3.267716in}{4.665994in}}{\pgfqpoint{3.263598in}{4.661876in}}%
\pgfpathcurveto{\pgfqpoint{3.259480in}{4.657758in}}{\pgfqpoint{3.257166in}{4.652172in}}{\pgfqpoint{3.257166in}{4.646348in}}%
\pgfpathcurveto{\pgfqpoint{3.257166in}{4.640524in}}{\pgfqpoint{3.259480in}{4.634938in}}{\pgfqpoint{3.263598in}{4.630820in}}%
\pgfpathcurveto{\pgfqpoint{3.267716in}{4.626702in}}{\pgfqpoint{3.273302in}{4.624388in}}{\pgfqpoint{3.279126in}{4.624388in}}%
\pgfpathlineto{\pgfqpoint{3.279126in}{4.624388in}}%
\pgfpathclose%
\pgfusepath{stroke,fill}%
\end{pgfscope}%
\begin{pgfscope}%
\pgfpathrectangle{\pgfqpoint{1.414736in}{0.548769in}}{\pgfqpoint{3.497747in}{7.253186in}}%
\pgfusepath{clip}%
\pgfsetbuttcap%
\pgfsetroundjoin%
\definecolor{currentfill}{rgb}{0.206756,0.371758,0.553117}%
\pgfsetfillcolor{currentfill}%
\pgfsetlinewidth{1.003750pt}%
\definecolor{currentstroke}{rgb}{0.206756,0.371758,0.553117}%
\pgfsetstrokecolor{currentstroke}%
\pgfsetdash{}{0pt}%
\pgfpathmoveto{\pgfqpoint{4.165244in}{5.095374in}}%
\pgfpathcurveto{\pgfqpoint{4.171068in}{5.095374in}}{\pgfqpoint{4.176655in}{5.097688in}}{\pgfqpoint{4.180773in}{5.101806in}}%
\pgfpathcurveto{\pgfqpoint{4.184891in}{5.105924in}}{\pgfqpoint{4.187205in}{5.111510in}}{\pgfqpoint{4.187205in}{5.117334in}}%
\pgfpathcurveto{\pgfqpoint{4.187205in}{5.123158in}}{\pgfqpoint{4.184891in}{5.128744in}}{\pgfqpoint{4.180773in}{5.132862in}}%
\pgfpathcurveto{\pgfqpoint{4.176655in}{5.136980in}}{\pgfqpoint{4.171068in}{5.139294in}}{\pgfqpoint{4.165244in}{5.139294in}}%
\pgfpathcurveto{\pgfqpoint{4.159420in}{5.139294in}}{\pgfqpoint{4.153834in}{5.136980in}}{\pgfqpoint{4.149716in}{5.132862in}}%
\pgfpathcurveto{\pgfqpoint{4.145598in}{5.128744in}}{\pgfqpoint{4.143284in}{5.123158in}}{\pgfqpoint{4.143284in}{5.117334in}}%
\pgfpathcurveto{\pgfqpoint{4.143284in}{5.111510in}}{\pgfqpoint{4.145598in}{5.105924in}}{\pgfqpoint{4.149716in}{5.101806in}}%
\pgfpathcurveto{\pgfqpoint{4.153834in}{5.097688in}}{\pgfqpoint{4.159420in}{5.095374in}}{\pgfqpoint{4.165244in}{5.095374in}}%
\pgfpathlineto{\pgfqpoint{4.165244in}{5.095374in}}%
\pgfpathclose%
\pgfusepath{stroke,fill}%
\end{pgfscope}%
\begin{pgfscope}%
\pgfpathrectangle{\pgfqpoint{1.414736in}{0.548769in}}{\pgfqpoint{3.497747in}{7.253186in}}%
\pgfusepath{clip}%
\pgfsetbuttcap%
\pgfsetroundjoin%
\definecolor{currentfill}{rgb}{0.172719,0.448791,0.557885}%
\pgfsetfillcolor{currentfill}%
\pgfsetlinewidth{1.003750pt}%
\definecolor{currentstroke}{rgb}{0.172719,0.448791,0.557885}%
\pgfsetstrokecolor{currentstroke}%
\pgfsetdash{}{0pt}%
\pgfpathmoveto{\pgfqpoint{3.714126in}{5.566360in}}%
\pgfpathcurveto{\pgfqpoint{3.719950in}{5.566360in}}{\pgfqpoint{3.725536in}{5.568674in}}{\pgfqpoint{3.729654in}{5.572792in}}%
\pgfpathcurveto{\pgfqpoint{3.733772in}{5.576910in}}{\pgfqpoint{3.736086in}{5.582496in}}{\pgfqpoint{3.736086in}{5.588320in}}%
\pgfpathcurveto{\pgfqpoint{3.736086in}{5.594144in}}{\pgfqpoint{3.733772in}{5.599730in}}{\pgfqpoint{3.729654in}{5.603848in}}%
\pgfpathcurveto{\pgfqpoint{3.725536in}{5.607967in}}{\pgfqpoint{3.719950in}{5.610280in}}{\pgfqpoint{3.714126in}{5.610280in}}%
\pgfpathcurveto{\pgfqpoint{3.708302in}{5.610280in}}{\pgfqpoint{3.702716in}{5.607967in}}{\pgfqpoint{3.698598in}{5.603848in}}%
\pgfpathcurveto{\pgfqpoint{3.694479in}{5.599730in}}{\pgfqpoint{3.692166in}{5.594144in}}{\pgfqpoint{3.692166in}{5.588320in}}%
\pgfpathcurveto{\pgfqpoint{3.692166in}{5.582496in}}{\pgfqpoint{3.694479in}{5.576910in}}{\pgfqpoint{3.698598in}{5.572792in}}%
\pgfpathcurveto{\pgfqpoint{3.702716in}{5.568674in}}{\pgfqpoint{3.708302in}{5.566360in}}{\pgfqpoint{3.714126in}{5.566360in}}%
\pgfpathlineto{\pgfqpoint{3.714126in}{5.566360in}}%
\pgfpathclose%
\pgfusepath{stroke,fill}%
\end{pgfscope}%
\begin{pgfscope}%
\pgfpathrectangle{\pgfqpoint{1.414736in}{0.548769in}}{\pgfqpoint{3.497747in}{7.253186in}}%
\pgfusepath{clip}%
\pgfsetbuttcap%
\pgfsetroundjoin%
\definecolor{currentfill}{rgb}{0.477504,0.821444,0.318195}%
\pgfsetfillcolor{currentfill}%
\pgfsetlinewidth{1.003750pt}%
\definecolor{currentstroke}{rgb}{0.477504,0.821444,0.318195}%
\pgfsetstrokecolor{currentstroke}%
\pgfsetdash{}{0pt}%
\pgfpathmoveto{\pgfqpoint{3.413458in}{6.037346in}}%
\pgfpathcurveto{\pgfqpoint{3.419282in}{6.037346in}}{\pgfqpoint{3.424868in}{6.039660in}}{\pgfqpoint{3.428986in}{6.043778in}}%
\pgfpathcurveto{\pgfqpoint{3.433104in}{6.047896in}}{\pgfqpoint{3.435418in}{6.053482in}}{\pgfqpoint{3.435418in}{6.059306in}}%
\pgfpathcurveto{\pgfqpoint{3.435418in}{6.065130in}}{\pgfqpoint{3.433104in}{6.070716in}}{\pgfqpoint{3.428986in}{6.074835in}}%
\pgfpathcurveto{\pgfqpoint{3.424868in}{6.078953in}}{\pgfqpoint{3.419282in}{6.081267in}}{\pgfqpoint{3.413458in}{6.081267in}}%
\pgfpathcurveto{\pgfqpoint{3.407634in}{6.081267in}}{\pgfqpoint{3.402048in}{6.078953in}}{\pgfqpoint{3.397930in}{6.074835in}}%
\pgfpathcurveto{\pgfqpoint{3.393812in}{6.070716in}}{\pgfqpoint{3.391498in}{6.065130in}}{\pgfqpoint{3.391498in}{6.059306in}}%
\pgfpathcurveto{\pgfqpoint{3.391498in}{6.053482in}}{\pgfqpoint{3.393812in}{6.047896in}}{\pgfqpoint{3.397930in}{6.043778in}}%
\pgfpathcurveto{\pgfqpoint{3.402048in}{6.039660in}}{\pgfqpoint{3.407634in}{6.037346in}}{\pgfqpoint{3.413458in}{6.037346in}}%
\pgfpathlineto{\pgfqpoint{3.413458in}{6.037346in}}%
\pgfpathclose%
\pgfusepath{stroke,fill}%
\end{pgfscope}%
\begin{pgfscope}%
\pgfpathrectangle{\pgfqpoint{1.414736in}{0.548769in}}{\pgfqpoint{3.497747in}{7.253186in}}%
\pgfusepath{clip}%
\pgfsetbuttcap%
\pgfsetroundjoin%
\definecolor{currentfill}{rgb}{0.993248,0.906157,0.143936}%
\pgfsetfillcolor{currentfill}%
\pgfsetlinewidth{1.003750pt}%
\definecolor{currentstroke}{rgb}{0.993248,0.906157,0.143936}%
\pgfsetstrokecolor{currentstroke}%
\pgfsetdash{}{0pt}%
\pgfpathmoveto{\pgfqpoint{3.226157in}{6.508332in}}%
\pgfpathcurveto{\pgfqpoint{3.231981in}{6.508332in}}{\pgfqpoint{3.237567in}{6.510646in}}{\pgfqpoint{3.241685in}{6.514764in}}%
\pgfpathcurveto{\pgfqpoint{3.245804in}{6.518882in}}{\pgfqpoint{3.248117in}{6.524468in}}{\pgfqpoint{3.248117in}{6.530292in}}%
\pgfpathcurveto{\pgfqpoint{3.248117in}{6.536116in}}{\pgfqpoint{3.245804in}{6.541702in}}{\pgfqpoint{3.241685in}{6.545821in}}%
\pgfpathcurveto{\pgfqpoint{3.237567in}{6.549939in}}{\pgfqpoint{3.231981in}{6.552253in}}{\pgfqpoint{3.226157in}{6.552253in}}%
\pgfpathcurveto{\pgfqpoint{3.220333in}{6.552253in}}{\pgfqpoint{3.214747in}{6.549939in}}{\pgfqpoint{3.210629in}{6.545821in}}%
\pgfpathcurveto{\pgfqpoint{3.206511in}{6.541702in}}{\pgfqpoint{3.204197in}{6.536116in}}{\pgfqpoint{3.204197in}{6.530292in}}%
\pgfpathcurveto{\pgfqpoint{3.204197in}{6.524468in}}{\pgfqpoint{3.206511in}{6.518882in}}{\pgfqpoint{3.210629in}{6.514764in}}%
\pgfpathcurveto{\pgfqpoint{3.214747in}{6.510646in}}{\pgfqpoint{3.220333in}{6.508332in}}{\pgfqpoint{3.226157in}{6.508332in}}%
\pgfpathlineto{\pgfqpoint{3.226157in}{6.508332in}}%
\pgfpathclose%
\pgfusepath{stroke,fill}%
\end{pgfscope}%
\begin{pgfscope}%
\pgfpathrectangle{\pgfqpoint{1.414736in}{0.548769in}}{\pgfqpoint{3.497747in}{7.253186in}}%
\pgfusepath{clip}%
\pgfsetbuttcap%
\pgfsetroundjoin%
\definecolor{currentfill}{rgb}{0.430983,0.808473,0.346476}%
\pgfsetfillcolor{currentfill}%
\pgfsetlinewidth{1.003750pt}%
\definecolor{currentstroke}{rgb}{0.430983,0.808473,0.346476}%
\pgfsetstrokecolor{currentstroke}%
\pgfsetdash{}{0pt}%
\pgfpathmoveto{\pgfqpoint{3.497739in}{6.979318in}}%
\pgfpathcurveto{\pgfqpoint{3.503563in}{6.979318in}}{\pgfqpoint{3.509149in}{6.981632in}}{\pgfqpoint{3.513267in}{6.985750in}}%
\pgfpathcurveto{\pgfqpoint{3.517385in}{6.989868in}}{\pgfqpoint{3.519699in}{6.995455in}}{\pgfqpoint{3.519699in}{7.001278in}}%
\pgfpathcurveto{\pgfqpoint{3.519699in}{7.007102in}}{\pgfqpoint{3.517385in}{7.012689in}}{\pgfqpoint{3.513267in}{7.016807in}}%
\pgfpathcurveto{\pgfqpoint{3.509149in}{7.020925in}}{\pgfqpoint{3.503563in}{7.023239in}}{\pgfqpoint{3.497739in}{7.023239in}}%
\pgfpathcurveto{\pgfqpoint{3.491915in}{7.023239in}}{\pgfqpoint{3.486329in}{7.020925in}}{\pgfqpoint{3.482211in}{7.016807in}}%
\pgfpathcurveto{\pgfqpoint{3.478093in}{7.012689in}}{\pgfqpoint{3.475779in}{7.007102in}}{\pgfqpoint{3.475779in}{7.001278in}}%
\pgfpathcurveto{\pgfqpoint{3.475779in}{6.995455in}}{\pgfqpoint{3.478093in}{6.989868in}}{\pgfqpoint{3.482211in}{6.985750in}}%
\pgfpathcurveto{\pgfqpoint{3.486329in}{6.981632in}}{\pgfqpoint{3.491915in}{6.979318in}}{\pgfqpoint{3.497739in}{6.979318in}}%
\pgfpathlineto{\pgfqpoint{3.497739in}{6.979318in}}%
\pgfpathclose%
\pgfusepath{stroke,fill}%
\end{pgfscope}%
\begin{pgfscope}%
\pgfpathrectangle{\pgfqpoint{1.414736in}{0.548769in}}{\pgfqpoint{3.497747in}{7.253186in}}%
\pgfusepath{clip}%
\pgfsetbuttcap%
\pgfsetroundjoin%
\definecolor{currentfill}{rgb}{0.699415,0.867117,0.175971}%
\pgfsetfillcolor{currentfill}%
\pgfsetlinewidth{1.003750pt}%
\definecolor{currentstroke}{rgb}{0.699415,0.867117,0.175971}%
\pgfsetstrokecolor{currentstroke}%
\pgfsetdash{}{0pt}%
\pgfpathmoveto{\pgfqpoint{3.738633in}{7.450304in}}%
\pgfpathcurveto{\pgfqpoint{3.744457in}{7.450304in}}{\pgfqpoint{3.750043in}{7.452618in}}{\pgfqpoint{3.754161in}{7.456736in}}%
\pgfpathcurveto{\pgfqpoint{3.758279in}{7.460854in}}{\pgfqpoint{3.760593in}{7.466441in}}{\pgfqpoint{3.760593in}{7.472265in}}%
\pgfpathcurveto{\pgfqpoint{3.760593in}{7.478089in}}{\pgfqpoint{3.758279in}{7.483675in}}{\pgfqpoint{3.754161in}{7.487793in}}%
\pgfpathcurveto{\pgfqpoint{3.750043in}{7.491911in}}{\pgfqpoint{3.744457in}{7.494225in}}{\pgfqpoint{3.738633in}{7.494225in}}%
\pgfpathcurveto{\pgfqpoint{3.732809in}{7.494225in}}{\pgfqpoint{3.727223in}{7.491911in}}{\pgfqpoint{3.723105in}{7.487793in}}%
\pgfpathcurveto{\pgfqpoint{3.718987in}{7.483675in}}{\pgfqpoint{3.716673in}{7.478089in}}{\pgfqpoint{3.716673in}{7.472265in}}%
\pgfpathcurveto{\pgfqpoint{3.716673in}{7.466441in}}{\pgfqpoint{3.718987in}{7.460854in}}{\pgfqpoint{3.723105in}{7.456736in}}%
\pgfpathcurveto{\pgfqpoint{3.727223in}{7.452618in}}{\pgfqpoint{3.732809in}{7.450304in}}{\pgfqpoint{3.738633in}{7.450304in}}%
\pgfpathlineto{\pgfqpoint{3.738633in}{7.450304in}}%
\pgfpathclose%
\pgfusepath{stroke,fill}%
\end{pgfscope}%
\begin{pgfscope}%
\pgfpathrectangle{\pgfqpoint{1.414736in}{0.548769in}}{\pgfqpoint{3.497747in}{7.253186in}}%
\pgfusepath{clip}%
\pgfsetbuttcap%
\pgfsetroundjoin%
\definecolor{currentfill}{rgb}{0.974417,0.903590,0.130215}%
\pgfsetfillcolor{currentfill}%
\pgfsetlinewidth{1.003750pt}%
\definecolor{currentstroke}{rgb}{0.974417,0.903590,0.130215}%
\pgfsetstrokecolor{currentstroke}%
\pgfsetdash{}{0pt}%
\pgfpathmoveto{\pgfqpoint{3.479368in}{0.856499in}}%
\pgfpathcurveto{\pgfqpoint{3.485192in}{0.856499in}}{\pgfqpoint{3.490778in}{0.858813in}}{\pgfqpoint{3.494896in}{0.862931in}}%
\pgfpathcurveto{\pgfqpoint{3.499014in}{0.867049in}}{\pgfqpoint{3.501328in}{0.872635in}}{\pgfqpoint{3.501328in}{0.878459in}}%
\pgfpathcurveto{\pgfqpoint{3.501328in}{0.884283in}}{\pgfqpoint{3.499014in}{0.889869in}}{\pgfqpoint{3.494896in}{0.893987in}}%
\pgfpathcurveto{\pgfqpoint{3.490778in}{0.898105in}}{\pgfqpoint{3.485192in}{0.900419in}}{\pgfqpoint{3.479368in}{0.900419in}}%
\pgfpathcurveto{\pgfqpoint{3.473544in}{0.900419in}}{\pgfqpoint{3.467957in}{0.898105in}}{\pgfqpoint{3.463839in}{0.893987in}}%
\pgfpathcurveto{\pgfqpoint{3.459721in}{0.889869in}}{\pgfqpoint{3.457407in}{0.884283in}}{\pgfqpoint{3.457407in}{0.878459in}}%
\pgfpathcurveto{\pgfqpoint{3.457407in}{0.872635in}}{\pgfqpoint{3.459721in}{0.867049in}}{\pgfqpoint{3.463839in}{0.862931in}}%
\pgfpathcurveto{\pgfqpoint{3.467957in}{0.858813in}}{\pgfqpoint{3.473544in}{0.856499in}}{\pgfqpoint{3.479368in}{0.856499in}}%
\pgfpathlineto{\pgfqpoint{3.479368in}{0.856499in}}%
\pgfpathclose%
\pgfusepath{stroke,fill}%
\end{pgfscope}%
\begin{pgfscope}%
\pgfpathrectangle{\pgfqpoint{1.414736in}{0.548769in}}{\pgfqpoint{3.497747in}{7.253186in}}%
\pgfusepath{clip}%
\pgfsetbuttcap%
\pgfsetroundjoin%
\definecolor{currentfill}{rgb}{0.468053,0.818921,0.323998}%
\pgfsetfillcolor{currentfill}%
\pgfsetlinewidth{1.003750pt}%
\definecolor{currentstroke}{rgb}{0.468053,0.818921,0.323998}%
\pgfsetstrokecolor{currentstroke}%
\pgfsetdash{}{0pt}%
\pgfpathmoveto{\pgfqpoint{3.599304in}{1.327485in}}%
\pgfpathcurveto{\pgfqpoint{3.605128in}{1.327485in}}{\pgfqpoint{3.610714in}{1.329799in}}{\pgfqpoint{3.614833in}{1.333917in}}%
\pgfpathcurveto{\pgfqpoint{3.618951in}{1.338035in}}{\pgfqpoint{3.621265in}{1.343621in}}{\pgfqpoint{3.621265in}{1.349445in}}%
\pgfpathcurveto{\pgfqpoint{3.621265in}{1.355269in}}{\pgfqpoint{3.618951in}{1.360855in}}{\pgfqpoint{3.614833in}{1.364973in}}%
\pgfpathcurveto{\pgfqpoint{3.610714in}{1.369092in}}{\pgfqpoint{3.605128in}{1.371405in}}{\pgfqpoint{3.599304in}{1.371405in}}%
\pgfpathcurveto{\pgfqpoint{3.593480in}{1.371405in}}{\pgfqpoint{3.587894in}{1.369092in}}{\pgfqpoint{3.583776in}{1.364973in}}%
\pgfpathcurveto{\pgfqpoint{3.579658in}{1.360855in}}{\pgfqpoint{3.577344in}{1.355269in}}{\pgfqpoint{3.577344in}{1.349445in}}%
\pgfpathcurveto{\pgfqpoint{3.577344in}{1.343621in}}{\pgfqpoint{3.579658in}{1.338035in}}{\pgfqpoint{3.583776in}{1.333917in}}%
\pgfpathcurveto{\pgfqpoint{3.587894in}{1.329799in}}{\pgfqpoint{3.593480in}{1.327485in}}{\pgfqpoint{3.599304in}{1.327485in}}%
\pgfpathlineto{\pgfqpoint{3.599304in}{1.327485in}}%
\pgfpathclose%
\pgfusepath{stroke,fill}%
\end{pgfscope}%
\begin{pgfscope}%
\pgfpathrectangle{\pgfqpoint{1.414736in}{0.548769in}}{\pgfqpoint{3.497747in}{7.253186in}}%
\pgfusepath{clip}%
\pgfsetbuttcap%
\pgfsetroundjoin%
\definecolor{currentfill}{rgb}{0.449368,0.813768,0.335384}%
\pgfsetfillcolor{currentfill}%
\pgfsetlinewidth{1.003750pt}%
\definecolor{currentstroke}{rgb}{0.449368,0.813768,0.335384}%
\pgfsetstrokecolor{currentstroke}%
\pgfsetdash{}{0pt}%
\pgfpathmoveto{\pgfqpoint{3.482299in}{1.798471in}}%
\pgfpathcurveto{\pgfqpoint{3.488123in}{1.798471in}}{\pgfqpoint{3.493709in}{1.800785in}}{\pgfqpoint{3.497827in}{1.804903in}}%
\pgfpathcurveto{\pgfqpoint{3.501946in}{1.809021in}}{\pgfqpoint{3.504259in}{1.814607in}}{\pgfqpoint{3.504259in}{1.820431in}}%
\pgfpathcurveto{\pgfqpoint{3.504259in}{1.826255in}}{\pgfqpoint{3.501946in}{1.831841in}}{\pgfqpoint{3.497827in}{1.835960in}}%
\pgfpathcurveto{\pgfqpoint{3.493709in}{1.840078in}}{\pgfqpoint{3.488123in}{1.842392in}}{\pgfqpoint{3.482299in}{1.842392in}}%
\pgfpathcurveto{\pgfqpoint{3.476475in}{1.842392in}}{\pgfqpoint{3.470889in}{1.840078in}}{\pgfqpoint{3.466771in}{1.835960in}}%
\pgfpathcurveto{\pgfqpoint{3.462653in}{1.831841in}}{\pgfqpoint{3.460339in}{1.826255in}}{\pgfqpoint{3.460339in}{1.820431in}}%
\pgfpathcurveto{\pgfqpoint{3.460339in}{1.814607in}}{\pgfqpoint{3.462653in}{1.809021in}}{\pgfqpoint{3.466771in}{1.804903in}}%
\pgfpathcurveto{\pgfqpoint{3.470889in}{1.800785in}}{\pgfqpoint{3.476475in}{1.798471in}}{\pgfqpoint{3.482299in}{1.798471in}}%
\pgfpathlineto{\pgfqpoint{3.482299in}{1.798471in}}%
\pgfpathclose%
\pgfusepath{stroke,fill}%
\end{pgfscope}%
\begin{pgfscope}%
\pgfpathrectangle{\pgfqpoint{1.414736in}{0.548769in}}{\pgfqpoint{3.497747in}{7.253186in}}%
\pgfusepath{clip}%
\pgfsetbuttcap%
\pgfsetroundjoin%
\definecolor{currentfill}{rgb}{0.916242,0.896091,0.100717}%
\pgfsetfillcolor{currentfill}%
\pgfsetlinewidth{1.003750pt}%
\definecolor{currentstroke}{rgb}{0.916242,0.896091,0.100717}%
\pgfsetstrokecolor{currentstroke}%
\pgfsetdash{}{0pt}%
\pgfpathmoveto{\pgfqpoint{3.495973in}{2.269457in}}%
\pgfpathcurveto{\pgfqpoint{3.501797in}{2.269457in}}{\pgfqpoint{3.507383in}{2.271771in}}{\pgfqpoint{3.511501in}{2.275889in}}%
\pgfpathcurveto{\pgfqpoint{3.515619in}{2.280007in}}{\pgfqpoint{3.517933in}{2.285593in}}{\pgfqpoint{3.517933in}{2.291417in}}%
\pgfpathcurveto{\pgfqpoint{3.517933in}{2.297241in}}{\pgfqpoint{3.515619in}{2.302828in}}{\pgfqpoint{3.511501in}{2.306946in}}%
\pgfpathcurveto{\pgfqpoint{3.507383in}{2.311064in}}{\pgfqpoint{3.501797in}{2.313378in}}{\pgfqpoint{3.495973in}{2.313378in}}%
\pgfpathcurveto{\pgfqpoint{3.490149in}{2.313378in}}{\pgfqpoint{3.484563in}{2.311064in}}{\pgfqpoint{3.480445in}{2.306946in}}%
\pgfpathcurveto{\pgfqpoint{3.476326in}{2.302828in}}{\pgfqpoint{3.474013in}{2.297241in}}{\pgfqpoint{3.474013in}{2.291417in}}%
\pgfpathcurveto{\pgfqpoint{3.474013in}{2.285593in}}{\pgfqpoint{3.476326in}{2.280007in}}{\pgfqpoint{3.480445in}{2.275889in}}%
\pgfpathcurveto{\pgfqpoint{3.484563in}{2.271771in}}{\pgfqpoint{3.490149in}{2.269457in}}{\pgfqpoint{3.495973in}{2.269457in}}%
\pgfpathlineto{\pgfqpoint{3.495973in}{2.269457in}}%
\pgfpathclose%
\pgfusepath{stroke,fill}%
\end{pgfscope}%
\begin{pgfscope}%
\pgfpathrectangle{\pgfqpoint{1.414736in}{0.548769in}}{\pgfqpoint{3.497747in}{7.253186in}}%
\pgfusepath{clip}%
\pgfsetbuttcap%
\pgfsetroundjoin%
\definecolor{currentfill}{rgb}{0.585678,0.846661,0.249897}%
\pgfsetfillcolor{currentfill}%
\pgfsetlinewidth{1.003750pt}%
\definecolor{currentstroke}{rgb}{0.585678,0.846661,0.249897}%
\pgfsetstrokecolor{currentstroke}%
\pgfsetdash{}{0pt}%
\pgfpathmoveto{\pgfqpoint{3.443016in}{2.740443in}}%
\pgfpathcurveto{\pgfqpoint{3.448840in}{2.740443in}}{\pgfqpoint{3.454427in}{2.742757in}}{\pgfqpoint{3.458545in}{2.746875in}}%
\pgfpathcurveto{\pgfqpoint{3.462663in}{2.750993in}}{\pgfqpoint{3.464977in}{2.756580in}}{\pgfqpoint{3.464977in}{2.762404in}}%
\pgfpathcurveto{\pgfqpoint{3.464977in}{2.768227in}}{\pgfqpoint{3.462663in}{2.773814in}}{\pgfqpoint{3.458545in}{2.777932in}}%
\pgfpathcurveto{\pgfqpoint{3.454427in}{2.782050in}}{\pgfqpoint{3.448840in}{2.784364in}}{\pgfqpoint{3.443016in}{2.784364in}}%
\pgfpathcurveto{\pgfqpoint{3.437193in}{2.784364in}}{\pgfqpoint{3.431606in}{2.782050in}}{\pgfqpoint{3.427488in}{2.777932in}}%
\pgfpathcurveto{\pgfqpoint{3.423370in}{2.773814in}}{\pgfqpoint{3.421056in}{2.768227in}}{\pgfqpoint{3.421056in}{2.762404in}}%
\pgfpathcurveto{\pgfqpoint{3.421056in}{2.756580in}}{\pgfqpoint{3.423370in}{2.750993in}}{\pgfqpoint{3.427488in}{2.746875in}}%
\pgfpathcurveto{\pgfqpoint{3.431606in}{2.742757in}}{\pgfqpoint{3.437193in}{2.740443in}}{\pgfqpoint{3.443016in}{2.740443in}}%
\pgfpathlineto{\pgfqpoint{3.443016in}{2.740443in}}%
\pgfpathclose%
\pgfusepath{stroke,fill}%
\end{pgfscope}%
\begin{pgfscope}%
\pgfpathrectangle{\pgfqpoint{1.414736in}{0.548769in}}{\pgfqpoint{3.497747in}{7.253186in}}%
\pgfusepath{clip}%
\pgfsetbuttcap%
\pgfsetroundjoin%
\definecolor{currentfill}{rgb}{0.678489,0.863742,0.189503}%
\pgfsetfillcolor{currentfill}%
\pgfsetlinewidth{1.003750pt}%
\definecolor{currentstroke}{rgb}{0.678489,0.863742,0.189503}%
\pgfsetstrokecolor{currentstroke}%
\pgfsetdash{}{0pt}%
\pgfpathmoveto{\pgfqpoint{3.523172in}{3.211429in}}%
\pgfpathcurveto{\pgfqpoint{3.528996in}{3.211429in}}{\pgfqpoint{3.534583in}{3.213743in}}{\pgfqpoint{3.538701in}{3.217861in}}%
\pgfpathcurveto{\pgfqpoint{3.542819in}{3.221980in}}{\pgfqpoint{3.545133in}{3.227566in}}{\pgfqpoint{3.545133in}{3.233390in}}%
\pgfpathcurveto{\pgfqpoint{3.545133in}{3.239214in}}{\pgfqpoint{3.542819in}{3.244800in}}{\pgfqpoint{3.538701in}{3.248918in}}%
\pgfpathcurveto{\pgfqpoint{3.534583in}{3.253036in}}{\pgfqpoint{3.528996in}{3.255350in}}{\pgfqpoint{3.523172in}{3.255350in}}%
\pgfpathcurveto{\pgfqpoint{3.517349in}{3.255350in}}{\pgfqpoint{3.511762in}{3.253036in}}{\pgfqpoint{3.507644in}{3.248918in}}%
\pgfpathcurveto{\pgfqpoint{3.503526in}{3.244800in}}{\pgfqpoint{3.501212in}{3.239214in}}{\pgfqpoint{3.501212in}{3.233390in}}%
\pgfpathcurveto{\pgfqpoint{3.501212in}{3.227566in}}{\pgfqpoint{3.503526in}{3.221980in}}{\pgfqpoint{3.507644in}{3.217861in}}%
\pgfpathcurveto{\pgfqpoint{3.511762in}{3.213743in}}{\pgfqpoint{3.517349in}{3.211429in}}{\pgfqpoint{3.523172in}{3.211429in}}%
\pgfpathlineto{\pgfqpoint{3.523172in}{3.211429in}}%
\pgfpathclose%
\pgfusepath{stroke,fill}%
\end{pgfscope}%
\begin{pgfscope}%
\pgfpathrectangle{\pgfqpoint{1.414736in}{0.548769in}}{\pgfqpoint{3.497747in}{7.253186in}}%
\pgfusepath{clip}%
\pgfsetbuttcap%
\pgfsetroundjoin%
\definecolor{currentfill}{rgb}{0.793760,0.880678,0.120005}%
\pgfsetfillcolor{currentfill}%
\pgfsetlinewidth{1.003750pt}%
\definecolor{currentstroke}{rgb}{0.793760,0.880678,0.120005}%
\pgfsetstrokecolor{currentstroke}%
\pgfsetdash{}{0pt}%
\pgfpathmoveto{\pgfqpoint{3.543748in}{3.682415in}}%
\pgfpathcurveto{\pgfqpoint{3.549572in}{3.682415in}}{\pgfqpoint{3.555158in}{3.684729in}}{\pgfqpoint{3.559276in}{3.688847in}}%
\pgfpathcurveto{\pgfqpoint{3.563394in}{3.692966in}}{\pgfqpoint{3.565708in}{3.698552in}}{\pgfqpoint{3.565708in}{3.704376in}}%
\pgfpathcurveto{\pgfqpoint{3.565708in}{3.710200in}}{\pgfqpoint{3.563394in}{3.715786in}}{\pgfqpoint{3.559276in}{3.719904in}}%
\pgfpathcurveto{\pgfqpoint{3.555158in}{3.724022in}}{\pgfqpoint{3.549572in}{3.726336in}}{\pgfqpoint{3.543748in}{3.726336in}}%
\pgfpathcurveto{\pgfqpoint{3.537924in}{3.726336in}}{\pgfqpoint{3.532338in}{3.724022in}}{\pgfqpoint{3.528220in}{3.719904in}}%
\pgfpathcurveto{\pgfqpoint{3.524101in}{3.715786in}}{\pgfqpoint{3.521788in}{3.710200in}}{\pgfqpoint{3.521788in}{3.704376in}}%
\pgfpathcurveto{\pgfqpoint{3.521788in}{3.698552in}}{\pgfqpoint{3.524101in}{3.692966in}}{\pgfqpoint{3.528220in}{3.688847in}}%
\pgfpathcurveto{\pgfqpoint{3.532338in}{3.684729in}}{\pgfqpoint{3.537924in}{3.682415in}}{\pgfqpoint{3.543748in}{3.682415in}}%
\pgfpathlineto{\pgfqpoint{3.543748in}{3.682415in}}%
\pgfpathclose%
\pgfusepath{stroke,fill}%
\end{pgfscope}%
\begin{pgfscope}%
\pgfpathrectangle{\pgfqpoint{1.414736in}{0.548769in}}{\pgfqpoint{3.497747in}{7.253186in}}%
\pgfusepath{clip}%
\pgfsetbuttcap%
\pgfsetroundjoin%
\definecolor{currentfill}{rgb}{0.545524,0.838039,0.275626}%
\pgfsetfillcolor{currentfill}%
\pgfsetlinewidth{1.003750pt}%
\definecolor{currentstroke}{rgb}{0.545524,0.838039,0.275626}%
\pgfsetstrokecolor{currentstroke}%
\pgfsetdash{}{0pt}%
\pgfpathmoveto{\pgfqpoint{3.826766in}{4.153402in}}%
\pgfpathcurveto{\pgfqpoint{3.832590in}{4.153402in}}{\pgfqpoint{3.838176in}{4.155715in}}{\pgfqpoint{3.842295in}{4.159834in}}%
\pgfpathcurveto{\pgfqpoint{3.846413in}{4.163952in}}{\pgfqpoint{3.848727in}{4.169538in}}{\pgfqpoint{3.848727in}{4.175362in}}%
\pgfpathcurveto{\pgfqpoint{3.848727in}{4.181186in}}{\pgfqpoint{3.846413in}{4.186772in}}{\pgfqpoint{3.842295in}{4.190890in}}%
\pgfpathcurveto{\pgfqpoint{3.838176in}{4.195008in}}{\pgfqpoint{3.832590in}{4.197322in}}{\pgfqpoint{3.826766in}{4.197322in}}%
\pgfpathcurveto{\pgfqpoint{3.820942in}{4.197322in}}{\pgfqpoint{3.815356in}{4.195008in}}{\pgfqpoint{3.811238in}{4.190890in}}%
\pgfpathcurveto{\pgfqpoint{3.807120in}{4.186772in}}{\pgfqpoint{3.804806in}{4.181186in}}{\pgfqpoint{3.804806in}{4.175362in}}%
\pgfpathcurveto{\pgfqpoint{3.804806in}{4.169538in}}{\pgfqpoint{3.807120in}{4.163952in}}{\pgfqpoint{3.811238in}{4.159834in}}%
\pgfpathcurveto{\pgfqpoint{3.815356in}{4.155715in}}{\pgfqpoint{3.820942in}{4.153402in}}{\pgfqpoint{3.826766in}{4.153402in}}%
\pgfpathlineto{\pgfqpoint{3.826766in}{4.153402in}}%
\pgfpathclose%
\pgfusepath{stroke,fill}%
\end{pgfscope}%
\begin{pgfscope}%
\pgfpathrectangle{\pgfqpoint{1.414736in}{0.548769in}}{\pgfqpoint{3.497747in}{7.253186in}}%
\pgfusepath{clip}%
\pgfsetbuttcap%
\pgfsetroundjoin%
\definecolor{currentfill}{rgb}{0.668054,0.861999,0.196293}%
\pgfsetfillcolor{currentfill}%
\pgfsetlinewidth{1.003750pt}%
\definecolor{currentstroke}{rgb}{0.668054,0.861999,0.196293}%
\pgfsetstrokecolor{currentstroke}%
\pgfsetdash{}{0pt}%
\pgfpathmoveto{\pgfqpoint{3.389142in}{4.624388in}}%
\pgfpathcurveto{\pgfqpoint{3.394965in}{4.624388in}}{\pgfqpoint{3.400552in}{4.626702in}}{\pgfqpoint{3.404670in}{4.630820in}}%
\pgfpathcurveto{\pgfqpoint{3.408788in}{4.634938in}}{\pgfqpoint{3.411102in}{4.640524in}}{\pgfqpoint{3.411102in}{4.646348in}}%
\pgfpathcurveto{\pgfqpoint{3.411102in}{4.652172in}}{\pgfqpoint{3.408788in}{4.657758in}}{\pgfqpoint{3.404670in}{4.661876in}}%
\pgfpathcurveto{\pgfqpoint{3.400552in}{4.665994in}}{\pgfqpoint{3.394965in}{4.668308in}}{\pgfqpoint{3.389142in}{4.668308in}}%
\pgfpathcurveto{\pgfqpoint{3.383318in}{4.668308in}}{\pgfqpoint{3.377731in}{4.665994in}}{\pgfqpoint{3.373613in}{4.661876in}}%
\pgfpathcurveto{\pgfqpoint{3.369495in}{4.657758in}}{\pgfqpoint{3.367181in}{4.652172in}}{\pgfqpoint{3.367181in}{4.646348in}}%
\pgfpathcurveto{\pgfqpoint{3.367181in}{4.640524in}}{\pgfqpoint{3.369495in}{4.634938in}}{\pgfqpoint{3.373613in}{4.630820in}}%
\pgfpathcurveto{\pgfqpoint{3.377731in}{4.626702in}}{\pgfqpoint{3.383318in}{4.624388in}}{\pgfqpoint{3.389142in}{4.624388in}}%
\pgfpathlineto{\pgfqpoint{3.389142in}{4.624388in}}%
\pgfpathclose%
\pgfusepath{stroke,fill}%
\end{pgfscope}%
\begin{pgfscope}%
\pgfpathrectangle{\pgfqpoint{1.414736in}{0.548769in}}{\pgfqpoint{3.497747in}{7.253186in}}%
\pgfusepath{clip}%
\pgfsetbuttcap%
\pgfsetroundjoin%
\definecolor{currentfill}{rgb}{0.208623,0.367752,0.552675}%
\pgfsetfillcolor{currentfill}%
\pgfsetlinewidth{1.003750pt}%
\definecolor{currentstroke}{rgb}{0.208623,0.367752,0.552675}%
\pgfsetstrokecolor{currentstroke}%
\pgfsetdash{}{0pt}%
\pgfpathmoveto{\pgfqpoint{3.375893in}{5.095374in}}%
\pgfpathcurveto{\pgfqpoint{3.381717in}{5.095374in}}{\pgfqpoint{3.387303in}{5.097688in}}{\pgfqpoint{3.391421in}{5.101806in}}%
\pgfpathcurveto{\pgfqpoint{3.395539in}{5.105924in}}{\pgfqpoint{3.397853in}{5.111510in}}{\pgfqpoint{3.397853in}{5.117334in}}%
\pgfpathcurveto{\pgfqpoint{3.397853in}{5.123158in}}{\pgfqpoint{3.395539in}{5.128744in}}{\pgfqpoint{3.391421in}{5.132862in}}%
\pgfpathcurveto{\pgfqpoint{3.387303in}{5.136980in}}{\pgfqpoint{3.381717in}{5.139294in}}{\pgfqpoint{3.375893in}{5.139294in}}%
\pgfpathcurveto{\pgfqpoint{3.370069in}{5.139294in}}{\pgfqpoint{3.364483in}{5.136980in}}{\pgfqpoint{3.360365in}{5.132862in}}%
\pgfpathcurveto{\pgfqpoint{3.356246in}{5.128744in}}{\pgfqpoint{3.353933in}{5.123158in}}{\pgfqpoint{3.353933in}{5.117334in}}%
\pgfpathcurveto{\pgfqpoint{3.353933in}{5.111510in}}{\pgfqpoint{3.356246in}{5.105924in}}{\pgfqpoint{3.360365in}{5.101806in}}%
\pgfpathcurveto{\pgfqpoint{3.364483in}{5.097688in}}{\pgfqpoint{3.370069in}{5.095374in}}{\pgfqpoint{3.375893in}{5.095374in}}%
\pgfpathlineto{\pgfqpoint{3.375893in}{5.095374in}}%
\pgfpathclose%
\pgfusepath{stroke,fill}%
\end{pgfscope}%
\begin{pgfscope}%
\pgfpathrectangle{\pgfqpoint{1.414736in}{0.548769in}}{\pgfqpoint{3.497747in}{7.253186in}}%
\pgfusepath{clip}%
\pgfsetbuttcap%
\pgfsetroundjoin%
\definecolor{currentfill}{rgb}{0.188923,0.410910,0.556326}%
\pgfsetfillcolor{currentfill}%
\pgfsetlinewidth{1.003750pt}%
\definecolor{currentstroke}{rgb}{0.188923,0.410910,0.556326}%
\pgfsetstrokecolor{currentstroke}%
\pgfsetdash{}{0pt}%
\pgfpathmoveto{\pgfqpoint{3.685359in}{5.566360in}}%
\pgfpathcurveto{\pgfqpoint{3.691183in}{5.566360in}}{\pgfqpoint{3.696769in}{5.568674in}}{\pgfqpoint{3.700887in}{5.572792in}}%
\pgfpathcurveto{\pgfqpoint{3.705005in}{5.576910in}}{\pgfqpoint{3.707319in}{5.582496in}}{\pgfqpoint{3.707319in}{5.588320in}}%
\pgfpathcurveto{\pgfqpoint{3.707319in}{5.594144in}}{\pgfqpoint{3.705005in}{5.599730in}}{\pgfqpoint{3.700887in}{5.603848in}}%
\pgfpathcurveto{\pgfqpoint{3.696769in}{5.607967in}}{\pgfqpoint{3.691183in}{5.610280in}}{\pgfqpoint{3.685359in}{5.610280in}}%
\pgfpathcurveto{\pgfqpoint{3.679535in}{5.610280in}}{\pgfqpoint{3.673949in}{5.607967in}}{\pgfqpoint{3.669831in}{5.603848in}}%
\pgfpathcurveto{\pgfqpoint{3.665713in}{5.599730in}}{\pgfqpoint{3.663399in}{5.594144in}}{\pgfqpoint{3.663399in}{5.588320in}}%
\pgfpathcurveto{\pgfqpoint{3.663399in}{5.582496in}}{\pgfqpoint{3.665713in}{5.576910in}}{\pgfqpoint{3.669831in}{5.572792in}}%
\pgfpathcurveto{\pgfqpoint{3.673949in}{5.568674in}}{\pgfqpoint{3.679535in}{5.566360in}}{\pgfqpoint{3.685359in}{5.566360in}}%
\pgfpathlineto{\pgfqpoint{3.685359in}{5.566360in}}%
\pgfpathclose%
\pgfusepath{stroke,fill}%
\end{pgfscope}%
\begin{pgfscope}%
\pgfpathrectangle{\pgfqpoint{1.414736in}{0.548769in}}{\pgfqpoint{3.497747in}{7.253186in}}%
\pgfusepath{clip}%
\pgfsetbuttcap%
\pgfsetroundjoin%
\definecolor{currentfill}{rgb}{0.119699,0.618490,0.536347}%
\pgfsetfillcolor{currentfill}%
\pgfsetlinewidth{1.003750pt}%
\definecolor{currentstroke}{rgb}{0.119699,0.618490,0.536347}%
\pgfsetstrokecolor{currentstroke}%
\pgfsetdash{}{0pt}%
\pgfpathmoveto{\pgfqpoint{3.765685in}{6.037346in}}%
\pgfpathcurveto{\pgfqpoint{3.771509in}{6.037346in}}{\pgfqpoint{3.777095in}{6.039660in}}{\pgfqpoint{3.781213in}{6.043778in}}%
\pgfpathcurveto{\pgfqpoint{3.785331in}{6.047896in}}{\pgfqpoint{3.787645in}{6.053482in}}{\pgfqpoint{3.787645in}{6.059306in}}%
\pgfpathcurveto{\pgfqpoint{3.787645in}{6.065130in}}{\pgfqpoint{3.785331in}{6.070716in}}{\pgfqpoint{3.781213in}{6.074835in}}%
\pgfpathcurveto{\pgfqpoint{3.777095in}{6.078953in}}{\pgfqpoint{3.771509in}{6.081267in}}{\pgfqpoint{3.765685in}{6.081267in}}%
\pgfpathcurveto{\pgfqpoint{3.759861in}{6.081267in}}{\pgfqpoint{3.754275in}{6.078953in}}{\pgfqpoint{3.750156in}{6.074835in}}%
\pgfpathcurveto{\pgfqpoint{3.746038in}{6.070716in}}{\pgfqpoint{3.743724in}{6.065130in}}{\pgfqpoint{3.743724in}{6.059306in}}%
\pgfpathcurveto{\pgfqpoint{3.743724in}{6.053482in}}{\pgfqpoint{3.746038in}{6.047896in}}{\pgfqpoint{3.750156in}{6.043778in}}%
\pgfpathcurveto{\pgfqpoint{3.754275in}{6.039660in}}{\pgfqpoint{3.759861in}{6.037346in}}{\pgfqpoint{3.765685in}{6.037346in}}%
\pgfpathlineto{\pgfqpoint{3.765685in}{6.037346in}}%
\pgfpathclose%
\pgfusepath{stroke,fill}%
\end{pgfscope}%
\begin{pgfscope}%
\pgfpathrectangle{\pgfqpoint{1.414736in}{0.548769in}}{\pgfqpoint{3.497747in}{7.253186in}}%
\pgfusepath{clip}%
\pgfsetbuttcap%
\pgfsetroundjoin%
\definecolor{currentfill}{rgb}{0.288921,0.758394,0.428426}%
\pgfsetfillcolor{currentfill}%
\pgfsetlinewidth{1.003750pt}%
\definecolor{currentstroke}{rgb}{0.288921,0.758394,0.428426}%
\pgfsetstrokecolor{currentstroke}%
\pgfsetdash{}{0pt}%
\pgfpathmoveto{\pgfqpoint{3.723507in}{6.508332in}}%
\pgfpathcurveto{\pgfqpoint{3.729331in}{6.508332in}}{\pgfqpoint{3.734917in}{6.510646in}}{\pgfqpoint{3.739035in}{6.514764in}}%
\pgfpathcurveto{\pgfqpoint{3.743153in}{6.518882in}}{\pgfqpoint{3.745467in}{6.524468in}}{\pgfqpoint{3.745467in}{6.530292in}}%
\pgfpathcurveto{\pgfqpoint{3.745467in}{6.536116in}}{\pgfqpoint{3.743153in}{6.541702in}}{\pgfqpoint{3.739035in}{6.545821in}}%
\pgfpathcurveto{\pgfqpoint{3.734917in}{6.549939in}}{\pgfqpoint{3.729331in}{6.552253in}}{\pgfqpoint{3.723507in}{6.552253in}}%
\pgfpathcurveto{\pgfqpoint{3.717683in}{6.552253in}}{\pgfqpoint{3.712097in}{6.549939in}}{\pgfqpoint{3.707979in}{6.545821in}}%
\pgfpathcurveto{\pgfqpoint{3.703860in}{6.541702in}}{\pgfqpoint{3.701546in}{6.536116in}}{\pgfqpoint{3.701546in}{6.530292in}}%
\pgfpathcurveto{\pgfqpoint{3.701546in}{6.524468in}}{\pgfqpoint{3.703860in}{6.518882in}}{\pgfqpoint{3.707979in}{6.514764in}}%
\pgfpathcurveto{\pgfqpoint{3.712097in}{6.510646in}}{\pgfqpoint{3.717683in}{6.508332in}}{\pgfqpoint{3.723507in}{6.508332in}}%
\pgfpathlineto{\pgfqpoint{3.723507in}{6.508332in}}%
\pgfpathclose%
\pgfusepath{stroke,fill}%
\end{pgfscope}%
\begin{pgfscope}%
\pgfpathrectangle{\pgfqpoint{1.414736in}{0.548769in}}{\pgfqpoint{3.497747in}{7.253186in}}%
\pgfusepath{clip}%
\pgfsetbuttcap%
\pgfsetroundjoin%
\definecolor{currentfill}{rgb}{0.404001,0.800275,0.362552}%
\pgfsetfillcolor{currentfill}%
\pgfsetlinewidth{1.003750pt}%
\definecolor{currentstroke}{rgb}{0.404001,0.800275,0.362552}%
\pgfsetstrokecolor{currentstroke}%
\pgfsetdash{}{0pt}%
\pgfpathmoveto{\pgfqpoint{4.304138in}{6.979318in}}%
\pgfpathcurveto{\pgfqpoint{4.309962in}{6.979318in}}{\pgfqpoint{4.315548in}{6.981632in}}{\pgfqpoint{4.319666in}{6.985750in}}%
\pgfpathcurveto{\pgfqpoint{4.323784in}{6.989868in}}{\pgfqpoint{4.326098in}{6.995455in}}{\pgfqpoint{4.326098in}{7.001278in}}%
\pgfpathcurveto{\pgfqpoint{4.326098in}{7.007102in}}{\pgfqpoint{4.323784in}{7.012689in}}{\pgfqpoint{4.319666in}{7.016807in}}%
\pgfpathcurveto{\pgfqpoint{4.315548in}{7.020925in}}{\pgfqpoint{4.309962in}{7.023239in}}{\pgfqpoint{4.304138in}{7.023239in}}%
\pgfpathcurveto{\pgfqpoint{4.298314in}{7.023239in}}{\pgfqpoint{4.292728in}{7.020925in}}{\pgfqpoint{4.288610in}{7.016807in}}%
\pgfpathcurveto{\pgfqpoint{4.284491in}{7.012689in}}{\pgfqpoint{4.282178in}{7.007102in}}{\pgfqpoint{4.282178in}{7.001278in}}%
\pgfpathcurveto{\pgfqpoint{4.282178in}{6.995455in}}{\pgfqpoint{4.284491in}{6.989868in}}{\pgfqpoint{4.288610in}{6.985750in}}%
\pgfpathcurveto{\pgfqpoint{4.292728in}{6.981632in}}{\pgfqpoint{4.298314in}{6.979318in}}{\pgfqpoint{4.304138in}{6.979318in}}%
\pgfpathlineto{\pgfqpoint{4.304138in}{6.979318in}}%
\pgfpathclose%
\pgfusepath{stroke,fill}%
\end{pgfscope}%
\begin{pgfscope}%
\pgfpathrectangle{\pgfqpoint{1.414736in}{0.548769in}}{\pgfqpoint{3.497747in}{7.253186in}}%
\pgfusepath{clip}%
\pgfsetbuttcap%
\pgfsetroundjoin%
\definecolor{currentfill}{rgb}{0.626579,0.854645,0.223353}%
\pgfsetfillcolor{currentfill}%
\pgfsetlinewidth{1.003750pt}%
\definecolor{currentstroke}{rgb}{0.626579,0.854645,0.223353}%
\pgfsetstrokecolor{currentstroke}%
\pgfsetdash{}{0pt}%
\pgfpathmoveto{\pgfqpoint{3.568468in}{7.450304in}}%
\pgfpathcurveto{\pgfqpoint{3.574292in}{7.450304in}}{\pgfqpoint{3.579878in}{7.452618in}}{\pgfqpoint{3.583996in}{7.456736in}}%
\pgfpathcurveto{\pgfqpoint{3.588115in}{7.460854in}}{\pgfqpoint{3.590428in}{7.466441in}}{\pgfqpoint{3.590428in}{7.472265in}}%
\pgfpathcurveto{\pgfqpoint{3.590428in}{7.478089in}}{\pgfqpoint{3.588115in}{7.483675in}}{\pgfqpoint{3.583996in}{7.487793in}}%
\pgfpathcurveto{\pgfqpoint{3.579878in}{7.491911in}}{\pgfqpoint{3.574292in}{7.494225in}}{\pgfqpoint{3.568468in}{7.494225in}}%
\pgfpathcurveto{\pgfqpoint{3.562644in}{7.494225in}}{\pgfqpoint{3.557058in}{7.491911in}}{\pgfqpoint{3.552940in}{7.487793in}}%
\pgfpathcurveto{\pgfqpoint{3.548822in}{7.483675in}}{\pgfqpoint{3.546508in}{7.478089in}}{\pgfqpoint{3.546508in}{7.472265in}}%
\pgfpathcurveto{\pgfqpoint{3.546508in}{7.466441in}}{\pgfqpoint{3.548822in}{7.460854in}}{\pgfqpoint{3.552940in}{7.456736in}}%
\pgfpathcurveto{\pgfqpoint{3.557058in}{7.452618in}}{\pgfqpoint{3.562644in}{7.450304in}}{\pgfqpoint{3.568468in}{7.450304in}}%
\pgfpathlineto{\pgfqpoint{3.568468in}{7.450304in}}%
\pgfpathclose%
\pgfusepath{stroke,fill}%
\end{pgfscope}%
\begin{pgfscope}%
\pgfpathrectangle{\pgfqpoint{1.414736in}{0.548769in}}{\pgfqpoint{3.497747in}{7.253186in}}%
\pgfusepath{clip}%
\pgfsetbuttcap%
\pgfsetroundjoin%
\definecolor{currentfill}{rgb}{0.993248,0.906157,0.143936}%
\pgfsetfillcolor{currentfill}%
\pgfsetlinewidth{1.003750pt}%
\definecolor{currentstroke}{rgb}{0.993248,0.906157,0.143936}%
\pgfsetstrokecolor{currentstroke}%
\pgfsetdash{}{0pt}%
\pgfpathmoveto{\pgfqpoint{3.757484in}{0.856499in}}%
\pgfpathcurveto{\pgfqpoint{3.763308in}{0.856499in}}{\pgfqpoint{3.768895in}{0.858813in}}{\pgfqpoint{3.773013in}{0.862931in}}%
\pgfpathcurveto{\pgfqpoint{3.777131in}{0.867049in}}{\pgfqpoint{3.779445in}{0.872635in}}{\pgfqpoint{3.779445in}{0.878459in}}%
\pgfpathcurveto{\pgfqpoint{3.779445in}{0.884283in}}{\pgfqpoint{3.777131in}{0.889869in}}{\pgfqpoint{3.773013in}{0.893987in}}%
\pgfpathcurveto{\pgfqpoint{3.768895in}{0.898105in}}{\pgfqpoint{3.763308in}{0.900419in}}{\pgfqpoint{3.757484in}{0.900419in}}%
\pgfpathcurveto{\pgfqpoint{3.751661in}{0.900419in}}{\pgfqpoint{3.746074in}{0.898105in}}{\pgfqpoint{3.741956in}{0.893987in}}%
\pgfpathcurveto{\pgfqpoint{3.737838in}{0.889869in}}{\pgfqpoint{3.735524in}{0.884283in}}{\pgfqpoint{3.735524in}{0.878459in}}%
\pgfpathcurveto{\pgfqpoint{3.735524in}{0.872635in}}{\pgfqpoint{3.737838in}{0.867049in}}{\pgfqpoint{3.741956in}{0.862931in}}%
\pgfpathcurveto{\pgfqpoint{3.746074in}{0.858813in}}{\pgfqpoint{3.751661in}{0.856499in}}{\pgfqpoint{3.757484in}{0.856499in}}%
\pgfpathlineto{\pgfqpoint{3.757484in}{0.856499in}}%
\pgfpathclose%
\pgfusepath{stroke,fill}%
\end{pgfscope}%
\begin{pgfscope}%
\pgfpathrectangle{\pgfqpoint{1.414736in}{0.548769in}}{\pgfqpoint{3.497747in}{7.253186in}}%
\pgfusepath{clip}%
\pgfsetbuttcap%
\pgfsetroundjoin%
\definecolor{currentfill}{rgb}{0.741388,0.873449,0.149561}%
\pgfsetfillcolor{currentfill}%
\pgfsetlinewidth{1.003750pt}%
\definecolor{currentstroke}{rgb}{0.741388,0.873449,0.149561}%
\pgfsetstrokecolor{currentstroke}%
\pgfsetdash{}{0pt}%
\pgfpathmoveto{\pgfqpoint{3.446162in}{1.327485in}}%
\pgfpathcurveto{\pgfqpoint{3.451986in}{1.327485in}}{\pgfqpoint{3.457572in}{1.329799in}}{\pgfqpoint{3.461690in}{1.333917in}}%
\pgfpathcurveto{\pgfqpoint{3.465808in}{1.338035in}}{\pgfqpoint{3.468122in}{1.343621in}}{\pgfqpoint{3.468122in}{1.349445in}}%
\pgfpathcurveto{\pgfqpoint{3.468122in}{1.355269in}}{\pgfqpoint{3.465808in}{1.360855in}}{\pgfqpoint{3.461690in}{1.364973in}}%
\pgfpathcurveto{\pgfqpoint{3.457572in}{1.369092in}}{\pgfqpoint{3.451986in}{1.371405in}}{\pgfqpoint{3.446162in}{1.371405in}}%
\pgfpathcurveto{\pgfqpoint{3.440338in}{1.371405in}}{\pgfqpoint{3.434752in}{1.369092in}}{\pgfqpoint{3.430634in}{1.364973in}}%
\pgfpathcurveto{\pgfqpoint{3.426515in}{1.360855in}}{\pgfqpoint{3.424202in}{1.355269in}}{\pgfqpoint{3.424202in}{1.349445in}}%
\pgfpathcurveto{\pgfqpoint{3.424202in}{1.343621in}}{\pgfqpoint{3.426515in}{1.338035in}}{\pgfqpoint{3.430634in}{1.333917in}}%
\pgfpathcurveto{\pgfqpoint{3.434752in}{1.329799in}}{\pgfqpoint{3.440338in}{1.327485in}}{\pgfqpoint{3.446162in}{1.327485in}}%
\pgfpathlineto{\pgfqpoint{3.446162in}{1.327485in}}%
\pgfpathclose%
\pgfusepath{stroke,fill}%
\end{pgfscope}%
\begin{pgfscope}%
\pgfpathrectangle{\pgfqpoint{1.414736in}{0.548769in}}{\pgfqpoint{3.497747in}{7.253186in}}%
\pgfusepath{clip}%
\pgfsetbuttcap%
\pgfsetroundjoin%
\definecolor{currentfill}{rgb}{0.699415,0.867117,0.175971}%
\pgfsetfillcolor{currentfill}%
\pgfsetlinewidth{1.003750pt}%
\definecolor{currentstroke}{rgb}{0.699415,0.867117,0.175971}%
\pgfsetstrokecolor{currentstroke}%
\pgfsetdash{}{0pt}%
\pgfpathmoveto{\pgfqpoint{3.700481in}{1.798471in}}%
\pgfpathcurveto{\pgfqpoint{3.706305in}{1.798471in}}{\pgfqpoint{3.711891in}{1.800785in}}{\pgfqpoint{3.716010in}{1.804903in}}%
\pgfpathcurveto{\pgfqpoint{3.720128in}{1.809021in}}{\pgfqpoint{3.722442in}{1.814607in}}{\pgfqpoint{3.722442in}{1.820431in}}%
\pgfpathcurveto{\pgfqpoint{3.722442in}{1.826255in}}{\pgfqpoint{3.720128in}{1.831841in}}{\pgfqpoint{3.716010in}{1.835960in}}%
\pgfpathcurveto{\pgfqpoint{3.711891in}{1.840078in}}{\pgfqpoint{3.706305in}{1.842392in}}{\pgfqpoint{3.700481in}{1.842392in}}%
\pgfpathcurveto{\pgfqpoint{3.694657in}{1.842392in}}{\pgfqpoint{3.689071in}{1.840078in}}{\pgfqpoint{3.684953in}{1.835960in}}%
\pgfpathcurveto{\pgfqpoint{3.680835in}{1.831841in}}{\pgfqpoint{3.678521in}{1.826255in}}{\pgfqpoint{3.678521in}{1.820431in}}%
\pgfpathcurveto{\pgfqpoint{3.678521in}{1.814607in}}{\pgfqpoint{3.680835in}{1.809021in}}{\pgfqpoint{3.684953in}{1.804903in}}%
\pgfpathcurveto{\pgfqpoint{3.689071in}{1.800785in}}{\pgfqpoint{3.694657in}{1.798471in}}{\pgfqpoint{3.700481in}{1.798471in}}%
\pgfpathlineto{\pgfqpoint{3.700481in}{1.798471in}}%
\pgfpathclose%
\pgfusepath{stroke,fill}%
\end{pgfscope}%
\begin{pgfscope}%
\pgfpathrectangle{\pgfqpoint{1.414736in}{0.548769in}}{\pgfqpoint{3.497747in}{7.253186in}}%
\pgfusepath{clip}%
\pgfsetbuttcap%
\pgfsetroundjoin%
\definecolor{currentfill}{rgb}{0.993248,0.906157,0.143936}%
\pgfsetfillcolor{currentfill}%
\pgfsetlinewidth{1.003750pt}%
\definecolor{currentstroke}{rgb}{0.993248,0.906157,0.143936}%
\pgfsetstrokecolor{currentstroke}%
\pgfsetdash{}{0pt}%
\pgfpathmoveto{\pgfqpoint{3.440788in}{2.269457in}}%
\pgfpathcurveto{\pgfqpoint{3.446612in}{2.269457in}}{\pgfqpoint{3.452198in}{2.271771in}}{\pgfqpoint{3.456316in}{2.275889in}}%
\pgfpathcurveto{\pgfqpoint{3.460435in}{2.280007in}}{\pgfqpoint{3.462748in}{2.285593in}}{\pgfqpoint{3.462748in}{2.291417in}}%
\pgfpathcurveto{\pgfqpoint{3.462748in}{2.297241in}}{\pgfqpoint{3.460435in}{2.302828in}}{\pgfqpoint{3.456316in}{2.306946in}}%
\pgfpathcurveto{\pgfqpoint{3.452198in}{2.311064in}}{\pgfqpoint{3.446612in}{2.313378in}}{\pgfqpoint{3.440788in}{2.313378in}}%
\pgfpathcurveto{\pgfqpoint{3.434964in}{2.313378in}}{\pgfqpoint{3.429378in}{2.311064in}}{\pgfqpoint{3.425260in}{2.306946in}}%
\pgfpathcurveto{\pgfqpoint{3.421142in}{2.302828in}}{\pgfqpoint{3.418828in}{2.297241in}}{\pgfqpoint{3.418828in}{2.291417in}}%
\pgfpathcurveto{\pgfqpoint{3.418828in}{2.285593in}}{\pgfqpoint{3.421142in}{2.280007in}}{\pgfqpoint{3.425260in}{2.275889in}}%
\pgfpathcurveto{\pgfqpoint{3.429378in}{2.271771in}}{\pgfqpoint{3.434964in}{2.269457in}}{\pgfqpoint{3.440788in}{2.269457in}}%
\pgfpathlineto{\pgfqpoint{3.440788in}{2.269457in}}%
\pgfpathclose%
\pgfusepath{stroke,fill}%
\end{pgfscope}%
\begin{pgfscope}%
\pgfpathrectangle{\pgfqpoint{1.414736in}{0.548769in}}{\pgfqpoint{3.497747in}{7.253186in}}%
\pgfusepath{clip}%
\pgfsetbuttcap%
\pgfsetroundjoin%
\definecolor{currentfill}{rgb}{0.274149,0.751988,0.436601}%
\pgfsetfillcolor{currentfill}%
\pgfsetlinewidth{1.003750pt}%
\definecolor{currentstroke}{rgb}{0.274149,0.751988,0.436601}%
\pgfsetstrokecolor{currentstroke}%
\pgfsetdash{}{0pt}%
\pgfpathmoveto{\pgfqpoint{3.733094in}{2.740443in}}%
\pgfpathcurveto{\pgfqpoint{3.738918in}{2.740443in}}{\pgfqpoint{3.744504in}{2.742757in}}{\pgfqpoint{3.748622in}{2.746875in}}%
\pgfpathcurveto{\pgfqpoint{3.752740in}{2.750993in}}{\pgfqpoint{3.755054in}{2.756580in}}{\pgfqpoint{3.755054in}{2.762404in}}%
\pgfpathcurveto{\pgfqpoint{3.755054in}{2.768227in}}{\pgfqpoint{3.752740in}{2.773814in}}{\pgfqpoint{3.748622in}{2.777932in}}%
\pgfpathcurveto{\pgfqpoint{3.744504in}{2.782050in}}{\pgfqpoint{3.738918in}{2.784364in}}{\pgfqpoint{3.733094in}{2.784364in}}%
\pgfpathcurveto{\pgfqpoint{3.727270in}{2.784364in}}{\pgfqpoint{3.721684in}{2.782050in}}{\pgfqpoint{3.717566in}{2.777932in}}%
\pgfpathcurveto{\pgfqpoint{3.713447in}{2.773814in}}{\pgfqpoint{3.711134in}{2.768227in}}{\pgfqpoint{3.711134in}{2.762404in}}%
\pgfpathcurveto{\pgfqpoint{3.711134in}{2.756580in}}{\pgfqpoint{3.713447in}{2.750993in}}{\pgfqpoint{3.717566in}{2.746875in}}%
\pgfpathcurveto{\pgfqpoint{3.721684in}{2.742757in}}{\pgfqpoint{3.727270in}{2.740443in}}{\pgfqpoint{3.733094in}{2.740443in}}%
\pgfpathlineto{\pgfqpoint{3.733094in}{2.740443in}}%
\pgfpathclose%
\pgfusepath{stroke,fill}%
\end{pgfscope}%
\begin{pgfscope}%
\pgfpathrectangle{\pgfqpoint{1.414736in}{0.548769in}}{\pgfqpoint{3.497747in}{7.253186in}}%
\pgfusepath{clip}%
\pgfsetbuttcap%
\pgfsetroundjoin%
\definecolor{currentfill}{rgb}{0.855810,0.888601,0.097452}%
\pgfsetfillcolor{currentfill}%
\pgfsetlinewidth{1.003750pt}%
\definecolor{currentstroke}{rgb}{0.855810,0.888601,0.097452}%
\pgfsetstrokecolor{currentstroke}%
\pgfsetdash{}{0pt}%
\pgfpathmoveto{\pgfqpoint{3.444977in}{3.211429in}}%
\pgfpathcurveto{\pgfqpoint{3.450801in}{3.211429in}}{\pgfqpoint{3.456387in}{3.213743in}}{\pgfqpoint{3.460505in}{3.217861in}}%
\pgfpathcurveto{\pgfqpoint{3.464623in}{3.221980in}}{\pgfqpoint{3.466937in}{3.227566in}}{\pgfqpoint{3.466937in}{3.233390in}}%
\pgfpathcurveto{\pgfqpoint{3.466937in}{3.239214in}}{\pgfqpoint{3.464623in}{3.244800in}}{\pgfqpoint{3.460505in}{3.248918in}}%
\pgfpathcurveto{\pgfqpoint{3.456387in}{3.253036in}}{\pgfqpoint{3.450801in}{3.255350in}}{\pgfqpoint{3.444977in}{3.255350in}}%
\pgfpathcurveto{\pgfqpoint{3.439153in}{3.255350in}}{\pgfqpoint{3.433567in}{3.253036in}}{\pgfqpoint{3.429449in}{3.248918in}}%
\pgfpathcurveto{\pgfqpoint{3.425331in}{3.244800in}}{\pgfqpoint{3.423017in}{3.239214in}}{\pgfqpoint{3.423017in}{3.233390in}}%
\pgfpathcurveto{\pgfqpoint{3.423017in}{3.227566in}}{\pgfqpoint{3.425331in}{3.221980in}}{\pgfqpoint{3.429449in}{3.217861in}}%
\pgfpathcurveto{\pgfqpoint{3.433567in}{3.213743in}}{\pgfqpoint{3.439153in}{3.211429in}}{\pgfqpoint{3.444977in}{3.211429in}}%
\pgfpathlineto{\pgfqpoint{3.444977in}{3.211429in}}%
\pgfpathclose%
\pgfusepath{stroke,fill}%
\end{pgfscope}%
\begin{pgfscope}%
\pgfpathrectangle{\pgfqpoint{1.414736in}{0.548769in}}{\pgfqpoint{3.497747in}{7.253186in}}%
\pgfusepath{clip}%
\pgfsetbuttcap%
\pgfsetroundjoin%
\definecolor{currentfill}{rgb}{0.993248,0.906157,0.143936}%
\pgfsetfillcolor{currentfill}%
\pgfsetlinewidth{1.003750pt}%
\definecolor{currentstroke}{rgb}{0.993248,0.906157,0.143936}%
\pgfsetstrokecolor{currentstroke}%
\pgfsetdash{}{0pt}%
\pgfpathmoveto{\pgfqpoint{3.535652in}{3.682415in}}%
\pgfpathcurveto{\pgfqpoint{3.541476in}{3.682415in}}{\pgfqpoint{3.547062in}{3.684729in}}{\pgfqpoint{3.551180in}{3.688847in}}%
\pgfpathcurveto{\pgfqpoint{3.555298in}{3.692966in}}{\pgfqpoint{3.557612in}{3.698552in}}{\pgfqpoint{3.557612in}{3.704376in}}%
\pgfpathcurveto{\pgfqpoint{3.557612in}{3.710200in}}{\pgfqpoint{3.555298in}{3.715786in}}{\pgfqpoint{3.551180in}{3.719904in}}%
\pgfpathcurveto{\pgfqpoint{3.547062in}{3.724022in}}{\pgfqpoint{3.541476in}{3.726336in}}{\pgfqpoint{3.535652in}{3.726336in}}%
\pgfpathcurveto{\pgfqpoint{3.529828in}{3.726336in}}{\pgfqpoint{3.524242in}{3.724022in}}{\pgfqpoint{3.520123in}{3.719904in}}%
\pgfpathcurveto{\pgfqpoint{3.516005in}{3.715786in}}{\pgfqpoint{3.513691in}{3.710200in}}{\pgfqpoint{3.513691in}{3.704376in}}%
\pgfpathcurveto{\pgfqpoint{3.513691in}{3.698552in}}{\pgfqpoint{3.516005in}{3.692966in}}{\pgfqpoint{3.520123in}{3.688847in}}%
\pgfpathcurveto{\pgfqpoint{3.524242in}{3.684729in}}{\pgfqpoint{3.529828in}{3.682415in}}{\pgfqpoint{3.535652in}{3.682415in}}%
\pgfpathlineto{\pgfqpoint{3.535652in}{3.682415in}}%
\pgfpathclose%
\pgfusepath{stroke,fill}%
\end{pgfscope}%
\begin{pgfscope}%
\pgfpathrectangle{\pgfqpoint{1.414736in}{0.548769in}}{\pgfqpoint{3.497747in}{7.253186in}}%
\pgfusepath{clip}%
\pgfsetbuttcap%
\pgfsetroundjoin%
\definecolor{currentfill}{rgb}{0.793760,0.880678,0.120005}%
\pgfsetfillcolor{currentfill}%
\pgfsetlinewidth{1.003750pt}%
\definecolor{currentstroke}{rgb}{0.793760,0.880678,0.120005}%
\pgfsetstrokecolor{currentstroke}%
\pgfsetdash{}{0pt}%
\pgfpathmoveto{\pgfqpoint{3.812572in}{4.153402in}}%
\pgfpathcurveto{\pgfqpoint{3.818396in}{4.153402in}}{\pgfqpoint{3.823982in}{4.155715in}}{\pgfqpoint{3.828100in}{4.159834in}}%
\pgfpathcurveto{\pgfqpoint{3.832218in}{4.163952in}}{\pgfqpoint{3.834532in}{4.169538in}}{\pgfqpoint{3.834532in}{4.175362in}}%
\pgfpathcurveto{\pgfqpoint{3.834532in}{4.181186in}}{\pgfqpoint{3.832218in}{4.186772in}}{\pgfqpoint{3.828100in}{4.190890in}}%
\pgfpathcurveto{\pgfqpoint{3.823982in}{4.195008in}}{\pgfqpoint{3.818396in}{4.197322in}}{\pgfqpoint{3.812572in}{4.197322in}}%
\pgfpathcurveto{\pgfqpoint{3.806748in}{4.197322in}}{\pgfqpoint{3.801162in}{4.195008in}}{\pgfqpoint{3.797044in}{4.190890in}}%
\pgfpathcurveto{\pgfqpoint{3.792925in}{4.186772in}}{\pgfqpoint{3.790612in}{4.181186in}}{\pgfqpoint{3.790612in}{4.175362in}}%
\pgfpathcurveto{\pgfqpoint{3.790612in}{4.169538in}}{\pgfqpoint{3.792925in}{4.163952in}}{\pgfqpoint{3.797044in}{4.159834in}}%
\pgfpathcurveto{\pgfqpoint{3.801162in}{4.155715in}}{\pgfqpoint{3.806748in}{4.153402in}}{\pgfqpoint{3.812572in}{4.153402in}}%
\pgfpathlineto{\pgfqpoint{3.812572in}{4.153402in}}%
\pgfpathclose%
\pgfusepath{stroke,fill}%
\end{pgfscope}%
\begin{pgfscope}%
\pgfpathrectangle{\pgfqpoint{1.414736in}{0.548769in}}{\pgfqpoint{3.497747in}{7.253186in}}%
\pgfusepath{clip}%
\pgfsetbuttcap%
\pgfsetroundjoin%
\definecolor{currentfill}{rgb}{0.855810,0.888601,0.097452}%
\pgfsetfillcolor{currentfill}%
\pgfsetlinewidth{1.003750pt}%
\definecolor{currentstroke}{rgb}{0.855810,0.888601,0.097452}%
\pgfsetstrokecolor{currentstroke}%
\pgfsetdash{}{0pt}%
\pgfpathmoveto{\pgfqpoint{3.868690in}{4.624388in}}%
\pgfpathcurveto{\pgfqpoint{3.874514in}{4.624388in}}{\pgfqpoint{3.880100in}{4.626702in}}{\pgfqpoint{3.884218in}{4.630820in}}%
\pgfpathcurveto{\pgfqpoint{3.888337in}{4.634938in}}{\pgfqpoint{3.890650in}{4.640524in}}{\pgfqpoint{3.890650in}{4.646348in}}%
\pgfpathcurveto{\pgfqpoint{3.890650in}{4.652172in}}{\pgfqpoint{3.888337in}{4.657758in}}{\pgfqpoint{3.884218in}{4.661876in}}%
\pgfpathcurveto{\pgfqpoint{3.880100in}{4.665994in}}{\pgfqpoint{3.874514in}{4.668308in}}{\pgfqpoint{3.868690in}{4.668308in}}%
\pgfpathcurveto{\pgfqpoint{3.862866in}{4.668308in}}{\pgfqpoint{3.857280in}{4.665994in}}{\pgfqpoint{3.853162in}{4.661876in}}%
\pgfpathcurveto{\pgfqpoint{3.849044in}{4.657758in}}{\pgfqpoint{3.846730in}{4.652172in}}{\pgfqpoint{3.846730in}{4.646348in}}%
\pgfpathcurveto{\pgfqpoint{3.846730in}{4.640524in}}{\pgfqpoint{3.849044in}{4.634938in}}{\pgfqpoint{3.853162in}{4.630820in}}%
\pgfpathcurveto{\pgfqpoint{3.857280in}{4.626702in}}{\pgfqpoint{3.862866in}{4.624388in}}{\pgfqpoint{3.868690in}{4.624388in}}%
\pgfpathlineto{\pgfqpoint{3.868690in}{4.624388in}}%
\pgfpathclose%
\pgfusepath{stroke,fill}%
\end{pgfscope}%
\begin{pgfscope}%
\pgfpathrectangle{\pgfqpoint{1.414736in}{0.548769in}}{\pgfqpoint{3.497747in}{7.253186in}}%
\pgfusepath{clip}%
\pgfsetbuttcap%
\pgfsetroundjoin%
\definecolor{currentfill}{rgb}{0.125394,0.574318,0.549086}%
\pgfsetfillcolor{currentfill}%
\pgfsetlinewidth{1.003750pt}%
\definecolor{currentstroke}{rgb}{0.125394,0.574318,0.549086}%
\pgfsetstrokecolor{currentstroke}%
\pgfsetdash{}{0pt}%
\pgfpathmoveto{\pgfqpoint{3.300628in}{5.095374in}}%
\pgfpathcurveto{\pgfqpoint{3.306452in}{5.095374in}}{\pgfqpoint{3.312038in}{5.097688in}}{\pgfqpoint{3.316156in}{5.101806in}}%
\pgfpathcurveto{\pgfqpoint{3.320274in}{5.105924in}}{\pgfqpoint{3.322588in}{5.111510in}}{\pgfqpoint{3.322588in}{5.117334in}}%
\pgfpathcurveto{\pgfqpoint{3.322588in}{5.123158in}}{\pgfqpoint{3.320274in}{5.128744in}}{\pgfqpoint{3.316156in}{5.132862in}}%
\pgfpathcurveto{\pgfqpoint{3.312038in}{5.136980in}}{\pgfqpoint{3.306452in}{5.139294in}}{\pgfqpoint{3.300628in}{5.139294in}}%
\pgfpathcurveto{\pgfqpoint{3.294804in}{5.139294in}}{\pgfqpoint{3.289218in}{5.136980in}}{\pgfqpoint{3.285100in}{5.132862in}}%
\pgfpathcurveto{\pgfqpoint{3.280982in}{5.128744in}}{\pgfqpoint{3.278668in}{5.123158in}}{\pgfqpoint{3.278668in}{5.117334in}}%
\pgfpathcurveto{\pgfqpoint{3.278668in}{5.111510in}}{\pgfqpoint{3.280982in}{5.105924in}}{\pgfqpoint{3.285100in}{5.101806in}}%
\pgfpathcurveto{\pgfqpoint{3.289218in}{5.097688in}}{\pgfqpoint{3.294804in}{5.095374in}}{\pgfqpoint{3.300628in}{5.095374in}}%
\pgfpathlineto{\pgfqpoint{3.300628in}{5.095374in}}%
\pgfpathclose%
\pgfusepath{stroke,fill}%
\end{pgfscope}%
\begin{pgfscope}%
\pgfpathrectangle{\pgfqpoint{1.414736in}{0.548769in}}{\pgfqpoint{3.497747in}{7.253186in}}%
\pgfusepath{clip}%
\pgfsetbuttcap%
\pgfsetroundjoin%
\definecolor{currentfill}{rgb}{0.150476,0.504369,0.557430}%
\pgfsetfillcolor{currentfill}%
\pgfsetlinewidth{1.003750pt}%
\definecolor{currentstroke}{rgb}{0.150476,0.504369,0.557430}%
\pgfsetstrokecolor{currentstroke}%
\pgfsetdash{}{0pt}%
\pgfpathmoveto{\pgfqpoint{3.870863in}{5.566360in}}%
\pgfpathcurveto{\pgfqpoint{3.876687in}{5.566360in}}{\pgfqpoint{3.882273in}{5.568674in}}{\pgfqpoint{3.886391in}{5.572792in}}%
\pgfpathcurveto{\pgfqpoint{3.890509in}{5.576910in}}{\pgfqpoint{3.892823in}{5.582496in}}{\pgfqpoint{3.892823in}{5.588320in}}%
\pgfpathcurveto{\pgfqpoint{3.892823in}{5.594144in}}{\pgfqpoint{3.890509in}{5.599730in}}{\pgfqpoint{3.886391in}{5.603848in}}%
\pgfpathcurveto{\pgfqpoint{3.882273in}{5.607967in}}{\pgfqpoint{3.876687in}{5.610280in}}{\pgfqpoint{3.870863in}{5.610280in}}%
\pgfpathcurveto{\pgfqpoint{3.865039in}{5.610280in}}{\pgfqpoint{3.859453in}{5.607967in}}{\pgfqpoint{3.855334in}{5.603848in}}%
\pgfpathcurveto{\pgfqpoint{3.851216in}{5.599730in}}{\pgfqpoint{3.848902in}{5.594144in}}{\pgfqpoint{3.848902in}{5.588320in}}%
\pgfpathcurveto{\pgfqpoint{3.848902in}{5.582496in}}{\pgfqpoint{3.851216in}{5.576910in}}{\pgfqpoint{3.855334in}{5.572792in}}%
\pgfpathcurveto{\pgfqpoint{3.859453in}{5.568674in}}{\pgfqpoint{3.865039in}{5.566360in}}{\pgfqpoint{3.870863in}{5.566360in}}%
\pgfpathlineto{\pgfqpoint{3.870863in}{5.566360in}}%
\pgfpathclose%
\pgfusepath{stroke,fill}%
\end{pgfscope}%
\begin{pgfscope}%
\pgfpathrectangle{\pgfqpoint{1.414736in}{0.548769in}}{\pgfqpoint{3.497747in}{7.253186in}}%
\pgfusepath{clip}%
\pgfsetbuttcap%
\pgfsetroundjoin%
\definecolor{currentfill}{rgb}{0.412913,0.803041,0.357269}%
\pgfsetfillcolor{currentfill}%
\pgfsetlinewidth{1.003750pt}%
\definecolor{currentstroke}{rgb}{0.412913,0.803041,0.357269}%
\pgfsetstrokecolor{currentstroke}%
\pgfsetdash{}{0pt}%
\pgfpathmoveto{\pgfqpoint{3.388951in}{6.037346in}}%
\pgfpathcurveto{\pgfqpoint{3.394775in}{6.037346in}}{\pgfqpoint{3.400361in}{6.039660in}}{\pgfqpoint{3.404480in}{6.043778in}}%
\pgfpathcurveto{\pgfqpoint{3.408598in}{6.047896in}}{\pgfqpoint{3.410912in}{6.053482in}}{\pgfqpoint{3.410912in}{6.059306in}}%
\pgfpathcurveto{\pgfqpoint{3.410912in}{6.065130in}}{\pgfqpoint{3.408598in}{6.070716in}}{\pgfqpoint{3.404480in}{6.074835in}}%
\pgfpathcurveto{\pgfqpoint{3.400361in}{6.078953in}}{\pgfqpoint{3.394775in}{6.081267in}}{\pgfqpoint{3.388951in}{6.081267in}}%
\pgfpathcurveto{\pgfqpoint{3.383127in}{6.081267in}}{\pgfqpoint{3.377541in}{6.078953in}}{\pgfqpoint{3.373423in}{6.074835in}}%
\pgfpathcurveto{\pgfqpoint{3.369305in}{6.070716in}}{\pgfqpoint{3.366991in}{6.065130in}}{\pgfqpoint{3.366991in}{6.059306in}}%
\pgfpathcurveto{\pgfqpoint{3.366991in}{6.053482in}}{\pgfqpoint{3.369305in}{6.047896in}}{\pgfqpoint{3.373423in}{6.043778in}}%
\pgfpathcurveto{\pgfqpoint{3.377541in}{6.039660in}}{\pgfqpoint{3.383127in}{6.037346in}}{\pgfqpoint{3.388951in}{6.037346in}}%
\pgfpathlineto{\pgfqpoint{3.388951in}{6.037346in}}%
\pgfpathclose%
\pgfusepath{stroke,fill}%
\end{pgfscope}%
\begin{pgfscope}%
\pgfpathrectangle{\pgfqpoint{1.414736in}{0.548769in}}{\pgfqpoint{3.497747in}{7.253186in}}%
\pgfusepath{clip}%
\pgfsetbuttcap%
\pgfsetroundjoin%
\definecolor{currentfill}{rgb}{0.974417,0.903590,0.130215}%
\pgfsetfillcolor{currentfill}%
\pgfsetlinewidth{1.003750pt}%
\definecolor{currentstroke}{rgb}{0.974417,0.903590,0.130215}%
\pgfsetstrokecolor{currentstroke}%
\pgfsetdash{}{0pt}%
\pgfpathmoveto{\pgfqpoint{3.997787in}{6.508332in}}%
\pgfpathcurveto{\pgfqpoint{4.003611in}{6.508332in}}{\pgfqpoint{4.009197in}{6.510646in}}{\pgfqpoint{4.013315in}{6.514764in}}%
\pgfpathcurveto{\pgfqpoint{4.017433in}{6.518882in}}{\pgfqpoint{4.019747in}{6.524468in}}{\pgfqpoint{4.019747in}{6.530292in}}%
\pgfpathcurveto{\pgfqpoint{4.019747in}{6.536116in}}{\pgfqpoint{4.017433in}{6.541702in}}{\pgfqpoint{4.013315in}{6.545821in}}%
\pgfpathcurveto{\pgfqpoint{4.009197in}{6.549939in}}{\pgfqpoint{4.003611in}{6.552253in}}{\pgfqpoint{3.997787in}{6.552253in}}%
\pgfpathcurveto{\pgfqpoint{3.991963in}{6.552253in}}{\pgfqpoint{3.986377in}{6.549939in}}{\pgfqpoint{3.982259in}{6.545821in}}%
\pgfpathcurveto{\pgfqpoint{3.978141in}{6.541702in}}{\pgfqpoint{3.975827in}{6.536116in}}{\pgfqpoint{3.975827in}{6.530292in}}%
\pgfpathcurveto{\pgfqpoint{3.975827in}{6.524468in}}{\pgfqpoint{3.978141in}{6.518882in}}{\pgfqpoint{3.982259in}{6.514764in}}%
\pgfpathcurveto{\pgfqpoint{3.986377in}{6.510646in}}{\pgfqpoint{3.991963in}{6.508332in}}{\pgfqpoint{3.997787in}{6.508332in}}%
\pgfpathlineto{\pgfqpoint{3.997787in}{6.508332in}}%
\pgfpathclose%
\pgfusepath{stroke,fill}%
\end{pgfscope}%
\begin{pgfscope}%
\pgfpathrectangle{\pgfqpoint{1.414736in}{0.548769in}}{\pgfqpoint{3.497747in}{7.253186in}}%
\pgfusepath{clip}%
\pgfsetbuttcap%
\pgfsetroundjoin%
\definecolor{currentfill}{rgb}{0.699415,0.867117,0.175971}%
\pgfsetfillcolor{currentfill}%
\pgfsetlinewidth{1.003750pt}%
\definecolor{currentstroke}{rgb}{0.699415,0.867117,0.175971}%
\pgfsetstrokecolor{currentstroke}%
\pgfsetdash{}{0pt}%
\pgfpathmoveto{\pgfqpoint{3.838420in}{6.979318in}}%
\pgfpathcurveto{\pgfqpoint{3.844244in}{6.979318in}}{\pgfqpoint{3.849830in}{6.981632in}}{\pgfqpoint{3.853948in}{6.985750in}}%
\pgfpathcurveto{\pgfqpoint{3.858066in}{6.989868in}}{\pgfqpoint{3.860380in}{6.995455in}}{\pgfqpoint{3.860380in}{7.001278in}}%
\pgfpathcurveto{\pgfqpoint{3.860380in}{7.007102in}}{\pgfqpoint{3.858066in}{7.012689in}}{\pgfqpoint{3.853948in}{7.016807in}}%
\pgfpathcurveto{\pgfqpoint{3.849830in}{7.020925in}}{\pgfqpoint{3.844244in}{7.023239in}}{\pgfqpoint{3.838420in}{7.023239in}}%
\pgfpathcurveto{\pgfqpoint{3.832596in}{7.023239in}}{\pgfqpoint{3.827010in}{7.020925in}}{\pgfqpoint{3.822891in}{7.016807in}}%
\pgfpathcurveto{\pgfqpoint{3.818773in}{7.012689in}}{\pgfqpoint{3.816459in}{7.007102in}}{\pgfqpoint{3.816459in}{7.001278in}}%
\pgfpathcurveto{\pgfqpoint{3.816459in}{6.995455in}}{\pgfqpoint{3.818773in}{6.989868in}}{\pgfqpoint{3.822891in}{6.985750in}}%
\pgfpathcurveto{\pgfqpoint{3.827010in}{6.981632in}}{\pgfqpoint{3.832596in}{6.979318in}}{\pgfqpoint{3.838420in}{6.979318in}}%
\pgfpathlineto{\pgfqpoint{3.838420in}{6.979318in}}%
\pgfpathclose%
\pgfusepath{stroke,fill}%
\end{pgfscope}%
\begin{pgfscope}%
\pgfpathrectangle{\pgfqpoint{1.414736in}{0.548769in}}{\pgfqpoint{3.497747in}{7.253186in}}%
\pgfusepath{clip}%
\pgfsetbuttcap%
\pgfsetroundjoin%
\definecolor{currentfill}{rgb}{0.783315,0.879285,0.125405}%
\pgfsetfillcolor{currentfill}%
\pgfsetlinewidth{1.003750pt}%
\definecolor{currentstroke}{rgb}{0.783315,0.879285,0.125405}%
\pgfsetstrokecolor{currentstroke}%
\pgfsetdash{}{0pt}%
\pgfpathmoveto{\pgfqpoint{3.578395in}{7.450304in}}%
\pgfpathcurveto{\pgfqpoint{3.584219in}{7.450304in}}{\pgfqpoint{3.589805in}{7.452618in}}{\pgfqpoint{3.593923in}{7.456736in}}%
\pgfpathcurveto{\pgfqpoint{3.598041in}{7.460854in}}{\pgfqpoint{3.600355in}{7.466441in}}{\pgfqpoint{3.600355in}{7.472265in}}%
\pgfpathcurveto{\pgfqpoint{3.600355in}{7.478089in}}{\pgfqpoint{3.598041in}{7.483675in}}{\pgfqpoint{3.593923in}{7.487793in}}%
\pgfpathcurveto{\pgfqpoint{3.589805in}{7.491911in}}{\pgfqpoint{3.584219in}{7.494225in}}{\pgfqpoint{3.578395in}{7.494225in}}%
\pgfpathcurveto{\pgfqpoint{3.572571in}{7.494225in}}{\pgfqpoint{3.566985in}{7.491911in}}{\pgfqpoint{3.562867in}{7.487793in}}%
\pgfpathcurveto{\pgfqpoint{3.558749in}{7.483675in}}{\pgfqpoint{3.556435in}{7.478089in}}{\pgfqpoint{3.556435in}{7.472265in}}%
\pgfpathcurveto{\pgfqpoint{3.556435in}{7.466441in}}{\pgfqpoint{3.558749in}{7.460854in}}{\pgfqpoint{3.562867in}{7.456736in}}%
\pgfpathcurveto{\pgfqpoint{3.566985in}{7.452618in}}{\pgfqpoint{3.572571in}{7.450304in}}{\pgfqpoint{3.578395in}{7.450304in}}%
\pgfpathlineto{\pgfqpoint{3.578395in}{7.450304in}}%
\pgfpathclose%
\pgfusepath{stroke,fill}%
\end{pgfscope}%
\begin{pgfscope}%
\pgfpathrectangle{\pgfqpoint{1.414736in}{0.548769in}}{\pgfqpoint{3.497747in}{7.253186in}}%
\pgfusepath{clip}%
\pgfsetbuttcap%
\pgfsetroundjoin%
\definecolor{currentfill}{rgb}{0.955300,0.901065,0.118128}%
\pgfsetfillcolor{currentfill}%
\pgfsetlinewidth{1.003750pt}%
\definecolor{currentstroke}{rgb}{0.955300,0.901065,0.118128}%
\pgfsetstrokecolor{currentstroke}%
\pgfsetdash{}{0pt}%
\pgfpathmoveto{\pgfqpoint{3.734912in}{0.856499in}}%
\pgfpathcurveto{\pgfqpoint{3.740736in}{0.856499in}}{\pgfqpoint{3.746322in}{0.858813in}}{\pgfqpoint{3.750440in}{0.862931in}}%
\pgfpathcurveto{\pgfqpoint{3.754558in}{0.867049in}}{\pgfqpoint{3.756872in}{0.872635in}}{\pgfqpoint{3.756872in}{0.878459in}}%
\pgfpathcurveto{\pgfqpoint{3.756872in}{0.884283in}}{\pgfqpoint{3.754558in}{0.889869in}}{\pgfqpoint{3.750440in}{0.893987in}}%
\pgfpathcurveto{\pgfqpoint{3.746322in}{0.898105in}}{\pgfqpoint{3.740736in}{0.900419in}}{\pgfqpoint{3.734912in}{0.900419in}}%
\pgfpathcurveto{\pgfqpoint{3.729088in}{0.900419in}}{\pgfqpoint{3.723502in}{0.898105in}}{\pgfqpoint{3.719384in}{0.893987in}}%
\pgfpathcurveto{\pgfqpoint{3.715265in}{0.889869in}}{\pgfqpoint{3.712952in}{0.884283in}}{\pgfqpoint{3.712952in}{0.878459in}}%
\pgfpathcurveto{\pgfqpoint{3.712952in}{0.872635in}}{\pgfqpoint{3.715265in}{0.867049in}}{\pgfqpoint{3.719384in}{0.862931in}}%
\pgfpathcurveto{\pgfqpoint{3.723502in}{0.858813in}}{\pgfqpoint{3.729088in}{0.856499in}}{\pgfqpoint{3.734912in}{0.856499in}}%
\pgfpathlineto{\pgfqpoint{3.734912in}{0.856499in}}%
\pgfpathclose%
\pgfusepath{stroke,fill}%
\end{pgfscope}%
\begin{pgfscope}%
\pgfpathrectangle{\pgfqpoint{1.414736in}{0.548769in}}{\pgfqpoint{3.497747in}{7.253186in}}%
\pgfusepath{clip}%
\pgfsetbuttcap%
\pgfsetroundjoin%
\definecolor{currentfill}{rgb}{0.741388,0.873449,0.149561}%
\pgfsetfillcolor{currentfill}%
\pgfsetlinewidth{1.003750pt}%
\definecolor{currentstroke}{rgb}{0.741388,0.873449,0.149561}%
\pgfsetstrokecolor{currentstroke}%
\pgfsetdash{}{0pt}%
\pgfpathmoveto{\pgfqpoint{3.436657in}{1.327485in}}%
\pgfpathcurveto{\pgfqpoint{3.442481in}{1.327485in}}{\pgfqpoint{3.448068in}{1.329799in}}{\pgfqpoint{3.452186in}{1.333917in}}%
\pgfpathcurveto{\pgfqpoint{3.456304in}{1.338035in}}{\pgfqpoint{3.458618in}{1.343621in}}{\pgfqpoint{3.458618in}{1.349445in}}%
\pgfpathcurveto{\pgfqpoint{3.458618in}{1.355269in}}{\pgfqpoint{3.456304in}{1.360855in}}{\pgfqpoint{3.452186in}{1.364973in}}%
\pgfpathcurveto{\pgfqpoint{3.448068in}{1.369092in}}{\pgfqpoint{3.442481in}{1.371405in}}{\pgfqpoint{3.436657in}{1.371405in}}%
\pgfpathcurveto{\pgfqpoint{3.430834in}{1.371405in}}{\pgfqpoint{3.425247in}{1.369092in}}{\pgfqpoint{3.421129in}{1.364973in}}%
\pgfpathcurveto{\pgfqpoint{3.417011in}{1.360855in}}{\pgfqpoint{3.414697in}{1.355269in}}{\pgfqpoint{3.414697in}{1.349445in}}%
\pgfpathcurveto{\pgfqpoint{3.414697in}{1.343621in}}{\pgfqpoint{3.417011in}{1.338035in}}{\pgfqpoint{3.421129in}{1.333917in}}%
\pgfpathcurveto{\pgfqpoint{3.425247in}{1.329799in}}{\pgfqpoint{3.430834in}{1.327485in}}{\pgfqpoint{3.436657in}{1.327485in}}%
\pgfpathlineto{\pgfqpoint{3.436657in}{1.327485in}}%
\pgfpathclose%
\pgfusepath{stroke,fill}%
\end{pgfscope}%
\begin{pgfscope}%
\pgfpathrectangle{\pgfqpoint{1.414736in}{0.548769in}}{\pgfqpoint{3.497747in}{7.253186in}}%
\pgfusepath{clip}%
\pgfsetbuttcap%
\pgfsetroundjoin%
\definecolor{currentfill}{rgb}{0.699415,0.867117,0.175971}%
\pgfsetfillcolor{currentfill}%
\pgfsetlinewidth{1.003750pt}%
\definecolor{currentstroke}{rgb}{0.699415,0.867117,0.175971}%
\pgfsetstrokecolor{currentstroke}%
\pgfsetdash{}{0pt}%
\pgfpathmoveto{\pgfqpoint{3.715169in}{1.798471in}}%
\pgfpathcurveto{\pgfqpoint{3.720993in}{1.798471in}}{\pgfqpoint{3.726579in}{1.800785in}}{\pgfqpoint{3.730697in}{1.804903in}}%
\pgfpathcurveto{\pgfqpoint{3.734815in}{1.809021in}}{\pgfqpoint{3.737129in}{1.814607in}}{\pgfqpoint{3.737129in}{1.820431in}}%
\pgfpathcurveto{\pgfqpoint{3.737129in}{1.826255in}}{\pgfqpoint{3.734815in}{1.831841in}}{\pgfqpoint{3.730697in}{1.835960in}}%
\pgfpathcurveto{\pgfqpoint{3.726579in}{1.840078in}}{\pgfqpoint{3.720993in}{1.842392in}}{\pgfqpoint{3.715169in}{1.842392in}}%
\pgfpathcurveto{\pgfqpoint{3.709345in}{1.842392in}}{\pgfqpoint{3.703759in}{1.840078in}}{\pgfqpoint{3.699640in}{1.835960in}}%
\pgfpathcurveto{\pgfqpoint{3.695522in}{1.831841in}}{\pgfqpoint{3.693208in}{1.826255in}}{\pgfqpoint{3.693208in}{1.820431in}}%
\pgfpathcurveto{\pgfqpoint{3.693208in}{1.814607in}}{\pgfqpoint{3.695522in}{1.809021in}}{\pgfqpoint{3.699640in}{1.804903in}}%
\pgfpathcurveto{\pgfqpoint{3.703759in}{1.800785in}}{\pgfqpoint{3.709345in}{1.798471in}}{\pgfqpoint{3.715169in}{1.798471in}}%
\pgfpathlineto{\pgfqpoint{3.715169in}{1.798471in}}%
\pgfpathclose%
\pgfusepath{stroke,fill}%
\end{pgfscope}%
\begin{pgfscope}%
\pgfpathrectangle{\pgfqpoint{1.414736in}{0.548769in}}{\pgfqpoint{3.497747in}{7.253186in}}%
\pgfusepath{clip}%
\pgfsetbuttcap%
\pgfsetroundjoin%
\definecolor{currentfill}{rgb}{0.993248,0.906157,0.143936}%
\pgfsetfillcolor{currentfill}%
\pgfsetlinewidth{1.003750pt}%
\definecolor{currentstroke}{rgb}{0.993248,0.906157,0.143936}%
\pgfsetstrokecolor{currentstroke}%
\pgfsetdash{}{0pt}%
\pgfpathmoveto{\pgfqpoint{3.578509in}{2.269457in}}%
\pgfpathcurveto{\pgfqpoint{3.584333in}{2.269457in}}{\pgfqpoint{3.589919in}{2.271771in}}{\pgfqpoint{3.594037in}{2.275889in}}%
\pgfpathcurveto{\pgfqpoint{3.598156in}{2.280007in}}{\pgfqpoint{3.600469in}{2.285593in}}{\pgfqpoint{3.600469in}{2.291417in}}%
\pgfpathcurveto{\pgfqpoint{3.600469in}{2.297241in}}{\pgfqpoint{3.598156in}{2.302828in}}{\pgfqpoint{3.594037in}{2.306946in}}%
\pgfpathcurveto{\pgfqpoint{3.589919in}{2.311064in}}{\pgfqpoint{3.584333in}{2.313378in}}{\pgfqpoint{3.578509in}{2.313378in}}%
\pgfpathcurveto{\pgfqpoint{3.572685in}{2.313378in}}{\pgfqpoint{3.567099in}{2.311064in}}{\pgfqpoint{3.562981in}{2.306946in}}%
\pgfpathcurveto{\pgfqpoint{3.558863in}{2.302828in}}{\pgfqpoint{3.556549in}{2.297241in}}{\pgfqpoint{3.556549in}{2.291417in}}%
\pgfpathcurveto{\pgfqpoint{3.556549in}{2.285593in}}{\pgfqpoint{3.558863in}{2.280007in}}{\pgfqpoint{3.562981in}{2.275889in}}%
\pgfpathcurveto{\pgfqpoint{3.567099in}{2.271771in}}{\pgfqpoint{3.572685in}{2.269457in}}{\pgfqpoint{3.578509in}{2.269457in}}%
\pgfpathlineto{\pgfqpoint{3.578509in}{2.269457in}}%
\pgfpathclose%
\pgfusepath{stroke,fill}%
\end{pgfscope}%
\begin{pgfscope}%
\pgfpathrectangle{\pgfqpoint{1.414736in}{0.548769in}}{\pgfqpoint{3.497747in}{7.253186in}}%
\pgfusepath{clip}%
\pgfsetbuttcap%
\pgfsetroundjoin%
\definecolor{currentfill}{rgb}{0.288921,0.758394,0.428426}%
\pgfsetfillcolor{currentfill}%
\pgfsetlinewidth{1.003750pt}%
\definecolor{currentstroke}{rgb}{0.288921,0.758394,0.428426}%
\pgfsetstrokecolor{currentstroke}%
\pgfsetdash{}{0pt}%
\pgfpathmoveto{\pgfqpoint{3.686121in}{2.740443in}}%
\pgfpathcurveto{\pgfqpoint{3.691945in}{2.740443in}}{\pgfqpoint{3.697531in}{2.742757in}}{\pgfqpoint{3.701649in}{2.746875in}}%
\pgfpathcurveto{\pgfqpoint{3.705768in}{2.750993in}}{\pgfqpoint{3.708081in}{2.756580in}}{\pgfqpoint{3.708081in}{2.762404in}}%
\pgfpathcurveto{\pgfqpoint{3.708081in}{2.768227in}}{\pgfqpoint{3.705768in}{2.773814in}}{\pgfqpoint{3.701649in}{2.777932in}}%
\pgfpathcurveto{\pgfqpoint{3.697531in}{2.782050in}}{\pgfqpoint{3.691945in}{2.784364in}}{\pgfqpoint{3.686121in}{2.784364in}}%
\pgfpathcurveto{\pgfqpoint{3.680297in}{2.784364in}}{\pgfqpoint{3.674711in}{2.782050in}}{\pgfqpoint{3.670593in}{2.777932in}}%
\pgfpathcurveto{\pgfqpoint{3.666475in}{2.773814in}}{\pgfqpoint{3.664161in}{2.768227in}}{\pgfqpoint{3.664161in}{2.762404in}}%
\pgfpathcurveto{\pgfqpoint{3.664161in}{2.756580in}}{\pgfqpoint{3.666475in}{2.750993in}}{\pgfqpoint{3.670593in}{2.746875in}}%
\pgfpathcurveto{\pgfqpoint{3.674711in}{2.742757in}}{\pgfqpoint{3.680297in}{2.740443in}}{\pgfqpoint{3.686121in}{2.740443in}}%
\pgfpathlineto{\pgfqpoint{3.686121in}{2.740443in}}%
\pgfpathclose%
\pgfusepath{stroke,fill}%
\end{pgfscope}%
\begin{pgfscope}%
\pgfpathrectangle{\pgfqpoint{1.414736in}{0.548769in}}{\pgfqpoint{3.497747in}{7.253186in}}%
\pgfusepath{clip}%
\pgfsetbuttcap%
\pgfsetroundjoin%
\definecolor{currentfill}{rgb}{0.783315,0.879285,0.125405}%
\pgfsetfillcolor{currentfill}%
\pgfsetlinewidth{1.003750pt}%
\definecolor{currentstroke}{rgb}{0.783315,0.879285,0.125405}%
\pgfsetstrokecolor{currentstroke}%
\pgfsetdash{}{0pt}%
\pgfpathmoveto{\pgfqpoint{3.444702in}{3.211429in}}%
\pgfpathcurveto{\pgfqpoint{3.450526in}{3.211429in}}{\pgfqpoint{3.456112in}{3.213743in}}{\pgfqpoint{3.460230in}{3.217861in}}%
\pgfpathcurveto{\pgfqpoint{3.464349in}{3.221980in}}{\pgfqpoint{3.466662in}{3.227566in}}{\pgfqpoint{3.466662in}{3.233390in}}%
\pgfpathcurveto{\pgfqpoint{3.466662in}{3.239214in}}{\pgfqpoint{3.464349in}{3.244800in}}{\pgfqpoint{3.460230in}{3.248918in}}%
\pgfpathcurveto{\pgfqpoint{3.456112in}{3.253036in}}{\pgfqpoint{3.450526in}{3.255350in}}{\pgfqpoint{3.444702in}{3.255350in}}%
\pgfpathcurveto{\pgfqpoint{3.438878in}{3.255350in}}{\pgfqpoint{3.433292in}{3.253036in}}{\pgfqpoint{3.429174in}{3.248918in}}%
\pgfpathcurveto{\pgfqpoint{3.425056in}{3.244800in}}{\pgfqpoint{3.422742in}{3.239214in}}{\pgfqpoint{3.422742in}{3.233390in}}%
\pgfpathcurveto{\pgfqpoint{3.422742in}{3.227566in}}{\pgfqpoint{3.425056in}{3.221980in}}{\pgfqpoint{3.429174in}{3.217861in}}%
\pgfpathcurveto{\pgfqpoint{3.433292in}{3.213743in}}{\pgfqpoint{3.438878in}{3.211429in}}{\pgfqpoint{3.444702in}{3.211429in}}%
\pgfpathlineto{\pgfqpoint{3.444702in}{3.211429in}}%
\pgfpathclose%
\pgfusepath{stroke,fill}%
\end{pgfscope}%
\begin{pgfscope}%
\pgfpathrectangle{\pgfqpoint{1.414736in}{0.548769in}}{\pgfqpoint{3.497747in}{7.253186in}}%
\pgfusepath{clip}%
\pgfsetbuttcap%
\pgfsetroundjoin%
\definecolor{currentfill}{rgb}{0.993248,0.906157,0.143936}%
\pgfsetfillcolor{currentfill}%
\pgfsetlinewidth{1.003750pt}%
\definecolor{currentstroke}{rgb}{0.993248,0.906157,0.143936}%
\pgfsetstrokecolor{currentstroke}%
\pgfsetdash{}{0pt}%
\pgfpathmoveto{\pgfqpoint{3.364838in}{3.682415in}}%
\pgfpathcurveto{\pgfqpoint{3.370662in}{3.682415in}}{\pgfqpoint{3.376248in}{3.684729in}}{\pgfqpoint{3.380367in}{3.688847in}}%
\pgfpathcurveto{\pgfqpoint{3.384485in}{3.692966in}}{\pgfqpoint{3.386799in}{3.698552in}}{\pgfqpoint{3.386799in}{3.704376in}}%
\pgfpathcurveto{\pgfqpoint{3.386799in}{3.710200in}}{\pgfqpoint{3.384485in}{3.715786in}}{\pgfqpoint{3.380367in}{3.719904in}}%
\pgfpathcurveto{\pgfqpoint{3.376248in}{3.724022in}}{\pgfqpoint{3.370662in}{3.726336in}}{\pgfqpoint{3.364838in}{3.726336in}}%
\pgfpathcurveto{\pgfqpoint{3.359014in}{3.726336in}}{\pgfqpoint{3.353428in}{3.724022in}}{\pgfqpoint{3.349310in}{3.719904in}}%
\pgfpathcurveto{\pgfqpoint{3.345192in}{3.715786in}}{\pgfqpoint{3.342878in}{3.710200in}}{\pgfqpoint{3.342878in}{3.704376in}}%
\pgfpathcurveto{\pgfqpoint{3.342878in}{3.698552in}}{\pgfqpoint{3.345192in}{3.692966in}}{\pgfqpoint{3.349310in}{3.688847in}}%
\pgfpathcurveto{\pgfqpoint{3.353428in}{3.684729in}}{\pgfqpoint{3.359014in}{3.682415in}}{\pgfqpoint{3.364838in}{3.682415in}}%
\pgfpathlineto{\pgfqpoint{3.364838in}{3.682415in}}%
\pgfpathclose%
\pgfusepath{stroke,fill}%
\end{pgfscope}%
\begin{pgfscope}%
\pgfpathrectangle{\pgfqpoint{1.414736in}{0.548769in}}{\pgfqpoint{3.497747in}{7.253186in}}%
\pgfusepath{clip}%
\pgfsetbuttcap%
\pgfsetroundjoin%
\definecolor{currentfill}{rgb}{0.935904,0.898570,0.108131}%
\pgfsetfillcolor{currentfill}%
\pgfsetlinewidth{1.003750pt}%
\definecolor{currentstroke}{rgb}{0.935904,0.898570,0.108131}%
\pgfsetstrokecolor{currentstroke}%
\pgfsetdash{}{0pt}%
\pgfpathmoveto{\pgfqpoint{3.794610in}{4.153402in}}%
\pgfpathcurveto{\pgfqpoint{3.800434in}{4.153402in}}{\pgfqpoint{3.806020in}{4.155715in}}{\pgfqpoint{3.810138in}{4.159834in}}%
\pgfpathcurveto{\pgfqpoint{3.814256in}{4.163952in}}{\pgfqpoint{3.816570in}{4.169538in}}{\pgfqpoint{3.816570in}{4.175362in}}%
\pgfpathcurveto{\pgfqpoint{3.816570in}{4.181186in}}{\pgfqpoint{3.814256in}{4.186772in}}{\pgfqpoint{3.810138in}{4.190890in}}%
\pgfpathcurveto{\pgfqpoint{3.806020in}{4.195008in}}{\pgfqpoint{3.800434in}{4.197322in}}{\pgfqpoint{3.794610in}{4.197322in}}%
\pgfpathcurveto{\pgfqpoint{3.788786in}{4.197322in}}{\pgfqpoint{3.783200in}{4.195008in}}{\pgfqpoint{3.779082in}{4.190890in}}%
\pgfpathcurveto{\pgfqpoint{3.774964in}{4.186772in}}{\pgfqpoint{3.772650in}{4.181186in}}{\pgfqpoint{3.772650in}{4.175362in}}%
\pgfpathcurveto{\pgfqpoint{3.772650in}{4.169538in}}{\pgfqpoint{3.774964in}{4.163952in}}{\pgfqpoint{3.779082in}{4.159834in}}%
\pgfpathcurveto{\pgfqpoint{3.783200in}{4.155715in}}{\pgfqpoint{3.788786in}{4.153402in}}{\pgfqpoint{3.794610in}{4.153402in}}%
\pgfpathlineto{\pgfqpoint{3.794610in}{4.153402in}}%
\pgfpathclose%
\pgfusepath{stroke,fill}%
\end{pgfscope}%
\begin{pgfscope}%
\pgfpathrectangle{\pgfqpoint{1.414736in}{0.548769in}}{\pgfqpoint{3.497747in}{7.253186in}}%
\pgfusepath{clip}%
\pgfsetbuttcap%
\pgfsetroundjoin%
\definecolor{currentfill}{rgb}{0.741388,0.873449,0.149561}%
\pgfsetfillcolor{currentfill}%
\pgfsetlinewidth{1.003750pt}%
\definecolor{currentstroke}{rgb}{0.741388,0.873449,0.149561}%
\pgfsetstrokecolor{currentstroke}%
\pgfsetdash{}{0pt}%
\pgfpathmoveto{\pgfqpoint{3.614774in}{4.624388in}}%
\pgfpathcurveto{\pgfqpoint{3.620598in}{4.624388in}}{\pgfqpoint{3.626184in}{4.626702in}}{\pgfqpoint{3.630302in}{4.630820in}}%
\pgfpathcurveto{\pgfqpoint{3.634421in}{4.634938in}}{\pgfqpoint{3.636734in}{4.640524in}}{\pgfqpoint{3.636734in}{4.646348in}}%
\pgfpathcurveto{\pgfqpoint{3.636734in}{4.652172in}}{\pgfqpoint{3.634421in}{4.657758in}}{\pgfqpoint{3.630302in}{4.661876in}}%
\pgfpathcurveto{\pgfqpoint{3.626184in}{4.665994in}}{\pgfqpoint{3.620598in}{4.668308in}}{\pgfqpoint{3.614774in}{4.668308in}}%
\pgfpathcurveto{\pgfqpoint{3.608950in}{4.668308in}}{\pgfqpoint{3.603364in}{4.665994in}}{\pgfqpoint{3.599246in}{4.661876in}}%
\pgfpathcurveto{\pgfqpoint{3.595128in}{4.657758in}}{\pgfqpoint{3.592814in}{4.652172in}}{\pgfqpoint{3.592814in}{4.646348in}}%
\pgfpathcurveto{\pgfqpoint{3.592814in}{4.640524in}}{\pgfqpoint{3.595128in}{4.634938in}}{\pgfqpoint{3.599246in}{4.630820in}}%
\pgfpathcurveto{\pgfqpoint{3.603364in}{4.626702in}}{\pgfqpoint{3.608950in}{4.624388in}}{\pgfqpoint{3.614774in}{4.624388in}}%
\pgfpathlineto{\pgfqpoint{3.614774in}{4.624388in}}%
\pgfpathclose%
\pgfusepath{stroke,fill}%
\end{pgfscope}%
\begin{pgfscope}%
\pgfpathrectangle{\pgfqpoint{1.414736in}{0.548769in}}{\pgfqpoint{3.497747in}{7.253186in}}%
\pgfusepath{clip}%
\pgfsetbuttcap%
\pgfsetroundjoin%
\definecolor{currentfill}{rgb}{0.157729,0.485932,0.558013}%
\pgfsetfillcolor{currentfill}%
\pgfsetlinewidth{1.003750pt}%
\definecolor{currentstroke}{rgb}{0.157729,0.485932,0.558013}%
\pgfsetstrokecolor{currentstroke}%
\pgfsetdash{}{0pt}%
\pgfpathmoveto{\pgfqpoint{3.500616in}{5.095374in}}%
\pgfpathcurveto{\pgfqpoint{3.506440in}{5.095374in}}{\pgfqpoint{3.512026in}{5.097688in}}{\pgfqpoint{3.516144in}{5.101806in}}%
\pgfpathcurveto{\pgfqpoint{3.520262in}{5.105924in}}{\pgfqpoint{3.522576in}{5.111510in}}{\pgfqpoint{3.522576in}{5.117334in}}%
\pgfpathcurveto{\pgfqpoint{3.522576in}{5.123158in}}{\pgfqpoint{3.520262in}{5.128744in}}{\pgfqpoint{3.516144in}{5.132862in}}%
\pgfpathcurveto{\pgfqpoint{3.512026in}{5.136980in}}{\pgfqpoint{3.506440in}{5.139294in}}{\pgfqpoint{3.500616in}{5.139294in}}%
\pgfpathcurveto{\pgfqpoint{3.494792in}{5.139294in}}{\pgfqpoint{3.489206in}{5.136980in}}{\pgfqpoint{3.485088in}{5.132862in}}%
\pgfpathcurveto{\pgfqpoint{3.480970in}{5.128744in}}{\pgfqpoint{3.478656in}{5.123158in}}{\pgfqpoint{3.478656in}{5.117334in}}%
\pgfpathcurveto{\pgfqpoint{3.478656in}{5.111510in}}{\pgfqpoint{3.480970in}{5.105924in}}{\pgfqpoint{3.485088in}{5.101806in}}%
\pgfpathcurveto{\pgfqpoint{3.489206in}{5.097688in}}{\pgfqpoint{3.494792in}{5.095374in}}{\pgfqpoint{3.500616in}{5.095374in}}%
\pgfpathlineto{\pgfqpoint{3.500616in}{5.095374in}}%
\pgfpathclose%
\pgfusepath{stroke,fill}%
\end{pgfscope}%
\begin{pgfscope}%
\pgfpathrectangle{\pgfqpoint{1.414736in}{0.548769in}}{\pgfqpoint{3.497747in}{7.253186in}}%
\pgfusepath{clip}%
\pgfsetbuttcap%
\pgfsetroundjoin%
\definecolor{currentfill}{rgb}{0.204903,0.375746,0.553533}%
\pgfsetfillcolor{currentfill}%
\pgfsetlinewidth{1.003750pt}%
\definecolor{currentstroke}{rgb}{0.204903,0.375746,0.553533}%
\pgfsetstrokecolor{currentstroke}%
\pgfsetdash{}{0pt}%
\pgfpathmoveto{\pgfqpoint{3.971359in}{5.566360in}}%
\pgfpathcurveto{\pgfqpoint{3.977183in}{5.566360in}}{\pgfqpoint{3.982769in}{5.568674in}}{\pgfqpoint{3.986887in}{5.572792in}}%
\pgfpathcurveto{\pgfqpoint{3.991005in}{5.576910in}}{\pgfqpoint{3.993319in}{5.582496in}}{\pgfqpoint{3.993319in}{5.588320in}}%
\pgfpathcurveto{\pgfqpoint{3.993319in}{5.594144in}}{\pgfqpoint{3.991005in}{5.599730in}}{\pgfqpoint{3.986887in}{5.603848in}}%
\pgfpathcurveto{\pgfqpoint{3.982769in}{5.607967in}}{\pgfqpoint{3.977183in}{5.610280in}}{\pgfqpoint{3.971359in}{5.610280in}}%
\pgfpathcurveto{\pgfqpoint{3.965535in}{5.610280in}}{\pgfqpoint{3.959949in}{5.607967in}}{\pgfqpoint{3.955831in}{5.603848in}}%
\pgfpathcurveto{\pgfqpoint{3.951713in}{5.599730in}}{\pgfqpoint{3.949399in}{5.594144in}}{\pgfqpoint{3.949399in}{5.588320in}}%
\pgfpathcurveto{\pgfqpoint{3.949399in}{5.582496in}}{\pgfqpoint{3.951713in}{5.576910in}}{\pgfqpoint{3.955831in}{5.572792in}}%
\pgfpathcurveto{\pgfqpoint{3.959949in}{5.568674in}}{\pgfqpoint{3.965535in}{5.566360in}}{\pgfqpoint{3.971359in}{5.566360in}}%
\pgfpathlineto{\pgfqpoint{3.971359in}{5.566360in}}%
\pgfpathclose%
\pgfusepath{stroke,fill}%
\end{pgfscope}%
\begin{pgfscope}%
\pgfpathrectangle{\pgfqpoint{1.414736in}{0.548769in}}{\pgfqpoint{3.497747in}{7.253186in}}%
\pgfusepath{clip}%
\pgfsetbuttcap%
\pgfsetroundjoin%
\definecolor{currentfill}{rgb}{0.274149,0.751988,0.436601}%
\pgfsetfillcolor{currentfill}%
\pgfsetlinewidth{1.003750pt}%
\definecolor{currentstroke}{rgb}{0.274149,0.751988,0.436601}%
\pgfsetstrokecolor{currentstroke}%
\pgfsetdash{}{0pt}%
\pgfpathmoveto{\pgfqpoint{3.378071in}{6.037346in}}%
\pgfpathcurveto{\pgfqpoint{3.383895in}{6.037346in}}{\pgfqpoint{3.389481in}{6.039660in}}{\pgfqpoint{3.393599in}{6.043778in}}%
\pgfpathcurveto{\pgfqpoint{3.397717in}{6.047896in}}{\pgfqpoint{3.400031in}{6.053482in}}{\pgfqpoint{3.400031in}{6.059306in}}%
\pgfpathcurveto{\pgfqpoint{3.400031in}{6.065130in}}{\pgfqpoint{3.397717in}{6.070716in}}{\pgfqpoint{3.393599in}{6.074835in}}%
\pgfpathcurveto{\pgfqpoint{3.389481in}{6.078953in}}{\pgfqpoint{3.383895in}{6.081267in}}{\pgfqpoint{3.378071in}{6.081267in}}%
\pgfpathcurveto{\pgfqpoint{3.372247in}{6.081267in}}{\pgfqpoint{3.366661in}{6.078953in}}{\pgfqpoint{3.362542in}{6.074835in}}%
\pgfpathcurveto{\pgfqpoint{3.358424in}{6.070716in}}{\pgfqpoint{3.356110in}{6.065130in}}{\pgfqpoint{3.356110in}{6.059306in}}%
\pgfpathcurveto{\pgfqpoint{3.356110in}{6.053482in}}{\pgfqpoint{3.358424in}{6.047896in}}{\pgfqpoint{3.362542in}{6.043778in}}%
\pgfpathcurveto{\pgfqpoint{3.366661in}{6.039660in}}{\pgfqpoint{3.372247in}{6.037346in}}{\pgfqpoint{3.378071in}{6.037346in}}%
\pgfpathlineto{\pgfqpoint{3.378071in}{6.037346in}}%
\pgfpathclose%
\pgfusepath{stroke,fill}%
\end{pgfscope}%
\begin{pgfscope}%
\pgfpathrectangle{\pgfqpoint{1.414736in}{0.548769in}}{\pgfqpoint{3.497747in}{7.253186in}}%
\pgfusepath{clip}%
\pgfsetbuttcap%
\pgfsetroundjoin%
\definecolor{currentfill}{rgb}{0.993248,0.906157,0.143936}%
\pgfsetfillcolor{currentfill}%
\pgfsetlinewidth{1.003750pt}%
\definecolor{currentstroke}{rgb}{0.993248,0.906157,0.143936}%
\pgfsetstrokecolor{currentstroke}%
\pgfsetdash{}{0pt}%
\pgfpathmoveto{\pgfqpoint{3.710205in}{6.508332in}}%
\pgfpathcurveto{\pgfqpoint{3.716029in}{6.508332in}}{\pgfqpoint{3.721615in}{6.510646in}}{\pgfqpoint{3.725733in}{6.514764in}}%
\pgfpathcurveto{\pgfqpoint{3.729851in}{6.518882in}}{\pgfqpoint{3.732165in}{6.524468in}}{\pgfqpoint{3.732165in}{6.530292in}}%
\pgfpathcurveto{\pgfqpoint{3.732165in}{6.536116in}}{\pgfqpoint{3.729851in}{6.541702in}}{\pgfqpoint{3.725733in}{6.545821in}}%
\pgfpathcurveto{\pgfqpoint{3.721615in}{6.549939in}}{\pgfqpoint{3.716029in}{6.552253in}}{\pgfqpoint{3.710205in}{6.552253in}}%
\pgfpathcurveto{\pgfqpoint{3.704381in}{6.552253in}}{\pgfqpoint{3.698795in}{6.549939in}}{\pgfqpoint{3.694677in}{6.545821in}}%
\pgfpathcurveto{\pgfqpoint{3.690558in}{6.541702in}}{\pgfqpoint{3.688245in}{6.536116in}}{\pgfqpoint{3.688245in}{6.530292in}}%
\pgfpathcurveto{\pgfqpoint{3.688245in}{6.524468in}}{\pgfqpoint{3.690558in}{6.518882in}}{\pgfqpoint{3.694677in}{6.514764in}}%
\pgfpathcurveto{\pgfqpoint{3.698795in}{6.510646in}}{\pgfqpoint{3.704381in}{6.508332in}}{\pgfqpoint{3.710205in}{6.508332in}}%
\pgfpathlineto{\pgfqpoint{3.710205in}{6.508332in}}%
\pgfpathclose%
\pgfusepath{stroke,fill}%
\end{pgfscope}%
\begin{pgfscope}%
\pgfpathrectangle{\pgfqpoint{1.414736in}{0.548769in}}{\pgfqpoint{3.497747in}{7.253186in}}%
\pgfusepath{clip}%
\pgfsetbuttcap%
\pgfsetroundjoin%
\definecolor{currentfill}{rgb}{0.993248,0.906157,0.143936}%
\pgfsetfillcolor{currentfill}%
\pgfsetlinewidth{1.003750pt}%
\definecolor{currentstroke}{rgb}{0.993248,0.906157,0.143936}%
\pgfsetstrokecolor{currentstroke}%
\pgfsetdash{}{0pt}%
\pgfpathmoveto{\pgfqpoint{3.641267in}{6.979318in}}%
\pgfpathcurveto{\pgfqpoint{3.647091in}{6.979318in}}{\pgfqpoint{3.652677in}{6.981632in}}{\pgfqpoint{3.656795in}{6.985750in}}%
\pgfpathcurveto{\pgfqpoint{3.660913in}{6.989868in}}{\pgfqpoint{3.663227in}{6.995455in}}{\pgfqpoint{3.663227in}{7.001278in}}%
\pgfpathcurveto{\pgfqpoint{3.663227in}{7.007102in}}{\pgfqpoint{3.660913in}{7.012689in}}{\pgfqpoint{3.656795in}{7.016807in}}%
\pgfpathcurveto{\pgfqpoint{3.652677in}{7.020925in}}{\pgfqpoint{3.647091in}{7.023239in}}{\pgfqpoint{3.641267in}{7.023239in}}%
\pgfpathcurveto{\pgfqpoint{3.635443in}{7.023239in}}{\pgfqpoint{3.629857in}{7.020925in}}{\pgfqpoint{3.625738in}{7.016807in}}%
\pgfpathcurveto{\pgfqpoint{3.621620in}{7.012689in}}{\pgfqpoint{3.619306in}{7.007102in}}{\pgfqpoint{3.619306in}{7.001278in}}%
\pgfpathcurveto{\pgfqpoint{3.619306in}{6.995455in}}{\pgfqpoint{3.621620in}{6.989868in}}{\pgfqpoint{3.625738in}{6.985750in}}%
\pgfpathcurveto{\pgfqpoint{3.629857in}{6.981632in}}{\pgfqpoint{3.635443in}{6.979318in}}{\pgfqpoint{3.641267in}{6.979318in}}%
\pgfpathlineto{\pgfqpoint{3.641267in}{6.979318in}}%
\pgfpathclose%
\pgfusepath{stroke,fill}%
\end{pgfscope}%
\begin{pgfscope}%
\pgfpathrectangle{\pgfqpoint{1.414736in}{0.548769in}}{\pgfqpoint{3.497747in}{7.253186in}}%
\pgfusepath{clip}%
\pgfsetbuttcap%
\pgfsetroundjoin%
\definecolor{currentfill}{rgb}{0.412913,0.803041,0.357269}%
\pgfsetfillcolor{currentfill}%
\pgfsetlinewidth{1.003750pt}%
\definecolor{currentstroke}{rgb}{0.412913,0.803041,0.357269}%
\pgfsetstrokecolor{currentstroke}%
\pgfsetdash{}{0pt}%
\pgfpathmoveto{\pgfqpoint{4.492190in}{7.450304in}}%
\pgfpathcurveto{\pgfqpoint{4.498014in}{7.450304in}}{\pgfqpoint{4.503600in}{7.452618in}}{\pgfqpoint{4.507718in}{7.456736in}}%
\pgfpathcurveto{\pgfqpoint{4.511836in}{7.460854in}}{\pgfqpoint{4.514150in}{7.466441in}}{\pgfqpoint{4.514150in}{7.472265in}}%
\pgfpathcurveto{\pgfqpoint{4.514150in}{7.478089in}}{\pgfqpoint{4.511836in}{7.483675in}}{\pgfqpoint{4.507718in}{7.487793in}}%
\pgfpathcurveto{\pgfqpoint{4.503600in}{7.491911in}}{\pgfqpoint{4.498014in}{7.494225in}}{\pgfqpoint{4.492190in}{7.494225in}}%
\pgfpathcurveto{\pgfqpoint{4.486366in}{7.494225in}}{\pgfqpoint{4.480779in}{7.491911in}}{\pgfqpoint{4.476661in}{7.487793in}}%
\pgfpathcurveto{\pgfqpoint{4.472543in}{7.483675in}}{\pgfqpoint{4.470229in}{7.478089in}}{\pgfqpoint{4.470229in}{7.472265in}}%
\pgfpathcurveto{\pgfqpoint{4.470229in}{7.466441in}}{\pgfqpoint{4.472543in}{7.460854in}}{\pgfqpoint{4.476661in}{7.456736in}}%
\pgfpathcurveto{\pgfqpoint{4.480779in}{7.452618in}}{\pgfqpoint{4.486366in}{7.450304in}}{\pgfqpoint{4.492190in}{7.450304in}}%
\pgfpathlineto{\pgfqpoint{4.492190in}{7.450304in}}%
\pgfpathclose%
\pgfusepath{stroke,fill}%
\end{pgfscope}%
\begin{pgfscope}%
\pgfpathrectangle{\pgfqpoint{1.414736in}{0.548769in}}{\pgfqpoint{3.497747in}{7.253186in}}%
\pgfusepath{clip}%
\pgfsetbuttcap%
\pgfsetroundjoin%
\definecolor{currentfill}{rgb}{0.855810,0.888601,0.097452}%
\pgfsetfillcolor{currentfill}%
\pgfsetlinewidth{1.003750pt}%
\definecolor{currentstroke}{rgb}{0.855810,0.888601,0.097452}%
\pgfsetstrokecolor{currentstroke}%
\pgfsetdash{}{0pt}%
\pgfpathmoveto{\pgfqpoint{3.404719in}{0.856499in}}%
\pgfpathcurveto{\pgfqpoint{3.410543in}{0.856499in}}{\pgfqpoint{3.416129in}{0.858813in}}{\pgfqpoint{3.420247in}{0.862931in}}%
\pgfpathcurveto{\pgfqpoint{3.424366in}{0.867049in}}{\pgfqpoint{3.426679in}{0.872635in}}{\pgfqpoint{3.426679in}{0.878459in}}%
\pgfpathcurveto{\pgfqpoint{3.426679in}{0.884283in}}{\pgfqpoint{3.424366in}{0.889869in}}{\pgfqpoint{3.420247in}{0.893987in}}%
\pgfpathcurveto{\pgfqpoint{3.416129in}{0.898105in}}{\pgfqpoint{3.410543in}{0.900419in}}{\pgfqpoint{3.404719in}{0.900419in}}%
\pgfpathcurveto{\pgfqpoint{3.398895in}{0.900419in}}{\pgfqpoint{3.393309in}{0.898105in}}{\pgfqpoint{3.389191in}{0.893987in}}%
\pgfpathcurveto{\pgfqpoint{3.385073in}{0.889869in}}{\pgfqpoint{3.382759in}{0.884283in}}{\pgfqpoint{3.382759in}{0.878459in}}%
\pgfpathcurveto{\pgfqpoint{3.382759in}{0.872635in}}{\pgfqpoint{3.385073in}{0.867049in}}{\pgfqpoint{3.389191in}{0.862931in}}%
\pgfpathcurveto{\pgfqpoint{3.393309in}{0.858813in}}{\pgfqpoint{3.398895in}{0.856499in}}{\pgfqpoint{3.404719in}{0.856499in}}%
\pgfpathlineto{\pgfqpoint{3.404719in}{0.856499in}}%
\pgfpathclose%
\pgfusepath{stroke,fill}%
\end{pgfscope}%
\begin{pgfscope}%
\pgfpathrectangle{\pgfqpoint{1.414736in}{0.548769in}}{\pgfqpoint{3.497747in}{7.253186in}}%
\pgfusepath{clip}%
\pgfsetbuttcap%
\pgfsetroundjoin%
\definecolor{currentfill}{rgb}{0.421908,0.805774,0.351910}%
\pgfsetfillcolor{currentfill}%
\pgfsetlinewidth{1.003750pt}%
\definecolor{currentstroke}{rgb}{0.421908,0.805774,0.351910}%
\pgfsetstrokecolor{currentstroke}%
\pgfsetdash{}{0pt}%
\pgfpathmoveto{\pgfqpoint{3.732501in}{1.327485in}}%
\pgfpathcurveto{\pgfqpoint{3.738325in}{1.327485in}}{\pgfqpoint{3.743911in}{1.329799in}}{\pgfqpoint{3.748029in}{1.333917in}}%
\pgfpathcurveto{\pgfqpoint{3.752147in}{1.338035in}}{\pgfqpoint{3.754461in}{1.343621in}}{\pgfqpoint{3.754461in}{1.349445in}}%
\pgfpathcurveto{\pgfqpoint{3.754461in}{1.355269in}}{\pgfqpoint{3.752147in}{1.360855in}}{\pgfqpoint{3.748029in}{1.364973in}}%
\pgfpathcurveto{\pgfqpoint{3.743911in}{1.369092in}}{\pgfqpoint{3.738325in}{1.371405in}}{\pgfqpoint{3.732501in}{1.371405in}}%
\pgfpathcurveto{\pgfqpoint{3.726677in}{1.371405in}}{\pgfqpoint{3.721091in}{1.369092in}}{\pgfqpoint{3.716973in}{1.364973in}}%
\pgfpathcurveto{\pgfqpoint{3.712854in}{1.360855in}}{\pgfqpoint{3.710541in}{1.355269in}}{\pgfqpoint{3.710541in}{1.349445in}}%
\pgfpathcurveto{\pgfqpoint{3.710541in}{1.343621in}}{\pgfqpoint{3.712854in}{1.338035in}}{\pgfqpoint{3.716973in}{1.333917in}}%
\pgfpathcurveto{\pgfqpoint{3.721091in}{1.329799in}}{\pgfqpoint{3.726677in}{1.327485in}}{\pgfqpoint{3.732501in}{1.327485in}}%
\pgfpathlineto{\pgfqpoint{3.732501in}{1.327485in}}%
\pgfpathclose%
\pgfusepath{stroke,fill}%
\end{pgfscope}%
\begin{pgfscope}%
\pgfpathrectangle{\pgfqpoint{1.414736in}{0.548769in}}{\pgfqpoint{3.497747in}{7.253186in}}%
\pgfusepath{clip}%
\pgfsetbuttcap%
\pgfsetroundjoin%
\definecolor{currentfill}{rgb}{0.699415,0.867117,0.175971}%
\pgfsetfillcolor{currentfill}%
\pgfsetlinewidth{1.003750pt}%
\definecolor{currentstroke}{rgb}{0.699415,0.867117,0.175971}%
\pgfsetstrokecolor{currentstroke}%
\pgfsetdash{}{0pt}%
\pgfpathmoveto{\pgfqpoint{3.799805in}{1.798471in}}%
\pgfpathcurveto{\pgfqpoint{3.805628in}{1.798471in}}{\pgfqpoint{3.811215in}{1.800785in}}{\pgfqpoint{3.815333in}{1.804903in}}%
\pgfpathcurveto{\pgfqpoint{3.819451in}{1.809021in}}{\pgfqpoint{3.821765in}{1.814607in}}{\pgfqpoint{3.821765in}{1.820431in}}%
\pgfpathcurveto{\pgfqpoint{3.821765in}{1.826255in}}{\pgfqpoint{3.819451in}{1.831841in}}{\pgfqpoint{3.815333in}{1.835960in}}%
\pgfpathcurveto{\pgfqpoint{3.811215in}{1.840078in}}{\pgfqpoint{3.805628in}{1.842392in}}{\pgfqpoint{3.799805in}{1.842392in}}%
\pgfpathcurveto{\pgfqpoint{3.793981in}{1.842392in}}{\pgfqpoint{3.788394in}{1.840078in}}{\pgfqpoint{3.784276in}{1.835960in}}%
\pgfpathcurveto{\pgfqpoint{3.780158in}{1.831841in}}{\pgfqpoint{3.777844in}{1.826255in}}{\pgfqpoint{3.777844in}{1.820431in}}%
\pgfpathcurveto{\pgfqpoint{3.777844in}{1.814607in}}{\pgfqpoint{3.780158in}{1.809021in}}{\pgfqpoint{3.784276in}{1.804903in}}%
\pgfpathcurveto{\pgfqpoint{3.788394in}{1.800785in}}{\pgfqpoint{3.793981in}{1.798471in}}{\pgfqpoint{3.799805in}{1.798471in}}%
\pgfpathlineto{\pgfqpoint{3.799805in}{1.798471in}}%
\pgfpathclose%
\pgfusepath{stroke,fill}%
\end{pgfscope}%
\begin{pgfscope}%
\pgfpathrectangle{\pgfqpoint{1.414736in}{0.548769in}}{\pgfqpoint{3.497747in}{7.253186in}}%
\pgfusepath{clip}%
\pgfsetbuttcap%
\pgfsetroundjoin%
\definecolor{currentfill}{rgb}{0.993248,0.906157,0.143936}%
\pgfsetfillcolor{currentfill}%
\pgfsetlinewidth{1.003750pt}%
\definecolor{currentstroke}{rgb}{0.993248,0.906157,0.143936}%
\pgfsetstrokecolor{currentstroke}%
\pgfsetdash{}{0pt}%
\pgfpathmoveto{\pgfqpoint{3.857278in}{2.269457in}}%
\pgfpathcurveto{\pgfqpoint{3.863102in}{2.269457in}}{\pgfqpoint{3.868688in}{2.271771in}}{\pgfqpoint{3.872806in}{2.275889in}}%
\pgfpathcurveto{\pgfqpoint{3.876924in}{2.280007in}}{\pgfqpoint{3.879238in}{2.285593in}}{\pgfqpoint{3.879238in}{2.291417in}}%
\pgfpathcurveto{\pgfqpoint{3.879238in}{2.297241in}}{\pgfqpoint{3.876924in}{2.302828in}}{\pgfqpoint{3.872806in}{2.306946in}}%
\pgfpathcurveto{\pgfqpoint{3.868688in}{2.311064in}}{\pgfqpoint{3.863102in}{2.313378in}}{\pgfqpoint{3.857278in}{2.313378in}}%
\pgfpathcurveto{\pgfqpoint{3.851454in}{2.313378in}}{\pgfqpoint{3.845868in}{2.311064in}}{\pgfqpoint{3.841749in}{2.306946in}}%
\pgfpathcurveto{\pgfqpoint{3.837631in}{2.302828in}}{\pgfqpoint{3.835317in}{2.297241in}}{\pgfqpoint{3.835317in}{2.291417in}}%
\pgfpathcurveto{\pgfqpoint{3.835317in}{2.285593in}}{\pgfqpoint{3.837631in}{2.280007in}}{\pgfqpoint{3.841749in}{2.275889in}}%
\pgfpathcurveto{\pgfqpoint{3.845868in}{2.271771in}}{\pgfqpoint{3.851454in}{2.269457in}}{\pgfqpoint{3.857278in}{2.269457in}}%
\pgfpathlineto{\pgfqpoint{3.857278in}{2.269457in}}%
\pgfpathclose%
\pgfusepath{stroke,fill}%
\end{pgfscope}%
\begin{pgfscope}%
\pgfpathrectangle{\pgfqpoint{1.414736in}{0.548769in}}{\pgfqpoint{3.497747in}{7.253186in}}%
\pgfusepath{clip}%
\pgfsetbuttcap%
\pgfsetroundjoin%
\definecolor{currentfill}{rgb}{0.288921,0.758394,0.428426}%
\pgfsetfillcolor{currentfill}%
\pgfsetlinewidth{1.003750pt}%
\definecolor{currentstroke}{rgb}{0.288921,0.758394,0.428426}%
\pgfsetstrokecolor{currentstroke}%
\pgfsetdash{}{0pt}%
\pgfpathmoveto{\pgfqpoint{3.720922in}{2.740443in}}%
\pgfpathcurveto{\pgfqpoint{3.726745in}{2.740443in}}{\pgfqpoint{3.732332in}{2.742757in}}{\pgfqpoint{3.736450in}{2.746875in}}%
\pgfpathcurveto{\pgfqpoint{3.740568in}{2.750993in}}{\pgfqpoint{3.742882in}{2.756580in}}{\pgfqpoint{3.742882in}{2.762404in}}%
\pgfpathcurveto{\pgfqpoint{3.742882in}{2.768227in}}{\pgfqpoint{3.740568in}{2.773814in}}{\pgfqpoint{3.736450in}{2.777932in}}%
\pgfpathcurveto{\pgfqpoint{3.732332in}{2.782050in}}{\pgfqpoint{3.726745in}{2.784364in}}{\pgfqpoint{3.720922in}{2.784364in}}%
\pgfpathcurveto{\pgfqpoint{3.715098in}{2.784364in}}{\pgfqpoint{3.709511in}{2.782050in}}{\pgfqpoint{3.705393in}{2.777932in}}%
\pgfpathcurveto{\pgfqpoint{3.701275in}{2.773814in}}{\pgfqpoint{3.698961in}{2.768227in}}{\pgfqpoint{3.698961in}{2.762404in}}%
\pgfpathcurveto{\pgfqpoint{3.698961in}{2.756580in}}{\pgfqpoint{3.701275in}{2.750993in}}{\pgfqpoint{3.705393in}{2.746875in}}%
\pgfpathcurveto{\pgfqpoint{3.709511in}{2.742757in}}{\pgfqpoint{3.715098in}{2.740443in}}{\pgfqpoint{3.720922in}{2.740443in}}%
\pgfpathlineto{\pgfqpoint{3.720922in}{2.740443in}}%
\pgfpathclose%
\pgfusepath{stroke,fill}%
\end{pgfscope}%
\begin{pgfscope}%
\pgfpathrectangle{\pgfqpoint{1.414736in}{0.548769in}}{\pgfqpoint{3.497747in}{7.253186in}}%
\pgfusepath{clip}%
\pgfsetbuttcap%
\pgfsetroundjoin%
\definecolor{currentfill}{rgb}{0.926106,0.897330,0.104071}%
\pgfsetfillcolor{currentfill}%
\pgfsetlinewidth{1.003750pt}%
\definecolor{currentstroke}{rgb}{0.926106,0.897330,0.104071}%
\pgfsetstrokecolor{currentstroke}%
\pgfsetdash{}{0pt}%
\pgfpathmoveto{\pgfqpoint{3.515121in}{3.211429in}}%
\pgfpathcurveto{\pgfqpoint{3.520945in}{3.211429in}}{\pgfqpoint{3.526531in}{3.213743in}}{\pgfqpoint{3.530649in}{3.217861in}}%
\pgfpathcurveto{\pgfqpoint{3.534768in}{3.221980in}}{\pgfqpoint{3.537081in}{3.227566in}}{\pgfqpoint{3.537081in}{3.233390in}}%
\pgfpathcurveto{\pgfqpoint{3.537081in}{3.239214in}}{\pgfqpoint{3.534768in}{3.244800in}}{\pgfqpoint{3.530649in}{3.248918in}}%
\pgfpathcurveto{\pgfqpoint{3.526531in}{3.253036in}}{\pgfqpoint{3.520945in}{3.255350in}}{\pgfqpoint{3.515121in}{3.255350in}}%
\pgfpathcurveto{\pgfqpoint{3.509297in}{3.255350in}}{\pgfqpoint{3.503711in}{3.253036in}}{\pgfqpoint{3.499593in}{3.248918in}}%
\pgfpathcurveto{\pgfqpoint{3.495475in}{3.244800in}}{\pgfqpoint{3.493161in}{3.239214in}}{\pgfqpoint{3.493161in}{3.233390in}}%
\pgfpathcurveto{\pgfqpoint{3.493161in}{3.227566in}}{\pgfqpoint{3.495475in}{3.221980in}}{\pgfqpoint{3.499593in}{3.217861in}}%
\pgfpathcurveto{\pgfqpoint{3.503711in}{3.213743in}}{\pgfqpoint{3.509297in}{3.211429in}}{\pgfqpoint{3.515121in}{3.211429in}}%
\pgfpathlineto{\pgfqpoint{3.515121in}{3.211429in}}%
\pgfpathclose%
\pgfusepath{stroke,fill}%
\end{pgfscope}%
\begin{pgfscope}%
\pgfpathrectangle{\pgfqpoint{1.414736in}{0.548769in}}{\pgfqpoint{3.497747in}{7.253186in}}%
\pgfusepath{clip}%
\pgfsetbuttcap%
\pgfsetroundjoin%
\definecolor{currentfill}{rgb}{0.993248,0.906157,0.143936}%
\pgfsetfillcolor{currentfill}%
\pgfsetlinewidth{1.003750pt}%
\definecolor{currentstroke}{rgb}{0.993248,0.906157,0.143936}%
\pgfsetstrokecolor{currentstroke}%
\pgfsetdash{}{0pt}%
\pgfpathmoveto{\pgfqpoint{3.260378in}{3.682415in}}%
\pgfpathcurveto{\pgfqpoint{3.266201in}{3.682415in}}{\pgfqpoint{3.271788in}{3.684729in}}{\pgfqpoint{3.275906in}{3.688847in}}%
\pgfpathcurveto{\pgfqpoint{3.280024in}{3.692966in}}{\pgfqpoint{3.282338in}{3.698552in}}{\pgfqpoint{3.282338in}{3.704376in}}%
\pgfpathcurveto{\pgfqpoint{3.282338in}{3.710200in}}{\pgfqpoint{3.280024in}{3.715786in}}{\pgfqpoint{3.275906in}{3.719904in}}%
\pgfpathcurveto{\pgfqpoint{3.271788in}{3.724022in}}{\pgfqpoint{3.266201in}{3.726336in}}{\pgfqpoint{3.260378in}{3.726336in}}%
\pgfpathcurveto{\pgfqpoint{3.254554in}{3.726336in}}{\pgfqpoint{3.248967in}{3.724022in}}{\pgfqpoint{3.244849in}{3.719904in}}%
\pgfpathcurveto{\pgfqpoint{3.240731in}{3.715786in}}{\pgfqpoint{3.238417in}{3.710200in}}{\pgfqpoint{3.238417in}{3.704376in}}%
\pgfpathcurveto{\pgfqpoint{3.238417in}{3.698552in}}{\pgfqpoint{3.240731in}{3.692966in}}{\pgfqpoint{3.244849in}{3.688847in}}%
\pgfpathcurveto{\pgfqpoint{3.248967in}{3.684729in}}{\pgfqpoint{3.254554in}{3.682415in}}{\pgfqpoint{3.260378in}{3.682415in}}%
\pgfpathlineto{\pgfqpoint{3.260378in}{3.682415in}}%
\pgfpathclose%
\pgfusepath{stroke,fill}%
\end{pgfscope}%
\begin{pgfscope}%
\pgfpathrectangle{\pgfqpoint{1.414736in}{0.548769in}}{\pgfqpoint{3.497747in}{7.253186in}}%
\pgfusepath{clip}%
\pgfsetbuttcap%
\pgfsetroundjoin%
\definecolor{currentfill}{rgb}{0.993248,0.906157,0.143936}%
\pgfsetfillcolor{currentfill}%
\pgfsetlinewidth{1.003750pt}%
\definecolor{currentstroke}{rgb}{0.993248,0.906157,0.143936}%
\pgfsetstrokecolor{currentstroke}%
\pgfsetdash{}{0pt}%
\pgfpathmoveto{\pgfqpoint{3.776019in}{4.153402in}}%
\pgfpathcurveto{\pgfqpoint{3.781843in}{4.153402in}}{\pgfqpoint{3.787429in}{4.155715in}}{\pgfqpoint{3.791547in}{4.159834in}}%
\pgfpathcurveto{\pgfqpoint{3.795665in}{4.163952in}}{\pgfqpoint{3.797979in}{4.169538in}}{\pgfqpoint{3.797979in}{4.175362in}}%
\pgfpathcurveto{\pgfqpoint{3.797979in}{4.181186in}}{\pgfqpoint{3.795665in}{4.186772in}}{\pgfqpoint{3.791547in}{4.190890in}}%
\pgfpathcurveto{\pgfqpoint{3.787429in}{4.195008in}}{\pgfqpoint{3.781843in}{4.197322in}}{\pgfqpoint{3.776019in}{4.197322in}}%
\pgfpathcurveto{\pgfqpoint{3.770195in}{4.197322in}}{\pgfqpoint{3.764609in}{4.195008in}}{\pgfqpoint{3.760491in}{4.190890in}}%
\pgfpathcurveto{\pgfqpoint{3.756372in}{4.186772in}}{\pgfqpoint{3.754059in}{4.181186in}}{\pgfqpoint{3.754059in}{4.175362in}}%
\pgfpathcurveto{\pgfqpoint{3.754059in}{4.169538in}}{\pgfqpoint{3.756372in}{4.163952in}}{\pgfqpoint{3.760491in}{4.159834in}}%
\pgfpathcurveto{\pgfqpoint{3.764609in}{4.155715in}}{\pgfqpoint{3.770195in}{4.153402in}}{\pgfqpoint{3.776019in}{4.153402in}}%
\pgfpathlineto{\pgfqpoint{3.776019in}{4.153402in}}%
\pgfpathclose%
\pgfusepath{stroke,fill}%
\end{pgfscope}%
\begin{pgfscope}%
\pgfpathrectangle{\pgfqpoint{1.414736in}{0.548769in}}{\pgfqpoint{3.497747in}{7.253186in}}%
\pgfusepath{clip}%
\pgfsetbuttcap%
\pgfsetroundjoin%
\definecolor{currentfill}{rgb}{0.616293,0.852709,0.230052}%
\pgfsetfillcolor{currentfill}%
\pgfsetlinewidth{1.003750pt}%
\definecolor{currentstroke}{rgb}{0.616293,0.852709,0.230052}%
\pgfsetstrokecolor{currentstroke}%
\pgfsetdash{}{0pt}%
\pgfpathmoveto{\pgfqpoint{3.106785in}{4.624388in}}%
\pgfpathcurveto{\pgfqpoint{3.112609in}{4.624388in}}{\pgfqpoint{3.118196in}{4.626702in}}{\pgfqpoint{3.122314in}{4.630820in}}%
\pgfpathcurveto{\pgfqpoint{3.126432in}{4.634938in}}{\pgfqpoint{3.128746in}{4.640524in}}{\pgfqpoint{3.128746in}{4.646348in}}%
\pgfpathcurveto{\pgfqpoint{3.128746in}{4.652172in}}{\pgfqpoint{3.126432in}{4.657758in}}{\pgfqpoint{3.122314in}{4.661876in}}%
\pgfpathcurveto{\pgfqpoint{3.118196in}{4.665994in}}{\pgfqpoint{3.112609in}{4.668308in}}{\pgfqpoint{3.106785in}{4.668308in}}%
\pgfpathcurveto{\pgfqpoint{3.100961in}{4.668308in}}{\pgfqpoint{3.095375in}{4.665994in}}{\pgfqpoint{3.091257in}{4.661876in}}%
\pgfpathcurveto{\pgfqpoint{3.087139in}{4.657758in}}{\pgfqpoint{3.084825in}{4.652172in}}{\pgfqpoint{3.084825in}{4.646348in}}%
\pgfpathcurveto{\pgfqpoint{3.084825in}{4.640524in}}{\pgfqpoint{3.087139in}{4.634938in}}{\pgfqpoint{3.091257in}{4.630820in}}%
\pgfpathcurveto{\pgfqpoint{3.095375in}{4.626702in}}{\pgfqpoint{3.100961in}{4.624388in}}{\pgfqpoint{3.106785in}{4.624388in}}%
\pgfpathlineto{\pgfqpoint{3.106785in}{4.624388in}}%
\pgfpathclose%
\pgfusepath{stroke,fill}%
\end{pgfscope}%
\begin{pgfscope}%
\pgfpathrectangle{\pgfqpoint{1.414736in}{0.548769in}}{\pgfqpoint{3.497747in}{7.253186in}}%
\pgfusepath{clip}%
\pgfsetbuttcap%
\pgfsetroundjoin%
\definecolor{currentfill}{rgb}{0.120092,0.600104,0.542530}%
\pgfsetfillcolor{currentfill}%
\pgfsetlinewidth{1.003750pt}%
\definecolor{currentstroke}{rgb}{0.120092,0.600104,0.542530}%
\pgfsetstrokecolor{currentstroke}%
\pgfsetdash{}{0pt}%
\pgfpathmoveto{\pgfqpoint{3.853416in}{5.095374in}}%
\pgfpathcurveto{\pgfqpoint{3.859240in}{5.095374in}}{\pgfqpoint{3.864826in}{5.097688in}}{\pgfqpoint{3.868944in}{5.101806in}}%
\pgfpathcurveto{\pgfqpoint{3.873062in}{5.105924in}}{\pgfqpoint{3.875376in}{5.111510in}}{\pgfqpoint{3.875376in}{5.117334in}}%
\pgfpathcurveto{\pgfqpoint{3.875376in}{5.123158in}}{\pgfqpoint{3.873062in}{5.128744in}}{\pgfqpoint{3.868944in}{5.132862in}}%
\pgfpathcurveto{\pgfqpoint{3.864826in}{5.136980in}}{\pgfqpoint{3.859240in}{5.139294in}}{\pgfqpoint{3.853416in}{5.139294in}}%
\pgfpathcurveto{\pgfqpoint{3.847592in}{5.139294in}}{\pgfqpoint{3.842006in}{5.136980in}}{\pgfqpoint{3.837888in}{5.132862in}}%
\pgfpathcurveto{\pgfqpoint{3.833770in}{5.128744in}}{\pgfqpoint{3.831456in}{5.123158in}}{\pgfqpoint{3.831456in}{5.117334in}}%
\pgfpathcurveto{\pgfqpoint{3.831456in}{5.111510in}}{\pgfqpoint{3.833770in}{5.105924in}}{\pgfqpoint{3.837888in}{5.101806in}}%
\pgfpathcurveto{\pgfqpoint{3.842006in}{5.097688in}}{\pgfqpoint{3.847592in}{5.095374in}}{\pgfqpoint{3.853416in}{5.095374in}}%
\pgfpathlineto{\pgfqpoint{3.853416in}{5.095374in}}%
\pgfpathclose%
\pgfusepath{stroke,fill}%
\end{pgfscope}%
\begin{pgfscope}%
\pgfpathrectangle{\pgfqpoint{1.414736in}{0.548769in}}{\pgfqpoint{3.497747in}{7.253186in}}%
\pgfusepath{clip}%
\pgfsetbuttcap%
\pgfsetroundjoin%
\definecolor{currentfill}{rgb}{0.157729,0.485932,0.558013}%
\pgfsetfillcolor{currentfill}%
\pgfsetlinewidth{1.003750pt}%
\definecolor{currentstroke}{rgb}{0.157729,0.485932,0.558013}%
\pgfsetstrokecolor{currentstroke}%
\pgfsetdash{}{0pt}%
\pgfpathmoveto{\pgfqpoint{3.403306in}{5.566360in}}%
\pgfpathcurveto{\pgfqpoint{3.409130in}{5.566360in}}{\pgfqpoint{3.414716in}{5.568674in}}{\pgfqpoint{3.418834in}{5.572792in}}%
\pgfpathcurveto{\pgfqpoint{3.422952in}{5.576910in}}{\pgfqpoint{3.425266in}{5.582496in}}{\pgfqpoint{3.425266in}{5.588320in}}%
\pgfpathcurveto{\pgfqpoint{3.425266in}{5.594144in}}{\pgfqpoint{3.422952in}{5.599730in}}{\pgfqpoint{3.418834in}{5.603848in}}%
\pgfpathcurveto{\pgfqpoint{3.414716in}{5.607967in}}{\pgfqpoint{3.409130in}{5.610280in}}{\pgfqpoint{3.403306in}{5.610280in}}%
\pgfpathcurveto{\pgfqpoint{3.397482in}{5.610280in}}{\pgfqpoint{3.391896in}{5.607967in}}{\pgfqpoint{3.387778in}{5.603848in}}%
\pgfpathcurveto{\pgfqpoint{3.383660in}{5.599730in}}{\pgfqpoint{3.381346in}{5.594144in}}{\pgfqpoint{3.381346in}{5.588320in}}%
\pgfpathcurveto{\pgfqpoint{3.381346in}{5.582496in}}{\pgfqpoint{3.383660in}{5.576910in}}{\pgfqpoint{3.387778in}{5.572792in}}%
\pgfpathcurveto{\pgfqpoint{3.391896in}{5.568674in}}{\pgfqpoint{3.397482in}{5.566360in}}{\pgfqpoint{3.403306in}{5.566360in}}%
\pgfpathlineto{\pgfqpoint{3.403306in}{5.566360in}}%
\pgfpathclose%
\pgfusepath{stroke,fill}%
\end{pgfscope}%
\begin{pgfscope}%
\pgfpathrectangle{\pgfqpoint{1.414736in}{0.548769in}}{\pgfqpoint{3.497747in}{7.253186in}}%
\pgfusepath{clip}%
\pgfsetbuttcap%
\pgfsetroundjoin%
\definecolor{currentfill}{rgb}{0.458674,0.816363,0.329727}%
\pgfsetfillcolor{currentfill}%
\pgfsetlinewidth{1.003750pt}%
\definecolor{currentstroke}{rgb}{0.458674,0.816363,0.329727}%
\pgfsetstrokecolor{currentstroke}%
\pgfsetdash{}{0pt}%
\pgfpathmoveto{\pgfqpoint{3.371843in}{6.037346in}}%
\pgfpathcurveto{\pgfqpoint{3.377666in}{6.037346in}}{\pgfqpoint{3.383253in}{6.039660in}}{\pgfqpoint{3.387371in}{6.043778in}}%
\pgfpathcurveto{\pgfqpoint{3.391489in}{6.047896in}}{\pgfqpoint{3.393803in}{6.053482in}}{\pgfqpoint{3.393803in}{6.059306in}}%
\pgfpathcurveto{\pgfqpoint{3.393803in}{6.065130in}}{\pgfqpoint{3.391489in}{6.070716in}}{\pgfqpoint{3.387371in}{6.074835in}}%
\pgfpathcurveto{\pgfqpoint{3.383253in}{6.078953in}}{\pgfqpoint{3.377666in}{6.081267in}}{\pgfqpoint{3.371843in}{6.081267in}}%
\pgfpathcurveto{\pgfqpoint{3.366019in}{6.081267in}}{\pgfqpoint{3.360432in}{6.078953in}}{\pgfqpoint{3.356314in}{6.074835in}}%
\pgfpathcurveto{\pgfqpoint{3.352196in}{6.070716in}}{\pgfqpoint{3.349882in}{6.065130in}}{\pgfqpoint{3.349882in}{6.059306in}}%
\pgfpathcurveto{\pgfqpoint{3.349882in}{6.053482in}}{\pgfqpoint{3.352196in}{6.047896in}}{\pgfqpoint{3.356314in}{6.043778in}}%
\pgfpathcurveto{\pgfqpoint{3.360432in}{6.039660in}}{\pgfqpoint{3.366019in}{6.037346in}}{\pgfqpoint{3.371843in}{6.037346in}}%
\pgfpathlineto{\pgfqpoint{3.371843in}{6.037346in}}%
\pgfpathclose%
\pgfusepath{stroke,fill}%
\end{pgfscope}%
\begin{pgfscope}%
\pgfpathrectangle{\pgfqpoint{1.414736in}{0.548769in}}{\pgfqpoint{3.497747in}{7.253186in}}%
\pgfusepath{clip}%
\pgfsetbuttcap%
\pgfsetroundjoin%
\definecolor{currentfill}{rgb}{0.993248,0.906157,0.143936}%
\pgfsetfillcolor{currentfill}%
\pgfsetlinewidth{1.003750pt}%
\definecolor{currentstroke}{rgb}{0.993248,0.906157,0.143936}%
\pgfsetstrokecolor{currentstroke}%
\pgfsetdash{}{0pt}%
\pgfpathmoveto{\pgfqpoint{3.102011in}{6.508332in}}%
\pgfpathcurveto{\pgfqpoint{3.107835in}{6.508332in}}{\pgfqpoint{3.113422in}{6.510646in}}{\pgfqpoint{3.117540in}{6.514764in}}%
\pgfpathcurveto{\pgfqpoint{3.121658in}{6.518882in}}{\pgfqpoint{3.123972in}{6.524468in}}{\pgfqpoint{3.123972in}{6.530292in}}%
\pgfpathcurveto{\pgfqpoint{3.123972in}{6.536116in}}{\pgfqpoint{3.121658in}{6.541702in}}{\pgfqpoint{3.117540in}{6.545821in}}%
\pgfpathcurveto{\pgfqpoint{3.113422in}{6.549939in}}{\pgfqpoint{3.107835in}{6.552253in}}{\pgfqpoint{3.102011in}{6.552253in}}%
\pgfpathcurveto{\pgfqpoint{3.096188in}{6.552253in}}{\pgfqpoint{3.090601in}{6.549939in}}{\pgfqpoint{3.086483in}{6.545821in}}%
\pgfpathcurveto{\pgfqpoint{3.082365in}{6.541702in}}{\pgfqpoint{3.080051in}{6.536116in}}{\pgfqpoint{3.080051in}{6.530292in}}%
\pgfpathcurveto{\pgfqpoint{3.080051in}{6.524468in}}{\pgfqpoint{3.082365in}{6.518882in}}{\pgfqpoint{3.086483in}{6.514764in}}%
\pgfpathcurveto{\pgfqpoint{3.090601in}{6.510646in}}{\pgfqpoint{3.096188in}{6.508332in}}{\pgfqpoint{3.102011in}{6.508332in}}%
\pgfpathlineto{\pgfqpoint{3.102011in}{6.508332in}}%
\pgfpathclose%
\pgfusepath{stroke,fill}%
\end{pgfscope}%
\begin{pgfscope}%
\pgfpathrectangle{\pgfqpoint{1.414736in}{0.548769in}}{\pgfqpoint{3.497747in}{7.253186in}}%
\pgfusepath{clip}%
\pgfsetbuttcap%
\pgfsetroundjoin%
\definecolor{currentfill}{rgb}{0.845561,0.887322,0.099702}%
\pgfsetfillcolor{currentfill}%
\pgfsetlinewidth{1.003750pt}%
\definecolor{currentstroke}{rgb}{0.845561,0.887322,0.099702}%
\pgfsetstrokecolor{currentstroke}%
\pgfsetdash{}{0pt}%
\pgfpathmoveto{\pgfqpoint{4.101478in}{6.979318in}}%
\pgfpathcurveto{\pgfqpoint{4.107302in}{6.979318in}}{\pgfqpoint{4.112888in}{6.981632in}}{\pgfqpoint{4.117006in}{6.985750in}}%
\pgfpathcurveto{\pgfqpoint{4.121124in}{6.989868in}}{\pgfqpoint{4.123438in}{6.995455in}}{\pgfqpoint{4.123438in}{7.001278in}}%
\pgfpathcurveto{\pgfqpoint{4.123438in}{7.007102in}}{\pgfqpoint{4.121124in}{7.012689in}}{\pgfqpoint{4.117006in}{7.016807in}}%
\pgfpathcurveto{\pgfqpoint{4.112888in}{7.020925in}}{\pgfqpoint{4.107302in}{7.023239in}}{\pgfqpoint{4.101478in}{7.023239in}}%
\pgfpathcurveto{\pgfqpoint{4.095654in}{7.023239in}}{\pgfqpoint{4.090068in}{7.020925in}}{\pgfqpoint{4.085950in}{7.016807in}}%
\pgfpathcurveto{\pgfqpoint{4.081831in}{7.012689in}}{\pgfqpoint{4.079518in}{7.007102in}}{\pgfqpoint{4.079518in}{7.001278in}}%
\pgfpathcurveto{\pgfqpoint{4.079518in}{6.995455in}}{\pgfqpoint{4.081831in}{6.989868in}}{\pgfqpoint{4.085950in}{6.985750in}}%
\pgfpathcurveto{\pgfqpoint{4.090068in}{6.981632in}}{\pgfqpoint{4.095654in}{6.979318in}}{\pgfqpoint{4.101478in}{6.979318in}}%
\pgfpathlineto{\pgfqpoint{4.101478in}{6.979318in}}%
\pgfpathclose%
\pgfusepath{stroke,fill}%
\end{pgfscope}%
\begin{pgfscope}%
\pgfpathrectangle{\pgfqpoint{1.414736in}{0.548769in}}{\pgfqpoint{3.497747in}{7.253186in}}%
\pgfusepath{clip}%
\pgfsetbuttcap%
\pgfsetroundjoin%
\definecolor{currentfill}{rgb}{0.916242,0.896091,0.100717}%
\pgfsetfillcolor{currentfill}%
\pgfsetlinewidth{1.003750pt}%
\definecolor{currentstroke}{rgb}{0.916242,0.896091,0.100717}%
\pgfsetstrokecolor{currentstroke}%
\pgfsetdash{}{0pt}%
\pgfpathmoveto{\pgfqpoint{3.798682in}{7.450304in}}%
\pgfpathcurveto{\pgfqpoint{3.804506in}{7.450304in}}{\pgfqpoint{3.810092in}{7.452618in}}{\pgfqpoint{3.814210in}{7.456736in}}%
\pgfpathcurveto{\pgfqpoint{3.818328in}{7.460854in}}{\pgfqpoint{3.820642in}{7.466441in}}{\pgfqpoint{3.820642in}{7.472265in}}%
\pgfpathcurveto{\pgfqpoint{3.820642in}{7.478089in}}{\pgfqpoint{3.818328in}{7.483675in}}{\pgfqpoint{3.814210in}{7.487793in}}%
\pgfpathcurveto{\pgfqpoint{3.810092in}{7.491911in}}{\pgfqpoint{3.804506in}{7.494225in}}{\pgfqpoint{3.798682in}{7.494225in}}%
\pgfpathcurveto{\pgfqpoint{3.792858in}{7.494225in}}{\pgfqpoint{3.787272in}{7.491911in}}{\pgfqpoint{3.783154in}{7.487793in}}%
\pgfpathcurveto{\pgfqpoint{3.779036in}{7.483675in}}{\pgfqpoint{3.776722in}{7.478089in}}{\pgfqpoint{3.776722in}{7.472265in}}%
\pgfpathcurveto{\pgfqpoint{3.776722in}{7.466441in}}{\pgfqpoint{3.779036in}{7.460854in}}{\pgfqpoint{3.783154in}{7.456736in}}%
\pgfpathcurveto{\pgfqpoint{3.787272in}{7.452618in}}{\pgfqpoint{3.792858in}{7.450304in}}{\pgfqpoint{3.798682in}{7.450304in}}%
\pgfpathlineto{\pgfqpoint{3.798682in}{7.450304in}}%
\pgfpathclose%
\pgfusepath{stroke,fill}%
\end{pgfscope}%
\begin{pgfscope}%
\pgfpathrectangle{\pgfqpoint{1.414736in}{0.548769in}}{\pgfqpoint{3.497747in}{7.253186in}}%
\pgfusepath{clip}%
\pgfsetbuttcap%
\pgfsetroundjoin%
\definecolor{currentfill}{rgb}{0.935904,0.898570,0.108131}%
\pgfsetfillcolor{currentfill}%
\pgfsetlinewidth{1.003750pt}%
\definecolor{currentstroke}{rgb}{0.935904,0.898570,0.108131}%
\pgfsetstrokecolor{currentstroke}%
\pgfsetdash{}{0pt}%
\pgfpathmoveto{\pgfqpoint{3.432316in}{0.856499in}}%
\pgfpathcurveto{\pgfqpoint{3.438140in}{0.856499in}}{\pgfqpoint{3.443726in}{0.858813in}}{\pgfqpoint{3.447844in}{0.862931in}}%
\pgfpathcurveto{\pgfqpoint{3.451962in}{0.867049in}}{\pgfqpoint{3.454276in}{0.872635in}}{\pgfqpoint{3.454276in}{0.878459in}}%
\pgfpathcurveto{\pgfqpoint{3.454276in}{0.884283in}}{\pgfqpoint{3.451962in}{0.889869in}}{\pgfqpoint{3.447844in}{0.893987in}}%
\pgfpathcurveto{\pgfqpoint{3.443726in}{0.898105in}}{\pgfqpoint{3.438140in}{0.900419in}}{\pgfqpoint{3.432316in}{0.900419in}}%
\pgfpathcurveto{\pgfqpoint{3.426492in}{0.900419in}}{\pgfqpoint{3.420905in}{0.898105in}}{\pgfqpoint{3.416787in}{0.893987in}}%
\pgfpathcurveto{\pgfqpoint{3.412669in}{0.889869in}}{\pgfqpoint{3.410355in}{0.884283in}}{\pgfqpoint{3.410355in}{0.878459in}}%
\pgfpathcurveto{\pgfqpoint{3.410355in}{0.872635in}}{\pgfqpoint{3.412669in}{0.867049in}}{\pgfqpoint{3.416787in}{0.862931in}}%
\pgfpathcurveto{\pgfqpoint{3.420905in}{0.858813in}}{\pgfqpoint{3.426492in}{0.856499in}}{\pgfqpoint{3.432316in}{0.856499in}}%
\pgfpathlineto{\pgfqpoint{3.432316in}{0.856499in}}%
\pgfpathclose%
\pgfusepath{stroke,fill}%
\end{pgfscope}%
\begin{pgfscope}%
\pgfpathrectangle{\pgfqpoint{1.414736in}{0.548769in}}{\pgfqpoint{3.497747in}{7.253186in}}%
\pgfusepath{clip}%
\pgfsetbuttcap%
\pgfsetroundjoin%
\definecolor{currentfill}{rgb}{0.449368,0.813768,0.335384}%
\pgfsetfillcolor{currentfill}%
\pgfsetlinewidth{1.003750pt}%
\definecolor{currentstroke}{rgb}{0.449368,0.813768,0.335384}%
\pgfsetstrokecolor{currentstroke}%
\pgfsetdash{}{0pt}%
\pgfpathmoveto{\pgfqpoint{3.674225in}{1.327485in}}%
\pgfpathcurveto{\pgfqpoint{3.680049in}{1.327485in}}{\pgfqpoint{3.685635in}{1.329799in}}{\pgfqpoint{3.689753in}{1.333917in}}%
\pgfpathcurveto{\pgfqpoint{3.693871in}{1.338035in}}{\pgfqpoint{3.696185in}{1.343621in}}{\pgfqpoint{3.696185in}{1.349445in}}%
\pgfpathcurveto{\pgfqpoint{3.696185in}{1.355269in}}{\pgfqpoint{3.693871in}{1.360855in}}{\pgfqpoint{3.689753in}{1.364973in}}%
\pgfpathcurveto{\pgfqpoint{3.685635in}{1.369092in}}{\pgfqpoint{3.680049in}{1.371405in}}{\pgfqpoint{3.674225in}{1.371405in}}%
\pgfpathcurveto{\pgfqpoint{3.668401in}{1.371405in}}{\pgfqpoint{3.662815in}{1.369092in}}{\pgfqpoint{3.658697in}{1.364973in}}%
\pgfpathcurveto{\pgfqpoint{3.654578in}{1.360855in}}{\pgfqpoint{3.652265in}{1.355269in}}{\pgfqpoint{3.652265in}{1.349445in}}%
\pgfpathcurveto{\pgfqpoint{3.652265in}{1.343621in}}{\pgfqpoint{3.654578in}{1.338035in}}{\pgfqpoint{3.658697in}{1.333917in}}%
\pgfpathcurveto{\pgfqpoint{3.662815in}{1.329799in}}{\pgfqpoint{3.668401in}{1.327485in}}{\pgfqpoint{3.674225in}{1.327485in}}%
\pgfpathlineto{\pgfqpoint{3.674225in}{1.327485in}}%
\pgfpathclose%
\pgfusepath{stroke,fill}%
\end{pgfscope}%
\begin{pgfscope}%
\pgfpathrectangle{\pgfqpoint{1.414736in}{0.548769in}}{\pgfqpoint{3.497747in}{7.253186in}}%
\pgfusepath{clip}%
\pgfsetbuttcap%
\pgfsetroundjoin%
\definecolor{currentfill}{rgb}{0.699415,0.867117,0.175971}%
\pgfsetfillcolor{currentfill}%
\pgfsetlinewidth{1.003750pt}%
\definecolor{currentstroke}{rgb}{0.699415,0.867117,0.175971}%
\pgfsetstrokecolor{currentstroke}%
\pgfsetdash{}{0pt}%
\pgfpathmoveto{\pgfqpoint{3.674964in}{1.798471in}}%
\pgfpathcurveto{\pgfqpoint{3.680788in}{1.798471in}}{\pgfqpoint{3.686374in}{1.800785in}}{\pgfqpoint{3.690493in}{1.804903in}}%
\pgfpathcurveto{\pgfqpoint{3.694611in}{1.809021in}}{\pgfqpoint{3.696925in}{1.814607in}}{\pgfqpoint{3.696925in}{1.820431in}}%
\pgfpathcurveto{\pgfqpoint{3.696925in}{1.826255in}}{\pgfqpoint{3.694611in}{1.831841in}}{\pgfqpoint{3.690493in}{1.835960in}}%
\pgfpathcurveto{\pgfqpoint{3.686374in}{1.840078in}}{\pgfqpoint{3.680788in}{1.842392in}}{\pgfqpoint{3.674964in}{1.842392in}}%
\pgfpathcurveto{\pgfqpoint{3.669140in}{1.842392in}}{\pgfqpoint{3.663554in}{1.840078in}}{\pgfqpoint{3.659436in}{1.835960in}}%
\pgfpathcurveto{\pgfqpoint{3.655318in}{1.831841in}}{\pgfqpoint{3.653004in}{1.826255in}}{\pgfqpoint{3.653004in}{1.820431in}}%
\pgfpathcurveto{\pgfqpoint{3.653004in}{1.814607in}}{\pgfqpoint{3.655318in}{1.809021in}}{\pgfqpoint{3.659436in}{1.804903in}}%
\pgfpathcurveto{\pgfqpoint{3.663554in}{1.800785in}}{\pgfqpoint{3.669140in}{1.798471in}}{\pgfqpoint{3.674964in}{1.798471in}}%
\pgfpathlineto{\pgfqpoint{3.674964in}{1.798471in}}%
\pgfpathclose%
\pgfusepath{stroke,fill}%
\end{pgfscope}%
\begin{pgfscope}%
\pgfpathrectangle{\pgfqpoint{1.414736in}{0.548769in}}{\pgfqpoint{3.497747in}{7.253186in}}%
\pgfusepath{clip}%
\pgfsetbuttcap%
\pgfsetroundjoin%
\definecolor{currentfill}{rgb}{0.993248,0.906157,0.143936}%
\pgfsetfillcolor{currentfill}%
\pgfsetlinewidth{1.003750pt}%
\definecolor{currentstroke}{rgb}{0.993248,0.906157,0.143936}%
\pgfsetstrokecolor{currentstroke}%
\pgfsetdash{}{0pt}%
\pgfpathmoveto{\pgfqpoint{3.402016in}{2.269457in}}%
\pgfpathcurveto{\pgfqpoint{3.407840in}{2.269457in}}{\pgfqpoint{3.413426in}{2.271771in}}{\pgfqpoint{3.417544in}{2.275889in}}%
\pgfpathcurveto{\pgfqpoint{3.421662in}{2.280007in}}{\pgfqpoint{3.423976in}{2.285593in}}{\pgfqpoint{3.423976in}{2.291417in}}%
\pgfpathcurveto{\pgfqpoint{3.423976in}{2.297241in}}{\pgfqpoint{3.421662in}{2.302828in}}{\pgfqpoint{3.417544in}{2.306946in}}%
\pgfpathcurveto{\pgfqpoint{3.413426in}{2.311064in}}{\pgfqpoint{3.407840in}{2.313378in}}{\pgfqpoint{3.402016in}{2.313378in}}%
\pgfpathcurveto{\pgfqpoint{3.396192in}{2.313378in}}{\pgfqpoint{3.390606in}{2.311064in}}{\pgfqpoint{3.386487in}{2.306946in}}%
\pgfpathcurveto{\pgfqpoint{3.382369in}{2.302828in}}{\pgfqpoint{3.380055in}{2.297241in}}{\pgfqpoint{3.380055in}{2.291417in}}%
\pgfpathcurveto{\pgfqpoint{3.380055in}{2.285593in}}{\pgfqpoint{3.382369in}{2.280007in}}{\pgfqpoint{3.386487in}{2.275889in}}%
\pgfpathcurveto{\pgfqpoint{3.390606in}{2.271771in}}{\pgfqpoint{3.396192in}{2.269457in}}{\pgfqpoint{3.402016in}{2.269457in}}%
\pgfpathlineto{\pgfqpoint{3.402016in}{2.269457in}}%
\pgfpathclose%
\pgfusepath{stroke,fill}%
\end{pgfscope}%
\begin{pgfscope}%
\pgfpathrectangle{\pgfqpoint{1.414736in}{0.548769in}}{\pgfqpoint{3.497747in}{7.253186in}}%
\pgfusepath{clip}%
\pgfsetbuttcap%
\pgfsetroundjoin%
\definecolor{currentfill}{rgb}{0.468053,0.818921,0.323998}%
\pgfsetfillcolor{currentfill}%
\pgfsetlinewidth{1.003750pt}%
\definecolor{currentstroke}{rgb}{0.468053,0.818921,0.323998}%
\pgfsetstrokecolor{currentstroke}%
\pgfsetdash{}{0pt}%
\pgfpathmoveto{\pgfqpoint{3.610139in}{2.740443in}}%
\pgfpathcurveto{\pgfqpoint{3.615963in}{2.740443in}}{\pgfqpoint{3.621549in}{2.742757in}}{\pgfqpoint{3.625667in}{2.746875in}}%
\pgfpathcurveto{\pgfqpoint{3.629785in}{2.750993in}}{\pgfqpoint{3.632099in}{2.756580in}}{\pgfqpoint{3.632099in}{2.762404in}}%
\pgfpathcurveto{\pgfqpoint{3.632099in}{2.768227in}}{\pgfqpoint{3.629785in}{2.773814in}}{\pgfqpoint{3.625667in}{2.777932in}}%
\pgfpathcurveto{\pgfqpoint{3.621549in}{2.782050in}}{\pgfqpoint{3.615963in}{2.784364in}}{\pgfqpoint{3.610139in}{2.784364in}}%
\pgfpathcurveto{\pgfqpoint{3.604315in}{2.784364in}}{\pgfqpoint{3.598729in}{2.782050in}}{\pgfqpoint{3.594611in}{2.777932in}}%
\pgfpathcurveto{\pgfqpoint{3.590493in}{2.773814in}}{\pgfqpoint{3.588179in}{2.768227in}}{\pgfqpoint{3.588179in}{2.762404in}}%
\pgfpathcurveto{\pgfqpoint{3.588179in}{2.756580in}}{\pgfqpoint{3.590493in}{2.750993in}}{\pgfqpoint{3.594611in}{2.746875in}}%
\pgfpathcurveto{\pgfqpoint{3.598729in}{2.742757in}}{\pgfqpoint{3.604315in}{2.740443in}}{\pgfqpoint{3.610139in}{2.740443in}}%
\pgfpathlineto{\pgfqpoint{3.610139in}{2.740443in}}%
\pgfpathclose%
\pgfusepath{stroke,fill}%
\end{pgfscope}%
\begin{pgfscope}%
\pgfpathrectangle{\pgfqpoint{1.414736in}{0.548769in}}{\pgfqpoint{3.497747in}{7.253186in}}%
\pgfusepath{clip}%
\pgfsetbuttcap%
\pgfsetroundjoin%
\definecolor{currentfill}{rgb}{0.741388,0.873449,0.149561}%
\pgfsetfillcolor{currentfill}%
\pgfsetlinewidth{1.003750pt}%
\definecolor{currentstroke}{rgb}{0.741388,0.873449,0.149561}%
\pgfsetstrokecolor{currentstroke}%
\pgfsetdash{}{0pt}%
\pgfpathmoveto{\pgfqpoint{3.427609in}{3.211429in}}%
\pgfpathcurveto{\pgfqpoint{3.433433in}{3.211429in}}{\pgfqpoint{3.439019in}{3.213743in}}{\pgfqpoint{3.443137in}{3.217861in}}%
\pgfpathcurveto{\pgfqpoint{3.447255in}{3.221980in}}{\pgfqpoint{3.449569in}{3.227566in}}{\pgfqpoint{3.449569in}{3.233390in}}%
\pgfpathcurveto{\pgfqpoint{3.449569in}{3.239214in}}{\pgfqpoint{3.447255in}{3.244800in}}{\pgfqpoint{3.443137in}{3.248918in}}%
\pgfpathcurveto{\pgfqpoint{3.439019in}{3.253036in}}{\pgfqpoint{3.433433in}{3.255350in}}{\pgfqpoint{3.427609in}{3.255350in}}%
\pgfpathcurveto{\pgfqpoint{3.421785in}{3.255350in}}{\pgfqpoint{3.416199in}{3.253036in}}{\pgfqpoint{3.412080in}{3.248918in}}%
\pgfpathcurveto{\pgfqpoint{3.407962in}{3.244800in}}{\pgfqpoint{3.405648in}{3.239214in}}{\pgfqpoint{3.405648in}{3.233390in}}%
\pgfpathcurveto{\pgfqpoint{3.405648in}{3.227566in}}{\pgfqpoint{3.407962in}{3.221980in}}{\pgfqpoint{3.412080in}{3.217861in}}%
\pgfpathcurveto{\pgfqpoint{3.416199in}{3.213743in}}{\pgfqpoint{3.421785in}{3.211429in}}{\pgfqpoint{3.427609in}{3.211429in}}%
\pgfpathlineto{\pgfqpoint{3.427609in}{3.211429in}}%
\pgfpathclose%
\pgfusepath{stroke,fill}%
\end{pgfscope}%
\begin{pgfscope}%
\pgfpathrectangle{\pgfqpoint{1.414736in}{0.548769in}}{\pgfqpoint{3.497747in}{7.253186in}}%
\pgfusepath{clip}%
\pgfsetbuttcap%
\pgfsetroundjoin%
\definecolor{currentfill}{rgb}{0.993248,0.906157,0.143936}%
\pgfsetfillcolor{currentfill}%
\pgfsetlinewidth{1.003750pt}%
\definecolor{currentstroke}{rgb}{0.993248,0.906157,0.143936}%
\pgfsetstrokecolor{currentstroke}%
\pgfsetdash{}{0pt}%
\pgfpathmoveto{\pgfqpoint{2.851355in}{3.682415in}}%
\pgfpathcurveto{\pgfqpoint{2.857179in}{3.682415in}}{\pgfqpoint{2.862765in}{3.684729in}}{\pgfqpoint{2.866884in}{3.688847in}}%
\pgfpathcurveto{\pgfqpoint{2.871002in}{3.692966in}}{\pgfqpoint{2.873316in}{3.698552in}}{\pgfqpoint{2.873316in}{3.704376in}}%
\pgfpathcurveto{\pgfqpoint{2.873316in}{3.710200in}}{\pgfqpoint{2.871002in}{3.715786in}}{\pgfqpoint{2.866884in}{3.719904in}}%
\pgfpathcurveto{\pgfqpoint{2.862765in}{3.724022in}}{\pgfqpoint{2.857179in}{3.726336in}}{\pgfqpoint{2.851355in}{3.726336in}}%
\pgfpathcurveto{\pgfqpoint{2.845531in}{3.726336in}}{\pgfqpoint{2.839945in}{3.724022in}}{\pgfqpoint{2.835827in}{3.719904in}}%
\pgfpathcurveto{\pgfqpoint{2.831709in}{3.715786in}}{\pgfqpoint{2.829395in}{3.710200in}}{\pgfqpoint{2.829395in}{3.704376in}}%
\pgfpathcurveto{\pgfqpoint{2.829395in}{3.698552in}}{\pgfqpoint{2.831709in}{3.692966in}}{\pgfqpoint{2.835827in}{3.688847in}}%
\pgfpathcurveto{\pgfqpoint{2.839945in}{3.684729in}}{\pgfqpoint{2.845531in}{3.682415in}}{\pgfqpoint{2.851355in}{3.682415in}}%
\pgfpathlineto{\pgfqpoint{2.851355in}{3.682415in}}%
\pgfpathclose%
\pgfusepath{stroke,fill}%
\end{pgfscope}%
\begin{pgfscope}%
\pgfpathrectangle{\pgfqpoint{1.414736in}{0.548769in}}{\pgfqpoint{3.497747in}{7.253186in}}%
\pgfusepath{clip}%
\pgfsetbuttcap%
\pgfsetroundjoin%
\definecolor{currentfill}{rgb}{0.993248,0.906157,0.143936}%
\pgfsetfillcolor{currentfill}%
\pgfsetlinewidth{1.003750pt}%
\definecolor{currentstroke}{rgb}{0.993248,0.906157,0.143936}%
\pgfsetstrokecolor{currentstroke}%
\pgfsetdash{}{0pt}%
\pgfpathmoveto{\pgfqpoint{3.718575in}{4.153402in}}%
\pgfpathcurveto{\pgfqpoint{3.724399in}{4.153402in}}{\pgfqpoint{3.729985in}{4.155715in}}{\pgfqpoint{3.734103in}{4.159834in}}%
\pgfpathcurveto{\pgfqpoint{3.738221in}{4.163952in}}{\pgfqpoint{3.740535in}{4.169538in}}{\pgfqpoint{3.740535in}{4.175362in}}%
\pgfpathcurveto{\pgfqpoint{3.740535in}{4.181186in}}{\pgfqpoint{3.738221in}{4.186772in}}{\pgfqpoint{3.734103in}{4.190890in}}%
\pgfpathcurveto{\pgfqpoint{3.729985in}{4.195008in}}{\pgfqpoint{3.724399in}{4.197322in}}{\pgfqpoint{3.718575in}{4.197322in}}%
\pgfpathcurveto{\pgfqpoint{3.712751in}{4.197322in}}{\pgfqpoint{3.707165in}{4.195008in}}{\pgfqpoint{3.703046in}{4.190890in}}%
\pgfpathcurveto{\pgfqpoint{3.698928in}{4.186772in}}{\pgfqpoint{3.696614in}{4.181186in}}{\pgfqpoint{3.696614in}{4.175362in}}%
\pgfpathcurveto{\pgfqpoint{3.696614in}{4.169538in}}{\pgfqpoint{3.698928in}{4.163952in}}{\pgfqpoint{3.703046in}{4.159834in}}%
\pgfpathcurveto{\pgfqpoint{3.707165in}{4.155715in}}{\pgfqpoint{3.712751in}{4.153402in}}{\pgfqpoint{3.718575in}{4.153402in}}%
\pgfpathlineto{\pgfqpoint{3.718575in}{4.153402in}}%
\pgfpathclose%
\pgfusepath{stroke,fill}%
\end{pgfscope}%
\begin{pgfscope}%
\pgfpathrectangle{\pgfqpoint{1.414736in}{0.548769in}}{\pgfqpoint{3.497747in}{7.253186in}}%
\pgfusepath{clip}%
\pgfsetbuttcap%
\pgfsetroundjoin%
\definecolor{currentfill}{rgb}{0.668054,0.861999,0.196293}%
\pgfsetfillcolor{currentfill}%
\pgfsetlinewidth{1.003750pt}%
\definecolor{currentstroke}{rgb}{0.668054,0.861999,0.196293}%
\pgfsetstrokecolor{currentstroke}%
\pgfsetdash{}{0pt}%
\pgfpathmoveto{\pgfqpoint{3.245148in}{4.624388in}}%
\pgfpathcurveto{\pgfqpoint{3.250972in}{4.624388in}}{\pgfqpoint{3.256558in}{4.626702in}}{\pgfqpoint{3.260676in}{4.630820in}}%
\pgfpathcurveto{\pgfqpoint{3.264794in}{4.634938in}}{\pgfqpoint{3.267108in}{4.640524in}}{\pgfqpoint{3.267108in}{4.646348in}}%
\pgfpathcurveto{\pgfqpoint{3.267108in}{4.652172in}}{\pgfqpoint{3.264794in}{4.657758in}}{\pgfqpoint{3.260676in}{4.661876in}}%
\pgfpathcurveto{\pgfqpoint{3.256558in}{4.665994in}}{\pgfqpoint{3.250972in}{4.668308in}}{\pgfqpoint{3.245148in}{4.668308in}}%
\pgfpathcurveto{\pgfqpoint{3.239324in}{4.668308in}}{\pgfqpoint{3.233738in}{4.665994in}}{\pgfqpoint{3.229620in}{4.661876in}}%
\pgfpathcurveto{\pgfqpoint{3.225501in}{4.657758in}}{\pgfqpoint{3.223188in}{4.652172in}}{\pgfqpoint{3.223188in}{4.646348in}}%
\pgfpathcurveto{\pgfqpoint{3.223188in}{4.640524in}}{\pgfqpoint{3.225501in}{4.634938in}}{\pgfqpoint{3.229620in}{4.630820in}}%
\pgfpathcurveto{\pgfqpoint{3.233738in}{4.626702in}}{\pgfqpoint{3.239324in}{4.624388in}}{\pgfqpoint{3.245148in}{4.624388in}}%
\pgfpathlineto{\pgfqpoint{3.245148in}{4.624388in}}%
\pgfpathclose%
\pgfusepath{stroke,fill}%
\end{pgfscope}%
\begin{pgfscope}%
\pgfpathrectangle{\pgfqpoint{1.414736in}{0.548769in}}{\pgfqpoint{3.497747in}{7.253186in}}%
\pgfusepath{clip}%
\pgfsetbuttcap%
\pgfsetroundjoin%
\definecolor{currentfill}{rgb}{0.169646,0.456262,0.558030}%
\pgfsetfillcolor{currentfill}%
\pgfsetlinewidth{1.003750pt}%
\definecolor{currentstroke}{rgb}{0.169646,0.456262,0.558030}%
\pgfsetstrokecolor{currentstroke}%
\pgfsetdash{}{0pt}%
\pgfpathmoveto{\pgfqpoint{3.838085in}{5.095374in}}%
\pgfpathcurveto{\pgfqpoint{3.843909in}{5.095374in}}{\pgfqpoint{3.849495in}{5.097688in}}{\pgfqpoint{3.853613in}{5.101806in}}%
\pgfpathcurveto{\pgfqpoint{3.857731in}{5.105924in}}{\pgfqpoint{3.860045in}{5.111510in}}{\pgfqpoint{3.860045in}{5.117334in}}%
\pgfpathcurveto{\pgfqpoint{3.860045in}{5.123158in}}{\pgfqpoint{3.857731in}{5.128744in}}{\pgfqpoint{3.853613in}{5.132862in}}%
\pgfpathcurveto{\pgfqpoint{3.849495in}{5.136980in}}{\pgfqpoint{3.843909in}{5.139294in}}{\pgfqpoint{3.838085in}{5.139294in}}%
\pgfpathcurveto{\pgfqpoint{3.832261in}{5.139294in}}{\pgfqpoint{3.826675in}{5.136980in}}{\pgfqpoint{3.822556in}{5.132862in}}%
\pgfpathcurveto{\pgfqpoint{3.818438in}{5.128744in}}{\pgfqpoint{3.816124in}{5.123158in}}{\pgfqpoint{3.816124in}{5.117334in}}%
\pgfpathcurveto{\pgfqpoint{3.816124in}{5.111510in}}{\pgfqpoint{3.818438in}{5.105924in}}{\pgfqpoint{3.822556in}{5.101806in}}%
\pgfpathcurveto{\pgfqpoint{3.826675in}{5.097688in}}{\pgfqpoint{3.832261in}{5.095374in}}{\pgfqpoint{3.838085in}{5.095374in}}%
\pgfpathlineto{\pgfqpoint{3.838085in}{5.095374in}}%
\pgfpathclose%
\pgfusepath{stroke,fill}%
\end{pgfscope}%
\begin{pgfscope}%
\pgfpathrectangle{\pgfqpoint{1.414736in}{0.548769in}}{\pgfqpoint{3.497747in}{7.253186in}}%
\pgfusepath{clip}%
\pgfsetbuttcap%
\pgfsetroundjoin%
\definecolor{currentfill}{rgb}{0.151918,0.500685,0.557587}%
\pgfsetfillcolor{currentfill}%
\pgfsetlinewidth{1.003750pt}%
\definecolor{currentstroke}{rgb}{0.151918,0.500685,0.557587}%
\pgfsetstrokecolor{currentstroke}%
\pgfsetdash{}{0pt}%
\pgfpathmoveto{\pgfqpoint{3.649071in}{5.566360in}}%
\pgfpathcurveto{\pgfqpoint{3.654895in}{5.566360in}}{\pgfqpoint{3.660481in}{5.568674in}}{\pgfqpoint{3.664599in}{5.572792in}}%
\pgfpathcurveto{\pgfqpoint{3.668718in}{5.576910in}}{\pgfqpoint{3.671031in}{5.582496in}}{\pgfqpoint{3.671031in}{5.588320in}}%
\pgfpathcurveto{\pgfqpoint{3.671031in}{5.594144in}}{\pgfqpoint{3.668718in}{5.599730in}}{\pgfqpoint{3.664599in}{5.603848in}}%
\pgfpathcurveto{\pgfqpoint{3.660481in}{5.607967in}}{\pgfqpoint{3.654895in}{5.610280in}}{\pgfqpoint{3.649071in}{5.610280in}}%
\pgfpathcurveto{\pgfqpoint{3.643247in}{5.610280in}}{\pgfqpoint{3.637661in}{5.607967in}}{\pgfqpoint{3.633543in}{5.603848in}}%
\pgfpathcurveto{\pgfqpoint{3.629425in}{5.599730in}}{\pgfqpoint{3.627111in}{5.594144in}}{\pgfqpoint{3.627111in}{5.588320in}}%
\pgfpathcurveto{\pgfqpoint{3.627111in}{5.582496in}}{\pgfqpoint{3.629425in}{5.576910in}}{\pgfqpoint{3.633543in}{5.572792in}}%
\pgfpathcurveto{\pgfqpoint{3.637661in}{5.568674in}}{\pgfqpoint{3.643247in}{5.566360in}}{\pgfqpoint{3.649071in}{5.566360in}}%
\pgfpathlineto{\pgfqpoint{3.649071in}{5.566360in}}%
\pgfpathclose%
\pgfusepath{stroke,fill}%
\end{pgfscope}%
\begin{pgfscope}%
\pgfpathrectangle{\pgfqpoint{1.414736in}{0.548769in}}{\pgfqpoint{3.497747in}{7.253186in}}%
\pgfusepath{clip}%
\pgfsetbuttcap%
\pgfsetroundjoin%
\definecolor{currentfill}{rgb}{0.515992,0.831158,0.294279}%
\pgfsetfillcolor{currentfill}%
\pgfsetlinewidth{1.003750pt}%
\definecolor{currentstroke}{rgb}{0.515992,0.831158,0.294279}%
\pgfsetstrokecolor{currentstroke}%
\pgfsetdash{}{0pt}%
\pgfpathmoveto{\pgfqpoint{3.190529in}{6.037346in}}%
\pgfpathcurveto{\pgfqpoint{3.196353in}{6.037346in}}{\pgfqpoint{3.201939in}{6.039660in}}{\pgfqpoint{3.206057in}{6.043778in}}%
\pgfpathcurveto{\pgfqpoint{3.210175in}{6.047896in}}{\pgfqpoint{3.212489in}{6.053482in}}{\pgfqpoint{3.212489in}{6.059306in}}%
\pgfpathcurveto{\pgfqpoint{3.212489in}{6.065130in}}{\pgfqpoint{3.210175in}{6.070716in}}{\pgfqpoint{3.206057in}{6.074835in}}%
\pgfpathcurveto{\pgfqpoint{3.201939in}{6.078953in}}{\pgfqpoint{3.196353in}{6.081267in}}{\pgfqpoint{3.190529in}{6.081267in}}%
\pgfpathcurveto{\pgfqpoint{3.184705in}{6.081267in}}{\pgfqpoint{3.179119in}{6.078953in}}{\pgfqpoint{3.175001in}{6.074835in}}%
\pgfpathcurveto{\pgfqpoint{3.170883in}{6.070716in}}{\pgfqpoint{3.168569in}{6.065130in}}{\pgfqpoint{3.168569in}{6.059306in}}%
\pgfpathcurveto{\pgfqpoint{3.168569in}{6.053482in}}{\pgfqpoint{3.170883in}{6.047896in}}{\pgfqpoint{3.175001in}{6.043778in}}%
\pgfpathcurveto{\pgfqpoint{3.179119in}{6.039660in}}{\pgfqpoint{3.184705in}{6.037346in}}{\pgfqpoint{3.190529in}{6.037346in}}%
\pgfpathlineto{\pgfqpoint{3.190529in}{6.037346in}}%
\pgfpathclose%
\pgfusepath{stroke,fill}%
\end{pgfscope}%
\begin{pgfscope}%
\pgfpathrectangle{\pgfqpoint{1.414736in}{0.548769in}}{\pgfqpoint{3.497747in}{7.253186in}}%
\pgfusepath{clip}%
\pgfsetbuttcap%
\pgfsetroundjoin%
\definecolor{currentfill}{rgb}{0.993248,0.906157,0.143936}%
\pgfsetfillcolor{currentfill}%
\pgfsetlinewidth{1.003750pt}%
\definecolor{currentstroke}{rgb}{0.993248,0.906157,0.143936}%
\pgfsetstrokecolor{currentstroke}%
\pgfsetdash{}{0pt}%
\pgfpathmoveto{\pgfqpoint{3.451248in}{6.508332in}}%
\pgfpathcurveto{\pgfqpoint{3.457072in}{6.508332in}}{\pgfqpoint{3.462658in}{6.510646in}}{\pgfqpoint{3.466776in}{6.514764in}}%
\pgfpathcurveto{\pgfqpoint{3.470894in}{6.518882in}}{\pgfqpoint{3.473208in}{6.524468in}}{\pgfqpoint{3.473208in}{6.530292in}}%
\pgfpathcurveto{\pgfqpoint{3.473208in}{6.536116in}}{\pgfqpoint{3.470894in}{6.541702in}}{\pgfqpoint{3.466776in}{6.545821in}}%
\pgfpathcurveto{\pgfqpoint{3.462658in}{6.549939in}}{\pgfqpoint{3.457072in}{6.552253in}}{\pgfqpoint{3.451248in}{6.552253in}}%
\pgfpathcurveto{\pgfqpoint{3.445424in}{6.552253in}}{\pgfqpoint{3.439838in}{6.549939in}}{\pgfqpoint{3.435720in}{6.545821in}}%
\pgfpathcurveto{\pgfqpoint{3.431601in}{6.541702in}}{\pgfqpoint{3.429288in}{6.536116in}}{\pgfqpoint{3.429288in}{6.530292in}}%
\pgfpathcurveto{\pgfqpoint{3.429288in}{6.524468in}}{\pgfqpoint{3.431601in}{6.518882in}}{\pgfqpoint{3.435720in}{6.514764in}}%
\pgfpathcurveto{\pgfqpoint{3.439838in}{6.510646in}}{\pgfqpoint{3.445424in}{6.508332in}}{\pgfqpoint{3.451248in}{6.508332in}}%
\pgfpathlineto{\pgfqpoint{3.451248in}{6.508332in}}%
\pgfpathclose%
\pgfusepath{stroke,fill}%
\end{pgfscope}%
\begin{pgfscope}%
\pgfpathrectangle{\pgfqpoint{1.414736in}{0.548769in}}{\pgfqpoint{3.497747in}{7.253186in}}%
\pgfusepath{clip}%
\pgfsetbuttcap%
\pgfsetroundjoin%
\definecolor{currentfill}{rgb}{0.440137,0.811138,0.340967}%
\pgfsetfillcolor{currentfill}%
\pgfsetlinewidth{1.003750pt}%
\definecolor{currentstroke}{rgb}{0.440137,0.811138,0.340967}%
\pgfsetstrokecolor{currentstroke}%
\pgfsetdash{}{0pt}%
\pgfpathmoveto{\pgfqpoint{4.608532in}{6.979318in}}%
\pgfpathcurveto{\pgfqpoint{4.614356in}{6.979318in}}{\pgfqpoint{4.619942in}{6.981632in}}{\pgfqpoint{4.624060in}{6.985750in}}%
\pgfpathcurveto{\pgfqpoint{4.628178in}{6.989868in}}{\pgfqpoint{4.630492in}{6.995455in}}{\pgfqpoint{4.630492in}{7.001278in}}%
\pgfpathcurveto{\pgfqpoint{4.630492in}{7.007102in}}{\pgfqpoint{4.628178in}{7.012689in}}{\pgfqpoint{4.624060in}{7.016807in}}%
\pgfpathcurveto{\pgfqpoint{4.619942in}{7.020925in}}{\pgfqpoint{4.614356in}{7.023239in}}{\pgfqpoint{4.608532in}{7.023239in}}%
\pgfpathcurveto{\pgfqpoint{4.602708in}{7.023239in}}{\pgfqpoint{4.597122in}{7.020925in}}{\pgfqpoint{4.593003in}{7.016807in}}%
\pgfpathcurveto{\pgfqpoint{4.588885in}{7.012689in}}{\pgfqpoint{4.586571in}{7.007102in}}{\pgfqpoint{4.586571in}{7.001278in}}%
\pgfpathcurveto{\pgfqpoint{4.586571in}{6.995455in}}{\pgfqpoint{4.588885in}{6.989868in}}{\pgfqpoint{4.593003in}{6.985750in}}%
\pgfpathcurveto{\pgfqpoint{4.597122in}{6.981632in}}{\pgfqpoint{4.602708in}{6.979318in}}{\pgfqpoint{4.608532in}{6.979318in}}%
\pgfpathlineto{\pgfqpoint{4.608532in}{6.979318in}}%
\pgfpathclose%
\pgfusepath{stroke,fill}%
\end{pgfscope}%
\begin{pgfscope}%
\pgfpathrectangle{\pgfqpoint{1.414736in}{0.548769in}}{\pgfqpoint{3.497747in}{7.253186in}}%
\pgfusepath{clip}%
\pgfsetbuttcap%
\pgfsetroundjoin%
\definecolor{currentfill}{rgb}{0.866013,0.889868,0.095953}%
\pgfsetfillcolor{currentfill}%
\pgfsetlinewidth{1.003750pt}%
\definecolor{currentstroke}{rgb}{0.866013,0.889868,0.095953}%
\pgfsetstrokecolor{currentstroke}%
\pgfsetdash{}{0pt}%
\pgfpathmoveto{\pgfqpoint{3.697130in}{7.450304in}}%
\pgfpathcurveto{\pgfqpoint{3.702954in}{7.450304in}}{\pgfqpoint{3.708540in}{7.452618in}}{\pgfqpoint{3.712658in}{7.456736in}}%
\pgfpathcurveto{\pgfqpoint{3.716777in}{7.460854in}}{\pgfqpoint{3.719090in}{7.466441in}}{\pgfqpoint{3.719090in}{7.472265in}}%
\pgfpathcurveto{\pgfqpoint{3.719090in}{7.478089in}}{\pgfqpoint{3.716777in}{7.483675in}}{\pgfqpoint{3.712658in}{7.487793in}}%
\pgfpathcurveto{\pgfqpoint{3.708540in}{7.491911in}}{\pgfqpoint{3.702954in}{7.494225in}}{\pgfqpoint{3.697130in}{7.494225in}}%
\pgfpathcurveto{\pgfqpoint{3.691306in}{7.494225in}}{\pgfqpoint{3.685720in}{7.491911in}}{\pgfqpoint{3.681602in}{7.487793in}}%
\pgfpathcurveto{\pgfqpoint{3.677484in}{7.483675in}}{\pgfqpoint{3.675170in}{7.478089in}}{\pgfqpoint{3.675170in}{7.472265in}}%
\pgfpathcurveto{\pgfqpoint{3.675170in}{7.466441in}}{\pgfqpoint{3.677484in}{7.460854in}}{\pgfqpoint{3.681602in}{7.456736in}}%
\pgfpathcurveto{\pgfqpoint{3.685720in}{7.452618in}}{\pgfqpoint{3.691306in}{7.450304in}}{\pgfqpoint{3.697130in}{7.450304in}}%
\pgfpathlineto{\pgfqpoint{3.697130in}{7.450304in}}%
\pgfpathclose%
\pgfusepath{stroke,fill}%
\end{pgfscope}%
\begin{pgfscope}%
\pgfpathrectangle{\pgfqpoint{1.414736in}{0.548769in}}{\pgfqpoint{3.497747in}{7.253186in}}%
\pgfusepath{clip}%
\pgfsetbuttcap%
\pgfsetroundjoin%
\definecolor{currentfill}{rgb}{0.993248,0.906157,0.143936}%
\pgfsetfillcolor{currentfill}%
\pgfsetlinewidth{1.003750pt}%
\definecolor{currentstroke}{rgb}{0.993248,0.906157,0.143936}%
\pgfsetstrokecolor{currentstroke}%
\pgfsetdash{}{0pt}%
\pgfpathmoveto{\pgfqpoint{3.444714in}{0.856499in}}%
\pgfpathcurveto{\pgfqpoint{3.450538in}{0.856499in}}{\pgfqpoint{3.456124in}{0.858813in}}{\pgfqpoint{3.460242in}{0.862931in}}%
\pgfpathcurveto{\pgfqpoint{3.464360in}{0.867049in}}{\pgfqpoint{3.466674in}{0.872635in}}{\pgfqpoint{3.466674in}{0.878459in}}%
\pgfpathcurveto{\pgfqpoint{3.466674in}{0.884283in}}{\pgfqpoint{3.464360in}{0.889869in}}{\pgfqpoint{3.460242in}{0.893987in}}%
\pgfpathcurveto{\pgfqpoint{3.456124in}{0.898105in}}{\pgfqpoint{3.450538in}{0.900419in}}{\pgfqpoint{3.444714in}{0.900419in}}%
\pgfpathcurveto{\pgfqpoint{3.438890in}{0.900419in}}{\pgfqpoint{3.433304in}{0.898105in}}{\pgfqpoint{3.429186in}{0.893987in}}%
\pgfpathcurveto{\pgfqpoint{3.425068in}{0.889869in}}{\pgfqpoint{3.422754in}{0.884283in}}{\pgfqpoint{3.422754in}{0.878459in}}%
\pgfpathcurveto{\pgfqpoint{3.422754in}{0.872635in}}{\pgfqpoint{3.425068in}{0.867049in}}{\pgfqpoint{3.429186in}{0.862931in}}%
\pgfpathcurveto{\pgfqpoint{3.433304in}{0.858813in}}{\pgfqpoint{3.438890in}{0.856499in}}{\pgfqpoint{3.444714in}{0.856499in}}%
\pgfpathlineto{\pgfqpoint{3.444714in}{0.856499in}}%
\pgfpathclose%
\pgfusepath{stroke,fill}%
\end{pgfscope}%
\begin{pgfscope}%
\pgfpathrectangle{\pgfqpoint{1.414736in}{0.548769in}}{\pgfqpoint{3.497747in}{7.253186in}}%
\pgfusepath{clip}%
\pgfsetbuttcap%
\pgfsetroundjoin%
\definecolor{currentfill}{rgb}{0.555484,0.840254,0.269281}%
\pgfsetfillcolor{currentfill}%
\pgfsetlinewidth{1.003750pt}%
\definecolor{currentstroke}{rgb}{0.555484,0.840254,0.269281}%
\pgfsetstrokecolor{currentstroke}%
\pgfsetdash{}{0pt}%
\pgfpathmoveto{\pgfqpoint{3.636549in}{1.327485in}}%
\pgfpathcurveto{\pgfqpoint{3.642373in}{1.327485in}}{\pgfqpoint{3.647959in}{1.329799in}}{\pgfqpoint{3.652077in}{1.333917in}}%
\pgfpathcurveto{\pgfqpoint{3.656195in}{1.338035in}}{\pgfqpoint{3.658509in}{1.343621in}}{\pgfqpoint{3.658509in}{1.349445in}}%
\pgfpathcurveto{\pgfqpoint{3.658509in}{1.355269in}}{\pgfqpoint{3.656195in}{1.360855in}}{\pgfqpoint{3.652077in}{1.364973in}}%
\pgfpathcurveto{\pgfqpoint{3.647959in}{1.369092in}}{\pgfqpoint{3.642373in}{1.371405in}}{\pgfqpoint{3.636549in}{1.371405in}}%
\pgfpathcurveto{\pgfqpoint{3.630725in}{1.371405in}}{\pgfqpoint{3.625139in}{1.369092in}}{\pgfqpoint{3.621021in}{1.364973in}}%
\pgfpathcurveto{\pgfqpoint{3.616903in}{1.360855in}}{\pgfqpoint{3.614589in}{1.355269in}}{\pgfqpoint{3.614589in}{1.349445in}}%
\pgfpathcurveto{\pgfqpoint{3.614589in}{1.343621in}}{\pgfqpoint{3.616903in}{1.338035in}}{\pgfqpoint{3.621021in}{1.333917in}}%
\pgfpathcurveto{\pgfqpoint{3.625139in}{1.329799in}}{\pgfqpoint{3.630725in}{1.327485in}}{\pgfqpoint{3.636549in}{1.327485in}}%
\pgfpathlineto{\pgfqpoint{3.636549in}{1.327485in}}%
\pgfpathclose%
\pgfusepath{stroke,fill}%
\end{pgfscope}%
\begin{pgfscope}%
\pgfpathrectangle{\pgfqpoint{1.414736in}{0.548769in}}{\pgfqpoint{3.497747in}{7.253186in}}%
\pgfusepath{clip}%
\pgfsetbuttcap%
\pgfsetroundjoin%
\definecolor{currentfill}{rgb}{0.699415,0.867117,0.175971}%
\pgfsetfillcolor{currentfill}%
\pgfsetlinewidth{1.003750pt}%
\definecolor{currentstroke}{rgb}{0.699415,0.867117,0.175971}%
\pgfsetstrokecolor{currentstroke}%
\pgfsetdash{}{0pt}%
\pgfpathmoveto{\pgfqpoint{3.759863in}{1.798471in}}%
\pgfpathcurveto{\pgfqpoint{3.765687in}{1.798471in}}{\pgfqpoint{3.771273in}{1.800785in}}{\pgfqpoint{3.775391in}{1.804903in}}%
\pgfpathcurveto{\pgfqpoint{3.779510in}{1.809021in}}{\pgfqpoint{3.781823in}{1.814607in}}{\pgfqpoint{3.781823in}{1.820431in}}%
\pgfpathcurveto{\pgfqpoint{3.781823in}{1.826255in}}{\pgfqpoint{3.779510in}{1.831841in}}{\pgfqpoint{3.775391in}{1.835960in}}%
\pgfpathcurveto{\pgfqpoint{3.771273in}{1.840078in}}{\pgfqpoint{3.765687in}{1.842392in}}{\pgfqpoint{3.759863in}{1.842392in}}%
\pgfpathcurveto{\pgfqpoint{3.754039in}{1.842392in}}{\pgfqpoint{3.748453in}{1.840078in}}{\pgfqpoint{3.744335in}{1.835960in}}%
\pgfpathcurveto{\pgfqpoint{3.740217in}{1.831841in}}{\pgfqpoint{3.737903in}{1.826255in}}{\pgfqpoint{3.737903in}{1.820431in}}%
\pgfpathcurveto{\pgfqpoint{3.737903in}{1.814607in}}{\pgfqpoint{3.740217in}{1.809021in}}{\pgfqpoint{3.744335in}{1.804903in}}%
\pgfpathcurveto{\pgfqpoint{3.748453in}{1.800785in}}{\pgfqpoint{3.754039in}{1.798471in}}{\pgfqpoint{3.759863in}{1.798471in}}%
\pgfpathlineto{\pgfqpoint{3.759863in}{1.798471in}}%
\pgfpathclose%
\pgfusepath{stroke,fill}%
\end{pgfscope}%
\begin{pgfscope}%
\pgfpathrectangle{\pgfqpoint{1.414736in}{0.548769in}}{\pgfqpoint{3.497747in}{7.253186in}}%
\pgfusepath{clip}%
\pgfsetbuttcap%
\pgfsetroundjoin%
\definecolor{currentfill}{rgb}{0.993248,0.906157,0.143936}%
\pgfsetfillcolor{currentfill}%
\pgfsetlinewidth{1.003750pt}%
\definecolor{currentstroke}{rgb}{0.993248,0.906157,0.143936}%
\pgfsetstrokecolor{currentstroke}%
\pgfsetdash{}{0pt}%
\pgfpathmoveto{\pgfqpoint{3.856335in}{2.269457in}}%
\pgfpathcurveto{\pgfqpoint{3.862159in}{2.269457in}}{\pgfqpoint{3.867745in}{2.271771in}}{\pgfqpoint{3.871864in}{2.275889in}}%
\pgfpathcurveto{\pgfqpoint{3.875982in}{2.280007in}}{\pgfqpoint{3.878296in}{2.285593in}}{\pgfqpoint{3.878296in}{2.291417in}}%
\pgfpathcurveto{\pgfqpoint{3.878296in}{2.297241in}}{\pgfqpoint{3.875982in}{2.302828in}}{\pgfqpoint{3.871864in}{2.306946in}}%
\pgfpathcurveto{\pgfqpoint{3.867745in}{2.311064in}}{\pgfqpoint{3.862159in}{2.313378in}}{\pgfqpoint{3.856335in}{2.313378in}}%
\pgfpathcurveto{\pgfqpoint{3.850511in}{2.313378in}}{\pgfqpoint{3.844925in}{2.311064in}}{\pgfqpoint{3.840807in}{2.306946in}}%
\pgfpathcurveto{\pgfqpoint{3.836689in}{2.302828in}}{\pgfqpoint{3.834375in}{2.297241in}}{\pgfqpoint{3.834375in}{2.291417in}}%
\pgfpathcurveto{\pgfqpoint{3.834375in}{2.285593in}}{\pgfqpoint{3.836689in}{2.280007in}}{\pgfqpoint{3.840807in}{2.275889in}}%
\pgfpathcurveto{\pgfqpoint{3.844925in}{2.271771in}}{\pgfqpoint{3.850511in}{2.269457in}}{\pgfqpoint{3.856335in}{2.269457in}}%
\pgfpathlineto{\pgfqpoint{3.856335in}{2.269457in}}%
\pgfpathclose%
\pgfusepath{stroke,fill}%
\end{pgfscope}%
\begin{pgfscope}%
\pgfpathrectangle{\pgfqpoint{1.414736in}{0.548769in}}{\pgfqpoint{3.497747in}{7.253186in}}%
\pgfusepath{clip}%
\pgfsetbuttcap%
\pgfsetroundjoin%
\definecolor{currentfill}{rgb}{0.202219,0.715272,0.476084}%
\pgfsetfillcolor{currentfill}%
\pgfsetlinewidth{1.003750pt}%
\definecolor{currentstroke}{rgb}{0.202219,0.715272,0.476084}%
\pgfsetstrokecolor{currentstroke}%
\pgfsetdash{}{0pt}%
\pgfpathmoveto{\pgfqpoint{3.950451in}{2.740443in}}%
\pgfpathcurveto{\pgfqpoint{3.956275in}{2.740443in}}{\pgfqpoint{3.961861in}{2.742757in}}{\pgfqpoint{3.965979in}{2.746875in}}%
\pgfpathcurveto{\pgfqpoint{3.970097in}{2.750993in}}{\pgfqpoint{3.972411in}{2.756580in}}{\pgfqpoint{3.972411in}{2.762404in}}%
\pgfpathcurveto{\pgfqpoint{3.972411in}{2.768227in}}{\pgfqpoint{3.970097in}{2.773814in}}{\pgfqpoint{3.965979in}{2.777932in}}%
\pgfpathcurveto{\pgfqpoint{3.961861in}{2.782050in}}{\pgfqpoint{3.956275in}{2.784364in}}{\pgfqpoint{3.950451in}{2.784364in}}%
\pgfpathcurveto{\pgfqpoint{3.944627in}{2.784364in}}{\pgfqpoint{3.939041in}{2.782050in}}{\pgfqpoint{3.934923in}{2.777932in}}%
\pgfpathcurveto{\pgfqpoint{3.930804in}{2.773814in}}{\pgfqpoint{3.928491in}{2.768227in}}{\pgfqpoint{3.928491in}{2.762404in}}%
\pgfpathcurveto{\pgfqpoint{3.928491in}{2.756580in}}{\pgfqpoint{3.930804in}{2.750993in}}{\pgfqpoint{3.934923in}{2.746875in}}%
\pgfpathcurveto{\pgfqpoint{3.939041in}{2.742757in}}{\pgfqpoint{3.944627in}{2.740443in}}{\pgfqpoint{3.950451in}{2.740443in}}%
\pgfpathlineto{\pgfqpoint{3.950451in}{2.740443in}}%
\pgfpathclose%
\pgfusepath{stroke,fill}%
\end{pgfscope}%
\begin{pgfscope}%
\pgfpathrectangle{\pgfqpoint{1.414736in}{0.548769in}}{\pgfqpoint{3.497747in}{7.253186in}}%
\pgfusepath{clip}%
\pgfsetbuttcap%
\pgfsetroundjoin%
\definecolor{currentfill}{rgb}{0.606045,0.850733,0.236712}%
\pgfsetfillcolor{currentfill}%
\pgfsetlinewidth{1.003750pt}%
\definecolor{currentstroke}{rgb}{0.606045,0.850733,0.236712}%
\pgfsetstrokecolor{currentstroke}%
\pgfsetdash{}{0pt}%
\pgfpathmoveto{\pgfqpoint{3.484985in}{3.211429in}}%
\pgfpathcurveto{\pgfqpoint{3.490809in}{3.211429in}}{\pgfqpoint{3.496395in}{3.213743in}}{\pgfqpoint{3.500513in}{3.217861in}}%
\pgfpathcurveto{\pgfqpoint{3.504632in}{3.221980in}}{\pgfqpoint{3.506945in}{3.227566in}}{\pgfqpoint{3.506945in}{3.233390in}}%
\pgfpathcurveto{\pgfqpoint{3.506945in}{3.239214in}}{\pgfqpoint{3.504632in}{3.244800in}}{\pgfqpoint{3.500513in}{3.248918in}}%
\pgfpathcurveto{\pgfqpoint{3.496395in}{3.253036in}}{\pgfqpoint{3.490809in}{3.255350in}}{\pgfqpoint{3.484985in}{3.255350in}}%
\pgfpathcurveto{\pgfqpoint{3.479161in}{3.255350in}}{\pgfqpoint{3.473575in}{3.253036in}}{\pgfqpoint{3.469457in}{3.248918in}}%
\pgfpathcurveto{\pgfqpoint{3.465339in}{3.244800in}}{\pgfqpoint{3.463025in}{3.239214in}}{\pgfqpoint{3.463025in}{3.233390in}}%
\pgfpathcurveto{\pgfqpoint{3.463025in}{3.227566in}}{\pgfqpoint{3.465339in}{3.221980in}}{\pgfqpoint{3.469457in}{3.217861in}}%
\pgfpathcurveto{\pgfqpoint{3.473575in}{3.213743in}}{\pgfqpoint{3.479161in}{3.211429in}}{\pgfqpoint{3.484985in}{3.211429in}}%
\pgfpathlineto{\pgfqpoint{3.484985in}{3.211429in}}%
\pgfpathclose%
\pgfusepath{stroke,fill}%
\end{pgfscope}%
\begin{pgfscope}%
\pgfpathrectangle{\pgfqpoint{1.414736in}{0.548769in}}{\pgfqpoint{3.497747in}{7.253186in}}%
\pgfusepath{clip}%
\pgfsetbuttcap%
\pgfsetroundjoin%
\definecolor{currentfill}{rgb}{0.814576,0.883393,0.110347}%
\pgfsetfillcolor{currentfill}%
\pgfsetlinewidth{1.003750pt}%
\definecolor{currentstroke}{rgb}{0.814576,0.883393,0.110347}%
\pgfsetstrokecolor{currentstroke}%
\pgfsetdash{}{0pt}%
\pgfpathmoveto{\pgfqpoint{3.516957in}{3.682415in}}%
\pgfpathcurveto{\pgfqpoint{3.522781in}{3.682415in}}{\pgfqpoint{3.528367in}{3.684729in}}{\pgfqpoint{3.532486in}{3.688847in}}%
\pgfpathcurveto{\pgfqpoint{3.536604in}{3.692966in}}{\pgfqpoint{3.538918in}{3.698552in}}{\pgfqpoint{3.538918in}{3.704376in}}%
\pgfpathcurveto{\pgfqpoint{3.538918in}{3.710200in}}{\pgfqpoint{3.536604in}{3.715786in}}{\pgfqpoint{3.532486in}{3.719904in}}%
\pgfpathcurveto{\pgfqpoint{3.528367in}{3.724022in}}{\pgfqpoint{3.522781in}{3.726336in}}{\pgfqpoint{3.516957in}{3.726336in}}%
\pgfpathcurveto{\pgfqpoint{3.511133in}{3.726336in}}{\pgfqpoint{3.505547in}{3.724022in}}{\pgfqpoint{3.501429in}{3.719904in}}%
\pgfpathcurveto{\pgfqpoint{3.497311in}{3.715786in}}{\pgfqpoint{3.494997in}{3.710200in}}{\pgfqpoint{3.494997in}{3.704376in}}%
\pgfpathcurveto{\pgfqpoint{3.494997in}{3.698552in}}{\pgfqpoint{3.497311in}{3.692966in}}{\pgfqpoint{3.501429in}{3.688847in}}%
\pgfpathcurveto{\pgfqpoint{3.505547in}{3.684729in}}{\pgfqpoint{3.511133in}{3.682415in}}{\pgfqpoint{3.516957in}{3.682415in}}%
\pgfpathlineto{\pgfqpoint{3.516957in}{3.682415in}}%
\pgfpathclose%
\pgfusepath{stroke,fill}%
\end{pgfscope}%
\begin{pgfscope}%
\pgfpathrectangle{\pgfqpoint{1.414736in}{0.548769in}}{\pgfqpoint{3.497747in}{7.253186in}}%
\pgfusepath{clip}%
\pgfsetbuttcap%
\pgfsetroundjoin%
\definecolor{currentfill}{rgb}{0.595839,0.848717,0.243329}%
\pgfsetfillcolor{currentfill}%
\pgfsetlinewidth{1.003750pt}%
\definecolor{currentstroke}{rgb}{0.595839,0.848717,0.243329}%
\pgfsetstrokecolor{currentstroke}%
\pgfsetdash{}{0pt}%
\pgfpathmoveto{\pgfqpoint{3.489337in}{4.153402in}}%
\pgfpathcurveto{\pgfqpoint{3.495161in}{4.153402in}}{\pgfqpoint{3.500747in}{4.155715in}}{\pgfqpoint{3.504865in}{4.159834in}}%
\pgfpathcurveto{\pgfqpoint{3.508983in}{4.163952in}}{\pgfqpoint{3.511297in}{4.169538in}}{\pgfqpoint{3.511297in}{4.175362in}}%
\pgfpathcurveto{\pgfqpoint{3.511297in}{4.181186in}}{\pgfqpoint{3.508983in}{4.186772in}}{\pgfqpoint{3.504865in}{4.190890in}}%
\pgfpathcurveto{\pgfqpoint{3.500747in}{4.195008in}}{\pgfqpoint{3.495161in}{4.197322in}}{\pgfqpoint{3.489337in}{4.197322in}}%
\pgfpathcurveto{\pgfqpoint{3.483513in}{4.197322in}}{\pgfqpoint{3.477927in}{4.195008in}}{\pgfqpoint{3.473809in}{4.190890in}}%
\pgfpathcurveto{\pgfqpoint{3.469690in}{4.186772in}}{\pgfqpoint{3.467377in}{4.181186in}}{\pgfqpoint{3.467377in}{4.175362in}}%
\pgfpathcurveto{\pgfqpoint{3.467377in}{4.169538in}}{\pgfqpoint{3.469690in}{4.163952in}}{\pgfqpoint{3.473809in}{4.159834in}}%
\pgfpathcurveto{\pgfqpoint{3.477927in}{4.155715in}}{\pgfqpoint{3.483513in}{4.153402in}}{\pgfqpoint{3.489337in}{4.153402in}}%
\pgfpathlineto{\pgfqpoint{3.489337in}{4.153402in}}%
\pgfpathclose%
\pgfusepath{stroke,fill}%
\end{pgfscope}%
\begin{pgfscope}%
\pgfpathrectangle{\pgfqpoint{1.414736in}{0.548769in}}{\pgfqpoint{3.497747in}{7.253186in}}%
\pgfusepath{clip}%
\pgfsetbuttcap%
\pgfsetroundjoin%
\definecolor{currentfill}{rgb}{0.896320,0.893616,0.096335}%
\pgfsetfillcolor{currentfill}%
\pgfsetlinewidth{1.003750pt}%
\definecolor{currentstroke}{rgb}{0.896320,0.893616,0.096335}%
\pgfsetstrokecolor{currentstroke}%
\pgfsetdash{}{0pt}%
\pgfpathmoveto{\pgfqpoint{3.462924in}{4.624388in}}%
\pgfpathcurveto{\pgfqpoint{3.468748in}{4.624388in}}{\pgfqpoint{3.474334in}{4.626702in}}{\pgfqpoint{3.478452in}{4.630820in}}%
\pgfpathcurveto{\pgfqpoint{3.482570in}{4.634938in}}{\pgfqpoint{3.484884in}{4.640524in}}{\pgfqpoint{3.484884in}{4.646348in}}%
\pgfpathcurveto{\pgfqpoint{3.484884in}{4.652172in}}{\pgfqpoint{3.482570in}{4.657758in}}{\pgfqpoint{3.478452in}{4.661876in}}%
\pgfpathcurveto{\pgfqpoint{3.474334in}{4.665994in}}{\pgfqpoint{3.468748in}{4.668308in}}{\pgfqpoint{3.462924in}{4.668308in}}%
\pgfpathcurveto{\pgfqpoint{3.457100in}{4.668308in}}{\pgfqpoint{3.451514in}{4.665994in}}{\pgfqpoint{3.447396in}{4.661876in}}%
\pgfpathcurveto{\pgfqpoint{3.443277in}{4.657758in}}{\pgfqpoint{3.440964in}{4.652172in}}{\pgfqpoint{3.440964in}{4.646348in}}%
\pgfpathcurveto{\pgfqpoint{3.440964in}{4.640524in}}{\pgfqpoint{3.443277in}{4.634938in}}{\pgfqpoint{3.447396in}{4.630820in}}%
\pgfpathcurveto{\pgfqpoint{3.451514in}{4.626702in}}{\pgfqpoint{3.457100in}{4.624388in}}{\pgfqpoint{3.462924in}{4.624388in}}%
\pgfpathlineto{\pgfqpoint{3.462924in}{4.624388in}}%
\pgfpathclose%
\pgfusepath{stroke,fill}%
\end{pgfscope}%
\begin{pgfscope}%
\pgfpathrectangle{\pgfqpoint{1.414736in}{0.548769in}}{\pgfqpoint{3.497747in}{7.253186in}}%
\pgfusepath{clip}%
\pgfsetbuttcap%
\pgfsetroundjoin%
\definecolor{currentfill}{rgb}{0.122606,0.585371,0.546557}%
\pgfsetfillcolor{currentfill}%
\pgfsetlinewidth{1.003750pt}%
\definecolor{currentstroke}{rgb}{0.122606,0.585371,0.546557}%
\pgfsetstrokecolor{currentstroke}%
\pgfsetdash{}{0pt}%
\pgfpathmoveto{\pgfqpoint{3.948664in}{5.095374in}}%
\pgfpathcurveto{\pgfqpoint{3.954488in}{5.095374in}}{\pgfqpoint{3.960074in}{5.097688in}}{\pgfqpoint{3.964192in}{5.101806in}}%
\pgfpathcurveto{\pgfqpoint{3.968310in}{5.105924in}}{\pgfqpoint{3.970624in}{5.111510in}}{\pgfqpoint{3.970624in}{5.117334in}}%
\pgfpathcurveto{\pgfqpoint{3.970624in}{5.123158in}}{\pgfqpoint{3.968310in}{5.128744in}}{\pgfqpoint{3.964192in}{5.132862in}}%
\pgfpathcurveto{\pgfqpoint{3.960074in}{5.136980in}}{\pgfqpoint{3.954488in}{5.139294in}}{\pgfqpoint{3.948664in}{5.139294in}}%
\pgfpathcurveto{\pgfqpoint{3.942840in}{5.139294in}}{\pgfqpoint{3.937254in}{5.136980in}}{\pgfqpoint{3.933136in}{5.132862in}}%
\pgfpathcurveto{\pgfqpoint{3.929017in}{5.128744in}}{\pgfqpoint{3.926704in}{5.123158in}}{\pgfqpoint{3.926704in}{5.117334in}}%
\pgfpathcurveto{\pgfqpoint{3.926704in}{5.111510in}}{\pgfqpoint{3.929017in}{5.105924in}}{\pgfqpoint{3.933136in}{5.101806in}}%
\pgfpathcurveto{\pgfqpoint{3.937254in}{5.097688in}}{\pgfqpoint{3.942840in}{5.095374in}}{\pgfqpoint{3.948664in}{5.095374in}}%
\pgfpathlineto{\pgfqpoint{3.948664in}{5.095374in}}%
\pgfpathclose%
\pgfusepath{stroke,fill}%
\end{pgfscope}%
\begin{pgfscope}%
\pgfpathrectangle{\pgfqpoint{1.414736in}{0.548769in}}{\pgfqpoint{3.497747in}{7.253186in}}%
\pgfusepath{clip}%
\pgfsetbuttcap%
\pgfsetroundjoin%
\definecolor{currentfill}{rgb}{0.131172,0.555899,0.552459}%
\pgfsetfillcolor{currentfill}%
\pgfsetlinewidth{1.003750pt}%
\definecolor{currentstroke}{rgb}{0.131172,0.555899,0.552459}%
\pgfsetstrokecolor{currentstroke}%
\pgfsetdash{}{0pt}%
\pgfpathmoveto{\pgfqpoint{3.670159in}{5.566360in}}%
\pgfpathcurveto{\pgfqpoint{3.675983in}{5.566360in}}{\pgfqpoint{3.681570in}{5.568674in}}{\pgfqpoint{3.685688in}{5.572792in}}%
\pgfpathcurveto{\pgfqpoint{3.689806in}{5.576910in}}{\pgfqpoint{3.692120in}{5.582496in}}{\pgfqpoint{3.692120in}{5.588320in}}%
\pgfpathcurveto{\pgfqpoint{3.692120in}{5.594144in}}{\pgfqpoint{3.689806in}{5.599730in}}{\pgfqpoint{3.685688in}{5.603848in}}%
\pgfpathcurveto{\pgfqpoint{3.681570in}{5.607967in}}{\pgfqpoint{3.675983in}{5.610280in}}{\pgfqpoint{3.670159in}{5.610280in}}%
\pgfpathcurveto{\pgfqpoint{3.664335in}{5.610280in}}{\pgfqpoint{3.658749in}{5.607967in}}{\pgfqpoint{3.654631in}{5.603848in}}%
\pgfpathcurveto{\pgfqpoint{3.650513in}{5.599730in}}{\pgfqpoint{3.648199in}{5.594144in}}{\pgfqpoint{3.648199in}{5.588320in}}%
\pgfpathcurveto{\pgfqpoint{3.648199in}{5.582496in}}{\pgfqpoint{3.650513in}{5.576910in}}{\pgfqpoint{3.654631in}{5.572792in}}%
\pgfpathcurveto{\pgfqpoint{3.658749in}{5.568674in}}{\pgfqpoint{3.664335in}{5.566360in}}{\pgfqpoint{3.670159in}{5.566360in}}%
\pgfpathlineto{\pgfqpoint{3.670159in}{5.566360in}}%
\pgfpathclose%
\pgfusepath{stroke,fill}%
\end{pgfscope}%
\begin{pgfscope}%
\pgfpathrectangle{\pgfqpoint{1.414736in}{0.548769in}}{\pgfqpoint{3.497747in}{7.253186in}}%
\pgfusepath{clip}%
\pgfsetbuttcap%
\pgfsetroundjoin%
\definecolor{currentfill}{rgb}{0.252899,0.742211,0.448284}%
\pgfsetfillcolor{currentfill}%
\pgfsetlinewidth{1.003750pt}%
\definecolor{currentstroke}{rgb}{0.252899,0.742211,0.448284}%
\pgfsetstrokecolor{currentstroke}%
\pgfsetdash{}{0pt}%
\pgfpathmoveto{\pgfqpoint{3.569202in}{6.037346in}}%
\pgfpathcurveto{\pgfqpoint{3.575026in}{6.037346in}}{\pgfqpoint{3.580612in}{6.039660in}}{\pgfqpoint{3.584730in}{6.043778in}}%
\pgfpathcurveto{\pgfqpoint{3.588848in}{6.047896in}}{\pgfqpoint{3.591162in}{6.053482in}}{\pgfqpoint{3.591162in}{6.059306in}}%
\pgfpathcurveto{\pgfqpoint{3.591162in}{6.065130in}}{\pgfqpoint{3.588848in}{6.070716in}}{\pgfqpoint{3.584730in}{6.074835in}}%
\pgfpathcurveto{\pgfqpoint{3.580612in}{6.078953in}}{\pgfqpoint{3.575026in}{6.081267in}}{\pgfqpoint{3.569202in}{6.081267in}}%
\pgfpathcurveto{\pgfqpoint{3.563378in}{6.081267in}}{\pgfqpoint{3.557792in}{6.078953in}}{\pgfqpoint{3.553673in}{6.074835in}}%
\pgfpathcurveto{\pgfqpoint{3.549555in}{6.070716in}}{\pgfqpoint{3.547241in}{6.065130in}}{\pgfqpoint{3.547241in}{6.059306in}}%
\pgfpathcurveto{\pgfqpoint{3.547241in}{6.053482in}}{\pgfqpoint{3.549555in}{6.047896in}}{\pgfqpoint{3.553673in}{6.043778in}}%
\pgfpathcurveto{\pgfqpoint{3.557792in}{6.039660in}}{\pgfqpoint{3.563378in}{6.037346in}}{\pgfqpoint{3.569202in}{6.037346in}}%
\pgfpathlineto{\pgfqpoint{3.569202in}{6.037346in}}%
\pgfpathclose%
\pgfusepath{stroke,fill}%
\end{pgfscope}%
\begin{pgfscope}%
\pgfpathrectangle{\pgfqpoint{1.414736in}{0.548769in}}{\pgfqpoint{3.497747in}{7.253186in}}%
\pgfusepath{clip}%
\pgfsetbuttcap%
\pgfsetroundjoin%
\definecolor{currentfill}{rgb}{0.993248,0.906157,0.143936}%
\pgfsetfillcolor{currentfill}%
\pgfsetlinewidth{1.003750pt}%
\definecolor{currentstroke}{rgb}{0.993248,0.906157,0.143936}%
\pgfsetstrokecolor{currentstroke}%
\pgfsetdash{}{0pt}%
\pgfpathmoveto{\pgfqpoint{3.685499in}{6.508332in}}%
\pgfpathcurveto{\pgfqpoint{3.691323in}{6.508332in}}{\pgfqpoint{3.696909in}{6.510646in}}{\pgfqpoint{3.701027in}{6.514764in}}%
\pgfpathcurveto{\pgfqpoint{3.705145in}{6.518882in}}{\pgfqpoint{3.707459in}{6.524468in}}{\pgfqpoint{3.707459in}{6.530292in}}%
\pgfpathcurveto{\pgfqpoint{3.707459in}{6.536116in}}{\pgfqpoint{3.705145in}{6.541702in}}{\pgfqpoint{3.701027in}{6.545821in}}%
\pgfpathcurveto{\pgfqpoint{3.696909in}{6.549939in}}{\pgfqpoint{3.691323in}{6.552253in}}{\pgfqpoint{3.685499in}{6.552253in}}%
\pgfpathcurveto{\pgfqpoint{3.679675in}{6.552253in}}{\pgfqpoint{3.674089in}{6.549939in}}{\pgfqpoint{3.669971in}{6.545821in}}%
\pgfpathcurveto{\pgfqpoint{3.665852in}{6.541702in}}{\pgfqpoint{3.663539in}{6.536116in}}{\pgfqpoint{3.663539in}{6.530292in}}%
\pgfpathcurveto{\pgfqpoint{3.663539in}{6.524468in}}{\pgfqpoint{3.665852in}{6.518882in}}{\pgfqpoint{3.669971in}{6.514764in}}%
\pgfpathcurveto{\pgfqpoint{3.674089in}{6.510646in}}{\pgfqpoint{3.679675in}{6.508332in}}{\pgfqpoint{3.685499in}{6.508332in}}%
\pgfpathlineto{\pgfqpoint{3.685499in}{6.508332in}}%
\pgfpathclose%
\pgfusepath{stroke,fill}%
\end{pgfscope}%
\begin{pgfscope}%
\pgfpathrectangle{\pgfqpoint{1.414736in}{0.548769in}}{\pgfqpoint{3.497747in}{7.253186in}}%
\pgfusepath{clip}%
\pgfsetbuttcap%
\pgfsetroundjoin%
\definecolor{currentfill}{rgb}{0.983868,0.904867,0.136897}%
\pgfsetfillcolor{currentfill}%
\pgfsetlinewidth{1.003750pt}%
\definecolor{currentstroke}{rgb}{0.983868,0.904867,0.136897}%
\pgfsetstrokecolor{currentstroke}%
\pgfsetdash{}{0pt}%
\pgfpathmoveto{\pgfqpoint{3.656517in}{6.979318in}}%
\pgfpathcurveto{\pgfqpoint{3.662341in}{6.979318in}}{\pgfqpoint{3.667928in}{6.981632in}}{\pgfqpoint{3.672046in}{6.985750in}}%
\pgfpathcurveto{\pgfqpoint{3.676164in}{6.989868in}}{\pgfqpoint{3.678478in}{6.995455in}}{\pgfqpoint{3.678478in}{7.001278in}}%
\pgfpathcurveto{\pgfqpoint{3.678478in}{7.007102in}}{\pgfqpoint{3.676164in}{7.012689in}}{\pgfqpoint{3.672046in}{7.016807in}}%
\pgfpathcurveto{\pgfqpoint{3.667928in}{7.020925in}}{\pgfqpoint{3.662341in}{7.023239in}}{\pgfqpoint{3.656517in}{7.023239in}}%
\pgfpathcurveto{\pgfqpoint{3.650694in}{7.023239in}}{\pgfqpoint{3.645107in}{7.020925in}}{\pgfqpoint{3.640989in}{7.016807in}}%
\pgfpathcurveto{\pgfqpoint{3.636871in}{7.012689in}}{\pgfqpoint{3.634557in}{7.007102in}}{\pgfqpoint{3.634557in}{7.001278in}}%
\pgfpathcurveto{\pgfqpoint{3.634557in}{6.995455in}}{\pgfqpoint{3.636871in}{6.989868in}}{\pgfqpoint{3.640989in}{6.985750in}}%
\pgfpathcurveto{\pgfqpoint{3.645107in}{6.981632in}}{\pgfqpoint{3.650694in}{6.979318in}}{\pgfqpoint{3.656517in}{6.979318in}}%
\pgfpathlineto{\pgfqpoint{3.656517in}{6.979318in}}%
\pgfpathclose%
\pgfusepath{stroke,fill}%
\end{pgfscope}%
\begin{pgfscope}%
\pgfpathrectangle{\pgfqpoint{1.414736in}{0.548769in}}{\pgfqpoint{3.497747in}{7.253186in}}%
\pgfusepath{clip}%
\pgfsetbuttcap%
\pgfsetroundjoin%
\definecolor{currentfill}{rgb}{0.487026,0.823929,0.312321}%
\pgfsetfillcolor{currentfill}%
\pgfsetlinewidth{1.003750pt}%
\definecolor{currentstroke}{rgb}{0.487026,0.823929,0.312321}%
\pgfsetstrokecolor{currentstroke}%
\pgfsetdash{}{0pt}%
\pgfpathmoveto{\pgfqpoint{3.710156in}{7.450304in}}%
\pgfpathcurveto{\pgfqpoint{3.715980in}{7.450304in}}{\pgfqpoint{3.721566in}{7.452618in}}{\pgfqpoint{3.725685in}{7.456736in}}%
\pgfpathcurveto{\pgfqpoint{3.729803in}{7.460854in}}{\pgfqpoint{3.732117in}{7.466441in}}{\pgfqpoint{3.732117in}{7.472265in}}%
\pgfpathcurveto{\pgfqpoint{3.732117in}{7.478089in}}{\pgfqpoint{3.729803in}{7.483675in}}{\pgfqpoint{3.725685in}{7.487793in}}%
\pgfpathcurveto{\pgfqpoint{3.721566in}{7.491911in}}{\pgfqpoint{3.715980in}{7.494225in}}{\pgfqpoint{3.710156in}{7.494225in}}%
\pgfpathcurveto{\pgfqpoint{3.704332in}{7.494225in}}{\pgfqpoint{3.698746in}{7.491911in}}{\pgfqpoint{3.694628in}{7.487793in}}%
\pgfpathcurveto{\pgfqpoint{3.690510in}{7.483675in}}{\pgfqpoint{3.688196in}{7.478089in}}{\pgfqpoint{3.688196in}{7.472265in}}%
\pgfpathcurveto{\pgfqpoint{3.688196in}{7.466441in}}{\pgfqpoint{3.690510in}{7.460854in}}{\pgfqpoint{3.694628in}{7.456736in}}%
\pgfpathcurveto{\pgfqpoint{3.698746in}{7.452618in}}{\pgfqpoint{3.704332in}{7.450304in}}{\pgfqpoint{3.710156in}{7.450304in}}%
\pgfpathlineto{\pgfqpoint{3.710156in}{7.450304in}}%
\pgfpathclose%
\pgfusepath{stroke,fill}%
\end{pgfscope}%
\begin{pgfscope}%
\pgfpathrectangle{\pgfqpoint{1.414736in}{0.548769in}}{\pgfqpoint{3.497747in}{7.253186in}}%
\pgfusepath{clip}%
\pgfsetbuttcap%
\pgfsetroundjoin%
\definecolor{currentfill}{rgb}{0.945636,0.899815,0.112838}%
\pgfsetfillcolor{currentfill}%
\pgfsetlinewidth{1.003750pt}%
\definecolor{currentstroke}{rgb}{0.945636,0.899815,0.112838}%
\pgfsetstrokecolor{currentstroke}%
\pgfsetdash{}{0pt}%
\pgfpathmoveto{\pgfqpoint{3.502481in}{0.856499in}}%
\pgfpathcurveto{\pgfqpoint{3.508305in}{0.856499in}}{\pgfqpoint{3.513891in}{0.858813in}}{\pgfqpoint{3.518009in}{0.862931in}}%
\pgfpathcurveto{\pgfqpoint{3.522128in}{0.867049in}}{\pgfqpoint{3.524441in}{0.872635in}}{\pgfqpoint{3.524441in}{0.878459in}}%
\pgfpathcurveto{\pgfqpoint{3.524441in}{0.884283in}}{\pgfqpoint{3.522128in}{0.889869in}}{\pgfqpoint{3.518009in}{0.893987in}}%
\pgfpathcurveto{\pgfqpoint{3.513891in}{0.898105in}}{\pgfqpoint{3.508305in}{0.900419in}}{\pgfqpoint{3.502481in}{0.900419in}}%
\pgfpathcurveto{\pgfqpoint{3.496657in}{0.900419in}}{\pgfqpoint{3.491071in}{0.898105in}}{\pgfqpoint{3.486953in}{0.893987in}}%
\pgfpathcurveto{\pgfqpoint{3.482835in}{0.889869in}}{\pgfqpoint{3.480521in}{0.884283in}}{\pgfqpoint{3.480521in}{0.878459in}}%
\pgfpathcurveto{\pgfqpoint{3.480521in}{0.872635in}}{\pgfqpoint{3.482835in}{0.867049in}}{\pgfqpoint{3.486953in}{0.862931in}}%
\pgfpathcurveto{\pgfqpoint{3.491071in}{0.858813in}}{\pgfqpoint{3.496657in}{0.856499in}}{\pgfqpoint{3.502481in}{0.856499in}}%
\pgfpathlineto{\pgfqpoint{3.502481in}{0.856499in}}%
\pgfpathclose%
\pgfusepath{stroke,fill}%
\end{pgfscope}%
\begin{pgfscope}%
\pgfpathrectangle{\pgfqpoint{1.414736in}{0.548769in}}{\pgfqpoint{3.497747in}{7.253186in}}%
\pgfusepath{clip}%
\pgfsetbuttcap%
\pgfsetroundjoin%
\definecolor{currentfill}{rgb}{0.369214,0.788888,0.382914}%
\pgfsetfillcolor{currentfill}%
\pgfsetlinewidth{1.003750pt}%
\definecolor{currentstroke}{rgb}{0.369214,0.788888,0.382914}%
\pgfsetstrokecolor{currentstroke}%
\pgfsetdash{}{0pt}%
\pgfpathmoveto{\pgfqpoint{3.727204in}{1.327485in}}%
\pgfpathcurveto{\pgfqpoint{3.733027in}{1.327485in}}{\pgfqpoint{3.738614in}{1.329799in}}{\pgfqpoint{3.742732in}{1.333917in}}%
\pgfpathcurveto{\pgfqpoint{3.746850in}{1.338035in}}{\pgfqpoint{3.749164in}{1.343621in}}{\pgfqpoint{3.749164in}{1.349445in}}%
\pgfpathcurveto{\pgfqpoint{3.749164in}{1.355269in}}{\pgfqpoint{3.746850in}{1.360855in}}{\pgfqpoint{3.742732in}{1.364973in}}%
\pgfpathcurveto{\pgfqpoint{3.738614in}{1.369092in}}{\pgfqpoint{3.733027in}{1.371405in}}{\pgfqpoint{3.727204in}{1.371405in}}%
\pgfpathcurveto{\pgfqpoint{3.721380in}{1.371405in}}{\pgfqpoint{3.715793in}{1.369092in}}{\pgfqpoint{3.711675in}{1.364973in}}%
\pgfpathcurveto{\pgfqpoint{3.707557in}{1.360855in}}{\pgfqpoint{3.705243in}{1.355269in}}{\pgfqpoint{3.705243in}{1.349445in}}%
\pgfpathcurveto{\pgfqpoint{3.705243in}{1.343621in}}{\pgfqpoint{3.707557in}{1.338035in}}{\pgfqpoint{3.711675in}{1.333917in}}%
\pgfpathcurveto{\pgfqpoint{3.715793in}{1.329799in}}{\pgfqpoint{3.721380in}{1.327485in}}{\pgfqpoint{3.727204in}{1.327485in}}%
\pgfpathlineto{\pgfqpoint{3.727204in}{1.327485in}}%
\pgfpathclose%
\pgfusepath{stroke,fill}%
\end{pgfscope}%
\begin{pgfscope}%
\pgfpathrectangle{\pgfqpoint{1.414736in}{0.548769in}}{\pgfqpoint{3.497747in}{7.253186in}}%
\pgfusepath{clip}%
\pgfsetbuttcap%
\pgfsetroundjoin%
\definecolor{currentfill}{rgb}{0.699415,0.867117,0.175971}%
\pgfsetfillcolor{currentfill}%
\pgfsetlinewidth{1.003750pt}%
\definecolor{currentstroke}{rgb}{0.699415,0.867117,0.175971}%
\pgfsetstrokecolor{currentstroke}%
\pgfsetdash{}{0pt}%
\pgfpathmoveto{\pgfqpoint{3.675167in}{1.798471in}}%
\pgfpathcurveto{\pgfqpoint{3.680991in}{1.798471in}}{\pgfqpoint{3.686577in}{1.800785in}}{\pgfqpoint{3.690695in}{1.804903in}}%
\pgfpathcurveto{\pgfqpoint{3.694813in}{1.809021in}}{\pgfqpoint{3.697127in}{1.814607in}}{\pgfqpoint{3.697127in}{1.820431in}}%
\pgfpathcurveto{\pgfqpoint{3.697127in}{1.826255in}}{\pgfqpoint{3.694813in}{1.831841in}}{\pgfqpoint{3.690695in}{1.835960in}}%
\pgfpathcurveto{\pgfqpoint{3.686577in}{1.840078in}}{\pgfqpoint{3.680991in}{1.842392in}}{\pgfqpoint{3.675167in}{1.842392in}}%
\pgfpathcurveto{\pgfqpoint{3.669343in}{1.842392in}}{\pgfqpoint{3.663757in}{1.840078in}}{\pgfqpoint{3.659639in}{1.835960in}}%
\pgfpathcurveto{\pgfqpoint{3.655520in}{1.831841in}}{\pgfqpoint{3.653207in}{1.826255in}}{\pgfqpoint{3.653207in}{1.820431in}}%
\pgfpathcurveto{\pgfqpoint{3.653207in}{1.814607in}}{\pgfqpoint{3.655520in}{1.809021in}}{\pgfqpoint{3.659639in}{1.804903in}}%
\pgfpathcurveto{\pgfqpoint{3.663757in}{1.800785in}}{\pgfqpoint{3.669343in}{1.798471in}}{\pgfqpoint{3.675167in}{1.798471in}}%
\pgfpathlineto{\pgfqpoint{3.675167in}{1.798471in}}%
\pgfpathclose%
\pgfusepath{stroke,fill}%
\end{pgfscope}%
\begin{pgfscope}%
\pgfpathrectangle{\pgfqpoint{1.414736in}{0.548769in}}{\pgfqpoint{3.497747in}{7.253186in}}%
\pgfusepath{clip}%
\pgfsetbuttcap%
\pgfsetroundjoin%
\definecolor{currentfill}{rgb}{0.804182,0.882046,0.114965}%
\pgfsetfillcolor{currentfill}%
\pgfsetlinewidth{1.003750pt}%
\definecolor{currentstroke}{rgb}{0.804182,0.882046,0.114965}%
\pgfsetstrokecolor{currentstroke}%
\pgfsetdash{}{0pt}%
\pgfpathmoveto{\pgfqpoint{3.399595in}{2.269457in}}%
\pgfpathcurveto{\pgfqpoint{3.405419in}{2.269457in}}{\pgfqpoint{3.411005in}{2.271771in}}{\pgfqpoint{3.415123in}{2.275889in}}%
\pgfpathcurveto{\pgfqpoint{3.419241in}{2.280007in}}{\pgfqpoint{3.421555in}{2.285593in}}{\pgfqpoint{3.421555in}{2.291417in}}%
\pgfpathcurveto{\pgfqpoint{3.421555in}{2.297241in}}{\pgfqpoint{3.419241in}{2.302828in}}{\pgfqpoint{3.415123in}{2.306946in}}%
\pgfpathcurveto{\pgfqpoint{3.411005in}{2.311064in}}{\pgfqpoint{3.405419in}{2.313378in}}{\pgfqpoint{3.399595in}{2.313378in}}%
\pgfpathcurveto{\pgfqpoint{3.393771in}{2.313378in}}{\pgfqpoint{3.388185in}{2.311064in}}{\pgfqpoint{3.384067in}{2.306946in}}%
\pgfpathcurveto{\pgfqpoint{3.379949in}{2.302828in}}{\pgfqpoint{3.377635in}{2.297241in}}{\pgfqpoint{3.377635in}{2.291417in}}%
\pgfpathcurveto{\pgfqpoint{3.377635in}{2.285593in}}{\pgfqpoint{3.379949in}{2.280007in}}{\pgfqpoint{3.384067in}{2.275889in}}%
\pgfpathcurveto{\pgfqpoint{3.388185in}{2.271771in}}{\pgfqpoint{3.393771in}{2.269457in}}{\pgfqpoint{3.399595in}{2.269457in}}%
\pgfpathlineto{\pgfqpoint{3.399595in}{2.269457in}}%
\pgfpathclose%
\pgfusepath{stroke,fill}%
\end{pgfscope}%
\begin{pgfscope}%
\pgfpathrectangle{\pgfqpoint{1.414736in}{0.548769in}}{\pgfqpoint{3.497747in}{7.253186in}}%
\pgfusepath{clip}%
\pgfsetbuttcap%
\pgfsetroundjoin%
\definecolor{currentfill}{rgb}{0.252899,0.742211,0.448284}%
\pgfsetfillcolor{currentfill}%
\pgfsetlinewidth{1.003750pt}%
\definecolor{currentstroke}{rgb}{0.252899,0.742211,0.448284}%
\pgfsetstrokecolor{currentstroke}%
\pgfsetdash{}{0pt}%
\pgfpathmoveto{\pgfqpoint{3.792572in}{2.740443in}}%
\pgfpathcurveto{\pgfqpoint{3.798396in}{2.740443in}}{\pgfqpoint{3.803982in}{2.742757in}}{\pgfqpoint{3.808100in}{2.746875in}}%
\pgfpathcurveto{\pgfqpoint{3.812218in}{2.750993in}}{\pgfqpoint{3.814532in}{2.756580in}}{\pgfqpoint{3.814532in}{2.762404in}}%
\pgfpathcurveto{\pgfqpoint{3.814532in}{2.768227in}}{\pgfqpoint{3.812218in}{2.773814in}}{\pgfqpoint{3.808100in}{2.777932in}}%
\pgfpathcurveto{\pgfqpoint{3.803982in}{2.782050in}}{\pgfqpoint{3.798396in}{2.784364in}}{\pgfqpoint{3.792572in}{2.784364in}}%
\pgfpathcurveto{\pgfqpoint{3.786748in}{2.784364in}}{\pgfqpoint{3.781162in}{2.782050in}}{\pgfqpoint{3.777044in}{2.777932in}}%
\pgfpathcurveto{\pgfqpoint{3.772926in}{2.773814in}}{\pgfqpoint{3.770612in}{2.768227in}}{\pgfqpoint{3.770612in}{2.762404in}}%
\pgfpathcurveto{\pgfqpoint{3.770612in}{2.756580in}}{\pgfqpoint{3.772926in}{2.750993in}}{\pgfqpoint{3.777044in}{2.746875in}}%
\pgfpathcurveto{\pgfqpoint{3.781162in}{2.742757in}}{\pgfqpoint{3.786748in}{2.740443in}}{\pgfqpoint{3.792572in}{2.740443in}}%
\pgfpathlineto{\pgfqpoint{3.792572in}{2.740443in}}%
\pgfpathclose%
\pgfusepath{stroke,fill}%
\end{pgfscope}%
\begin{pgfscope}%
\pgfpathrectangle{\pgfqpoint{1.414736in}{0.548769in}}{\pgfqpoint{3.497747in}{7.253186in}}%
\pgfusepath{clip}%
\pgfsetbuttcap%
\pgfsetroundjoin%
\definecolor{currentfill}{rgb}{0.595839,0.848717,0.243329}%
\pgfsetfillcolor{currentfill}%
\pgfsetlinewidth{1.003750pt}%
\definecolor{currentstroke}{rgb}{0.595839,0.848717,0.243329}%
\pgfsetstrokecolor{currentstroke}%
\pgfsetdash{}{0pt}%
\pgfpathmoveto{\pgfqpoint{3.356853in}{3.211429in}}%
\pgfpathcurveto{\pgfqpoint{3.362676in}{3.211429in}}{\pgfqpoint{3.368263in}{3.213743in}}{\pgfqpoint{3.372381in}{3.217861in}}%
\pgfpathcurveto{\pgfqpoint{3.376499in}{3.221980in}}{\pgfqpoint{3.378813in}{3.227566in}}{\pgfqpoint{3.378813in}{3.233390in}}%
\pgfpathcurveto{\pgfqpoint{3.378813in}{3.239214in}}{\pgfqpoint{3.376499in}{3.244800in}}{\pgfqpoint{3.372381in}{3.248918in}}%
\pgfpathcurveto{\pgfqpoint{3.368263in}{3.253036in}}{\pgfqpoint{3.362676in}{3.255350in}}{\pgfqpoint{3.356853in}{3.255350in}}%
\pgfpathcurveto{\pgfqpoint{3.351029in}{3.255350in}}{\pgfqpoint{3.345442in}{3.253036in}}{\pgfqpoint{3.341324in}{3.248918in}}%
\pgfpathcurveto{\pgfqpoint{3.337206in}{3.244800in}}{\pgfqpoint{3.334892in}{3.239214in}}{\pgfqpoint{3.334892in}{3.233390in}}%
\pgfpathcurveto{\pgfqpoint{3.334892in}{3.227566in}}{\pgfqpoint{3.337206in}{3.221980in}}{\pgfqpoint{3.341324in}{3.217861in}}%
\pgfpathcurveto{\pgfqpoint{3.345442in}{3.213743in}}{\pgfqpoint{3.351029in}{3.211429in}}{\pgfqpoint{3.356853in}{3.211429in}}%
\pgfpathlineto{\pgfqpoint{3.356853in}{3.211429in}}%
\pgfpathclose%
\pgfusepath{stroke,fill}%
\end{pgfscope}%
\begin{pgfscope}%
\pgfpathrectangle{\pgfqpoint{1.414736in}{0.548769in}}{\pgfqpoint{3.497747in}{7.253186in}}%
\pgfusepath{clip}%
\pgfsetbuttcap%
\pgfsetroundjoin%
\definecolor{currentfill}{rgb}{0.824940,0.884720,0.106217}%
\pgfsetfillcolor{currentfill}%
\pgfsetlinewidth{1.003750pt}%
\definecolor{currentstroke}{rgb}{0.824940,0.884720,0.106217}%
\pgfsetstrokecolor{currentstroke}%
\pgfsetdash{}{0pt}%
\pgfpathmoveto{\pgfqpoint{3.771057in}{3.682415in}}%
\pgfpathcurveto{\pgfqpoint{3.776881in}{3.682415in}}{\pgfqpoint{3.782467in}{3.684729in}}{\pgfqpoint{3.786585in}{3.688847in}}%
\pgfpathcurveto{\pgfqpoint{3.790703in}{3.692966in}}{\pgfqpoint{3.793017in}{3.698552in}}{\pgfqpoint{3.793017in}{3.704376in}}%
\pgfpathcurveto{\pgfqpoint{3.793017in}{3.710200in}}{\pgfqpoint{3.790703in}{3.715786in}}{\pgfqpoint{3.786585in}{3.719904in}}%
\pgfpathcurveto{\pgfqpoint{3.782467in}{3.724022in}}{\pgfqpoint{3.776881in}{3.726336in}}{\pgfqpoint{3.771057in}{3.726336in}}%
\pgfpathcurveto{\pgfqpoint{3.765233in}{3.726336in}}{\pgfqpoint{3.759647in}{3.724022in}}{\pgfqpoint{3.755529in}{3.719904in}}%
\pgfpathcurveto{\pgfqpoint{3.751411in}{3.715786in}}{\pgfqpoint{3.749097in}{3.710200in}}{\pgfqpoint{3.749097in}{3.704376in}}%
\pgfpathcurveto{\pgfqpoint{3.749097in}{3.698552in}}{\pgfqpoint{3.751411in}{3.692966in}}{\pgfqpoint{3.755529in}{3.688847in}}%
\pgfpathcurveto{\pgfqpoint{3.759647in}{3.684729in}}{\pgfqpoint{3.765233in}{3.682415in}}{\pgfqpoint{3.771057in}{3.682415in}}%
\pgfpathlineto{\pgfqpoint{3.771057in}{3.682415in}}%
\pgfpathclose%
\pgfusepath{stroke,fill}%
\end{pgfscope}%
\begin{pgfscope}%
\pgfpathrectangle{\pgfqpoint{1.414736in}{0.548769in}}{\pgfqpoint{3.497747in}{7.253186in}}%
\pgfusepath{clip}%
\pgfsetbuttcap%
\pgfsetroundjoin%
\definecolor{currentfill}{rgb}{0.468053,0.818921,0.323998}%
\pgfsetfillcolor{currentfill}%
\pgfsetlinewidth{1.003750pt}%
\definecolor{currentstroke}{rgb}{0.468053,0.818921,0.323998}%
\pgfsetstrokecolor{currentstroke}%
\pgfsetdash{}{0pt}%
\pgfpathmoveto{\pgfqpoint{3.517177in}{4.153402in}}%
\pgfpathcurveto{\pgfqpoint{3.523001in}{4.153402in}}{\pgfqpoint{3.528587in}{4.155715in}}{\pgfqpoint{3.532705in}{4.159834in}}%
\pgfpathcurveto{\pgfqpoint{3.536823in}{4.163952in}}{\pgfqpoint{3.539137in}{4.169538in}}{\pgfqpoint{3.539137in}{4.175362in}}%
\pgfpathcurveto{\pgfqpoint{3.539137in}{4.181186in}}{\pgfqpoint{3.536823in}{4.186772in}}{\pgfqpoint{3.532705in}{4.190890in}}%
\pgfpathcurveto{\pgfqpoint{3.528587in}{4.195008in}}{\pgfqpoint{3.523001in}{4.197322in}}{\pgfqpoint{3.517177in}{4.197322in}}%
\pgfpathcurveto{\pgfqpoint{3.511353in}{4.197322in}}{\pgfqpoint{3.505767in}{4.195008in}}{\pgfqpoint{3.501649in}{4.190890in}}%
\pgfpathcurveto{\pgfqpoint{3.497531in}{4.186772in}}{\pgfqpoint{3.495217in}{4.181186in}}{\pgfqpoint{3.495217in}{4.175362in}}%
\pgfpathcurveto{\pgfqpoint{3.495217in}{4.169538in}}{\pgfqpoint{3.497531in}{4.163952in}}{\pgfqpoint{3.501649in}{4.159834in}}%
\pgfpathcurveto{\pgfqpoint{3.505767in}{4.155715in}}{\pgfqpoint{3.511353in}{4.153402in}}{\pgfqpoint{3.517177in}{4.153402in}}%
\pgfpathlineto{\pgfqpoint{3.517177in}{4.153402in}}%
\pgfpathclose%
\pgfusepath{stroke,fill}%
\end{pgfscope}%
\begin{pgfscope}%
\pgfpathrectangle{\pgfqpoint{1.414736in}{0.548769in}}{\pgfqpoint{3.497747in}{7.253186in}}%
\pgfusepath{clip}%
\pgfsetbuttcap%
\pgfsetroundjoin%
\definecolor{currentfill}{rgb}{0.678489,0.863742,0.189503}%
\pgfsetfillcolor{currentfill}%
\pgfsetlinewidth{1.003750pt}%
\definecolor{currentstroke}{rgb}{0.678489,0.863742,0.189503}%
\pgfsetstrokecolor{currentstroke}%
\pgfsetdash{}{0pt}%
\pgfpathmoveto{\pgfqpoint{3.589625in}{4.624388in}}%
\pgfpathcurveto{\pgfqpoint{3.595448in}{4.624388in}}{\pgfqpoint{3.601035in}{4.626702in}}{\pgfqpoint{3.605153in}{4.630820in}}%
\pgfpathcurveto{\pgfqpoint{3.609271in}{4.634938in}}{\pgfqpoint{3.611585in}{4.640524in}}{\pgfqpoint{3.611585in}{4.646348in}}%
\pgfpathcurveto{\pgfqpoint{3.611585in}{4.652172in}}{\pgfqpoint{3.609271in}{4.657758in}}{\pgfqpoint{3.605153in}{4.661876in}}%
\pgfpathcurveto{\pgfqpoint{3.601035in}{4.665994in}}{\pgfqpoint{3.595448in}{4.668308in}}{\pgfqpoint{3.589625in}{4.668308in}}%
\pgfpathcurveto{\pgfqpoint{3.583801in}{4.668308in}}{\pgfqpoint{3.578214in}{4.665994in}}{\pgfqpoint{3.574096in}{4.661876in}}%
\pgfpathcurveto{\pgfqpoint{3.569978in}{4.657758in}}{\pgfqpoint{3.567664in}{4.652172in}}{\pgfqpoint{3.567664in}{4.646348in}}%
\pgfpathcurveto{\pgfqpoint{3.567664in}{4.640524in}}{\pgfqpoint{3.569978in}{4.634938in}}{\pgfqpoint{3.574096in}{4.630820in}}%
\pgfpathcurveto{\pgfqpoint{3.578214in}{4.626702in}}{\pgfqpoint{3.583801in}{4.624388in}}{\pgfqpoint{3.589625in}{4.624388in}}%
\pgfpathlineto{\pgfqpoint{3.589625in}{4.624388in}}%
\pgfpathclose%
\pgfusepath{stroke,fill}%
\end{pgfscope}%
\begin{pgfscope}%
\pgfpathrectangle{\pgfqpoint{1.414736in}{0.548769in}}{\pgfqpoint{3.497747in}{7.253186in}}%
\pgfusepath{clip}%
\pgfsetbuttcap%
\pgfsetroundjoin%
\definecolor{currentfill}{rgb}{0.274952,0.037752,0.364543}%
\pgfsetfillcolor{currentfill}%
\pgfsetlinewidth{1.003750pt}%
\definecolor{currentstroke}{rgb}{0.274952,0.037752,0.364543}%
\pgfsetstrokecolor{currentstroke}%
\pgfsetdash{}{0pt}%
\pgfpathmoveto{\pgfqpoint{3.471220in}{5.095374in}}%
\pgfpathcurveto{\pgfqpoint{3.477044in}{5.095374in}}{\pgfqpoint{3.482631in}{5.097688in}}{\pgfqpoint{3.486749in}{5.101806in}}%
\pgfpathcurveto{\pgfqpoint{3.490867in}{5.105924in}}{\pgfqpoint{3.493181in}{5.111510in}}{\pgfqpoint{3.493181in}{5.117334in}}%
\pgfpathcurveto{\pgfqpoint{3.493181in}{5.123158in}}{\pgfqpoint{3.490867in}{5.128744in}}{\pgfqpoint{3.486749in}{5.132862in}}%
\pgfpathcurveto{\pgfqpoint{3.482631in}{5.136980in}}{\pgfqpoint{3.477044in}{5.139294in}}{\pgfqpoint{3.471220in}{5.139294in}}%
\pgfpathcurveto{\pgfqpoint{3.465397in}{5.139294in}}{\pgfqpoint{3.459810in}{5.136980in}}{\pgfqpoint{3.455692in}{5.132862in}}%
\pgfpathcurveto{\pgfqpoint{3.451574in}{5.128744in}}{\pgfqpoint{3.449260in}{5.123158in}}{\pgfqpoint{3.449260in}{5.117334in}}%
\pgfpathcurveto{\pgfqpoint{3.449260in}{5.111510in}}{\pgfqpoint{3.451574in}{5.105924in}}{\pgfqpoint{3.455692in}{5.101806in}}%
\pgfpathcurveto{\pgfqpoint{3.459810in}{5.097688in}}{\pgfqpoint{3.465397in}{5.095374in}}{\pgfqpoint{3.471220in}{5.095374in}}%
\pgfpathlineto{\pgfqpoint{3.471220in}{5.095374in}}%
\pgfpathclose%
\pgfusepath{stroke,fill}%
\end{pgfscope}%
\begin{pgfscope}%
\pgfpathrectangle{\pgfqpoint{1.414736in}{0.548769in}}{\pgfqpoint{3.497747in}{7.253186in}}%
\pgfusepath{clip}%
\pgfsetbuttcap%
\pgfsetroundjoin%
\definecolor{currentfill}{rgb}{0.241237,0.296485,0.539709}%
\pgfsetfillcolor{currentfill}%
\pgfsetlinewidth{1.003750pt}%
\definecolor{currentstroke}{rgb}{0.241237,0.296485,0.539709}%
\pgfsetstrokecolor{currentstroke}%
\pgfsetdash{}{0pt}%
\pgfpathmoveto{\pgfqpoint{3.808188in}{5.566360in}}%
\pgfpathcurveto{\pgfqpoint{3.814012in}{5.566360in}}{\pgfqpoint{3.819598in}{5.568674in}}{\pgfqpoint{3.823716in}{5.572792in}}%
\pgfpathcurveto{\pgfqpoint{3.827834in}{5.576910in}}{\pgfqpoint{3.830148in}{5.582496in}}{\pgfqpoint{3.830148in}{5.588320in}}%
\pgfpathcurveto{\pgfqpoint{3.830148in}{5.594144in}}{\pgfqpoint{3.827834in}{5.599730in}}{\pgfqpoint{3.823716in}{5.603848in}}%
\pgfpathcurveto{\pgfqpoint{3.819598in}{5.607967in}}{\pgfqpoint{3.814012in}{5.610280in}}{\pgfqpoint{3.808188in}{5.610280in}}%
\pgfpathcurveto{\pgfqpoint{3.802364in}{5.610280in}}{\pgfqpoint{3.796778in}{5.607967in}}{\pgfqpoint{3.792660in}{5.603848in}}%
\pgfpathcurveto{\pgfqpoint{3.788542in}{5.599730in}}{\pgfqpoint{3.786228in}{5.594144in}}{\pgfqpoint{3.786228in}{5.588320in}}%
\pgfpathcurveto{\pgfqpoint{3.786228in}{5.582496in}}{\pgfqpoint{3.788542in}{5.576910in}}{\pgfqpoint{3.792660in}{5.572792in}}%
\pgfpathcurveto{\pgfqpoint{3.796778in}{5.568674in}}{\pgfqpoint{3.802364in}{5.566360in}}{\pgfqpoint{3.808188in}{5.566360in}}%
\pgfpathlineto{\pgfqpoint{3.808188in}{5.566360in}}%
\pgfpathclose%
\pgfusepath{stroke,fill}%
\end{pgfscope}%
\begin{pgfscope}%
\pgfpathrectangle{\pgfqpoint{1.414736in}{0.548769in}}{\pgfqpoint{3.497747in}{7.253186in}}%
\pgfusepath{clip}%
\pgfsetbuttcap%
\pgfsetroundjoin%
\definecolor{currentfill}{rgb}{0.147607,0.511733,0.557049}%
\pgfsetfillcolor{currentfill}%
\pgfsetlinewidth{1.003750pt}%
\definecolor{currentstroke}{rgb}{0.147607,0.511733,0.557049}%
\pgfsetstrokecolor{currentstroke}%
\pgfsetdash{}{0pt}%
\pgfpathmoveto{\pgfqpoint{3.652121in}{6.037346in}}%
\pgfpathcurveto{\pgfqpoint{3.657945in}{6.037346in}}{\pgfqpoint{3.663531in}{6.039660in}}{\pgfqpoint{3.667649in}{6.043778in}}%
\pgfpathcurveto{\pgfqpoint{3.671767in}{6.047896in}}{\pgfqpoint{3.674081in}{6.053482in}}{\pgfqpoint{3.674081in}{6.059306in}}%
\pgfpathcurveto{\pgfqpoint{3.674081in}{6.065130in}}{\pgfqpoint{3.671767in}{6.070716in}}{\pgfqpoint{3.667649in}{6.074835in}}%
\pgfpathcurveto{\pgfqpoint{3.663531in}{6.078953in}}{\pgfqpoint{3.657945in}{6.081267in}}{\pgfqpoint{3.652121in}{6.081267in}}%
\pgfpathcurveto{\pgfqpoint{3.646297in}{6.081267in}}{\pgfqpoint{3.640711in}{6.078953in}}{\pgfqpoint{3.636592in}{6.074835in}}%
\pgfpathcurveto{\pgfqpoint{3.632474in}{6.070716in}}{\pgfqpoint{3.630160in}{6.065130in}}{\pgfqpoint{3.630160in}{6.059306in}}%
\pgfpathcurveto{\pgfqpoint{3.630160in}{6.053482in}}{\pgfqpoint{3.632474in}{6.047896in}}{\pgfqpoint{3.636592in}{6.043778in}}%
\pgfpathcurveto{\pgfqpoint{3.640711in}{6.039660in}}{\pgfqpoint{3.646297in}{6.037346in}}{\pgfqpoint{3.652121in}{6.037346in}}%
\pgfpathlineto{\pgfqpoint{3.652121in}{6.037346in}}%
\pgfpathclose%
\pgfusepath{stroke,fill}%
\end{pgfscope}%
\begin{pgfscope}%
\pgfpathrectangle{\pgfqpoint{1.414736in}{0.548769in}}{\pgfqpoint{3.497747in}{7.253186in}}%
\pgfusepath{clip}%
\pgfsetbuttcap%
\pgfsetroundjoin%
\definecolor{currentfill}{rgb}{0.814576,0.883393,0.110347}%
\pgfsetfillcolor{currentfill}%
\pgfsetlinewidth{1.003750pt}%
\definecolor{currentstroke}{rgb}{0.814576,0.883393,0.110347}%
\pgfsetstrokecolor{currentstroke}%
\pgfsetdash{}{0pt}%
\pgfpathmoveto{\pgfqpoint{3.846890in}{6.508332in}}%
\pgfpathcurveto{\pgfqpoint{3.852714in}{6.508332in}}{\pgfqpoint{3.858300in}{6.510646in}}{\pgfqpoint{3.862418in}{6.514764in}}%
\pgfpathcurveto{\pgfqpoint{3.866536in}{6.518882in}}{\pgfqpoint{3.868850in}{6.524468in}}{\pgfqpoint{3.868850in}{6.530292in}}%
\pgfpathcurveto{\pgfqpoint{3.868850in}{6.536116in}}{\pgfqpoint{3.866536in}{6.541702in}}{\pgfqpoint{3.862418in}{6.545821in}}%
\pgfpathcurveto{\pgfqpoint{3.858300in}{6.549939in}}{\pgfqpoint{3.852714in}{6.552253in}}{\pgfqpoint{3.846890in}{6.552253in}}%
\pgfpathcurveto{\pgfqpoint{3.841066in}{6.552253in}}{\pgfqpoint{3.835480in}{6.549939in}}{\pgfqpoint{3.831362in}{6.545821in}}%
\pgfpathcurveto{\pgfqpoint{3.827244in}{6.541702in}}{\pgfqpoint{3.824930in}{6.536116in}}{\pgfqpoint{3.824930in}{6.530292in}}%
\pgfpathcurveto{\pgfqpoint{3.824930in}{6.524468in}}{\pgfqpoint{3.827244in}{6.518882in}}{\pgfqpoint{3.831362in}{6.514764in}}%
\pgfpathcurveto{\pgfqpoint{3.835480in}{6.510646in}}{\pgfqpoint{3.841066in}{6.508332in}}{\pgfqpoint{3.846890in}{6.508332in}}%
\pgfpathlineto{\pgfqpoint{3.846890in}{6.508332in}}%
\pgfpathclose%
\pgfusepath{stroke,fill}%
\end{pgfscope}%
\begin{pgfscope}%
\pgfpathrectangle{\pgfqpoint{1.414736in}{0.548769in}}{\pgfqpoint{3.497747in}{7.253186in}}%
\pgfusepath{clip}%
\pgfsetbuttcap%
\pgfsetroundjoin%
\definecolor{currentfill}{rgb}{0.304148,0.764704,0.419943}%
\pgfsetfillcolor{currentfill}%
\pgfsetlinewidth{1.003750pt}%
\definecolor{currentstroke}{rgb}{0.304148,0.764704,0.419943}%
\pgfsetstrokecolor{currentstroke}%
\pgfsetdash{}{0pt}%
\pgfpathmoveto{\pgfqpoint{3.936751in}{6.979318in}}%
\pgfpathcurveto{\pgfqpoint{3.942575in}{6.979318in}}{\pgfqpoint{3.948162in}{6.981632in}}{\pgfqpoint{3.952280in}{6.985750in}}%
\pgfpathcurveto{\pgfqpoint{3.956398in}{6.989868in}}{\pgfqpoint{3.958712in}{6.995455in}}{\pgfqpoint{3.958712in}{7.001278in}}%
\pgfpathcurveto{\pgfqpoint{3.958712in}{7.007102in}}{\pgfqpoint{3.956398in}{7.012689in}}{\pgfqpoint{3.952280in}{7.016807in}}%
\pgfpathcurveto{\pgfqpoint{3.948162in}{7.020925in}}{\pgfqpoint{3.942575in}{7.023239in}}{\pgfqpoint{3.936751in}{7.023239in}}%
\pgfpathcurveto{\pgfqpoint{3.930928in}{7.023239in}}{\pgfqpoint{3.925341in}{7.020925in}}{\pgfqpoint{3.921223in}{7.016807in}}%
\pgfpathcurveto{\pgfqpoint{3.917105in}{7.012689in}}{\pgfqpoint{3.914791in}{7.007102in}}{\pgfqpoint{3.914791in}{7.001278in}}%
\pgfpathcurveto{\pgfqpoint{3.914791in}{6.995455in}}{\pgfqpoint{3.917105in}{6.989868in}}{\pgfqpoint{3.921223in}{6.985750in}}%
\pgfpathcurveto{\pgfqpoint{3.925341in}{6.981632in}}{\pgfqpoint{3.930928in}{6.979318in}}{\pgfqpoint{3.936751in}{6.979318in}}%
\pgfpathlineto{\pgfqpoint{3.936751in}{6.979318in}}%
\pgfpathclose%
\pgfusepath{stroke,fill}%
\end{pgfscope}%
\begin{pgfscope}%
\pgfpathrectangle{\pgfqpoint{1.414736in}{0.548769in}}{\pgfqpoint{3.497747in}{7.253186in}}%
\pgfusepath{clip}%
\pgfsetbuttcap%
\pgfsetroundjoin%
\definecolor{currentfill}{rgb}{0.208030,0.718701,0.472873}%
\pgfsetfillcolor{currentfill}%
\pgfsetlinewidth{1.003750pt}%
\definecolor{currentstroke}{rgb}{0.208030,0.718701,0.472873}%
\pgfsetstrokecolor{currentstroke}%
\pgfsetdash{}{0pt}%
\pgfpathmoveto{\pgfqpoint{3.815927in}{7.450304in}}%
\pgfpathcurveto{\pgfqpoint{3.821751in}{7.450304in}}{\pgfqpoint{3.827337in}{7.452618in}}{\pgfqpoint{3.831455in}{7.456736in}}%
\pgfpathcurveto{\pgfqpoint{3.835573in}{7.460854in}}{\pgfqpoint{3.837887in}{7.466441in}}{\pgfqpoint{3.837887in}{7.472265in}}%
\pgfpathcurveto{\pgfqpoint{3.837887in}{7.478089in}}{\pgfqpoint{3.835573in}{7.483675in}}{\pgfqpoint{3.831455in}{7.487793in}}%
\pgfpathcurveto{\pgfqpoint{3.827337in}{7.491911in}}{\pgfqpoint{3.821751in}{7.494225in}}{\pgfqpoint{3.815927in}{7.494225in}}%
\pgfpathcurveto{\pgfqpoint{3.810103in}{7.494225in}}{\pgfqpoint{3.804517in}{7.491911in}}{\pgfqpoint{3.800398in}{7.487793in}}%
\pgfpathcurveto{\pgfqpoint{3.796280in}{7.483675in}}{\pgfqpoint{3.793966in}{7.478089in}}{\pgfqpoint{3.793966in}{7.472265in}}%
\pgfpathcurveto{\pgfqpoint{3.793966in}{7.466441in}}{\pgfqpoint{3.796280in}{7.460854in}}{\pgfqpoint{3.800398in}{7.456736in}}%
\pgfpathcurveto{\pgfqpoint{3.804517in}{7.452618in}}{\pgfqpoint{3.810103in}{7.450304in}}{\pgfqpoint{3.815927in}{7.450304in}}%
\pgfpathlineto{\pgfqpoint{3.815927in}{7.450304in}}%
\pgfpathclose%
\pgfusepath{stroke,fill}%
\end{pgfscope}%
\begin{pgfscope}%
\pgfpathrectangle{\pgfqpoint{1.414736in}{0.548769in}}{\pgfqpoint{3.497747in}{7.253186in}}%
\pgfusepath{clip}%
\pgfsetbuttcap%
\pgfsetroundjoin%
\definecolor{currentfill}{rgb}{0.955300,0.901065,0.118128}%
\pgfsetfillcolor{currentfill}%
\pgfsetlinewidth{1.003750pt}%
\definecolor{currentstroke}{rgb}{0.955300,0.901065,0.118128}%
\pgfsetstrokecolor{currentstroke}%
\pgfsetdash{}{0pt}%
\pgfpathmoveto{\pgfqpoint{3.666353in}{0.856499in}}%
\pgfpathcurveto{\pgfqpoint{3.672177in}{0.856499in}}{\pgfqpoint{3.677763in}{0.858813in}}{\pgfqpoint{3.681881in}{0.862931in}}%
\pgfpathcurveto{\pgfqpoint{3.685999in}{0.867049in}}{\pgfqpoint{3.688313in}{0.872635in}}{\pgfqpoint{3.688313in}{0.878459in}}%
\pgfpathcurveto{\pgfqpoint{3.688313in}{0.884283in}}{\pgfqpoint{3.685999in}{0.889869in}}{\pgfqpoint{3.681881in}{0.893987in}}%
\pgfpathcurveto{\pgfqpoint{3.677763in}{0.898105in}}{\pgfqpoint{3.672177in}{0.900419in}}{\pgfqpoint{3.666353in}{0.900419in}}%
\pgfpathcurveto{\pgfqpoint{3.660529in}{0.900419in}}{\pgfqpoint{3.654942in}{0.898105in}}{\pgfqpoint{3.650824in}{0.893987in}}%
\pgfpathcurveto{\pgfqpoint{3.646706in}{0.889869in}}{\pgfqpoint{3.644392in}{0.884283in}}{\pgfqpoint{3.644392in}{0.878459in}}%
\pgfpathcurveto{\pgfqpoint{3.644392in}{0.872635in}}{\pgfqpoint{3.646706in}{0.867049in}}{\pgfqpoint{3.650824in}{0.862931in}}%
\pgfpathcurveto{\pgfqpoint{3.654942in}{0.858813in}}{\pgfqpoint{3.660529in}{0.856499in}}{\pgfqpoint{3.666353in}{0.856499in}}%
\pgfpathlineto{\pgfqpoint{3.666353in}{0.856499in}}%
\pgfpathclose%
\pgfusepath{stroke,fill}%
\end{pgfscope}%
\begin{pgfscope}%
\pgfpathrectangle{\pgfqpoint{1.414736in}{0.548769in}}{\pgfqpoint{3.497747in}{7.253186in}}%
\pgfusepath{clip}%
\pgfsetbuttcap%
\pgfsetroundjoin%
\definecolor{currentfill}{rgb}{0.762373,0.876424,0.137064}%
\pgfsetfillcolor{currentfill}%
\pgfsetlinewidth{1.003750pt}%
\definecolor{currentstroke}{rgb}{0.762373,0.876424,0.137064}%
\pgfsetstrokecolor{currentstroke}%
\pgfsetdash{}{0pt}%
\pgfpathmoveto{\pgfqpoint{3.389050in}{1.327485in}}%
\pgfpathcurveto{\pgfqpoint{3.394874in}{1.327485in}}{\pgfqpoint{3.400460in}{1.329799in}}{\pgfqpoint{3.404578in}{1.333917in}}%
\pgfpathcurveto{\pgfqpoint{3.408696in}{1.338035in}}{\pgfqpoint{3.411010in}{1.343621in}}{\pgfqpoint{3.411010in}{1.349445in}}%
\pgfpathcurveto{\pgfqpoint{3.411010in}{1.355269in}}{\pgfqpoint{3.408696in}{1.360855in}}{\pgfqpoint{3.404578in}{1.364973in}}%
\pgfpathcurveto{\pgfqpoint{3.400460in}{1.369092in}}{\pgfqpoint{3.394874in}{1.371405in}}{\pgfqpoint{3.389050in}{1.371405in}}%
\pgfpathcurveto{\pgfqpoint{3.383226in}{1.371405in}}{\pgfqpoint{3.377640in}{1.369092in}}{\pgfqpoint{3.373521in}{1.364973in}}%
\pgfpathcurveto{\pgfqpoint{3.369403in}{1.360855in}}{\pgfqpoint{3.367089in}{1.355269in}}{\pgfqpoint{3.367089in}{1.349445in}}%
\pgfpathcurveto{\pgfqpoint{3.367089in}{1.343621in}}{\pgfqpoint{3.369403in}{1.338035in}}{\pgfqpoint{3.373521in}{1.333917in}}%
\pgfpathcurveto{\pgfqpoint{3.377640in}{1.329799in}}{\pgfqpoint{3.383226in}{1.327485in}}{\pgfqpoint{3.389050in}{1.327485in}}%
\pgfpathlineto{\pgfqpoint{3.389050in}{1.327485in}}%
\pgfpathclose%
\pgfusepath{stroke,fill}%
\end{pgfscope}%
\begin{pgfscope}%
\pgfpathrectangle{\pgfqpoint{1.414736in}{0.548769in}}{\pgfqpoint{3.497747in}{7.253186in}}%
\pgfusepath{clip}%
\pgfsetbuttcap%
\pgfsetroundjoin%
\definecolor{currentfill}{rgb}{0.699415,0.867117,0.175971}%
\pgfsetfillcolor{currentfill}%
\pgfsetlinewidth{1.003750pt}%
\definecolor{currentstroke}{rgb}{0.699415,0.867117,0.175971}%
\pgfsetstrokecolor{currentstroke}%
\pgfsetdash{}{0pt}%
\pgfpathmoveto{\pgfqpoint{3.690025in}{1.798471in}}%
\pgfpathcurveto{\pgfqpoint{3.695849in}{1.798471in}}{\pgfqpoint{3.701436in}{1.800785in}}{\pgfqpoint{3.705554in}{1.804903in}}%
\pgfpathcurveto{\pgfqpoint{3.709672in}{1.809021in}}{\pgfqpoint{3.711986in}{1.814607in}}{\pgfqpoint{3.711986in}{1.820431in}}%
\pgfpathcurveto{\pgfqpoint{3.711986in}{1.826255in}}{\pgfqpoint{3.709672in}{1.831841in}}{\pgfqpoint{3.705554in}{1.835960in}}%
\pgfpathcurveto{\pgfqpoint{3.701436in}{1.840078in}}{\pgfqpoint{3.695849in}{1.842392in}}{\pgfqpoint{3.690025in}{1.842392in}}%
\pgfpathcurveto{\pgfqpoint{3.684202in}{1.842392in}}{\pgfqpoint{3.678615in}{1.840078in}}{\pgfqpoint{3.674497in}{1.835960in}}%
\pgfpathcurveto{\pgfqpoint{3.670379in}{1.831841in}}{\pgfqpoint{3.668065in}{1.826255in}}{\pgfqpoint{3.668065in}{1.820431in}}%
\pgfpathcurveto{\pgfqpoint{3.668065in}{1.814607in}}{\pgfqpoint{3.670379in}{1.809021in}}{\pgfqpoint{3.674497in}{1.804903in}}%
\pgfpathcurveto{\pgfqpoint{3.678615in}{1.800785in}}{\pgfqpoint{3.684202in}{1.798471in}}{\pgfqpoint{3.690025in}{1.798471in}}%
\pgfpathlineto{\pgfqpoint{3.690025in}{1.798471in}}%
\pgfpathclose%
\pgfusepath{stroke,fill}%
\end{pgfscope}%
\begin{pgfscope}%
\pgfpathrectangle{\pgfqpoint{1.414736in}{0.548769in}}{\pgfqpoint{3.497747in}{7.253186in}}%
\pgfusepath{clip}%
\pgfsetbuttcap%
\pgfsetroundjoin%
\definecolor{currentfill}{rgb}{0.993248,0.906157,0.143936}%
\pgfsetfillcolor{currentfill}%
\pgfsetlinewidth{1.003750pt}%
\definecolor{currentstroke}{rgb}{0.993248,0.906157,0.143936}%
\pgfsetstrokecolor{currentstroke}%
\pgfsetdash{}{0pt}%
\pgfpathmoveto{\pgfqpoint{3.552734in}{2.269457in}}%
\pgfpathcurveto{\pgfqpoint{3.558557in}{2.269457in}}{\pgfqpoint{3.564144in}{2.271771in}}{\pgfqpoint{3.568262in}{2.275889in}}%
\pgfpathcurveto{\pgfqpoint{3.572380in}{2.280007in}}{\pgfqpoint{3.574694in}{2.285593in}}{\pgfqpoint{3.574694in}{2.291417in}}%
\pgfpathcurveto{\pgfqpoint{3.574694in}{2.297241in}}{\pgfqpoint{3.572380in}{2.302828in}}{\pgfqpoint{3.568262in}{2.306946in}}%
\pgfpathcurveto{\pgfqpoint{3.564144in}{2.311064in}}{\pgfqpoint{3.558557in}{2.313378in}}{\pgfqpoint{3.552734in}{2.313378in}}%
\pgfpathcurveto{\pgfqpoint{3.546910in}{2.313378in}}{\pgfqpoint{3.541323in}{2.311064in}}{\pgfqpoint{3.537205in}{2.306946in}}%
\pgfpathcurveto{\pgfqpoint{3.533087in}{2.302828in}}{\pgfqpoint{3.530773in}{2.297241in}}{\pgfqpoint{3.530773in}{2.291417in}}%
\pgfpathcurveto{\pgfqpoint{3.530773in}{2.285593in}}{\pgfqpoint{3.533087in}{2.280007in}}{\pgfqpoint{3.537205in}{2.275889in}}%
\pgfpathcurveto{\pgfqpoint{3.541323in}{2.271771in}}{\pgfqpoint{3.546910in}{2.269457in}}{\pgfqpoint{3.552734in}{2.269457in}}%
\pgfpathlineto{\pgfqpoint{3.552734in}{2.269457in}}%
\pgfpathclose%
\pgfusepath{stroke,fill}%
\end{pgfscope}%
\begin{pgfscope}%
\pgfpathrectangle{\pgfqpoint{1.414736in}{0.548769in}}{\pgfqpoint{3.497747in}{7.253186in}}%
\pgfusepath{clip}%
\pgfsetbuttcap%
\pgfsetroundjoin%
\definecolor{currentfill}{rgb}{0.266941,0.748751,0.440573}%
\pgfsetfillcolor{currentfill}%
\pgfsetlinewidth{1.003750pt}%
\definecolor{currentstroke}{rgb}{0.266941,0.748751,0.440573}%
\pgfsetstrokecolor{currentstroke}%
\pgfsetdash{}{0pt}%
\pgfpathmoveto{\pgfqpoint{3.696499in}{2.740443in}}%
\pgfpathcurveto{\pgfqpoint{3.702323in}{2.740443in}}{\pgfqpoint{3.707909in}{2.742757in}}{\pgfqpoint{3.712027in}{2.746875in}}%
\pgfpathcurveto{\pgfqpoint{3.716145in}{2.750993in}}{\pgfqpoint{3.718459in}{2.756580in}}{\pgfqpoint{3.718459in}{2.762404in}}%
\pgfpathcurveto{\pgfqpoint{3.718459in}{2.768227in}}{\pgfqpoint{3.716145in}{2.773814in}}{\pgfqpoint{3.712027in}{2.777932in}}%
\pgfpathcurveto{\pgfqpoint{3.707909in}{2.782050in}}{\pgfqpoint{3.702323in}{2.784364in}}{\pgfqpoint{3.696499in}{2.784364in}}%
\pgfpathcurveto{\pgfqpoint{3.690675in}{2.784364in}}{\pgfqpoint{3.685089in}{2.782050in}}{\pgfqpoint{3.680971in}{2.777932in}}%
\pgfpathcurveto{\pgfqpoint{3.676853in}{2.773814in}}{\pgfqpoint{3.674539in}{2.768227in}}{\pgfqpoint{3.674539in}{2.762404in}}%
\pgfpathcurveto{\pgfqpoint{3.674539in}{2.756580in}}{\pgfqpoint{3.676853in}{2.750993in}}{\pgfqpoint{3.680971in}{2.746875in}}%
\pgfpathcurveto{\pgfqpoint{3.685089in}{2.742757in}}{\pgfqpoint{3.690675in}{2.740443in}}{\pgfqpoint{3.696499in}{2.740443in}}%
\pgfpathlineto{\pgfqpoint{3.696499in}{2.740443in}}%
\pgfpathclose%
\pgfusepath{stroke,fill}%
\end{pgfscope}%
\begin{pgfscope}%
\pgfpathrectangle{\pgfqpoint{1.414736in}{0.548769in}}{\pgfqpoint{3.497747in}{7.253186in}}%
\pgfusepath{clip}%
\pgfsetbuttcap%
\pgfsetroundjoin%
\definecolor{currentfill}{rgb}{0.751884,0.874951,0.143228}%
\pgfsetfillcolor{currentfill}%
\pgfsetlinewidth{1.003750pt}%
\definecolor{currentstroke}{rgb}{0.751884,0.874951,0.143228}%
\pgfsetstrokecolor{currentstroke}%
\pgfsetdash{}{0pt}%
\pgfpathmoveto{\pgfqpoint{3.444754in}{3.211429in}}%
\pgfpathcurveto{\pgfqpoint{3.450578in}{3.211429in}}{\pgfqpoint{3.456164in}{3.213743in}}{\pgfqpoint{3.460282in}{3.217861in}}%
\pgfpathcurveto{\pgfqpoint{3.464400in}{3.221980in}}{\pgfqpoint{3.466714in}{3.227566in}}{\pgfqpoint{3.466714in}{3.233390in}}%
\pgfpathcurveto{\pgfqpoint{3.466714in}{3.239214in}}{\pgfqpoint{3.464400in}{3.244800in}}{\pgfqpoint{3.460282in}{3.248918in}}%
\pgfpathcurveto{\pgfqpoint{3.456164in}{3.253036in}}{\pgfqpoint{3.450578in}{3.255350in}}{\pgfqpoint{3.444754in}{3.255350in}}%
\pgfpathcurveto{\pgfqpoint{3.438930in}{3.255350in}}{\pgfqpoint{3.433344in}{3.253036in}}{\pgfqpoint{3.429226in}{3.248918in}}%
\pgfpathcurveto{\pgfqpoint{3.425108in}{3.244800in}}{\pgfqpoint{3.422794in}{3.239214in}}{\pgfqpoint{3.422794in}{3.233390in}}%
\pgfpathcurveto{\pgfqpoint{3.422794in}{3.227566in}}{\pgfqpoint{3.425108in}{3.221980in}}{\pgfqpoint{3.429226in}{3.217861in}}%
\pgfpathcurveto{\pgfqpoint{3.433344in}{3.213743in}}{\pgfqpoint{3.438930in}{3.211429in}}{\pgfqpoint{3.444754in}{3.211429in}}%
\pgfpathlineto{\pgfqpoint{3.444754in}{3.211429in}}%
\pgfpathclose%
\pgfusepath{stroke,fill}%
\end{pgfscope}%
\begin{pgfscope}%
\pgfpathrectangle{\pgfqpoint{1.414736in}{0.548769in}}{\pgfqpoint{3.497747in}{7.253186in}}%
\pgfusepath{clip}%
\pgfsetbuttcap%
\pgfsetroundjoin%
\definecolor{currentfill}{rgb}{0.993248,0.906157,0.143936}%
\pgfsetfillcolor{currentfill}%
\pgfsetlinewidth{1.003750pt}%
\definecolor{currentstroke}{rgb}{0.993248,0.906157,0.143936}%
\pgfsetstrokecolor{currentstroke}%
\pgfsetdash{}{0pt}%
\pgfpathmoveto{\pgfqpoint{3.427073in}{3.682415in}}%
\pgfpathcurveto{\pgfqpoint{3.432897in}{3.682415in}}{\pgfqpoint{3.438483in}{3.684729in}}{\pgfqpoint{3.442601in}{3.688847in}}%
\pgfpathcurveto{\pgfqpoint{3.446719in}{3.692966in}}{\pgfqpoint{3.449033in}{3.698552in}}{\pgfqpoint{3.449033in}{3.704376in}}%
\pgfpathcurveto{\pgfqpoint{3.449033in}{3.710200in}}{\pgfqpoint{3.446719in}{3.715786in}}{\pgfqpoint{3.442601in}{3.719904in}}%
\pgfpathcurveto{\pgfqpoint{3.438483in}{3.724022in}}{\pgfqpoint{3.432897in}{3.726336in}}{\pgfqpoint{3.427073in}{3.726336in}}%
\pgfpathcurveto{\pgfqpoint{3.421249in}{3.726336in}}{\pgfqpoint{3.415663in}{3.724022in}}{\pgfqpoint{3.411545in}{3.719904in}}%
\pgfpathcurveto{\pgfqpoint{3.407426in}{3.715786in}}{\pgfqpoint{3.405113in}{3.710200in}}{\pgfqpoint{3.405113in}{3.704376in}}%
\pgfpathcurveto{\pgfqpoint{3.405113in}{3.698552in}}{\pgfqpoint{3.407426in}{3.692966in}}{\pgfqpoint{3.411545in}{3.688847in}}%
\pgfpathcurveto{\pgfqpoint{3.415663in}{3.684729in}}{\pgfqpoint{3.421249in}{3.682415in}}{\pgfqpoint{3.427073in}{3.682415in}}%
\pgfpathlineto{\pgfqpoint{3.427073in}{3.682415in}}%
\pgfpathclose%
\pgfusepath{stroke,fill}%
\end{pgfscope}%
\begin{pgfscope}%
\pgfpathrectangle{\pgfqpoint{1.414736in}{0.548769in}}{\pgfqpoint{3.497747in}{7.253186in}}%
\pgfusepath{clip}%
\pgfsetbuttcap%
\pgfsetroundjoin%
\definecolor{currentfill}{rgb}{0.955300,0.901065,0.118128}%
\pgfsetfillcolor{currentfill}%
\pgfsetlinewidth{1.003750pt}%
\definecolor{currentstroke}{rgb}{0.955300,0.901065,0.118128}%
\pgfsetstrokecolor{currentstroke}%
\pgfsetdash{}{0pt}%
\pgfpathmoveto{\pgfqpoint{3.946403in}{4.153402in}}%
\pgfpathcurveto{\pgfqpoint{3.952227in}{4.153402in}}{\pgfqpoint{3.957813in}{4.155715in}}{\pgfqpoint{3.961931in}{4.159834in}}%
\pgfpathcurveto{\pgfqpoint{3.966049in}{4.163952in}}{\pgfqpoint{3.968363in}{4.169538in}}{\pgfqpoint{3.968363in}{4.175362in}}%
\pgfpathcurveto{\pgfqpoint{3.968363in}{4.181186in}}{\pgfqpoint{3.966049in}{4.186772in}}{\pgfqpoint{3.961931in}{4.190890in}}%
\pgfpathcurveto{\pgfqpoint{3.957813in}{4.195008in}}{\pgfqpoint{3.952227in}{4.197322in}}{\pgfqpoint{3.946403in}{4.197322in}}%
\pgfpathcurveto{\pgfqpoint{3.940579in}{4.197322in}}{\pgfqpoint{3.934992in}{4.195008in}}{\pgfqpoint{3.930874in}{4.190890in}}%
\pgfpathcurveto{\pgfqpoint{3.926756in}{4.186772in}}{\pgfqpoint{3.924442in}{4.181186in}}{\pgfqpoint{3.924442in}{4.175362in}}%
\pgfpathcurveto{\pgfqpoint{3.924442in}{4.169538in}}{\pgfqpoint{3.926756in}{4.163952in}}{\pgfqpoint{3.930874in}{4.159834in}}%
\pgfpathcurveto{\pgfqpoint{3.934992in}{4.155715in}}{\pgfqpoint{3.940579in}{4.153402in}}{\pgfqpoint{3.946403in}{4.153402in}}%
\pgfpathlineto{\pgfqpoint{3.946403in}{4.153402in}}%
\pgfpathclose%
\pgfusepath{stroke,fill}%
\end{pgfscope}%
\begin{pgfscope}%
\pgfpathrectangle{\pgfqpoint{1.414736in}{0.548769in}}{\pgfqpoint{3.497747in}{7.253186in}}%
\pgfusepath{clip}%
\pgfsetbuttcap%
\pgfsetroundjoin%
\definecolor{currentfill}{rgb}{0.772852,0.877868,0.131109}%
\pgfsetfillcolor{currentfill}%
\pgfsetlinewidth{1.003750pt}%
\definecolor{currentstroke}{rgb}{0.772852,0.877868,0.131109}%
\pgfsetstrokecolor{currentstroke}%
\pgfsetdash{}{0pt}%
\pgfpathmoveto{\pgfqpoint{3.659104in}{4.624388in}}%
\pgfpathcurveto{\pgfqpoint{3.664927in}{4.624388in}}{\pgfqpoint{3.670514in}{4.626702in}}{\pgfqpoint{3.674632in}{4.630820in}}%
\pgfpathcurveto{\pgfqpoint{3.678750in}{4.634938in}}{\pgfqpoint{3.681064in}{4.640524in}}{\pgfqpoint{3.681064in}{4.646348in}}%
\pgfpathcurveto{\pgfqpoint{3.681064in}{4.652172in}}{\pgfqpoint{3.678750in}{4.657758in}}{\pgfqpoint{3.674632in}{4.661876in}}%
\pgfpathcurveto{\pgfqpoint{3.670514in}{4.665994in}}{\pgfqpoint{3.664927in}{4.668308in}}{\pgfqpoint{3.659104in}{4.668308in}}%
\pgfpathcurveto{\pgfqpoint{3.653280in}{4.668308in}}{\pgfqpoint{3.647693in}{4.665994in}}{\pgfqpoint{3.643575in}{4.661876in}}%
\pgfpathcurveto{\pgfqpoint{3.639457in}{4.657758in}}{\pgfqpoint{3.637143in}{4.652172in}}{\pgfqpoint{3.637143in}{4.646348in}}%
\pgfpathcurveto{\pgfqpoint{3.637143in}{4.640524in}}{\pgfqpoint{3.639457in}{4.634938in}}{\pgfqpoint{3.643575in}{4.630820in}}%
\pgfpathcurveto{\pgfqpoint{3.647693in}{4.626702in}}{\pgfqpoint{3.653280in}{4.624388in}}{\pgfqpoint{3.659104in}{4.624388in}}%
\pgfpathlineto{\pgfqpoint{3.659104in}{4.624388in}}%
\pgfpathclose%
\pgfusepath{stroke,fill}%
\end{pgfscope}%
\begin{pgfscope}%
\pgfpathrectangle{\pgfqpoint{1.414736in}{0.548769in}}{\pgfqpoint{3.497747in}{7.253186in}}%
\pgfusepath{clip}%
\pgfsetbuttcap%
\pgfsetroundjoin%
\definecolor{currentfill}{rgb}{0.153364,0.497000,0.557724}%
\pgfsetfillcolor{currentfill}%
\pgfsetlinewidth{1.003750pt}%
\definecolor{currentstroke}{rgb}{0.153364,0.497000,0.557724}%
\pgfsetstrokecolor{currentstroke}%
\pgfsetdash{}{0pt}%
\pgfpathmoveto{\pgfqpoint{3.578211in}{5.095374in}}%
\pgfpathcurveto{\pgfqpoint{3.584035in}{5.095374in}}{\pgfqpoint{3.589622in}{5.097688in}}{\pgfqpoint{3.593740in}{5.101806in}}%
\pgfpathcurveto{\pgfqpoint{3.597858in}{5.105924in}}{\pgfqpoint{3.600172in}{5.111510in}}{\pgfqpoint{3.600172in}{5.117334in}}%
\pgfpathcurveto{\pgfqpoint{3.600172in}{5.123158in}}{\pgfqpoint{3.597858in}{5.128744in}}{\pgfqpoint{3.593740in}{5.132862in}}%
\pgfpathcurveto{\pgfqpoint{3.589622in}{5.136980in}}{\pgfqpoint{3.584035in}{5.139294in}}{\pgfqpoint{3.578211in}{5.139294in}}%
\pgfpathcurveto{\pgfqpoint{3.572388in}{5.139294in}}{\pgfqpoint{3.566801in}{5.136980in}}{\pgfqpoint{3.562683in}{5.132862in}}%
\pgfpathcurveto{\pgfqpoint{3.558565in}{5.128744in}}{\pgfqpoint{3.556251in}{5.123158in}}{\pgfqpoint{3.556251in}{5.117334in}}%
\pgfpathcurveto{\pgfqpoint{3.556251in}{5.111510in}}{\pgfqpoint{3.558565in}{5.105924in}}{\pgfqpoint{3.562683in}{5.101806in}}%
\pgfpathcurveto{\pgfqpoint{3.566801in}{5.097688in}}{\pgfqpoint{3.572388in}{5.095374in}}{\pgfqpoint{3.578211in}{5.095374in}}%
\pgfpathlineto{\pgfqpoint{3.578211in}{5.095374in}}%
\pgfpathclose%
\pgfusepath{stroke,fill}%
\end{pgfscope}%
\begin{pgfscope}%
\pgfpathrectangle{\pgfqpoint{1.414736in}{0.548769in}}{\pgfqpoint{3.497747in}{7.253186in}}%
\pgfusepath{clip}%
\pgfsetbuttcap%
\pgfsetroundjoin%
\definecolor{currentfill}{rgb}{0.199430,0.387607,0.554642}%
\pgfsetfillcolor{currentfill}%
\pgfsetlinewidth{1.003750pt}%
\definecolor{currentstroke}{rgb}{0.199430,0.387607,0.554642}%
\pgfsetstrokecolor{currentstroke}%
\pgfsetdash{}{0pt}%
\pgfpathmoveto{\pgfqpoint{3.940391in}{5.566360in}}%
\pgfpathcurveto{\pgfqpoint{3.946215in}{5.566360in}}{\pgfqpoint{3.951801in}{5.568674in}}{\pgfqpoint{3.955920in}{5.572792in}}%
\pgfpathcurveto{\pgfqpoint{3.960038in}{5.576910in}}{\pgfqpoint{3.962352in}{5.582496in}}{\pgfqpoint{3.962352in}{5.588320in}}%
\pgfpathcurveto{\pgfqpoint{3.962352in}{5.594144in}}{\pgfqpoint{3.960038in}{5.599730in}}{\pgfqpoint{3.955920in}{5.603848in}}%
\pgfpathcurveto{\pgfqpoint{3.951801in}{5.607967in}}{\pgfqpoint{3.946215in}{5.610280in}}{\pgfqpoint{3.940391in}{5.610280in}}%
\pgfpathcurveto{\pgfqpoint{3.934567in}{5.610280in}}{\pgfqpoint{3.928981in}{5.607967in}}{\pgfqpoint{3.924863in}{5.603848in}}%
\pgfpathcurveto{\pgfqpoint{3.920745in}{5.599730in}}{\pgfqpoint{3.918431in}{5.594144in}}{\pgfqpoint{3.918431in}{5.588320in}}%
\pgfpathcurveto{\pgfqpoint{3.918431in}{5.582496in}}{\pgfqpoint{3.920745in}{5.576910in}}{\pgfqpoint{3.924863in}{5.572792in}}%
\pgfpathcurveto{\pgfqpoint{3.928981in}{5.568674in}}{\pgfqpoint{3.934567in}{5.566360in}}{\pgfqpoint{3.940391in}{5.566360in}}%
\pgfpathlineto{\pgfqpoint{3.940391in}{5.566360in}}%
\pgfpathclose%
\pgfusepath{stroke,fill}%
\end{pgfscope}%
\begin{pgfscope}%
\pgfpathrectangle{\pgfqpoint{1.414736in}{0.548769in}}{\pgfqpoint{3.497747in}{7.253186in}}%
\pgfusepath{clip}%
\pgfsetbuttcap%
\pgfsetroundjoin%
\definecolor{currentfill}{rgb}{0.311925,0.767822,0.415586}%
\pgfsetfillcolor{currentfill}%
\pgfsetlinewidth{1.003750pt}%
\definecolor{currentstroke}{rgb}{0.311925,0.767822,0.415586}%
\pgfsetstrokecolor{currentstroke}%
\pgfsetdash{}{0pt}%
\pgfpathmoveto{\pgfqpoint{3.237539in}{6.037346in}}%
\pgfpathcurveto{\pgfqpoint{3.243363in}{6.037346in}}{\pgfqpoint{3.248949in}{6.039660in}}{\pgfqpoint{3.253068in}{6.043778in}}%
\pgfpathcurveto{\pgfqpoint{3.257186in}{6.047896in}}{\pgfqpoint{3.259500in}{6.053482in}}{\pgfqpoint{3.259500in}{6.059306in}}%
\pgfpathcurveto{\pgfqpoint{3.259500in}{6.065130in}}{\pgfqpoint{3.257186in}{6.070716in}}{\pgfqpoint{3.253068in}{6.074835in}}%
\pgfpathcurveto{\pgfqpoint{3.248949in}{6.078953in}}{\pgfqpoint{3.243363in}{6.081267in}}{\pgfqpoint{3.237539in}{6.081267in}}%
\pgfpathcurveto{\pgfqpoint{3.231715in}{6.081267in}}{\pgfqpoint{3.226129in}{6.078953in}}{\pgfqpoint{3.222011in}{6.074835in}}%
\pgfpathcurveto{\pgfqpoint{3.217893in}{6.070716in}}{\pgfqpoint{3.215579in}{6.065130in}}{\pgfqpoint{3.215579in}{6.059306in}}%
\pgfpathcurveto{\pgfqpoint{3.215579in}{6.053482in}}{\pgfqpoint{3.217893in}{6.047896in}}{\pgfqpoint{3.222011in}{6.043778in}}%
\pgfpathcurveto{\pgfqpoint{3.226129in}{6.039660in}}{\pgfqpoint{3.231715in}{6.037346in}}{\pgfqpoint{3.237539in}{6.037346in}}%
\pgfpathlineto{\pgfqpoint{3.237539in}{6.037346in}}%
\pgfpathclose%
\pgfusepath{stroke,fill}%
\end{pgfscope}%
\begin{pgfscope}%
\pgfpathrectangle{\pgfqpoint{1.414736in}{0.548769in}}{\pgfqpoint{3.497747in}{7.253186in}}%
\pgfusepath{clip}%
\pgfsetbuttcap%
\pgfsetroundjoin%
\definecolor{currentfill}{rgb}{0.993248,0.906157,0.143936}%
\pgfsetfillcolor{currentfill}%
\pgfsetlinewidth{1.003750pt}%
\definecolor{currentstroke}{rgb}{0.993248,0.906157,0.143936}%
\pgfsetstrokecolor{currentstroke}%
\pgfsetdash{}{0pt}%
\pgfpathmoveto{\pgfqpoint{3.756720in}{6.508332in}}%
\pgfpathcurveto{\pgfqpoint{3.762544in}{6.508332in}}{\pgfqpoint{3.768130in}{6.510646in}}{\pgfqpoint{3.772248in}{6.514764in}}%
\pgfpathcurveto{\pgfqpoint{3.776366in}{6.518882in}}{\pgfqpoint{3.778680in}{6.524468in}}{\pgfqpoint{3.778680in}{6.530292in}}%
\pgfpathcurveto{\pgfqpoint{3.778680in}{6.536116in}}{\pgfqpoint{3.776366in}{6.541702in}}{\pgfqpoint{3.772248in}{6.545821in}}%
\pgfpathcurveto{\pgfqpoint{3.768130in}{6.549939in}}{\pgfqpoint{3.762544in}{6.552253in}}{\pgfqpoint{3.756720in}{6.552253in}}%
\pgfpathcurveto{\pgfqpoint{3.750896in}{6.552253in}}{\pgfqpoint{3.745310in}{6.549939in}}{\pgfqpoint{3.741191in}{6.545821in}}%
\pgfpathcurveto{\pgfqpoint{3.737073in}{6.541702in}}{\pgfqpoint{3.734759in}{6.536116in}}{\pgfqpoint{3.734759in}{6.530292in}}%
\pgfpathcurveto{\pgfqpoint{3.734759in}{6.524468in}}{\pgfqpoint{3.737073in}{6.518882in}}{\pgfqpoint{3.741191in}{6.514764in}}%
\pgfpathcurveto{\pgfqpoint{3.745310in}{6.510646in}}{\pgfqpoint{3.750896in}{6.508332in}}{\pgfqpoint{3.756720in}{6.508332in}}%
\pgfpathlineto{\pgfqpoint{3.756720in}{6.508332in}}%
\pgfpathclose%
\pgfusepath{stroke,fill}%
\end{pgfscope}%
\begin{pgfscope}%
\pgfpathrectangle{\pgfqpoint{1.414736in}{0.548769in}}{\pgfqpoint{3.497747in}{7.253186in}}%
\pgfusepath{clip}%
\pgfsetbuttcap%
\pgfsetroundjoin%
\definecolor{currentfill}{rgb}{0.993248,0.906157,0.143936}%
\pgfsetfillcolor{currentfill}%
\pgfsetlinewidth{1.003750pt}%
\definecolor{currentstroke}{rgb}{0.993248,0.906157,0.143936}%
\pgfsetstrokecolor{currentstroke}%
\pgfsetdash{}{0pt}%
\pgfpathmoveto{\pgfqpoint{3.644316in}{6.979318in}}%
\pgfpathcurveto{\pgfqpoint{3.650140in}{6.979318in}}{\pgfqpoint{3.655727in}{6.981632in}}{\pgfqpoint{3.659845in}{6.985750in}}%
\pgfpathcurveto{\pgfqpoint{3.663963in}{6.989868in}}{\pgfqpoint{3.666277in}{6.995455in}}{\pgfqpoint{3.666277in}{7.001278in}}%
\pgfpathcurveto{\pgfqpoint{3.666277in}{7.007102in}}{\pgfqpoint{3.663963in}{7.012689in}}{\pgfqpoint{3.659845in}{7.016807in}}%
\pgfpathcurveto{\pgfqpoint{3.655727in}{7.020925in}}{\pgfqpoint{3.650140in}{7.023239in}}{\pgfqpoint{3.644316in}{7.023239in}}%
\pgfpathcurveto{\pgfqpoint{3.638493in}{7.023239in}}{\pgfqpoint{3.632906in}{7.020925in}}{\pgfqpoint{3.628788in}{7.016807in}}%
\pgfpathcurveto{\pgfqpoint{3.624670in}{7.012689in}}{\pgfqpoint{3.622356in}{7.007102in}}{\pgfqpoint{3.622356in}{7.001278in}}%
\pgfpathcurveto{\pgfqpoint{3.622356in}{6.995455in}}{\pgfqpoint{3.624670in}{6.989868in}}{\pgfqpoint{3.628788in}{6.985750in}}%
\pgfpathcurveto{\pgfqpoint{3.632906in}{6.981632in}}{\pgfqpoint{3.638493in}{6.979318in}}{\pgfqpoint{3.644316in}{6.979318in}}%
\pgfpathlineto{\pgfqpoint{3.644316in}{6.979318in}}%
\pgfpathclose%
\pgfusepath{stroke,fill}%
\end{pgfscope}%
\begin{pgfscope}%
\pgfpathrectangle{\pgfqpoint{1.414736in}{0.548769in}}{\pgfqpoint{3.497747in}{7.253186in}}%
\pgfusepath{clip}%
\pgfsetbuttcap%
\pgfsetroundjoin%
\definecolor{currentfill}{rgb}{0.468053,0.818921,0.323998}%
\pgfsetfillcolor{currentfill}%
\pgfsetlinewidth{1.003750pt}%
\definecolor{currentstroke}{rgb}{0.468053,0.818921,0.323998}%
\pgfsetstrokecolor{currentstroke}%
\pgfsetdash{}{0pt}%
\pgfpathmoveto{\pgfqpoint{4.324761in}{7.450304in}}%
\pgfpathcurveto{\pgfqpoint{4.330585in}{7.450304in}}{\pgfqpoint{4.336171in}{7.452618in}}{\pgfqpoint{4.340289in}{7.456736in}}%
\pgfpathcurveto{\pgfqpoint{4.344407in}{7.460854in}}{\pgfqpoint{4.346721in}{7.466441in}}{\pgfqpoint{4.346721in}{7.472265in}}%
\pgfpathcurveto{\pgfqpoint{4.346721in}{7.478089in}}{\pgfqpoint{4.344407in}{7.483675in}}{\pgfqpoint{4.340289in}{7.487793in}}%
\pgfpathcurveto{\pgfqpoint{4.336171in}{7.491911in}}{\pgfqpoint{4.330585in}{7.494225in}}{\pgfqpoint{4.324761in}{7.494225in}}%
\pgfpathcurveto{\pgfqpoint{4.318937in}{7.494225in}}{\pgfqpoint{4.313351in}{7.491911in}}{\pgfqpoint{4.309233in}{7.487793in}}%
\pgfpathcurveto{\pgfqpoint{4.305114in}{7.483675in}}{\pgfqpoint{4.302801in}{7.478089in}}{\pgfqpoint{4.302801in}{7.472265in}}%
\pgfpathcurveto{\pgfqpoint{4.302801in}{7.466441in}}{\pgfqpoint{4.305114in}{7.460854in}}{\pgfqpoint{4.309233in}{7.456736in}}%
\pgfpathcurveto{\pgfqpoint{4.313351in}{7.452618in}}{\pgfqpoint{4.318937in}{7.450304in}}{\pgfqpoint{4.324761in}{7.450304in}}%
\pgfpathlineto{\pgfqpoint{4.324761in}{7.450304in}}%
\pgfpathclose%
\pgfusepath{stroke,fill}%
\end{pgfscope}%
\begin{pgfscope}%
\pgfpathrectangle{\pgfqpoint{1.414736in}{0.548769in}}{\pgfqpoint{3.497747in}{7.253186in}}%
\pgfusepath{clip}%
\pgfsetbuttcap%
\pgfsetroundjoin%
\definecolor{currentfill}{rgb}{0.487026,0.823929,0.312321}%
\pgfsetfillcolor{currentfill}%
\pgfsetlinewidth{1.003750pt}%
\definecolor{currentstroke}{rgb}{0.487026,0.823929,0.312321}%
\pgfsetstrokecolor{currentstroke}%
\pgfsetdash{}{0pt}%
\pgfpathmoveto{\pgfqpoint{3.309341in}{0.856499in}}%
\pgfpathcurveto{\pgfqpoint{3.315165in}{0.856499in}}{\pgfqpoint{3.320751in}{0.858813in}}{\pgfqpoint{3.324869in}{0.862931in}}%
\pgfpathcurveto{\pgfqpoint{3.328988in}{0.867049in}}{\pgfqpoint{3.331301in}{0.872635in}}{\pgfqpoint{3.331301in}{0.878459in}}%
\pgfpathcurveto{\pgfqpoint{3.331301in}{0.884283in}}{\pgfqpoint{3.328988in}{0.889869in}}{\pgfqpoint{3.324869in}{0.893987in}}%
\pgfpathcurveto{\pgfqpoint{3.320751in}{0.898105in}}{\pgfqpoint{3.315165in}{0.900419in}}{\pgfqpoint{3.309341in}{0.900419in}}%
\pgfpathcurveto{\pgfqpoint{3.303517in}{0.900419in}}{\pgfqpoint{3.297931in}{0.898105in}}{\pgfqpoint{3.293813in}{0.893987in}}%
\pgfpathcurveto{\pgfqpoint{3.289695in}{0.889869in}}{\pgfqpoint{3.287381in}{0.884283in}}{\pgfqpoint{3.287381in}{0.878459in}}%
\pgfpathcurveto{\pgfqpoint{3.287381in}{0.872635in}}{\pgfqpoint{3.289695in}{0.867049in}}{\pgfqpoint{3.293813in}{0.862931in}}%
\pgfpathcurveto{\pgfqpoint{3.297931in}{0.858813in}}{\pgfqpoint{3.303517in}{0.856499in}}{\pgfqpoint{3.309341in}{0.856499in}}%
\pgfpathlineto{\pgfqpoint{3.309341in}{0.856499in}}%
\pgfpathclose%
\pgfusepath{stroke,fill}%
\end{pgfscope}%
\begin{pgfscope}%
\pgfpathrectangle{\pgfqpoint{1.414736in}{0.548769in}}{\pgfqpoint{3.497747in}{7.253186in}}%
\pgfusepath{clip}%
\pgfsetbuttcap%
\pgfsetroundjoin%
\definecolor{currentfill}{rgb}{0.208030,0.718701,0.472873}%
\pgfsetfillcolor{currentfill}%
\pgfsetlinewidth{1.003750pt}%
\definecolor{currentstroke}{rgb}{0.208030,0.718701,0.472873}%
\pgfsetstrokecolor{currentstroke}%
\pgfsetdash{}{0pt}%
\pgfpathmoveto{\pgfqpoint{3.709735in}{1.327485in}}%
\pgfpathcurveto{\pgfqpoint{3.715559in}{1.327485in}}{\pgfqpoint{3.721145in}{1.329799in}}{\pgfqpoint{3.725263in}{1.333917in}}%
\pgfpathcurveto{\pgfqpoint{3.729381in}{1.338035in}}{\pgfqpoint{3.731695in}{1.343621in}}{\pgfqpoint{3.731695in}{1.349445in}}%
\pgfpathcurveto{\pgfqpoint{3.731695in}{1.355269in}}{\pgfqpoint{3.729381in}{1.360855in}}{\pgfqpoint{3.725263in}{1.364973in}}%
\pgfpathcurveto{\pgfqpoint{3.721145in}{1.369092in}}{\pgfqpoint{3.715559in}{1.371405in}}{\pgfqpoint{3.709735in}{1.371405in}}%
\pgfpathcurveto{\pgfqpoint{3.703911in}{1.371405in}}{\pgfqpoint{3.698325in}{1.369092in}}{\pgfqpoint{3.694206in}{1.364973in}}%
\pgfpathcurveto{\pgfqpoint{3.690088in}{1.360855in}}{\pgfqpoint{3.687774in}{1.355269in}}{\pgfqpoint{3.687774in}{1.349445in}}%
\pgfpathcurveto{\pgfqpoint{3.687774in}{1.343621in}}{\pgfqpoint{3.690088in}{1.338035in}}{\pgfqpoint{3.694206in}{1.333917in}}%
\pgfpathcurveto{\pgfqpoint{3.698325in}{1.329799in}}{\pgfqpoint{3.703911in}{1.327485in}}{\pgfqpoint{3.709735in}{1.327485in}}%
\pgfpathlineto{\pgfqpoint{3.709735in}{1.327485in}}%
\pgfpathclose%
\pgfusepath{stroke,fill}%
\end{pgfscope}%
\begin{pgfscope}%
\pgfpathrectangle{\pgfqpoint{1.414736in}{0.548769in}}{\pgfqpoint{3.497747in}{7.253186in}}%
\pgfusepath{clip}%
\pgfsetbuttcap%
\pgfsetroundjoin%
\definecolor{currentfill}{rgb}{0.804182,0.882046,0.114965}%
\pgfsetfillcolor{currentfill}%
\pgfsetlinewidth{1.003750pt}%
\definecolor{currentstroke}{rgb}{0.804182,0.882046,0.114965}%
\pgfsetstrokecolor{currentstroke}%
\pgfsetdash{}{0pt}%
\pgfpathmoveto{\pgfqpoint{3.896374in}{1.798471in}}%
\pgfpathcurveto{\pgfqpoint{3.902198in}{1.798471in}}{\pgfqpoint{3.907784in}{1.800785in}}{\pgfqpoint{3.911902in}{1.804903in}}%
\pgfpathcurveto{\pgfqpoint{3.916020in}{1.809021in}}{\pgfqpoint{3.918334in}{1.814607in}}{\pgfqpoint{3.918334in}{1.820431in}}%
\pgfpathcurveto{\pgfqpoint{3.918334in}{1.826255in}}{\pgfqpoint{3.916020in}{1.831841in}}{\pgfqpoint{3.911902in}{1.835960in}}%
\pgfpathcurveto{\pgfqpoint{3.907784in}{1.840078in}}{\pgfqpoint{3.902198in}{1.842392in}}{\pgfqpoint{3.896374in}{1.842392in}}%
\pgfpathcurveto{\pgfqpoint{3.890550in}{1.842392in}}{\pgfqpoint{3.884964in}{1.840078in}}{\pgfqpoint{3.880845in}{1.835960in}}%
\pgfpathcurveto{\pgfqpoint{3.876727in}{1.831841in}}{\pgfqpoint{3.874413in}{1.826255in}}{\pgfqpoint{3.874413in}{1.820431in}}%
\pgfpathcurveto{\pgfqpoint{3.874413in}{1.814607in}}{\pgfqpoint{3.876727in}{1.809021in}}{\pgfqpoint{3.880845in}{1.804903in}}%
\pgfpathcurveto{\pgfqpoint{3.884964in}{1.800785in}}{\pgfqpoint{3.890550in}{1.798471in}}{\pgfqpoint{3.896374in}{1.798471in}}%
\pgfpathlineto{\pgfqpoint{3.896374in}{1.798471in}}%
\pgfpathclose%
\pgfusepath{stroke,fill}%
\end{pgfscope}%
\begin{pgfscope}%
\pgfpathrectangle{\pgfqpoint{1.414736in}{0.548769in}}{\pgfqpoint{3.497747in}{7.253186in}}%
\pgfusepath{clip}%
\pgfsetbuttcap%
\pgfsetroundjoin%
\definecolor{currentfill}{rgb}{0.668054,0.861999,0.196293}%
\pgfsetfillcolor{currentfill}%
\pgfsetlinewidth{1.003750pt}%
\definecolor{currentstroke}{rgb}{0.668054,0.861999,0.196293}%
\pgfsetstrokecolor{currentstroke}%
\pgfsetdash{}{0pt}%
\pgfpathmoveto{\pgfqpoint{3.794862in}{2.269457in}}%
\pgfpathcurveto{\pgfqpoint{3.800686in}{2.269457in}}{\pgfqpoint{3.806272in}{2.271771in}}{\pgfqpoint{3.810390in}{2.275889in}}%
\pgfpathcurveto{\pgfqpoint{3.814508in}{2.280007in}}{\pgfqpoint{3.816822in}{2.285593in}}{\pgfqpoint{3.816822in}{2.291417in}}%
\pgfpathcurveto{\pgfqpoint{3.816822in}{2.297241in}}{\pgfqpoint{3.814508in}{2.302828in}}{\pgfqpoint{3.810390in}{2.306946in}}%
\pgfpathcurveto{\pgfqpoint{3.806272in}{2.311064in}}{\pgfqpoint{3.800686in}{2.313378in}}{\pgfqpoint{3.794862in}{2.313378in}}%
\pgfpathcurveto{\pgfqpoint{3.789038in}{2.313378in}}{\pgfqpoint{3.783452in}{2.311064in}}{\pgfqpoint{3.779334in}{2.306946in}}%
\pgfpathcurveto{\pgfqpoint{3.775216in}{2.302828in}}{\pgfqpoint{3.772902in}{2.297241in}}{\pgfqpoint{3.772902in}{2.291417in}}%
\pgfpathcurveto{\pgfqpoint{3.772902in}{2.285593in}}{\pgfqpoint{3.775216in}{2.280007in}}{\pgfqpoint{3.779334in}{2.275889in}}%
\pgfpathcurveto{\pgfqpoint{3.783452in}{2.271771in}}{\pgfqpoint{3.789038in}{2.269457in}}{\pgfqpoint{3.794862in}{2.269457in}}%
\pgfpathlineto{\pgfqpoint{3.794862in}{2.269457in}}%
\pgfpathclose%
\pgfusepath{stroke,fill}%
\end{pgfscope}%
\begin{pgfscope}%
\pgfpathrectangle{\pgfqpoint{1.414736in}{0.548769in}}{\pgfqpoint{3.497747in}{7.253186in}}%
\pgfusepath{clip}%
\pgfsetbuttcap%
\pgfsetroundjoin%
\definecolor{currentfill}{rgb}{0.352360,0.783011,0.392636}%
\pgfsetfillcolor{currentfill}%
\pgfsetlinewidth{1.003750pt}%
\definecolor{currentstroke}{rgb}{0.352360,0.783011,0.392636}%
\pgfsetstrokecolor{currentstroke}%
\pgfsetdash{}{0pt}%
\pgfpathmoveto{\pgfqpoint{3.166118in}{2.740443in}}%
\pgfpathcurveto{\pgfqpoint{3.171942in}{2.740443in}}{\pgfqpoint{3.177528in}{2.742757in}}{\pgfqpoint{3.181647in}{2.746875in}}%
\pgfpathcurveto{\pgfqpoint{3.185765in}{2.750993in}}{\pgfqpoint{3.188079in}{2.756580in}}{\pgfqpoint{3.188079in}{2.762404in}}%
\pgfpathcurveto{\pgfqpoint{3.188079in}{2.768227in}}{\pgfqpoint{3.185765in}{2.773814in}}{\pgfqpoint{3.181647in}{2.777932in}}%
\pgfpathcurveto{\pgfqpoint{3.177528in}{2.782050in}}{\pgfqpoint{3.171942in}{2.784364in}}{\pgfqpoint{3.166118in}{2.784364in}}%
\pgfpathcurveto{\pgfqpoint{3.160294in}{2.784364in}}{\pgfqpoint{3.154708in}{2.782050in}}{\pgfqpoint{3.150590in}{2.777932in}}%
\pgfpathcurveto{\pgfqpoint{3.146472in}{2.773814in}}{\pgfqpoint{3.144158in}{2.768227in}}{\pgfqpoint{3.144158in}{2.762404in}}%
\pgfpathcurveto{\pgfqpoint{3.144158in}{2.756580in}}{\pgfqpoint{3.146472in}{2.750993in}}{\pgfqpoint{3.150590in}{2.746875in}}%
\pgfpathcurveto{\pgfqpoint{3.154708in}{2.742757in}}{\pgfqpoint{3.160294in}{2.740443in}}{\pgfqpoint{3.166118in}{2.740443in}}%
\pgfpathlineto{\pgfqpoint{3.166118in}{2.740443in}}%
\pgfpathclose%
\pgfusepath{stroke,fill}%
\end{pgfscope}%
\begin{pgfscope}%
\pgfpathrectangle{\pgfqpoint{1.414736in}{0.548769in}}{\pgfqpoint{3.497747in}{7.253186in}}%
\pgfusepath{clip}%
\pgfsetbuttcap%
\pgfsetroundjoin%
\definecolor{currentfill}{rgb}{0.525776,0.833491,0.288127}%
\pgfsetfillcolor{currentfill}%
\pgfsetlinewidth{1.003750pt}%
\definecolor{currentstroke}{rgb}{0.525776,0.833491,0.288127}%
\pgfsetstrokecolor{currentstroke}%
\pgfsetdash{}{0pt}%
\pgfpathmoveto{\pgfqpoint{3.796674in}{3.211429in}}%
\pgfpathcurveto{\pgfqpoint{3.802498in}{3.211429in}}{\pgfqpoint{3.808084in}{3.213743in}}{\pgfqpoint{3.812202in}{3.217861in}}%
\pgfpathcurveto{\pgfqpoint{3.816320in}{3.221980in}}{\pgfqpoint{3.818634in}{3.227566in}}{\pgfqpoint{3.818634in}{3.233390in}}%
\pgfpathcurveto{\pgfqpoint{3.818634in}{3.239214in}}{\pgfqpoint{3.816320in}{3.244800in}}{\pgfqpoint{3.812202in}{3.248918in}}%
\pgfpathcurveto{\pgfqpoint{3.808084in}{3.253036in}}{\pgfqpoint{3.802498in}{3.255350in}}{\pgfqpoint{3.796674in}{3.255350in}}%
\pgfpathcurveto{\pgfqpoint{3.790850in}{3.255350in}}{\pgfqpoint{3.785263in}{3.253036in}}{\pgfqpoint{3.781145in}{3.248918in}}%
\pgfpathcurveto{\pgfqpoint{3.777027in}{3.244800in}}{\pgfqpoint{3.774713in}{3.239214in}}{\pgfqpoint{3.774713in}{3.233390in}}%
\pgfpathcurveto{\pgfqpoint{3.774713in}{3.227566in}}{\pgfqpoint{3.777027in}{3.221980in}}{\pgfqpoint{3.781145in}{3.217861in}}%
\pgfpathcurveto{\pgfqpoint{3.785263in}{3.213743in}}{\pgfqpoint{3.790850in}{3.211429in}}{\pgfqpoint{3.796674in}{3.211429in}}%
\pgfpathlineto{\pgfqpoint{3.796674in}{3.211429in}}%
\pgfpathclose%
\pgfusepath{stroke,fill}%
\end{pgfscope}%
\begin{pgfscope}%
\pgfpathrectangle{\pgfqpoint{1.414736in}{0.548769in}}{\pgfqpoint{3.497747in}{7.253186in}}%
\pgfusepath{clip}%
\pgfsetbuttcap%
\pgfsetroundjoin%
\definecolor{currentfill}{rgb}{0.935904,0.898570,0.108131}%
\pgfsetfillcolor{currentfill}%
\pgfsetlinewidth{1.003750pt}%
\definecolor{currentstroke}{rgb}{0.935904,0.898570,0.108131}%
\pgfsetstrokecolor{currentstroke}%
\pgfsetdash{}{0pt}%
\pgfpathmoveto{\pgfqpoint{4.065650in}{3.682415in}}%
\pgfpathcurveto{\pgfqpoint{4.071474in}{3.682415in}}{\pgfqpoint{4.077060in}{3.684729in}}{\pgfqpoint{4.081178in}{3.688847in}}%
\pgfpathcurveto{\pgfqpoint{4.085296in}{3.692966in}}{\pgfqpoint{4.087610in}{3.698552in}}{\pgfqpoint{4.087610in}{3.704376in}}%
\pgfpathcurveto{\pgfqpoint{4.087610in}{3.710200in}}{\pgfqpoint{4.085296in}{3.715786in}}{\pgfqpoint{4.081178in}{3.719904in}}%
\pgfpathcurveto{\pgfqpoint{4.077060in}{3.724022in}}{\pgfqpoint{4.071474in}{3.726336in}}{\pgfqpoint{4.065650in}{3.726336in}}%
\pgfpathcurveto{\pgfqpoint{4.059826in}{3.726336in}}{\pgfqpoint{4.054240in}{3.724022in}}{\pgfqpoint{4.050121in}{3.719904in}}%
\pgfpathcurveto{\pgfqpoint{4.046003in}{3.715786in}}{\pgfqpoint{4.043689in}{3.710200in}}{\pgfqpoint{4.043689in}{3.704376in}}%
\pgfpathcurveto{\pgfqpoint{4.043689in}{3.698552in}}{\pgfqpoint{4.046003in}{3.692966in}}{\pgfqpoint{4.050121in}{3.688847in}}%
\pgfpathcurveto{\pgfqpoint{4.054240in}{3.684729in}}{\pgfqpoint{4.059826in}{3.682415in}}{\pgfqpoint{4.065650in}{3.682415in}}%
\pgfpathlineto{\pgfqpoint{4.065650in}{3.682415in}}%
\pgfpathclose%
\pgfusepath{stroke,fill}%
\end{pgfscope}%
\begin{pgfscope}%
\pgfpathrectangle{\pgfqpoint{1.414736in}{0.548769in}}{\pgfqpoint{3.497747in}{7.253186in}}%
\pgfusepath{clip}%
\pgfsetbuttcap%
\pgfsetroundjoin%
\definecolor{currentfill}{rgb}{0.993248,0.906157,0.143936}%
\pgfsetfillcolor{currentfill}%
\pgfsetlinewidth{1.003750pt}%
\definecolor{currentstroke}{rgb}{0.993248,0.906157,0.143936}%
\pgfsetstrokecolor{currentstroke}%
\pgfsetdash{}{0pt}%
\pgfpathmoveto{\pgfqpoint{3.491033in}{4.153402in}}%
\pgfpathcurveto{\pgfqpoint{3.496857in}{4.153402in}}{\pgfqpoint{3.502443in}{4.155715in}}{\pgfqpoint{3.506562in}{4.159834in}}%
\pgfpathcurveto{\pgfqpoint{3.510680in}{4.163952in}}{\pgfqpoint{3.512994in}{4.169538in}}{\pgfqpoint{3.512994in}{4.175362in}}%
\pgfpathcurveto{\pgfqpoint{3.512994in}{4.181186in}}{\pgfqpoint{3.510680in}{4.186772in}}{\pgfqpoint{3.506562in}{4.190890in}}%
\pgfpathcurveto{\pgfqpoint{3.502443in}{4.195008in}}{\pgfqpoint{3.496857in}{4.197322in}}{\pgfqpoint{3.491033in}{4.197322in}}%
\pgfpathcurveto{\pgfqpoint{3.485209in}{4.197322in}}{\pgfqpoint{3.479623in}{4.195008in}}{\pgfqpoint{3.475505in}{4.190890in}}%
\pgfpathcurveto{\pgfqpoint{3.471387in}{4.186772in}}{\pgfqpoint{3.469073in}{4.181186in}}{\pgfqpoint{3.469073in}{4.175362in}}%
\pgfpathcurveto{\pgfqpoint{3.469073in}{4.169538in}}{\pgfqpoint{3.471387in}{4.163952in}}{\pgfqpoint{3.475505in}{4.159834in}}%
\pgfpathcurveto{\pgfqpoint{3.479623in}{4.155715in}}{\pgfqpoint{3.485209in}{4.153402in}}{\pgfqpoint{3.491033in}{4.153402in}}%
\pgfpathlineto{\pgfqpoint{3.491033in}{4.153402in}}%
\pgfpathclose%
\pgfusepath{stroke,fill}%
\end{pgfscope}%
\begin{pgfscope}%
\pgfpathrectangle{\pgfqpoint{1.414736in}{0.548769in}}{\pgfqpoint{3.497747in}{7.253186in}}%
\pgfusepath{clip}%
\pgfsetbuttcap%
\pgfsetroundjoin%
\definecolor{currentfill}{rgb}{0.319809,0.770914,0.411152}%
\pgfsetfillcolor{currentfill}%
\pgfsetlinewidth{1.003750pt}%
\definecolor{currentstroke}{rgb}{0.319809,0.770914,0.411152}%
\pgfsetstrokecolor{currentstroke}%
\pgfsetdash{}{0pt}%
\pgfpathmoveto{\pgfqpoint{3.418738in}{4.624388in}}%
\pgfpathcurveto{\pgfqpoint{3.424562in}{4.624388in}}{\pgfqpoint{3.430148in}{4.626702in}}{\pgfqpoint{3.434266in}{4.630820in}}%
\pgfpathcurveto{\pgfqpoint{3.438384in}{4.634938in}}{\pgfqpoint{3.440698in}{4.640524in}}{\pgfqpoint{3.440698in}{4.646348in}}%
\pgfpathcurveto{\pgfqpoint{3.440698in}{4.652172in}}{\pgfqpoint{3.438384in}{4.657758in}}{\pgfqpoint{3.434266in}{4.661876in}}%
\pgfpathcurveto{\pgfqpoint{3.430148in}{4.665994in}}{\pgfqpoint{3.424562in}{4.668308in}}{\pgfqpoint{3.418738in}{4.668308in}}%
\pgfpathcurveto{\pgfqpoint{3.412914in}{4.668308in}}{\pgfqpoint{3.407328in}{4.665994in}}{\pgfqpoint{3.403209in}{4.661876in}}%
\pgfpathcurveto{\pgfqpoint{3.399091in}{4.657758in}}{\pgfqpoint{3.396777in}{4.652172in}}{\pgfqpoint{3.396777in}{4.646348in}}%
\pgfpathcurveto{\pgfqpoint{3.396777in}{4.640524in}}{\pgfqpoint{3.399091in}{4.634938in}}{\pgfqpoint{3.403209in}{4.630820in}}%
\pgfpathcurveto{\pgfqpoint{3.407328in}{4.626702in}}{\pgfqpoint{3.412914in}{4.624388in}}{\pgfqpoint{3.418738in}{4.624388in}}%
\pgfpathlineto{\pgfqpoint{3.418738in}{4.624388in}}%
\pgfpathclose%
\pgfusepath{stroke,fill}%
\end{pgfscope}%
\begin{pgfscope}%
\pgfpathrectangle{\pgfqpoint{1.414736in}{0.548769in}}{\pgfqpoint{3.497747in}{7.253186in}}%
\pgfusepath{clip}%
\pgfsetbuttcap%
\pgfsetroundjoin%
\definecolor{currentfill}{rgb}{0.265145,0.232956,0.516599}%
\pgfsetfillcolor{currentfill}%
\pgfsetlinewidth{1.003750pt}%
\definecolor{currentstroke}{rgb}{0.265145,0.232956,0.516599}%
\pgfsetstrokecolor{currentstroke}%
\pgfsetdash{}{0pt}%
\pgfpathmoveto{\pgfqpoint{3.343136in}{5.095374in}}%
\pgfpathcurveto{\pgfqpoint{3.348960in}{5.095374in}}{\pgfqpoint{3.354546in}{5.097688in}}{\pgfqpoint{3.358664in}{5.101806in}}%
\pgfpathcurveto{\pgfqpoint{3.362782in}{5.105924in}}{\pgfqpoint{3.365096in}{5.111510in}}{\pgfqpoint{3.365096in}{5.117334in}}%
\pgfpathcurveto{\pgfqpoint{3.365096in}{5.123158in}}{\pgfqpoint{3.362782in}{5.128744in}}{\pgfqpoint{3.358664in}{5.132862in}}%
\pgfpathcurveto{\pgfqpoint{3.354546in}{5.136980in}}{\pgfqpoint{3.348960in}{5.139294in}}{\pgfqpoint{3.343136in}{5.139294in}}%
\pgfpathcurveto{\pgfqpoint{3.337312in}{5.139294in}}{\pgfqpoint{3.331726in}{5.136980in}}{\pgfqpoint{3.327608in}{5.132862in}}%
\pgfpathcurveto{\pgfqpoint{3.323490in}{5.128744in}}{\pgfqpoint{3.321176in}{5.123158in}}{\pgfqpoint{3.321176in}{5.117334in}}%
\pgfpathcurveto{\pgfqpoint{3.321176in}{5.111510in}}{\pgfqpoint{3.323490in}{5.105924in}}{\pgfqpoint{3.327608in}{5.101806in}}%
\pgfpathcurveto{\pgfqpoint{3.331726in}{5.097688in}}{\pgfqpoint{3.337312in}{5.095374in}}{\pgfqpoint{3.343136in}{5.095374in}}%
\pgfpathlineto{\pgfqpoint{3.343136in}{5.095374in}}%
\pgfpathclose%
\pgfusepath{stroke,fill}%
\end{pgfscope}%
\begin{pgfscope}%
\pgfpathrectangle{\pgfqpoint{1.414736in}{0.548769in}}{\pgfqpoint{3.497747in}{7.253186in}}%
\pgfusepath{clip}%
\pgfsetbuttcap%
\pgfsetroundjoin%
\definecolor{currentfill}{rgb}{0.216210,0.351535,0.550627}%
\pgfsetfillcolor{currentfill}%
\pgfsetlinewidth{1.003750pt}%
\definecolor{currentstroke}{rgb}{0.216210,0.351535,0.550627}%
\pgfsetstrokecolor{currentstroke}%
\pgfsetdash{}{0pt}%
\pgfpathmoveto{\pgfqpoint{3.926425in}{5.566360in}}%
\pgfpathcurveto{\pgfqpoint{3.932249in}{5.566360in}}{\pgfqpoint{3.937835in}{5.568674in}}{\pgfqpoint{3.941953in}{5.572792in}}%
\pgfpathcurveto{\pgfqpoint{3.946071in}{5.576910in}}{\pgfqpoint{3.948385in}{5.582496in}}{\pgfqpoint{3.948385in}{5.588320in}}%
\pgfpathcurveto{\pgfqpoint{3.948385in}{5.594144in}}{\pgfqpoint{3.946071in}{5.599730in}}{\pgfqpoint{3.941953in}{5.603848in}}%
\pgfpathcurveto{\pgfqpoint{3.937835in}{5.607967in}}{\pgfqpoint{3.932249in}{5.610280in}}{\pgfqpoint{3.926425in}{5.610280in}}%
\pgfpathcurveto{\pgfqpoint{3.920601in}{5.610280in}}{\pgfqpoint{3.915015in}{5.607967in}}{\pgfqpoint{3.910897in}{5.603848in}}%
\pgfpathcurveto{\pgfqpoint{3.906779in}{5.599730in}}{\pgfqpoint{3.904465in}{5.594144in}}{\pgfqpoint{3.904465in}{5.588320in}}%
\pgfpathcurveto{\pgfqpoint{3.904465in}{5.582496in}}{\pgfqpoint{3.906779in}{5.576910in}}{\pgfqpoint{3.910897in}{5.572792in}}%
\pgfpathcurveto{\pgfqpoint{3.915015in}{5.568674in}}{\pgfqpoint{3.920601in}{5.566360in}}{\pgfqpoint{3.926425in}{5.566360in}}%
\pgfpathlineto{\pgfqpoint{3.926425in}{5.566360in}}%
\pgfpathclose%
\pgfusepath{stroke,fill}%
\end{pgfscope}%
\begin{pgfscope}%
\pgfpathrectangle{\pgfqpoint{1.414736in}{0.548769in}}{\pgfqpoint{3.497747in}{7.253186in}}%
\pgfusepath{clip}%
\pgfsetbuttcap%
\pgfsetroundjoin%
\definecolor{currentfill}{rgb}{0.119512,0.607464,0.540218}%
\pgfsetfillcolor{currentfill}%
\pgfsetlinewidth{1.003750pt}%
\definecolor{currentstroke}{rgb}{0.119512,0.607464,0.540218}%
\pgfsetstrokecolor{currentstroke}%
\pgfsetdash{}{0pt}%
\pgfpathmoveto{\pgfqpoint{3.745369in}{6.037346in}}%
\pgfpathcurveto{\pgfqpoint{3.751193in}{6.037346in}}{\pgfqpoint{3.756779in}{6.039660in}}{\pgfqpoint{3.760897in}{6.043778in}}%
\pgfpathcurveto{\pgfqpoint{3.765015in}{6.047896in}}{\pgfqpoint{3.767329in}{6.053482in}}{\pgfqpoint{3.767329in}{6.059306in}}%
\pgfpathcurveto{\pgfqpoint{3.767329in}{6.065130in}}{\pgfqpoint{3.765015in}{6.070716in}}{\pgfqpoint{3.760897in}{6.074835in}}%
\pgfpathcurveto{\pgfqpoint{3.756779in}{6.078953in}}{\pgfqpoint{3.751193in}{6.081267in}}{\pgfqpoint{3.745369in}{6.081267in}}%
\pgfpathcurveto{\pgfqpoint{3.739545in}{6.081267in}}{\pgfqpoint{3.733959in}{6.078953in}}{\pgfqpoint{3.729841in}{6.074835in}}%
\pgfpathcurveto{\pgfqpoint{3.725723in}{6.070716in}}{\pgfqpoint{3.723409in}{6.065130in}}{\pgfqpoint{3.723409in}{6.059306in}}%
\pgfpathcurveto{\pgfqpoint{3.723409in}{6.053482in}}{\pgfqpoint{3.725723in}{6.047896in}}{\pgfqpoint{3.729841in}{6.043778in}}%
\pgfpathcurveto{\pgfqpoint{3.733959in}{6.039660in}}{\pgfqpoint{3.739545in}{6.037346in}}{\pgfqpoint{3.745369in}{6.037346in}}%
\pgfpathlineto{\pgfqpoint{3.745369in}{6.037346in}}%
\pgfpathclose%
\pgfusepath{stroke,fill}%
\end{pgfscope}%
\begin{pgfscope}%
\pgfpathrectangle{\pgfqpoint{1.414736in}{0.548769in}}{\pgfqpoint{3.497747in}{7.253186in}}%
\pgfusepath{clip}%
\pgfsetbuttcap%
\pgfsetroundjoin%
\definecolor{currentfill}{rgb}{0.993248,0.906157,0.143936}%
\pgfsetfillcolor{currentfill}%
\pgfsetlinewidth{1.003750pt}%
\definecolor{currentstroke}{rgb}{0.993248,0.906157,0.143936}%
\pgfsetstrokecolor{currentstroke}%
\pgfsetdash{}{0pt}%
\pgfpathmoveto{\pgfqpoint{3.002858in}{6.508332in}}%
\pgfpathcurveto{\pgfqpoint{3.008682in}{6.508332in}}{\pgfqpoint{3.014268in}{6.510646in}}{\pgfqpoint{3.018386in}{6.514764in}}%
\pgfpathcurveto{\pgfqpoint{3.022504in}{6.518882in}}{\pgfqpoint{3.024818in}{6.524468in}}{\pgfqpoint{3.024818in}{6.530292in}}%
\pgfpathcurveto{\pgfqpoint{3.024818in}{6.536116in}}{\pgfqpoint{3.022504in}{6.541702in}}{\pgfqpoint{3.018386in}{6.545821in}}%
\pgfpathcurveto{\pgfqpoint{3.014268in}{6.549939in}}{\pgfqpoint{3.008682in}{6.552253in}}{\pgfqpoint{3.002858in}{6.552253in}}%
\pgfpathcurveto{\pgfqpoint{2.997034in}{6.552253in}}{\pgfqpoint{2.991448in}{6.549939in}}{\pgfqpoint{2.987329in}{6.545821in}}%
\pgfpathcurveto{\pgfqpoint{2.983211in}{6.541702in}}{\pgfqpoint{2.980897in}{6.536116in}}{\pgfqpoint{2.980897in}{6.530292in}}%
\pgfpathcurveto{\pgfqpoint{2.980897in}{6.524468in}}{\pgfqpoint{2.983211in}{6.518882in}}{\pgfqpoint{2.987329in}{6.514764in}}%
\pgfpathcurveto{\pgfqpoint{2.991448in}{6.510646in}}{\pgfqpoint{2.997034in}{6.508332in}}{\pgfqpoint{3.002858in}{6.508332in}}%
\pgfpathlineto{\pgfqpoint{3.002858in}{6.508332in}}%
\pgfpathclose%
\pgfusepath{stroke,fill}%
\end{pgfscope}%
\begin{pgfscope}%
\pgfpathrectangle{\pgfqpoint{1.414736in}{0.548769in}}{\pgfqpoint{3.497747in}{7.253186in}}%
\pgfusepath{clip}%
\pgfsetbuttcap%
\pgfsetroundjoin%
\definecolor{currentfill}{rgb}{0.993248,0.906157,0.143936}%
\pgfsetfillcolor{currentfill}%
\pgfsetlinewidth{1.003750pt}%
\definecolor{currentstroke}{rgb}{0.993248,0.906157,0.143936}%
\pgfsetstrokecolor{currentstroke}%
\pgfsetdash{}{0pt}%
\pgfpathmoveto{\pgfqpoint{3.223582in}{6.979318in}}%
\pgfpathcurveto{\pgfqpoint{3.229406in}{6.979318in}}{\pgfqpoint{3.234992in}{6.981632in}}{\pgfqpoint{3.239110in}{6.985750in}}%
\pgfpathcurveto{\pgfqpoint{3.243228in}{6.989868in}}{\pgfqpoint{3.245542in}{6.995455in}}{\pgfqpoint{3.245542in}{7.001278in}}%
\pgfpathcurveto{\pgfqpoint{3.245542in}{7.007102in}}{\pgfqpoint{3.243228in}{7.012689in}}{\pgfqpoint{3.239110in}{7.016807in}}%
\pgfpathcurveto{\pgfqpoint{3.234992in}{7.020925in}}{\pgfqpoint{3.229406in}{7.023239in}}{\pgfqpoint{3.223582in}{7.023239in}}%
\pgfpathcurveto{\pgfqpoint{3.217758in}{7.023239in}}{\pgfqpoint{3.212172in}{7.020925in}}{\pgfqpoint{3.208053in}{7.016807in}}%
\pgfpathcurveto{\pgfqpoint{3.203935in}{7.012689in}}{\pgfqpoint{3.201621in}{7.007102in}}{\pgfqpoint{3.201621in}{7.001278in}}%
\pgfpathcurveto{\pgfqpoint{3.201621in}{6.995455in}}{\pgfqpoint{3.203935in}{6.989868in}}{\pgfqpoint{3.208053in}{6.985750in}}%
\pgfpathcurveto{\pgfqpoint{3.212172in}{6.981632in}}{\pgfqpoint{3.217758in}{6.979318in}}{\pgfqpoint{3.223582in}{6.979318in}}%
\pgfpathlineto{\pgfqpoint{3.223582in}{6.979318in}}%
\pgfpathclose%
\pgfusepath{stroke,fill}%
\end{pgfscope}%
\begin{pgfscope}%
\pgfpathrectangle{\pgfqpoint{1.414736in}{0.548769in}}{\pgfqpoint{3.497747in}{7.253186in}}%
\pgfusepath{clip}%
\pgfsetbuttcap%
\pgfsetroundjoin%
\definecolor{currentfill}{rgb}{0.288921,0.758394,0.428426}%
\pgfsetfillcolor{currentfill}%
\pgfsetlinewidth{1.003750pt}%
\definecolor{currentstroke}{rgb}{0.288921,0.758394,0.428426}%
\pgfsetstrokecolor{currentstroke}%
\pgfsetdash{}{0pt}%
\pgfpathmoveto{\pgfqpoint{4.327101in}{7.450304in}}%
\pgfpathcurveto{\pgfqpoint{4.332925in}{7.450304in}}{\pgfqpoint{4.338511in}{7.452618in}}{\pgfqpoint{4.342629in}{7.456736in}}%
\pgfpathcurveto{\pgfqpoint{4.346748in}{7.460854in}}{\pgfqpoint{4.349061in}{7.466441in}}{\pgfqpoint{4.349061in}{7.472265in}}%
\pgfpathcurveto{\pgfqpoint{4.349061in}{7.478089in}}{\pgfqpoint{4.346748in}{7.483675in}}{\pgfqpoint{4.342629in}{7.487793in}}%
\pgfpathcurveto{\pgfqpoint{4.338511in}{7.491911in}}{\pgfqpoint{4.332925in}{7.494225in}}{\pgfqpoint{4.327101in}{7.494225in}}%
\pgfpathcurveto{\pgfqpoint{4.321277in}{7.494225in}}{\pgfqpoint{4.315691in}{7.491911in}}{\pgfqpoint{4.311573in}{7.487793in}}%
\pgfpathcurveto{\pgfqpoint{4.307455in}{7.483675in}}{\pgfqpoint{4.305141in}{7.478089in}}{\pgfqpoint{4.305141in}{7.472265in}}%
\pgfpathcurveto{\pgfqpoint{4.305141in}{7.466441in}}{\pgfqpoint{4.307455in}{7.460854in}}{\pgfqpoint{4.311573in}{7.456736in}}%
\pgfpathcurveto{\pgfqpoint{4.315691in}{7.452618in}}{\pgfqpoint{4.321277in}{7.450304in}}{\pgfqpoint{4.327101in}{7.450304in}}%
\pgfpathlineto{\pgfqpoint{4.327101in}{7.450304in}}%
\pgfpathclose%
\pgfusepath{stroke,fill}%
\end{pgfscope}%
\begin{pgfscope}%
\pgfpathrectangle{\pgfqpoint{1.414736in}{0.548769in}}{\pgfqpoint{3.497747in}{7.253186in}}%
\pgfusepath{clip}%
\pgfsetbuttcap%
\pgfsetroundjoin%
\definecolor{currentfill}{rgb}{0.783315,0.879285,0.125405}%
\pgfsetfillcolor{currentfill}%
\pgfsetlinewidth{1.003750pt}%
\definecolor{currentstroke}{rgb}{0.783315,0.879285,0.125405}%
\pgfsetstrokecolor{currentstroke}%
\pgfsetdash{}{0pt}%
\pgfpathmoveto{\pgfqpoint{3.387642in}{0.856499in}}%
\pgfpathcurveto{\pgfqpoint{3.393466in}{0.856499in}}{\pgfqpoint{3.399053in}{0.858813in}}{\pgfqpoint{3.403171in}{0.862931in}}%
\pgfpathcurveto{\pgfqpoint{3.407289in}{0.867049in}}{\pgfqpoint{3.409603in}{0.872635in}}{\pgfqpoint{3.409603in}{0.878459in}}%
\pgfpathcurveto{\pgfqpoint{3.409603in}{0.884283in}}{\pgfqpoint{3.407289in}{0.889869in}}{\pgfqpoint{3.403171in}{0.893987in}}%
\pgfpathcurveto{\pgfqpoint{3.399053in}{0.898105in}}{\pgfqpoint{3.393466in}{0.900419in}}{\pgfqpoint{3.387642in}{0.900419in}}%
\pgfpathcurveto{\pgfqpoint{3.381819in}{0.900419in}}{\pgfqpoint{3.376232in}{0.898105in}}{\pgfqpoint{3.372114in}{0.893987in}}%
\pgfpathcurveto{\pgfqpoint{3.367996in}{0.889869in}}{\pgfqpoint{3.365682in}{0.884283in}}{\pgfqpoint{3.365682in}{0.878459in}}%
\pgfpathcurveto{\pgfqpoint{3.365682in}{0.872635in}}{\pgfqpoint{3.367996in}{0.867049in}}{\pgfqpoint{3.372114in}{0.862931in}}%
\pgfpathcurveto{\pgfqpoint{3.376232in}{0.858813in}}{\pgfqpoint{3.381819in}{0.856499in}}{\pgfqpoint{3.387642in}{0.856499in}}%
\pgfpathlineto{\pgfqpoint{3.387642in}{0.856499in}}%
\pgfpathclose%
\pgfusepath{stroke,fill}%
\end{pgfscope}%
\begin{pgfscope}%
\pgfpathrectangle{\pgfqpoint{1.414736in}{0.548769in}}{\pgfqpoint{3.497747in}{7.253186in}}%
\pgfusepath{clip}%
\pgfsetbuttcap%
\pgfsetroundjoin%
\definecolor{currentfill}{rgb}{0.166383,0.690856,0.496502}%
\pgfsetfillcolor{currentfill}%
\pgfsetlinewidth{1.003750pt}%
\definecolor{currentstroke}{rgb}{0.166383,0.690856,0.496502}%
\pgfsetstrokecolor{currentstroke}%
\pgfsetdash{}{0pt}%
\pgfpathmoveto{\pgfqpoint{3.979453in}{1.327485in}}%
\pgfpathcurveto{\pgfqpoint{3.985277in}{1.327485in}}{\pgfqpoint{3.990864in}{1.329799in}}{\pgfqpoint{3.994982in}{1.333917in}}%
\pgfpathcurveto{\pgfqpoint{3.999100in}{1.338035in}}{\pgfqpoint{4.001414in}{1.343621in}}{\pgfqpoint{4.001414in}{1.349445in}}%
\pgfpathcurveto{\pgfqpoint{4.001414in}{1.355269in}}{\pgfqpoint{3.999100in}{1.360855in}}{\pgfqpoint{3.994982in}{1.364973in}}%
\pgfpathcurveto{\pgfqpoint{3.990864in}{1.369092in}}{\pgfqpoint{3.985277in}{1.371405in}}{\pgfqpoint{3.979453in}{1.371405in}}%
\pgfpathcurveto{\pgfqpoint{3.973629in}{1.371405in}}{\pgfqpoint{3.968043in}{1.369092in}}{\pgfqpoint{3.963925in}{1.364973in}}%
\pgfpathcurveto{\pgfqpoint{3.959807in}{1.360855in}}{\pgfqpoint{3.957493in}{1.355269in}}{\pgfqpoint{3.957493in}{1.349445in}}%
\pgfpathcurveto{\pgfqpoint{3.957493in}{1.343621in}}{\pgfqpoint{3.959807in}{1.338035in}}{\pgfqpoint{3.963925in}{1.333917in}}%
\pgfpathcurveto{\pgfqpoint{3.968043in}{1.329799in}}{\pgfqpoint{3.973629in}{1.327485in}}{\pgfqpoint{3.979453in}{1.327485in}}%
\pgfpathlineto{\pgfqpoint{3.979453in}{1.327485in}}%
\pgfpathclose%
\pgfusepath{stroke,fill}%
\end{pgfscope}%
\begin{pgfscope}%
\pgfpathrectangle{\pgfqpoint{1.414736in}{0.548769in}}{\pgfqpoint{3.497747in}{7.253186in}}%
\pgfusepath{clip}%
\pgfsetbuttcap%
\pgfsetroundjoin%
\definecolor{currentfill}{rgb}{0.804182,0.882046,0.114965}%
\pgfsetfillcolor{currentfill}%
\pgfsetlinewidth{1.003750pt}%
\definecolor{currentstroke}{rgb}{0.804182,0.882046,0.114965}%
\pgfsetstrokecolor{currentstroke}%
\pgfsetdash{}{0pt}%
\pgfpathmoveto{\pgfqpoint{3.755566in}{1.798471in}}%
\pgfpathcurveto{\pgfqpoint{3.761390in}{1.798471in}}{\pgfqpoint{3.766976in}{1.800785in}}{\pgfqpoint{3.771094in}{1.804903in}}%
\pgfpathcurveto{\pgfqpoint{3.775213in}{1.809021in}}{\pgfqpoint{3.777526in}{1.814607in}}{\pgfqpoint{3.777526in}{1.820431in}}%
\pgfpathcurveto{\pgfqpoint{3.777526in}{1.826255in}}{\pgfqpoint{3.775213in}{1.831841in}}{\pgfqpoint{3.771094in}{1.835960in}}%
\pgfpathcurveto{\pgfqpoint{3.766976in}{1.840078in}}{\pgfqpoint{3.761390in}{1.842392in}}{\pgfqpoint{3.755566in}{1.842392in}}%
\pgfpathcurveto{\pgfqpoint{3.749742in}{1.842392in}}{\pgfqpoint{3.744156in}{1.840078in}}{\pgfqpoint{3.740038in}{1.835960in}}%
\pgfpathcurveto{\pgfqpoint{3.735920in}{1.831841in}}{\pgfqpoint{3.733606in}{1.826255in}}{\pgfqpoint{3.733606in}{1.820431in}}%
\pgfpathcurveto{\pgfqpoint{3.733606in}{1.814607in}}{\pgfqpoint{3.735920in}{1.809021in}}{\pgfqpoint{3.740038in}{1.804903in}}%
\pgfpathcurveto{\pgfqpoint{3.744156in}{1.800785in}}{\pgfqpoint{3.749742in}{1.798471in}}{\pgfqpoint{3.755566in}{1.798471in}}%
\pgfpathlineto{\pgfqpoint{3.755566in}{1.798471in}}%
\pgfpathclose%
\pgfusepath{stroke,fill}%
\end{pgfscope}%
\begin{pgfscope}%
\pgfpathrectangle{\pgfqpoint{1.414736in}{0.548769in}}{\pgfqpoint{3.497747in}{7.253186in}}%
\pgfusepath{clip}%
\pgfsetbuttcap%
\pgfsetroundjoin%
\definecolor{currentfill}{rgb}{0.595839,0.848717,0.243329}%
\pgfsetfillcolor{currentfill}%
\pgfsetlinewidth{1.003750pt}%
\definecolor{currentstroke}{rgb}{0.595839,0.848717,0.243329}%
\pgfsetstrokecolor{currentstroke}%
\pgfsetdash{}{0pt}%
\pgfpathmoveto{\pgfqpoint{3.630970in}{2.269457in}}%
\pgfpathcurveto{\pgfqpoint{3.636794in}{2.269457in}}{\pgfqpoint{3.642380in}{2.271771in}}{\pgfqpoint{3.646498in}{2.275889in}}%
\pgfpathcurveto{\pgfqpoint{3.650617in}{2.280007in}}{\pgfqpoint{3.652930in}{2.285593in}}{\pgfqpoint{3.652930in}{2.291417in}}%
\pgfpathcurveto{\pgfqpoint{3.652930in}{2.297241in}}{\pgfqpoint{3.650617in}{2.302828in}}{\pgfqpoint{3.646498in}{2.306946in}}%
\pgfpathcurveto{\pgfqpoint{3.642380in}{2.311064in}}{\pgfqpoint{3.636794in}{2.313378in}}{\pgfqpoint{3.630970in}{2.313378in}}%
\pgfpathcurveto{\pgfqpoint{3.625146in}{2.313378in}}{\pgfqpoint{3.619560in}{2.311064in}}{\pgfqpoint{3.615442in}{2.306946in}}%
\pgfpathcurveto{\pgfqpoint{3.611324in}{2.302828in}}{\pgfqpoint{3.609010in}{2.297241in}}{\pgfqpoint{3.609010in}{2.291417in}}%
\pgfpathcurveto{\pgfqpoint{3.609010in}{2.285593in}}{\pgfqpoint{3.611324in}{2.280007in}}{\pgfqpoint{3.615442in}{2.275889in}}%
\pgfpathcurveto{\pgfqpoint{3.619560in}{2.271771in}}{\pgfqpoint{3.625146in}{2.269457in}}{\pgfqpoint{3.630970in}{2.269457in}}%
\pgfpathlineto{\pgfqpoint{3.630970in}{2.269457in}}%
\pgfpathclose%
\pgfusepath{stroke,fill}%
\end{pgfscope}%
\begin{pgfscope}%
\pgfpathrectangle{\pgfqpoint{1.414736in}{0.548769in}}{\pgfqpoint{3.497747in}{7.253186in}}%
\pgfusepath{clip}%
\pgfsetbuttcap%
\pgfsetroundjoin%
\definecolor{currentfill}{rgb}{0.440137,0.811138,0.340967}%
\pgfsetfillcolor{currentfill}%
\pgfsetlinewidth{1.003750pt}%
\definecolor{currentstroke}{rgb}{0.440137,0.811138,0.340967}%
\pgfsetstrokecolor{currentstroke}%
\pgfsetdash{}{0pt}%
\pgfpathmoveto{\pgfqpoint{3.671422in}{2.740443in}}%
\pgfpathcurveto{\pgfqpoint{3.677246in}{2.740443in}}{\pgfqpoint{3.682832in}{2.742757in}}{\pgfqpoint{3.686950in}{2.746875in}}%
\pgfpathcurveto{\pgfqpoint{3.691068in}{2.750993in}}{\pgfqpoint{3.693382in}{2.756580in}}{\pgfqpoint{3.693382in}{2.762404in}}%
\pgfpathcurveto{\pgfqpoint{3.693382in}{2.768227in}}{\pgfqpoint{3.691068in}{2.773814in}}{\pgfqpoint{3.686950in}{2.777932in}}%
\pgfpathcurveto{\pgfqpoint{3.682832in}{2.782050in}}{\pgfqpoint{3.677246in}{2.784364in}}{\pgfqpoint{3.671422in}{2.784364in}}%
\pgfpathcurveto{\pgfqpoint{3.665598in}{2.784364in}}{\pgfqpoint{3.660012in}{2.782050in}}{\pgfqpoint{3.655893in}{2.777932in}}%
\pgfpathcurveto{\pgfqpoint{3.651775in}{2.773814in}}{\pgfqpoint{3.649461in}{2.768227in}}{\pgfqpoint{3.649461in}{2.762404in}}%
\pgfpathcurveto{\pgfqpoint{3.649461in}{2.756580in}}{\pgfqpoint{3.651775in}{2.750993in}}{\pgfqpoint{3.655893in}{2.746875in}}%
\pgfpathcurveto{\pgfqpoint{3.660012in}{2.742757in}}{\pgfqpoint{3.665598in}{2.740443in}}{\pgfqpoint{3.671422in}{2.740443in}}%
\pgfpathlineto{\pgfqpoint{3.671422in}{2.740443in}}%
\pgfpathclose%
\pgfusepath{stroke,fill}%
\end{pgfscope}%
\begin{pgfscope}%
\pgfpathrectangle{\pgfqpoint{1.414736in}{0.548769in}}{\pgfqpoint{3.497747in}{7.253186in}}%
\pgfusepath{clip}%
\pgfsetbuttcap%
\pgfsetroundjoin%
\definecolor{currentfill}{rgb}{0.352360,0.783011,0.392636}%
\pgfsetfillcolor{currentfill}%
\pgfsetlinewidth{1.003750pt}%
\definecolor{currentstroke}{rgb}{0.352360,0.783011,0.392636}%
\pgfsetstrokecolor{currentstroke}%
\pgfsetdash{}{0pt}%
\pgfpathmoveto{\pgfqpoint{4.134245in}{3.211429in}}%
\pgfpathcurveto{\pgfqpoint{4.140069in}{3.211429in}}{\pgfqpoint{4.145655in}{3.213743in}}{\pgfqpoint{4.149773in}{3.217861in}}%
\pgfpathcurveto{\pgfqpoint{4.153891in}{3.221980in}}{\pgfqpoint{4.156205in}{3.227566in}}{\pgfqpoint{4.156205in}{3.233390in}}%
\pgfpathcurveto{\pgfqpoint{4.156205in}{3.239214in}}{\pgfqpoint{4.153891in}{3.244800in}}{\pgfqpoint{4.149773in}{3.248918in}}%
\pgfpathcurveto{\pgfqpoint{4.145655in}{3.253036in}}{\pgfqpoint{4.140069in}{3.255350in}}{\pgfqpoint{4.134245in}{3.255350in}}%
\pgfpathcurveto{\pgfqpoint{4.128421in}{3.255350in}}{\pgfqpoint{4.122835in}{3.253036in}}{\pgfqpoint{4.118717in}{3.248918in}}%
\pgfpathcurveto{\pgfqpoint{4.114598in}{3.244800in}}{\pgfqpoint{4.112285in}{3.239214in}}{\pgfqpoint{4.112285in}{3.233390in}}%
\pgfpathcurveto{\pgfqpoint{4.112285in}{3.227566in}}{\pgfqpoint{4.114598in}{3.221980in}}{\pgfqpoint{4.118717in}{3.217861in}}%
\pgfpathcurveto{\pgfqpoint{4.122835in}{3.213743in}}{\pgfqpoint{4.128421in}{3.211429in}}{\pgfqpoint{4.134245in}{3.211429in}}%
\pgfpathlineto{\pgfqpoint{4.134245in}{3.211429in}}%
\pgfpathclose%
\pgfusepath{stroke,fill}%
\end{pgfscope}%
\begin{pgfscope}%
\pgfpathrectangle{\pgfqpoint{1.414736in}{0.548769in}}{\pgfqpoint{3.497747in}{7.253186in}}%
\pgfusepath{clip}%
\pgfsetbuttcap%
\pgfsetroundjoin%
\definecolor{currentfill}{rgb}{0.993248,0.906157,0.143936}%
\pgfsetfillcolor{currentfill}%
\pgfsetlinewidth{1.003750pt}%
\definecolor{currentstroke}{rgb}{0.993248,0.906157,0.143936}%
\pgfsetstrokecolor{currentstroke}%
\pgfsetdash{}{0pt}%
\pgfpathmoveto{\pgfqpoint{3.545800in}{3.682415in}}%
\pgfpathcurveto{\pgfqpoint{3.551624in}{3.682415in}}{\pgfqpoint{3.557211in}{3.684729in}}{\pgfqpoint{3.561329in}{3.688847in}}%
\pgfpathcurveto{\pgfqpoint{3.565447in}{3.692966in}}{\pgfqpoint{3.567761in}{3.698552in}}{\pgfqpoint{3.567761in}{3.704376in}}%
\pgfpathcurveto{\pgfqpoint{3.567761in}{3.710200in}}{\pgfqpoint{3.565447in}{3.715786in}}{\pgfqpoint{3.561329in}{3.719904in}}%
\pgfpathcurveto{\pgfqpoint{3.557211in}{3.724022in}}{\pgfqpoint{3.551624in}{3.726336in}}{\pgfqpoint{3.545800in}{3.726336in}}%
\pgfpathcurveto{\pgfqpoint{3.539977in}{3.726336in}}{\pgfqpoint{3.534390in}{3.724022in}}{\pgfqpoint{3.530272in}{3.719904in}}%
\pgfpathcurveto{\pgfqpoint{3.526154in}{3.715786in}}{\pgfqpoint{3.523840in}{3.710200in}}{\pgfqpoint{3.523840in}{3.704376in}}%
\pgfpathcurveto{\pgfqpoint{3.523840in}{3.698552in}}{\pgfqpoint{3.526154in}{3.692966in}}{\pgfqpoint{3.530272in}{3.688847in}}%
\pgfpathcurveto{\pgfqpoint{3.534390in}{3.684729in}}{\pgfqpoint{3.539977in}{3.682415in}}{\pgfqpoint{3.545800in}{3.682415in}}%
\pgfpathlineto{\pgfqpoint{3.545800in}{3.682415in}}%
\pgfpathclose%
\pgfusepath{stroke,fill}%
\end{pgfscope}%
\begin{pgfscope}%
\pgfpathrectangle{\pgfqpoint{1.414736in}{0.548769in}}{\pgfqpoint{3.497747in}{7.253186in}}%
\pgfusepath{clip}%
\pgfsetbuttcap%
\pgfsetroundjoin%
\definecolor{currentfill}{rgb}{0.993248,0.906157,0.143936}%
\pgfsetfillcolor{currentfill}%
\pgfsetlinewidth{1.003750pt}%
\definecolor{currentstroke}{rgb}{0.993248,0.906157,0.143936}%
\pgfsetstrokecolor{currentstroke}%
\pgfsetdash{}{0pt}%
\pgfpathmoveto{\pgfqpoint{3.761811in}{4.153402in}}%
\pgfpathcurveto{\pgfqpoint{3.767635in}{4.153402in}}{\pgfqpoint{3.773221in}{4.155715in}}{\pgfqpoint{3.777339in}{4.159834in}}%
\pgfpathcurveto{\pgfqpoint{3.781457in}{4.163952in}}{\pgfqpoint{3.783771in}{4.169538in}}{\pgfqpoint{3.783771in}{4.175362in}}%
\pgfpathcurveto{\pgfqpoint{3.783771in}{4.181186in}}{\pgfqpoint{3.781457in}{4.186772in}}{\pgfqpoint{3.777339in}{4.190890in}}%
\pgfpathcurveto{\pgfqpoint{3.773221in}{4.195008in}}{\pgfqpoint{3.767635in}{4.197322in}}{\pgfqpoint{3.761811in}{4.197322in}}%
\pgfpathcurveto{\pgfqpoint{3.755987in}{4.197322in}}{\pgfqpoint{3.750401in}{4.195008in}}{\pgfqpoint{3.746283in}{4.190890in}}%
\pgfpathcurveto{\pgfqpoint{3.742165in}{4.186772in}}{\pgfqpoint{3.739851in}{4.181186in}}{\pgfqpoint{3.739851in}{4.175362in}}%
\pgfpathcurveto{\pgfqpoint{3.739851in}{4.169538in}}{\pgfqpoint{3.742165in}{4.163952in}}{\pgfqpoint{3.746283in}{4.159834in}}%
\pgfpathcurveto{\pgfqpoint{3.750401in}{4.155715in}}{\pgfqpoint{3.755987in}{4.153402in}}{\pgfqpoint{3.761811in}{4.153402in}}%
\pgfpathlineto{\pgfqpoint{3.761811in}{4.153402in}}%
\pgfpathclose%
\pgfusepath{stroke,fill}%
\end{pgfscope}%
\begin{pgfscope}%
\pgfpathrectangle{\pgfqpoint{1.414736in}{0.548769in}}{\pgfqpoint{3.497747in}{7.253186in}}%
\pgfusepath{clip}%
\pgfsetbuttcap%
\pgfsetroundjoin%
\definecolor{currentfill}{rgb}{0.595839,0.848717,0.243329}%
\pgfsetfillcolor{currentfill}%
\pgfsetlinewidth{1.003750pt}%
\definecolor{currentstroke}{rgb}{0.595839,0.848717,0.243329}%
\pgfsetstrokecolor{currentstroke}%
\pgfsetdash{}{0pt}%
\pgfpathmoveto{\pgfqpoint{3.057757in}{4.624388in}}%
\pgfpathcurveto{\pgfqpoint{3.063581in}{4.624388in}}{\pgfqpoint{3.069168in}{4.626702in}}{\pgfqpoint{3.073286in}{4.630820in}}%
\pgfpathcurveto{\pgfqpoint{3.077404in}{4.634938in}}{\pgfqpoint{3.079718in}{4.640524in}}{\pgfqpoint{3.079718in}{4.646348in}}%
\pgfpathcurveto{\pgfqpoint{3.079718in}{4.652172in}}{\pgfqpoint{3.077404in}{4.657758in}}{\pgfqpoint{3.073286in}{4.661876in}}%
\pgfpathcurveto{\pgfqpoint{3.069168in}{4.665994in}}{\pgfqpoint{3.063581in}{4.668308in}}{\pgfqpoint{3.057757in}{4.668308in}}%
\pgfpathcurveto{\pgfqpoint{3.051934in}{4.668308in}}{\pgfqpoint{3.046347in}{4.665994in}}{\pgfqpoint{3.042229in}{4.661876in}}%
\pgfpathcurveto{\pgfqpoint{3.038111in}{4.657758in}}{\pgfqpoint{3.035797in}{4.652172in}}{\pgfqpoint{3.035797in}{4.646348in}}%
\pgfpathcurveto{\pgfqpoint{3.035797in}{4.640524in}}{\pgfqpoint{3.038111in}{4.634938in}}{\pgfqpoint{3.042229in}{4.630820in}}%
\pgfpathcurveto{\pgfqpoint{3.046347in}{4.626702in}}{\pgfqpoint{3.051934in}{4.624388in}}{\pgfqpoint{3.057757in}{4.624388in}}%
\pgfpathlineto{\pgfqpoint{3.057757in}{4.624388in}}%
\pgfpathclose%
\pgfusepath{stroke,fill}%
\end{pgfscope}%
\begin{pgfscope}%
\pgfpathrectangle{\pgfqpoint{1.414736in}{0.548769in}}{\pgfqpoint{3.497747in}{7.253186in}}%
\pgfusepath{clip}%
\pgfsetbuttcap%
\pgfsetroundjoin%
\definecolor{currentfill}{rgb}{0.280255,0.165693,0.476498}%
\pgfsetfillcolor{currentfill}%
\pgfsetlinewidth{1.003750pt}%
\definecolor{currentstroke}{rgb}{0.280255,0.165693,0.476498}%
\pgfsetstrokecolor{currentstroke}%
\pgfsetdash{}{0pt}%
\pgfpathmoveto{\pgfqpoint{2.719126in}{5.095374in}}%
\pgfpathcurveto{\pgfqpoint{2.724950in}{5.095374in}}{\pgfqpoint{2.730537in}{5.097688in}}{\pgfqpoint{2.734655in}{5.101806in}}%
\pgfpathcurveto{\pgfqpoint{2.738773in}{5.105924in}}{\pgfqpoint{2.741087in}{5.111510in}}{\pgfqpoint{2.741087in}{5.117334in}}%
\pgfpathcurveto{\pgfqpoint{2.741087in}{5.123158in}}{\pgfqpoint{2.738773in}{5.128744in}}{\pgfqpoint{2.734655in}{5.132862in}}%
\pgfpathcurveto{\pgfqpoint{2.730537in}{5.136980in}}{\pgfqpoint{2.724950in}{5.139294in}}{\pgfqpoint{2.719126in}{5.139294in}}%
\pgfpathcurveto{\pgfqpoint{2.713303in}{5.139294in}}{\pgfqpoint{2.707716in}{5.136980in}}{\pgfqpoint{2.703598in}{5.132862in}}%
\pgfpathcurveto{\pgfqpoint{2.699480in}{5.128744in}}{\pgfqpoint{2.697166in}{5.123158in}}{\pgfqpoint{2.697166in}{5.117334in}}%
\pgfpathcurveto{\pgfqpoint{2.697166in}{5.111510in}}{\pgfqpoint{2.699480in}{5.105924in}}{\pgfqpoint{2.703598in}{5.101806in}}%
\pgfpathcurveto{\pgfqpoint{2.707716in}{5.097688in}}{\pgfqpoint{2.713303in}{5.095374in}}{\pgfqpoint{2.719126in}{5.095374in}}%
\pgfpathlineto{\pgfqpoint{2.719126in}{5.095374in}}%
\pgfpathclose%
\pgfusepath{stroke,fill}%
\end{pgfscope}%
\begin{pgfscope}%
\pgfpathrectangle{\pgfqpoint{1.414736in}{0.548769in}}{\pgfqpoint{3.497747in}{7.253186in}}%
\pgfusepath{clip}%
\pgfsetbuttcap%
\pgfsetroundjoin%
\definecolor{currentfill}{rgb}{0.218130,0.347432,0.550038}%
\pgfsetfillcolor{currentfill}%
\pgfsetlinewidth{1.003750pt}%
\definecolor{currentstroke}{rgb}{0.218130,0.347432,0.550038}%
\pgfsetstrokecolor{currentstroke}%
\pgfsetdash{}{0pt}%
\pgfpathmoveto{\pgfqpoint{4.107096in}{5.566360in}}%
\pgfpathcurveto{\pgfqpoint{4.112920in}{5.566360in}}{\pgfqpoint{4.118506in}{5.568674in}}{\pgfqpoint{4.122624in}{5.572792in}}%
\pgfpathcurveto{\pgfqpoint{4.126742in}{5.576910in}}{\pgfqpoint{4.129056in}{5.582496in}}{\pgfqpoint{4.129056in}{5.588320in}}%
\pgfpathcurveto{\pgfqpoint{4.129056in}{5.594144in}}{\pgfqpoint{4.126742in}{5.599730in}}{\pgfqpoint{4.122624in}{5.603848in}}%
\pgfpathcurveto{\pgfqpoint{4.118506in}{5.607967in}}{\pgfqpoint{4.112920in}{5.610280in}}{\pgfqpoint{4.107096in}{5.610280in}}%
\pgfpathcurveto{\pgfqpoint{4.101272in}{5.610280in}}{\pgfqpoint{4.095686in}{5.607967in}}{\pgfqpoint{4.091567in}{5.603848in}}%
\pgfpathcurveto{\pgfqpoint{4.087449in}{5.599730in}}{\pgfqpoint{4.085135in}{5.594144in}}{\pgfqpoint{4.085135in}{5.588320in}}%
\pgfpathcurveto{\pgfqpoint{4.085135in}{5.582496in}}{\pgfqpoint{4.087449in}{5.576910in}}{\pgfqpoint{4.091567in}{5.572792in}}%
\pgfpathcurveto{\pgfqpoint{4.095686in}{5.568674in}}{\pgfqpoint{4.101272in}{5.566360in}}{\pgfqpoint{4.107096in}{5.566360in}}%
\pgfpathlineto{\pgfqpoint{4.107096in}{5.566360in}}%
\pgfpathclose%
\pgfusepath{stroke,fill}%
\end{pgfscope}%
\begin{pgfscope}%
\pgfpathrectangle{\pgfqpoint{1.414736in}{0.548769in}}{\pgfqpoint{3.497747in}{7.253186in}}%
\pgfusepath{clip}%
\pgfsetbuttcap%
\pgfsetroundjoin%
\definecolor{currentfill}{rgb}{0.126326,0.644107,0.525311}%
\pgfsetfillcolor{currentfill}%
\pgfsetlinewidth{1.003750pt}%
\definecolor{currentstroke}{rgb}{0.126326,0.644107,0.525311}%
\pgfsetstrokecolor{currentstroke}%
\pgfsetdash{}{0pt}%
\pgfpathmoveto{\pgfqpoint{4.134948in}{6.037346in}}%
\pgfpathcurveto{\pgfqpoint{4.140772in}{6.037346in}}{\pgfqpoint{4.146358in}{6.039660in}}{\pgfqpoint{4.150477in}{6.043778in}}%
\pgfpathcurveto{\pgfqpoint{4.154595in}{6.047896in}}{\pgfqpoint{4.156909in}{6.053482in}}{\pgfqpoint{4.156909in}{6.059306in}}%
\pgfpathcurveto{\pgfqpoint{4.156909in}{6.065130in}}{\pgfqpoint{4.154595in}{6.070716in}}{\pgfqpoint{4.150477in}{6.074835in}}%
\pgfpathcurveto{\pgfqpoint{4.146358in}{6.078953in}}{\pgfqpoint{4.140772in}{6.081267in}}{\pgfqpoint{4.134948in}{6.081267in}}%
\pgfpathcurveto{\pgfqpoint{4.129124in}{6.081267in}}{\pgfqpoint{4.123538in}{6.078953in}}{\pgfqpoint{4.119420in}{6.074835in}}%
\pgfpathcurveto{\pgfqpoint{4.115302in}{6.070716in}}{\pgfqpoint{4.112988in}{6.065130in}}{\pgfqpoint{4.112988in}{6.059306in}}%
\pgfpathcurveto{\pgfqpoint{4.112988in}{6.053482in}}{\pgfqpoint{4.115302in}{6.047896in}}{\pgfqpoint{4.119420in}{6.043778in}}%
\pgfpathcurveto{\pgfqpoint{4.123538in}{6.039660in}}{\pgfqpoint{4.129124in}{6.037346in}}{\pgfqpoint{4.134948in}{6.037346in}}%
\pgfpathlineto{\pgfqpoint{4.134948in}{6.037346in}}%
\pgfpathclose%
\pgfusepath{stroke,fill}%
\end{pgfscope}%
\begin{pgfscope}%
\pgfpathrectangle{\pgfqpoint{1.414736in}{0.548769in}}{\pgfqpoint{3.497747in}{7.253186in}}%
\pgfusepath{clip}%
\pgfsetbuttcap%
\pgfsetroundjoin%
\definecolor{currentfill}{rgb}{0.993248,0.906157,0.143936}%
\pgfsetfillcolor{currentfill}%
\pgfsetlinewidth{1.003750pt}%
\definecolor{currentstroke}{rgb}{0.993248,0.906157,0.143936}%
\pgfsetstrokecolor{currentstroke}%
\pgfsetdash{}{0pt}%
\pgfpathmoveto{\pgfqpoint{1.767145in}{6.508332in}}%
\pgfpathcurveto{\pgfqpoint{1.772969in}{6.508332in}}{\pgfqpoint{1.778555in}{6.510646in}}{\pgfqpoint{1.782673in}{6.514764in}}%
\pgfpathcurveto{\pgfqpoint{1.786791in}{6.518882in}}{\pgfqpoint{1.789105in}{6.524468in}}{\pgfqpoint{1.789105in}{6.530292in}}%
\pgfpathcurveto{\pgfqpoint{1.789105in}{6.536116in}}{\pgfqpoint{1.786791in}{6.541702in}}{\pgfqpoint{1.782673in}{6.545821in}}%
\pgfpathcurveto{\pgfqpoint{1.778555in}{6.549939in}}{\pgfqpoint{1.772969in}{6.552253in}}{\pgfqpoint{1.767145in}{6.552253in}}%
\pgfpathcurveto{\pgfqpoint{1.761321in}{6.552253in}}{\pgfqpoint{1.755735in}{6.549939in}}{\pgfqpoint{1.751617in}{6.545821in}}%
\pgfpathcurveto{\pgfqpoint{1.747499in}{6.541702in}}{\pgfqpoint{1.745185in}{6.536116in}}{\pgfqpoint{1.745185in}{6.530292in}}%
\pgfpathcurveto{\pgfqpoint{1.745185in}{6.524468in}}{\pgfqpoint{1.747499in}{6.518882in}}{\pgfqpoint{1.751617in}{6.514764in}}%
\pgfpathcurveto{\pgfqpoint{1.755735in}{6.510646in}}{\pgfqpoint{1.761321in}{6.508332in}}{\pgfqpoint{1.767145in}{6.508332in}}%
\pgfpathlineto{\pgfqpoint{1.767145in}{6.508332in}}%
\pgfpathclose%
\pgfusepath{stroke,fill}%
\end{pgfscope}%
\begin{pgfscope}%
\pgfpathrectangle{\pgfqpoint{1.414736in}{0.548769in}}{\pgfqpoint{3.497747in}{7.253186in}}%
\pgfusepath{clip}%
\pgfsetbuttcap%
\pgfsetroundjoin%
\definecolor{currentfill}{rgb}{0.876168,0.891125,0.095250}%
\pgfsetfillcolor{currentfill}%
\pgfsetlinewidth{1.003750pt}%
\definecolor{currentstroke}{rgb}{0.876168,0.891125,0.095250}%
\pgfsetstrokecolor{currentstroke}%
\pgfsetdash{}{0pt}%
\pgfpathmoveto{\pgfqpoint{2.872835in}{6.979318in}}%
\pgfpathcurveto{\pgfqpoint{2.878659in}{6.979318in}}{\pgfqpoint{2.884246in}{6.981632in}}{\pgfqpoint{2.888364in}{6.985750in}}%
\pgfpathcurveto{\pgfqpoint{2.892482in}{6.989868in}}{\pgfqpoint{2.894796in}{6.995455in}}{\pgfqpoint{2.894796in}{7.001278in}}%
\pgfpathcurveto{\pgfqpoint{2.894796in}{7.007102in}}{\pgfqpoint{2.892482in}{7.012689in}}{\pgfqpoint{2.888364in}{7.016807in}}%
\pgfpathcurveto{\pgfqpoint{2.884246in}{7.020925in}}{\pgfqpoint{2.878659in}{7.023239in}}{\pgfqpoint{2.872835in}{7.023239in}}%
\pgfpathcurveto{\pgfqpoint{2.867012in}{7.023239in}}{\pgfqpoint{2.861425in}{7.020925in}}{\pgfqpoint{2.857307in}{7.016807in}}%
\pgfpathcurveto{\pgfqpoint{2.853189in}{7.012689in}}{\pgfqpoint{2.850875in}{7.007102in}}{\pgfqpoint{2.850875in}{7.001278in}}%
\pgfpathcurveto{\pgfqpoint{2.850875in}{6.995455in}}{\pgfqpoint{2.853189in}{6.989868in}}{\pgfqpoint{2.857307in}{6.985750in}}%
\pgfpathcurveto{\pgfqpoint{2.861425in}{6.981632in}}{\pgfqpoint{2.867012in}{6.979318in}}{\pgfqpoint{2.872835in}{6.979318in}}%
\pgfpathlineto{\pgfqpoint{2.872835in}{6.979318in}}%
\pgfpathclose%
\pgfusepath{stroke,fill}%
\end{pgfscope}%
\begin{pgfscope}%
\pgfpathrectangle{\pgfqpoint{1.414736in}{0.548769in}}{\pgfqpoint{3.497747in}{7.253186in}}%
\pgfusepath{clip}%
\pgfsetbuttcap%
\pgfsetroundjoin%
\definecolor{currentfill}{rgb}{0.741388,0.873449,0.149561}%
\pgfsetfillcolor{currentfill}%
\pgfsetlinewidth{1.003750pt}%
\definecolor{currentstroke}{rgb}{0.741388,0.873449,0.149561}%
\pgfsetstrokecolor{currentstroke}%
\pgfsetdash{}{0pt}%
\pgfpathmoveto{\pgfqpoint{4.406744in}{7.450304in}}%
\pgfpathcurveto{\pgfqpoint{4.412568in}{7.450304in}}{\pgfqpoint{4.418154in}{7.452618in}}{\pgfqpoint{4.422272in}{7.456736in}}%
\pgfpathcurveto{\pgfqpoint{4.426390in}{7.460854in}}{\pgfqpoint{4.428704in}{7.466441in}}{\pgfqpoint{4.428704in}{7.472265in}}%
\pgfpathcurveto{\pgfqpoint{4.428704in}{7.478089in}}{\pgfqpoint{4.426390in}{7.483675in}}{\pgfqpoint{4.422272in}{7.487793in}}%
\pgfpathcurveto{\pgfqpoint{4.418154in}{7.491911in}}{\pgfqpoint{4.412568in}{7.494225in}}{\pgfqpoint{4.406744in}{7.494225in}}%
\pgfpathcurveto{\pgfqpoint{4.400920in}{7.494225in}}{\pgfqpoint{4.395334in}{7.491911in}}{\pgfqpoint{4.391216in}{7.487793in}}%
\pgfpathcurveto{\pgfqpoint{4.387098in}{7.483675in}}{\pgfqpoint{4.384784in}{7.478089in}}{\pgfqpoint{4.384784in}{7.472265in}}%
\pgfpathcurveto{\pgfqpoint{4.384784in}{7.466441in}}{\pgfqpoint{4.387098in}{7.460854in}}{\pgfqpoint{4.391216in}{7.456736in}}%
\pgfpathcurveto{\pgfqpoint{4.395334in}{7.452618in}}{\pgfqpoint{4.400920in}{7.450304in}}{\pgfqpoint{4.406744in}{7.450304in}}%
\pgfpathlineto{\pgfqpoint{4.406744in}{7.450304in}}%
\pgfpathclose%
\pgfusepath{stroke,fill}%
\end{pgfscope}%
\begin{pgfscope}%
\pgfpathrectangle{\pgfqpoint{1.414736in}{0.548769in}}{\pgfqpoint{3.497747in}{7.253186in}}%
\pgfusepath{clip}%
\pgfsetbuttcap%
\pgfsetroundjoin%
\definecolor{currentfill}{rgb}{0.814576,0.883393,0.110347}%
\pgfsetfillcolor{currentfill}%
\pgfsetlinewidth{1.003750pt}%
\definecolor{currentstroke}{rgb}{0.814576,0.883393,0.110347}%
\pgfsetstrokecolor{currentstroke}%
\pgfsetdash{}{0pt}%
\pgfpathmoveto{\pgfqpoint{3.403882in}{0.856499in}}%
\pgfpathcurveto{\pgfqpoint{3.409706in}{0.856499in}}{\pgfqpoint{3.415292in}{0.858813in}}{\pgfqpoint{3.419410in}{0.862931in}}%
\pgfpathcurveto{\pgfqpoint{3.423528in}{0.867049in}}{\pgfqpoint{3.425842in}{0.872635in}}{\pgfqpoint{3.425842in}{0.878459in}}%
\pgfpathcurveto{\pgfqpoint{3.425842in}{0.884283in}}{\pgfqpoint{3.423528in}{0.889869in}}{\pgfqpoint{3.419410in}{0.893987in}}%
\pgfpathcurveto{\pgfqpoint{3.415292in}{0.898105in}}{\pgfqpoint{3.409706in}{0.900419in}}{\pgfqpoint{3.403882in}{0.900419in}}%
\pgfpathcurveto{\pgfqpoint{3.398058in}{0.900419in}}{\pgfqpoint{3.392472in}{0.898105in}}{\pgfqpoint{3.388353in}{0.893987in}}%
\pgfpathcurveto{\pgfqpoint{3.384235in}{0.889869in}}{\pgfqpoint{3.381921in}{0.884283in}}{\pgfqpoint{3.381921in}{0.878459in}}%
\pgfpathcurveto{\pgfqpoint{3.381921in}{0.872635in}}{\pgfqpoint{3.384235in}{0.867049in}}{\pgfqpoint{3.388353in}{0.862931in}}%
\pgfpathcurveto{\pgfqpoint{3.392472in}{0.858813in}}{\pgfqpoint{3.398058in}{0.856499in}}{\pgfqpoint{3.403882in}{0.856499in}}%
\pgfpathlineto{\pgfqpoint{3.403882in}{0.856499in}}%
\pgfpathclose%
\pgfusepath{stroke,fill}%
\end{pgfscope}%
\begin{pgfscope}%
\pgfpathrectangle{\pgfqpoint{1.414736in}{0.548769in}}{\pgfqpoint{3.497747in}{7.253186in}}%
\pgfusepath{clip}%
\pgfsetbuttcap%
\pgfsetroundjoin%
\definecolor{currentfill}{rgb}{0.202219,0.715272,0.476084}%
\pgfsetfillcolor{currentfill}%
\pgfsetlinewidth{1.003750pt}%
\definecolor{currentstroke}{rgb}{0.202219,0.715272,0.476084}%
\pgfsetstrokecolor{currentstroke}%
\pgfsetdash{}{0pt}%
\pgfpathmoveto{\pgfqpoint{3.977090in}{1.327485in}}%
\pgfpathcurveto{\pgfqpoint{3.982913in}{1.327485in}}{\pgfqpoint{3.988500in}{1.329799in}}{\pgfqpoint{3.992618in}{1.333917in}}%
\pgfpathcurveto{\pgfqpoint{3.996736in}{1.338035in}}{\pgfqpoint{3.999050in}{1.343621in}}{\pgfqpoint{3.999050in}{1.349445in}}%
\pgfpathcurveto{\pgfqpoint{3.999050in}{1.355269in}}{\pgfqpoint{3.996736in}{1.360855in}}{\pgfqpoint{3.992618in}{1.364973in}}%
\pgfpathcurveto{\pgfqpoint{3.988500in}{1.369092in}}{\pgfqpoint{3.982913in}{1.371405in}}{\pgfqpoint{3.977090in}{1.371405in}}%
\pgfpathcurveto{\pgfqpoint{3.971266in}{1.371405in}}{\pgfqpoint{3.965679in}{1.369092in}}{\pgfqpoint{3.961561in}{1.364973in}}%
\pgfpathcurveto{\pgfqpoint{3.957443in}{1.360855in}}{\pgfqpoint{3.955129in}{1.355269in}}{\pgfqpoint{3.955129in}{1.349445in}}%
\pgfpathcurveto{\pgfqpoint{3.955129in}{1.343621in}}{\pgfqpoint{3.957443in}{1.338035in}}{\pgfqpoint{3.961561in}{1.333917in}}%
\pgfpathcurveto{\pgfqpoint{3.965679in}{1.329799in}}{\pgfqpoint{3.971266in}{1.327485in}}{\pgfqpoint{3.977090in}{1.327485in}}%
\pgfpathlineto{\pgfqpoint{3.977090in}{1.327485in}}%
\pgfpathclose%
\pgfusepath{stroke,fill}%
\end{pgfscope}%
\begin{pgfscope}%
\pgfpathrectangle{\pgfqpoint{1.414736in}{0.548769in}}{\pgfqpoint{3.497747in}{7.253186in}}%
\pgfusepath{clip}%
\pgfsetbuttcap%
\pgfsetroundjoin%
\definecolor{currentfill}{rgb}{0.804182,0.882046,0.114965}%
\pgfsetfillcolor{currentfill}%
\pgfsetlinewidth{1.003750pt}%
\definecolor{currentstroke}{rgb}{0.804182,0.882046,0.114965}%
\pgfsetstrokecolor{currentstroke}%
\pgfsetdash{}{0pt}%
\pgfpathmoveto{\pgfqpoint{3.410647in}{1.798471in}}%
\pgfpathcurveto{\pgfqpoint{3.416471in}{1.798471in}}{\pgfqpoint{3.422057in}{1.800785in}}{\pgfqpoint{3.426175in}{1.804903in}}%
\pgfpathcurveto{\pgfqpoint{3.430293in}{1.809021in}}{\pgfqpoint{3.432607in}{1.814607in}}{\pgfqpoint{3.432607in}{1.820431in}}%
\pgfpathcurveto{\pgfqpoint{3.432607in}{1.826255in}}{\pgfqpoint{3.430293in}{1.831841in}}{\pgfqpoint{3.426175in}{1.835960in}}%
\pgfpathcurveto{\pgfqpoint{3.422057in}{1.840078in}}{\pgfqpoint{3.416471in}{1.842392in}}{\pgfqpoint{3.410647in}{1.842392in}}%
\pgfpathcurveto{\pgfqpoint{3.404823in}{1.842392in}}{\pgfqpoint{3.399237in}{1.840078in}}{\pgfqpoint{3.395119in}{1.835960in}}%
\pgfpathcurveto{\pgfqpoint{3.391000in}{1.831841in}}{\pgfqpoint{3.388687in}{1.826255in}}{\pgfqpoint{3.388687in}{1.820431in}}%
\pgfpathcurveto{\pgfqpoint{3.388687in}{1.814607in}}{\pgfqpoint{3.391000in}{1.809021in}}{\pgfqpoint{3.395119in}{1.804903in}}%
\pgfpathcurveto{\pgfqpoint{3.399237in}{1.800785in}}{\pgfqpoint{3.404823in}{1.798471in}}{\pgfqpoint{3.410647in}{1.798471in}}%
\pgfpathlineto{\pgfqpoint{3.410647in}{1.798471in}}%
\pgfpathclose%
\pgfusepath{stroke,fill}%
\end{pgfscope}%
\begin{pgfscope}%
\pgfpathrectangle{\pgfqpoint{1.414736in}{0.548769in}}{\pgfqpoint{3.497747in}{7.253186in}}%
\pgfusepath{clip}%
\pgfsetbuttcap%
\pgfsetroundjoin%
\definecolor{currentfill}{rgb}{0.845561,0.887322,0.099702}%
\pgfsetfillcolor{currentfill}%
\pgfsetlinewidth{1.003750pt}%
\definecolor{currentstroke}{rgb}{0.845561,0.887322,0.099702}%
\pgfsetstrokecolor{currentstroke}%
\pgfsetdash{}{0pt}%
\pgfpathmoveto{\pgfqpoint{4.151890in}{2.269457in}}%
\pgfpathcurveto{\pgfqpoint{4.157714in}{2.269457in}}{\pgfqpoint{4.163300in}{2.271771in}}{\pgfqpoint{4.167418in}{2.275889in}}%
\pgfpathcurveto{\pgfqpoint{4.171536in}{2.280007in}}{\pgfqpoint{4.173850in}{2.285593in}}{\pgfqpoint{4.173850in}{2.291417in}}%
\pgfpathcurveto{\pgfqpoint{4.173850in}{2.297241in}}{\pgfqpoint{4.171536in}{2.302828in}}{\pgfqpoint{4.167418in}{2.306946in}}%
\pgfpathcurveto{\pgfqpoint{4.163300in}{2.311064in}}{\pgfqpoint{4.157714in}{2.313378in}}{\pgfqpoint{4.151890in}{2.313378in}}%
\pgfpathcurveto{\pgfqpoint{4.146066in}{2.313378in}}{\pgfqpoint{4.140480in}{2.311064in}}{\pgfqpoint{4.136361in}{2.306946in}}%
\pgfpathcurveto{\pgfqpoint{4.132243in}{2.302828in}}{\pgfqpoint{4.129929in}{2.297241in}}{\pgfqpoint{4.129929in}{2.291417in}}%
\pgfpathcurveto{\pgfqpoint{4.129929in}{2.285593in}}{\pgfqpoint{4.132243in}{2.280007in}}{\pgfqpoint{4.136361in}{2.275889in}}%
\pgfpathcurveto{\pgfqpoint{4.140480in}{2.271771in}}{\pgfqpoint{4.146066in}{2.269457in}}{\pgfqpoint{4.151890in}{2.269457in}}%
\pgfpathlineto{\pgfqpoint{4.151890in}{2.269457in}}%
\pgfpathclose%
\pgfusepath{stroke,fill}%
\end{pgfscope}%
\begin{pgfscope}%
\pgfpathrectangle{\pgfqpoint{1.414736in}{0.548769in}}{\pgfqpoint{3.497747in}{7.253186in}}%
\pgfusepath{clip}%
\pgfsetbuttcap%
\pgfsetroundjoin%
\definecolor{currentfill}{rgb}{0.657642,0.860219,0.203082}%
\pgfsetfillcolor{currentfill}%
\pgfsetlinewidth{1.003750pt}%
\definecolor{currentstroke}{rgb}{0.657642,0.860219,0.203082}%
\pgfsetstrokecolor{currentstroke}%
\pgfsetdash{}{0pt}%
\pgfpathmoveto{\pgfqpoint{3.277768in}{2.740443in}}%
\pgfpathcurveto{\pgfqpoint{3.283592in}{2.740443in}}{\pgfqpoint{3.289178in}{2.742757in}}{\pgfqpoint{3.293296in}{2.746875in}}%
\pgfpathcurveto{\pgfqpoint{3.297414in}{2.750993in}}{\pgfqpoint{3.299728in}{2.756580in}}{\pgfqpoint{3.299728in}{2.762404in}}%
\pgfpathcurveto{\pgfqpoint{3.299728in}{2.768227in}}{\pgfqpoint{3.297414in}{2.773814in}}{\pgfqpoint{3.293296in}{2.777932in}}%
\pgfpathcurveto{\pgfqpoint{3.289178in}{2.782050in}}{\pgfqpoint{3.283592in}{2.784364in}}{\pgfqpoint{3.277768in}{2.784364in}}%
\pgfpathcurveto{\pgfqpoint{3.271944in}{2.784364in}}{\pgfqpoint{3.266358in}{2.782050in}}{\pgfqpoint{3.262240in}{2.777932in}}%
\pgfpathcurveto{\pgfqpoint{3.258121in}{2.773814in}}{\pgfqpoint{3.255808in}{2.768227in}}{\pgfqpoint{3.255808in}{2.762404in}}%
\pgfpathcurveto{\pgfqpoint{3.255808in}{2.756580in}}{\pgfqpoint{3.258121in}{2.750993in}}{\pgfqpoint{3.262240in}{2.746875in}}%
\pgfpathcurveto{\pgfqpoint{3.266358in}{2.742757in}}{\pgfqpoint{3.271944in}{2.740443in}}{\pgfqpoint{3.277768in}{2.740443in}}%
\pgfpathlineto{\pgfqpoint{3.277768in}{2.740443in}}%
\pgfpathclose%
\pgfusepath{stroke,fill}%
\end{pgfscope}%
\begin{pgfscope}%
\pgfpathrectangle{\pgfqpoint{1.414736in}{0.548769in}}{\pgfqpoint{3.497747in}{7.253186in}}%
\pgfusepath{clip}%
\pgfsetbuttcap%
\pgfsetroundjoin%
\definecolor{currentfill}{rgb}{0.335885,0.777018,0.402049}%
\pgfsetfillcolor{currentfill}%
\pgfsetlinewidth{1.003750pt}%
\definecolor{currentstroke}{rgb}{0.335885,0.777018,0.402049}%
\pgfsetstrokecolor{currentstroke}%
\pgfsetdash{}{0pt}%
\pgfpathmoveto{\pgfqpoint{4.321733in}{3.211429in}}%
\pgfpathcurveto{\pgfqpoint{4.327557in}{3.211429in}}{\pgfqpoint{4.333143in}{3.213743in}}{\pgfqpoint{4.337262in}{3.217861in}}%
\pgfpathcurveto{\pgfqpoint{4.341380in}{3.221980in}}{\pgfqpoint{4.343694in}{3.227566in}}{\pgfqpoint{4.343694in}{3.233390in}}%
\pgfpathcurveto{\pgfqpoint{4.343694in}{3.239214in}}{\pgfqpoint{4.341380in}{3.244800in}}{\pgfqpoint{4.337262in}{3.248918in}}%
\pgfpathcurveto{\pgfqpoint{4.333143in}{3.253036in}}{\pgfqpoint{4.327557in}{3.255350in}}{\pgfqpoint{4.321733in}{3.255350in}}%
\pgfpathcurveto{\pgfqpoint{4.315909in}{3.255350in}}{\pgfqpoint{4.310323in}{3.253036in}}{\pgfqpoint{4.306205in}{3.248918in}}%
\pgfpathcurveto{\pgfqpoint{4.302087in}{3.244800in}}{\pgfqpoint{4.299773in}{3.239214in}}{\pgfqpoint{4.299773in}{3.233390in}}%
\pgfpathcurveto{\pgfqpoint{4.299773in}{3.227566in}}{\pgfqpoint{4.302087in}{3.221980in}}{\pgfqpoint{4.306205in}{3.217861in}}%
\pgfpathcurveto{\pgfqpoint{4.310323in}{3.213743in}}{\pgfqpoint{4.315909in}{3.211429in}}{\pgfqpoint{4.321733in}{3.211429in}}%
\pgfpathlineto{\pgfqpoint{4.321733in}{3.211429in}}%
\pgfpathclose%
\pgfusepath{stroke,fill}%
\end{pgfscope}%
\begin{pgfscope}%
\pgfpathrectangle{\pgfqpoint{1.414736in}{0.548769in}}{\pgfqpoint{3.497747in}{7.253186in}}%
\pgfusepath{clip}%
\pgfsetbuttcap%
\pgfsetroundjoin%
\definecolor{currentfill}{rgb}{0.993248,0.906157,0.143936}%
\pgfsetfillcolor{currentfill}%
\pgfsetlinewidth{1.003750pt}%
\definecolor{currentstroke}{rgb}{0.993248,0.906157,0.143936}%
\pgfsetstrokecolor{currentstroke}%
\pgfsetdash{}{0pt}%
\pgfpathmoveto{\pgfqpoint{3.153807in}{3.682415in}}%
\pgfpathcurveto{\pgfqpoint{3.159631in}{3.682415in}}{\pgfqpoint{3.165218in}{3.684729in}}{\pgfqpoint{3.169336in}{3.688847in}}%
\pgfpathcurveto{\pgfqpoint{3.173454in}{3.692966in}}{\pgfqpoint{3.175768in}{3.698552in}}{\pgfqpoint{3.175768in}{3.704376in}}%
\pgfpathcurveto{\pgfqpoint{3.175768in}{3.710200in}}{\pgfqpoint{3.173454in}{3.715786in}}{\pgfqpoint{3.169336in}{3.719904in}}%
\pgfpathcurveto{\pgfqpoint{3.165218in}{3.724022in}}{\pgfqpoint{3.159631in}{3.726336in}}{\pgfqpoint{3.153807in}{3.726336in}}%
\pgfpathcurveto{\pgfqpoint{3.147984in}{3.726336in}}{\pgfqpoint{3.142397in}{3.724022in}}{\pgfqpoint{3.138279in}{3.719904in}}%
\pgfpathcurveto{\pgfqpoint{3.134161in}{3.715786in}}{\pgfqpoint{3.131847in}{3.710200in}}{\pgfqpoint{3.131847in}{3.704376in}}%
\pgfpathcurveto{\pgfqpoint{3.131847in}{3.698552in}}{\pgfqpoint{3.134161in}{3.692966in}}{\pgfqpoint{3.138279in}{3.688847in}}%
\pgfpathcurveto{\pgfqpoint{3.142397in}{3.684729in}}{\pgfqpoint{3.147984in}{3.682415in}}{\pgfqpoint{3.153807in}{3.682415in}}%
\pgfpathlineto{\pgfqpoint{3.153807in}{3.682415in}}%
\pgfpathclose%
\pgfusepath{stroke,fill}%
\end{pgfscope}%
\begin{pgfscope}%
\pgfpathrectangle{\pgfqpoint{1.414736in}{0.548769in}}{\pgfqpoint{3.497747in}{7.253186in}}%
\pgfusepath{clip}%
\pgfsetbuttcap%
\pgfsetroundjoin%
\definecolor{currentfill}{rgb}{0.993248,0.906157,0.143936}%
\pgfsetfillcolor{currentfill}%
\pgfsetlinewidth{1.003750pt}%
\definecolor{currentstroke}{rgb}{0.993248,0.906157,0.143936}%
\pgfsetstrokecolor{currentstroke}%
\pgfsetdash{}{0pt}%
\pgfpathmoveto{\pgfqpoint{4.003799in}{4.153402in}}%
\pgfpathcurveto{\pgfqpoint{4.009623in}{4.153402in}}{\pgfqpoint{4.015209in}{4.155715in}}{\pgfqpoint{4.019327in}{4.159834in}}%
\pgfpathcurveto{\pgfqpoint{4.023446in}{4.163952in}}{\pgfqpoint{4.025759in}{4.169538in}}{\pgfqpoint{4.025759in}{4.175362in}}%
\pgfpathcurveto{\pgfqpoint{4.025759in}{4.181186in}}{\pgfqpoint{4.023446in}{4.186772in}}{\pgfqpoint{4.019327in}{4.190890in}}%
\pgfpathcurveto{\pgfqpoint{4.015209in}{4.195008in}}{\pgfqpoint{4.009623in}{4.197322in}}{\pgfqpoint{4.003799in}{4.197322in}}%
\pgfpathcurveto{\pgfqpoint{3.997975in}{4.197322in}}{\pgfqpoint{3.992389in}{4.195008in}}{\pgfqpoint{3.988271in}{4.190890in}}%
\pgfpathcurveto{\pgfqpoint{3.984153in}{4.186772in}}{\pgfqpoint{3.981839in}{4.181186in}}{\pgfqpoint{3.981839in}{4.175362in}}%
\pgfpathcurveto{\pgfqpoint{3.981839in}{4.169538in}}{\pgfqpoint{3.984153in}{4.163952in}}{\pgfqpoint{3.988271in}{4.159834in}}%
\pgfpathcurveto{\pgfqpoint{3.992389in}{4.155715in}}{\pgfqpoint{3.997975in}{4.153402in}}{\pgfqpoint{4.003799in}{4.153402in}}%
\pgfpathlineto{\pgfqpoint{4.003799in}{4.153402in}}%
\pgfpathclose%
\pgfusepath{stroke,fill}%
\end{pgfscope}%
\begin{pgfscope}%
\pgfpathrectangle{\pgfqpoint{1.414736in}{0.548769in}}{\pgfqpoint{3.497747in}{7.253186in}}%
\pgfusepath{clip}%
\pgfsetbuttcap%
\pgfsetroundjoin%
\definecolor{currentfill}{rgb}{0.535621,0.835785,0.281908}%
\pgfsetfillcolor{currentfill}%
\pgfsetlinewidth{1.003750pt}%
\definecolor{currentstroke}{rgb}{0.535621,0.835785,0.281908}%
\pgfsetstrokecolor{currentstroke}%
\pgfsetdash{}{0pt}%
\pgfpathmoveto{\pgfqpoint{3.029278in}{4.624388in}}%
\pgfpathcurveto{\pgfqpoint{3.035102in}{4.624388in}}{\pgfqpoint{3.040688in}{4.626702in}}{\pgfqpoint{3.044807in}{4.630820in}}%
\pgfpathcurveto{\pgfqpoint{3.048925in}{4.634938in}}{\pgfqpoint{3.051239in}{4.640524in}}{\pgfqpoint{3.051239in}{4.646348in}}%
\pgfpathcurveto{\pgfqpoint{3.051239in}{4.652172in}}{\pgfqpoint{3.048925in}{4.657758in}}{\pgfqpoint{3.044807in}{4.661876in}}%
\pgfpathcurveto{\pgfqpoint{3.040688in}{4.665994in}}{\pgfqpoint{3.035102in}{4.668308in}}{\pgfqpoint{3.029278in}{4.668308in}}%
\pgfpathcurveto{\pgfqpoint{3.023454in}{4.668308in}}{\pgfqpoint{3.017868in}{4.665994in}}{\pgfqpoint{3.013750in}{4.661876in}}%
\pgfpathcurveto{\pgfqpoint{3.009632in}{4.657758in}}{\pgfqpoint{3.007318in}{4.652172in}}{\pgfqpoint{3.007318in}{4.646348in}}%
\pgfpathcurveto{\pgfqpoint{3.007318in}{4.640524in}}{\pgfqpoint{3.009632in}{4.634938in}}{\pgfqpoint{3.013750in}{4.630820in}}%
\pgfpathcurveto{\pgfqpoint{3.017868in}{4.626702in}}{\pgfqpoint{3.023454in}{4.624388in}}{\pgfqpoint{3.029278in}{4.624388in}}%
\pgfpathlineto{\pgfqpoint{3.029278in}{4.624388in}}%
\pgfpathclose%
\pgfusepath{stroke,fill}%
\end{pgfscope}%
\begin{pgfscope}%
\pgfpathrectangle{\pgfqpoint{1.414736in}{0.548769in}}{\pgfqpoint{3.497747in}{7.253186in}}%
\pgfusepath{clip}%
\pgfsetbuttcap%
\pgfsetroundjoin%
\definecolor{currentfill}{rgb}{0.175841,0.441290,0.557685}%
\pgfsetfillcolor{currentfill}%
\pgfsetlinewidth{1.003750pt}%
\definecolor{currentstroke}{rgb}{0.175841,0.441290,0.557685}%
\pgfsetstrokecolor{currentstroke}%
\pgfsetdash{}{0pt}%
\pgfpathmoveto{\pgfqpoint{3.641703in}{5.095374in}}%
\pgfpathcurveto{\pgfqpoint{3.647527in}{5.095374in}}{\pgfqpoint{3.653113in}{5.097688in}}{\pgfqpoint{3.657231in}{5.101806in}}%
\pgfpathcurveto{\pgfqpoint{3.661349in}{5.105924in}}{\pgfqpoint{3.663663in}{5.111510in}}{\pgfqpoint{3.663663in}{5.117334in}}%
\pgfpathcurveto{\pgfqpoint{3.663663in}{5.123158in}}{\pgfqpoint{3.661349in}{5.128744in}}{\pgfqpoint{3.657231in}{5.132862in}}%
\pgfpathcurveto{\pgfqpoint{3.653113in}{5.136980in}}{\pgfqpoint{3.647527in}{5.139294in}}{\pgfqpoint{3.641703in}{5.139294in}}%
\pgfpathcurveto{\pgfqpoint{3.635879in}{5.139294in}}{\pgfqpoint{3.630293in}{5.136980in}}{\pgfqpoint{3.626175in}{5.132862in}}%
\pgfpathcurveto{\pgfqpoint{3.622056in}{5.128744in}}{\pgfqpoint{3.619743in}{5.123158in}}{\pgfqpoint{3.619743in}{5.117334in}}%
\pgfpathcurveto{\pgfqpoint{3.619743in}{5.111510in}}{\pgfqpoint{3.622056in}{5.105924in}}{\pgfqpoint{3.626175in}{5.101806in}}%
\pgfpathcurveto{\pgfqpoint{3.630293in}{5.097688in}}{\pgfqpoint{3.635879in}{5.095374in}}{\pgfqpoint{3.641703in}{5.095374in}}%
\pgfpathlineto{\pgfqpoint{3.641703in}{5.095374in}}%
\pgfpathclose%
\pgfusepath{stroke,fill}%
\end{pgfscope}%
\begin{pgfscope}%
\pgfpathrectangle{\pgfqpoint{1.414736in}{0.548769in}}{\pgfqpoint{3.497747in}{7.253186in}}%
\pgfusepath{clip}%
\pgfsetbuttcap%
\pgfsetroundjoin%
\definecolor{currentfill}{rgb}{0.154815,0.493313,0.557840}%
\pgfsetfillcolor{currentfill}%
\pgfsetlinewidth{1.003750pt}%
\definecolor{currentstroke}{rgb}{0.154815,0.493313,0.557840}%
\pgfsetstrokecolor{currentstroke}%
\pgfsetdash{}{0pt}%
\pgfpathmoveto{\pgfqpoint{3.940313in}{5.566360in}}%
\pgfpathcurveto{\pgfqpoint{3.946137in}{5.566360in}}{\pgfqpoint{3.951723in}{5.568674in}}{\pgfqpoint{3.955841in}{5.572792in}}%
\pgfpathcurveto{\pgfqpoint{3.959960in}{5.576910in}}{\pgfqpoint{3.962274in}{5.582496in}}{\pgfqpoint{3.962274in}{5.588320in}}%
\pgfpathcurveto{\pgfqpoint{3.962274in}{5.594144in}}{\pgfqpoint{3.959960in}{5.599730in}}{\pgfqpoint{3.955841in}{5.603848in}}%
\pgfpathcurveto{\pgfqpoint{3.951723in}{5.607967in}}{\pgfqpoint{3.946137in}{5.610280in}}{\pgfqpoint{3.940313in}{5.610280in}}%
\pgfpathcurveto{\pgfqpoint{3.934489in}{5.610280in}}{\pgfqpoint{3.928903in}{5.607967in}}{\pgfqpoint{3.924785in}{5.603848in}}%
\pgfpathcurveto{\pgfqpoint{3.920667in}{5.599730in}}{\pgfqpoint{3.918353in}{5.594144in}}{\pgfqpoint{3.918353in}{5.588320in}}%
\pgfpathcurveto{\pgfqpoint{3.918353in}{5.582496in}}{\pgfqpoint{3.920667in}{5.576910in}}{\pgfqpoint{3.924785in}{5.572792in}}%
\pgfpathcurveto{\pgfqpoint{3.928903in}{5.568674in}}{\pgfqpoint{3.934489in}{5.566360in}}{\pgfqpoint{3.940313in}{5.566360in}}%
\pgfpathlineto{\pgfqpoint{3.940313in}{5.566360in}}%
\pgfpathclose%
\pgfusepath{stroke,fill}%
\end{pgfscope}%
\begin{pgfscope}%
\pgfpathrectangle{\pgfqpoint{1.414736in}{0.548769in}}{\pgfqpoint{3.497747in}{7.253186in}}%
\pgfusepath{clip}%
\pgfsetbuttcap%
\pgfsetroundjoin%
\definecolor{currentfill}{rgb}{0.126326,0.644107,0.525311}%
\pgfsetfillcolor{currentfill}%
\pgfsetlinewidth{1.003750pt}%
\definecolor{currentstroke}{rgb}{0.126326,0.644107,0.525311}%
\pgfsetstrokecolor{currentstroke}%
\pgfsetdash{}{0pt}%
\pgfpathmoveto{\pgfqpoint{3.409160in}{6.037346in}}%
\pgfpathcurveto{\pgfqpoint{3.414984in}{6.037346in}}{\pgfqpoint{3.420571in}{6.039660in}}{\pgfqpoint{3.424689in}{6.043778in}}%
\pgfpathcurveto{\pgfqpoint{3.428807in}{6.047896in}}{\pgfqpoint{3.431121in}{6.053482in}}{\pgfqpoint{3.431121in}{6.059306in}}%
\pgfpathcurveto{\pgfqpoint{3.431121in}{6.065130in}}{\pgfqpoint{3.428807in}{6.070716in}}{\pgfqpoint{3.424689in}{6.074835in}}%
\pgfpathcurveto{\pgfqpoint{3.420571in}{6.078953in}}{\pgfqpoint{3.414984in}{6.081267in}}{\pgfqpoint{3.409160in}{6.081267in}}%
\pgfpathcurveto{\pgfqpoint{3.403337in}{6.081267in}}{\pgfqpoint{3.397750in}{6.078953in}}{\pgfqpoint{3.393632in}{6.074835in}}%
\pgfpathcurveto{\pgfqpoint{3.389514in}{6.070716in}}{\pgfqpoint{3.387200in}{6.065130in}}{\pgfqpoint{3.387200in}{6.059306in}}%
\pgfpathcurveto{\pgfqpoint{3.387200in}{6.053482in}}{\pgfqpoint{3.389514in}{6.047896in}}{\pgfqpoint{3.393632in}{6.043778in}}%
\pgfpathcurveto{\pgfqpoint{3.397750in}{6.039660in}}{\pgfqpoint{3.403337in}{6.037346in}}{\pgfqpoint{3.409160in}{6.037346in}}%
\pgfpathlineto{\pgfqpoint{3.409160in}{6.037346in}}%
\pgfpathclose%
\pgfusepath{stroke,fill}%
\end{pgfscope}%
\begin{pgfscope}%
\pgfpathrectangle{\pgfqpoint{1.414736in}{0.548769in}}{\pgfqpoint{3.497747in}{7.253186in}}%
\pgfusepath{clip}%
\pgfsetbuttcap%
\pgfsetroundjoin%
\definecolor{currentfill}{rgb}{0.993248,0.906157,0.143936}%
\pgfsetfillcolor{currentfill}%
\pgfsetlinewidth{1.003750pt}%
\definecolor{currentstroke}{rgb}{0.993248,0.906157,0.143936}%
\pgfsetstrokecolor{currentstroke}%
\pgfsetdash{}{0pt}%
\pgfpathmoveto{\pgfqpoint{3.383655in}{6.508332in}}%
\pgfpathcurveto{\pgfqpoint{3.389479in}{6.508332in}}{\pgfqpoint{3.395065in}{6.510646in}}{\pgfqpoint{3.399183in}{6.514764in}}%
\pgfpathcurveto{\pgfqpoint{3.403301in}{6.518882in}}{\pgfqpoint{3.405615in}{6.524468in}}{\pgfqpoint{3.405615in}{6.530292in}}%
\pgfpathcurveto{\pgfqpoint{3.405615in}{6.536116in}}{\pgfqpoint{3.403301in}{6.541702in}}{\pgfqpoint{3.399183in}{6.545821in}}%
\pgfpathcurveto{\pgfqpoint{3.395065in}{6.549939in}}{\pgfqpoint{3.389479in}{6.552253in}}{\pgfqpoint{3.383655in}{6.552253in}}%
\pgfpathcurveto{\pgfqpoint{3.377831in}{6.552253in}}{\pgfqpoint{3.372245in}{6.549939in}}{\pgfqpoint{3.368126in}{6.545821in}}%
\pgfpathcurveto{\pgfqpoint{3.364008in}{6.541702in}}{\pgfqpoint{3.361694in}{6.536116in}}{\pgfqpoint{3.361694in}{6.530292in}}%
\pgfpathcurveto{\pgfqpoint{3.361694in}{6.524468in}}{\pgfqpoint{3.364008in}{6.518882in}}{\pgfqpoint{3.368126in}{6.514764in}}%
\pgfpathcurveto{\pgfqpoint{3.372245in}{6.510646in}}{\pgfqpoint{3.377831in}{6.508332in}}{\pgfqpoint{3.383655in}{6.508332in}}%
\pgfpathlineto{\pgfqpoint{3.383655in}{6.508332in}}%
\pgfpathclose%
\pgfusepath{stroke,fill}%
\end{pgfscope}%
\begin{pgfscope}%
\pgfpathrectangle{\pgfqpoint{1.414736in}{0.548769in}}{\pgfqpoint{3.497747in}{7.253186in}}%
\pgfusepath{clip}%
\pgfsetbuttcap%
\pgfsetroundjoin%
\definecolor{currentfill}{rgb}{0.668054,0.861999,0.196293}%
\pgfsetfillcolor{currentfill}%
\pgfsetlinewidth{1.003750pt}%
\definecolor{currentstroke}{rgb}{0.668054,0.861999,0.196293}%
\pgfsetstrokecolor{currentstroke}%
\pgfsetdash{}{0pt}%
\pgfpathmoveto{\pgfqpoint{3.148465in}{6.979318in}}%
\pgfpathcurveto{\pgfqpoint{3.154289in}{6.979318in}}{\pgfqpoint{3.159875in}{6.981632in}}{\pgfqpoint{3.163993in}{6.985750in}}%
\pgfpathcurveto{\pgfqpoint{3.168111in}{6.989868in}}{\pgfqpoint{3.170425in}{6.995455in}}{\pgfqpoint{3.170425in}{7.001278in}}%
\pgfpathcurveto{\pgfqpoint{3.170425in}{7.007102in}}{\pgfqpoint{3.168111in}{7.012689in}}{\pgfqpoint{3.163993in}{7.016807in}}%
\pgfpathcurveto{\pgfqpoint{3.159875in}{7.020925in}}{\pgfqpoint{3.154289in}{7.023239in}}{\pgfqpoint{3.148465in}{7.023239in}}%
\pgfpathcurveto{\pgfqpoint{3.142641in}{7.023239in}}{\pgfqpoint{3.137055in}{7.020925in}}{\pgfqpoint{3.132937in}{7.016807in}}%
\pgfpathcurveto{\pgfqpoint{3.128819in}{7.012689in}}{\pgfqpoint{3.126505in}{7.007102in}}{\pgfqpoint{3.126505in}{7.001278in}}%
\pgfpathcurveto{\pgfqpoint{3.126505in}{6.995455in}}{\pgfqpoint{3.128819in}{6.989868in}}{\pgfqpoint{3.132937in}{6.985750in}}%
\pgfpathcurveto{\pgfqpoint{3.137055in}{6.981632in}}{\pgfqpoint{3.142641in}{6.979318in}}{\pgfqpoint{3.148465in}{6.979318in}}%
\pgfpathlineto{\pgfqpoint{3.148465in}{6.979318in}}%
\pgfpathclose%
\pgfusepath{stroke,fill}%
\end{pgfscope}%
\begin{pgfscope}%
\pgfpathrectangle{\pgfqpoint{1.414736in}{0.548769in}}{\pgfqpoint{3.497747in}{7.253186in}}%
\pgfusepath{clip}%
\pgfsetbuttcap%
\pgfsetroundjoin%
\definecolor{currentfill}{rgb}{0.993248,0.906157,0.143936}%
\pgfsetfillcolor{currentfill}%
\pgfsetlinewidth{1.003750pt}%
\definecolor{currentstroke}{rgb}{0.993248,0.906157,0.143936}%
\pgfsetstrokecolor{currentstroke}%
\pgfsetdash{}{0pt}%
\pgfpathmoveto{\pgfqpoint{4.549937in}{7.450304in}}%
\pgfpathcurveto{\pgfqpoint{4.555761in}{7.450304in}}{\pgfqpoint{4.561347in}{7.452618in}}{\pgfqpoint{4.565465in}{7.456736in}}%
\pgfpathcurveto{\pgfqpoint{4.569584in}{7.460854in}}{\pgfqpoint{4.571898in}{7.466441in}}{\pgfqpoint{4.571898in}{7.472265in}}%
\pgfpathcurveto{\pgfqpoint{4.571898in}{7.478089in}}{\pgfqpoint{4.569584in}{7.483675in}}{\pgfqpoint{4.565465in}{7.487793in}}%
\pgfpathcurveto{\pgfqpoint{4.561347in}{7.491911in}}{\pgfqpoint{4.555761in}{7.494225in}}{\pgfqpoint{4.549937in}{7.494225in}}%
\pgfpathcurveto{\pgfqpoint{4.544113in}{7.494225in}}{\pgfqpoint{4.538527in}{7.491911in}}{\pgfqpoint{4.534409in}{7.487793in}}%
\pgfpathcurveto{\pgfqpoint{4.530291in}{7.483675in}}{\pgfqpoint{4.527977in}{7.478089in}}{\pgfqpoint{4.527977in}{7.472265in}}%
\pgfpathcurveto{\pgfqpoint{4.527977in}{7.466441in}}{\pgfqpoint{4.530291in}{7.460854in}}{\pgfqpoint{4.534409in}{7.456736in}}%
\pgfpathcurveto{\pgfqpoint{4.538527in}{7.452618in}}{\pgfqpoint{4.544113in}{7.450304in}}{\pgfqpoint{4.549937in}{7.450304in}}%
\pgfpathlineto{\pgfqpoint{4.549937in}{7.450304in}}%
\pgfpathclose%
\pgfusepath{stroke,fill}%
\end{pgfscope}%
\begin{pgfscope}%
\pgfpathrectangle{\pgfqpoint{1.414736in}{0.548769in}}{\pgfqpoint{3.497747in}{7.253186in}}%
\pgfusepath{clip}%
\pgfsetbuttcap%
\pgfsetroundjoin%
\definecolor{currentfill}{rgb}{0.730889,0.871916,0.156029}%
\pgfsetfillcolor{currentfill}%
\pgfsetlinewidth{1.003750pt}%
\definecolor{currentstroke}{rgb}{0.730889,0.871916,0.156029}%
\pgfsetstrokecolor{currentstroke}%
\pgfsetdash{}{0pt}%
\pgfpathmoveto{\pgfqpoint{3.573316in}{0.856499in}}%
\pgfpathcurveto{\pgfqpoint{3.579140in}{0.856499in}}{\pgfqpoint{3.584726in}{0.858813in}}{\pgfqpoint{3.588844in}{0.862931in}}%
\pgfpathcurveto{\pgfqpoint{3.592963in}{0.867049in}}{\pgfqpoint{3.595276in}{0.872635in}}{\pgfqpoint{3.595276in}{0.878459in}}%
\pgfpathcurveto{\pgfqpoint{3.595276in}{0.884283in}}{\pgfqpoint{3.592963in}{0.889869in}}{\pgfqpoint{3.588844in}{0.893987in}}%
\pgfpathcurveto{\pgfqpoint{3.584726in}{0.898105in}}{\pgfqpoint{3.579140in}{0.900419in}}{\pgfqpoint{3.573316in}{0.900419in}}%
\pgfpathcurveto{\pgfqpoint{3.567492in}{0.900419in}}{\pgfqpoint{3.561906in}{0.898105in}}{\pgfqpoint{3.557788in}{0.893987in}}%
\pgfpathcurveto{\pgfqpoint{3.553670in}{0.889869in}}{\pgfqpoint{3.551356in}{0.884283in}}{\pgfqpoint{3.551356in}{0.878459in}}%
\pgfpathcurveto{\pgfqpoint{3.551356in}{0.872635in}}{\pgfqpoint{3.553670in}{0.867049in}}{\pgfqpoint{3.557788in}{0.862931in}}%
\pgfpathcurveto{\pgfqpoint{3.561906in}{0.858813in}}{\pgfqpoint{3.567492in}{0.856499in}}{\pgfqpoint{3.573316in}{0.856499in}}%
\pgfpathlineto{\pgfqpoint{3.573316in}{0.856499in}}%
\pgfpathclose%
\pgfusepath{stroke,fill}%
\end{pgfscope}%
\begin{pgfscope}%
\pgfpathrectangle{\pgfqpoint{1.414736in}{0.548769in}}{\pgfqpoint{3.497747in}{7.253186in}}%
\pgfusepath{clip}%
\pgfsetbuttcap%
\pgfsetroundjoin%
\definecolor{currentfill}{rgb}{0.369214,0.788888,0.382914}%
\pgfsetfillcolor{currentfill}%
\pgfsetlinewidth{1.003750pt}%
\definecolor{currentstroke}{rgb}{0.369214,0.788888,0.382914}%
\pgfsetstrokecolor{currentstroke}%
\pgfsetdash{}{0pt}%
\pgfpathmoveto{\pgfqpoint{3.807591in}{1.327485in}}%
\pgfpathcurveto{\pgfqpoint{3.813414in}{1.327485in}}{\pgfqpoint{3.819001in}{1.329799in}}{\pgfqpoint{3.823119in}{1.333917in}}%
\pgfpathcurveto{\pgfqpoint{3.827237in}{1.338035in}}{\pgfqpoint{3.829551in}{1.343621in}}{\pgfqpoint{3.829551in}{1.349445in}}%
\pgfpathcurveto{\pgfqpoint{3.829551in}{1.355269in}}{\pgfqpoint{3.827237in}{1.360855in}}{\pgfqpoint{3.823119in}{1.364973in}}%
\pgfpathcurveto{\pgfqpoint{3.819001in}{1.369092in}}{\pgfqpoint{3.813414in}{1.371405in}}{\pgfqpoint{3.807591in}{1.371405in}}%
\pgfpathcurveto{\pgfqpoint{3.801767in}{1.371405in}}{\pgfqpoint{3.796180in}{1.369092in}}{\pgfqpoint{3.792062in}{1.364973in}}%
\pgfpathcurveto{\pgfqpoint{3.787944in}{1.360855in}}{\pgfqpoint{3.785630in}{1.355269in}}{\pgfqpoint{3.785630in}{1.349445in}}%
\pgfpathcurveto{\pgfqpoint{3.785630in}{1.343621in}}{\pgfqpoint{3.787944in}{1.338035in}}{\pgfqpoint{3.792062in}{1.333917in}}%
\pgfpathcurveto{\pgfqpoint{3.796180in}{1.329799in}}{\pgfqpoint{3.801767in}{1.327485in}}{\pgfqpoint{3.807591in}{1.327485in}}%
\pgfpathlineto{\pgfqpoint{3.807591in}{1.327485in}}%
\pgfpathclose%
\pgfusepath{stroke,fill}%
\end{pgfscope}%
\begin{pgfscope}%
\pgfpathrectangle{\pgfqpoint{1.414736in}{0.548769in}}{\pgfqpoint{3.497747in}{7.253186in}}%
\pgfusepath{clip}%
\pgfsetbuttcap%
\pgfsetroundjoin%
\definecolor{currentfill}{rgb}{0.699415,0.867117,0.175971}%
\pgfsetfillcolor{currentfill}%
\pgfsetlinewidth{1.003750pt}%
\definecolor{currentstroke}{rgb}{0.699415,0.867117,0.175971}%
\pgfsetstrokecolor{currentstroke}%
\pgfsetdash{}{0pt}%
\pgfpathmoveto{\pgfqpoint{3.734093in}{1.798471in}}%
\pgfpathcurveto{\pgfqpoint{3.739917in}{1.798471in}}{\pgfqpoint{3.745503in}{1.800785in}}{\pgfqpoint{3.749621in}{1.804903in}}%
\pgfpathcurveto{\pgfqpoint{3.753739in}{1.809021in}}{\pgfqpoint{3.756053in}{1.814607in}}{\pgfqpoint{3.756053in}{1.820431in}}%
\pgfpathcurveto{\pgfqpoint{3.756053in}{1.826255in}}{\pgfqpoint{3.753739in}{1.831841in}}{\pgfqpoint{3.749621in}{1.835960in}}%
\pgfpathcurveto{\pgfqpoint{3.745503in}{1.840078in}}{\pgfqpoint{3.739917in}{1.842392in}}{\pgfqpoint{3.734093in}{1.842392in}}%
\pgfpathcurveto{\pgfqpoint{3.728269in}{1.842392in}}{\pgfqpoint{3.722683in}{1.840078in}}{\pgfqpoint{3.718565in}{1.835960in}}%
\pgfpathcurveto{\pgfqpoint{3.714447in}{1.831841in}}{\pgfqpoint{3.712133in}{1.826255in}}{\pgfqpoint{3.712133in}{1.820431in}}%
\pgfpathcurveto{\pgfqpoint{3.712133in}{1.814607in}}{\pgfqpoint{3.714447in}{1.809021in}}{\pgfqpoint{3.718565in}{1.804903in}}%
\pgfpathcurveto{\pgfqpoint{3.722683in}{1.800785in}}{\pgfqpoint{3.728269in}{1.798471in}}{\pgfqpoint{3.734093in}{1.798471in}}%
\pgfpathlineto{\pgfqpoint{3.734093in}{1.798471in}}%
\pgfpathclose%
\pgfusepath{stroke,fill}%
\end{pgfscope}%
\begin{pgfscope}%
\pgfpathrectangle{\pgfqpoint{1.414736in}{0.548769in}}{\pgfqpoint{3.497747in}{7.253186in}}%
\pgfusepath{clip}%
\pgfsetbuttcap%
\pgfsetroundjoin%
\definecolor{currentfill}{rgb}{0.140210,0.665859,0.513427}%
\pgfsetfillcolor{currentfill}%
\pgfsetlinewidth{1.003750pt}%
\definecolor{currentstroke}{rgb}{0.140210,0.665859,0.513427}%
\pgfsetstrokecolor{currentstroke}%
\pgfsetdash{}{0pt}%
\pgfpathmoveto{\pgfqpoint{3.295926in}{2.269457in}}%
\pgfpathcurveto{\pgfqpoint{3.301750in}{2.269457in}}{\pgfqpoint{3.307336in}{2.271771in}}{\pgfqpoint{3.311455in}{2.275889in}}%
\pgfpathcurveto{\pgfqpoint{3.315573in}{2.280007in}}{\pgfqpoint{3.317887in}{2.285593in}}{\pgfqpoint{3.317887in}{2.291417in}}%
\pgfpathcurveto{\pgfqpoint{3.317887in}{2.297241in}}{\pgfqpoint{3.315573in}{2.302828in}}{\pgfqpoint{3.311455in}{2.306946in}}%
\pgfpathcurveto{\pgfqpoint{3.307336in}{2.311064in}}{\pgfqpoint{3.301750in}{2.313378in}}{\pgfqpoint{3.295926in}{2.313378in}}%
\pgfpathcurveto{\pgfqpoint{3.290102in}{2.313378in}}{\pgfqpoint{3.284516in}{2.311064in}}{\pgfqpoint{3.280398in}{2.306946in}}%
\pgfpathcurveto{\pgfqpoint{3.276280in}{2.302828in}}{\pgfqpoint{3.273966in}{2.297241in}}{\pgfqpoint{3.273966in}{2.291417in}}%
\pgfpathcurveto{\pgfqpoint{3.273966in}{2.285593in}}{\pgfqpoint{3.276280in}{2.280007in}}{\pgfqpoint{3.280398in}{2.275889in}}%
\pgfpathcurveto{\pgfqpoint{3.284516in}{2.271771in}}{\pgfqpoint{3.290102in}{2.269457in}}{\pgfqpoint{3.295926in}{2.269457in}}%
\pgfpathlineto{\pgfqpoint{3.295926in}{2.269457in}}%
\pgfpathclose%
\pgfusepath{stroke,fill}%
\end{pgfscope}%
\begin{pgfscope}%
\pgfpathrectangle{\pgfqpoint{1.414736in}{0.548769in}}{\pgfqpoint{3.497747in}{7.253186in}}%
\pgfusepath{clip}%
\pgfsetbuttcap%
\pgfsetroundjoin%
\definecolor{currentfill}{rgb}{0.360741,0.785964,0.387814}%
\pgfsetfillcolor{currentfill}%
\pgfsetlinewidth{1.003750pt}%
\definecolor{currentstroke}{rgb}{0.360741,0.785964,0.387814}%
\pgfsetstrokecolor{currentstroke}%
\pgfsetdash{}{0pt}%
\pgfpathmoveto{\pgfqpoint{3.869374in}{2.740443in}}%
\pgfpathcurveto{\pgfqpoint{3.875198in}{2.740443in}}{\pgfqpoint{3.880784in}{2.742757in}}{\pgfqpoint{3.884902in}{2.746875in}}%
\pgfpathcurveto{\pgfqpoint{3.889020in}{2.750993in}}{\pgfqpoint{3.891334in}{2.756580in}}{\pgfqpoint{3.891334in}{2.762404in}}%
\pgfpathcurveto{\pgfqpoint{3.891334in}{2.768227in}}{\pgfqpoint{3.889020in}{2.773814in}}{\pgfqpoint{3.884902in}{2.777932in}}%
\pgfpathcurveto{\pgfqpoint{3.880784in}{2.782050in}}{\pgfqpoint{3.875198in}{2.784364in}}{\pgfqpoint{3.869374in}{2.784364in}}%
\pgfpathcurveto{\pgfqpoint{3.863550in}{2.784364in}}{\pgfqpoint{3.857963in}{2.782050in}}{\pgfqpoint{3.853845in}{2.777932in}}%
\pgfpathcurveto{\pgfqpoint{3.849727in}{2.773814in}}{\pgfqpoint{3.847413in}{2.768227in}}{\pgfqpoint{3.847413in}{2.762404in}}%
\pgfpathcurveto{\pgfqpoint{3.847413in}{2.756580in}}{\pgfqpoint{3.849727in}{2.750993in}}{\pgfqpoint{3.853845in}{2.746875in}}%
\pgfpathcurveto{\pgfqpoint{3.857963in}{2.742757in}}{\pgfqpoint{3.863550in}{2.740443in}}{\pgfqpoint{3.869374in}{2.740443in}}%
\pgfpathlineto{\pgfqpoint{3.869374in}{2.740443in}}%
\pgfpathclose%
\pgfusepath{stroke,fill}%
\end{pgfscope}%
\begin{pgfscope}%
\pgfpathrectangle{\pgfqpoint{1.414736in}{0.548769in}}{\pgfqpoint{3.497747in}{7.253186in}}%
\pgfusepath{clip}%
\pgfsetbuttcap%
\pgfsetroundjoin%
\definecolor{currentfill}{rgb}{0.626579,0.854645,0.223353}%
\pgfsetfillcolor{currentfill}%
\pgfsetlinewidth{1.003750pt}%
\definecolor{currentstroke}{rgb}{0.626579,0.854645,0.223353}%
\pgfsetstrokecolor{currentstroke}%
\pgfsetdash{}{0pt}%
\pgfpathmoveto{\pgfqpoint{3.889161in}{3.211429in}}%
\pgfpathcurveto{\pgfqpoint{3.894984in}{3.211429in}}{\pgfqpoint{3.900571in}{3.213743in}}{\pgfqpoint{3.904689in}{3.217861in}}%
\pgfpathcurveto{\pgfqpoint{3.908807in}{3.221980in}}{\pgfqpoint{3.911121in}{3.227566in}}{\pgfqpoint{3.911121in}{3.233390in}}%
\pgfpathcurveto{\pgfqpoint{3.911121in}{3.239214in}}{\pgfqpoint{3.908807in}{3.244800in}}{\pgfqpoint{3.904689in}{3.248918in}}%
\pgfpathcurveto{\pgfqpoint{3.900571in}{3.253036in}}{\pgfqpoint{3.894984in}{3.255350in}}{\pgfqpoint{3.889161in}{3.255350in}}%
\pgfpathcurveto{\pgfqpoint{3.883337in}{3.255350in}}{\pgfqpoint{3.877750in}{3.253036in}}{\pgfqpoint{3.873632in}{3.248918in}}%
\pgfpathcurveto{\pgfqpoint{3.869514in}{3.244800in}}{\pgfqpoint{3.867200in}{3.239214in}}{\pgfqpoint{3.867200in}{3.233390in}}%
\pgfpathcurveto{\pgfqpoint{3.867200in}{3.227566in}}{\pgfqpoint{3.869514in}{3.221980in}}{\pgfqpoint{3.873632in}{3.217861in}}%
\pgfpathcurveto{\pgfqpoint{3.877750in}{3.213743in}}{\pgfqpoint{3.883337in}{3.211429in}}{\pgfqpoint{3.889161in}{3.211429in}}%
\pgfpathlineto{\pgfqpoint{3.889161in}{3.211429in}}%
\pgfpathclose%
\pgfusepath{stroke,fill}%
\end{pgfscope}%
\begin{pgfscope}%
\pgfpathrectangle{\pgfqpoint{1.414736in}{0.548769in}}{\pgfqpoint{3.497747in}{7.253186in}}%
\pgfusepath{clip}%
\pgfsetbuttcap%
\pgfsetroundjoin%
\definecolor{currentfill}{rgb}{0.983868,0.904867,0.136897}%
\pgfsetfillcolor{currentfill}%
\pgfsetlinewidth{1.003750pt}%
\definecolor{currentstroke}{rgb}{0.983868,0.904867,0.136897}%
\pgfsetstrokecolor{currentstroke}%
\pgfsetdash{}{0pt}%
\pgfpathmoveto{\pgfqpoint{3.134413in}{3.682415in}}%
\pgfpathcurveto{\pgfqpoint{3.140237in}{3.682415in}}{\pgfqpoint{3.145823in}{3.684729in}}{\pgfqpoint{3.149941in}{3.688847in}}%
\pgfpathcurveto{\pgfqpoint{3.154059in}{3.692966in}}{\pgfqpoint{3.156373in}{3.698552in}}{\pgfqpoint{3.156373in}{3.704376in}}%
\pgfpathcurveto{\pgfqpoint{3.156373in}{3.710200in}}{\pgfqpoint{3.154059in}{3.715786in}}{\pgfqpoint{3.149941in}{3.719904in}}%
\pgfpathcurveto{\pgfqpoint{3.145823in}{3.724022in}}{\pgfqpoint{3.140237in}{3.726336in}}{\pgfqpoint{3.134413in}{3.726336in}}%
\pgfpathcurveto{\pgfqpoint{3.128589in}{3.726336in}}{\pgfqpoint{3.123002in}{3.724022in}}{\pgfqpoint{3.118884in}{3.719904in}}%
\pgfpathcurveto{\pgfqpoint{3.114766in}{3.715786in}}{\pgfqpoint{3.112452in}{3.710200in}}{\pgfqpoint{3.112452in}{3.704376in}}%
\pgfpathcurveto{\pgfqpoint{3.112452in}{3.698552in}}{\pgfqpoint{3.114766in}{3.692966in}}{\pgfqpoint{3.118884in}{3.688847in}}%
\pgfpathcurveto{\pgfqpoint{3.123002in}{3.684729in}}{\pgfqpoint{3.128589in}{3.682415in}}{\pgfqpoint{3.134413in}{3.682415in}}%
\pgfpathlineto{\pgfqpoint{3.134413in}{3.682415in}}%
\pgfpathclose%
\pgfusepath{stroke,fill}%
\end{pgfscope}%
\begin{pgfscope}%
\pgfpathrectangle{\pgfqpoint{1.414736in}{0.548769in}}{\pgfqpoint{3.497747in}{7.253186in}}%
\pgfusepath{clip}%
\pgfsetbuttcap%
\pgfsetroundjoin%
\definecolor{currentfill}{rgb}{0.993248,0.906157,0.143936}%
\pgfsetfillcolor{currentfill}%
\pgfsetlinewidth{1.003750pt}%
\definecolor{currentstroke}{rgb}{0.993248,0.906157,0.143936}%
\pgfsetstrokecolor{currentstroke}%
\pgfsetdash{}{0pt}%
\pgfpathmoveto{\pgfqpoint{4.261476in}{4.153402in}}%
\pgfpathcurveto{\pgfqpoint{4.267300in}{4.153402in}}{\pgfqpoint{4.272886in}{4.155715in}}{\pgfqpoint{4.277004in}{4.159834in}}%
\pgfpathcurveto{\pgfqpoint{4.281123in}{4.163952in}}{\pgfqpoint{4.283436in}{4.169538in}}{\pgfqpoint{4.283436in}{4.175362in}}%
\pgfpathcurveto{\pgfqpoint{4.283436in}{4.181186in}}{\pgfqpoint{4.281123in}{4.186772in}}{\pgfqpoint{4.277004in}{4.190890in}}%
\pgfpathcurveto{\pgfqpoint{4.272886in}{4.195008in}}{\pgfqpoint{4.267300in}{4.197322in}}{\pgfqpoint{4.261476in}{4.197322in}}%
\pgfpathcurveto{\pgfqpoint{4.255652in}{4.197322in}}{\pgfqpoint{4.250066in}{4.195008in}}{\pgfqpoint{4.245948in}{4.190890in}}%
\pgfpathcurveto{\pgfqpoint{4.241830in}{4.186772in}}{\pgfqpoint{4.239516in}{4.181186in}}{\pgfqpoint{4.239516in}{4.175362in}}%
\pgfpathcurveto{\pgfqpoint{4.239516in}{4.169538in}}{\pgfqpoint{4.241830in}{4.163952in}}{\pgfqpoint{4.245948in}{4.159834in}}%
\pgfpathcurveto{\pgfqpoint{4.250066in}{4.155715in}}{\pgfqpoint{4.255652in}{4.153402in}}{\pgfqpoint{4.261476in}{4.153402in}}%
\pgfpathlineto{\pgfqpoint{4.261476in}{4.153402in}}%
\pgfpathclose%
\pgfusepath{stroke,fill}%
\end{pgfscope}%
\begin{pgfscope}%
\pgfpathrectangle{\pgfqpoint{1.414736in}{0.548769in}}{\pgfqpoint{3.497747in}{7.253186in}}%
\pgfusepath{clip}%
\pgfsetbuttcap%
\pgfsetroundjoin%
\definecolor{currentfill}{rgb}{0.688944,0.865448,0.182725}%
\pgfsetfillcolor{currentfill}%
\pgfsetlinewidth{1.003750pt}%
\definecolor{currentstroke}{rgb}{0.688944,0.865448,0.182725}%
\pgfsetstrokecolor{currentstroke}%
\pgfsetdash{}{0pt}%
\pgfpathmoveto{\pgfqpoint{3.040250in}{4.624388in}}%
\pgfpathcurveto{\pgfqpoint{3.046074in}{4.624388in}}{\pgfqpoint{3.051660in}{4.626702in}}{\pgfqpoint{3.055779in}{4.630820in}}%
\pgfpathcurveto{\pgfqpoint{3.059897in}{4.634938in}}{\pgfqpoint{3.062211in}{4.640524in}}{\pgfqpoint{3.062211in}{4.646348in}}%
\pgfpathcurveto{\pgfqpoint{3.062211in}{4.652172in}}{\pgfqpoint{3.059897in}{4.657758in}}{\pgfqpoint{3.055779in}{4.661876in}}%
\pgfpathcurveto{\pgfqpoint{3.051660in}{4.665994in}}{\pgfqpoint{3.046074in}{4.668308in}}{\pgfqpoint{3.040250in}{4.668308in}}%
\pgfpathcurveto{\pgfqpoint{3.034426in}{4.668308in}}{\pgfqpoint{3.028840in}{4.665994in}}{\pgfqpoint{3.024722in}{4.661876in}}%
\pgfpathcurveto{\pgfqpoint{3.020604in}{4.657758in}}{\pgfqpoint{3.018290in}{4.652172in}}{\pgfqpoint{3.018290in}{4.646348in}}%
\pgfpathcurveto{\pgfqpoint{3.018290in}{4.640524in}}{\pgfqpoint{3.020604in}{4.634938in}}{\pgfqpoint{3.024722in}{4.630820in}}%
\pgfpathcurveto{\pgfqpoint{3.028840in}{4.626702in}}{\pgfqpoint{3.034426in}{4.624388in}}{\pgfqpoint{3.040250in}{4.624388in}}%
\pgfpathlineto{\pgfqpoint{3.040250in}{4.624388in}}%
\pgfpathclose%
\pgfusepath{stroke,fill}%
\end{pgfscope}%
\begin{pgfscope}%
\pgfpathrectangle{\pgfqpoint{1.414736in}{0.548769in}}{\pgfqpoint{3.497747in}{7.253186in}}%
\pgfusepath{clip}%
\pgfsetbuttcap%
\pgfsetroundjoin%
\definecolor{currentfill}{rgb}{0.283197,0.115680,0.436115}%
\pgfsetfillcolor{currentfill}%
\pgfsetlinewidth{1.003750pt}%
\definecolor{currentstroke}{rgb}{0.283197,0.115680,0.436115}%
\pgfsetstrokecolor{currentstroke}%
\pgfsetdash{}{0pt}%
\pgfpathmoveto{\pgfqpoint{3.130114in}{5.095374in}}%
\pgfpathcurveto{\pgfqpoint{3.135938in}{5.095374in}}{\pgfqpoint{3.141524in}{5.097688in}}{\pgfqpoint{3.145643in}{5.101806in}}%
\pgfpathcurveto{\pgfqpoint{3.149761in}{5.105924in}}{\pgfqpoint{3.152075in}{5.111510in}}{\pgfqpoint{3.152075in}{5.117334in}}%
\pgfpathcurveto{\pgfqpoint{3.152075in}{5.123158in}}{\pgfqpoint{3.149761in}{5.128744in}}{\pgfqpoint{3.145643in}{5.132862in}}%
\pgfpathcurveto{\pgfqpoint{3.141524in}{5.136980in}}{\pgfqpoint{3.135938in}{5.139294in}}{\pgfqpoint{3.130114in}{5.139294in}}%
\pgfpathcurveto{\pgfqpoint{3.124290in}{5.139294in}}{\pgfqpoint{3.118704in}{5.136980in}}{\pgfqpoint{3.114586in}{5.132862in}}%
\pgfpathcurveto{\pgfqpoint{3.110468in}{5.128744in}}{\pgfqpoint{3.108154in}{5.123158in}}{\pgfqpoint{3.108154in}{5.117334in}}%
\pgfpathcurveto{\pgfqpoint{3.108154in}{5.111510in}}{\pgfqpoint{3.110468in}{5.105924in}}{\pgfqpoint{3.114586in}{5.101806in}}%
\pgfpathcurveto{\pgfqpoint{3.118704in}{5.097688in}}{\pgfqpoint{3.124290in}{5.095374in}}{\pgfqpoint{3.130114in}{5.095374in}}%
\pgfpathlineto{\pgfqpoint{3.130114in}{5.095374in}}%
\pgfpathclose%
\pgfusepath{stroke,fill}%
\end{pgfscope}%
\begin{pgfscope}%
\pgfpathrectangle{\pgfqpoint{1.414736in}{0.548769in}}{\pgfqpoint{3.497747in}{7.253186in}}%
\pgfusepath{clip}%
\pgfsetbuttcap%
\pgfsetroundjoin%
\definecolor{currentfill}{rgb}{0.206756,0.371758,0.553117}%
\pgfsetfillcolor{currentfill}%
\pgfsetlinewidth{1.003750pt}%
\definecolor{currentstroke}{rgb}{0.206756,0.371758,0.553117}%
\pgfsetstrokecolor{currentstroke}%
\pgfsetdash{}{0pt}%
\pgfpathmoveto{\pgfqpoint{3.778160in}{5.566360in}}%
\pgfpathcurveto{\pgfqpoint{3.783984in}{5.566360in}}{\pgfqpoint{3.789570in}{5.568674in}}{\pgfqpoint{3.793688in}{5.572792in}}%
\pgfpathcurveto{\pgfqpoint{3.797807in}{5.576910in}}{\pgfqpoint{3.800120in}{5.582496in}}{\pgfqpoint{3.800120in}{5.588320in}}%
\pgfpathcurveto{\pgfqpoint{3.800120in}{5.594144in}}{\pgfqpoint{3.797807in}{5.599730in}}{\pgfqpoint{3.793688in}{5.603848in}}%
\pgfpathcurveto{\pgfqpoint{3.789570in}{5.607967in}}{\pgfqpoint{3.783984in}{5.610280in}}{\pgfqpoint{3.778160in}{5.610280in}}%
\pgfpathcurveto{\pgfqpoint{3.772336in}{5.610280in}}{\pgfqpoint{3.766750in}{5.607967in}}{\pgfqpoint{3.762632in}{5.603848in}}%
\pgfpathcurveto{\pgfqpoint{3.758514in}{5.599730in}}{\pgfqpoint{3.756200in}{5.594144in}}{\pgfqpoint{3.756200in}{5.588320in}}%
\pgfpathcurveto{\pgfqpoint{3.756200in}{5.582496in}}{\pgfqpoint{3.758514in}{5.576910in}}{\pgfqpoint{3.762632in}{5.572792in}}%
\pgfpathcurveto{\pgfqpoint{3.766750in}{5.568674in}}{\pgfqpoint{3.772336in}{5.566360in}}{\pgfqpoint{3.778160in}{5.566360in}}%
\pgfpathlineto{\pgfqpoint{3.778160in}{5.566360in}}%
\pgfpathclose%
\pgfusepath{stroke,fill}%
\end{pgfscope}%
\begin{pgfscope}%
\pgfpathrectangle{\pgfqpoint{1.414736in}{0.548769in}}{\pgfqpoint{3.497747in}{7.253186in}}%
\pgfusepath{clip}%
\pgfsetbuttcap%
\pgfsetroundjoin%
\definecolor{currentfill}{rgb}{0.227802,0.326594,0.546532}%
\pgfsetfillcolor{currentfill}%
\pgfsetlinewidth{1.003750pt}%
\definecolor{currentstroke}{rgb}{0.227802,0.326594,0.546532}%
\pgfsetstrokecolor{currentstroke}%
\pgfsetdash{}{0pt}%
\pgfpathmoveto{\pgfqpoint{4.477401in}{6.037346in}}%
\pgfpathcurveto{\pgfqpoint{4.483225in}{6.037346in}}{\pgfqpoint{4.488811in}{6.039660in}}{\pgfqpoint{4.492929in}{6.043778in}}%
\pgfpathcurveto{\pgfqpoint{4.497048in}{6.047896in}}{\pgfqpoint{4.499361in}{6.053482in}}{\pgfqpoint{4.499361in}{6.059306in}}%
\pgfpathcurveto{\pgfqpoint{4.499361in}{6.065130in}}{\pgfqpoint{4.497048in}{6.070716in}}{\pgfqpoint{4.492929in}{6.074835in}}%
\pgfpathcurveto{\pgfqpoint{4.488811in}{6.078953in}}{\pgfqpoint{4.483225in}{6.081267in}}{\pgfqpoint{4.477401in}{6.081267in}}%
\pgfpathcurveto{\pgfqpoint{4.471577in}{6.081267in}}{\pgfqpoint{4.465991in}{6.078953in}}{\pgfqpoint{4.461873in}{6.074835in}}%
\pgfpathcurveto{\pgfqpoint{4.457755in}{6.070716in}}{\pgfqpoint{4.455441in}{6.065130in}}{\pgfqpoint{4.455441in}{6.059306in}}%
\pgfpathcurveto{\pgfqpoint{4.455441in}{6.053482in}}{\pgfqpoint{4.457755in}{6.047896in}}{\pgfqpoint{4.461873in}{6.043778in}}%
\pgfpathcurveto{\pgfqpoint{4.465991in}{6.039660in}}{\pgfqpoint{4.471577in}{6.037346in}}{\pgfqpoint{4.477401in}{6.037346in}}%
\pgfpathlineto{\pgfqpoint{4.477401in}{6.037346in}}%
\pgfpathclose%
\pgfusepath{stroke,fill}%
\end{pgfscope}%
\begin{pgfscope}%
\pgfpathrectangle{\pgfqpoint{1.414736in}{0.548769in}}{\pgfqpoint{3.497747in}{7.253186in}}%
\pgfusepath{clip}%
\pgfsetbuttcap%
\pgfsetroundjoin%
\definecolor{currentfill}{rgb}{0.993248,0.906157,0.143936}%
\pgfsetfillcolor{currentfill}%
\pgfsetlinewidth{1.003750pt}%
\definecolor{currentstroke}{rgb}{0.993248,0.906157,0.143936}%
\pgfsetstrokecolor{currentstroke}%
\pgfsetdash{}{0pt}%
\pgfpathmoveto{\pgfqpoint{3.445994in}{6.508332in}}%
\pgfpathcurveto{\pgfqpoint{3.451818in}{6.508332in}}{\pgfqpoint{3.457404in}{6.510646in}}{\pgfqpoint{3.461522in}{6.514764in}}%
\pgfpathcurveto{\pgfqpoint{3.465640in}{6.518882in}}{\pgfqpoint{3.467954in}{6.524468in}}{\pgfqpoint{3.467954in}{6.530292in}}%
\pgfpathcurveto{\pgfqpoint{3.467954in}{6.536116in}}{\pgfqpoint{3.465640in}{6.541702in}}{\pgfqpoint{3.461522in}{6.545821in}}%
\pgfpathcurveto{\pgfqpoint{3.457404in}{6.549939in}}{\pgfqpoint{3.451818in}{6.552253in}}{\pgfqpoint{3.445994in}{6.552253in}}%
\pgfpathcurveto{\pgfqpoint{3.440170in}{6.552253in}}{\pgfqpoint{3.434584in}{6.549939in}}{\pgfqpoint{3.430466in}{6.545821in}}%
\pgfpathcurveto{\pgfqpoint{3.426348in}{6.541702in}}{\pgfqpoint{3.424034in}{6.536116in}}{\pgfqpoint{3.424034in}{6.530292in}}%
\pgfpathcurveto{\pgfqpoint{3.424034in}{6.524468in}}{\pgfqpoint{3.426348in}{6.518882in}}{\pgfqpoint{3.430466in}{6.514764in}}%
\pgfpathcurveto{\pgfqpoint{3.434584in}{6.510646in}}{\pgfqpoint{3.440170in}{6.508332in}}{\pgfqpoint{3.445994in}{6.508332in}}%
\pgfpathlineto{\pgfqpoint{3.445994in}{6.508332in}}%
\pgfpathclose%
\pgfusepath{stroke,fill}%
\end{pgfscope}%
\begin{pgfscope}%
\pgfpathrectangle{\pgfqpoint{1.414736in}{0.548769in}}{\pgfqpoint{3.497747in}{7.253186in}}%
\pgfusepath{clip}%
\pgfsetbuttcap%
\pgfsetroundjoin%
\definecolor{currentfill}{rgb}{0.430983,0.808473,0.346476}%
\pgfsetfillcolor{currentfill}%
\pgfsetlinewidth{1.003750pt}%
\definecolor{currentstroke}{rgb}{0.430983,0.808473,0.346476}%
\pgfsetstrokecolor{currentstroke}%
\pgfsetdash{}{0pt}%
\pgfpathmoveto{\pgfqpoint{2.203162in}{6.979318in}}%
\pgfpathcurveto{\pgfqpoint{2.208986in}{6.979318in}}{\pgfqpoint{2.214572in}{6.981632in}}{\pgfqpoint{2.218690in}{6.985750in}}%
\pgfpathcurveto{\pgfqpoint{2.222808in}{6.989868in}}{\pgfqpoint{2.225122in}{6.995455in}}{\pgfqpoint{2.225122in}{7.001278in}}%
\pgfpathcurveto{\pgfqpoint{2.225122in}{7.007102in}}{\pgfqpoint{2.222808in}{7.012689in}}{\pgfqpoint{2.218690in}{7.016807in}}%
\pgfpathcurveto{\pgfqpoint{2.214572in}{7.020925in}}{\pgfqpoint{2.208986in}{7.023239in}}{\pgfqpoint{2.203162in}{7.023239in}}%
\pgfpathcurveto{\pgfqpoint{2.197338in}{7.023239in}}{\pgfqpoint{2.191752in}{7.020925in}}{\pgfqpoint{2.187634in}{7.016807in}}%
\pgfpathcurveto{\pgfqpoint{2.183516in}{7.012689in}}{\pgfqpoint{2.181202in}{7.007102in}}{\pgfqpoint{2.181202in}{7.001278in}}%
\pgfpathcurveto{\pgfqpoint{2.181202in}{6.995455in}}{\pgfqpoint{2.183516in}{6.989868in}}{\pgfqpoint{2.187634in}{6.985750in}}%
\pgfpathcurveto{\pgfqpoint{2.191752in}{6.981632in}}{\pgfqpoint{2.197338in}{6.979318in}}{\pgfqpoint{2.203162in}{6.979318in}}%
\pgfpathlineto{\pgfqpoint{2.203162in}{6.979318in}}%
\pgfpathclose%
\pgfusepath{stroke,fill}%
\end{pgfscope}%
\begin{pgfscope}%
\pgfpathrectangle{\pgfqpoint{1.414736in}{0.548769in}}{\pgfqpoint{3.497747in}{7.253186in}}%
\pgfusepath{clip}%
\pgfsetbuttcap%
\pgfsetroundjoin%
\definecolor{currentfill}{rgb}{0.626579,0.854645,0.223353}%
\pgfsetfillcolor{currentfill}%
\pgfsetlinewidth{1.003750pt}%
\definecolor{currentstroke}{rgb}{0.626579,0.854645,0.223353}%
\pgfsetstrokecolor{currentstroke}%
\pgfsetdash{}{0pt}%
\pgfpathmoveto{\pgfqpoint{4.224841in}{7.450304in}}%
\pgfpathcurveto{\pgfqpoint{4.230665in}{7.450304in}}{\pgfqpoint{4.236251in}{7.452618in}}{\pgfqpoint{4.240369in}{7.456736in}}%
\pgfpathcurveto{\pgfqpoint{4.244487in}{7.460854in}}{\pgfqpoint{4.246801in}{7.466441in}}{\pgfqpoint{4.246801in}{7.472265in}}%
\pgfpathcurveto{\pgfqpoint{4.246801in}{7.478089in}}{\pgfqpoint{4.244487in}{7.483675in}}{\pgfqpoint{4.240369in}{7.487793in}}%
\pgfpathcurveto{\pgfqpoint{4.236251in}{7.491911in}}{\pgfqpoint{4.230665in}{7.494225in}}{\pgfqpoint{4.224841in}{7.494225in}}%
\pgfpathcurveto{\pgfqpoint{4.219017in}{7.494225in}}{\pgfqpoint{4.213431in}{7.491911in}}{\pgfqpoint{4.209313in}{7.487793in}}%
\pgfpathcurveto{\pgfqpoint{4.205195in}{7.483675in}}{\pgfqpoint{4.202881in}{7.478089in}}{\pgfqpoint{4.202881in}{7.472265in}}%
\pgfpathcurveto{\pgfqpoint{4.202881in}{7.466441in}}{\pgfqpoint{4.205195in}{7.460854in}}{\pgfqpoint{4.209313in}{7.456736in}}%
\pgfpathcurveto{\pgfqpoint{4.213431in}{7.452618in}}{\pgfqpoint{4.219017in}{7.450304in}}{\pgfqpoint{4.224841in}{7.450304in}}%
\pgfpathlineto{\pgfqpoint{4.224841in}{7.450304in}}%
\pgfpathclose%
\pgfusepath{stroke,fill}%
\end{pgfscope}%
\begin{pgfscope}%
\pgfpathrectangle{\pgfqpoint{1.414736in}{0.548769in}}{\pgfqpoint{3.497747in}{7.253186in}}%
\pgfusepath{clip}%
\pgfsetbuttcap%
\pgfsetroundjoin%
\definecolor{currentfill}{rgb}{0.993248,0.906157,0.143936}%
\pgfsetfillcolor{currentfill}%
\pgfsetlinewidth{1.003750pt}%
\definecolor{currentstroke}{rgb}{0.993248,0.906157,0.143936}%
\pgfsetstrokecolor{currentstroke}%
\pgfsetdash{}{0pt}%
\pgfpathmoveto{\pgfqpoint{3.416016in}{0.856499in}}%
\pgfpathcurveto{\pgfqpoint{3.421840in}{0.856499in}}{\pgfqpoint{3.427426in}{0.858813in}}{\pgfqpoint{3.431544in}{0.862931in}}%
\pgfpathcurveto{\pgfqpoint{3.435663in}{0.867049in}}{\pgfqpoint{3.437976in}{0.872635in}}{\pgfqpoint{3.437976in}{0.878459in}}%
\pgfpathcurveto{\pgfqpoint{3.437976in}{0.884283in}}{\pgfqpoint{3.435663in}{0.889869in}}{\pgfqpoint{3.431544in}{0.893987in}}%
\pgfpathcurveto{\pgfqpoint{3.427426in}{0.898105in}}{\pgfqpoint{3.421840in}{0.900419in}}{\pgfqpoint{3.416016in}{0.900419in}}%
\pgfpathcurveto{\pgfqpoint{3.410192in}{0.900419in}}{\pgfqpoint{3.404606in}{0.898105in}}{\pgfqpoint{3.400488in}{0.893987in}}%
\pgfpathcurveto{\pgfqpoint{3.396370in}{0.889869in}}{\pgfqpoint{3.394056in}{0.884283in}}{\pgfqpoint{3.394056in}{0.878459in}}%
\pgfpathcurveto{\pgfqpoint{3.394056in}{0.872635in}}{\pgfqpoint{3.396370in}{0.867049in}}{\pgfqpoint{3.400488in}{0.862931in}}%
\pgfpathcurveto{\pgfqpoint{3.404606in}{0.858813in}}{\pgfqpoint{3.410192in}{0.856499in}}{\pgfqpoint{3.416016in}{0.856499in}}%
\pgfpathlineto{\pgfqpoint{3.416016in}{0.856499in}}%
\pgfpathclose%
\pgfusepath{stroke,fill}%
\end{pgfscope}%
\begin{pgfscope}%
\pgfpathrectangle{\pgfqpoint{1.414736in}{0.548769in}}{\pgfqpoint{3.497747in}{7.253186in}}%
\pgfusepath{clip}%
\pgfsetbuttcap%
\pgfsetroundjoin%
\definecolor{currentfill}{rgb}{0.616293,0.852709,0.230052}%
\pgfsetfillcolor{currentfill}%
\pgfsetlinewidth{1.003750pt}%
\definecolor{currentstroke}{rgb}{0.616293,0.852709,0.230052}%
\pgfsetstrokecolor{currentstroke}%
\pgfsetdash{}{0pt}%
\pgfpathmoveto{\pgfqpoint{3.489122in}{1.327485in}}%
\pgfpathcurveto{\pgfqpoint{3.494946in}{1.327485in}}{\pgfqpoint{3.500532in}{1.329799in}}{\pgfqpoint{3.504650in}{1.333917in}}%
\pgfpathcurveto{\pgfqpoint{3.508769in}{1.338035in}}{\pgfqpoint{3.511082in}{1.343621in}}{\pgfqpoint{3.511082in}{1.349445in}}%
\pgfpathcurveto{\pgfqpoint{3.511082in}{1.355269in}}{\pgfqpoint{3.508769in}{1.360855in}}{\pgfqpoint{3.504650in}{1.364973in}}%
\pgfpathcurveto{\pgfqpoint{3.500532in}{1.369092in}}{\pgfqpoint{3.494946in}{1.371405in}}{\pgfqpoint{3.489122in}{1.371405in}}%
\pgfpathcurveto{\pgfqpoint{3.483298in}{1.371405in}}{\pgfqpoint{3.477712in}{1.369092in}}{\pgfqpoint{3.473594in}{1.364973in}}%
\pgfpathcurveto{\pgfqpoint{3.469476in}{1.360855in}}{\pgfqpoint{3.467162in}{1.355269in}}{\pgfqpoint{3.467162in}{1.349445in}}%
\pgfpathcurveto{\pgfqpoint{3.467162in}{1.343621in}}{\pgfqpoint{3.469476in}{1.338035in}}{\pgfqpoint{3.473594in}{1.333917in}}%
\pgfpathcurveto{\pgfqpoint{3.477712in}{1.329799in}}{\pgfqpoint{3.483298in}{1.327485in}}{\pgfqpoint{3.489122in}{1.327485in}}%
\pgfpathlineto{\pgfqpoint{3.489122in}{1.327485in}}%
\pgfpathclose%
\pgfusepath{stroke,fill}%
\end{pgfscope}%
\begin{pgfscope}%
\pgfpathrectangle{\pgfqpoint{1.414736in}{0.548769in}}{\pgfqpoint{3.497747in}{7.253186in}}%
\pgfusepath{clip}%
\pgfsetbuttcap%
\pgfsetroundjoin%
\definecolor{currentfill}{rgb}{0.395174,0.797475,0.367757}%
\pgfsetfillcolor{currentfill}%
\pgfsetlinewidth{1.003750pt}%
\definecolor{currentstroke}{rgb}{0.395174,0.797475,0.367757}%
\pgfsetstrokecolor{currentstroke}%
\pgfsetdash{}{0pt}%
\pgfpathmoveto{\pgfqpoint{3.198033in}{1.798471in}}%
\pgfpathcurveto{\pgfqpoint{3.203857in}{1.798471in}}{\pgfqpoint{3.209443in}{1.800785in}}{\pgfqpoint{3.213561in}{1.804903in}}%
\pgfpathcurveto{\pgfqpoint{3.217679in}{1.809021in}}{\pgfqpoint{3.219993in}{1.814607in}}{\pgfqpoint{3.219993in}{1.820431in}}%
\pgfpathcurveto{\pgfqpoint{3.219993in}{1.826255in}}{\pgfqpoint{3.217679in}{1.831841in}}{\pgfqpoint{3.213561in}{1.835960in}}%
\pgfpathcurveto{\pgfqpoint{3.209443in}{1.840078in}}{\pgfqpoint{3.203857in}{1.842392in}}{\pgfqpoint{3.198033in}{1.842392in}}%
\pgfpathcurveto{\pgfqpoint{3.192209in}{1.842392in}}{\pgfqpoint{3.186623in}{1.840078in}}{\pgfqpoint{3.182505in}{1.835960in}}%
\pgfpathcurveto{\pgfqpoint{3.178386in}{1.831841in}}{\pgfqpoint{3.176073in}{1.826255in}}{\pgfqpoint{3.176073in}{1.820431in}}%
\pgfpathcurveto{\pgfqpoint{3.176073in}{1.814607in}}{\pgfqpoint{3.178386in}{1.809021in}}{\pgfqpoint{3.182505in}{1.804903in}}%
\pgfpathcurveto{\pgfqpoint{3.186623in}{1.800785in}}{\pgfqpoint{3.192209in}{1.798471in}}{\pgfqpoint{3.198033in}{1.798471in}}%
\pgfpathlineto{\pgfqpoint{3.198033in}{1.798471in}}%
\pgfpathclose%
\pgfusepath{stroke,fill}%
\end{pgfscope}%
\begin{pgfscope}%
\pgfpathrectangle{\pgfqpoint{1.414736in}{0.548769in}}{\pgfqpoint{3.497747in}{7.253186in}}%
\pgfusepath{clip}%
\pgfsetbuttcap%
\pgfsetroundjoin%
\definecolor{currentfill}{rgb}{0.993248,0.906157,0.143936}%
\pgfsetfillcolor{currentfill}%
\pgfsetlinewidth{1.003750pt}%
\definecolor{currentstroke}{rgb}{0.993248,0.906157,0.143936}%
\pgfsetstrokecolor{currentstroke}%
\pgfsetdash{}{0pt}%
\pgfpathmoveto{\pgfqpoint{3.554984in}{2.269457in}}%
\pgfpathcurveto{\pgfqpoint{3.560808in}{2.269457in}}{\pgfqpoint{3.566394in}{2.271771in}}{\pgfqpoint{3.570512in}{2.275889in}}%
\pgfpathcurveto{\pgfqpoint{3.574630in}{2.280007in}}{\pgfqpoint{3.576944in}{2.285593in}}{\pgfqpoint{3.576944in}{2.291417in}}%
\pgfpathcurveto{\pgfqpoint{3.576944in}{2.297241in}}{\pgfqpoint{3.574630in}{2.302828in}}{\pgfqpoint{3.570512in}{2.306946in}}%
\pgfpathcurveto{\pgfqpoint{3.566394in}{2.311064in}}{\pgfqpoint{3.560808in}{2.313378in}}{\pgfqpoint{3.554984in}{2.313378in}}%
\pgfpathcurveto{\pgfqpoint{3.549160in}{2.313378in}}{\pgfqpoint{3.543574in}{2.311064in}}{\pgfqpoint{3.539456in}{2.306946in}}%
\pgfpathcurveto{\pgfqpoint{3.535338in}{2.302828in}}{\pgfqpoint{3.533024in}{2.297241in}}{\pgfqpoint{3.533024in}{2.291417in}}%
\pgfpathcurveto{\pgfqpoint{3.533024in}{2.285593in}}{\pgfqpoint{3.535338in}{2.280007in}}{\pgfqpoint{3.539456in}{2.275889in}}%
\pgfpathcurveto{\pgfqpoint{3.543574in}{2.271771in}}{\pgfqpoint{3.549160in}{2.269457in}}{\pgfqpoint{3.554984in}{2.269457in}}%
\pgfpathlineto{\pgfqpoint{3.554984in}{2.269457in}}%
\pgfpathclose%
\pgfusepath{stroke,fill}%
\end{pgfscope}%
\begin{pgfscope}%
\pgfpathrectangle{\pgfqpoint{1.414736in}{0.548769in}}{\pgfqpoint{3.497747in}{7.253186in}}%
\pgfusepath{clip}%
\pgfsetbuttcap%
\pgfsetroundjoin%
\definecolor{currentfill}{rgb}{0.274149,0.751988,0.436601}%
\pgfsetfillcolor{currentfill}%
\pgfsetlinewidth{1.003750pt}%
\definecolor{currentstroke}{rgb}{0.274149,0.751988,0.436601}%
\pgfsetstrokecolor{currentstroke}%
\pgfsetdash{}{0pt}%
\pgfpathmoveto{\pgfqpoint{3.866174in}{2.740443in}}%
\pgfpathcurveto{\pgfqpoint{3.871998in}{2.740443in}}{\pgfqpoint{3.877584in}{2.742757in}}{\pgfqpoint{3.881702in}{2.746875in}}%
\pgfpathcurveto{\pgfqpoint{3.885820in}{2.750993in}}{\pgfqpoint{3.888134in}{2.756580in}}{\pgfqpoint{3.888134in}{2.762404in}}%
\pgfpathcurveto{\pgfqpoint{3.888134in}{2.768227in}}{\pgfqpoint{3.885820in}{2.773814in}}{\pgfqpoint{3.881702in}{2.777932in}}%
\pgfpathcurveto{\pgfqpoint{3.877584in}{2.782050in}}{\pgfqpoint{3.871998in}{2.784364in}}{\pgfqpoint{3.866174in}{2.784364in}}%
\pgfpathcurveto{\pgfqpoint{3.860350in}{2.784364in}}{\pgfqpoint{3.854764in}{2.782050in}}{\pgfqpoint{3.850646in}{2.777932in}}%
\pgfpathcurveto{\pgfqpoint{3.846528in}{2.773814in}}{\pgfqpoint{3.844214in}{2.768227in}}{\pgfqpoint{3.844214in}{2.762404in}}%
\pgfpathcurveto{\pgfqpoint{3.844214in}{2.756580in}}{\pgfqpoint{3.846528in}{2.750993in}}{\pgfqpoint{3.850646in}{2.746875in}}%
\pgfpathcurveto{\pgfqpoint{3.854764in}{2.742757in}}{\pgfqpoint{3.860350in}{2.740443in}}{\pgfqpoint{3.866174in}{2.740443in}}%
\pgfpathlineto{\pgfqpoint{3.866174in}{2.740443in}}%
\pgfpathclose%
\pgfusepath{stroke,fill}%
\end{pgfscope}%
\begin{pgfscope}%
\pgfpathrectangle{\pgfqpoint{1.414736in}{0.548769in}}{\pgfqpoint{3.497747in}{7.253186in}}%
\pgfusepath{clip}%
\pgfsetbuttcap%
\pgfsetroundjoin%
\definecolor{currentfill}{rgb}{0.595839,0.848717,0.243329}%
\pgfsetfillcolor{currentfill}%
\pgfsetlinewidth{1.003750pt}%
\definecolor{currentstroke}{rgb}{0.595839,0.848717,0.243329}%
\pgfsetstrokecolor{currentstroke}%
\pgfsetdash{}{0pt}%
\pgfpathmoveto{\pgfqpoint{3.414612in}{3.211429in}}%
\pgfpathcurveto{\pgfqpoint{3.420436in}{3.211429in}}{\pgfqpoint{3.426022in}{3.213743in}}{\pgfqpoint{3.430140in}{3.217861in}}%
\pgfpathcurveto{\pgfqpoint{3.434258in}{3.221980in}}{\pgfqpoint{3.436572in}{3.227566in}}{\pgfqpoint{3.436572in}{3.233390in}}%
\pgfpathcurveto{\pgfqpoint{3.436572in}{3.239214in}}{\pgfqpoint{3.434258in}{3.244800in}}{\pgfqpoint{3.430140in}{3.248918in}}%
\pgfpathcurveto{\pgfqpoint{3.426022in}{3.253036in}}{\pgfqpoint{3.420436in}{3.255350in}}{\pgfqpoint{3.414612in}{3.255350in}}%
\pgfpathcurveto{\pgfqpoint{3.408788in}{3.255350in}}{\pgfqpoint{3.403202in}{3.253036in}}{\pgfqpoint{3.399084in}{3.248918in}}%
\pgfpathcurveto{\pgfqpoint{3.394966in}{3.244800in}}{\pgfqpoint{3.392652in}{3.239214in}}{\pgfqpoint{3.392652in}{3.233390in}}%
\pgfpathcurveto{\pgfqpoint{3.392652in}{3.227566in}}{\pgfqpoint{3.394966in}{3.221980in}}{\pgfqpoint{3.399084in}{3.217861in}}%
\pgfpathcurveto{\pgfqpoint{3.403202in}{3.213743in}}{\pgfqpoint{3.408788in}{3.211429in}}{\pgfqpoint{3.414612in}{3.211429in}}%
\pgfpathlineto{\pgfqpoint{3.414612in}{3.211429in}}%
\pgfpathclose%
\pgfusepath{stroke,fill}%
\end{pgfscope}%
\begin{pgfscope}%
\pgfpathrectangle{\pgfqpoint{1.414736in}{0.548769in}}{\pgfqpoint{3.497747in}{7.253186in}}%
\pgfusepath{clip}%
\pgfsetbuttcap%
\pgfsetroundjoin%
\definecolor{currentfill}{rgb}{0.993248,0.906157,0.143936}%
\pgfsetfillcolor{currentfill}%
\pgfsetlinewidth{1.003750pt}%
\definecolor{currentstroke}{rgb}{0.993248,0.906157,0.143936}%
\pgfsetstrokecolor{currentstroke}%
\pgfsetdash{}{0pt}%
\pgfpathmoveto{\pgfqpoint{4.010885in}{3.682415in}}%
\pgfpathcurveto{\pgfqpoint{4.016709in}{3.682415in}}{\pgfqpoint{4.022295in}{3.684729in}}{\pgfqpoint{4.026414in}{3.688847in}}%
\pgfpathcurveto{\pgfqpoint{4.030532in}{3.692966in}}{\pgfqpoint{4.032846in}{3.698552in}}{\pgfqpoint{4.032846in}{3.704376in}}%
\pgfpathcurveto{\pgfqpoint{4.032846in}{3.710200in}}{\pgfqpoint{4.030532in}{3.715786in}}{\pgfqpoint{4.026414in}{3.719904in}}%
\pgfpathcurveto{\pgfqpoint{4.022295in}{3.724022in}}{\pgfqpoint{4.016709in}{3.726336in}}{\pgfqpoint{4.010885in}{3.726336in}}%
\pgfpathcurveto{\pgfqpoint{4.005061in}{3.726336in}}{\pgfqpoint{3.999475in}{3.724022in}}{\pgfqpoint{3.995357in}{3.719904in}}%
\pgfpathcurveto{\pgfqpoint{3.991239in}{3.715786in}}{\pgfqpoint{3.988925in}{3.710200in}}{\pgfqpoint{3.988925in}{3.704376in}}%
\pgfpathcurveto{\pgfqpoint{3.988925in}{3.698552in}}{\pgfqpoint{3.991239in}{3.692966in}}{\pgfqpoint{3.995357in}{3.688847in}}%
\pgfpathcurveto{\pgfqpoint{3.999475in}{3.684729in}}{\pgfqpoint{4.005061in}{3.682415in}}{\pgfqpoint{4.010885in}{3.682415in}}%
\pgfpathlineto{\pgfqpoint{4.010885in}{3.682415in}}%
\pgfpathclose%
\pgfusepath{stroke,fill}%
\end{pgfscope}%
\begin{pgfscope}%
\pgfpathrectangle{\pgfqpoint{1.414736in}{0.548769in}}{\pgfqpoint{3.497747in}{7.253186in}}%
\pgfusepath{clip}%
\pgfsetbuttcap%
\pgfsetroundjoin%
\definecolor{currentfill}{rgb}{0.993248,0.906157,0.143936}%
\pgfsetfillcolor{currentfill}%
\pgfsetlinewidth{1.003750pt}%
\definecolor{currentstroke}{rgb}{0.993248,0.906157,0.143936}%
\pgfsetstrokecolor{currentstroke}%
\pgfsetdash{}{0pt}%
\pgfpathmoveto{\pgfqpoint{3.691433in}{4.153402in}}%
\pgfpathcurveto{\pgfqpoint{3.697257in}{4.153402in}}{\pgfqpoint{3.702844in}{4.155715in}}{\pgfqpoint{3.706962in}{4.159834in}}%
\pgfpathcurveto{\pgfqpoint{3.711080in}{4.163952in}}{\pgfqpoint{3.713394in}{4.169538in}}{\pgfqpoint{3.713394in}{4.175362in}}%
\pgfpathcurveto{\pgfqpoint{3.713394in}{4.181186in}}{\pgfqpoint{3.711080in}{4.186772in}}{\pgfqpoint{3.706962in}{4.190890in}}%
\pgfpathcurveto{\pgfqpoint{3.702844in}{4.195008in}}{\pgfqpoint{3.697257in}{4.197322in}}{\pgfqpoint{3.691433in}{4.197322in}}%
\pgfpathcurveto{\pgfqpoint{3.685609in}{4.197322in}}{\pgfqpoint{3.680023in}{4.195008in}}{\pgfqpoint{3.675905in}{4.190890in}}%
\pgfpathcurveto{\pgfqpoint{3.671787in}{4.186772in}}{\pgfqpoint{3.669473in}{4.181186in}}{\pgfqpoint{3.669473in}{4.175362in}}%
\pgfpathcurveto{\pgfqpoint{3.669473in}{4.169538in}}{\pgfqpoint{3.671787in}{4.163952in}}{\pgfqpoint{3.675905in}{4.159834in}}%
\pgfpathcurveto{\pgfqpoint{3.680023in}{4.155715in}}{\pgfqpoint{3.685609in}{4.153402in}}{\pgfqpoint{3.691433in}{4.153402in}}%
\pgfpathlineto{\pgfqpoint{3.691433in}{4.153402in}}%
\pgfpathclose%
\pgfusepath{stroke,fill}%
\end{pgfscope}%
\begin{pgfscope}%
\pgfpathrectangle{\pgfqpoint{1.414736in}{0.548769in}}{\pgfqpoint{3.497747in}{7.253186in}}%
\pgfusepath{clip}%
\pgfsetbuttcap%
\pgfsetroundjoin%
\definecolor{currentfill}{rgb}{0.896320,0.893616,0.096335}%
\pgfsetfillcolor{currentfill}%
\pgfsetlinewidth{1.003750pt}%
\definecolor{currentstroke}{rgb}{0.896320,0.893616,0.096335}%
\pgfsetstrokecolor{currentstroke}%
\pgfsetdash{}{0pt}%
\pgfpathmoveto{\pgfqpoint{3.581616in}{4.624388in}}%
\pgfpathcurveto{\pgfqpoint{3.587440in}{4.624388in}}{\pgfqpoint{3.593026in}{4.626702in}}{\pgfqpoint{3.597144in}{4.630820in}}%
\pgfpathcurveto{\pgfqpoint{3.601262in}{4.634938in}}{\pgfqpoint{3.603576in}{4.640524in}}{\pgfqpoint{3.603576in}{4.646348in}}%
\pgfpathcurveto{\pgfqpoint{3.603576in}{4.652172in}}{\pgfqpoint{3.601262in}{4.657758in}}{\pgfqpoint{3.597144in}{4.661876in}}%
\pgfpathcurveto{\pgfqpoint{3.593026in}{4.665994in}}{\pgfqpoint{3.587440in}{4.668308in}}{\pgfqpoint{3.581616in}{4.668308in}}%
\pgfpathcurveto{\pgfqpoint{3.575792in}{4.668308in}}{\pgfqpoint{3.570206in}{4.665994in}}{\pgfqpoint{3.566088in}{4.661876in}}%
\pgfpathcurveto{\pgfqpoint{3.561969in}{4.657758in}}{\pgfqpoint{3.559656in}{4.652172in}}{\pgfqpoint{3.559656in}{4.646348in}}%
\pgfpathcurveto{\pgfqpoint{3.559656in}{4.640524in}}{\pgfqpoint{3.561969in}{4.634938in}}{\pgfqpoint{3.566088in}{4.630820in}}%
\pgfpathcurveto{\pgfqpoint{3.570206in}{4.626702in}}{\pgfqpoint{3.575792in}{4.624388in}}{\pgfqpoint{3.581616in}{4.624388in}}%
\pgfpathlineto{\pgfqpoint{3.581616in}{4.624388in}}%
\pgfpathclose%
\pgfusepath{stroke,fill}%
\end{pgfscope}%
\begin{pgfscope}%
\pgfpathrectangle{\pgfqpoint{1.414736in}{0.548769in}}{\pgfqpoint{3.497747in}{7.253186in}}%
\pgfusepath{clip}%
\pgfsetbuttcap%
\pgfsetroundjoin%
\definecolor{currentfill}{rgb}{0.144759,0.519093,0.556572}%
\pgfsetfillcolor{currentfill}%
\pgfsetlinewidth{1.003750pt}%
\definecolor{currentstroke}{rgb}{0.144759,0.519093,0.556572}%
\pgfsetstrokecolor{currentstroke}%
\pgfsetdash{}{0pt}%
\pgfpathmoveto{\pgfqpoint{4.185756in}{5.095374in}}%
\pgfpathcurveto{\pgfqpoint{4.191580in}{5.095374in}}{\pgfqpoint{4.197166in}{5.097688in}}{\pgfqpoint{4.201284in}{5.101806in}}%
\pgfpathcurveto{\pgfqpoint{4.205402in}{5.105924in}}{\pgfqpoint{4.207716in}{5.111510in}}{\pgfqpoint{4.207716in}{5.117334in}}%
\pgfpathcurveto{\pgfqpoint{4.207716in}{5.123158in}}{\pgfqpoint{4.205402in}{5.128744in}}{\pgfqpoint{4.201284in}{5.132862in}}%
\pgfpathcurveto{\pgfqpoint{4.197166in}{5.136980in}}{\pgfqpoint{4.191580in}{5.139294in}}{\pgfqpoint{4.185756in}{5.139294in}}%
\pgfpathcurveto{\pgfqpoint{4.179932in}{5.139294in}}{\pgfqpoint{4.174346in}{5.136980in}}{\pgfqpoint{4.170228in}{5.132862in}}%
\pgfpathcurveto{\pgfqpoint{4.166109in}{5.128744in}}{\pgfqpoint{4.163796in}{5.123158in}}{\pgfqpoint{4.163796in}{5.117334in}}%
\pgfpathcurveto{\pgfqpoint{4.163796in}{5.111510in}}{\pgfqpoint{4.166109in}{5.105924in}}{\pgfqpoint{4.170228in}{5.101806in}}%
\pgfpathcurveto{\pgfqpoint{4.174346in}{5.097688in}}{\pgfqpoint{4.179932in}{5.095374in}}{\pgfqpoint{4.185756in}{5.095374in}}%
\pgfpathlineto{\pgfqpoint{4.185756in}{5.095374in}}%
\pgfpathclose%
\pgfusepath{stroke,fill}%
\end{pgfscope}%
\begin{pgfscope}%
\pgfpathrectangle{\pgfqpoint{1.414736in}{0.548769in}}{\pgfqpoint{3.497747in}{7.253186in}}%
\pgfusepath{clip}%
\pgfsetbuttcap%
\pgfsetroundjoin%
\definecolor{currentfill}{rgb}{0.159194,0.482237,0.558073}%
\pgfsetfillcolor{currentfill}%
\pgfsetlinewidth{1.003750pt}%
\definecolor{currentstroke}{rgb}{0.159194,0.482237,0.558073}%
\pgfsetstrokecolor{currentstroke}%
\pgfsetdash{}{0pt}%
\pgfpathmoveto{\pgfqpoint{3.775517in}{5.566360in}}%
\pgfpathcurveto{\pgfqpoint{3.781341in}{5.566360in}}{\pgfqpoint{3.786927in}{5.568674in}}{\pgfqpoint{3.791045in}{5.572792in}}%
\pgfpathcurveto{\pgfqpoint{3.795163in}{5.576910in}}{\pgfqpoint{3.797477in}{5.582496in}}{\pgfqpoint{3.797477in}{5.588320in}}%
\pgfpathcurveto{\pgfqpoint{3.797477in}{5.594144in}}{\pgfqpoint{3.795163in}{5.599730in}}{\pgfqpoint{3.791045in}{5.603848in}}%
\pgfpathcurveto{\pgfqpoint{3.786927in}{5.607967in}}{\pgfqpoint{3.781341in}{5.610280in}}{\pgfqpoint{3.775517in}{5.610280in}}%
\pgfpathcurveto{\pgfqpoint{3.769693in}{5.610280in}}{\pgfqpoint{3.764107in}{5.607967in}}{\pgfqpoint{3.759989in}{5.603848in}}%
\pgfpathcurveto{\pgfqpoint{3.755871in}{5.599730in}}{\pgfqpoint{3.753557in}{5.594144in}}{\pgfqpoint{3.753557in}{5.588320in}}%
\pgfpathcurveto{\pgfqpoint{3.753557in}{5.582496in}}{\pgfqpoint{3.755871in}{5.576910in}}{\pgfqpoint{3.759989in}{5.572792in}}%
\pgfpathcurveto{\pgfqpoint{3.764107in}{5.568674in}}{\pgfqpoint{3.769693in}{5.566360in}}{\pgfqpoint{3.775517in}{5.566360in}}%
\pgfpathlineto{\pgfqpoint{3.775517in}{5.566360in}}%
\pgfpathclose%
\pgfusepath{stroke,fill}%
\end{pgfscope}%
\begin{pgfscope}%
\pgfpathrectangle{\pgfqpoint{1.414736in}{0.548769in}}{\pgfqpoint{3.497747in}{7.253186in}}%
\pgfusepath{clip}%
\pgfsetbuttcap%
\pgfsetroundjoin%
\definecolor{currentfill}{rgb}{0.440137,0.811138,0.340967}%
\pgfsetfillcolor{currentfill}%
\pgfsetlinewidth{1.003750pt}%
\definecolor{currentstroke}{rgb}{0.440137,0.811138,0.340967}%
\pgfsetstrokecolor{currentstroke}%
\pgfsetdash{}{0pt}%
\pgfpathmoveto{\pgfqpoint{3.052549in}{6.037346in}}%
\pgfpathcurveto{\pgfqpoint{3.058373in}{6.037346in}}{\pgfqpoint{3.063960in}{6.039660in}}{\pgfqpoint{3.068078in}{6.043778in}}%
\pgfpathcurveto{\pgfqpoint{3.072196in}{6.047896in}}{\pgfqpoint{3.074510in}{6.053482in}}{\pgfqpoint{3.074510in}{6.059306in}}%
\pgfpathcurveto{\pgfqpoint{3.074510in}{6.065130in}}{\pgfqpoint{3.072196in}{6.070716in}}{\pgfqpoint{3.068078in}{6.074835in}}%
\pgfpathcurveto{\pgfqpoint{3.063960in}{6.078953in}}{\pgfqpoint{3.058373in}{6.081267in}}{\pgfqpoint{3.052549in}{6.081267in}}%
\pgfpathcurveto{\pgfqpoint{3.046726in}{6.081267in}}{\pgfqpoint{3.041139in}{6.078953in}}{\pgfqpoint{3.037021in}{6.074835in}}%
\pgfpathcurveto{\pgfqpoint{3.032903in}{6.070716in}}{\pgfqpoint{3.030589in}{6.065130in}}{\pgfqpoint{3.030589in}{6.059306in}}%
\pgfpathcurveto{\pgfqpoint{3.030589in}{6.053482in}}{\pgfqpoint{3.032903in}{6.047896in}}{\pgfqpoint{3.037021in}{6.043778in}}%
\pgfpathcurveto{\pgfqpoint{3.041139in}{6.039660in}}{\pgfqpoint{3.046726in}{6.037346in}}{\pgfqpoint{3.052549in}{6.037346in}}%
\pgfpathlineto{\pgfqpoint{3.052549in}{6.037346in}}%
\pgfpathclose%
\pgfusepath{stroke,fill}%
\end{pgfscope}%
\begin{pgfscope}%
\pgfpathrectangle{\pgfqpoint{1.414736in}{0.548769in}}{\pgfqpoint{3.497747in}{7.253186in}}%
\pgfusepath{clip}%
\pgfsetbuttcap%
\pgfsetroundjoin%
\definecolor{currentfill}{rgb}{0.993248,0.906157,0.143936}%
\pgfsetfillcolor{currentfill}%
\pgfsetlinewidth{1.003750pt}%
\definecolor{currentstroke}{rgb}{0.993248,0.906157,0.143936}%
\pgfsetstrokecolor{currentstroke}%
\pgfsetdash{}{0pt}%
\pgfpathmoveto{\pgfqpoint{3.983580in}{6.508332in}}%
\pgfpathcurveto{\pgfqpoint{3.989404in}{6.508332in}}{\pgfqpoint{3.994990in}{6.510646in}}{\pgfqpoint{3.999109in}{6.514764in}}%
\pgfpathcurveto{\pgfqpoint{4.003227in}{6.518882in}}{\pgfqpoint{4.005541in}{6.524468in}}{\pgfqpoint{4.005541in}{6.530292in}}%
\pgfpathcurveto{\pgfqpoint{4.005541in}{6.536116in}}{\pgfqpoint{4.003227in}{6.541702in}}{\pgfqpoint{3.999109in}{6.545821in}}%
\pgfpathcurveto{\pgfqpoint{3.994990in}{6.549939in}}{\pgfqpoint{3.989404in}{6.552253in}}{\pgfqpoint{3.983580in}{6.552253in}}%
\pgfpathcurveto{\pgfqpoint{3.977756in}{6.552253in}}{\pgfqpoint{3.972170in}{6.549939in}}{\pgfqpoint{3.968052in}{6.545821in}}%
\pgfpathcurveto{\pgfqpoint{3.963934in}{6.541702in}}{\pgfqpoint{3.961620in}{6.536116in}}{\pgfqpoint{3.961620in}{6.530292in}}%
\pgfpathcurveto{\pgfqpoint{3.961620in}{6.524468in}}{\pgfqpoint{3.963934in}{6.518882in}}{\pgfqpoint{3.968052in}{6.514764in}}%
\pgfpathcurveto{\pgfqpoint{3.972170in}{6.510646in}}{\pgfqpoint{3.977756in}{6.508332in}}{\pgfqpoint{3.983580in}{6.508332in}}%
\pgfpathlineto{\pgfqpoint{3.983580in}{6.508332in}}%
\pgfpathclose%
\pgfusepath{stroke,fill}%
\end{pgfscope}%
\begin{pgfscope}%
\pgfpathrectangle{\pgfqpoint{1.414736in}{0.548769in}}{\pgfqpoint{3.497747in}{7.253186in}}%
\pgfusepath{clip}%
\pgfsetbuttcap%
\pgfsetroundjoin%
\definecolor{currentfill}{rgb}{0.146616,0.673050,0.508936}%
\pgfsetfillcolor{currentfill}%
\pgfsetlinewidth{1.003750pt}%
\definecolor{currentstroke}{rgb}{0.146616,0.673050,0.508936}%
\pgfsetstrokecolor{currentstroke}%
\pgfsetdash{}{0pt}%
\pgfpathmoveto{\pgfqpoint{3.650382in}{6.979318in}}%
\pgfpathcurveto{\pgfqpoint{3.656206in}{6.979318in}}{\pgfqpoint{3.661792in}{6.981632in}}{\pgfqpoint{3.665910in}{6.985750in}}%
\pgfpathcurveto{\pgfqpoint{3.670028in}{6.989868in}}{\pgfqpoint{3.672342in}{6.995455in}}{\pgfqpoint{3.672342in}{7.001278in}}%
\pgfpathcurveto{\pgfqpoint{3.672342in}{7.007102in}}{\pgfqpoint{3.670028in}{7.012689in}}{\pgfqpoint{3.665910in}{7.016807in}}%
\pgfpathcurveto{\pgfqpoint{3.661792in}{7.020925in}}{\pgfqpoint{3.656206in}{7.023239in}}{\pgfqpoint{3.650382in}{7.023239in}}%
\pgfpathcurveto{\pgfqpoint{3.644558in}{7.023239in}}{\pgfqpoint{3.638972in}{7.020925in}}{\pgfqpoint{3.634853in}{7.016807in}}%
\pgfpathcurveto{\pgfqpoint{3.630735in}{7.012689in}}{\pgfqpoint{3.628421in}{7.007102in}}{\pgfqpoint{3.628421in}{7.001278in}}%
\pgfpathcurveto{\pgfqpoint{3.628421in}{6.995455in}}{\pgfqpoint{3.630735in}{6.989868in}}{\pgfqpoint{3.634853in}{6.985750in}}%
\pgfpathcurveto{\pgfqpoint{3.638972in}{6.981632in}}{\pgfqpoint{3.644558in}{6.979318in}}{\pgfqpoint{3.650382in}{6.979318in}}%
\pgfpathlineto{\pgfqpoint{3.650382in}{6.979318in}}%
\pgfpathclose%
\pgfusepath{stroke,fill}%
\end{pgfscope}%
\begin{pgfscope}%
\pgfpathrectangle{\pgfqpoint{1.414736in}{0.548769in}}{\pgfqpoint{3.497747in}{7.253186in}}%
\pgfusepath{clip}%
\pgfsetbuttcap%
\pgfsetroundjoin%
\definecolor{currentfill}{rgb}{0.906311,0.894855,0.098125}%
\pgfsetfillcolor{currentfill}%
\pgfsetlinewidth{1.003750pt}%
\definecolor{currentstroke}{rgb}{0.906311,0.894855,0.098125}%
\pgfsetstrokecolor{currentstroke}%
\pgfsetdash{}{0pt}%
\pgfpathmoveto{\pgfqpoint{3.804665in}{7.450304in}}%
\pgfpathcurveto{\pgfqpoint{3.810489in}{7.450304in}}{\pgfqpoint{3.816075in}{7.452618in}}{\pgfqpoint{3.820193in}{7.456736in}}%
\pgfpathcurveto{\pgfqpoint{3.824312in}{7.460854in}}{\pgfqpoint{3.826625in}{7.466441in}}{\pgfqpoint{3.826625in}{7.472265in}}%
\pgfpathcurveto{\pgfqpoint{3.826625in}{7.478089in}}{\pgfqpoint{3.824312in}{7.483675in}}{\pgfqpoint{3.820193in}{7.487793in}}%
\pgfpathcurveto{\pgfqpoint{3.816075in}{7.491911in}}{\pgfqpoint{3.810489in}{7.494225in}}{\pgfqpoint{3.804665in}{7.494225in}}%
\pgfpathcurveto{\pgfqpoint{3.798841in}{7.494225in}}{\pgfqpoint{3.793255in}{7.491911in}}{\pgfqpoint{3.789137in}{7.487793in}}%
\pgfpathcurveto{\pgfqpoint{3.785019in}{7.483675in}}{\pgfqpoint{3.782705in}{7.478089in}}{\pgfqpoint{3.782705in}{7.472265in}}%
\pgfpathcurveto{\pgfqpoint{3.782705in}{7.466441in}}{\pgfqpoint{3.785019in}{7.460854in}}{\pgfqpoint{3.789137in}{7.456736in}}%
\pgfpathcurveto{\pgfqpoint{3.793255in}{7.452618in}}{\pgfqpoint{3.798841in}{7.450304in}}{\pgfqpoint{3.804665in}{7.450304in}}%
\pgfpathlineto{\pgfqpoint{3.804665in}{7.450304in}}%
\pgfpathclose%
\pgfusepath{stroke,fill}%
\end{pgfscope}%
\begin{pgfscope}%
\pgfpathrectangle{\pgfqpoint{1.414736in}{0.548769in}}{\pgfqpoint{3.497747in}{7.253186in}}%
\pgfusepath{clip}%
\pgfsetbuttcap%
\pgfsetroundjoin%
\definecolor{currentfill}{rgb}{0.751884,0.874951,0.143228}%
\pgfsetfillcolor{currentfill}%
\pgfsetlinewidth{1.003750pt}%
\definecolor{currentstroke}{rgb}{0.751884,0.874951,0.143228}%
\pgfsetstrokecolor{currentstroke}%
\pgfsetdash{}{0pt}%
\pgfpathmoveto{\pgfqpoint{3.734803in}{0.856499in}}%
\pgfpathcurveto{\pgfqpoint{3.740627in}{0.856499in}}{\pgfqpoint{3.746213in}{0.858813in}}{\pgfqpoint{3.750331in}{0.862931in}}%
\pgfpathcurveto{\pgfqpoint{3.754449in}{0.867049in}}{\pgfqpoint{3.756763in}{0.872635in}}{\pgfqpoint{3.756763in}{0.878459in}}%
\pgfpathcurveto{\pgfqpoint{3.756763in}{0.884283in}}{\pgfqpoint{3.754449in}{0.889869in}}{\pgfqpoint{3.750331in}{0.893987in}}%
\pgfpathcurveto{\pgfqpoint{3.746213in}{0.898105in}}{\pgfqpoint{3.740627in}{0.900419in}}{\pgfqpoint{3.734803in}{0.900419in}}%
\pgfpathcurveto{\pgfqpoint{3.728979in}{0.900419in}}{\pgfqpoint{3.723393in}{0.898105in}}{\pgfqpoint{3.719274in}{0.893987in}}%
\pgfpathcurveto{\pgfqpoint{3.715156in}{0.889869in}}{\pgfqpoint{3.712842in}{0.884283in}}{\pgfqpoint{3.712842in}{0.878459in}}%
\pgfpathcurveto{\pgfqpoint{3.712842in}{0.872635in}}{\pgfqpoint{3.715156in}{0.867049in}}{\pgfqpoint{3.719274in}{0.862931in}}%
\pgfpathcurveto{\pgfqpoint{3.723393in}{0.858813in}}{\pgfqpoint{3.728979in}{0.856499in}}{\pgfqpoint{3.734803in}{0.856499in}}%
\pgfpathlineto{\pgfqpoint{3.734803in}{0.856499in}}%
\pgfpathclose%
\pgfusepath{stroke,fill}%
\end{pgfscope}%
\begin{pgfscope}%
\pgfpathrectangle{\pgfqpoint{1.414736in}{0.548769in}}{\pgfqpoint{3.497747in}{7.253186in}}%
\pgfusepath{clip}%
\pgfsetbuttcap%
\pgfsetroundjoin%
\definecolor{currentfill}{rgb}{0.595839,0.848717,0.243329}%
\pgfsetfillcolor{currentfill}%
\pgfsetlinewidth{1.003750pt}%
\definecolor{currentstroke}{rgb}{0.595839,0.848717,0.243329}%
\pgfsetstrokecolor{currentstroke}%
\pgfsetdash{}{0pt}%
\pgfpathmoveto{\pgfqpoint{3.437593in}{1.327485in}}%
\pgfpathcurveto{\pgfqpoint{3.443417in}{1.327485in}}{\pgfqpoint{3.449003in}{1.329799in}}{\pgfqpoint{3.453121in}{1.333917in}}%
\pgfpathcurveto{\pgfqpoint{3.457239in}{1.338035in}}{\pgfqpoint{3.459553in}{1.343621in}}{\pgfqpoint{3.459553in}{1.349445in}}%
\pgfpathcurveto{\pgfqpoint{3.459553in}{1.355269in}}{\pgfqpoint{3.457239in}{1.360855in}}{\pgfqpoint{3.453121in}{1.364973in}}%
\pgfpathcurveto{\pgfqpoint{3.449003in}{1.369092in}}{\pgfqpoint{3.443417in}{1.371405in}}{\pgfqpoint{3.437593in}{1.371405in}}%
\pgfpathcurveto{\pgfqpoint{3.431769in}{1.371405in}}{\pgfqpoint{3.426183in}{1.369092in}}{\pgfqpoint{3.422065in}{1.364973in}}%
\pgfpathcurveto{\pgfqpoint{3.417946in}{1.360855in}}{\pgfqpoint{3.415633in}{1.355269in}}{\pgfqpoint{3.415633in}{1.349445in}}%
\pgfpathcurveto{\pgfqpoint{3.415633in}{1.343621in}}{\pgfqpoint{3.417946in}{1.338035in}}{\pgfqpoint{3.422065in}{1.333917in}}%
\pgfpathcurveto{\pgfqpoint{3.426183in}{1.329799in}}{\pgfqpoint{3.431769in}{1.327485in}}{\pgfqpoint{3.437593in}{1.327485in}}%
\pgfpathlineto{\pgfqpoint{3.437593in}{1.327485in}}%
\pgfpathclose%
\pgfusepath{stroke,fill}%
\end{pgfscope}%
\begin{pgfscope}%
\pgfpathrectangle{\pgfqpoint{1.414736in}{0.548769in}}{\pgfqpoint{3.497747in}{7.253186in}}%
\pgfusepath{clip}%
\pgfsetbuttcap%
\pgfsetroundjoin%
\definecolor{currentfill}{rgb}{0.699415,0.867117,0.175971}%
\pgfsetfillcolor{currentfill}%
\pgfsetlinewidth{1.003750pt}%
\definecolor{currentstroke}{rgb}{0.699415,0.867117,0.175971}%
\pgfsetstrokecolor{currentstroke}%
\pgfsetdash{}{0pt}%
\pgfpathmoveto{\pgfqpoint{3.754411in}{1.798471in}}%
\pgfpathcurveto{\pgfqpoint{3.760235in}{1.798471in}}{\pgfqpoint{3.765821in}{1.800785in}}{\pgfqpoint{3.769940in}{1.804903in}}%
\pgfpathcurveto{\pgfqpoint{3.774058in}{1.809021in}}{\pgfqpoint{3.776372in}{1.814607in}}{\pgfqpoint{3.776372in}{1.820431in}}%
\pgfpathcurveto{\pgfqpoint{3.776372in}{1.826255in}}{\pgfqpoint{3.774058in}{1.831841in}}{\pgfqpoint{3.769940in}{1.835960in}}%
\pgfpathcurveto{\pgfqpoint{3.765821in}{1.840078in}}{\pgfqpoint{3.760235in}{1.842392in}}{\pgfqpoint{3.754411in}{1.842392in}}%
\pgfpathcurveto{\pgfqpoint{3.748587in}{1.842392in}}{\pgfqpoint{3.743001in}{1.840078in}}{\pgfqpoint{3.738883in}{1.835960in}}%
\pgfpathcurveto{\pgfqpoint{3.734765in}{1.831841in}}{\pgfqpoint{3.732451in}{1.826255in}}{\pgfqpoint{3.732451in}{1.820431in}}%
\pgfpathcurveto{\pgfqpoint{3.732451in}{1.814607in}}{\pgfqpoint{3.734765in}{1.809021in}}{\pgfqpoint{3.738883in}{1.804903in}}%
\pgfpathcurveto{\pgfqpoint{3.743001in}{1.800785in}}{\pgfqpoint{3.748587in}{1.798471in}}{\pgfqpoint{3.754411in}{1.798471in}}%
\pgfpathlineto{\pgfqpoint{3.754411in}{1.798471in}}%
\pgfpathclose%
\pgfusepath{stroke,fill}%
\end{pgfscope}%
\begin{pgfscope}%
\pgfpathrectangle{\pgfqpoint{1.414736in}{0.548769in}}{\pgfqpoint{3.497747in}{7.253186in}}%
\pgfusepath{clip}%
\pgfsetbuttcap%
\pgfsetroundjoin%
\definecolor{currentfill}{rgb}{0.993248,0.906157,0.143936}%
\pgfsetfillcolor{currentfill}%
\pgfsetlinewidth{1.003750pt}%
\definecolor{currentstroke}{rgb}{0.993248,0.906157,0.143936}%
\pgfsetstrokecolor{currentstroke}%
\pgfsetdash{}{0pt}%
\pgfpathmoveto{\pgfqpoint{3.816695in}{2.269457in}}%
\pgfpathcurveto{\pgfqpoint{3.822519in}{2.269457in}}{\pgfqpoint{3.828106in}{2.271771in}}{\pgfqpoint{3.832224in}{2.275889in}}%
\pgfpathcurveto{\pgfqpoint{3.836342in}{2.280007in}}{\pgfqpoint{3.838656in}{2.285593in}}{\pgfqpoint{3.838656in}{2.291417in}}%
\pgfpathcurveto{\pgfqpoint{3.838656in}{2.297241in}}{\pgfqpoint{3.836342in}{2.302828in}}{\pgfqpoint{3.832224in}{2.306946in}}%
\pgfpathcurveto{\pgfqpoint{3.828106in}{2.311064in}}{\pgfqpoint{3.822519in}{2.313378in}}{\pgfqpoint{3.816695in}{2.313378in}}%
\pgfpathcurveto{\pgfqpoint{3.810872in}{2.313378in}}{\pgfqpoint{3.805285in}{2.311064in}}{\pgfqpoint{3.801167in}{2.306946in}}%
\pgfpathcurveto{\pgfqpoint{3.797049in}{2.302828in}}{\pgfqpoint{3.794735in}{2.297241in}}{\pgfqpoint{3.794735in}{2.291417in}}%
\pgfpathcurveto{\pgfqpoint{3.794735in}{2.285593in}}{\pgfqpoint{3.797049in}{2.280007in}}{\pgfqpoint{3.801167in}{2.275889in}}%
\pgfpathcurveto{\pgfqpoint{3.805285in}{2.271771in}}{\pgfqpoint{3.810872in}{2.269457in}}{\pgfqpoint{3.816695in}{2.269457in}}%
\pgfpathlineto{\pgfqpoint{3.816695in}{2.269457in}}%
\pgfpathclose%
\pgfusepath{stroke,fill}%
\end{pgfscope}%
\begin{pgfscope}%
\pgfpathrectangle{\pgfqpoint{1.414736in}{0.548769in}}{\pgfqpoint{3.497747in}{7.253186in}}%
\pgfusepath{clip}%
\pgfsetbuttcap%
\pgfsetroundjoin%
\definecolor{currentfill}{rgb}{0.252899,0.742211,0.448284}%
\pgfsetfillcolor{currentfill}%
\pgfsetlinewidth{1.003750pt}%
\definecolor{currentstroke}{rgb}{0.252899,0.742211,0.448284}%
\pgfsetstrokecolor{currentstroke}%
\pgfsetdash{}{0pt}%
\pgfpathmoveto{\pgfqpoint{3.848296in}{2.740443in}}%
\pgfpathcurveto{\pgfqpoint{3.854120in}{2.740443in}}{\pgfqpoint{3.859706in}{2.742757in}}{\pgfqpoint{3.863824in}{2.746875in}}%
\pgfpathcurveto{\pgfqpoint{3.867942in}{2.750993in}}{\pgfqpoint{3.870256in}{2.756580in}}{\pgfqpoint{3.870256in}{2.762404in}}%
\pgfpathcurveto{\pgfqpoint{3.870256in}{2.768227in}}{\pgfqpoint{3.867942in}{2.773814in}}{\pgfqpoint{3.863824in}{2.777932in}}%
\pgfpathcurveto{\pgfqpoint{3.859706in}{2.782050in}}{\pgfqpoint{3.854120in}{2.784364in}}{\pgfqpoint{3.848296in}{2.784364in}}%
\pgfpathcurveto{\pgfqpoint{3.842472in}{2.784364in}}{\pgfqpoint{3.836886in}{2.782050in}}{\pgfqpoint{3.832768in}{2.777932in}}%
\pgfpathcurveto{\pgfqpoint{3.828650in}{2.773814in}}{\pgfqpoint{3.826336in}{2.768227in}}{\pgfqpoint{3.826336in}{2.762404in}}%
\pgfpathcurveto{\pgfqpoint{3.826336in}{2.756580in}}{\pgfqpoint{3.828650in}{2.750993in}}{\pgfqpoint{3.832768in}{2.746875in}}%
\pgfpathcurveto{\pgfqpoint{3.836886in}{2.742757in}}{\pgfqpoint{3.842472in}{2.740443in}}{\pgfqpoint{3.848296in}{2.740443in}}%
\pgfpathlineto{\pgfqpoint{3.848296in}{2.740443in}}%
\pgfpathclose%
\pgfusepath{stroke,fill}%
\end{pgfscope}%
\begin{pgfscope}%
\pgfpathrectangle{\pgfqpoint{1.414736in}{0.548769in}}{\pgfqpoint{3.497747in}{7.253186in}}%
\pgfusepath{clip}%
\pgfsetbuttcap%
\pgfsetroundjoin%
\definecolor{currentfill}{rgb}{0.955300,0.901065,0.118128}%
\pgfsetfillcolor{currentfill}%
\pgfsetlinewidth{1.003750pt}%
\definecolor{currentstroke}{rgb}{0.955300,0.901065,0.118128}%
\pgfsetstrokecolor{currentstroke}%
\pgfsetdash{}{0pt}%
\pgfpathmoveto{\pgfqpoint{3.432139in}{3.211429in}}%
\pgfpathcurveto{\pgfqpoint{3.437962in}{3.211429in}}{\pgfqpoint{3.443549in}{3.213743in}}{\pgfqpoint{3.447667in}{3.217861in}}%
\pgfpathcurveto{\pgfqpoint{3.451785in}{3.221980in}}{\pgfqpoint{3.454099in}{3.227566in}}{\pgfqpoint{3.454099in}{3.233390in}}%
\pgfpathcurveto{\pgfqpoint{3.454099in}{3.239214in}}{\pgfqpoint{3.451785in}{3.244800in}}{\pgfqpoint{3.447667in}{3.248918in}}%
\pgfpathcurveto{\pgfqpoint{3.443549in}{3.253036in}}{\pgfqpoint{3.437962in}{3.255350in}}{\pgfqpoint{3.432139in}{3.255350in}}%
\pgfpathcurveto{\pgfqpoint{3.426315in}{3.255350in}}{\pgfqpoint{3.420728in}{3.253036in}}{\pgfqpoint{3.416610in}{3.248918in}}%
\pgfpathcurveto{\pgfqpoint{3.412492in}{3.244800in}}{\pgfqpoint{3.410178in}{3.239214in}}{\pgfqpoint{3.410178in}{3.233390in}}%
\pgfpathcurveto{\pgfqpoint{3.410178in}{3.227566in}}{\pgfqpoint{3.412492in}{3.221980in}}{\pgfqpoint{3.416610in}{3.217861in}}%
\pgfpathcurveto{\pgfqpoint{3.420728in}{3.213743in}}{\pgfqpoint{3.426315in}{3.211429in}}{\pgfqpoint{3.432139in}{3.211429in}}%
\pgfpathlineto{\pgfqpoint{3.432139in}{3.211429in}}%
\pgfpathclose%
\pgfusepath{stroke,fill}%
\end{pgfscope}%
\begin{pgfscope}%
\pgfpathrectangle{\pgfqpoint{1.414736in}{0.548769in}}{\pgfqpoint{3.497747in}{7.253186in}}%
\pgfusepath{clip}%
\pgfsetbuttcap%
\pgfsetroundjoin%
\definecolor{currentfill}{rgb}{0.855810,0.888601,0.097452}%
\pgfsetfillcolor{currentfill}%
\pgfsetlinewidth{1.003750pt}%
\definecolor{currentstroke}{rgb}{0.855810,0.888601,0.097452}%
\pgfsetstrokecolor{currentstroke}%
\pgfsetdash{}{0pt}%
\pgfpathmoveto{\pgfqpoint{3.474318in}{3.682415in}}%
\pgfpathcurveto{\pgfqpoint{3.480142in}{3.682415in}}{\pgfqpoint{3.485728in}{3.684729in}}{\pgfqpoint{3.489847in}{3.688847in}}%
\pgfpathcurveto{\pgfqpoint{3.493965in}{3.692966in}}{\pgfqpoint{3.496279in}{3.698552in}}{\pgfqpoint{3.496279in}{3.704376in}}%
\pgfpathcurveto{\pgfqpoint{3.496279in}{3.710200in}}{\pgfqpoint{3.493965in}{3.715786in}}{\pgfqpoint{3.489847in}{3.719904in}}%
\pgfpathcurveto{\pgfqpoint{3.485728in}{3.724022in}}{\pgfqpoint{3.480142in}{3.726336in}}{\pgfqpoint{3.474318in}{3.726336in}}%
\pgfpathcurveto{\pgfqpoint{3.468494in}{3.726336in}}{\pgfqpoint{3.462908in}{3.724022in}}{\pgfqpoint{3.458790in}{3.719904in}}%
\pgfpathcurveto{\pgfqpoint{3.454672in}{3.715786in}}{\pgfqpoint{3.452358in}{3.710200in}}{\pgfqpoint{3.452358in}{3.704376in}}%
\pgfpathcurveto{\pgfqpoint{3.452358in}{3.698552in}}{\pgfqpoint{3.454672in}{3.692966in}}{\pgfqpoint{3.458790in}{3.688847in}}%
\pgfpathcurveto{\pgfqpoint{3.462908in}{3.684729in}}{\pgfqpoint{3.468494in}{3.682415in}}{\pgfqpoint{3.474318in}{3.682415in}}%
\pgfpathlineto{\pgfqpoint{3.474318in}{3.682415in}}%
\pgfpathclose%
\pgfusepath{stroke,fill}%
\end{pgfscope}%
\begin{pgfscope}%
\pgfpathrectangle{\pgfqpoint{1.414736in}{0.548769in}}{\pgfqpoint{3.497747in}{7.253186in}}%
\pgfusepath{clip}%
\pgfsetbuttcap%
\pgfsetroundjoin%
\definecolor{currentfill}{rgb}{0.477504,0.821444,0.318195}%
\pgfsetfillcolor{currentfill}%
\pgfsetlinewidth{1.003750pt}%
\definecolor{currentstroke}{rgb}{0.477504,0.821444,0.318195}%
\pgfsetstrokecolor{currentstroke}%
\pgfsetdash{}{0pt}%
\pgfpathmoveto{\pgfqpoint{3.398763in}{4.153402in}}%
\pgfpathcurveto{\pgfqpoint{3.404587in}{4.153402in}}{\pgfqpoint{3.410174in}{4.155715in}}{\pgfqpoint{3.414292in}{4.159834in}}%
\pgfpathcurveto{\pgfqpoint{3.418410in}{4.163952in}}{\pgfqpoint{3.420724in}{4.169538in}}{\pgfqpoint{3.420724in}{4.175362in}}%
\pgfpathcurveto{\pgfqpoint{3.420724in}{4.181186in}}{\pgfqpoint{3.418410in}{4.186772in}}{\pgfqpoint{3.414292in}{4.190890in}}%
\pgfpathcurveto{\pgfqpoint{3.410174in}{4.195008in}}{\pgfqpoint{3.404587in}{4.197322in}}{\pgfqpoint{3.398763in}{4.197322in}}%
\pgfpathcurveto{\pgfqpoint{3.392939in}{4.197322in}}{\pgfqpoint{3.387353in}{4.195008in}}{\pgfqpoint{3.383235in}{4.190890in}}%
\pgfpathcurveto{\pgfqpoint{3.379117in}{4.186772in}}{\pgfqpoint{3.376803in}{4.181186in}}{\pgfqpoint{3.376803in}{4.175362in}}%
\pgfpathcurveto{\pgfqpoint{3.376803in}{4.169538in}}{\pgfqpoint{3.379117in}{4.163952in}}{\pgfqpoint{3.383235in}{4.159834in}}%
\pgfpathcurveto{\pgfqpoint{3.387353in}{4.155715in}}{\pgfqpoint{3.392939in}{4.153402in}}{\pgfqpoint{3.398763in}{4.153402in}}%
\pgfpathlineto{\pgfqpoint{3.398763in}{4.153402in}}%
\pgfpathclose%
\pgfusepath{stroke,fill}%
\end{pgfscope}%
\begin{pgfscope}%
\pgfpathrectangle{\pgfqpoint{1.414736in}{0.548769in}}{\pgfqpoint{3.497747in}{7.253186in}}%
\pgfusepath{clip}%
\pgfsetbuttcap%
\pgfsetroundjoin%
\definecolor{currentfill}{rgb}{0.804182,0.882046,0.114965}%
\pgfsetfillcolor{currentfill}%
\pgfsetlinewidth{1.003750pt}%
\definecolor{currentstroke}{rgb}{0.804182,0.882046,0.114965}%
\pgfsetstrokecolor{currentstroke}%
\pgfsetdash{}{0pt}%
\pgfpathmoveto{\pgfqpoint{3.803367in}{4.624388in}}%
\pgfpathcurveto{\pgfqpoint{3.809191in}{4.624388in}}{\pgfqpoint{3.814777in}{4.626702in}}{\pgfqpoint{3.818895in}{4.630820in}}%
\pgfpathcurveto{\pgfqpoint{3.823013in}{4.634938in}}{\pgfqpoint{3.825327in}{4.640524in}}{\pgfqpoint{3.825327in}{4.646348in}}%
\pgfpathcurveto{\pgfqpoint{3.825327in}{4.652172in}}{\pgfqpoint{3.823013in}{4.657758in}}{\pgfqpoint{3.818895in}{4.661876in}}%
\pgfpathcurveto{\pgfqpoint{3.814777in}{4.665994in}}{\pgfqpoint{3.809191in}{4.668308in}}{\pgfqpoint{3.803367in}{4.668308in}}%
\pgfpathcurveto{\pgfqpoint{3.797543in}{4.668308in}}{\pgfqpoint{3.791957in}{4.665994in}}{\pgfqpoint{3.787839in}{4.661876in}}%
\pgfpathcurveto{\pgfqpoint{3.783721in}{4.657758in}}{\pgfqpoint{3.781407in}{4.652172in}}{\pgfqpoint{3.781407in}{4.646348in}}%
\pgfpathcurveto{\pgfqpoint{3.781407in}{4.640524in}}{\pgfqpoint{3.783721in}{4.634938in}}{\pgfqpoint{3.787839in}{4.630820in}}%
\pgfpathcurveto{\pgfqpoint{3.791957in}{4.626702in}}{\pgfqpoint{3.797543in}{4.624388in}}{\pgfqpoint{3.803367in}{4.624388in}}%
\pgfpathlineto{\pgfqpoint{3.803367in}{4.624388in}}%
\pgfpathclose%
\pgfusepath{stroke,fill}%
\end{pgfscope}%
\begin{pgfscope}%
\pgfpathrectangle{\pgfqpoint{1.414736in}{0.548769in}}{\pgfqpoint{3.497747in}{7.253186in}}%
\pgfusepath{clip}%
\pgfsetbuttcap%
\pgfsetroundjoin%
\definecolor{currentfill}{rgb}{0.180653,0.701402,0.488189}%
\pgfsetfillcolor{currentfill}%
\pgfsetlinewidth{1.003750pt}%
\definecolor{currentstroke}{rgb}{0.180653,0.701402,0.488189}%
\pgfsetstrokecolor{currentstroke}%
\pgfsetdash{}{0pt}%
\pgfpathmoveto{\pgfqpoint{3.992165in}{5.095374in}}%
\pgfpathcurveto{\pgfqpoint{3.997989in}{5.095374in}}{\pgfqpoint{4.003575in}{5.097688in}}{\pgfqpoint{4.007693in}{5.101806in}}%
\pgfpathcurveto{\pgfqpoint{4.011811in}{5.105924in}}{\pgfqpoint{4.014125in}{5.111510in}}{\pgfqpoint{4.014125in}{5.117334in}}%
\pgfpathcurveto{\pgfqpoint{4.014125in}{5.123158in}}{\pgfqpoint{4.011811in}{5.128744in}}{\pgfqpoint{4.007693in}{5.132862in}}%
\pgfpathcurveto{\pgfqpoint{4.003575in}{5.136980in}}{\pgfqpoint{3.997989in}{5.139294in}}{\pgfqpoint{3.992165in}{5.139294in}}%
\pgfpathcurveto{\pgfqpoint{3.986341in}{5.139294in}}{\pgfqpoint{3.980755in}{5.136980in}}{\pgfqpoint{3.976637in}{5.132862in}}%
\pgfpathcurveto{\pgfqpoint{3.972519in}{5.128744in}}{\pgfqpoint{3.970205in}{5.123158in}}{\pgfqpoint{3.970205in}{5.117334in}}%
\pgfpathcurveto{\pgfqpoint{3.970205in}{5.111510in}}{\pgfqpoint{3.972519in}{5.105924in}}{\pgfqpoint{3.976637in}{5.101806in}}%
\pgfpathcurveto{\pgfqpoint{3.980755in}{5.097688in}}{\pgfqpoint{3.986341in}{5.095374in}}{\pgfqpoint{3.992165in}{5.095374in}}%
\pgfpathlineto{\pgfqpoint{3.992165in}{5.095374in}}%
\pgfpathclose%
\pgfusepath{stroke,fill}%
\end{pgfscope}%
\begin{pgfscope}%
\pgfpathrectangle{\pgfqpoint{1.414736in}{0.548769in}}{\pgfqpoint{3.497747in}{7.253186in}}%
\pgfusepath{clip}%
\pgfsetbuttcap%
\pgfsetroundjoin%
\definecolor{currentfill}{rgb}{0.120565,0.596422,0.543611}%
\pgfsetfillcolor{currentfill}%
\pgfsetlinewidth{1.003750pt}%
\definecolor{currentstroke}{rgb}{0.120565,0.596422,0.543611}%
\pgfsetstrokecolor{currentstroke}%
\pgfsetdash{}{0pt}%
\pgfpathmoveto{\pgfqpoint{3.175104in}{5.566360in}}%
\pgfpathcurveto{\pgfqpoint{3.180928in}{5.566360in}}{\pgfqpoint{3.186515in}{5.568674in}}{\pgfqpoint{3.190633in}{5.572792in}}%
\pgfpathcurveto{\pgfqpoint{3.194751in}{5.576910in}}{\pgfqpoint{3.197065in}{5.582496in}}{\pgfqpoint{3.197065in}{5.588320in}}%
\pgfpathcurveto{\pgfqpoint{3.197065in}{5.594144in}}{\pgfqpoint{3.194751in}{5.599730in}}{\pgfqpoint{3.190633in}{5.603848in}}%
\pgfpathcurveto{\pgfqpoint{3.186515in}{5.607967in}}{\pgfqpoint{3.180928in}{5.610280in}}{\pgfqpoint{3.175104in}{5.610280in}}%
\pgfpathcurveto{\pgfqpoint{3.169281in}{5.610280in}}{\pgfqpoint{3.163694in}{5.607967in}}{\pgfqpoint{3.159576in}{5.603848in}}%
\pgfpathcurveto{\pgfqpoint{3.155458in}{5.599730in}}{\pgfqpoint{3.153144in}{5.594144in}}{\pgfqpoint{3.153144in}{5.588320in}}%
\pgfpathcurveto{\pgfqpoint{3.153144in}{5.582496in}}{\pgfqpoint{3.155458in}{5.576910in}}{\pgfqpoint{3.159576in}{5.572792in}}%
\pgfpathcurveto{\pgfqpoint{3.163694in}{5.568674in}}{\pgfqpoint{3.169281in}{5.566360in}}{\pgfqpoint{3.175104in}{5.566360in}}%
\pgfpathlineto{\pgfqpoint{3.175104in}{5.566360in}}%
\pgfpathclose%
\pgfusepath{stroke,fill}%
\end{pgfscope}%
\begin{pgfscope}%
\pgfpathrectangle{\pgfqpoint{1.414736in}{0.548769in}}{\pgfqpoint{3.497747in}{7.253186in}}%
\pgfusepath{clip}%
\pgfsetbuttcap%
\pgfsetroundjoin%
\definecolor{currentfill}{rgb}{0.232815,0.732247,0.459277}%
\pgfsetfillcolor{currentfill}%
\pgfsetlinewidth{1.003750pt}%
\definecolor{currentstroke}{rgb}{0.232815,0.732247,0.459277}%
\pgfsetstrokecolor{currentstroke}%
\pgfsetdash{}{0pt}%
\pgfpathmoveto{\pgfqpoint{3.901767in}{6.037346in}}%
\pgfpathcurveto{\pgfqpoint{3.907591in}{6.037346in}}{\pgfqpoint{3.913177in}{6.039660in}}{\pgfqpoint{3.917295in}{6.043778in}}%
\pgfpathcurveto{\pgfqpoint{3.921413in}{6.047896in}}{\pgfqpoint{3.923727in}{6.053482in}}{\pgfqpoint{3.923727in}{6.059306in}}%
\pgfpathcurveto{\pgfqpoint{3.923727in}{6.065130in}}{\pgfqpoint{3.921413in}{6.070716in}}{\pgfqpoint{3.917295in}{6.074835in}}%
\pgfpathcurveto{\pgfqpoint{3.913177in}{6.078953in}}{\pgfqpoint{3.907591in}{6.081267in}}{\pgfqpoint{3.901767in}{6.081267in}}%
\pgfpathcurveto{\pgfqpoint{3.895943in}{6.081267in}}{\pgfqpoint{3.890357in}{6.078953in}}{\pgfqpoint{3.886238in}{6.074835in}}%
\pgfpathcurveto{\pgfqpoint{3.882120in}{6.070716in}}{\pgfqpoint{3.879806in}{6.065130in}}{\pgfqpoint{3.879806in}{6.059306in}}%
\pgfpathcurveto{\pgfqpoint{3.879806in}{6.053482in}}{\pgfqpoint{3.882120in}{6.047896in}}{\pgfqpoint{3.886238in}{6.043778in}}%
\pgfpathcurveto{\pgfqpoint{3.890357in}{6.039660in}}{\pgfqpoint{3.895943in}{6.037346in}}{\pgfqpoint{3.901767in}{6.037346in}}%
\pgfpathlineto{\pgfqpoint{3.901767in}{6.037346in}}%
\pgfpathclose%
\pgfusepath{stroke,fill}%
\end{pgfscope}%
\begin{pgfscope}%
\pgfpathrectangle{\pgfqpoint{1.414736in}{0.548769in}}{\pgfqpoint{3.497747in}{7.253186in}}%
\pgfusepath{clip}%
\pgfsetbuttcap%
\pgfsetroundjoin%
\definecolor{currentfill}{rgb}{0.545524,0.838039,0.275626}%
\pgfsetfillcolor{currentfill}%
\pgfsetlinewidth{1.003750pt}%
\definecolor{currentstroke}{rgb}{0.545524,0.838039,0.275626}%
\pgfsetstrokecolor{currentstroke}%
\pgfsetdash{}{0pt}%
\pgfpathmoveto{\pgfqpoint{3.738712in}{6.508332in}}%
\pgfpathcurveto{\pgfqpoint{3.744536in}{6.508332in}}{\pgfqpoint{3.750122in}{6.510646in}}{\pgfqpoint{3.754241in}{6.514764in}}%
\pgfpathcurveto{\pgfqpoint{3.758359in}{6.518882in}}{\pgfqpoint{3.760673in}{6.524468in}}{\pgfqpoint{3.760673in}{6.530292in}}%
\pgfpathcurveto{\pgfqpoint{3.760673in}{6.536116in}}{\pgfqpoint{3.758359in}{6.541702in}}{\pgfqpoint{3.754241in}{6.545821in}}%
\pgfpathcurveto{\pgfqpoint{3.750122in}{6.549939in}}{\pgfqpoint{3.744536in}{6.552253in}}{\pgfqpoint{3.738712in}{6.552253in}}%
\pgfpathcurveto{\pgfqpoint{3.732888in}{6.552253in}}{\pgfqpoint{3.727302in}{6.549939in}}{\pgfqpoint{3.723184in}{6.545821in}}%
\pgfpathcurveto{\pgfqpoint{3.719066in}{6.541702in}}{\pgfqpoint{3.716752in}{6.536116in}}{\pgfqpoint{3.716752in}{6.530292in}}%
\pgfpathcurveto{\pgfqpoint{3.716752in}{6.524468in}}{\pgfqpoint{3.719066in}{6.518882in}}{\pgfqpoint{3.723184in}{6.514764in}}%
\pgfpathcurveto{\pgfqpoint{3.727302in}{6.510646in}}{\pgfqpoint{3.732888in}{6.508332in}}{\pgfqpoint{3.738712in}{6.508332in}}%
\pgfpathlineto{\pgfqpoint{3.738712in}{6.508332in}}%
\pgfpathclose%
\pgfusepath{stroke,fill}%
\end{pgfscope}%
\begin{pgfscope}%
\pgfpathrectangle{\pgfqpoint{1.414736in}{0.548769in}}{\pgfqpoint{3.497747in}{7.253186in}}%
\pgfusepath{clip}%
\pgfsetbuttcap%
\pgfsetroundjoin%
\definecolor{currentfill}{rgb}{0.585678,0.846661,0.249897}%
\pgfsetfillcolor{currentfill}%
\pgfsetlinewidth{1.003750pt}%
\definecolor{currentstroke}{rgb}{0.585678,0.846661,0.249897}%
\pgfsetstrokecolor{currentstroke}%
\pgfsetdash{}{0pt}%
\pgfpathmoveto{\pgfqpoint{3.859282in}{6.979318in}}%
\pgfpathcurveto{\pgfqpoint{3.865106in}{6.979318in}}{\pgfqpoint{3.870692in}{6.981632in}}{\pgfqpoint{3.874810in}{6.985750in}}%
\pgfpathcurveto{\pgfqpoint{3.878928in}{6.989868in}}{\pgfqpoint{3.881242in}{6.995455in}}{\pgfqpoint{3.881242in}{7.001278in}}%
\pgfpathcurveto{\pgfqpoint{3.881242in}{7.007102in}}{\pgfqpoint{3.878928in}{7.012689in}}{\pgfqpoint{3.874810in}{7.016807in}}%
\pgfpathcurveto{\pgfqpoint{3.870692in}{7.020925in}}{\pgfqpoint{3.865106in}{7.023239in}}{\pgfqpoint{3.859282in}{7.023239in}}%
\pgfpathcurveto{\pgfqpoint{3.853458in}{7.023239in}}{\pgfqpoint{3.847872in}{7.020925in}}{\pgfqpoint{3.843754in}{7.016807in}}%
\pgfpathcurveto{\pgfqpoint{3.839636in}{7.012689in}}{\pgfqpoint{3.837322in}{7.007102in}}{\pgfqpoint{3.837322in}{7.001278in}}%
\pgfpathcurveto{\pgfqpoint{3.837322in}{6.995455in}}{\pgfqpoint{3.839636in}{6.989868in}}{\pgfqpoint{3.843754in}{6.985750in}}%
\pgfpathcurveto{\pgfqpoint{3.847872in}{6.981632in}}{\pgfqpoint{3.853458in}{6.979318in}}{\pgfqpoint{3.859282in}{6.979318in}}%
\pgfpathlineto{\pgfqpoint{3.859282in}{6.979318in}}%
\pgfpathclose%
\pgfusepath{stroke,fill}%
\end{pgfscope}%
\begin{pgfscope}%
\pgfpathrectangle{\pgfqpoint{1.414736in}{0.548769in}}{\pgfqpoint{3.497747in}{7.253186in}}%
\pgfusepath{clip}%
\pgfsetbuttcap%
\pgfsetroundjoin%
\definecolor{currentfill}{rgb}{0.412913,0.803041,0.357269}%
\pgfsetfillcolor{currentfill}%
\pgfsetlinewidth{1.003750pt}%
\definecolor{currentstroke}{rgb}{0.412913,0.803041,0.357269}%
\pgfsetstrokecolor{currentstroke}%
\pgfsetdash{}{0pt}%
\pgfpathmoveto{\pgfqpoint{3.121564in}{7.450304in}}%
\pgfpathcurveto{\pgfqpoint{3.127388in}{7.450304in}}{\pgfqpoint{3.132974in}{7.452618in}}{\pgfqpoint{3.137093in}{7.456736in}}%
\pgfpathcurveto{\pgfqpoint{3.141211in}{7.460854in}}{\pgfqpoint{3.143525in}{7.466441in}}{\pgfqpoint{3.143525in}{7.472265in}}%
\pgfpathcurveto{\pgfqpoint{3.143525in}{7.478089in}}{\pgfqpoint{3.141211in}{7.483675in}}{\pgfqpoint{3.137093in}{7.487793in}}%
\pgfpathcurveto{\pgfqpoint{3.132974in}{7.491911in}}{\pgfqpoint{3.127388in}{7.494225in}}{\pgfqpoint{3.121564in}{7.494225in}}%
\pgfpathcurveto{\pgfqpoint{3.115740in}{7.494225in}}{\pgfqpoint{3.110154in}{7.491911in}}{\pgfqpoint{3.106036in}{7.487793in}}%
\pgfpathcurveto{\pgfqpoint{3.101918in}{7.483675in}}{\pgfqpoint{3.099604in}{7.478089in}}{\pgfqpoint{3.099604in}{7.472265in}}%
\pgfpathcurveto{\pgfqpoint{3.099604in}{7.466441in}}{\pgfqpoint{3.101918in}{7.460854in}}{\pgfqpoint{3.106036in}{7.456736in}}%
\pgfpathcurveto{\pgfqpoint{3.110154in}{7.452618in}}{\pgfqpoint{3.115740in}{7.450304in}}{\pgfqpoint{3.121564in}{7.450304in}}%
\pgfpathlineto{\pgfqpoint{3.121564in}{7.450304in}}%
\pgfpathclose%
\pgfusepath{stroke,fill}%
\end{pgfscope}%
\begin{pgfscope}%
\pgfpathrectangle{\pgfqpoint{1.414736in}{0.548769in}}{\pgfqpoint{3.497747in}{7.253186in}}%
\pgfusepath{clip}%
\pgfsetbuttcap%
\pgfsetroundjoin%
\definecolor{currentfill}{rgb}{0.993248,0.906157,0.143936}%
\pgfsetfillcolor{currentfill}%
\pgfsetlinewidth{1.003750pt}%
\definecolor{currentstroke}{rgb}{0.993248,0.906157,0.143936}%
\pgfsetstrokecolor{currentstroke}%
\pgfsetdash{}{0pt}%
\pgfpathmoveto{\pgfqpoint{3.455720in}{0.856499in}}%
\pgfpathcurveto{\pgfqpoint{3.461544in}{0.856499in}}{\pgfqpoint{3.467130in}{0.858813in}}{\pgfqpoint{3.471248in}{0.862931in}}%
\pgfpathcurveto{\pgfqpoint{3.475366in}{0.867049in}}{\pgfqpoint{3.477680in}{0.872635in}}{\pgfqpoint{3.477680in}{0.878459in}}%
\pgfpathcurveto{\pgfqpoint{3.477680in}{0.884283in}}{\pgfqpoint{3.475366in}{0.889869in}}{\pgfqpoint{3.471248in}{0.893987in}}%
\pgfpathcurveto{\pgfqpoint{3.467130in}{0.898105in}}{\pgfqpoint{3.461544in}{0.900419in}}{\pgfqpoint{3.455720in}{0.900419in}}%
\pgfpathcurveto{\pgfqpoint{3.449896in}{0.900419in}}{\pgfqpoint{3.444310in}{0.898105in}}{\pgfqpoint{3.440191in}{0.893987in}}%
\pgfpathcurveto{\pgfqpoint{3.436073in}{0.889869in}}{\pgfqpoint{3.433759in}{0.884283in}}{\pgfqpoint{3.433759in}{0.878459in}}%
\pgfpathcurveto{\pgfqpoint{3.433759in}{0.872635in}}{\pgfqpoint{3.436073in}{0.867049in}}{\pgfqpoint{3.440191in}{0.862931in}}%
\pgfpathcurveto{\pgfqpoint{3.444310in}{0.858813in}}{\pgfqpoint{3.449896in}{0.856499in}}{\pgfqpoint{3.455720in}{0.856499in}}%
\pgfpathlineto{\pgfqpoint{3.455720in}{0.856499in}}%
\pgfpathclose%
\pgfusepath{stroke,fill}%
\end{pgfscope}%
\begin{pgfscope}%
\pgfpathrectangle{\pgfqpoint{1.414736in}{0.548769in}}{\pgfqpoint{3.497747in}{7.253186in}}%
\pgfusepath{clip}%
\pgfsetbuttcap%
\pgfsetroundjoin%
\definecolor{currentfill}{rgb}{0.487026,0.823929,0.312321}%
\pgfsetfillcolor{currentfill}%
\pgfsetlinewidth{1.003750pt}%
\definecolor{currentstroke}{rgb}{0.487026,0.823929,0.312321}%
\pgfsetstrokecolor{currentstroke}%
\pgfsetdash{}{0pt}%
\pgfpathmoveto{\pgfqpoint{3.645668in}{1.327485in}}%
\pgfpathcurveto{\pgfqpoint{3.651492in}{1.327485in}}{\pgfqpoint{3.657078in}{1.329799in}}{\pgfqpoint{3.661196in}{1.333917in}}%
\pgfpathcurveto{\pgfqpoint{3.665314in}{1.338035in}}{\pgfqpoint{3.667628in}{1.343621in}}{\pgfqpoint{3.667628in}{1.349445in}}%
\pgfpathcurveto{\pgfqpoint{3.667628in}{1.355269in}}{\pgfqpoint{3.665314in}{1.360855in}}{\pgfqpoint{3.661196in}{1.364973in}}%
\pgfpathcurveto{\pgfqpoint{3.657078in}{1.369092in}}{\pgfqpoint{3.651492in}{1.371405in}}{\pgfqpoint{3.645668in}{1.371405in}}%
\pgfpathcurveto{\pgfqpoint{3.639844in}{1.371405in}}{\pgfqpoint{3.634258in}{1.369092in}}{\pgfqpoint{3.630140in}{1.364973in}}%
\pgfpathcurveto{\pgfqpoint{3.626022in}{1.360855in}}{\pgfqpoint{3.623708in}{1.355269in}}{\pgfqpoint{3.623708in}{1.349445in}}%
\pgfpathcurveto{\pgfqpoint{3.623708in}{1.343621in}}{\pgfqpoint{3.626022in}{1.338035in}}{\pgfqpoint{3.630140in}{1.333917in}}%
\pgfpathcurveto{\pgfqpoint{3.634258in}{1.329799in}}{\pgfqpoint{3.639844in}{1.327485in}}{\pgfqpoint{3.645668in}{1.327485in}}%
\pgfpathlineto{\pgfqpoint{3.645668in}{1.327485in}}%
\pgfpathclose%
\pgfusepath{stroke,fill}%
\end{pgfscope}%
\begin{pgfscope}%
\pgfpathrectangle{\pgfqpoint{1.414736in}{0.548769in}}{\pgfqpoint{3.497747in}{7.253186in}}%
\pgfusepath{clip}%
\pgfsetbuttcap%
\pgfsetroundjoin%
\definecolor{currentfill}{rgb}{0.430983,0.808473,0.346476}%
\pgfsetfillcolor{currentfill}%
\pgfsetlinewidth{1.003750pt}%
\definecolor{currentstroke}{rgb}{0.430983,0.808473,0.346476}%
\pgfsetstrokecolor{currentstroke}%
\pgfsetdash{}{0pt}%
\pgfpathmoveto{\pgfqpoint{3.410550in}{1.798471in}}%
\pgfpathcurveto{\pgfqpoint{3.416374in}{1.798471in}}{\pgfqpoint{3.421960in}{1.800785in}}{\pgfqpoint{3.426079in}{1.804903in}}%
\pgfpathcurveto{\pgfqpoint{3.430197in}{1.809021in}}{\pgfqpoint{3.432511in}{1.814607in}}{\pgfqpoint{3.432511in}{1.820431in}}%
\pgfpathcurveto{\pgfqpoint{3.432511in}{1.826255in}}{\pgfqpoint{3.430197in}{1.831841in}}{\pgfqpoint{3.426079in}{1.835960in}}%
\pgfpathcurveto{\pgfqpoint{3.421960in}{1.840078in}}{\pgfqpoint{3.416374in}{1.842392in}}{\pgfqpoint{3.410550in}{1.842392in}}%
\pgfpathcurveto{\pgfqpoint{3.404726in}{1.842392in}}{\pgfqpoint{3.399140in}{1.840078in}}{\pgfqpoint{3.395022in}{1.835960in}}%
\pgfpathcurveto{\pgfqpoint{3.390904in}{1.831841in}}{\pgfqpoint{3.388590in}{1.826255in}}{\pgfqpoint{3.388590in}{1.820431in}}%
\pgfpathcurveto{\pgfqpoint{3.388590in}{1.814607in}}{\pgfqpoint{3.390904in}{1.809021in}}{\pgfqpoint{3.395022in}{1.804903in}}%
\pgfpathcurveto{\pgfqpoint{3.399140in}{1.800785in}}{\pgfqpoint{3.404726in}{1.798471in}}{\pgfqpoint{3.410550in}{1.798471in}}%
\pgfpathlineto{\pgfqpoint{3.410550in}{1.798471in}}%
\pgfpathclose%
\pgfusepath{stroke,fill}%
\end{pgfscope}%
\begin{pgfscope}%
\pgfpathrectangle{\pgfqpoint{1.414736in}{0.548769in}}{\pgfqpoint{3.497747in}{7.253186in}}%
\pgfusepath{clip}%
\pgfsetbuttcap%
\pgfsetroundjoin%
\definecolor{currentfill}{rgb}{0.974417,0.903590,0.130215}%
\pgfsetfillcolor{currentfill}%
\pgfsetlinewidth{1.003750pt}%
\definecolor{currentstroke}{rgb}{0.974417,0.903590,0.130215}%
\pgfsetstrokecolor{currentstroke}%
\pgfsetdash{}{0pt}%
\pgfpathmoveto{\pgfqpoint{3.577125in}{2.269457in}}%
\pgfpathcurveto{\pgfqpoint{3.582949in}{2.269457in}}{\pgfqpoint{3.588535in}{2.271771in}}{\pgfqpoint{3.592653in}{2.275889in}}%
\pgfpathcurveto{\pgfqpoint{3.596771in}{2.280007in}}{\pgfqpoint{3.599085in}{2.285593in}}{\pgfqpoint{3.599085in}{2.291417in}}%
\pgfpathcurveto{\pgfqpoint{3.599085in}{2.297241in}}{\pgfqpoint{3.596771in}{2.302828in}}{\pgfqpoint{3.592653in}{2.306946in}}%
\pgfpathcurveto{\pgfqpoint{3.588535in}{2.311064in}}{\pgfqpoint{3.582949in}{2.313378in}}{\pgfqpoint{3.577125in}{2.313378in}}%
\pgfpathcurveto{\pgfqpoint{3.571301in}{2.313378in}}{\pgfqpoint{3.565715in}{2.311064in}}{\pgfqpoint{3.561596in}{2.306946in}}%
\pgfpathcurveto{\pgfqpoint{3.557478in}{2.302828in}}{\pgfqpoint{3.555164in}{2.297241in}}{\pgfqpoint{3.555164in}{2.291417in}}%
\pgfpathcurveto{\pgfqpoint{3.555164in}{2.285593in}}{\pgfqpoint{3.557478in}{2.280007in}}{\pgfqpoint{3.561596in}{2.275889in}}%
\pgfpathcurveto{\pgfqpoint{3.565715in}{2.271771in}}{\pgfqpoint{3.571301in}{2.269457in}}{\pgfqpoint{3.577125in}{2.269457in}}%
\pgfpathlineto{\pgfqpoint{3.577125in}{2.269457in}}%
\pgfpathclose%
\pgfusepath{stroke,fill}%
\end{pgfscope}%
\begin{pgfscope}%
\pgfpathrectangle{\pgfqpoint{1.414736in}{0.548769in}}{\pgfqpoint{3.497747in}{7.253186in}}%
\pgfusepath{clip}%
\pgfsetbuttcap%
\pgfsetroundjoin%
\definecolor{currentfill}{rgb}{0.545524,0.838039,0.275626}%
\pgfsetfillcolor{currentfill}%
\pgfsetlinewidth{1.003750pt}%
\definecolor{currentstroke}{rgb}{0.545524,0.838039,0.275626}%
\pgfsetstrokecolor{currentstroke}%
\pgfsetdash{}{0pt}%
\pgfpathmoveto{\pgfqpoint{3.563173in}{2.740443in}}%
\pgfpathcurveto{\pgfqpoint{3.568997in}{2.740443in}}{\pgfqpoint{3.574583in}{2.742757in}}{\pgfqpoint{3.578701in}{2.746875in}}%
\pgfpathcurveto{\pgfqpoint{3.582819in}{2.750993in}}{\pgfqpoint{3.585133in}{2.756580in}}{\pgfqpoint{3.585133in}{2.762404in}}%
\pgfpathcurveto{\pgfqpoint{3.585133in}{2.768227in}}{\pgfqpoint{3.582819in}{2.773814in}}{\pgfqpoint{3.578701in}{2.777932in}}%
\pgfpathcurveto{\pgfqpoint{3.574583in}{2.782050in}}{\pgfqpoint{3.568997in}{2.784364in}}{\pgfqpoint{3.563173in}{2.784364in}}%
\pgfpathcurveto{\pgfqpoint{3.557349in}{2.784364in}}{\pgfqpoint{3.551763in}{2.782050in}}{\pgfqpoint{3.547645in}{2.777932in}}%
\pgfpathcurveto{\pgfqpoint{3.543527in}{2.773814in}}{\pgfqpoint{3.541213in}{2.768227in}}{\pgfqpoint{3.541213in}{2.762404in}}%
\pgfpathcurveto{\pgfqpoint{3.541213in}{2.756580in}}{\pgfqpoint{3.543527in}{2.750993in}}{\pgfqpoint{3.547645in}{2.746875in}}%
\pgfpathcurveto{\pgfqpoint{3.551763in}{2.742757in}}{\pgfqpoint{3.557349in}{2.740443in}}{\pgfqpoint{3.563173in}{2.740443in}}%
\pgfpathlineto{\pgfqpoint{3.563173in}{2.740443in}}%
\pgfpathclose%
\pgfusepath{stroke,fill}%
\end{pgfscope}%
\begin{pgfscope}%
\pgfpathrectangle{\pgfqpoint{1.414736in}{0.548769in}}{\pgfqpoint{3.497747in}{7.253186in}}%
\pgfusepath{clip}%
\pgfsetbuttcap%
\pgfsetroundjoin%
\definecolor{currentfill}{rgb}{0.730889,0.871916,0.156029}%
\pgfsetfillcolor{currentfill}%
\pgfsetlinewidth{1.003750pt}%
\definecolor{currentstroke}{rgb}{0.730889,0.871916,0.156029}%
\pgfsetstrokecolor{currentstroke}%
\pgfsetdash{}{0pt}%
\pgfpathmoveto{\pgfqpoint{3.395241in}{3.211429in}}%
\pgfpathcurveto{\pgfqpoint{3.401065in}{3.211429in}}{\pgfqpoint{3.406651in}{3.213743in}}{\pgfqpoint{3.410769in}{3.217861in}}%
\pgfpathcurveto{\pgfqpoint{3.414887in}{3.221980in}}{\pgfqpoint{3.417201in}{3.227566in}}{\pgfqpoint{3.417201in}{3.233390in}}%
\pgfpathcurveto{\pgfqpoint{3.417201in}{3.239214in}}{\pgfqpoint{3.414887in}{3.244800in}}{\pgfqpoint{3.410769in}{3.248918in}}%
\pgfpathcurveto{\pgfqpoint{3.406651in}{3.253036in}}{\pgfqpoint{3.401065in}{3.255350in}}{\pgfqpoint{3.395241in}{3.255350in}}%
\pgfpathcurveto{\pgfqpoint{3.389417in}{3.255350in}}{\pgfqpoint{3.383831in}{3.253036in}}{\pgfqpoint{3.379713in}{3.248918in}}%
\pgfpathcurveto{\pgfqpoint{3.375594in}{3.244800in}}{\pgfqpoint{3.373281in}{3.239214in}}{\pgfqpoint{3.373281in}{3.233390in}}%
\pgfpathcurveto{\pgfqpoint{3.373281in}{3.227566in}}{\pgfqpoint{3.375594in}{3.221980in}}{\pgfqpoint{3.379713in}{3.217861in}}%
\pgfpathcurveto{\pgfqpoint{3.383831in}{3.213743in}}{\pgfqpoint{3.389417in}{3.211429in}}{\pgfqpoint{3.395241in}{3.211429in}}%
\pgfpathlineto{\pgfqpoint{3.395241in}{3.211429in}}%
\pgfpathclose%
\pgfusepath{stroke,fill}%
\end{pgfscope}%
\begin{pgfscope}%
\pgfpathrectangle{\pgfqpoint{1.414736in}{0.548769in}}{\pgfqpoint{3.497747in}{7.253186in}}%
\pgfusepath{clip}%
\pgfsetbuttcap%
\pgfsetroundjoin%
\definecolor{currentfill}{rgb}{0.886271,0.892374,0.095374}%
\pgfsetfillcolor{currentfill}%
\pgfsetlinewidth{1.003750pt}%
\definecolor{currentstroke}{rgb}{0.886271,0.892374,0.095374}%
\pgfsetstrokecolor{currentstroke}%
\pgfsetdash{}{0pt}%
\pgfpathmoveto{\pgfqpoint{3.555830in}{3.682415in}}%
\pgfpathcurveto{\pgfqpoint{3.561654in}{3.682415in}}{\pgfqpoint{3.567240in}{3.684729in}}{\pgfqpoint{3.571358in}{3.688847in}}%
\pgfpathcurveto{\pgfqpoint{3.575476in}{3.692966in}}{\pgfqpoint{3.577790in}{3.698552in}}{\pgfqpoint{3.577790in}{3.704376in}}%
\pgfpathcurveto{\pgfqpoint{3.577790in}{3.710200in}}{\pgfqpoint{3.575476in}{3.715786in}}{\pgfqpoint{3.571358in}{3.719904in}}%
\pgfpathcurveto{\pgfqpoint{3.567240in}{3.724022in}}{\pgfqpoint{3.561654in}{3.726336in}}{\pgfqpoint{3.555830in}{3.726336in}}%
\pgfpathcurveto{\pgfqpoint{3.550006in}{3.726336in}}{\pgfqpoint{3.544420in}{3.724022in}}{\pgfqpoint{3.540302in}{3.719904in}}%
\pgfpathcurveto{\pgfqpoint{3.536184in}{3.715786in}}{\pgfqpoint{3.533870in}{3.710200in}}{\pgfqpoint{3.533870in}{3.704376in}}%
\pgfpathcurveto{\pgfqpoint{3.533870in}{3.698552in}}{\pgfqpoint{3.536184in}{3.692966in}}{\pgfqpoint{3.540302in}{3.688847in}}%
\pgfpathcurveto{\pgfqpoint{3.544420in}{3.684729in}}{\pgfqpoint{3.550006in}{3.682415in}}{\pgfqpoint{3.555830in}{3.682415in}}%
\pgfpathlineto{\pgfqpoint{3.555830in}{3.682415in}}%
\pgfpathclose%
\pgfusepath{stroke,fill}%
\end{pgfscope}%
\begin{pgfscope}%
\pgfpathrectangle{\pgfqpoint{1.414736in}{0.548769in}}{\pgfqpoint{3.497747in}{7.253186in}}%
\pgfusepath{clip}%
\pgfsetbuttcap%
\pgfsetroundjoin%
\definecolor{currentfill}{rgb}{0.730889,0.871916,0.156029}%
\pgfsetfillcolor{currentfill}%
\pgfsetlinewidth{1.003750pt}%
\definecolor{currentstroke}{rgb}{0.730889,0.871916,0.156029}%
\pgfsetstrokecolor{currentstroke}%
\pgfsetdash{}{0pt}%
\pgfpathmoveto{\pgfqpoint{3.456353in}{4.153402in}}%
\pgfpathcurveto{\pgfqpoint{3.462177in}{4.153402in}}{\pgfqpoint{3.467763in}{4.155715in}}{\pgfqpoint{3.471881in}{4.159834in}}%
\pgfpathcurveto{\pgfqpoint{3.476000in}{4.163952in}}{\pgfqpoint{3.478313in}{4.169538in}}{\pgfqpoint{3.478313in}{4.175362in}}%
\pgfpathcurveto{\pgfqpoint{3.478313in}{4.181186in}}{\pgfqpoint{3.476000in}{4.186772in}}{\pgfqpoint{3.471881in}{4.190890in}}%
\pgfpathcurveto{\pgfqpoint{3.467763in}{4.195008in}}{\pgfqpoint{3.462177in}{4.197322in}}{\pgfqpoint{3.456353in}{4.197322in}}%
\pgfpathcurveto{\pgfqpoint{3.450529in}{4.197322in}}{\pgfqpoint{3.444943in}{4.195008in}}{\pgfqpoint{3.440825in}{4.190890in}}%
\pgfpathcurveto{\pgfqpoint{3.436707in}{4.186772in}}{\pgfqpoint{3.434393in}{4.181186in}}{\pgfqpoint{3.434393in}{4.175362in}}%
\pgfpathcurveto{\pgfqpoint{3.434393in}{4.169538in}}{\pgfqpoint{3.436707in}{4.163952in}}{\pgfqpoint{3.440825in}{4.159834in}}%
\pgfpathcurveto{\pgfqpoint{3.444943in}{4.155715in}}{\pgfqpoint{3.450529in}{4.153402in}}{\pgfqpoint{3.456353in}{4.153402in}}%
\pgfpathlineto{\pgfqpoint{3.456353in}{4.153402in}}%
\pgfpathclose%
\pgfusepath{stroke,fill}%
\end{pgfscope}%
\begin{pgfscope}%
\pgfpathrectangle{\pgfqpoint{1.414736in}{0.548769in}}{\pgfqpoint{3.497747in}{7.253186in}}%
\pgfusepath{clip}%
\pgfsetbuttcap%
\pgfsetroundjoin%
\definecolor{currentfill}{rgb}{0.699415,0.867117,0.175971}%
\pgfsetfillcolor{currentfill}%
\pgfsetlinewidth{1.003750pt}%
\definecolor{currentstroke}{rgb}{0.699415,0.867117,0.175971}%
\pgfsetstrokecolor{currentstroke}%
\pgfsetdash{}{0pt}%
\pgfpathmoveto{\pgfqpoint{3.424475in}{4.624388in}}%
\pgfpathcurveto{\pgfqpoint{3.430299in}{4.624388in}}{\pgfqpoint{3.435886in}{4.626702in}}{\pgfqpoint{3.440004in}{4.630820in}}%
\pgfpathcurveto{\pgfqpoint{3.444122in}{4.634938in}}{\pgfqpoint{3.446436in}{4.640524in}}{\pgfqpoint{3.446436in}{4.646348in}}%
\pgfpathcurveto{\pgfqpoint{3.446436in}{4.652172in}}{\pgfqpoint{3.444122in}{4.657758in}}{\pgfqpoint{3.440004in}{4.661876in}}%
\pgfpathcurveto{\pgfqpoint{3.435886in}{4.665994in}}{\pgfqpoint{3.430299in}{4.668308in}}{\pgfqpoint{3.424475in}{4.668308in}}%
\pgfpathcurveto{\pgfqpoint{3.418652in}{4.668308in}}{\pgfqpoint{3.413065in}{4.665994in}}{\pgfqpoint{3.408947in}{4.661876in}}%
\pgfpathcurveto{\pgfqpoint{3.404829in}{4.657758in}}{\pgfqpoint{3.402515in}{4.652172in}}{\pgfqpoint{3.402515in}{4.646348in}}%
\pgfpathcurveto{\pgfqpoint{3.402515in}{4.640524in}}{\pgfqpoint{3.404829in}{4.634938in}}{\pgfqpoint{3.408947in}{4.630820in}}%
\pgfpathcurveto{\pgfqpoint{3.413065in}{4.626702in}}{\pgfqpoint{3.418652in}{4.624388in}}{\pgfqpoint{3.424475in}{4.624388in}}%
\pgfpathlineto{\pgfqpoint{3.424475in}{4.624388in}}%
\pgfpathclose%
\pgfusepath{stroke,fill}%
\end{pgfscope}%
\begin{pgfscope}%
\pgfpathrectangle{\pgfqpoint{1.414736in}{0.548769in}}{\pgfqpoint{3.497747in}{7.253186in}}%
\pgfusepath{clip}%
\pgfsetbuttcap%
\pgfsetroundjoin%
\definecolor{currentfill}{rgb}{0.258965,0.251537,0.524736}%
\pgfsetfillcolor{currentfill}%
\pgfsetlinewidth{1.003750pt}%
\definecolor{currentstroke}{rgb}{0.258965,0.251537,0.524736}%
\pgfsetstrokecolor{currentstroke}%
\pgfsetdash{}{0pt}%
\pgfpathmoveto{\pgfqpoint{3.502851in}{5.095374in}}%
\pgfpathcurveto{\pgfqpoint{3.508675in}{5.095374in}}{\pgfqpoint{3.514261in}{5.097688in}}{\pgfqpoint{3.518379in}{5.101806in}}%
\pgfpathcurveto{\pgfqpoint{3.522497in}{5.105924in}}{\pgfqpoint{3.524811in}{5.111510in}}{\pgfqpoint{3.524811in}{5.117334in}}%
\pgfpathcurveto{\pgfqpoint{3.524811in}{5.123158in}}{\pgfqpoint{3.522497in}{5.128744in}}{\pgfqpoint{3.518379in}{5.132862in}}%
\pgfpathcurveto{\pgfqpoint{3.514261in}{5.136980in}}{\pgfqpoint{3.508675in}{5.139294in}}{\pgfqpoint{3.502851in}{5.139294in}}%
\pgfpathcurveto{\pgfqpoint{3.497027in}{5.139294in}}{\pgfqpoint{3.491440in}{5.136980in}}{\pgfqpoint{3.487322in}{5.132862in}}%
\pgfpathcurveto{\pgfqpoint{3.483204in}{5.128744in}}{\pgfqpoint{3.480890in}{5.123158in}}{\pgfqpoint{3.480890in}{5.117334in}}%
\pgfpathcurveto{\pgfqpoint{3.480890in}{5.111510in}}{\pgfqpoint{3.483204in}{5.105924in}}{\pgfqpoint{3.487322in}{5.101806in}}%
\pgfpathcurveto{\pgfqpoint{3.491440in}{5.097688in}}{\pgfqpoint{3.497027in}{5.095374in}}{\pgfqpoint{3.502851in}{5.095374in}}%
\pgfpathlineto{\pgfqpoint{3.502851in}{5.095374in}}%
\pgfpathclose%
\pgfusepath{stroke,fill}%
\end{pgfscope}%
\begin{pgfscope}%
\pgfpathrectangle{\pgfqpoint{1.414736in}{0.548769in}}{\pgfqpoint{3.497747in}{7.253186in}}%
\pgfusepath{clip}%
\pgfsetbuttcap%
\pgfsetroundjoin%
\definecolor{currentfill}{rgb}{0.204903,0.375746,0.553533}%
\pgfsetfillcolor{currentfill}%
\pgfsetlinewidth{1.003750pt}%
\definecolor{currentstroke}{rgb}{0.204903,0.375746,0.553533}%
\pgfsetstrokecolor{currentstroke}%
\pgfsetdash{}{0pt}%
\pgfpathmoveto{\pgfqpoint{3.595655in}{5.566360in}}%
\pgfpathcurveto{\pgfqpoint{3.601479in}{5.566360in}}{\pgfqpoint{3.607065in}{5.568674in}}{\pgfqpoint{3.611183in}{5.572792in}}%
\pgfpathcurveto{\pgfqpoint{3.615301in}{5.576910in}}{\pgfqpoint{3.617615in}{5.582496in}}{\pgfqpoint{3.617615in}{5.588320in}}%
\pgfpathcurveto{\pgfqpoint{3.617615in}{5.594144in}}{\pgfqpoint{3.615301in}{5.599730in}}{\pgfqpoint{3.611183in}{5.603848in}}%
\pgfpathcurveto{\pgfqpoint{3.607065in}{5.607967in}}{\pgfqpoint{3.601479in}{5.610280in}}{\pgfqpoint{3.595655in}{5.610280in}}%
\pgfpathcurveto{\pgfqpoint{3.589831in}{5.610280in}}{\pgfqpoint{3.584245in}{5.607967in}}{\pgfqpoint{3.580127in}{5.603848in}}%
\pgfpathcurveto{\pgfqpoint{3.576009in}{5.599730in}}{\pgfqpoint{3.573695in}{5.594144in}}{\pgfqpoint{3.573695in}{5.588320in}}%
\pgfpathcurveto{\pgfqpoint{3.573695in}{5.582496in}}{\pgfqpoint{3.576009in}{5.576910in}}{\pgfqpoint{3.580127in}{5.572792in}}%
\pgfpathcurveto{\pgfqpoint{3.584245in}{5.568674in}}{\pgfqpoint{3.589831in}{5.566360in}}{\pgfqpoint{3.595655in}{5.566360in}}%
\pgfpathlineto{\pgfqpoint{3.595655in}{5.566360in}}%
\pgfpathclose%
\pgfusepath{stroke,fill}%
\end{pgfscope}%
\begin{pgfscope}%
\pgfpathrectangle{\pgfqpoint{1.414736in}{0.548769in}}{\pgfqpoint{3.497747in}{7.253186in}}%
\pgfusepath{clip}%
\pgfsetbuttcap%
\pgfsetroundjoin%
\definecolor{currentfill}{rgb}{0.259857,0.745492,0.444467}%
\pgfsetfillcolor{currentfill}%
\pgfsetlinewidth{1.003750pt}%
\definecolor{currentstroke}{rgb}{0.259857,0.745492,0.444467}%
\pgfsetstrokecolor{currentstroke}%
\pgfsetdash{}{0pt}%
\pgfpathmoveto{\pgfqpoint{3.383642in}{6.037346in}}%
\pgfpathcurveto{\pgfqpoint{3.389466in}{6.037346in}}{\pgfqpoint{3.395052in}{6.039660in}}{\pgfqpoint{3.399170in}{6.043778in}}%
\pgfpathcurveto{\pgfqpoint{3.403289in}{6.047896in}}{\pgfqpoint{3.405602in}{6.053482in}}{\pgfqpoint{3.405602in}{6.059306in}}%
\pgfpathcurveto{\pgfqpoint{3.405602in}{6.065130in}}{\pgfqpoint{3.403289in}{6.070716in}}{\pgfqpoint{3.399170in}{6.074835in}}%
\pgfpathcurveto{\pgfqpoint{3.395052in}{6.078953in}}{\pgfqpoint{3.389466in}{6.081267in}}{\pgfqpoint{3.383642in}{6.081267in}}%
\pgfpathcurveto{\pgfqpoint{3.377818in}{6.081267in}}{\pgfqpoint{3.372232in}{6.078953in}}{\pgfqpoint{3.368114in}{6.074835in}}%
\pgfpathcurveto{\pgfqpoint{3.363996in}{6.070716in}}{\pgfqpoint{3.361682in}{6.065130in}}{\pgfqpoint{3.361682in}{6.059306in}}%
\pgfpathcurveto{\pgfqpoint{3.361682in}{6.053482in}}{\pgfqpoint{3.363996in}{6.047896in}}{\pgfqpoint{3.368114in}{6.043778in}}%
\pgfpathcurveto{\pgfqpoint{3.372232in}{6.039660in}}{\pgfqpoint{3.377818in}{6.037346in}}{\pgfqpoint{3.383642in}{6.037346in}}%
\pgfpathlineto{\pgfqpoint{3.383642in}{6.037346in}}%
\pgfpathclose%
\pgfusepath{stroke,fill}%
\end{pgfscope}%
\begin{pgfscope}%
\pgfpathrectangle{\pgfqpoint{1.414736in}{0.548769in}}{\pgfqpoint{3.497747in}{7.253186in}}%
\pgfusepath{clip}%
\pgfsetbuttcap%
\pgfsetroundjoin%
\definecolor{currentfill}{rgb}{0.657642,0.860219,0.203082}%
\pgfsetfillcolor{currentfill}%
\pgfsetlinewidth{1.003750pt}%
\definecolor{currentstroke}{rgb}{0.657642,0.860219,0.203082}%
\pgfsetstrokecolor{currentstroke}%
\pgfsetdash{}{0pt}%
\pgfpathmoveto{\pgfqpoint{3.971467in}{6.508332in}}%
\pgfpathcurveto{\pgfqpoint{3.977291in}{6.508332in}}{\pgfqpoint{3.982877in}{6.510646in}}{\pgfqpoint{3.986995in}{6.514764in}}%
\pgfpathcurveto{\pgfqpoint{3.991113in}{6.518882in}}{\pgfqpoint{3.993427in}{6.524468in}}{\pgfqpoint{3.993427in}{6.530292in}}%
\pgfpathcurveto{\pgfqpoint{3.993427in}{6.536116in}}{\pgfqpoint{3.991113in}{6.541702in}}{\pgfqpoint{3.986995in}{6.545821in}}%
\pgfpathcurveto{\pgfqpoint{3.982877in}{6.549939in}}{\pgfqpoint{3.977291in}{6.552253in}}{\pgfqpoint{3.971467in}{6.552253in}}%
\pgfpathcurveto{\pgfqpoint{3.965643in}{6.552253in}}{\pgfqpoint{3.960057in}{6.549939in}}{\pgfqpoint{3.955938in}{6.545821in}}%
\pgfpathcurveto{\pgfqpoint{3.951820in}{6.541702in}}{\pgfqpoint{3.949506in}{6.536116in}}{\pgfqpoint{3.949506in}{6.530292in}}%
\pgfpathcurveto{\pgfqpoint{3.949506in}{6.524468in}}{\pgfqpoint{3.951820in}{6.518882in}}{\pgfqpoint{3.955938in}{6.514764in}}%
\pgfpathcurveto{\pgfqpoint{3.960057in}{6.510646in}}{\pgfqpoint{3.965643in}{6.508332in}}{\pgfqpoint{3.971467in}{6.508332in}}%
\pgfpathlineto{\pgfqpoint{3.971467in}{6.508332in}}%
\pgfpathclose%
\pgfusepath{stroke,fill}%
\end{pgfscope}%
\begin{pgfscope}%
\pgfpathrectangle{\pgfqpoint{1.414736in}{0.548769in}}{\pgfqpoint{3.497747in}{7.253186in}}%
\pgfusepath{clip}%
\pgfsetbuttcap%
\pgfsetroundjoin%
\definecolor{currentfill}{rgb}{0.296479,0.761561,0.424223}%
\pgfsetfillcolor{currentfill}%
\pgfsetlinewidth{1.003750pt}%
\definecolor{currentstroke}{rgb}{0.296479,0.761561,0.424223}%
\pgfsetstrokecolor{currentstroke}%
\pgfsetdash{}{0pt}%
\pgfpathmoveto{\pgfqpoint{3.568077in}{6.979318in}}%
\pgfpathcurveto{\pgfqpoint{3.573901in}{6.979318in}}{\pgfqpoint{3.579487in}{6.981632in}}{\pgfqpoint{3.583605in}{6.985750in}}%
\pgfpathcurveto{\pgfqpoint{3.587723in}{6.989868in}}{\pgfqpoint{3.590037in}{6.995455in}}{\pgfqpoint{3.590037in}{7.001278in}}%
\pgfpathcurveto{\pgfqpoint{3.590037in}{7.007102in}}{\pgfqpoint{3.587723in}{7.012689in}}{\pgfqpoint{3.583605in}{7.016807in}}%
\pgfpathcurveto{\pgfqpoint{3.579487in}{7.020925in}}{\pgfqpoint{3.573901in}{7.023239in}}{\pgfqpoint{3.568077in}{7.023239in}}%
\pgfpathcurveto{\pgfqpoint{3.562253in}{7.023239in}}{\pgfqpoint{3.556667in}{7.020925in}}{\pgfqpoint{3.552548in}{7.016807in}}%
\pgfpathcurveto{\pgfqpoint{3.548430in}{7.012689in}}{\pgfqpoint{3.546116in}{7.007102in}}{\pgfqpoint{3.546116in}{7.001278in}}%
\pgfpathcurveto{\pgfqpoint{3.546116in}{6.995455in}}{\pgfqpoint{3.548430in}{6.989868in}}{\pgfqpoint{3.552548in}{6.985750in}}%
\pgfpathcurveto{\pgfqpoint{3.556667in}{6.981632in}}{\pgfqpoint{3.562253in}{6.979318in}}{\pgfqpoint{3.568077in}{6.979318in}}%
\pgfpathlineto{\pgfqpoint{3.568077in}{6.979318in}}%
\pgfpathclose%
\pgfusepath{stroke,fill}%
\end{pgfscope}%
\begin{pgfscope}%
\pgfpathrectangle{\pgfqpoint{1.414736in}{0.548769in}}{\pgfqpoint{3.497747in}{7.253186in}}%
\pgfusepath{clip}%
\pgfsetbuttcap%
\pgfsetroundjoin%
\definecolor{currentfill}{rgb}{0.741388,0.873449,0.149561}%
\pgfsetfillcolor{currentfill}%
\pgfsetlinewidth{1.003750pt}%
\definecolor{currentstroke}{rgb}{0.741388,0.873449,0.149561}%
\pgfsetstrokecolor{currentstroke}%
\pgfsetdash{}{0pt}%
\pgfpathmoveto{\pgfqpoint{3.861071in}{7.450304in}}%
\pgfpathcurveto{\pgfqpoint{3.866895in}{7.450304in}}{\pgfqpoint{3.872482in}{7.452618in}}{\pgfqpoint{3.876600in}{7.456736in}}%
\pgfpathcurveto{\pgfqpoint{3.880718in}{7.460854in}}{\pgfqpoint{3.883032in}{7.466441in}}{\pgfqpoint{3.883032in}{7.472265in}}%
\pgfpathcurveto{\pgfqpoint{3.883032in}{7.478089in}}{\pgfqpoint{3.880718in}{7.483675in}}{\pgfqpoint{3.876600in}{7.487793in}}%
\pgfpathcurveto{\pgfqpoint{3.872482in}{7.491911in}}{\pgfqpoint{3.866895in}{7.494225in}}{\pgfqpoint{3.861071in}{7.494225in}}%
\pgfpathcurveto{\pgfqpoint{3.855247in}{7.494225in}}{\pgfqpoint{3.849661in}{7.491911in}}{\pgfqpoint{3.845543in}{7.487793in}}%
\pgfpathcurveto{\pgfqpoint{3.841425in}{7.483675in}}{\pgfqpoint{3.839111in}{7.478089in}}{\pgfqpoint{3.839111in}{7.472265in}}%
\pgfpathcurveto{\pgfqpoint{3.839111in}{7.466441in}}{\pgfqpoint{3.841425in}{7.460854in}}{\pgfqpoint{3.845543in}{7.456736in}}%
\pgfpathcurveto{\pgfqpoint{3.849661in}{7.452618in}}{\pgfqpoint{3.855247in}{7.450304in}}{\pgfqpoint{3.861071in}{7.450304in}}%
\pgfpathlineto{\pgfqpoint{3.861071in}{7.450304in}}%
\pgfpathclose%
\pgfusepath{stroke,fill}%
\end{pgfscope}%
\begin{pgfscope}%
\pgfpathrectangle{\pgfqpoint{1.414736in}{0.548769in}}{\pgfqpoint{3.497747in}{7.253186in}}%
\pgfusepath{clip}%
\pgfsetbuttcap%
\pgfsetroundjoin%
\definecolor{currentfill}{rgb}{0.993248,0.906157,0.143936}%
\pgfsetfillcolor{currentfill}%
\pgfsetlinewidth{1.003750pt}%
\definecolor{currentstroke}{rgb}{0.993248,0.906157,0.143936}%
\pgfsetstrokecolor{currentstroke}%
\pgfsetdash{}{0pt}%
\pgfpathmoveto{\pgfqpoint{3.761453in}{0.856499in}}%
\pgfpathcurveto{\pgfqpoint{3.767277in}{0.856499in}}{\pgfqpoint{3.772863in}{0.858813in}}{\pgfqpoint{3.776982in}{0.862931in}}%
\pgfpathcurveto{\pgfqpoint{3.781100in}{0.867049in}}{\pgfqpoint{3.783414in}{0.872635in}}{\pgfqpoint{3.783414in}{0.878459in}}%
\pgfpathcurveto{\pgfqpoint{3.783414in}{0.884283in}}{\pgfqpoint{3.781100in}{0.889869in}}{\pgfqpoint{3.776982in}{0.893987in}}%
\pgfpathcurveto{\pgfqpoint{3.772863in}{0.898105in}}{\pgfqpoint{3.767277in}{0.900419in}}{\pgfqpoint{3.761453in}{0.900419in}}%
\pgfpathcurveto{\pgfqpoint{3.755629in}{0.900419in}}{\pgfqpoint{3.750043in}{0.898105in}}{\pgfqpoint{3.745925in}{0.893987in}}%
\pgfpathcurveto{\pgfqpoint{3.741807in}{0.889869in}}{\pgfqpoint{3.739493in}{0.884283in}}{\pgfqpoint{3.739493in}{0.878459in}}%
\pgfpathcurveto{\pgfqpoint{3.739493in}{0.872635in}}{\pgfqpoint{3.741807in}{0.867049in}}{\pgfqpoint{3.745925in}{0.862931in}}%
\pgfpathcurveto{\pgfqpoint{3.750043in}{0.858813in}}{\pgfqpoint{3.755629in}{0.856499in}}{\pgfqpoint{3.761453in}{0.856499in}}%
\pgfpathlineto{\pgfqpoint{3.761453in}{0.856499in}}%
\pgfpathclose%
\pgfusepath{stroke,fill}%
\end{pgfscope}%
\begin{pgfscope}%
\pgfpathrectangle{\pgfqpoint{1.414736in}{0.548769in}}{\pgfqpoint{3.497747in}{7.253186in}}%
\pgfusepath{clip}%
\pgfsetbuttcap%
\pgfsetroundjoin%
\definecolor{currentfill}{rgb}{0.741388,0.873449,0.149561}%
\pgfsetfillcolor{currentfill}%
\pgfsetlinewidth{1.003750pt}%
\definecolor{currentstroke}{rgb}{0.741388,0.873449,0.149561}%
\pgfsetstrokecolor{currentstroke}%
\pgfsetdash{}{0pt}%
\pgfpathmoveto{\pgfqpoint{3.388600in}{1.327485in}}%
\pgfpathcurveto{\pgfqpoint{3.394424in}{1.327485in}}{\pgfqpoint{3.400010in}{1.329799in}}{\pgfqpoint{3.404129in}{1.333917in}}%
\pgfpathcurveto{\pgfqpoint{3.408247in}{1.338035in}}{\pgfqpoint{3.410561in}{1.343621in}}{\pgfqpoint{3.410561in}{1.349445in}}%
\pgfpathcurveto{\pgfqpoint{3.410561in}{1.355269in}}{\pgfqpoint{3.408247in}{1.360855in}}{\pgfqpoint{3.404129in}{1.364973in}}%
\pgfpathcurveto{\pgfqpoint{3.400010in}{1.369092in}}{\pgfqpoint{3.394424in}{1.371405in}}{\pgfqpoint{3.388600in}{1.371405in}}%
\pgfpathcurveto{\pgfqpoint{3.382776in}{1.371405in}}{\pgfqpoint{3.377190in}{1.369092in}}{\pgfqpoint{3.373072in}{1.364973in}}%
\pgfpathcurveto{\pgfqpoint{3.368954in}{1.360855in}}{\pgfqpoint{3.366640in}{1.355269in}}{\pgfqpoint{3.366640in}{1.349445in}}%
\pgfpathcurveto{\pgfqpoint{3.366640in}{1.343621in}}{\pgfqpoint{3.368954in}{1.338035in}}{\pgfqpoint{3.373072in}{1.333917in}}%
\pgfpathcurveto{\pgfqpoint{3.377190in}{1.329799in}}{\pgfqpoint{3.382776in}{1.327485in}}{\pgfqpoint{3.388600in}{1.327485in}}%
\pgfpathlineto{\pgfqpoint{3.388600in}{1.327485in}}%
\pgfpathclose%
\pgfusepath{stroke,fill}%
\end{pgfscope}%
\begin{pgfscope}%
\pgfpathrectangle{\pgfqpoint{1.414736in}{0.548769in}}{\pgfqpoint{3.497747in}{7.253186in}}%
\pgfusepath{clip}%
\pgfsetbuttcap%
\pgfsetroundjoin%
\definecolor{currentfill}{rgb}{0.699415,0.867117,0.175971}%
\pgfsetfillcolor{currentfill}%
\pgfsetlinewidth{1.003750pt}%
\definecolor{currentstroke}{rgb}{0.699415,0.867117,0.175971}%
\pgfsetstrokecolor{currentstroke}%
\pgfsetdash{}{0pt}%
\pgfpathmoveto{\pgfqpoint{3.823497in}{1.798471in}}%
\pgfpathcurveto{\pgfqpoint{3.829321in}{1.798471in}}{\pgfqpoint{3.834907in}{1.800785in}}{\pgfqpoint{3.839026in}{1.804903in}}%
\pgfpathcurveto{\pgfqpoint{3.843144in}{1.809021in}}{\pgfqpoint{3.845458in}{1.814607in}}{\pgfqpoint{3.845458in}{1.820431in}}%
\pgfpathcurveto{\pgfqpoint{3.845458in}{1.826255in}}{\pgfqpoint{3.843144in}{1.831841in}}{\pgfqpoint{3.839026in}{1.835960in}}%
\pgfpathcurveto{\pgfqpoint{3.834907in}{1.840078in}}{\pgfqpoint{3.829321in}{1.842392in}}{\pgfqpoint{3.823497in}{1.842392in}}%
\pgfpathcurveto{\pgfqpoint{3.817673in}{1.842392in}}{\pgfqpoint{3.812087in}{1.840078in}}{\pgfqpoint{3.807969in}{1.835960in}}%
\pgfpathcurveto{\pgfqpoint{3.803851in}{1.831841in}}{\pgfqpoint{3.801537in}{1.826255in}}{\pgfqpoint{3.801537in}{1.820431in}}%
\pgfpathcurveto{\pgfqpoint{3.801537in}{1.814607in}}{\pgfqpoint{3.803851in}{1.809021in}}{\pgfqpoint{3.807969in}{1.804903in}}%
\pgfpathcurveto{\pgfqpoint{3.812087in}{1.800785in}}{\pgfqpoint{3.817673in}{1.798471in}}{\pgfqpoint{3.823497in}{1.798471in}}%
\pgfpathlineto{\pgfqpoint{3.823497in}{1.798471in}}%
\pgfpathclose%
\pgfusepath{stroke,fill}%
\end{pgfscope}%
\begin{pgfscope}%
\pgfpathrectangle{\pgfqpoint{1.414736in}{0.548769in}}{\pgfqpoint{3.497747in}{7.253186in}}%
\pgfusepath{clip}%
\pgfsetbuttcap%
\pgfsetroundjoin%
\definecolor{currentfill}{rgb}{0.993248,0.906157,0.143936}%
\pgfsetfillcolor{currentfill}%
\pgfsetlinewidth{1.003750pt}%
\definecolor{currentstroke}{rgb}{0.993248,0.906157,0.143936}%
\pgfsetstrokecolor{currentstroke}%
\pgfsetdash{}{0pt}%
\pgfpathmoveto{\pgfqpoint{3.430045in}{2.269457in}}%
\pgfpathcurveto{\pgfqpoint{3.435868in}{2.269457in}}{\pgfqpoint{3.441455in}{2.271771in}}{\pgfqpoint{3.445573in}{2.275889in}}%
\pgfpathcurveto{\pgfqpoint{3.449691in}{2.280007in}}{\pgfqpoint{3.452005in}{2.285593in}}{\pgfqpoint{3.452005in}{2.291417in}}%
\pgfpathcurveto{\pgfqpoint{3.452005in}{2.297241in}}{\pgfqpoint{3.449691in}{2.302828in}}{\pgfqpoint{3.445573in}{2.306946in}}%
\pgfpathcurveto{\pgfqpoint{3.441455in}{2.311064in}}{\pgfqpoint{3.435868in}{2.313378in}}{\pgfqpoint{3.430045in}{2.313378in}}%
\pgfpathcurveto{\pgfqpoint{3.424221in}{2.313378in}}{\pgfqpoint{3.418634in}{2.311064in}}{\pgfqpoint{3.414516in}{2.306946in}}%
\pgfpathcurveto{\pgfqpoint{3.410398in}{2.302828in}}{\pgfqpoint{3.408084in}{2.297241in}}{\pgfqpoint{3.408084in}{2.291417in}}%
\pgfpathcurveto{\pgfqpoint{3.408084in}{2.285593in}}{\pgfqpoint{3.410398in}{2.280007in}}{\pgfqpoint{3.414516in}{2.275889in}}%
\pgfpathcurveto{\pgfqpoint{3.418634in}{2.271771in}}{\pgfqpoint{3.424221in}{2.269457in}}{\pgfqpoint{3.430045in}{2.269457in}}%
\pgfpathlineto{\pgfqpoint{3.430045in}{2.269457in}}%
\pgfpathclose%
\pgfusepath{stroke,fill}%
\end{pgfscope}%
\begin{pgfscope}%
\pgfpathrectangle{\pgfqpoint{1.414736in}{0.548769in}}{\pgfqpoint{3.497747in}{7.253186in}}%
\pgfusepath{clip}%
\pgfsetbuttcap%
\pgfsetroundjoin%
\definecolor{currentfill}{rgb}{0.214000,0.722114,0.469588}%
\pgfsetfillcolor{currentfill}%
\pgfsetlinewidth{1.003750pt}%
\definecolor{currentstroke}{rgb}{0.214000,0.722114,0.469588}%
\pgfsetstrokecolor{currentstroke}%
\pgfsetdash{}{0pt}%
\pgfpathmoveto{\pgfqpoint{4.045539in}{2.740443in}}%
\pgfpathcurveto{\pgfqpoint{4.051363in}{2.740443in}}{\pgfqpoint{4.056950in}{2.742757in}}{\pgfqpoint{4.061068in}{2.746875in}}%
\pgfpathcurveto{\pgfqpoint{4.065186in}{2.750993in}}{\pgfqpoint{4.067500in}{2.756580in}}{\pgfqpoint{4.067500in}{2.762404in}}%
\pgfpathcurveto{\pgfqpoint{4.067500in}{2.768227in}}{\pgfqpoint{4.065186in}{2.773814in}}{\pgfqpoint{4.061068in}{2.777932in}}%
\pgfpathcurveto{\pgfqpoint{4.056950in}{2.782050in}}{\pgfqpoint{4.051363in}{2.784364in}}{\pgfqpoint{4.045539in}{2.784364in}}%
\pgfpathcurveto{\pgfqpoint{4.039716in}{2.784364in}}{\pgfqpoint{4.034129in}{2.782050in}}{\pgfqpoint{4.030011in}{2.777932in}}%
\pgfpathcurveto{\pgfqpoint{4.025893in}{2.773814in}}{\pgfqpoint{4.023579in}{2.768227in}}{\pgfqpoint{4.023579in}{2.762404in}}%
\pgfpathcurveto{\pgfqpoint{4.023579in}{2.756580in}}{\pgfqpoint{4.025893in}{2.750993in}}{\pgfqpoint{4.030011in}{2.746875in}}%
\pgfpathcurveto{\pgfqpoint{4.034129in}{2.742757in}}{\pgfqpoint{4.039716in}{2.740443in}}{\pgfqpoint{4.045539in}{2.740443in}}%
\pgfpathlineto{\pgfqpoint{4.045539in}{2.740443in}}%
\pgfpathclose%
\pgfusepath{stroke,fill}%
\end{pgfscope}%
\begin{pgfscope}%
\pgfpathrectangle{\pgfqpoint{1.414736in}{0.548769in}}{\pgfqpoint{3.497747in}{7.253186in}}%
\pgfusepath{clip}%
\pgfsetbuttcap%
\pgfsetroundjoin%
\definecolor{currentfill}{rgb}{0.575563,0.844566,0.256415}%
\pgfsetfillcolor{currentfill}%
\pgfsetlinewidth{1.003750pt}%
\definecolor{currentstroke}{rgb}{0.575563,0.844566,0.256415}%
\pgfsetstrokecolor{currentstroke}%
\pgfsetdash{}{0pt}%
\pgfpathmoveto{\pgfqpoint{3.310179in}{3.211429in}}%
\pgfpathcurveto{\pgfqpoint{3.316003in}{3.211429in}}{\pgfqpoint{3.321589in}{3.213743in}}{\pgfqpoint{3.325707in}{3.217861in}}%
\pgfpathcurveto{\pgfqpoint{3.329826in}{3.221980in}}{\pgfqpoint{3.332139in}{3.227566in}}{\pgfqpoint{3.332139in}{3.233390in}}%
\pgfpathcurveto{\pgfqpoint{3.332139in}{3.239214in}}{\pgfqpoint{3.329826in}{3.244800in}}{\pgfqpoint{3.325707in}{3.248918in}}%
\pgfpathcurveto{\pgfqpoint{3.321589in}{3.253036in}}{\pgfqpoint{3.316003in}{3.255350in}}{\pgfqpoint{3.310179in}{3.255350in}}%
\pgfpathcurveto{\pgfqpoint{3.304355in}{3.255350in}}{\pgfqpoint{3.298769in}{3.253036in}}{\pgfqpoint{3.294651in}{3.248918in}}%
\pgfpathcurveto{\pgfqpoint{3.290533in}{3.244800in}}{\pgfqpoint{3.288219in}{3.239214in}}{\pgfqpoint{3.288219in}{3.233390in}}%
\pgfpathcurveto{\pgfqpoint{3.288219in}{3.227566in}}{\pgfqpoint{3.290533in}{3.221980in}}{\pgfqpoint{3.294651in}{3.217861in}}%
\pgfpathcurveto{\pgfqpoint{3.298769in}{3.213743in}}{\pgfqpoint{3.304355in}{3.211429in}}{\pgfqpoint{3.310179in}{3.211429in}}%
\pgfpathlineto{\pgfqpoint{3.310179in}{3.211429in}}%
\pgfpathclose%
\pgfusepath{stroke,fill}%
\end{pgfscope}%
\begin{pgfscope}%
\pgfpathrectangle{\pgfqpoint{1.414736in}{0.548769in}}{\pgfqpoint{3.497747in}{7.253186in}}%
\pgfusepath{clip}%
\pgfsetbuttcap%
\pgfsetroundjoin%
\definecolor{currentfill}{rgb}{0.993248,0.906157,0.143936}%
\pgfsetfillcolor{currentfill}%
\pgfsetlinewidth{1.003750pt}%
\definecolor{currentstroke}{rgb}{0.993248,0.906157,0.143936}%
\pgfsetstrokecolor{currentstroke}%
\pgfsetdash{}{0pt}%
\pgfpathmoveto{\pgfqpoint{3.687592in}{3.682415in}}%
\pgfpathcurveto{\pgfqpoint{3.693416in}{3.682415in}}{\pgfqpoint{3.699002in}{3.684729in}}{\pgfqpoint{3.703120in}{3.688847in}}%
\pgfpathcurveto{\pgfqpoint{3.707238in}{3.692966in}}{\pgfqpoint{3.709552in}{3.698552in}}{\pgfqpoint{3.709552in}{3.704376in}}%
\pgfpathcurveto{\pgfqpoint{3.709552in}{3.710200in}}{\pgfqpoint{3.707238in}{3.715786in}}{\pgfqpoint{3.703120in}{3.719904in}}%
\pgfpathcurveto{\pgfqpoint{3.699002in}{3.724022in}}{\pgfqpoint{3.693416in}{3.726336in}}{\pgfqpoint{3.687592in}{3.726336in}}%
\pgfpathcurveto{\pgfqpoint{3.681768in}{3.726336in}}{\pgfqpoint{3.676182in}{3.724022in}}{\pgfqpoint{3.672063in}{3.719904in}}%
\pgfpathcurveto{\pgfqpoint{3.667945in}{3.715786in}}{\pgfqpoint{3.665631in}{3.710200in}}{\pgfqpoint{3.665631in}{3.704376in}}%
\pgfpathcurveto{\pgfqpoint{3.665631in}{3.698552in}}{\pgfqpoint{3.667945in}{3.692966in}}{\pgfqpoint{3.672063in}{3.688847in}}%
\pgfpathcurveto{\pgfqpoint{3.676182in}{3.684729in}}{\pgfqpoint{3.681768in}{3.682415in}}{\pgfqpoint{3.687592in}{3.682415in}}%
\pgfpathlineto{\pgfqpoint{3.687592in}{3.682415in}}%
\pgfpathclose%
\pgfusepath{stroke,fill}%
\end{pgfscope}%
\begin{pgfscope}%
\pgfpathrectangle{\pgfqpoint{1.414736in}{0.548769in}}{\pgfqpoint{3.497747in}{7.253186in}}%
\pgfusepath{clip}%
\pgfsetbuttcap%
\pgfsetroundjoin%
\definecolor{currentfill}{rgb}{0.993248,0.906157,0.143936}%
\pgfsetfillcolor{currentfill}%
\pgfsetlinewidth{1.003750pt}%
\definecolor{currentstroke}{rgb}{0.993248,0.906157,0.143936}%
\pgfsetstrokecolor{currentstroke}%
\pgfsetdash{}{0pt}%
\pgfpathmoveto{\pgfqpoint{3.425195in}{4.153402in}}%
\pgfpathcurveto{\pgfqpoint{3.431019in}{4.153402in}}{\pgfqpoint{3.436605in}{4.155715in}}{\pgfqpoint{3.440723in}{4.159834in}}%
\pgfpathcurveto{\pgfqpoint{3.444842in}{4.163952in}}{\pgfqpoint{3.447155in}{4.169538in}}{\pgfqpoint{3.447155in}{4.175362in}}%
\pgfpathcurveto{\pgfqpoint{3.447155in}{4.181186in}}{\pgfqpoint{3.444842in}{4.186772in}}{\pgfqpoint{3.440723in}{4.190890in}}%
\pgfpathcurveto{\pgfqpoint{3.436605in}{4.195008in}}{\pgfqpoint{3.431019in}{4.197322in}}{\pgfqpoint{3.425195in}{4.197322in}}%
\pgfpathcurveto{\pgfqpoint{3.419371in}{4.197322in}}{\pgfqpoint{3.413785in}{4.195008in}}{\pgfqpoint{3.409667in}{4.190890in}}%
\pgfpathcurveto{\pgfqpoint{3.405549in}{4.186772in}}{\pgfqpoint{3.403235in}{4.181186in}}{\pgfqpoint{3.403235in}{4.175362in}}%
\pgfpathcurveto{\pgfqpoint{3.403235in}{4.169538in}}{\pgfqpoint{3.405549in}{4.163952in}}{\pgfqpoint{3.409667in}{4.159834in}}%
\pgfpathcurveto{\pgfqpoint{3.413785in}{4.155715in}}{\pgfqpoint{3.419371in}{4.153402in}}{\pgfqpoint{3.425195in}{4.153402in}}%
\pgfpathlineto{\pgfqpoint{3.425195in}{4.153402in}}%
\pgfpathclose%
\pgfusepath{stroke,fill}%
\end{pgfscope}%
\begin{pgfscope}%
\pgfpathrectangle{\pgfqpoint{1.414736in}{0.548769in}}{\pgfqpoint{3.497747in}{7.253186in}}%
\pgfusepath{clip}%
\pgfsetbuttcap%
\pgfsetroundjoin%
\definecolor{currentfill}{rgb}{0.916242,0.896091,0.100717}%
\pgfsetfillcolor{currentfill}%
\pgfsetlinewidth{1.003750pt}%
\definecolor{currentstroke}{rgb}{0.916242,0.896091,0.100717}%
\pgfsetstrokecolor{currentstroke}%
\pgfsetdash{}{0pt}%
\pgfpathmoveto{\pgfqpoint{4.275074in}{4.624388in}}%
\pgfpathcurveto{\pgfqpoint{4.280898in}{4.624388in}}{\pgfqpoint{4.286484in}{4.626702in}}{\pgfqpoint{4.290602in}{4.630820in}}%
\pgfpathcurveto{\pgfqpoint{4.294720in}{4.634938in}}{\pgfqpoint{4.297034in}{4.640524in}}{\pgfqpoint{4.297034in}{4.646348in}}%
\pgfpathcurveto{\pgfqpoint{4.297034in}{4.652172in}}{\pgfqpoint{4.294720in}{4.657758in}}{\pgfqpoint{4.290602in}{4.661876in}}%
\pgfpathcurveto{\pgfqpoint{4.286484in}{4.665994in}}{\pgfqpoint{4.280898in}{4.668308in}}{\pgfqpoint{4.275074in}{4.668308in}}%
\pgfpathcurveto{\pgfqpoint{4.269250in}{4.668308in}}{\pgfqpoint{4.263664in}{4.665994in}}{\pgfqpoint{4.259546in}{4.661876in}}%
\pgfpathcurveto{\pgfqpoint{4.255428in}{4.657758in}}{\pgfqpoint{4.253114in}{4.652172in}}{\pgfqpoint{4.253114in}{4.646348in}}%
\pgfpathcurveto{\pgfqpoint{4.253114in}{4.640524in}}{\pgfqpoint{4.255428in}{4.634938in}}{\pgfqpoint{4.259546in}{4.630820in}}%
\pgfpathcurveto{\pgfqpoint{4.263664in}{4.626702in}}{\pgfqpoint{4.269250in}{4.624388in}}{\pgfqpoint{4.275074in}{4.624388in}}%
\pgfpathlineto{\pgfqpoint{4.275074in}{4.624388in}}%
\pgfpathclose%
\pgfusepath{stroke,fill}%
\end{pgfscope}%
\begin{pgfscope}%
\pgfpathrectangle{\pgfqpoint{1.414736in}{0.548769in}}{\pgfqpoint{3.497747in}{7.253186in}}%
\pgfusepath{clip}%
\pgfsetbuttcap%
\pgfsetroundjoin%
\definecolor{currentfill}{rgb}{0.288921,0.758394,0.428426}%
\pgfsetfillcolor{currentfill}%
\pgfsetlinewidth{1.003750pt}%
\definecolor{currentstroke}{rgb}{0.288921,0.758394,0.428426}%
\pgfsetstrokecolor{currentstroke}%
\pgfsetdash{}{0pt}%
\pgfpathmoveto{\pgfqpoint{3.427911in}{5.095374in}}%
\pgfpathcurveto{\pgfqpoint{3.433735in}{5.095374in}}{\pgfqpoint{3.439321in}{5.097688in}}{\pgfqpoint{3.443439in}{5.101806in}}%
\pgfpathcurveto{\pgfqpoint{3.447557in}{5.105924in}}{\pgfqpoint{3.449871in}{5.111510in}}{\pgfqpoint{3.449871in}{5.117334in}}%
\pgfpathcurveto{\pgfqpoint{3.449871in}{5.123158in}}{\pgfqpoint{3.447557in}{5.128744in}}{\pgfqpoint{3.443439in}{5.132862in}}%
\pgfpathcurveto{\pgfqpoint{3.439321in}{5.136980in}}{\pgfqpoint{3.433735in}{5.139294in}}{\pgfqpoint{3.427911in}{5.139294in}}%
\pgfpathcurveto{\pgfqpoint{3.422087in}{5.139294in}}{\pgfqpoint{3.416501in}{5.136980in}}{\pgfqpoint{3.412383in}{5.132862in}}%
\pgfpathcurveto{\pgfqpoint{3.408265in}{5.128744in}}{\pgfqpoint{3.405951in}{5.123158in}}{\pgfqpoint{3.405951in}{5.117334in}}%
\pgfpathcurveto{\pgfqpoint{3.405951in}{5.111510in}}{\pgfqpoint{3.408265in}{5.105924in}}{\pgfqpoint{3.412383in}{5.101806in}}%
\pgfpathcurveto{\pgfqpoint{3.416501in}{5.097688in}}{\pgfqpoint{3.422087in}{5.095374in}}{\pgfqpoint{3.427911in}{5.095374in}}%
\pgfpathlineto{\pgfqpoint{3.427911in}{5.095374in}}%
\pgfpathclose%
\pgfusepath{stroke,fill}%
\end{pgfscope}%
\begin{pgfscope}%
\pgfpathrectangle{\pgfqpoint{1.414736in}{0.548769in}}{\pgfqpoint{3.497747in}{7.253186in}}%
\pgfusepath{clip}%
\pgfsetbuttcap%
\pgfsetroundjoin%
\definecolor{currentfill}{rgb}{0.174274,0.445044,0.557792}%
\pgfsetfillcolor{currentfill}%
\pgfsetlinewidth{1.003750pt}%
\definecolor{currentstroke}{rgb}{0.174274,0.445044,0.557792}%
\pgfsetstrokecolor{currentstroke}%
\pgfsetdash{}{0pt}%
\pgfpathmoveto{\pgfqpoint{3.718277in}{5.566360in}}%
\pgfpathcurveto{\pgfqpoint{3.724101in}{5.566360in}}{\pgfqpoint{3.729687in}{5.568674in}}{\pgfqpoint{3.733805in}{5.572792in}}%
\pgfpathcurveto{\pgfqpoint{3.737923in}{5.576910in}}{\pgfqpoint{3.740237in}{5.582496in}}{\pgfqpoint{3.740237in}{5.588320in}}%
\pgfpathcurveto{\pgfqpoint{3.740237in}{5.594144in}}{\pgfqpoint{3.737923in}{5.599730in}}{\pgfqpoint{3.733805in}{5.603848in}}%
\pgfpathcurveto{\pgfqpoint{3.729687in}{5.607967in}}{\pgfqpoint{3.724101in}{5.610280in}}{\pgfqpoint{3.718277in}{5.610280in}}%
\pgfpathcurveto{\pgfqpoint{3.712453in}{5.610280in}}{\pgfqpoint{3.706867in}{5.607967in}}{\pgfqpoint{3.702749in}{5.603848in}}%
\pgfpathcurveto{\pgfqpoint{3.698631in}{5.599730in}}{\pgfqpoint{3.696317in}{5.594144in}}{\pgfqpoint{3.696317in}{5.588320in}}%
\pgfpathcurveto{\pgfqpoint{3.696317in}{5.582496in}}{\pgfqpoint{3.698631in}{5.576910in}}{\pgfqpoint{3.702749in}{5.572792in}}%
\pgfpathcurveto{\pgfqpoint{3.706867in}{5.568674in}}{\pgfqpoint{3.712453in}{5.566360in}}{\pgfqpoint{3.718277in}{5.566360in}}%
\pgfpathlineto{\pgfqpoint{3.718277in}{5.566360in}}%
\pgfpathclose%
\pgfusepath{stroke,fill}%
\end{pgfscope}%
\begin{pgfscope}%
\pgfpathrectangle{\pgfqpoint{1.414736in}{0.548769in}}{\pgfqpoint{3.497747in}{7.253186in}}%
\pgfusepath{clip}%
\pgfsetbuttcap%
\pgfsetroundjoin%
\definecolor{currentfill}{rgb}{0.377779,0.791781,0.377939}%
\pgfsetfillcolor{currentfill}%
\pgfsetlinewidth{1.003750pt}%
\definecolor{currentstroke}{rgb}{0.377779,0.791781,0.377939}%
\pgfsetstrokecolor{currentstroke}%
\pgfsetdash{}{0pt}%
\pgfpathmoveto{\pgfqpoint{3.118203in}{6.037346in}}%
\pgfpathcurveto{\pgfqpoint{3.124027in}{6.037346in}}{\pgfqpoint{3.129613in}{6.039660in}}{\pgfqpoint{3.133732in}{6.043778in}}%
\pgfpathcurveto{\pgfqpoint{3.137850in}{6.047896in}}{\pgfqpoint{3.140164in}{6.053482in}}{\pgfqpoint{3.140164in}{6.059306in}}%
\pgfpathcurveto{\pgfqpoint{3.140164in}{6.065130in}}{\pgfqpoint{3.137850in}{6.070716in}}{\pgfqpoint{3.133732in}{6.074835in}}%
\pgfpathcurveto{\pgfqpoint{3.129613in}{6.078953in}}{\pgfqpoint{3.124027in}{6.081267in}}{\pgfqpoint{3.118203in}{6.081267in}}%
\pgfpathcurveto{\pgfqpoint{3.112379in}{6.081267in}}{\pgfqpoint{3.106793in}{6.078953in}}{\pgfqpoint{3.102675in}{6.074835in}}%
\pgfpathcurveto{\pgfqpoint{3.098557in}{6.070716in}}{\pgfqpoint{3.096243in}{6.065130in}}{\pgfqpoint{3.096243in}{6.059306in}}%
\pgfpathcurveto{\pgfqpoint{3.096243in}{6.053482in}}{\pgfqpoint{3.098557in}{6.047896in}}{\pgfqpoint{3.102675in}{6.043778in}}%
\pgfpathcurveto{\pgfqpoint{3.106793in}{6.039660in}}{\pgfqpoint{3.112379in}{6.037346in}}{\pgfqpoint{3.118203in}{6.037346in}}%
\pgfpathlineto{\pgfqpoint{3.118203in}{6.037346in}}%
\pgfpathclose%
\pgfusepath{stroke,fill}%
\end{pgfscope}%
\begin{pgfscope}%
\pgfpathrectangle{\pgfqpoint{1.414736in}{0.548769in}}{\pgfqpoint{3.497747in}{7.253186in}}%
\pgfusepath{clip}%
\pgfsetbuttcap%
\pgfsetroundjoin%
\definecolor{currentfill}{rgb}{0.993248,0.906157,0.143936}%
\pgfsetfillcolor{currentfill}%
\pgfsetlinewidth{1.003750pt}%
\definecolor{currentstroke}{rgb}{0.993248,0.906157,0.143936}%
\pgfsetstrokecolor{currentstroke}%
\pgfsetdash{}{0pt}%
\pgfpathmoveto{\pgfqpoint{3.096355in}{6.508332in}}%
\pgfpathcurveto{\pgfqpoint{3.102179in}{6.508332in}}{\pgfqpoint{3.107765in}{6.510646in}}{\pgfqpoint{3.111883in}{6.514764in}}%
\pgfpathcurveto{\pgfqpoint{3.116001in}{6.518882in}}{\pgfqpoint{3.118315in}{6.524468in}}{\pgfqpoint{3.118315in}{6.530292in}}%
\pgfpathcurveto{\pgfqpoint{3.118315in}{6.536116in}}{\pgfqpoint{3.116001in}{6.541702in}}{\pgfqpoint{3.111883in}{6.545821in}}%
\pgfpathcurveto{\pgfqpoint{3.107765in}{6.549939in}}{\pgfqpoint{3.102179in}{6.552253in}}{\pgfqpoint{3.096355in}{6.552253in}}%
\pgfpathcurveto{\pgfqpoint{3.090531in}{6.552253in}}{\pgfqpoint{3.084945in}{6.549939in}}{\pgfqpoint{3.080827in}{6.545821in}}%
\pgfpathcurveto{\pgfqpoint{3.076708in}{6.541702in}}{\pgfqpoint{3.074395in}{6.536116in}}{\pgfqpoint{3.074395in}{6.530292in}}%
\pgfpathcurveto{\pgfqpoint{3.074395in}{6.524468in}}{\pgfqpoint{3.076708in}{6.518882in}}{\pgfqpoint{3.080827in}{6.514764in}}%
\pgfpathcurveto{\pgfqpoint{3.084945in}{6.510646in}}{\pgfqpoint{3.090531in}{6.508332in}}{\pgfqpoint{3.096355in}{6.508332in}}%
\pgfpathlineto{\pgfqpoint{3.096355in}{6.508332in}}%
\pgfpathclose%
\pgfusepath{stroke,fill}%
\end{pgfscope}%
\begin{pgfscope}%
\pgfpathrectangle{\pgfqpoint{1.414736in}{0.548769in}}{\pgfqpoint{3.497747in}{7.253186in}}%
\pgfusepath{clip}%
\pgfsetbuttcap%
\pgfsetroundjoin%
\definecolor{currentfill}{rgb}{0.886271,0.892374,0.095374}%
\pgfsetfillcolor{currentfill}%
\pgfsetlinewidth{1.003750pt}%
\definecolor{currentstroke}{rgb}{0.886271,0.892374,0.095374}%
\pgfsetstrokecolor{currentstroke}%
\pgfsetdash{}{0pt}%
\pgfpathmoveto{\pgfqpoint{2.631410in}{6.979318in}}%
\pgfpathcurveto{\pgfqpoint{2.637234in}{6.979318in}}{\pgfqpoint{2.642820in}{6.981632in}}{\pgfqpoint{2.646938in}{6.985750in}}%
\pgfpathcurveto{\pgfqpoint{2.651056in}{6.989868in}}{\pgfqpoint{2.653370in}{6.995455in}}{\pgfqpoint{2.653370in}{7.001278in}}%
\pgfpathcurveto{\pgfqpoint{2.653370in}{7.007102in}}{\pgfqpoint{2.651056in}{7.012689in}}{\pgfqpoint{2.646938in}{7.016807in}}%
\pgfpathcurveto{\pgfqpoint{2.642820in}{7.020925in}}{\pgfqpoint{2.637234in}{7.023239in}}{\pgfqpoint{2.631410in}{7.023239in}}%
\pgfpathcurveto{\pgfqpoint{2.625586in}{7.023239in}}{\pgfqpoint{2.620000in}{7.020925in}}{\pgfqpoint{2.615882in}{7.016807in}}%
\pgfpathcurveto{\pgfqpoint{2.611764in}{7.012689in}}{\pgfqpoint{2.609450in}{7.007102in}}{\pgfqpoint{2.609450in}{7.001278in}}%
\pgfpathcurveto{\pgfqpoint{2.609450in}{6.995455in}}{\pgfqpoint{2.611764in}{6.989868in}}{\pgfqpoint{2.615882in}{6.985750in}}%
\pgfpathcurveto{\pgfqpoint{2.620000in}{6.981632in}}{\pgfqpoint{2.625586in}{6.979318in}}{\pgfqpoint{2.631410in}{6.979318in}}%
\pgfpathlineto{\pgfqpoint{2.631410in}{6.979318in}}%
\pgfpathclose%
\pgfusepath{stroke,fill}%
\end{pgfscope}%
\begin{pgfscope}%
\pgfpathrectangle{\pgfqpoint{1.414736in}{0.548769in}}{\pgfqpoint{3.497747in}{7.253186in}}%
\pgfusepath{clip}%
\pgfsetbuttcap%
\pgfsetroundjoin%
\definecolor{currentfill}{rgb}{0.916242,0.896091,0.100717}%
\pgfsetfillcolor{currentfill}%
\pgfsetlinewidth{1.003750pt}%
\definecolor{currentstroke}{rgb}{0.916242,0.896091,0.100717}%
\pgfsetstrokecolor{currentstroke}%
\pgfsetdash{}{0pt}%
\pgfpathmoveto{\pgfqpoint{3.286290in}{7.450304in}}%
\pgfpathcurveto{\pgfqpoint{3.292114in}{7.450304in}}{\pgfqpoint{3.297700in}{7.452618in}}{\pgfqpoint{3.301818in}{7.456736in}}%
\pgfpathcurveto{\pgfqpoint{3.305936in}{7.460854in}}{\pgfqpoint{3.308250in}{7.466441in}}{\pgfqpoint{3.308250in}{7.472265in}}%
\pgfpathcurveto{\pgfqpoint{3.308250in}{7.478089in}}{\pgfqpoint{3.305936in}{7.483675in}}{\pgfqpoint{3.301818in}{7.487793in}}%
\pgfpathcurveto{\pgfqpoint{3.297700in}{7.491911in}}{\pgfqpoint{3.292114in}{7.494225in}}{\pgfqpoint{3.286290in}{7.494225in}}%
\pgfpathcurveto{\pgfqpoint{3.280466in}{7.494225in}}{\pgfqpoint{3.274880in}{7.491911in}}{\pgfqpoint{3.270761in}{7.487793in}}%
\pgfpathcurveto{\pgfqpoint{3.266643in}{7.483675in}}{\pgfqpoint{3.264329in}{7.478089in}}{\pgfqpoint{3.264329in}{7.472265in}}%
\pgfpathcurveto{\pgfqpoint{3.264329in}{7.466441in}}{\pgfqpoint{3.266643in}{7.460854in}}{\pgfqpoint{3.270761in}{7.456736in}}%
\pgfpathcurveto{\pgfqpoint{3.274880in}{7.452618in}}{\pgfqpoint{3.280466in}{7.450304in}}{\pgfqpoint{3.286290in}{7.450304in}}%
\pgfpathlineto{\pgfqpoint{3.286290in}{7.450304in}}%
\pgfpathclose%
\pgfusepath{stroke,fill}%
\end{pgfscope}%
\begin{pgfscope}%
\pgfpathrectangle{\pgfqpoint{1.414736in}{0.548769in}}{\pgfqpoint{3.497747in}{7.253186in}}%
\pgfusepath{clip}%
\pgfsetbuttcap%
\pgfsetroundjoin%
\definecolor{currentfill}{rgb}{0.993248,0.906157,0.143936}%
\pgfsetfillcolor{currentfill}%
\pgfsetlinewidth{1.003750pt}%
\definecolor{currentstroke}{rgb}{0.993248,0.906157,0.143936}%
\pgfsetstrokecolor{currentstroke}%
\pgfsetdash{}{0pt}%
\pgfpathmoveto{\pgfqpoint{3.653745in}{0.856499in}}%
\pgfpathcurveto{\pgfqpoint{3.659569in}{0.856499in}}{\pgfqpoint{3.665155in}{0.858813in}}{\pgfqpoint{3.669273in}{0.862931in}}%
\pgfpathcurveto{\pgfqpoint{3.673391in}{0.867049in}}{\pgfqpoint{3.675705in}{0.872635in}}{\pgfqpoint{3.675705in}{0.878459in}}%
\pgfpathcurveto{\pgfqpoint{3.675705in}{0.884283in}}{\pgfqpoint{3.673391in}{0.889869in}}{\pgfqpoint{3.669273in}{0.893987in}}%
\pgfpathcurveto{\pgfqpoint{3.665155in}{0.898105in}}{\pgfqpoint{3.659569in}{0.900419in}}{\pgfqpoint{3.653745in}{0.900419in}}%
\pgfpathcurveto{\pgfqpoint{3.647921in}{0.900419in}}{\pgfqpoint{3.642335in}{0.898105in}}{\pgfqpoint{3.638217in}{0.893987in}}%
\pgfpathcurveto{\pgfqpoint{3.634098in}{0.889869in}}{\pgfqpoint{3.631785in}{0.884283in}}{\pgfqpoint{3.631785in}{0.878459in}}%
\pgfpathcurveto{\pgfqpoint{3.631785in}{0.872635in}}{\pgfqpoint{3.634098in}{0.867049in}}{\pgfqpoint{3.638217in}{0.862931in}}%
\pgfpathcurveto{\pgfqpoint{3.642335in}{0.858813in}}{\pgfqpoint{3.647921in}{0.856499in}}{\pgfqpoint{3.653745in}{0.856499in}}%
\pgfpathlineto{\pgfqpoint{3.653745in}{0.856499in}}%
\pgfpathclose%
\pgfusepath{stroke,fill}%
\end{pgfscope}%
\begin{pgfscope}%
\pgfpathrectangle{\pgfqpoint{1.414736in}{0.548769in}}{\pgfqpoint{3.497747in}{7.253186in}}%
\pgfusepath{clip}%
\pgfsetbuttcap%
\pgfsetroundjoin%
\definecolor{currentfill}{rgb}{0.515992,0.831158,0.294279}%
\pgfsetfillcolor{currentfill}%
\pgfsetlinewidth{1.003750pt}%
\definecolor{currentstroke}{rgb}{0.515992,0.831158,0.294279}%
\pgfsetstrokecolor{currentstroke}%
\pgfsetdash{}{0pt}%
\pgfpathmoveto{\pgfqpoint{3.557076in}{1.327485in}}%
\pgfpathcurveto{\pgfqpoint{3.562899in}{1.327485in}}{\pgfqpoint{3.568486in}{1.329799in}}{\pgfqpoint{3.572604in}{1.333917in}}%
\pgfpathcurveto{\pgfqpoint{3.576722in}{1.338035in}}{\pgfqpoint{3.579036in}{1.343621in}}{\pgfqpoint{3.579036in}{1.349445in}}%
\pgfpathcurveto{\pgfqpoint{3.579036in}{1.355269in}}{\pgfqpoint{3.576722in}{1.360855in}}{\pgfqpoint{3.572604in}{1.364973in}}%
\pgfpathcurveto{\pgfqpoint{3.568486in}{1.369092in}}{\pgfqpoint{3.562899in}{1.371405in}}{\pgfqpoint{3.557076in}{1.371405in}}%
\pgfpathcurveto{\pgfqpoint{3.551252in}{1.371405in}}{\pgfqpoint{3.545665in}{1.369092in}}{\pgfqpoint{3.541547in}{1.364973in}}%
\pgfpathcurveto{\pgfqpoint{3.537429in}{1.360855in}}{\pgfqpoint{3.535115in}{1.355269in}}{\pgfqpoint{3.535115in}{1.349445in}}%
\pgfpathcurveto{\pgfqpoint{3.535115in}{1.343621in}}{\pgfqpoint{3.537429in}{1.338035in}}{\pgfqpoint{3.541547in}{1.333917in}}%
\pgfpathcurveto{\pgfqpoint{3.545665in}{1.329799in}}{\pgfqpoint{3.551252in}{1.327485in}}{\pgfqpoint{3.557076in}{1.327485in}}%
\pgfpathlineto{\pgfqpoint{3.557076in}{1.327485in}}%
\pgfpathclose%
\pgfusepath{stroke,fill}%
\end{pgfscope}%
\begin{pgfscope}%
\pgfpathrectangle{\pgfqpoint{1.414736in}{0.548769in}}{\pgfqpoint{3.497747in}{7.253186in}}%
\pgfusepath{clip}%
\pgfsetbuttcap%
\pgfsetroundjoin%
\definecolor{currentfill}{rgb}{0.430983,0.808473,0.346476}%
\pgfsetfillcolor{currentfill}%
\pgfsetlinewidth{1.003750pt}%
\definecolor{currentstroke}{rgb}{0.430983,0.808473,0.346476}%
\pgfsetstrokecolor{currentstroke}%
\pgfsetdash{}{0pt}%
\pgfpathmoveto{\pgfqpoint{3.350784in}{1.798471in}}%
\pgfpathcurveto{\pgfqpoint{3.356608in}{1.798471in}}{\pgfqpoint{3.362194in}{1.800785in}}{\pgfqpoint{3.366313in}{1.804903in}}%
\pgfpathcurveto{\pgfqpoint{3.370431in}{1.809021in}}{\pgfqpoint{3.372745in}{1.814607in}}{\pgfqpoint{3.372745in}{1.820431in}}%
\pgfpathcurveto{\pgfqpoint{3.372745in}{1.826255in}}{\pgfqpoint{3.370431in}{1.831841in}}{\pgfqpoint{3.366313in}{1.835960in}}%
\pgfpathcurveto{\pgfqpoint{3.362194in}{1.840078in}}{\pgfqpoint{3.356608in}{1.842392in}}{\pgfqpoint{3.350784in}{1.842392in}}%
\pgfpathcurveto{\pgfqpoint{3.344960in}{1.842392in}}{\pgfqpoint{3.339374in}{1.840078in}}{\pgfqpoint{3.335256in}{1.835960in}}%
\pgfpathcurveto{\pgfqpoint{3.331138in}{1.831841in}}{\pgfqpoint{3.328824in}{1.826255in}}{\pgfqpoint{3.328824in}{1.820431in}}%
\pgfpathcurveto{\pgfqpoint{3.328824in}{1.814607in}}{\pgfqpoint{3.331138in}{1.809021in}}{\pgfqpoint{3.335256in}{1.804903in}}%
\pgfpathcurveto{\pgfqpoint{3.339374in}{1.800785in}}{\pgfqpoint{3.344960in}{1.798471in}}{\pgfqpoint{3.350784in}{1.798471in}}%
\pgfpathlineto{\pgfqpoint{3.350784in}{1.798471in}}%
\pgfpathclose%
\pgfusepath{stroke,fill}%
\end{pgfscope}%
\begin{pgfscope}%
\pgfpathrectangle{\pgfqpoint{1.414736in}{0.548769in}}{\pgfqpoint{3.497747in}{7.253186in}}%
\pgfusepath{clip}%
\pgfsetbuttcap%
\pgfsetroundjoin%
\definecolor{currentfill}{rgb}{0.993248,0.906157,0.143936}%
\pgfsetfillcolor{currentfill}%
\pgfsetlinewidth{1.003750pt}%
\definecolor{currentstroke}{rgb}{0.993248,0.906157,0.143936}%
\pgfsetstrokecolor{currentstroke}%
\pgfsetdash{}{0pt}%
\pgfpathmoveto{\pgfqpoint{3.080849in}{2.269457in}}%
\pgfpathcurveto{\pgfqpoint{3.086673in}{2.269457in}}{\pgfqpoint{3.092259in}{2.271771in}}{\pgfqpoint{3.096377in}{2.275889in}}%
\pgfpathcurveto{\pgfqpoint{3.100495in}{2.280007in}}{\pgfqpoint{3.102809in}{2.285593in}}{\pgfqpoint{3.102809in}{2.291417in}}%
\pgfpathcurveto{\pgfqpoint{3.102809in}{2.297241in}}{\pgfqpoint{3.100495in}{2.302828in}}{\pgfqpoint{3.096377in}{2.306946in}}%
\pgfpathcurveto{\pgfqpoint{3.092259in}{2.311064in}}{\pgfqpoint{3.086673in}{2.313378in}}{\pgfqpoint{3.080849in}{2.313378in}}%
\pgfpathcurveto{\pgfqpoint{3.075025in}{2.313378in}}{\pgfqpoint{3.069439in}{2.311064in}}{\pgfqpoint{3.065320in}{2.306946in}}%
\pgfpathcurveto{\pgfqpoint{3.061202in}{2.302828in}}{\pgfqpoint{3.058888in}{2.297241in}}{\pgfqpoint{3.058888in}{2.291417in}}%
\pgfpathcurveto{\pgfqpoint{3.058888in}{2.285593in}}{\pgfqpoint{3.061202in}{2.280007in}}{\pgfqpoint{3.065320in}{2.275889in}}%
\pgfpathcurveto{\pgfqpoint{3.069439in}{2.271771in}}{\pgfqpoint{3.075025in}{2.269457in}}{\pgfqpoint{3.080849in}{2.269457in}}%
\pgfpathlineto{\pgfqpoint{3.080849in}{2.269457in}}%
\pgfpathclose%
\pgfusepath{stroke,fill}%
\end{pgfscope}%
\begin{pgfscope}%
\pgfpathrectangle{\pgfqpoint{1.414736in}{0.548769in}}{\pgfqpoint{3.497747in}{7.253186in}}%
\pgfusepath{clip}%
\pgfsetbuttcap%
\pgfsetroundjoin%
\definecolor{currentfill}{rgb}{0.232815,0.732247,0.459277}%
\pgfsetfillcolor{currentfill}%
\pgfsetlinewidth{1.003750pt}%
\definecolor{currentstroke}{rgb}{0.232815,0.732247,0.459277}%
\pgfsetstrokecolor{currentstroke}%
\pgfsetdash{}{0pt}%
\pgfpathmoveto{\pgfqpoint{3.954902in}{2.740443in}}%
\pgfpathcurveto{\pgfqpoint{3.960726in}{2.740443in}}{\pgfqpoint{3.966313in}{2.742757in}}{\pgfqpoint{3.970431in}{2.746875in}}%
\pgfpathcurveto{\pgfqpoint{3.974549in}{2.750993in}}{\pgfqpoint{3.976863in}{2.756580in}}{\pgfqpoint{3.976863in}{2.762404in}}%
\pgfpathcurveto{\pgfqpoint{3.976863in}{2.768227in}}{\pgfqpoint{3.974549in}{2.773814in}}{\pgfqpoint{3.970431in}{2.777932in}}%
\pgfpathcurveto{\pgfqpoint{3.966313in}{2.782050in}}{\pgfqpoint{3.960726in}{2.784364in}}{\pgfqpoint{3.954902in}{2.784364in}}%
\pgfpathcurveto{\pgfqpoint{3.949078in}{2.784364in}}{\pgfqpoint{3.943492in}{2.782050in}}{\pgfqpoint{3.939374in}{2.777932in}}%
\pgfpathcurveto{\pgfqpoint{3.935256in}{2.773814in}}{\pgfqpoint{3.932942in}{2.768227in}}{\pgfqpoint{3.932942in}{2.762404in}}%
\pgfpathcurveto{\pgfqpoint{3.932942in}{2.756580in}}{\pgfqpoint{3.935256in}{2.750993in}}{\pgfqpoint{3.939374in}{2.746875in}}%
\pgfpathcurveto{\pgfqpoint{3.943492in}{2.742757in}}{\pgfqpoint{3.949078in}{2.740443in}}{\pgfqpoint{3.954902in}{2.740443in}}%
\pgfpathlineto{\pgfqpoint{3.954902in}{2.740443in}}%
\pgfpathclose%
\pgfusepath{stroke,fill}%
\end{pgfscope}%
\begin{pgfscope}%
\pgfpathrectangle{\pgfqpoint{1.414736in}{0.548769in}}{\pgfqpoint{3.497747in}{7.253186in}}%
\pgfusepath{clip}%
\pgfsetbuttcap%
\pgfsetroundjoin%
\definecolor{currentfill}{rgb}{0.606045,0.850733,0.236712}%
\pgfsetfillcolor{currentfill}%
\pgfsetlinewidth{1.003750pt}%
\definecolor{currentstroke}{rgb}{0.606045,0.850733,0.236712}%
\pgfsetstrokecolor{currentstroke}%
\pgfsetdash{}{0pt}%
\pgfpathmoveto{\pgfqpoint{3.263158in}{3.211429in}}%
\pgfpathcurveto{\pgfqpoint{3.268982in}{3.211429in}}{\pgfqpoint{3.274569in}{3.213743in}}{\pgfqpoint{3.278687in}{3.217861in}}%
\pgfpathcurveto{\pgfqpoint{3.282805in}{3.221980in}}{\pgfqpoint{3.285119in}{3.227566in}}{\pgfqpoint{3.285119in}{3.233390in}}%
\pgfpathcurveto{\pgfqpoint{3.285119in}{3.239214in}}{\pgfqpoint{3.282805in}{3.244800in}}{\pgfqpoint{3.278687in}{3.248918in}}%
\pgfpathcurveto{\pgfqpoint{3.274569in}{3.253036in}}{\pgfqpoint{3.268982in}{3.255350in}}{\pgfqpoint{3.263158in}{3.255350in}}%
\pgfpathcurveto{\pgfqpoint{3.257334in}{3.255350in}}{\pgfqpoint{3.251748in}{3.253036in}}{\pgfqpoint{3.247630in}{3.248918in}}%
\pgfpathcurveto{\pgfqpoint{3.243512in}{3.244800in}}{\pgfqpoint{3.241198in}{3.239214in}}{\pgfqpoint{3.241198in}{3.233390in}}%
\pgfpathcurveto{\pgfqpoint{3.241198in}{3.227566in}}{\pgfqpoint{3.243512in}{3.221980in}}{\pgfqpoint{3.247630in}{3.217861in}}%
\pgfpathcurveto{\pgfqpoint{3.251748in}{3.213743in}}{\pgfqpoint{3.257334in}{3.211429in}}{\pgfqpoint{3.263158in}{3.211429in}}%
\pgfpathlineto{\pgfqpoint{3.263158in}{3.211429in}}%
\pgfpathclose%
\pgfusepath{stroke,fill}%
\end{pgfscope}%
\begin{pgfscope}%
\pgfpathrectangle{\pgfqpoint{1.414736in}{0.548769in}}{\pgfqpoint{3.497747in}{7.253186in}}%
\pgfusepath{clip}%
\pgfsetbuttcap%
\pgfsetroundjoin%
\definecolor{currentfill}{rgb}{0.793760,0.880678,0.120005}%
\pgfsetfillcolor{currentfill}%
\pgfsetlinewidth{1.003750pt}%
\definecolor{currentstroke}{rgb}{0.793760,0.880678,0.120005}%
\pgfsetstrokecolor{currentstroke}%
\pgfsetdash{}{0pt}%
\pgfpathmoveto{\pgfqpoint{3.821071in}{3.682415in}}%
\pgfpathcurveto{\pgfqpoint{3.826895in}{3.682415in}}{\pgfqpoint{3.832481in}{3.684729in}}{\pgfqpoint{3.836599in}{3.688847in}}%
\pgfpathcurveto{\pgfqpoint{3.840717in}{3.692966in}}{\pgfqpoint{3.843031in}{3.698552in}}{\pgfqpoint{3.843031in}{3.704376in}}%
\pgfpathcurveto{\pgfqpoint{3.843031in}{3.710200in}}{\pgfqpoint{3.840717in}{3.715786in}}{\pgfqpoint{3.836599in}{3.719904in}}%
\pgfpathcurveto{\pgfqpoint{3.832481in}{3.724022in}}{\pgfqpoint{3.826895in}{3.726336in}}{\pgfqpoint{3.821071in}{3.726336in}}%
\pgfpathcurveto{\pgfqpoint{3.815247in}{3.726336in}}{\pgfqpoint{3.809661in}{3.724022in}}{\pgfqpoint{3.805543in}{3.719904in}}%
\pgfpathcurveto{\pgfqpoint{3.801425in}{3.715786in}}{\pgfqpoint{3.799111in}{3.710200in}}{\pgfqpoint{3.799111in}{3.704376in}}%
\pgfpathcurveto{\pgfqpoint{3.799111in}{3.698552in}}{\pgfqpoint{3.801425in}{3.692966in}}{\pgfqpoint{3.805543in}{3.688847in}}%
\pgfpathcurveto{\pgfqpoint{3.809661in}{3.684729in}}{\pgfqpoint{3.815247in}{3.682415in}}{\pgfqpoint{3.821071in}{3.682415in}}%
\pgfpathlineto{\pgfqpoint{3.821071in}{3.682415in}}%
\pgfpathclose%
\pgfusepath{stroke,fill}%
\end{pgfscope}%
\begin{pgfscope}%
\pgfpathrectangle{\pgfqpoint{1.414736in}{0.548769in}}{\pgfqpoint{3.497747in}{7.253186in}}%
\pgfusepath{clip}%
\pgfsetbuttcap%
\pgfsetroundjoin%
\definecolor{currentfill}{rgb}{0.993248,0.906157,0.143936}%
\pgfsetfillcolor{currentfill}%
\pgfsetlinewidth{1.003750pt}%
\definecolor{currentstroke}{rgb}{0.993248,0.906157,0.143936}%
\pgfsetstrokecolor{currentstroke}%
\pgfsetdash{}{0pt}%
\pgfpathmoveto{\pgfqpoint{3.366188in}{4.153402in}}%
\pgfpathcurveto{\pgfqpoint{3.372012in}{4.153402in}}{\pgfqpoint{3.377599in}{4.155715in}}{\pgfqpoint{3.381717in}{4.159834in}}%
\pgfpathcurveto{\pgfqpoint{3.385835in}{4.163952in}}{\pgfqpoint{3.388149in}{4.169538in}}{\pgfqpoint{3.388149in}{4.175362in}}%
\pgfpathcurveto{\pgfqpoint{3.388149in}{4.181186in}}{\pgfqpoint{3.385835in}{4.186772in}}{\pgfqpoint{3.381717in}{4.190890in}}%
\pgfpathcurveto{\pgfqpoint{3.377599in}{4.195008in}}{\pgfqpoint{3.372012in}{4.197322in}}{\pgfqpoint{3.366188in}{4.197322in}}%
\pgfpathcurveto{\pgfqpoint{3.360365in}{4.197322in}}{\pgfqpoint{3.354778in}{4.195008in}}{\pgfqpoint{3.350660in}{4.190890in}}%
\pgfpathcurveto{\pgfqpoint{3.346542in}{4.186772in}}{\pgfqpoint{3.344228in}{4.181186in}}{\pgfqpoint{3.344228in}{4.175362in}}%
\pgfpathcurveto{\pgfqpoint{3.344228in}{4.169538in}}{\pgfqpoint{3.346542in}{4.163952in}}{\pgfqpoint{3.350660in}{4.159834in}}%
\pgfpathcurveto{\pgfqpoint{3.354778in}{4.155715in}}{\pgfqpoint{3.360365in}{4.153402in}}{\pgfqpoint{3.366188in}{4.153402in}}%
\pgfpathlineto{\pgfqpoint{3.366188in}{4.153402in}}%
\pgfpathclose%
\pgfusepath{stroke,fill}%
\end{pgfscope}%
\begin{pgfscope}%
\pgfpathrectangle{\pgfqpoint{1.414736in}{0.548769in}}{\pgfqpoint{3.497747in}{7.253186in}}%
\pgfusepath{clip}%
\pgfsetbuttcap%
\pgfsetroundjoin%
\definecolor{currentfill}{rgb}{0.824940,0.884720,0.106217}%
\pgfsetfillcolor{currentfill}%
\pgfsetlinewidth{1.003750pt}%
\definecolor{currentstroke}{rgb}{0.824940,0.884720,0.106217}%
\pgfsetstrokecolor{currentstroke}%
\pgfsetdash{}{0pt}%
\pgfpathmoveto{\pgfqpoint{3.727884in}{4.624388in}}%
\pgfpathcurveto{\pgfqpoint{3.733708in}{4.624388in}}{\pgfqpoint{3.739294in}{4.626702in}}{\pgfqpoint{3.743412in}{4.630820in}}%
\pgfpathcurveto{\pgfqpoint{3.747531in}{4.634938in}}{\pgfqpoint{3.749844in}{4.640524in}}{\pgfqpoint{3.749844in}{4.646348in}}%
\pgfpathcurveto{\pgfqpoint{3.749844in}{4.652172in}}{\pgfqpoint{3.747531in}{4.657758in}}{\pgfqpoint{3.743412in}{4.661876in}}%
\pgfpathcurveto{\pgfqpoint{3.739294in}{4.665994in}}{\pgfqpoint{3.733708in}{4.668308in}}{\pgfqpoint{3.727884in}{4.668308in}}%
\pgfpathcurveto{\pgfqpoint{3.722060in}{4.668308in}}{\pgfqpoint{3.716474in}{4.665994in}}{\pgfqpoint{3.712356in}{4.661876in}}%
\pgfpathcurveto{\pgfqpoint{3.708238in}{4.657758in}}{\pgfqpoint{3.705924in}{4.652172in}}{\pgfqpoint{3.705924in}{4.646348in}}%
\pgfpathcurveto{\pgfqpoint{3.705924in}{4.640524in}}{\pgfqpoint{3.708238in}{4.634938in}}{\pgfqpoint{3.712356in}{4.630820in}}%
\pgfpathcurveto{\pgfqpoint{3.716474in}{4.626702in}}{\pgfqpoint{3.722060in}{4.624388in}}{\pgfqpoint{3.727884in}{4.624388in}}%
\pgfpathlineto{\pgfqpoint{3.727884in}{4.624388in}}%
\pgfpathclose%
\pgfusepath{stroke,fill}%
\end{pgfscope}%
\begin{pgfscope}%
\pgfpathrectangle{\pgfqpoint{1.414736in}{0.548769in}}{\pgfqpoint{3.497747in}{7.253186in}}%
\pgfusepath{clip}%
\pgfsetbuttcap%
\pgfsetroundjoin%
\definecolor{currentfill}{rgb}{0.281446,0.084320,0.407414}%
\pgfsetfillcolor{currentfill}%
\pgfsetlinewidth{1.003750pt}%
\definecolor{currentstroke}{rgb}{0.281446,0.084320,0.407414}%
\pgfsetstrokecolor{currentstroke}%
\pgfsetdash{}{0pt}%
\pgfpathmoveto{\pgfqpoint{3.669991in}{5.095374in}}%
\pgfpathcurveto{\pgfqpoint{3.675815in}{5.095374in}}{\pgfqpoint{3.681401in}{5.097688in}}{\pgfqpoint{3.685520in}{5.101806in}}%
\pgfpathcurveto{\pgfqpoint{3.689638in}{5.105924in}}{\pgfqpoint{3.691952in}{5.111510in}}{\pgfqpoint{3.691952in}{5.117334in}}%
\pgfpathcurveto{\pgfqpoint{3.691952in}{5.123158in}}{\pgfqpoint{3.689638in}{5.128744in}}{\pgfqpoint{3.685520in}{5.132862in}}%
\pgfpathcurveto{\pgfqpoint{3.681401in}{5.136980in}}{\pgfqpoint{3.675815in}{5.139294in}}{\pgfqpoint{3.669991in}{5.139294in}}%
\pgfpathcurveto{\pgfqpoint{3.664167in}{5.139294in}}{\pgfqpoint{3.658581in}{5.136980in}}{\pgfqpoint{3.654463in}{5.132862in}}%
\pgfpathcurveto{\pgfqpoint{3.650345in}{5.128744in}}{\pgfqpoint{3.648031in}{5.123158in}}{\pgfqpoint{3.648031in}{5.117334in}}%
\pgfpathcurveto{\pgfqpoint{3.648031in}{5.111510in}}{\pgfqpoint{3.650345in}{5.105924in}}{\pgfqpoint{3.654463in}{5.101806in}}%
\pgfpathcurveto{\pgfqpoint{3.658581in}{5.097688in}}{\pgfqpoint{3.664167in}{5.095374in}}{\pgfqpoint{3.669991in}{5.095374in}}%
\pgfpathlineto{\pgfqpoint{3.669991in}{5.095374in}}%
\pgfpathclose%
\pgfusepath{stroke,fill}%
\end{pgfscope}%
\begin{pgfscope}%
\pgfpathrectangle{\pgfqpoint{1.414736in}{0.548769in}}{\pgfqpoint{3.497747in}{7.253186in}}%
\pgfusepath{clip}%
\pgfsetbuttcap%
\pgfsetroundjoin%
\definecolor{currentfill}{rgb}{0.252194,0.269783,0.531579}%
\pgfsetfillcolor{currentfill}%
\pgfsetlinewidth{1.003750pt}%
\definecolor{currentstroke}{rgb}{0.252194,0.269783,0.531579}%
\pgfsetstrokecolor{currentstroke}%
\pgfsetdash{}{0pt}%
\pgfpathmoveto{\pgfqpoint{3.946046in}{5.566360in}}%
\pgfpathcurveto{\pgfqpoint{3.951870in}{5.566360in}}{\pgfqpoint{3.957456in}{5.568674in}}{\pgfqpoint{3.961575in}{5.572792in}}%
\pgfpathcurveto{\pgfqpoint{3.965693in}{5.576910in}}{\pgfqpoint{3.968007in}{5.582496in}}{\pgfqpoint{3.968007in}{5.588320in}}%
\pgfpathcurveto{\pgfqpoint{3.968007in}{5.594144in}}{\pgfqpoint{3.965693in}{5.599730in}}{\pgfqpoint{3.961575in}{5.603848in}}%
\pgfpathcurveto{\pgfqpoint{3.957456in}{5.607967in}}{\pgfqpoint{3.951870in}{5.610280in}}{\pgfqpoint{3.946046in}{5.610280in}}%
\pgfpathcurveto{\pgfqpoint{3.940222in}{5.610280in}}{\pgfqpoint{3.934636in}{5.607967in}}{\pgfqpoint{3.930518in}{5.603848in}}%
\pgfpathcurveto{\pgfqpoint{3.926400in}{5.599730in}}{\pgfqpoint{3.924086in}{5.594144in}}{\pgfqpoint{3.924086in}{5.588320in}}%
\pgfpathcurveto{\pgfqpoint{3.924086in}{5.582496in}}{\pgfqpoint{3.926400in}{5.576910in}}{\pgfqpoint{3.930518in}{5.572792in}}%
\pgfpathcurveto{\pgfqpoint{3.934636in}{5.568674in}}{\pgfqpoint{3.940222in}{5.566360in}}{\pgfqpoint{3.946046in}{5.566360in}}%
\pgfpathlineto{\pgfqpoint{3.946046in}{5.566360in}}%
\pgfpathclose%
\pgfusepath{stroke,fill}%
\end{pgfscope}%
\begin{pgfscope}%
\pgfpathrectangle{\pgfqpoint{1.414736in}{0.548769in}}{\pgfqpoint{3.497747in}{7.253186in}}%
\pgfusepath{clip}%
\pgfsetbuttcap%
\pgfsetroundjoin%
\definecolor{currentfill}{rgb}{0.125394,0.574318,0.549086}%
\pgfsetfillcolor{currentfill}%
\pgfsetlinewidth{1.003750pt}%
\definecolor{currentstroke}{rgb}{0.125394,0.574318,0.549086}%
\pgfsetstrokecolor{currentstroke}%
\pgfsetdash{}{0pt}%
\pgfpathmoveto{\pgfqpoint{3.726604in}{6.037346in}}%
\pgfpathcurveto{\pgfqpoint{3.732427in}{6.037346in}}{\pgfqpoint{3.738014in}{6.039660in}}{\pgfqpoint{3.742132in}{6.043778in}}%
\pgfpathcurveto{\pgfqpoint{3.746250in}{6.047896in}}{\pgfqpoint{3.748564in}{6.053482in}}{\pgfqpoint{3.748564in}{6.059306in}}%
\pgfpathcurveto{\pgfqpoint{3.748564in}{6.065130in}}{\pgfqpoint{3.746250in}{6.070716in}}{\pgfqpoint{3.742132in}{6.074835in}}%
\pgfpathcurveto{\pgfqpoint{3.738014in}{6.078953in}}{\pgfqpoint{3.732427in}{6.081267in}}{\pgfqpoint{3.726604in}{6.081267in}}%
\pgfpathcurveto{\pgfqpoint{3.720780in}{6.081267in}}{\pgfqpoint{3.715193in}{6.078953in}}{\pgfqpoint{3.711075in}{6.074835in}}%
\pgfpathcurveto{\pgfqpoint{3.706957in}{6.070716in}}{\pgfqpoint{3.704643in}{6.065130in}}{\pgfqpoint{3.704643in}{6.059306in}}%
\pgfpathcurveto{\pgfqpoint{3.704643in}{6.053482in}}{\pgfqpoint{3.706957in}{6.047896in}}{\pgfqpoint{3.711075in}{6.043778in}}%
\pgfpathcurveto{\pgfqpoint{3.715193in}{6.039660in}}{\pgfqpoint{3.720780in}{6.037346in}}{\pgfqpoint{3.726604in}{6.037346in}}%
\pgfpathlineto{\pgfqpoint{3.726604in}{6.037346in}}%
\pgfpathclose%
\pgfusepath{stroke,fill}%
\end{pgfscope}%
\begin{pgfscope}%
\pgfpathrectangle{\pgfqpoint{1.414736in}{0.548769in}}{\pgfqpoint{3.497747in}{7.253186in}}%
\pgfusepath{clip}%
\pgfsetbuttcap%
\pgfsetroundjoin%
\definecolor{currentfill}{rgb}{0.824940,0.884720,0.106217}%
\pgfsetfillcolor{currentfill}%
\pgfsetlinewidth{1.003750pt}%
\definecolor{currentstroke}{rgb}{0.824940,0.884720,0.106217}%
\pgfsetstrokecolor{currentstroke}%
\pgfsetdash{}{0pt}%
\pgfpathmoveto{\pgfqpoint{3.489334in}{6.508332in}}%
\pgfpathcurveto{\pgfqpoint{3.495158in}{6.508332in}}{\pgfqpoint{3.500744in}{6.510646in}}{\pgfqpoint{3.504862in}{6.514764in}}%
\pgfpathcurveto{\pgfqpoint{3.508980in}{6.518882in}}{\pgfqpoint{3.511294in}{6.524468in}}{\pgfqpoint{3.511294in}{6.530292in}}%
\pgfpathcurveto{\pgfqpoint{3.511294in}{6.536116in}}{\pgfqpoint{3.508980in}{6.541702in}}{\pgfqpoint{3.504862in}{6.545821in}}%
\pgfpathcurveto{\pgfqpoint{3.500744in}{6.549939in}}{\pgfqpoint{3.495158in}{6.552253in}}{\pgfqpoint{3.489334in}{6.552253in}}%
\pgfpathcurveto{\pgfqpoint{3.483510in}{6.552253in}}{\pgfqpoint{3.477924in}{6.549939in}}{\pgfqpoint{3.473806in}{6.545821in}}%
\pgfpathcurveto{\pgfqpoint{3.469688in}{6.541702in}}{\pgfqpoint{3.467374in}{6.536116in}}{\pgfqpoint{3.467374in}{6.530292in}}%
\pgfpathcurveto{\pgfqpoint{3.467374in}{6.524468in}}{\pgfqpoint{3.469688in}{6.518882in}}{\pgfqpoint{3.473806in}{6.514764in}}%
\pgfpathcurveto{\pgfqpoint{3.477924in}{6.510646in}}{\pgfqpoint{3.483510in}{6.508332in}}{\pgfqpoint{3.489334in}{6.508332in}}%
\pgfpathlineto{\pgfqpoint{3.489334in}{6.508332in}}%
\pgfpathclose%
\pgfusepath{stroke,fill}%
\end{pgfscope}%
\begin{pgfscope}%
\pgfpathrectangle{\pgfqpoint{1.414736in}{0.548769in}}{\pgfqpoint{3.497747in}{7.253186in}}%
\pgfusepath{clip}%
\pgfsetbuttcap%
\pgfsetroundjoin%
\definecolor{currentfill}{rgb}{0.993248,0.906157,0.143936}%
\pgfsetfillcolor{currentfill}%
\pgfsetlinewidth{1.003750pt}%
\definecolor{currentstroke}{rgb}{0.993248,0.906157,0.143936}%
\pgfsetstrokecolor{currentstroke}%
\pgfsetdash{}{0pt}%
\pgfpathmoveto{\pgfqpoint{2.871177in}{6.979318in}}%
\pgfpathcurveto{\pgfqpoint{2.877001in}{6.979318in}}{\pgfqpoint{2.882587in}{6.981632in}}{\pgfqpoint{2.886705in}{6.985750in}}%
\pgfpathcurveto{\pgfqpoint{2.890824in}{6.989868in}}{\pgfqpoint{2.893137in}{6.995455in}}{\pgfqpoint{2.893137in}{7.001278in}}%
\pgfpathcurveto{\pgfqpoint{2.893137in}{7.007102in}}{\pgfqpoint{2.890824in}{7.012689in}}{\pgfqpoint{2.886705in}{7.016807in}}%
\pgfpathcurveto{\pgfqpoint{2.882587in}{7.020925in}}{\pgfqpoint{2.877001in}{7.023239in}}{\pgfqpoint{2.871177in}{7.023239in}}%
\pgfpathcurveto{\pgfqpoint{2.865353in}{7.023239in}}{\pgfqpoint{2.859767in}{7.020925in}}{\pgfqpoint{2.855649in}{7.016807in}}%
\pgfpathcurveto{\pgfqpoint{2.851531in}{7.012689in}}{\pgfqpoint{2.849217in}{7.007102in}}{\pgfqpoint{2.849217in}{7.001278in}}%
\pgfpathcurveto{\pgfqpoint{2.849217in}{6.995455in}}{\pgfqpoint{2.851531in}{6.989868in}}{\pgfqpoint{2.855649in}{6.985750in}}%
\pgfpathcurveto{\pgfqpoint{2.859767in}{6.981632in}}{\pgfqpoint{2.865353in}{6.979318in}}{\pgfqpoint{2.871177in}{6.979318in}}%
\pgfpathlineto{\pgfqpoint{2.871177in}{6.979318in}}%
\pgfpathclose%
\pgfusepath{stroke,fill}%
\end{pgfscope}%
\begin{pgfscope}%
\pgfpathrectangle{\pgfqpoint{1.414736in}{0.548769in}}{\pgfqpoint{3.497747in}{7.253186in}}%
\pgfusepath{clip}%
\pgfsetbuttcap%
\pgfsetroundjoin%
\definecolor{currentfill}{rgb}{0.845561,0.887322,0.099702}%
\pgfsetfillcolor{currentfill}%
\pgfsetlinewidth{1.003750pt}%
\definecolor{currentstroke}{rgb}{0.845561,0.887322,0.099702}%
\pgfsetstrokecolor{currentstroke}%
\pgfsetdash{}{0pt}%
\pgfpathmoveto{\pgfqpoint{4.147917in}{7.450304in}}%
\pgfpathcurveto{\pgfqpoint{4.153741in}{7.450304in}}{\pgfqpoint{4.159327in}{7.452618in}}{\pgfqpoint{4.163446in}{7.456736in}}%
\pgfpathcurveto{\pgfqpoint{4.167564in}{7.460854in}}{\pgfqpoint{4.169878in}{7.466441in}}{\pgfqpoint{4.169878in}{7.472265in}}%
\pgfpathcurveto{\pgfqpoint{4.169878in}{7.478089in}}{\pgfqpoint{4.167564in}{7.483675in}}{\pgfqpoint{4.163446in}{7.487793in}}%
\pgfpathcurveto{\pgfqpoint{4.159327in}{7.491911in}}{\pgfqpoint{4.153741in}{7.494225in}}{\pgfqpoint{4.147917in}{7.494225in}}%
\pgfpathcurveto{\pgfqpoint{4.142093in}{7.494225in}}{\pgfqpoint{4.136507in}{7.491911in}}{\pgfqpoint{4.132389in}{7.487793in}}%
\pgfpathcurveto{\pgfqpoint{4.128271in}{7.483675in}}{\pgfqpoint{4.125957in}{7.478089in}}{\pgfqpoint{4.125957in}{7.472265in}}%
\pgfpathcurveto{\pgfqpoint{4.125957in}{7.466441in}}{\pgfqpoint{4.128271in}{7.460854in}}{\pgfqpoint{4.132389in}{7.456736in}}%
\pgfpathcurveto{\pgfqpoint{4.136507in}{7.452618in}}{\pgfqpoint{4.142093in}{7.450304in}}{\pgfqpoint{4.147917in}{7.450304in}}%
\pgfpathlineto{\pgfqpoint{4.147917in}{7.450304in}}%
\pgfpathclose%
\pgfusepath{stroke,fill}%
\end{pgfscope}%
\begin{pgfscope}%
\pgfpathrectangle{\pgfqpoint{1.414736in}{0.548769in}}{\pgfqpoint{3.497747in}{7.253186in}}%
\pgfusepath{clip}%
\pgfsetbuttcap%
\pgfsetroundjoin%
\definecolor{currentfill}{rgb}{0.855810,0.888601,0.097452}%
\pgfsetfillcolor{currentfill}%
\pgfsetlinewidth{1.003750pt}%
\definecolor{currentstroke}{rgb}{0.855810,0.888601,0.097452}%
\pgfsetstrokecolor{currentstroke}%
\pgfsetdash{}{0pt}%
\pgfpathmoveto{\pgfqpoint{3.650491in}{0.856499in}}%
\pgfpathcurveto{\pgfqpoint{3.656315in}{0.856499in}}{\pgfqpoint{3.661901in}{0.858813in}}{\pgfqpoint{3.666019in}{0.862931in}}%
\pgfpathcurveto{\pgfqpoint{3.670137in}{0.867049in}}{\pgfqpoint{3.672451in}{0.872635in}}{\pgfqpoint{3.672451in}{0.878459in}}%
\pgfpathcurveto{\pgfqpoint{3.672451in}{0.884283in}}{\pgfqpoint{3.670137in}{0.889869in}}{\pgfqpoint{3.666019in}{0.893987in}}%
\pgfpathcurveto{\pgfqpoint{3.661901in}{0.898105in}}{\pgfqpoint{3.656315in}{0.900419in}}{\pgfqpoint{3.650491in}{0.900419in}}%
\pgfpathcurveto{\pgfqpoint{3.644667in}{0.900419in}}{\pgfqpoint{3.639081in}{0.898105in}}{\pgfqpoint{3.634963in}{0.893987in}}%
\pgfpathcurveto{\pgfqpoint{3.630845in}{0.889869in}}{\pgfqpoint{3.628531in}{0.884283in}}{\pgfqpoint{3.628531in}{0.878459in}}%
\pgfpathcurveto{\pgfqpoint{3.628531in}{0.872635in}}{\pgfqpoint{3.630845in}{0.867049in}}{\pgfqpoint{3.634963in}{0.862931in}}%
\pgfpathcurveto{\pgfqpoint{3.639081in}{0.858813in}}{\pgfqpoint{3.644667in}{0.856499in}}{\pgfqpoint{3.650491in}{0.856499in}}%
\pgfpathlineto{\pgfqpoint{3.650491in}{0.856499in}}%
\pgfpathclose%
\pgfusepath{stroke,fill}%
\end{pgfscope}%
\begin{pgfscope}%
\pgfpathrectangle{\pgfqpoint{1.414736in}{0.548769in}}{\pgfqpoint{3.497747in}{7.253186in}}%
\pgfusepath{clip}%
\pgfsetbuttcap%
\pgfsetroundjoin%
\definecolor{currentfill}{rgb}{0.585678,0.846661,0.249897}%
\pgfsetfillcolor{currentfill}%
\pgfsetlinewidth{1.003750pt}%
\definecolor{currentstroke}{rgb}{0.585678,0.846661,0.249897}%
\pgfsetstrokecolor{currentstroke}%
\pgfsetdash{}{0pt}%
\pgfpathmoveto{\pgfqpoint{3.446668in}{1.327485in}}%
\pgfpathcurveto{\pgfqpoint{3.452492in}{1.327485in}}{\pgfqpoint{3.458078in}{1.329799in}}{\pgfqpoint{3.462196in}{1.333917in}}%
\pgfpathcurveto{\pgfqpoint{3.466315in}{1.338035in}}{\pgfqpoint{3.468628in}{1.343621in}}{\pgfqpoint{3.468628in}{1.349445in}}%
\pgfpathcurveto{\pgfqpoint{3.468628in}{1.355269in}}{\pgfqpoint{3.466315in}{1.360855in}}{\pgfqpoint{3.462196in}{1.364973in}}%
\pgfpathcurveto{\pgfqpoint{3.458078in}{1.369092in}}{\pgfqpoint{3.452492in}{1.371405in}}{\pgfqpoint{3.446668in}{1.371405in}}%
\pgfpathcurveto{\pgfqpoint{3.440844in}{1.371405in}}{\pgfqpoint{3.435258in}{1.369092in}}{\pgfqpoint{3.431140in}{1.364973in}}%
\pgfpathcurveto{\pgfqpoint{3.427022in}{1.360855in}}{\pgfqpoint{3.424708in}{1.355269in}}{\pgfqpoint{3.424708in}{1.349445in}}%
\pgfpathcurveto{\pgfqpoint{3.424708in}{1.343621in}}{\pgfqpoint{3.427022in}{1.338035in}}{\pgfqpoint{3.431140in}{1.333917in}}%
\pgfpathcurveto{\pgfqpoint{3.435258in}{1.329799in}}{\pgfqpoint{3.440844in}{1.327485in}}{\pgfqpoint{3.446668in}{1.327485in}}%
\pgfpathlineto{\pgfqpoint{3.446668in}{1.327485in}}%
\pgfpathclose%
\pgfusepath{stroke,fill}%
\end{pgfscope}%
\begin{pgfscope}%
\pgfpathrectangle{\pgfqpoint{1.414736in}{0.548769in}}{\pgfqpoint{3.497747in}{7.253186in}}%
\pgfusepath{clip}%
\pgfsetbuttcap%
\pgfsetroundjoin%
\definecolor{currentfill}{rgb}{0.804182,0.882046,0.114965}%
\pgfsetfillcolor{currentfill}%
\pgfsetlinewidth{1.003750pt}%
\definecolor{currentstroke}{rgb}{0.804182,0.882046,0.114965}%
\pgfsetstrokecolor{currentstroke}%
\pgfsetdash{}{0pt}%
\pgfpathmoveto{\pgfqpoint{3.923788in}{1.798471in}}%
\pgfpathcurveto{\pgfqpoint{3.929612in}{1.798471in}}{\pgfqpoint{3.935198in}{1.800785in}}{\pgfqpoint{3.939316in}{1.804903in}}%
\pgfpathcurveto{\pgfqpoint{3.943434in}{1.809021in}}{\pgfqpoint{3.945748in}{1.814607in}}{\pgfqpoint{3.945748in}{1.820431in}}%
\pgfpathcurveto{\pgfqpoint{3.945748in}{1.826255in}}{\pgfqpoint{3.943434in}{1.831841in}}{\pgfqpoint{3.939316in}{1.835960in}}%
\pgfpathcurveto{\pgfqpoint{3.935198in}{1.840078in}}{\pgfqpoint{3.929612in}{1.842392in}}{\pgfqpoint{3.923788in}{1.842392in}}%
\pgfpathcurveto{\pgfqpoint{3.917964in}{1.842392in}}{\pgfqpoint{3.912378in}{1.840078in}}{\pgfqpoint{3.908260in}{1.835960in}}%
\pgfpathcurveto{\pgfqpoint{3.904142in}{1.831841in}}{\pgfqpoint{3.901828in}{1.826255in}}{\pgfqpoint{3.901828in}{1.820431in}}%
\pgfpathcurveto{\pgfqpoint{3.901828in}{1.814607in}}{\pgfqpoint{3.904142in}{1.809021in}}{\pgfqpoint{3.908260in}{1.804903in}}%
\pgfpathcurveto{\pgfqpoint{3.912378in}{1.800785in}}{\pgfqpoint{3.917964in}{1.798471in}}{\pgfqpoint{3.923788in}{1.798471in}}%
\pgfpathlineto{\pgfqpoint{3.923788in}{1.798471in}}%
\pgfpathclose%
\pgfusepath{stroke,fill}%
\end{pgfscope}%
\begin{pgfscope}%
\pgfpathrectangle{\pgfqpoint{1.414736in}{0.548769in}}{\pgfqpoint{3.497747in}{7.253186in}}%
\pgfusepath{clip}%
\pgfsetbuttcap%
\pgfsetroundjoin%
\definecolor{currentfill}{rgb}{0.993248,0.906157,0.143936}%
\pgfsetfillcolor{currentfill}%
\pgfsetlinewidth{1.003750pt}%
\definecolor{currentstroke}{rgb}{0.993248,0.906157,0.143936}%
\pgfsetstrokecolor{currentstroke}%
\pgfsetdash{}{0pt}%
\pgfpathmoveto{\pgfqpoint{3.692717in}{2.269457in}}%
\pgfpathcurveto{\pgfqpoint{3.698541in}{2.269457in}}{\pgfqpoint{3.704127in}{2.271771in}}{\pgfqpoint{3.708246in}{2.275889in}}%
\pgfpathcurveto{\pgfqpoint{3.712364in}{2.280007in}}{\pgfqpoint{3.714678in}{2.285593in}}{\pgfqpoint{3.714678in}{2.291417in}}%
\pgfpathcurveto{\pgfqpoint{3.714678in}{2.297241in}}{\pgfqpoint{3.712364in}{2.302828in}}{\pgfqpoint{3.708246in}{2.306946in}}%
\pgfpathcurveto{\pgfqpoint{3.704127in}{2.311064in}}{\pgfqpoint{3.698541in}{2.313378in}}{\pgfqpoint{3.692717in}{2.313378in}}%
\pgfpathcurveto{\pgfqpoint{3.686893in}{2.313378in}}{\pgfqpoint{3.681307in}{2.311064in}}{\pgfqpoint{3.677189in}{2.306946in}}%
\pgfpathcurveto{\pgfqpoint{3.673071in}{2.302828in}}{\pgfqpoint{3.670757in}{2.297241in}}{\pgfqpoint{3.670757in}{2.291417in}}%
\pgfpathcurveto{\pgfqpoint{3.670757in}{2.285593in}}{\pgfqpoint{3.673071in}{2.280007in}}{\pgfqpoint{3.677189in}{2.275889in}}%
\pgfpathcurveto{\pgfqpoint{3.681307in}{2.271771in}}{\pgfqpoint{3.686893in}{2.269457in}}{\pgfqpoint{3.692717in}{2.269457in}}%
\pgfpathlineto{\pgfqpoint{3.692717in}{2.269457in}}%
\pgfpathclose%
\pgfusepath{stroke,fill}%
\end{pgfscope}%
\begin{pgfscope}%
\pgfpathrectangle{\pgfqpoint{1.414736in}{0.548769in}}{\pgfqpoint{3.497747in}{7.253186in}}%
\pgfusepath{clip}%
\pgfsetbuttcap%
\pgfsetroundjoin%
\definecolor{currentfill}{rgb}{0.296479,0.761561,0.424223}%
\pgfsetfillcolor{currentfill}%
\pgfsetlinewidth{1.003750pt}%
\definecolor{currentstroke}{rgb}{0.296479,0.761561,0.424223}%
\pgfsetstrokecolor{currentstroke}%
\pgfsetdash{}{0pt}%
\pgfpathmoveto{\pgfqpoint{3.728648in}{2.740443in}}%
\pgfpathcurveto{\pgfqpoint{3.734472in}{2.740443in}}{\pgfqpoint{3.740058in}{2.742757in}}{\pgfqpoint{3.744176in}{2.746875in}}%
\pgfpathcurveto{\pgfqpoint{3.748295in}{2.750993in}}{\pgfqpoint{3.750608in}{2.756580in}}{\pgfqpoint{3.750608in}{2.762404in}}%
\pgfpathcurveto{\pgfqpoint{3.750608in}{2.768227in}}{\pgfqpoint{3.748295in}{2.773814in}}{\pgfqpoint{3.744176in}{2.777932in}}%
\pgfpathcurveto{\pgfqpoint{3.740058in}{2.782050in}}{\pgfqpoint{3.734472in}{2.784364in}}{\pgfqpoint{3.728648in}{2.784364in}}%
\pgfpathcurveto{\pgfqpoint{3.722824in}{2.784364in}}{\pgfqpoint{3.717238in}{2.782050in}}{\pgfqpoint{3.713120in}{2.777932in}}%
\pgfpathcurveto{\pgfqpoint{3.709002in}{2.773814in}}{\pgfqpoint{3.706688in}{2.768227in}}{\pgfqpoint{3.706688in}{2.762404in}}%
\pgfpathcurveto{\pgfqpoint{3.706688in}{2.756580in}}{\pgfqpoint{3.709002in}{2.750993in}}{\pgfqpoint{3.713120in}{2.746875in}}%
\pgfpathcurveto{\pgfqpoint{3.717238in}{2.742757in}}{\pgfqpoint{3.722824in}{2.740443in}}{\pgfqpoint{3.728648in}{2.740443in}}%
\pgfpathlineto{\pgfqpoint{3.728648in}{2.740443in}}%
\pgfpathclose%
\pgfusepath{stroke,fill}%
\end{pgfscope}%
\begin{pgfscope}%
\pgfpathrectangle{\pgfqpoint{1.414736in}{0.548769in}}{\pgfqpoint{3.497747in}{7.253186in}}%
\pgfusepath{clip}%
\pgfsetbuttcap%
\pgfsetroundjoin%
\definecolor{currentfill}{rgb}{0.720391,0.870350,0.162603}%
\pgfsetfillcolor{currentfill}%
\pgfsetlinewidth{1.003750pt}%
\definecolor{currentstroke}{rgb}{0.720391,0.870350,0.162603}%
\pgfsetstrokecolor{currentstroke}%
\pgfsetdash{}{0pt}%
\pgfpathmoveto{\pgfqpoint{3.102304in}{3.211429in}}%
\pgfpathcurveto{\pgfqpoint{3.108128in}{3.211429in}}{\pgfqpoint{3.113714in}{3.213743in}}{\pgfqpoint{3.117833in}{3.217861in}}%
\pgfpathcurveto{\pgfqpoint{3.121951in}{3.221980in}}{\pgfqpoint{3.124265in}{3.227566in}}{\pgfqpoint{3.124265in}{3.233390in}}%
\pgfpathcurveto{\pgfqpoint{3.124265in}{3.239214in}}{\pgfqpoint{3.121951in}{3.244800in}}{\pgfqpoint{3.117833in}{3.248918in}}%
\pgfpathcurveto{\pgfqpoint{3.113714in}{3.253036in}}{\pgfqpoint{3.108128in}{3.255350in}}{\pgfqpoint{3.102304in}{3.255350in}}%
\pgfpathcurveto{\pgfqpoint{3.096480in}{3.255350in}}{\pgfqpoint{3.090894in}{3.253036in}}{\pgfqpoint{3.086776in}{3.248918in}}%
\pgfpathcurveto{\pgfqpoint{3.082658in}{3.244800in}}{\pgfqpoint{3.080344in}{3.239214in}}{\pgfqpoint{3.080344in}{3.233390in}}%
\pgfpathcurveto{\pgfqpoint{3.080344in}{3.227566in}}{\pgfqpoint{3.082658in}{3.221980in}}{\pgfqpoint{3.086776in}{3.217861in}}%
\pgfpathcurveto{\pgfqpoint{3.090894in}{3.213743in}}{\pgfqpoint{3.096480in}{3.211429in}}{\pgfqpoint{3.102304in}{3.211429in}}%
\pgfpathlineto{\pgfqpoint{3.102304in}{3.211429in}}%
\pgfpathclose%
\pgfusepath{stroke,fill}%
\end{pgfscope}%
\begin{pgfscope}%
\pgfpathrectangle{\pgfqpoint{1.414736in}{0.548769in}}{\pgfqpoint{3.497747in}{7.253186in}}%
\pgfusepath{clip}%
\pgfsetbuttcap%
\pgfsetroundjoin%
\definecolor{currentfill}{rgb}{0.964894,0.902323,0.123941}%
\pgfsetfillcolor{currentfill}%
\pgfsetlinewidth{1.003750pt}%
\definecolor{currentstroke}{rgb}{0.964894,0.902323,0.123941}%
\pgfsetstrokecolor{currentstroke}%
\pgfsetdash{}{0pt}%
\pgfpathmoveto{\pgfqpoint{3.263244in}{3.682415in}}%
\pgfpathcurveto{\pgfqpoint{3.269068in}{3.682415in}}{\pgfqpoint{3.274654in}{3.684729in}}{\pgfqpoint{3.278772in}{3.688847in}}%
\pgfpathcurveto{\pgfqpoint{3.282890in}{3.692966in}}{\pgfqpoint{3.285204in}{3.698552in}}{\pgfqpoint{3.285204in}{3.704376in}}%
\pgfpathcurveto{\pgfqpoint{3.285204in}{3.710200in}}{\pgfqpoint{3.282890in}{3.715786in}}{\pgfqpoint{3.278772in}{3.719904in}}%
\pgfpathcurveto{\pgfqpoint{3.274654in}{3.724022in}}{\pgfqpoint{3.269068in}{3.726336in}}{\pgfqpoint{3.263244in}{3.726336in}}%
\pgfpathcurveto{\pgfqpoint{3.257420in}{3.726336in}}{\pgfqpoint{3.251834in}{3.724022in}}{\pgfqpoint{3.247716in}{3.719904in}}%
\pgfpathcurveto{\pgfqpoint{3.243597in}{3.715786in}}{\pgfqpoint{3.241284in}{3.710200in}}{\pgfqpoint{3.241284in}{3.704376in}}%
\pgfpathcurveto{\pgfqpoint{3.241284in}{3.698552in}}{\pgfqpoint{3.243597in}{3.692966in}}{\pgfqpoint{3.247716in}{3.688847in}}%
\pgfpathcurveto{\pgfqpoint{3.251834in}{3.684729in}}{\pgfqpoint{3.257420in}{3.682415in}}{\pgfqpoint{3.263244in}{3.682415in}}%
\pgfpathlineto{\pgfqpoint{3.263244in}{3.682415in}}%
\pgfpathclose%
\pgfusepath{stroke,fill}%
\end{pgfscope}%
\begin{pgfscope}%
\pgfpathrectangle{\pgfqpoint{1.414736in}{0.548769in}}{\pgfqpoint{3.497747in}{7.253186in}}%
\pgfusepath{clip}%
\pgfsetbuttcap%
\pgfsetroundjoin%
\definecolor{currentfill}{rgb}{0.974417,0.903590,0.130215}%
\pgfsetfillcolor{currentfill}%
\pgfsetlinewidth{1.003750pt}%
\definecolor{currentstroke}{rgb}{0.974417,0.903590,0.130215}%
\pgfsetstrokecolor{currentstroke}%
\pgfsetdash{}{0pt}%
\pgfpathmoveto{\pgfqpoint{3.825694in}{4.153402in}}%
\pgfpathcurveto{\pgfqpoint{3.831518in}{4.153402in}}{\pgfqpoint{3.837105in}{4.155715in}}{\pgfqpoint{3.841223in}{4.159834in}}%
\pgfpathcurveto{\pgfqpoint{3.845341in}{4.163952in}}{\pgfqpoint{3.847655in}{4.169538in}}{\pgfqpoint{3.847655in}{4.175362in}}%
\pgfpathcurveto{\pgfqpoint{3.847655in}{4.181186in}}{\pgfqpoint{3.845341in}{4.186772in}}{\pgfqpoint{3.841223in}{4.190890in}}%
\pgfpathcurveto{\pgfqpoint{3.837105in}{4.195008in}}{\pgfqpoint{3.831518in}{4.197322in}}{\pgfqpoint{3.825694in}{4.197322in}}%
\pgfpathcurveto{\pgfqpoint{3.819871in}{4.197322in}}{\pgfqpoint{3.814284in}{4.195008in}}{\pgfqpoint{3.810166in}{4.190890in}}%
\pgfpathcurveto{\pgfqpoint{3.806048in}{4.186772in}}{\pgfqpoint{3.803734in}{4.181186in}}{\pgfqpoint{3.803734in}{4.175362in}}%
\pgfpathcurveto{\pgfqpoint{3.803734in}{4.169538in}}{\pgfqpoint{3.806048in}{4.163952in}}{\pgfqpoint{3.810166in}{4.159834in}}%
\pgfpathcurveto{\pgfqpoint{3.814284in}{4.155715in}}{\pgfqpoint{3.819871in}{4.153402in}}{\pgfqpoint{3.825694in}{4.153402in}}%
\pgfpathlineto{\pgfqpoint{3.825694in}{4.153402in}}%
\pgfpathclose%
\pgfusepath{stroke,fill}%
\end{pgfscope}%
\begin{pgfscope}%
\pgfpathrectangle{\pgfqpoint{1.414736in}{0.548769in}}{\pgfqpoint{3.497747in}{7.253186in}}%
\pgfusepath{clip}%
\pgfsetbuttcap%
\pgfsetroundjoin%
\definecolor{currentfill}{rgb}{0.647257,0.858400,0.209861}%
\pgfsetfillcolor{currentfill}%
\pgfsetlinewidth{1.003750pt}%
\definecolor{currentstroke}{rgb}{0.647257,0.858400,0.209861}%
\pgfsetstrokecolor{currentstroke}%
\pgfsetdash{}{0pt}%
\pgfpathmoveto{\pgfqpoint{3.300773in}{4.624388in}}%
\pgfpathcurveto{\pgfqpoint{3.306597in}{4.624388in}}{\pgfqpoint{3.312183in}{4.626702in}}{\pgfqpoint{3.316301in}{4.630820in}}%
\pgfpathcurveto{\pgfqpoint{3.320419in}{4.634938in}}{\pgfqpoint{3.322733in}{4.640524in}}{\pgfqpoint{3.322733in}{4.646348in}}%
\pgfpathcurveto{\pgfqpoint{3.322733in}{4.652172in}}{\pgfqpoint{3.320419in}{4.657758in}}{\pgfqpoint{3.316301in}{4.661876in}}%
\pgfpathcurveto{\pgfqpoint{3.312183in}{4.665994in}}{\pgfqpoint{3.306597in}{4.668308in}}{\pgfqpoint{3.300773in}{4.668308in}}%
\pgfpathcurveto{\pgfqpoint{3.294949in}{4.668308in}}{\pgfqpoint{3.289363in}{4.665994in}}{\pgfqpoint{3.285244in}{4.661876in}}%
\pgfpathcurveto{\pgfqpoint{3.281126in}{4.657758in}}{\pgfqpoint{3.278812in}{4.652172in}}{\pgfqpoint{3.278812in}{4.646348in}}%
\pgfpathcurveto{\pgfqpoint{3.278812in}{4.640524in}}{\pgfqpoint{3.281126in}{4.634938in}}{\pgfqpoint{3.285244in}{4.630820in}}%
\pgfpathcurveto{\pgfqpoint{3.289363in}{4.626702in}}{\pgfqpoint{3.294949in}{4.624388in}}{\pgfqpoint{3.300773in}{4.624388in}}%
\pgfpathlineto{\pgfqpoint{3.300773in}{4.624388in}}%
\pgfpathclose%
\pgfusepath{stroke,fill}%
\end{pgfscope}%
\begin{pgfscope}%
\pgfpathrectangle{\pgfqpoint{1.414736in}{0.548769in}}{\pgfqpoint{3.497747in}{7.253186in}}%
\pgfusepath{clip}%
\pgfsetbuttcap%
\pgfsetroundjoin%
\definecolor{currentfill}{rgb}{0.201239,0.383670,0.554294}%
\pgfsetfillcolor{currentfill}%
\pgfsetlinewidth{1.003750pt}%
\definecolor{currentstroke}{rgb}{0.201239,0.383670,0.554294}%
\pgfsetstrokecolor{currentstroke}%
\pgfsetdash{}{0pt}%
\pgfpathmoveto{\pgfqpoint{3.223930in}{5.095374in}}%
\pgfpathcurveto{\pgfqpoint{3.229754in}{5.095374in}}{\pgfqpoint{3.235340in}{5.097688in}}{\pgfqpoint{3.239458in}{5.101806in}}%
\pgfpathcurveto{\pgfqpoint{3.243576in}{5.105924in}}{\pgfqpoint{3.245890in}{5.111510in}}{\pgfqpoint{3.245890in}{5.117334in}}%
\pgfpathcurveto{\pgfqpoint{3.245890in}{5.123158in}}{\pgfqpoint{3.243576in}{5.128744in}}{\pgfqpoint{3.239458in}{5.132862in}}%
\pgfpathcurveto{\pgfqpoint{3.235340in}{5.136980in}}{\pgfqpoint{3.229754in}{5.139294in}}{\pgfqpoint{3.223930in}{5.139294in}}%
\pgfpathcurveto{\pgfqpoint{3.218106in}{5.139294in}}{\pgfqpoint{3.212520in}{5.136980in}}{\pgfqpoint{3.208401in}{5.132862in}}%
\pgfpathcurveto{\pgfqpoint{3.204283in}{5.128744in}}{\pgfqpoint{3.201969in}{5.123158in}}{\pgfqpoint{3.201969in}{5.117334in}}%
\pgfpathcurveto{\pgfqpoint{3.201969in}{5.111510in}}{\pgfqpoint{3.204283in}{5.105924in}}{\pgfqpoint{3.208401in}{5.101806in}}%
\pgfpathcurveto{\pgfqpoint{3.212520in}{5.097688in}}{\pgfqpoint{3.218106in}{5.095374in}}{\pgfqpoint{3.223930in}{5.095374in}}%
\pgfpathlineto{\pgfqpoint{3.223930in}{5.095374in}}%
\pgfpathclose%
\pgfusepath{stroke,fill}%
\end{pgfscope}%
\begin{pgfscope}%
\pgfpathrectangle{\pgfqpoint{1.414736in}{0.548769in}}{\pgfqpoint{3.497747in}{7.253186in}}%
\pgfusepath{clip}%
\pgfsetbuttcap%
\pgfsetroundjoin%
\definecolor{currentfill}{rgb}{0.225863,0.330805,0.547314}%
\pgfsetfillcolor{currentfill}%
\pgfsetlinewidth{1.003750pt}%
\definecolor{currentstroke}{rgb}{0.225863,0.330805,0.547314}%
\pgfsetstrokecolor{currentstroke}%
\pgfsetdash{}{0pt}%
\pgfpathmoveto{\pgfqpoint{3.968606in}{5.566360in}}%
\pgfpathcurveto{\pgfqpoint{3.974430in}{5.566360in}}{\pgfqpoint{3.980016in}{5.568674in}}{\pgfqpoint{3.984135in}{5.572792in}}%
\pgfpathcurveto{\pgfqpoint{3.988253in}{5.576910in}}{\pgfqpoint{3.990567in}{5.582496in}}{\pgfqpoint{3.990567in}{5.588320in}}%
\pgfpathcurveto{\pgfqpoint{3.990567in}{5.594144in}}{\pgfqpoint{3.988253in}{5.599730in}}{\pgfqpoint{3.984135in}{5.603848in}}%
\pgfpathcurveto{\pgfqpoint{3.980016in}{5.607967in}}{\pgfqpoint{3.974430in}{5.610280in}}{\pgfqpoint{3.968606in}{5.610280in}}%
\pgfpathcurveto{\pgfqpoint{3.962782in}{5.610280in}}{\pgfqpoint{3.957196in}{5.607967in}}{\pgfqpoint{3.953078in}{5.603848in}}%
\pgfpathcurveto{\pgfqpoint{3.948960in}{5.599730in}}{\pgfqpoint{3.946646in}{5.594144in}}{\pgfqpoint{3.946646in}{5.588320in}}%
\pgfpathcurveto{\pgfqpoint{3.946646in}{5.582496in}}{\pgfqpoint{3.948960in}{5.576910in}}{\pgfqpoint{3.953078in}{5.572792in}}%
\pgfpathcurveto{\pgfqpoint{3.957196in}{5.568674in}}{\pgfqpoint{3.962782in}{5.566360in}}{\pgfqpoint{3.968606in}{5.566360in}}%
\pgfpathlineto{\pgfqpoint{3.968606in}{5.566360in}}%
\pgfpathclose%
\pgfusepath{stroke,fill}%
\end{pgfscope}%
\begin{pgfscope}%
\pgfpathrectangle{\pgfqpoint{1.414736in}{0.548769in}}{\pgfqpoint{3.497747in}{7.253186in}}%
\pgfusepath{clip}%
\pgfsetbuttcap%
\pgfsetroundjoin%
\definecolor{currentfill}{rgb}{0.162016,0.687316,0.499129}%
\pgfsetfillcolor{currentfill}%
\pgfsetlinewidth{1.003750pt}%
\definecolor{currentstroke}{rgb}{0.162016,0.687316,0.499129}%
\pgfsetstrokecolor{currentstroke}%
\pgfsetdash{}{0pt}%
\pgfpathmoveto{\pgfqpoint{3.227194in}{6.037346in}}%
\pgfpathcurveto{\pgfqpoint{3.233018in}{6.037346in}}{\pgfqpoint{3.238605in}{6.039660in}}{\pgfqpoint{3.242723in}{6.043778in}}%
\pgfpathcurveto{\pgfqpoint{3.246841in}{6.047896in}}{\pgfqpoint{3.249155in}{6.053482in}}{\pgfqpoint{3.249155in}{6.059306in}}%
\pgfpathcurveto{\pgfqpoint{3.249155in}{6.065130in}}{\pgfqpoint{3.246841in}{6.070716in}}{\pgfqpoint{3.242723in}{6.074835in}}%
\pgfpathcurveto{\pgfqpoint{3.238605in}{6.078953in}}{\pgfqpoint{3.233018in}{6.081267in}}{\pgfqpoint{3.227194in}{6.081267in}}%
\pgfpathcurveto{\pgfqpoint{3.221370in}{6.081267in}}{\pgfqpoint{3.215784in}{6.078953in}}{\pgfqpoint{3.211666in}{6.074835in}}%
\pgfpathcurveto{\pgfqpoint{3.207548in}{6.070716in}}{\pgfqpoint{3.205234in}{6.065130in}}{\pgfqpoint{3.205234in}{6.059306in}}%
\pgfpathcurveto{\pgfqpoint{3.205234in}{6.053482in}}{\pgfqpoint{3.207548in}{6.047896in}}{\pgfqpoint{3.211666in}{6.043778in}}%
\pgfpathcurveto{\pgfqpoint{3.215784in}{6.039660in}}{\pgfqpoint{3.221370in}{6.037346in}}{\pgfqpoint{3.227194in}{6.037346in}}%
\pgfpathlineto{\pgfqpoint{3.227194in}{6.037346in}}%
\pgfpathclose%
\pgfusepath{stroke,fill}%
\end{pgfscope}%
\begin{pgfscope}%
\pgfpathrectangle{\pgfqpoint{1.414736in}{0.548769in}}{\pgfqpoint{3.497747in}{7.253186in}}%
\pgfusepath{clip}%
\pgfsetbuttcap%
\pgfsetroundjoin%
\definecolor{currentfill}{rgb}{0.993248,0.906157,0.143936}%
\pgfsetfillcolor{currentfill}%
\pgfsetlinewidth{1.003750pt}%
\definecolor{currentstroke}{rgb}{0.993248,0.906157,0.143936}%
\pgfsetstrokecolor{currentstroke}%
\pgfsetdash{}{0pt}%
\pgfpathmoveto{\pgfqpoint{2.924570in}{6.508332in}}%
\pgfpathcurveto{\pgfqpoint{2.930394in}{6.508332in}}{\pgfqpoint{2.935980in}{6.510646in}}{\pgfqpoint{2.940098in}{6.514764in}}%
\pgfpathcurveto{\pgfqpoint{2.944217in}{6.518882in}}{\pgfqpoint{2.946531in}{6.524468in}}{\pgfqpoint{2.946531in}{6.530292in}}%
\pgfpathcurveto{\pgfqpoint{2.946531in}{6.536116in}}{\pgfqpoint{2.944217in}{6.541702in}}{\pgfqpoint{2.940098in}{6.545821in}}%
\pgfpathcurveto{\pgfqpoint{2.935980in}{6.549939in}}{\pgfqpoint{2.930394in}{6.552253in}}{\pgfqpoint{2.924570in}{6.552253in}}%
\pgfpathcurveto{\pgfqpoint{2.918746in}{6.552253in}}{\pgfqpoint{2.913160in}{6.549939in}}{\pgfqpoint{2.909042in}{6.545821in}}%
\pgfpathcurveto{\pgfqpoint{2.904924in}{6.541702in}}{\pgfqpoint{2.902610in}{6.536116in}}{\pgfqpoint{2.902610in}{6.530292in}}%
\pgfpathcurveto{\pgfqpoint{2.902610in}{6.524468in}}{\pgfqpoint{2.904924in}{6.518882in}}{\pgfqpoint{2.909042in}{6.514764in}}%
\pgfpathcurveto{\pgfqpoint{2.913160in}{6.510646in}}{\pgfqpoint{2.918746in}{6.508332in}}{\pgfqpoint{2.924570in}{6.508332in}}%
\pgfpathlineto{\pgfqpoint{2.924570in}{6.508332in}}%
\pgfpathclose%
\pgfusepath{stroke,fill}%
\end{pgfscope}%
\begin{pgfscope}%
\pgfpathrectangle{\pgfqpoint{1.414736in}{0.548769in}}{\pgfqpoint{3.497747in}{7.253186in}}%
\pgfusepath{clip}%
\pgfsetbuttcap%
\pgfsetroundjoin%
\definecolor{currentfill}{rgb}{0.783315,0.879285,0.125405}%
\pgfsetfillcolor{currentfill}%
\pgfsetlinewidth{1.003750pt}%
\definecolor{currentstroke}{rgb}{0.783315,0.879285,0.125405}%
\pgfsetstrokecolor{currentstroke}%
\pgfsetdash{}{0pt}%
\pgfpathmoveto{\pgfqpoint{3.408360in}{6.979318in}}%
\pgfpathcurveto{\pgfqpoint{3.414184in}{6.979318in}}{\pgfqpoint{3.419770in}{6.981632in}}{\pgfqpoint{3.423888in}{6.985750in}}%
\pgfpathcurveto{\pgfqpoint{3.428006in}{6.989868in}}{\pgfqpoint{3.430320in}{6.995455in}}{\pgfqpoint{3.430320in}{7.001278in}}%
\pgfpathcurveto{\pgfqpoint{3.430320in}{7.007102in}}{\pgfqpoint{3.428006in}{7.012689in}}{\pgfqpoint{3.423888in}{7.016807in}}%
\pgfpathcurveto{\pgfqpoint{3.419770in}{7.020925in}}{\pgfqpoint{3.414184in}{7.023239in}}{\pgfqpoint{3.408360in}{7.023239in}}%
\pgfpathcurveto{\pgfqpoint{3.402536in}{7.023239in}}{\pgfqpoint{3.396950in}{7.020925in}}{\pgfqpoint{3.392832in}{7.016807in}}%
\pgfpathcurveto{\pgfqpoint{3.388714in}{7.012689in}}{\pgfqpoint{3.386400in}{7.007102in}}{\pgfqpoint{3.386400in}{7.001278in}}%
\pgfpathcurveto{\pgfqpoint{3.386400in}{6.995455in}}{\pgfqpoint{3.388714in}{6.989868in}}{\pgfqpoint{3.392832in}{6.985750in}}%
\pgfpathcurveto{\pgfqpoint{3.396950in}{6.981632in}}{\pgfqpoint{3.402536in}{6.979318in}}{\pgfqpoint{3.408360in}{6.979318in}}%
\pgfpathlineto{\pgfqpoint{3.408360in}{6.979318in}}%
\pgfpathclose%
\pgfusepath{stroke,fill}%
\end{pgfscope}%
\begin{pgfscope}%
\pgfpathrectangle{\pgfqpoint{1.414736in}{0.548769in}}{\pgfqpoint{3.497747in}{7.253186in}}%
\pgfusepath{clip}%
\pgfsetbuttcap%
\pgfsetroundjoin%
\definecolor{currentfill}{rgb}{0.855810,0.888601,0.097452}%
\pgfsetfillcolor{currentfill}%
\pgfsetlinewidth{1.003750pt}%
\definecolor{currentstroke}{rgb}{0.855810,0.888601,0.097452}%
\pgfsetstrokecolor{currentstroke}%
\pgfsetdash{}{0pt}%
\pgfpathmoveto{\pgfqpoint{4.257176in}{7.450304in}}%
\pgfpathcurveto{\pgfqpoint{4.263000in}{7.450304in}}{\pgfqpoint{4.268587in}{7.452618in}}{\pgfqpoint{4.272705in}{7.456736in}}%
\pgfpathcurveto{\pgfqpoint{4.276823in}{7.460854in}}{\pgfqpoint{4.279137in}{7.466441in}}{\pgfqpoint{4.279137in}{7.472265in}}%
\pgfpathcurveto{\pgfqpoint{4.279137in}{7.478089in}}{\pgfqpoint{4.276823in}{7.483675in}}{\pgfqpoint{4.272705in}{7.487793in}}%
\pgfpathcurveto{\pgfqpoint{4.268587in}{7.491911in}}{\pgfqpoint{4.263000in}{7.494225in}}{\pgfqpoint{4.257176in}{7.494225in}}%
\pgfpathcurveto{\pgfqpoint{4.251352in}{7.494225in}}{\pgfqpoint{4.245766in}{7.491911in}}{\pgfqpoint{4.241648in}{7.487793in}}%
\pgfpathcurveto{\pgfqpoint{4.237530in}{7.483675in}}{\pgfqpoint{4.235216in}{7.478089in}}{\pgfqpoint{4.235216in}{7.472265in}}%
\pgfpathcurveto{\pgfqpoint{4.235216in}{7.466441in}}{\pgfqpoint{4.237530in}{7.460854in}}{\pgfqpoint{4.241648in}{7.456736in}}%
\pgfpathcurveto{\pgfqpoint{4.245766in}{7.452618in}}{\pgfqpoint{4.251352in}{7.450304in}}{\pgfqpoint{4.257176in}{7.450304in}}%
\pgfpathlineto{\pgfqpoint{4.257176in}{7.450304in}}%
\pgfpathclose%
\pgfusepath{stroke,fill}%
\end{pgfscope}%
\begin{pgfscope}%
\pgfpathrectangle{\pgfqpoint{1.414736in}{0.548769in}}{\pgfqpoint{3.497747in}{7.253186in}}%
\pgfusepath{clip}%
\pgfsetbuttcap%
\pgfsetroundjoin%
\definecolor{currentfill}{rgb}{0.993248,0.906157,0.143936}%
\pgfsetfillcolor{currentfill}%
\pgfsetlinewidth{1.003750pt}%
\definecolor{currentstroke}{rgb}{0.993248,0.906157,0.143936}%
\pgfsetstrokecolor{currentstroke}%
\pgfsetdash{}{0pt}%
\pgfpathmoveto{\pgfqpoint{3.775945in}{0.856499in}}%
\pgfpathcurveto{\pgfqpoint{3.781769in}{0.856499in}}{\pgfqpoint{3.787355in}{0.858813in}}{\pgfqpoint{3.791473in}{0.862931in}}%
\pgfpathcurveto{\pgfqpoint{3.795591in}{0.867049in}}{\pgfqpoint{3.797905in}{0.872635in}}{\pgfqpoint{3.797905in}{0.878459in}}%
\pgfpathcurveto{\pgfqpoint{3.797905in}{0.884283in}}{\pgfqpoint{3.795591in}{0.889869in}}{\pgfqpoint{3.791473in}{0.893987in}}%
\pgfpathcurveto{\pgfqpoint{3.787355in}{0.898105in}}{\pgfqpoint{3.781769in}{0.900419in}}{\pgfqpoint{3.775945in}{0.900419in}}%
\pgfpathcurveto{\pgfqpoint{3.770121in}{0.900419in}}{\pgfqpoint{3.764535in}{0.898105in}}{\pgfqpoint{3.760417in}{0.893987in}}%
\pgfpathcurveto{\pgfqpoint{3.756299in}{0.889869in}}{\pgfqpoint{3.753985in}{0.884283in}}{\pgfqpoint{3.753985in}{0.878459in}}%
\pgfpathcurveto{\pgfqpoint{3.753985in}{0.872635in}}{\pgfqpoint{3.756299in}{0.867049in}}{\pgfqpoint{3.760417in}{0.862931in}}%
\pgfpathcurveto{\pgfqpoint{3.764535in}{0.858813in}}{\pgfqpoint{3.770121in}{0.856499in}}{\pgfqpoint{3.775945in}{0.856499in}}%
\pgfpathlineto{\pgfqpoint{3.775945in}{0.856499in}}%
\pgfpathclose%
\pgfusepath{stroke,fill}%
\end{pgfscope}%
\begin{pgfscope}%
\pgfpathrectangle{\pgfqpoint{1.414736in}{0.548769in}}{\pgfqpoint{3.497747in}{7.253186in}}%
\pgfusepath{clip}%
\pgfsetbuttcap%
\pgfsetroundjoin%
\definecolor{currentfill}{rgb}{0.804182,0.882046,0.114965}%
\pgfsetfillcolor{currentfill}%
\pgfsetlinewidth{1.003750pt}%
\definecolor{currentstroke}{rgb}{0.804182,0.882046,0.114965}%
\pgfsetstrokecolor{currentstroke}%
\pgfsetdash{}{0pt}%
\pgfpathmoveto{\pgfqpoint{3.393736in}{1.327485in}}%
\pgfpathcurveto{\pgfqpoint{3.399560in}{1.327485in}}{\pgfqpoint{3.405146in}{1.329799in}}{\pgfqpoint{3.409264in}{1.333917in}}%
\pgfpathcurveto{\pgfqpoint{3.413382in}{1.338035in}}{\pgfqpoint{3.415696in}{1.343621in}}{\pgfqpoint{3.415696in}{1.349445in}}%
\pgfpathcurveto{\pgfqpoint{3.415696in}{1.355269in}}{\pgfqpoint{3.413382in}{1.360855in}}{\pgfqpoint{3.409264in}{1.364973in}}%
\pgfpathcurveto{\pgfqpoint{3.405146in}{1.369092in}}{\pgfqpoint{3.399560in}{1.371405in}}{\pgfqpoint{3.393736in}{1.371405in}}%
\pgfpathcurveto{\pgfqpoint{3.387912in}{1.371405in}}{\pgfqpoint{3.382326in}{1.369092in}}{\pgfqpoint{3.378207in}{1.364973in}}%
\pgfpathcurveto{\pgfqpoint{3.374089in}{1.360855in}}{\pgfqpoint{3.371775in}{1.355269in}}{\pgfqpoint{3.371775in}{1.349445in}}%
\pgfpathcurveto{\pgfqpoint{3.371775in}{1.343621in}}{\pgfqpoint{3.374089in}{1.338035in}}{\pgfqpoint{3.378207in}{1.333917in}}%
\pgfpathcurveto{\pgfqpoint{3.382326in}{1.329799in}}{\pgfqpoint{3.387912in}{1.327485in}}{\pgfqpoint{3.393736in}{1.327485in}}%
\pgfpathlineto{\pgfqpoint{3.393736in}{1.327485in}}%
\pgfpathclose%
\pgfusepath{stroke,fill}%
\end{pgfscope}%
\begin{pgfscope}%
\pgfpathrectangle{\pgfqpoint{1.414736in}{0.548769in}}{\pgfqpoint{3.497747in}{7.253186in}}%
\pgfusepath{clip}%
\pgfsetbuttcap%
\pgfsetroundjoin%
\definecolor{currentfill}{rgb}{0.699415,0.867117,0.175971}%
\pgfsetfillcolor{currentfill}%
\pgfsetlinewidth{1.003750pt}%
\definecolor{currentstroke}{rgb}{0.699415,0.867117,0.175971}%
\pgfsetstrokecolor{currentstroke}%
\pgfsetdash{}{0pt}%
\pgfpathmoveto{\pgfqpoint{3.703715in}{1.798471in}}%
\pgfpathcurveto{\pgfqpoint{3.709539in}{1.798471in}}{\pgfqpoint{3.715125in}{1.800785in}}{\pgfqpoint{3.719243in}{1.804903in}}%
\pgfpathcurveto{\pgfqpoint{3.723361in}{1.809021in}}{\pgfqpoint{3.725675in}{1.814607in}}{\pgfqpoint{3.725675in}{1.820431in}}%
\pgfpathcurveto{\pgfqpoint{3.725675in}{1.826255in}}{\pgfqpoint{3.723361in}{1.831841in}}{\pgfqpoint{3.719243in}{1.835960in}}%
\pgfpathcurveto{\pgfqpoint{3.715125in}{1.840078in}}{\pgfqpoint{3.709539in}{1.842392in}}{\pgfqpoint{3.703715in}{1.842392in}}%
\pgfpathcurveto{\pgfqpoint{3.697891in}{1.842392in}}{\pgfqpoint{3.692305in}{1.840078in}}{\pgfqpoint{3.688187in}{1.835960in}}%
\pgfpathcurveto{\pgfqpoint{3.684068in}{1.831841in}}{\pgfqpoint{3.681755in}{1.826255in}}{\pgfqpoint{3.681755in}{1.820431in}}%
\pgfpathcurveto{\pgfqpoint{3.681755in}{1.814607in}}{\pgfqpoint{3.684068in}{1.809021in}}{\pgfqpoint{3.688187in}{1.804903in}}%
\pgfpathcurveto{\pgfqpoint{3.692305in}{1.800785in}}{\pgfqpoint{3.697891in}{1.798471in}}{\pgfqpoint{3.703715in}{1.798471in}}%
\pgfpathlineto{\pgfqpoint{3.703715in}{1.798471in}}%
\pgfpathclose%
\pgfusepath{stroke,fill}%
\end{pgfscope}%
\begin{pgfscope}%
\pgfpathrectangle{\pgfqpoint{1.414736in}{0.548769in}}{\pgfqpoint{3.497747in}{7.253186in}}%
\pgfusepath{clip}%
\pgfsetbuttcap%
\pgfsetroundjoin%
\definecolor{currentfill}{rgb}{0.993248,0.906157,0.143936}%
\pgfsetfillcolor{currentfill}%
\pgfsetlinewidth{1.003750pt}%
\definecolor{currentstroke}{rgb}{0.993248,0.906157,0.143936}%
\pgfsetstrokecolor{currentstroke}%
\pgfsetdash{}{0pt}%
\pgfpathmoveto{\pgfqpoint{3.527371in}{2.269457in}}%
\pgfpathcurveto{\pgfqpoint{3.533195in}{2.269457in}}{\pgfqpoint{3.538781in}{2.271771in}}{\pgfqpoint{3.542899in}{2.275889in}}%
\pgfpathcurveto{\pgfqpoint{3.547017in}{2.280007in}}{\pgfqpoint{3.549331in}{2.285593in}}{\pgfqpoint{3.549331in}{2.291417in}}%
\pgfpathcurveto{\pgfqpoint{3.549331in}{2.297241in}}{\pgfqpoint{3.547017in}{2.302828in}}{\pgfqpoint{3.542899in}{2.306946in}}%
\pgfpathcurveto{\pgfqpoint{3.538781in}{2.311064in}}{\pgfqpoint{3.533195in}{2.313378in}}{\pgfqpoint{3.527371in}{2.313378in}}%
\pgfpathcurveto{\pgfqpoint{3.521547in}{2.313378in}}{\pgfqpoint{3.515961in}{2.311064in}}{\pgfqpoint{3.511843in}{2.306946in}}%
\pgfpathcurveto{\pgfqpoint{3.507725in}{2.302828in}}{\pgfqpoint{3.505411in}{2.297241in}}{\pgfqpoint{3.505411in}{2.291417in}}%
\pgfpathcurveto{\pgfqpoint{3.505411in}{2.285593in}}{\pgfqpoint{3.507725in}{2.280007in}}{\pgfqpoint{3.511843in}{2.275889in}}%
\pgfpathcurveto{\pgfqpoint{3.515961in}{2.271771in}}{\pgfqpoint{3.521547in}{2.269457in}}{\pgfqpoint{3.527371in}{2.269457in}}%
\pgfpathlineto{\pgfqpoint{3.527371in}{2.269457in}}%
\pgfpathclose%
\pgfusepath{stroke,fill}%
\end{pgfscope}%
\begin{pgfscope}%
\pgfpathrectangle{\pgfqpoint{1.414736in}{0.548769in}}{\pgfqpoint{3.497747in}{7.253186in}}%
\pgfusepath{clip}%
\pgfsetbuttcap%
\pgfsetroundjoin%
\definecolor{currentfill}{rgb}{0.266941,0.748751,0.440573}%
\pgfsetfillcolor{currentfill}%
\pgfsetlinewidth{1.003750pt}%
\definecolor{currentstroke}{rgb}{0.266941,0.748751,0.440573}%
\pgfsetstrokecolor{currentstroke}%
\pgfsetdash{}{0pt}%
\pgfpathmoveto{\pgfqpoint{3.738168in}{2.740443in}}%
\pgfpathcurveto{\pgfqpoint{3.743992in}{2.740443in}}{\pgfqpoint{3.749578in}{2.742757in}}{\pgfqpoint{3.753697in}{2.746875in}}%
\pgfpathcurveto{\pgfqpoint{3.757815in}{2.750993in}}{\pgfqpoint{3.760129in}{2.756580in}}{\pgfqpoint{3.760129in}{2.762404in}}%
\pgfpathcurveto{\pgfqpoint{3.760129in}{2.768227in}}{\pgfqpoint{3.757815in}{2.773814in}}{\pgfqpoint{3.753697in}{2.777932in}}%
\pgfpathcurveto{\pgfqpoint{3.749578in}{2.782050in}}{\pgfqpoint{3.743992in}{2.784364in}}{\pgfqpoint{3.738168in}{2.784364in}}%
\pgfpathcurveto{\pgfqpoint{3.732344in}{2.784364in}}{\pgfqpoint{3.726758in}{2.782050in}}{\pgfqpoint{3.722640in}{2.777932in}}%
\pgfpathcurveto{\pgfqpoint{3.718522in}{2.773814in}}{\pgfqpoint{3.716208in}{2.768227in}}{\pgfqpoint{3.716208in}{2.762404in}}%
\pgfpathcurveto{\pgfqpoint{3.716208in}{2.756580in}}{\pgfqpoint{3.718522in}{2.750993in}}{\pgfqpoint{3.722640in}{2.746875in}}%
\pgfpathcurveto{\pgfqpoint{3.726758in}{2.742757in}}{\pgfqpoint{3.732344in}{2.740443in}}{\pgfqpoint{3.738168in}{2.740443in}}%
\pgfpathlineto{\pgfqpoint{3.738168in}{2.740443in}}%
\pgfpathclose%
\pgfusepath{stroke,fill}%
\end{pgfscope}%
\begin{pgfscope}%
\pgfpathrectangle{\pgfqpoint{1.414736in}{0.548769in}}{\pgfqpoint{3.497747in}{7.253186in}}%
\pgfusepath{clip}%
\pgfsetbuttcap%
\pgfsetroundjoin%
\definecolor{currentfill}{rgb}{0.772852,0.877868,0.131109}%
\pgfsetfillcolor{currentfill}%
\pgfsetlinewidth{1.003750pt}%
\definecolor{currentstroke}{rgb}{0.772852,0.877868,0.131109}%
\pgfsetstrokecolor{currentstroke}%
\pgfsetdash{}{0pt}%
\pgfpathmoveto{\pgfqpoint{3.462003in}{3.211429in}}%
\pgfpathcurveto{\pgfqpoint{3.467827in}{3.211429in}}{\pgfqpoint{3.473413in}{3.213743in}}{\pgfqpoint{3.477531in}{3.217861in}}%
\pgfpathcurveto{\pgfqpoint{3.481649in}{3.221980in}}{\pgfqpoint{3.483963in}{3.227566in}}{\pgfqpoint{3.483963in}{3.233390in}}%
\pgfpathcurveto{\pgfqpoint{3.483963in}{3.239214in}}{\pgfqpoint{3.481649in}{3.244800in}}{\pgfqpoint{3.477531in}{3.248918in}}%
\pgfpathcurveto{\pgfqpoint{3.473413in}{3.253036in}}{\pgfqpoint{3.467827in}{3.255350in}}{\pgfqpoint{3.462003in}{3.255350in}}%
\pgfpathcurveto{\pgfqpoint{3.456179in}{3.255350in}}{\pgfqpoint{3.450593in}{3.253036in}}{\pgfqpoint{3.446474in}{3.248918in}}%
\pgfpathcurveto{\pgfqpoint{3.442356in}{3.244800in}}{\pgfqpoint{3.440042in}{3.239214in}}{\pgfqpoint{3.440042in}{3.233390in}}%
\pgfpathcurveto{\pgfqpoint{3.440042in}{3.227566in}}{\pgfqpoint{3.442356in}{3.221980in}}{\pgfqpoint{3.446474in}{3.217861in}}%
\pgfpathcurveto{\pgfqpoint{3.450593in}{3.213743in}}{\pgfqpoint{3.456179in}{3.211429in}}{\pgfqpoint{3.462003in}{3.211429in}}%
\pgfpathlineto{\pgfqpoint{3.462003in}{3.211429in}}%
\pgfpathclose%
\pgfusepath{stroke,fill}%
\end{pgfscope}%
\begin{pgfscope}%
\pgfpathrectangle{\pgfqpoint{1.414736in}{0.548769in}}{\pgfqpoint{3.497747in}{7.253186in}}%
\pgfusepath{clip}%
\pgfsetbuttcap%
\pgfsetroundjoin%
\definecolor{currentfill}{rgb}{0.993248,0.906157,0.143936}%
\pgfsetfillcolor{currentfill}%
\pgfsetlinewidth{1.003750pt}%
\definecolor{currentstroke}{rgb}{0.993248,0.906157,0.143936}%
\pgfsetstrokecolor{currentstroke}%
\pgfsetdash{}{0pt}%
\pgfpathmoveto{\pgfqpoint{3.366505in}{3.682415in}}%
\pgfpathcurveto{\pgfqpoint{3.372329in}{3.682415in}}{\pgfqpoint{3.377915in}{3.684729in}}{\pgfqpoint{3.382033in}{3.688847in}}%
\pgfpathcurveto{\pgfqpoint{3.386151in}{3.692966in}}{\pgfqpoint{3.388465in}{3.698552in}}{\pgfqpoint{3.388465in}{3.704376in}}%
\pgfpathcurveto{\pgfqpoint{3.388465in}{3.710200in}}{\pgfqpoint{3.386151in}{3.715786in}}{\pgfqpoint{3.382033in}{3.719904in}}%
\pgfpathcurveto{\pgfqpoint{3.377915in}{3.724022in}}{\pgfqpoint{3.372329in}{3.726336in}}{\pgfqpoint{3.366505in}{3.726336in}}%
\pgfpathcurveto{\pgfqpoint{3.360681in}{3.726336in}}{\pgfqpoint{3.355095in}{3.724022in}}{\pgfqpoint{3.350977in}{3.719904in}}%
\pgfpathcurveto{\pgfqpoint{3.346859in}{3.715786in}}{\pgfqpoint{3.344545in}{3.710200in}}{\pgfqpoint{3.344545in}{3.704376in}}%
\pgfpathcurveto{\pgfqpoint{3.344545in}{3.698552in}}{\pgfqpoint{3.346859in}{3.692966in}}{\pgfqpoint{3.350977in}{3.688847in}}%
\pgfpathcurveto{\pgfqpoint{3.355095in}{3.684729in}}{\pgfqpoint{3.360681in}{3.682415in}}{\pgfqpoint{3.366505in}{3.682415in}}%
\pgfpathlineto{\pgfqpoint{3.366505in}{3.682415in}}%
\pgfpathclose%
\pgfusepath{stroke,fill}%
\end{pgfscope}%
\begin{pgfscope}%
\pgfpathrectangle{\pgfqpoint{1.414736in}{0.548769in}}{\pgfqpoint{3.497747in}{7.253186in}}%
\pgfusepath{clip}%
\pgfsetbuttcap%
\pgfsetroundjoin%
\definecolor{currentfill}{rgb}{0.709898,0.868751,0.169257}%
\pgfsetfillcolor{currentfill}%
\pgfsetlinewidth{1.003750pt}%
\definecolor{currentstroke}{rgb}{0.709898,0.868751,0.169257}%
\pgfsetstrokecolor{currentstroke}%
\pgfsetdash{}{0pt}%
\pgfpathmoveto{\pgfqpoint{3.777241in}{4.153402in}}%
\pgfpathcurveto{\pgfqpoint{3.783065in}{4.153402in}}{\pgfqpoint{3.788651in}{4.155715in}}{\pgfqpoint{3.792769in}{4.159834in}}%
\pgfpathcurveto{\pgfqpoint{3.796887in}{4.163952in}}{\pgfqpoint{3.799201in}{4.169538in}}{\pgfqpoint{3.799201in}{4.175362in}}%
\pgfpathcurveto{\pgfqpoint{3.799201in}{4.181186in}}{\pgfqpoint{3.796887in}{4.186772in}}{\pgfqpoint{3.792769in}{4.190890in}}%
\pgfpathcurveto{\pgfqpoint{3.788651in}{4.195008in}}{\pgfqpoint{3.783065in}{4.197322in}}{\pgfqpoint{3.777241in}{4.197322in}}%
\pgfpathcurveto{\pgfqpoint{3.771417in}{4.197322in}}{\pgfqpoint{3.765831in}{4.195008in}}{\pgfqpoint{3.761713in}{4.190890in}}%
\pgfpathcurveto{\pgfqpoint{3.757594in}{4.186772in}}{\pgfqpoint{3.755281in}{4.181186in}}{\pgfqpoint{3.755281in}{4.175362in}}%
\pgfpathcurveto{\pgfqpoint{3.755281in}{4.169538in}}{\pgfqpoint{3.757594in}{4.163952in}}{\pgfqpoint{3.761713in}{4.159834in}}%
\pgfpathcurveto{\pgfqpoint{3.765831in}{4.155715in}}{\pgfqpoint{3.771417in}{4.153402in}}{\pgfqpoint{3.777241in}{4.153402in}}%
\pgfpathlineto{\pgfqpoint{3.777241in}{4.153402in}}%
\pgfpathclose%
\pgfusepath{stroke,fill}%
\end{pgfscope}%
\begin{pgfscope}%
\pgfpathrectangle{\pgfqpoint{1.414736in}{0.548769in}}{\pgfqpoint{3.497747in}{7.253186in}}%
\pgfusepath{clip}%
\pgfsetbuttcap%
\pgfsetroundjoin%
\definecolor{currentfill}{rgb}{0.866013,0.889868,0.095953}%
\pgfsetfillcolor{currentfill}%
\pgfsetlinewidth{1.003750pt}%
\definecolor{currentstroke}{rgb}{0.866013,0.889868,0.095953}%
\pgfsetstrokecolor{currentstroke}%
\pgfsetdash{}{0pt}%
\pgfpathmoveto{\pgfqpoint{3.780052in}{4.624388in}}%
\pgfpathcurveto{\pgfqpoint{3.785876in}{4.624388in}}{\pgfqpoint{3.791462in}{4.626702in}}{\pgfqpoint{3.795581in}{4.630820in}}%
\pgfpathcurveto{\pgfqpoint{3.799699in}{4.634938in}}{\pgfqpoint{3.802013in}{4.640524in}}{\pgfqpoint{3.802013in}{4.646348in}}%
\pgfpathcurveto{\pgfqpoint{3.802013in}{4.652172in}}{\pgfqpoint{3.799699in}{4.657758in}}{\pgfqpoint{3.795581in}{4.661876in}}%
\pgfpathcurveto{\pgfqpoint{3.791462in}{4.665994in}}{\pgfqpoint{3.785876in}{4.668308in}}{\pgfqpoint{3.780052in}{4.668308in}}%
\pgfpathcurveto{\pgfqpoint{3.774228in}{4.668308in}}{\pgfqpoint{3.768642in}{4.665994in}}{\pgfqpoint{3.764524in}{4.661876in}}%
\pgfpathcurveto{\pgfqpoint{3.760406in}{4.657758in}}{\pgfqpoint{3.758092in}{4.652172in}}{\pgfqpoint{3.758092in}{4.646348in}}%
\pgfpathcurveto{\pgfqpoint{3.758092in}{4.640524in}}{\pgfqpoint{3.760406in}{4.634938in}}{\pgfqpoint{3.764524in}{4.630820in}}%
\pgfpathcurveto{\pgfqpoint{3.768642in}{4.626702in}}{\pgfqpoint{3.774228in}{4.624388in}}{\pgfqpoint{3.780052in}{4.624388in}}%
\pgfpathlineto{\pgfqpoint{3.780052in}{4.624388in}}%
\pgfpathclose%
\pgfusepath{stroke,fill}%
\end{pgfscope}%
\begin{pgfscope}%
\pgfpathrectangle{\pgfqpoint{1.414736in}{0.548769in}}{\pgfqpoint{3.497747in}{7.253186in}}%
\pgfusepath{clip}%
\pgfsetbuttcap%
\pgfsetroundjoin%
\definecolor{currentfill}{rgb}{0.132268,0.655014,0.519661}%
\pgfsetfillcolor{currentfill}%
\pgfsetlinewidth{1.003750pt}%
\definecolor{currentstroke}{rgb}{0.132268,0.655014,0.519661}%
\pgfsetstrokecolor{currentstroke}%
\pgfsetdash{}{0pt}%
\pgfpathmoveto{\pgfqpoint{3.479213in}{5.095374in}}%
\pgfpathcurveto{\pgfqpoint{3.485037in}{5.095374in}}{\pgfqpoint{3.490623in}{5.097688in}}{\pgfqpoint{3.494741in}{5.101806in}}%
\pgfpathcurveto{\pgfqpoint{3.498859in}{5.105924in}}{\pgfqpoint{3.501173in}{5.111510in}}{\pgfqpoint{3.501173in}{5.117334in}}%
\pgfpathcurveto{\pgfqpoint{3.501173in}{5.123158in}}{\pgfqpoint{3.498859in}{5.128744in}}{\pgfqpoint{3.494741in}{5.132862in}}%
\pgfpathcurveto{\pgfqpoint{3.490623in}{5.136980in}}{\pgfqpoint{3.485037in}{5.139294in}}{\pgfqpoint{3.479213in}{5.139294in}}%
\pgfpathcurveto{\pgfqpoint{3.473389in}{5.139294in}}{\pgfqpoint{3.467803in}{5.136980in}}{\pgfqpoint{3.463685in}{5.132862in}}%
\pgfpathcurveto{\pgfqpoint{3.459567in}{5.128744in}}{\pgfqpoint{3.457253in}{5.123158in}}{\pgfqpoint{3.457253in}{5.117334in}}%
\pgfpathcurveto{\pgfqpoint{3.457253in}{5.111510in}}{\pgfqpoint{3.459567in}{5.105924in}}{\pgfqpoint{3.463685in}{5.101806in}}%
\pgfpathcurveto{\pgfqpoint{3.467803in}{5.097688in}}{\pgfqpoint{3.473389in}{5.095374in}}{\pgfqpoint{3.479213in}{5.095374in}}%
\pgfpathlineto{\pgfqpoint{3.479213in}{5.095374in}}%
\pgfpathclose%
\pgfusepath{stroke,fill}%
\end{pgfscope}%
\begin{pgfscope}%
\pgfpathrectangle{\pgfqpoint{1.414736in}{0.548769in}}{\pgfqpoint{3.497747in}{7.253186in}}%
\pgfusepath{clip}%
\pgfsetbuttcap%
\pgfsetroundjoin%
\definecolor{currentfill}{rgb}{0.133743,0.548535,0.553541}%
\pgfsetfillcolor{currentfill}%
\pgfsetlinewidth{1.003750pt}%
\definecolor{currentstroke}{rgb}{0.133743,0.548535,0.553541}%
\pgfsetstrokecolor{currentstroke}%
\pgfsetdash{}{0pt}%
\pgfpathmoveto{\pgfqpoint{3.890341in}{5.566360in}}%
\pgfpathcurveto{\pgfqpoint{3.896165in}{5.566360in}}{\pgfqpoint{3.901751in}{5.568674in}}{\pgfqpoint{3.905869in}{5.572792in}}%
\pgfpathcurveto{\pgfqpoint{3.909987in}{5.576910in}}{\pgfqpoint{3.912301in}{5.582496in}}{\pgfqpoint{3.912301in}{5.588320in}}%
\pgfpathcurveto{\pgfqpoint{3.912301in}{5.594144in}}{\pgfqpoint{3.909987in}{5.599730in}}{\pgfqpoint{3.905869in}{5.603848in}}%
\pgfpathcurveto{\pgfqpoint{3.901751in}{5.607967in}}{\pgfqpoint{3.896165in}{5.610280in}}{\pgfqpoint{3.890341in}{5.610280in}}%
\pgfpathcurveto{\pgfqpoint{3.884517in}{5.610280in}}{\pgfqpoint{3.878931in}{5.607967in}}{\pgfqpoint{3.874812in}{5.603848in}}%
\pgfpathcurveto{\pgfqpoint{3.870694in}{5.599730in}}{\pgfqpoint{3.868380in}{5.594144in}}{\pgfqpoint{3.868380in}{5.588320in}}%
\pgfpathcurveto{\pgfqpoint{3.868380in}{5.582496in}}{\pgfqpoint{3.870694in}{5.576910in}}{\pgfqpoint{3.874812in}{5.572792in}}%
\pgfpathcurveto{\pgfqpoint{3.878931in}{5.568674in}}{\pgfqpoint{3.884517in}{5.566360in}}{\pgfqpoint{3.890341in}{5.566360in}}%
\pgfpathlineto{\pgfqpoint{3.890341in}{5.566360in}}%
\pgfpathclose%
\pgfusepath{stroke,fill}%
\end{pgfscope}%
\begin{pgfscope}%
\pgfpathrectangle{\pgfqpoint{1.414736in}{0.548769in}}{\pgfqpoint{3.497747in}{7.253186in}}%
\pgfusepath{clip}%
\pgfsetbuttcap%
\pgfsetroundjoin%
\definecolor{currentfill}{rgb}{0.404001,0.800275,0.362552}%
\pgfsetfillcolor{currentfill}%
\pgfsetlinewidth{1.003750pt}%
\definecolor{currentstroke}{rgb}{0.404001,0.800275,0.362552}%
\pgfsetstrokecolor{currentstroke}%
\pgfsetdash{}{0pt}%
\pgfpathmoveto{\pgfqpoint{3.270365in}{6.037346in}}%
\pgfpathcurveto{\pgfqpoint{3.276189in}{6.037346in}}{\pgfqpoint{3.281775in}{6.039660in}}{\pgfqpoint{3.285893in}{6.043778in}}%
\pgfpathcurveto{\pgfqpoint{3.290011in}{6.047896in}}{\pgfqpoint{3.292325in}{6.053482in}}{\pgfqpoint{3.292325in}{6.059306in}}%
\pgfpathcurveto{\pgfqpoint{3.292325in}{6.065130in}}{\pgfqpoint{3.290011in}{6.070716in}}{\pgfqpoint{3.285893in}{6.074835in}}%
\pgfpathcurveto{\pgfqpoint{3.281775in}{6.078953in}}{\pgfqpoint{3.276189in}{6.081267in}}{\pgfqpoint{3.270365in}{6.081267in}}%
\pgfpathcurveto{\pgfqpoint{3.264541in}{6.081267in}}{\pgfqpoint{3.258955in}{6.078953in}}{\pgfqpoint{3.254837in}{6.074835in}}%
\pgfpathcurveto{\pgfqpoint{3.250719in}{6.070716in}}{\pgfqpoint{3.248405in}{6.065130in}}{\pgfqpoint{3.248405in}{6.059306in}}%
\pgfpathcurveto{\pgfqpoint{3.248405in}{6.053482in}}{\pgfqpoint{3.250719in}{6.047896in}}{\pgfqpoint{3.254837in}{6.043778in}}%
\pgfpathcurveto{\pgfqpoint{3.258955in}{6.039660in}}{\pgfqpoint{3.264541in}{6.037346in}}{\pgfqpoint{3.270365in}{6.037346in}}%
\pgfpathlineto{\pgfqpoint{3.270365in}{6.037346in}}%
\pgfpathclose%
\pgfusepath{stroke,fill}%
\end{pgfscope}%
\begin{pgfscope}%
\pgfpathrectangle{\pgfqpoint{1.414736in}{0.548769in}}{\pgfqpoint{3.497747in}{7.253186in}}%
\pgfusepath{clip}%
\pgfsetbuttcap%
\pgfsetroundjoin%
\definecolor{currentfill}{rgb}{0.993248,0.906157,0.143936}%
\pgfsetfillcolor{currentfill}%
\pgfsetlinewidth{1.003750pt}%
\definecolor{currentstroke}{rgb}{0.993248,0.906157,0.143936}%
\pgfsetstrokecolor{currentstroke}%
\pgfsetdash{}{0pt}%
\pgfpathmoveto{\pgfqpoint{3.935381in}{6.508332in}}%
\pgfpathcurveto{\pgfqpoint{3.941205in}{6.508332in}}{\pgfqpoint{3.946791in}{6.510646in}}{\pgfqpoint{3.950909in}{6.514764in}}%
\pgfpathcurveto{\pgfqpoint{3.955027in}{6.518882in}}{\pgfqpoint{3.957341in}{6.524468in}}{\pgfqpoint{3.957341in}{6.530292in}}%
\pgfpathcurveto{\pgfqpoint{3.957341in}{6.536116in}}{\pgfqpoint{3.955027in}{6.541702in}}{\pgfqpoint{3.950909in}{6.545821in}}%
\pgfpathcurveto{\pgfqpoint{3.946791in}{6.549939in}}{\pgfqpoint{3.941205in}{6.552253in}}{\pgfqpoint{3.935381in}{6.552253in}}%
\pgfpathcurveto{\pgfqpoint{3.929557in}{6.552253in}}{\pgfqpoint{3.923970in}{6.549939in}}{\pgfqpoint{3.919852in}{6.545821in}}%
\pgfpathcurveto{\pgfqpoint{3.915734in}{6.541702in}}{\pgfqpoint{3.913420in}{6.536116in}}{\pgfqpoint{3.913420in}{6.530292in}}%
\pgfpathcurveto{\pgfqpoint{3.913420in}{6.524468in}}{\pgfqpoint{3.915734in}{6.518882in}}{\pgfqpoint{3.919852in}{6.514764in}}%
\pgfpathcurveto{\pgfqpoint{3.923970in}{6.510646in}}{\pgfqpoint{3.929557in}{6.508332in}}{\pgfqpoint{3.935381in}{6.508332in}}%
\pgfpathlineto{\pgfqpoint{3.935381in}{6.508332in}}%
\pgfpathclose%
\pgfusepath{stroke,fill}%
\end{pgfscope}%
\begin{pgfscope}%
\pgfpathrectangle{\pgfqpoint{1.414736in}{0.548769in}}{\pgfqpoint{3.497747in}{7.253186in}}%
\pgfusepath{clip}%
\pgfsetbuttcap%
\pgfsetroundjoin%
\definecolor{currentfill}{rgb}{0.993248,0.906157,0.143936}%
\pgfsetfillcolor{currentfill}%
\pgfsetlinewidth{1.003750pt}%
\definecolor{currentstroke}{rgb}{0.993248,0.906157,0.143936}%
\pgfsetstrokecolor{currentstroke}%
\pgfsetdash{}{0pt}%
\pgfpathmoveto{\pgfqpoint{4.163369in}{6.979318in}}%
\pgfpathcurveto{\pgfqpoint{4.169193in}{6.979318in}}{\pgfqpoint{4.174779in}{6.981632in}}{\pgfqpoint{4.178897in}{6.985750in}}%
\pgfpathcurveto{\pgfqpoint{4.183015in}{6.989868in}}{\pgfqpoint{4.185329in}{6.995455in}}{\pgfqpoint{4.185329in}{7.001278in}}%
\pgfpathcurveto{\pgfqpoint{4.185329in}{7.007102in}}{\pgfqpoint{4.183015in}{7.012689in}}{\pgfqpoint{4.178897in}{7.016807in}}%
\pgfpathcurveto{\pgfqpoint{4.174779in}{7.020925in}}{\pgfqpoint{4.169193in}{7.023239in}}{\pgfqpoint{4.163369in}{7.023239in}}%
\pgfpathcurveto{\pgfqpoint{4.157545in}{7.023239in}}{\pgfqpoint{4.151959in}{7.020925in}}{\pgfqpoint{4.147841in}{7.016807in}}%
\pgfpathcurveto{\pgfqpoint{4.143722in}{7.012689in}}{\pgfqpoint{4.141408in}{7.007102in}}{\pgfqpoint{4.141408in}{7.001278in}}%
\pgfpathcurveto{\pgfqpoint{4.141408in}{6.995455in}}{\pgfqpoint{4.143722in}{6.989868in}}{\pgfqpoint{4.147841in}{6.985750in}}%
\pgfpathcurveto{\pgfqpoint{4.151959in}{6.981632in}}{\pgfqpoint{4.157545in}{6.979318in}}{\pgfqpoint{4.163369in}{6.979318in}}%
\pgfpathlineto{\pgfqpoint{4.163369in}{6.979318in}}%
\pgfpathclose%
\pgfusepath{stroke,fill}%
\end{pgfscope}%
\begin{pgfscope}%
\pgfpathrectangle{\pgfqpoint{1.414736in}{0.548769in}}{\pgfqpoint{3.497747in}{7.253186in}}%
\pgfusepath{clip}%
\pgfsetbuttcap%
\pgfsetroundjoin%
\definecolor{currentfill}{rgb}{0.814576,0.883393,0.110347}%
\pgfsetfillcolor{currentfill}%
\pgfsetlinewidth{1.003750pt}%
\definecolor{currentstroke}{rgb}{0.814576,0.883393,0.110347}%
\pgfsetstrokecolor{currentstroke}%
\pgfsetdash{}{0pt}%
\pgfpathmoveto{\pgfqpoint{3.670877in}{7.450304in}}%
\pgfpathcurveto{\pgfqpoint{3.676701in}{7.450304in}}{\pgfqpoint{3.682287in}{7.452618in}}{\pgfqpoint{3.686406in}{7.456736in}}%
\pgfpathcurveto{\pgfqpoint{3.690524in}{7.460854in}}{\pgfqpoint{3.692838in}{7.466441in}}{\pgfqpoint{3.692838in}{7.472265in}}%
\pgfpathcurveto{\pgfqpoint{3.692838in}{7.478089in}}{\pgfqpoint{3.690524in}{7.483675in}}{\pgfqpoint{3.686406in}{7.487793in}}%
\pgfpathcurveto{\pgfqpoint{3.682287in}{7.491911in}}{\pgfqpoint{3.676701in}{7.494225in}}{\pgfqpoint{3.670877in}{7.494225in}}%
\pgfpathcurveto{\pgfqpoint{3.665053in}{7.494225in}}{\pgfqpoint{3.659467in}{7.491911in}}{\pgfqpoint{3.655349in}{7.487793in}}%
\pgfpathcurveto{\pgfqpoint{3.651231in}{7.483675in}}{\pgfqpoint{3.648917in}{7.478089in}}{\pgfqpoint{3.648917in}{7.472265in}}%
\pgfpathcurveto{\pgfqpoint{3.648917in}{7.466441in}}{\pgfqpoint{3.651231in}{7.460854in}}{\pgfqpoint{3.655349in}{7.456736in}}%
\pgfpathcurveto{\pgfqpoint{3.659467in}{7.452618in}}{\pgfqpoint{3.665053in}{7.450304in}}{\pgfqpoint{3.670877in}{7.450304in}}%
\pgfpathlineto{\pgfqpoint{3.670877in}{7.450304in}}%
\pgfpathclose%
\pgfusepath{stroke,fill}%
\end{pgfscope}%
\begin{pgfscope}%
\pgfpathrectangle{\pgfqpoint{1.414736in}{0.548769in}}{\pgfqpoint{3.497747in}{7.253186in}}%
\pgfusepath{clip}%
\pgfsetbuttcap%
\pgfsetroundjoin%
\definecolor{currentfill}{rgb}{0.993248,0.906157,0.143936}%
\pgfsetfillcolor{currentfill}%
\pgfsetlinewidth{1.003750pt}%
\definecolor{currentstroke}{rgb}{0.993248,0.906157,0.143936}%
\pgfsetstrokecolor{currentstroke}%
\pgfsetdash{}{0pt}%
\pgfpathmoveto{\pgfqpoint{3.948314in}{0.856499in}}%
\pgfpathcurveto{\pgfqpoint{3.954137in}{0.856499in}}{\pgfqpoint{3.959724in}{0.858813in}}{\pgfqpoint{3.963842in}{0.862931in}}%
\pgfpathcurveto{\pgfqpoint{3.967960in}{0.867049in}}{\pgfqpoint{3.970274in}{0.872635in}}{\pgfqpoint{3.970274in}{0.878459in}}%
\pgfpathcurveto{\pgfqpoint{3.970274in}{0.884283in}}{\pgfqpoint{3.967960in}{0.889869in}}{\pgfqpoint{3.963842in}{0.893987in}}%
\pgfpathcurveto{\pgfqpoint{3.959724in}{0.898105in}}{\pgfqpoint{3.954137in}{0.900419in}}{\pgfqpoint{3.948314in}{0.900419in}}%
\pgfpathcurveto{\pgfqpoint{3.942490in}{0.900419in}}{\pgfqpoint{3.936903in}{0.898105in}}{\pgfqpoint{3.932785in}{0.893987in}}%
\pgfpathcurveto{\pgfqpoint{3.928667in}{0.889869in}}{\pgfqpoint{3.926353in}{0.884283in}}{\pgfqpoint{3.926353in}{0.878459in}}%
\pgfpathcurveto{\pgfqpoint{3.926353in}{0.872635in}}{\pgfqpoint{3.928667in}{0.867049in}}{\pgfqpoint{3.932785in}{0.862931in}}%
\pgfpathcurveto{\pgfqpoint{3.936903in}{0.858813in}}{\pgfqpoint{3.942490in}{0.856499in}}{\pgfqpoint{3.948314in}{0.856499in}}%
\pgfpathlineto{\pgfqpoint{3.948314in}{0.856499in}}%
\pgfpathclose%
\pgfusepath{stroke,fill}%
\end{pgfscope}%
\begin{pgfscope}%
\pgfpathrectangle{\pgfqpoint{1.414736in}{0.548769in}}{\pgfqpoint{3.497747in}{7.253186in}}%
\pgfusepath{clip}%
\pgfsetbuttcap%
\pgfsetroundjoin%
\definecolor{currentfill}{rgb}{0.751884,0.874951,0.143228}%
\pgfsetfillcolor{currentfill}%
\pgfsetlinewidth{1.003750pt}%
\definecolor{currentstroke}{rgb}{0.751884,0.874951,0.143228}%
\pgfsetstrokecolor{currentstroke}%
\pgfsetdash{}{0pt}%
\pgfpathmoveto{\pgfqpoint{3.465769in}{1.327485in}}%
\pgfpathcurveto{\pgfqpoint{3.471593in}{1.327485in}}{\pgfqpoint{3.477179in}{1.329799in}}{\pgfqpoint{3.481298in}{1.333917in}}%
\pgfpathcurveto{\pgfqpoint{3.485416in}{1.338035in}}{\pgfqpoint{3.487730in}{1.343621in}}{\pgfqpoint{3.487730in}{1.349445in}}%
\pgfpathcurveto{\pgfqpoint{3.487730in}{1.355269in}}{\pgfqpoint{3.485416in}{1.360855in}}{\pgfqpoint{3.481298in}{1.364973in}}%
\pgfpathcurveto{\pgfqpoint{3.477179in}{1.369092in}}{\pgfqpoint{3.471593in}{1.371405in}}{\pgfqpoint{3.465769in}{1.371405in}}%
\pgfpathcurveto{\pgfqpoint{3.459945in}{1.371405in}}{\pgfqpoint{3.454359in}{1.369092in}}{\pgfqpoint{3.450241in}{1.364973in}}%
\pgfpathcurveto{\pgfqpoint{3.446123in}{1.360855in}}{\pgfqpoint{3.443809in}{1.355269in}}{\pgfqpoint{3.443809in}{1.349445in}}%
\pgfpathcurveto{\pgfqpoint{3.443809in}{1.343621in}}{\pgfqpoint{3.446123in}{1.338035in}}{\pgfqpoint{3.450241in}{1.333917in}}%
\pgfpathcurveto{\pgfqpoint{3.454359in}{1.329799in}}{\pgfqpoint{3.459945in}{1.327485in}}{\pgfqpoint{3.465769in}{1.327485in}}%
\pgfpathlineto{\pgfqpoint{3.465769in}{1.327485in}}%
\pgfpathclose%
\pgfusepath{stroke,fill}%
\end{pgfscope}%
\begin{pgfscope}%
\pgfpathrectangle{\pgfqpoint{1.414736in}{0.548769in}}{\pgfqpoint{3.497747in}{7.253186in}}%
\pgfusepath{clip}%
\pgfsetbuttcap%
\pgfsetroundjoin%
\definecolor{currentfill}{rgb}{0.412913,0.803041,0.357269}%
\pgfsetfillcolor{currentfill}%
\pgfsetlinewidth{1.003750pt}%
\definecolor{currentstroke}{rgb}{0.412913,0.803041,0.357269}%
\pgfsetstrokecolor{currentstroke}%
\pgfsetdash{}{0pt}%
\pgfpathmoveto{\pgfqpoint{3.389139in}{1.798471in}}%
\pgfpathcurveto{\pgfqpoint{3.394963in}{1.798471in}}{\pgfqpoint{3.400549in}{1.800785in}}{\pgfqpoint{3.404668in}{1.804903in}}%
\pgfpathcurveto{\pgfqpoint{3.408786in}{1.809021in}}{\pgfqpoint{3.411100in}{1.814607in}}{\pgfqpoint{3.411100in}{1.820431in}}%
\pgfpathcurveto{\pgfqpoint{3.411100in}{1.826255in}}{\pgfqpoint{3.408786in}{1.831841in}}{\pgfqpoint{3.404668in}{1.835960in}}%
\pgfpathcurveto{\pgfqpoint{3.400549in}{1.840078in}}{\pgfqpoint{3.394963in}{1.842392in}}{\pgfqpoint{3.389139in}{1.842392in}}%
\pgfpathcurveto{\pgfqpoint{3.383315in}{1.842392in}}{\pgfqpoint{3.377729in}{1.840078in}}{\pgfqpoint{3.373611in}{1.835960in}}%
\pgfpathcurveto{\pgfqpoint{3.369493in}{1.831841in}}{\pgfqpoint{3.367179in}{1.826255in}}{\pgfqpoint{3.367179in}{1.820431in}}%
\pgfpathcurveto{\pgfqpoint{3.367179in}{1.814607in}}{\pgfqpoint{3.369493in}{1.809021in}}{\pgfqpoint{3.373611in}{1.804903in}}%
\pgfpathcurveto{\pgfqpoint{3.377729in}{1.800785in}}{\pgfqpoint{3.383315in}{1.798471in}}{\pgfqpoint{3.389139in}{1.798471in}}%
\pgfpathlineto{\pgfqpoint{3.389139in}{1.798471in}}%
\pgfpathclose%
\pgfusepath{stroke,fill}%
\end{pgfscope}%
\begin{pgfscope}%
\pgfpathrectangle{\pgfqpoint{1.414736in}{0.548769in}}{\pgfqpoint{3.497747in}{7.253186in}}%
\pgfusepath{clip}%
\pgfsetbuttcap%
\pgfsetroundjoin%
\definecolor{currentfill}{rgb}{0.993248,0.906157,0.143936}%
\pgfsetfillcolor{currentfill}%
\pgfsetlinewidth{1.003750pt}%
\definecolor{currentstroke}{rgb}{0.993248,0.906157,0.143936}%
\pgfsetstrokecolor{currentstroke}%
\pgfsetdash{}{0pt}%
\pgfpathmoveto{\pgfqpoint{3.222283in}{2.269457in}}%
\pgfpathcurveto{\pgfqpoint{3.228107in}{2.269457in}}{\pgfqpoint{3.233693in}{2.271771in}}{\pgfqpoint{3.237811in}{2.275889in}}%
\pgfpathcurveto{\pgfqpoint{3.241929in}{2.280007in}}{\pgfqpoint{3.244243in}{2.285593in}}{\pgfqpoint{3.244243in}{2.291417in}}%
\pgfpathcurveto{\pgfqpoint{3.244243in}{2.297241in}}{\pgfqpoint{3.241929in}{2.302828in}}{\pgfqpoint{3.237811in}{2.306946in}}%
\pgfpathcurveto{\pgfqpoint{3.233693in}{2.311064in}}{\pgfqpoint{3.228107in}{2.313378in}}{\pgfqpoint{3.222283in}{2.313378in}}%
\pgfpathcurveto{\pgfqpoint{3.216459in}{2.313378in}}{\pgfqpoint{3.210873in}{2.311064in}}{\pgfqpoint{3.206755in}{2.306946in}}%
\pgfpathcurveto{\pgfqpoint{3.202637in}{2.302828in}}{\pgfqpoint{3.200323in}{2.297241in}}{\pgfqpoint{3.200323in}{2.291417in}}%
\pgfpathcurveto{\pgfqpoint{3.200323in}{2.285593in}}{\pgfqpoint{3.202637in}{2.280007in}}{\pgfqpoint{3.206755in}{2.275889in}}%
\pgfpathcurveto{\pgfqpoint{3.210873in}{2.271771in}}{\pgfqpoint{3.216459in}{2.269457in}}{\pgfqpoint{3.222283in}{2.269457in}}%
\pgfpathlineto{\pgfqpoint{3.222283in}{2.269457in}}%
\pgfpathclose%
\pgfusepath{stroke,fill}%
\end{pgfscope}%
\begin{pgfscope}%
\pgfpathrectangle{\pgfqpoint{1.414736in}{0.548769in}}{\pgfqpoint{3.497747in}{7.253186in}}%
\pgfusepath{clip}%
\pgfsetbuttcap%
\pgfsetroundjoin%
\definecolor{currentfill}{rgb}{0.395174,0.797475,0.367757}%
\pgfsetfillcolor{currentfill}%
\pgfsetlinewidth{1.003750pt}%
\definecolor{currentstroke}{rgb}{0.395174,0.797475,0.367757}%
\pgfsetstrokecolor{currentstroke}%
\pgfsetdash{}{0pt}%
\pgfpathmoveto{\pgfqpoint{3.657615in}{2.740443in}}%
\pgfpathcurveto{\pgfqpoint{3.663439in}{2.740443in}}{\pgfqpoint{3.669025in}{2.742757in}}{\pgfqpoint{3.673143in}{2.746875in}}%
\pgfpathcurveto{\pgfqpoint{3.677261in}{2.750993in}}{\pgfqpoint{3.679575in}{2.756580in}}{\pgfqpoint{3.679575in}{2.762404in}}%
\pgfpathcurveto{\pgfqpoint{3.679575in}{2.768227in}}{\pgfqpoint{3.677261in}{2.773814in}}{\pgfqpoint{3.673143in}{2.777932in}}%
\pgfpathcurveto{\pgfqpoint{3.669025in}{2.782050in}}{\pgfqpoint{3.663439in}{2.784364in}}{\pgfqpoint{3.657615in}{2.784364in}}%
\pgfpathcurveto{\pgfqpoint{3.651791in}{2.784364in}}{\pgfqpoint{3.646205in}{2.782050in}}{\pgfqpoint{3.642087in}{2.777932in}}%
\pgfpathcurveto{\pgfqpoint{3.637969in}{2.773814in}}{\pgfqpoint{3.635655in}{2.768227in}}{\pgfqpoint{3.635655in}{2.762404in}}%
\pgfpathcurveto{\pgfqpoint{3.635655in}{2.756580in}}{\pgfqpoint{3.637969in}{2.750993in}}{\pgfqpoint{3.642087in}{2.746875in}}%
\pgfpathcurveto{\pgfqpoint{3.646205in}{2.742757in}}{\pgfqpoint{3.651791in}{2.740443in}}{\pgfqpoint{3.657615in}{2.740443in}}%
\pgfpathlineto{\pgfqpoint{3.657615in}{2.740443in}}%
\pgfpathclose%
\pgfusepath{stroke,fill}%
\end{pgfscope}%
\begin{pgfscope}%
\pgfpathrectangle{\pgfqpoint{1.414736in}{0.548769in}}{\pgfqpoint{3.497747in}{7.253186in}}%
\pgfusepath{clip}%
\pgfsetbuttcap%
\pgfsetroundjoin%
\definecolor{currentfill}{rgb}{0.616293,0.852709,0.230052}%
\pgfsetfillcolor{currentfill}%
\pgfsetlinewidth{1.003750pt}%
\definecolor{currentstroke}{rgb}{0.616293,0.852709,0.230052}%
\pgfsetstrokecolor{currentstroke}%
\pgfsetdash{}{0pt}%
\pgfpathmoveto{\pgfqpoint{3.432337in}{3.211429in}}%
\pgfpathcurveto{\pgfqpoint{3.438161in}{3.211429in}}{\pgfqpoint{3.443747in}{3.213743in}}{\pgfqpoint{3.447866in}{3.217861in}}%
\pgfpathcurveto{\pgfqpoint{3.451984in}{3.221980in}}{\pgfqpoint{3.454298in}{3.227566in}}{\pgfqpoint{3.454298in}{3.233390in}}%
\pgfpathcurveto{\pgfqpoint{3.454298in}{3.239214in}}{\pgfqpoint{3.451984in}{3.244800in}}{\pgfqpoint{3.447866in}{3.248918in}}%
\pgfpathcurveto{\pgfqpoint{3.443747in}{3.253036in}}{\pgfqpoint{3.438161in}{3.255350in}}{\pgfqpoint{3.432337in}{3.255350in}}%
\pgfpathcurveto{\pgfqpoint{3.426513in}{3.255350in}}{\pgfqpoint{3.420927in}{3.253036in}}{\pgfqpoint{3.416809in}{3.248918in}}%
\pgfpathcurveto{\pgfqpoint{3.412691in}{3.244800in}}{\pgfqpoint{3.410377in}{3.239214in}}{\pgfqpoint{3.410377in}{3.233390in}}%
\pgfpathcurveto{\pgfqpoint{3.410377in}{3.227566in}}{\pgfqpoint{3.412691in}{3.221980in}}{\pgfqpoint{3.416809in}{3.217861in}}%
\pgfpathcurveto{\pgfqpoint{3.420927in}{3.213743in}}{\pgfqpoint{3.426513in}{3.211429in}}{\pgfqpoint{3.432337in}{3.211429in}}%
\pgfpathlineto{\pgfqpoint{3.432337in}{3.211429in}}%
\pgfpathclose%
\pgfusepath{stroke,fill}%
\end{pgfscope}%
\begin{pgfscope}%
\pgfpathrectangle{\pgfqpoint{1.414736in}{0.548769in}}{\pgfqpoint{3.497747in}{7.253186in}}%
\pgfusepath{clip}%
\pgfsetbuttcap%
\pgfsetroundjoin%
\definecolor{currentfill}{rgb}{0.565498,0.842430,0.262877}%
\pgfsetfillcolor{currentfill}%
\pgfsetlinewidth{1.003750pt}%
\definecolor{currentstroke}{rgb}{0.565498,0.842430,0.262877}%
\pgfsetstrokecolor{currentstroke}%
\pgfsetdash{}{0pt}%
\pgfpathmoveto{\pgfqpoint{4.018851in}{3.682415in}}%
\pgfpathcurveto{\pgfqpoint{4.024675in}{3.682415in}}{\pgfqpoint{4.030261in}{3.684729in}}{\pgfqpoint{4.034379in}{3.688847in}}%
\pgfpathcurveto{\pgfqpoint{4.038497in}{3.692966in}}{\pgfqpoint{4.040811in}{3.698552in}}{\pgfqpoint{4.040811in}{3.704376in}}%
\pgfpathcurveto{\pgfqpoint{4.040811in}{3.710200in}}{\pgfqpoint{4.038497in}{3.715786in}}{\pgfqpoint{4.034379in}{3.719904in}}%
\pgfpathcurveto{\pgfqpoint{4.030261in}{3.724022in}}{\pgfqpoint{4.024675in}{3.726336in}}{\pgfqpoint{4.018851in}{3.726336in}}%
\pgfpathcurveto{\pgfqpoint{4.013027in}{3.726336in}}{\pgfqpoint{4.007441in}{3.724022in}}{\pgfqpoint{4.003322in}{3.719904in}}%
\pgfpathcurveto{\pgfqpoint{3.999204in}{3.715786in}}{\pgfqpoint{3.996890in}{3.710200in}}{\pgfqpoint{3.996890in}{3.704376in}}%
\pgfpathcurveto{\pgfqpoint{3.996890in}{3.698552in}}{\pgfqpoint{3.999204in}{3.692966in}}{\pgfqpoint{4.003322in}{3.688847in}}%
\pgfpathcurveto{\pgfqpoint{4.007441in}{3.684729in}}{\pgfqpoint{4.013027in}{3.682415in}}{\pgfqpoint{4.018851in}{3.682415in}}%
\pgfpathlineto{\pgfqpoint{4.018851in}{3.682415in}}%
\pgfpathclose%
\pgfusepath{stroke,fill}%
\end{pgfscope}%
\begin{pgfscope}%
\pgfpathrectangle{\pgfqpoint{1.414736in}{0.548769in}}{\pgfqpoint{3.497747in}{7.253186in}}%
\pgfusepath{clip}%
\pgfsetbuttcap%
\pgfsetroundjoin%
\definecolor{currentfill}{rgb}{0.993248,0.906157,0.143936}%
\pgfsetfillcolor{currentfill}%
\pgfsetlinewidth{1.003750pt}%
\definecolor{currentstroke}{rgb}{0.993248,0.906157,0.143936}%
\pgfsetstrokecolor{currentstroke}%
\pgfsetdash{}{0pt}%
\pgfpathmoveto{\pgfqpoint{3.186163in}{4.153402in}}%
\pgfpathcurveto{\pgfqpoint{3.191987in}{4.153402in}}{\pgfqpoint{3.197573in}{4.155715in}}{\pgfqpoint{3.201691in}{4.159834in}}%
\pgfpathcurveto{\pgfqpoint{3.205809in}{4.163952in}}{\pgfqpoint{3.208123in}{4.169538in}}{\pgfqpoint{3.208123in}{4.175362in}}%
\pgfpathcurveto{\pgfqpoint{3.208123in}{4.181186in}}{\pgfqpoint{3.205809in}{4.186772in}}{\pgfqpoint{3.201691in}{4.190890in}}%
\pgfpathcurveto{\pgfqpoint{3.197573in}{4.195008in}}{\pgfqpoint{3.191987in}{4.197322in}}{\pgfqpoint{3.186163in}{4.197322in}}%
\pgfpathcurveto{\pgfqpoint{3.180339in}{4.197322in}}{\pgfqpoint{3.174753in}{4.195008in}}{\pgfqpoint{3.170635in}{4.190890in}}%
\pgfpathcurveto{\pgfqpoint{3.166516in}{4.186772in}}{\pgfqpoint{3.164203in}{4.181186in}}{\pgfqpoint{3.164203in}{4.175362in}}%
\pgfpathcurveto{\pgfqpoint{3.164203in}{4.169538in}}{\pgfqpoint{3.166516in}{4.163952in}}{\pgfqpoint{3.170635in}{4.159834in}}%
\pgfpathcurveto{\pgfqpoint{3.174753in}{4.155715in}}{\pgfqpoint{3.180339in}{4.153402in}}{\pgfqpoint{3.186163in}{4.153402in}}%
\pgfpathlineto{\pgfqpoint{3.186163in}{4.153402in}}%
\pgfpathclose%
\pgfusepath{stroke,fill}%
\end{pgfscope}%
\begin{pgfscope}%
\pgfpathrectangle{\pgfqpoint{1.414736in}{0.548769in}}{\pgfqpoint{3.497747in}{7.253186in}}%
\pgfusepath{clip}%
\pgfsetbuttcap%
\pgfsetroundjoin%
\definecolor{currentfill}{rgb}{0.993248,0.906157,0.143936}%
\pgfsetfillcolor{currentfill}%
\pgfsetlinewidth{1.003750pt}%
\definecolor{currentstroke}{rgb}{0.993248,0.906157,0.143936}%
\pgfsetstrokecolor{currentstroke}%
\pgfsetdash{}{0pt}%
\pgfpathmoveto{\pgfqpoint{4.250001in}{4.624388in}}%
\pgfpathcurveto{\pgfqpoint{4.255825in}{4.624388in}}{\pgfqpoint{4.261411in}{4.626702in}}{\pgfqpoint{4.265529in}{4.630820in}}%
\pgfpathcurveto{\pgfqpoint{4.269647in}{4.634938in}}{\pgfqpoint{4.271961in}{4.640524in}}{\pgfqpoint{4.271961in}{4.646348in}}%
\pgfpathcurveto{\pgfqpoint{4.271961in}{4.652172in}}{\pgfqpoint{4.269647in}{4.657758in}}{\pgfqpoint{4.265529in}{4.661876in}}%
\pgfpathcurveto{\pgfqpoint{4.261411in}{4.665994in}}{\pgfqpoint{4.255825in}{4.668308in}}{\pgfqpoint{4.250001in}{4.668308in}}%
\pgfpathcurveto{\pgfqpoint{4.244177in}{4.668308in}}{\pgfqpoint{4.238591in}{4.665994in}}{\pgfqpoint{4.234473in}{4.661876in}}%
\pgfpathcurveto{\pgfqpoint{4.230354in}{4.657758in}}{\pgfqpoint{4.228041in}{4.652172in}}{\pgfqpoint{4.228041in}{4.646348in}}%
\pgfpathcurveto{\pgfqpoint{4.228041in}{4.640524in}}{\pgfqpoint{4.230354in}{4.634938in}}{\pgfqpoint{4.234473in}{4.630820in}}%
\pgfpathcurveto{\pgfqpoint{4.238591in}{4.626702in}}{\pgfqpoint{4.244177in}{4.624388in}}{\pgfqpoint{4.250001in}{4.624388in}}%
\pgfpathlineto{\pgfqpoint{4.250001in}{4.624388in}}%
\pgfpathclose%
\pgfusepath{stroke,fill}%
\end{pgfscope}%
\begin{pgfscope}%
\pgfpathrectangle{\pgfqpoint{1.414736in}{0.548769in}}{\pgfqpoint{3.497747in}{7.253186in}}%
\pgfusepath{clip}%
\pgfsetbuttcap%
\pgfsetroundjoin%
\definecolor{currentfill}{rgb}{0.232815,0.732247,0.459277}%
\pgfsetfillcolor{currentfill}%
\pgfsetlinewidth{1.003750pt}%
\definecolor{currentstroke}{rgb}{0.232815,0.732247,0.459277}%
\pgfsetstrokecolor{currentstroke}%
\pgfsetdash{}{0pt}%
\pgfpathmoveto{\pgfqpoint{4.411994in}{5.095374in}}%
\pgfpathcurveto{\pgfqpoint{4.417818in}{5.095374in}}{\pgfqpoint{4.423404in}{5.097688in}}{\pgfqpoint{4.427522in}{5.101806in}}%
\pgfpathcurveto{\pgfqpoint{4.431640in}{5.105924in}}{\pgfqpoint{4.433954in}{5.111510in}}{\pgfqpoint{4.433954in}{5.117334in}}%
\pgfpathcurveto{\pgfqpoint{4.433954in}{5.123158in}}{\pgfqpoint{4.431640in}{5.128744in}}{\pgfqpoint{4.427522in}{5.132862in}}%
\pgfpathcurveto{\pgfqpoint{4.423404in}{5.136980in}}{\pgfqpoint{4.417818in}{5.139294in}}{\pgfqpoint{4.411994in}{5.139294in}}%
\pgfpathcurveto{\pgfqpoint{4.406170in}{5.139294in}}{\pgfqpoint{4.400584in}{5.136980in}}{\pgfqpoint{4.396465in}{5.132862in}}%
\pgfpathcurveto{\pgfqpoint{4.392347in}{5.128744in}}{\pgfqpoint{4.390033in}{5.123158in}}{\pgfqpoint{4.390033in}{5.117334in}}%
\pgfpathcurveto{\pgfqpoint{4.390033in}{5.111510in}}{\pgfqpoint{4.392347in}{5.105924in}}{\pgfqpoint{4.396465in}{5.101806in}}%
\pgfpathcurveto{\pgfqpoint{4.400584in}{5.097688in}}{\pgfqpoint{4.406170in}{5.095374in}}{\pgfqpoint{4.411994in}{5.095374in}}%
\pgfpathlineto{\pgfqpoint{4.411994in}{5.095374in}}%
\pgfpathclose%
\pgfusepath{stroke,fill}%
\end{pgfscope}%
\begin{pgfscope}%
\pgfpathrectangle{\pgfqpoint{1.414736in}{0.548769in}}{\pgfqpoint{3.497747in}{7.253186in}}%
\pgfusepath{clip}%
\pgfsetbuttcap%
\pgfsetroundjoin%
\definecolor{currentfill}{rgb}{0.121148,0.592739,0.544641}%
\pgfsetfillcolor{currentfill}%
\pgfsetlinewidth{1.003750pt}%
\definecolor{currentstroke}{rgb}{0.121148,0.592739,0.544641}%
\pgfsetstrokecolor{currentstroke}%
\pgfsetdash{}{0pt}%
\pgfpathmoveto{\pgfqpoint{3.197235in}{5.566360in}}%
\pgfpathcurveto{\pgfqpoint{3.203059in}{5.566360in}}{\pgfqpoint{3.208645in}{5.568674in}}{\pgfqpoint{3.212763in}{5.572792in}}%
\pgfpathcurveto{\pgfqpoint{3.216882in}{5.576910in}}{\pgfqpoint{3.219195in}{5.582496in}}{\pgfqpoint{3.219195in}{5.588320in}}%
\pgfpathcurveto{\pgfqpoint{3.219195in}{5.594144in}}{\pgfqpoint{3.216882in}{5.599730in}}{\pgfqpoint{3.212763in}{5.603848in}}%
\pgfpathcurveto{\pgfqpoint{3.208645in}{5.607967in}}{\pgfqpoint{3.203059in}{5.610280in}}{\pgfqpoint{3.197235in}{5.610280in}}%
\pgfpathcurveto{\pgfqpoint{3.191411in}{5.610280in}}{\pgfqpoint{3.185825in}{5.607967in}}{\pgfqpoint{3.181707in}{5.603848in}}%
\pgfpathcurveto{\pgfqpoint{3.177589in}{5.599730in}}{\pgfqpoint{3.175275in}{5.594144in}}{\pgfqpoint{3.175275in}{5.588320in}}%
\pgfpathcurveto{\pgfqpoint{3.175275in}{5.582496in}}{\pgfqpoint{3.177589in}{5.576910in}}{\pgfqpoint{3.181707in}{5.572792in}}%
\pgfpathcurveto{\pgfqpoint{3.185825in}{5.568674in}}{\pgfqpoint{3.191411in}{5.566360in}}{\pgfqpoint{3.197235in}{5.566360in}}%
\pgfpathlineto{\pgfqpoint{3.197235in}{5.566360in}}%
\pgfpathclose%
\pgfusepath{stroke,fill}%
\end{pgfscope}%
\begin{pgfscope}%
\pgfpathrectangle{\pgfqpoint{1.414736in}{0.548769in}}{\pgfqpoint{3.497747in}{7.253186in}}%
\pgfusepath{clip}%
\pgfsetbuttcap%
\pgfsetroundjoin%
\definecolor{currentfill}{rgb}{0.814576,0.883393,0.110347}%
\pgfsetfillcolor{currentfill}%
\pgfsetlinewidth{1.003750pt}%
\definecolor{currentstroke}{rgb}{0.814576,0.883393,0.110347}%
\pgfsetstrokecolor{currentstroke}%
\pgfsetdash{}{0pt}%
\pgfpathmoveto{\pgfqpoint{2.849303in}{6.037346in}}%
\pgfpathcurveto{\pgfqpoint{2.855127in}{6.037346in}}{\pgfqpoint{2.860713in}{6.039660in}}{\pgfqpoint{2.864831in}{6.043778in}}%
\pgfpathcurveto{\pgfqpoint{2.868950in}{6.047896in}}{\pgfqpoint{2.871263in}{6.053482in}}{\pgfqpoint{2.871263in}{6.059306in}}%
\pgfpathcurveto{\pgfqpoint{2.871263in}{6.065130in}}{\pgfqpoint{2.868950in}{6.070716in}}{\pgfqpoint{2.864831in}{6.074835in}}%
\pgfpathcurveto{\pgfqpoint{2.860713in}{6.078953in}}{\pgfqpoint{2.855127in}{6.081267in}}{\pgfqpoint{2.849303in}{6.081267in}}%
\pgfpathcurveto{\pgfqpoint{2.843479in}{6.081267in}}{\pgfqpoint{2.837893in}{6.078953in}}{\pgfqpoint{2.833775in}{6.074835in}}%
\pgfpathcurveto{\pgfqpoint{2.829657in}{6.070716in}}{\pgfqpoint{2.827343in}{6.065130in}}{\pgfqpoint{2.827343in}{6.059306in}}%
\pgfpathcurveto{\pgfqpoint{2.827343in}{6.053482in}}{\pgfqpoint{2.829657in}{6.047896in}}{\pgfqpoint{2.833775in}{6.043778in}}%
\pgfpathcurveto{\pgfqpoint{2.837893in}{6.039660in}}{\pgfqpoint{2.843479in}{6.037346in}}{\pgfqpoint{2.849303in}{6.037346in}}%
\pgfpathlineto{\pgfqpoint{2.849303in}{6.037346in}}%
\pgfpathclose%
\pgfusepath{stroke,fill}%
\end{pgfscope}%
\begin{pgfscope}%
\pgfpathrectangle{\pgfqpoint{1.414736in}{0.548769in}}{\pgfqpoint{3.497747in}{7.253186in}}%
\pgfusepath{clip}%
\pgfsetbuttcap%
\pgfsetroundjoin%
\definecolor{currentfill}{rgb}{0.993248,0.906157,0.143936}%
\pgfsetfillcolor{currentfill}%
\pgfsetlinewidth{1.003750pt}%
\definecolor{currentstroke}{rgb}{0.993248,0.906157,0.143936}%
\pgfsetstrokecolor{currentstroke}%
\pgfsetdash{}{0pt}%
\pgfpathmoveto{\pgfqpoint{3.645878in}{6.508332in}}%
\pgfpathcurveto{\pgfqpoint{3.651702in}{6.508332in}}{\pgfqpoint{3.657288in}{6.510646in}}{\pgfqpoint{3.661406in}{6.514764in}}%
\pgfpathcurveto{\pgfqpoint{3.665524in}{6.518882in}}{\pgfqpoint{3.667838in}{6.524468in}}{\pgfqpoint{3.667838in}{6.530292in}}%
\pgfpathcurveto{\pgfqpoint{3.667838in}{6.536116in}}{\pgfqpoint{3.665524in}{6.541702in}}{\pgfqpoint{3.661406in}{6.545821in}}%
\pgfpathcurveto{\pgfqpoint{3.657288in}{6.549939in}}{\pgfqpoint{3.651702in}{6.552253in}}{\pgfqpoint{3.645878in}{6.552253in}}%
\pgfpathcurveto{\pgfqpoint{3.640054in}{6.552253in}}{\pgfqpoint{3.634468in}{6.549939in}}{\pgfqpoint{3.630350in}{6.545821in}}%
\pgfpathcurveto{\pgfqpoint{3.626232in}{6.541702in}}{\pgfqpoint{3.623918in}{6.536116in}}{\pgfqpoint{3.623918in}{6.530292in}}%
\pgfpathcurveto{\pgfqpoint{3.623918in}{6.524468in}}{\pgfqpoint{3.626232in}{6.518882in}}{\pgfqpoint{3.630350in}{6.514764in}}%
\pgfpathcurveto{\pgfqpoint{3.634468in}{6.510646in}}{\pgfqpoint{3.640054in}{6.508332in}}{\pgfqpoint{3.645878in}{6.508332in}}%
\pgfpathlineto{\pgfqpoint{3.645878in}{6.508332in}}%
\pgfpathclose%
\pgfusepath{stroke,fill}%
\end{pgfscope}%
\begin{pgfscope}%
\pgfpathrectangle{\pgfqpoint{1.414736in}{0.548769in}}{\pgfqpoint{3.497747in}{7.253186in}}%
\pgfusepath{clip}%
\pgfsetbuttcap%
\pgfsetroundjoin%
\definecolor{currentfill}{rgb}{0.835270,0.886029,0.102646}%
\pgfsetfillcolor{currentfill}%
\pgfsetlinewidth{1.003750pt}%
\definecolor{currentstroke}{rgb}{0.835270,0.886029,0.102646}%
\pgfsetstrokecolor{currentstroke}%
\pgfsetdash{}{0pt}%
\pgfpathmoveto{\pgfqpoint{2.879580in}{6.979318in}}%
\pgfpathcurveto{\pgfqpoint{2.885404in}{6.979318in}}{\pgfqpoint{2.890990in}{6.981632in}}{\pgfqpoint{2.895109in}{6.985750in}}%
\pgfpathcurveto{\pgfqpoint{2.899227in}{6.989868in}}{\pgfqpoint{2.901541in}{6.995455in}}{\pgfqpoint{2.901541in}{7.001278in}}%
\pgfpathcurveto{\pgfqpoint{2.901541in}{7.007102in}}{\pgfqpoint{2.899227in}{7.012689in}}{\pgfqpoint{2.895109in}{7.016807in}}%
\pgfpathcurveto{\pgfqpoint{2.890990in}{7.020925in}}{\pgfqpoint{2.885404in}{7.023239in}}{\pgfqpoint{2.879580in}{7.023239in}}%
\pgfpathcurveto{\pgfqpoint{2.873756in}{7.023239in}}{\pgfqpoint{2.868170in}{7.020925in}}{\pgfqpoint{2.864052in}{7.016807in}}%
\pgfpathcurveto{\pgfqpoint{2.859934in}{7.012689in}}{\pgfqpoint{2.857620in}{7.007102in}}{\pgfqpoint{2.857620in}{7.001278in}}%
\pgfpathcurveto{\pgfqpoint{2.857620in}{6.995455in}}{\pgfqpoint{2.859934in}{6.989868in}}{\pgfqpoint{2.864052in}{6.985750in}}%
\pgfpathcurveto{\pgfqpoint{2.868170in}{6.981632in}}{\pgfqpoint{2.873756in}{6.979318in}}{\pgfqpoint{2.879580in}{6.979318in}}%
\pgfpathlineto{\pgfqpoint{2.879580in}{6.979318in}}%
\pgfpathclose%
\pgfusepath{stroke,fill}%
\end{pgfscope}%
\begin{pgfscope}%
\pgfpathrectangle{\pgfqpoint{1.414736in}{0.548769in}}{\pgfqpoint{3.497747in}{7.253186in}}%
\pgfusepath{clip}%
\pgfsetbuttcap%
\pgfsetroundjoin%
\definecolor{currentfill}{rgb}{0.993248,0.906157,0.143936}%
\pgfsetfillcolor{currentfill}%
\pgfsetlinewidth{1.003750pt}%
\definecolor{currentstroke}{rgb}{0.993248,0.906157,0.143936}%
\pgfsetstrokecolor{currentstroke}%
\pgfsetdash{}{0pt}%
\pgfpathmoveto{\pgfqpoint{3.293687in}{7.450304in}}%
\pgfpathcurveto{\pgfqpoint{3.299511in}{7.450304in}}{\pgfqpoint{3.305097in}{7.452618in}}{\pgfqpoint{3.309215in}{7.456736in}}%
\pgfpathcurveto{\pgfqpoint{3.313333in}{7.460854in}}{\pgfqpoint{3.315647in}{7.466441in}}{\pgfqpoint{3.315647in}{7.472265in}}%
\pgfpathcurveto{\pgfqpoint{3.315647in}{7.478089in}}{\pgfqpoint{3.313333in}{7.483675in}}{\pgfqpoint{3.309215in}{7.487793in}}%
\pgfpathcurveto{\pgfqpoint{3.305097in}{7.491911in}}{\pgfqpoint{3.299511in}{7.494225in}}{\pgfqpoint{3.293687in}{7.494225in}}%
\pgfpathcurveto{\pgfqpoint{3.287863in}{7.494225in}}{\pgfqpoint{3.282277in}{7.491911in}}{\pgfqpoint{3.278159in}{7.487793in}}%
\pgfpathcurveto{\pgfqpoint{3.274041in}{7.483675in}}{\pgfqpoint{3.271727in}{7.478089in}}{\pgfqpoint{3.271727in}{7.472265in}}%
\pgfpathcurveto{\pgfqpoint{3.271727in}{7.466441in}}{\pgfqpoint{3.274041in}{7.460854in}}{\pgfqpoint{3.278159in}{7.456736in}}%
\pgfpathcurveto{\pgfqpoint{3.282277in}{7.452618in}}{\pgfqpoint{3.287863in}{7.450304in}}{\pgfqpoint{3.293687in}{7.450304in}}%
\pgfpathlineto{\pgfqpoint{3.293687in}{7.450304in}}%
\pgfpathclose%
\pgfusepath{stroke,fill}%
\end{pgfscope}%
\begin{pgfscope}%
\pgfpathrectangle{\pgfqpoint{1.414736in}{0.548769in}}{\pgfqpoint{3.497747in}{7.253186in}}%
\pgfusepath{clip}%
\pgfsetbuttcap%
\pgfsetroundjoin%
\definecolor{currentfill}{rgb}{0.983868,0.904867,0.136897}%
\pgfsetfillcolor{currentfill}%
\pgfsetlinewidth{1.003750pt}%
\definecolor{currentstroke}{rgb}{0.983868,0.904867,0.136897}%
\pgfsetstrokecolor{currentstroke}%
\pgfsetdash{}{0pt}%
\pgfpathmoveto{\pgfqpoint{3.415435in}{0.856499in}}%
\pgfpathcurveto{\pgfqpoint{3.421259in}{0.856499in}}{\pgfqpoint{3.426846in}{0.858813in}}{\pgfqpoint{3.430964in}{0.862931in}}%
\pgfpathcurveto{\pgfqpoint{3.435082in}{0.867049in}}{\pgfqpoint{3.437396in}{0.872635in}}{\pgfqpoint{3.437396in}{0.878459in}}%
\pgfpathcurveto{\pgfqpoint{3.437396in}{0.884283in}}{\pgfqpoint{3.435082in}{0.889869in}}{\pgfqpoint{3.430964in}{0.893987in}}%
\pgfpathcurveto{\pgfqpoint{3.426846in}{0.898105in}}{\pgfqpoint{3.421259in}{0.900419in}}{\pgfqpoint{3.415435in}{0.900419in}}%
\pgfpathcurveto{\pgfqpoint{3.409612in}{0.900419in}}{\pgfqpoint{3.404025in}{0.898105in}}{\pgfqpoint{3.399907in}{0.893987in}}%
\pgfpathcurveto{\pgfqpoint{3.395789in}{0.889869in}}{\pgfqpoint{3.393475in}{0.884283in}}{\pgfqpoint{3.393475in}{0.878459in}}%
\pgfpathcurveto{\pgfqpoint{3.393475in}{0.872635in}}{\pgfqpoint{3.395789in}{0.867049in}}{\pgfqpoint{3.399907in}{0.862931in}}%
\pgfpathcurveto{\pgfqpoint{3.404025in}{0.858813in}}{\pgfqpoint{3.409612in}{0.856499in}}{\pgfqpoint{3.415435in}{0.856499in}}%
\pgfpathlineto{\pgfqpoint{3.415435in}{0.856499in}}%
\pgfpathclose%
\pgfusepath{stroke,fill}%
\end{pgfscope}%
\begin{pgfscope}%
\pgfpathrectangle{\pgfqpoint{1.414736in}{0.548769in}}{\pgfqpoint{3.497747in}{7.253186in}}%
\pgfusepath{clip}%
\pgfsetbuttcap%
\pgfsetroundjoin%
\definecolor{currentfill}{rgb}{0.412913,0.803041,0.357269}%
\pgfsetfillcolor{currentfill}%
\pgfsetlinewidth{1.003750pt}%
\definecolor{currentstroke}{rgb}{0.412913,0.803041,0.357269}%
\pgfsetstrokecolor{currentstroke}%
\pgfsetdash{}{0pt}%
\pgfpathmoveto{\pgfqpoint{3.784933in}{1.327485in}}%
\pgfpathcurveto{\pgfqpoint{3.790757in}{1.327485in}}{\pgfqpoint{3.796343in}{1.329799in}}{\pgfqpoint{3.800461in}{1.333917in}}%
\pgfpathcurveto{\pgfqpoint{3.804579in}{1.338035in}}{\pgfqpoint{3.806893in}{1.343621in}}{\pgfqpoint{3.806893in}{1.349445in}}%
\pgfpathcurveto{\pgfqpoint{3.806893in}{1.355269in}}{\pgfqpoint{3.804579in}{1.360855in}}{\pgfqpoint{3.800461in}{1.364973in}}%
\pgfpathcurveto{\pgfqpoint{3.796343in}{1.369092in}}{\pgfqpoint{3.790757in}{1.371405in}}{\pgfqpoint{3.784933in}{1.371405in}}%
\pgfpathcurveto{\pgfqpoint{3.779109in}{1.371405in}}{\pgfqpoint{3.773523in}{1.369092in}}{\pgfqpoint{3.769404in}{1.364973in}}%
\pgfpathcurveto{\pgfqpoint{3.765286in}{1.360855in}}{\pgfqpoint{3.762972in}{1.355269in}}{\pgfqpoint{3.762972in}{1.349445in}}%
\pgfpathcurveto{\pgfqpoint{3.762972in}{1.343621in}}{\pgfqpoint{3.765286in}{1.338035in}}{\pgfqpoint{3.769404in}{1.333917in}}%
\pgfpathcurveto{\pgfqpoint{3.773523in}{1.329799in}}{\pgfqpoint{3.779109in}{1.327485in}}{\pgfqpoint{3.784933in}{1.327485in}}%
\pgfpathlineto{\pgfqpoint{3.784933in}{1.327485in}}%
\pgfpathclose%
\pgfusepath{stroke,fill}%
\end{pgfscope}%
\begin{pgfscope}%
\pgfpathrectangle{\pgfqpoint{1.414736in}{0.548769in}}{\pgfqpoint{3.497747in}{7.253186in}}%
\pgfusepath{clip}%
\pgfsetbuttcap%
\pgfsetroundjoin%
\definecolor{currentfill}{rgb}{0.449368,0.813768,0.335384}%
\pgfsetfillcolor{currentfill}%
\pgfsetlinewidth{1.003750pt}%
\definecolor{currentstroke}{rgb}{0.449368,0.813768,0.335384}%
\pgfsetstrokecolor{currentstroke}%
\pgfsetdash{}{0pt}%
\pgfpathmoveto{\pgfqpoint{3.355077in}{1.798471in}}%
\pgfpathcurveto{\pgfqpoint{3.360901in}{1.798471in}}{\pgfqpoint{3.366487in}{1.800785in}}{\pgfqpoint{3.370605in}{1.804903in}}%
\pgfpathcurveto{\pgfqpoint{3.374723in}{1.809021in}}{\pgfqpoint{3.377037in}{1.814607in}}{\pgfqpoint{3.377037in}{1.820431in}}%
\pgfpathcurveto{\pgfqpoint{3.377037in}{1.826255in}}{\pgfqpoint{3.374723in}{1.831841in}}{\pgfqpoint{3.370605in}{1.835960in}}%
\pgfpathcurveto{\pgfqpoint{3.366487in}{1.840078in}}{\pgfqpoint{3.360901in}{1.842392in}}{\pgfqpoint{3.355077in}{1.842392in}}%
\pgfpathcurveto{\pgfqpoint{3.349253in}{1.842392in}}{\pgfqpoint{3.343667in}{1.840078in}}{\pgfqpoint{3.339549in}{1.835960in}}%
\pgfpathcurveto{\pgfqpoint{3.335431in}{1.831841in}}{\pgfqpoint{3.333117in}{1.826255in}}{\pgfqpoint{3.333117in}{1.820431in}}%
\pgfpathcurveto{\pgfqpoint{3.333117in}{1.814607in}}{\pgfqpoint{3.335431in}{1.809021in}}{\pgfqpoint{3.339549in}{1.804903in}}%
\pgfpathcurveto{\pgfqpoint{3.343667in}{1.800785in}}{\pgfqpoint{3.349253in}{1.798471in}}{\pgfqpoint{3.355077in}{1.798471in}}%
\pgfpathlineto{\pgfqpoint{3.355077in}{1.798471in}}%
\pgfpathclose%
\pgfusepath{stroke,fill}%
\end{pgfscope}%
\begin{pgfscope}%
\pgfpathrectangle{\pgfqpoint{1.414736in}{0.548769in}}{\pgfqpoint{3.497747in}{7.253186in}}%
\pgfusepath{clip}%
\pgfsetbuttcap%
\pgfsetroundjoin%
\definecolor{currentfill}{rgb}{0.824940,0.884720,0.106217}%
\pgfsetfillcolor{currentfill}%
\pgfsetlinewidth{1.003750pt}%
\definecolor{currentstroke}{rgb}{0.824940,0.884720,0.106217}%
\pgfsetstrokecolor{currentstroke}%
\pgfsetdash{}{0pt}%
\pgfpathmoveto{\pgfqpoint{3.383906in}{2.269457in}}%
\pgfpathcurveto{\pgfqpoint{3.389730in}{2.269457in}}{\pgfqpoint{3.395316in}{2.271771in}}{\pgfqpoint{3.399435in}{2.275889in}}%
\pgfpathcurveto{\pgfqpoint{3.403553in}{2.280007in}}{\pgfqpoint{3.405867in}{2.285593in}}{\pgfqpoint{3.405867in}{2.291417in}}%
\pgfpathcurveto{\pgfqpoint{3.405867in}{2.297241in}}{\pgfqpoint{3.403553in}{2.302828in}}{\pgfqpoint{3.399435in}{2.306946in}}%
\pgfpathcurveto{\pgfqpoint{3.395316in}{2.311064in}}{\pgfqpoint{3.389730in}{2.313378in}}{\pgfqpoint{3.383906in}{2.313378in}}%
\pgfpathcurveto{\pgfqpoint{3.378082in}{2.313378in}}{\pgfqpoint{3.372496in}{2.311064in}}{\pgfqpoint{3.368378in}{2.306946in}}%
\pgfpathcurveto{\pgfqpoint{3.364260in}{2.302828in}}{\pgfqpoint{3.361946in}{2.297241in}}{\pgfqpoint{3.361946in}{2.291417in}}%
\pgfpathcurveto{\pgfqpoint{3.361946in}{2.285593in}}{\pgfqpoint{3.364260in}{2.280007in}}{\pgfqpoint{3.368378in}{2.275889in}}%
\pgfpathcurveto{\pgfqpoint{3.372496in}{2.271771in}}{\pgfqpoint{3.378082in}{2.269457in}}{\pgfqpoint{3.383906in}{2.269457in}}%
\pgfpathlineto{\pgfqpoint{3.383906in}{2.269457in}}%
\pgfpathclose%
\pgfusepath{stroke,fill}%
\end{pgfscope}%
\begin{pgfscope}%
\pgfpathrectangle{\pgfqpoint{1.414736in}{0.548769in}}{\pgfqpoint{3.497747in}{7.253186in}}%
\pgfusepath{clip}%
\pgfsetbuttcap%
\pgfsetroundjoin%
\definecolor{currentfill}{rgb}{0.208030,0.718701,0.472873}%
\pgfsetfillcolor{currentfill}%
\pgfsetlinewidth{1.003750pt}%
\definecolor{currentstroke}{rgb}{0.208030,0.718701,0.472873}%
\pgfsetstrokecolor{currentstroke}%
\pgfsetdash{}{0pt}%
\pgfpathmoveto{\pgfqpoint{4.000063in}{2.740443in}}%
\pgfpathcurveto{\pgfqpoint{4.005887in}{2.740443in}}{\pgfqpoint{4.011473in}{2.742757in}}{\pgfqpoint{4.015591in}{2.746875in}}%
\pgfpathcurveto{\pgfqpoint{4.019709in}{2.750993in}}{\pgfqpoint{4.022023in}{2.756580in}}{\pgfqpoint{4.022023in}{2.762404in}}%
\pgfpathcurveto{\pgfqpoint{4.022023in}{2.768227in}}{\pgfqpoint{4.019709in}{2.773814in}}{\pgfqpoint{4.015591in}{2.777932in}}%
\pgfpathcurveto{\pgfqpoint{4.011473in}{2.782050in}}{\pgfqpoint{4.005887in}{2.784364in}}{\pgfqpoint{4.000063in}{2.784364in}}%
\pgfpathcurveto{\pgfqpoint{3.994239in}{2.784364in}}{\pgfqpoint{3.988653in}{2.782050in}}{\pgfqpoint{3.984535in}{2.777932in}}%
\pgfpathcurveto{\pgfqpoint{3.980416in}{2.773814in}}{\pgfqpoint{3.978103in}{2.768227in}}{\pgfqpoint{3.978103in}{2.762404in}}%
\pgfpathcurveto{\pgfqpoint{3.978103in}{2.756580in}}{\pgfqpoint{3.980416in}{2.750993in}}{\pgfqpoint{3.984535in}{2.746875in}}%
\pgfpathcurveto{\pgfqpoint{3.988653in}{2.742757in}}{\pgfqpoint{3.994239in}{2.740443in}}{\pgfqpoint{4.000063in}{2.740443in}}%
\pgfpathlineto{\pgfqpoint{4.000063in}{2.740443in}}%
\pgfpathclose%
\pgfusepath{stroke,fill}%
\end{pgfscope}%
\begin{pgfscope}%
\pgfpathrectangle{\pgfqpoint{1.414736in}{0.548769in}}{\pgfqpoint{3.497747in}{7.253186in}}%
\pgfusepath{clip}%
\pgfsetbuttcap%
\pgfsetroundjoin%
\definecolor{currentfill}{rgb}{0.525776,0.833491,0.288127}%
\pgfsetfillcolor{currentfill}%
\pgfsetlinewidth{1.003750pt}%
\definecolor{currentstroke}{rgb}{0.525776,0.833491,0.288127}%
\pgfsetstrokecolor{currentstroke}%
\pgfsetdash{}{0pt}%
\pgfpathmoveto{\pgfqpoint{3.388646in}{3.211429in}}%
\pgfpathcurveto{\pgfqpoint{3.394470in}{3.211429in}}{\pgfqpoint{3.400056in}{3.213743in}}{\pgfqpoint{3.404175in}{3.217861in}}%
\pgfpathcurveto{\pgfqpoint{3.408293in}{3.221980in}}{\pgfqpoint{3.410607in}{3.227566in}}{\pgfqpoint{3.410607in}{3.233390in}}%
\pgfpathcurveto{\pgfqpoint{3.410607in}{3.239214in}}{\pgfqpoint{3.408293in}{3.244800in}}{\pgfqpoint{3.404175in}{3.248918in}}%
\pgfpathcurveto{\pgfqpoint{3.400056in}{3.253036in}}{\pgfqpoint{3.394470in}{3.255350in}}{\pgfqpoint{3.388646in}{3.255350in}}%
\pgfpathcurveto{\pgfqpoint{3.382822in}{3.255350in}}{\pgfqpoint{3.377236in}{3.253036in}}{\pgfqpoint{3.373118in}{3.248918in}}%
\pgfpathcurveto{\pgfqpoint{3.369000in}{3.244800in}}{\pgfqpoint{3.366686in}{3.239214in}}{\pgfqpoint{3.366686in}{3.233390in}}%
\pgfpathcurveto{\pgfqpoint{3.366686in}{3.227566in}}{\pgfqpoint{3.369000in}{3.221980in}}{\pgfqpoint{3.373118in}{3.217861in}}%
\pgfpathcurveto{\pgfqpoint{3.377236in}{3.213743in}}{\pgfqpoint{3.382822in}{3.211429in}}{\pgfqpoint{3.388646in}{3.211429in}}%
\pgfpathlineto{\pgfqpoint{3.388646in}{3.211429in}}%
\pgfpathclose%
\pgfusepath{stroke,fill}%
\end{pgfscope}%
\begin{pgfscope}%
\pgfpathrectangle{\pgfqpoint{1.414736in}{0.548769in}}{\pgfqpoint{3.497747in}{7.253186in}}%
\pgfusepath{clip}%
\pgfsetbuttcap%
\pgfsetroundjoin%
\definecolor{currentfill}{rgb}{0.606045,0.850733,0.236712}%
\pgfsetfillcolor{currentfill}%
\pgfsetlinewidth{1.003750pt}%
\definecolor{currentstroke}{rgb}{0.606045,0.850733,0.236712}%
\pgfsetstrokecolor{currentstroke}%
\pgfsetdash{}{0pt}%
\pgfpathmoveto{\pgfqpoint{3.707989in}{3.682415in}}%
\pgfpathcurveto{\pgfqpoint{3.713813in}{3.682415in}}{\pgfqpoint{3.719399in}{3.684729in}}{\pgfqpoint{3.723518in}{3.688847in}}%
\pgfpathcurveto{\pgfqpoint{3.727636in}{3.692966in}}{\pgfqpoint{3.729950in}{3.698552in}}{\pgfqpoint{3.729950in}{3.704376in}}%
\pgfpathcurveto{\pgfqpoint{3.729950in}{3.710200in}}{\pgfqpoint{3.727636in}{3.715786in}}{\pgfqpoint{3.723518in}{3.719904in}}%
\pgfpathcurveto{\pgfqpoint{3.719399in}{3.724022in}}{\pgfqpoint{3.713813in}{3.726336in}}{\pgfqpoint{3.707989in}{3.726336in}}%
\pgfpathcurveto{\pgfqpoint{3.702165in}{3.726336in}}{\pgfqpoint{3.696579in}{3.724022in}}{\pgfqpoint{3.692461in}{3.719904in}}%
\pgfpathcurveto{\pgfqpoint{3.688343in}{3.715786in}}{\pgfqpoint{3.686029in}{3.710200in}}{\pgfqpoint{3.686029in}{3.704376in}}%
\pgfpathcurveto{\pgfqpoint{3.686029in}{3.698552in}}{\pgfqpoint{3.688343in}{3.692966in}}{\pgfqpoint{3.692461in}{3.688847in}}%
\pgfpathcurveto{\pgfqpoint{3.696579in}{3.684729in}}{\pgfqpoint{3.702165in}{3.682415in}}{\pgfqpoint{3.707989in}{3.682415in}}%
\pgfpathlineto{\pgfqpoint{3.707989in}{3.682415in}}%
\pgfpathclose%
\pgfusepath{stroke,fill}%
\end{pgfscope}%
\begin{pgfscope}%
\pgfpathrectangle{\pgfqpoint{1.414736in}{0.548769in}}{\pgfqpoint{3.497747in}{7.253186in}}%
\pgfusepath{clip}%
\pgfsetbuttcap%
\pgfsetroundjoin%
\definecolor{currentfill}{rgb}{0.974417,0.903590,0.130215}%
\pgfsetfillcolor{currentfill}%
\pgfsetlinewidth{1.003750pt}%
\definecolor{currentstroke}{rgb}{0.974417,0.903590,0.130215}%
\pgfsetstrokecolor{currentstroke}%
\pgfsetdash{}{0pt}%
\pgfpathmoveto{\pgfqpoint{3.559090in}{4.153402in}}%
\pgfpathcurveto{\pgfqpoint{3.564914in}{4.153402in}}{\pgfqpoint{3.570500in}{4.155715in}}{\pgfqpoint{3.574618in}{4.159834in}}%
\pgfpathcurveto{\pgfqpoint{3.578736in}{4.163952in}}{\pgfqpoint{3.581050in}{4.169538in}}{\pgfqpoint{3.581050in}{4.175362in}}%
\pgfpathcurveto{\pgfqpoint{3.581050in}{4.181186in}}{\pgfqpoint{3.578736in}{4.186772in}}{\pgfqpoint{3.574618in}{4.190890in}}%
\pgfpathcurveto{\pgfqpoint{3.570500in}{4.195008in}}{\pgfqpoint{3.564914in}{4.197322in}}{\pgfqpoint{3.559090in}{4.197322in}}%
\pgfpathcurveto{\pgfqpoint{3.553266in}{4.197322in}}{\pgfqpoint{3.547680in}{4.195008in}}{\pgfqpoint{3.543561in}{4.190890in}}%
\pgfpathcurveto{\pgfqpoint{3.539443in}{4.186772in}}{\pgfqpoint{3.537129in}{4.181186in}}{\pgfqpoint{3.537129in}{4.175362in}}%
\pgfpathcurveto{\pgfqpoint{3.537129in}{4.169538in}}{\pgfqpoint{3.539443in}{4.163952in}}{\pgfqpoint{3.543561in}{4.159834in}}%
\pgfpathcurveto{\pgfqpoint{3.547680in}{4.155715in}}{\pgfqpoint{3.553266in}{4.153402in}}{\pgfqpoint{3.559090in}{4.153402in}}%
\pgfpathlineto{\pgfqpoint{3.559090in}{4.153402in}}%
\pgfpathclose%
\pgfusepath{stroke,fill}%
\end{pgfscope}%
\begin{pgfscope}%
\pgfpathrectangle{\pgfqpoint{1.414736in}{0.548769in}}{\pgfqpoint{3.497747in}{7.253186in}}%
\pgfusepath{clip}%
\pgfsetbuttcap%
\pgfsetroundjoin%
\definecolor{currentfill}{rgb}{0.657642,0.860219,0.203082}%
\pgfsetfillcolor{currentfill}%
\pgfsetlinewidth{1.003750pt}%
\definecolor{currentstroke}{rgb}{0.657642,0.860219,0.203082}%
\pgfsetstrokecolor{currentstroke}%
\pgfsetdash{}{0pt}%
\pgfpathmoveto{\pgfqpoint{3.089784in}{4.624388in}}%
\pgfpathcurveto{\pgfqpoint{3.095608in}{4.624388in}}{\pgfqpoint{3.101194in}{4.626702in}}{\pgfqpoint{3.105312in}{4.630820in}}%
\pgfpathcurveto{\pgfqpoint{3.109430in}{4.634938in}}{\pgfqpoint{3.111744in}{4.640524in}}{\pgfqpoint{3.111744in}{4.646348in}}%
\pgfpathcurveto{\pgfqpoint{3.111744in}{4.652172in}}{\pgfqpoint{3.109430in}{4.657758in}}{\pgfqpoint{3.105312in}{4.661876in}}%
\pgfpathcurveto{\pgfqpoint{3.101194in}{4.665994in}}{\pgfqpoint{3.095608in}{4.668308in}}{\pgfqpoint{3.089784in}{4.668308in}}%
\pgfpathcurveto{\pgfqpoint{3.083960in}{4.668308in}}{\pgfqpoint{3.078374in}{4.665994in}}{\pgfqpoint{3.074255in}{4.661876in}}%
\pgfpathcurveto{\pgfqpoint{3.070137in}{4.657758in}}{\pgfqpoint{3.067823in}{4.652172in}}{\pgfqpoint{3.067823in}{4.646348in}}%
\pgfpathcurveto{\pgfqpoint{3.067823in}{4.640524in}}{\pgfqpoint{3.070137in}{4.634938in}}{\pgfqpoint{3.074255in}{4.630820in}}%
\pgfpathcurveto{\pgfqpoint{3.078374in}{4.626702in}}{\pgfqpoint{3.083960in}{4.624388in}}{\pgfqpoint{3.089784in}{4.624388in}}%
\pgfpathlineto{\pgfqpoint{3.089784in}{4.624388in}}%
\pgfpathclose%
\pgfusepath{stroke,fill}%
\end{pgfscope}%
\begin{pgfscope}%
\pgfpathrectangle{\pgfqpoint{1.414736in}{0.548769in}}{\pgfqpoint{3.497747in}{7.253186in}}%
\pgfusepath{clip}%
\pgfsetbuttcap%
\pgfsetroundjoin%
\definecolor{currentfill}{rgb}{0.252194,0.269783,0.531579}%
\pgfsetfillcolor{currentfill}%
\pgfsetlinewidth{1.003750pt}%
\definecolor{currentstroke}{rgb}{0.252194,0.269783,0.531579}%
\pgfsetstrokecolor{currentstroke}%
\pgfsetdash{}{0pt}%
\pgfpathmoveto{\pgfqpoint{3.277010in}{5.095374in}}%
\pgfpathcurveto{\pgfqpoint{3.282834in}{5.095374in}}{\pgfqpoint{3.288421in}{5.097688in}}{\pgfqpoint{3.292539in}{5.101806in}}%
\pgfpathcurveto{\pgfqpoint{3.296657in}{5.105924in}}{\pgfqpoint{3.298971in}{5.111510in}}{\pgfqpoint{3.298971in}{5.117334in}}%
\pgfpathcurveto{\pgfqpoint{3.298971in}{5.123158in}}{\pgfqpoint{3.296657in}{5.128744in}}{\pgfqpoint{3.292539in}{5.132862in}}%
\pgfpathcurveto{\pgfqpoint{3.288421in}{5.136980in}}{\pgfqpoint{3.282834in}{5.139294in}}{\pgfqpoint{3.277010in}{5.139294in}}%
\pgfpathcurveto{\pgfqpoint{3.271186in}{5.139294in}}{\pgfqpoint{3.265600in}{5.136980in}}{\pgfqpoint{3.261482in}{5.132862in}}%
\pgfpathcurveto{\pgfqpoint{3.257364in}{5.128744in}}{\pgfqpoint{3.255050in}{5.123158in}}{\pgfqpoint{3.255050in}{5.117334in}}%
\pgfpathcurveto{\pgfqpoint{3.255050in}{5.111510in}}{\pgfqpoint{3.257364in}{5.105924in}}{\pgfqpoint{3.261482in}{5.101806in}}%
\pgfpathcurveto{\pgfqpoint{3.265600in}{5.097688in}}{\pgfqpoint{3.271186in}{5.095374in}}{\pgfqpoint{3.277010in}{5.095374in}}%
\pgfpathlineto{\pgfqpoint{3.277010in}{5.095374in}}%
\pgfpathclose%
\pgfusepath{stroke,fill}%
\end{pgfscope}%
\begin{pgfscope}%
\pgfpathrectangle{\pgfqpoint{1.414736in}{0.548769in}}{\pgfqpoint{3.497747in}{7.253186in}}%
\pgfusepath{clip}%
\pgfsetbuttcap%
\pgfsetroundjoin%
\definecolor{currentfill}{rgb}{0.151918,0.500685,0.557587}%
\pgfsetfillcolor{currentfill}%
\pgfsetlinewidth{1.003750pt}%
\definecolor{currentstroke}{rgb}{0.151918,0.500685,0.557587}%
\pgfsetstrokecolor{currentstroke}%
\pgfsetdash{}{0pt}%
\pgfpathmoveto{\pgfqpoint{3.521231in}{5.566360in}}%
\pgfpathcurveto{\pgfqpoint{3.527055in}{5.566360in}}{\pgfqpoint{3.532641in}{5.568674in}}{\pgfqpoint{3.536759in}{5.572792in}}%
\pgfpathcurveto{\pgfqpoint{3.540877in}{5.576910in}}{\pgfqpoint{3.543191in}{5.582496in}}{\pgfqpoint{3.543191in}{5.588320in}}%
\pgfpathcurveto{\pgfqpoint{3.543191in}{5.594144in}}{\pgfqpoint{3.540877in}{5.599730in}}{\pgfqpoint{3.536759in}{5.603848in}}%
\pgfpathcurveto{\pgfqpoint{3.532641in}{5.607967in}}{\pgfqpoint{3.527055in}{5.610280in}}{\pgfqpoint{3.521231in}{5.610280in}}%
\pgfpathcurveto{\pgfqpoint{3.515407in}{5.610280in}}{\pgfqpoint{3.509821in}{5.607967in}}{\pgfqpoint{3.505702in}{5.603848in}}%
\pgfpathcurveto{\pgfqpoint{3.501584in}{5.599730in}}{\pgfqpoint{3.499270in}{5.594144in}}{\pgfqpoint{3.499270in}{5.588320in}}%
\pgfpathcurveto{\pgfqpoint{3.499270in}{5.582496in}}{\pgfqpoint{3.501584in}{5.576910in}}{\pgfqpoint{3.505702in}{5.572792in}}%
\pgfpathcurveto{\pgfqpoint{3.509821in}{5.568674in}}{\pgfqpoint{3.515407in}{5.566360in}}{\pgfqpoint{3.521231in}{5.566360in}}%
\pgfpathlineto{\pgfqpoint{3.521231in}{5.566360in}}%
\pgfpathclose%
\pgfusepath{stroke,fill}%
\end{pgfscope}%
\begin{pgfscope}%
\pgfpathrectangle{\pgfqpoint{1.414736in}{0.548769in}}{\pgfqpoint{3.497747in}{7.253186in}}%
\pgfusepath{clip}%
\pgfsetbuttcap%
\pgfsetroundjoin%
\definecolor{currentfill}{rgb}{0.130067,0.651384,0.521608}%
\pgfsetfillcolor{currentfill}%
\pgfsetlinewidth{1.003750pt}%
\definecolor{currentstroke}{rgb}{0.130067,0.651384,0.521608}%
\pgfsetstrokecolor{currentstroke}%
\pgfsetdash{}{0pt}%
\pgfpathmoveto{\pgfqpoint{4.108793in}{6.037346in}}%
\pgfpathcurveto{\pgfqpoint{4.114617in}{6.037346in}}{\pgfqpoint{4.120203in}{6.039660in}}{\pgfqpoint{4.124321in}{6.043778in}}%
\pgfpathcurveto{\pgfqpoint{4.128439in}{6.047896in}}{\pgfqpoint{4.130753in}{6.053482in}}{\pgfqpoint{4.130753in}{6.059306in}}%
\pgfpathcurveto{\pgfqpoint{4.130753in}{6.065130in}}{\pgfqpoint{4.128439in}{6.070716in}}{\pgfqpoint{4.124321in}{6.074835in}}%
\pgfpathcurveto{\pgfqpoint{4.120203in}{6.078953in}}{\pgfqpoint{4.114617in}{6.081267in}}{\pgfqpoint{4.108793in}{6.081267in}}%
\pgfpathcurveto{\pgfqpoint{4.102969in}{6.081267in}}{\pgfqpoint{4.097383in}{6.078953in}}{\pgfqpoint{4.093265in}{6.074835in}}%
\pgfpathcurveto{\pgfqpoint{4.089147in}{6.070716in}}{\pgfqpoint{4.086833in}{6.065130in}}{\pgfqpoint{4.086833in}{6.059306in}}%
\pgfpathcurveto{\pgfqpoint{4.086833in}{6.053482in}}{\pgfqpoint{4.089147in}{6.047896in}}{\pgfqpoint{4.093265in}{6.043778in}}%
\pgfpathcurveto{\pgfqpoint{4.097383in}{6.039660in}}{\pgfqpoint{4.102969in}{6.037346in}}{\pgfqpoint{4.108793in}{6.037346in}}%
\pgfpathlineto{\pgfqpoint{4.108793in}{6.037346in}}%
\pgfpathclose%
\pgfusepath{stroke,fill}%
\end{pgfscope}%
\begin{pgfscope}%
\pgfpathrectangle{\pgfqpoint{1.414736in}{0.548769in}}{\pgfqpoint{3.497747in}{7.253186in}}%
\pgfusepath{clip}%
\pgfsetbuttcap%
\pgfsetroundjoin%
\definecolor{currentfill}{rgb}{0.993248,0.906157,0.143936}%
\pgfsetfillcolor{currentfill}%
\pgfsetlinewidth{1.003750pt}%
\definecolor{currentstroke}{rgb}{0.993248,0.906157,0.143936}%
\pgfsetstrokecolor{currentstroke}%
\pgfsetdash{}{0pt}%
\pgfpathmoveto{\pgfqpoint{3.913904in}{6.508332in}}%
\pgfpathcurveto{\pgfqpoint{3.919728in}{6.508332in}}{\pgfqpoint{3.925314in}{6.510646in}}{\pgfqpoint{3.929432in}{6.514764in}}%
\pgfpathcurveto{\pgfqpoint{3.933550in}{6.518882in}}{\pgfqpoint{3.935864in}{6.524468in}}{\pgfqpoint{3.935864in}{6.530292in}}%
\pgfpathcurveto{\pgfqpoint{3.935864in}{6.536116in}}{\pgfqpoint{3.933550in}{6.541702in}}{\pgfqpoint{3.929432in}{6.545821in}}%
\pgfpathcurveto{\pgfqpoint{3.925314in}{6.549939in}}{\pgfqpoint{3.919728in}{6.552253in}}{\pgfqpoint{3.913904in}{6.552253in}}%
\pgfpathcurveto{\pgfqpoint{3.908080in}{6.552253in}}{\pgfqpoint{3.902494in}{6.549939in}}{\pgfqpoint{3.898376in}{6.545821in}}%
\pgfpathcurveto{\pgfqpoint{3.894258in}{6.541702in}}{\pgfqpoint{3.891944in}{6.536116in}}{\pgfqpoint{3.891944in}{6.530292in}}%
\pgfpathcurveto{\pgfqpoint{3.891944in}{6.524468in}}{\pgfqpoint{3.894258in}{6.518882in}}{\pgfqpoint{3.898376in}{6.514764in}}%
\pgfpathcurveto{\pgfqpoint{3.902494in}{6.510646in}}{\pgfqpoint{3.908080in}{6.508332in}}{\pgfqpoint{3.913904in}{6.508332in}}%
\pgfpathlineto{\pgfqpoint{3.913904in}{6.508332in}}%
\pgfpathclose%
\pgfusepath{stroke,fill}%
\end{pgfscope}%
\begin{pgfscope}%
\pgfpathrectangle{\pgfqpoint{1.414736in}{0.548769in}}{\pgfqpoint{3.497747in}{7.253186in}}%
\pgfusepath{clip}%
\pgfsetbuttcap%
\pgfsetroundjoin%
\definecolor{currentfill}{rgb}{0.657642,0.860219,0.203082}%
\pgfsetfillcolor{currentfill}%
\pgfsetlinewidth{1.003750pt}%
\definecolor{currentstroke}{rgb}{0.657642,0.860219,0.203082}%
\pgfsetstrokecolor{currentstroke}%
\pgfsetdash{}{0pt}%
\pgfpathmoveto{\pgfqpoint{3.894589in}{6.979318in}}%
\pgfpathcurveto{\pgfqpoint{3.900413in}{6.979318in}}{\pgfqpoint{3.905999in}{6.981632in}}{\pgfqpoint{3.910117in}{6.985750in}}%
\pgfpathcurveto{\pgfqpoint{3.914236in}{6.989868in}}{\pgfqpoint{3.916549in}{6.995455in}}{\pgfqpoint{3.916549in}{7.001278in}}%
\pgfpathcurveto{\pgfqpoint{3.916549in}{7.007102in}}{\pgfqpoint{3.914236in}{7.012689in}}{\pgfqpoint{3.910117in}{7.016807in}}%
\pgfpathcurveto{\pgfqpoint{3.905999in}{7.020925in}}{\pgfqpoint{3.900413in}{7.023239in}}{\pgfqpoint{3.894589in}{7.023239in}}%
\pgfpathcurveto{\pgfqpoint{3.888765in}{7.023239in}}{\pgfqpoint{3.883179in}{7.020925in}}{\pgfqpoint{3.879061in}{7.016807in}}%
\pgfpathcurveto{\pgfqpoint{3.874943in}{7.012689in}}{\pgfqpoint{3.872629in}{7.007102in}}{\pgfqpoint{3.872629in}{7.001278in}}%
\pgfpathcurveto{\pgfqpoint{3.872629in}{6.995455in}}{\pgfqpoint{3.874943in}{6.989868in}}{\pgfqpoint{3.879061in}{6.985750in}}%
\pgfpathcurveto{\pgfqpoint{3.883179in}{6.981632in}}{\pgfqpoint{3.888765in}{6.979318in}}{\pgfqpoint{3.894589in}{6.979318in}}%
\pgfpathlineto{\pgfqpoint{3.894589in}{6.979318in}}%
\pgfpathclose%
\pgfusepath{stroke,fill}%
\end{pgfscope}%
\begin{pgfscope}%
\pgfpathrectangle{\pgfqpoint{1.414736in}{0.548769in}}{\pgfqpoint{3.497747in}{7.253186in}}%
\pgfusepath{clip}%
\pgfsetbuttcap%
\pgfsetroundjoin%
\definecolor{currentfill}{rgb}{0.259857,0.745492,0.444467}%
\pgfsetfillcolor{currentfill}%
\pgfsetlinewidth{1.003750pt}%
\definecolor{currentstroke}{rgb}{0.259857,0.745492,0.444467}%
\pgfsetstrokecolor{currentstroke}%
\pgfsetdash{}{0pt}%
\pgfpathmoveto{\pgfqpoint{3.862969in}{7.450304in}}%
\pgfpathcurveto{\pgfqpoint{3.868793in}{7.450304in}}{\pgfqpoint{3.874379in}{7.452618in}}{\pgfqpoint{3.878497in}{7.456736in}}%
\pgfpathcurveto{\pgfqpoint{3.882616in}{7.460854in}}{\pgfqpoint{3.884929in}{7.466441in}}{\pgfqpoint{3.884929in}{7.472265in}}%
\pgfpathcurveto{\pgfqpoint{3.884929in}{7.478089in}}{\pgfqpoint{3.882616in}{7.483675in}}{\pgfqpoint{3.878497in}{7.487793in}}%
\pgfpathcurveto{\pgfqpoint{3.874379in}{7.491911in}}{\pgfqpoint{3.868793in}{7.494225in}}{\pgfqpoint{3.862969in}{7.494225in}}%
\pgfpathcurveto{\pgfqpoint{3.857145in}{7.494225in}}{\pgfqpoint{3.851559in}{7.491911in}}{\pgfqpoint{3.847441in}{7.487793in}}%
\pgfpathcurveto{\pgfqpoint{3.843323in}{7.483675in}}{\pgfqpoint{3.841009in}{7.478089in}}{\pgfqpoint{3.841009in}{7.472265in}}%
\pgfpathcurveto{\pgfqpoint{3.841009in}{7.466441in}}{\pgfqpoint{3.843323in}{7.460854in}}{\pgfqpoint{3.847441in}{7.456736in}}%
\pgfpathcurveto{\pgfqpoint{3.851559in}{7.452618in}}{\pgfqpoint{3.857145in}{7.450304in}}{\pgfqpoint{3.862969in}{7.450304in}}%
\pgfpathlineto{\pgfqpoint{3.862969in}{7.450304in}}%
\pgfpathclose%
\pgfusepath{stroke,fill}%
\end{pgfscope}%
\begin{pgfscope}%
\pgfpathrectangle{\pgfqpoint{1.414736in}{0.548769in}}{\pgfqpoint{3.497747in}{7.253186in}}%
\pgfusepath{clip}%
\pgfsetbuttcap%
\pgfsetroundjoin%
\definecolor{currentfill}{rgb}{0.993248,0.906157,0.143936}%
\pgfsetfillcolor{currentfill}%
\pgfsetlinewidth{1.003750pt}%
\definecolor{currentstroke}{rgb}{0.993248,0.906157,0.143936}%
\pgfsetstrokecolor{currentstroke}%
\pgfsetdash{}{0pt}%
\pgfpathmoveto{\pgfqpoint{3.682713in}{0.856499in}}%
\pgfpathcurveto{\pgfqpoint{3.688537in}{0.856499in}}{\pgfqpoint{3.694123in}{0.858813in}}{\pgfqpoint{3.698241in}{0.862931in}}%
\pgfpathcurveto{\pgfqpoint{3.702359in}{0.867049in}}{\pgfqpoint{3.704673in}{0.872635in}}{\pgfqpoint{3.704673in}{0.878459in}}%
\pgfpathcurveto{\pgfqpoint{3.704673in}{0.884283in}}{\pgfqpoint{3.702359in}{0.889869in}}{\pgfqpoint{3.698241in}{0.893987in}}%
\pgfpathcurveto{\pgfqpoint{3.694123in}{0.898105in}}{\pgfqpoint{3.688537in}{0.900419in}}{\pgfqpoint{3.682713in}{0.900419in}}%
\pgfpathcurveto{\pgfqpoint{3.676889in}{0.900419in}}{\pgfqpoint{3.671303in}{0.898105in}}{\pgfqpoint{3.667185in}{0.893987in}}%
\pgfpathcurveto{\pgfqpoint{3.663066in}{0.889869in}}{\pgfqpoint{3.660753in}{0.884283in}}{\pgfqpoint{3.660753in}{0.878459in}}%
\pgfpathcurveto{\pgfqpoint{3.660753in}{0.872635in}}{\pgfqpoint{3.663066in}{0.867049in}}{\pgfqpoint{3.667185in}{0.862931in}}%
\pgfpathcurveto{\pgfqpoint{3.671303in}{0.858813in}}{\pgfqpoint{3.676889in}{0.856499in}}{\pgfqpoint{3.682713in}{0.856499in}}%
\pgfpathlineto{\pgfqpoint{3.682713in}{0.856499in}}%
\pgfpathclose%
\pgfusepath{stroke,fill}%
\end{pgfscope}%
\begin{pgfscope}%
\pgfpathrectangle{\pgfqpoint{1.414736in}{0.548769in}}{\pgfqpoint{3.497747in}{7.253186in}}%
\pgfusepath{clip}%
\pgfsetbuttcap%
\pgfsetroundjoin%
\definecolor{currentfill}{rgb}{0.616293,0.852709,0.230052}%
\pgfsetfillcolor{currentfill}%
\pgfsetlinewidth{1.003750pt}%
\definecolor{currentstroke}{rgb}{0.616293,0.852709,0.230052}%
\pgfsetstrokecolor{currentstroke}%
\pgfsetdash{}{0pt}%
\pgfpathmoveto{\pgfqpoint{3.566333in}{1.327485in}}%
\pgfpathcurveto{\pgfqpoint{3.572157in}{1.327485in}}{\pgfqpoint{3.577743in}{1.329799in}}{\pgfqpoint{3.581861in}{1.333917in}}%
\pgfpathcurveto{\pgfqpoint{3.585979in}{1.338035in}}{\pgfqpoint{3.588293in}{1.343621in}}{\pgfqpoint{3.588293in}{1.349445in}}%
\pgfpathcurveto{\pgfqpoint{3.588293in}{1.355269in}}{\pgfqpoint{3.585979in}{1.360855in}}{\pgfqpoint{3.581861in}{1.364973in}}%
\pgfpathcurveto{\pgfqpoint{3.577743in}{1.369092in}}{\pgfqpoint{3.572157in}{1.371405in}}{\pgfqpoint{3.566333in}{1.371405in}}%
\pgfpathcurveto{\pgfqpoint{3.560509in}{1.371405in}}{\pgfqpoint{3.554923in}{1.369092in}}{\pgfqpoint{3.550805in}{1.364973in}}%
\pgfpathcurveto{\pgfqpoint{3.546686in}{1.360855in}}{\pgfqpoint{3.544373in}{1.355269in}}{\pgfqpoint{3.544373in}{1.349445in}}%
\pgfpathcurveto{\pgfqpoint{3.544373in}{1.343621in}}{\pgfqpoint{3.546686in}{1.338035in}}{\pgfqpoint{3.550805in}{1.333917in}}%
\pgfpathcurveto{\pgfqpoint{3.554923in}{1.329799in}}{\pgfqpoint{3.560509in}{1.327485in}}{\pgfqpoint{3.566333in}{1.327485in}}%
\pgfpathlineto{\pgfqpoint{3.566333in}{1.327485in}}%
\pgfpathclose%
\pgfusepath{stroke,fill}%
\end{pgfscope}%
\begin{pgfscope}%
\pgfpathrectangle{\pgfqpoint{1.414736in}{0.548769in}}{\pgfqpoint{3.497747in}{7.253186in}}%
\pgfusepath{clip}%
\pgfsetbuttcap%
\pgfsetroundjoin%
\definecolor{currentfill}{rgb}{0.449368,0.813768,0.335384}%
\pgfsetfillcolor{currentfill}%
\pgfsetlinewidth{1.003750pt}%
\definecolor{currentstroke}{rgb}{0.449368,0.813768,0.335384}%
\pgfsetstrokecolor{currentstroke}%
\pgfsetdash{}{0pt}%
\pgfpathmoveto{\pgfqpoint{3.420600in}{1.798471in}}%
\pgfpathcurveto{\pgfqpoint{3.426424in}{1.798471in}}{\pgfqpoint{3.432010in}{1.800785in}}{\pgfqpoint{3.436128in}{1.804903in}}%
\pgfpathcurveto{\pgfqpoint{3.440246in}{1.809021in}}{\pgfqpoint{3.442560in}{1.814607in}}{\pgfqpoint{3.442560in}{1.820431in}}%
\pgfpathcurveto{\pgfqpoint{3.442560in}{1.826255in}}{\pgfqpoint{3.440246in}{1.831841in}}{\pgfqpoint{3.436128in}{1.835960in}}%
\pgfpathcurveto{\pgfqpoint{3.432010in}{1.840078in}}{\pgfqpoint{3.426424in}{1.842392in}}{\pgfqpoint{3.420600in}{1.842392in}}%
\pgfpathcurveto{\pgfqpoint{3.414776in}{1.842392in}}{\pgfqpoint{3.409190in}{1.840078in}}{\pgfqpoint{3.405072in}{1.835960in}}%
\pgfpathcurveto{\pgfqpoint{3.400954in}{1.831841in}}{\pgfqpoint{3.398640in}{1.826255in}}{\pgfqpoint{3.398640in}{1.820431in}}%
\pgfpathcurveto{\pgfqpoint{3.398640in}{1.814607in}}{\pgfqpoint{3.400954in}{1.809021in}}{\pgfqpoint{3.405072in}{1.804903in}}%
\pgfpathcurveto{\pgfqpoint{3.409190in}{1.800785in}}{\pgfqpoint{3.414776in}{1.798471in}}{\pgfqpoint{3.420600in}{1.798471in}}%
\pgfpathlineto{\pgfqpoint{3.420600in}{1.798471in}}%
\pgfpathclose%
\pgfusepath{stroke,fill}%
\end{pgfscope}%
\begin{pgfscope}%
\pgfpathrectangle{\pgfqpoint{1.414736in}{0.548769in}}{\pgfqpoint{3.497747in}{7.253186in}}%
\pgfusepath{clip}%
\pgfsetbuttcap%
\pgfsetroundjoin%
\definecolor{currentfill}{rgb}{0.993248,0.906157,0.143936}%
\pgfsetfillcolor{currentfill}%
\pgfsetlinewidth{1.003750pt}%
\definecolor{currentstroke}{rgb}{0.993248,0.906157,0.143936}%
\pgfsetstrokecolor{currentstroke}%
\pgfsetdash{}{0pt}%
\pgfpathmoveto{\pgfqpoint{3.250295in}{2.269457in}}%
\pgfpathcurveto{\pgfqpoint{3.256118in}{2.269457in}}{\pgfqpoint{3.261705in}{2.271771in}}{\pgfqpoint{3.265823in}{2.275889in}}%
\pgfpathcurveto{\pgfqpoint{3.269941in}{2.280007in}}{\pgfqpoint{3.272255in}{2.285593in}}{\pgfqpoint{3.272255in}{2.291417in}}%
\pgfpathcurveto{\pgfqpoint{3.272255in}{2.297241in}}{\pgfqpoint{3.269941in}{2.302828in}}{\pgfqpoint{3.265823in}{2.306946in}}%
\pgfpathcurveto{\pgfqpoint{3.261705in}{2.311064in}}{\pgfqpoint{3.256118in}{2.313378in}}{\pgfqpoint{3.250295in}{2.313378in}}%
\pgfpathcurveto{\pgfqpoint{3.244471in}{2.313378in}}{\pgfqpoint{3.238884in}{2.311064in}}{\pgfqpoint{3.234766in}{2.306946in}}%
\pgfpathcurveto{\pgfqpoint{3.230648in}{2.302828in}}{\pgfqpoint{3.228334in}{2.297241in}}{\pgfqpoint{3.228334in}{2.291417in}}%
\pgfpathcurveto{\pgfqpoint{3.228334in}{2.285593in}}{\pgfqpoint{3.230648in}{2.280007in}}{\pgfqpoint{3.234766in}{2.275889in}}%
\pgfpathcurveto{\pgfqpoint{3.238884in}{2.271771in}}{\pgfqpoint{3.244471in}{2.269457in}}{\pgfqpoint{3.250295in}{2.269457in}}%
\pgfpathlineto{\pgfqpoint{3.250295in}{2.269457in}}%
\pgfpathclose%
\pgfusepath{stroke,fill}%
\end{pgfscope}%
\begin{pgfscope}%
\pgfpathrectangle{\pgfqpoint{1.414736in}{0.548769in}}{\pgfqpoint{3.497747in}{7.253186in}}%
\pgfusepath{clip}%
\pgfsetbuttcap%
\pgfsetroundjoin%
\definecolor{currentfill}{rgb}{0.266941,0.748751,0.440573}%
\pgfsetfillcolor{currentfill}%
\pgfsetlinewidth{1.003750pt}%
\definecolor{currentstroke}{rgb}{0.266941,0.748751,0.440573}%
\pgfsetstrokecolor{currentstroke}%
\pgfsetdash{}{0pt}%
\pgfpathmoveto{\pgfqpoint{3.872610in}{2.740443in}}%
\pgfpathcurveto{\pgfqpoint{3.878434in}{2.740443in}}{\pgfqpoint{3.884020in}{2.742757in}}{\pgfqpoint{3.888138in}{2.746875in}}%
\pgfpathcurveto{\pgfqpoint{3.892257in}{2.750993in}}{\pgfqpoint{3.894570in}{2.756580in}}{\pgfqpoint{3.894570in}{2.762404in}}%
\pgfpathcurveto{\pgfqpoint{3.894570in}{2.768227in}}{\pgfqpoint{3.892257in}{2.773814in}}{\pgfqpoint{3.888138in}{2.777932in}}%
\pgfpathcurveto{\pgfqpoint{3.884020in}{2.782050in}}{\pgfqpoint{3.878434in}{2.784364in}}{\pgfqpoint{3.872610in}{2.784364in}}%
\pgfpathcurveto{\pgfqpoint{3.866786in}{2.784364in}}{\pgfqpoint{3.861200in}{2.782050in}}{\pgfqpoint{3.857082in}{2.777932in}}%
\pgfpathcurveto{\pgfqpoint{3.852964in}{2.773814in}}{\pgfqpoint{3.850650in}{2.768227in}}{\pgfqpoint{3.850650in}{2.762404in}}%
\pgfpathcurveto{\pgfqpoint{3.850650in}{2.756580in}}{\pgfqpoint{3.852964in}{2.750993in}}{\pgfqpoint{3.857082in}{2.746875in}}%
\pgfpathcurveto{\pgfqpoint{3.861200in}{2.742757in}}{\pgfqpoint{3.866786in}{2.740443in}}{\pgfqpoint{3.872610in}{2.740443in}}%
\pgfpathlineto{\pgfqpoint{3.872610in}{2.740443in}}%
\pgfpathclose%
\pgfusepath{stroke,fill}%
\end{pgfscope}%
\begin{pgfscope}%
\pgfpathrectangle{\pgfqpoint{1.414736in}{0.548769in}}{\pgfqpoint{3.497747in}{7.253186in}}%
\pgfusepath{clip}%
\pgfsetbuttcap%
\pgfsetroundjoin%
\definecolor{currentfill}{rgb}{0.657642,0.860219,0.203082}%
\pgfsetfillcolor{currentfill}%
\pgfsetlinewidth{1.003750pt}%
\definecolor{currentstroke}{rgb}{0.657642,0.860219,0.203082}%
\pgfsetstrokecolor{currentstroke}%
\pgfsetdash{}{0pt}%
\pgfpathmoveto{\pgfqpoint{3.329228in}{3.211429in}}%
\pgfpathcurveto{\pgfqpoint{3.335052in}{3.211429in}}{\pgfqpoint{3.340638in}{3.213743in}}{\pgfqpoint{3.344756in}{3.217861in}}%
\pgfpathcurveto{\pgfqpoint{3.348874in}{3.221980in}}{\pgfqpoint{3.351188in}{3.227566in}}{\pgfqpoint{3.351188in}{3.233390in}}%
\pgfpathcurveto{\pgfqpoint{3.351188in}{3.239214in}}{\pgfqpoint{3.348874in}{3.244800in}}{\pgfqpoint{3.344756in}{3.248918in}}%
\pgfpathcurveto{\pgfqpoint{3.340638in}{3.253036in}}{\pgfqpoint{3.335052in}{3.255350in}}{\pgfqpoint{3.329228in}{3.255350in}}%
\pgfpathcurveto{\pgfqpoint{3.323404in}{3.255350in}}{\pgfqpoint{3.317818in}{3.253036in}}{\pgfqpoint{3.313700in}{3.248918in}}%
\pgfpathcurveto{\pgfqpoint{3.309582in}{3.244800in}}{\pgfqpoint{3.307268in}{3.239214in}}{\pgfqpoint{3.307268in}{3.233390in}}%
\pgfpathcurveto{\pgfqpoint{3.307268in}{3.227566in}}{\pgfqpoint{3.309582in}{3.221980in}}{\pgfqpoint{3.313700in}{3.217861in}}%
\pgfpathcurveto{\pgfqpoint{3.317818in}{3.213743in}}{\pgfqpoint{3.323404in}{3.211429in}}{\pgfqpoint{3.329228in}{3.211429in}}%
\pgfpathlineto{\pgfqpoint{3.329228in}{3.211429in}}%
\pgfpathclose%
\pgfusepath{stroke,fill}%
\end{pgfscope}%
\begin{pgfscope}%
\pgfpathrectangle{\pgfqpoint{1.414736in}{0.548769in}}{\pgfqpoint{3.497747in}{7.253186in}}%
\pgfusepath{clip}%
\pgfsetbuttcap%
\pgfsetroundjoin%
\definecolor{currentfill}{rgb}{0.845561,0.887322,0.099702}%
\pgfsetfillcolor{currentfill}%
\pgfsetlinewidth{1.003750pt}%
\definecolor{currentstroke}{rgb}{0.845561,0.887322,0.099702}%
\pgfsetstrokecolor{currentstroke}%
\pgfsetdash{}{0pt}%
\pgfpathmoveto{\pgfqpoint{3.349016in}{3.682415in}}%
\pgfpathcurveto{\pgfqpoint{3.354840in}{3.682415in}}{\pgfqpoint{3.360426in}{3.684729in}}{\pgfqpoint{3.364544in}{3.688847in}}%
\pgfpathcurveto{\pgfqpoint{3.368662in}{3.692966in}}{\pgfqpoint{3.370976in}{3.698552in}}{\pgfqpoint{3.370976in}{3.704376in}}%
\pgfpathcurveto{\pgfqpoint{3.370976in}{3.710200in}}{\pgfqpoint{3.368662in}{3.715786in}}{\pgfqpoint{3.364544in}{3.719904in}}%
\pgfpathcurveto{\pgfqpoint{3.360426in}{3.724022in}}{\pgfqpoint{3.354840in}{3.726336in}}{\pgfqpoint{3.349016in}{3.726336in}}%
\pgfpathcurveto{\pgfqpoint{3.343192in}{3.726336in}}{\pgfqpoint{3.337606in}{3.724022in}}{\pgfqpoint{3.333488in}{3.719904in}}%
\pgfpathcurveto{\pgfqpoint{3.329370in}{3.715786in}}{\pgfqpoint{3.327056in}{3.710200in}}{\pgfqpoint{3.327056in}{3.704376in}}%
\pgfpathcurveto{\pgfqpoint{3.327056in}{3.698552in}}{\pgfqpoint{3.329370in}{3.692966in}}{\pgfqpoint{3.333488in}{3.688847in}}%
\pgfpathcurveto{\pgfqpoint{3.337606in}{3.684729in}}{\pgfqpoint{3.343192in}{3.682415in}}{\pgfqpoint{3.349016in}{3.682415in}}%
\pgfpathlineto{\pgfqpoint{3.349016in}{3.682415in}}%
\pgfpathclose%
\pgfusepath{stroke,fill}%
\end{pgfscope}%
\begin{pgfscope}%
\pgfpathrectangle{\pgfqpoint{1.414736in}{0.548769in}}{\pgfqpoint{3.497747in}{7.253186in}}%
\pgfusepath{clip}%
\pgfsetbuttcap%
\pgfsetroundjoin%
\definecolor{currentfill}{rgb}{0.993248,0.906157,0.143936}%
\pgfsetfillcolor{currentfill}%
\pgfsetlinewidth{1.003750pt}%
\definecolor{currentstroke}{rgb}{0.993248,0.906157,0.143936}%
\pgfsetstrokecolor{currentstroke}%
\pgfsetdash{}{0pt}%
\pgfpathmoveto{\pgfqpoint{3.075813in}{4.153402in}}%
\pgfpathcurveto{\pgfqpoint{3.081637in}{4.153402in}}{\pgfqpoint{3.087223in}{4.155715in}}{\pgfqpoint{3.091341in}{4.159834in}}%
\pgfpathcurveto{\pgfqpoint{3.095460in}{4.163952in}}{\pgfqpoint{3.097773in}{4.169538in}}{\pgfqpoint{3.097773in}{4.175362in}}%
\pgfpathcurveto{\pgfqpoint{3.097773in}{4.181186in}}{\pgfqpoint{3.095460in}{4.186772in}}{\pgfqpoint{3.091341in}{4.190890in}}%
\pgfpathcurveto{\pgfqpoint{3.087223in}{4.195008in}}{\pgfqpoint{3.081637in}{4.197322in}}{\pgfqpoint{3.075813in}{4.197322in}}%
\pgfpathcurveto{\pgfqpoint{3.069989in}{4.197322in}}{\pgfqpoint{3.064403in}{4.195008in}}{\pgfqpoint{3.060285in}{4.190890in}}%
\pgfpathcurveto{\pgfqpoint{3.056167in}{4.186772in}}{\pgfqpoint{3.053853in}{4.181186in}}{\pgfqpoint{3.053853in}{4.175362in}}%
\pgfpathcurveto{\pgfqpoint{3.053853in}{4.169538in}}{\pgfqpoint{3.056167in}{4.163952in}}{\pgfqpoint{3.060285in}{4.159834in}}%
\pgfpathcurveto{\pgfqpoint{3.064403in}{4.155715in}}{\pgfqpoint{3.069989in}{4.153402in}}{\pgfqpoint{3.075813in}{4.153402in}}%
\pgfpathlineto{\pgfqpoint{3.075813in}{4.153402in}}%
\pgfpathclose%
\pgfusepath{stroke,fill}%
\end{pgfscope}%
\begin{pgfscope}%
\pgfpathrectangle{\pgfqpoint{1.414736in}{0.548769in}}{\pgfqpoint{3.497747in}{7.253186in}}%
\pgfusepath{clip}%
\pgfsetbuttcap%
\pgfsetroundjoin%
\definecolor{currentfill}{rgb}{0.886271,0.892374,0.095374}%
\pgfsetfillcolor{currentfill}%
\pgfsetlinewidth{1.003750pt}%
\definecolor{currentstroke}{rgb}{0.886271,0.892374,0.095374}%
\pgfsetstrokecolor{currentstroke}%
\pgfsetdash{}{0pt}%
\pgfpathmoveto{\pgfqpoint{3.738621in}{4.624388in}}%
\pgfpathcurveto{\pgfqpoint{3.744445in}{4.624388in}}{\pgfqpoint{3.750031in}{4.626702in}}{\pgfqpoint{3.754149in}{4.630820in}}%
\pgfpathcurveto{\pgfqpoint{3.758267in}{4.634938in}}{\pgfqpoint{3.760581in}{4.640524in}}{\pgfqpoint{3.760581in}{4.646348in}}%
\pgfpathcurveto{\pgfqpoint{3.760581in}{4.652172in}}{\pgfqpoint{3.758267in}{4.657758in}}{\pgfqpoint{3.754149in}{4.661876in}}%
\pgfpathcurveto{\pgfqpoint{3.750031in}{4.665994in}}{\pgfqpoint{3.744445in}{4.668308in}}{\pgfqpoint{3.738621in}{4.668308in}}%
\pgfpathcurveto{\pgfqpoint{3.732797in}{4.668308in}}{\pgfqpoint{3.727211in}{4.665994in}}{\pgfqpoint{3.723092in}{4.661876in}}%
\pgfpathcurveto{\pgfqpoint{3.718974in}{4.657758in}}{\pgfqpoint{3.716660in}{4.652172in}}{\pgfqpoint{3.716660in}{4.646348in}}%
\pgfpathcurveto{\pgfqpoint{3.716660in}{4.640524in}}{\pgfqpoint{3.718974in}{4.634938in}}{\pgfqpoint{3.723092in}{4.630820in}}%
\pgfpathcurveto{\pgfqpoint{3.727211in}{4.626702in}}{\pgfqpoint{3.732797in}{4.624388in}}{\pgfqpoint{3.738621in}{4.624388in}}%
\pgfpathlineto{\pgfqpoint{3.738621in}{4.624388in}}%
\pgfpathclose%
\pgfusepath{stroke,fill}%
\end{pgfscope}%
\begin{pgfscope}%
\pgfpathrectangle{\pgfqpoint{1.414736in}{0.548769in}}{\pgfqpoint{3.497747in}{7.253186in}}%
\pgfusepath{clip}%
\pgfsetbuttcap%
\pgfsetroundjoin%
\definecolor{currentfill}{rgb}{0.159194,0.482237,0.558073}%
\pgfsetfillcolor{currentfill}%
\pgfsetlinewidth{1.003750pt}%
\definecolor{currentstroke}{rgb}{0.159194,0.482237,0.558073}%
\pgfsetstrokecolor{currentstroke}%
\pgfsetdash{}{0pt}%
\pgfpathmoveto{\pgfqpoint{3.627326in}{5.095374in}}%
\pgfpathcurveto{\pgfqpoint{3.633150in}{5.095374in}}{\pgfqpoint{3.638736in}{5.097688in}}{\pgfqpoint{3.642855in}{5.101806in}}%
\pgfpathcurveto{\pgfqpoint{3.646973in}{5.105924in}}{\pgfqpoint{3.649287in}{5.111510in}}{\pgfqpoint{3.649287in}{5.117334in}}%
\pgfpathcurveto{\pgfqpoint{3.649287in}{5.123158in}}{\pgfqpoint{3.646973in}{5.128744in}}{\pgfqpoint{3.642855in}{5.132862in}}%
\pgfpathcurveto{\pgfqpoint{3.638736in}{5.136980in}}{\pgfqpoint{3.633150in}{5.139294in}}{\pgfqpoint{3.627326in}{5.139294in}}%
\pgfpathcurveto{\pgfqpoint{3.621502in}{5.139294in}}{\pgfqpoint{3.615916in}{5.136980in}}{\pgfqpoint{3.611798in}{5.132862in}}%
\pgfpathcurveto{\pgfqpoint{3.607680in}{5.128744in}}{\pgfqpoint{3.605366in}{5.123158in}}{\pgfqpoint{3.605366in}{5.117334in}}%
\pgfpathcurveto{\pgfqpoint{3.605366in}{5.111510in}}{\pgfqpoint{3.607680in}{5.105924in}}{\pgfqpoint{3.611798in}{5.101806in}}%
\pgfpathcurveto{\pgfqpoint{3.615916in}{5.097688in}}{\pgfqpoint{3.621502in}{5.095374in}}{\pgfqpoint{3.627326in}{5.095374in}}%
\pgfpathlineto{\pgfqpoint{3.627326in}{5.095374in}}%
\pgfpathclose%
\pgfusepath{stroke,fill}%
\end{pgfscope}%
\begin{pgfscope}%
\pgfpathrectangle{\pgfqpoint{1.414736in}{0.548769in}}{\pgfqpoint{3.497747in}{7.253186in}}%
\pgfusepath{clip}%
\pgfsetbuttcap%
\pgfsetroundjoin%
\definecolor{currentfill}{rgb}{0.175841,0.441290,0.557685}%
\pgfsetfillcolor{currentfill}%
\pgfsetlinewidth{1.003750pt}%
\definecolor{currentstroke}{rgb}{0.175841,0.441290,0.557685}%
\pgfsetstrokecolor{currentstroke}%
\pgfsetdash{}{0pt}%
\pgfpathmoveto{\pgfqpoint{3.803694in}{5.566360in}}%
\pgfpathcurveto{\pgfqpoint{3.809518in}{5.566360in}}{\pgfqpoint{3.815104in}{5.568674in}}{\pgfqpoint{3.819222in}{5.572792in}}%
\pgfpathcurveto{\pgfqpoint{3.823341in}{5.576910in}}{\pgfqpoint{3.825655in}{5.582496in}}{\pgfqpoint{3.825655in}{5.588320in}}%
\pgfpathcurveto{\pgfqpoint{3.825655in}{5.594144in}}{\pgfqpoint{3.823341in}{5.599730in}}{\pgfqpoint{3.819222in}{5.603848in}}%
\pgfpathcurveto{\pgfqpoint{3.815104in}{5.607967in}}{\pgfqpoint{3.809518in}{5.610280in}}{\pgfqpoint{3.803694in}{5.610280in}}%
\pgfpathcurveto{\pgfqpoint{3.797870in}{5.610280in}}{\pgfqpoint{3.792284in}{5.607967in}}{\pgfqpoint{3.788166in}{5.603848in}}%
\pgfpathcurveto{\pgfqpoint{3.784048in}{5.599730in}}{\pgfqpoint{3.781734in}{5.594144in}}{\pgfqpoint{3.781734in}{5.588320in}}%
\pgfpathcurveto{\pgfqpoint{3.781734in}{5.582496in}}{\pgfqpoint{3.784048in}{5.576910in}}{\pgfqpoint{3.788166in}{5.572792in}}%
\pgfpathcurveto{\pgfqpoint{3.792284in}{5.568674in}}{\pgfqpoint{3.797870in}{5.566360in}}{\pgfqpoint{3.803694in}{5.566360in}}%
\pgfpathlineto{\pgfqpoint{3.803694in}{5.566360in}}%
\pgfpathclose%
\pgfusepath{stroke,fill}%
\end{pgfscope}%
\begin{pgfscope}%
\pgfpathrectangle{\pgfqpoint{1.414736in}{0.548769in}}{\pgfqpoint{3.497747in}{7.253186in}}%
\pgfusepath{clip}%
\pgfsetbuttcap%
\pgfsetroundjoin%
\definecolor{currentfill}{rgb}{0.143303,0.669459,0.511215}%
\pgfsetfillcolor{currentfill}%
\pgfsetlinewidth{1.003750pt}%
\definecolor{currentstroke}{rgb}{0.143303,0.669459,0.511215}%
\pgfsetstrokecolor{currentstroke}%
\pgfsetdash{}{0pt}%
\pgfpathmoveto{\pgfqpoint{3.436763in}{6.037346in}}%
\pgfpathcurveto{\pgfqpoint{3.442587in}{6.037346in}}{\pgfqpoint{3.448173in}{6.039660in}}{\pgfqpoint{3.452291in}{6.043778in}}%
\pgfpathcurveto{\pgfqpoint{3.456409in}{6.047896in}}{\pgfqpoint{3.458723in}{6.053482in}}{\pgfqpoint{3.458723in}{6.059306in}}%
\pgfpathcurveto{\pgfqpoint{3.458723in}{6.065130in}}{\pgfqpoint{3.456409in}{6.070716in}}{\pgfqpoint{3.452291in}{6.074835in}}%
\pgfpathcurveto{\pgfqpoint{3.448173in}{6.078953in}}{\pgfqpoint{3.442587in}{6.081267in}}{\pgfqpoint{3.436763in}{6.081267in}}%
\pgfpathcurveto{\pgfqpoint{3.430939in}{6.081267in}}{\pgfqpoint{3.425353in}{6.078953in}}{\pgfqpoint{3.421235in}{6.074835in}}%
\pgfpathcurveto{\pgfqpoint{3.417116in}{6.070716in}}{\pgfqpoint{3.414803in}{6.065130in}}{\pgfqpoint{3.414803in}{6.059306in}}%
\pgfpathcurveto{\pgfqpoint{3.414803in}{6.053482in}}{\pgfqpoint{3.417116in}{6.047896in}}{\pgfqpoint{3.421235in}{6.043778in}}%
\pgfpathcurveto{\pgfqpoint{3.425353in}{6.039660in}}{\pgfqpoint{3.430939in}{6.037346in}}{\pgfqpoint{3.436763in}{6.037346in}}%
\pgfpathlineto{\pgfqpoint{3.436763in}{6.037346in}}%
\pgfpathclose%
\pgfusepath{stroke,fill}%
\end{pgfscope}%
\begin{pgfscope}%
\pgfpathrectangle{\pgfqpoint{1.414736in}{0.548769in}}{\pgfqpoint{3.497747in}{7.253186in}}%
\pgfusepath{clip}%
\pgfsetbuttcap%
\pgfsetroundjoin%
\definecolor{currentfill}{rgb}{0.772852,0.877868,0.131109}%
\pgfsetfillcolor{currentfill}%
\pgfsetlinewidth{1.003750pt}%
\definecolor{currentstroke}{rgb}{0.772852,0.877868,0.131109}%
\pgfsetstrokecolor{currentstroke}%
\pgfsetdash{}{0pt}%
\pgfpathmoveto{\pgfqpoint{3.936148in}{6.508332in}}%
\pgfpathcurveto{\pgfqpoint{3.941972in}{6.508332in}}{\pgfqpoint{3.947558in}{6.510646in}}{\pgfqpoint{3.951676in}{6.514764in}}%
\pgfpathcurveto{\pgfqpoint{3.955794in}{6.518882in}}{\pgfqpoint{3.958108in}{6.524468in}}{\pgfqpoint{3.958108in}{6.530292in}}%
\pgfpathcurveto{\pgfqpoint{3.958108in}{6.536116in}}{\pgfqpoint{3.955794in}{6.541702in}}{\pgfqpoint{3.951676in}{6.545821in}}%
\pgfpathcurveto{\pgfqpoint{3.947558in}{6.549939in}}{\pgfqpoint{3.941972in}{6.552253in}}{\pgfqpoint{3.936148in}{6.552253in}}%
\pgfpathcurveto{\pgfqpoint{3.930324in}{6.552253in}}{\pgfqpoint{3.924737in}{6.549939in}}{\pgfqpoint{3.920619in}{6.545821in}}%
\pgfpathcurveto{\pgfqpoint{3.916501in}{6.541702in}}{\pgfqpoint{3.914187in}{6.536116in}}{\pgfqpoint{3.914187in}{6.530292in}}%
\pgfpathcurveto{\pgfqpoint{3.914187in}{6.524468in}}{\pgfqpoint{3.916501in}{6.518882in}}{\pgfqpoint{3.920619in}{6.514764in}}%
\pgfpathcurveto{\pgfqpoint{3.924737in}{6.510646in}}{\pgfqpoint{3.930324in}{6.508332in}}{\pgfqpoint{3.936148in}{6.508332in}}%
\pgfpathlineto{\pgfqpoint{3.936148in}{6.508332in}}%
\pgfpathclose%
\pgfusepath{stroke,fill}%
\end{pgfscope}%
\begin{pgfscope}%
\pgfpathrectangle{\pgfqpoint{1.414736in}{0.548769in}}{\pgfqpoint{3.497747in}{7.253186in}}%
\pgfusepath{clip}%
\pgfsetbuttcap%
\pgfsetroundjoin%
\definecolor{currentfill}{rgb}{0.993248,0.906157,0.143936}%
\pgfsetfillcolor{currentfill}%
\pgfsetlinewidth{1.003750pt}%
\definecolor{currentstroke}{rgb}{0.993248,0.906157,0.143936}%
\pgfsetstrokecolor{currentstroke}%
\pgfsetdash{}{0pt}%
\pgfpathmoveto{\pgfqpoint{3.918427in}{6.979318in}}%
\pgfpathcurveto{\pgfqpoint{3.924251in}{6.979318in}}{\pgfqpoint{3.929837in}{6.981632in}}{\pgfqpoint{3.933955in}{6.985750in}}%
\pgfpathcurveto{\pgfqpoint{3.938073in}{6.989868in}}{\pgfqpoint{3.940387in}{6.995455in}}{\pgfqpoint{3.940387in}{7.001278in}}%
\pgfpathcurveto{\pgfqpoint{3.940387in}{7.007102in}}{\pgfqpoint{3.938073in}{7.012689in}}{\pgfqpoint{3.933955in}{7.016807in}}%
\pgfpathcurveto{\pgfqpoint{3.929837in}{7.020925in}}{\pgfqpoint{3.924251in}{7.023239in}}{\pgfqpoint{3.918427in}{7.023239in}}%
\pgfpathcurveto{\pgfqpoint{3.912603in}{7.023239in}}{\pgfqpoint{3.907017in}{7.020925in}}{\pgfqpoint{3.902899in}{7.016807in}}%
\pgfpathcurveto{\pgfqpoint{3.898781in}{7.012689in}}{\pgfqpoint{3.896467in}{7.007102in}}{\pgfqpoint{3.896467in}{7.001278in}}%
\pgfpathcurveto{\pgfqpoint{3.896467in}{6.995455in}}{\pgfqpoint{3.898781in}{6.989868in}}{\pgfqpoint{3.902899in}{6.985750in}}%
\pgfpathcurveto{\pgfqpoint{3.907017in}{6.981632in}}{\pgfqpoint{3.912603in}{6.979318in}}{\pgfqpoint{3.918427in}{6.979318in}}%
\pgfpathlineto{\pgfqpoint{3.918427in}{6.979318in}}%
\pgfpathclose%
\pgfusepath{stroke,fill}%
\end{pgfscope}%
\begin{pgfscope}%
\pgfpathrectangle{\pgfqpoint{1.414736in}{0.548769in}}{\pgfqpoint{3.497747in}{7.253186in}}%
\pgfusepath{clip}%
\pgfsetbuttcap%
\pgfsetroundjoin%
\definecolor{currentfill}{rgb}{0.525776,0.833491,0.288127}%
\pgfsetfillcolor{currentfill}%
\pgfsetlinewidth{1.003750pt}%
\definecolor{currentstroke}{rgb}{0.525776,0.833491,0.288127}%
\pgfsetstrokecolor{currentstroke}%
\pgfsetdash{}{0pt}%
\pgfpathmoveto{\pgfqpoint{4.027018in}{7.450304in}}%
\pgfpathcurveto{\pgfqpoint{4.032842in}{7.450304in}}{\pgfqpoint{4.038428in}{7.452618in}}{\pgfqpoint{4.042546in}{7.456736in}}%
\pgfpathcurveto{\pgfqpoint{4.046665in}{7.460854in}}{\pgfqpoint{4.048978in}{7.466441in}}{\pgfqpoint{4.048978in}{7.472265in}}%
\pgfpathcurveto{\pgfqpoint{4.048978in}{7.478089in}}{\pgfqpoint{4.046665in}{7.483675in}}{\pgfqpoint{4.042546in}{7.487793in}}%
\pgfpathcurveto{\pgfqpoint{4.038428in}{7.491911in}}{\pgfqpoint{4.032842in}{7.494225in}}{\pgfqpoint{4.027018in}{7.494225in}}%
\pgfpathcurveto{\pgfqpoint{4.021194in}{7.494225in}}{\pgfqpoint{4.015608in}{7.491911in}}{\pgfqpoint{4.011490in}{7.487793in}}%
\pgfpathcurveto{\pgfqpoint{4.007372in}{7.483675in}}{\pgfqpoint{4.005058in}{7.478089in}}{\pgfqpoint{4.005058in}{7.472265in}}%
\pgfpathcurveto{\pgfqpoint{4.005058in}{7.466441in}}{\pgfqpoint{4.007372in}{7.460854in}}{\pgfqpoint{4.011490in}{7.456736in}}%
\pgfpathcurveto{\pgfqpoint{4.015608in}{7.452618in}}{\pgfqpoint{4.021194in}{7.450304in}}{\pgfqpoint{4.027018in}{7.450304in}}%
\pgfpathlineto{\pgfqpoint{4.027018in}{7.450304in}}%
\pgfpathclose%
\pgfusepath{stroke,fill}%
\end{pgfscope}%
\begin{pgfscope}%
\pgfpathrectangle{\pgfqpoint{1.414736in}{0.548769in}}{\pgfqpoint{3.497747in}{7.253186in}}%
\pgfusepath{clip}%
\pgfsetbuttcap%
\pgfsetroundjoin%
\definecolor{currentfill}{rgb}{0.993248,0.906157,0.143936}%
\pgfsetfillcolor{currentfill}%
\pgfsetlinewidth{1.003750pt}%
\definecolor{currentstroke}{rgb}{0.993248,0.906157,0.143936}%
\pgfsetstrokecolor{currentstroke}%
\pgfsetdash{}{0pt}%
\pgfpathmoveto{\pgfqpoint{3.884313in}{0.856499in}}%
\pgfpathcurveto{\pgfqpoint{3.890137in}{0.856499in}}{\pgfqpoint{3.895723in}{0.858813in}}{\pgfqpoint{3.899841in}{0.862931in}}%
\pgfpathcurveto{\pgfqpoint{3.903959in}{0.867049in}}{\pgfqpoint{3.906273in}{0.872635in}}{\pgfqpoint{3.906273in}{0.878459in}}%
\pgfpathcurveto{\pgfqpoint{3.906273in}{0.884283in}}{\pgfqpoint{3.903959in}{0.889869in}}{\pgfqpoint{3.899841in}{0.893987in}}%
\pgfpathcurveto{\pgfqpoint{3.895723in}{0.898105in}}{\pgfqpoint{3.890137in}{0.900419in}}{\pgfqpoint{3.884313in}{0.900419in}}%
\pgfpathcurveto{\pgfqpoint{3.878489in}{0.900419in}}{\pgfqpoint{3.872903in}{0.898105in}}{\pgfqpoint{3.868785in}{0.893987in}}%
\pgfpathcurveto{\pgfqpoint{3.864667in}{0.889869in}}{\pgfqpoint{3.862353in}{0.884283in}}{\pgfqpoint{3.862353in}{0.878459in}}%
\pgfpathcurveto{\pgfqpoint{3.862353in}{0.872635in}}{\pgfqpoint{3.864667in}{0.867049in}}{\pgfqpoint{3.868785in}{0.862931in}}%
\pgfpathcurveto{\pgfqpoint{3.872903in}{0.858813in}}{\pgfqpoint{3.878489in}{0.856499in}}{\pgfqpoint{3.884313in}{0.856499in}}%
\pgfpathlineto{\pgfqpoint{3.884313in}{0.856499in}}%
\pgfpathclose%
\pgfusepath{stroke,fill}%
\end{pgfscope}%
\begin{pgfscope}%
\pgfpathrectangle{\pgfqpoint{1.414736in}{0.548769in}}{\pgfqpoint{3.497747in}{7.253186in}}%
\pgfusepath{clip}%
\pgfsetbuttcap%
\pgfsetroundjoin%
\definecolor{currentfill}{rgb}{0.720391,0.870350,0.162603}%
\pgfsetfillcolor{currentfill}%
\pgfsetlinewidth{1.003750pt}%
\definecolor{currentstroke}{rgb}{0.720391,0.870350,0.162603}%
\pgfsetstrokecolor{currentstroke}%
\pgfsetdash{}{0pt}%
\pgfpathmoveto{\pgfqpoint{3.385058in}{1.327485in}}%
\pgfpathcurveto{\pgfqpoint{3.390882in}{1.327485in}}{\pgfqpoint{3.396468in}{1.329799in}}{\pgfqpoint{3.400587in}{1.333917in}}%
\pgfpathcurveto{\pgfqpoint{3.404705in}{1.338035in}}{\pgfqpoint{3.407019in}{1.343621in}}{\pgfqpoint{3.407019in}{1.349445in}}%
\pgfpathcurveto{\pgfqpoint{3.407019in}{1.355269in}}{\pgfqpoint{3.404705in}{1.360855in}}{\pgfqpoint{3.400587in}{1.364973in}}%
\pgfpathcurveto{\pgfqpoint{3.396468in}{1.369092in}}{\pgfqpoint{3.390882in}{1.371405in}}{\pgfqpoint{3.385058in}{1.371405in}}%
\pgfpathcurveto{\pgfqpoint{3.379234in}{1.371405in}}{\pgfqpoint{3.373648in}{1.369092in}}{\pgfqpoint{3.369530in}{1.364973in}}%
\pgfpathcurveto{\pgfqpoint{3.365412in}{1.360855in}}{\pgfqpoint{3.363098in}{1.355269in}}{\pgfqpoint{3.363098in}{1.349445in}}%
\pgfpathcurveto{\pgfqpoint{3.363098in}{1.343621in}}{\pgfqpoint{3.365412in}{1.338035in}}{\pgfqpoint{3.369530in}{1.333917in}}%
\pgfpathcurveto{\pgfqpoint{3.373648in}{1.329799in}}{\pgfqpoint{3.379234in}{1.327485in}}{\pgfqpoint{3.385058in}{1.327485in}}%
\pgfpathlineto{\pgfqpoint{3.385058in}{1.327485in}}%
\pgfpathclose%
\pgfusepath{stroke,fill}%
\end{pgfscope}%
\begin{pgfscope}%
\pgfpathrectangle{\pgfqpoint{1.414736in}{0.548769in}}{\pgfqpoint{3.497747in}{7.253186in}}%
\pgfusepath{clip}%
\pgfsetbuttcap%
\pgfsetroundjoin%
\definecolor{currentfill}{rgb}{0.412913,0.803041,0.357269}%
\pgfsetfillcolor{currentfill}%
\pgfsetlinewidth{1.003750pt}%
\definecolor{currentstroke}{rgb}{0.412913,0.803041,0.357269}%
\pgfsetstrokecolor{currentstroke}%
\pgfsetdash{}{0pt}%
\pgfpathmoveto{\pgfqpoint{3.311147in}{1.798471in}}%
\pgfpathcurveto{\pgfqpoint{3.316971in}{1.798471in}}{\pgfqpoint{3.322557in}{1.800785in}}{\pgfqpoint{3.326675in}{1.804903in}}%
\pgfpathcurveto{\pgfqpoint{3.330793in}{1.809021in}}{\pgfqpoint{3.333107in}{1.814607in}}{\pgfqpoint{3.333107in}{1.820431in}}%
\pgfpathcurveto{\pgfqpoint{3.333107in}{1.826255in}}{\pgfqpoint{3.330793in}{1.831841in}}{\pgfqpoint{3.326675in}{1.835960in}}%
\pgfpathcurveto{\pgfqpoint{3.322557in}{1.840078in}}{\pgfqpoint{3.316971in}{1.842392in}}{\pgfqpoint{3.311147in}{1.842392in}}%
\pgfpathcurveto{\pgfqpoint{3.305323in}{1.842392in}}{\pgfqpoint{3.299737in}{1.840078in}}{\pgfqpoint{3.295619in}{1.835960in}}%
\pgfpathcurveto{\pgfqpoint{3.291501in}{1.831841in}}{\pgfqpoint{3.289187in}{1.826255in}}{\pgfqpoint{3.289187in}{1.820431in}}%
\pgfpathcurveto{\pgfqpoint{3.289187in}{1.814607in}}{\pgfqpoint{3.291501in}{1.809021in}}{\pgfqpoint{3.295619in}{1.804903in}}%
\pgfpathcurveto{\pgfqpoint{3.299737in}{1.800785in}}{\pgfqpoint{3.305323in}{1.798471in}}{\pgfqpoint{3.311147in}{1.798471in}}%
\pgfpathlineto{\pgfqpoint{3.311147in}{1.798471in}}%
\pgfpathclose%
\pgfusepath{stroke,fill}%
\end{pgfscope}%
\begin{pgfscope}%
\pgfpathrectangle{\pgfqpoint{1.414736in}{0.548769in}}{\pgfqpoint{3.497747in}{7.253186in}}%
\pgfusepath{clip}%
\pgfsetbuttcap%
\pgfsetroundjoin%
\definecolor{currentfill}{rgb}{0.993248,0.906157,0.143936}%
\pgfsetfillcolor{currentfill}%
\pgfsetlinewidth{1.003750pt}%
\definecolor{currentstroke}{rgb}{0.993248,0.906157,0.143936}%
\pgfsetstrokecolor{currentstroke}%
\pgfsetdash{}{0pt}%
\pgfpathmoveto{\pgfqpoint{3.155453in}{2.269457in}}%
\pgfpathcurveto{\pgfqpoint{3.161277in}{2.269457in}}{\pgfqpoint{3.166863in}{2.271771in}}{\pgfqpoint{3.170981in}{2.275889in}}%
\pgfpathcurveto{\pgfqpoint{3.175099in}{2.280007in}}{\pgfqpoint{3.177413in}{2.285593in}}{\pgfqpoint{3.177413in}{2.291417in}}%
\pgfpathcurveto{\pgfqpoint{3.177413in}{2.297241in}}{\pgfqpoint{3.175099in}{2.302828in}}{\pgfqpoint{3.170981in}{2.306946in}}%
\pgfpathcurveto{\pgfqpoint{3.166863in}{2.311064in}}{\pgfqpoint{3.161277in}{2.313378in}}{\pgfqpoint{3.155453in}{2.313378in}}%
\pgfpathcurveto{\pgfqpoint{3.149629in}{2.313378in}}{\pgfqpoint{3.144043in}{2.311064in}}{\pgfqpoint{3.139925in}{2.306946in}}%
\pgfpathcurveto{\pgfqpoint{3.135807in}{2.302828in}}{\pgfqpoint{3.133493in}{2.297241in}}{\pgfqpoint{3.133493in}{2.291417in}}%
\pgfpathcurveto{\pgfqpoint{3.133493in}{2.285593in}}{\pgfqpoint{3.135807in}{2.280007in}}{\pgfqpoint{3.139925in}{2.275889in}}%
\pgfpathcurveto{\pgfqpoint{3.144043in}{2.271771in}}{\pgfqpoint{3.149629in}{2.269457in}}{\pgfqpoint{3.155453in}{2.269457in}}%
\pgfpathlineto{\pgfqpoint{3.155453in}{2.269457in}}%
\pgfpathclose%
\pgfusepath{stroke,fill}%
\end{pgfscope}%
\begin{pgfscope}%
\pgfpathrectangle{\pgfqpoint{1.414736in}{0.548769in}}{\pgfqpoint{3.497747in}{7.253186in}}%
\pgfusepath{clip}%
\pgfsetbuttcap%
\pgfsetroundjoin%
\definecolor{currentfill}{rgb}{0.525776,0.833491,0.288127}%
\pgfsetfillcolor{currentfill}%
\pgfsetlinewidth{1.003750pt}%
\definecolor{currentstroke}{rgb}{0.525776,0.833491,0.288127}%
\pgfsetstrokecolor{currentstroke}%
\pgfsetdash{}{0pt}%
\pgfpathmoveto{\pgfqpoint{3.420107in}{2.740443in}}%
\pgfpathcurveto{\pgfqpoint{3.425931in}{2.740443in}}{\pgfqpoint{3.431517in}{2.742757in}}{\pgfqpoint{3.435635in}{2.746875in}}%
\pgfpathcurveto{\pgfqpoint{3.439753in}{2.750993in}}{\pgfqpoint{3.442067in}{2.756580in}}{\pgfqpoint{3.442067in}{2.762404in}}%
\pgfpathcurveto{\pgfqpoint{3.442067in}{2.768227in}}{\pgfqpoint{3.439753in}{2.773814in}}{\pgfqpoint{3.435635in}{2.777932in}}%
\pgfpathcurveto{\pgfqpoint{3.431517in}{2.782050in}}{\pgfqpoint{3.425931in}{2.784364in}}{\pgfqpoint{3.420107in}{2.784364in}}%
\pgfpathcurveto{\pgfqpoint{3.414283in}{2.784364in}}{\pgfqpoint{3.408697in}{2.782050in}}{\pgfqpoint{3.404579in}{2.777932in}}%
\pgfpathcurveto{\pgfqpoint{3.400461in}{2.773814in}}{\pgfqpoint{3.398147in}{2.768227in}}{\pgfqpoint{3.398147in}{2.762404in}}%
\pgfpathcurveto{\pgfqpoint{3.398147in}{2.756580in}}{\pgfqpoint{3.400461in}{2.750993in}}{\pgfqpoint{3.404579in}{2.746875in}}%
\pgfpathcurveto{\pgfqpoint{3.408697in}{2.742757in}}{\pgfqpoint{3.414283in}{2.740443in}}{\pgfqpoint{3.420107in}{2.740443in}}%
\pgfpathlineto{\pgfqpoint{3.420107in}{2.740443in}}%
\pgfpathclose%
\pgfusepath{stroke,fill}%
\end{pgfscope}%
\begin{pgfscope}%
\pgfpathrectangle{\pgfqpoint{1.414736in}{0.548769in}}{\pgfqpoint{3.497747in}{7.253186in}}%
\pgfusepath{clip}%
\pgfsetbuttcap%
\pgfsetroundjoin%
\definecolor{currentfill}{rgb}{0.657642,0.860219,0.203082}%
\pgfsetfillcolor{currentfill}%
\pgfsetlinewidth{1.003750pt}%
\definecolor{currentstroke}{rgb}{0.657642,0.860219,0.203082}%
\pgfsetstrokecolor{currentstroke}%
\pgfsetdash{}{0pt}%
\pgfpathmoveto{\pgfqpoint{3.347701in}{3.211429in}}%
\pgfpathcurveto{\pgfqpoint{3.353525in}{3.211429in}}{\pgfqpoint{3.359111in}{3.213743in}}{\pgfqpoint{3.363229in}{3.217861in}}%
\pgfpathcurveto{\pgfqpoint{3.367348in}{3.221980in}}{\pgfqpoint{3.369661in}{3.227566in}}{\pgfqpoint{3.369661in}{3.233390in}}%
\pgfpathcurveto{\pgfqpoint{3.369661in}{3.239214in}}{\pgfqpoint{3.367348in}{3.244800in}}{\pgfqpoint{3.363229in}{3.248918in}}%
\pgfpathcurveto{\pgfqpoint{3.359111in}{3.253036in}}{\pgfqpoint{3.353525in}{3.255350in}}{\pgfqpoint{3.347701in}{3.255350in}}%
\pgfpathcurveto{\pgfqpoint{3.341877in}{3.255350in}}{\pgfqpoint{3.336291in}{3.253036in}}{\pgfqpoint{3.332173in}{3.248918in}}%
\pgfpathcurveto{\pgfqpoint{3.328055in}{3.244800in}}{\pgfqpoint{3.325741in}{3.239214in}}{\pgfqpoint{3.325741in}{3.233390in}}%
\pgfpathcurveto{\pgfqpoint{3.325741in}{3.227566in}}{\pgfqpoint{3.328055in}{3.221980in}}{\pgfqpoint{3.332173in}{3.217861in}}%
\pgfpathcurveto{\pgfqpoint{3.336291in}{3.213743in}}{\pgfqpoint{3.341877in}{3.211429in}}{\pgfqpoint{3.347701in}{3.211429in}}%
\pgfpathlineto{\pgfqpoint{3.347701in}{3.211429in}}%
\pgfpathclose%
\pgfusepath{stroke,fill}%
\end{pgfscope}%
\begin{pgfscope}%
\pgfpathrectangle{\pgfqpoint{1.414736in}{0.548769in}}{\pgfqpoint{3.497747in}{7.253186in}}%
\pgfusepath{clip}%
\pgfsetbuttcap%
\pgfsetroundjoin%
\definecolor{currentfill}{rgb}{0.964894,0.902323,0.123941}%
\pgfsetfillcolor{currentfill}%
\pgfsetlinewidth{1.003750pt}%
\definecolor{currentstroke}{rgb}{0.964894,0.902323,0.123941}%
\pgfsetstrokecolor{currentstroke}%
\pgfsetdash{}{0pt}%
\pgfpathmoveto{\pgfqpoint{3.029193in}{3.682415in}}%
\pgfpathcurveto{\pgfqpoint{3.035017in}{3.682415in}}{\pgfqpoint{3.040603in}{3.684729in}}{\pgfqpoint{3.044721in}{3.688847in}}%
\pgfpathcurveto{\pgfqpoint{3.048839in}{3.692966in}}{\pgfqpoint{3.051153in}{3.698552in}}{\pgfqpoint{3.051153in}{3.704376in}}%
\pgfpathcurveto{\pgfqpoint{3.051153in}{3.710200in}}{\pgfqpoint{3.048839in}{3.715786in}}{\pgfqpoint{3.044721in}{3.719904in}}%
\pgfpathcurveto{\pgfqpoint{3.040603in}{3.724022in}}{\pgfqpoint{3.035017in}{3.726336in}}{\pgfqpoint{3.029193in}{3.726336in}}%
\pgfpathcurveto{\pgfqpoint{3.023369in}{3.726336in}}{\pgfqpoint{3.017783in}{3.724022in}}{\pgfqpoint{3.013664in}{3.719904in}}%
\pgfpathcurveto{\pgfqpoint{3.009546in}{3.715786in}}{\pgfqpoint{3.007232in}{3.710200in}}{\pgfqpoint{3.007232in}{3.704376in}}%
\pgfpathcurveto{\pgfqpoint{3.007232in}{3.698552in}}{\pgfqpoint{3.009546in}{3.692966in}}{\pgfqpoint{3.013664in}{3.688847in}}%
\pgfpathcurveto{\pgfqpoint{3.017783in}{3.684729in}}{\pgfqpoint{3.023369in}{3.682415in}}{\pgfqpoint{3.029193in}{3.682415in}}%
\pgfpathlineto{\pgfqpoint{3.029193in}{3.682415in}}%
\pgfpathclose%
\pgfusepath{stroke,fill}%
\end{pgfscope}%
\begin{pgfscope}%
\pgfpathrectangle{\pgfqpoint{1.414736in}{0.548769in}}{\pgfqpoint{3.497747in}{7.253186in}}%
\pgfusepath{clip}%
\pgfsetbuttcap%
\pgfsetroundjoin%
\definecolor{currentfill}{rgb}{0.983868,0.904867,0.136897}%
\pgfsetfillcolor{currentfill}%
\pgfsetlinewidth{1.003750pt}%
\definecolor{currentstroke}{rgb}{0.983868,0.904867,0.136897}%
\pgfsetstrokecolor{currentstroke}%
\pgfsetdash{}{0pt}%
\pgfpathmoveto{\pgfqpoint{2.792220in}{4.153402in}}%
\pgfpathcurveto{\pgfqpoint{2.798044in}{4.153402in}}{\pgfqpoint{2.803630in}{4.155715in}}{\pgfqpoint{2.807748in}{4.159834in}}%
\pgfpathcurveto{\pgfqpoint{2.811866in}{4.163952in}}{\pgfqpoint{2.814180in}{4.169538in}}{\pgfqpoint{2.814180in}{4.175362in}}%
\pgfpathcurveto{\pgfqpoint{2.814180in}{4.181186in}}{\pgfqpoint{2.811866in}{4.186772in}}{\pgfqpoint{2.807748in}{4.190890in}}%
\pgfpathcurveto{\pgfqpoint{2.803630in}{4.195008in}}{\pgfqpoint{2.798044in}{4.197322in}}{\pgfqpoint{2.792220in}{4.197322in}}%
\pgfpathcurveto{\pgfqpoint{2.786396in}{4.197322in}}{\pgfqpoint{2.780810in}{4.195008in}}{\pgfqpoint{2.776692in}{4.190890in}}%
\pgfpathcurveto{\pgfqpoint{2.772573in}{4.186772in}}{\pgfqpoint{2.770260in}{4.181186in}}{\pgfqpoint{2.770260in}{4.175362in}}%
\pgfpathcurveto{\pgfqpoint{2.770260in}{4.169538in}}{\pgfqpoint{2.772573in}{4.163952in}}{\pgfqpoint{2.776692in}{4.159834in}}%
\pgfpathcurveto{\pgfqpoint{2.780810in}{4.155715in}}{\pgfqpoint{2.786396in}{4.153402in}}{\pgfqpoint{2.792220in}{4.153402in}}%
\pgfpathlineto{\pgfqpoint{2.792220in}{4.153402in}}%
\pgfpathclose%
\pgfusepath{stroke,fill}%
\end{pgfscope}%
\begin{pgfscope}%
\pgfpathrectangle{\pgfqpoint{1.414736in}{0.548769in}}{\pgfqpoint{3.497747in}{7.253186in}}%
\pgfusepath{clip}%
\pgfsetbuttcap%
\pgfsetroundjoin%
\definecolor{currentfill}{rgb}{0.964894,0.902323,0.123941}%
\pgfsetfillcolor{currentfill}%
\pgfsetlinewidth{1.003750pt}%
\definecolor{currentstroke}{rgb}{0.964894,0.902323,0.123941}%
\pgfsetstrokecolor{currentstroke}%
\pgfsetdash{}{0pt}%
\pgfpathmoveto{\pgfqpoint{4.251527in}{4.624388in}}%
\pgfpathcurveto{\pgfqpoint{4.257351in}{4.624388in}}{\pgfqpoint{4.262937in}{4.626702in}}{\pgfqpoint{4.267055in}{4.630820in}}%
\pgfpathcurveto{\pgfqpoint{4.271173in}{4.634938in}}{\pgfqpoint{4.273487in}{4.640524in}}{\pgfqpoint{4.273487in}{4.646348in}}%
\pgfpathcurveto{\pgfqpoint{4.273487in}{4.652172in}}{\pgfqpoint{4.271173in}{4.657758in}}{\pgfqpoint{4.267055in}{4.661876in}}%
\pgfpathcurveto{\pgfqpoint{4.262937in}{4.665994in}}{\pgfqpoint{4.257351in}{4.668308in}}{\pgfqpoint{4.251527in}{4.668308in}}%
\pgfpathcurveto{\pgfqpoint{4.245703in}{4.668308in}}{\pgfqpoint{4.240117in}{4.665994in}}{\pgfqpoint{4.235998in}{4.661876in}}%
\pgfpathcurveto{\pgfqpoint{4.231880in}{4.657758in}}{\pgfqpoint{4.229566in}{4.652172in}}{\pgfqpoint{4.229566in}{4.646348in}}%
\pgfpathcurveto{\pgfqpoint{4.229566in}{4.640524in}}{\pgfqpoint{4.231880in}{4.634938in}}{\pgfqpoint{4.235998in}{4.630820in}}%
\pgfpathcurveto{\pgfqpoint{4.240117in}{4.626702in}}{\pgfqpoint{4.245703in}{4.624388in}}{\pgfqpoint{4.251527in}{4.624388in}}%
\pgfpathlineto{\pgfqpoint{4.251527in}{4.624388in}}%
\pgfpathclose%
\pgfusepath{stroke,fill}%
\end{pgfscope}%
\begin{pgfscope}%
\pgfpathrectangle{\pgfqpoint{1.414736in}{0.548769in}}{\pgfqpoint{3.497747in}{7.253186in}}%
\pgfusepath{clip}%
\pgfsetbuttcap%
\pgfsetroundjoin%
\definecolor{currentfill}{rgb}{0.214000,0.722114,0.469588}%
\pgfsetfillcolor{currentfill}%
\pgfsetlinewidth{1.003750pt}%
\definecolor{currentstroke}{rgb}{0.214000,0.722114,0.469588}%
\pgfsetstrokecolor{currentstroke}%
\pgfsetdash{}{0pt}%
\pgfpathmoveto{\pgfqpoint{3.804960in}{5.095374in}}%
\pgfpathcurveto{\pgfqpoint{3.810784in}{5.095374in}}{\pgfqpoint{3.816370in}{5.097688in}}{\pgfqpoint{3.820488in}{5.101806in}}%
\pgfpathcurveto{\pgfqpoint{3.824606in}{5.105924in}}{\pgfqpoint{3.826920in}{5.111510in}}{\pgfqpoint{3.826920in}{5.117334in}}%
\pgfpathcurveto{\pgfqpoint{3.826920in}{5.123158in}}{\pgfqpoint{3.824606in}{5.128744in}}{\pgfqpoint{3.820488in}{5.132862in}}%
\pgfpathcurveto{\pgfqpoint{3.816370in}{5.136980in}}{\pgfqpoint{3.810784in}{5.139294in}}{\pgfqpoint{3.804960in}{5.139294in}}%
\pgfpathcurveto{\pgfqpoint{3.799136in}{5.139294in}}{\pgfqpoint{3.793550in}{5.136980in}}{\pgfqpoint{3.789431in}{5.132862in}}%
\pgfpathcurveto{\pgfqpoint{3.785313in}{5.128744in}}{\pgfqpoint{3.782999in}{5.123158in}}{\pgfqpoint{3.782999in}{5.117334in}}%
\pgfpathcurveto{\pgfqpoint{3.782999in}{5.111510in}}{\pgfqpoint{3.785313in}{5.105924in}}{\pgfqpoint{3.789431in}{5.101806in}}%
\pgfpathcurveto{\pgfqpoint{3.793550in}{5.097688in}}{\pgfqpoint{3.799136in}{5.095374in}}{\pgfqpoint{3.804960in}{5.095374in}}%
\pgfpathlineto{\pgfqpoint{3.804960in}{5.095374in}}%
\pgfpathclose%
\pgfusepath{stroke,fill}%
\end{pgfscope}%
\begin{pgfscope}%
\pgfpathrectangle{\pgfqpoint{1.414736in}{0.548769in}}{\pgfqpoint{3.497747in}{7.253186in}}%
\pgfusepath{clip}%
\pgfsetbuttcap%
\pgfsetroundjoin%
\definecolor{currentfill}{rgb}{0.123463,0.581687,0.547445}%
\pgfsetfillcolor{currentfill}%
\pgfsetlinewidth{1.003750pt}%
\definecolor{currentstroke}{rgb}{0.123463,0.581687,0.547445}%
\pgfsetstrokecolor{currentstroke}%
\pgfsetdash{}{0pt}%
\pgfpathmoveto{\pgfqpoint{3.737641in}{5.566360in}}%
\pgfpathcurveto{\pgfqpoint{3.743465in}{5.566360in}}{\pgfqpoint{3.749051in}{5.568674in}}{\pgfqpoint{3.753169in}{5.572792in}}%
\pgfpathcurveto{\pgfqpoint{3.757287in}{5.576910in}}{\pgfqpoint{3.759601in}{5.582496in}}{\pgfqpoint{3.759601in}{5.588320in}}%
\pgfpathcurveto{\pgfqpoint{3.759601in}{5.594144in}}{\pgfqpoint{3.757287in}{5.599730in}}{\pgfqpoint{3.753169in}{5.603848in}}%
\pgfpathcurveto{\pgfqpoint{3.749051in}{5.607967in}}{\pgfqpoint{3.743465in}{5.610280in}}{\pgfqpoint{3.737641in}{5.610280in}}%
\pgfpathcurveto{\pgfqpoint{3.731817in}{5.610280in}}{\pgfqpoint{3.726231in}{5.607967in}}{\pgfqpoint{3.722113in}{5.603848in}}%
\pgfpathcurveto{\pgfqpoint{3.717995in}{5.599730in}}{\pgfqpoint{3.715681in}{5.594144in}}{\pgfqpoint{3.715681in}{5.588320in}}%
\pgfpathcurveto{\pgfqpoint{3.715681in}{5.582496in}}{\pgfqpoint{3.717995in}{5.576910in}}{\pgfqpoint{3.722113in}{5.572792in}}%
\pgfpathcurveto{\pgfqpoint{3.726231in}{5.568674in}}{\pgfqpoint{3.731817in}{5.566360in}}{\pgfqpoint{3.737641in}{5.566360in}}%
\pgfpathlineto{\pgfqpoint{3.737641in}{5.566360in}}%
\pgfpathclose%
\pgfusepath{stroke,fill}%
\end{pgfscope}%
\begin{pgfscope}%
\pgfpathrectangle{\pgfqpoint{1.414736in}{0.548769in}}{\pgfqpoint{3.497747in}{7.253186in}}%
\pgfusepath{clip}%
\pgfsetbuttcap%
\pgfsetroundjoin%
\definecolor{currentfill}{rgb}{0.814576,0.883393,0.110347}%
\pgfsetfillcolor{currentfill}%
\pgfsetlinewidth{1.003750pt}%
\definecolor{currentstroke}{rgb}{0.814576,0.883393,0.110347}%
\pgfsetstrokecolor{currentstroke}%
\pgfsetdash{}{0pt}%
\pgfpathmoveto{\pgfqpoint{2.846948in}{6.037346in}}%
\pgfpathcurveto{\pgfqpoint{2.852772in}{6.037346in}}{\pgfqpoint{2.858358in}{6.039660in}}{\pgfqpoint{2.862476in}{6.043778in}}%
\pgfpathcurveto{\pgfqpoint{2.866594in}{6.047896in}}{\pgfqpoint{2.868908in}{6.053482in}}{\pgfqpoint{2.868908in}{6.059306in}}%
\pgfpathcurveto{\pgfqpoint{2.868908in}{6.065130in}}{\pgfqpoint{2.866594in}{6.070716in}}{\pgfqpoint{2.862476in}{6.074835in}}%
\pgfpathcurveto{\pgfqpoint{2.858358in}{6.078953in}}{\pgfqpoint{2.852772in}{6.081267in}}{\pgfqpoint{2.846948in}{6.081267in}}%
\pgfpathcurveto{\pgfqpoint{2.841124in}{6.081267in}}{\pgfqpoint{2.835538in}{6.078953in}}{\pgfqpoint{2.831420in}{6.074835in}}%
\pgfpathcurveto{\pgfqpoint{2.827302in}{6.070716in}}{\pgfqpoint{2.824988in}{6.065130in}}{\pgfqpoint{2.824988in}{6.059306in}}%
\pgfpathcurveto{\pgfqpoint{2.824988in}{6.053482in}}{\pgfqpoint{2.827302in}{6.047896in}}{\pgfqpoint{2.831420in}{6.043778in}}%
\pgfpathcurveto{\pgfqpoint{2.835538in}{6.039660in}}{\pgfqpoint{2.841124in}{6.037346in}}{\pgfqpoint{2.846948in}{6.037346in}}%
\pgfpathlineto{\pgfqpoint{2.846948in}{6.037346in}}%
\pgfpathclose%
\pgfusepath{stroke,fill}%
\end{pgfscope}%
\begin{pgfscope}%
\pgfpathrectangle{\pgfqpoint{1.414736in}{0.548769in}}{\pgfqpoint{3.497747in}{7.253186in}}%
\pgfusepath{clip}%
\pgfsetbuttcap%
\pgfsetroundjoin%
\definecolor{currentfill}{rgb}{0.993248,0.906157,0.143936}%
\pgfsetfillcolor{currentfill}%
\pgfsetlinewidth{1.003750pt}%
\definecolor{currentstroke}{rgb}{0.993248,0.906157,0.143936}%
\pgfsetstrokecolor{currentstroke}%
\pgfsetdash{}{0pt}%
\pgfpathmoveto{\pgfqpoint{3.210635in}{6.508332in}}%
\pgfpathcurveto{\pgfqpoint{3.216459in}{6.508332in}}{\pgfqpoint{3.222045in}{6.510646in}}{\pgfqpoint{3.226163in}{6.514764in}}%
\pgfpathcurveto{\pgfqpoint{3.230281in}{6.518882in}}{\pgfqpoint{3.232595in}{6.524468in}}{\pgfqpoint{3.232595in}{6.530292in}}%
\pgfpathcurveto{\pgfqpoint{3.232595in}{6.536116in}}{\pgfqpoint{3.230281in}{6.541702in}}{\pgfqpoint{3.226163in}{6.545821in}}%
\pgfpathcurveto{\pgfqpoint{3.222045in}{6.549939in}}{\pgfqpoint{3.216459in}{6.552253in}}{\pgfqpoint{3.210635in}{6.552253in}}%
\pgfpathcurveto{\pgfqpoint{3.204811in}{6.552253in}}{\pgfqpoint{3.199225in}{6.549939in}}{\pgfqpoint{3.195107in}{6.545821in}}%
\pgfpathcurveto{\pgfqpoint{3.190989in}{6.541702in}}{\pgfqpoint{3.188675in}{6.536116in}}{\pgfqpoint{3.188675in}{6.530292in}}%
\pgfpathcurveto{\pgfqpoint{3.188675in}{6.524468in}}{\pgfqpoint{3.190989in}{6.518882in}}{\pgfqpoint{3.195107in}{6.514764in}}%
\pgfpathcurveto{\pgfqpoint{3.199225in}{6.510646in}}{\pgfqpoint{3.204811in}{6.508332in}}{\pgfqpoint{3.210635in}{6.508332in}}%
\pgfpathlineto{\pgfqpoint{3.210635in}{6.508332in}}%
\pgfpathclose%
\pgfusepath{stroke,fill}%
\end{pgfscope}%
\begin{pgfscope}%
\pgfpathrectangle{\pgfqpoint{1.414736in}{0.548769in}}{\pgfqpoint{3.497747in}{7.253186in}}%
\pgfusepath{clip}%
\pgfsetbuttcap%
\pgfsetroundjoin%
\definecolor{currentfill}{rgb}{0.886271,0.892374,0.095374}%
\pgfsetfillcolor{currentfill}%
\pgfsetlinewidth{1.003750pt}%
\definecolor{currentstroke}{rgb}{0.886271,0.892374,0.095374}%
\pgfsetstrokecolor{currentstroke}%
\pgfsetdash{}{0pt}%
\pgfpathmoveto{\pgfqpoint{3.567737in}{6.979318in}}%
\pgfpathcurveto{\pgfqpoint{3.573561in}{6.979318in}}{\pgfqpoint{3.579147in}{6.981632in}}{\pgfqpoint{3.583266in}{6.985750in}}%
\pgfpathcurveto{\pgfqpoint{3.587384in}{6.989868in}}{\pgfqpoint{3.589698in}{6.995455in}}{\pgfqpoint{3.589698in}{7.001278in}}%
\pgfpathcurveto{\pgfqpoint{3.589698in}{7.007102in}}{\pgfqpoint{3.587384in}{7.012689in}}{\pgfqpoint{3.583266in}{7.016807in}}%
\pgfpathcurveto{\pgfqpoint{3.579147in}{7.020925in}}{\pgfqpoint{3.573561in}{7.023239in}}{\pgfqpoint{3.567737in}{7.023239in}}%
\pgfpathcurveto{\pgfqpoint{3.561913in}{7.023239in}}{\pgfqpoint{3.556327in}{7.020925in}}{\pgfqpoint{3.552209in}{7.016807in}}%
\pgfpathcurveto{\pgfqpoint{3.548091in}{7.012689in}}{\pgfqpoint{3.545777in}{7.007102in}}{\pgfqpoint{3.545777in}{7.001278in}}%
\pgfpathcurveto{\pgfqpoint{3.545777in}{6.995455in}}{\pgfqpoint{3.548091in}{6.989868in}}{\pgfqpoint{3.552209in}{6.985750in}}%
\pgfpathcurveto{\pgfqpoint{3.556327in}{6.981632in}}{\pgfqpoint{3.561913in}{6.979318in}}{\pgfqpoint{3.567737in}{6.979318in}}%
\pgfpathlineto{\pgfqpoint{3.567737in}{6.979318in}}%
\pgfpathclose%
\pgfusepath{stroke,fill}%
\end{pgfscope}%
\begin{pgfscope}%
\pgfpathrectangle{\pgfqpoint{1.414736in}{0.548769in}}{\pgfqpoint{3.497747in}{7.253186in}}%
\pgfusepath{clip}%
\pgfsetbuttcap%
\pgfsetroundjoin%
\definecolor{currentfill}{rgb}{0.835270,0.886029,0.102646}%
\pgfsetfillcolor{currentfill}%
\pgfsetlinewidth{1.003750pt}%
\definecolor{currentstroke}{rgb}{0.835270,0.886029,0.102646}%
\pgfsetstrokecolor{currentstroke}%
\pgfsetdash{}{0pt}%
\pgfpathmoveto{\pgfqpoint{4.533718in}{7.450304in}}%
\pgfpathcurveto{\pgfqpoint{4.539542in}{7.450304in}}{\pgfqpoint{4.545128in}{7.452618in}}{\pgfqpoint{4.549246in}{7.456736in}}%
\pgfpathcurveto{\pgfqpoint{4.553364in}{7.460854in}}{\pgfqpoint{4.555678in}{7.466441in}}{\pgfqpoint{4.555678in}{7.472265in}}%
\pgfpathcurveto{\pgfqpoint{4.555678in}{7.478089in}}{\pgfqpoint{4.553364in}{7.483675in}}{\pgfqpoint{4.549246in}{7.487793in}}%
\pgfpathcurveto{\pgfqpoint{4.545128in}{7.491911in}}{\pgfqpoint{4.539542in}{7.494225in}}{\pgfqpoint{4.533718in}{7.494225in}}%
\pgfpathcurveto{\pgfqpoint{4.527894in}{7.494225in}}{\pgfqpoint{4.522308in}{7.491911in}}{\pgfqpoint{4.518189in}{7.487793in}}%
\pgfpathcurveto{\pgfqpoint{4.514071in}{7.483675in}}{\pgfqpoint{4.511757in}{7.478089in}}{\pgfqpoint{4.511757in}{7.472265in}}%
\pgfpathcurveto{\pgfqpoint{4.511757in}{7.466441in}}{\pgfqpoint{4.514071in}{7.460854in}}{\pgfqpoint{4.518189in}{7.456736in}}%
\pgfpathcurveto{\pgfqpoint{4.522308in}{7.452618in}}{\pgfqpoint{4.527894in}{7.450304in}}{\pgfqpoint{4.533718in}{7.450304in}}%
\pgfpathlineto{\pgfqpoint{4.533718in}{7.450304in}}%
\pgfpathclose%
\pgfusepath{stroke,fill}%
\end{pgfscope}%
\begin{pgfscope}%
\pgfpathrectangle{\pgfqpoint{1.414736in}{0.548769in}}{\pgfqpoint{3.497747in}{7.253186in}}%
\pgfusepath{clip}%
\pgfsetbuttcap%
\pgfsetroundjoin%
\definecolor{currentfill}{rgb}{0.616293,0.852709,0.230052}%
\pgfsetfillcolor{currentfill}%
\pgfsetlinewidth{1.003750pt}%
\definecolor{currentstroke}{rgb}{0.616293,0.852709,0.230052}%
\pgfsetstrokecolor{currentstroke}%
\pgfsetdash{}{0pt}%
\pgfpathmoveto{\pgfqpoint{3.317229in}{0.856499in}}%
\pgfpathcurveto{\pgfqpoint{3.323052in}{0.856499in}}{\pgfqpoint{3.328639in}{0.858813in}}{\pgfqpoint{3.332757in}{0.862931in}}%
\pgfpathcurveto{\pgfqpoint{3.336875in}{0.867049in}}{\pgfqpoint{3.339189in}{0.872635in}}{\pgfqpoint{3.339189in}{0.878459in}}%
\pgfpathcurveto{\pgfqpoint{3.339189in}{0.884283in}}{\pgfqpoint{3.336875in}{0.889869in}}{\pgfqpoint{3.332757in}{0.893987in}}%
\pgfpathcurveto{\pgfqpoint{3.328639in}{0.898105in}}{\pgfqpoint{3.323052in}{0.900419in}}{\pgfqpoint{3.317229in}{0.900419in}}%
\pgfpathcurveto{\pgfqpoint{3.311405in}{0.900419in}}{\pgfqpoint{3.305818in}{0.898105in}}{\pgfqpoint{3.301700in}{0.893987in}}%
\pgfpathcurveto{\pgfqpoint{3.297582in}{0.889869in}}{\pgfqpoint{3.295268in}{0.884283in}}{\pgfqpoint{3.295268in}{0.878459in}}%
\pgfpathcurveto{\pgfqpoint{3.295268in}{0.872635in}}{\pgfqpoint{3.297582in}{0.867049in}}{\pgfqpoint{3.301700in}{0.862931in}}%
\pgfpathcurveto{\pgfqpoint{3.305818in}{0.858813in}}{\pgfqpoint{3.311405in}{0.856499in}}{\pgfqpoint{3.317229in}{0.856499in}}%
\pgfpathlineto{\pgfqpoint{3.317229in}{0.856499in}}%
\pgfpathclose%
\pgfusepath{stroke,fill}%
\end{pgfscope}%
\begin{pgfscope}%
\pgfpathrectangle{\pgfqpoint{1.414736in}{0.548769in}}{\pgfqpoint{3.497747in}{7.253186in}}%
\pgfusepath{clip}%
\pgfsetbuttcap%
\pgfsetroundjoin%
\definecolor{currentfill}{rgb}{0.274149,0.751988,0.436601}%
\pgfsetfillcolor{currentfill}%
\pgfsetlinewidth{1.003750pt}%
\definecolor{currentstroke}{rgb}{0.274149,0.751988,0.436601}%
\pgfsetstrokecolor{currentstroke}%
\pgfsetdash{}{0pt}%
\pgfpathmoveto{\pgfqpoint{3.677646in}{1.327485in}}%
\pgfpathcurveto{\pgfqpoint{3.683470in}{1.327485in}}{\pgfqpoint{3.689056in}{1.329799in}}{\pgfqpoint{3.693175in}{1.333917in}}%
\pgfpathcurveto{\pgfqpoint{3.697293in}{1.338035in}}{\pgfqpoint{3.699607in}{1.343621in}}{\pgfqpoint{3.699607in}{1.349445in}}%
\pgfpathcurveto{\pgfqpoint{3.699607in}{1.355269in}}{\pgfqpoint{3.697293in}{1.360855in}}{\pgfqpoint{3.693175in}{1.364973in}}%
\pgfpathcurveto{\pgfqpoint{3.689056in}{1.369092in}}{\pgfqpoint{3.683470in}{1.371405in}}{\pgfqpoint{3.677646in}{1.371405in}}%
\pgfpathcurveto{\pgfqpoint{3.671822in}{1.371405in}}{\pgfqpoint{3.666236in}{1.369092in}}{\pgfqpoint{3.662118in}{1.364973in}}%
\pgfpathcurveto{\pgfqpoint{3.658000in}{1.360855in}}{\pgfqpoint{3.655686in}{1.355269in}}{\pgfqpoint{3.655686in}{1.349445in}}%
\pgfpathcurveto{\pgfqpoint{3.655686in}{1.343621in}}{\pgfqpoint{3.658000in}{1.338035in}}{\pgfqpoint{3.662118in}{1.333917in}}%
\pgfpathcurveto{\pgfqpoint{3.666236in}{1.329799in}}{\pgfqpoint{3.671822in}{1.327485in}}{\pgfqpoint{3.677646in}{1.327485in}}%
\pgfpathlineto{\pgfqpoint{3.677646in}{1.327485in}}%
\pgfpathclose%
\pgfusepath{stroke,fill}%
\end{pgfscope}%
\begin{pgfscope}%
\pgfpathrectangle{\pgfqpoint{1.414736in}{0.548769in}}{\pgfqpoint{3.497747in}{7.253186in}}%
\pgfusepath{clip}%
\pgfsetbuttcap%
\pgfsetroundjoin%
\definecolor{currentfill}{rgb}{0.855810,0.888601,0.097452}%
\pgfsetfillcolor{currentfill}%
\pgfsetlinewidth{1.003750pt}%
\definecolor{currentstroke}{rgb}{0.855810,0.888601,0.097452}%
\pgfsetstrokecolor{currentstroke}%
\pgfsetdash{}{0pt}%
\pgfpathmoveto{\pgfqpoint{3.796943in}{1.798471in}}%
\pgfpathcurveto{\pgfqpoint{3.802767in}{1.798471in}}{\pgfqpoint{3.808353in}{1.800785in}}{\pgfqpoint{3.812471in}{1.804903in}}%
\pgfpathcurveto{\pgfqpoint{3.816589in}{1.809021in}}{\pgfqpoint{3.818903in}{1.814607in}}{\pgfqpoint{3.818903in}{1.820431in}}%
\pgfpathcurveto{\pgfqpoint{3.818903in}{1.826255in}}{\pgfqpoint{3.816589in}{1.831841in}}{\pgfqpoint{3.812471in}{1.835960in}}%
\pgfpathcurveto{\pgfqpoint{3.808353in}{1.840078in}}{\pgfqpoint{3.802767in}{1.842392in}}{\pgfqpoint{3.796943in}{1.842392in}}%
\pgfpathcurveto{\pgfqpoint{3.791119in}{1.842392in}}{\pgfqpoint{3.785533in}{1.840078in}}{\pgfqpoint{3.781415in}{1.835960in}}%
\pgfpathcurveto{\pgfqpoint{3.777296in}{1.831841in}}{\pgfqpoint{3.774983in}{1.826255in}}{\pgfqpoint{3.774983in}{1.820431in}}%
\pgfpathcurveto{\pgfqpoint{3.774983in}{1.814607in}}{\pgfqpoint{3.777296in}{1.809021in}}{\pgfqpoint{3.781415in}{1.804903in}}%
\pgfpathcurveto{\pgfqpoint{3.785533in}{1.800785in}}{\pgfqpoint{3.791119in}{1.798471in}}{\pgfqpoint{3.796943in}{1.798471in}}%
\pgfpathlineto{\pgfqpoint{3.796943in}{1.798471in}}%
\pgfpathclose%
\pgfusepath{stroke,fill}%
\end{pgfscope}%
\begin{pgfscope}%
\pgfpathrectangle{\pgfqpoint{1.414736in}{0.548769in}}{\pgfqpoint{3.497747in}{7.253186in}}%
\pgfusepath{clip}%
\pgfsetbuttcap%
\pgfsetroundjoin%
\definecolor{currentfill}{rgb}{0.993248,0.906157,0.143936}%
\pgfsetfillcolor{currentfill}%
\pgfsetlinewidth{1.003750pt}%
\definecolor{currentstroke}{rgb}{0.993248,0.906157,0.143936}%
\pgfsetstrokecolor{currentstroke}%
\pgfsetdash{}{0pt}%
\pgfpathmoveto{\pgfqpoint{3.845570in}{2.269457in}}%
\pgfpathcurveto{\pgfqpoint{3.851394in}{2.269457in}}{\pgfqpoint{3.856980in}{2.271771in}}{\pgfqpoint{3.861098in}{2.275889in}}%
\pgfpathcurveto{\pgfqpoint{3.865216in}{2.280007in}}{\pgfqpoint{3.867530in}{2.285593in}}{\pgfqpoint{3.867530in}{2.291417in}}%
\pgfpathcurveto{\pgfqpoint{3.867530in}{2.297241in}}{\pgfqpoint{3.865216in}{2.302828in}}{\pgfqpoint{3.861098in}{2.306946in}}%
\pgfpathcurveto{\pgfqpoint{3.856980in}{2.311064in}}{\pgfqpoint{3.851394in}{2.313378in}}{\pgfqpoint{3.845570in}{2.313378in}}%
\pgfpathcurveto{\pgfqpoint{3.839746in}{2.313378in}}{\pgfqpoint{3.834160in}{2.311064in}}{\pgfqpoint{3.830042in}{2.306946in}}%
\pgfpathcurveto{\pgfqpoint{3.825923in}{2.302828in}}{\pgfqpoint{3.823610in}{2.297241in}}{\pgfqpoint{3.823610in}{2.291417in}}%
\pgfpathcurveto{\pgfqpoint{3.823610in}{2.285593in}}{\pgfqpoint{3.825923in}{2.280007in}}{\pgfqpoint{3.830042in}{2.275889in}}%
\pgfpathcurveto{\pgfqpoint{3.834160in}{2.271771in}}{\pgfqpoint{3.839746in}{2.269457in}}{\pgfqpoint{3.845570in}{2.269457in}}%
\pgfpathlineto{\pgfqpoint{3.845570in}{2.269457in}}%
\pgfpathclose%
\pgfusepath{stroke,fill}%
\end{pgfscope}%
\begin{pgfscope}%
\pgfpathrectangle{\pgfqpoint{1.414736in}{0.548769in}}{\pgfqpoint{3.497747in}{7.253186in}}%
\pgfusepath{clip}%
\pgfsetbuttcap%
\pgfsetroundjoin%
\definecolor{currentfill}{rgb}{0.616293,0.852709,0.230052}%
\pgfsetfillcolor{currentfill}%
\pgfsetlinewidth{1.003750pt}%
\definecolor{currentstroke}{rgb}{0.616293,0.852709,0.230052}%
\pgfsetstrokecolor{currentstroke}%
\pgfsetdash{}{0pt}%
\pgfpathmoveto{\pgfqpoint{3.131652in}{2.740443in}}%
\pgfpathcurveto{\pgfqpoint{3.137476in}{2.740443in}}{\pgfqpoint{3.143062in}{2.742757in}}{\pgfqpoint{3.147180in}{2.746875in}}%
\pgfpathcurveto{\pgfqpoint{3.151298in}{2.750993in}}{\pgfqpoint{3.153612in}{2.756580in}}{\pgfqpoint{3.153612in}{2.762404in}}%
\pgfpathcurveto{\pgfqpoint{3.153612in}{2.768227in}}{\pgfqpoint{3.151298in}{2.773814in}}{\pgfqpoint{3.147180in}{2.777932in}}%
\pgfpathcurveto{\pgfqpoint{3.143062in}{2.782050in}}{\pgfqpoint{3.137476in}{2.784364in}}{\pgfqpoint{3.131652in}{2.784364in}}%
\pgfpathcurveto{\pgfqpoint{3.125828in}{2.784364in}}{\pgfqpoint{3.120242in}{2.782050in}}{\pgfqpoint{3.116124in}{2.777932in}}%
\pgfpathcurveto{\pgfqpoint{3.112006in}{2.773814in}}{\pgfqpoint{3.109692in}{2.768227in}}{\pgfqpoint{3.109692in}{2.762404in}}%
\pgfpathcurveto{\pgfqpoint{3.109692in}{2.756580in}}{\pgfqpoint{3.112006in}{2.750993in}}{\pgfqpoint{3.116124in}{2.746875in}}%
\pgfpathcurveto{\pgfqpoint{3.120242in}{2.742757in}}{\pgfqpoint{3.125828in}{2.740443in}}{\pgfqpoint{3.131652in}{2.740443in}}%
\pgfpathlineto{\pgfqpoint{3.131652in}{2.740443in}}%
\pgfpathclose%
\pgfusepath{stroke,fill}%
\end{pgfscope}%
\begin{pgfscope}%
\pgfpathrectangle{\pgfqpoint{1.414736in}{0.548769in}}{\pgfqpoint{3.497747in}{7.253186in}}%
\pgfusepath{clip}%
\pgfsetbuttcap%
\pgfsetroundjoin%
\definecolor{currentfill}{rgb}{0.430983,0.808473,0.346476}%
\pgfsetfillcolor{currentfill}%
\pgfsetlinewidth{1.003750pt}%
\definecolor{currentstroke}{rgb}{0.430983,0.808473,0.346476}%
\pgfsetstrokecolor{currentstroke}%
\pgfsetdash{}{0pt}%
\pgfpathmoveto{\pgfqpoint{4.205923in}{3.211429in}}%
\pgfpathcurveto{\pgfqpoint{4.211747in}{3.211429in}}{\pgfqpoint{4.217333in}{3.213743in}}{\pgfqpoint{4.221451in}{3.217861in}}%
\pgfpathcurveto{\pgfqpoint{4.225570in}{3.221980in}}{\pgfqpoint{4.227883in}{3.227566in}}{\pgfqpoint{4.227883in}{3.233390in}}%
\pgfpathcurveto{\pgfqpoint{4.227883in}{3.239214in}}{\pgfqpoint{4.225570in}{3.244800in}}{\pgfqpoint{4.221451in}{3.248918in}}%
\pgfpathcurveto{\pgfqpoint{4.217333in}{3.253036in}}{\pgfqpoint{4.211747in}{3.255350in}}{\pgfqpoint{4.205923in}{3.255350in}}%
\pgfpathcurveto{\pgfqpoint{4.200099in}{3.255350in}}{\pgfqpoint{4.194513in}{3.253036in}}{\pgfqpoint{4.190395in}{3.248918in}}%
\pgfpathcurveto{\pgfqpoint{4.186277in}{3.244800in}}{\pgfqpoint{4.183963in}{3.239214in}}{\pgfqpoint{4.183963in}{3.233390in}}%
\pgfpathcurveto{\pgfqpoint{4.183963in}{3.227566in}}{\pgfqpoint{4.186277in}{3.221980in}}{\pgfqpoint{4.190395in}{3.217861in}}%
\pgfpathcurveto{\pgfqpoint{4.194513in}{3.213743in}}{\pgfqpoint{4.200099in}{3.211429in}}{\pgfqpoint{4.205923in}{3.211429in}}%
\pgfpathlineto{\pgfqpoint{4.205923in}{3.211429in}}%
\pgfpathclose%
\pgfusepath{stroke,fill}%
\end{pgfscope}%
\begin{pgfscope}%
\pgfpathrectangle{\pgfqpoint{1.414736in}{0.548769in}}{\pgfqpoint{3.497747in}{7.253186in}}%
\pgfusepath{clip}%
\pgfsetbuttcap%
\pgfsetroundjoin%
\definecolor{currentfill}{rgb}{0.993248,0.906157,0.143936}%
\pgfsetfillcolor{currentfill}%
\pgfsetlinewidth{1.003750pt}%
\definecolor{currentstroke}{rgb}{0.993248,0.906157,0.143936}%
\pgfsetstrokecolor{currentstroke}%
\pgfsetdash{}{0pt}%
\pgfpathmoveto{\pgfqpoint{3.418301in}{3.682415in}}%
\pgfpathcurveto{\pgfqpoint{3.424125in}{3.682415in}}{\pgfqpoint{3.429711in}{3.684729in}}{\pgfqpoint{3.433830in}{3.688847in}}%
\pgfpathcurveto{\pgfqpoint{3.437948in}{3.692966in}}{\pgfqpoint{3.440262in}{3.698552in}}{\pgfqpoint{3.440262in}{3.704376in}}%
\pgfpathcurveto{\pgfqpoint{3.440262in}{3.710200in}}{\pgfqpoint{3.437948in}{3.715786in}}{\pgfqpoint{3.433830in}{3.719904in}}%
\pgfpathcurveto{\pgfqpoint{3.429711in}{3.724022in}}{\pgfqpoint{3.424125in}{3.726336in}}{\pgfqpoint{3.418301in}{3.726336in}}%
\pgfpathcurveto{\pgfqpoint{3.412477in}{3.726336in}}{\pgfqpoint{3.406891in}{3.724022in}}{\pgfqpoint{3.402773in}{3.719904in}}%
\pgfpathcurveto{\pgfqpoint{3.398655in}{3.715786in}}{\pgfqpoint{3.396341in}{3.710200in}}{\pgfqpoint{3.396341in}{3.704376in}}%
\pgfpathcurveto{\pgfqpoint{3.396341in}{3.698552in}}{\pgfqpoint{3.398655in}{3.692966in}}{\pgfqpoint{3.402773in}{3.688847in}}%
\pgfpathcurveto{\pgfqpoint{3.406891in}{3.684729in}}{\pgfqpoint{3.412477in}{3.682415in}}{\pgfqpoint{3.418301in}{3.682415in}}%
\pgfpathlineto{\pgfqpoint{3.418301in}{3.682415in}}%
\pgfpathclose%
\pgfusepath{stroke,fill}%
\end{pgfscope}%
\begin{pgfscope}%
\pgfpathrectangle{\pgfqpoint{1.414736in}{0.548769in}}{\pgfqpoint{3.497747in}{7.253186in}}%
\pgfusepath{clip}%
\pgfsetbuttcap%
\pgfsetroundjoin%
\definecolor{currentfill}{rgb}{0.993248,0.906157,0.143936}%
\pgfsetfillcolor{currentfill}%
\pgfsetlinewidth{1.003750pt}%
\definecolor{currentstroke}{rgb}{0.993248,0.906157,0.143936}%
\pgfsetstrokecolor{currentstroke}%
\pgfsetdash{}{0pt}%
\pgfpathmoveto{\pgfqpoint{3.896067in}{4.153402in}}%
\pgfpathcurveto{\pgfqpoint{3.901891in}{4.153402in}}{\pgfqpoint{3.907477in}{4.155715in}}{\pgfqpoint{3.911595in}{4.159834in}}%
\pgfpathcurveto{\pgfqpoint{3.915714in}{4.163952in}}{\pgfqpoint{3.918027in}{4.169538in}}{\pgfqpoint{3.918027in}{4.175362in}}%
\pgfpathcurveto{\pgfqpoint{3.918027in}{4.181186in}}{\pgfqpoint{3.915714in}{4.186772in}}{\pgfqpoint{3.911595in}{4.190890in}}%
\pgfpathcurveto{\pgfqpoint{3.907477in}{4.195008in}}{\pgfqpoint{3.901891in}{4.197322in}}{\pgfqpoint{3.896067in}{4.197322in}}%
\pgfpathcurveto{\pgfqpoint{3.890243in}{4.197322in}}{\pgfqpoint{3.884657in}{4.195008in}}{\pgfqpoint{3.880539in}{4.190890in}}%
\pgfpathcurveto{\pgfqpoint{3.876421in}{4.186772in}}{\pgfqpoint{3.874107in}{4.181186in}}{\pgfqpoint{3.874107in}{4.175362in}}%
\pgfpathcurveto{\pgfqpoint{3.874107in}{4.169538in}}{\pgfqpoint{3.876421in}{4.163952in}}{\pgfqpoint{3.880539in}{4.159834in}}%
\pgfpathcurveto{\pgfqpoint{3.884657in}{4.155715in}}{\pgfqpoint{3.890243in}{4.153402in}}{\pgfqpoint{3.896067in}{4.153402in}}%
\pgfpathlineto{\pgfqpoint{3.896067in}{4.153402in}}%
\pgfpathclose%
\pgfusepath{stroke,fill}%
\end{pgfscope}%
\begin{pgfscope}%
\pgfpathrectangle{\pgfqpoint{1.414736in}{0.548769in}}{\pgfqpoint{3.497747in}{7.253186in}}%
\pgfusepath{clip}%
\pgfsetbuttcap%
\pgfsetroundjoin%
\definecolor{currentfill}{rgb}{0.506271,0.828786,0.300362}%
\pgfsetfillcolor{currentfill}%
\pgfsetlinewidth{1.003750pt}%
\definecolor{currentstroke}{rgb}{0.506271,0.828786,0.300362}%
\pgfsetstrokecolor{currentstroke}%
\pgfsetdash{}{0pt}%
\pgfpathmoveto{\pgfqpoint{3.582663in}{4.624388in}}%
\pgfpathcurveto{\pgfqpoint{3.588487in}{4.624388in}}{\pgfqpoint{3.594073in}{4.626702in}}{\pgfqpoint{3.598192in}{4.630820in}}%
\pgfpathcurveto{\pgfqpoint{3.602310in}{4.634938in}}{\pgfqpoint{3.604624in}{4.640524in}}{\pgfqpoint{3.604624in}{4.646348in}}%
\pgfpathcurveto{\pgfqpoint{3.604624in}{4.652172in}}{\pgfqpoint{3.602310in}{4.657758in}}{\pgfqpoint{3.598192in}{4.661876in}}%
\pgfpathcurveto{\pgfqpoint{3.594073in}{4.665994in}}{\pgfqpoint{3.588487in}{4.668308in}}{\pgfqpoint{3.582663in}{4.668308in}}%
\pgfpathcurveto{\pgfqpoint{3.576839in}{4.668308in}}{\pgfqpoint{3.571253in}{4.665994in}}{\pgfqpoint{3.567135in}{4.661876in}}%
\pgfpathcurveto{\pgfqpoint{3.563017in}{4.657758in}}{\pgfqpoint{3.560703in}{4.652172in}}{\pgfqpoint{3.560703in}{4.646348in}}%
\pgfpathcurveto{\pgfqpoint{3.560703in}{4.640524in}}{\pgfqpoint{3.563017in}{4.634938in}}{\pgfqpoint{3.567135in}{4.630820in}}%
\pgfpathcurveto{\pgfqpoint{3.571253in}{4.626702in}}{\pgfqpoint{3.576839in}{4.624388in}}{\pgfqpoint{3.582663in}{4.624388in}}%
\pgfpathlineto{\pgfqpoint{3.582663in}{4.624388in}}%
\pgfpathclose%
\pgfusepath{stroke,fill}%
\end{pgfscope}%
\begin{pgfscope}%
\pgfpathrectangle{\pgfqpoint{1.414736in}{0.548769in}}{\pgfqpoint{3.497747in}{7.253186in}}%
\pgfusepath{clip}%
\pgfsetbuttcap%
\pgfsetroundjoin%
\definecolor{currentfill}{rgb}{0.149039,0.508051,0.557250}%
\pgfsetfillcolor{currentfill}%
\pgfsetlinewidth{1.003750pt}%
\definecolor{currentstroke}{rgb}{0.149039,0.508051,0.557250}%
\pgfsetstrokecolor{currentstroke}%
\pgfsetdash{}{0pt}%
\pgfpathmoveto{\pgfqpoint{3.511192in}{5.095374in}}%
\pgfpathcurveto{\pgfqpoint{3.517016in}{5.095374in}}{\pgfqpoint{3.522603in}{5.097688in}}{\pgfqpoint{3.526721in}{5.101806in}}%
\pgfpathcurveto{\pgfqpoint{3.530839in}{5.105924in}}{\pgfqpoint{3.533153in}{5.111510in}}{\pgfqpoint{3.533153in}{5.117334in}}%
\pgfpathcurveto{\pgfqpoint{3.533153in}{5.123158in}}{\pgfqpoint{3.530839in}{5.128744in}}{\pgfqpoint{3.526721in}{5.132862in}}%
\pgfpathcurveto{\pgfqpoint{3.522603in}{5.136980in}}{\pgfqpoint{3.517016in}{5.139294in}}{\pgfqpoint{3.511192in}{5.139294in}}%
\pgfpathcurveto{\pgfqpoint{3.505368in}{5.139294in}}{\pgfqpoint{3.499782in}{5.136980in}}{\pgfqpoint{3.495664in}{5.132862in}}%
\pgfpathcurveto{\pgfqpoint{3.491546in}{5.128744in}}{\pgfqpoint{3.489232in}{5.123158in}}{\pgfqpoint{3.489232in}{5.117334in}}%
\pgfpathcurveto{\pgfqpoint{3.489232in}{5.111510in}}{\pgfqpoint{3.491546in}{5.105924in}}{\pgfqpoint{3.495664in}{5.101806in}}%
\pgfpathcurveto{\pgfqpoint{3.499782in}{5.097688in}}{\pgfqpoint{3.505368in}{5.095374in}}{\pgfqpoint{3.511192in}{5.095374in}}%
\pgfpathlineto{\pgfqpoint{3.511192in}{5.095374in}}%
\pgfpathclose%
\pgfusepath{stroke,fill}%
\end{pgfscope}%
\begin{pgfscope}%
\pgfpathrectangle{\pgfqpoint{1.414736in}{0.548769in}}{\pgfqpoint{3.497747in}{7.253186in}}%
\pgfusepath{clip}%
\pgfsetbuttcap%
\pgfsetroundjoin%
\definecolor{currentfill}{rgb}{0.180629,0.429975,0.557282}%
\pgfsetfillcolor{currentfill}%
\pgfsetlinewidth{1.003750pt}%
\definecolor{currentstroke}{rgb}{0.180629,0.429975,0.557282}%
\pgfsetstrokecolor{currentstroke}%
\pgfsetdash{}{0pt}%
\pgfpathmoveto{\pgfqpoint{3.729507in}{5.566360in}}%
\pgfpathcurveto{\pgfqpoint{3.735331in}{5.566360in}}{\pgfqpoint{3.740917in}{5.568674in}}{\pgfqpoint{3.745035in}{5.572792in}}%
\pgfpathcurveto{\pgfqpoint{3.749153in}{5.576910in}}{\pgfqpoint{3.751467in}{5.582496in}}{\pgfqpoint{3.751467in}{5.588320in}}%
\pgfpathcurveto{\pgfqpoint{3.751467in}{5.594144in}}{\pgfqpoint{3.749153in}{5.599730in}}{\pgfqpoint{3.745035in}{5.603848in}}%
\pgfpathcurveto{\pgfqpoint{3.740917in}{5.607967in}}{\pgfqpoint{3.735331in}{5.610280in}}{\pgfqpoint{3.729507in}{5.610280in}}%
\pgfpathcurveto{\pgfqpoint{3.723683in}{5.610280in}}{\pgfqpoint{3.718097in}{5.607967in}}{\pgfqpoint{3.713979in}{5.603848in}}%
\pgfpathcurveto{\pgfqpoint{3.709861in}{5.599730in}}{\pgfqpoint{3.707547in}{5.594144in}}{\pgfqpoint{3.707547in}{5.588320in}}%
\pgfpathcurveto{\pgfqpoint{3.707547in}{5.582496in}}{\pgfqpoint{3.709861in}{5.576910in}}{\pgfqpoint{3.713979in}{5.572792in}}%
\pgfpathcurveto{\pgfqpoint{3.718097in}{5.568674in}}{\pgfqpoint{3.723683in}{5.566360in}}{\pgfqpoint{3.729507in}{5.566360in}}%
\pgfpathlineto{\pgfqpoint{3.729507in}{5.566360in}}%
\pgfpathclose%
\pgfusepath{stroke,fill}%
\end{pgfscope}%
\begin{pgfscope}%
\pgfpathrectangle{\pgfqpoint{1.414736in}{0.548769in}}{\pgfqpoint{3.497747in}{7.253186in}}%
\pgfusepath{clip}%
\pgfsetbuttcap%
\pgfsetroundjoin%
\definecolor{currentfill}{rgb}{0.288921,0.758394,0.428426}%
\pgfsetfillcolor{currentfill}%
\pgfsetlinewidth{1.003750pt}%
\definecolor{currentstroke}{rgb}{0.288921,0.758394,0.428426}%
\pgfsetstrokecolor{currentstroke}%
\pgfsetdash{}{0pt}%
\pgfpathmoveto{\pgfqpoint{3.689543in}{6.037346in}}%
\pgfpathcurveto{\pgfqpoint{3.695367in}{6.037346in}}{\pgfqpoint{3.700953in}{6.039660in}}{\pgfqpoint{3.705071in}{6.043778in}}%
\pgfpathcurveto{\pgfqpoint{3.709189in}{6.047896in}}{\pgfqpoint{3.711503in}{6.053482in}}{\pgfqpoint{3.711503in}{6.059306in}}%
\pgfpathcurveto{\pgfqpoint{3.711503in}{6.065130in}}{\pgfqpoint{3.709189in}{6.070716in}}{\pgfqpoint{3.705071in}{6.074835in}}%
\pgfpathcurveto{\pgfqpoint{3.700953in}{6.078953in}}{\pgfqpoint{3.695367in}{6.081267in}}{\pgfqpoint{3.689543in}{6.081267in}}%
\pgfpathcurveto{\pgfqpoint{3.683719in}{6.081267in}}{\pgfqpoint{3.678133in}{6.078953in}}{\pgfqpoint{3.674015in}{6.074835in}}%
\pgfpathcurveto{\pgfqpoint{3.669896in}{6.070716in}}{\pgfqpoint{3.667583in}{6.065130in}}{\pgfqpoint{3.667583in}{6.059306in}}%
\pgfpathcurveto{\pgfqpoint{3.667583in}{6.053482in}}{\pgfqpoint{3.669896in}{6.047896in}}{\pgfqpoint{3.674015in}{6.043778in}}%
\pgfpathcurveto{\pgfqpoint{3.678133in}{6.039660in}}{\pgfqpoint{3.683719in}{6.037346in}}{\pgfqpoint{3.689543in}{6.037346in}}%
\pgfpathlineto{\pgfqpoint{3.689543in}{6.037346in}}%
\pgfpathclose%
\pgfusepath{stroke,fill}%
\end{pgfscope}%
\begin{pgfscope}%
\pgfpathrectangle{\pgfqpoint{1.414736in}{0.548769in}}{\pgfqpoint{3.497747in}{7.253186in}}%
\pgfusepath{clip}%
\pgfsetbuttcap%
\pgfsetroundjoin%
\definecolor{currentfill}{rgb}{0.993248,0.906157,0.143936}%
\pgfsetfillcolor{currentfill}%
\pgfsetlinewidth{1.003750pt}%
\definecolor{currentstroke}{rgb}{0.993248,0.906157,0.143936}%
\pgfsetstrokecolor{currentstroke}%
\pgfsetdash{}{0pt}%
\pgfpathmoveto{\pgfqpoint{3.485800in}{6.508332in}}%
\pgfpathcurveto{\pgfqpoint{3.491624in}{6.508332in}}{\pgfqpoint{3.497210in}{6.510646in}}{\pgfqpoint{3.501328in}{6.514764in}}%
\pgfpathcurveto{\pgfqpoint{3.505446in}{6.518882in}}{\pgfqpoint{3.507760in}{6.524468in}}{\pgfqpoint{3.507760in}{6.530292in}}%
\pgfpathcurveto{\pgfqpoint{3.507760in}{6.536116in}}{\pgfqpoint{3.505446in}{6.541702in}}{\pgfqpoint{3.501328in}{6.545821in}}%
\pgfpathcurveto{\pgfqpoint{3.497210in}{6.549939in}}{\pgfqpoint{3.491624in}{6.552253in}}{\pgfqpoint{3.485800in}{6.552253in}}%
\pgfpathcurveto{\pgfqpoint{3.479976in}{6.552253in}}{\pgfqpoint{3.474390in}{6.549939in}}{\pgfqpoint{3.470272in}{6.545821in}}%
\pgfpathcurveto{\pgfqpoint{3.466153in}{6.541702in}}{\pgfqpoint{3.463840in}{6.536116in}}{\pgfqpoint{3.463840in}{6.530292in}}%
\pgfpathcurveto{\pgfqpoint{3.463840in}{6.524468in}}{\pgfqpoint{3.466153in}{6.518882in}}{\pgfqpoint{3.470272in}{6.514764in}}%
\pgfpathcurveto{\pgfqpoint{3.474390in}{6.510646in}}{\pgfqpoint{3.479976in}{6.508332in}}{\pgfqpoint{3.485800in}{6.508332in}}%
\pgfpathlineto{\pgfqpoint{3.485800in}{6.508332in}}%
\pgfpathclose%
\pgfusepath{stroke,fill}%
\end{pgfscope}%
\begin{pgfscope}%
\pgfpathrectangle{\pgfqpoint{1.414736in}{0.548769in}}{\pgfqpoint{3.497747in}{7.253186in}}%
\pgfusepath{clip}%
\pgfsetbuttcap%
\pgfsetroundjoin%
\definecolor{currentfill}{rgb}{0.964894,0.902323,0.123941}%
\pgfsetfillcolor{currentfill}%
\pgfsetlinewidth{1.003750pt}%
\definecolor{currentstroke}{rgb}{0.964894,0.902323,0.123941}%
\pgfsetstrokecolor{currentstroke}%
\pgfsetdash{}{0pt}%
\pgfpathmoveto{\pgfqpoint{3.567820in}{6.979318in}}%
\pgfpathcurveto{\pgfqpoint{3.573644in}{6.979318in}}{\pgfqpoint{3.579230in}{6.981632in}}{\pgfqpoint{3.583348in}{6.985750in}}%
\pgfpathcurveto{\pgfqpoint{3.587466in}{6.989868in}}{\pgfqpoint{3.589780in}{6.995455in}}{\pgfqpoint{3.589780in}{7.001278in}}%
\pgfpathcurveto{\pgfqpoint{3.589780in}{7.007102in}}{\pgfqpoint{3.587466in}{7.012689in}}{\pgfqpoint{3.583348in}{7.016807in}}%
\pgfpathcurveto{\pgfqpoint{3.579230in}{7.020925in}}{\pgfqpoint{3.573644in}{7.023239in}}{\pgfqpoint{3.567820in}{7.023239in}}%
\pgfpathcurveto{\pgfqpoint{3.561996in}{7.023239in}}{\pgfqpoint{3.556410in}{7.020925in}}{\pgfqpoint{3.552292in}{7.016807in}}%
\pgfpathcurveto{\pgfqpoint{3.548174in}{7.012689in}}{\pgfqpoint{3.545860in}{7.007102in}}{\pgfqpoint{3.545860in}{7.001278in}}%
\pgfpathcurveto{\pgfqpoint{3.545860in}{6.995455in}}{\pgfqpoint{3.548174in}{6.989868in}}{\pgfqpoint{3.552292in}{6.985750in}}%
\pgfpathcurveto{\pgfqpoint{3.556410in}{6.981632in}}{\pgfqpoint{3.561996in}{6.979318in}}{\pgfqpoint{3.567820in}{6.979318in}}%
\pgfpathlineto{\pgfqpoint{3.567820in}{6.979318in}}%
\pgfpathclose%
\pgfusepath{stroke,fill}%
\end{pgfscope}%
\begin{pgfscope}%
\pgfpathrectangle{\pgfqpoint{1.414736in}{0.548769in}}{\pgfqpoint{3.497747in}{7.253186in}}%
\pgfusepath{clip}%
\pgfsetbuttcap%
\pgfsetroundjoin%
\definecolor{currentfill}{rgb}{0.468053,0.818921,0.323998}%
\pgfsetfillcolor{currentfill}%
\pgfsetlinewidth{1.003750pt}%
\definecolor{currentstroke}{rgb}{0.468053,0.818921,0.323998}%
\pgfsetstrokecolor{currentstroke}%
\pgfsetdash{}{0pt}%
\pgfpathmoveto{\pgfqpoint{3.710845in}{7.450304in}}%
\pgfpathcurveto{\pgfqpoint{3.716669in}{7.450304in}}{\pgfqpoint{3.722256in}{7.452618in}}{\pgfqpoint{3.726374in}{7.456736in}}%
\pgfpathcurveto{\pgfqpoint{3.730492in}{7.460854in}}{\pgfqpoint{3.732806in}{7.466441in}}{\pgfqpoint{3.732806in}{7.472265in}}%
\pgfpathcurveto{\pgfqpoint{3.732806in}{7.478089in}}{\pgfqpoint{3.730492in}{7.483675in}}{\pgfqpoint{3.726374in}{7.487793in}}%
\pgfpathcurveto{\pgfqpoint{3.722256in}{7.491911in}}{\pgfqpoint{3.716669in}{7.494225in}}{\pgfqpoint{3.710845in}{7.494225in}}%
\pgfpathcurveto{\pgfqpoint{3.705022in}{7.494225in}}{\pgfqpoint{3.699435in}{7.491911in}}{\pgfqpoint{3.695317in}{7.487793in}}%
\pgfpathcurveto{\pgfqpoint{3.691199in}{7.483675in}}{\pgfqpoint{3.688885in}{7.478089in}}{\pgfqpoint{3.688885in}{7.472265in}}%
\pgfpathcurveto{\pgfqpoint{3.688885in}{7.466441in}}{\pgfqpoint{3.691199in}{7.460854in}}{\pgfqpoint{3.695317in}{7.456736in}}%
\pgfpathcurveto{\pgfqpoint{3.699435in}{7.452618in}}{\pgfqpoint{3.705022in}{7.450304in}}{\pgfqpoint{3.710845in}{7.450304in}}%
\pgfpathlineto{\pgfqpoint{3.710845in}{7.450304in}}%
\pgfpathclose%
\pgfusepath{stroke,fill}%
\end{pgfscope}%
\begin{pgfscope}%
\pgfpathrectangle{\pgfqpoint{1.414736in}{0.548769in}}{\pgfqpoint{3.497747in}{7.253186in}}%
\pgfusepath{clip}%
\pgfsetbuttcap%
\pgfsetroundjoin%
\definecolor{currentfill}{rgb}{0.993248,0.906157,0.143936}%
\pgfsetfillcolor{currentfill}%
\pgfsetlinewidth{1.003750pt}%
\definecolor{currentstroke}{rgb}{0.993248,0.906157,0.143936}%
\pgfsetstrokecolor{currentstroke}%
\pgfsetdash{}{0pt}%
\pgfpathmoveto{\pgfqpoint{3.439460in}{0.856499in}}%
\pgfpathcurveto{\pgfqpoint{3.445284in}{0.856499in}}{\pgfqpoint{3.450870in}{0.858813in}}{\pgfqpoint{3.454988in}{0.862931in}}%
\pgfpathcurveto{\pgfqpoint{3.459106in}{0.867049in}}{\pgfqpoint{3.461420in}{0.872635in}}{\pgfqpoint{3.461420in}{0.878459in}}%
\pgfpathcurveto{\pgfqpoint{3.461420in}{0.884283in}}{\pgfqpoint{3.459106in}{0.889869in}}{\pgfqpoint{3.454988in}{0.893987in}}%
\pgfpathcurveto{\pgfqpoint{3.450870in}{0.898105in}}{\pgfqpoint{3.445284in}{0.900419in}}{\pgfqpoint{3.439460in}{0.900419in}}%
\pgfpathcurveto{\pgfqpoint{3.433636in}{0.900419in}}{\pgfqpoint{3.428050in}{0.898105in}}{\pgfqpoint{3.423931in}{0.893987in}}%
\pgfpathcurveto{\pgfqpoint{3.419813in}{0.889869in}}{\pgfqpoint{3.417499in}{0.884283in}}{\pgfqpoint{3.417499in}{0.878459in}}%
\pgfpathcurveto{\pgfqpoint{3.417499in}{0.872635in}}{\pgfqpoint{3.419813in}{0.867049in}}{\pgfqpoint{3.423931in}{0.862931in}}%
\pgfpathcurveto{\pgfqpoint{3.428050in}{0.858813in}}{\pgfqpoint{3.433636in}{0.856499in}}{\pgfqpoint{3.439460in}{0.856499in}}%
\pgfpathlineto{\pgfqpoint{3.439460in}{0.856499in}}%
\pgfpathclose%
\pgfusepath{stroke,fill}%
\end{pgfscope}%
\begin{pgfscope}%
\pgfpathrectangle{\pgfqpoint{1.414736in}{0.548769in}}{\pgfqpoint{3.497747in}{7.253186in}}%
\pgfusepath{clip}%
\pgfsetbuttcap%
\pgfsetroundjoin%
\definecolor{currentfill}{rgb}{0.404001,0.800275,0.362552}%
\pgfsetfillcolor{currentfill}%
\pgfsetlinewidth{1.003750pt}%
\definecolor{currentstroke}{rgb}{0.404001,0.800275,0.362552}%
\pgfsetstrokecolor{currentstroke}%
\pgfsetdash{}{0pt}%
\pgfpathmoveto{\pgfqpoint{3.745987in}{1.327485in}}%
\pgfpathcurveto{\pgfqpoint{3.751811in}{1.327485in}}{\pgfqpoint{3.757397in}{1.329799in}}{\pgfqpoint{3.761515in}{1.333917in}}%
\pgfpathcurveto{\pgfqpoint{3.765633in}{1.338035in}}{\pgfqpoint{3.767947in}{1.343621in}}{\pgfqpoint{3.767947in}{1.349445in}}%
\pgfpathcurveto{\pgfqpoint{3.767947in}{1.355269in}}{\pgfqpoint{3.765633in}{1.360855in}}{\pgfqpoint{3.761515in}{1.364973in}}%
\pgfpathcurveto{\pgfqpoint{3.757397in}{1.369092in}}{\pgfqpoint{3.751811in}{1.371405in}}{\pgfqpoint{3.745987in}{1.371405in}}%
\pgfpathcurveto{\pgfqpoint{3.740163in}{1.371405in}}{\pgfqpoint{3.734577in}{1.369092in}}{\pgfqpoint{3.730459in}{1.364973in}}%
\pgfpathcurveto{\pgfqpoint{3.726341in}{1.360855in}}{\pgfqpoint{3.724027in}{1.355269in}}{\pgfqpoint{3.724027in}{1.349445in}}%
\pgfpathcurveto{\pgfqpoint{3.724027in}{1.343621in}}{\pgfqpoint{3.726341in}{1.338035in}}{\pgfqpoint{3.730459in}{1.333917in}}%
\pgfpathcurveto{\pgfqpoint{3.734577in}{1.329799in}}{\pgfqpoint{3.740163in}{1.327485in}}{\pgfqpoint{3.745987in}{1.327485in}}%
\pgfpathlineto{\pgfqpoint{3.745987in}{1.327485in}}%
\pgfpathclose%
\pgfusepath{stroke,fill}%
\end{pgfscope}%
\begin{pgfscope}%
\pgfpathrectangle{\pgfqpoint{1.414736in}{0.548769in}}{\pgfqpoint{3.497747in}{7.253186in}}%
\pgfusepath{clip}%
\pgfsetbuttcap%
\pgfsetroundjoin%
\definecolor{currentfill}{rgb}{0.449368,0.813768,0.335384}%
\pgfsetfillcolor{currentfill}%
\pgfsetlinewidth{1.003750pt}%
\definecolor{currentstroke}{rgb}{0.449368,0.813768,0.335384}%
\pgfsetstrokecolor{currentstroke}%
\pgfsetdash{}{0pt}%
\pgfpathmoveto{\pgfqpoint{3.512546in}{1.798471in}}%
\pgfpathcurveto{\pgfqpoint{3.518370in}{1.798471in}}{\pgfqpoint{3.523956in}{1.800785in}}{\pgfqpoint{3.528074in}{1.804903in}}%
\pgfpathcurveto{\pgfqpoint{3.532192in}{1.809021in}}{\pgfqpoint{3.534506in}{1.814607in}}{\pgfqpoint{3.534506in}{1.820431in}}%
\pgfpathcurveto{\pgfqpoint{3.534506in}{1.826255in}}{\pgfqpoint{3.532192in}{1.831841in}}{\pgfqpoint{3.528074in}{1.835960in}}%
\pgfpathcurveto{\pgfqpoint{3.523956in}{1.840078in}}{\pgfqpoint{3.518370in}{1.842392in}}{\pgfqpoint{3.512546in}{1.842392in}}%
\pgfpathcurveto{\pgfqpoint{3.506722in}{1.842392in}}{\pgfqpoint{3.501136in}{1.840078in}}{\pgfqpoint{3.497018in}{1.835960in}}%
\pgfpathcurveto{\pgfqpoint{3.492899in}{1.831841in}}{\pgfqpoint{3.490586in}{1.826255in}}{\pgfqpoint{3.490586in}{1.820431in}}%
\pgfpathcurveto{\pgfqpoint{3.490586in}{1.814607in}}{\pgfqpoint{3.492899in}{1.809021in}}{\pgfqpoint{3.497018in}{1.804903in}}%
\pgfpathcurveto{\pgfqpoint{3.501136in}{1.800785in}}{\pgfqpoint{3.506722in}{1.798471in}}{\pgfqpoint{3.512546in}{1.798471in}}%
\pgfpathlineto{\pgfqpoint{3.512546in}{1.798471in}}%
\pgfpathclose%
\pgfusepath{stroke,fill}%
\end{pgfscope}%
\begin{pgfscope}%
\pgfpathrectangle{\pgfqpoint{1.414736in}{0.548769in}}{\pgfqpoint{3.497747in}{7.253186in}}%
\pgfusepath{clip}%
\pgfsetbuttcap%
\pgfsetroundjoin%
\definecolor{currentfill}{rgb}{0.993248,0.906157,0.143936}%
\pgfsetfillcolor{currentfill}%
\pgfsetlinewidth{1.003750pt}%
\definecolor{currentstroke}{rgb}{0.993248,0.906157,0.143936}%
\pgfsetstrokecolor{currentstroke}%
\pgfsetdash{}{0pt}%
\pgfpathmoveto{\pgfqpoint{3.506277in}{2.269457in}}%
\pgfpathcurveto{\pgfqpoint{3.512101in}{2.269457in}}{\pgfqpoint{3.517687in}{2.271771in}}{\pgfqpoint{3.521805in}{2.275889in}}%
\pgfpathcurveto{\pgfqpoint{3.525923in}{2.280007in}}{\pgfqpoint{3.528237in}{2.285593in}}{\pgfqpoint{3.528237in}{2.291417in}}%
\pgfpathcurveto{\pgfqpoint{3.528237in}{2.297241in}}{\pgfqpoint{3.525923in}{2.302828in}}{\pgfqpoint{3.521805in}{2.306946in}}%
\pgfpathcurveto{\pgfqpoint{3.517687in}{2.311064in}}{\pgfqpoint{3.512101in}{2.313378in}}{\pgfqpoint{3.506277in}{2.313378in}}%
\pgfpathcurveto{\pgfqpoint{3.500453in}{2.313378in}}{\pgfqpoint{3.494867in}{2.311064in}}{\pgfqpoint{3.490748in}{2.306946in}}%
\pgfpathcurveto{\pgfqpoint{3.486630in}{2.302828in}}{\pgfqpoint{3.484316in}{2.297241in}}{\pgfqpoint{3.484316in}{2.291417in}}%
\pgfpathcurveto{\pgfqpoint{3.484316in}{2.285593in}}{\pgfqpoint{3.486630in}{2.280007in}}{\pgfqpoint{3.490748in}{2.275889in}}%
\pgfpathcurveto{\pgfqpoint{3.494867in}{2.271771in}}{\pgfqpoint{3.500453in}{2.269457in}}{\pgfqpoint{3.506277in}{2.269457in}}%
\pgfpathlineto{\pgfqpoint{3.506277in}{2.269457in}}%
\pgfpathclose%
\pgfusepath{stroke,fill}%
\end{pgfscope}%
\begin{pgfscope}%
\pgfpathrectangle{\pgfqpoint{1.414736in}{0.548769in}}{\pgfqpoint{3.497747in}{7.253186in}}%
\pgfusepath{clip}%
\pgfsetbuttcap%
\pgfsetroundjoin%
\definecolor{currentfill}{rgb}{0.440137,0.811138,0.340967}%
\pgfsetfillcolor{currentfill}%
\pgfsetlinewidth{1.003750pt}%
\definecolor{currentstroke}{rgb}{0.440137,0.811138,0.340967}%
\pgfsetstrokecolor{currentstroke}%
\pgfsetdash{}{0pt}%
\pgfpathmoveto{\pgfqpoint{3.629560in}{2.740443in}}%
\pgfpathcurveto{\pgfqpoint{3.635384in}{2.740443in}}{\pgfqpoint{3.640971in}{2.742757in}}{\pgfqpoint{3.645089in}{2.746875in}}%
\pgfpathcurveto{\pgfqpoint{3.649207in}{2.750993in}}{\pgfqpoint{3.651521in}{2.756580in}}{\pgfqpoint{3.651521in}{2.762404in}}%
\pgfpathcurveto{\pgfqpoint{3.651521in}{2.768227in}}{\pgfqpoint{3.649207in}{2.773814in}}{\pgfqpoint{3.645089in}{2.777932in}}%
\pgfpathcurveto{\pgfqpoint{3.640971in}{2.782050in}}{\pgfqpoint{3.635384in}{2.784364in}}{\pgfqpoint{3.629560in}{2.784364in}}%
\pgfpathcurveto{\pgfqpoint{3.623737in}{2.784364in}}{\pgfqpoint{3.618150in}{2.782050in}}{\pgfqpoint{3.614032in}{2.777932in}}%
\pgfpathcurveto{\pgfqpoint{3.609914in}{2.773814in}}{\pgfqpoint{3.607600in}{2.768227in}}{\pgfqpoint{3.607600in}{2.762404in}}%
\pgfpathcurveto{\pgfqpoint{3.607600in}{2.756580in}}{\pgfqpoint{3.609914in}{2.750993in}}{\pgfqpoint{3.614032in}{2.746875in}}%
\pgfpathcurveto{\pgfqpoint{3.618150in}{2.742757in}}{\pgfqpoint{3.623737in}{2.740443in}}{\pgfqpoint{3.629560in}{2.740443in}}%
\pgfpathlineto{\pgfqpoint{3.629560in}{2.740443in}}%
\pgfpathclose%
\pgfusepath{stroke,fill}%
\end{pgfscope}%
\begin{pgfscope}%
\pgfpathrectangle{\pgfqpoint{1.414736in}{0.548769in}}{\pgfqpoint{3.497747in}{7.253186in}}%
\pgfusepath{clip}%
\pgfsetbuttcap%
\pgfsetroundjoin%
\definecolor{currentfill}{rgb}{0.855810,0.888601,0.097452}%
\pgfsetfillcolor{currentfill}%
\pgfsetlinewidth{1.003750pt}%
\definecolor{currentstroke}{rgb}{0.855810,0.888601,0.097452}%
\pgfsetstrokecolor{currentstroke}%
\pgfsetdash{}{0pt}%
\pgfpathmoveto{\pgfqpoint{3.483480in}{3.211429in}}%
\pgfpathcurveto{\pgfqpoint{3.489304in}{3.211429in}}{\pgfqpoint{3.494890in}{3.213743in}}{\pgfqpoint{3.499009in}{3.217861in}}%
\pgfpathcurveto{\pgfqpoint{3.503127in}{3.221980in}}{\pgfqpoint{3.505441in}{3.227566in}}{\pgfqpoint{3.505441in}{3.233390in}}%
\pgfpathcurveto{\pgfqpoint{3.505441in}{3.239214in}}{\pgfqpoint{3.503127in}{3.244800in}}{\pgfqpoint{3.499009in}{3.248918in}}%
\pgfpathcurveto{\pgfqpoint{3.494890in}{3.253036in}}{\pgfqpoint{3.489304in}{3.255350in}}{\pgfqpoint{3.483480in}{3.255350in}}%
\pgfpathcurveto{\pgfqpoint{3.477656in}{3.255350in}}{\pgfqpoint{3.472070in}{3.253036in}}{\pgfqpoint{3.467952in}{3.248918in}}%
\pgfpathcurveto{\pgfqpoint{3.463834in}{3.244800in}}{\pgfqpoint{3.461520in}{3.239214in}}{\pgfqpoint{3.461520in}{3.233390in}}%
\pgfpathcurveto{\pgfqpoint{3.461520in}{3.227566in}}{\pgfqpoint{3.463834in}{3.221980in}}{\pgfqpoint{3.467952in}{3.217861in}}%
\pgfpathcurveto{\pgfqpoint{3.472070in}{3.213743in}}{\pgfqpoint{3.477656in}{3.211429in}}{\pgfqpoint{3.483480in}{3.211429in}}%
\pgfpathlineto{\pgfqpoint{3.483480in}{3.211429in}}%
\pgfpathclose%
\pgfusepath{stroke,fill}%
\end{pgfscope}%
\begin{pgfscope}%
\pgfpathrectangle{\pgfqpoint{1.414736in}{0.548769in}}{\pgfqpoint{3.497747in}{7.253186in}}%
\pgfusepath{clip}%
\pgfsetbuttcap%
\pgfsetroundjoin%
\definecolor{currentfill}{rgb}{0.993248,0.906157,0.143936}%
\pgfsetfillcolor{currentfill}%
\pgfsetlinewidth{1.003750pt}%
\definecolor{currentstroke}{rgb}{0.993248,0.906157,0.143936}%
\pgfsetstrokecolor{currentstroke}%
\pgfsetdash{}{0pt}%
\pgfpathmoveto{\pgfqpoint{3.048730in}{3.682415in}}%
\pgfpathcurveto{\pgfqpoint{3.054554in}{3.682415in}}{\pgfqpoint{3.060141in}{3.684729in}}{\pgfqpoint{3.064259in}{3.688847in}}%
\pgfpathcurveto{\pgfqpoint{3.068377in}{3.692966in}}{\pgfqpoint{3.070691in}{3.698552in}}{\pgfqpoint{3.070691in}{3.704376in}}%
\pgfpathcurveto{\pgfqpoint{3.070691in}{3.710200in}}{\pgfqpoint{3.068377in}{3.715786in}}{\pgfqpoint{3.064259in}{3.719904in}}%
\pgfpathcurveto{\pgfqpoint{3.060141in}{3.724022in}}{\pgfqpoint{3.054554in}{3.726336in}}{\pgfqpoint{3.048730in}{3.726336in}}%
\pgfpathcurveto{\pgfqpoint{3.042907in}{3.726336in}}{\pgfqpoint{3.037320in}{3.724022in}}{\pgfqpoint{3.033202in}{3.719904in}}%
\pgfpathcurveto{\pgfqpoint{3.029084in}{3.715786in}}{\pgfqpoint{3.026770in}{3.710200in}}{\pgfqpoint{3.026770in}{3.704376in}}%
\pgfpathcurveto{\pgfqpoint{3.026770in}{3.698552in}}{\pgfqpoint{3.029084in}{3.692966in}}{\pgfqpoint{3.033202in}{3.688847in}}%
\pgfpathcurveto{\pgfqpoint{3.037320in}{3.684729in}}{\pgfqpoint{3.042907in}{3.682415in}}{\pgfqpoint{3.048730in}{3.682415in}}%
\pgfpathlineto{\pgfqpoint{3.048730in}{3.682415in}}%
\pgfpathclose%
\pgfusepath{stroke,fill}%
\end{pgfscope}%
\begin{pgfscope}%
\pgfpathrectangle{\pgfqpoint{1.414736in}{0.548769in}}{\pgfqpoint{3.497747in}{7.253186in}}%
\pgfusepath{clip}%
\pgfsetbuttcap%
\pgfsetroundjoin%
\definecolor{currentfill}{rgb}{0.377779,0.791781,0.377939}%
\pgfsetfillcolor{currentfill}%
\pgfsetlinewidth{1.003750pt}%
\definecolor{currentstroke}{rgb}{0.377779,0.791781,0.377939}%
\pgfsetstrokecolor{currentstroke}%
\pgfsetdash{}{0pt}%
\pgfpathmoveto{\pgfqpoint{3.533495in}{4.153402in}}%
\pgfpathcurveto{\pgfqpoint{3.539319in}{4.153402in}}{\pgfqpoint{3.544905in}{4.155715in}}{\pgfqpoint{3.549023in}{4.159834in}}%
\pgfpathcurveto{\pgfqpoint{3.553141in}{4.163952in}}{\pgfqpoint{3.555455in}{4.169538in}}{\pgfqpoint{3.555455in}{4.175362in}}%
\pgfpathcurveto{\pgfqpoint{3.555455in}{4.181186in}}{\pgfqpoint{3.553141in}{4.186772in}}{\pgfqpoint{3.549023in}{4.190890in}}%
\pgfpathcurveto{\pgfqpoint{3.544905in}{4.195008in}}{\pgfqpoint{3.539319in}{4.197322in}}{\pgfqpoint{3.533495in}{4.197322in}}%
\pgfpathcurveto{\pgfqpoint{3.527671in}{4.197322in}}{\pgfqpoint{3.522084in}{4.195008in}}{\pgfqpoint{3.517966in}{4.190890in}}%
\pgfpathcurveto{\pgfqpoint{3.513848in}{4.186772in}}{\pgfqpoint{3.511534in}{4.181186in}}{\pgfqpoint{3.511534in}{4.175362in}}%
\pgfpathcurveto{\pgfqpoint{3.511534in}{4.169538in}}{\pgfqpoint{3.513848in}{4.163952in}}{\pgfqpoint{3.517966in}{4.159834in}}%
\pgfpathcurveto{\pgfqpoint{3.522084in}{4.155715in}}{\pgfqpoint{3.527671in}{4.153402in}}{\pgfqpoint{3.533495in}{4.153402in}}%
\pgfpathlineto{\pgfqpoint{3.533495in}{4.153402in}}%
\pgfpathclose%
\pgfusepath{stroke,fill}%
\end{pgfscope}%
\begin{pgfscope}%
\pgfpathrectangle{\pgfqpoint{1.414736in}{0.548769in}}{\pgfqpoint{3.497747in}{7.253186in}}%
\pgfusepath{clip}%
\pgfsetbuttcap%
\pgfsetroundjoin%
\definecolor{currentfill}{rgb}{0.772852,0.877868,0.131109}%
\pgfsetfillcolor{currentfill}%
\pgfsetlinewidth{1.003750pt}%
\definecolor{currentstroke}{rgb}{0.772852,0.877868,0.131109}%
\pgfsetstrokecolor{currentstroke}%
\pgfsetdash{}{0pt}%
\pgfpathmoveto{\pgfqpoint{3.515128in}{4.624388in}}%
\pgfpathcurveto{\pgfqpoint{3.520952in}{4.624388in}}{\pgfqpoint{3.526538in}{4.626702in}}{\pgfqpoint{3.530656in}{4.630820in}}%
\pgfpathcurveto{\pgfqpoint{3.534774in}{4.634938in}}{\pgfqpoint{3.537088in}{4.640524in}}{\pgfqpoint{3.537088in}{4.646348in}}%
\pgfpathcurveto{\pgfqpoint{3.537088in}{4.652172in}}{\pgfqpoint{3.534774in}{4.657758in}}{\pgfqpoint{3.530656in}{4.661876in}}%
\pgfpathcurveto{\pgfqpoint{3.526538in}{4.665994in}}{\pgfqpoint{3.520952in}{4.668308in}}{\pgfqpoint{3.515128in}{4.668308in}}%
\pgfpathcurveto{\pgfqpoint{3.509304in}{4.668308in}}{\pgfqpoint{3.503718in}{4.665994in}}{\pgfqpoint{3.499600in}{4.661876in}}%
\pgfpathcurveto{\pgfqpoint{3.495482in}{4.657758in}}{\pgfqpoint{3.493168in}{4.652172in}}{\pgfqpoint{3.493168in}{4.646348in}}%
\pgfpathcurveto{\pgfqpoint{3.493168in}{4.640524in}}{\pgfqpoint{3.495482in}{4.634938in}}{\pgfqpoint{3.499600in}{4.630820in}}%
\pgfpathcurveto{\pgfqpoint{3.503718in}{4.626702in}}{\pgfqpoint{3.509304in}{4.624388in}}{\pgfqpoint{3.515128in}{4.624388in}}%
\pgfpathlineto{\pgfqpoint{3.515128in}{4.624388in}}%
\pgfpathclose%
\pgfusepath{stroke,fill}%
\end{pgfscope}%
\begin{pgfscope}%
\pgfpathrectangle{\pgfqpoint{1.414736in}{0.548769in}}{\pgfqpoint{3.497747in}{7.253186in}}%
\pgfusepath{clip}%
\pgfsetbuttcap%
\pgfsetroundjoin%
\definecolor{currentfill}{rgb}{0.282623,0.140926,0.457517}%
\pgfsetfillcolor{currentfill}%
\pgfsetlinewidth{1.003750pt}%
\definecolor{currentstroke}{rgb}{0.282623,0.140926,0.457517}%
\pgfsetstrokecolor{currentstroke}%
\pgfsetdash{}{0pt}%
\pgfpathmoveto{\pgfqpoint{3.723964in}{5.095374in}}%
\pgfpathcurveto{\pgfqpoint{3.729788in}{5.095374in}}{\pgfqpoint{3.735374in}{5.097688in}}{\pgfqpoint{3.739492in}{5.101806in}}%
\pgfpathcurveto{\pgfqpoint{3.743611in}{5.105924in}}{\pgfqpoint{3.745924in}{5.111510in}}{\pgfqpoint{3.745924in}{5.117334in}}%
\pgfpathcurveto{\pgfqpoint{3.745924in}{5.123158in}}{\pgfqpoint{3.743611in}{5.128744in}}{\pgfqpoint{3.739492in}{5.132862in}}%
\pgfpathcurveto{\pgfqpoint{3.735374in}{5.136980in}}{\pgfqpoint{3.729788in}{5.139294in}}{\pgfqpoint{3.723964in}{5.139294in}}%
\pgfpathcurveto{\pgfqpoint{3.718140in}{5.139294in}}{\pgfqpoint{3.712554in}{5.136980in}}{\pgfqpoint{3.708436in}{5.132862in}}%
\pgfpathcurveto{\pgfqpoint{3.704318in}{5.128744in}}{\pgfqpoint{3.702004in}{5.123158in}}{\pgfqpoint{3.702004in}{5.117334in}}%
\pgfpathcurveto{\pgfqpoint{3.702004in}{5.111510in}}{\pgfqpoint{3.704318in}{5.105924in}}{\pgfqpoint{3.708436in}{5.101806in}}%
\pgfpathcurveto{\pgfqpoint{3.712554in}{5.097688in}}{\pgfqpoint{3.718140in}{5.095374in}}{\pgfqpoint{3.723964in}{5.095374in}}%
\pgfpathlineto{\pgfqpoint{3.723964in}{5.095374in}}%
\pgfpathclose%
\pgfusepath{stroke,fill}%
\end{pgfscope}%
\begin{pgfscope}%
\pgfpathrectangle{\pgfqpoint{1.414736in}{0.548769in}}{\pgfqpoint{3.497747in}{7.253186in}}%
\pgfusepath{clip}%
\pgfsetbuttcap%
\pgfsetroundjoin%
\definecolor{currentfill}{rgb}{0.218130,0.347432,0.550038}%
\pgfsetfillcolor{currentfill}%
\pgfsetlinewidth{1.003750pt}%
\definecolor{currentstroke}{rgb}{0.218130,0.347432,0.550038}%
\pgfsetstrokecolor{currentstroke}%
\pgfsetdash{}{0pt}%
\pgfpathmoveto{\pgfqpoint{3.784661in}{5.566360in}}%
\pgfpathcurveto{\pgfqpoint{3.790485in}{5.566360in}}{\pgfqpoint{3.796071in}{5.568674in}}{\pgfqpoint{3.800189in}{5.572792in}}%
\pgfpathcurveto{\pgfqpoint{3.804307in}{5.576910in}}{\pgfqpoint{3.806621in}{5.582496in}}{\pgfqpoint{3.806621in}{5.588320in}}%
\pgfpathcurveto{\pgfqpoint{3.806621in}{5.594144in}}{\pgfqpoint{3.804307in}{5.599730in}}{\pgfqpoint{3.800189in}{5.603848in}}%
\pgfpathcurveto{\pgfqpoint{3.796071in}{5.607967in}}{\pgfqpoint{3.790485in}{5.610280in}}{\pgfqpoint{3.784661in}{5.610280in}}%
\pgfpathcurveto{\pgfqpoint{3.778837in}{5.610280in}}{\pgfqpoint{3.773251in}{5.607967in}}{\pgfqpoint{3.769133in}{5.603848in}}%
\pgfpathcurveto{\pgfqpoint{3.765015in}{5.599730in}}{\pgfqpoint{3.762701in}{5.594144in}}{\pgfqpoint{3.762701in}{5.588320in}}%
\pgfpathcurveto{\pgfqpoint{3.762701in}{5.582496in}}{\pgfqpoint{3.765015in}{5.576910in}}{\pgfqpoint{3.769133in}{5.572792in}}%
\pgfpathcurveto{\pgfqpoint{3.773251in}{5.568674in}}{\pgfqpoint{3.778837in}{5.566360in}}{\pgfqpoint{3.784661in}{5.566360in}}%
\pgfpathlineto{\pgfqpoint{3.784661in}{5.566360in}}%
\pgfpathclose%
\pgfusepath{stroke,fill}%
\end{pgfscope}%
\begin{pgfscope}%
\pgfpathrectangle{\pgfqpoint{1.414736in}{0.548769in}}{\pgfqpoint{3.497747in}{7.253186in}}%
\pgfusepath{clip}%
\pgfsetbuttcap%
\pgfsetroundjoin%
\definecolor{currentfill}{rgb}{0.124780,0.640461,0.527068}%
\pgfsetfillcolor{currentfill}%
\pgfsetlinewidth{1.003750pt}%
\definecolor{currentstroke}{rgb}{0.124780,0.640461,0.527068}%
\pgfsetstrokecolor{currentstroke}%
\pgfsetdash{}{0pt}%
\pgfpathmoveto{\pgfqpoint{3.600827in}{6.037346in}}%
\pgfpathcurveto{\pgfqpoint{3.606651in}{6.037346in}}{\pgfqpoint{3.612237in}{6.039660in}}{\pgfqpoint{3.616355in}{6.043778in}}%
\pgfpathcurveto{\pgfqpoint{3.620473in}{6.047896in}}{\pgfqpoint{3.622787in}{6.053482in}}{\pgfqpoint{3.622787in}{6.059306in}}%
\pgfpathcurveto{\pgfqpoint{3.622787in}{6.065130in}}{\pgfqpoint{3.620473in}{6.070716in}}{\pgfqpoint{3.616355in}{6.074835in}}%
\pgfpathcurveto{\pgfqpoint{3.612237in}{6.078953in}}{\pgfqpoint{3.606651in}{6.081267in}}{\pgfqpoint{3.600827in}{6.081267in}}%
\pgfpathcurveto{\pgfqpoint{3.595003in}{6.081267in}}{\pgfqpoint{3.589417in}{6.078953in}}{\pgfqpoint{3.585299in}{6.074835in}}%
\pgfpathcurveto{\pgfqpoint{3.581181in}{6.070716in}}{\pgfqpoint{3.578867in}{6.065130in}}{\pgfqpoint{3.578867in}{6.059306in}}%
\pgfpathcurveto{\pgfqpoint{3.578867in}{6.053482in}}{\pgfqpoint{3.581181in}{6.047896in}}{\pgfqpoint{3.585299in}{6.043778in}}%
\pgfpathcurveto{\pgfqpoint{3.589417in}{6.039660in}}{\pgfqpoint{3.595003in}{6.037346in}}{\pgfqpoint{3.600827in}{6.037346in}}%
\pgfpathlineto{\pgfqpoint{3.600827in}{6.037346in}}%
\pgfpathclose%
\pgfusepath{stroke,fill}%
\end{pgfscope}%
\begin{pgfscope}%
\pgfpathrectangle{\pgfqpoint{1.414736in}{0.548769in}}{\pgfqpoint{3.497747in}{7.253186in}}%
\pgfusepath{clip}%
\pgfsetbuttcap%
\pgfsetroundjoin%
\definecolor{currentfill}{rgb}{0.793760,0.880678,0.120005}%
\pgfsetfillcolor{currentfill}%
\pgfsetlinewidth{1.003750pt}%
\definecolor{currentstroke}{rgb}{0.793760,0.880678,0.120005}%
\pgfsetstrokecolor{currentstroke}%
\pgfsetdash{}{0pt}%
\pgfpathmoveto{\pgfqpoint{3.757116in}{6.508332in}}%
\pgfpathcurveto{\pgfqpoint{3.762940in}{6.508332in}}{\pgfqpoint{3.768526in}{6.510646in}}{\pgfqpoint{3.772644in}{6.514764in}}%
\pgfpathcurveto{\pgfqpoint{3.776762in}{6.518882in}}{\pgfqpoint{3.779076in}{6.524468in}}{\pgfqpoint{3.779076in}{6.530292in}}%
\pgfpathcurveto{\pgfqpoint{3.779076in}{6.536116in}}{\pgfqpoint{3.776762in}{6.541702in}}{\pgfqpoint{3.772644in}{6.545821in}}%
\pgfpathcurveto{\pgfqpoint{3.768526in}{6.549939in}}{\pgfqpoint{3.762940in}{6.552253in}}{\pgfqpoint{3.757116in}{6.552253in}}%
\pgfpathcurveto{\pgfqpoint{3.751292in}{6.552253in}}{\pgfqpoint{3.745706in}{6.549939in}}{\pgfqpoint{3.741588in}{6.545821in}}%
\pgfpathcurveto{\pgfqpoint{3.737469in}{6.541702in}}{\pgfqpoint{3.735156in}{6.536116in}}{\pgfqpoint{3.735156in}{6.530292in}}%
\pgfpathcurveto{\pgfqpoint{3.735156in}{6.524468in}}{\pgfqpoint{3.737469in}{6.518882in}}{\pgfqpoint{3.741588in}{6.514764in}}%
\pgfpathcurveto{\pgfqpoint{3.745706in}{6.510646in}}{\pgfqpoint{3.751292in}{6.508332in}}{\pgfqpoint{3.757116in}{6.508332in}}%
\pgfpathlineto{\pgfqpoint{3.757116in}{6.508332in}}%
\pgfpathclose%
\pgfusepath{stroke,fill}%
\end{pgfscope}%
\begin{pgfscope}%
\pgfpathrectangle{\pgfqpoint{1.414736in}{0.548769in}}{\pgfqpoint{3.497747in}{7.253186in}}%
\pgfusepath{clip}%
\pgfsetbuttcap%
\pgfsetroundjoin%
\definecolor{currentfill}{rgb}{0.344074,0.780029,0.397381}%
\pgfsetfillcolor{currentfill}%
\pgfsetlinewidth{1.003750pt}%
\definecolor{currentstroke}{rgb}{0.344074,0.780029,0.397381}%
\pgfsetstrokecolor{currentstroke}%
\pgfsetdash{}{0pt}%
\pgfpathmoveto{\pgfqpoint{3.976245in}{6.979318in}}%
\pgfpathcurveto{\pgfqpoint{3.982069in}{6.979318in}}{\pgfqpoint{3.987655in}{6.981632in}}{\pgfqpoint{3.991773in}{6.985750in}}%
\pgfpathcurveto{\pgfqpoint{3.995891in}{6.989868in}}{\pgfqpoint{3.998205in}{6.995455in}}{\pgfqpoint{3.998205in}{7.001278in}}%
\pgfpathcurveto{\pgfqpoint{3.998205in}{7.007102in}}{\pgfqpoint{3.995891in}{7.012689in}}{\pgfqpoint{3.991773in}{7.016807in}}%
\pgfpathcurveto{\pgfqpoint{3.987655in}{7.020925in}}{\pgfqpoint{3.982069in}{7.023239in}}{\pgfqpoint{3.976245in}{7.023239in}}%
\pgfpathcurveto{\pgfqpoint{3.970421in}{7.023239in}}{\pgfqpoint{3.964835in}{7.020925in}}{\pgfqpoint{3.960717in}{7.016807in}}%
\pgfpathcurveto{\pgfqpoint{3.956599in}{7.012689in}}{\pgfqpoint{3.954285in}{7.007102in}}{\pgfqpoint{3.954285in}{7.001278in}}%
\pgfpathcurveto{\pgfqpoint{3.954285in}{6.995455in}}{\pgfqpoint{3.956599in}{6.989868in}}{\pgfqpoint{3.960717in}{6.985750in}}%
\pgfpathcurveto{\pgfqpoint{3.964835in}{6.981632in}}{\pgfqpoint{3.970421in}{6.979318in}}{\pgfqpoint{3.976245in}{6.979318in}}%
\pgfpathlineto{\pgfqpoint{3.976245in}{6.979318in}}%
\pgfpathclose%
\pgfusepath{stroke,fill}%
\end{pgfscope}%
\begin{pgfscope}%
\pgfpathrectangle{\pgfqpoint{1.414736in}{0.548769in}}{\pgfqpoint{3.497747in}{7.253186in}}%
\pgfusepath{clip}%
\pgfsetbuttcap%
\pgfsetroundjoin%
\definecolor{currentfill}{rgb}{0.335885,0.777018,0.402049}%
\pgfsetfillcolor{currentfill}%
\pgfsetlinewidth{1.003750pt}%
\definecolor{currentstroke}{rgb}{0.335885,0.777018,0.402049}%
\pgfsetstrokecolor{currentstroke}%
\pgfsetdash{}{0pt}%
\pgfpathmoveto{\pgfqpoint{4.183319in}{7.450304in}}%
\pgfpathcurveto{\pgfqpoint{4.189143in}{7.450304in}}{\pgfqpoint{4.194729in}{7.452618in}}{\pgfqpoint{4.198847in}{7.456736in}}%
\pgfpathcurveto{\pgfqpoint{4.202966in}{7.460854in}}{\pgfqpoint{4.205279in}{7.466441in}}{\pgfqpoint{4.205279in}{7.472265in}}%
\pgfpathcurveto{\pgfqpoint{4.205279in}{7.478089in}}{\pgfqpoint{4.202966in}{7.483675in}}{\pgfqpoint{4.198847in}{7.487793in}}%
\pgfpathcurveto{\pgfqpoint{4.194729in}{7.491911in}}{\pgfqpoint{4.189143in}{7.494225in}}{\pgfqpoint{4.183319in}{7.494225in}}%
\pgfpathcurveto{\pgfqpoint{4.177495in}{7.494225in}}{\pgfqpoint{4.171909in}{7.491911in}}{\pgfqpoint{4.167791in}{7.487793in}}%
\pgfpathcurveto{\pgfqpoint{4.163673in}{7.483675in}}{\pgfqpoint{4.161359in}{7.478089in}}{\pgfqpoint{4.161359in}{7.472265in}}%
\pgfpathcurveto{\pgfqpoint{4.161359in}{7.466441in}}{\pgfqpoint{4.163673in}{7.460854in}}{\pgfqpoint{4.167791in}{7.456736in}}%
\pgfpathcurveto{\pgfqpoint{4.171909in}{7.452618in}}{\pgfqpoint{4.177495in}{7.450304in}}{\pgfqpoint{4.183319in}{7.450304in}}%
\pgfpathlineto{\pgfqpoint{4.183319in}{7.450304in}}%
\pgfpathclose%
\pgfusepath{stroke,fill}%
\end{pgfscope}%
\begin{pgfscope}%
\pgfsetbuttcap%
\pgfsetroundjoin%
\definecolor{currentfill}{rgb}{0.000000,0.000000,0.000000}%
\pgfsetfillcolor{currentfill}%
\pgfsetlinewidth{0.803000pt}%
\definecolor{currentstroke}{rgb}{0.000000,0.000000,0.000000}%
\pgfsetstrokecolor{currentstroke}%
\pgfsetdash{}{0pt}%
\pgfsys@defobject{currentmarker}{\pgfqpoint{0.000000in}{-0.048611in}}{\pgfqpoint{0.000000in}{0.000000in}}{%
\pgfpathmoveto{\pgfqpoint{0.000000in}{0.000000in}}%
\pgfpathlineto{\pgfqpoint{0.000000in}{-0.048611in}}%
\pgfusepath{stroke,fill}%
}%
\begin{pgfscope}%
\pgfsys@transformshift{1.873364in}{0.548769in}%
\pgfsys@useobject{currentmarker}{}%
\end{pgfscope}%
\end{pgfscope}%
\begin{pgfscope}%
\definecolor{textcolor}{rgb}{0.000000,0.000000,0.000000}%
\pgfsetstrokecolor{textcolor}%
\pgfsetfillcolor{textcolor}%
\pgftext[x=1.873364in,y=0.451547in,,top]{\color{textcolor}{\rmfamily\fontsize{10.000000}{12.000000}\selectfont\catcode`\^=\active\def^{\ifmmode\sp\else\^{}\fi}\catcode`\%=\active\def%{\%}$\mathdefault{\ensuremath{-}0.15}$}}%
\end{pgfscope}%
\begin{pgfscope}%
\pgfsetbuttcap%
\pgfsetroundjoin%
\definecolor{currentfill}{rgb}{0.000000,0.000000,0.000000}%
\pgfsetfillcolor{currentfill}%
\pgfsetlinewidth{0.803000pt}%
\definecolor{currentstroke}{rgb}{0.000000,0.000000,0.000000}%
\pgfsetstrokecolor{currentstroke}%
\pgfsetdash{}{0pt}%
\pgfsys@defobject{currentmarker}{\pgfqpoint{0.000000in}{-0.048611in}}{\pgfqpoint{0.000000in}{0.000000in}}{%
\pgfpathmoveto{\pgfqpoint{0.000000in}{0.000000in}}%
\pgfpathlineto{\pgfqpoint{0.000000in}{-0.048611in}}%
\pgfusepath{stroke,fill}%
}%
\begin{pgfscope}%
\pgfsys@transformshift{2.440236in}{0.548769in}%
\pgfsys@useobject{currentmarker}{}%
\end{pgfscope}%
\end{pgfscope}%
\begin{pgfscope}%
\definecolor{textcolor}{rgb}{0.000000,0.000000,0.000000}%
\pgfsetstrokecolor{textcolor}%
\pgfsetfillcolor{textcolor}%
\pgftext[x=2.440236in,y=0.451547in,,top]{\color{textcolor}{\rmfamily\fontsize{10.000000}{12.000000}\selectfont\catcode`\^=\active\def^{\ifmmode\sp\else\^{}\fi}\catcode`\%=\active\def%{\%}$\mathdefault{\ensuremath{-}0.10}$}}%
\end{pgfscope}%
\begin{pgfscope}%
\pgfsetbuttcap%
\pgfsetroundjoin%
\definecolor{currentfill}{rgb}{0.000000,0.000000,0.000000}%
\pgfsetfillcolor{currentfill}%
\pgfsetlinewidth{0.803000pt}%
\definecolor{currentstroke}{rgb}{0.000000,0.000000,0.000000}%
\pgfsetstrokecolor{currentstroke}%
\pgfsetdash{}{0pt}%
\pgfsys@defobject{currentmarker}{\pgfqpoint{0.000000in}{-0.048611in}}{\pgfqpoint{0.000000in}{0.000000in}}{%
\pgfpathmoveto{\pgfqpoint{0.000000in}{0.000000in}}%
\pgfpathlineto{\pgfqpoint{0.000000in}{-0.048611in}}%
\pgfusepath{stroke,fill}%
}%
\begin{pgfscope}%
\pgfsys@transformshift{3.007107in}{0.548769in}%
\pgfsys@useobject{currentmarker}{}%
\end{pgfscope}%
\end{pgfscope}%
\begin{pgfscope}%
\definecolor{textcolor}{rgb}{0.000000,0.000000,0.000000}%
\pgfsetstrokecolor{textcolor}%
\pgfsetfillcolor{textcolor}%
\pgftext[x=3.007107in,y=0.451547in,,top]{\color{textcolor}{\rmfamily\fontsize{10.000000}{12.000000}\selectfont\catcode`\^=\active\def^{\ifmmode\sp\else\^{}\fi}\catcode`\%=\active\def%{\%}$\mathdefault{\ensuremath{-}0.05}$}}%
\end{pgfscope}%
\begin{pgfscope}%
\pgfsetbuttcap%
\pgfsetroundjoin%
\definecolor{currentfill}{rgb}{0.000000,0.000000,0.000000}%
\pgfsetfillcolor{currentfill}%
\pgfsetlinewidth{0.803000pt}%
\definecolor{currentstroke}{rgb}{0.000000,0.000000,0.000000}%
\pgfsetstrokecolor{currentstroke}%
\pgfsetdash{}{0pt}%
\pgfsys@defobject{currentmarker}{\pgfqpoint{0.000000in}{-0.048611in}}{\pgfqpoint{0.000000in}{0.000000in}}{%
\pgfpathmoveto{\pgfqpoint{0.000000in}{0.000000in}}%
\pgfpathlineto{\pgfqpoint{0.000000in}{-0.048611in}}%
\pgfusepath{stroke,fill}%
}%
\begin{pgfscope}%
\pgfsys@transformshift{3.573979in}{0.548769in}%
\pgfsys@useobject{currentmarker}{}%
\end{pgfscope}%
\end{pgfscope}%
\begin{pgfscope}%
\definecolor{textcolor}{rgb}{0.000000,0.000000,0.000000}%
\pgfsetstrokecolor{textcolor}%
\pgfsetfillcolor{textcolor}%
\pgftext[x=3.573979in,y=0.451547in,,top]{\color{textcolor}{\rmfamily\fontsize{10.000000}{12.000000}\selectfont\catcode`\^=\active\def^{\ifmmode\sp\else\^{}\fi}\catcode`\%=\active\def%{\%}$\mathdefault{0.00}$}}%
\end{pgfscope}%
\begin{pgfscope}%
\pgfsetbuttcap%
\pgfsetroundjoin%
\definecolor{currentfill}{rgb}{0.000000,0.000000,0.000000}%
\pgfsetfillcolor{currentfill}%
\pgfsetlinewidth{0.803000pt}%
\definecolor{currentstroke}{rgb}{0.000000,0.000000,0.000000}%
\pgfsetstrokecolor{currentstroke}%
\pgfsetdash{}{0pt}%
\pgfsys@defobject{currentmarker}{\pgfqpoint{0.000000in}{-0.048611in}}{\pgfqpoint{0.000000in}{0.000000in}}{%
\pgfpathmoveto{\pgfqpoint{0.000000in}{0.000000in}}%
\pgfpathlineto{\pgfqpoint{0.000000in}{-0.048611in}}%
\pgfusepath{stroke,fill}%
}%
\begin{pgfscope}%
\pgfsys@transformshift{4.140850in}{0.548769in}%
\pgfsys@useobject{currentmarker}{}%
\end{pgfscope}%
\end{pgfscope}%
\begin{pgfscope}%
\definecolor{textcolor}{rgb}{0.000000,0.000000,0.000000}%
\pgfsetstrokecolor{textcolor}%
\pgfsetfillcolor{textcolor}%
\pgftext[x=4.140850in,y=0.451547in,,top]{\color{textcolor}{\rmfamily\fontsize{10.000000}{12.000000}\selectfont\catcode`\^=\active\def^{\ifmmode\sp\else\^{}\fi}\catcode`\%=\active\def%{\%}$\mathdefault{0.05}$}}%
\end{pgfscope}%
\begin{pgfscope}%
\pgfsetbuttcap%
\pgfsetroundjoin%
\definecolor{currentfill}{rgb}{0.000000,0.000000,0.000000}%
\pgfsetfillcolor{currentfill}%
\pgfsetlinewidth{0.803000pt}%
\definecolor{currentstroke}{rgb}{0.000000,0.000000,0.000000}%
\pgfsetstrokecolor{currentstroke}%
\pgfsetdash{}{0pt}%
\pgfsys@defobject{currentmarker}{\pgfqpoint{0.000000in}{-0.048611in}}{\pgfqpoint{0.000000in}{0.000000in}}{%
\pgfpathmoveto{\pgfqpoint{0.000000in}{0.000000in}}%
\pgfpathlineto{\pgfqpoint{0.000000in}{-0.048611in}}%
\pgfusepath{stroke,fill}%
}%
\begin{pgfscope}%
\pgfsys@transformshift{4.707722in}{0.548769in}%
\pgfsys@useobject{currentmarker}{}%
\end{pgfscope}%
\end{pgfscope}%
\begin{pgfscope}%
\definecolor{textcolor}{rgb}{0.000000,0.000000,0.000000}%
\pgfsetstrokecolor{textcolor}%
\pgfsetfillcolor{textcolor}%
\pgftext[x=4.707722in,y=0.451547in,,top]{\color{textcolor}{\rmfamily\fontsize{10.000000}{12.000000}\selectfont\catcode`\^=\active\def^{\ifmmode\sp\else\^{}\fi}\catcode`\%=\active\def%{\%}$\mathdefault{0.10}$}}%
\end{pgfscope}%
\begin{pgfscope}%
\definecolor{textcolor}{rgb}{0.000000,0.000000,0.000000}%
\pgfsetstrokecolor{textcolor}%
\pgfsetfillcolor{textcolor}%
\pgftext[x=3.163610in,y=0.272534in,,top]{\color{textcolor}{\rmfamily\fontsize{10.000000}{12.000000}\selectfont\catcode`\^=\active\def^{\ifmmode\sp\else\^{}\fi}\catcode`\%=\active\def%{\%}Shapley Value}}%
\end{pgfscope}%
\begin{pgfscope}%
\pgfsetbuttcap%
\pgfsetroundjoin%
\definecolor{currentfill}{rgb}{0.000000,0.000000,0.000000}%
\pgfsetfillcolor{currentfill}%
\pgfsetlinewidth{0.803000pt}%
\definecolor{currentstroke}{rgb}{0.000000,0.000000,0.000000}%
\pgfsetstrokecolor{currentstroke}%
\pgfsetdash{}{0pt}%
\pgfsys@defobject{currentmarker}{\pgfqpoint{-0.048611in}{0.000000in}}{\pgfqpoint{-0.000000in}{0.000000in}}{%
\pgfpathmoveto{\pgfqpoint{-0.000000in}{0.000000in}}%
\pgfpathlineto{\pgfqpoint{-0.048611in}{0.000000in}}%
\pgfusepath{stroke,fill}%
}%
\begin{pgfscope}%
\pgfsys@transformshift{1.414736in}{0.878459in}%
\pgfsys@useobject{currentmarker}{}%
\end{pgfscope}%
\end{pgfscope}%
\begin{pgfscope}%
\definecolor{textcolor}{rgb}{0.000000,0.000000,0.000000}%
\pgfsetstrokecolor{textcolor}%
\pgfsetfillcolor{textcolor}%
\pgftext[x=0.629627in, y=0.830234in, left, base]{\color{textcolor}{\rmfamily\fontsize{10.000000}{12.000000}\selectfont\catcode`\^=\active\def^{\ifmmode\sp\else\^{}\fi}\catcode`\%=\active\def%{\%}vel x sword}}%
\end{pgfscope}%
\begin{pgfscope}%
\pgfsetbuttcap%
\pgfsetroundjoin%
\definecolor{currentfill}{rgb}{0.000000,0.000000,0.000000}%
\pgfsetfillcolor{currentfill}%
\pgfsetlinewidth{0.803000pt}%
\definecolor{currentstroke}{rgb}{0.000000,0.000000,0.000000}%
\pgfsetstrokecolor{currentstroke}%
\pgfsetdash{}{0pt}%
\pgfsys@defobject{currentmarker}{\pgfqpoint{-0.048611in}{0.000000in}}{\pgfqpoint{-0.000000in}{0.000000in}}{%
\pgfpathmoveto{\pgfqpoint{-0.000000in}{0.000000in}}%
\pgfpathlineto{\pgfqpoint{-0.048611in}{0.000000in}}%
\pgfusepath{stroke,fill}%
}%
\begin{pgfscope}%
\pgfsys@transformshift{1.414736in}{1.349445in}%
\pgfsys@useobject{currentmarker}{}%
\end{pgfscope}%
\end{pgfscope}%
\begin{pgfscope}%
\definecolor{textcolor}{rgb}{0.000000,0.000000,0.000000}%
\pgfsetstrokecolor{textcolor}%
\pgfsetfillcolor{textcolor}%
\pgftext[x=0.393130in, y=1.301220in, left, base]{\color{textcolor}{\rmfamily\fontsize{10.000000}{12.000000}\selectfont\catcode`\^=\active\def^{\ifmmode\sp\else\^{}\fi}\catcode`\%=\active\def%{\%}rel pos y sword}}%
\end{pgfscope}%
\begin{pgfscope}%
\pgfsetbuttcap%
\pgfsetroundjoin%
\definecolor{currentfill}{rgb}{0.000000,0.000000,0.000000}%
\pgfsetfillcolor{currentfill}%
\pgfsetlinewidth{0.803000pt}%
\definecolor{currentstroke}{rgb}{0.000000,0.000000,0.000000}%
\pgfsetstrokecolor{currentstroke}%
\pgfsetdash{}{0pt}%
\pgfsys@defobject{currentmarker}{\pgfqpoint{-0.048611in}{0.000000in}}{\pgfqpoint{-0.000000in}{0.000000in}}{%
\pgfpathmoveto{\pgfqpoint{-0.000000in}{0.000000in}}%
\pgfpathlineto{\pgfqpoint{-0.048611in}{0.000000in}}%
\pgfusepath{stroke,fill}%
}%
\begin{pgfscope}%
\pgfsys@transformshift{1.414736in}{1.820431in}%
\pgfsys@useobject{currentmarker}{}%
\end{pgfscope}%
\end{pgfscope}%
\begin{pgfscope}%
\definecolor{textcolor}{rgb}{0.000000,0.000000,0.000000}%
\pgfsetstrokecolor{textcolor}%
\pgfsetfillcolor{textcolor}%
\pgftext[x=0.784334in, y=1.772206in, left, base]{\color{textcolor}{\rmfamily\fontsize{10.000000}{12.000000}\selectfont\catcode`\^=\active\def^{\ifmmode\sp\else\^{}\fi}\catcode`\%=\active\def%{\%}vel x self}}%
\end{pgfscope}%
\begin{pgfscope}%
\pgfsetbuttcap%
\pgfsetroundjoin%
\definecolor{currentfill}{rgb}{0.000000,0.000000,0.000000}%
\pgfsetfillcolor{currentfill}%
\pgfsetlinewidth{0.803000pt}%
\definecolor{currentstroke}{rgb}{0.000000,0.000000,0.000000}%
\pgfsetstrokecolor{currentstroke}%
\pgfsetdash{}{0pt}%
\pgfsys@defobject{currentmarker}{\pgfqpoint{-0.048611in}{0.000000in}}{\pgfqpoint{-0.000000in}{0.000000in}}{%
\pgfpathmoveto{\pgfqpoint{-0.000000in}{0.000000in}}%
\pgfpathlineto{\pgfqpoint{-0.048611in}{0.000000in}}%
\pgfusepath{stroke,fill}%
}%
\begin{pgfscope}%
\pgfsys@transformshift{1.414736in}{2.291417in}%
\pgfsys@useobject{currentmarker}{}%
\end{pgfscope}%
\end{pgfscope}%
\begin{pgfscope}%
\definecolor{textcolor}{rgb}{0.000000,0.000000,0.000000}%
\pgfsetstrokecolor{textcolor}%
\pgfsetfillcolor{textcolor}%
\pgftext[x=0.702544in, y=2.243192in, left, base]{\color{textcolor}{\rmfamily\fontsize{10.000000}{12.000000}\selectfont\catcode`\^=\active\def^{\ifmmode\sp\else\^{}\fi}\catcode`\%=\active\def%{\%}dist arrow}}%
\end{pgfscope}%
\begin{pgfscope}%
\pgfsetbuttcap%
\pgfsetroundjoin%
\definecolor{currentfill}{rgb}{0.000000,0.000000,0.000000}%
\pgfsetfillcolor{currentfill}%
\pgfsetlinewidth{0.803000pt}%
\definecolor{currentstroke}{rgb}{0.000000,0.000000,0.000000}%
\pgfsetstrokecolor{currentstroke}%
\pgfsetdash{}{0pt}%
\pgfsys@defobject{currentmarker}{\pgfqpoint{-0.048611in}{0.000000in}}{\pgfqpoint{-0.000000in}{0.000000in}}{%
\pgfpathmoveto{\pgfqpoint{-0.000000in}{0.000000in}}%
\pgfpathlineto{\pgfqpoint{-0.048611in}{0.000000in}}%
\pgfusepath{stroke,fill}%
}%
\begin{pgfscope}%
\pgfsys@transformshift{1.414736in}{2.762404in}%
\pgfsys@useobject{currentmarker}{}%
\end{pgfscope}%
\end{pgfscope}%
\begin{pgfscope}%
\definecolor{textcolor}{rgb}{0.000000,0.000000,0.000000}%
\pgfsetstrokecolor{textcolor}%
\pgfsetfillcolor{textcolor}%
\pgftext[x=0.359564in, y=2.714178in, left, base]{\color{textcolor}{\rmfamily\fontsize{10.000000}{12.000000}\selectfont\catcode`\^=\active\def^{\ifmmode\sp\else\^{}\fi}\catcode`\%=\active\def%{\%}rel pos y knight}}%
\end{pgfscope}%
\begin{pgfscope}%
\pgfsetbuttcap%
\pgfsetroundjoin%
\definecolor{currentfill}{rgb}{0.000000,0.000000,0.000000}%
\pgfsetfillcolor{currentfill}%
\pgfsetlinewidth{0.803000pt}%
\definecolor{currentstroke}{rgb}{0.000000,0.000000,0.000000}%
\pgfsetstrokecolor{currentstroke}%
\pgfsetdash{}{0pt}%
\pgfsys@defobject{currentmarker}{\pgfqpoint{-0.048611in}{0.000000in}}{\pgfqpoint{-0.000000in}{0.000000in}}{%
\pgfpathmoveto{\pgfqpoint{-0.000000in}{0.000000in}}%
\pgfpathlineto{\pgfqpoint{-0.048611in}{0.000000in}}%
\pgfusepath{stroke,fill}%
}%
\begin{pgfscope}%
\pgfsys@transformshift{1.414736in}{3.233390in}%
\pgfsys@useobject{currentmarker}{}%
\end{pgfscope}%
\end{pgfscope}%
\begin{pgfscope}%
\definecolor{textcolor}{rgb}{0.000000,0.000000,0.000000}%
\pgfsetstrokecolor{textcolor}%
\pgfsetfillcolor{textcolor}%
\pgftext[x=0.606865in, y=3.185164in, left, base]{\color{textcolor}{\rmfamily\fontsize{10.000000}{12.000000}\selectfont\catcode`\^=\active\def^{\ifmmode\sp\else\^{}\fi}\catcode`\%=\active\def%{\%}vel y archer}}%
\end{pgfscope}%
\begin{pgfscope}%
\pgfsetbuttcap%
\pgfsetroundjoin%
\definecolor{currentfill}{rgb}{0.000000,0.000000,0.000000}%
\pgfsetfillcolor{currentfill}%
\pgfsetlinewidth{0.803000pt}%
\definecolor{currentstroke}{rgb}{0.000000,0.000000,0.000000}%
\pgfsetstrokecolor{currentstroke}%
\pgfsetdash{}{0pt}%
\pgfsys@defobject{currentmarker}{\pgfqpoint{-0.048611in}{0.000000in}}{\pgfqpoint{-0.000000in}{0.000000in}}{%
\pgfpathmoveto{\pgfqpoint{-0.000000in}{0.000000in}}%
\pgfpathlineto{\pgfqpoint{-0.048611in}{0.000000in}}%
\pgfusepath{stroke,fill}%
}%
\begin{pgfscope}%
\pgfsys@transformshift{1.414736in}{3.704376in}%
\pgfsys@useobject{currentmarker}{}%
\end{pgfscope}%
\end{pgfscope}%
\begin{pgfscope}%
\definecolor{textcolor}{rgb}{0.000000,0.000000,0.000000}%
\pgfsetstrokecolor{textcolor}%
\pgfsetfillcolor{textcolor}%
\pgftext[x=0.626155in, y=3.656150in, left, base]{\color{textcolor}{\rmfamily\fontsize{10.000000}{12.000000}\selectfont\catcode`\^=\active\def^{\ifmmode\sp\else\^{}\fi}\catcode`\%=\active\def%{\%}dist zombie}}%
\end{pgfscope}%
\begin{pgfscope}%
\pgfsetbuttcap%
\pgfsetroundjoin%
\definecolor{currentfill}{rgb}{0.000000,0.000000,0.000000}%
\pgfsetfillcolor{currentfill}%
\pgfsetlinewidth{0.803000pt}%
\definecolor{currentstroke}{rgb}{0.000000,0.000000,0.000000}%
\pgfsetstrokecolor{currentstroke}%
\pgfsetdash{}{0pt}%
\pgfsys@defobject{currentmarker}{\pgfqpoint{-0.048611in}{0.000000in}}{\pgfqpoint{-0.000000in}{0.000000in}}{%
\pgfpathmoveto{\pgfqpoint{-0.000000in}{0.000000in}}%
\pgfpathlineto{\pgfqpoint{-0.048611in}{0.000000in}}%
\pgfusepath{stroke,fill}%
}%
\begin{pgfscope}%
\pgfsys@transformshift{1.414736in}{4.175362in}%
\pgfsys@useobject{currentmarker}{}%
\end{pgfscope}%
\end{pgfscope}%
\begin{pgfscope}%
\definecolor{textcolor}{rgb}{0.000000,0.000000,0.000000}%
\pgfsetstrokecolor{textcolor}%
\pgfsetfillcolor{textcolor}%
\pgftext[x=0.324842in, y=4.127137in, left, base]{\color{textcolor}{\rmfamily\fontsize{10.000000}{12.000000}\selectfont\catcode`\^=\active\def^{\ifmmode\sp\else\^{}\fi}\catcode`\%=\active\def%{\%}rel pos y zombie}}%
\end{pgfscope}%
\begin{pgfscope}%
\pgfsetbuttcap%
\pgfsetroundjoin%
\definecolor{currentfill}{rgb}{0.000000,0.000000,0.000000}%
\pgfsetfillcolor{currentfill}%
\pgfsetlinewidth{0.803000pt}%
\definecolor{currentstroke}{rgb}{0.000000,0.000000,0.000000}%
\pgfsetstrokecolor{currentstroke}%
\pgfsetdash{}{0pt}%
\pgfsys@defobject{currentmarker}{\pgfqpoint{-0.048611in}{0.000000in}}{\pgfqpoint{-0.000000in}{0.000000in}}{%
\pgfpathmoveto{\pgfqpoint{-0.000000in}{0.000000in}}%
\pgfpathlineto{\pgfqpoint{-0.048611in}{0.000000in}}%
\pgfusepath{stroke,fill}%
}%
\begin{pgfscope}%
\pgfsys@transformshift{1.414736in}{4.646348in}%
\pgfsys@useobject{currentmarker}{}%
\end{pgfscope}%
\end{pgfscope}%
\begin{pgfscope}%
\definecolor{textcolor}{rgb}{0.000000,0.000000,0.000000}%
\pgfsetstrokecolor{textcolor}%
\pgfsetfillcolor{textcolor}%
\pgftext[x=0.629627in, y=4.598123in, left, base]{\color{textcolor}{\rmfamily\fontsize{10.000000}{12.000000}\selectfont\catcode`\^=\active\def^{\ifmmode\sp\else\^{}\fi}\catcode`\%=\active\def%{\%}vel y sword}}%
\end{pgfscope}%
\begin{pgfscope}%
\pgfsetbuttcap%
\pgfsetroundjoin%
\definecolor{currentfill}{rgb}{0.000000,0.000000,0.000000}%
\pgfsetfillcolor{currentfill}%
\pgfsetlinewidth{0.803000pt}%
\definecolor{currentstroke}{rgb}{0.000000,0.000000,0.000000}%
\pgfsetstrokecolor{currentstroke}%
\pgfsetdash{}{0pt}%
\pgfsys@defobject{currentmarker}{\pgfqpoint{-0.048611in}{0.000000in}}{\pgfqpoint{-0.000000in}{0.000000in}}{%
\pgfpathmoveto{\pgfqpoint{-0.000000in}{0.000000in}}%
\pgfpathlineto{\pgfqpoint{-0.048611in}{0.000000in}}%
\pgfusepath{stroke,fill}%
}%
\begin{pgfscope}%
\pgfsys@transformshift{1.414736in}{5.117334in}%
\pgfsys@useobject{currentmarker}{}%
\end{pgfscope}%
\end{pgfscope}%
\begin{pgfscope}%
\definecolor{textcolor}{rgb}{0.000000,0.000000,0.000000}%
\pgfsetstrokecolor{textcolor}%
\pgfsetfillcolor{textcolor}%
\pgftext[x=0.637729in, y=5.069109in, left, base]{\color{textcolor}{\rmfamily\fontsize{10.000000}{12.000000}\selectfont\catcode`\^=\active\def^{\ifmmode\sp\else\^{}\fi}\catcode`\%=\active\def%{\%}vel y arrow}}%
\end{pgfscope}%
\begin{pgfscope}%
\pgfsetbuttcap%
\pgfsetroundjoin%
\definecolor{currentfill}{rgb}{0.000000,0.000000,0.000000}%
\pgfsetfillcolor{currentfill}%
\pgfsetlinewidth{0.803000pt}%
\definecolor{currentstroke}{rgb}{0.000000,0.000000,0.000000}%
\pgfsetstrokecolor{currentstroke}%
\pgfsetdash{}{0pt}%
\pgfsys@defobject{currentmarker}{\pgfqpoint{-0.048611in}{0.000000in}}{\pgfqpoint{-0.000000in}{0.000000in}}{%
\pgfpathmoveto{\pgfqpoint{-0.000000in}{0.000000in}}%
\pgfpathlineto{\pgfqpoint{-0.048611in}{0.000000in}}%
\pgfusepath{stroke,fill}%
}%
\begin{pgfscope}%
\pgfsys@transformshift{1.414736in}{5.588320in}%
\pgfsys@useobject{currentmarker}{}%
\end{pgfscope}%
\end{pgfscope}%
\begin{pgfscope}%
\definecolor{textcolor}{rgb}{0.000000,0.000000,0.000000}%
\pgfsetstrokecolor{textcolor}%
\pgfsetfillcolor{textcolor}%
\pgftext[x=0.401232in, y=5.540095in, left, base]{\color{textcolor}{\rmfamily\fontsize{10.000000}{12.000000}\selectfont\catcode`\^=\active\def^{\ifmmode\sp\else\^{}\fi}\catcode`\%=\active\def%{\%}rel pos x arrow}}%
\end{pgfscope}%
\begin{pgfscope}%
\pgfsetbuttcap%
\pgfsetroundjoin%
\definecolor{currentfill}{rgb}{0.000000,0.000000,0.000000}%
\pgfsetfillcolor{currentfill}%
\pgfsetlinewidth{0.803000pt}%
\definecolor{currentstroke}{rgb}{0.000000,0.000000,0.000000}%
\pgfsetstrokecolor{currentstroke}%
\pgfsetdash{}{0pt}%
\pgfsys@defobject{currentmarker}{\pgfqpoint{-0.048611in}{0.000000in}}{\pgfqpoint{-0.000000in}{0.000000in}}{%
\pgfpathmoveto{\pgfqpoint{-0.000000in}{0.000000in}}%
\pgfpathlineto{\pgfqpoint{-0.048611in}{0.000000in}}%
\pgfusepath{stroke,fill}%
}%
\begin{pgfscope}%
\pgfsys@transformshift{1.414736in}{6.059306in}%
\pgfsys@useobject{currentmarker}{}%
\end{pgfscope}%
\end{pgfscope}%
\begin{pgfscope}%
\definecolor{textcolor}{rgb}{0.000000,0.000000,0.000000}%
\pgfsetstrokecolor{textcolor}%
\pgfsetfillcolor{textcolor}%
\pgftext[x=0.637729in, y=6.011081in, left, base]{\color{textcolor}{\rmfamily\fontsize{10.000000}{12.000000}\selectfont\catcode`\^=\active\def^{\ifmmode\sp\else\^{}\fi}\catcode`\%=\active\def%{\%}vel x arrow}}%
\end{pgfscope}%
\begin{pgfscope}%
\pgfsetbuttcap%
\pgfsetroundjoin%
\definecolor{currentfill}{rgb}{0.000000,0.000000,0.000000}%
\pgfsetfillcolor{currentfill}%
\pgfsetlinewidth{0.803000pt}%
\definecolor{currentstroke}{rgb}{0.000000,0.000000,0.000000}%
\pgfsetstrokecolor{currentstroke}%
\pgfsetdash{}{0pt}%
\pgfsys@defobject{currentmarker}{\pgfqpoint{-0.048611in}{0.000000in}}{\pgfqpoint{-0.000000in}{0.000000in}}{%
\pgfpathmoveto{\pgfqpoint{-0.000000in}{0.000000in}}%
\pgfpathlineto{\pgfqpoint{-0.048611in}{0.000000in}}%
\pgfusepath{stroke,fill}%
}%
\begin{pgfscope}%
\pgfsys@transformshift{1.414736in}{6.530292in}%
\pgfsys@useobject{currentmarker}{}%
\end{pgfscope}%
\end{pgfscope}%
\begin{pgfscope}%
\definecolor{textcolor}{rgb}{0.000000,0.000000,0.000000}%
\pgfsetstrokecolor{textcolor}%
\pgfsetfillcolor{textcolor}%
\pgftext[x=0.561340in, y=6.482067in, left, base]{\color{textcolor}{\rmfamily\fontsize{10.000000}{12.000000}\selectfont\catcode`\^=\active\def^{\ifmmode\sp\else\^{}\fi}\catcode`\%=\active\def%{\%}vel x zombie}}%
\end{pgfscope}%
\begin{pgfscope}%
\pgfsetbuttcap%
\pgfsetroundjoin%
\definecolor{currentfill}{rgb}{0.000000,0.000000,0.000000}%
\pgfsetfillcolor{currentfill}%
\pgfsetlinewidth{0.803000pt}%
\definecolor{currentstroke}{rgb}{0.000000,0.000000,0.000000}%
\pgfsetstrokecolor{currentstroke}%
\pgfsetdash{}{0pt}%
\pgfsys@defobject{currentmarker}{\pgfqpoint{-0.048611in}{0.000000in}}{\pgfqpoint{-0.000000in}{0.000000in}}{%
\pgfpathmoveto{\pgfqpoint{-0.000000in}{0.000000in}}%
\pgfpathlineto{\pgfqpoint{-0.048611in}{0.000000in}}%
\pgfusepath{stroke,fill}%
}%
\begin{pgfscope}%
\pgfsys@transformshift{1.414736in}{7.001278in}%
\pgfsys@useobject{currentmarker}{}%
\end{pgfscope}%
\end{pgfscope}%
\begin{pgfscope}%
\definecolor{textcolor}{rgb}{0.000000,0.000000,0.000000}%
\pgfsetstrokecolor{textcolor}%
\pgfsetfillcolor{textcolor}%
\pgftext[x=0.561340in, y=6.953053in, left, base]{\color{textcolor}{\rmfamily\fontsize{10.000000}{12.000000}\selectfont\catcode`\^=\active\def^{\ifmmode\sp\else\^{}\fi}\catcode`\%=\active\def%{\%}vel y zombie}}%
\end{pgfscope}%
\begin{pgfscope}%
\pgfsetbuttcap%
\pgfsetroundjoin%
\definecolor{currentfill}{rgb}{0.000000,0.000000,0.000000}%
\pgfsetfillcolor{currentfill}%
\pgfsetlinewidth{0.803000pt}%
\definecolor{currentstroke}{rgb}{0.000000,0.000000,0.000000}%
\pgfsetstrokecolor{currentstroke}%
\pgfsetdash{}{0pt}%
\pgfsys@defobject{currentmarker}{\pgfqpoint{-0.048611in}{0.000000in}}{\pgfqpoint{-0.000000in}{0.000000in}}{%
\pgfpathmoveto{\pgfqpoint{-0.000000in}{0.000000in}}%
\pgfpathlineto{\pgfqpoint{-0.048611in}{0.000000in}}%
\pgfusepath{stroke,fill}%
}%
\begin{pgfscope}%
\pgfsys@transformshift{1.414736in}{7.472265in}%
\pgfsys@useobject{currentmarker}{}%
\end{pgfscope}%
\end{pgfscope}%
\begin{pgfscope}%
\definecolor{textcolor}{rgb}{0.000000,0.000000,0.000000}%
\pgfsetstrokecolor{textcolor}%
\pgfsetfillcolor{textcolor}%
\pgftext[x=0.324842in, y=7.424039in, left, base]{\color{textcolor}{\rmfamily\fontsize{10.000000}{12.000000}\selectfont\catcode`\^=\active\def^{\ifmmode\sp\else\^{}\fi}\catcode`\%=\active\def%{\%}rel pos x zombie}}%
\end{pgfscope}%
\begin{pgfscope}%
\definecolor{textcolor}{rgb}{0.000000,0.000000,0.000000}%
\pgfsetstrokecolor{textcolor}%
\pgfsetfillcolor{textcolor}%
\pgftext[x=0.269287in,y=4.175362in,,bottom,rotate=90.000000]{\color{textcolor}{\rmfamily\fontsize{10.000000}{12.000000}\selectfont\catcode`\^=\active\def^{\ifmmode\sp\else\^{}\fi}\catcode`\%=\active\def%{\%}Feature}}%
\end{pgfscope}%
\begin{pgfscope}%
\pgfpathrectangle{\pgfqpoint{1.414736in}{0.548769in}}{\pgfqpoint{3.497747in}{7.253186in}}%
\pgfusepath{clip}%
\pgfsetbuttcap%
\pgfsetroundjoin%
\pgfsetlinewidth{0.501875pt}%
\definecolor{currentstroke}{rgb}{0.501961,0.501961,0.501961}%
\pgfsetstrokecolor{currentstroke}%
\pgfsetdash{{1.850000pt}{0.800000pt}}{0.000000pt}%
\pgfpathmoveto{\pgfqpoint{1.414736in}{0.878459in}}%
\pgfpathlineto{\pgfqpoint{4.912483in}{0.878459in}}%
\pgfusepath{stroke}%
\end{pgfscope}%
\begin{pgfscope}%
\pgfpathrectangle{\pgfqpoint{1.414736in}{0.548769in}}{\pgfqpoint{3.497747in}{7.253186in}}%
\pgfusepath{clip}%
\pgfsetbuttcap%
\pgfsetroundjoin%
\pgfsetlinewidth{0.501875pt}%
\definecolor{currentstroke}{rgb}{0.501961,0.501961,0.501961}%
\pgfsetstrokecolor{currentstroke}%
\pgfsetdash{{1.850000pt}{0.800000pt}}{0.000000pt}%
\pgfpathmoveto{\pgfqpoint{1.414736in}{1.349445in}}%
\pgfpathlineto{\pgfqpoint{4.912483in}{1.349445in}}%
\pgfusepath{stroke}%
\end{pgfscope}%
\begin{pgfscope}%
\pgfpathrectangle{\pgfqpoint{1.414736in}{0.548769in}}{\pgfqpoint{3.497747in}{7.253186in}}%
\pgfusepath{clip}%
\pgfsetbuttcap%
\pgfsetroundjoin%
\pgfsetlinewidth{0.501875pt}%
\definecolor{currentstroke}{rgb}{0.501961,0.501961,0.501961}%
\pgfsetstrokecolor{currentstroke}%
\pgfsetdash{{1.850000pt}{0.800000pt}}{0.000000pt}%
\pgfpathmoveto{\pgfqpoint{1.414736in}{1.820431in}}%
\pgfpathlineto{\pgfqpoint{4.912483in}{1.820431in}}%
\pgfusepath{stroke}%
\end{pgfscope}%
\begin{pgfscope}%
\pgfpathrectangle{\pgfqpoint{1.414736in}{0.548769in}}{\pgfqpoint{3.497747in}{7.253186in}}%
\pgfusepath{clip}%
\pgfsetbuttcap%
\pgfsetroundjoin%
\pgfsetlinewidth{0.501875pt}%
\definecolor{currentstroke}{rgb}{0.501961,0.501961,0.501961}%
\pgfsetstrokecolor{currentstroke}%
\pgfsetdash{{1.850000pt}{0.800000pt}}{0.000000pt}%
\pgfpathmoveto{\pgfqpoint{1.414736in}{2.291417in}}%
\pgfpathlineto{\pgfqpoint{4.912483in}{2.291417in}}%
\pgfusepath{stroke}%
\end{pgfscope}%
\begin{pgfscope}%
\pgfpathrectangle{\pgfqpoint{1.414736in}{0.548769in}}{\pgfqpoint{3.497747in}{7.253186in}}%
\pgfusepath{clip}%
\pgfsetbuttcap%
\pgfsetroundjoin%
\pgfsetlinewidth{0.501875pt}%
\definecolor{currentstroke}{rgb}{0.501961,0.501961,0.501961}%
\pgfsetstrokecolor{currentstroke}%
\pgfsetdash{{1.850000pt}{0.800000pt}}{0.000000pt}%
\pgfpathmoveto{\pgfqpoint{1.414736in}{2.762404in}}%
\pgfpathlineto{\pgfqpoint{4.912483in}{2.762404in}}%
\pgfusepath{stroke}%
\end{pgfscope}%
\begin{pgfscope}%
\pgfpathrectangle{\pgfqpoint{1.414736in}{0.548769in}}{\pgfqpoint{3.497747in}{7.253186in}}%
\pgfusepath{clip}%
\pgfsetbuttcap%
\pgfsetroundjoin%
\pgfsetlinewidth{0.501875pt}%
\definecolor{currentstroke}{rgb}{0.501961,0.501961,0.501961}%
\pgfsetstrokecolor{currentstroke}%
\pgfsetdash{{1.850000pt}{0.800000pt}}{0.000000pt}%
\pgfpathmoveto{\pgfqpoint{1.414736in}{3.233390in}}%
\pgfpathlineto{\pgfqpoint{4.912483in}{3.233390in}}%
\pgfusepath{stroke}%
\end{pgfscope}%
\begin{pgfscope}%
\pgfpathrectangle{\pgfqpoint{1.414736in}{0.548769in}}{\pgfqpoint{3.497747in}{7.253186in}}%
\pgfusepath{clip}%
\pgfsetbuttcap%
\pgfsetroundjoin%
\pgfsetlinewidth{0.501875pt}%
\definecolor{currentstroke}{rgb}{0.501961,0.501961,0.501961}%
\pgfsetstrokecolor{currentstroke}%
\pgfsetdash{{1.850000pt}{0.800000pt}}{0.000000pt}%
\pgfpathmoveto{\pgfqpoint{1.414736in}{3.704376in}}%
\pgfpathlineto{\pgfqpoint{4.912483in}{3.704376in}}%
\pgfusepath{stroke}%
\end{pgfscope}%
\begin{pgfscope}%
\pgfpathrectangle{\pgfqpoint{1.414736in}{0.548769in}}{\pgfqpoint{3.497747in}{7.253186in}}%
\pgfusepath{clip}%
\pgfsetbuttcap%
\pgfsetroundjoin%
\pgfsetlinewidth{0.501875pt}%
\definecolor{currentstroke}{rgb}{0.501961,0.501961,0.501961}%
\pgfsetstrokecolor{currentstroke}%
\pgfsetdash{{1.850000pt}{0.800000pt}}{0.000000pt}%
\pgfpathmoveto{\pgfqpoint{1.414736in}{4.175362in}}%
\pgfpathlineto{\pgfqpoint{4.912483in}{4.175362in}}%
\pgfusepath{stroke}%
\end{pgfscope}%
\begin{pgfscope}%
\pgfpathrectangle{\pgfqpoint{1.414736in}{0.548769in}}{\pgfqpoint{3.497747in}{7.253186in}}%
\pgfusepath{clip}%
\pgfsetbuttcap%
\pgfsetroundjoin%
\pgfsetlinewidth{0.501875pt}%
\definecolor{currentstroke}{rgb}{0.501961,0.501961,0.501961}%
\pgfsetstrokecolor{currentstroke}%
\pgfsetdash{{1.850000pt}{0.800000pt}}{0.000000pt}%
\pgfpathmoveto{\pgfqpoint{1.414736in}{4.646348in}}%
\pgfpathlineto{\pgfqpoint{4.912483in}{4.646348in}}%
\pgfusepath{stroke}%
\end{pgfscope}%
\begin{pgfscope}%
\pgfpathrectangle{\pgfqpoint{1.414736in}{0.548769in}}{\pgfqpoint{3.497747in}{7.253186in}}%
\pgfusepath{clip}%
\pgfsetbuttcap%
\pgfsetroundjoin%
\pgfsetlinewidth{0.501875pt}%
\definecolor{currentstroke}{rgb}{0.501961,0.501961,0.501961}%
\pgfsetstrokecolor{currentstroke}%
\pgfsetdash{{1.850000pt}{0.800000pt}}{0.000000pt}%
\pgfpathmoveto{\pgfqpoint{1.414736in}{5.117334in}}%
\pgfpathlineto{\pgfqpoint{4.912483in}{5.117334in}}%
\pgfusepath{stroke}%
\end{pgfscope}%
\begin{pgfscope}%
\pgfpathrectangle{\pgfqpoint{1.414736in}{0.548769in}}{\pgfqpoint{3.497747in}{7.253186in}}%
\pgfusepath{clip}%
\pgfsetbuttcap%
\pgfsetroundjoin%
\pgfsetlinewidth{0.501875pt}%
\definecolor{currentstroke}{rgb}{0.501961,0.501961,0.501961}%
\pgfsetstrokecolor{currentstroke}%
\pgfsetdash{{1.850000pt}{0.800000pt}}{0.000000pt}%
\pgfpathmoveto{\pgfqpoint{1.414736in}{5.588320in}}%
\pgfpathlineto{\pgfqpoint{4.912483in}{5.588320in}}%
\pgfusepath{stroke}%
\end{pgfscope}%
\begin{pgfscope}%
\pgfpathrectangle{\pgfqpoint{1.414736in}{0.548769in}}{\pgfqpoint{3.497747in}{7.253186in}}%
\pgfusepath{clip}%
\pgfsetbuttcap%
\pgfsetroundjoin%
\pgfsetlinewidth{0.501875pt}%
\definecolor{currentstroke}{rgb}{0.501961,0.501961,0.501961}%
\pgfsetstrokecolor{currentstroke}%
\pgfsetdash{{1.850000pt}{0.800000pt}}{0.000000pt}%
\pgfpathmoveto{\pgfqpoint{1.414736in}{6.059306in}}%
\pgfpathlineto{\pgfqpoint{4.912483in}{6.059306in}}%
\pgfusepath{stroke}%
\end{pgfscope}%
\begin{pgfscope}%
\pgfpathrectangle{\pgfqpoint{1.414736in}{0.548769in}}{\pgfqpoint{3.497747in}{7.253186in}}%
\pgfusepath{clip}%
\pgfsetbuttcap%
\pgfsetroundjoin%
\pgfsetlinewidth{0.501875pt}%
\definecolor{currentstroke}{rgb}{0.501961,0.501961,0.501961}%
\pgfsetstrokecolor{currentstroke}%
\pgfsetdash{{1.850000pt}{0.800000pt}}{0.000000pt}%
\pgfpathmoveto{\pgfqpoint{1.414736in}{6.530292in}}%
\pgfpathlineto{\pgfqpoint{4.912483in}{6.530292in}}%
\pgfusepath{stroke}%
\end{pgfscope}%
\begin{pgfscope}%
\pgfpathrectangle{\pgfqpoint{1.414736in}{0.548769in}}{\pgfqpoint{3.497747in}{7.253186in}}%
\pgfusepath{clip}%
\pgfsetbuttcap%
\pgfsetroundjoin%
\pgfsetlinewidth{0.501875pt}%
\definecolor{currentstroke}{rgb}{0.501961,0.501961,0.501961}%
\pgfsetstrokecolor{currentstroke}%
\pgfsetdash{{1.850000pt}{0.800000pt}}{0.000000pt}%
\pgfpathmoveto{\pgfqpoint{1.414736in}{7.001278in}}%
\pgfpathlineto{\pgfqpoint{4.912483in}{7.001278in}}%
\pgfusepath{stroke}%
\end{pgfscope}%
\begin{pgfscope}%
\pgfpathrectangle{\pgfqpoint{1.414736in}{0.548769in}}{\pgfqpoint{3.497747in}{7.253186in}}%
\pgfusepath{clip}%
\pgfsetbuttcap%
\pgfsetroundjoin%
\pgfsetlinewidth{0.501875pt}%
\definecolor{currentstroke}{rgb}{0.501961,0.501961,0.501961}%
\pgfsetstrokecolor{currentstroke}%
\pgfsetdash{{1.850000pt}{0.800000pt}}{0.000000pt}%
\pgfpathmoveto{\pgfqpoint{1.414736in}{7.472265in}}%
\pgfpathlineto{\pgfqpoint{4.912483in}{7.472265in}}%
\pgfusepath{stroke}%
\end{pgfscope}%
\begin{pgfscope}%
\pgfpathrectangle{\pgfqpoint{1.414736in}{0.548769in}}{\pgfqpoint{3.497747in}{7.253186in}}%
\pgfusepath{clip}%
\pgfsetrectcap%
\pgfsetroundjoin%
\pgfsetlinewidth{0.501875pt}%
\definecolor{currentstroke}{rgb}{0.501961,0.501961,0.501961}%
\pgfsetstrokecolor{currentstroke}%
\pgfsetdash{}{0pt}%
\pgfpathmoveto{\pgfqpoint{3.573979in}{0.548769in}}%
\pgfpathlineto{\pgfqpoint{3.573979in}{7.801955in}}%
\pgfusepath{stroke}%
\end{pgfscope}%
\begin{pgfscope}%
\pgfsetrectcap%
\pgfsetmiterjoin%
\pgfsetlinewidth{0.803000pt}%
\definecolor{currentstroke}{rgb}{0.000000,0.000000,0.000000}%
\pgfsetstrokecolor{currentstroke}%
\pgfsetdash{}{0pt}%
\pgfpathmoveto{\pgfqpoint{1.414736in}{0.548769in}}%
\pgfpathlineto{\pgfqpoint{1.414736in}{7.801955in}}%
\pgfusepath{stroke}%
\end{pgfscope}%
\begin{pgfscope}%
\pgfsetrectcap%
\pgfsetmiterjoin%
\pgfsetlinewidth{0.803000pt}%
\definecolor{currentstroke}{rgb}{0.000000,0.000000,0.000000}%
\pgfsetstrokecolor{currentstroke}%
\pgfsetdash{}{0pt}%
\pgfpathmoveto{\pgfqpoint{4.912483in}{0.548769in}}%
\pgfpathlineto{\pgfqpoint{4.912483in}{7.801955in}}%
\pgfusepath{stroke}%
\end{pgfscope}%
\begin{pgfscope}%
\pgfsetrectcap%
\pgfsetmiterjoin%
\pgfsetlinewidth{0.803000pt}%
\definecolor{currentstroke}{rgb}{0.000000,0.000000,0.000000}%
\pgfsetstrokecolor{currentstroke}%
\pgfsetdash{}{0pt}%
\pgfpathmoveto{\pgfqpoint{1.414736in}{0.548769in}}%
\pgfpathlineto{\pgfqpoint{4.912483in}{0.548769in}}%
\pgfusepath{stroke}%
\end{pgfscope}%
\begin{pgfscope}%
\pgfsetrectcap%
\pgfsetmiterjoin%
\pgfsetlinewidth{0.803000pt}%
\definecolor{currentstroke}{rgb}{0.000000,0.000000,0.000000}%
\pgfsetstrokecolor{currentstroke}%
\pgfsetdash{}{0pt}%
\pgfpathmoveto{\pgfqpoint{1.414736in}{7.801955in}}%
\pgfpathlineto{\pgfqpoint{4.912483in}{7.801955in}}%
\pgfusepath{stroke}%
\end{pgfscope}%
\begin{pgfscope}%
\pgfsetbuttcap%
\pgfsetmiterjoin%
\definecolor{currentfill}{rgb}{1.000000,1.000000,1.000000}%
\pgfsetfillcolor{currentfill}%
\pgfsetlinewidth{0.000000pt}%
\definecolor{currentstroke}{rgb}{0.000000,0.000000,0.000000}%
\pgfsetstrokecolor{currentstroke}%
\pgfsetstrokeopacity{0.000000}%
\pgfsetdash{}{0pt}%
\pgfpathmoveto{\pgfqpoint{5.131092in}{0.548769in}}%
\pgfpathlineto{\pgfqpoint{5.493752in}{0.548769in}}%
\pgfpathlineto{\pgfqpoint{5.493752in}{7.801955in}}%
\pgfpathlineto{\pgfqpoint{5.131092in}{7.801955in}}%
\pgfpathlineto{\pgfqpoint{5.131092in}{0.548769in}}%
\pgfpathclose%
\pgfusepath{fill}%
\end{pgfscope}%
\begin{pgfscope}%
\pgfsys@transformshift{5.130000in}{0.550000in}%
\pgftext[left,bottom]{\includegraphics[interpolate=true,width=0.360000in,height=7.250000in]{0_none_none_shap-img0.png}}%
\end{pgfscope}%
\begin{pgfscope}%
\pgfsetbuttcap%
\pgfsetroundjoin%
\definecolor{currentfill}{rgb}{0.000000,0.000000,0.000000}%
\pgfsetfillcolor{currentfill}%
\pgfsetlinewidth{0.803000pt}%
\definecolor{currentstroke}{rgb}{0.000000,0.000000,0.000000}%
\pgfsetstrokecolor{currentstroke}%
\pgfsetdash{}{0pt}%
\pgfsys@defobject{currentmarker}{\pgfqpoint{0.000000in}{0.000000in}}{\pgfqpoint{0.048611in}{0.000000in}}{%
\pgfpathmoveto{\pgfqpoint{0.000000in}{0.000000in}}%
\pgfpathlineto{\pgfqpoint{0.048611in}{0.000000in}}%
\pgfusepath{stroke,fill}%
}%
\begin{pgfscope}%
\pgfsys@transformshift{5.493752in}{0.548769in}%
\pgfsys@useobject{currentmarker}{}%
\end{pgfscope}%
\end{pgfscope}%
\begin{pgfscope}%
\definecolor{textcolor}{rgb}{0.000000,0.000000,0.000000}%
\pgfsetstrokecolor{textcolor}%
\pgfsetfillcolor{textcolor}%
\pgftext[x=5.590974in, y=0.500544in, left, base]{\color{textcolor}{\rmfamily\fontsize{10.000000}{12.000000}\selectfont\catcode`\^=\active\def^{\ifmmode\sp\else\^{}\fi}\catcode`\%=\active\def%{\%}$\mathdefault{0.0}$}}%
\end{pgfscope}%
\begin{pgfscope}%
\pgfsetbuttcap%
\pgfsetroundjoin%
\definecolor{currentfill}{rgb}{0.000000,0.000000,0.000000}%
\pgfsetfillcolor{currentfill}%
\pgfsetlinewidth{0.803000pt}%
\definecolor{currentstroke}{rgb}{0.000000,0.000000,0.000000}%
\pgfsetstrokecolor{currentstroke}%
\pgfsetdash{}{0pt}%
\pgfsys@defobject{currentmarker}{\pgfqpoint{0.000000in}{0.000000in}}{\pgfqpoint{0.048611in}{0.000000in}}{%
\pgfpathmoveto{\pgfqpoint{0.000000in}{0.000000in}}%
\pgfpathlineto{\pgfqpoint{0.048611in}{0.000000in}}%
\pgfusepath{stroke,fill}%
}%
\begin{pgfscope}%
\pgfsys@transformshift{5.493752in}{1.999406in}%
\pgfsys@useobject{currentmarker}{}%
\end{pgfscope}%
\end{pgfscope}%
\begin{pgfscope}%
\definecolor{textcolor}{rgb}{0.000000,0.000000,0.000000}%
\pgfsetstrokecolor{textcolor}%
\pgfsetfillcolor{textcolor}%
\pgftext[x=5.590974in, y=1.951181in, left, base]{\color{textcolor}{\rmfamily\fontsize{10.000000}{12.000000}\selectfont\catcode`\^=\active\def^{\ifmmode\sp\else\^{}\fi}\catcode`\%=\active\def%{\%}$\mathdefault{0.2}$}}%
\end{pgfscope}%
\begin{pgfscope}%
\pgfsetbuttcap%
\pgfsetroundjoin%
\definecolor{currentfill}{rgb}{0.000000,0.000000,0.000000}%
\pgfsetfillcolor{currentfill}%
\pgfsetlinewidth{0.803000pt}%
\definecolor{currentstroke}{rgb}{0.000000,0.000000,0.000000}%
\pgfsetstrokecolor{currentstroke}%
\pgfsetdash{}{0pt}%
\pgfsys@defobject{currentmarker}{\pgfqpoint{0.000000in}{0.000000in}}{\pgfqpoint{0.048611in}{0.000000in}}{%
\pgfpathmoveto{\pgfqpoint{0.000000in}{0.000000in}}%
\pgfpathlineto{\pgfqpoint{0.048611in}{0.000000in}}%
\pgfusepath{stroke,fill}%
}%
\begin{pgfscope}%
\pgfsys@transformshift{5.493752in}{3.450043in}%
\pgfsys@useobject{currentmarker}{}%
\end{pgfscope}%
\end{pgfscope}%
\begin{pgfscope}%
\definecolor{textcolor}{rgb}{0.000000,0.000000,0.000000}%
\pgfsetstrokecolor{textcolor}%
\pgfsetfillcolor{textcolor}%
\pgftext[x=5.590974in, y=3.401818in, left, base]{\color{textcolor}{\rmfamily\fontsize{10.000000}{12.000000}\selectfont\catcode`\^=\active\def^{\ifmmode\sp\else\^{}\fi}\catcode`\%=\active\def%{\%}$\mathdefault{0.4}$}}%
\end{pgfscope}%
\begin{pgfscope}%
\pgfsetbuttcap%
\pgfsetroundjoin%
\definecolor{currentfill}{rgb}{0.000000,0.000000,0.000000}%
\pgfsetfillcolor{currentfill}%
\pgfsetlinewidth{0.803000pt}%
\definecolor{currentstroke}{rgb}{0.000000,0.000000,0.000000}%
\pgfsetstrokecolor{currentstroke}%
\pgfsetdash{}{0pt}%
\pgfsys@defobject{currentmarker}{\pgfqpoint{0.000000in}{0.000000in}}{\pgfqpoint{0.048611in}{0.000000in}}{%
\pgfpathmoveto{\pgfqpoint{0.000000in}{0.000000in}}%
\pgfpathlineto{\pgfqpoint{0.048611in}{0.000000in}}%
\pgfusepath{stroke,fill}%
}%
\begin{pgfscope}%
\pgfsys@transformshift{5.493752in}{4.900680in}%
\pgfsys@useobject{currentmarker}{}%
\end{pgfscope}%
\end{pgfscope}%
\begin{pgfscope}%
\definecolor{textcolor}{rgb}{0.000000,0.000000,0.000000}%
\pgfsetstrokecolor{textcolor}%
\pgfsetfillcolor{textcolor}%
\pgftext[x=5.590974in, y=4.852455in, left, base]{\color{textcolor}{\rmfamily\fontsize{10.000000}{12.000000}\selectfont\catcode`\^=\active\def^{\ifmmode\sp\else\^{}\fi}\catcode`\%=\active\def%{\%}$\mathdefault{0.6}$}}%
\end{pgfscope}%
\begin{pgfscope}%
\pgfsetbuttcap%
\pgfsetroundjoin%
\definecolor{currentfill}{rgb}{0.000000,0.000000,0.000000}%
\pgfsetfillcolor{currentfill}%
\pgfsetlinewidth{0.803000pt}%
\definecolor{currentstroke}{rgb}{0.000000,0.000000,0.000000}%
\pgfsetstrokecolor{currentstroke}%
\pgfsetdash{}{0pt}%
\pgfsys@defobject{currentmarker}{\pgfqpoint{0.000000in}{0.000000in}}{\pgfqpoint{0.048611in}{0.000000in}}{%
\pgfpathmoveto{\pgfqpoint{0.000000in}{0.000000in}}%
\pgfpathlineto{\pgfqpoint{0.048611in}{0.000000in}}%
\pgfusepath{stroke,fill}%
}%
\begin{pgfscope}%
\pgfsys@transformshift{5.493752in}{6.351318in}%
\pgfsys@useobject{currentmarker}{}%
\end{pgfscope}%
\end{pgfscope}%
\begin{pgfscope}%
\definecolor{textcolor}{rgb}{0.000000,0.000000,0.000000}%
\pgfsetstrokecolor{textcolor}%
\pgfsetfillcolor{textcolor}%
\pgftext[x=5.590974in, y=6.303092in, left, base]{\color{textcolor}{\rmfamily\fontsize{10.000000}{12.000000}\selectfont\catcode`\^=\active\def^{\ifmmode\sp\else\^{}\fi}\catcode`\%=\active\def%{\%}$\mathdefault{0.8}$}}%
\end{pgfscope}%
\begin{pgfscope}%
\pgfsetbuttcap%
\pgfsetroundjoin%
\definecolor{currentfill}{rgb}{0.000000,0.000000,0.000000}%
\pgfsetfillcolor{currentfill}%
\pgfsetlinewidth{0.803000pt}%
\definecolor{currentstroke}{rgb}{0.000000,0.000000,0.000000}%
\pgfsetstrokecolor{currentstroke}%
\pgfsetdash{}{0pt}%
\pgfsys@defobject{currentmarker}{\pgfqpoint{0.000000in}{0.000000in}}{\pgfqpoint{0.048611in}{0.000000in}}{%
\pgfpathmoveto{\pgfqpoint{0.000000in}{0.000000in}}%
\pgfpathlineto{\pgfqpoint{0.048611in}{0.000000in}}%
\pgfusepath{stroke,fill}%
}%
\begin{pgfscope}%
\pgfsys@transformshift{5.493752in}{7.801955in}%
\pgfsys@useobject{currentmarker}{}%
\end{pgfscope}%
\end{pgfscope}%
\begin{pgfscope}%
\definecolor{textcolor}{rgb}{0.000000,0.000000,0.000000}%
\pgfsetstrokecolor{textcolor}%
\pgfsetfillcolor{textcolor}%
\pgftext[x=5.590974in, y=7.753730in, left, base]{\color{textcolor}{\rmfamily\fontsize{10.000000}{12.000000}\selectfont\catcode`\^=\active\def^{\ifmmode\sp\else\^{}\fi}\catcode`\%=\active\def%{\%}$\mathdefault{1.0}$}}%
\end{pgfscope}%
\begin{pgfscope}%
\definecolor{textcolor}{rgb}{0.000000,0.000000,0.000000}%
\pgfsetstrokecolor{textcolor}%
\pgfsetfillcolor{textcolor}%
\pgftext[x=5.823999in,y=4.175362in,,top,rotate=90.000000]{\color{textcolor}{\rmfamily\fontsize{10.000000}{12.000000}\selectfont\catcode`\^=\active\def^{\ifmmode\sp\else\^{}\fi}\catcode`\%=\active\def%{\%}Feature value}}%
\end{pgfscope}%
\begin{pgfscope}%
\pgfsetrectcap%
\pgfsetmiterjoin%
\pgfsetlinewidth{0.803000pt}%
\definecolor{currentstroke}{rgb}{0.000000,0.000000,0.000000}%
\pgfsetstrokecolor{currentstroke}%
\pgfsetdash{}{0pt}%
\pgfpathmoveto{\pgfqpoint{5.131092in}{0.548769in}}%
\pgfpathlineto{\pgfqpoint{5.312422in}{0.548769in}}%
\pgfpathlineto{\pgfqpoint{5.493752in}{0.548769in}}%
\pgfpathlineto{\pgfqpoint{5.493752in}{7.801955in}}%
\pgfpathlineto{\pgfqpoint{5.312422in}{7.801955in}}%
\pgfpathlineto{\pgfqpoint{5.131092in}{7.801955in}}%
\pgfpathlineto{\pgfqpoint{5.131092in}{0.548769in}}%
\pgfpathclose%
\pgfusepath{stroke}%
\end{pgfscope}%
\end{pgfpicture}%
\makeatother%
\endgroup%

\end{figure}

\section{Experiment 3}
\subsection{Setup}
\subsection{Results}
\subsection{Analysis}

\chapter{Discussion}

\section{Conclusion}

\section{Limitations}

\section{Future work}


\printbibliography{}
\end{document}
